% Auto-generated from data/segments.json + layout.json (+ transforms) — do not edit by hand.
% volume: 38Abhi10
% mode: reading
% segments: 5741
% transform_rules: 0
% toc_marks: 150

% Volume layout defaults (layout.json ``layout``).
\csromanlayoutapply{1}{1.5}{21.6pt}{5pt}{6.3pt}{65pt}{2.5em}%

% Reading physical-page layout (layout.json ``page_layout_reading_mode``).
\csromanreadingpagelayoutvolume{1}{1.5}{21.6pt}{5pt}{6.3pt}{65pt}{2.5em}%
\csromanreadingpagelayoutenable

\pitaka{\textbf{อภิธมฺมปิฏก}}
\csromanheader{2}{\textbf{ธมฺมานุโลม}}
\csromanmatikaanchor{mtk.0}
\csromantocmarknopage{chapter}{ปฏฺฐานปาฬิ}{mtk.0}
\bookmark[dest={mtk.0},level=0]{ปฏฺฐานปาฬิ}
\gambhira{\textbf{ทุกปฏฺฐานปาฬิ}}
\namakkaram{นโม ตสฺส ภควโต อรหโต สมฺมาสมฺพุทฺธสฺส.}
\csromanmatikaanchor{mtk.1}
\csromanmatikaanchor{mtk.2}
\csromantocmarknopage{section}{ตติยภาค}{mtk.1}
\bookmark[dest={mtk.1},level=1]{ตติยภาค}
\csromantocmarknopage{chapter}{1. เหตุโคจฺฉก}{mtk.2}
\bookmark[dest={mtk.2},level=0]{1. เหตุโคจฺฉก}
\chapterhead{1. \textbf{เหตุ}โคจฺฉก}
\csromanmatikaanchor{mtk.3}
\csromantocmarkref{section}{1. เหตุทุก}{mtk.3}
\bookmark[dest={mtk.3},level=1]{1. เหตุทุก}
\csromanheader{1}{1. \textbf{เหตุ}ทุก 1. ปฏิจฺจวาร}
\csromanheader{2}{1. ปจฺจยานุโลม 1. วิภงฺควาร}
\csromanheader{2}{\textbf{เหตุ}}
\vspace{\tipitakaparskip}
\csromancenter{\textbf{เหตุ}ํ ธมฺมํ ปฏิจฺจ เหตุ ธมฺโม อุปฺปชฺชติ เหตุปจฺจยา.\csromanspacer{} อโลภํ ปฏิจฺจ อโทโส อโมโห.\csromanspacer{} อโทสํ ปฏิจฺจ อโลโภ อโมโห.\csromanspacer{} อโมหํ ปฏิจฺจ อโลโภ อโทโส.\csromanspacer{} โลภํ ปฏิจฺจ โมโห.\csromanspacer{} โมหํ ปฏิจฺจ โลโภ.\csromanspacer{} โทสํ ปฏิจฺจ โมโห.\csromanspacer{} โมหํ ปฏิจฺจ โทโส.\csromanspacer{} ปฏิสนฺธิกฺขเณ ฯเปฯ. (1)}
\csromancenter{\textbf{เหตุ}ํ ธมฺมํ ปฏิจฺจ นเหตุ ธมฺโม อุปฺปชฺชติ เหตุปจฺจยา.\csromanspacer{} \textbf{เหตุ}ํ ธมฺมํ ปฏิจฺจ สมฺปยุตฺตกา ขนฺธา จิตฺตสมุฏฺฐานญฺจ รูปํ.\csromanspacer{} ปฏิสนฺธิกฺขเณ ฯเปฯ. (2)}
\csromancenter{\textbf{เหตุ}ํ ธมฺมํ ปฏิจฺจ เหตุ จ นเหตุ จ ธมฺมา อุปฺปชฺชนฺติ เหตุปจฺจยา.\csromanspacer{} อโลภํ ปฏิจฺจ อโทโส อโมโห สมฺปยุตฺตกา จ ขนฺธา จิตฺตสมุฏฺฐานญฺจ รูปํ. (จกฺกํ.) โลภํ ปฏิจฺจ โมโห สมฺปยุตฺตกา จ ขนฺธา จิตฺตสมุฏฺฐานญฺจ รูปํ ฯเปฯ.\csromanspacer{} ปฏิสนฺธิกฺขเณ ฯเปฯ. (3)}
\proseitem{2}{\csromanfolio{2}นเหตุํ ธมฺมํ ปฏิจฺจ นเหตุ ธมฺโม อุปฺปชฺชติ เหตุปจฺจยา.\csromanspacer{} นเหตุํ เอกํ ขนฺธํ ปฏิจฺจ ตโย ขนฺธา จิตฺตสมุฏฺฐานญฺจ รูปํ ฯเปฯ เทฺว ขนฺเธ ปฏิจฺจ เทฺว ขนฺธา จิตฺตสมุฏฺฐานญฺจ รูปํ.\csromanspacer{} ปฏิสนฺธิกฺขเณ ฯเปฯ.\csromanspacer{} ขนฺเธ ปฏิจฺจ วตฺถุ, วตฺถุํ ปฏิจฺจ ขนฺธา.\csromanspacer{} เอกํ มหาภูตํ ฯเปฯ. (1)}

\prose{นเหตุํ ธมฺมํ ปฏิจฺจ เหตุ ธมฺโม อุปฺปชฺชติ เหตุปจฺจยา.\csromanspacer{} นเหตู ขนฺเธ ปฏิจฺจ เหตู.\csromanspacer{} ปฏิสนฺธิกฺขเณ ฯเปฯ.\csromanspacer{} วตฺถุํ ปฏิจฺจ เหตู.}

\prose{นเหตุํ ธมฺมํ ปฏิจฺจ เหตุ จ นเหตุ จ ธมฺมา อุปฺปชฺชนฺติ เหตุปจฺจยา.\csromanspacer{} นเหตุํ เอกํ ขนฺธํ ปฏิจฺจ ตโย ขนฺธา เหตุ จ จิตฺตสมุฏฺฐานญฺจ รูปํ ฯเปฯ เทฺว ขนฺเธ ปฏิจฺจ เทฺว ขนฺธา เหตุ จ จิตฺตสมุฏฺฐานญฺจ รูปํ.\csromanspacer{} ปฏิสนฺธิกฺขเณ ฯเปฯ.\csromanspacer{} ปฏิสนฺธิกฺขเณ วตฺถุํ ปฏิจฺจ เหตู สมฺปยุตฺตกา จ ขนฺธา. (3)}

\proseitem{3}{เหตุญฺจ นเหตุญฺจ ธมฺมํ ปฏิจฺจ เหตุ ธมฺโม อุปฺปชฺชติ เหตุปจฺจยา.\csromanspacer{} อโลภญฺจ สมฺปยุตฺตเก จ ขนฺเธ ปฏิจฺจ อโทโส อโมโห. (จกฺกํ.) โลภญฺจ สมฺปยุตฺตเก จ ขนฺเธ ปฏิจฺจ โมโห.\csromanspacer{} โทสญฺจ สมฺปยุตฺตเก จ ขนฺเธ ปฏิจฺจ โมโห ฯเปฯ.\csromanspacer{} ปฏิสนฺธิกฺขเณ ฯเปฯ.\csromanspacer{} อโลภญฺจ วตฺถุญฺจ ปฏิจฺจ อโทโส อโมโห ฯเปฯ. (1)}

\prose{เหตุญฺจ นเหตุญฺจ ธมฺมํ ปฏิจฺจ นเหตุ ธมฺโม อุปฺปชฺชติ เหตุปจฺจยา.\csromanspacer{} นเหตุํ เอกํ ขนฺธญฺจ เหตุญฺจ ปฏิจฺจ ตโย ขนฺธา จิตฺตสมุฏฺฐานญฺจ รูปํ ฯเปฯ เทฺว ขนฺเธ จ เหตุญฺจ ปฏิจฺจ เทฺว ขนฺธา จิตฺตสมุฏฺฐานญฺจ รูปํ.\csromanspacer{} ปฏิสนฺธิกฺขเณ ฯเปฯ.\csromanspacer{} ปฏิสนฺธิกฺขเณ วตฺถุญฺจ เหตุญฺจ ปฏิจฺจ สมฺปยุตฺตกา ขนฺธา. (2) ปปฺ}

\prose{เหตุญฺจ นเหตุญฺจ ธมฺมํ ปฏิจฺจ เหตุ จ นเหตุ จ ธมฺมา อุปฺปชฺชนฺติ เหตุปจฺจยา.\csromanspacer{} นเหตุํ เอกํ ขนฺธญฺจ อโลภญฺจ ปฏิจฺจ ตโย ขนฺธา อโทโส อโมโห จ จิตฺตสมุฏฺฐานญฺจ รูปํ ฯเปฯ เทฺว ขนฺเธ จ อโลภญฺจ ปฏิจฺจ เทฺว ขนฺธา อโทโส อโมโห จ จิตฺตสมุฏฺฐานญฺจ รูปํ. (จกฺกํ.) นเหตุํ เอกํ ขนฺธญฺจ โลภญฺจ ปฏิจฺจ ตโย ขนฺธา โมโห จ จิตฺตสมุฏฺฐานญฺจ รูปํ ฯเปฯ เทฺว ขนฺเธ ฯเปฯ.\csromanspacer{} ปฏิสนฺธิกฺขเณ ฯเปฯ.\csromanspacer{} วตฺถุญฺจ อโลภญฺจ ปฏิจฺจ อโทโส อโมโห สมฺปยุตฺตกา จ ขนฺธา. (3)}

\csromanheader{2}{\textbf{อารมฺมณาทิ}}
\proseitem{4}{\csromanfolio{3}เหตุํ ธมฺมํ ปฏิจฺจ เหตุ ธมฺโม อุปฺปชฺชติ อารมฺมณปจฺจยา. (รูปํ ฉฑฺเฑตฺวา อรูเปเยว นว ปญฺหา.) อธิปติปจฺจยา. (ปฏิสนฺธิ นตฺถิ.\csromanspacer{} ปริปุณฺณํ.) เอกํ มหาภูตํ ปฏิจฺจ ฯเปฯ มหาภูเต ปฏิจฺจ จิตฺตสมุฏฺฐานํ รูปํ อุปาทารูปํ. (อิมํ นานํ.) อนนฺตรปจฺจยา.\csromanspacer{} สมนนฺตรปจฺจยา.\csromanspacer{} สหชาตปจฺจยา. (สพฺเพ มหาภูตา, ยาว อสญฺญสตฺตา.) อญฺญมญฺญปจฺจยา.\csromanspacer{} นิสฺสยปจฺจยา.\csromanspacer{} อุปนิสฺสยปจฺจยา.\csromanspacer{} ปุเรชาตปจฺจยา.\csromanspacer{} อาเสวนปจฺจยา. (ทฺวีสุปิ ปฏิสนฺธิ นตฺถิ.) กมฺมปจฺจยา.\csromanspacer{} วิปากปจฺจยา. (สํขิตฺตํ.) อวิคตปจฺจยา.}

\csromanheader{2}{1. ปจฺจยานุโลม 2. สงฺขฺยาวาร}
\proseitem{5}{เหตุยา นว, อารมฺมเณ นว, (สพฺพตฺถ นว.) อวิคเต นว. (เอวํ คเณตพฺพํ.)}

\vspace{\tipitakaparskip}
\csromancenter{อนุโลมํ.}
\csromanheader{2}{2. ปจฺจยปจฺจนีย 1. วิภงฺควาร}
\csromanheader{2}{\textbf{นเหตุ}}
\vspace{\tipitakaparskip}
\csromancenter{\textbf{นเหตุ}ํ ธมฺมํ ปฏิจฺจ นเหตุ ธมฺโม อุปฺปชฺชติ นเหตุปจฺจยา.\csromanspacer{} อเหตุกํ นเหตุํ เอกํ ขนฺธํ ปฏิจฺจ ตโย ขนฺธา จิตฺตสมุฏฺฐานญฺจ รูปํ ฯเปฯ เทฺว ขนฺเธ ฯเปฯ.\csromanspacer{} อเหตุกปฏิสนฺธิกฺขเณ ฯเปฯ.\csromanspacer{} ขนฺเธ ปฏิจฺจ วตฺถุ, วตฺถุํ ปฏิจฺจ ขนฺธา.\csromanspacer{} เอกํ มหาภูตํ ฯเปฯ.\csromanspacer{} พาหิรํ.\csromanspacer{} อาหารสมุฏฺฐานํ.\csromanspacer{} อุตุสมุฏฺฐานํ.\csromanspacer{} อสญฺญสตฺตานํ ฯเปฯ. (1)}
\csromancenter{\textbf{นเหตุ}ํ ธมฺมํ ปฏิจฺจ เหตุ ธมฺโม อุปฺปชฺชติ นเหตุปจฺจยา.\csromanspacer{} วิจิกิจฺฉาสหคเต อุทฺธจฺจสหคเต ขนฺเธ ปฏิจฺจ วิจิกิจฺฉาสหคโต อุทฺธจฺจสหคโต โมโห. (2)}
\csromanheader{2}{\textbf{น-อารมฺมณาทิ}}
\proseitem{7}{เหตุํ ธมฺมํ ปฏิจฺจ นเหตุ ธมฺโม อุปฺปชฺชติ น-อารมฺมณปจฺจยา.\csromanspacer{} เหตุํ ปฏิจฺจ จิตฺตสมุฏฺฐานํ รูปํ.\csromanspacer{} ปฏิสนฺธิกฺขเณ ฯเปฯ. (1)}

\prose{\csromanfolio{4}นเหตุํ ธมฺมํ ปฏิจฺจ นเหตุ ธมฺโม อุปฺปชฺชติ น-อารมฺมณปจฺจยา.\csromanspacer{} นเหตู ขนฺเธ ปฏิจฺจ จิตฺตสมุฏฺฐานํ รูปํ.\csromanspacer{} ปฏิสนฺธิกฺขเณ ฯเปฯ.\csromanspacer{} สพฺเพ มหาภูตา ฯเปฯ. (1)}

\prose{เหตุญฺจ นเหตุญฺจ ธมฺมํ ปฏิจฺจ นเหตุ ธมฺโม อุปฺปชฺชติ น-อารมฺมณปจฺจยา.\csromanspacer{} เหตุญฺจ นเหตุญฺจ ขนฺเธ ปฏิจฺจ จิตฺตสมุฏฺฐานํ รูปํ.\csromanspacer{} ปฏิสนฺธิกฺขเณ ฯเปฯ น-อธิปติปจฺจยา. (ปริปุณฺณํ.) น-อนนฺตรปจฺจยา.\csromanspacer{} นสมนนฺตรปจฺจยา.\csromanspacer{} น-อญฺญมญฺญปจฺจยา.\csromanspacer{} น-อุปนิสฺสยปจฺจยา.}

\csromanheader{2}{\textbf{นปุเรชาต}}
\proseitem{8}{เหตุํ ธมฺมํ ปฏิจฺจ เหตุ ธมฺโม อุปฺปชฺชติ นปุเรชาตปจฺจยา.\csromanspacer{} อรูเป อโลภํ ปฏิจฺจ อโทโส อโมโห. (จกฺกํ.) โลภํ ปฏิจฺจ โมโห, โมหํ ปฏิจฺจ โลโภ.\csromanspacer{} ปฏิสนฺธิกฺขเณ ฯเปฯ. (1)}

\prose{เหตุํ ธมฺมํ ปฏิจฺจ นเหตุ ธมฺโม อุปฺปชฺชติ นปุเรชาตปจฺจยา.\csromanspacer{} อรูเป เหตุํ ปฏิจฺจ สมฺปยุตฺตกา ขนฺธา.\csromanspacer{} เหตุํ ปฏิจฺจ จิตฺตสมุฏฺฐานํ รูปํ.\csromanspacer{} ปฏิสนฺธิกฺขเณ ฯเปฯ. (2)}

\prose{เหตุํ ธมฺมํ ปฏิจฺจ เหตุ จ นเหตุ จ ธมฺมา อุปฺปชฺชนฺติ นปุเรชาตปจฺจยา.\csromanspacer{} อรูเป อโลภํ ปฏิจฺจ อโทโส อโมโห สมฺปยุตฺตกา จ ขนฺธา. (จกฺกํ.) โลภํ ปฏิจฺจ โมโห สมฺปยุตฺตกา จ ขนฺธา. (จกฺกํ.) ปฏิสนฺธิกฺขเณ ฯเปฯ. (3)}

\proseitem{9}{นเหตุํ ธมฺมํ ปฏิจฺจ นเหตุ ธมฺโม อุปฺปชฺชติ นปุเรชาตปจฺจยา.\csromanspacer{} อรูเป นเหตุํ เอกํ ขนฺธํ ปฏิจฺจ ตโย ขนฺธา ฯเปฯ.\csromanspacer{} นเหตู ขนฺเธ ปฏิจฺจ จิตฺตสมุฏฺฐานํ รูปํ.\csromanspacer{} ปฏิสนฺธิกฺขเณ ฯเปฯ.\csromanspacer{} เอกํ มหาภูตํ ฯเปฯ. (1)}

\prose{นเหตุํ ธมฺมํ ปฏิจฺจ เหตุ ธมฺโม อุปฺปชฺชติ นปุเรชาตปจฺจยา.\csromanspacer{} อรูเป นเหตู ขนฺเธ ปฏิจฺจ เหตู.\csromanspacer{} ปฏิสนฺธิกฺขเณ ฯเปฯ. (2)}

\prose{นเหตุํ ธมฺมํ ปฏิจฺจ เหตุ จ นเหตุ จ ธมฺมา อุปฺปชฺชนฺติ นปุเรชาตปจฺจยา.\csromanspacer{} อรูเป นเหตุํ เอกํ ขนฺธํ ปฏิจฺจ ตโย ขนฺธา เหตุ จ ฯเปฯ เทฺว ขนฺเธ ฯเปฯ.\csromanspacer{} ปฏิสนฺธิกฺขเณ ฯเปฯ. (3)}

\proseitem{10}{เหตุญฺจ นเหตุญฺจ ธมฺมํ ปฏิจฺจ เหตุ ธมฺโม อุปฺปชฺชติ นปุเรชาตปจฺจยา.\csromanspacer{} อรูเป อโลภญฺจ สมฺปยุตฺตเก จ ขนฺเธ ปฏิจฺจ อโทโส \csromanfolio{5}อโมโห. (จกฺกํ.) อรูเป โลภญฺจ สมฺปยุตฺตเก จ ขนฺเธ ปฏิจฺจ โมโห. (จกฺกํ.) ปฏิสนฺธิกฺขเณ. (1)}

\prose{เหตุญฺจ นเหตุญฺจ ธมฺมํ ปฏิจฺจ นเหตุ ธมฺโม อุปฺปชฺชติ นปุเรชาตปจฺจยา.\csromanspacer{} อรูเป นเหตุํ เอกํ ขนฺธญฺจ เหตุญฺจ ปฏิจฺจ ตโย ขนฺธา ฯเปฯ เทฺว ขนฺเธ ฯเปฯ.\csromanspacer{} นเหตู ขนฺเธ จ เหตุญฺจ ปฏิจฺจ จิตฺตสมุฏฺฐานํ รูปํ.\csromanspacer{} เหตุญฺจ มหาภูเต จ ปฏิจฺจ จิตฺตสมุฏฺฐานํ รูปํ.\csromanspacer{} ปฏิสนฺธิกฺขเณ ฯเปฯ. (2)}

\prose{เหตุญฺจ นเหตุญฺจ ธมฺมํ ปฏิจฺจ เหตุ จ นเหตุ จ ธมฺมา อุปฺปชฺชนฺติ นปุเรชาตปจฺจยา.\csromanspacer{} อรูเป นเหตุํ เอกํ ขนฺธญฺจ อโลภญฺจ ปฏิจฺจ ตโย ขนฺธา อโทโส อโมโห จ ฯเปฯ เทฺว ขนฺเธ ฯเปฯ. (จกฺกํ.) นเหตุํ เอกํ ขนฺธญฺจ โลภญฺจ ปฏิจฺจ ตโย ขนฺธา โมโห จ. (จกฺกํ.) ปฏิสนฺธิกฺขเณ ฯเปฯ. (3)}

\csromanheader{2}{\textbf{นปจฺฉาชาตาทิ}}
\proseitem{11}{เหตุํ ธมฺมํ ปฏิจฺจ เหตุ ธมฺโม อุปฺปชฺชติ นปจฺฉาชาตปจฺจยา.\csromanspacer{} น-อาเสวนปจฺจยา.}

\csromanheader{2}{\textbf{นกมฺมาทิ}}
\proseitem{12}{เหตุํ ธมฺมํ ปฏิจฺจ นเหตุ ธมฺโม อุปฺปชฺชติ นกมฺมปจฺจยา.\csromanspacer{} เหตุํ ปฏิจฺจ สมฺปยุตฺตกา เจตนา. (1)}

\prose{นเหตุํ ธมฺมํ ปฏิจฺจ นเหตุ ธมฺโม อุปฺปชฺชติ นกมฺมปจฺจยา.\csromanspacer{} นเหตู ขนฺเธ ปฏิจฺจ สมฺปยุตฺตกา เจตนา.\csromanspacer{} พาหิรํ.\csromanspacer{} อาหารสมุฏฺฐานํ.\csromanspacer{} อุตุสมุฏฺฐานํ ฯเปฯ. (1)}

\prose{เหตุญฺจ นเหตุญฺจ ธมฺมํ ปฏิจฺจ นเหตุ ธมฺโม อุปฺปชฺชติ นกมฺมปจฺจยา.\csromanspacer{} เหตุญฺจ สมฺปยุตฺตเก จ ขนฺเธ ปฏิจฺจ สมฺปยุตฺตกา เจตนา. (1)}

\prose{เหตุํ ธมฺมํ ปฏิจฺจ เหตุ ธมฺโม อุปฺปชฺชติ นวิปากปจฺจยา.\csromanspacer{} นว.}

\csromanheader{2}{\textbf{น-อาหาราทิ}}
\proseitem{13}{นเหตุํ ธมฺมํ ปฏิจฺจ นเหตุ ธมฺโม อุปฺปชฺชติ น-อาหารปจฺจยา.\csromanspacer{} พาหิรํ.\csromanspacer{} อุตุสมุฏฺฐานํ.\csromanspacer{} อสญฺญสตฺตานํ ฯเปฯ น-อินฺทฺริยปจฺจยา.\csromanspacer{} พาหิรํ. \csromanfolio{6}อาหารสมุฏฺฐานํ.\csromanspacer{} อุตุสมุฏฺฐานํ ฯเปฯ.\csromanspacer{} อสญฺญสตฺตานํ มหาภูเต ปฏิจฺจ รูปชีวิตินฺทฺริยํ.\csromanspacer{} นฌานปจฺจยา.\csromanspacer{} ปญฺจวิญฺญาณํ ฯเปฯ.\csromanspacer{} พาหิรํ.\csromanspacer{} อาหารสมุฏฺฐานํ.\csromanspacer{} อุตุสมุฏฺฐานํ.\csromanspacer{} อสญฺญสตฺตานํ ฯเปฯ น- มคฺคปจฺจยา.\csromanspacer{} อเหตุกํ น-เหตุํ เอกํ ขนฺธํ ปฏิจฺจ ฯเปฯ.\csromanspacer{} พาหิรํ.\csromanspacer{} อาหารสมุฏฺฐานํ.\csromanspacer{} อุตุสมุฏฺฐานํ.\csromanspacer{} อสญฺญสตฺตานํ ฯเปฯ น สมฺปยุตฺตปจฺจยา.\csromanspacer{} นวิปฺปยุตฺตปจฺจยา. (นปุเรชาตสทิสํ.\csromanspacer{} อรูปปญฺหาเยว.) โนนตฺถิปจฺจยา.\csromanspacer{} โนวิคตปจฺจยา.}

\csromanheader{2}{2. ปจฺจยปจฺจนีย 2. สงฺขฺยาวาร}
\csromanheader{2}{\textbf{สุทฺธ}}
\proseitem{14}{นเหตุยา เทฺว, น-อารมฺมเณ ตีณิ, น-อธิปติยา นว, น-อนนฺตเร ตีณิ, นสมนนฺตเร ตีณิ, น-อญฺญมญฺเญ ตีณิ, น-อุปนิสฺสเย ตีณิ, นปุเรชาเต นว, นปจฺฉาชาเต นว, น-อาเสวเน นว, นกมฺเม ตีณิ, นวิปาเก นว, น-อาหาเร เอกํ, น-อินฺทฺริเย เอกํ, นฌาเน เอกํ, นมคฺเค เอกํ, นสมฺปยุตฺเต ตีณิ, นวิปฺปยุตฺเต นว, โน นตฺถิยา ตีณิ, โน วิคเต ตีณิ. (เอวํ คเณตพฺพํ.)}

\vspace{\tipitakaparskip}
\csromancenter{ปจฺจนียํ.}
\csromanheader{2}{3. ปจฺจยานุโลมปจฺจนีย}
\csromanheader{2}{\textbf{เหตุทุก}}
\proseitem{15}{เหตุปจฺจยา น-อารมฺมเณ ตีณิ, น-อธิปติยา นว, น-อนนฺตเร ตีณิ ฯเปฯ น-อุปนิสฺสเย ตีณิ, นปุเรชาเต นว, นปจฺฉาชาเต นว, น-อาเสวเน นว, นกมฺเม ตีณิ, นวิปาเก นว, นสมฺปยุตฺเต ตีณิ, นวิปฺปยุตฺเต นว, โนนตฺถิยา ตีณิ, โนวิคเต ตีณิ.}

\vspace{\tipitakaparskip}
\csromancenter{อนุโลมปจฺจนียํ.}
\csromanheader{2}{4. ปจฺจยปจฺจนียานุโลม}
\csromanheader{2}{\textbf{นเหตุทุก}}
\proseitem{16}{\csromanfolio{7}นเหตุปจฺจยา อารมฺมเณ เทฺว, อนนฺตเร เทฺว ฯเปฯ กมฺเม เทฺว, วิปาเก เอกํ, อาหาเร เทฺว, อินฺทฺริเย เทฺว, ฌาเน เทฺว, มคฺเค เอกํ, สมฺปยุตฺเต เทฺว ฯเปฯ อวิคเต เทฺว.\csromanspacer{} ปจฺจนียานุโลมํ.}

\prose{(สหชาตวาโรปิ ปจฺจยวาโรปิ นิสฺสยวาโรปิ ปฏิจฺจวารสทิสาเยว ปญฺหา.\csromanspacer{} มหาภูเตสุ นิฏฺฐิเตสุ “วตฺถุํ ปจฺจยา”ติ กาตพฺพา.\csromanspacer{} ปญฺจายตนานิ อนุโลเมปิ ปจฺจนีเยปิ ยถา ลพฺภนฺติ, ตถา กาตพฺพา.\csromanspacer{} สํสฏฺฐวาโรปิ สมฺปยุตฺตวาโรปิ ปริปุณฺโณ.\csromanspacer{} รูปํ นตฺถิ, อรูปเมว.)}

\csromanheader{2}{1. \textbf{เหตุ}ทุก 7. ปญฺหาวาร}
\csromanheader{2}{1. ปจฺจยานุโลม 1. วิภงฺควาร}
\csromanheader{2}{\textbf{เหตุ}}
\vspace{\tipitakaparskip}
\csromancenter{\textbf{เหตุ} ธมฺโม เหตุสฺส ธมฺมสฺส เหตุปจฺจเยน ปจฺจโย.\csromanspacer{} อโลโภ อโทสสฺส อโมหสฺส เหตุปจฺจเยน ปจฺจโย. (จกฺกํ.) โลโภ โมหสฺส เหตุปจฺจเยน ปจฺจโย.\csromanspacer{} โทโส โมหสฺส เหตุปจฺจเยน ปจฺจโย.\csromanspacer{} ปฏิสนฺธิกฺขเณ ฯเปฯ. (1)}
\csromancenter{\textbf{เหตุ} ธมฺโม นเหตุสฺส ธมฺมสฺส เหตุปจฺจเยน ปจฺจโย.\csromanspacer{} เหตู สมฺปยุตฺตกานํ ขนฺธานํ จิตฺตสมุฏฺฐานานญฺจ รูปานํ เหตุปจฺจเยน ปจฺจโย.\csromanspacer{} ปฏิสนฺธิกฺขเณ ฯเปฯ. (2)}
\csromancenter{\textbf{เหตุ} ธมฺโม เหตุสฺส จ นเหตุสฺส จ ธมฺมสฺส เหตุปจฺจเยน ปจฺจโย.\csromanspacer{} อโลโภ อโทสสฺส อโมหสฺส สมฺปยุตฺตกานญฺจ ขนฺธานํ จิตฺตสมุฏฺฐานานญฺจ รูปานํ เหตุปจฺจเยน ปจฺจโย. (จกฺกํ.) โลโภ โมหสฺส ฯเปฯ.\csromanspacer{} ปฏิสนฺธิกฺขเณ ฯเปฯ. (3)}
\csromanheader{2}{\textbf{อารมฺมณ}}
\proseitem{18}{\csromanfolio{8}เหตุ ธมฺโม เหตุสฺส ธมฺมสฺส อารมฺมณปจฺจเยน ปจฺจโย.\csromanspacer{} เหตุํ อารพฺภ เหตู อุปฺปชฺชนฺติ. (1)}

\prose{เหตุ ธมฺโม นเหตุสฺส ธมฺมสฺส อารมฺมณปจฺจเยน ปจฺจโย.\csromanspacer{} เหตุํ อารพฺภ นเหตู ขนฺธา อุปฺปชฺชนฺติ. (2)}

\prose{เหตุ ธมฺโม เหตุสฺส จ นเหตุสฺส จ ธมฺมสฺส อารมฺมณปจฺจเยน ปจฺจโย.\csromanspacer{} เหตุํ อารพฺภ เหตู จ สมฺปยุตฺตกา จ ขนฺธา อุปฺปชฺชนฺติ. (3)}

\proseitem{19}{นเหตุ ธมฺโม นเหตุสฺส ธมฺมสฺส อารมฺมณปจฺจเยน ปจฺจโย.\csromanspacer{} ทานํ ทตฺวา, สีลํ ฯเปฯ อุโปสถกมฺมํ กตฺวา ตํ ปจฺจเวกฺขติ.\csromanspacer{} ปุพฺเพ สุจิณฺณานิ ปจฺจเวกฺขติ.\csromanspacer{} ฌานา วุฏฺฐหิตฺวา ฯเปฯ.\csromanspacer{} อริยา มคฺคา วุฏฺฐหิตฺวา มคฺคํ ปจฺจเวกฺขนฺติ, ผลํ ฯเปฯ นิพฺพานํ ฯเปฯ.\csromanspacer{} นิพฺพานํ โคตฺรภุสฺส, โวทานสฺส มคฺคสฺส, ผลสฺส, อาวชฺชนาย อารมฺมณปจฺจเยน ปจฺจโย.\csromanspacer{} อริยา นเหตู ปหีเน กิเลเส ปจฺจเวกฺขนฺติ.\csromanspacer{} วิกฺขมฺภิ เต กิเลเส ปจฺจเวกฺขนฺติ.\csromanspacer{} ปุพฺเพ สมุทาจิณฺเณ กิเลเส ชานนฺติ.\csromanspacer{} จกฺขุํ ฯเปฯ วตฺถุํ, นเหตู ขนฺเธ อนิจฺจโต ฯเปฯ โทมนสฺสํ อุปฺปชฺชติ.\csromanspacer{} ทิพฺเพน จกฺขุนา รูปํ ปสฺสติ.\csromanspacer{} ทิพฺพาย โสตธาตุยา สทฺทํ สุณาติ.\csromanspacer{} เจโตปริยญาเณน นเหตุจิตฺตสมงฺคิสฺส จิตฺตํ ชานาติ.\csromanspacer{} อากาสานญฺจายตนํ\footnote{อากาสานญฺจายตนกิริสํ (สฺยา) เอวมุปริปิ ตีสุ ฐาเนสุ.} วิญฺญาณญฺจายตนสฺส ฯเปฯ.\csromanspacer{} อากิญฺจญฺญายตนํ เนวสญฺญานาสญฺญายตนสฺส ฯเปฯ.\csromanspacer{} รูปายตนํ จกฺขุวิญฺญาณสฺส ฯเปฯ.\csromanspacer{} โผฏฺฐพฺพายตนํ กายวิญฺญาณสฺส ฯเปฯ.\csromanspacer{} นเหตู ขนฺธา อิทฺธิวิธญาณสฺส, เจโตปริยญาณสฺส, ปุพฺเพนิวาสานุสฺสติญาณสฺส, ยถากมฺมูปคญาณสฺส, อนาคตํสญาณสฺส อาวชฺชนาย อารมฺมณปจฺจเยน ปจฺจโย. (1)}

\prose{นเหตุ ธมฺโม เหตุสฺส ธมฺมสฺส อารมฺมณปจฺจเยน ปจฺจโย.\csromanspacer{} ทานํ ทตฺวา. (ปฐมคมนํเยว.\csromanspacer{} อาวชฺชนา นตฺถิ. “รูปายตนํ จกฺขุวิญฺญาณสฺส ฯเปฯ.\csromanspacer{} โผฏฺฐพฺพายตนํ กายวิญฺญาณสฺสา”ติ อิทํ นตฺถิ.)}

\prose{นเหตุ ธมฺโม เหตุสฺส จ นเหตุสฺส จ ธมฺมสฺส อารมฺมณปจฺจเยน ปจฺจโย.\csromanspacer{} ทานํ ทตฺวา, สีลํ ฯเปฯ อุโปสถกมฺมํ กตฺวา ตํ ปจฺจเวกฺขติ. \csromanfolio{9}ตํ อารพฺภ เหตู จ สมฺปยุตฺตกา จ ขนฺธา อุปฺปชฺชนฺติ. (ตตฺถ ตตฺถ ฐิเตน อิมํ กาตพฺพํ.\csromanspacer{} ทุติยคมนสทิสํ.) (3)}

\proseitem{20}{เหตุ จ นเหตุ จ ธมฺมา เหตุสฺส ธมฺมสฺส อารมฺมณปจฺจเยน ปจฺจโย.\csromanspacer{} เหตุญฺจ สมฺปยุตฺตเก จ ขนฺเธ อารพฺภ เหตู อุปฺปชฺชนฺติ. (1)}

\prose{เหตุ จ นเหตุ จ ธมฺมา นเหตุสฺส ธมฺมสฺส อารมฺมณปจฺจเยน ปจฺจโย.\csromanspacer{} เหตุญฺจ สมฺปยุตฺตเก จ ขนฺเธ อารพฺภ นเหตู ขนฺธา อุปฺปชฺชนฺติ. (2)}

\prose{เหตุ จ นเหตุ จ ธมฺมา เหตุสฺส จ นเหตุสฺส จ ธมฺมสฺส อารมฺมณปจฺจเยน ปจฺจโย.\csromanspacer{} เหตุญฺจ สมฺปยุตฺตเก จ ขนฺเธ อารพฺภ เหตู จ สมฺปยุตฺตกา จ ขนฺธา อุปฺปชฺชนฺติ. (3)}

\csromanheader{2}{\textbf{อธิปติ}}
\proseitem{21}{เหตุ ธมฺโม เหตุสฺส ธมฺมสฺส อธิปติปจฺจเยน ปจฺจโย.\csromanspacer{} \textbf{อารมฺมณาธิปติ}, สหชาตาธิปติ.\csromanspacer{}\csromanspacer{}\textbf{อารมฺมณาธิปติ\csromandash{}}เหตุํ ครุํ กตฺวา เหตู อุปฺปชฺชนฺติ.\csromanspacer{} \textbf{สหชาตาธิปติ}\csromandash{}เหตุ อธิปติ สมฺปยุตฺตกานํ เหตูนํ อธิปติปจฺจเยน ปจฺจโย. (1)}

\prose{เหตุ ธมฺโม นเหตุสฺส ธมฺมสฺส อธิปติปจฺจเยน ปจฺจโย.\csromanspacer{} \textbf{อารมฺมณาธิปติ}, สหชาตาธิปติ.\csromanspacer{}\csromanspacer{}\textbf{อารมฺมณาธิปติ\csromandash{}}เหตุํ ครุํ กตฺวา นเหตู ขนฺธา อุปฺปชฺชนฺติ.\csromanspacer{} \textbf{สหชาตาธิปติ}\csromandash{}เหตุ อธิปติ สมฺปยุตฺตกานํ ขนฺธานํ จิตฺตสมุฏฺฐานานญฺจ รูปานํ อธิปติปจฺจเยน ปจฺจโย. (2)}

\prose{เหตุ ธมฺโม เหตุสฺส จ นเหตุสฺส จ ธมฺมสฺส อธิปติปจฺจเยน ปจฺจโย.\csromanspacer{} \textbf{อารมฺมณาธิปติ}, สหชาตาธิปติ.\csromanspacer{}\csromanspacer{}\textbf{อารมฺมณาธิปติ\csromandash{}}เหตุํ ครุํ กตฺวา เหตู จ สมฺปยุตฺตกา จ ขนฺธา อุปฺปชฺชนฺติ.\csromanspacer{} \textbf{สหชาตาธิปติ}\csromandash{} เหตุ อธิปติ สมฺปยุตฺตกานํ ขนฺธานํ เหตูนญฺจ จิตฺตสมุฏฺฐานานญฺจ รูปานํ อธิปติปจฺจเยน ปจฺจโย. (3)}

\proseitem{22}{นเหตุ ธมฺโม นเหตุสฺส ธมฺมสฺส อธิปติปจฺจเยน ปจฺจโย.\csromanspacer{} \textbf{อารมฺมณาธิปติ}, สหชาตาธิปติ.\csromanspacer{}\csromanspacer{}\textbf{อารมฺมณาธิปติ\csromandash{}}ทานํ ทตฺวา. (วิตฺถาเรตพฺพํ.\csromanspacer{} ยาว นเหตู ขนฺธา.) \textbf{สหชาตาธิปติ}\csromandash{}นเหตุ อธิปติ สมฺปยุตฺตกานํ ขนฺธานํ จิตฺตสมุฏฺฐานานญฺจ รูปานํ อธิปติปจฺจเยน ปจฺจโย. (1)}

\prose{\csromanfolio{10}นเหตุ ธมฺโม เหตุสฺส ธมฺมสฺส อธิปติปจฺจเยน ปจฺจโย.\csromanspacer{} \textbf{อารมฺมณาธิปติ}, สหชาตาธิปติ.\csromanspacer{}\csromanspacer{}\textbf{อารมฺมณาธิปติ\csromandash{}}ทานํ ทตฺวา. (สํขิตฺตํ.\csromanspacer{} ยาว วตฺถุ นเหตู จ ขนฺธา, ตาว กาตพฺพํ.) \textbf{สหชาตาธิปติ\csromandash{}}นเหตุ อธิปติ สมฺปยุตฺตกานํ เหตูนํ อธิปติปจฺจเยน ปจฺจโย. (2)}

\prose{นเหตุ ธมฺโม เหตุสฺส จ นเหตุสฺส จ ธมฺมสฺส อธิปติปจฺจเยน ปจฺจโย.\csromanspacer{} \textbf{อารมฺมณาธิปติ}, สหชาตาธิปติ.\csromanspacer{}\csromanspacer{}\textbf{อารมฺมณาธิปติ\csromandash{}}ทานํ ทตฺวา, สีลํ ฯเปฯ อุโปสถกมฺมํ กตฺวา ตํ ครุํ กตฺวา ปจฺจเวกฺขติ.\csromanspacer{} ตํ ครุํ กตฺวา นเหตู ขนฺธา จ เหตู จ อุปฺปชฺชนฺติ.\csromanspacer{} ปุพฺเพ สุจิณฺณานิ. (ยาว วตฺถุ นเหตู ขนฺธา จ, ตาว กาตพฺพํ.) \textbf{สหชาตาธิปติ\csromandash{}}นเหตุ อธิปติ สมฺปยุตฺตกานํ ขนฺธานํ เหตูนญฺจ จิตฺตสมุฏฺฐานานญฺจ รูปานํ อธิปติปจฺจเยน ปจฺจโย. (3)}

\proseitem{23}{เหตุ จ นเหตุ จ ธมฺมา เหตุสฺส ธมฺมสฺส อธิปติปจฺจเยน ปจฺจโย.\csromanspacer{} \textbf{อารมฺมณาธิปติ\csromandash{}}เหตุญฺจ สมฺปยุตฺตเก จ ขนฺเธ ครุํ กตฺวา เหตู อุปฺปชฺชนฺติ. (1)}

\prose{เหตุ จ นเหตุ จ ธมฺมา นเหตุสฺส ธมฺมสฺส อธิปติปจฺจเยน ปจฺจโย.\csromanspacer{} \textbf{อารมฺมณาธิปติ\csromandash{}}เหตุญฺจ สมฺปยุตฺตเก จ ขนฺเธ ครุํ กตฺวา นเหตู ขนฺธา อุปฺปชฺชนฺติ. (2)}

\prose{เหตุ จ นเหตุ จ ธมฺมา เหตุสฺส จ นเหตุสฺส จ ธมฺมสฺส อธิปติปจฺจเยน ปจฺจโย.\csromanspacer{} \textbf{อารมฺมณาธิปติ\csromandash{}}เหตุญฺจ สมฺปยุตฺตเก จ ขนฺเธ ครุํ กตฺวา เหตู จ สมฺปยุตฺตกา จ ขนฺธา อุปฺปชฺชนฺติ. (3)}

\csromanheader{2}{\textbf{อนนฺตร}}
\proseitem{24}{เหตุ ธมฺโม เหตุสฺส ธมฺมสฺส อนนฺตรปจฺจเยน ปจฺจโย.\csromanspacer{} ปุริมา ปุริมา เหตู ปจฺฉิมานํ ปจฺฉิมานํ เหตูนํ อนนฺตรปจฺจเยน ปจฺจโย. (1)}

\prose{เหตุ ธมฺโม นเหตุสฺส ธมฺมสฺส อนนฺตรปจฺจเยน ปจฺจโย.\csromanspacer{} ปุริมา ปุริมา เหตู ปจฺฉิมานํ ปจฺฉิมานํ นเหตูนํ ขนฺธานํ อนนฺตรปจฺจเยน ปจฺจโย. (2)}

\prose{เหตุ ธมฺโม เหตุสฺส จ นเหตุสฺส จ ธมฺมสฺส อนนฺตรปจฺจเยน ปจฺจโย.\csromanspacer{} ปุริมา ปุริมา เหตู ปจฺฉิมานํ ปจฺฉิมานํ เหตูนํ สมฺปยุตฺตกานญฺจ ขนฺธานํ อนนฺตรปจฺจเยน ปจฺจโย. (3)}

\proseitem{25}{\csromanfolio{11}นเหตุ ธมฺโม นเหตุสฺส ธมฺมสฺส อนนฺตรปจฺจเยน ปจฺจโย.\csromanspacer{} ปุริมา ปุริมา นเหตู ขนฺธา ปจฺฉิมานํ ปจฺฉิมานํ นเหตูนํ ขนฺธานํ อนนฺตรปจฺจเยน ปจฺจโย.\csromanspacer{} อนุโลมํ โคตฺรภุสฺส. (สํขิตฺตํ.) เนวสญฺญานาสญฺญายตนํ ผลสมาปตฺติยา อนนฺตรปจฺจเยน ปจฺจโย. (1)}

\prose{นเหตุ ธมฺโม เหตุสฺส ธมฺมสฺส อนนฺตรปจฺจเยน ปจฺจโย ฯเปฯ.}

\prose{นเหตุ ธมฺโม เหตุสฺส จ นเหตุสฺส จ ธมฺมสฺส อนนฺตรปจฺจเยน ปจฺจโย. (นเหตุมูลกํ ตีณิปิ เอกสทิสํ.) (3)}

\proseitem{26}{เหตู จ นเหตู จ ธมฺมา เหตุสฺส ธมฺมสฺส อนนฺตรปจฺจเยน ปจฺจโย.\csromanspacer{} ปุริมา ปุริมา เหตู จ สมฺปยุตฺตกา จ ขนฺธา ปจฺฉิมานํ ปจฺฉิมานํ เหตูนํ อนนฺตรปจฺจเยน ปจฺจโย. (1)}

\prose{เหตู จ นเหตู จ ธมฺมา นเหตุสฺส ธมฺมสฺส อนนฺตรปจฺจเยน ปจฺจโย.\csromanspacer{} ปุริมา ปุริมา เหตู จ สมฺปยุตฺตกา จ ขนฺธา ปจฺฉิมานํ ปจฺฉิมานํ นเหตูนํ ขนฺธานํ อนนฺตรปจฺจเยน ปจฺจโย. (2)}

\prose{เหตู จ นเหตู จ ธมฺมา เหตุสฺส จ นเหตุสฺส จ ธมฺมสฺส อนนฺตรปจฺจเยน ปจฺจโย.\csromanspacer{} ปุริมา ปุริมา เหตู จ สมฺปยุตฺตกา จ ขนฺธา ปจฺฉิมานํ ปจฺฉิมานํ เหตูนํ สมฺปยุตฺตกานญฺจ ขนฺธานํ อนนฺตรปจฺจเยน ปจฺจโย. (3)}

\csromanheader{2}{\textbf{สมนนฺตราทิ}}
\proseitem{27}{เหตุ ธมฺโม เหตุสฺส ธมฺมสฺส สมนนฺตรปจฺจเยน ปจฺจโย. (อนนฺตรสทิสํ.) สหชาตปจฺจเยน ปจฺจโย.\csromanspacer{} อญฺญมญฺญปจฺจเยน ปจฺจโย. (อิเม เทฺวปิ ปฏิจฺจสทิสา.\csromanspacer{} นิสฺสยปจฺจโย ปจฺจยวาเร นิสฺสยปจฺจยสทิโส.)}

\csromanheader{2}{\textbf{อุปนิสฺสย}}
\proseitem{28}{เหตุ ธมฺโม เหตุสฺส ธมฺมสฺส อุปนิสฺสยปจฺจเยน ปจฺจโย.\csromanspacer{} อารมฺมณูปนิสฺสโย, อนนฺตรูปนิสฺสโย, ปกตูปนิสฺสโย ฯเปฯ.\csromanspacer{} ปกตูปนิสฺสโย\csromandash{}เหตู เหตูนํ อุปนิสฺสยปจฺจเยน ปจฺจโย. (1)}

\prose{เหตุ ธมฺโม นเหตุสฺส ธมฺมสฺส อุปนิสฺสยปจฺจเยน ปจฺจโย.\csromanspacer{} อารมฺมณูปนิสฺสโย, อนนฺตรูปนิสฺสโย, ปกตูปนิสฺสโย ฯเปฯ.\csromanspacer{} ปกตูปนิสฺสโย\csromandash{}เหตู นเหตูนํ ขนฺธานํ อุปนิสฺสยปจฺจเยน ปจฺจโย. (2)}

\prose{\csromanfolio{12}เหตุ ธมฺโม เหตุสฺส จ นเหตุสฺส จ ธมฺมสฺส อุปนิสฺสยปจฺจเยน ปจฺจโย.\csromanspacer{} อารมฺมณูปนิสฺสโย, อนนฺตรูปนิสฺสโย, ปกตูปนิสฺสโย ฯเปฯ.\csromanspacer{} ปกตูปนิสฺสโย\csromandash{}เหตู เหตูนํ สมฺปยุตฺตกานญฺจ ขนฺธานํ อุปนิสฺสยปจฺจเยน ปจฺจโย. (3)}

\proseitem{29}{นเหตุ ธมฺโม นเหตุสฺส ธมฺมสฺส อุปนิสฺสยปจฺจเยน ปจฺจโย.\csromanspacer{} อารมฺมณูปนิสฺสโย, อนนฺตรูปนิสฺสโย, ปกตูปนิสฺสโย ฯเปฯ.\csromanspacer{} ปกตูปนิสฺสโย\csromandash{}สทฺธํ อุปนิสฺสาย ทานํ เทติ ฯเปฯ สมาปตฺติํ อุปฺปาเทติ ฯเปฯ ทิฏฺฐิํ คณฺหาติ.\csromanspacer{} สีลํ ฯเปฯ เสนาสนํ อุปนิสฺสาย ทานํ เทติ ฯเปฯ สํฆํ ภินฺทติ.\csromanspacer{} สทฺธา ฯเปฯ เสนาสนํ สทฺธาย ฯเปฯ ผลสมาปตฺถิยา อุปนิสฺสยปจฺจเยน ปจฺจโย. (1)}

\prose{นเหตุ ธมฺโม เหตุสฺส ธมฺมสฺส อุปนิสฺสยปจฺจเยน ปจฺจโย.\csromanspacer{} อารมฺมณูปนิสฺสโย, อนนฺตรูปนิสฺสโย, ปกตูปนิสฺสโย ฯเปฯ.\csromanspacer{} ปกตูปนิสฺสโย\csromandash{}สทฺธํ ฯเปฯ เสนาสนํ อุปนิสฺสาย ทานํ เทติ ฯเปฯ สํฆํ ภินฺทติ.\csromanspacer{} สทฺธา ฯเปฯ เสนาสนํ สทฺธาย ฯเปฯ ปตฺถนาย มคฺคสฺส ผลสมาปตฺติยา อุปนิสฺสยปจฺจเยน ปจฺจโย. (2)}

\prose{นเหตุ ธมฺโม เหตุสฺส จ นเหตุสฺส จ ธมฺมสฺส อุปนิสฺสยปจฺจเยน ปจฺจโย.\csromanspacer{} อารมฺมณูปนิสฺสโย, อนนฺตรูปนิสฺสโย, ปกตูปนิสฺสโย ฯเปฯ.\csromanspacer{} ปกตูปนิสฺสโย. (ทุติย-อุปนิสฺสยสทิสํ.) (3)}

\proseitem{30}{เหตู จ นเหตู จ ธมฺมา เหตุสฺส ธมฺมสฺส อุปนิสฺสยปจฺจเยน ปจฺจโย.\csromanspacer{} อารมฺมณูปนิสฺสโย, อนนฺตรูปนิสฺสโย, ปกตูปนิสฺสโย ฯเปฯ.\csromanspacer{} ปกตูปนิสฺสโย\csromandash{}เหตู จ สมฺปยุตฺตกา จ ขนฺธา เหตูนํ อุปนิสฺสยปจฺจเยน ปจฺจโย. (1)}

\prose{เหตู จ นเหตู จ ธมฺมา นเหตุสฺส ธมฺมสฺส อุปนิสฺสยปจฺจเยน ปจฺจโย.\csromanspacer{} อารมฺมณูปนิสฺสโย, อนนฺตรูปนิสฺสโย, ปกตูปนิสฺสโย ฯเปฯ.\csromanspacer{} ปกตูปนิสฺสโย\csromandash{}เหตู จ สมฺปยุตฺตกา จ ขนฺธา เนเหตูนํ ขนฺธานํ อุปนิสฺสยปจฺจเยน ปจฺจโย. (2)}

\prose{เหตู จ นเหตู จ ธมฺมา เหตุสฺส จ นเหตุสฺส จ ธมฺมสฺส อุปนิสฺสยปจฺจเยน ปจฺจโย.\csromanspacer{} อารมฺมณูปนิสฺสโย, อนนฺตรูปนิสฺสโย, \csromanfolio{13}ปกตูปนิสฺสโย ฯเปฯ.\csromanspacer{} ปกตูปนิสฺสโย\csromandash{}เหตู จ สมฺปยุตฺตกา จ ขนฺธา เหตูนญฺจ สมฺปยุตฺตกานญฺจ ขนฺธานํ อุปนิสฺสยปจฺจเยน ปจฺจโย. (3)}

\csromanheader{2}{\textbf{ปุเรชาต}}
\proseitem{31}{นเหตุ ธมฺโม นเหตุสฺส ธมฺมสฺส ปุเรชาตปจฺจเยน ปจฺจโย.\csromanspacer{} \textbf{อารมฺมณปุเรชาตํ}, วตฺถุปุเรชาตํ.\csromanspacer{}\csromanspacer{}\textbf{อารมฺมณปุเรชาตํ\csromandash{}}จกฺขุํ ฯเปฯ วตฺถุํ อนิจฺจโต ฯเปฯ โทมนสฺสํ อุปฺปชฺชติ.\csromanspacer{} ทิพฺเพน จกฺขุนา รูปํ ปสฺสติ.\csromanspacer{} ทิพฺพาย โสตธาตุยา สทฺทํ สุณาติ.\csromanspacer{} รูปายตนํ จกฺขุวิญฺญาณสฺส ฯเปฯ.\csromanspacer{} โผฏฺฐพฺพายตนํ กายวิญฺญาณสฺส ฯเปฯ.\csromanspacer{} \textbf{วตฺถุปุเรชาตํ}\csromandash{}จกฺขายตนํ ฯเปฯ กายายตนํ ฯเปฯ.\csromanspacer{} วตฺถุ นเหตูนํ ขนฺธานํ ปุเรชาตปจฺจเยน ปจฺจโย. (1)}

\prose{นเหตุ ธมฺโม เหตุสฺส ธมฺมสฺส ปุเรชาตปจฺจเยน ปจฺจโย.\csromanspacer{} \textbf{อารมฺมณปุเรชาตํ}, วตฺถุปุเรชาตํ.\csromanspacer{}\csromanspacer{}\textbf{อารมฺมณปุเรชาตํ\csromandash{}}จกฺขุํ ฯเปฯ วตฺถุํ อนิจฺจโต ฯเปฯ โทมนสฺสํ อุปฺปชฺชติ.\csromanspacer{} ทิพฺเพน จกฺขุนา รูปํ ปสฺสติ.\csromanspacer{} ทิพฺพาย โสตธาตุยา สทฺทํ สุณาติ.\csromanspacer{} \textbf{วตฺถุปุเรชาตํ}\csromandash{}วตฺถุ เหตูนํ ปุเรชาตปจฺจเยน ปจฺจโย. (2)}

\prose{นเหตุ ธมฺโม เหตุสฺส จ นเหตุสฺส จ ธมฺมสฺส ปุเรชาตปจฺจเยน ปจฺจโย.\csromanspacer{} \textbf{อารมฺมณปุเรชาตํ}, วตฺถุปุเรชาตํ.\csromanspacer{}\csromanspacer{}\textbf{อารมฺมณปุเรชาตํ\csromandash{}} จกฺขุํ ฯเปฯ วตฺถุํ อนิจฺจโต ฯเปฯ โทมนสฺสํ อุปฺปชฺชติ.\csromanspacer{} \textbf{วตฺถุปุเรชาตํ}\csromandash{}วตฺถุ เหตูนํ สมฺปยุตฺตกานญฺจ ขนฺธานํ ปุเรชาตปจฺจเยน ปจฺจโย. (3)}

\csromanheader{2}{\textbf{ปจฺฉาชาตาทิ}}
\proseitem{32}{เหตุ ธมฺโม นเหตุสฺส ธมฺมสฺส ปจฺฉาชาตปจฺจเยน ปจฺจโย.\csromanspacer{} ปจฺฉาชาตา เหตู ปุเรชาตสฺส อิมสฺส กายสฺส ปจฺฉาชาตปจฺจเยน ปจฺจโย. (1)}

\prose{นเหตุ ธมฺโม นเหตุสฺส ธมฺมสฺส ปจฺฉาชาตปจฺจเยน ปจฺจโย.\csromanspacer{} ปจฺฉาชาตา นเหตู ขนฺธา ปุเรชาตสฺส อิมสฺส กายสฺส ปจฺฉาชาตปจฺจเยน ปจฺจโย. (1)}

\prose{เหตู จ นเหตู จ ธมฺมา นเหตุสฺส ธมฺมสฺส ปจฺฉาชาตปจฺจเยน.\csromanspacer{} ปจฺฉาชาตา เหตู จ สมฺปยุตฺตกา จ ขนฺธา ปุเรชาตสฺส อิมสฺส กายสฺส ปจฺฉาชาตปจฺจเยน ปจฺจโย. (1)}

\prose{\csromanfolio{14}เหตุ ธมฺโม เหตุสฺส ธมฺมสฺส อาเสวนปจฺจเยน ปจฺจโย. (อนนฺตรสทิสํ.)}

\csromanheader{2}{\textbf{กมฺม}}
\proseitem{33}{นเหตุ ธมฺโม นเหตุสฺส ธมฺมสฺส กมฺมปจฺจเยน ปจฺจโย.\csromanspacer{} \textbf{สหชาตา}, นานากฺขณิกา.\csromanspacer{}\csromanspacer{}\textbf{สหชาตา} นเหตุ เจตนา สมฺปยุตฺตกานํ ขนฺธานํ จิตฺตสมุฏฺฐานานญฺจ รูปานํ กมฺมปจฺจเยน ปจฺจโย.\csromanspacer{} ปฏิสนฺธิกฺขเณ ฯเปฯ.\csromanspacer{} \textbf{นานากฺขณิกา} นเหตุ เจตนา วิปากานํ ขนฺธานํ กฏตฺตา จ รูปานํ กมฺมปจฺจเยน ปจฺจโย. (1)}

\prose{นเหตุ ธมฺโม เหตุสฺส ธมฺมสฺส กมฺมปจฺจเยน ปจฺจโย.\csromanspacer{} \textbf{สหชาตา}, นานากฺขณิกา.\csromanspacer{}\csromanspacer{}\textbf{สหชาตา} นเหตุ เจตนา สมฺปยุตฺตกานํ เหตูนํ กมฺมปจฺจเยน ปจฺจโย.\csromanspacer{} ปฏิสนฺธิกฺขเณ ฯเปฯ.\csromanspacer{} \textbf{นานากฺขณิกา} นเหตุ เจตนา วิปากานํ เหตูนํ กมฺมปจฺจเยน ปจฺจโย. (2)}

\prose{นเหตุ ธมฺโม เหตุสฺส จ นเหตุสฺส จ ธมฺมสฺส กมฺมปจฺจเยน ปจฺจโย.\csromanspacer{} \textbf{สหชาตา}, นานากฺขณิกา.\csromanspacer{}\csromanspacer{}\textbf{สหชาตา} นเหตุ เจตนา สมฺปยุตฺตกานํ ขนฺธานํ เหตูนํ จิตฺตสมุฏฺฐานานญฺจ รูปานํ กมฺมปจฺจเยน ปจฺจโย.\csromanspacer{} ปฏิสนฺธิกฺขเณ ฯเปฯ.\csromanspacer{} \textbf{นานากฺขณิกา} นเหตุ เจตนา วิปากานํ ขนฺธานํ เหตูนํ กฏตฺตา จ รูปานํ กมฺมปจฺจเยน ปจฺจโย. (3)}

\csromanheader{2}{\textbf{วิปาก}}
\proseitem{34}{เหตุ ธมฺโม เหตุสฺส ธมฺมสฺส วิปากปจฺจเยน ปจฺจโย.\csromanspacer{} วิปาโก อโลโภ อโทสสฺส อโมหสฺส วิปากปจฺจเยน ปจฺจโย. (ปฏิจฺจวารสทิสํ.\csromanspacer{} \textbf{วิปาก}วิภงฺเค นว ปญฺหา.)}

\csromanheader{2}{\textbf{อาหาร}}
\proseitem{35}{นเหตุ ธมฺโม นเหตุสฺส ธมฺมสฺส อาหารปจฺจเยน ปจฺจโย.\csromanspacer{} นเหตู อาหารา สมฺปยุตฺตกานํ ขนฺธานํ จิตฺตสมุฏฺฐานานญฺจ รูปานํ อาหารปจฺจเยน ปจฺจโย.\csromanspacer{} ปฏิสนฺธิกฺขเณ ฯเปฯ.\csromanspacer{} กพฬีกาโร อาหาโร อิมสฺส กายสฺส อาหารปจฺจเยน ปจฺจโย. (1)}

\prose{นเหตุ ธมฺโม เหตุสฺส ธมฺมสฺส อาหารปจฺจเยน ปจฺจโย.\csromanspacer{} นเหตู อาหารา สมฺปยุตฺตกานํ เหตูนํ อาหารปจฺจเยน ปจฺจโย.\csromanspacer{} ปฏิสนฺธิกฺขเณ ฯเปฯ. (2)}

\prose{\csromanfolio{15}นเหตุ ธมฺโม เหตุสฺส จ น เหตุสฺส จ ธมฺมสฺส อาหารปจฺจเยน ปจฺจโย.\csromanspacer{} นเหตู อาหารา สมฺปยุตฺตกานํ ขนฺธานํ, เหตูนํ จิตฺตสมุฏฺฐานานญฺจ รูปานํ อาหารปจฺจเยน ปจฺจโย.\csromanspacer{} ปฏิสนฺธิกฺขเณ ฯเปฯ. (3)}

\csromanheader{2}{\textbf{อินฺทฺริย}}
\proseitem{36}{เหตุ ธมฺโม เหตุสฺส ธมฺมสฺส อินฺทฺริยปจฺจเยน ปจฺจโย ฯเปฯ. (เหตุมูลเก ตีณิ)}

\prose{นเหตุ ธมฺโม นเหตุสฺส ธมฺมสฺส อินฺทฺริยปจฺจเยน ปจฺจโย.\csromanspacer{} นเหตู อินฺทฺริยา สมฺปยุตฺตกานํ ขนฺธานํ จิตฺตสมุฏฺฐานานญฺจ รูปานํ อินฺทฺริยปจฺจเยน ปจฺจโย.\csromanspacer{} ปฏิสนฺธิกฺขเณ ฯเปฯ.\csromanspacer{} จกฺขุนฺทฺริยํ จกฺขุวิญฺญาณสฺส ฯเปฯ.\csromanspacer{} กายินฺทฺริยํ กายวิญฺญาณสฺส ฯเปฯ.\csromanspacer{} รูปชีวิตินฺทฺริยํ กฏตฺตารูปานํ อินฺทฺริยปจฺจเยน ปจฺจโย. (เอวํ อินฺทฺริยปจฺจยา วิตฺถาเรตพฺพา.\csromanspacer{} นว.)}

\csromanheader{2}{\textbf{ฌานาทิ}}
\proseitem{37}{นเหตุ ธมฺโม นเหตุสฺส ธมฺมสฺส ฌานปจฺจเยน ปจฺจโย.\csromanspacer{} ตีณิ.}

\prose{เหตุ ธมฺโม เหตุสฺส ธมฺมสฺส มคฺคปจฺจเยน ปจฺจโย.\csromanspacer{} สมฺปยุตฺตปจฺจเยน ปจฺจโย. (อิเมสุ ทฺวีสุ นว.)}

\csromanheader{2}{\textbf{วิปฺปยุตฺต}}
\proseitem{38}{เหตุ ธมฺโม นเหตุสฺส ธมฺมสฺส วิปฺปยุตฺตปจฺจเยน ปจฺจโย.\csromanspacer{} สหชาตํ, ปจฺฉาชาตํ.\csromanspacer{}\csromanspacer{}\textbf{สหชาตา} เหตู จิตฺตสมุฏฺฐานานํ รูปานํ วิปฺปยุตฺตปจฺจเยน ปจฺจโย.\csromanspacer{} ปฏิสนฺธิกฺขเณ เหตู กฏตฺตารูปานํ วิปฺปยุตฺตปจฺจเยน ปจฺจโย.\csromanspacer{} เหตู วตฺถุสฺส วิปฺปยุตฺตปจฺจเยน ปจฺจโย.\csromanspacer{} \textbf{ปจฺฉาชาตา} เหตู ปุเรชาตสฺส อิมสฺส กายสฺส วิปฺปยุตฺตปจฺจเยน ปจฺจโย. (1)}

\prose{นเหตุ ธมฺโม นเหตุสฺส ธมฺมสฺส วิปฺปยุตฺตปจฺจเยน ปจฺจโย.\csromanspacer{} สหชาตํ, ปุเรชาตํ, ปจฺฉาชาตํ.\csromanspacer{}\csromanspacer{}\textbf{สหชาตา} นเหตู ขนฺธา จิตฺตสมุฏฺฐานานํ รูปานํ วิปฺปยุตฺตปจฺจเยน ปจฺจโย.\csromanspacer{} ปฏิสนฺธิกฺขเณ นเหตู ขนฺธา กฏตฺตารูปานํ วิปฺปยุตฺตปจฺจเยน ปจฺจโย.\csromanspacer{} ขนฺธา วตฺถุสฺส ฯเปฯ วตฺถุ ขนฺธานํ วิปฺปยุตฺตปจฺจเยน ปจฺจโย.\csromanspacer{} \textbf{ปุเรชาตํ} จกฺขายตนํ จกฺขุวิญฺญาณสฺส ฯเปฯ.\csromanspacer{} กายายตนํ กายวิญฺญาณสฺส.\csromanspacer{} วตฺถุ นเหตูนํ ขนฺธานํ วิปฺปยุตฺตปจฺจเยน \csromanfolio{16}ปจฺจโย.\csromanspacer{} \textbf{ปจฺฉาชาตา} นเหตู ขนฺธา ปุเรชาตสฺส อิมสฺส กายสฺส วิปฺปยุตฺตปจฺจเยน ปจฺจโย. (1)}

\prose{นเหตุ ธมฺโม เหตุสฺส ธมฺมสฺส วิปฺปยุตฺตปจฺจเยน ปจฺจโย.\csromanspacer{} \textbf{สหชาตํ}, ปุเรชาตํ.\csromanspacer{}\csromanspacer{}\textbf{สหชาตํ} ปฏิสนฺธิกฺขเณ วตฺถุ เหตูนํ วิปฺปยุตฺตปจฺจเยน ปจฺจโย.\csromanspacer{} \textbf{ปุเรชาตํ} วตฺถุ เหตูนํ วิปฺปยุตฺตปจฺจเยน ปจฺจโย. (2)}

\prose{นเหตุ ธมฺโม เหตุสฺส จ นเหตุสฺส จ ธมฺมสฺส วิปฺปยุตฺตปจฺจเยน ปจฺจโย.\csromanspacer{} \textbf{สหชาตํ}, ปุเรชาตํ.\csromanspacer{}\csromanspacer{}\textbf{สหชาตํ} ปฏิสนฺธิกฺขเณ วตฺถุ เหตูนํ สมฺปยุตฺตกานญฺจ ขนฺธานํ วิปฺปยุตฺตปจฺจเยน ปจฺจโย.\csromanspacer{} \textbf{ปุเรชาตํ} วตฺถุ เหตูนํ สมฺปยุตฺตกานญฺจ ขนฺธานํ วิปฺปยุตฺตปจฺจเยน ปจฺจโย. (3)}

\prose{เหตู จ นเหตู จ ธมฺมา นเหตุสฺส ธมฺมสฺส วิปฺปยุตฺตปจฺจเยน ปจฺจโย.\csromanspacer{} \textbf{สหชาตํ}, ปจฺฉาชาตํ.\csromanspacer{}\csromanspacer{}\textbf{สหชาตา} เหตู จ สมฺปยุตฺตกา จ ขนฺธา จิตฺตสมุฏฺฐานานํ รูปานํ วิปฺปยุตฺตปจฺจเยน ปจฺจโย.\csromanspacer{} ปฏิสนฺธิกฺขเณ เหตู จ สมฺปยุตฺตกา จ ขนฺธา กฏตฺตารูปานํ วิปฺปยุตฺตปจฺจเยน ปจฺจโย.\csromanspacer{} \textbf{ปจฺฉาชาตา} เหตู จ สมฺปยุตฺตกา จ ขนฺธา ปุเรชาตสฺส อิมสฺส กายสฺส วิปฺปยุตฺตปจฺจเยน ปจฺจโย. (1)}

\csromanheader{2}{\textbf{อตฺถิ-อาทิ}}
\proseitem{39}{เหตุ ธมฺโม เหตุสฺส ธมฺมสฺส อตฺถิปจฺจเยน ปจฺจโย.\csromanspacer{} อโลโภ อโทสสฺส อโมหสฺส อตฺถิปจฺจเยน ปจฺจโย. (จกฺกํ.) โลโภ โมหสฺส อตฺถิปจฺจเยน ปจฺจโย. (จกฺกํ.) ปฏิสนฺธิกฺขเณ ฯเปฯ. (1)}

\prose{เหตุ ธมฺโม นเหตุสฺส ธมฺมสฺส อตฺถิปจฺจเยน ปจฺจโย.\csromanspacer{} \textbf{สหชาตํ}, ปจฺฉาชาตํ.\csromanspacer{}\csromanspacer{}\textbf{สหชาตา} เหตู สมฺปยุตฺตกานํ ขนฺธานํ จิตฺตสมุฏฺฐานานญฺจ รูปานํ อตฺถิปจฺจเยน ปจฺจโย.\csromanspacer{} ปฏิสนฺธิกฺขเณ ฯเปฯ.\csromanspacer{} \textbf{ปจฺฉาชาตา} เหตู ปุเรชาตสฺส อิมสฺส กายสฺส อตฺถิปจฺจเยน ปจฺจโย.}

\prose{เหตุ ธมฺโม เหตุสฺส จ นเหตุสฺส จ ธมฺมสฺส อตฺถิปจฺจเยน ปจฺจโย.\csromanspacer{} อโลโภ อโทสสฺส อโมหสฺส สมฺปยุตฺตกานํ ขนฺธานํ จิตฺตสมุฏฺฐานานญฺจ รูปานํ อตฺถิปจฺจเยน ปจฺจโย. (จกฺกํ) โลโภ โมหสฺส สมฺปยุตฺตกานํ ขนฺธานํ จิตฺตสมุฏฺฐานานญฺจ รูปานํ อตฺถิปจฺจเยน ปจฺจโย. (จกฺกํ.) ปฏิสนฺธิกฺขเณ ฯเปฯ. (3)}

\proseitem{40}{\csromanfolio{17}นเหตุ ธมฺโม นเหตุสฺส ธมฺมสฺส อตฺถิปจฺจเยน ปจฺจโย.\csromanspacer{} \textbf{สหชาตํ, ปุเรชาตํ}, ปจฺฉาชาตํ, อาหารํ, อินฺทฺริยํ.\csromanspacer{}\csromanspacer{}\textbf{สหชาโต} นเหตุ เอโก ขนฺโธ ติณฺณนฺนํ ขนฺธานํ จิตฺตสมุฏฺฐานานญฺจ รูปานํ อตฺถิปจฺจเยน ปจฺจโย ฯเปฯ เทฺว ขนฺธา ฯเปฯ.\csromanspacer{} ปฏิสนฺธิกฺขเณ นเหตุ เอโก ขนฺโธ ติณฺณนฺนํ ขนฺธานํ กฏตฺตา จ รูปานํ ฯเปฯ.\csromanspacer{} ขนฺธา วตฺถุสฺส อตฺถิปจฺจเยน ปจฺจโย.\csromanspacer{} วตฺถุ ขนฺธานํ อตฺถิปจฺจเยน ปจฺจโย.\csromanspacer{} เอกํ มหาภูตํ ฯเปฯ.\csromanspacer{} พาหิรํ.\csromanspacer{} อาหารสมุฏฺฐานํ.\csromanspacer{} อุตุสมุฏฺฐานํ.\csromanspacer{} อสญฺญสตฺตานํ ฯเปฯ.\csromanspacer{} \textbf{ปุเรชาตํ} จกฺขุํ ฯเปฯ.\csromanspacer{} วตฺถุํ ฯเปฯ.\csromanspacer{} ทิพฺเพน จกฺขุนา รูปํ ปสฺสติ.\csromanspacer{} ทิพฺพาย โสตธาตุยา สทฺทํ สุณาติ.\csromanspacer{} รูปายตนํ จกฺขุวิญฺญาณสฺส ฯเปฯ.\csromanspacer{} โผฏฺฐพฺพายตนํ กายวิญฺญาณสฺส.\csromanspacer{} จกฺขายตนํ จกฺขุวิญฺญาณสฺส ฯเปฯ.\csromanspacer{} กายายตนํ กายวิญฺญาณสฺส.\csromanspacer{} วตฺถุ นเหตูนํ ขนฺธานํ อตฺถิปจฺจเยน ปจฺจโย.\csromanspacer{} \textbf{ปจฺฉาชาตา} นเหตู ขนฺธา ปุเรชาตสฺส อิมสฺส กายสฺส.\csromanspacer{} อตฺถิปจฺจเยน ปจฺจโย.\csromanspacer{} \textbf{กพฬีกาโร อาหาโร} อิมสฺส กายสฺส อตฺถิปจฺจเยน ปจฺจโย.\csromanspacer{} \textbf{รูปชีวิตินฺทฺริยํ} กฏตฺตารูปานํ อตฺถิปจฺจเยน ปจฺจโย. (1)}

\prose{นเหตุ ธมฺโม เหตุสฺส ธมฺมสฺส อตฺถิปจฺจเยน ปจฺจโย.\csromanspacer{} \textbf{สหชาตํ, ปุเรชาตํ}.\csromanspacer{}\csromanspacer{}\textbf{สหชาตา} นเหตู ขนฺธา สมฺปยุตฺตกานํ เหตูนํ อตฺถิปจฺจเยน ปจฺจโย.\csromanspacer{} ปฏิสนฺธิกฺขเณ ฯเปฯ.\csromanspacer{} วตฺถุ เหตูนํ อตฺถิปจฺจเยน ปจฺจโย.\csromanspacer{} \textbf{ปุเรชาตํ} จกฺขุํ ฯเปฯ วตฺถุํ อนิจฺจโต ฯเปฯ โทมนสฺสํ อุปฺปชฺชติ.\csromanspacer{} ทิพฺเพน จกฺขุนา รูปํ ปสฺสติ.\csromanspacer{} ทิพฺพาย โสตธาตุยา สทฺทํ สุณาติ.\csromanspacer{} วตฺถุ เหตูนํ อตฺถิปจฺจเยน ปจฺจโย. (2)}

\prose{นเหตุ ธมฺโม เหตุสฺส จ นเหตุสฺส จ ธมฺมสฺส อตฺถิปจฺจเยน ปจฺจโย.\csromanspacer{} \textbf{สหชาตํ, ปุเรชาตํ}.\csromanspacer{}\csromanspacer{}\textbf{สหชาโต} นเหตุ เอโก ขนฺโธ ติณฺณนฺนํ ขนฺธานํ เหตูนํ จิตฺตสมุฏฺฐานานญฺจ รูปานํ อตฺถิปจฺจเยน ปจฺจโย.\csromanspacer{} ปฏิสนฺธิกฺขเณ ฯเปฯ.\csromanspacer{} วตฺถุ เหตูนํ สมฺปยุตฺตกานญฺจ ขนฺธานํ อตฺถิปจฺจเยน ปจฺจโย.\csromanspacer{} \textbf{ปุเรชาตํ} จกฺขุํ ฯเปฯ วตฺถุํ อนิจฺจโต ฯเปฯ โทมนสฺสํ อุปฺปชฺชติ.\csromanspacer{} ทิพฺเพน จกฺขุนา รูปํ ปสฺสติ.\csromanspacer{} ทิพฺพาย โสตธาตุยา สทฺทํ สุณาติ.\csromanspacer{} วตฺถุ เหตูนํ สมฺปยุตฺตกานญฺจ ขนฺธานํ อตฺถิปจฺจเยน ปจฺจโย. (3)}

\proseitem{41}{เหตุ จ นเหตุ จ ธมฺมา เหตุสฺส ธมฺมสฺส อตฺถิปจฺจเยน ปจฺจโย.\csromanspacer{} \textbf{สหชาตํ, ปุเรชาตํ}.\csromanspacer{}\csromanspacer{}\textbf{สหชาโต} อโลโภ จ สมฺปยุตฺตกา จ ขนฺธา อโทสสฺส อโมหสฺส อตฺถิปจฺจเยน ปจฺจโย. \csromanfolio{18}(จกฺกํ.) โลโภ จ สมฺปยุตฺตกา จ ขนฺธา โมหสฺส อตฺถิปจฺจเยน ปจฺจโย. (จกฺกํ.) ปฏิสนฺธิกฺขเณ ฯเปฯ.\csromanspacer{} อโลโภ จ วตฺถุ จ อโทสสฺส อโมหสฺส อตฺถิปจฺจเยน ปจฺจโย. (จกฺกํ.) (1)}

\prose{เหตุ จ นเหตุ จ ธมฺมา นเหตุสฺส ธมฺมสฺส อตฺถิปจฺจเยน ปจฺจโย.\csromanspacer{} สหชาตํ, ปจฺฉาชาตํ, อาหารํ, อินฺทฺริยํ.\csromanspacer{}\csromanspacer{}\textbf{สหชาโต} นเหตุ เอโก ขนฺโธ จ เหตู จ ติณฺณนฺนํ ขนฺธานํ จิตฺตสมุฏฺฐานานญฺจ รูปานํ อตฺถิปจฺจเยน ปจฺจโย ฯเปฯ เทฺว ขนฺธา ฯเปฯ.\csromanspacer{} ปฏิสนฺธิกฺขเณ ฯเปฯ.\csromanspacer{} ปฏิสนฺธิกฺขเณ เหตู จ วตฺถุ จ นเหตูนํ ขนฺธานํ อตฺถิปจฺจเยน ปจฺจโย.\csromanspacer{} สหชาตา เหตู จ มหาภูตา จ จิตฺตสมุฏฺฐานานํ รูปานํ อตฺถิปจฺจเยน ปจฺจโย.\csromanspacer{} \textbf{ปจฺฉาชาตา} เหตู จ กพฬีกาโร อาหาโร จ อิมสฺส กายสฺส อตฺถิปจฺจเยน ปจฺจโย.\csromanspacer{} \textbf{ปจฺฉาชาตา} เหตู จ รูปชีวิตินฺทฺริยญฺจ กฏตฺตารูปานํ อตฺถิปจฺจเยน ปจฺจโย. (2)}

\prose{เหตุ จ นเหตุ จ ธมฺมา เหตุสฺส จ นเหตุสฺส จ ธมฺมสฺส อตฺถิปจฺจเยน ปจฺจโย.\csromanspacer{} \textbf{สหชาตํ, ปุเรชาตํ}.\csromanspacer{}\csromanspacer{}\textbf{สหชาโต} นเหตุ เอโก ขนฺโธ จ อโลโภ จ ติณฺณนฺนํ ขนฺธานํ อโทสสฺส อโมหสฺส จิตฺตสมุฏฺฐานานญฺจ รูปานํ อตฺถิปจฺจเยน ปจฺจโย. (จกฺกํ.) นเหตุ เอโก ขนฺโธ จ โลโภ จ ติณฺณนฺนํ ขนฺธานํ โมหสฺส จิตฺตสมุฏฺฐานานญฺจ รูปานํ อตฺถิปจฺจเยน ปจฺจโย. (จกฺกํ.) ปฏิสนฺธิกฺขเณ นเหตุ เอโก ขนฺโธ จ อโลโภ จ. (จกฺกํ.) ปฏิสนฺธิกฺขเณ ฯเปฯ.\csromanspacer{} อโลโภ จ วตฺถุ จ อโทสสฺส อโมหสฺส สมฺปยุตฺตกานญฺจ ขนฺธานํ อตฺถิปจฺจเยน ปจฺจโย.\csromanspacer{} โลโภ จ วตฺถุ จ โมหสฺส สมฺปยุตฺตกานญฺจ ขนฺธานํ อตฺถิปจฺจเยน ปจฺจโย.}

\prose{นตฺถิปจฺจเยน ปจฺจโย.\csromanspacer{} วิคตปจฺจเยน ปจฺจโย.\csromanspacer{} อวิคตปจฺจเยน ปจฺจโย. (3)}

\csromanheader{2}{1. ปจฺจยานุโลม 2. สงฺขฺยาวาร}
\csromanheader{2}{\textbf{สุทฺธ}}
\proseitem{42}{เหตุยา ตีณิ, อารมฺมเณ นว, อธิปติยา นว, อนนฺตเร นว, สมนนฺตเร นว, สหชาเต นว, อญฺญมญฺเญ นว, นิสฺสเย นว, อุปนิสฺสเย นว, ปุเรชาเต ตีณิ, ปจฺฉาชาเต ตีณิ, อาเสวเน นว, กมฺเม ตีณิ, วิปาเก นว, อาหาเร ตีณิ, อินฺทฺริเย นว, ฌาเน ตีณิ, \csromanfolio{19}มคฺเค นว, สมฺปยุตฺเต นว, วิปฺปยุตฺเต ปญฺจ, อตฺถิยา นว, นตฺถิยา นว, วิคเต นว, อวิคเต นว. (เอวํ อนุมชฺชนฺเตน คเณตพฺพํ.)}

\vspace{\tipitakaparskip}
\csromancenter{อนุโลมํ.}
\csromanheader{2}{2. ปจฺจนียุทฺธาร}
\proseitem{43}{เหตุ ธมฺโม เหตุสฺส ธมฺมสฺส อารมฺมณปจฺจเยน ปจฺจโย.\csromanspacer{} สหชาตปจฺจเยน ปจฺจโย.\csromanspacer{} อุปนิสฺสยปจฺจเยน ปจฺจโย. (1)}

\prose{เหตุ ธมฺโม นเหตุสฺส ธมฺมสฺส อารมฺมณปจฺจเยน ปจฺจโย.\csromanspacer{} สหชาตปจฺจเยน ปจฺจโย.\csromanspacer{} อุปนิสฺสยปจฺจเยน ปจฺจโย.\csromanspacer{} ปจฺฉาชาตปจฺจเยน ปจฺจโย. (2)}

\prose{เหตุ ธมฺโม เหตุสฺส จ นเหตุสฺส จ ธมฺมสฺส อารมฺมณปจฺจเยน ปจฺจโย.\csromanspacer{} สหชาตปจฺจเยน ปจฺจโย.\csromanspacer{} อุปนิสฺสยปจฺจเยน ปจฺจโย. (3)}

\proseitem{44}{นเหตุ ธมฺโม นเหตุสฺส ธมฺมสฺส อารมฺมณปจฺจเยน ปจฺจโย.\csromanspacer{} สหชาตปจฺจเยน ปจฺจโย.\csromanspacer{} อุปนิสฺสยปจฺจเยน ปจฺจโย.\csromanspacer{} ปุเรชาตปจฺจเยน ปจฺจโย.\csromanspacer{} ปจฺฉาชาตปจฺจเยน ปจฺจโย.\csromanspacer{} กมฺมปจฺจเยน ปจฺจโย.\csromanspacer{} อาหารปจฺจเยน ปจฺจโย.\csromanspacer{} อินฺทฺริยปจฺจเยน ปจฺจโย. (1)}

\prose{นเหตุ ธมฺโม เหตุสฺส ธมฺมสฺส อารมฺมณปจฺจเยน ปจฺจโย.\csromanspacer{} สหชาตปจฺจเยน ปจฺจโย.\csromanspacer{} อุปนิสฺสยปจฺจเยน ปจฺจโย.\csromanspacer{} ปุเรชาตปจฺจเยน ปจฺจโย.\csromanspacer{} กมฺมปจฺจเยน ปจฺจโย. (2)}

\prose{นเหตุ ธมฺโม เหตุสฺส จ นเหตุสฺส จ ธมฺมสฺส อารมฺมณปจฺจเยน ปจฺจโย.\csromanspacer{} สหชาตปจฺจเยน ปจฺจโย.\csromanspacer{} อุปนิสฺสยปจฺจเยน ปจฺจโย.\csromanspacer{} ปุเรชาตปจฺจเยน ปจฺจโย.\csromanspacer{} กมฺมปจฺจเยน ปจฺจโย. (3)}

\proseitem{45}{เหตุ จ นเหตุ จ ธมฺมา เหตุสฺส ธมฺมสฺส อารมฺมณปจฺจเยน ปจฺจโย.\csromanspacer{} สหชาตปจฺจเยน ปจฺจโย.\csromanspacer{} อุปนิสฺสยปจฺจเยน ปจฺจโย. (1)}

\prose{เหตุ จ นเหตุ จ ธมฺมา นเหตุสฺส ธมฺมสฺส อารมฺมณปจฺจเยน ปจฺจโย.\csromanspacer{} สหชาตปจฺจเยน ปจฺจโย.\csromanspacer{} อุปนิสฺสยปจฺจเยน ปจฺจโย.\csromanspacer{} ปจฺฉาชาตปจฺจเยน ปจฺจโย. (2)}

\prose{\csromanfolio{20}เหตุ จ นเหตุ จ ธมฺมา เหตุสฺส จ นเหตุสฺส จ ธมฺมสฺส อารมฺมณปจฺจเยน ปจฺจโย.\csromanspacer{} สหชาตปจฺจเยน ปจฺจโย.\csromanspacer{} อุปนิสฺสยปจฺจเยน ปจฺจโย. (3)}

\csromanheader{2}{2. ปจฺจยปจฺจนีย 2. สงฺขฺยาวาร}
\proseitem{46}{นเหตุยา นว, น-อารมฺมเณ นว ฯเปฯ โน-อวิคเต นว. (เอวํ คเณตพฺพํ.)}

\vspace{\tipitakaparskip}
\csromancenter{ปจฺจนียํ.}
\csromanheader{2}{3. ปจฺจยานุโลมปจฺจนีย}
\csromanheader{2}{\textbf{เหตุทุก}}
\proseitem{47}{เหตุปจฺจยา น-อารมฺมเณ ตีณิ, น-อธิปติยา ตีณิ, น-อนนฺตเร ตีณิ, นสมนนฺตเร ตีณิ, น-อญฺญมญฺเญ เอกํ, น-อุปนิสฺสเย ตีณิ ฯเปฯ นมคฺเค ตีณิ, นสมฺปยุตฺเต เอกํ, นวิปฺปยุตฺเต ตีณิ, โนนตฺถิยา ตีณิ, โนวิคเต ตีณิ. (เอวํ คเณตพฺพํ.)}

\vspace{\tipitakaparskip}
\csromancenter{อนุโลมปจฺจนียํ.}
\csromanheader{2}{4. ปจฺจยปจฺจนียานุโลม}
\csromanheader{2}{\textbf{นเหตุทุก}}
\proseitem{48}{นเหตุปจฺจยา อารมฺมเณ นว, อธิปติยา นว, อนนฺตเร นว, สมนนฺตเร นว, สหชาเต ตีณิ, อญฺญมญฺเญ ตีณิ, นิสฺสเย ตีณิ, อุปนิสฺสเย นว, ปุเรชาเต ตีณิ, ปจฺฉาชาเต ตีณิ, อาเสวเน นว, กมฺเม ตีณิ, วิปาเก ตีณิ, อาหาเร ตีณิ, อินฺทฺริเย ตีณิ, ฌาเน ตีณิ, มคฺเค ตีณิ, สมฺปยุตฺเต ตีณิ, วิปฺปยุตฺเต ตีณิ, อตฺถิยา ตีณิ, นตฺถิยา นว, วิคเต นว, อวิคเต ตีณิ. (เอวํ คเณตพฺพํ.)}

\vspace{\tipitakaparskip}
\csromancenter{ปจฺจนียานุโลมํ.}
\nitthitam{เหตุทุกํ นิฏฺฐิตํ.}
\csromanmatikaanchor{mtk.4}
\csromantocmarkref{section}{2. สเหตุกทุก}{mtk.4}
\bookmark[dest={mtk.4},level=1]{2. สเหตุกทุก}
\csromanheader{1}{2. สเหตุทุก \textbf{2. สเหตุกทุก 1. ปฏิจฺจวาร}}
\csromanheader{2}{1. ปจฺจยานุโลม 1. วิภงฺควาร}
\csromanheader{2}{\textbf{เหตุ}}
\proseitem{49}{\csromanfolio{21}สเหตุกํ ธมฺมํ ปฏิจฺจ สเหตุโก ธมฺโม อุปฺปชฺชติ เหตุปจฺจยา.\csromanspacer{} สเหตุกํ เอกํ ขนฺธํ ปฏิจฺจ ตโย ขนฺธา ฯเปฯ เทฺว ขนฺเธ ฯเปฯ.\csromanspacer{} ปฏิสนฺธิกฺขเณ ฯเปฯ. (1)}

\prose{สเหตุกํ ธมฺมํ ปฏิจฺจ อเหตุโก ธมฺโม อุปฺปชฺชติ เหตุปจฺจยา.\csromanspacer{} สเหตุเก ขนฺเธ ปฏิจฺจ จิตฺตสมุฏฺฐานํ รูปํ.\csromanspacer{} ปฏิสนฺธิกฺขเณ ฯเปฯ.}

\prose{สเหตุกํ ธมฺมํ ปฏิจฺจ สเหตุโก จ อเหตุโก จ ธมฺมา อุปฺปชฺชนฺติ เหตุปจฺจยา.\csromanspacer{} สเหตุกํ เอกํ ขนฺธํ ปฏิจฺจ ตโย ขนฺธา จิตฺตสมุฏฺฐานญฺจ รูปํ ฯเปฯ เทฺว ขนฺเธ ฯเปฯ.\csromanspacer{} ปฏิสนฺธิกฺขเณ ฯเปฯ. (3)}

\proseitem{50}{อเหตุกํ ธมฺมํ ปฏิจฺจ อเหตุโก ธมฺโม อุปฺปชฺชติ เหตุปจฺจยา.\csromanspacer{} วิจิกิจฺฉาสหคตํ อุทฺธจฺจสหคตํ โมหํ ปฏิจฺจ จิตฺตสมุฏฺฐานํ รูปํ.\csromanspacer{} เอกํ มหาภูตํ ปฏิจฺจ ตโย มหาภูตา ฯเปฯ มหาภูเต ปฏิจฺจ จิตฺตสมุฏฺฐานํ รูปํ, กฏตฺตารูปํ, อุปาทารูปํ. (1)}

\prose{อเหตุกํ ธมฺมํ ปฏิจฺจ สเหตุโก ธมฺโม อุปฺปชฺชติ เหตุปจฺจยา.\csromanspacer{} วิจิกิจฺฉาสหคตํ อุทฺธจฺจสหคตํ โมหํ ปฏิจฺจ สมฺปยุตฺตกา ขนฺธา.\csromanspacer{} ปฏิสนฺธิกฺขเณ วตฺถุํ ปฏิจฺจ สเหตุกา ขนฺธา. (2)}

\prose{อเหตุกํ ธมฺมํ ปฏิจฺจ สเหตุโก จ อเหตุโก จ ธมฺมา อุปฺปชฺชนฺติ อเหตุปจฺจยา.\csromanspacer{} วิจิกิจฺฉาสหคตํ อุทฺธจฺจสหคตํ โมหํ ปฏิจฺจ สมฺปยุตฺตกา ขนฺธา จิตฺตสมุฏฺฐานญฺจ รูปํ.\csromanspacer{} ปฏิสนฺธิกฺขเณ วตฺถุํ ปฏิจฺจ สเหตุกา ขนฺธา.\csromanspacer{} มหาภูเต ปฏิจฺจ กฏตฺตารูปํ. (3)}

\proseitem{51}{สเหตุกญฺจ อเหตุกญฺจ ธมฺมํ ปฏิจฺจ สเหตุโก ธมฺโม อุปฺปชฺชติ เหตุปจฺจยา.\csromanspacer{} วิจิกิจฺฉาสหคตํ อุทฺธจฺจสหคตํ เอกํ ขนฺธญฺจ โมหญฺจ ปฏิจฺจ ตโย ขนฺธา ฯเปฯ เทฺว ขนฺเธ ฯเปฯ.\csromanspacer{} ปฏิสนฺธิกฺขเณ สเหตุกํ เอกํ ขนฺธญฺจ วตฺถุญฺจ ปฏิจฺจ ตโย ขนฺธา ฯเปฯ เทฺว ขนฺเธ ฯเปฯ. (1)}

\prose{\csromanfolio{22}สเหตุกญฺจ อเหตุกญฺจ ธมฺมํ ปฏิจฺจ อเหตุโก ธมฺโม อุปฺปชฺชติ เหตุปจฺจยา.\csromanspacer{} สเหตุเก ขนฺเธ จ มหาภูเต จ ปฏิจฺจ จิตฺตสมุฏฺฐานํ รูปํ.\csromanspacer{} วิจิกิจฺฉาสหคเต อุทฺธจฺจสหคเต ขนฺเธ จ โมหญฺจ ปฏิจฺจ จิตฺตสมุฏฺฐานํ รูปํ.\csromanspacer{} ปฏิสนฺธิกฺขเณ ฯเปฯ. (2)}

\prose{สเหตุกญฺจ อเหตุกญฺจ ธมฺมํ ปฏิจฺจ สเหตุโก จ อเหตุโก จ ธมฺมา อุปฺปชฺชนฺติ เหตุปจฺจยา.\csromanspacer{} วิจิกิจฺฉาสหคตํ อุทฺธจฺจสหคตํ เอกํ ขนฺธญฺจ โมหญฺจ ปฏิจฺจ ตโย ขนฺธา จิตฺตสมุฏฺฐานญฺจ รูปํ ฯเปฯ เทฺว ขนฺเธ ฯเปฯ.\csromanspacer{} ปฏิสนฺธิกฺขเณ สเหตุกํ เอกํ ขนฺธญฺจ วตฺถุญฺจ ปฏิจฺจ ตโย ขนฺธา ฯเปฯ เทฺว ขนฺเธ ฯเปฯ.\csromanspacer{} สเหตุเก ขนฺเธ จ มหาภูเต จ ปฏิจฺจ กฏตฺตารูปํ. (3)}

\csromanheader{2}{\textbf{อารมฺมณ}}
\proseitem{52}{สเหตุกํ ธมฺมํ ปฏิจฺจ สเหตุโก ธมฺโม อุปฺปชฺชติ อารมฺมณปจฺจยา.\csromanspacer{} สเหตุกํ เอกํ ขนฺธํ ปฏิจฺจ ตโย ขนฺธา ฯเปฯ เทฺว ขนฺเธ ฯเปฯ.\csromanspacer{} ปฏิสนฺธิกฺขเณ ฯเปฯ. (1)}

\prose{สเหตุกํ ธมฺมํ ปฏิจฺจ อเหตุโก ธมฺโม อุปฺปชฺชติ อารมฺมณปจฺจยา.\csromanspacer{} วิจิกิจฺฉาสหคเต อุทฺธจฺจสหคเต ขนฺเธ ปฏิจฺจ วิจิกิจฺฉาสหคโต อุทฺธจฺจสหคโต โมโห. (2)}

\prose{สเหตุกํ ธมฺมํ ปฏิจฺจ สเหตุโก จ อเหตุโก จ ธมฺมา อุปฺปชฺชนฺติ อารมฺมณปจฺจยา.\csromanspacer{} วิจิกิจฺฉาสหคตํ อุทฺธจฺจสหคตํ เอกํ ขนฺธํ ปฏิจฺจ ตโย ขนฺธา โมโห จ ฯเปฯ เทฺว ขนฺเธ ฯเปฯ. (3)}

\proseitem{53}{อเหตุกํ ธมฺมํ ปฏิจฺจ อเหตุโก ธมฺโม อุปฺปชฺชติ อารมฺมณปจฺจยา.\csromanspacer{} อเหตุกํ เอกํ ขนฺธํ ปฏิจฺจ ตโย ขนฺธา ฯเปฯ เทฺว ขนฺเธ ฯเปฯ.\csromanspacer{} ปฏิสนฺธิกฺขเณ วตฺถุํ ปฏิจฺจ อเหตุกา ขนฺธา. (1)}

\prose{อเหตุกํ ธมฺมํ ปฏิจฺจ สเหตุโก ธมฺโม อุปฺปชฺชติ อารมฺมณปจฺจยา.\csromanspacer{} วิจิกิจฺฉาสหคตํ อุทฺธจฺจสหคตํ โมหํ ปฏิจฺจ สมฺปยุตฺตกา ขนฺธา.\csromanspacer{} ปฏิสนฺธิกฺขเณ วตฺถุํ ปฏิจฺจ สเหตุกา ขนฺธา. (2)}

\prose{สเหตุกญฺจ อเหตุกญฺจ ธมฺมํ ปฏิจฺจ สเหตุโก ธมฺโม อุปฺปชฺชติ อารมฺมณปจฺจยา.\csromanspacer{} วิจิกิจฺฉาสหคตํ อุทฺธจฺจสหคตํ เอกํ ขนฺธญฺจ โมหญฺจ ปฏิจฺจ}

\vspace{\tipitakaparskip}
\csromancenter{สเหตุทุก ตโย ขนฺธา ฯเปฯ เทฺว ขนฺเธ ฯเปฯ.\csromanspacer{} ปฏิสนฺธิกฺขเณ สเหตุกํ เอกํ ขนฺธญฺจ วตฺถุญฺจ ปฏิจฺจ ตโย ขนฺธา ฯเปฯ เทฺว ขนฺเธ ฯเปฯ. (1)}
\csromanheader{2}{\textbf{อธิปติ}}
\proseitem{54}{\csromanfolio{23}สเหตุกํ ธมฺมํ ปฏิจฺจ สเหตุโก ธมฺโม อุปฺปชฺชติ อธิปติปจฺจยา.\csromanspacer{} สเหตุกํ เอกํ ขนฺธํ ปฏิจฺจ ตโย ขนฺธา ฯเปฯ เทฺว ขนฺเธ ฯเปฯ. (1)}

\prose{สเหตุกํ ธมฺมํ ปฏิจฺจ อเหตุโก ธมฺโม อุปฺปชฺชติ อธิปติปจฺจยา.\csromanspacer{} สเหตุเก ขนฺเธ ปฏิจฺจ จิตฺตสมุฏฺฐานํ รูปํ. (2)}

\prose{สเหตุกํ ธมฺมํ ปฏิจฺจ สเหตุโก จ อเหตุโก จ ธมฺมา อุปฺปชฺชนฺติ อธิปติปจฺจยา.\csromanspacer{} สเหตุกํ เอกํ ขนฺธํ ปฏิจฺจ ตโย ขนฺธา จิตฺตสมุฏฺฐานญฺจ รูปํ ฯเปฯ เทฺว ขนฺเธ ฯเปฯ. (3)}

\proseitem{55}{อเหตุกํ ธมฺมํ ปฏิจฺจ อเหตุโก ธมฺโม อุปฺปชฺชติ อธิปติปจฺจยา.\csromanspacer{} เอกํ มหาภูตํ ฯเปฯ มหาภูเต ปฏิจฺจ จิตฺตสมุฏฺฐานํ รูปํ, อุปาทารูปํ. (1)}

\prose{สเหตุกญฺจ อเหตุกญฺจ ธมฺมํ ปฏิจฺจ อเหตุโก ธมฺโม อุปฺปชฺชติ อธิปติปจฺจยา.\csromanspacer{} สเหตุเก ขนฺเธ จ มหาภูเต จ ปฏิจฺจ จิตฺตสมุฏฺฐานํ รูปํ. (1)}

\csromanheader{2}{\textbf{อนนฺตราทิ}}
\proseitem{56}{สเหตุกํ ธมฺมํ ปฏิจฺจ สเหตุโก ธมฺโม อุปฺปชฺชติ อนนฺตรปจฺจยา.\csromanspacer{} สมนนฺตรปจฺจยา.\csromanspacer{} สหชาตปจฺจยา.\csromanspacer{} สเหตุกํ เอกํ ขนฺธํ ปฏิจฺจ ตโย ขนฺธา ฯเปฯ.\csromanspacer{} ปฏิสนฺธิกฺขเณ ฯเปฯ. (1)}

\prose{สเหตุกํ ธมฺมํ ปฏิจฺจ อเหตุโก ธมฺโม อุปฺปชฺชติ สหชาตปจฺจยา.\csromanspacer{} สเหตุเก ขนฺเธ ปฏิจฺจ จิตฺตสมุฏฺฐานํ รูปํ วิจิกิจฺฉาสหคเต อุทฺธจฺจสหคเต ขนฺเธ ปฏิจฺจ วิจิกิจฺฉาสหคโต อุทฺธจฺจสหคโต โมโห จิตฺตสมุฏฺฐานญฺจ รูปํ.\csromanspacer{} ปฏิสนฺธิกฺขเณ ฯเปฯ.}

\prose{สเหตุกํ ธมฺมํ ปฏิจฺจ สเหตุโก จ อเหตุโก จ ธมฺมา อุปฺปชฺชนฺติ สหชาตปจฺจยา.\csromanspacer{} สเหตุกํ เอกํ ขนฺธํ ปฏิจฺจ ตโย ขนฺธา จิตฺตสมุฏฺฐานญฺจ รูปํ เทฺว ขนฺเธ ฯเปฯ.\csromanspacer{} วิจิกิจฺฉาสหคตํ อุทฺธจฺจสหคตํ เอกํ ขนฺธํ ปฏิจฺจ ตโย ขนฺธา โมโห จ จิตฺตสมุฏฺฐานญฺจ รูปํ ฯเปฯ เทฺว ขนฺเธ ฯเปฯ.\csromanspacer{} ปฏิสนฺธิกฺขเณ ฯเปฯ. (3)}

\proseitem{57}{\csromanfolio{24}อเหตุกํ ธมฺมํ ปฏิจฺจ อเหตุโก ธมฺโม อุปฺปชฺชติ สหชาตปจฺจยา.\csromanspacer{} อเหตุกํ เอกํ ขนฺธํ ปฏิจฺจ ตโย ขนฺธา จิตฺตสมุฏฺฐานญฺจ รูปํ ฯเปฯ เทฺว ขนฺเธ ฯเปฯ.\csromanspacer{} วิจิกิจฺฉาสหคตํ อุทฺธจฺจสหคตํ โมหํ ปฏิจฺจ จิตฺตสมุฏฺฐานํ รูปํ.\csromanspacer{} ปฏิสนฺธิกฺขเณ ฯเปฯ.\csromanspacer{} ขนฺเธ ปฏิจฺจ วตฺถุ, วตฺถุํ ปฏิจฺจ ขนฺธา.\csromanspacer{} เอกํ มหาภูตํ ฯเปฯ.\csromanspacer{} พาหิรํ.\csromanspacer{} อาหารสมุฏฺฐานํ.\csromanspacer{} อุตุสมุฏฺฐานํ.\csromanspacer{} อสญฺญสตฺตานํ เอกํ มหาภูตํ ฯเปฯ. (1)}

\prose{อเหตุกํ ธมฺมํ ปฏิจฺจ สเหตุโก ธมฺโม อุปฺปชฺชติ สหชาตปจฺจยา. (อิเม ปญฺจ ปญฺหา เหตุสทิสา, นินฺนานํ.) (2)}

\csromanheader{2}{\textbf{อญฺญมญฺญ}}
\proseitem{58}{สเหตุกํ ธมฺมํ ปฏิจฺจ สเหตุโก ธมฺโม อุปฺปชฺชติ อญฺญมญฺญปจฺจยา.\csromanspacer{} สเหตุกํ เอกํ ขนฺธํ ปฏิจฺจ ตโย ขนฺธา ฯเปฯ.\csromanspacer{} ปฏิสนฺธิกฺขเณ ฯเปฯ. (1)}

\prose{สเหตุกํ ธมฺมํ ปฏิจฺจ อเหตุโก ธมฺโม อุปฺปชฺชติ อญฺญมญฺญปจฺจยา.\csromanspacer{} วิจิกิจฺฉาสหคเต อุทฺธจฺจสหคเต ขนฺเธ ปฏิจฺจ วิจิกิจฺฉาสหคโต อุทฺธจฺจสหคโต โมโห.\csromanspacer{} ปฏิสนฺธิกฺขเณ สเหตุเก ขนฺเธ ปฏิจฺจ วตฺถุ. (2)}

\prose{สเหตุกํ ธมฺมํ ปฏิจฺจ สเหตุโก จ อเหตุโก จ ธมฺมา อุปฺปชฺชนฺติ อญฺญมญฺญปจฺจยา.\csromanspacer{} วิจิกิจฺฉาสหคตํ อุทฺธจฺจสหคตํ เอกํ ขนฺธํ ปฏิจฺจ ตโย ขนฺธา โมโห จ ฯเปฯ เทฺว ขนฺเธ ฯเปฯ.\csromanspacer{} ปฏิสนฺธิกฺขเณ สเหตุกํ เอกํ ขนฺธํ ปฏิจฺจ ตโย ขนฺธา วตฺถุ จ ฯเปฯ เทฺว ขนฺเธ ฯเปฯ. (3)}

\proseitem{59}{อเหตุกํ ธมฺมํ ปฏิจฺจ อเหตุโก ธมฺโม อุปฺปชฺชติ อญฺญมญฺญปจฺจยา.\csromanspacer{} อเหตุกํ เอกํ ขนฺธํ ปฏิจฺจ ตโย ขนฺธา ฯเปฯ เทฺว ขนฺเธ ฯเปฯ.\csromanspacer{} ปฏิสนฺธิกฺขเณ อเหตุกํ เอกํ ขนฺธํ ปฏิจฺจ ตโย ขนฺธา วตฺถุ จ ฯเปฯ เทฺว ขนฺเธ ฯเปฯ. (สํขิตฺตํ.\csromanspacer{} ยาว อสญฺญสตฺตา.) (1)}

\prose{อเหตุกํ ธมฺมํ ปฏิจฺจ สเหตุโก ธมฺโม อุปฺปชฺชติ อญฺญมญฺญปจฺจยา.\csromanspacer{} วิจิกิจฺฉาสหคตํ อุทฺธจฺจสหคตํ โมหํ ปฏิจฺจ สมฺปยุตฺตกา ขนฺธา.\csromanspacer{} ปฏิสนฺธิกฺขเณ วตฺถุํ ปฏิจฺจ สเหตุกา ขนฺธา. (2)}

\prose{สเหตุกญฺจ อเหตุกญฺจ ธมฺมํ ปฏิจฺจ สเหตุโก ธมฺโม อุปฺปชฺชติ อญฺญมญฺญปจฺจยา.\csromanspacer{} วิจิกิจฺฉาสหคตํ อุทฺธจฺจสหคตํ เอกํ ขนฺธญฺจ โมหญฺจ}

\vspace{\tipitakaparskip}
\csromancenter{สเหตุทุก ปฏิจฺจ ตโย ขนฺธา ฯเปฯ เทฺว ขนฺเธ ฯเปฯ.\csromanspacer{} ปฏิสนฺธิกฺขเณ สเหตุกํ เอกํ ขนฺธญฺจ วตฺถุญฺจ ปฏิจฺจ ตโย ขนฺธา ฯเปฯ เทฺว ขนฺเธ ฯเปฯ. (1)}
\csromanheader{2}{\textbf{นิสฺสยาทิ}}
\proseitem{60}{\csromanfolio{25}สเหตุกํ ธมฺมํ ปฏิจฺจ สเหตุโก ธมฺโม อุปฺปชฺชติ นิสฺสยปจฺจยา.\csromanspacer{} อุปนิสฺสยปจฺจยา.\csromanspacer{} ปุเรชาตปจฺจยา.\csromanspacer{} อาเสวนปจฺจยา.\csromanspacer{} กมฺมปจฺจยา.\csromanspacer{} วิปากปจฺจยา.\csromanspacer{} อาหารปจฺจยา.\csromanspacer{} อินฺทฺริยปจฺจยา.\csromanspacer{} ฌานปจฺจยา.\csromanspacer{} มคฺคปจฺจยา. (ฌานมฺปิ มคฺคมฺปิ สหชาตปจฺจสสทิสา.\csromanspacer{} พาหิรา มหาภูตา นตฺถิ.) สมฺปยุตฺตปจฺจยา.\csromanspacer{} วิปฺปยุตฺตปจฺจยา.\csromanspacer{} อตฺถิปจฺจยา.\csromanspacer{} นตฺถิปจฺจยา.\csromanspacer{} วิคตปจฺจยา.\csromanspacer{} อวิคตปจฺจยา.}

\csromanheader{2}{1. ปจฺจยานุโลม 2. สงฺขฺยาวาร}
\csromanheader{2}{\textbf{สุทฺธ}}
\proseitem{61}{เหตุยา นว, อารมฺมเณ ฉ, อธิปติยา ปญฺจ, อนนฺตเร ฉ, สมนนฺตเร ฉ, สหชาเต นว, อญฺญมญฺเญ ฉ, นิสฺสเย นว, อุปนิสฺสเย ฉ, ปุเรชาเต ฉ, อาเสวเน ฉ, กมฺเม นว, วิปาเก นว, อาหาเร นว, อินฺทฺริเย นว, ฌาเน นว, มคฺเค นว, สมฺปยุตฺเต ฉ, วิปฺปยุตฺเต นว, อตฺถิยา นว, นตฺถิยา ฉ, วิคเต ฉ, อวิคเต นว. (เอวํ คเณตพฺพํ.)}

\vspace{\tipitakaparskip}
\csromancenter{อนุโลมํ.}
\csromanheader{2}{2. ปจฺจยปจฺจนีย 1. วิภงฺควาร}
\csromanheader{2}{\textbf{นเหตุ}}
\proseitem{62}{สเหตุกํ ธมฺมํ ปฏิจฺจ อเหตุโก ธมฺโม อุปฺปชฺชติ นเหตุปจฺจยา.\csromanspacer{} วิจิกิจฺฉาสหคเต อุทฺธจฺจสหคเต ขนฺเธ ปฏิจฺจ วิจิกิจฺฉาสหคโต อุทฺธจฺจสหคโต โมโห. (1)}

\prose{อเหตุกํ ธมฺมํ ปฏิจฺจ อเหตุโก ธมฺโม อุปฺปชฺชติ นเหตุปจฺจยา อเหตุกํ เอกํ ขนฺธํ ปฏิจฺจ ตโย ขนฺธา จิตฺตสมุฏฺฐานํ รูปํ ฯเปฯ เทฺว ขนฺเธ ฯเปฯ.\csromanspacer{} ปฏิสนฺธิกฺขเณ ฯเปฯ. (1)}

\vspace{\tipitakaparskip}
\csromancenter{(สพฺพํ ยาว อสญฺญสตฺตา, ตาว กาตพฺพํ.)}
\csromanheader{2}{\textbf{น-อารมฺมณาทิ}}
\proseitem{63}{\csromanfolio{26}สเหตุกํ ธมฺมํ ปฏิจฺจ อเหตุโก ธมฺโม อุปฺปชฺชติ น-อารมฺมณปจฺจยา.\csromanspacer{} สเหตุเก ขนฺเธ ปฏิจฺจ จิตฺตสมุฏฺฐานํ รูปํ.\csromanspacer{} ปฏิสนฺธิกฺขเณ ฯเปฯ. (1)}

\prose{อเหตุกํ ธมฺมํ ปฏิจฺจ อเหตุโก ธมฺโม อุปฺปชฺชติ น-อารมฺมณปจฺจยา.\csromanspacer{} อเหตุเก ขนฺเธ ปฏิจฺจ จิตฺตสมุฏฺฐานํ รูปํ.\csromanspacer{} วิจิกิจฺฉาสหคตํ อุทฺธจฺจสหคตํ โมหํ ปฏิจฺจ จิตฺตสมุฏฺฐานํ รูปํ.\csromanspacer{} ปฏิสนฺธิกฺขเณ ฯเปฯ.\csromanspacer{} ขนฺเธ ปฏิจฺจ วตฺถุ.\csromanspacer{} เอกํ มหาภูตํ ฯเปฯ.\csromanspacer{} อสญฺญสตฺตานํ เอกํ ฯเปฯ. (1)}

\prose{สเหตุกญฺจ อเหตุกญฺจ ธมฺมํ ปฏิจฺจ อเหตุโก ธมฺโม อุปฺปชฺชติ น-อารมฺมณปจฺจยา.\csromanspacer{} สเหตุเก ขนฺเธ จ มหาภูเต จ ปฏิจฺจจิตฺตสมุฏฺฐานํ รูปํ.\csromanspacer{} วิจิกิจฺฉาสหคเต อุทฺธจฺจสหคเต ขนฺเธ จ โมหญฺจ ปฏิจฺจ จิตฺตสมุฏฺฐานํ รูปํ.\csromanspacer{} ปฏิสนฺธิกฺขเณ ฯเปฯ. (1)}

\prose{สเหตุกํ ธมฺมํ ปฏิจฺจ สเหตุโก ธมฺโม อุปฺปชฺชติ น-อธิปติปจฺจยา. (อนุโลมสหชาตสทิสา.) น-อนนฺตรปจฺจยา.\csromanspacer{} นสมนนฺตรปจฺจยา.\csromanspacer{} น-อญฺญมญฺญปจฺจยา.\csromanspacer{} น-อุปนิสฺสยปจฺจยา.\csromanspacer{} นปุเรชาตปจฺจยา.\csromanspacer{} อรูเป สเหตุกํ เอกํ ขนฺธํ ฯเปฯ.\csromanspacer{} ปฏิสนฺธิกฺขเณ ฯเปฯ. (1)}

\proseitem{64}{สเหตุกํ ธมฺมํ ปฏิจฺจ อเหตุโก ธมฺโม อุปฺปชฺชติ นปุเรชาตปจฺจยา.\csromanspacer{} อรูเป วิจิกิจฺฉาสหคเต อุทฺธจฺจสหคเต ขนฺเธ ปฏิจฺจ วิจิกิจฺฉาสหคโต อุทฺธจฺจสหคโต โมโห.\csromanspacer{} สเหตุเก ขนฺเธ ปฏิจฺจ จิตฺตสมุฏฺฐานํ รูปํ.\csromanspacer{} ปฏิสนฺธิกฺขเณ ฯเปฯ. (2)}

\prose{สเหตุกํ ธมฺมํ ปฏิจฺจ สเหตุโก จ อเหตุโก จ ธมฺมา อุปฺปชฺชนฺติ นปุเรชาตปจฺจยา.\csromanspacer{} อรูเป วิจิกิจฺฉาสหคตํ อุทฺธจฺจสหคตํ เอกํ ขนฺธํ ปฏิจฺจ ตโย ขนฺธา โมโห จ ฯเปฯ เทฺว ขนฺเธ ฯเปฯ.\csromanspacer{} ปฏิสนฺธิกฺขเณ ฯเปฯ. (3)}

\prose{อเหตุกํ ธมฺมํ ปฏิจฺจ อเหตุโก ธมฺโม อุปฺปชฺชติ นปุเรชาตปจฺจยา.\csromanspacer{} อรูเป อเหตุกํ เอกํ ขนฺธํ ฯเปฯ เทฺว ขนฺเธ ฯเปฯ.\csromanspacer{} อเหตุเก ขนฺเธ ปฏิจฺจ จิตฺตสมุฏฺฐานํ รูปํ.\csromanspacer{} วิจิกิจฺฉาสหคตํ อุทฺธจฺจสหคตํ โมหํ ปฏิจฺจ จิตฺตสมุฏฺฐานํ รูปํ.\csromanspacer{} ปฏิสนฺธิกฺขเณ ฯเปฯ. (ยาว อสญฺญสตฺตา, ตาว วิตฺถาโร.) (1)}

\csromanheader{2}{2. สเหตุทุก}
\prose{\csromanfolio{27}อเหตุกํ ธมฺมํ ปฏิจฺจ สเหตุโก ธมฺโม อุปฺปชฺชติ นปุเรชาตปจฺจยา.\csromanspacer{} อรูเป วิจิกิจฺฉาสหคตํ อุทฺธจฺจสหคตํ โมหํ ปฏิจฺจ สมฺปยุตฺตกา ขนฺธา.\csromanspacer{} ปฏิสนฺธิกฺขเณ วตฺถุํ ปฏิจฺจ สเหตุกา ขนฺธา. (2)}

\prose{อเหตุกํ ธมฺมํ ปฏิจฺจ สเหตุโก จ อเหตุโก จ ธมฺมา อุปฺปชฺชนฺติ นปุเรชาตปจฺจยา.\csromanspacer{} ปฏิสนฺธิกฺขเณ วตฺถุํ ปฏิจฺจ สเหตุกา ขนฺธา.\csromanspacer{} มหาภูเต ปฏิจฺจ กฏตฺตารูปํ. (3)}

\proseitem{65}{สเหตุกญฺจ อเหตุกญฺจ ธมฺมํ ปฏิจฺจ สเหตุโก ธมฺโม อุปฺปชฺชติ นปุเรชาตปจฺจยา.\csromanspacer{} อรูเป วิจิกิจฺฉาสหคตํ อุทฺธจฺจสหคตํ เอกํ ขนฺธญฺจ โมหญฺจ ปฏิจฺจ ตโย ขนฺธา ฯเปฯ เทฺว ขนฺเธ ฯเปฯ.\csromanspacer{} ปฏิสนฺธิกฺขเณ สเหตุกํ เอกํ ขนฺธญฺจ วตฺถุญฺจ ปฏิจฺจ ตโย ขนฺธา ฯเปฯ เทฺว ขนฺเธ ฯเปฯ. (1)}

\prose{สเหตุกญฺจ อเหตุกญฺจ ธมฺมํ ปฏิจฺจ อเหตุโก ธมฺโม อุปฺปชฺชติ นปุเรชาตปจฺจยา.\csromanspacer{} สเหตุเก ขนฺเธ จ มหาภูเต จ ปฏิจฺจ จิตฺตสมุฏฺฐานํ รูปํ.\csromanspacer{} วิจิกิจฺฉาสหคเต อุทฺธจฺจสหคเต ขนฺเธ จ โมหญฺจ ปฏิจฺจ จิตฺตสมุฏฺฐานํ รูปํ.\csromanspacer{} ปฏิสนฺธิกฺขเณ ฯเปฯ. (2)}

\prose{สเหตุกญฺจ อเหตุกญฺจ ธมฺมํ ปฏิจฺจ สเหตุโก จ อเหตุโก จ ธมฺมา อุปฺปชฺชนฺติ นปุเรชาตปจฺจยา.\csromanspacer{} ปฏิสนฺธิกฺขเณ สเหตุกํ เอกํ ขนฺธญฺจ วตฺถุญฺจ ปฏิจฺจ ตโย ขนฺธา ฯเปฯ เทฺว ขนฺเธ ฯเปฯ.\csromanspacer{} สเหตุเก ขนฺเธ จ มหาภูเต จ ปฏิจฺจ กฏตฺตารูปํ. (3)}

\csromanheader{2}{\textbf{นปจฺฉาชาตาทิ}}
\proseitem{66}{สเหตุกํ ธมฺมํ ปฏิจฺจ สเหตุโก ธมฺโม อุปฺปชฺชติ นปจฺฉาชาตปจฺจยา.\csromanspacer{} น-อาเสวนปจฺจยา.\csromanspacer{} นกมฺมปจฺจยา.\csromanspacer{} สเหตุเก ขนฺเธ ปฏิจฺจ สเหตุกา เจตนา. (1)}

\prose{อเหตุกํ ธมฺมํ ปฏิจฺจ อเหตุโก ธมฺโม อุปฺปชฺชติ นกมฺมปจฺจยา.\csromanspacer{} อเหตุเก ขนฺเธ ปฏิจฺจ อเหตุกา เจตนา.\csromanspacer{} พาหิรํ.\csromanspacer{} อาหารสมุฏฺฐานํ.\csromanspacer{} อุตุสมุฏฺฐานํ ฯเปฯ. (1)}

\prose{อเหตุกํ ธมฺมํ ปฏิจฺจ สเหตุโก ธมฺโม อุปฺปชฺชติ นกมฺมปจฺจยา.\csromanspacer{} วิจิกิจฺฉาสหคตํ อุทฺธจฺจสหคตํ โมหํ ปฏิจฺจ สมฺปยุตฺตกา เจตนา. (2)}

\prose{\csromanfolio{28}สเหตุกญฺจ อเหตุกญฺจ ธมฺมํ ปฏิจฺจ สเหตุโก ธมฺโม อุปฺปชฺชติ นกมฺมปจฺจยา.\csromanspacer{} วิจิกิจฺฉาสหคเต อุทฺธจฺจสหคเต ขนฺเธ จ โมหญฺจ ปฏิจฺจ สมฺปยุตฺตกา เจตนา.\csromanspacer{} นวิปากปจฺจยา. (ปฏิสนฺธิ นตฺถิ.)}

\csromanheader{2}{\textbf{น-อาหาราทิ}}
\proseitem{67}{อเหตุกํ ธมฺมํ ปฏิจฺจ อเหตุโก ธมฺโม อุปฺปชฺชติ น-อาหารปจฺจยา.\csromanspacer{} น-อินฺทฺริยปจฺจยา.\csromanspacer{} นฌานปจฺจยา.\csromanspacer{} นมคฺคปจฺจยา.\csromanspacer{} นสมฺปยุตฺตปจฺจยา.}

\csromanheader{2}{\textbf{นวิปฺปยุตฺตาทิ}}
\proseitem{68}{สเหตุกํ ธมฺมํ ปฏิจฺจ สเหตุโก ธมฺโม อุปฺปชฺชติ นวิปฺปยุตฺตปจฺจยา.\csromanspacer{} อรูเป สเหตุกํ เอกํ ขนฺธํ ฯเปฯ. (1)}

\prose{สเหตุกํ ธมฺมํ ปฏิจฺจ อเหตุโก ธมฺโม อุปฺปชฺชติ นวิปฺปยุตฺตปจฺจยา.\csromanspacer{} อรูเป วิจิกิจฺฉาสหคเต อุทฺธจฺจสหคเต ขนฺเธ ปฏิจฺจ วิจิกิจฺฉาสหคโต อุทฺธจฺจสหคโต โมโห. (2)}

\prose{สเหตุกํ ธมฺมํ ปฏิจฺจ สเหตุโก จ อเหตุโก จ ธมฺมา อุปฺปชฺชนฺติ นวิปฺปยุตฺตปจฺจยา.\csromanspacer{} อรูเป วิจิกิจฺฉาสหคตํ อุทฺธจฺจสหคตํ เอกํ ขนฺธํ ปฏิจฺจ ตโย ขนฺธา โมโห จ ฯเปฯ เทฺว ขนฺเธ ฯเปฯ. (3)}

\prose{อเหตุกํ ธมฺมํ ปฏิจฺจ อเหตุโก ธมฺโม อุปฺปชฺชติ นวิปฺปยุตฺตปจฺจยา.\csromanspacer{} อรูเป อเหตุกํ เอกํ ขนฺธํ ปฏิจฺจ ตโย ขนฺธา ฯเปฯ เทฺว ขนฺเธ ฯเปฯ. (1)}

\prose{อเหตุกํ ธมฺมํ ปฏิจฺจ สเหตุโก ธมฺโม อุปฺปชฺชติ นวิปฺปยุตฺตปจฺจยา.\csromanspacer{} อรูเป วิจิกิจฺฉาสหคตํ อุทฺธจฺจสหคตํ โมหํ ปฏิจฺจ สมฺปยุตฺตกา ขนฺธา. (2)}

\prose{สเหตุกญฺจ อเหตุกญฺจ ธมฺมํ ปฏิจฺจ สเหตุโก ธมฺโม อุปฺปชฺชติ นวิปฺปยุตฺตปจฺจยา.\csromanspacer{} อรูเป วิจิกิจฺฉาสหคตํ อุทฺธจฺจสหคตํ เอกํ ขนฺธญฺจ โมหญฺจ ปฏิจฺจ ตโย ขนฺธา ฯเปฯ เทฺว ขนฺเธ ฯเปฯ.\csromanspacer{} โนนตฺถิปจฺจยา.\csromanspacer{} โนวิคตปจฺจยา. (1)}

\csromanheader{2}{2. ปจฺจยปจฺจนีย 2. สงฺขฺยาวาร}
\csromanheader{2}{\textbf{สุทฺธ}}
\proseitem{69}{นเหตุยา เทฺว, น-อารมฺมเณ ตีณิ, น-อธิปติยา นว, น-อนนฺตเร ตีณิ, นสมนนฺตเร ตีณิ, น-อญฺญมญฺเญ ตีณิ, น-อุปนิสฺสเย}

\vspace{\tipitakaparskip}
\csromancenter{สเหตุทุก ตีณิ, นปุเรชาเต นว, นปจฺฉาชาเต นว, น-อาเสวเน นว, นกมฺเม จตฺตาริ, นวิปาเก นว, น-อาหาเร เอกํ, น-อินฺทฺริเย เอกํ, นฌาเน เอกํ, นมคฺเค เอกํ, นสมฺปยุตฺเต ตีณิ, นวิปฺปยุตฺเต ฉ, โนนตฺถิยา ตีณิ, โนวิคเต ตีณิ. (เอวํ คเณตพฺพํ.)}
\csromancenter{ปจฺจนียํ.}
\csromanheader{2}{3. ปจฺจยานุโลมปจฺจนีย}
\csromanheader{2}{\textbf{เหตุทุก}}
\proseitem{70}{\csromanfolio{29}เหตุปจฺจยา น-อารมฺมเณ ตีณิ, น-อธิปติยา นว, น-อนนฺตเร ตีณิ, นสมนนฺตเร ตีณิ, น-อญฺญมญฺเญ ตีณิ, น-อุปนิสฺสเย ตีณิ, นปุเรชาเต นว, นปจฺฉาชาเต นว, น-อาเสวเน นว, นกมฺเม ตีณิ, นวิปาเก นว, นสมฺปยุตฺเต ตีณิ, นวิปฺปยุตฺเต ตีณิ, โนนตฺถิยา ตีณิ, โนวิคเต ตีณิ. (เอวํ คเณตพฺพํ.)}

\vspace{\tipitakaparskip}
\csromancenter{อนุโลมปจฺจนียํ.}
\csromanheader{2}{4. ปจฺจยปจฺจนียานุโลม}
\csromanheader{2}{\textbf{นเหตุทุก}}
\proseitem{71}{นเหตุปจฺจยา อารมฺมเณ เทฺว, อนนฺตเร เทฺว, สมนนฺตเร เทฺว, (สพฺพตฺถ เทฺว,) วิปาเก เอกํ, อาหาเร เทฺว, อินฺทฺริเย เทฺว, ฌาเน เทฺว, มคฺเค เอกํ, สมฺปยุตฺเต เทฺว ฯเปฯ อวิคเต เทฺว. (เอวํ คเณตพฺพํ.)}

\vspace{\tipitakaparskip}
\csromancenter{ปจฺจนียานุโลมํ.}
\csromancenter{(สหชาตวาโร ปฏิจฺจวารสทิโส.)}
\csromanheader{2}{2. สเหตุกทุก 3. ปจฺจยวาร}
\csromanheader{2}{1. ปจฺจยานุโลม 1. วิภงฺควาร}
\csromanheader{2}{\textbf{เหตุ}}
\proseitem{72}{\csromanfolio{30}สเหตุกํ ธมฺมํ ปจฺจยา สเหตุโก ธมฺโม อุปฺปชฺชติ เหตุปจฺจยา. (สเหตุกมูลกํ ปฏิจฺจวารสทิสํ.)}

\prose{อเหตุกํ ธมฺมํ ปจฺจยา อเหตุโก ธมฺโม ฯเปฯ. (ปฏิจฺจวารสทิสํเยว.) (1)}

\prose{อเหตุกํ ธมฺมํ ปจฺจยา สเหตุโก ธมฺโม อุปฺปชฺชติ เหตุปจฺจยา.\csromanspacer{} วตฺถุํ ปจฺจยา สเหตุกา ขนฺธา.\csromanspacer{} วิจิกิจฺฉาสหคตํ อุทฺธจฺจสหคตํ โมหํ ปจฺจยา สมฺปยุตฺตกา ขนฺธา.\csromanspacer{} ปฏิสนฺธิกฺขเณ วตฺถุํ ปจฺจยา สเหตุกา ขนฺธา. (2)}

\prose{อเหตุกํ ธมฺมํ ปจฺจยา สเหตุโก จ อเหตุโก จ ธมฺมา อุปฺปชฺชนฺติ เหตุปจฺจยา.\csromanspacer{} วตฺถุํ ปจฺจยา สเหตุกา ขนฺธา.\csromanspacer{} มหาภูเต ปจฺจยา จิตฺตสมุฏฺฐานํ รูปํ.\csromanspacer{} วิจิกิจฺฉาสหคตํ อุทฺธจฺจสหคตํ โมหํ ปจฺจยา สมฺปยุตฺตกา ขนฺธา จิตฺตสมุฏฺฐานญฺจ รูปํ.\csromanspacer{} ปฏิสนฺธิกฺขเณ วตฺถุํ ฯเปฯ. (3)}

\proseitem{73}{สเหตุกญฺจ อเหตุกญฺจ ธมฺมํ ปจฺจยา สเหตุโก ธมฺโม อุปฺปชฺชติ เหตุปจฺจยา.\csromanspacer{} สเหตุกํ เอกํ ขนฺธญฺจ วตฺถุญฺจ ปจฺจยา ตโย ขนฺธา ฯเปฯ เทฺว ขนฺเธ ฯเปฯ.\csromanspacer{} วิจิกิจฺฉาสหคตํ อุทฺธจฺจสหคตํ เอกํ ขนฺธญฺจ โมหญฺจ ปจฺจยา ตโย ขนฺธา ฯเปฯ เทฺว ขนฺเธ ฯเปฯ.\csromanspacer{} ปฏิสนฺธิกฺขเณ ฯเปฯ. (1)}

\prose{สเหตุกญฺจ อเหตุกญฺจ ธมฺมํ ปจฺจยา อเหตุโก ธมฺโม อุปฺปชฺชติ เหตุปจฺจยา.\csromanspacer{} สเหตุเก ขนฺเธ จ มหาภูเต จ ปจฺจยา จิตฺตสมุฏฺฐานํ รูปํ.\csromanspacer{} วิจิกิจฺฉาสหคเต อุทฺธจฺจสหคเต ขนฺเธ จ โมหญฺจ ปจฺจยา จิตฺตสมุฏฺฐานํ รูปํ.\csromanspacer{} ปฏิสนฺธิกฺขเณ ฯเปฯ. (2)}

\prose{สเหตุกญฺจ อเหตุกญฺจ ธมฺมํ ปจฺจยา สเหตุโก จ อเหตุโก จ ธมฺมา อุปฺปชฺชนฺติ เหตุปจฺจยา.\csromanspacer{} สเหตุกํ เอกํ ขนฺธญฺจ วตฺถุญฺจ ปจฺจยา ตโย ขนฺธา ฯเปฯ เทฺว ขนฺเธ ฯเปฯ.\csromanspacer{} สเหตุเก ขนฺเธ จ มหาภูเต จ ปจฺจยา จิตฺตสมุฏฺฐานํ รูปํ.\csromanspacer{} วิจิกิจฺฉาสหคตํ อุทฺธจฺจสหคตํ เอกํ ขนฺธญฺจ}

\vspace{\tipitakaparskip}
\csromancenter{สเหตุทุก โมหญฺจ ปจฺจยา ตโย ขนฺธา จิตฺตสมุฏฺฐานญฺจ รูปํ ฯเปฯ เทฺว ขนฺเธ ฯเปฯ.\csromanspacer{} ปฏิสนฺธิกฺขเณ สเหตุกํ เอกํ ขนฺธญฺจ วตฺถุญฺจ ปจฺจยา ตโย ขนฺธา ฯเปฯ เทฺว ขนฺเธ ฯเปฯ.\csromanspacer{} สเหตุเก ขนฺเธ จ มหาภูเต จ ปจฺจยา กฏตฺตารูปํ. (3)}
\csromanheader{2}{\textbf{อารมฺมณ}}
\proseitem{74}{\csromanfolio{31}สเหตุกํ ธมฺมํ ปจฺจยา สเหตุโก ธมฺโม อุปฺปชฺชติ อารมฺมณปจฺจยา.\csromanspacer{} สเหตุกํ เอกํ ขนฺธํ ปจฺจยา ตโย ขนฺธา ฯเปฯ.\csromanspacer{} ปฏิสนฺธิกฺขเณ ฯเปฯ. (1)}

\prose{สเหตุกํ ธมฺมํ ปจฺจยา อเหตุโก ธมฺโม อุปฺปชฺชติ อารมฺมณปจฺจยา.\csromanspacer{} วิจิกิจฺฉาสหคเต อุทฺธจฺจสหคเต ขนฺเธ ปจฺจยา วิจิกิจฺฉาสหคโต อุทฺธจฺจสหคโต โมโห. (2)}

\prose{สเหตุกํ ธมฺมํ ปจฺจยา สเหตุโก จ อเหตุโก จ ธมฺมา อุปฺปชฺชนฺติ อารมฺมณปจฺจยา.\csromanspacer{} วิจิกิจฺฉาสหคตํ อุทฺธจฺจสหคตํ เอกํ ขนฺธํ ปจฺจยา ตโย ขนฺธา โมโห จ ฯเปฯ เทฺว ขนฺเธ ฯเปฯ. (3)}

\proseitem{75}{อเหตุกํ ธมฺมํ ปจฺจยา อเหตุโก ธมฺโม อุปฺปชฺชติ อารมฺมณปจฺจยา.\csromanspacer{} อเหตุกํ เอกํ ขนฺธํ ฯเปฯ เทฺว ขนฺเธ ฯเปฯ.\csromanspacer{} ปฏิสนฺธิกฺขเณ วตฺถุํ ปจฺจยา ขนฺธา.\csromanspacer{} จกฺขายตนํ ปจฺจยา จกฺขุวิญฺญาณํ ฯเปฯ.\csromanspacer{} กายายตนํ ปจฺจยา กายวิญฺญาณํ.\csromanspacer{} วตฺถุํ ปจฺจยา อเหตุกา ขนฺธา. (1)}

\prose{อเหตุกํ ธมฺมํ ปจฺจยา สเหตุโก ธมฺโม อุปฺปชฺชติ อารมฺมณปจฺจยา.\csromanspacer{} วตฺถุํ ปจฺจยา สเหตุกา ขนฺธา.\csromanspacer{} วิจิกิจฺฉาสหคตํ อุทฺธจฺจสหคตํ โมหํ ปจฺจยา สมฺปยุตฺตกา ขนฺธา.\csromanspacer{} ปฏิสนฺธิกฺขเณ ฯเปฯ. (2)}

\prose{อเหตุกํ ธมฺมํ ปจฺจยา สเหตุโก จ อเหตุโก จ ธมฺมา อุปฺปชฺชนฺติ อารมฺมณปจฺจยา.\csromanspacer{} วตฺถุํ ปจฺจยา วิจิกิจฺฉาสหคตา อุทฺธจฺจสหคตา ขนฺธา โมโห จ. (3)}

\proseitem{76}{สเหตุกญฺจ อเหตุกญฺจ ธมฺมํ ปจฺจยา สเหตุโก ธมฺโม อุปฺปชฺชติ อารมฺมณปจฺจยา.\csromanspacer{} สเหตุกํ เอกํ ขนฺธญฺจ วตฺถุญฺจ ปจฺจยา ตโย ขนฺธา ฯเปฯ.\csromanspacer{} วิจิกิจฺฉาสหคตํ อุทฺธจฺจสหคตํ เอกํ ขนฺธญฺจ โมหญฺจ ปจฺจยา ตโย ขนฺธา ฯเปฯ เทฺว ขนฺเธ ฯเปฯ.\csromanspacer{} ปฏิสนฺธิกฺขเณ ฯเปฯ. (1)}

\prose{\csromanfolio{32}สเหตุกญฺจ อเหตุกญฺจ ธมฺมํ ปจฺจยา อเหตุโก ธมฺโม อุปฺปชฺชติ อารมฺมณปจฺจยา.\csromanspacer{} วิจิกิจฺฉาสหคเต อุทฺธจฺจสหคเต ขนฺเธ จ วตฺถุญฺจ ปจฺจยา วิจิกิจฺฉาสหคโต อุทฺธจฺจสหคโต โมโห. (2)}

\prose{สเหตุกญฺจ อเหตุกญฺจ ธมฺมํ ปจฺจยา สเหตุโก จ อเหตุโก จ ธมฺมา อุปฺปชฺชนฺติ อารมฺมณปจฺจยา.\csromanspacer{} วิจิกิจฺฉาสหคตํ อุทฺธจฺจสหคตํ เอกํ ขนฺธญฺจ วตฺถุญฺจ ปจฺจยา ตโย ขนฺธา โมโห จ ฯเปฯ เทฺว ขนฺเธ ฯเปฯ. (3)}

\csromanheader{2}{\textbf{อธิปติ}}
\proseitem{77}{สเหตุกํ ธมฺมํ ปจฺจยา สเหตุโก ธมฺโม อุปฺปชฺชติ อธิปติ ปจฺจยา. (\textbf{อธิปติ}ยา นว ปญฺหา ปวตฺเตเยว.)}

\csromanheader{2}{\textbf{อนนฺตราทิ}}
\proseitem{78}{สเหตุกํ ธมฺมํ ปจฺจยา สเหตุโก ธมฺโม อุปฺปชฺชติ อนนฺตรปจฺจยา.\csromanspacer{} สมนนฺตรปจฺจยา.\csromanspacer{} สหชาตปจฺจยา.\csromanspacer{} ตีณิ. (ปฏิจฺจวารสทิสา.)}

\prose{อเหตุกํ ธมฺมํ ปจฺจยา อเหตุโก ธมฺโม อุปฺปชฺชติ สหชาตปจฺจยา.\csromanspacer{} อเหตุกํ เอกํ ขนฺธํ ปจฺจยา ตโย ขนฺธา จิตฺตสมุฏฺฐานญฺจ รูปํ ฯเปฯ เทฺว ขนฺเธ ฯเปฯ.\csromanspacer{} ปฏิสนฺธิกฺขเณ ฯเปฯ. (ยาว อสญฺญสตฺตา.) จกฺขายตนํ ปจฺจยา จกฺขุวิญฺญาณํ ฯเปฯ.\csromanspacer{} กายายตนํ ปจฺจยา กายวิญฺญาณํ.\csromanspacer{} วตฺถุํ ปจฺจยา อเหตุกา ขนฺธา. (1)}

\prose{อเหตุกํ ธมฺมํ ปจฺจยา สเหตุโก ธมฺโม อุปฺปชฺชติ สหชาตปจฺจยา.\csromanspacer{} วตฺถุํ ปจฺจยา สเหตุกา ขนฺธา.\csromanspacer{} วิจิกิจฺฉาสหคตํ อุทฺธจฺจสหคตํ โมหํ ปจฺจยา สมฺปยุตฺตกา ขนฺธา.\csromanspacer{} ปฏิสนฺธิกฺขเณ ฯเปฯ. (2)}

\prose{อเหตุกํ ธมฺมํ ปจฺจยา สเหตุโก จ อเหตุโก จ ธมฺมา อุปฺปชฺชนฺติ สหชาตปจฺจยา.\csromanspacer{} วตฺถุํ ปจฺจยา สเหตุกา ขนฺธา.\csromanspacer{} มหาภูเต ปจฺจยา จิตฺตสมุฏฺฐานํ รูปํ.\csromanspacer{} วิจิกิจฺฉาสหคตํ อุทฺธจฺจสหคตํ โมหํ ปจฺจยา สมฺปยุตฺตกา ขนฺธา จิตฺตสมุฏฺฐานญฺจ รูปํ.\csromanspacer{} ปฏิสนฺธิกฺขเณ วตฺถุํ ฯเปฯ. (3)}

\csromanheader{2}{2. สเหตุทุก}
\proseitem{79}{\csromanfolio{33}สเหตุกญฺจ อเหตุกญฺจ ธมฺมํ ปจฺจยา สเหตุโก ธมฺโม อุปฺปชฺชติ สหชาตปจฺจยา.\csromanspacer{} สเหตุกํ เอกํ ขนฺธญฺจ วตฺถุญฺจ ปจฺจยา ตโย ขนฺธา ฯเปฯ เทฺว ขนฺเธ ฯเปฯ.\csromanspacer{} วิจิกิจฺฉาสหคตํ อุทฺธจฺจสหคตํ เอกํ ขนฺธญฺจ โมหญฺจ ปจฺจยา ตโย ขนฺธา ฯเปฯ เทฺว ขนฺเธ ฯเปฯ.\csromanspacer{} ปฏิสนฺธิกฺขเณ ฯเปฯ. (1)}

\prose{สเหตุกญฺจ อเหตุกญฺจ ธมฺมํ ปจฺจยา อเหตุโก ธมฺโม อุปฺปชฺชติ สหชาตปจฺจยา.\csromanspacer{} สเหตุเก ขนฺเธ จ มหาภูเต จ ปจฺจยา จิตฺตสมุฏฺฐานํ รูปํ.\csromanspacer{} วิจิกิจฺฉาสหคเต อุทฺธจฺจสหคเต ขนฺเธ จ โมหญฺจ ปจฺจยา จิตฺตสมุฏฺฐานํ รูปํ.\csromanspacer{} ปฏิสนฺธิกฺขเณ ฯเปฯ.\csromanspacer{} วิจิกิจฺฉาสหคเต อุทฺธจฺจสหคเต ขนฺเธ จ วตฺถุญฺจ ปจฺจยา วิจิกิจฺฉาสหคโต อุทฺธจฺจสหคโต โมโห. (2)}

\prose{สเหตุกญฺจ อเหตุกญฺจ ธมฺมํ ปจฺจยา สเหตุโก จ อเหตุโก จ ธมฺมา อุปฺปชฺชนฺติ สหชาตปจฺจยา.\csromanspacer{} สเหตุกํ เอกํ ขนฺธญฺจ วตฺถุญฺจ ปจฺจยา ตโย ขนฺธา ฯเปฯ เทฺว ขนฺเธ ฯเปฯ.\csromanspacer{} สเหตุเก ขนฺเธ จ มหาภูเต จ ปจฺจยา จิตฺตสมุฏฺฐานํ รูปํ.\csromanspacer{} วิจิกิจฺฉาสหคตํ อุทฺธจฺจสหคตํ เอกํ ขนฺธญฺจ วตฺถุญฺจ ปจฺจยา ตโย ขนฺธา ฯเปฯ เทฺว ขนฺเธ ฯเปฯ.\csromanspacer{} สเหตุเก ขนฺเธ จ มหาภูเต จ ปจฺจยา จิตฺตสมุฏฺฐานํ รูปํ.\csromanspacer{} วิจิกิจฺฉาสหคตํ อุทฺธจฺจสหคตํ เอกํ ขนฺธญฺจ โมหญฺจ ปจฺจยา ตโย ขนฺธา จิตฺตสมุฏฺฐานญฺจ รูปํ ฯเปฯ เทฺว ขนฺเธ ฯเปฯ.\csromanspacer{} วิจิกิจฺฉาสหคตํ อุทฺธจฺจสหคตํ เอกํ ขนฺธญฺจ วตฺถุญฺจ ปจฺจยา ตโย ขนฺธา โมโห จ ฯเปฯ.\csromanspacer{} ปฏิสนฺธิกฺขเณ สเหตุกํ เอกํ ขนฺธญฺจ วตฺถุญฺจ ปจฺจยา ตโย ขนฺธา ฯเปฯ เทฺว ขนฺเธ ฯเปฯ.\csromanspacer{} สเหตุเก ขนฺเธ จ มหาภูเต จ ปจฺจยา กฏตฺตารูปํ. (3)}

\csromanheader{2}{\textbf{อญฺญมญฺญาทิ}}
\proseitem{80}{สเหตุกํ ธมฺมํ ปจฺจยา สเหตุโก ธมฺโม อุปฺปชฺชติ อญฺญมญฺญปจฺจยา ฯเปฯ อวิคตปจฺจยา.}

\csromanheader{2}{1. ปจฺจยานุโลม 2. สงฺขฺยาวาร}
\proseitem{81}{เหตุยา นว, อารมฺมเณ นว, อธิปติยา นว, (สพฺพตฺถ นว,) อวิคเต นว. (เอวํ คเณตพฺพํ.)}

\vspace{\tipitakaparskip}
\csromancenter{อนุโลมํ.}
\csromanheader{2}{2. ปจฺจยปจฺจนีย 1. วิภงฺควาร}
\csromanheader{2}{\textbf{นเหตุ}}
\proseitem{82}{\csromanfolio{34}สเหตุกํ ธมฺมํ ปจฺจยา อเหตุโก ธมฺโม อุปฺปชฺชติ นเหตุปจฺจยา.\csromanspacer{} วิจิกิจฺฉาสหคเต อุทฺธจฺจสหคเต ขนฺเธ ปจฺจยา วิจิกิจฺฉาสหคโต อุทฺธจฺจสหคโต โมโห. (1)}

\prose{อเหตุกํ ธมฺมํ ปจฺจยา อเหตุโก ธมฺโม อุปฺปชฺชติ นเหตุปจฺจยา.\csromanspacer{} อเหตุกํ เอกํ ขนฺธํ ฯเปฯ.\csromanspacer{} ปฏิสนฺธิกฺขเณ ฯเปฯ. (ยาว อสญฺญสตฺตา.) จกฺขายตนํ ปจฺจยา จกฺขุวิญฺญาณํ ฯเปฯ.\csromanspacer{} กายายตนํ ปจฺจยา กายวิญฺญาณํ.\csromanspacer{} วตฺถุํ ปจฺจยา อเหตุกา ขนฺธา โมโห จ. (1)}

\prose{สเหตุกญฺจ อเหตุกญฺจ ธมฺมํ ปจฺจยา อเหตุโก ธมฺโม อุปฺปชฺชติ นเหตุปจฺจยา.\csromanspacer{} วิจิกิจฺฉาสหคเต อุทฺธจฺจสหคเต ขนฺเธ จ วตฺถุญฺจ ปจฺจยา วิจิกิจฺฉาสหคโต อุทฺธจฺจสหคโต โมโห. (สํขิตฺตํ.) (1)}

\csromanheader{2}{2. ปจฺจยปจฺจนีย 2. สงฺขฺยาวาร}
\csromanheader{2}{\textbf{สุทฺธ}}
\vspace{\tipitakaparskip}
\csromancenter{\textbf{นเหตุ}ยา ตีณิ, น-อารมฺมเณ ตีณิ, น-อธิปติยา นว, น-อนนฺตเร ตีณิ, นสมนนฺตเร ตีณิ, น-อญฺญมญฺเญ ตีณิ, นนิสฺสเย ตีณิ, น-อุปนิสฺสเย ตีณิ, นปุเรชาเต นว, นปจฺฉาชาเต นว, น-อาเสวเน นว, นกมฺเม จตฺตาริ, นวิปาเก นว, น-อาหาเร เอกํ, น-อินฺทฺริเย เอกํ, นฌาเน เอกํ, นมคฺเค เอกํ, นสมฺปยุตฺเต ตีณิ, นวิปฺปยุตฺเต ฉ, โนนตฺถิยา ตีณิ, โนวิคเต ตีณิ. (เอวํ คเณตพฺพํ.)}
\csromancenter{ปจฺจนียํ.}
\csromanheader{2}{3. ปจฺจยานุโลมปจฺจนีย}
\csromanheader{2}{\textbf{เหตุทุก}}
\proseitem{84}{เหตุปจฺจยา น-อารมฺมเณ ตีณิ, น-อธิปติยา นว, น-อนนฺตเร ตีณิ, นสมนนฺตเร ตีณิ, น-อญฺญมญฺเญ ตีณิ, น-อุปนิสฺสเย ตีณิ,}

\vspace{\tipitakaparskip}
\csromancenter{สเหตุทุก นปุเรชาเต นว, นปจฺฉาชาเต นว, น-อาเสวเน นว, นกมฺเม ตีณิ, นวิปาเก นว, นสมฺปยุตฺเต ตีณิ, นวิปฺปยุตฺเต ตีณิ, โนนตฺถิยา ตีณิ, โนวิคเต ตีณิ. (เอวํ คเณตพฺพํ.)}
\csromancenter{อนุโลมปจฺจนียํ.}
\csromanheader{2}{4. ปจฺจยปจฺจนียานุโลม}
\csromanheader{2}{\textbf{นเหตุทุก}}
\proseitem{85}{\csromanfolio{35}นเหตุปจฺจยา อารมฺมเณ ตีณิ, อนนฺตเร ตีณิ ฯเปฯ มคฺเค ตีณิ, สมฺปยุตฺเต ตีณิ ฯเปฯ อวิคเต ตีณิ. (เอวํ คเณตพฺพํ.)}

\vspace{\tipitakaparskip}
\csromancenter{ปจฺจนียานุโลมํ.}
\csromancenter{(นิสฺสยวาโร ปจฺจยวารสทิโส.)}
\csromanheader{2}{2. สเหตุกทุก 5. สํสฏฺฐวาร}
\csromanheader{2}{1. ปจฺจยานุโลม 1. วิภงฺควาร}
\csromanheader{2}{\textbf{เหตุ}}
\proseitem{86}{สเหตุกํ ธมฺมํ สํสฏฺโฐ สเหตุโก ธมฺโม อุปฺปชฺชติ เหตุปจฺจยา.\csromanspacer{} สเหตุกํ เอกํ ขนฺธํ สํสฏฺฐา ตโย ขนฺธา ฯเปฯ เทฺว ขนฺเธ ฯเปฯ.\csromanspacer{} ปฏิสนฺธิกฺขเณ ฯเปฯ. (1)}

\prose{อเหตุกํ ธมฺมํ สํสฏฺโฐ สเหตุโก ธมฺโม อุปฺปชฺชติ เหตุปจฺจยา.\csromanspacer{} วิจิกิจฺฉาสหคตํ อุทฺธจฺจสหคตํ โมหํ สํสฏฺฐา วิจิกิจฺฉาสหคตา อุทฺธจฺจสหคตา ขนฺธา. (1)}

\prose{สเหตุกญฺจ อเหตุกญฺจ ธมฺมํ สํสฏฺโฐ สเหตุโก ธมฺโม อุปฺปชฺชติ เหตุปจฺจยา.\csromanspacer{} วิจิกิจฺฉาสหคตํ อุทฺธจฺจสหคตํ เอกํ ขนฺธญฺจ โมหญฺจ สํสฏฺฐา ตโย ขนฺธา ฯเปฯ เทฺว ขนฺเธ ฯเปฯ. (1)}

\csromanheader{2}{\textbf{อารมฺมณ}}
\proseitem{87}{\csromanfolio{36}สเหตุกํ ธมฺมํ สํสฏฺโฐ อเหตุโก ธมฺโม อุปฺปชฺชติ.\csromanspacer{} \textbf{อารมฺมณ}ปจฺจยา.\csromanspacer{} สเหตุกํ เอกํ ขนฺธํ สํสฏฺฐา ตโย ขนฺธา ฯเปฯ เทฺว ขนฺเธ ฯเปฯ.\csromanspacer{} ปฏิสนฺธิกฺขเณ ฯเปฯ. (1)}

\prose{สเหตุกํ ธมฺมํ สํสฏฺโฐ อเหตุโก ธมฺโม อุปฺปชฺชติ อารมฺมณปจฺจยา.\csromanspacer{} วิจิกิจฺฉาสหคเต อุทฺธจฺจสหคเต ขนฺเธ สํสฏฺโฐ วิจิกิจฺฉาสหคโต อุทฺธจฺจสหคโต โมโห. (2)}

\prose{สเหตุกํ ธมฺมํ สํสฏฺโฐ สเหตุโก จ อเหตุโก จ ธมฺมา อุปฺปชฺชนฺติ อารมฺมณปจฺจยา.\csromanspacer{} วิจิกิจฺฉาสหคตํ อุทฺธจฺจสหคตํ เอกํ ขนฺธํ สํสฏฺฐา ตโย ขนฺธา โมโห จ ฯเปฯ เทฺว ขนฺเธ ฯเปฯ. (3)}

\proseitem{88}{อเหตุกํ ธมฺมํ สํสฏฺโฐ อเหตุโก ธมฺโม อุปฺปชฺชติ อารมฺมณปจฺจยา.\csromanspacer{} อเหตุกํ เอกํ ขนฺธํ สํสฏฺฐา ตโย ขนฺธา ฯเปฯ เทฺว ขนฺเธ ฯเปฯ.\csromanspacer{} ปฏิสนฺธิกฺขเณ ฯเปฯ. (1)}

\prose{อเหตุกํ ธมฺมํ สํสฏฺโฐ สเหตุโก ธมฺโม อุปฺปชฺชติ อารมฺมณปจฺจยา.\csromanspacer{} วิจิกิจฺฉาสหคตํ อุทฺธจฺจสหคตํ โมหํ สํสฏฺฐา วิจิกิจฺฉาสหคตา อุทฺธจฺจสหคตา ขนฺธา. (2)}

\prose{สเหตุกญฺจ อเหตุกญฺจ ธมฺมํ สํสฏฺโฐ สเหตุโก ธมฺโม อุปฺปชฺชติ อารมฺมณปจฺจยา.\csromanspacer{} วิจิกิจฺฉาสหคตํ อุทฺธจฺจสหคตํ เอกํ ขนฺธญฺจ โมหญฺจ สํสฏฺฐา ตโย ขนฺธา ฯเปฯ เทฺว ขนฺเธ ฯเปฯ. (1)}

\csromanheader{2}{\textbf{อธิปติ}}
\proseitem{89}{สเหตุกํ ธมฺมํ สํสฏฺโฐ สเหตุโก ธมฺโม อุปฺปชฺชติ อธิปติปจฺจยา.\csromanspacer{} สเหตุกํ เอกํ ขนฺธํ สํสฏฺฐา ตโย ขนฺธา ฯเปฯ เทฺว ขนฺเธ ฯเปฯ. (1)}

\csromanheader{2}{\textbf{อนนฺตราทิ}}
\proseitem{90}{สเหตุกํ ธมฺมํ สํสฏฺโฐ สเหตุโก ธมฺโม อุปฺปชฺชติ อนนฺตรปจฺจยา.\csromanspacer{} สมนนฺตรปจฺจยา.\csromanspacer{} สหชาตปจฺจยา ฯเปฯ วิปากปจฺจยา.\csromanspacer{} วิปากํ สเหตุกํ เอกํ ขนฺธํ สํสฏฺฐา ตโย ขนฺธา ฯเปฯ เทฺว ขนฺเธ ฯเปฯ.\csromanspacer{} ปฏิสนฺธิกฺขเณ ฯเปฯ.}

\csromanheader{2}{2. สเหตุทุก}
\prose{\csromanfolio{37}อเหตุกํ ธมฺมํ สํสฏฺโฐ อเหตุโก ธมฺโม อุปฺปชฺชติ วิปากปจฺจยา.\csromanspacer{} วิปากํ อเหตุกํ เอกํ ขนฺธํ สํสฏฺฐา ตโย ขนฺธา ฯเปฯ เทฺว ขนฺเธ ฯเปฯ.\csromanspacer{} ปฏิสนฺธิกฺขเณ ฯเปฯ ฌานปจฺจยา ฯเปฯ อวิคตปจฺจยา.}

\csromanheader{2}{1. ปจฺจยานุโลม 2. สงฺขฺยาวาร}
\csromanheader{2}{\textbf{สุทฺธ}}
\proseitem{91}{เหตุยา ตีณิ, อารมฺมเณ ฉ, อธิปติยา เอกํ, อนนฺตเร ฉ, สมนนฺตเร ฉ, สหชาเต ฉ, อญฺญมญฺเญ ฉ, นิสฺสเย ฉ, อุปนิสฺสเย ฉ, ปุเรชาเต ฉ ฯเปฯ วิปาเก เทฺว, อาหาเร ฉ, อินฺทฺริเย ฉ, ฌาเน ฉ, มคฺเค ปญฺจ ฯเปฯ อวิคเต ฉ. (เอวํ คเณตพฺพํ.)}

\vspace{\tipitakaparskip}
\csromancenter{อนุโลมํ.}
\csromanheader{2}{2. ปจฺจยปจฺจนีย 1. วิภงฺควาร}
\csromanheader{2}{\textbf{นเหตุ}}
\proseitem{92}{สเหตุกํ ธมฺมํ สํสฏฺโฐ อเหตุโก ธมฺโม อุปฺปชฺชติ นเหตุปจฺจยา.\csromanspacer{} วิจิกิจฺฉาสหคเต อุทฺธจฺจสหคเต ขนฺเธ สํสฏฺโฐ วิจิกิจฺฉาสหคโต อุทฺธจฺจสหคโต โมโห. (1)}

\prose{อเหตุกํ ธมฺมํ สํสฏฺโฐ อเหตุโก ธมฺโม อุปฺปชฺชติ นเหตุปจฺจยา.\csromanspacer{} อเหตุกํ เอกํ ขนฺธํ สํสฏฺฐา ตโย ขนฺธา ฯเปฯ เทฺว ขนฺเธ ฯเปฯ.\csromanspacer{} ปฏิสนฺธิกฺขเณ ฯเปฯ. (สํขิตฺตํ.) (1)}

\csromanheader{2}{2. ปจฺจยปจฺจนีย 2. สงฺขฺยาวาร}
\csromanheader{2}{\textbf{สุทฺธ}}
\vspace{\tipitakaparskip}
\csromancenter{\textbf{นเหตุ}ยา เทฺว, น-อธิปติยา ฉ, นปุเรชาเต ฉ, นปจฺฉาชาเต ฉ, น-อาเสวเน ฉ, นกมฺเม จตฺตาริ, นวิปาเก ฉ, นฌาเน เอกํ, นมคฺเค เอกํ, นวิปฺปยุตฺเต ฉ. (เอวํ คเณตพฺพํ.)}
\csromancenter{ปจฺจนียํ.}
\csromanheader{2}{3. ปจฺจยานุโลมปจฺจนีย}
\csromanheader{2}{\textbf{เหตุทุก}}
\proseitem{94}{\csromanfolio{38}\textbf{เหตุ}ปจฺจยา น-อธิปติยา ตีณิ, นปุเรชาเต ตีณิ, นปจฺฉาชาเต ตีณิ, น-อาเสวเน ตีณิ, นกมฺเม ตีณิ, นวิปาเก ตีณิ, นวิปฺปยุตฺเต ตีณิ.}

\vspace{\tipitakaparskip}
\csromancenter{อนุโลมปจฺจนียํ.}
\csromanheader{2}{4. ปจฺจยปจฺจนียานุโลม}
\csromanheader{2}{\textbf{นเหตุทุก}}
\proseitem{95}{นเหตุปจฺจยา อารมฺมเณ เทฺว, อนนฺตเร เทฺว ฯเปฯ กมฺเม เทฺว, วิปาเก เอกํ, อาหาเร เทฺว ฯเปฯ มคฺเค เอกํ ฯเปฯ อวิคเต เทฺว. (เอวํ คเณตพฺพํ)}

\vspace{\tipitakaparskip}
\csromancenter{ปจฺจนียานุโลมํ.}
\csromancenter{(สมฺปยุตฺตวาโร สํสฏฺฐวารสทิโส.)}
\csromanheader{2}{2. สเหตุกทุก 7. ปญฺหาวาร}
\csromanheader{2}{1. ปจฺจยานุโลม 1. วิภงฺควาร}
\csromanheader{2}{\textbf{เหตุ}}
\proseitem{96}{สเหตุโก ธมฺโม สเหตุกสฺส ธมฺมสฺส เหตุปจฺจเยน ปจฺจโย.\csromanspacer{} สเหตุกา เหตู สมฺปยุตฺตกานํ ขนฺธานํ เหตุปจฺจเยน ปจฺจโย.\csromanspacer{} ปฏิสนฺธิกฺขเณ ฯเปฯ. (1)}

\prose{สเหตุโก ธมฺโม อเหตุกสฺส ธมฺมสฺส เหตุปจฺจเยน ปจฺจโย.\csromanspacer{} สเหตุกา เหตู จิตฺตสมุฏฺฐานานํ รูปานํ เหตุปจฺจเยน ปจฺจโย.\csromanspacer{} ปฏิสนฺธิกฺขเณ ฯเปฯ. (2)}

\csromanheader{2}{2. สเหตุทุก}
\prose{\csromanfolio{39}สเหตุโก ธมฺโม สเหตุกสฺส จ อเหตุกสฺส จ ธมฺมสฺส เหตุปจฺจเยน ปจฺจโย.\csromanspacer{} สเหตุกา เหตู สมฺปยุตฺตกานํ ขนฺธานํ จิตฺตสมุฏฺฐานานญฺจ รูปานํ เหตุปจฺจเยน ปจฺจโย.\csromanspacer{} ปฏิสนฺธิกฺขเณ ฯเปฯ. (3)}

\proseitem{97}{อเหตุโก ธมฺโม อเหตุกสฺส ธมฺมสฺส เหตุปจฺจเยน ปจฺจโย.\csromanspacer{} วิจิกิจฺฉาสหคโต อุทฺธจฺจสหคโต โมโห จิตฺตสมุฏฺฐานานํ รูปานํ เหตุปจฺจเยน ปจฺจโย. (1)}

\prose{อเหตุโก ธมฺโม สเหตุกสฺส ธมฺมสฺส เหตุปจฺจเยน ปจฺจโย.\csromanspacer{} วิจิกิจฺฉาสหคโต อุทฺธจฺจสหคโต โมโห สมฺปยุตฺตกานํ ขนฺธานํ เหตุปจฺจเยน ปจฺจโย. (2)}

\prose{อเหตุโก ธมฺโม สเหตุกสฺส จ อเหตุกสฺส จ ธมฺมสฺส เหตุปจฺจเยน ปจฺจโย.\csromanspacer{} วิจิกิจฺฉาสหคโต อุทฺธจฺจสหคโต โมโห สมฺปยุตฺตกานํ ขนฺธานํ จิตฺตสมุฏฺฐานญฺจ รูปานํ เหตุปจฺจเยน ปจฺจโย. (3)}

\csromanheader{2}{\textbf{อารมฺมณ}}
\proseitem{98}{สเหตุโก ธมฺโม สเหตุกสฺส ธมฺมสฺส อารมฺมณปจฺจเยน ปจฺจโย.\csromanspacer{} ทานํ ทตฺวา, สีลํ ฯเปฯ อุโปสถกมฺมํ กตฺวา ตํ ปจฺจเวกฺขติ, ปุพฺเพ สุจิณฺณานิ ปจฺจเวกฺขติ.\csromanspacer{} ฌานา วุฏฺฐหิตฺวา ฌานํ ปจฺจเวกฺขติ.\csromanspacer{} อริยา มคฺคา วุฏฺฐหิตฺวา มคฺคํ ปจฺจเวกฺขนฺติ, ผลํ ปจฺจเวกฺขนฺติ.\csromanspacer{} ปหีเน กิเลเส ฯเปฯ.\csromanspacer{} วิกฺขมฺภิเต กิเลเส ฯเปฯ.\csromanspacer{} ปุพฺเพ ฯเปฯ.\csromanspacer{} สเหตุเก ขนฺเธ อนิจฺจโต ฯเปฯ โทมนสฺสํ อุปฺปชฺชติ.\csromanspacer{} กุสลากุสเล นิรุทฺเธ สเหตุโก วิปาโก ตทารมฺมณตา อุปฺปชฺชติ.\csromanspacer{} เจโตปริยญาเณน สเหตุกจิตฺตสมงฺคิสฺส จิตฺตํ ชานนฺติ.\csromanspacer{} อากาสานญฺจายตนํ วิญฺญาณญฺจายตนสฺส ฯเปฯ.\csromanspacer{} อากิญฺจญฺญายตนํ เนวสญฺญานาสญฺญายตนสฺส ฯเปฯ.\csromanspacer{} สเหตุกา ขนฺธา อิทฺธิวิธญาณสฺส, เจโตปริยญาณสฺส, ปุพฺเพนิวาสานุสฺสติญาณสฺส, ยถากมฺมูปคญาณสฺส, อนาคตํสญาณสฺส อารมฺมณปจฺจเยน ปจฺจโย.\csromanspacer{} สเหตุเก ขนฺเธ อารพฺภ สเหตุกา ขนฺธา อุปฺปชฺชนฺติ. (1)}

\prose{สเหตุโก ธมฺโม อเหตุกสฺส ธมฺมสฺส อารมฺมณปจฺจเยน ปจฺจโย.\csromanspacer{} สเหตุเก ขนฺเธ อนิจฺจโต ฯเปฯ โทมนสฺสํ อุปฺปชฺชติ.\csromanspacer{} กุสลากุสเล นิรุทฺเธ อเหตุโก วิปาโก ตทารมฺมณตา อุปฺปชฺชติ.\csromanspacer{} สเหตุเก ขนฺเธ อารพฺภ อเหตุกา ขนฺธา จ โมโห จ อุปฺปชฺชนฺติ. (2)}

\prose{\csromanfolio{40}สเหตุโก ธมฺโม สเหตุกสฺส จ อเหตุกสฺส จ ธมฺมสฺส อารมฺมณปจฺจเยน ปจฺจโย.\csromanspacer{} สเหตุเก ขนฺเธ อารพฺภ วิจิกิจฺฉาสหคตา อุทฺธจฺจสหคตา ขนฺธา จ โมโห จ อุปฺปชฺชนฺติ. (3)}

\proseitem{99}{อเหตุโก ธมฺโม อเหตุกสฺส ธมฺมสฺส อารมฺมณปจฺจเยน ปจฺจโย.\csromanspacer{} นิพฺพานํ อาวชฺชนาย อารมฺมณปจฺจเยน ปจฺจโย.\csromanspacer{} จกฺขุํ ฯเปฯ.\csromanspacer{} วตฺถุํ.\csromanspacer{} อเหตุเก ขนฺเธ จ โมหญฺจ อนิจฺจโต ฯเปฯ โทมนสฺสํ อุปฺปชฺชติ.\csromanspacer{} กุสลากุสเล นิรุทฺเธ อเหตุโก วิปาโก ตทารมฺมณตา อุปฺปชฺชติ.\csromanspacer{} รูปายตนํ จกฺขุวิญฺญาณสฺส ฯเปฯ.\csromanspacer{} โผฏฺฐพฺพายตนํ กายวิญฺญาณสฺส ฯเปฯ.\csromanspacer{} อเหตุเก ขนฺเธ จ โมหญฺจ อารพฺภ อเหตุกา ขนฺธา จ โมโห จ อุปฺปชฺชนฺติ. (1)}

\prose{อเหตุโก ธมฺโม สเหตุกสฺส ธมฺมสฺส อารมฺมณปจฺจเยน ปจฺจโย.\csromanspacer{} อริยา นิพฺพานํ ปจฺจเวกฺขนฺติ.\csromanspacer{} นิพฺพานํ โคตฺรภุสฺส, โวทานสฺส, มคฺคสฺส, ผลสฺส อารมฺมณปจฺจเยน ปจฺจโย.\csromanspacer{} อริยา อเหตุเก ปหีเน กิเลเส ปจฺจเวกฺขนฺติ.\csromanspacer{} วิกฺขมฺภิเต กิเลเส ฯเปฯ.\csromanspacer{} ปุพฺเพ ฯเปฯ.\csromanspacer{} จกฺขุํ ฯเปฯ วตฺถุํ ฯเปฯ.\csromanspacer{} อเหตุเก ขนฺเธ จ โมหญฺจ อนิจฺจโต ฯเปฯ โทมนสฺสํ อุปฺปชฺชติ.\csromanspacer{} กุสลากุสเล นิรุทฺเธ สเหตุโก วิปาโก ตทารมฺมณตา อุปฺปชฺชติ.\csromanspacer{} ทิพฺเพน จกฺขุนา รูปํ ปสฺสติ.\csromanspacer{} ทิพฺพาย โสตธาตุยา สทฺทํ สุณาติ.\csromanspacer{} เจโตปริยญาเณน อเหตุกจิตฺตสมงฺคิสฺส จิตฺตํ ชานนฺติ.\csromanspacer{} อเหตุกา ขนฺธา อิทฺธิวิธญาณสฺส, เจโตปริยญาณสฺส, ปุพฺเพนิวาสานุสฺสติญาณสฺส, อนาคตํสญาณสฺส อารมฺมณปจฺจเยน ปจฺจโย.\csromanspacer{} อเหตุเก ขนฺเธ จ โมหญฺจ อารพฺภ สเหตุกา ขนฺธา อุปฺปชฺชนฺติ. (2)}

\prose{อเหตุโก ธมฺโม สเหตุกสฺส จ อเหตุกสฺส จ ธมฺมสฺส อารมฺมณปจฺจเยน ปจฺจโย.\csromanspacer{} จกฺขุํ อารพฺภ วิจิกิจฺฉาสหคตา อุทฺธจฺจสหคตา ขนฺธา จ โมโห จ อุปฺปชฺชนฺติ.\csromanspacer{} โสตํ ฯเปฯ.\csromanspacer{} วตฺถุํ.\csromanspacer{} อเหตุเก ขนฺเธ จ โมหญฺจ อารพฺภ วิจิกิจฺฉาสหคตา อุทฺธจฺจสหคตา ขนฺธา จ โมโห จ อุปฺปชฺชนฺติ. (3)}

\proseitem{100}{สเหตุโก จ อเหตุโก จ ธมฺมา สเหตุกสฺส ธมฺมสฺส อารมฺมณปจฺจเยน ปจฺจโย.\csromanspacer{} วิจิกิจฺฉาสหคเต อุทฺธจฺจสหคเต ขนฺเธ จ โมหญฺจ อารพฺภ สเหตุกา ขนฺธา อุปฺปชฺชนฺติ. (1)}

\csromanheader{2}{2. สเหตุทุก}
\prose{\csromanfolio{41}สเหตุโก จ อเหตุโก จ ธมฺมา อเหตุกสฺส ธมฺมสฺส อารมฺมณปจฺจเยน ปจฺจโย.\csromanspacer{} วิจิกิจฺฉาสหคเต อุทฺธจฺจสหคเต ขนฺเธ จ โมหญฺจ อารพฺภ อเหตุกา ขนฺธา จ โมโห จ อุปฺปชฺชนฺติ. (2)}

\prose{สเหตุโก จ อเหตุโก จ ธมฺมา สเหตุกสฺส จ อเหตุกสฺส จ ธมฺมสฺส อารมฺมณปจฺจเยน ปจฺจโย.\csromanspacer{} วิจิกิจฺฉาสหคเต อุทฺธจฺจสหคเต ขนฺเธ จ โมหญฺจ อารพฺภ วิจิกิจฺฉาสหคตา อุทฺธจฺจสหคตา ขนฺธา จ โมโห จ อุปฺปชฺชนฺติ. (3)}

\csromanheader{2}{\textbf{อธิปติ}}
\proseitem{101}{สเหตุโก ธมฺโม สเหตุกสฺส ธมฺมสฺส อธิปติปจฺจเยน ปจฺจโย.\csromanspacer{} \textbf{อารมฺมณาธิปติ}, สหชาธิปติ.\csromanspacer{}\csromanspacer{}\textbf{อารมฺมณาธิปติ}\csromandash{}ทานํ ทตฺวา, สีลํ ฯเปฯ อุโปสถกมฺมํ กตฺวา ตํ ครุํ กตฺวา ปจฺจเวกฺขติ.\csromanspacer{} ฌานา ฯเปฯ.\csromanspacer{} อริยา มคฺคา วุฏฺฐหิตฺวา มคฺคํ ครุํ กตฺวา ฯเปฯ.\csromanspacer{} ผลํ ฯเปฯ.\csromanspacer{} สเหตุเก ขนฺเธ ครุํ กตฺวา อสฺสาเทติ อภินนฺทติ, ตํ ครุํ กตฺวา ราโค อุปฺปชฺชติ, ทิฏฺฐิ อุปฺปชฺชติ.\csromanspacer{} \textbf{สหชาตาธิปติ}\csromandash{}สเหตุกาธิปติ สมฺปยุตฺตกานํ ขนฺธานํ อธิปติปจฺจเยน ปจฺจโย. (1)}

\prose{สเหตุโก ธมฺโม อเหตุกสฺส ธมฺมสฺส อธิปติปจฺจเยน ปจฺจโย.\csromanspacer{} \textbf{สหชาตาธิปติ}\csromandash{}สเหตุกาธิปติ จิตฺตสมุฏฺฐานานํ รูปานํ อธิปติปจฺจเยน ปจฺจโย. (2)}

\prose{สเหตุโก ธมฺโม สเหตุกสฺส จ อเหตุกสฺส จ ธมฺมสฺส อธิปติปจฺจเยน ปจฺจโย.\csromanspacer{} \textbf{สหชาตาธิปติ}\csromandash{}สเหตุกาธิปติ สมฺปยุตฺตกานํ ขนฺธานํ จิตฺตสมุฏฺฐานานญฺจ รูปานํ อธิปติปจฺจเยน ปจฺจโย. (3)}

\prose{อเหตุโก ธมฺโม สเหตุกสฺส ธมฺมสฺส อธิปติปจฺจเยน ปจฺจโย.\csromanspacer{} \textbf{อารมฺมณาธิปติ}\csromandash{}อริยา นิพฺพานํ ครุํ กตฺวา ปจฺจเวกฺขนฺติ.\csromanspacer{} นิพฺพานํ โคตฺรภุสฺส, โวทานสฺส, มคฺคสฺส, ผลสฺส อธิปติปจฺจเยน ปจฺจโย.\csromanspacer{} จกฺขุํ ฯเปฯ วตฺถุํ.\csromanspacer{} อเหตุเก ขนฺเธ ครุํ กตฺวา อสฺสาเทติ อภินนฺทติ, ตํ ครุํ กตฺวา ราโค อุปฺปชฺชติ, ทิฏฺฐิ อุปฺปชฺชติ. (1)}

\csromanheader{2}{\textbf{อนนฺตร}}
\proseitem{102}{\csromanfolio{42}สเหตุโก ธมฺโม สเหตุกสฺส ธมฺมสฺส อนนฺตรปจฺจเยน ปจฺจโย.\csromanspacer{} ปุริมา ปุริมา สเหตุกา ขนฺธา ปจฺฉิมานํ ปจฺฉิมานํ สเหตุกานํ ขนฺธานํ อนนฺตรปจฺจเยน ปจฺจโย.\csromanspacer{} อนุโลมํ โคตฺรภุสฺส.\csromanspacer{} อนุโลมํ โวทานสฺส.\csromanspacer{} โคตฺรภุ มคฺคสฺส.\csromanspacer{} โวทานํ มคฺคสฺส.\csromanspacer{} มคฺโค ผลสฺส.\csromanspacer{} ผลํ ผลสฺส.\csromanspacer{} อนุโลมํ ผลสมาปตฺติยา.\csromanspacer{} นิโรธา วุฏฺฐหนฺตสฺส เนวสญฺญานาสญฺญายตนํ ผลสมาปตฺติยา อนนฺตรปจฺจเยน ปจฺจโย. (1)}

\prose{สเหตุโก ธมฺโม อเหตุกสฺส ธมฺมสฺส อนนฺตรปจฺจเยน ปจฺจโย.\csromanspacer{} ปุริมา ปุริมา วิจิกิจฺฉาสหคตา อุทฺธจฺจสหคตา ขนฺธา ปจฺฉิมสฺส ปจฺฉิมสฺส วิจิกิจฺฉาสหคตสฺส อุทฺธจฺจสหคตสฺส โมหสฺส อนนฺตรปจฺจเยน ปจฺจโย.\csromanspacer{} สเหตุกํ จุติจิตฺตํ อเหตุกสฺส อุปปตฺติจิตฺตสฺส อนนฺตรปจฺจเยน ปจฺจโย.\csromanspacer{} สเหตุกํ ภวงฺคํ อาวชฺชนาย อนนฺตรปจฺจเยน ปจฺจโย.\csromanspacer{} สเหตุกํ ภวงฺคํ อเหตุกสฺส ภวงฺคสฺส อนนฺตรปจฺจเยน ปจฺจโย.\csromanspacer{} สเหตุกา ขนฺธา อเหตุกสฺส วุฏฺฐานสฺส อนนฺตรปจฺจเยน ปจฺจโย. (2)}

\prose{สเหตุโก ธมฺโม สเหตุกสฺส จ อเหตุกสฺส จ ธมฺมสฺส อนนฺตรปจฺจเยน ปจฺจโย.\csromanspacer{} ปุริมา ปุริมา วิจิกิจฺฉาสหคตา อุทฺธจฺจสหคตา ขนฺธา ปจฺฉิมานํ ปจฺฉิมานํ วิจิกิจฺฉาสหคตานํ อุทฺธจฺจสหคตานํ ขนฺธานํ โมหสฺส จ อนนฺตรปจฺจเยน ปจฺจโย. (3)}

\proseitem{103}{อเหตุโก ธมฺโม อเหตุกสฺส ธมฺมสฺส อนนฺตรปจฺจเยน ปจฺจโย.\csromanspacer{} ปุริโม ปุริโม วิจิกิจฺฉาสหคโต อุทฺธจฺจสหคโต โมโห ปจฺฉิมสฺส ปจฺฉิมสฺส วิจิกิจฺฉาสหคตสฺส อุทฺธจฺจสหคตสฺส โมหสฺส อนนฺตรปจฺจเยน ปจฺจโย.\csromanspacer{} ปุริมา ปุริมา อเหตุกา ขนฺธา ปจฺฉิมานํ ปจฺฉิมานํ อเหตุกานํ ขนฺธานํ อนนฺตรปจฺจเยน ปจฺจโย.\csromanspacer{} อาวชฺชนา ปญฺจนฺนํ วิญฺญาณานํ อนนฺตรปจฺจเยน ปจฺจโย. (1)}

\prose{อเหตุโก ธมฺโม สเหตุกสฺส ธมฺมสฺส อนนฺตรปจฺจเยน ปจฺจโย.\csromanspacer{} ปุริโม ปุริโม วิจิกิจฺฉาสหคโต อุทฺธจฺจสหคโต โมโห ปจฺฉิมานํ ปจฺฉิมานํ วิจิกิจฺฉาสหคตานํ อุทฺธจฺจสหคตานํ ขนฺธานํ อนนฺตรปจฺจเยน}

\vspace{\tipitakaparskip}
\csromancenter{สเหตุทุก ปจฺจโย.\csromanspacer{} อเหตุกํ จุติจิตฺตํ สเหตุกสฺส อุปปตฺติจิตฺตสฺส อนนฺตรปจฺจเยน ปจฺจโย.\csromanspacer{} อเหตุกํ ภวงฺคํ สเหตุกสฺส ภวงฺคสฺส อนนฺตรปจฺจเยน ปจฺจโย.\csromanspacer{} อาวชฺชนา สเหตุกานํ ขนฺธานํ อนนฺตรปจฺจเยน ปจฺจโย.\csromanspacer{} อเหตุกา ขนฺธา สเหตุกสฺส วุฏฺฐานสฺส อนนฺตรปจฺจเยน ปจฺจโย. (2)}
\prose{\csromanfolio{43}อเหตุโก ธมฺโม สเหตุกสฺส จ อเหตุกสฺส จ ธมฺมสฺส อนนฺตรปจฺจเยน ปจฺจโย.\csromanspacer{} ปุริโม ปุริโม วิจิกิจฺฉาสหคโต อุทฺธจฺจสหคโต โมโห ปจฺฉิมานํ ปจฺฉิมานํ วิจิกิจฺฉาสหคตานํ อุทฺธจฺจสหคตานํ ขนฺธานํ โมหสฺส จ อนนฺตรปจฺจเยน ปจฺจโย.\csromanspacer{} อาวชฺชนา วิจิกิจฺฉาสหคตานํ อุทฺธจฺจสหคตานํ ขนฺธานํ โมหสฺส จ อนนฺตรปจฺจเยน ปจฺจโย. (3)}

\proseitem{104}{สเหตุโก จ อเหตุโก จ ธมฺมา สเหตุกสฺส ธมฺมสฺส อนนฺตรปจฺจเยน ปจฺจโย.\csromanspacer{} ปุริมา ปุริมา วิจิกิจฺฉาสหคตา อุทฺธจฺจสหคตา ขนฺธา จ โมโห จ ปจฺฉิมานํ ปจฺฉิมานํ วิจิกิจฺฉาสหคตานํ อุทฺธจฺจสหคตานํ ขนฺธานํ อนนฺตรปจฺจเยน ปจฺจโย. (1)}

\prose{สเหตุโก จ อเหตุโก จ ธมฺมา อเหตุกสฺส ธมฺมสฺส อนนฺตรปจฺจเยน ปจฺจโย.\csromanspacer{} ปุริมา ปุริมา วิจิกิจฺฉาสหคตา อุทฺธจฺจสหคตา ขนฺธา จ โมโห จ ปจฺฉิมสฺส ปจฺฉิมสฺส วิจิกิจฺฉาสหคตสฺส อุทฺธจฺจสหคตสฺส โมหสฺส อนนฺตรปจฺจเยน ปจฺจโย.\csromanspacer{} วิจิกิจฺฉาสหคตา อุทฺธจฺจสหคตา ขนฺธา จ โมโห จ อเหตุกสฺส วุฏฺฐานสฺส อนนฺตรปจฺจเยน ปจฺจโย. (2)}

\prose{สเหตุโก จ อเหตุโก จ ธมฺมา สเหตุกสฺส จ อเหตุกสฺส จ ธมฺมสฺส อนนฺตรปจฺจเยน ปจฺจโย.\csromanspacer{} ปุริมา ปุริมา วิจิกิจฺฉาสหคตา อุทฺธจฺจสหคถา ขนฺธา จ โมโห จ ปจฺฉิมานํ ปจฺฉิมานํ วิจิกิจฺฉาสหคตานํ อุทฺธจฺจสหคตานํ ขนฺธานํ โมหสฺส จ อนนฺตรปจฺจเยน ปจฺจโย. (3)}

\csromanheader{2}{\textbf{สหชาตาทิ}}
\proseitem{105}{สเหตุโก ธมฺโม สเหตุกสฺส ธมฺมสฺส สหชาตปจฺจเยน ปจฺจโย. (ปฏิจฺจวาเร สหชาตสทิสํ, อิห ฆฏนา นตฺถิ.) \csromanfolio{44}อญฺญมญฺญปจฺจเยน ปจฺจโย. (ปฏิจฺจวารสทิสํ.) นิสฺสยปจฺจเยน ปจฺจโย. (ปฏิจฺจวาเร นิสฺสยปจฺจยสทิสํ, อิห ฆฏนา นตฺถิ.)}

\csromanheader{2}{\textbf{อุปนิสฺสย}}
\proseitem{106}{สเหตุโก ธมฺโม สเหตุกสฺส ธมฺมสฺส อุปนิสฺสยปจฺจเยน ปจฺจโย.\csromanspacer{} \textbf{อารมฺมณูปนิสฺสโย, อนนฺตรูปนิสฺสโย, ปกตูปนิสฺสโย ฯเปฯ.}\csromanspacer{} ปกตูปนิสฺสโย\csromandash{}สเหตุกา ขนฺธา สเหตุกานํ ขนฺธานํ อุปนิสฺสยปจฺจเยน ปจฺจโย. (1)}

\prose{สเหตุโก ธมฺโม อเหตุกสฺส ธมฺมสฺส อุปนิสฺสยปจฺจเยน ปจฺจโย.\csromanspacer{} \textbf{อนนฺตรูปนิสฺสโย, ปกตูปนิสฺสโย ฯเปฯ.}\csromanspacer{} ปกตูปนิสฺสโย\csromandash{} สเหตุกา ขนฺธา อเหตุกานํ ขนฺธานํ โมหสฺส จ อุปนิสฺสยปจฺจเยน ปจฺจโย. (2)}

\prose{สเหตุโก ธมฺโม สเหตุกสฺส จ อเหตุกสฺส จ ธมฺมสฺส อุปนิสฺสยปจฺจเยน ปจฺจโย.\csromanspacer{} \textbf{อนนฺตรูปนิสฺสโย, ปกตูปนิสฺสโย ฯเปฯ.}\csromanspacer{} ปกตูปนิสฺสโย\csromandash{}สเหตุกา ขนฺธา วิจิกิจฺฉาสหคตานํ อุทฺธจฺจสหคตานํ ขนฺธานํ โมหสฺส จ อุปนิสฺสยปจฺจเยน ปจฺจโย. (3)}

\proseitem{107}{อเหตุโก ธมฺโม อเหตุกสฺส ธมฺมสฺส อุปนิสฺสยปจฺจเยน ปจฺจโย.\csromanspacer{} \textbf{อนนฺตรูปนิสฺสโย, ปกตูปนิสฺสโย ฯเปฯ.}\csromanspacer{} ปกตูปนิสฺสโย\csromandash{} กายิกํ สุขํ กายิกสฺส สุขสฺส, กายิกสฺส ทุกฺขสฺส, โมหสฺส จ อุปนิสฺสยปจฺจเยน ปจฺจโย.\csromanspacer{} กายิกํ ทุกฺขํ.\csromanspacer{} อุตุ.\csromanspacer{} โภชนํ.\csromanspacer{} เสนาสนํ กายิกสฺส สุขสฺส, กายิกสฺส ทุกฺขสฺส, โมหสฺส จ อุปนิสฺสยปจฺจเยน ปจฺจโย.\csromanspacer{} โมโห กายิกสฺส สุขสฺส, กายิกสฺส ทุกฺขสฺส, โมหสฺส จ อุปนิสฺสยปจฺจเยน ปจฺจโย.\csromanspacer{} กายิกํ สุขํ.\csromanspacer{} กายิกํ ทุกฺขํ.\csromanspacer{} อุตุ.\csromanspacer{} โภชนํ เสนาสนํ โมโห จ กายิกสฺส สุขสฺส, กายิกสฺส ทุกฺขสฺส, โมหสฺส จ อุปนิสฺสยปจฺจเยน ปจฺจโย. (1)}

\prose{อเหตุโก ธมฺโม สเหตุกสฺส ธมฺมสฺส อุปนิสฺสยปจฺจเยน ปจฺจโย.\csromanspacer{} \textbf{อารมฺมณูปนิสฺสโย, อนนฺตรูปนิสฺสโย, ปกตูปนิสฺสโย ฯเปฯ.}\csromanspacer{} ปกตูปนิสฺสโย\csromandash{}กายิกํ สุขํ อุปนิสฺสาย ทานํ เทติ ฯเปฯ สํฆํ ภินฺทติ.\csromanspacer{} กายิกํ ทุกฺขํ.\csromanspacer{} อุตุํ.\csromanspacer{} โภชนํ.\csromanspacer{} เสนาสนํ.\csromanspacer{} โมหํ}

\vspace{\tipitakaparskip}
\csromancenter{สเหตุทุก อุปนิสฺสาย ทานํ เทติ ฯเปฯ สํฆํ ภินฺทติ.\csromanspacer{} กายิกํ สุขํ ฯเปฯ โมโห จ สทฺธาย ฯเปฯ ปญฺญาย.\csromanspacer{} ราคสฺส ฯเปฯ ปตฺถนาย.\csromanspacer{} มคฺคสฺส ผลสมาปตฺติยา อุปนิสฺสยปจฺจเยน ปจฺจโย. (2)}
\prose{\csromanfolio{45}อเหตุโก ธมฺโม สเหตุกสฺส จ อเหตุกสฺส จ ธมฺมสฺส อุปนิสฺสยปจฺจเยน ปจฺจโย.\csromanspacer{} \textbf{อนนฺตรูปนิสฺสโย, ปกตูปนิสฺสโย ฯเปฯ.}\csromanspacer{} ปกตูปนิสฺสโย\csromandash{}กายิกํ สุขํ โมโห จ วิจิกิจฺฉาสหคตานํ อุทฺธจฺจสหคตานํ ขนฺธานํ โมหสฺส จ อุปนิสฺสยปจฺจเยน ปจฺจโย. (3)}

\proseitem{108}{สเหตุโก จ อเหตุโก จ ธมฺมา สเหตุกสฺส ธมฺมสฺส อุปนิสฺสยปจฺจเยน ปจฺจโย.\csromanspacer{} \textbf{อนนฺตรูปนิสฺสโย, ปกตูปนิสฺสโย ฯเปฯ.}\csromanspacer{} ปกตูปนิสฺสโย\csromandash{}วิจิกิจฺฉาสหคตา อุทฺธจฺจสหคตา ขนฺธา จ โมโห จ สเหตุกานํ ขนฺธานํ อุปนิสฺสยปจฺจเยน ปจฺจโย. (1)}

\prose{สเหตุโก จ อเหตุโก จ ธมฺมา อเหตุกสฺส ธมฺมสฺส อุปนิสฺสยปจฺจเยน ปจฺจโย.\csromanspacer{} \textbf{อนนฺตรูปนิสฺสโย, ปกตูปนิสฺสโย ฯเปฯ.}\csromanspacer{} ปกตูปนิสฺสโย\csromandash{}วิจิกิจฺฉาสหคตา อุทฺธจฺจสหคตา ขนฺธา จ โมโห จ อเหตุกานํ ขนฺธานํ โมหสฺส จ อุปนิสฺสยปจฺจเยน ปจฺจโย. (2)}

\prose{สเหตุโก จ อเหตุโก จ ธมฺมา สเหตุกสฺส จ อเหตุกสฺส จ ธมฺมสฺส อุปนิสฺสยปจฺจเยน ปจฺจโย.\csromanspacer{} \textbf{อนนฺตรูปนิสฺสโย, ปกตูปนิสฺสโย ฯเปฯ.}\csromanspacer{} ปกตูปนิสฺสโย\csromandash{}วิจิกิจฺฉาสหคตา อุทฺธจฺจสหคตา ขนฺธา จ โมโห จ วิจิกิจฺฉาสหคตานํ อุทฺธจฺจสหคตานํ ขนฺธานํ โมหสฺส จ อุปนิสฺสยปจฺจเยน ปจฺจโย. (3)}

\csromanheader{2}{\textbf{ปุเรชาต}}
\proseitem{109}{อเหตุโก ธมฺโม อเหตุกสฺส ธมฺมสฺส ปุเรชาตปจฺจเยน ปจฺจโย.\csromanspacer{} \textbf{อารมฺมณปุเรชาตํ}, วตฺถุปุเรชาตํ.\csromanspacer{}\csromanspacer{}\textbf{อารมฺมณปุเรชาตํ}\csromandash{} จกฺขุํ ฯเปฯ วตฺถุํ.\csromanspacer{} อนิจฺจโต ฯเปฯ โทมนสฺสํ อุปฺปชฺชติ.\csromanspacer{} กุสลากุสเล นิรุทฺเธ อเหตุโก วิปาโก ตทารมฺมณตา อุปฺปชฺชติ.\csromanspacer{} รูปายตนํ จกฺขุวิญฺญาณสฺส ฯเปฯ.\csromanspacer{} โผฏฺฐพฺพายตนํ กายวิญฺญาณสฺส ปุเรชาตปจฺจเยน \csromanfolio{46}ปจฺจโย.\csromanspacer{} \textbf{วตฺถุปุเรชาตํ}\csromandash{}จกฺขายตนํ จกฺขุวิญฺญาณสฺส ฯเปฯ กายายตนํ กายวิญฺญาณสฺส.\csromanspacer{} ปุเรชาตํ วตฺถุ อเหตุกานํ ขนฺธานํ โมหสฺส จ ปุเรชาตปจฺจเยน ปจฺจโย. (1)}

\prose{อเหตุโก ธมฺโม สเหตุกสฺส ธมฺมสฺส ปุเรชาตปจฺจเยน ปจฺจโย.\csromanspacer{} อารมฺมณปุเรชาตํ, วตฺถุปุเรชาตํ.\csromanspacer{}\csromanspacer{}\textbf{อารมฺมณปุเรชาตํ\csromandash{}}จกฺขุํ ฯเปฯ วตฺถุํ อนิจฺจโต ฯเปฯ โทมนสฺสํ อุปฺปชฺชติ.\csromanspacer{} กุสลากุสเล นิรุทฺเธ สเหตุโก วิปาโก ตทารมฺมณตา อุปฺปชฺชติ.\csromanspacer{} ทิพฺเพน จกฺขุนา รูปํ ปสฺสติ.\csromanspacer{} ทิพฺพาย โสตธาตุยา สทฺทํ สุณาติ.\csromanspacer{} \textbf{วตฺถุปุเรชาตํ}\csromandash{}วตฺถุ สเหตุกานํ ขนฺธานํ ปุเรชาตปจฺจเยน ปจฺจโย. (2)}

\prose{อเหตุโก ธมฺโม สเหตุกสฺส จ อเหตุกสฺส จ ธมฺมสฺส ปุเรชาตปจฺจเยน ปจฺจโย.\csromanspacer{} อารมฺมณปุเรชาตํ, วตฺถุปุเรชาตํ.\csromanspacer{}\csromanspacer{}\textbf{อารมฺมณปุเรชาตํ\csromandash{}}จกฺขุํ ฯเปฯ วตฺถุํ อารพฺภ วิจิกิจฺฉาสหคตา อุทฺธจฺจสหคตา ขนฺธา จ โมโห จ อุปฺปชฺชนฺติ.\csromanspacer{} \textbf{วตฺถุปุเรชาตํ}\csromandash{}วตฺถุ วิจิกิจฺฉาสหคตานํ อุทฺธจฺจสหคตานํ ขนฺธานํ โมหสฺส จ ปุเรชาตปจฺจเยน ปจฺจโย.}

\csromanheader{2}{\textbf{ปจฺฉาชาตาทิ}}
\proseitem{110}{สเหตุโก ธมฺโม อเหตุกสฺส ธมฺมสฺส ปจฺฉาชาตปจฺจเยน ปจฺจโย.\csromanspacer{} ปจฺฉาชาตา สเหตุกา ขนฺธา ปุเรชาตสฺส อิมสฺส กายสฺส ปจฺฉาชาตปจฺจเยน ปจฺจโย. (1)}

\prose{อเหตุโก ธมฺโม อเหตุกสฺส ธมฺมสฺส ปจฺฉาชาตปจฺจเยน ปจฺจโย.\csromanspacer{} ปจฺฉาชาตา อเหตุกา ขนฺธา จ โมโห จ ปุเรชาตสฺส อิมสฺส กายสฺส ปจฺฉาชาตปจฺจเยน ปจฺจโย. (1)}

\prose{สเหตุโก จ อเหตุโก จ ธมฺมา อเหตุกสฺส ธมฺมสฺส ปจฺฉาชาตปจฺจเยน ปจฺจโย.\csromanspacer{} ปจฺฉาชาตา วิจิกิจฺฉาสหคตา อุทฺธจฺจสหคตา ขนฺธา จ โมโห จ ปุเรชาตสฺส อิมสฺส กายสฺส ปจฺฉาชาตปจฺจเยน ปจฺจโย. (1)}

\prose{สเหตุโก ธมฺโม สเหตุกสฺส ธมฺมสฺส อาเสวนปจฺจเยน ปจฺจโย. (อนนฺตรสทิสํ.\csromanspacer{} อาวชฺชนมฺปิ ภวงฺคมฺปิ นตฺถิ, อาเสวนปจฺจเย วชฺเชตพฺพา นวปิ.)}

\csromanheader{2}{2. สเหตุทุก \textbf{กมฺม}}
\proseitem{111}{\csromanfolio{47}สเหตุโก ธมฺโม สเหตุกสฺส ธมฺมสฺส กมฺมปจฺจเยน ปจฺจโย.\csromanspacer{} \textbf{สหชาตา}, นานากฺขณิกา.\csromanspacer{}\csromanspacer{}\textbf{สหชาตา} สเหตุกา เจตนา สมฺปยุตฺตกานํ ขนฺธานํ กมฺมปจฺจเยน ปจฺจโย.\csromanspacer{} ปฏิสนฺธิกฺขเณ ฯเปฯ.\csromanspacer{} \textbf{นานากฺขณิกา} สเหตุกา เจตนา วิปากานํ สเหตุกานํ ขนฺธานํ กมฺมปจฺจเยน ปจฺจโย. (1)}

\prose{สเหตุโก ธมฺโม อเหตุกสฺส ธมฺมสฺส กมฺมปจฺจเยน ปจฺจโย.\csromanspacer{} \textbf{สหชาตา}, นานากฺขณิกา.\csromanspacer{}\csromanspacer{}\textbf{สหชาตา} สเหตุกา เจตนา จิตฺตสมุฏฺฐานานํ รูปานํ กมฺมปจฺจเยน ปจฺจโย.\csromanspacer{} ปฏิสนฺธิกฺขเณ ฯเปฯ.\csromanspacer{} \textbf{นานากฺขณิกา} สเหตุกา เจตนา วิปากานํ อเหตุกานํ ขนฺธานํ กฏตฺตา จ รูปานํ กมฺมปจฺจเยน ปจฺจโย. (2)}

\prose{สเหตุโก ธมฺโม สเหตุกสฺส จ อเหตุกสฺส จ ธมฺมสฺส กมฺมปจฺจเยน ปจฺจโย.\csromanspacer{} \textbf{สหชาตา}, นานากฺขณิกา.\csromanspacer{}\csromanspacer{}\textbf{สหชาตา} สเหตุกา เจตนา สมฺปยุตฺตกานํ ขนฺธานํ จิตฺตสมุฏฺฐานานญฺจ รูปานํ กมฺมปจฺจเยน ปจฺจโย.\csromanspacer{} \textbf{นานากฺขณิกา} สเหตุกา เจตนา วิปากานํ สเหตุกานํ ขนฺธานํ กฏตฺตา จ รูปานํ กมฺมปจฺจเยน ปจฺจโย. (3)}

\prose{อเหตุโก ธมฺโม อเหตุกสฺส ธมฺมสฺส กมฺมปจฺจเยน ปจฺจโย.\csromanspacer{} อเหตุกา เจตนา สมฺปยุตฺตกานํ ขนฺธานํ จิตฺตสมุฏฺฐานานญฺจ รูปานํ กมฺมปจฺจเยน ปจฺจโย.\csromanspacer{} ปฏิสนฺธิกฺขเณ ฯเปฯ. (1)}

\csromanheader{2}{\textbf{วิปาก}}
\proseitem{112}{สเหตุโก ธมฺโม สเหตุกสฺส ธมฺมสฺส วิปากปจฺจเยน ปจฺจโย.\csromanspacer{} วิปาโก สเหตุโก เอโก ขนฺโธ ติณฺณนฺนํ ขนฺธานํ ฯเปฯ เทฺว ขนฺธา ฯเปฯ.\csromanspacer{} ปฏิสนฺธิกฺขเณ ฯเปฯ. (1)}

\prose{สเหตุโก ธมฺโม อเหตุกสฺส ธมฺมสฺส วิปากปจฺจเยน ปจฺจโย.\csromanspacer{} วิปากา สเหตุกา ขนฺธา จิตฺตสมุฏฺฐานานํ รูปานํ วิปากปจฺจเยน ปจฺจโย.\csromanspacer{} ปฏิสนฺธิกฺขเณ ฯเปฯ. (2)}

\prose{สเหตุโก ธมฺโม สเหตุกสฺส จ อเหตุกสฺส จ ธมฺมสฺส วิปากปจฺจเยน ปจฺจโย.\csromanspacer{} วิปาโก สเหตุโก เอโก ขนฺโธ ติณฺณนฺนํ ขนฺธานํ จิตฺตสมุฏฺฐานานญฺจ รูปานํ ฯเปฯ เทฺว ขนฺธา ฯเปฯ.\csromanspacer{} ปฏิสนฺธิกฺขเณ ฯเปฯ. (3)}

\prose{\csromanfolio{48}อเหตุโก ธมฺโม อเหตุกสฺส ธมฺมสฺส วิปากปจฺจเยน ปจฺจโย.\csromanspacer{} วิปาโก อเหตุโก เอโก ขนฺโธ ติณฺณนฺนํ ขนฺธานํ จิตฺตสมุฏฺฐานานญฺจ รูปานํ ฯเปฯ เทฺว ขนฺธา ฯเปฯ.\csromanspacer{} ปฏิสนฺธิกฺขเณ ฯเปฯ.\csromanspacer{} ขนฺธา วตฺถุสฺส วิปากปจฺจเยน ปจฺจโย. (1)}

\csromanheader{2}{\textbf{อาหาร}}
\proseitem{113}{สเหตุโก ธมฺโม สเหตุกสฺส ธมฺมสฺส อาหารปจฺจเยน ปจฺจโย.\csromanspacer{} ตีณิ.}

\prose{อเหตุโก ธมฺโม อเหตุกสฺส ธมฺมสฺส อาหารปจฺจเยน ปจฺจโย.\csromanspacer{} อเหตุกา อาหารา สมฺปยุตฺตกานํ ขนฺธานํ จิตฺตสมุฏฺฐานานญฺจ รูปานํ ฯเปฯ.\csromanspacer{} ปฏิสนฺธิกฺขเณ ฯเปฯ.\csromanspacer{} กพฬีกาโร อาหาโร อิมสฺส กายสฺส อาหารปจฺจเยน ปจฺจโย. (1)}

\csromanheader{2}{\textbf{อินฺทฺริยาทิ}}
\proseitem{114}{สเหตุโก ธมฺโม สเหตุกสฺส ธมฺมสฺส อินฺทฺริยปจฺจเยน ปจฺจโย.\csromanspacer{} ตีณิ.}

\prose{อเหตุโก ธมฺโม อเหตุกสฺส ธมฺมสฺส อินฺทฺริยปจฺจเยน ปจฺจโย.\csromanspacer{} อเหตุกา อินฺทฺริยา สมฺปยุตฺตกานํ ขนฺธานํ จิตฺตสมุฏฺฐานานญฺจ รูปานํ ฯเปฯ.\csromanspacer{} ปฏิสนฺธิกฺขเณ ฯเปฯ.\csromanspacer{} จกฺขุนฺทฺริยํ จกฺขุวิญฺญาณสฺส ฯเปฯ.\csromanspacer{} กายินฺทฺริยํ กายวิญฺญาณสฺส อินฺทฺริยปจฺจเยน ปจฺจโย.\csromanspacer{} รูปชีวิตินฺทฺริยํ กฏตฺตารูปานํ อินฺทฺริยปจฺจเยน ปจฺจโย. (1)}

\prose{สเหตุโก ธมฺโม สเหตุกสฺส ธมฺมสฺส ฌานปจฺจเยน ปจฺจโย.\csromanspacer{} ตีณิ.}

\prose{อเหตุโก ธมฺโม อเหตุกสฺส ธมฺมสฺส ฌานปจฺจเยน ปจฺจโย.\csromanspacer{} อเหตุกานิ ฌานงฺคานิ สมฺปยุตฺตกานํ ขนฺธานํ จิตฺตสมุฏฺฐานานญฺจ รูปานํ ฯเปฯ.\csromanspacer{} ปฏิสนฺธิกฺขเณ ฯเปฯ. (1)}

\prose{สเหตุโก ธมฺโม สเหตุกสฺส ธมฺมสฺส มคฺคปจฺจเยน ปจฺจโย.\csromanspacer{} ตีณิ.}

\prose{สเหตุโก ธมฺโม สเหตุกสฺส ธมฺมสฺส สมฺปยุตฺตปจฺจเยน ปจฺจโย. (ปฏิจฺจวาเร สมฺปยุตฺตสทิสา ฉ ปญฺหา.)}

\csromanheader{2}{2. สเหตุทุก \textbf{วิปฺปยุตฺต}}
\proseitem{115}{\csromanfolio{49}สเหตุโก ธมฺโม อเหตุกสฺส ธมฺมสฺส วิปฺปยุตฺตปจฺจเยน ปจฺจโย.\csromanspacer{} \textbf{สหชาตํ}, ปจฺฉาชาตํ.\csromanspacer{}\csromanspacer{}\textbf{สหชาตา} สเหตุกา ขนฺธา จิตฺตสมุฏฺฐานานํ รูปานํ วิปฺปยุตฺตปจฺจเยน ปจฺจโย.\csromanspacer{} ปฏิสนฺธิกฺขเณ สเหตุกา ขนฺธา กฏตฺตารูปานํ วิปฺปยุตฺตปจฺจเยน ปจฺจโย.\csromanspacer{} \textbf{ปจฺฉาชาตา} สเหตุกา ขนฺธา ปุเรชาตสฺส อิมสฺส กายสฺส วิปฺปยุตฺตปจฺจเยน ปจฺจโย. (1)}

\prose{อเหตุโก ธมฺโม อเหตุกสฺส ธมฺมสฺส วิปฺปยุตฺตปจฺจเยน ปจฺจโย.\csromanspacer{} \textbf{สหชาตํ}, ปุเรชาตํ, ปจฺฉาชาตํ.\csromanspacer{}\csromanspacer{}\textbf{สหชาตา} อเหตุกา ขนฺธา จิตฺตสมุฏฺฐานานํ รูปานํ วิปฺปยุตฺตปจฺจเยน ปจฺจโย.\csromanspacer{} ปฏิสนฺธิกฺขเณ อเหตุกา ขนฺธา กฏตฺตารูปานํ วิปฺปยุตฺตปจฺจเยน ปจฺจโย.\csromanspacer{} ขนฺธา วตฺถุสฺส วิปฺปยุตฺตปจฺจเยน ปจฺจโย.\csromanspacer{} วตฺถุ ขนฺธานํ วิปฺปยุตฺตปจฺจเยน ปจฺจโย.\csromanspacer{} \textbf{ปุเรชาตํ} จกฺขาย-ตนํ จกฺขุวิญฺญาณสฺส ฯเปฯ.\csromanspacer{} กายายตนํ กายวิญฺญาณสฺส.\csromanspacer{} วตฺถุ อเหตุกานํ ขนฺธานํ โมหสฺส จ วิปฺปยุตฺตปจฺจเยน ปจฺจโย.\csromanspacer{} \textbf{ปจฺฉาชาตา} อเหตุกา ขนฺธา จ โมโห จ ปุเรชาตสฺส อิมสฺส กายสฺส วิปฺปยุตฺตปจฺจเยน ปจฺจโย. (1)}

\prose{อเหตุโก ธมฺโม สเหตุกสฺส ธมฺมสฺส วิปฺปยุตฺตปจฺจเยน ปจฺจโย.\csromanspacer{} \textbf{สหชาตํ}, ปุเรชาตํ.\csromanspacer{}\csromanspacer{}\textbf{สหชาตํ} ปฏิสนฺธิกฺขเณ วตฺถุ สเหตุกานํ ขนฺธานํ วิปฺปยุตฺตปจฺจเยน ปจฺจโย.\csromanspacer{} \textbf{ปุเรชาตํ} วตฺถุ สเหตุกานํ ขนฺธานํ วิปฺปยุตฺตปจฺจเยน ปจฺจโย. (2)}

\prose{อเหตุโก ธมฺโม สเหตุกสฺส จ อเหตุกสฺส จ ธมฺมสฺส วิปฺปยุตฺตปจฺจเยน ปจฺจโย.\csromanspacer{} \textbf{ปุเรชาตํ} วตฺถุ วิจิกิจฺฉาสหคตานํ อุทฺธจฺจสหคตานํ ขนฺธานํ โมหสฺส จ วิปฺปยุตฺตปจฺจเยน ปจฺจโย. (3)}

\prose{สเหตุโก จ อเหตุโก จ ธมฺมา อเหตุกสฺส ธมฺมสฺส วิปฺปยุตฺตปจฺจเยน ปจฺจโย.\csromanspacer{} \textbf{สหชาตํ}, ปจฺฉาชาตํ.\csromanspacer{}\csromanspacer{}\textbf{สหชาตา} วิจิกิจฺฉาสหคตา อุทฺธจฺจสหคตา ขนฺธา จ โมโห จ จิตฺตสมุฏฺฐานานํ รูปานํ วิปฺปยุตฺตปจฺจเยน ปจฺจโย.\csromanspacer{} \textbf{ปจฺฉาชาตา} วิจิกิจฺฉาสหคตา อุทฺธจฺจสหคตา ขนฺธา จ โมโห จ ปุเรชาตสฺส อิมสฺส กายสฺส วิปฺปยุตฺตปจฺจเยน ปจฺจโย. (1)}

\csromanheader{2}{\textbf{อตฺถิ}}
\proseitem{116}{\csromanfolio{50}สเหตุโก ธมฺโม สเหตุกสฺส ธมฺมสฺส อตฺถิปจฺจเยน ปจฺจโย.\csromanspacer{} สเหตุโก เอโก ขนฺโธ ติณฺณนฺนํ ขนฺธานํ ฯเปฯ เทฺว ขนฺธา ฯเปฯ.\csromanspacer{} ปฏิสนฺธิกฺขเณ ฯเปฯ. (1)}

\prose{สเหตุโก ธมฺโม อเหตุกสฺส ธมฺมสฺส อตฺถิปจฺจเยน ปจฺจโย.\csromanspacer{} สหชาตํ, ปจฺฉาชาตํ.\csromanspacer{}\csromanspacer{}\textbf{สหชาตา} สเหตุกา ขนฺธา จิตฺตสมุฏฺฐานานํ รูปานํ อตฺถิปจฺจเยน ปจฺจโย.\csromanspacer{} วิจิกิจฺฉาสหคตา อุทฺธจฺจสหคตา ขนฺธา โมหสฺส จ จิตฺตสมุฏฺฐานานญฺจ รูปานํ อตฺถิปจฺจเยน ปจฺจโย.\csromanspacer{} ปฏิสนฺธิกฺขเณ ฯเปฯ.\csromanspacer{} \textbf{ปจฺฉาชาตา} สเหตุกา ขนฺธา ปุเรชาตสฺส อิมสฺส กายสฺส อตฺถิปจฺจเยน ปจฺจโย. (2)}

\prose{สเหตุโก ธมฺโม สเหตุกสฺส จ อเหตุกสฺส จ ธมฺมสฺส อตฺถิปจฺจเยน ปจฺจโย.\csromanspacer{} สเหตุโก เอโก ขนฺโธ ติณฺณนฺนํ ขนฺธานํ จิตฺตสมุฏฺฐานานญฺจ รูปานํ อตฺถิปจฺจเยน ปจฺจโย.\csromanspacer{} วิจิกิจฺฉาสหคโต อุทฺธจฺจสหคโต เอโก ขนฺโธ ติณฺณนฺนํ ขนฺธานํ โมหสฺส จ จิตฺตสมุฏฺฐานานญฺจ รูปานํ อตฺถิปจฺจเยน ปจฺจโย ฯเปฯ เทฺว ขนฺธา ฯเปฯ.\csromanspacer{} ปฏิสนฺธิกฺขเณ สเหตุโก ฯเปฯ. (3)}

\proseitem{117}{อเหตุโก ธมฺโม อเหตุกสฺส ธมฺมสฺส อตฺถิปจฺจเยน ปจฺจโย.\csromanspacer{} สหชาตํ, ปุเรชาตํ, ปจฺฉาชาตํ, อาหารํ, อินฺทฺริยํ.\csromanspacer{}\csromanspacer{}\textbf{สหชาโต} อเหตุโก เอโก ขนฺโธ ติณฺณนฺนํ ขนฺธานํ จิตฺตสมุฏฺฐานานญฺจ รูปานํ อตฺถิปจฺจเยน ปจฺจโย ฯเปฯ เทฺว ขนฺธา ฯเปฯ.\csromanspacer{} วิจิกิจฺฉาสหคโต อุทฺธจฺจสหคโต โมโห จิตฺตสมุฏฺฐานานํ รูปานํ อตฺถิปจฺจเยน ปจฺจโย.\csromanspacer{} ปฏิสนฺธิกฺขเณ ฯเปฯ. (ยาว อสญฺญสตฺตา กาตพฺพํ.) \textbf{ปุเรชาตํ} จกฺขุํ ฯเปฯ วตฺถุํ อนิจฺจโต ฯเปฯ โทมนสฺสํ อุปฺปชฺชติ.\csromanspacer{} กุสลากุสเล นิรุทฺเธ อเหตุโก วิปาโก ตทารมฺมณตา อุปฺปชฺชติ.\csromanspacer{} รูปายตนํ จกฺขุวิญฺญาณสฺส ฯเปฯ.\csromanspacer{} โผฏฺฐพฺพายตนํ กายวิญฺญาณสฺส ฯเปฯ.\csromanspacer{} จกฺขายตนํ จกฺขุวิญฺญาณสฺส ฯเปฯ.\csromanspacer{} กายายตนํ กายวิญฺญาณสฺส อตฺถิปจฺจเยน ปจฺจโย.\csromanspacer{} วตฺถุ อเหตุกานํ ขนฺธานํ โมหสฺส จ อตฺถิปจฺจเยน ปจฺจโย.\csromanspacer{} \textbf{ปจฺฉาชาตา} อเหตุกา ขนฺธา จ โมโห จ ปุเรชาตสฺส อิมสฺส กายสฺส อตฺถิปจฺจเยน ปจฺจโย.\csromanspacer{} \textbf{กพฬีกาโร}}

\vspace{\tipitakaparskip}
\csromancenter{สเหตุทุก \textbf{อาหาโร} อิมสฺส กายสฺส อตฺถิปจฺจเยน ปจฺจโย.\csromanspacer{} \textbf{รูปชีวิตินฺทฺริยํ} กฏตฺตารูปานํ อตฺถิปจฺจเยน ปจฺจโย. (1)}
\prose{\csromanfolio{51}อเหตุโก ธมฺโม สเหตุกสฺส ธมฺมสฺส อตฺถิปจฺจเยน ปจฺจโย.\csromanspacer{} สหชาตํ, ปุเรชาตํ.\csromanspacer{}\csromanspacer{}\textbf{สหชาโต} วิจิกิจฺฉาสหคโต อุทฺธจฺจสหคโต โมโห สมฺปยุตฺตกานํ ขนฺธานํ อตฺถิปจฺจเยน ปจฺจโย.\csromanspacer{} ปฏิสนฺธิกฺขเณ วตฺถุ สเหตุกานํ ขนฺธานํ อตฺถิปจฺจเยน ปจฺจโย.\csromanspacer{} \textbf{ปุเรชาตํ} จกฺขุํ ฯเปฯ วตฺถุํ อนิจฺจโต ฯเปฯ โทมนสฺสํ อุปฺปชฺชติ.\csromanspacer{} กุสลากุสเล นิรุทฺเธ สเหตุโก วิปาโก ตทารมฺมณตา อุปฺปชฺชติ.\csromanspacer{} ทิพฺเพน จกฺขุนา รูปํ ปสฺสติ.\csromanspacer{} ทิพฺพาย โสตธาตุยา สทฺทํ สุณาติ.\csromanspacer{} วตฺถุ สเหตุกานํ ขนฺธานํ อตฺถิปจฺจเยน ปจฺจโย. (2)}

\prose{อเหตุโก ธมฺโม สเหตุกสฺส จ อเหตุกสฺส จ ธมฺมสฺส อตฺถิปจฺจเยน ปจฺจโย.\csromanspacer{} สหชาตํ, ปุเรชาตํ.\csromanspacer{}\csromanspacer{}\textbf{สหชาโต} วิจิกิจฺฉาสหคโต อุทฺธจฺจสหคโต โมโห สมฺปยุตฺตกานํ ขนฺธานํ จิตฺตสมุฏฺฐานานญฺจ รูปานํ อตฺถิปจฺจเยน ปจฺจโย.\csromanspacer{} \textbf{ปุเรชาตํ} จกฺขุํ ฯเปฯ วตฺถุํ อารพฺภ วิจิกิจฺฉาสหคตา อุทฺธจฺจสหคตา ขนฺธา จ โมโห จ อุปฺปชฺชนฺติ.\csromanspacer{} วตฺถุ วิจิกิจฺฉาสหคตานํ อุทฺธจฺจสหคตานํ ขนฺธานํ โมหสฺส จ อตฺถิปจฺจเยน ปจฺจโย. (3)}

\proseitem{118}{สเหตุโก จ อเหตุโก จ ธมฺมา สเหตุกสฺส ธมฺมสฺส อตฺถิปจฺจเยน ปจฺจโย.\csromanspacer{} สหชาตํ, ปุเรชาตํ.\csromanspacer{}\csromanspacer{}\textbf{สหชาโต} วิจิกิจฺฉาสหคโต อุทฺธจฺจสหคโต เอโก ขนฺโธ จ โมโห จ ติณฺณนฺนํ ขนฺธานํ อตฺถิปจฺจเยน ปจฺจโย ฯเปฯ เทฺว ขนฺธา จ ฯเปฯ.\csromanspacer{} ปฏิสนฺธิกฺขเณ สเหตุโก เอโก ขนฺโธ จ วตฺถุ จ ติณฺณนฺนํ ขนฺธานํ อตฺถิปจฺจเยน ปจฺจโย ฯเปฯ เทฺว ขนฺธา จ ฯเปฯ.\csromanspacer{} \textbf{สหชาโต} สเหตุโก เอโก ขนฺโธ จ วตฺถุ จ ติณฺณนฺนํ ขนฺธานํ อตฺถิปจฺจเยน ปจฺจโย ฯเปฯ เทฺว ขนฺธา จ ฯเปฯ. (1)}

\prose{สเหตุโก จ อเหตุโก จ ธมฺมา อเหตุกสฺส ธมฺมสฺส อตฺถิปจฺจเยน ปจฺจโย.\csromanspacer{} สหชาตํ, ปุเรชาตํ, ปจฺฉาชาตํ, อาหารํ, อินฺทฺริยํ.\csromanspacer{}\csromanspacer{}\textbf{สหชาตา} สเหตุกา ขนฺธา จ มหาภูตา จ จิตฺตสมุฏฺฐานานํ รูปานํ อตฺถิปจฺจเยน ปจฺจโย.\csromanspacer{} วิจิกิจฺฉาสหคตา อุทฺธจฺจสหคตา ขนฺธา จ โมโห จ จิตฺตสมุฏฺฐานานํ รูปานํ อตฺถิปจฺจเยน ปจฺจโย.\csromanspacer{} ปฏิสนฺธิกฺขเณ สเหตุกา ขนฺธา จ มหาภูตา จ กฏตฺตารูปานํ \csromanfolio{52}อตฺถิปจฺจเยน ปจฺจโย.\csromanspacer{} \textbf{สหชาตา} วิจิกิจฺฉาสหคตา อุทฺธจฺจสหคตา ขนฺธา จ วตฺถุ จ โมหสฺส อตฺถิปจฺจเยน ปจฺจโย.\csromanspacer{} \textbf{ปจฺฉาชาตา} วิจิกิจฺฉาสหคตา อุทฺธจฺจสหคตา ขนฺธา จ โมโห จ ปุเรชาตสฺส อิมสฺส กายสฺส อตฺถิปจฺจเยน ปจฺจโย.\csromanspacer{} \textbf{ปจฺฉาชาตา} สเหตุกา ขนฺธา จ กพฬีกาโร อาหาโร จ อิมสฺส กายสฺส อตฺถิปจฺจเยน ปจฺจโย.\csromanspacer{} \textbf{ปจฺฉาชาตา} สเหตุกา ขนฺธา จ รูปชีวิตินฺทฺริยญฺจ กฏตฺตา รูปานํ อตฺถิปจฺจเยน ปจฺจโย. (2)}

\prose{สเหตุโก จ อเหตุโก จ ธมฺมา สเหตุกสฺส จ อเหตุกสฺส จ ธมฺมสฺส อตฺถิปจฺจเยน ปจฺจโย.\csromanspacer{} สหชาตํ, ปุเรชาตํ.\csromanspacer{}\csromanspacer{}\textbf{สหชาโต} วิจิกิจฺฉาสหคโต อุทฺธจฺจสหคโต เอโก ขนฺโธ จ โมโห จ ติณฺณนฺนํ ขนฺธานํ จิตฺตสมุฏฺฐานานญฺจ รูปานํ อตฺถิปจฺจเยน ปจฺจโย ฯเปฯ.\csromanspacer{} \textbf{สหชาโต} วิจิกิจฺฉาสหคโต อุทฺธจฺจสหคโต เอโก ขนฺโธ จ วตฺถุ จ ติณฺณนฺนํ ขนฺธานํ โมหสฺส จ อตฺถิปจฺจเยน ปจฺจโย ฯเปฯ เทฺว ขนฺธา จ ฯเปฯ. (3)}

\csromanheader{2}{1. ปจฺจยานุโลม 2. สงฺขฺยาวาร}
\csromanheader{2}{\textbf{สุทฺธ}}
\proseitem{119}{เหตุยา ฉ, อารมฺมเณ นว, อธิปติยา จตฺตาริ, อนนฺตเร นว, สมนนฺตเร นว, สหชาเต นว, อญฺญมญฺเญ ฉ, นิสฺสเย นว, อุปนิสฺสเย นว, ปุเรชาเต ตีณิ, ปจฺฉาชาเต ตีณิ อาเสวเน นว, กมฺเม จตฺตาริ, วิปาเก จตฺตาริ, อาหาเร จตฺตาริ, อินฺทฺริเย จตฺตาริ, ฌาเน จตฺตาริ, มคฺเค ตีณิ, สมฺปยุตฺเต ฉ, วิปฺปยุตฺเต ปญฺจ, อตฺถิยา นว, นตฺถิยา นว, วิคเต นว, อวิคเต นว. (เอวํ คเณตพฺพํ.)}

\vspace{\tipitakaparskip}
\csromancenter{อนุโลมํ.}
\csromanheader{2}{2. ปจฺจนียุทฺธาร}
\proseitem{120}{สเหตุโก ธมฺโม สเหตุกสฺส ธมฺมสฺส อารมฺมณปจฺจเยน ปจฺจโย.\csromanspacer{} สหชาตปจฺจเยน ปจฺจโย.\csromanspacer{} อุปนิสฺสยปจฺจเยน ปจฺจโย.\csromanspacer{} กมฺมปจฺจเยน ปจฺจโย. (1)}

\csromanheader{2}{2. สเหตุทุก}
\prose{\csromanfolio{53}สเหตุโก ธมฺโม อเหตุกสฺส ธมฺมสฺส อารมฺมณปจฺจเยน ปจฺจโย.\csromanspacer{} สหชาตปจฺจเยน ปจฺจโย.\csromanspacer{} อุปนิสฺสยปจฺจเยน ปจฺจโย.\csromanspacer{} ปจฺฉาชาตปจฺจเยน ปจฺจโย.\csromanspacer{} กมฺมปจฺจเยน ปจฺจโย. (2)}

\prose{สเหตุโก ธมฺโม สเหตุกสฺส จ อเหตุกสฺส จ ธมฺมสฺส อารมฺมณปจฺจเยน ปจฺจโย.\csromanspacer{} สหชาตปจฺจเยน ปจฺจโย.\csromanspacer{} อุปนิสฺสยปจฺจเยน ปจฺจโย.\csromanspacer{} กมฺมปจฺจเยน ปจฺจโย. (3)}

\proseitem{121}{อเหตุโก ธมฺโม อเหตุกสฺส ธมฺมสฺส อารมฺมณปจฺจเยน ปจฺจโย.\csromanspacer{} สหชาตปจฺจเยน ปจฺจโย.\csromanspacer{} อุปนิสฺสยปจฺจเยน ปจฺจโย.\csromanspacer{} ปุเรชาตปจฺจเยน ปจฺจโย.\csromanspacer{} ปจฺฉาชาตปจฺจเยน ปจฺจโย.\csromanspacer{} อาหารปจฺจเยน ปจฺจโย.\csromanspacer{} อินฺทฺริยปจฺจเยน ปจฺจโย. (1)}

\prose{อเหตุโก ธมฺโม สเหตุกสฺส ธมฺมสฺส อารมฺมณปจฺจเยน ปจฺจโย.\csromanspacer{} สหชาตปจฺจเยน ปจฺจโย.\csromanspacer{} อุปนิสฺสยปจฺจเยน ปจฺจโย.\csromanspacer{} ปุเรชาตปจฺจเยน ปจฺจโย. (2)}

\prose{อเหตุโก ธมฺโม สเหตุกสฺส จ อเหตุกสฺส จ ธมฺมสฺส อารมฺมณปจฺจเยน ปจฺจโย.\csromanspacer{} สหชาตปจฺจเยน ปจฺจโย.\csromanspacer{} อุปนิสฺสยปจฺจเยน ปจฺจโย.\csromanspacer{} ปุเรชาตปจฺจเยน ปจฺจโย. (3)}

\proseitem{122}{สเหตุโก จ อเหตุโก จ ธมฺมา สเหตุกสฺส ธมฺมสฺส อารมฺมณปจฺจเยน ปจฺจโย.\csromanspacer{} สหชาตปจฺจเยน ปจฺจโย.\csromanspacer{} อุปนิสฺสยปจฺจเยน ปจฺจโย. (1)}

\prose{สเหตุโก จ อเหตุโก จ ธมฺมา อเหตุกสฺส ธมฺมสฺส อารมฺมณปจฺจเยน ปจฺจโย.\csromanspacer{} สหชาตปจฺจเยน ปจฺจโย.\csromanspacer{} อุปนิสฺสยปจฺจเยน ปจฺจโย.\csromanspacer{} ปจฺฉาชาตปจฺจเยน ปจฺจโย. (2)}

\prose{สเหตุโก จ อเหตุโก จ ธมฺมา สเหตุกสฺส จ อเหตุกสฺส จ ธมฺมสฺส อารมฺมณปจฺจเยน ปจฺจโย.\csromanspacer{} สหชาตปจฺจเยน ปจฺจโย.\csromanspacer{} อุปนิสฺสยปจฺจเยน ปจฺจโย. (3)}

\csromanheader{2}{2. ปจฺจยปจฺจนีย 2. สงฺขฺยาวาร}
\proseitem{123}{\csromanfolio{54}นเหตุยา นว ฯเปฯ (สพฺพตฺถ นว) โน-อวิคเต นว. (เอวํ คเณตพฺพํ.)}

\vspace{\tipitakaparskip}
\csromancenter{ปจฺจนียํ.}
\csromanheader{2}{3. ปจฺจยานุโลมปจฺจนีย}
\csromanheader{2}{\textbf{เหตุทุก}}
\proseitem{124}{เหตุปจฺจยา น-อารมฺมเณ ฉ, น-อธิปติยา ฉ, น-อนนฺตเร ฉ, นสมนนฺตเร ฉ, น-อญฺญมญฺเญ เทฺว, น-อุปนิสฺสเย ฉ ฯเปฯ นมคฺเค ฉ, นสมฺปยุตฺเต เทฺว, นวิปฺปยุตฺเต เทฺว, โนนตฺถิยา ฉ, โนวิคเต ฉ. (เอวํ คเณตพฺพํ.)}

\vspace{\tipitakaparskip}
\csromancenter{อนุโลมปจฺจนียํ.}
\csromanheader{2}{4. ปจฺจยปจฺจนียานุโลม}
\csromanheader{2}{\textbf{นเหตุทุก}}
\proseitem{125}{นเหตุปจฺจยา อารมฺมเณ นว, อธิปติยา จตฺตาริ, อนนฺตเร นว, สมนนฺตเร นว, สหชาเต นว, อญฺญมญฺเญ ฉ, นิสฺสเย นว, อุปนิสฺสเย นว, ปุเรชาเต ตีณิ, ปจฺฉาชาเต ตีณิ, อาเสวเน นว, กมฺเม จตฺตาริ, วิปาเก จตฺตาริ, อาหาเร จตฺตาริ, อินฺทฺริเย จตฺตาริ, ฌาเน จตฺตาริ, มคฺเค ตีณิ, สมฺปยุตฺเต ฉ, วิปฺปยุตฺเต ปญฺจ อตฺถิยา นว, นตฺถิยา นว, วิคเต นว, อวิคเต นว. (เอวํ คเณตพฺพํ.)}

\vspace{\tipitakaparskip}
\csromancenter{ปจฺจนียานุโลมํ.}
\nitthitam{สเหตุกทุกํ นิฏฺฐิตํ.}
\csromanmatikaanchor{mtk.5}
\csromantocmarkref{section}{3. เหตุสมฺปยุตฺตทุก}{mtk.5}
\bookmark[dest={mtk.5},level=1]{3. เหตุสมฺปยุตฺตทุก}
\csromanheader{1}{4. \textbf{เหตุ}สเหตุทุก \textbf{3. เหตุสมฺปยุตฺตทุก 1. ปฏิจฺจวาร}}
\csromanheader{2}{\textbf{เหตุ}}
\vspace{\tipitakaparskip}
\csromancenter{\textbf{เหตุ}สมฺปยุตฺตํ ธมฺมํ ปฏิจฺจ เหตุสมฺปยุตฺโต ธมฺโม อุปฺปชฺชติ เหตุปจฺจยา.\csromanspacer{} \textbf{เหตุ}สมฺปยุตฺตํ เอกํ ขนฺธํ ปฏิจฺจ ตโย ขนฺธา ฯเปฯ เทฺว ขนฺเธ ฯเปฯ.\csromanspacer{} ปฏิสนฺธิกฺขเณ ฯเปฯ. (1)}
\csromancenter{\textbf{เหตุ}สมฺปยุตฺตํ ธมฺมํ ปฏิจฺจ เหตุวิปฺปยุตฺโต ธมฺโม อุปฺปชฺชติ เหตุปจฺจยา.\csromanspacer{} \textbf{เหตุ}สมฺปยุตฺเต ขนฺเธ ปฏิจฺจ จิตฺตสมุฏฺฐานํ รูปํ.\csromanspacer{} ปฏิสนฺธิกฺขเณ ฯเปฯ. (2)}
\prose{\csromanfolio{55}(อิมินา การเณน วิตฺถาเรตพฺพํ.\csromanspacer{} ยถา สเหตุกทุกํ.\csromanspacer{} นินฺนานากรณํ.)}

\nitthitam{เหตุสมฺปยุตฺตทุกํ นิฏฺฐิตํ.}
\csromanmatikaanchor{mtk.6}
\csromantocmarkref{section}{4. เหตุสเหตุกทุก}{mtk.6}
\bookmark[dest={mtk.6},level=1]{4. เหตุสเหตุกทุก}
\csromanheader{1}{4. \textbf{เหตุ}สเหตุกทุก 1. ปฏิจฺจวาร}
\csromanheader{2}{1. ปจฺจยานุโลม 1. วิภงฺควาร}
\csromanheader{2}{\textbf{เหตุ}}
\vspace{\tipitakaparskip}
\csromancenter{\textbf{เหตุ}ญฺเจว สเหตุกญฺจ ธมฺมํ ปฏิจฺจ เหตุ เจว สเหตุโก จ ธมฺโม อุปฺปชฺชติ เหตุปจฺจยา.\csromanspacer{} อโลภํ ปฏิจฺจ อโทโส อโมโห. (จกฺกํ.) โลภํ ปฏิจฺจ โมโห. (จกฺกํ.) ปฏิสนฺธิกฺขเณ อโลภํ ปฏิจฺจ อโทโส อโมโห. (จกฺกํ.) (1)}
\csromancenter{\textbf{เหตุ}ญฺเจว สเหตุกญฺจ ธมฺมํ ปฏิจฺจ สเหตุโก เจว น จ เหตุ ธมฺโม อุปฺปชฺชติ เหตุปจฺจยา.\csromanspacer{} \textbf{เหตุ}ํ ปฏิจฺจ สมฺปยุตฺตกา ขนฺธา ปฏิสนฺธิกฺขเณ ฯเปฯ. (2)}
\csromancenter{\textbf{เหตุ}ญฺเจว สเหตุกญฺจ ธมฺมํ ปฏิจฺจ เหตุ เจว สเหตุโก จ สเหตุโก เจว น จ เหตุ จ ธมฺมา อุปฺปชฺชนฺติ เหตุปจฺจยา.\csromanspacer{} อโลภํ ปฏิจฺจ อโทโส อโมโห สมฺปยุตฺตกา จ ขนฺธา. (จกฺกํ.) โลภํ ปฏิจฺจ โมโห สมฺปยุตฺตกา จ ขนฺธา. (จกฺกํ.) ปฏิสนฺธิกฺขเณ ฯเปฯ. (3)}
\proseitem{128}{\csromanfolio{56}สเหตุกญฺเจว น จ เหตุํ ธมฺมํ ปฏิจฺจ สเหตุโก เจว น จ เหตุ ธมฺโม อุปฺปชฺชติ เหตุปจฺจยา.\csromanspacer{} สเหตุกญฺเจว น จ เหตุํ เอกํ ขนฺธํ ปฏิจฺจ ตโย ขนฺธา ฯเปฯ เทฺว ขนฺเธ ฯเปฯ.\csromanspacer{} ปฏิสนฺธิกฺขเณ ฯเปฯ. (1)}

\prose{สเหตุกญฺเจว น จ เหตุํ ธมฺมํ ปฏิจฺจ เหตุ เจว สเหตุโก จ ธมฺโม อุปฺปชฺชติ เหตุปจฺจยา.\csromanspacer{} สเหตุเก เจว น จ เหตู ขนฺเธ ปฏิจฺจ เหตู.\csromanspacer{} ปฏิสนฺธิกฺขเณ ฯเปฯ. (2)}

\prose{สเหตุกญฺเจว น จ เหตุํ ธมฺมํ ปฏิจฺจ เหตุ เจว สเหตุโก จ สเหตุโก เจว น จ เหตุ จ ธมฺมา อุปฺปชฺชนฺติ เหตุปจฺจยา.\csromanspacer{} สเหตุกญฺเจว น จ เหตุํ เอกํ ขนฺธํ ปฏิจฺจ ตโย ขนฺธา เหตุ จ ฯเปฯ เทฺว ขนฺเธ ฯเปฯ.\csromanspacer{} ปฏิสนฺธิกฺขเณ ฯเปฯ. (3)}

\proseitem{129}{เหตุญฺเจว สเหตุกญฺจ สเหตุกญฺเจว น จ เหตุญฺจ ธมฺมํ ปฏิจฺจ เหตุ เจว สเหตุโก จ ธมฺโม อุปฺปชฺชติ เหตุปจฺจยา.\csromanspacer{} อโลภญฺจ สมฺปยุตฺตเก จ ขนฺเธ ปฏิจฺจ อโทโส อโมโห. (จกฺกํ.) โลภญฺจ สมฺปยุตฺตเก จ ขนฺเธ ปฏิจฺจ โมโห. (จกฺกํ.) ปฏิสนฺธิกฺขเณ ฯเปฯ. (1)}

\prose{เหตุญฺเจว สเหตุกญฺจ สเหตุกญฺเจว น จ เหตุญฺจ ธมฺมํ ปฏิจฺจ สเหตุโก เจว น จ เหตุ ธมฺโม อุปฺปชฺชติ เหตุปจฺจยา.\csromanspacer{} สเหตุกญฺเจว น จ เหตุํ เอกํ ขนฺธญฺจ เหตุญฺจ ปฏิจฺจ ตโย ขนฺธา ฯเปฯ เทฺว ขนฺเธ ฯเปฯ.\csromanspacer{} ปฏิสนฺธิกฺขเณ ฯเปฯ. (2)}

\prose{เหตุญฺเจว สเหตุกญฺจ สเหตุกญฺเจว น จ เหตุญฺจ ธมฺมํ ปฏิจฺจ เหตุ เจว สเหตุโก จ สเหตุโก เจว น จ เหตุ จ ธมฺมา อุปฺปชฺชนฺติ เหตุปจฺจยา.\csromanspacer{} สเหตุกญฺเจว น จ เหตุํ เอกํ ขนฺธญฺจ อโลภญฺจ ปฏิจฺจ ตโย ขนฺธา อโทโส อโมโห จ ฯเปฯ เทฺว ขนฺเธ ฯเปฯ.\csromanspacer{} ปฏิสนฺธิกฺขเณ ฯเปฯ. (3) (สํขิตฺตํ.\csromanspacer{} เอวํ วิตฺถาเรตพฺพํ.)}

\csromanheader{2}{1. ปจฺจยานุโลม 2. สงฺขฺยาวาร}
\csromanheader{2}{\textbf{สุทฺธ}}
\proseitem{130}{เหตุยา นว, อารมฺมเณ นว, อธิปติยา นว, อนนฺตเร นว ฯเปฯ (สพฺพตฺถ นว) อวิคเต นว. (เอวํ คเณตพฺพํ.)}

\vspace{\tipitakaparskip}
\csromancenter{อนุโลมํ.}
\csromanheader{2}{4. เหตุสเหตุทุก \textbf{2. ปจฺจยปจฺจนีย 1. วิภงฺควาร}}
\csromanheader{2}{\textbf{น-อธิปตฺยาทิ}}
\proseitem{131}{\csromanfolio{57}เหตุญฺเจว สเหตุกญฺจ ธมฺมํ ปฏิจฺจ เหตุ เจว สเหตุโก จ ธมฺโม อุปฺปชฺชติ น-อธิปติปจฺจยา.\csromanspacer{} อโลภํ ปฏิจฺจ อโทโส อโมโห. (จกฺกํ.) ปฏิสนฺธิกฺขเณ ฯเปฯ. (ปริปุณฺณํ นว.) นปุเรชาเต นว, นปจฺฉาชาเต นว, น-อาเสวเน นว.}

\csromanheader{2}{\textbf{นกมฺมาทิ}}
\proseitem{132}{เหตุญฺเจว สเหตุกญฺจ ธมฺมํ ปฏิจฺจ สเหตุโก เจว น จ เหตุ ธมฺโม อุปฺปชฺชติ นกมฺมปจฺจยา.\csromanspacer{} เหตุํ ปฏิจฺจ สมฺปยุตฺตกา เจตนา. (1)}

\prose{สเหตุกญฺเจว น จ เหตุํ ธมฺมํ ปฏิจฺจ สเหตุโก เจว น จ เหตุ ธมฺโม อุปฺปชฺชติ นกมฺมปจฺจยา.\csromanspacer{} สเหตุเก เจว น จ เหตู ขนฺเธ ปฏิจฺจ สมฺปยุตฺตกา เจตนา.\csromanspacer{} ปฏิสนฺธิกฺขเณ ฯเปฯ. (1)}

\prose{เหตุญฺเจว สเหตุกญฺจ สเหตุกญฺเจว น จ เหตุญฺจ ธมฺมํ ปฏิจฺจ สเหตุโก เจว น จ เหตุ ธมฺโม อุปฺปชฺชติ นกมฺมปจฺจยา.\csromanspacer{} เหตุญฺจ สมฺปยุตฺตเก จ ขนฺเธ ปฏิจฺจ สมฺปยุตฺตกา เจตนา.\csromanspacer{} นวิปากปจฺจยา.\csromanspacer{} นวิปฺปยุตฺตปจฺจยา.}

\csromanheader{2}{2. ปจฺจยปจฺจนีย 2. สงฺขฺยาวาร}
\csromanheader{2}{\textbf{สุทฺธ}}
\proseitem{133}{น-อธิปติยา นว, นปุเรชาเต นว, นปจฺฉาชาเต นว, น-อาเสวเน นว, นกมฺเม ตีณิ, นวิปาเก นว, นวิปฺปยุตฺเต นว. (เอวํ คเณตพฺพํ.)}

\vspace{\tipitakaparskip}
\csromancenter{ปจฺจนียํ.}
\csromanheader{2}{3. ปจฺจยานุโลมปจฺจนีย}
\csromanheader{2}{\textbf{เหตุทุก}}
\proseitem{134}{\csromanfolio{58}\textbf{เหตุ}ปจฺจยา น-อธิปติยา นว, นปุเรชาเต นว, นปจฺฉาชาเต นว, น-อาเสวเน นว, นกมฺเม ตีณิ, นวิปาเก นว, นวิปฺปยุตฺเต นว. (เอวํ คเณตพฺพํ.)}

\vspace{\tipitakaparskip}
\csromancenter{อนุโลมปจฺจนียํ.}
\csromanheader{2}{4. ปจฺจยปจฺจนียานุโลม}
\csromanheader{2}{\textbf{น-อธิปติทุก}}
\proseitem{135}{น-อธิปติปจฺจยา เหตุยา นว, อารมฺมเณ นว, อนนฺตเร นว ฯเปฯ อวิคเต นว. (เอวํ คเณตพฺพํ.) ปจฺจนียานุโลมํ.}

\prose{(สหชาตวาโรปิ ปจฺจยวาโรปิ นิสฺสยวาโรปิ สํสฏฺฐวาโรปิ สมฺปยุตฺตวาโรปิ ปฏิจฺจวารสทิสา.)}

\csromanheader{2}{4. \textbf{เหตุ}สเหตุกทุก 7. ปญฺหาวาร}
\csromanheader{2}{1. ปจฺจยานุโลม 1. วิภงฺควาร}
\csromanheader{2}{\textbf{เหตุ}}
\vspace{\tipitakaparskip}
\csromancenter{\textbf{เหตุ} เจว สเหตุโก จ ธมฺโม เหตุสฺส เจว สเหตุกสฺส จ ธมฺมสฺส เหตุปจฺจเยน ปจฺจโย.\csromanspacer{} อโลโภ อโทสสฺส อโมหสฺส เหตุปจฺจเยน ปจฺจโย. (ยถา ปฏิจฺจวารสทิสํ.) (1)}
\csromancenter{\textbf{เหตุ} เจว สเหตุโก จ ธมฺโม สเหตุกสฺส เจว น จ เหตุสฺส ธมฺมสฺส เหตุปจฺจเยน ปจฺจโย.\csromanspacer{} เหตู สมฺปยุตฺตกานํ ขนฺธานํ เหตุปจฺจเยน ปจฺจโย.\csromanspacer{} ปฏิสนฺธิกฺขเณ ฯเปฯ. (2)}
\csromanheader{2}{4. เหตุสเหตุทุก}
\prose{\csromanfolio{59}เหตุ เจว สเหตุโก จ ธมฺโม เหตุสฺส เจว สเหตุกสฺส จ สเหตุกสฺส เจว น จ เหตุสฺส จ ธมฺมสฺส เหตุปจฺจเยน ปจฺจโย.\csromanspacer{} อโลโภ อโทสสฺส อโมหสฺส สมฺปยุตฺตกานญฺจ ขนฺธานํ เหตุปจฺจเยน ปจฺจโย. (วิตฺถาเรตพฺพํ.) (3)}

\csromanheader{2}{\textbf{อารมฺมณ}}
\proseitem{137}{เหตุ เจว สเหตุโก จ ธมฺโม เหตุสฺส เจว สเหตุกสฺส จ ธมฺมสฺส อารมฺมณปจฺจเยน ปจฺจโย.\csromanspacer{} เหตุํ อารพฺภ เหตู อุปฺปชฺชนฺติ. (1)}

\prose{เหตุ เจว สเหตุโก จ ธมฺโม สเหตุกสฺส เจว น จ เหตุสฺส ธมฺมสฺส อารมฺมณปจฺจเยน ปจฺจโย.\csromanspacer{} เหตุํ อารพฺภ สเหตุกา เจว น จ เหตู ขนฺธา อุปฺปชฺชนฺติ. (2)}

\prose{เหตุ เจว สเหตุโก จ ธมฺโม เหตุสฺส เจว สเหตุกสฺส จ สเหตุกสฺส เจว น จ เหตุสฺส จ ธมฺมสฺส อารมฺมณปจฺจเยน ปจฺจโย.\csromanspacer{} เหตุํ อารพฺภ เหตู จ สมฺปยุตฺตกา จ ขนฺธา อุปฺปชฺชนฺติ. (3)}

\prose{สเหตุโก เจว น จ เหตุ ธมฺโม สเหตุกสฺส เจว น จ เหตุสฺส ธมฺมสฺส อารมฺมณปจฺจเยน ปจฺจโย.\csromanspacer{} ทานํ ทตฺวา, สีลํ ฯเปฯ อุโปสถกมฺมํ กตฺวา ตํ ปจฺจเวกฺขติ.\csromanspacer{} ปุพฺเพ สุจิณฺณานิ ปจฺจเวกฺขติ.\csromanspacer{} ฌานา วุฏฺฐหิตฺวา ฯเปฯ.\csromanspacer{} อริยา มคฺคา วุฏฺฐหิตฺวา มคฺคํ ปจฺจเวกฺขนฺติ, ผลํ ปจฺจเวกฺขนฺติ.\csromanspacer{} ปหีเน กิเลเส ฯเปฯ.\csromanspacer{} วิกฺขมฺภิเต กิเลเส ปจฺจเวกฺขนฺติ.\csromanspacer{} ปุพฺเพ สมุทาจิณฺเณ กิเลเส ชานนฺติ.\csromanspacer{} สเหตุเก เจว น จ เหตู ขนฺเธ อนิจฺจโต ฯเปฯ โทมนสฺสํ อุปฺปชฺชติ.\csromanspacer{} เจโตปริยญาเณน สเหตุกา เจว น จ เหตุจิตฺตสมงฺคิสฺส จิตฺตํ ชานาติ.\csromanspacer{} อากาสานญฺจายตนํ วิญฺญาณญฺจายตนสฺส ฯเปฯ.\csromanspacer{} อากิญฺจญฺญายตนํ เนวสญฺญานาสญฺญายตนสฺส ฯเปฯ.\csromanspacer{} สเหตุกา เจว น จ เหตู ขนฺธา อิทฺธิวิธญาณสฺส, เจโตปริยญาณสฺส, ปุพฺเพนิวาสานุสฺสติญาณสฺส, ยถากมฺมูปคญาณสฺส, อนาคตํสญาณสฺส อารมฺมณปจฺจเยน ปจฺจโย. (1)}

\prose{สเหตุโก เจว น จ เหตุ ธมฺโม เหตุสฺส เจว สเหตุกสฺส จ ธมฺมสฺส อารมฺมณปจฺจเยน ปจฺจโย.\csromanspacer{} ทานํ ทตฺวา. (ยถา ปฐมคมนํ.\csromanspacer{} เอวํ นินฺนานํ.) (2)}

\prose{\csromanfolio{60}สเหตุโก เจว น จ เหตุ ธมฺโม เหตุสฺส เจว สเหตุกสฺส จ สเหตุกสฺส เจว น จ เหตุสฺส จ ธมฺมสฺส อารมฺมณปจฺจเยน ปจฺจโย.\csromanspacer{} ทานํ ทตฺวา. (ยถา ปฐมคมนํ.\csromanspacer{} เอวํ นินฺนานํ.) (3)}

\proseitem{138}{เหตุ เจว สเหตุโก จ สเหตุโก เจว น จ เหตุ จ ธมฺมา เหตุสฺส เจว สเหตุกสฺส จ ธมฺมสฺส อารมฺมณปจฺจเยน ปจฺจโย.\csromanspacer{} เหตุญฺจ สมฺปยุตฺตเก จ ขนฺเธ อารพฺภ เหตู อุปฺปชฺชนฺติ. (1)}

\prose{เหตุ เจว สเหตุโก จ สเหตุโก เจว น จ เหตุ จ ธมฺมา สเหตุกสฺส เจว น จ เหตุสฺส ธมฺมสฺส อารมฺมณปจฺจเยน ปจฺจโย.\csromanspacer{} เหตุญฺจ สมฺปยุตฺตเก จ ขนฺเธ อารพฺภ สเหตุกา เจว น จ เหตู ขนฺธา อุปฺปชฺชนฺติ. (2)}

\prose{เหตุ เจว สเหตุโก จ สเหตุโก เจว น จ เหตุ จ ธมฺมา เหตุสฺส เจว สเหตุกสฺส จ สเหตุกสฺส เจว น จ เหตุสฺส จ ธมฺมสฺส อารมฺมณปจฺจเยน ปจฺจโย.\csromanspacer{} เหตุญฺจ สมฺปยุตฺตเก จ ขนฺเธ อารพฺภ เหตู จ สมฺปยุตฺตกา จ ขนฺธา อุปฺปชฺชนฺติ. (3)}

\csromanheader{2}{\textbf{อธิปติ}}
\proseitem{139}{เหตุ เจว สเหตุโก จ ธมฺโม เหตุสฺส เจว สเหตุกสฺส จ ธมฺมสฺส อธิปติปจฺจเยน ปจฺจโย.\csromanspacer{} \textbf{อารมฺมณาธิปติ}, สหชาตาธิปติ.\csromanspacer{}\csromanspacer{}\textbf{อารมฺมณาธิปติ}\csromandash{}เหตุํ ครุํ กตฺวา เหตู อุปฺปชฺชนฺติ.\csromanspacer{} \textbf{สหชาตาธิปติ\csromandash{}} เหตุ เจว สเหตุกาธิปติ สมฺปยุตฺตกานํ เหตูนํ อธิปติปจฺจเยน ปจฺจโย. (1)}

\prose{เหตุ เจว สเหตุโก จ ธมฺโม สเหตุกสฺส เจว น จ เหตุสฺส ธมฺมสฺส อธิปติปจฺจเยน ปจฺจโย.\csromanspacer{} \textbf{อารมฺมณาธิปติ}, สหชาตาธิปติ.\csromanspacer{}\csromanspacer{}\textbf{อารมฺมณาธิปติ}\csromandash{}เหตุํ ครุํ กตฺวา สเหตุกา เจว น จ เหตู ขนฺธา อุปฺปชฺชนฺติ.\csromanspacer{} \textbf{สหชาตาธิปติ\csromandash{}}เหตุ เจว สเหตุกาธิปติ สมฺปยุตฺตกานํ ขนฺธานํ อธิปติปจฺจเยน ปจฺจโย. (2)}

\prose{เหตุ เจว สเหตุโก จ ธมฺโม เหตุสฺส เจว สเหตุกสฺส จ สเหตุกสฺส เจว น จ เหตุสฺส จ ธมฺมสฺส อธิปติปจฺจเยน ปจฺจโย.\csromanspacer{} \textbf{อารมฺมณาธิปติ}, สหชาตาธิปติ.\csromanspacer{}\csromanspacer{}\textbf{อารมฺมณาธิปติ}\csromandash{}เหตุํ ครุํ}

\vspace{\tipitakaparskip}
\csromancenter{เหตุสเหตุทุก กตฺวา เหตู จ สมฺปยุตฺตกา จ ขนฺธา อุปฺปชฺชนฺติ.\csromanspacer{} \textbf{สหชาตาธิปติ\csromandash{}}เหตุ เจว สเหตุกาธิปติ สมฺปยุตฺตกานํ ขนฺธานํ เหตูนญฺจ อธิปติปจฺจเยน ปจฺจโย. (3)}
\proseitem{140}{\csromanfolio{61}สเหตุโก เจว น จ เหตุ ธมฺโม สเหตุกสฺส เจว น จ เหตุสฺส ธมฺมสฺส อธิปติปจฺจเยน ปจฺจโย.\csromanspacer{} \textbf{อารมฺมณาธิปติ}, สหชาตาธิปติ.\csromanspacer{}\csromanspacer{}\textbf{อารมฺมณาธิปติ\csromandash{}}ทานํ ทตฺวา, สีลํ ฯเปฯ อุโปสถกมฺมํ กตฺวา ตํ ครุํ กตฺวา ปจฺจเวกฺขติ.\csromanspacer{} ปุพฺเพ สุจิณฺณานิ ฯเปฯ.\csromanspacer{} ฌานา วุฏฺฐหิตฺวา ฌานํ ครุํ กตฺวา ปจฺจเวกฺขติ.\csromanspacer{} อริยา มคฺคา วุฏฺฐหิตฺวา มคฺคํ ครุํ กตฺวา ปจฺจเวกฺขนฺติ, ผลํ ครุํ กตฺวา ปจฺจเวกฺขนฺติ.\csromanspacer{} สเหตุเก เจว น จ เหตู ขนฺเธ ครุํ กตฺวา อสฺสาเทติ อภินนฺทติ, ตํ ครุํ กตฺวา ราโค อุปฺปชฺชติ, ทิฏฺฐิ อุปฺปชฺชติ.\csromanspacer{} \textbf{สหชาตาธิปติ\csromandash{}}สเหตุโก เจว น จ เหตุ อธิปติ สมฺปยุตฺตกานํ ขนฺธานํ อธิปติปจฺจเยน ปจฺจโย. (1)}

\prose{สเหตุโก เจว น จ เหตุ ธมฺโม เหตุสฺส เจว สเหตุกสฺส จ ธมฺมสฺส อธิปติปจฺจเยน ปจฺจโย.\csromanspacer{} \textbf{อารมฺมณาธิปติ}, สหชาตาธิปติ.\csromanspacer{}\csromanspacer{}\textbf{อารมฺมณาธิปติ\csromandash{}}ทานํ ทตฺวา. (ปฐมคมนํเยว.) \textbf{สหชาตาธิปติ\csromandash{}}สเหตุโก เจว น จ เหตุ อธิปติ สมฺปยุตฺตกานํ เหตูนํ อธิปติปจฺจเยน ปจฺจโย. (2)}

\prose{สเหตุโก เจว น จ เหตุ ธมฺโม เหตุสฺส เจว สเหตุกสฺส จ สเหตุกสฺส เจว น จ เหตุสฺส จ ธมฺมสฺส อธิปติปจฺจเยน ปจฺจโย.\csromanspacer{} \textbf{อารมฺมณาธิปติ}, สหชาตาธิปติ.\csromanspacer{}\csromanspacer{}\textbf{อารมฺมณาธิปติ\csromandash{}}ทานํ ทตฺวา. (ปฐมคมนํเยว.) \textbf{สหชาตาธิปติ\csromandash{}}สเหตุโก เจว น จ เหตุ อธิปติ สมฺปยุตฺตกานํ ขนฺธานํ เหตูนญฺจ อธิปติปจฺจเยน ปจฺจโย. (3)}

\proseitem{141}{เหตุ เจว สเหตุโก จ สเหตุโก เจว น จ เหตุ จ ธมฺมา เหตุสฺส เจว สเหตุกสฺส จ ธมฺมสฺส อธิปติปจฺจเยน ปจฺจโย.\csromanspacer{} \textbf{อารมฺมณาธิปติ\csromandash{}}เหตุญฺจ สมฺปยุตฺตเก จ ขนฺเธ ครุํ กตฺวา เหตู อุปฺปชฺชนฺติ. (1)}

\prose{เหตุ เจว สเหตุโก จ สเหตุโก เจว น จ เหตุ จ ธมฺมา สเหตุกสฺส เจว น จ เหตุสฺส ธมฺมสฺส อธิปติปจฺจเยน ปจฺจโย. \csromanfolio{62}\textbf{อารมฺมณาธิปติ}\csromandash{}เหตุญฺจ สมฺปยุตฺตเก จ ขนฺเธ ครุํ กตฺวา สเหตุกา เจว น จ เหตู ขนฺธา อุปฺปชฺชนฺติ. (2)}

\prose{เหตุ เจว สเหตุโก จ สเหตุโก เจว น จ เหตุ จ ธมฺมา เหตุสฺส เจว สเหตุกสฺส จ สเหตุกสฺส เจว น จ เหตุสฺส จ ธมฺมสฺส อธิปติปจฺจเยน ปจฺจโย.\csromanspacer{} \textbf{อารมฺมณาธิปติ}\csromandash{}เหตุญฺจ สมฺปยุตฺตเก จ ขนฺเธ ครุํ กตฺวา เหตู จ สมฺปยุตฺตกา จ ขนฺธา อุปฺปชฺชนฺติ. (3)}

\csromanheader{2}{\textbf{อนนฺตร}}
\proseitem{142}{เหตุ เจว สเหตุโก จ ธมฺโม เหตุสฺส เจว สเหตุกสฺส จ ธมฺมสฺส อนนฺตรปจฺจเยน ปจฺจโย.\csromanspacer{} ปุริมา ปุริมา เหตู ปจฺฉิมานํ ปจฺฉิมานํ เหตูนํ อนนฺตรปจฺจเยน ปจฺจโย. (1)}

\prose{เหตุ เจว สเหตุโก จ ธมฺโม สเหตุกสฺส เจว น จ เหตุสฺส ธมฺมสฺส อนนฺตรปจฺจเยน ปจฺจโย.\csromanspacer{} ปุริมา ปุริมา เหตู ปจฺฉิมานํ ปจฺฉิมานํ สเหตุกานญฺเจว น จ เหตูนํ ขนฺธานํ อนนฺตรปจฺจเยน ปจฺจโย. (2)}

\prose{เหตุ เจว สเหตุโก จ ธมฺโม เหตุสฺส เจว สเหตุกสฺส จ สเหตุกสฺส เจว น จ เหตุสฺส จ ธมฺมสฺส อนนฺตรปจฺจเยน ปจฺจโย.\csromanspacer{} ปุริมา ปุริมา เหตู ปจฺฉิมานํ ปจฺฉิมานํ เหตูนํ สมฺปยุตฺตกานญฺจ ขนฺธานํ อนนฺตรปจฺจเยน ปจฺจโย. (3)}

\proseitem{143}{สเหตุโก เจว น จ เหตุ ธมฺโม สเหตุกสฺส เจว น จ เหตุสฺส ธมฺมสฺส อนนฺตรปจฺจเยน ปจฺจโย.\csromanspacer{} ปุริมา ปุริมา สเหตุกา เจว น จ เหตู ขนฺธา ปจฺฉิมานํ ปจฺฉิมานํ สเหตุกานญฺเจว น จ เหตูนํ ขนฺธานํ อนนฺตรปจฺจเยน ปจฺจโย.\csromanspacer{} อนุโลมํ โคตฺรภุสฺส.\csromanspacer{} อนุโลมํ โวทานสฺส ฯเปฯ.\csromanspacer{} นิโรธา วุฏฺฐหนฺตสฺส.\csromanspacer{} เนวสญฺญานาสญฺญายตนํ ผลสมาปตฺติยา อนนฺตรปจฺจเยน ปจฺจโย. (1)}

\prose{สเหตุโก เจว น จ เหตุ ธมฺโม เหตุสฺส เจว สเหตุกสฺส จ ธมฺมสฺส อนนฺตรปจฺจเยน ปจฺจโย.\csromanspacer{} ปุริมา ปุริมา สเหตุกา เจว น จ เหตู ขนฺธา ปจฺฉิมานํ ปจฺฉิมานํ เหตูนํ อนนฺตรปจฺจเยน ปจฺจโย.\csromanspacer{} อนุโลมํ โคตฺรภุสฺส. (สํขิตฺตํ.) (2)}

\csromanheader{2}{4. เหตุสเหตุทุก}
\prose{\csromanfolio{63}สเหตุโก เจว น จ เหตุ ธมฺโม เหตุสฺส เจว สเหตุกสฺส จ สเหตุกสฺส เจว น จ เหตุสฺส จ ธมฺมสฺส อนนฺตรปจฺจเยน ปจฺจโย.\csromanspacer{} ปุริมา ปุริมา สเหตุกา เจว น จ เหตู ขนฺธา ปจฺฉิมานํ ปจฺฉิมานํ เหตูนํ สมฺปยุตฺตกานญฺจ ขนฺธานํ อนนฺตรปจฺจเยน ปจฺจโย.\csromanspacer{} อนุโลมํ โคตฺรภุสฺส ฯเปฯ. (3)}

\prose{(สเหตุโก เจว น จ เหตุมูลกํ.\csromanspacer{} ตีณิปิ เอกสทิสา.)}

\proseitem{144}{เหตุ เจว สเหตุโก จ สเหตุโก เจว น จ เหตุ จ ธมฺมา เหตุสฺส เจว สเหตุกสฺส จ ธมฺมสฺส อนนฺตรปจฺจเยน ปจฺจโย.\csromanspacer{} ปุริมา ปุริมา เหตู จ สมฺปยุตฺตกา จ ขนฺธา ปจฺฉิมานํ ปจฺฉิมานํ เหตูนํ อนนฺตรปจฺจเยน ปจฺจโย. (1)}

\prose{เหตุ เจว สเหตุโก จ สเหตุโก เจว น จ เหตุ จ ธมฺมา สเหตุกสฺส เจว น จ เหตุสฺส ธมฺมสฺส อนนฺตรปจฺจเยน ปจฺจโย.\csromanspacer{} ปุริมา ปุริมา เหตู จ สมฺปยุตฺตกา จ ขนฺธา ปจฺฉิมานํ ปจฺฉิมานํ สเหตุกานญฺเจว น จ เหตูนํ ขนฺธานํ อนนฺตรปจฺจเยน ปจฺจโย. (2)}

\prose{เหตุ เจว สเหตุโก จ สเหตุโก เจว น จ เหตุ จ ธมฺมา เหตุสฺส เจว สเหตุกสฺส จ สเหตุกสฺส เจว น จ เหตุสฺส จ ธมฺมสฺส อนนฺตรปจฺจเยน ปจฺจโย.\csromanspacer{} ปุริมา ปุริมา เหตู จ สมฺปยุตฺตกา จ ขนฺธา ปจฺฉิมานํ ปจฺฉิมานํ เหตูนํ สมฺปยุตฺตกานญฺจ ขนฺธานํ อนนฺตรปจฺจเยน ปจฺจโย. (3)}

\csromanheader{2}{\textbf{สหชาตาทิ}}
\proseitem{145}{เหตุ เจว สเหตุโก จ ธมฺโม เหตุสฺส เจว สเหตุกสฺส จ ธมฺมสฺส สหชาตปจฺจเยน ปจฺจโย.\csromanspacer{} อญฺญมญฺญปจฺจเยน ปจฺจโย.\csromanspacer{} นิสฺสยปจฺจเยน ปจฺจโย. (ตีณิปิ ปจฺจยา ปฏิจฺจวาเร เหตุสทิสา.)}

\csromanheader{2}{\textbf{อุปนิสฺสย}}
\proseitem{146}{เหตุ เจว สเหตุโก จ ธมฺโม เหตุสฺส เจว สเหตุกสฺส จ ธมฺมสฺส อุปนิสฺสยปจฺจเยน ปจฺจโย.\csromanspacer{} \textbf{อารมฺมณูปนิสฺสโย, อนนฺตรูปนิสฺสโย, ปกตูปนิสฺสโย ฯเปฯ.}\csromanspacer{} ปกตูปนิสฺสโย\csromandash{}เหตู เหตูนํ อุปนิสฺสยปจฺจเยน ปจฺจโย. (มูลํ กาตพฺพํ.) เหตู \csromanfolio{64}สเหตุกานญฺเจว น จ เหตูนํ ขนฺธานํ อุปนิสฺสยปจฺจเยน ปจฺจโย. (มูลํ กาตพฺพํ.) เหตู เหตูนํ สมฺปยุตฺตกานญฺจ ขนฺธานํ อุปนิสฺสยปจฺจเยน ปจฺจโย. (อิเมสํ ทฺวินฺนมฺปิ ปญฺหานํ มูลานิ ปุจฺฉิตพฺพานิ.)}

\prose{สเหตุโก เจว น จ เหตุ ธมฺโม สเหตุกสฺส เจว น จ เหตุสฺส จ ธมฺมสฺส อุปนิสฺสยปจฺจเยน ปจฺจโย.\csromanspacer{} \textbf{อารมฺมณูปนิสฺสโย, อนนฺตรูปนิสฺสโย, ปกตูปนิสฺสโย ฯเปฯ.}\csromanspacer{} ปกตูปนิสฺสโย\csromandash{}สทฺธํ อุปนิสฺสาย ทานํ เทติ ฯเปฯ สมาปตฺติํ อุปฺปาเทติ, มานํ ชปฺเปติ, ทิฏฺฐิํ คณฺหาติ, สีลํ ฯเปฯ ปตฺถนํ อุปนิสฺสาย ทานํ เทติ ฯเปฯ สํฆํ ภินฺทติ.\csromanspacer{} สทฺธา ฯเปฯ.\csromanspacer{} ปตฺถนา สทฺธาย ฯเปฯ ปตฺถนาย มคฺคสฺส ผลสมาปตฺติยา อุปนิสฺสยปจฺจเยน ปจฺจโย. (1)}

\prose{(สเหตุโก เจว น จ เหตุมูลเก อิมินากาเรน วิตฺถาเรตพฺพา.\csromanspacer{} อวเสสา เทฺว ปญฺหา.)}

\prose{เหตุ เจว สเหตุโก จ สเหตุโก เจว น จ เหตุ จ ธมฺมา เหตุสฺส เจว สเหตุกสฺส จ ธมฺมสฺส อุปนิสฺสยปจฺจเยน ปจฺจโย.\csromanspacer{} \textbf{อารมฺมณูปนิสฺสโย, อนนฺตรูปนิสฺสโย, ปกตูปนิสฺสโย ฯเปฯ.}\csromanspacer{} ปกตูปนิสฺสโย\csromandash{}เหตู จ สมฺปยุตฺตกา จ ขนฺธา เหตูนํ อุปนิสฺสยปจฺจเยน ปจฺจโย. (เทฺว มูลานิ ปุจฺฉิตพฺพานิ.) เหตู จ สมฺปยุตฺตกา จ ขนฺธา สเหตุกานญฺเจว น จ เหตูนํ ขนฺธานํ อุปนิสฺสยปจฺจเยน ปจฺจโย. (มูลํ ปุจฺฉิตพฺพํ.) เหตู จ สมฺปยุตฺตกา จ ขนฺธา เหตูนํ สมฺปยุตฺตกานญฺจ ขนฺธานํ อุปนิสฺสยปจฺจเยน ปจฺจโย. (1)}

\csromanheader{2}{\textbf{อาเสวน}}
\proseitem{147}{เหตุ เจว สเหตุโก จ ธมฺโม เหตุสฺส เจว สเหตุกสฺส จ ธมฺมสฺส อาเสวนปจฺจเยน ปจฺจโย. (อนนฺตรสทิสํ.)}

\csromanheader{2}{\textbf{กมฺม}}
\proseitem{148}{สเหตุโก เจว น จ เหตุ ธมฺโม สเหตุกสฺส เจว น จ เหตุสฺส ธมฺมสฺส กมฺมปจฺจเยน ปจฺจโย.\csromanspacer{} \textbf{สหชาตา}, นานากฺขณิกา.\csromanspacer{}\csromanspacer{}\textbf{สหชาตา} สเหตุกา เจว น จ เหตู เจตนา สมฺปยุตฺตกานํ ขนฺธานํ กมฺมปจฺจเยน ปจฺจโย.\csromanspacer{} ปฏิสนฺธิกฺขเณ ฯเปฯ.}

\vspace{\tipitakaparskip}
\csromancenter{เหตุสเหตุทุก \textbf{นานากฺขณิกา} สเหตุกา เจว น จ เหตู เจตนา วิปากานํ สเหตุกานญฺเจว น จ เหตูนํ ขนฺธานํ กมฺมปจฺจเยน ปจฺจโย. (1)}
\prose{\csromanfolio{65}สเหตุโก เจว น จ เหตุ ธมฺโม เหตุสฺส เจว สเหตุกสฺส จ ธมฺมสฺส กมฺมปจฺจเยน ปจฺจโย.\csromanspacer{} \textbf{สหชาตา}, นานากฺขณิกา.\csromanspacer{}\csromanspacer{}\textbf{สหชาตา} สเหตุกา เจว น จ เหตู เจตนา สมฺปยุตฺตกานํ เหตูนํ กมฺมปจฺจเยน ปจฺจโย.\csromanspacer{} ปฏิสนฺธิกฺขเณ ฯเปฯ.\csromanspacer{} \textbf{นานากฺขณิกา} สเหตุกา เจว น จ เหตู เจตนา วิปากานํ เหตูนํ กมฺมปจฺจเยน ปจฺจโย. (2)}

\prose{สเหตุโก เจว น จ เหตุ ธมฺโม เหตุสฺส เจว สเหตุกสฺส จ สเหตุกสฺส เจว น จ เหตุสฺส จ ธมฺมสฺส กมฺมปจฺจเยน ปจฺจโย.\csromanspacer{} \textbf{สหชาตา}, นานากฺขณิกา.\csromanspacer{}\csromanspacer{}\textbf{สหชาตา} สเหตุกา เจว น จ เหตู เจตนา สมฺปยุตฺตกานํ ขนฺธานํ เหตูนญฺจ กมฺมปจฺจเยน ปจฺจโย.\csromanspacer{} ปฏิสนฺธิกฺขเณ ฯเปฯ.\csromanspacer{} \textbf{นานากฺขณิกา} สเหตุกา เจว น จ เหตู เจตนา วิปากานํ ขนฺธานํ เหตูนญฺจ กมฺมปจฺจเยน ปจฺจโย. (3)}

\csromanheader{2}{\textbf{วิปาก}}
\proseitem{149}{เหตุ เจว สเหตุโก จ ธมฺโม เหตุสฺส เจว สเหตุกสฺส จ ธมฺมสฺส วิปากปจฺจเยน ปจฺจโย.\csromanspacer{} วิปาโก อโลโภ อโทสสฺส อโมหสฺส จ วิปากปจฺจเยน ปจฺจโย. (จกฺกํ.) ปฏิสนฺธิกฺขเณ อโลโภ. (ยถา เหตุปจฺจยา, เอวํ วิตฺถาเรตพฺพํ.\csromanspacer{} นวปิ ‘วิปากนฺ’ติ นิยาเมตพฺพํ.)}

\csromanheader{2}{\textbf{อาหาราทิ}}
\proseitem{150}{สเหตุโก เจว น จ เหตุ ธมฺโม สเหตุกสฺส เจว น จ เหตุสฺส ธมฺมสฺส อาหารปจฺจเยน ปจฺจโย.\csromanspacer{} ตีณิ.}

\prose{เหตุ เจว สเหตุโก จ ธมฺโม เหตุสฺส เจว สเหตุกสฺส จ ธมฺมสฺส อินฺทฺริยปจฺจเยน ปจฺจโย. (‘อินฺทฺริยนฺ’ติ นิยาเมตพฺพํ.\csromanspacer{} นวปิ ปริปุณฺณํ.)}

\prose{สเหตุโก เจว น จ เหตุ ธมฺโม สเหตุกสฺส เจว น จ เหตุสฺส ธมฺมสฺส ฌานปจฺจเยน ปจฺจโย.\csromanspacer{} ตีณิ.}

\prose{\csromanfolio{66}เหตุ เจว สเหตุโก จ ธมฺโม เหตุสฺส เจว สเหตุกสฺส จ ธมฺมสฺส มคฺคปจฺจเยน ปจฺจโย.\csromanspacer{} สมฺปยุตฺตปจฺจเยน ปจฺจโย.\csromanspacer{} อตฺถิปจฺจเยน ปจฺจโย.\csromanspacer{} นตฺถิปจฺจเยน ปจฺจโย.\csromanspacer{} วิคตปจฺจเยน ปจฺจโย.\csromanspacer{} อวิคตปจฺจเยน ปจฺจโย.}

\csromanheader{2}{1. ปจฺจยานุโลม 2. สงฺขฺยาวาร}
\csromanheader{2}{\textbf{สุทฺธ}}
\proseitem{151}{เหตุยา ตีณิ, อารมฺมเณ นว, อธิปติยา นว, อนนฺตเร นว, สมนนฺตเร นว, สหชาเต นว, อญฺญมญฺเญ นว, นิสฺสเย นว, อุปนิสฺสเย นว, อาเสวเน นว, กมฺเม ตีณิ, วิปาเก นว, อาหาเร ตีณิ, อินฺทฺริเย นว, ฌาเน ตีณิ, มคฺเค นว, สมฺปยุตฺเต นว, อตฺถิยา นว, นตฺถิยา นว, วิคเต นว, อวิคเต นว. (เอวํ คเณตพฺพํ.)}

\csromanheader{2}{อนุโลมํ.}
\csromansectionrule
\csromanheader{2}{2. ปจฺจนียุทฺธาร}
\proseitem{152}{เหตุ เจว สเหตุโก จ ธมฺโม เหตุสฺส เจว สเหตุกสฺส จ ธมฺมสฺส อารมฺมณปจฺจเยน ปจฺจโย.\csromanspacer{} สหชาตปจฺจเยน ปจฺจโย.\csromanspacer{} อุปนิสฺสยปจฺจเยน ปจฺจโย. (1)}

\prose{เหตุ เจว สเหตุโก จ ธมฺโม สเหตุกสฺส เจว น จ เหตุสฺส ธมฺมสฺส อารมฺมณปจฺจเยน ปจฺจโย.\csromanspacer{} สหชาตปจฺจเยน ปจฺจโย.\csromanspacer{} อุปนิสฺสยปจฺจเยน ปจฺจโย. (2)}

\prose{เหตุ เจว สเหตุโก จ ธมฺโม เหตุสฺส เจว สเหตุกสฺส จ สเหตุกสฺส เจว น จ เหตุสฺส จ ธมฺมสฺส อารมฺมณปจฺจเยน ปจฺจโย.\csromanspacer{} สหชาตปจฺจเยน ปจฺจโย.\csromanspacer{} อุปนิสฺสยปจฺจเยน ปจฺจโย. (3)}

\proseitem{153}{สเหตุโก เจว น จ เหตุ ธมฺโม สเหตุกสฺส เจว น จ เหตุสฺส ธมฺมสฺส อารมฺมณปจฺจเยน ปจฺจโย.\csromanspacer{} สหชาตปจฺจเยน ปจฺจโย.\csromanspacer{} อุปนิสฺสยปจฺจเยน ปจฺจโย.\csromanspacer{} กมฺมปจฺจเยน ปจฺจโย. (1)}

\csromanheader{2}{4. เหตุสเหตุทุก}
\prose{\csromanfolio{67}สเหตุโก เจว น จ เหตุ ธมฺโม เหตุสฺส เจว สเหตุกสฺส จ ธมฺมสฺส อารมฺมณปจฺจเยน ปจฺจโย.\csromanspacer{} สหชาตปจฺจเยน ปจฺจโย.\csromanspacer{} อุปนิสฺสยปจฺจเยน ปจฺจโย.\csromanspacer{} กมฺมปจฺจเยน ปจฺจโย. (2)}

\prose{สเหตุโก เจว น จ เหตุ ธมฺโม เหตุสฺส เจว สเหตุกสฺส จ สเหตุกสฺส เจว น จ เหตุสฺส จ ธมฺมสฺส อารมฺมณปจฺจเยน ปจฺจโย.\csromanspacer{} สหชาตปจฺจเยน ปจฺจโย.\csromanspacer{} อุปนิสฺสยปจฺจเยน ปจฺจโย.\csromanspacer{} กมฺมปจฺจเยน ปจฺจโย. (3)}

\proseitem{154}{เหตุ เจว สเหตุโก จ สเหตุโก เจว น จ เหตุ จ ธมฺมา เหตุสฺส เจว สเหตุกสฺส จ ธมฺมสฺส อารมฺมณปจฺจเยน ปจฺจโย.\csromanspacer{} สหชาตปจฺจเยน ปจฺจโย.\csromanspacer{} อุปนิสฺสยปจฺจเยน ปจฺจโย. (1)}

\prose{เหตุ เจว สเหตุโก จ สเหตุโก เจว น จ เหตุ จ ธมฺมา สเหตุกสฺส เจว น จ เหตุสฺส ธมฺมสฺส อารมฺมณปจฺจเยน ปจฺจโย.\csromanspacer{} สหชาตปจฺจเยน ปจฺจโย.\csromanspacer{} อุปนิสฺสยปจฺจเยน ปจฺจโย. (2)}

\prose{เหตุ เจว สเหตุโก จ สเหตุโก เจว น จ เหตุ จ ธมฺมา เหตุสฺส เจว สเหตุกสฺส จ สเหตุกสฺส เจว น จ เหตุสฺส จ ธมฺมสฺส อารมฺมณปจฺจเยน ปจฺจโย.\csromanspacer{} สหชาตปจฺจเยน ปจฺจโย.\csromanspacer{} อุปนิสฺสยปจฺจเยน ปจฺจโย. (3)}

\csromanheader{2}{2. ปจฺจยปจฺจนีย 2. สงฺขฺยาวาร}
\proseitem{155}{นเหตุยา นว. (สํขิตฺตํ.\csromanspacer{} สพฺพตฺถ นว.\csromanspacer{} เอวํ คเณตพฺพํ.)}

\vspace{\tipitakaparskip}
\csromancenter{ปจฺจนียํ.}
\csromanheader{2}{3. ปจฺจยานุโลมปจฺจนีย}
\csromanheader{2}{\textbf{เหตุทุก}}
\proseitem{156}{เหตุปจฺจยา น-อารมฺมเณ ตีณิ, น-อธิปติยา ตีณิ, น-อนนฺตเร ตีณิ, นสมนนฺตเร ตีณิ, น-อุปนิสฺสเย ตีณิ, (สํขิตฺตํ.\csromanspacer{} สพฺพตฺถ \csromanfolio{68}ตีณิ.) นมคฺเค ตีณิ, โนนตฺถิยา ตีณิ, โนวิคเต ตีณิ. (เอวํ คเณตพฺพํ.)}

\vspace{\tipitakaparskip}
\csromancenter{อนุโลมปจฺจนียํ.}
\csromanheader{2}{4. ปจฺจยปจฺจนียานุโลม}
\csromanheader{2}{\textbf{นเหตุทุก}}
\proseitem{157}{นเหตุปจฺจยา อารมฺมเณ นว, อธิปติยา นว, อนนฺตเร นว, สมนนฺตเร นว, สหชาเต ตีณิ, อญฺญมญฺเญ ตีณิ, นิสฺสเย ตีณิ, อุปนิสฺสเย นว, อาเสวเน นว, กมฺเม ตีณิ, วิปาเก ตีณิ, อาหาเร ตีณิ, อินฺทฺริเย ตีณิ, ฌาเน ตีณิ, มคฺเค ตีณิ, สมฺปยุตฺเต ตีณิ, อตฺถิยา ตีณิ, นตฺถิยา นว, วิคเต นว, อวิคเต ตีณิ. (เอวํ คเณตพฺพํ.)}

\vspace{\tipitakaparskip}
\csromancenter{ปจฺจนียานุโลมํ.}
\nitthitam{เหตุสเหตุกทุกํ นิฏฺฐิตํ.}
\csromanmatikaanchor{mtk.7}
\csromantocmarkref{section}{5. เหตุเหตุสมฺปยุตฺตทุก}{mtk.7}
\bookmark[dest={mtk.7},level=1]{5. เหตุเหตุสมฺปยุตฺตทุก}
\csromanheader{1}{5. เหตุเหตุสมฺปยุตฺตทุก 1. ปฏิจฺจวาร}
\proseitem{158}{เหตุญฺเจว เหตุสมฺปยุตฺตญฺจ ธมฺมํ ปฏิจฺจ เหตุ เจว เหตุสมฺปยุตฺโต จ ธมฺโม อุปฺปชฺชติ เหตุปจฺจยา.\csromanspacer{} อโลภํ ปฏิจฺจ อโทโส อโมโห. (จกฺกํ.) โลภํ ปฏิจฺจ โมโห. (จกฺกํ.) ปฏิสนฺธิกฺขเณ ฯเปฯ.}

\prose{(ยถา เหตุสเหตุกทุกํ, เอวํ วิตฺถาเรตพฺพํ, นินฺนานากรณํ.)}

\nitthitam{เหตุเหตุสมฺปยุตฺตทุกํ นิฏฺฐิตํ.}
\csromanmatikaanchor{mtk.8}
\csromantocmarkref{section}{6. นเหตุสเหตุกทุก}{mtk.8}
\bookmark[dest={mtk.8},level=1]{6. นเหตุสเหตุกทุก}
\csromanheader{1}{6. นเหตุสเหตุทุก \textbf{6. นเหตุสเหตุกทุก 1. ปฏิจฺจวาร}}
\csromanheader{2}{1. ปจฺจยานุโลม 1. วิภงฺควาร}
\csromanheader{2}{\textbf{เหตุ}}
\proseitem{159}{\csromanfolio{69}นเหตุํ สเหตุกํ ธมฺมํ ปฏิจฺจ นเหตุ สเหตุโก ธมฺโม อุปฺปชฺชติ เหตุปจฺจยา.\csromanspacer{} นเหตุํ สเหตุกํ เอกํ ขนฺธํ ปฏิจฺจ ตโย ขนฺธา ฯเปฯ เทฺว ขนฺเธ ฯเปฯ.\csromanspacer{} ปฏิสนฺธิกฺขเณ ฯเปฯ. (1)}

\prose{นเหตุํ สเหตุกํ ธมฺมํ ปฏิจฺจ นเหตุ อเหตุโก ธมฺโม อุปฺปชฺชติ เหตุปจฺจยา.\csromanspacer{} นเหตู สเหตุเก ขนฺเธ ปฏิจฺจ จิตฺตสมุฏฺฐานํ รูปํ.\csromanspacer{} ปฏิสนฺธิกฺขเณ ฯเปฯ. (2)}

\prose{นเหตุํ สเหตุกํ ธมฺมํ ปฏิจฺจ นเหตุ สเหตุโก จ นเหตุ อเหตุโก จ ธมฺมา อุปฺปชฺชนฺติ เหตุปจฺจยา.\csromanspacer{} นเหตุํ สเหตุกํ เอกํ ขนฺธํ ปฏิจฺจ ตโย ขนฺธา จิตฺตสมุฏฺฐานญฺจ รูปํ ฯเปฯ เทฺว ขนฺเธ ฯเปฯ.\csromanspacer{} ปฏิสนฺธิกฺขเณ ฯเปฯ. (3)}

\proseitem{160}{นเหตุํ อเหตุกํ ธมฺมํ ปฏิจฺจ นเหตุ อเหตุโก ธมฺโม อุปฺปชฺชติ เหตุปจฺจยา ฯเปฯ.\csromanspacer{} เอกํ มหาภูตํ ปฏิจฺจ ฯเปฯ มหาภูเต ปฏิจฺจ จิตฺตสมุฏฺฐานํ รูปํ, กฏตฺตารูปํ, อุปาทารูปํ. (1)}

\prose{นเหตุํ อเหตุกํ ธมฺมํ ปฏิจฺจ นเหตุ สเหตุโก ธมฺโม อุปฺปชฺชติ เหตุปจฺจยา.\csromanspacer{} ปฏิสนฺธิกฺขเณ วตฺถุํ ปฏิจฺจ นเหตู สเหตุกา ขนฺธา. (2)}

\prose{นเหตุํ อเหตุกํ ธมฺมํ ปฏิจฺจ นเหตุ สเหตุโก จ นเหตุ อเหตุโก จ ธมฺมา อุปฺปชฺชนฺติ เหตุปจฺจยา.\csromanspacer{} ปฏิสนฺธิกฺขเณ วตฺถุํ ปฏิจฺจ นเหตู สเหตุกา ขนฺธา.\csromanspacer{} มหาภูเต ปฏิจฺจ กฏตฺตารูปํ. (3)}

\proseitem{161}{นเหตุํ สเหตุกญฺจ นเหตุํ อเหตุกญฺจ ธมฺมํ ปฏิจฺจ นเหตุ สเหตุโก ธมฺโม อุปฺปชฺชติ เหตุปจฺจยา.\csromanspacer{} ปฏิสนฺธิกฺขเณ นเหตุํ สเหตุกํ เอกํ ขนฺธญฺจ วตฺถุญฺจ ปฏิจฺจ ตโย ขนฺธา ฯเปฯ เทฺว ขนฺเธ จ ฯเปฯ. (1)}

\prose{นเหตุํ สเหตุกญฺจ นเหตุํ อเหตุกญฺจ ธมฺมํ ปฏิจฺจ นเหตุ อเหตุโก ธมฺโม อุปฺปชฺชติ เหตุปจฺจยา.\csromanspacer{} นเหตู สเหตุเก ขนฺเธ จ มหาภูเต จ ปฏิจฺจ จิตฺตสมุฏฺฐานํ รูปํ.\csromanspacer{} ปฏิสนฺธิกฺขเณ ฯเปฯ. (2)}

\prose{\csromanfolio{70}นเหตุํ สเหตุกญฺจ นเหตุํ อเหตุกญฺจ ธมฺมํ ปฏิจฺจ นเหตุ สเหตุโก จ นเหตุ อเหตุโก จ ธมฺมา อุปฺปชฺชนฺติ เหตุปจฺจยา.\csromanspacer{} ปฏิสนฺธิกฺขเณ นเหตุํ สเหตุกํ เอกํ ขนฺธญฺจ วตฺถุญฺจ ปฏิจฺจ ตโย ขนฺธา ฯเปฯ เทฺว ขนฺเธ ฯเปฯ.\csromanspacer{} นเหตู สเหตุเก ขนฺเธ จ มหาภูเต จ ปฏิจฺจ กฏตฺตารูปํ. (3)}

\csromanheader{2}{\textbf{อารมฺมณ}}
\proseitem{162}{นเหตุํ สเหตุกํ ธมฺมํ ปฏิจฺจ นเหตุ สเหตุโก ธมฺโม อุปฺปชฺชติ อารมฺมณปจฺจยา.\csromanspacer{} นเหตุํ สเหตุกํ เอกํ ขนฺธํ ปฏิจฺจ ตโย ขนฺธา ฯเปฯ เทฺว ขนฺเธ ฯเปฯ.\csromanspacer{} ปฏิสนฺธิกฺขเณ ฯเปฯ. (1)}

\prose{นเหตุํ อเหตุกํ ธมฺมํ ปฏิจฺจ นเหตุ อเหตุโก ธมฺโม อุปฺปชฺชติ อารมฺมณปจฺจยา.\csromanspacer{} นเหตุํ อเหตุกํ เอกํ ขนฺธํ ปฏิจฺจ ตโย ขนฺธา ฯเปฯ.\csromanspacer{} ปฏิสนฺธิกฺขเณ ฯเปฯ. (2)}

\prose{นเหตุํ อเหตุกํ ธมฺมํ ปฏิจฺจ นเหตุ สเหตุโก ธมฺโม อุปฺปชฺชติ อารมฺมณปจฺจยา.\csromanspacer{} ปฏิสนฺธิกฺขเณ วตฺถุํ ปฏิจฺจ นเหตู สเหตุกา ขนฺธา. (3)}

\prose{นเหตุํ สเหตุกญฺจ นเหตุํ อเหตุกญฺจ ธมฺมํ ปฏิจฺจ นเหตุ สเหตุโก ธมฺโม อุปฺปชฺชติ อารมฺมณปจฺจยา.\csromanspacer{} ปฏิสนฺธิกฺขเณ นเหตุํ สเหตุกํ เอกํ ขนฺธญฺจ วตฺถุญฺจ ปฏิจฺจ ตโย ขนฺธา ฯเปฯ เทฺว ขนฺเธ ฯเปฯ. (สํขิตฺตํ.\csromanspacer{} เอวํ วิภชิตพฺพํ.)}

\csromanheader{2}{1. ปจฺจยานุโลม 2. สงฺขฺยาวาร}
\csromanheader{2}{\textbf{สุทฺธ}}
\proseitem{163}{เหตุยา นว, อารมฺมเณ จตฺตาริ, อธิปติยา ปญฺจ, อนนฺตเร จตฺตาริ, สมนนฺตเร จตฺตาริ, สหชาเต นว, อญฺญมญฺเญ ฉ, นิสฺสเย นว, อุปนิสฺสเย จตฺตาริ, ปุเรชาเต เทฺว, อาเสวเน เทฺว, กมฺเม นว, วิปาเก นว, อาหาเร นว, (สํขิตฺตํ.\csromanspacer{} สพฺพตฺถ นว.) สมฺปยุตฺเต จตฺตาริ, วิปฺปยุตฺเต นว, อตฺถิยา นว, นตฺถิยา จตฺตาริ, วิคเต จตฺตาริ, อวิคเต นว. (เอวํ คเณตพฺพํ.)}

\vspace{\tipitakaparskip}
\csromancenter{อนุโลมํ.}
\csromanheader{2}{6. \textbf{นเหตุ}สเหตุทุก \textbf{2. ปจฺจยปจฺจนีย 1. วิภงฺควาร}}
\csromanheader{2}{\textbf{นเหตุ}}
\vspace{\tipitakaparskip}
\csromancenter{\textbf{นเหตุ}ํ อเหตุกํ ธมฺมํ ปฏิจฺจ นเหตุ อเหตุโก ธมฺโม อุปฺปชฺชติ นเหตุปจฺจยา.\csromanspacer{} \textbf{นเหตุ}ํ อเหตุกํ เอกํ ขนฺธํ ปฏิจฺจ ตโย ขนฺธา จิตฺตสมุฏฺฐานญฺจ รูปํ ฯเปฯ เทฺว ขนฺเธ ฯเปฯ.\csromanspacer{} ปฏิสนฺธิกฺขเณ. (ยาว อสญฺญสตฺตา โมโห นตฺถิ.) (1)}
\csromanheader{2}{\textbf{น-อารมฺมณ}}
\vspace{\tipitakaparskip}
\csromancenter{\textbf{นเหตุ}ํ สเหตุกํ ธมฺมํ ปฏิจฺจ นเหตุ อเหตุโก ธมฺโม อุปฺปชฺชติ น-อารมฺมณปจฺจยา.\csromanspacer{} นเหตู สเหตุเก ขนฺเธ ปฏิจฺจ จิตฺตสมุฏฺฐานํ รูปํ.\csromanspacer{} ปฏิสนฺธิกฺขเณ ฯเปฯ. (1)}
\csromancenter{\textbf{นเหตุ}ํ อเหตุกํ ธมฺมํ ปฏิจฺจ นเหตุ อเหตุโก ธมฺโม อุปฺปชฺชติ น-อารมฺมณปจฺจยา.\csromanspacer{} นเหตู อเหตุเก ขนฺเธ ปฏิจฺจ จิตฺตสมุฏฺฐานํ รูปํ.\csromanspacer{} ปฏิสนฺธิกฺขเณ ฯเปฯ. (ยาว อสญฺญสตฺตา.) (1)}
\csromancenter{\textbf{นเหตุ}ํ สเหตุกญฺจ นเหตุํ อเหตุกญฺจ ธมฺมํ ปฏิจฺจ นเหตุ อเหตุโก ธมฺโม อุปฺปชฺชติ น-อารมฺมณปจฺจยา.\csromanspacer{} นเหตู สเหตุเก ขนฺเธ จ มหาภูเต จ ปฏิจฺจ จิตฺตสมุฏฺฐานํ รูปํ.\csromanspacer{} ปฏิสนฺธิกฺขเณ ฯเปฯ. (สํขิตฺตํ.)}
\csromanheader{2}{2. ปจฺจยปจฺจนีย 2. สงฺขฺยาวาร}
\csromanheader{2}{\textbf{สุทฺธ}}
\vspace{\tipitakaparskip}
\csromancenter{\textbf{นเหตุ}ยา เอกํ, น-อารมฺมเณ ตีณิ, น-อธิปติยา นว, น-อนนฺตเร ตีณิ, นสมนนฺตเร ตีณิ, น-อญฺญมญฺเญ ตีณิ, น-อุปนิสฺสเย ตีณิ, นปุเรชาเต นว, นปจฺฉาชาเต นว, น-อาเสวเน นว, นกมฺเม เทฺว, นวิปาเก ปญฺจ, น-อาหาเร เอกํ, น-อินฺทฺริเย เอกํ, นฌาเน เอกํ, นมคฺเค เอกํ, นสมฺปยุตฺเต ตีณิ, นวิปฺปยุตฺเต เทฺว, โนนตฺถิยา ตีณิ, โนวิคเต ตีณิ. (เอวํ คเณตพฺพํ.)}
\csromancenter{ปจฺจนียํ.}
\csromanheader{2}{3. ปจฺจยานุโลมปจฺจนีย}
\csromanheader{2}{\textbf{เหตุทุก}}
\proseitem{167}{\csromanfolio{72}เหตุปจฺจยา น-อารมฺมเณ ตีณิ, น-อธิปติยา นว, น-อนนฺตเร นว, นสมนนฺตเร นว, น-อญฺญมญฺเญ นว, น-อุปนิสฺสเย ตีณิ, นปุเรชาเต นว, นปจฺฉาชาเต นว, น-อาเสวเน นว, นกมฺเม เอกํ, นวิปาเก ปญฺจ, นสมฺปยุตฺเต ตีณิ, นวิปฺปยุตฺเต เอกํ, โนนตฺถิยา ตีณิ, โนวิคเต ตีณิ. (เอวํ คเณตพฺพํ.)}

\vspace{\tipitakaparskip}
\csromancenter{อนุโลมปจฺจนียํ.}
\csromanheader{2}{4. ปจฺจยปจฺจนียานุโลม}
\csromanheader{2}{\textbf{นเหตุทุก}}
\proseitem{168}{นเหตุปจฺจยา อารมฺมเณ เอกํ ฯเปฯ อาหาเร เอกํ ฯเปฯ ฌาเน เอกํ, สมฺปยุตฺเต เอกํ, วิปฺปยุตฺเต เอกํ, ฯเปฯ วิคเต เอกํ, อวิคเต เอกํ. (เอวํ คเณตพฺพํ.}

\vspace{\tipitakaparskip}
\csromancenter{ปจฺจนียานุโลมํ.}
\csromancenter{(สหชาตวาเรปิ เอวํ คเณตพฺพํ.)}
\csromanheader{2}{6. นเหตุสเหตุกทุก 3. ปจฺจยวาร}
\csromanheader{2}{1. ปจฺจยานุโลมาทิ}
\proseitem{169}{นเหตุํ สเหตุกํ ธมฺมํ ปจฺจยา นเหตุ สเหตุโก ธมฺโม อุปฺปชฺชติ เหตุปจฺจยา.\csromanspacer{} ตีณิ.}

\prose{นเหตุํ อเหตุกํ ธมฺมํ ปจฺจยา นเหตุ อเหตุโก ธมฺโม อุปฺปชฺชติ เหตุปจฺจยา.\csromanspacer{} เอกํ มหาภูตํ ปจฺจยา ตโย มหาภูตา.\csromanspacer{} มหาภูเต ปจฺจยา จิตฺตสมุฏฺฐานํ รูปํ, กฏตฺตารูปํ, อุปาทารูปํ. (1)}

\csromanheader{2}{6. นเหตุสเหตุทุก}
\prose{\csromanfolio{73}นเหตุํ อเหตุกํ ธมฺมํ ปจฺจยา นเหตุ สเหตุโก ธมฺโม อุปฺปชฺชติ เหตุปจฺจยา.\csromanspacer{} วตฺถุํ ปจฺจยา นเหตู สเหตุกา ขนฺธา.\csromanspacer{} ปฏิสนฺธิกฺขเณ ฯเปฯ. (2)}

\prose{นเหตุํ อเหตุกํ ธมฺมํ ปจฺจยา นเหตุ สเหตุโก จ นเหตุ อเหตุโก จ ธมฺมา อุปฺปชฺชนฺติ เหตุปจฺจยา.\csromanspacer{} วตฺถุํ ปจฺจยา นเหตู สเหตุกา ขนฺธา.\csromanspacer{} มหาภูเต ปจฺจยา จิตฺตสมุฏฺฐานํ รูปํ.\csromanspacer{} ปฏิสนฺธิกฺขเณ ฯเปฯ. (3)}

\prose{นเหตุํ สเหตุกญฺจ นเหตุํ อเหตุกญฺจ ธมฺมํ ปจฺจยา นเหตุ สเหตุโก ธมฺโม อุปฺปชฺชติ เหตุปจฺจยา. (ฆฏนา ตีณิ, ปวตฺติปฏิสนฺธิ ปริปุณฺณํ.\csromanspacer{} สํขิตฺตํ.)}

\proseitem{170}{เหตุยา นว, อารมฺมเณ จตฺตาริ ฯเปฯ อญฺญมญฺเญ ฉ ฯเปฯ ปุเรชาเต, อาเสวเน จตฺตาริ ฯเปฯ อวิคเต นว. (เอวํ คเณตพฺพํ.)}

\vspace{\tipitakaparskip}
\csromancenter{อนุโลมํ.}
\csromancenter{นเหตุยา เอกํ, น-อารมฺมเณ ตีณิ ฯเปฯ โนวิคเต ตีณิ.}
\csromancenter{ปจฺจนียํ.}
\csromancenter{(นิสฺสยวาโร ปจฺจยวารสทิโส.)}
\csromanheader{2}{6. นเหตุสเหตุกทุก 5. สํสฏฺฐวาร}
\csromanheader{2}{1. ปจฺจยานุโลมาทิ}
\proseitem{172}{นเหตุํ สเหตุกํ ธมฺมํ สํสฏฺโฐ นเหตุ สเหตุโก ธมฺโม อุปฺปชฺชติ เหตุปจฺจยา.\csromanspacer{} นเหตุํ สเหตุกํ เอกํ ขนฺธํ ฯเปฯ.\csromanspacer{} ปฏิสนฺธิกฺขเณ ฯเปฯ.}

\proseitem{173}{เหตุยา เอกํ, อารมฺมเณ เทฺว, อธิปติยา เอกํ, อนนฺตเร เทฺว, (สพฺพตฺถ เทฺว.) มคฺเค เอกํ ฯเปฯ อวิคเต เทฺว.}

\vspace{\tipitakaparskip}
\csromancenter{อนุโลมํ.}
\proseitem{174}{\csromanfolio{74}นเหตุํ อเหตุกํ ธมฺมํ สํสฏฺโฐ นเหตุ อเหตุโก ธมฺโม อุปฺปชฺชติ นเหตุปจฺจยา.\csromanspacer{} นเหตุํ อเหตุกํ เอกํ ขนฺธํ ฯเปฯ.\csromanspacer{} ปฏิสนฺธิกฺขเณ ฯเปฯ.}

\proseitem{175}{นเหตุยา เอกํ, น-อธิปติยา เทฺว, นปุเรชาเต เทฺว, นปจฺฉาชาเต เทฺว, น-อาเสวเน เทฺว, นกมฺเม เทฺว, นวิปาเก เทฺว, นฌาเน เอกํ, นมคฺเค เอกํ, นวิปฺปยุตฺเต เทฺว.\csromanspacer{} ปจฺจนียํ. (เอวํ อวเสสาปิ เทฺว คณนา คเณตพฺพา.) (สมฺปยุตฺตวาโร สํสฏฺฐวารสทิโส.)}

\csromanheader{2}{6. นเหตุสเหตุกทุก 7. ปญฺหาวาร}
\csromanheader{2}{1. ปจฺจยานุโลม 1. วิภงฺควาร}
\csromanheader{2}{\textbf{อารมฺมณ}}
\proseitem{176}{นเหตุ สเหตุโก ธมฺโม นเหตุสเหตุกสฺส ธมฺมสฺส อารมฺมณปจฺจเยน ปจฺจโย.\csromanspacer{} ทานํ ทตฺวา, สีลํ ฯเปฯ อุโปสถกมฺมํ กตฺวา ตํ ปจฺจเวกฺขติ.\csromanspacer{} ปุพฺเพ สุจิณฺณานิ ปจฺจเวกฺขติ.\csromanspacer{} ฌานํ ฯเปฯ.\csromanspacer{} อริยา มคฺคา วุฏฺฐหิตฺวา มคฺคํ ปจฺจเวกฺขนฺติ, ผลํ ปจฺจเวกฺขนฺติ.\csromanspacer{} ปหีเน กิเลเส ฯเปฯ.\csromanspacer{} วิกฺขมฺภิเต กิเลเส ฯเปฯ.\csromanspacer{} ปุพฺเพ ฯเปฯ.\csromanspacer{} นเหตู สเหตุเก ขนฺเธ อนิจฺจโต ฯเปฯ โทมนสฺสํ อุปฺปชฺชติ.\csromanspacer{} กุสลากุสเล นิรุทฺเธ นเหตุ สเหตุโก วิปาโก ตทารมฺมณตา อุปฺปชฺชติ.\csromanspacer{} เจโตปริยญาเณน นเหตุสเหตุกจิตฺตสมงฺคิสฺส จิตฺตํ ชานาติ.\csromanspacer{} อากาสานญฺจายตนํ ฯเปฯ.\csromanspacer{} อากิญฺจญฺญายตนํ ฯเปฯ.\csromanspacer{} นเหตู สเหตุกา ขนฺธา อิทฺธิวิธญาณสฺส, เจโตปริยญาณสฺส, ปุพฺเพนิวาสานุสฺสติญาณสฺส, ยถากมฺมูปคญาณสฺส, อนาคตํสญาณสฺส, อารมฺมณปจฺจเยน ปจฺจโย.\csromanspacer{} นเหตู สเหตุเก ขนฺเธ อารพฺภ นเหตู สเหตุกา ขนฺธา อุปฺปชฺชนฺติ. (1)}

\prose{นเหตุ สเหตุโก ธมฺโม นเหตุ-อเหตุกสฺส ธมฺมสฺส อารมฺมณปจฺจเยน ปจฺจโย.\csromanspacer{} นเหตู สเหตุเก ขนฺเธ อนิจฺจโต ฯเปฯ}

\vspace{\tipitakaparskip}
\csromancenter{นเหตุสเหตุทุก โทมนสฺสํ อุปฺปชฺชติ.\csromanspacer{} กุสลากุสเล นิรุทฺเธ นเหตุ อเหตุโก วิปาโก ตทารมฺมณตา อุปฺปชฺชติ.\csromanspacer{} นเหตู สเหตุเก ขนฺเธ อารพฺภ นเหตู อเหตุกา ขนฺธา อุปฺปชฺชนฺติ. (2)}
\proseitem{177}{\csromanfolio{75}นเหตุ อเหตุโก ธมฺโม นเหตุ-อเหตุกสฺส ธมฺมสฺส อารมฺมณปจฺจเยน ปจฺจโย.\csromanspacer{} นิพฺพานํ อาวชฺชนาย อารมฺมณปจฺจเยน ปจฺจโย.\csromanspacer{} จกฺขุํ ฯเปฯ วตฺถุํ.\csromanspacer{} นเหตู อเหตุเก ขนฺเธ อนิจฺจโต ฯเปฯ โทมนสฺสํ อุปฺปชฺชติ.\csromanspacer{} กุสลากุสเล นิรุทฺเธ นเหตุ อเหตุโก วิปาโก ตทารมฺมณตา อุปฺปชฺชติ.\csromanspacer{} รูปายตนํ จกฺขุวิญฺญาณสฺส ฯเปฯ.\csromanspacer{} โผฏฺฐพฺพายตนํ กายวิญฺญาณสฺส ฯเปฯ.\csromanspacer{} นเหตู อเหตุเก ขนฺเธ อารพฺภ นเหตู อเหตุกา ขนฺธา อุปฺปชฺชนฺติ. (1)}

\prose{นเหตุ อเหตุโก ธมฺโม นเหตุสเหตุกสฺส ธมฺมสฺส อารมฺมณปจฺจเยน ปจฺจโย.\csromanspacer{} อริยา นิพฺพานํ ปจฺจเวกฺขนฺติ.\csromanspacer{} นิพฺพานํ โคตฺรภุสฺส, โวทานสฺส, มคฺคสฺส, ผลสฺส อารมฺมณปจฺจเยน ปจฺจโย.\csromanspacer{} จกฺขุํ ฯเปฯ วตฺถุํ.\csromanspacer{} นเหตู อเหตุเก ขนฺเธ อนิจฺจโต ฯเปฯ โทมนสฺสํ อุปฺปชฺชติ.\csromanspacer{} กุสลากุสเล นิรุทฺเธ นเหตุ สเหตุโก วิปาโก ตทารมฺมณตา อุปฺปชฺชติ.\csromanspacer{} ทิพฺเพน จกฺขุนา รูปํ ปสฺสติ.\csromanspacer{} ทิพฺพาย โสตธาตุยา สทฺทํ สุณาติ.\csromanspacer{} เจโตปริยญาเณน นเหตุ-อเหตุกจิตฺตสมงฺคิสฺส จิตฺตํ ชานาติ.\csromanspacer{} นเหตู อเหตุกา ขนฺธา อิทฺธิวิธญาณสฺส, เจโตปริยญาณสฺส, ปุพฺเพนิวาสานุสฺสติญาณสฺส, อนาคตํสญาณสฺส, อารมฺมณปจฺจเยน ปจฺจโย.\csromanspacer{} นเหตู อเหตุเก ขนฺเธ อารพฺภ นเหตู สเหตุกา ขนฺธา อุปฺปชฺชนฺติ. (2)}

\csromanheader{2}{\textbf{อธิปติ}}
\proseitem{178}{นเหตุ สเหตุโก ธมฺโม นเหตุสเหตุกสฺส ธมฺมสฺส อธิปติปจฺจเยน ปจฺจโย.\csromanspacer{} \textbf{อารมฺมณาธิปติ}, สหชาตาธิปติ.\csromanspacer{}\csromanspacer{}\textbf{อารมฺมณาธิปติ}\csromandash{}ทานํ ทตฺวา, สีลํ สมาทิยิตฺวา, อุโปสถกมฺมํ กตฺวา ตํ ครุํ กตฺวา ปจฺจเวกฺขติ.\csromanspacer{} ปุพฺเพ สุจิณฺณานิ ครุํ กตฺวา ปจฺจเวกฺขติ.\csromanspacer{} ฌานํ ฯเปฯ.\csromanspacer{} อริยา มคฺคา วุฏฺฐหิตฺวา มคฺคํ ครุํ กตฺวา ปจฺจเวกฺขนฺติ, ผลํ ครุํ กตฺวา ปจฺจเวกฺขนฺติ.\csromanspacer{} นเหตู สเหตุเก ขนฺเธ ครุํ กตฺวา อสฺสาเทติ อภินนฺทติ, ตํ ครุํ กตฺวา ราโค อุปฺปชฺชติ, ทิฏฺฐิ อุปฺปชฺชติ.\csromanspacer{} \textbf{สหชาตาธิปติ}\csromandash{} \csromanfolio{76}นเหตุสเหตุกาธิปติ สมฺปยุตฺตกานํ ขนฺธานํ อธิปติปจฺจเยน ปจฺจโย. (1)}

\prose{นเหตุ สเหตุโก ธมฺโม นเหตุ-อเหตุกสฺส ธมฺมสฺส อธิปติปจฺจเยน ปจฺจโย.\csromanspacer{} \textbf{สหชาตาธิปติ\csromandash{}}นเหตุ สเหตุกาธิปติ จิตฺตสมุฏฺฐานานํ รูปานํ อธิปติปจฺจเยน ปจฺจโย. (2)}

\prose{นเหตุ สเหตุโก ธมฺโม นเหตุสเหตุกสฺส จ นเหตุ-อเหตุกสฺส จ ธมฺมสฺส อธิปติปจฺจเยน ปจฺจโย.\csromanspacer{} \textbf{สหชาตาธิปติ\csromandash{}}นเหตุ สเหตุกาธิปติ สมฺปยุตฺตกานํ ขนฺธานํ จิตฺตสมุฏฺฐานานญฺจ รูปานํ อธิปติปจฺจเยน ปจฺจโย. (3)}

\prose{นเหตุ อเหตุโก ธมฺโม นเหตุสเหตุกสฺส ธมฺมสฺส อธิปติ ปจฺจเยน ปจฺจโย.\csromanspacer{} \textbf{อารมฺมณาธิปติ}\csromandash{}อริยา นิพฺพานํ ครุํ กตฺวา ปจฺจเวกฺขนฺติ.\csromanspacer{} นิพฺพานํ โคตฺรภุสฺส, โวทานสฺส, มคฺคสฺส, ผลสฺส อธิปติปจฺจเยน ปจฺจโย.\csromanspacer{} จกฺขุํ ฯเปฯ วตฺถุํ.\csromanspacer{} นเหตู อเหตุเก ขนฺเธ ครุํ กตฺวา อสฺสาเทติ อภินนฺทติ, ตํ ครุํ กตฺวา ราโค อุปฺปชฺชติ, ทิฏฺฐิ อุปฺปชฺชติ. (1)}

\csromanheader{2}{\textbf{อนนฺตร}}
\proseitem{179}{นเหตุ สเหตุโก ธมฺโม นเหตุสเหตุกสฺส ธมฺมสฺส อนนฺตรปจฺจเยน ปจฺจโย.\csromanspacer{} ปุริมา ปุริมา นเหตู สเหตุกา ขนฺธา ปจฺฉิมานํ ปจฺฉิมานํ นเหตุสเหตุกานํ ขนฺธานํ อนนฺตรปจฺจเยน ปจฺจโย.\csromanspacer{} อนุโลมํ โคตฺรภุสฺส ฯเปฯ.\csromanspacer{} เนวสญฺญานาสญฺญายตนํ ผลสมาปตฺติยา อนนฺตรปจฺจเยน ปจฺจโย. (1)}

\prose{นเหตุ สเหตุโก ธมฺโม นเหตุ-อเหตุกสฺส ธมฺมสฺส อนนฺตรปจฺจเยน ปจฺจโย.\csromanspacer{} นเหตุ สเหตุกํ จุติจิตฺตํ นเหตุ-อเหตุกสฺส อุปปตฺติจิตฺตสฺส อนนฺตรปจฺจเยน ปจฺจโย.\csromanspacer{} นเหตุ สเหตุกํ ภวงฺคํ อาวชฺชนาย.\csromanspacer{} นเหตุ สเหตุกํ ภวงฺคํ นเหตุ-อเหตุกสฺส ภวงฺคสฺส.\csromanspacer{} นเหตู สเหตุกา ขนฺธา นเหตุ-อเหตุกสฺส วุฏฺฐานสฺส อนนฺตรปจฺจเยน ปจฺจโย. (2)}

\prose{นเหตุ อเหตุโก ธมฺโม นเหตุ-อเหตุกสฺส ธมฺมสฺส อนนฺตรปจฺจเยน ปจฺจโย.\csromanspacer{} ปุริมา ปุริมา นเหตู อเหตุกา ขนฺธา ปจฺฉิมานํ}

\vspace{\tipitakaparskip}
\csromancenter{นเหตุสเหตุทุก ปจฺฉิมานํ นเหตุ-อเหตุกานํ ขนฺธานํ อนนฺตรปจฺจเยน ปจฺจโย.\csromanspacer{} อาวชฺชนา ปญฺจนฺนํ วิญฺญาณานํ อนนฺตรปจฺจเยน ปจฺจโย. (1)}
\prose{\csromanfolio{77}นเหตุ อเหตุโก ธมฺโม นเหตุสเหตุกสฺส ธมฺมสฺส อนนฺตรปจฺจเยน ปจฺจโย.\csromanspacer{} นเหตุ อเหตุกํ จุติจิตฺตํ นเหตุสเหตุกสฺส อุปปตฺติจิตฺตสฺส อนนฺตรปจฺจเยน ปจฺจโย.\csromanspacer{} อาวชฺชนา นเหตุสเหตุกานํ ขนฺธานํ อนนฺตรปจฺจเยน ปจฺจโย.\csromanspacer{} นเหตู อเหตุกา ขนฺธา นเหตุสเหตุกสฺส วุฏฺฐานสฺส อนนฺตรปจฺจเยน ปจฺจโย. (2)}

\csromanheader{2}{\textbf{สมนนฺตราทิ}}
\proseitem{180}{นเหตุ สเหตุโก ธมฺโม นเหตุสเหตุกสฺส ธมฺมสฺส สมนนฺตรปจฺจเยน ปจฺจโย.\csromanspacer{} สหชาตปจฺจเยน ปจฺจโย. (อิห ฆฏนา นตฺถิ.\csromanspacer{} สตฺต ปญฺหา.) อญฺญมญฺญปจฺจเยน ปจฺจโย.\csromanspacer{} ฉ ปญฺหา.\csromanspacer{} นิสฺสยปจฺจเยน ปจฺจโย. (ปวตฺติปฏิสนฺธิ สตฺต ปญฺหา.\csromanspacer{} อิห ฆฏนา นตฺถิ.)}

\csromanheader{2}{\textbf{อุปนิสฺสย}}
\proseitem{181}{นเหตุ สเหตุโก ธมฺโม นเหตุสเหตุกสฺส ธมฺมสฺส อุปนิสฺสยปจฺจเยน ปจฺจโย.\csromanspacer{} \textbf{อารมฺมณูปนิสฺสโย, อนนฺตรูปนิสฺสโย, ปกตูปนิสฺสโย ฯเปฯ.}\csromanspacer{} ปกตูปนิสฺสโย\csromandash{}สทฺธํ อุปนิสฺสาย ทานํ เทติ ฯเปฯ มานํ ชปฺเปติ.\csromanspacer{} ทิฏฺฐิํ คณฺหาติ.\csromanspacer{} สีลํ ฯเปฯ.\csromanspacer{} ปตฺถนํ อุปนิสฺสาย ทานํ เทติ ฯเปฯ สํฆํ ภินฺทติ.\csromanspacer{} สทฺธา ฯเปฯ.\csromanspacer{} ปตฺถนา สทฺธาย ฯเปฯ ปตฺถนาย มคฺคสฺส ผลสมาปตฺติยา อุปนิสฺสยปจฺจเยน ปจฺจโย. (1)}

\prose{นเหตุ สเหตุโก ธมฺโม นเหตุ-อเหตุกสฺส ธมฺมสฺส อุปนิสฺสยปจฺจเยน ปจฺจโย.\csromanspacer{} \textbf{อนนฺตรูปนิสฺสโย, ปกตูปนิสฺสโย ฯเปฯ.}\csromanspacer{} ปกตูปนิสฺสโย\csromandash{}สทฺธา กายิกสฺส สุขสฺส, กายิกสฺส ทุกฺขสฺส อุปนิสฺสยปจฺจเยน ปจฺจโย.\csromanspacer{} สีลํ ฯเปฯ.\csromanspacer{} ปตฺถนา กายิกสฺส สุขสฺส, กายิกสฺส ทุกฺขสฺส อุปนิสฺสยปจฺจเยน ปจฺจโย.\csromanspacer{} สทฺธา ฯเปฯ.\csromanspacer{} ปตฺถนา กายิกสฺส สุขสฺส, กายิกสฺส ทุกฺขสฺส อุปนิสฺสยปจฺจเยน ปจฺจโย. (2)}

\proseitem{182}{นเหตุ อเหตุโก ธมฺโม นเหตุ-อเหตุกสฺส ธมฺมสฺส อุปนิสฺสยปจฺจเยน ปจฺจโย.\csromanspacer{} \textbf{อนนฺตรูปนิสฺสโย, ปกตูปนิสฺสโย ฯเปฯ.}\csromanspacer{} ปกตูปนิสฺสโยกายิกํ สุขํ กายิกสฺส สุขสฺส, กายิกสฺส \csromanfolio{78}ทุกฺขสฺส อุปนิสฺสยปจฺจเยน ปจฺจโย.\csromanspacer{} กายิกํ ทุกฺขํ.\csromanspacer{} อุตุ.\csromanspacer{} โภชนํ.\csromanspacer{} เสนาสนํ กายิกสฺส สุขสฺส, กายิกสฺส ทุกฺขสฺส อุปนิสฺสยปจฺจเยน ปจฺจโย.\csromanspacer{} กายิกํ สุขํ.\csromanspacer{} กายิกํ ทุกฺขํ.\csromanspacer{} อุตุ.\csromanspacer{} โภชนํ.\csromanspacer{} เสนาสนํ กายิกสฺส สุขสฺส, กายิกสฺส ทุกฺขสฺส อุปนิสฺสยปจฺจเยน ปจฺจโย. (1)}

\prose{นเหตุ อเหตุโก ธมฺโม นเหตุสเหตุกสฺส ธมฺมสฺส อุปนิสฺสยปจฺจเยน ปจฺจโย.\csromanspacer{} \textbf{อารมฺมณูปนิสฺสโย, อนนฺตรูปนิสฺสโย, ปกตูปนิสฺสโย ฯเปฯ.}\csromanspacer{} ปกตูปนิสฺสโย\csromandash{}กายิกํ สุขํ อุปนิสฺสาย ทานํ เทติ ฯเปฯ สํฆํ ภินฺทติ.\csromanspacer{} กายิกํ ทุกฺขํ.\csromanspacer{} อุตุํ.\csromanspacer{} โภชนํ.\csromanspacer{} เสนาสนํ อุปนิสฺสาย ทานํ เทติ ฯเปฯ สํฆํ ภินฺทติ.\csromanspacer{} กายิกํ สุขํ ฯเปฯ.\csromanspacer{} เสนาสนํ สทฺธาย ฯเปฯ ปตฺถนาย มคฺคสฺส ผลสมาปตฺติยา อุปนิสฺสยปจฺจเยน ปจฺจโย. (2)}

\csromanheader{2}{\textbf{ปุเรชาต}}
\proseitem{183}{นเหตุ อเหตุโก ธมฺโม นเหตุ-อเหตุกสฺส ธมฺมสฺส ปุเรชาตปจฺจเยน ปจฺจโย.\csromanspacer{} อารมฺมณปุเรชาตํ, วตฺถุปุเรชาตํ.\csromanspacer{}\csromanspacer{}\textbf{อารมฺมณปุเรชาตํ\csromandash{}}จกฺขุํ ฯเปฯ วตฺถุํ อนิจฺจโต ฯเปฯ โทมนสฺสํ อุปฺปชฺชติ.\csromanspacer{} กุสลากุสเล นิรุทฺเธ นเหตุ อเหตุโก วิปาโก ตทารมฺมณตา อุปฺปชฺชติ.\csromanspacer{} รูปายตนํ จกฺขุวิญฺญาณสฺส ปุเรชาตปจฺจเยน ปจฺจโย ฯเปฯ.\csromanspacer{} โผฏฺฐพฺพายตนํ กายวิญฺญาณสฺส ปุเรชาตปจฺจเยน ปจฺจโย.\csromanspacer{} \textbf{วตฺถุปุเรชาตํ}\csromandash{}จกฺขายตนํ จกฺขุวิญฺญาณสฺส ฯเปฯ กายายตนํ กายวิญฺญาณสฺส.\csromanspacer{} วตฺถุ นเหตุ-อเหตุกานํ ขนฺธานํ ปุเรชาตปจฺจเยน ปจฺจโย. (1)}

\prose{นเหตุ อเหตุโก ธมฺโม นเหตุสเหตุกสฺส ธมฺมสฺส ปุเรชาตปจฺจเยน ปจฺจโย.\csromanspacer{} อารมฺมณปุเรชาตํ, วตฺถุปุเรชาตํ.\csromanspacer{}\csromanspacer{}\textbf{อารมฺมณปุเรชาตํ\csromandash{}}จกฺขุํ ฯเปฯ วตฺถุํ อนิจฺจโต ฯเปฯ โทมนสฺสํ อุปฺปชฺชติ.\csromanspacer{} กุสลากุสเล นิรุทฺเธ นเหตุ สเหตุโก วิปาโก ตทารมฺมณตา อุปฺปชฺชติ.\csromanspacer{} ทิพฺเพน จกฺขุนา รูปํ ปสฺสติ.\csromanspacer{} ทิพฺพาย โสตธาตุยา สทฺทํ สุณาติ.\csromanspacer{} \textbf{วตฺถุปุเรชาตํ}\csromandash{}วตฺถุ นเหตุสเหตุกานํ ขนฺธานํ ปุเรชาตปจฺจเยน ปจฺจโย. (2)}

\csromanheader{2}{6. นเหตุสเหตุทุก \textbf{ปจฺฉาชาต}}
\proseitem{184}{\csromanfolio{79}นเหตุ สเหตุโก ธมฺโม นเหตุ-อเหตุกสฺส ธมฺมสฺส ปจฺฉาชาตปจฺจเยน ปจฺจโย.\csromanspacer{} ปจฺฉาชาตา นเหตู สเหตุกา ขนฺธา ปุเรชาตสฺส อิมสฺส กายสฺส ปจฺฉาชาตปจฺจเยน ปจฺจโย. (1)}

\prose{นเหตุ อเหตุโก ธมฺโม นเหตุ-อเหตุกสฺส ธมฺมสฺส ปจฺฉาชาตปจฺจเยน ปจฺจโย.\csromanspacer{} ปจฺฉาชาตา นเหตู อเหตุกา ขนฺธา ปุเรชาตสฺส อิมสฺส กายสฺส ปจฺฉาชาตปจฺจเยน ปจฺจโย. (1)}

\csromanheader{2}{\textbf{อาเสวน}}
\proseitem{185}{นเหตุ สเหตุโก ธมฺโม นเหตุสเหตุกสฺส ธมฺมสฺส อาเสวนปจฺจเยน ปจฺจโย.\csromanspacer{} ปุริมา ปุริมา นเหตู สเหตุกา ขนฺธา ปจฺฉิมานํ ปจฺฉิมานํ นเหตุสเหตุกานํ ขนฺธานํ อาเสวนปจฺจเยน ปจฺจโย.\csromanspacer{} อนุโลมํ โคตฺรภุสฺส.\csromanspacer{} อนุโลมํ โวทานสฺส.\csromanspacer{} โคตฺรภุ มคฺคสฺส.\csromanspacer{} โวทานํ มคฺคสฺส อาเสวนปจฺจเยน ปจฺจโย. (1)}

\prose{นเหตุ อเหตุโก ธมฺโม นเหตุ-อเหตุกสฺส ธมฺมสฺส อาเสวนปจฺจเยน ปจฺจโย.\csromanspacer{} ปุริมา ปุริมา นเหตู อเหตุกา ขนฺธา ปจฺฉิมานํ ปจฺฉิมานํ นเหตุ-อเหตุกานํ ขนฺธานํ อาเสวนปจฺจเยน ปจฺจโย. (1)}

\csromanheader{2}{\textbf{กมฺม}}
\proseitem{186}{นเหตุ สเหตุโก ธมฺโม นเหตุสเหตุกสฺส ธมฺมสฺส กมฺมปจฺจเยน ปจฺจโย.\csromanspacer{} \textbf{สหชาตา}, นานากฺขณิกา.\csromanspacer{}\csromanspacer{}\textbf{สหชาตา} นเหตุ สเหตุกา เจตนา สมฺปยุตฺตกานํ ขนฺธานํ กมฺมปจฺจเยน ปจฺจโย.\csromanspacer{} ปฏิสนฺธิกฺขเณ ฯเปฯ.\csromanspacer{} \textbf{นานากฺขณิกา} นเหตุ สเหตุกา เจตนา วิปากานํ นเหตุสเหตุกานํ ขนฺธานํ กมฺมปจฺจเยน ปจฺจโย. (1)}

\prose{นเหตุ สเหตุโก ธมฺโม นเหตุ-อเหตุกสฺส ธมฺมสฺส กมฺมปจฺจเยน ปจฺจโย.\csromanspacer{} \textbf{สหชาตา}, นานากฺขณิกา.\csromanspacer{}\csromanspacer{}\textbf{สหชาตา} นเหตุ สเหตุกา เจตนา จิตฺตสมุฏฺฐานานํ รูปานํ กมฺมปจฺจเยน ปจฺจโย.\csromanspacer{} ปฏิสนฺธิกฺขเณ ฯเปฯ.\csromanspacer{} \textbf{นานากฺขณิกา} นเหตุ สเหตุกา เจตนา วิปากานํ นเหตุ-อเหตุกานํ ขนฺธานํ กฏตฺตา จ รูปานํ กมฺมปจฺจเยน ปจฺจโย. (1)}

\prose{\csromanfolio{80}นเหตุ สเหตุโก ธมฺโม นเหตุสเหตุกสฺส จ นเหตุ-อเหตุกสฺส จ ธมฺมสฺส กมฺมปจฺจเยน ปจฺจโย.\csromanspacer{} \textbf{สหชาตา}, นานากฺขณิกา.\csromanspacer{}\csromanspacer{}\textbf{สหชาตา} นเหตุ สเหตุกา เจตนา สมฺปยุตฺตกานํ ขนฺธานํ จิตฺตสมุฏฺฐานานญฺจ รูปานํ กมฺมปจฺจเยน ปจฺจโย.\csromanspacer{} ปฏิสนฺธิกฺขเณ ฯเปฯ.\csromanspacer{} \textbf{นานากฺขณิกา} นเหตุ สเหตุกา เจตนา วิปากานํ นเหตุสเหตุกานํ ขนฺธานํ กฏตฺตา จ รูปานํ กมฺมปจฺจเยน ปจฺจโย. (1)}

\prose{นเหตุ อเหตุโก ธมฺโม นเหตุ-อเหตุกสฺส ธมฺมสฺส กมฺมปจฺจเยน ปจฺจโย.\csromanspacer{} \textbf{สหชาตา} นเหตุ อเหตุกา เจตนา สมฺปยุตฺตกานํ ขนฺธานํ จิตฺตสมุฏฺฐานานญฺจ รูปานํ กมฺมปจฺจเยน ปจฺจโย.\csromanspacer{} ปฏิสนฺธิกฺขเณ นเหตุ อเหตุกา เจตนา สมฺปยุตฺตกานํ ขนฺธานํ กฏตฺตา จ รูปานํ กมฺมปจฺจเยน ปจฺจโย. (1)}

\csromanheader{2}{\textbf{วิปาก}}
\proseitem{187}{นเหตุ สเหตุโก ธมฺโม นเหตุสเหตุกสฺส ธมฺมสฺส วิปากปจฺจเยน ปจฺจโย.\csromanspacer{} ตีณิ.}

\prose{นเหตุ อเหตุโก ธมฺโม นเหตุ-อเหตุกสฺส ธมฺมสฺส วิปากปจฺจเยน ปจฺจโย.\csromanspacer{} เอกํ.}

\csromanheader{2}{\textbf{อาหาร}}
\proseitem{188}{นเหตุ สเหตุโก ธมฺโม นเหตุสเหตุกสฺส ธมฺมสฺส อาหารปจฺจเยน ปจฺจโย.\csromanspacer{} ตีณิ.}

\prose{นเหตุ อเหตุโก ธมฺโม นเหตุ-อเหตุกสฺส ธมฺมสฺส อาหารปจฺจเยน ปจฺจโย.\csromanspacer{} นเหตู อเหตุกา อาหารา สมฺปยุตฺตกานํ ขนฺธานํ จิตฺตสมุฏฺฐานานญฺจ รูปานํ อาหารปจฺจเยน ปจฺจโย.\csromanspacer{} ปฏิสนฺธิกฺขเณ ฯเปฯ.\csromanspacer{} กพฬีกาโร อาหาโร อิมสฺส กายสฺส อาหารปจฺจเยน ปจฺจโย. (1)}

\csromanheader{2}{\textbf{อินฺทฺริย}}
\proseitem{189}{นเหตุ สเหตุโก ธมฺโม นเหตุสเหตุกสฺส ธมฺมสฺส อินฺทฺริยปจฺจเยน ปจฺจโย.\csromanspacer{} ตีณิ.}

\csromanheader{2}{6. นเหตุสเหตุทุก}
\prose{\csromanfolio{81}นเหตุ อเหตุโก ธมฺโม นเหตุ-อเหตุกสฺส ธมฺมสฺส อินฺทฺริยปจฺจเยน ปจฺจโย.\csromanspacer{} นเหตู อเหตุกา อินฺทฺริยา สมฺปยุตฺตกานํ ขนฺธานํ จิตฺตสมุฏฺฐานานญฺจ รูปานํ อินฺทฺริยปจฺจเยน ปจฺจโย.\csromanspacer{} ปฏิสนฺธิกฺขเณ ฯเปฯ.\csromanspacer{} รูปชีวิตินฺทฺริยํ กฏตฺตารูปานํ อินฺทฺริยปจฺจเยน ปจฺจโย.}

\csromanheader{2}{\textbf{ฌานาทิ}}
\proseitem{190}{นเหตุ สเหตุโก ธมฺโม นเหตุสเหตุกสฺส ธมฺมสฺส ฌานปจฺจเยน ปจฺจโย ฯเปฯ. (จตฺตาริปิ กาตพฺพานิ.) มคฺคปจฺจเยน ปจฺจโย.\csromanspacer{} ตีณิ.}

\csromanheader{2}{\textbf{สมฺปยุตฺต}}
\proseitem{191}{นเหตุ สเหตุโก ธมฺโม นเหตุสเหตุกสฺส ธมฺมสฺส สมฺปยุตฺตปจฺจเยน ปจฺจโย.\csromanspacer{} นเหตุ สเหตุโก เอโก ขนฺโธ ติณฺณนฺนํ ขนฺธานํ ฯเปฯ.\csromanspacer{} ปฏิสนฺธิกฺขเณ ฯเปฯ. (1)}

\prose{นเหตุ อเหตุโก ธมฺโม นเหตุ-อเหตุกสฺส ธมฺมสฺส สมฺปยุตฺตปจฺจเยน ปจฺจโย.\csromanspacer{} นเหตุ อเหตุโก เอโก ขนฺโธ ติณฺณนฺนํ ขนฺธานํ ฯเปฯ.\csromanspacer{} ปฏิสนฺธิกฺขเณ ฯเปฯ. (2)}

\csromanheader{2}{\textbf{วิปฺปยุตฺต}}
\proseitem{192}{นเหตุ สเหตุโก ธมฺโม นเหตุ-อเหตุกสฺส ธมฺมสฺส วิปฺปยุตฺตปจฺจเยน ปจฺจโย.\csromanspacer{} สหชาตํ, ปจฺฉาชาตํ.\csromanspacer{}\csromanspacer{}\textbf{สหชาตา} นเหตู สเหตุกา ขนฺธา จิตฺตสมุฏฺฐานานํ รูปานํ วิปฺปยุตฺตปจฺจเยน ปจฺจโย.\csromanspacer{} ปฏิสนฺธิกฺขเณ ฯเปฯ.\csromanspacer{} \textbf{ปจฺฉาชาตา} นเหตู สเหตุกา ขนฺธา ปุเรชาตสฺส อิมสฺส กายสฺส วิปฺปยุตฺตปจฺจเยน ปจฺจโย. (1)}

\prose{นเหตุ อเหตุโก ธมฺโม นเหตุ-อเหตุกสฺส ธมฺมสฺส วิปฺปยุตฺตปจฺจเยน ปจฺจโย.\csromanspacer{} สหชาตํ, ปุเรชาตํ, ปจฺฉาชาตํ.\csromanspacer{}\csromanspacer{}\textbf{สหชาตา} นเหตู อเหตุกา ขนฺธา จิตฺตสมุฏฺฐานานํ รูปานํ วิปฺปยุตฺตปจฺจเยน ปจฺจโย.\csromanspacer{} ปฏิสนฺธิกฺขเณ ฯเปฯ.\csromanspacer{} ขนฺธา วตฺถุสฺส วิปฺปยุตฺตปจฺจเยน ปจฺจโย.\csromanspacer{} วตฺถุ ขนฺธานํ วิปฺปยุตฺตปจฺจเยน ปจฺจโย.\csromanspacer{} \textbf{ปุเรชาตํ} จกฺขายตนํ จกฺขุวิญฺญาณสฺส ฯเปฯ กายายตนํ กายวิญฺญาณสฺส ฯเปฯ วตฺถุ นเหตุสเหตุกานํ ขนฺธานํ วิปฺปยุตฺตปจฺจเยน ปจฺจโย.\csromanspacer{} \textbf{ปจฺฉาชาตา} นเหตู \csromanfolio{82}อเหตุกา ขนฺธา ปุเรชาตสฺส อิมสฺส กายสฺส วิปฺปยุตฺตปจฺจเยน ปจฺจโย. (1)}

\prose{นเหตุ อเหตุโก ธมฺโม นเหตุสเหตุกสฺส ธมฺมสฺส วิปฺปยุตฺตปจฺจเยน ปจฺจโย.\csromanspacer{} \textbf{สหชาตํ}, ปุเรชาตํ.\csromanspacer{}\csromanspacer{}\textbf{สหชาตํ} ปฏิสนฺธิกฺขเณ วตฺถุ นเหตุสเหตุกานํ ขนฺธานํ วิปฺปยุตฺตปจฺจเยน ปจฺจโย.\csromanspacer{} \textbf{ปุเรชาตํ} วตฺถุ นเหตุสเหตุกานํ ขนฺธานํ วิปฺปยุตฺตปจฺจเยน ปจฺจโย. (2)}

\csromanheader{2}{\textbf{อตฺถิ}}
\proseitem{193}{นเหตุ สเหตุโก ธมฺโม นเหตุสเหตุกสฺส ธมฺมสฺส อตฺถิปจฺจเยน ปจฺจโย.\csromanspacer{} นเหตุ สเหตุโก เอโก ขนฺโธ ติณฺณนฺนํ ขนฺธานํ ฯเปฯ.\csromanspacer{} ปฏิสนฺธิกฺขเณ ฯเปฯ. (1)}

\prose{นเหตุ สเหตุโก ธมฺโม นเหตุ-อเหตุกสฺส ธมฺมสฺส อตฺถิปจฺจเยน ปจฺจโย.\csromanspacer{} \textbf{สหชาตํ, ปจฺฉาชาตํ} ฯเปฯ. (2)}

\prose{นเหตุ สเหตุโก ธมฺโม นเหตุสเหตุกสฺส จ นเหตุ-อเหตุกสฺส จ ธมฺมสฺส อตฺถิปจฺจเยน ปจฺจโย.\csromanspacer{} นเหตุ สเหตุโก เอโก ขนฺโธ ติณฺณนฺนํ ขนฺธานํ จิตฺตสมุฏฺฐานานญฺจ รูปานํ อตฺถิปจฺจเยน ปจฺจโย ฯเปฯ.\csromanspacer{} ปฏิสนฺธิกฺขเณ ฯเปฯ. (3)}

\prose{นเหตุ อเหตุโก ธมฺโม นเหตุ-อเหตุกสฺส ธมฺมสฺส อตฺถิปจฺจเยน ปจฺจโย.\csromanspacer{} \textbf{สหชาตํ}, ปุเรชาตํ, ปจฺฉาชาตํ, อาหารํ, อินฺทฺริยํ.\csromanspacer{}\csromanspacer{}\textbf{สหชาโต} นเหตุ อเหตุโก เอโก ขนฺโธ ติณฺณนฺนํ ขนฺธานํ จิตฺตสมุฏฺฐานานญฺจ รูปานํ อตฺถิปจฺจเยน ปจฺจโย ฯเปฯ. (ยาว อสญฺญสตฺตา.) \textbf{ปุเรชาตํ} จกฺขุํ ฯเปฯ วตฺถุํ อนิจฺจโต ฯเปฯ โทมนสฺสํ อุปฺปชฺชติ.\csromanspacer{} กุสลากุสเล นิรุทฺเธ นเหตุ อเหตุโก วิปาโก ตทารมฺมณตา อุปฺปชฺชติ.\csromanspacer{} รูปายตนํ จกฺขุวิญฺญาณสฺส ฯเปฯ.\csromanspacer{} โผฏฺฐพฺพายตนํ กายวิญฺญาณสฺส ฯเปฯ.\csromanspacer{} จกฺขายตนํ ฯเปฯ.\csromanspacer{} กายายตนํ ฯเปฯ วตฺถุ นเหตุ-อเหตุกานํ ขนฺธานํ อตฺถิปจฺจเยน ปจฺจโย.\csromanspacer{} \textbf{ปจฺฉาชาตา} นเหตู อเหตุกา ขนฺธา ปุเรชาตสฺส ฯเปฯ.\csromanspacer{} \textbf{กพฬีกาโร อาหาโร} อิมสฺส กายสฺส ฯเปฯ.\csromanspacer{} \textbf{รูปชีวิตินฺทฺริยํ} กฏตฺตารูปานํ อตฺถิปจฺจเยน ปจฺจโย. (1)}

\prose{นเหตุ อเหตุโก ธมฺโม นเหตุสเหตุกสฺส ธมฺมสฺส อตฺถิปจฺจเยน ปจฺจโย.\csromanspacer{} \textbf{สหชาตํ}, ปุเรชาตํ.\csromanspacer{}\csromanspacer{}\textbf{สหชาตํ} ปฏิสนฺธิกฺขเณ}

\vspace{\tipitakaparskip}
\csromancenter{นเหตุสเหตุทุก วตฺถุ นเหตุสเหตุกานํ ขนฺธานํ อตฺถิปจฺจเยน ปจฺจโย.\csromanspacer{} \textbf{ปุเรชาตํ} จกฺขุํ ฯเปฯ วตฺถุํ อนิจฺจโต ฯเปฯ โทมนสฺสํ อุปฺปชฺชติ.\csromanspacer{} กุสลากุสเล นิรุทฺเธ นเหตุ สเหตุโก วิปาโก ตทารมฺมณตา อุปฺปชฺชติ.\csromanspacer{} ทิพฺเพน จกฺขุนา รูปํ ปสฺสติ.\csromanspacer{} ทิพฺพาย โสตธาตุยา สทฺทํ สุณาติ.\csromanspacer{} \textbf{วตฺถุปุเรชาตํ} วตฺถุ นเหตุสเหตุกานํ ขนฺธานํ อตฺถิปจฺจเยน ปจฺจโย. (2)}
\proseitem{194}{\csromanfolio{83}นเหตุ สเหตุโก จ นเหตุ อเหตุโก จ ธมฺมา นเหตุสเหตุกสฺส ธมฺมสฺส อตฺถิปจฺจเยน ปจฺจโย.\csromanspacer{} \textbf{สหชาตํ, ปุเรชาตํ}.\csromanspacer{}\csromanspacer{}สหชาโต นเหตุ สเหตุโก เอโก ขนฺโธ จ วตฺถุ จ ติณฺณนฺนํ ขนฺธานํ ฯเปฯ.\csromanspacer{} ปฏิสนฺธิกฺขเณ ฯเปฯ.\csromanspacer{} นเหตุ สเหตุโก เอโก ขนฺโธ จ วตฺถุ จ ติณฺณนฺนํ ขนฺธานํ ฯเปฯ. (1)}

\prose{นเหตุ สเหตุโก จ นเหตุ อเหตุโก จ ธมฺมา นเหตุ-อเหตุกสฺส ธมฺมสฺส อตฺถิปจฺจเยน ปจฺจโย.\csromanspacer{} สหชาตํ, ปจฺฉาชาตํ, อาหารํ, อินฺทฺริยํ.\csromanspacer{}\csromanspacer{}\textbf{สหชาตา} นเหตู สเหตุกา ขนฺธา จ มหาภูตา จ จิตฺตสมุฏฺฐานานํ รูปานํ อตฺถิปจฺจเยน ปจฺจโย.\csromanspacer{} ปฏิสนฺธิกฺขเณ ฯเปฯ.\csromanspacer{} \textbf{ปจฺฉาชาตา} นเหตู สเหตุกา ขนฺธา จ กพฬีกาโร อาหาโร จ อิมสฺส กายสฺส อตฺถิปจฺจเยน ปจฺจโย.\csromanspacer{} \textbf{ปจฺฉาชาตา} นเหตู สเหตุกา ขนฺธา จ รูปชีวิตินฺทฺริยญฺจ กฏตฺตารูปานํ อตฺถิปจฺจเยน ปจฺจโย ฯเปฯ. (2)}

\csromanheader{2}{1. ปจฺจยานุโลม 2. สงฺขฺยาวาร}
\csromanheader{2}{\textbf{สุทฺธ}}
\proseitem{195}{อารมฺมเณ จตฺตาริ, อธิปติยา จตฺตาริ, อนนฺตเร จตฺตาริ, สมนนฺตเร จตฺตาริ, สหชาเต สตฺต, อญฺญมญฺเญ ฉ, นิสฺสเย สตฺต, อุปนิสฺสเย จตฺตาริ, ปุเรชาเต เทฺว, ปจฺฉาชาเต เทฺว, อาเสวเน เทฺว, กมฺเม จตฺตาริ, วิปาเก จตฺตาริ, อาหาเร จตฺตาริ, อินฺทฺริเย จตฺตาริ, ฌาเน จตฺตาริ, มคฺเค ตีณิ, สมฺปยุตฺเต เทฺว, วิปฺปยุตฺเต ตีณิ, อตฺถิยา สตฺต, นตฺถิยา จตฺตาริ, วิคเต จตฺตาริ, อวิคเต สตฺต. (เอวํ คเณตพฺพํ)}

\vspace{\tipitakaparskip}
\csromancenter{อนุโลมํ.}
\csromanheader{2}{2. ปจฺจนียุทฺธาร}
\proseitem{196}{\csromanfolio{84}นเหตุ สเหตุโก ธมฺโม นเหตุสเหตุกสฺส ธมฺมสฺส อารมฺมณปจฺจเยน ปจฺจโย.\csromanspacer{} สหชาตปจฺจเยน ปจฺจโย.\csromanspacer{} อุปนิสฺสยปจฺจเยน ปจฺจโย.\csromanspacer{} กมฺมปจฺจเยน ปจฺจโย. (1)}

\prose{นเหตุ สเหตุโก ธมฺโม นเหตุ-อเหตุกสฺส ธมฺมสฺส อารมฺมณปจฺจเยน ปจฺจโย.\csromanspacer{} สหชาตปจฺจเยน ปจฺจโย.\csromanspacer{} อุปนิสฺสยปจฺจเยน ปจฺจโย.\csromanspacer{} ปจฺฉาชาตปจฺจเยน ปจฺจโย.\csromanspacer{} กมฺมปจฺจเยน ปจฺจโย. (2)}

\prose{นเหตุ สเหตุโก ธมฺโม นเหตุสเหตุกสฺส จ นเหตุ-อเหตุกสฺส จ ธมฺมสฺส สหชาตปจฺจเยน ปจฺจโย.\csromanspacer{} กมฺมปจฺจเยน ปจฺจโย. (3)}

\proseitem{197}{นเหตุ อเหตุโก ธมฺโม นเหตุ-อเหตุกสฺส ธมฺมสฺส อารมฺมณปจฺจเยน ปจฺจโย.\csromanspacer{} สหชาตปจฺจเยน ปจฺจโย.\csromanspacer{} อุปนิสฺสยปจฺจเยน ปจฺจโย.\csromanspacer{} ปุเรชาตปจฺจเยน ปจฺจโย.\csromanspacer{} ปจฺฉาชาตปจฺจเยน ปจฺจโย.\csromanspacer{} อาหารปจฺจเยน ปจฺจโย.\csromanspacer{} อินฺทฺริยปจฺจเยน ปจฺจโย. (1)}

\prose{นเหตุ อเหตุโก ธมฺโม นเหตุสเหตุกสฺส ธมฺมสฺส อารมฺมณปจฺจเยน ปจฺจโย.\csromanspacer{} สหชาตปจฺจเยน ปจฺจโย.\csromanspacer{} อุปนิสฺสยปจฺจเยน ปจฺจโย.\csromanspacer{} ปุเรชาตปจฺจเยน ปจฺจโย. (2)}

\proseitem{198}{นเหตุ สเหตุโก จ นเหตุ อเหตุโก จ ธมฺมา นเหตุสเหตุกสฺส ธมฺมสฺส สหชาตํ.\csromanspacer{} ปุเรชาตํ. (1)}

\prose{นเหตุ สเหตุโก จ นเหตุ อเหตุโก จ ธมฺมา นเหตุ-อเหตุกสฺส ธมฺมสฺส สหชาตํ.\csromanspacer{} ปจฺฉาชาตํ.\csromanspacer{} อาหารํ.\csromanspacer{} อินฺทฺริยํ. (2)}

\csromanheader{2}{2. ปจฺจยปจฺจนีย 2. สงฺขฺยาวาร}
\csromanheader{2}{\textbf{สุทฺธ}}
\proseitem{199}{นเหตุยา สตฺต, น-อารมฺมเณ สตฺต. (สํขิตฺตํ, สพฺพตฺถ สตฺต, นสหชาเต ฉ, น-อญฺญมญฺเญ ฉ, นนิสฺสเย ฉ, (สพฺพตฺถ สตฺต.)}

\vspace{\tipitakaparskip}
\csromancenter{นเหตุสเหตุทุก นสมฺปยุตฺเต ฉ, นวิปฺปยุตฺเต ปญฺจ, โน-อตฺถิยา ปญฺจ, โนนตฺถิยา สตฺต, โนวิคเต สตฺต, โน-อวิคเต ปญฺจ. (เอวํ คเณตพฺพํ.)}
\csromancenter{ปจฺจนียํ.}
\csromanheader{2}{3. ปจฺจยานุโลมปจฺจนีย}
\csromanheader{2}{\textbf{อารมฺมณทุก}}
\proseitem{200}{\csromanfolio{85}อารมฺมณปจฺจยา นเหตุยา จตฺตาริ, น-อธิปติยา จตฺตาริ, น-อนนฺตเร จตฺตาริ, (สพฺพตฺถ จตฺตาริ.) โนนตฺถิยา จตฺตาริ, โนวิคเต จตฺตาริ, โน-อวิคเต จตฺตาริ. (เอวํ คเณตพฺพํ.)}

\vspace{\tipitakaparskip}
\csromancenter{อนุโลมปจฺจนียํ.}
\csromanheader{2}{4. ปจฺจยปจฺจนียานุโลม}
\csromanheader{2}{\textbf{นเหตุทุก}}
\vspace{\tipitakaparskip}
\csromancenter{นเหตุปจฺจยา อารมฺมเณ จตฺตาริ, อธิปติยา จตฺตาริ ฯเปฯ อวิคเต สตฺต.}
\csromancenter{ปจฺจนียานุโลมํ.}
\nitthitam{นเหตุสเหตุกทุกํ นิฏฺฐิตํ.}
\nitthitam{เหตุโคจฺฉกํ นิฏฺฐิตํ.}
\csromanmatikaanchor{mtk.9}
\csromantocmarknopage{chapter}{2. จูฬนฺตรทุก}{mtk.9}
\bookmark[dest={mtk.9},level=0]{2. จูฬนฺตรทุก}
\chapterheadpageread{2. จูฬนฺตรทุก}
\csromanmatikaanchor{mtk.10}
\csromantocmarkref{section}{7. สปฺปจฺจยทุก}{mtk.10}
\bookmark[dest={mtk.10},level=1]{7. สปฺปจฺจยทุก}
\csromanheader{1}{7. สปฺปจฺจยทุก 1. ปฏิจฺจวาร}
\csromanheader{2}{1-4. ปจฺจยานุโลมาทิ}
\proseitem{1}{\csromanfolio{86}สปฺปจฺจยํ ธมฺมํ ปฏิจฺจ สปฺปจฺจโย ธมฺโม อุปฺปชฺชติ เหตุปจฺจยา.\csromanspacer{} สปฺปจฺจยํ เอกํ ขนฺธํ ปฏิจฺจ ตโย ขนฺธา จิตฺตสมุฏฺฐานญฺจ รูปํ ฯเปฯ เทฺว ขนฺเธ ฯเปฯ.\csromanspacer{} ปฏิสนฺธิกฺขเณ ฯเปฯ.\csromanspacer{} ขนฺเธ ปฏิจฺจ วตฺถุ, วตฺถุํ ปฏิจฺจ ขนฺธา.\csromanspacer{} เอกํ มหาภูตํ ฯเปฯ มหาภูเต ปฏิจฺจ จิตฺตสมุฏฺฐานํ รูปํ, กฏตฺตารูปํ, อุปาทารูปํ.}

\prose{สปฺปจฺจยํ ธมฺมํ ปฏิจฺจ สปฺปจฺจโย ธมฺโม อุปฺปชฺชติ อารมฺมณปจฺจยา ฯเปฯ อวิคตปจฺจยา.}

\proseitem{2}{เหตุยา เอกํ, อารมฺมเณ เอกํ ฯเปฯ อวิคเต เอกํ.}

\vspace{\tipitakaparskip}
\csromancenter{อนุโลมํ.}
\proseitem{3}{สปฺปจฺจยํ ธมฺมํ ปฏิจฺจ สปฺปจฺจโย ธมฺโม อุปฺปชฺชติ นเหตุปจฺจยา.\csromanspacer{} อเหตุกํ สปฺปจฺจยํ เอกํ ขนฺธํ ปฏิจฺจ ตโย ขนฺธา จิตฺตสมุฏฺฐานญฺจ รูปํ.\csromanspacer{} เทฺว ขนฺเธ ฯเปฯ.\csromanspacer{} อเหตุกปฏิสนฺธิกฺขเณ ฯเปฯ. (ยาว อสญฺญสตฺตา.) วิจิกิจฺฉาสเหคเต อุทฺธจฺจสหคเต ขนฺเธ ปฏิจฺจ วิจิกิจฺฉาสหคโต อุจฺฉจฺจสหคโต โมโห. (สํขิตฺตํ.)}

\proseitem{4}{นเหตุยา เอกํ, น-อารมฺมเณ เอกํ, น-อธิปติยา เอกํ ฯเปฯ โนวิคเต เอกํ.\csromanspacer{} ปจฺจนียํ.}

\proseitem{5}{เหตุปจฺจยา น-อารมฺมเณ เอกํ, น-อธิปติยา เอกํ ฯเปฯ โนวิคเต เอกํ.\csromanspacer{} อนุโลมปจฺจนียํ.}

\proseitem{6}{นเหตุปจฺจยา อารมฺมเณ เอกํ, อนนฺตเร เอกํ ฯเปฯ อวิคเต เอกํ.}

\vspace{\tipitakaparskip}
\csromancenter{ปจฺจนียานุโลมํ.}
\csromancenter{(สหชาตวาโร ปฏิจฺจวารสทิโส.)}
\csromanheader{2}{7. สปฺปจฺจยทุก \textbf{7. สปฺปจฺจยทุก 3. ปจฺจยวาร}}
\csromanheader{2}{1. ปจฺจยานุโลม 1.วิภงฺควาร}
\proseitem{7}{\csromanfolio{87}สปฺปจฺจยํ ธมฺมํ ปจฺจยา สปฺปจฺจโย ธมฺโม อุปฺปชฺชติ เหตุปจฺจยา.\csromanspacer{} สปฺปจฺจยํ เอกํ ขนฺธํ ปจฺจยา ตโย ขนฺธา จิตฺตสมุฏฺฐานญฺจ รูปํ ฯเปฯ.\csromanspacer{} ปฏิสนฺธิกฺขเณ ฯเปฯ.\csromanspacer{} ขนฺเธ ปจฺจยา วตฺถุ, วตฺถุํ ปจฺจยา ขนฺธา.\csromanspacer{} เอกํ มหาภูตํ ฯเปฯ.\csromanspacer{} วตฺถุํ ปจฺจยา สปฺปจฺจยา ขนฺธา.}

\prose{สปฺปจฺจยํ ธมฺมํ ปจฺจยา สปฺปจฺจโย ธมฺโม อุปฺปชฺชติ อารมฺมณปจฺจยา. (สํขิตฺตํ.) (เอวํ ปจฺจยวาโรปิ นิสฺสยวาโรปิ สํสฏฺฐวาโรปิ สมฺปยุตฺตวาโรปิ วิตฺถาเรตพฺโพ, สพฺพตฺถ เอกาเยว ปญฺหา.)}

\csromanheader{2}{7. สปฺปจฺจยทุก 7. ปญฺหาวาร}
\csromanheader{2}{1. ปจฺจยานุโลม 1. วิภงฺควาร}
\csromanheader{2}{\textbf{เหตุ}}
\proseitem{8}{สปฺปจฺจโย ธมฺโม สปฺปจฺจยสฺส ธมฺมสฺส เหตุปจฺจเยน ปจฺจโย.\csromanspacer{} สปฺปจฺจยา เหตู สมฺปยุตฺตกานํ ขนฺธานํ จิตฺตสมุฏฺฐานานญฺจ รูปานํ เหตุปจฺจเยน ปจฺจโย.\csromanspacer{} ปฏิสนฺธิกฺขเณ ฯเปฯ.}

\csromanheader{2}{\textbf{อารมฺมณ}}
\proseitem{9}{สปฺปจฺจโย ธมฺโม สปฺปจฺจยสฺส ธมฺมสฺส อารมฺมณปจฺจเยน ปจฺจโย.\csromanspacer{} ทานํ ทตฺวา, สีลํ สมาทิยิตฺวา, อุโปสถกมฺมํ กตฺวา ตํ ปจฺจเวกฺขติ.\csromanspacer{} ปุพฺเพ สุจิณฺณานิ ฯเปฯ.\csromanspacer{} ฌานา วุฏฺฐหิตฺวา ฯเปฯ.\csromanspacer{} อริยา มคฺคา วุฏฺฐหิตฺวา มคฺคํ ปจฺจเวกฺขนฺติ, ผลํ ปจฺจเวกฺขนฺติ.\csromanspacer{} ปหีเน กิเลเส ฯเปฯ.\csromanspacer{} วิกฺขมฺภิเต กิเลเส ฯเปฯ.\csromanspacer{} ปุพฺเพ สมุทาจิณฺเณ กิเลเส ชานนฺติ.\csromanspacer{} จกฺขุํ ฯเปฯ วตฺถุํ.\csromanspacer{} สปฺปจฺจเย ขนฺเธ อนิจฺจโต ฯเปฯ โทมนสฺสํ อุปฺปชฺชติ.\csromanspacer{} ทิพฺเพน จกฺขุนา รูปํ ปสฺสติ.\csromanspacer{} ทิพฺพาย โสตธาตุยา สทฺทํ สุณาติ.\csromanspacer{} เจโตปริญาเณน \csromanfolio{88}สปฺปจฺจยจิตฺตสมงฺคิสฺส จิตฺตํ ชานาติ.\csromanspacer{} อากาสานญฺจายตนํ วิญฺญาณญฺจายตนสฺส ฯเปฯ.\csromanspacer{} อากิญฺจญฺญายตนํ เนวสญฺญานาสญฺญายตนสฺส ฯเปฯ.\csromanspacer{} รูปายตนํ จกฺขุวิญฺญาณสฺส ฯเปฯ.\csromanspacer{} โผฏฺฐพฺพายตนํ กายวิญฺญาณสฺส ฯเปฯ.\csromanspacer{} สปฺปจฺจยา ขนฺธา อิทฺธิวิธญาณสฺส, เจโตปริยญาณสฺส, ปุพฺเพนิวาสานุสฺสติญาณสฺส, ยถากมฺมูปคญาณสฺส, อนาคตํสญาณสฺส, อาวชฺชนาย อารมฺมณปจฺจเยน ปจฺจโย. (1)}

\prose{อปฺปจฺจโย ธมฺโม สปฺปจฺจยสฺส ธมฺมสฺส อารมฺมณปจฺจเยน ปจฺจโย.\csromanspacer{} อริยา นิพฺพานํ ปจฺจเวกฺขนฺติ.\csromanspacer{} นิพฺพานํ โคตฺรภุสฺส, โวทานสฺส, มคฺคสฺส, ผลสฺส, อาวชฺชนาย อารมฺมณปจฺจเยน ปจฺจโย. (1)}

\csromanheader{2}{\textbf{อธิปติ}}
\proseitem{10}{สปฺปจฺจโย ธมฺโม สปฺปจฺจยสฺส ธมฺมสฺส อธิปติปจฺจเยน ปจฺจโย.\csromanspacer{} \textbf{อารมฺมณาธิปติ}, สหชาตาธิปติ.\csromanspacer{}\csromanspacer{}\textbf{อารมฺมณาธิปติ}\csromandash{}ทานํ ทตฺวา, สีลํ ฯเปฯ อุโปสถกมฺมํ กตฺวา ตํ ครุํ กตฺวา ปจฺจเวกฺขติ.\csromanspacer{} ปุพฺเพ ฯเปฯ.\csromanspacer{} ฌานา ฯเปฯ.\csromanspacer{} อริยา มคฺคา วุฏฺฐหิตฺวา มคฺคํ ครุํ กตฺวา ปจฺจเวกฺขนฺติ.\csromanspacer{} ผลํ ครุํ กตฺวา ฯเปฯ.\csromanspacer{} จกฺขุํ ฯเปฯ วตฺถุํ.\csromanspacer{} สปฺปจฺจเย ขนฺเธ ครุํ กตฺวา อสฺสาเทติ อภินนฺทติ, ตํ ครุํ กตฺวา ราโค อุปฺปชฺชติ.\csromanspacer{} ทิฏฺฐิ อุปฺปชฺชติ.\csromanspacer{} \textbf{สหชาตาธิปติ}\csromandash{}สปฺปจฺจยาธิปติ สมฺปยุตฺตกานํ ขนฺธานํ จิตฺตสมุฏฺฐานานญฺจ รูปานํ อธิปติปจฺจเยน ปจฺจโย. (1)}

\prose{อปฺปจฺจโย ธมฺโม สปฺปจฺจยสฺส ธมฺมสฺส อธิปติปจฺจเยน ปจฺจโย.\csromanspacer{} \textbf{อารมฺมณาธิปติ}\csromandash{}อริยา นิพฺพานํ ครุํ กตฺวา ปจฺจเวกฺขนฺติ.\csromanspacer{} นิพฺพานํ โคตฺรภุสฺส, โวทานสฺส, มคฺคสฺส, ผลสฺส อธิปติปจฺจเยน ปจฺจโย. (1)}

\csromanheader{2}{\textbf{อนนฺตราทิ}}
\proseitem{11}{สปฺปจฺจโย ธมฺโม สปฺปจฺจยสฺส ธมฺมสฺส อนนฺตรปจฺจเยน ปจฺจโย ฯเปฯ อุปนิสฺสยปจฺจเยน ปจฺจโย ฯเปฯ. (เทฺว ปญฺหา อุปนิสฺสยมูลํ.) ปุเรชาตปจฺจเยน ปจฺจโย ฯเปฯ อวิคตปจฺจเยน ปจฺจโย. (สพฺพตฺถ เอกาเยว ปญฺหา.)}

\csromanheader{2}{7. สปฺปจฺจยทุก \textbf{1. ปจฺจยานุโลม 2. สงฺขฺยาวาร}}
\csromanheader{2}{\textbf{สุทฺธ}}
\proseitem{12}{\csromanfolio{89}เหตุยา เอกํ, อารมฺมเณ เทฺว, อธิปติยา เทฺว, อนนฺตเร เอกํ, สมนนฺตเร เอกํ, สหชาเต เอกํ, อญฺญมญฺเญ เอกํ, นิสฺสเย เอกํ, อุปนิสฺสเย เทฺว, ปุเรชาเต เอกํ, (สพฺพตฺถ เอกํ.) อวิคเต เอกํ. (เอวํ คเณตพฺพํ.)}

\vspace{\tipitakaparskip}
\csromancenter{อนุโลมํ.}
\csromanheader{2}{2. ปจฺจนียุทฺธาร}
\proseitem{13}{สปฺปจฺจโย ธมฺโม สปฺปจฺจยสฺส ธมฺมสฺส อารมฺมณปจฺจเยน ปจฺจโย.\csromanspacer{} สหชาตปจฺจเยน ปจฺจโย.\csromanspacer{} อุปนิสฺสยปจฺจเยน ปจฺจโย.\csromanspacer{} ปุเรชาตปจฺจเยน ปจฺจโย.\csromanspacer{} ปจฺฉาชาตปจฺจเยน ปจฺจโย.\csromanspacer{} กมฺมปจฺจเยน ปจฺจโย.\csromanspacer{} อาหารปจฺจเยน ปจฺจโย.\csromanspacer{} อินฺทฺริยปจฺจเยน ปจฺจโย. (1)}

\prose{อปฺปจฺจโย ธมฺโม สปฺปจฺจยสฺส ธมฺมสฺส อารมฺมณปจฺจเยน ปจฺจโย.\csromanspacer{} อุปนิสฺสยปจฺจเยน ปจฺจโย. (1)}

\csromanheader{2}{2. ปจฺจยปจฺจนีย 2. สงฺขฺยาวาร}
\proseitem{14}{นเหตุยา เทฺว, น-อารมฺมเณ เอกํ, น-อธิปติยา เทฺว, น-อนนฺตเร เทฺว, นสมนนฺตเร เทฺว ฯเปฯ น-อุปนิสฺสเย เทฺว, นปุเรชาเต เทฺว ฯเปฯ โนวิคเต เทฺว, โน-อวิคเต เทฺว. (เอวํ คเณตพฺพํ.)}

\vspace{\tipitakaparskip}
\csromancenter{ปจฺจนียํ.}
\csromanheader{2}{3. ปจฺจยานุโลมปจฺจนีย}
\csromanheader{2}{\textbf{เหตุทุก}}
\proseitem{15}{เหตุปจฺจยา น-อารมฺมเณ เอกํ, น-อธิปติยา เอกํ, น-อนนฺตเร เอกํ, นสมนนฺตเร เอกํ, น-อญฺญมญฺเญ เอกํ, น-อุปนิสฺสเย เอกํ ฯเปฯ นสมฺปยุตฺเต เอกํ, นวิปฺปยุตฺเต เอกํ, โนนตฺถิยา เอกํ, โนวิคเต เอกํ. (เอวํ คเณตพฺพํ.)}

\vspace{\tipitakaparskip}
\csromancenter{อนุโลมปจฺจนียํ.}
\csromanheader{2}{4. ปจฺจยปจฺจนียานุโลม}
\csromanheader{2}{\textbf{นเหตุทุก}}
\proseitem{16}{\csromanfolio{90}นเหตุปจฺจยา อารมฺมเณ เทฺว, อธิปติยา เทฺว, อนนฺตเร เอกํ ฯเปฯ อุปนิสฺสเย เทฺว, ปุเรชาเต เอกํ ฯเปฯ อวิคเต เอกํ. (เอวํ คเณตพฺพํ.)}

\vspace{\tipitakaparskip}
\csromancenter{ปจฺจนียานุโลมํ.}
\nitthitam{สปฺปจฺจยทุกํ นิฏฺฐิตํ.}
\csromanmatikaanchor{mtk.11}
\csromantocmarkref{section}{8. สงฺขตทุก}{mtk.11}
\bookmark[dest={mtk.11},level=1]{8. สงฺขตทุก}
\csromanheader{1}{8. สงฺขตทุก 1. ปฏิจฺจวาร}
\csromanheader{2}{\textbf{เหตุ}}
\proseitem{17}{สงฺขตํ ธมฺมํ ปฏิจฺจ สงฺขโต ธมฺโม อุปฺปชฺชติ เหตุปจฺจยา.\csromanspacer{} สงฺขตํ เอกํ ขนฺธํ ปฏิจฺจ ตโย ขนฺธา จิตฺตสมุฏฺฐานญฺจ รูปํ ฯเปฯ.\csromanspacer{} ปฏิสนฺธิกฺขเณ ฯเปฯ.\csromanspacer{} ขนฺเธ ปฏิจฺจ วตฺถุ, วตฺถุํ ปฏิจฺจ ขนฺธา.\csromanspacer{} เอกํ มหาภูตํ ฯเปฯ มหาภูเต ปฏิจฺจ จิตฺตสมุฏฺฐานํ รูปํ, กฏตฺตารูปํ, อุปาทารูปํ.}

\prose{(อิมํ ทุกํ ยถา สปฺปจฺจยทุกํ, เอวํ คเณตพฺพํ, นินฺนานากรณํ.)}

\nitthitam{สงฺขตทุกํ นิฏฺฐิตํ.}
\csromanmatikaanchor{mtk.12}
\csromantocmarkref{section}{9. สนิทสฺสนทุก}{mtk.12}
\bookmark[dest={mtk.12},level=1]{9. สนิทสฺสนทุก}
\csromanheader{1}{9. สนิทสฺสนทุก 1. ปฏิจฺจวาร}
\csromanheader{2}{1. ปจฺจยานุโลม 1. วิภงฺควาร}
\csromanheader{2}{\textbf{เหตุ}}
\proseitem{18}{อนิทสฺสนํ ธมฺมํ ปฏิจฺจ อนิทสฺสโน ธมฺโม อุปฺปชฺชติ เหตุปจฺจยา.\csromanspacer{} อนิทสฺสนํ เอกํ ขนฺธํ ปฏิจฺจ ตโย ขนฺธา อนิทสฺสนํ จิตฺตสมุฏฺฐานญฺจ รูปํ ฯเปฯ เทฺว ขนฺเธ ฯเปฯ.\csromanspacer{} ปฏิสนฺธิกฺขเณ อนิทสฺสนํ เอกํ ขนฺธํ ปฏิจฺจ ตโย ขนฺธา อนิทสฺสนํ กฏตฺตา จ รูปํ ฯเปฯ เทฺว ขนฺเธ ฯเปฯ.\csromanspacer{} ขนฺเธ ปฏิจฺจ วตฺถุ, วตฺถุํ ปฏิจฺจ ขนฺธา.\csromanspacer{} เอกํ มหาภูตํ ฯเปฯ มหาภูเต ปฏิจฺจ อนิทสฺสนํ จิตฺตสมุฏฺฐานํ รูปํ, กฏตฺตารูปํ, อุปาทารูปํ. (1)}

\csromanheader{2}{9. สนิทสฺสนทุก}
\prose{\csromanfolio{91}อนิทสฺสนํ ธมฺมํ ปฏิจฺจ สนิทสฺสโน ธมฺโม อุปฺปชฺชติ เหตุปจฺจยา.\csromanspacer{} อนิทสฺสเน ขนฺเธ ปฏิจฺจ สนิทสฺสนํ จิตฺตสมุฏฺฐานํ รูปํ.\csromanspacer{} ปฏิสนฺธิกฺขเณ ฯเปฯ มหาภูเต ปฏิจฺจ สนิทสฺสนํ จิตฺตสมุฏฺฐานํ รูปํ, กฏตฺตารูปํ, อุปาทารูปํ. (2)}

\prose{อนิทสฺสนํ ธมฺมํ ปฏิจฺจ สนิทสฺสโน จ อนิทสฺสโน จ ธมฺมา อุปฺปชฺชนฺติ เหตุปจฺจยา.\csromanspacer{} อนิทสฺสนํ เอกํ ขนฺธํ ปฏิจฺจ ตโย ขนฺธา สนิทสฺสนญฺจ อนิทสฺสนญฺจ จิตฺตสมุฏฺฐานญฺจ รูปํ ฯเปฯ เทฺว ขนฺเธ ฯเปฯ.\csromanspacer{} ปฏิสนฺธิกฺขเณ ฯเปฯ มหาภูเต ปฏิจฺจ สนิทสฺสนญฺจ อนิทสฺสนญฺจ จิตฺตสมุฏฺฐานญฺจ รูปํ, กฏตฺตารูปํ, อุปาทารูปํ. (3)}

\csromanheader{2}{\textbf{อารมฺมณ}}
\proseitem{19}{อนิทสฺสนํ ธมฺมํ ปฏิจฺจ อนิทสฺสโน ธมฺโม อุปฺปชฺชติ อารมฺมณ ปจฺจโย.\csromanspacer{} อนิทสฺสนํ เอกํ ขนฺธํ ปฏิจฺจ ตโย ขนฺธา ฯเปฯ เทฺว ขนฺเธ ฯเปฯ.\csromanspacer{} ปฏิสนฺธิกฺขเณ ฯเปฯ.\csromanspacer{} วตฺถุํ ปฏิจฺจ ขนฺธา. (1)}

\csromanheader{2}{\textbf{อธิปติ}}
\proseitem{20}{อนิทสฺสนํ ธมฺมํ ปฏิจฺจ อนิทสฺสโน ธมฺโม อุปฺปชฺชติ อธิปติปจฺจยา.\csromanspacer{} อนิทสฺสนํ เอกํ ขนฺธํ ปฏิจฺจ ตโย ขนฺธา อนิทสฺสนํ จิตฺตสมุฏฺฐานญฺจ รูปํ ฯเปฯ เทฺว ขนฺเธ ฯเปฯ.\csromanspacer{} เอกํ มหาภูตํ ปฏิจฺจ ตโย มหาภูตา ฯเปฯ เทฺว มหาภูเต ฯเปฯ มหาภูเต ปฏิจฺจ อนิทสฺสนํ จิตฺตสมุฏฺฐานํ รูปํ, อุปาทารูปํ. (1)}

\prose{อนิทสฺสนํ ธมฺมํ ปฏิจฺจ สนิทสฺสโน ธมฺโม อุปฺปชฺชติ อธิปติปจฺจยา.\csromanspacer{} อนิทสฺสเน ขนฺเธ ปฏิจฺจ สนิทสฺสนํ จิตฺตสมุฏฺฐานํ รูปํ.\csromanspacer{} มหาภูเต ปฏิจฺจ สนิทสฺสนํ จิตฺตสมุฏฺฐานํ รูปํ. (2)}

\prose{อนิทสฺสนํ ธมฺมํ ปฏิจฺจ สนิทสฺสโน จ อนิทสฺสโน จ ธมฺมา อุปฺปชฺชนฺติ อธิปติปจฺจยา.\csromanspacer{} อนิทสฺสนํ เอกํ ขนฺธํ ปฏิจฺจ ตโย ขนฺธา สนิทสฺสนญฺจ อนิทสฺสนญฺจ จิตฺตสมุฏฺฐานํ รูปํ ฯเปฯ เทฺว ขนฺเธ ฯเปฯ มหาภูเต ปฏิจฺจ สนิทสฺสนญฺจ อนิทสฺสนญฺจ จิตฺตสมุฏฺฐานํ รูปํ, อุปาทารูปํ. (3) (สํขิตฺตํ.\csromanspacer{} สพฺเพ กาตพฺพา.)}

\csromanheader{2}{1. ปจฺจยานุโลม 2. สงฺขฺยาวาร}
\csromanheader{2}{\textbf{สุทฺธ}}
\proseitem{21}{\csromanfolio{92}เหตุยา ตีณิ, อารมฺมเณ เอกํ, อธิปติยา ตีณิ, อนนฺตเร เอกํ, สมนนฺตเร เอกํ, สหชาเต ตีณิ, อญฺญมญฺเญ เอกํ, นิสฺสเย ตีณิ, อุปนิสฺสเย เอกํ, ปุเรชาเต เอกํ, อาเสวเน เอกํ, กมฺเม ตีณิ, วิปาเก ตีณิ, (สพฺพตฺถ ตีณิ.) มคฺเค ตีณิ, สมฺปยุตฺเต เอกํ, วิปฺปยุตฺเต ตีณิ, อตฺถิยา ตีณิ, นตฺถิยา เอกํ, วิคเต เอกํ, อวิคเต ตีณิ. (เอวํ คเณตพฺพํ.)}

\csromanheader{2}{อนุโลมํ}
\csromanheader{2}{2. ปจฺจยปจฺจนีย 1. วิภงฺควาร}
\csromanheader{2}{\textbf{นเหตุ}}
\proseitem{22}{อนิทสฺสนํ ธมฺมํ ปฏิจฺจ อนิทสฺสโน ธมฺโม อุปฺปชฺชติ นเหตุปจฺจยา.\csromanspacer{} อเหตุกํ อนิทสฺสนํ เอกํ ขนฺธํ ปฏิจฺจ ตโย ขนฺธา อนิทสฺสนํ จิตฺตสมุฏฺฐานญฺจ รูปํ ฯเปฯ เทฺว ขนฺเธ ฯเปฯ.\csromanspacer{} อเหตุกปฏิสนฺธิกฺขเณ ฯเปฯ.\csromanspacer{} ขนฺเธ ปฏิจฺจ วตฺถุ, วตฺถุํ ปฏิจฺจ ขนฺธา.\csromanspacer{} เอกํ มหาภูตํ ฯเปฯ มหาภูเต ปฏิจฺจ อนิทสฺสนํ จิตฺตสมุฏฺฐานํ รูปํ, กฏตฺตารูปํ, อุปาทารูปํ.\csromanspacer{} พาหิรํ.\csromanspacer{} อาหารสมุฏฺฐานํ.\csromanspacer{} อุตุสมุฏฺฐานํ.\csromanspacer{} อสญฺญสตฺตานํ ฯเปฯ.\csromanspacer{} วิจิกิจฺฉาสหคเต อุทฺธจฺจสหคเต ขนฺเธ ปฏิจฺจ วิจิกิจฺฉาสหคโต อุทฺธจฺจสหคโต โมโห. (1)}

\prose{อนิทสฺสนํ ธมฺมํ ปฏิจฺจ สนิทสฺสโน ธมฺโม อุปฺปชฺชติ นเหตุปจฺจยา.\csromanspacer{} อเหตุเก อนิทสฺสเน ขนฺเธ ปฏิจฺจ สนิทสฺสนํ จิตฺตสมุฏฺฐานํ รูปํ.\csromanspacer{} ปฏิสนฺธิกฺขเณ ฯเปฯ มหาภูเต ปฏิจฺจ สนิทสฺสนํ จิตฺตสมุฏฺฐานํ รูปํ, กฏตฺตารูปํ, อุปาทารูปํ.\csromanspacer{} พาหิเร.\csromanspacer{} อาหารสมุฏฺฐาเน.\csromanspacer{} อุตุสมุฏฺฐาเน.\csromanspacer{} อสญฺญสตฺตานํ มหาภูเต ปฏิจฺจ สนิทสฺสนํ กฏตฺตารูปํ, อุปาทารูปํ. (2)}

\prose{อนิทสฺสนํ ธมฺมํ ปฏิจฺจ สนิทสฺสโน จ อนิทสฺสโน จ ธมฺมา อุปฺปชฺชนฺติ นเหตุปจฺจยา.\csromanspacer{} อเหตุกํ อนิทสฺสนํ เอกํ ขนฺธํ ปฏิจฺจ ตโย ขนฺธา}

\vspace{\tipitakaparskip}
\csromancenter{สนิทสฺสนทุก สนิทสฺสนญฺจ อนิทสฺสนญฺจ จิตฺตสมุฏฺฐานํ รูปํ ฯเปฯ เทฺว ขนฺเธ ฯเปฯ.\csromanspacer{} ปฏิสนฺธิกฺขเณ ฯเปฯ มหาภูเต ปฏิจฺจ ฯเปฯ.\csromanspacer{} พาหิเร.\csromanspacer{} อาหารสมุฏฺฐาเน.\csromanspacer{} อุตุสมุฏฺฐาเน.\csromanspacer{} อสญฺญสตฺตานํ มหาภูเต ปฏิจฺจ สนิทสฺสนญฺจ อนิทสฺสนญฺจ กฏตฺตารูปํ, อุปาทารูปํ. (3) (เอวํ สพฺเพ กาตพฺพา.)}
\csromanheader{2}{2. ปจฺจยปจฺจนีย 2. สงฺขฺยาวาร}
\csromanheader{2}{\textbf{สุทฺธ}}
\proseitem{23}{\csromanfolio{93}นเหตุยา ตีณิ น-อารมฺมเณ ตีณิ น-อธิปติยา ตีณิ, น-อนนฺตเร ตีณิ, นสมนนฺตเร ตีณิ น-อญฺญมญฺเญ ตีณิ, น-อุปนิสฺสเย ตีณิ นปุเรชาเต ตีณิ, นปจฺฉาชาเต ตีณิ, น-อาเสวเน ตีณิ, นกมฺเม ตีณิ, นวิปาเก ตีณิ, น-อาหาเร ตีณิ, น-อินฺทฺริเย ตีณิ, นฌาเน ตีณิ, นมคฺเค ตีณิ, นสมฺปยุตฺเต ตีณิ, นวิปฺปยุตฺเต ตีณิ, โนนตฺถิยา ตีณิ, โนวิคเต ตีณิ. (เอวํ คเณตพฺพํ.)}

\vspace{\tipitakaparskip}
\csromancenter{ปจฺจนียํ.}
\csromanheader{2}{3. ปจฺจยานุโลมปจฺจนีย}
\csromanheader{2}{\textbf{เหตุทุก}}
\proseitem{24}{เหตุปจฺจยา น-อารมฺมเณ ตีณิ, น-อธิปติยา ตีณิ, (สพฺพตฺถ ตีณิ.) นกมฺเม เอกํ, นวิปาเก ตีณิ, นสมฺปยุตฺเต ตีณิ, นวิปฺปยุตฺเต เอกํ, โนนตฺถิยา ตีณิ, โนวิคเต ตีณิ.}

\vspace{\tipitakaparskip}
\csromancenter{อนุโลมปจฺจนียํ.}
\csromanheader{2}{4. ปจฺจยปจฺจนียานุโลม}
\csromanheader{2}{\textbf{นเหตุทุก}}
\proseitem{25}{นเหตุปจฺจยา อารมฺมเณ เอกํ, อนนฺตเร เอกํ, สมนนฺตเร เอกํ, สหชาเต ตีณิ, อญฺญมญฺเญ เอกํ, นิสฺสเย ตีณิ, อุปนิสฺสเย \csromanfolio{94}เอกํ, ปุเรชาเต เอกํ, อาเสวเน เอกํ, กมฺเม ตีณิ ฯเปฯ ฌาเน ตีณิ, มคฺเค เอกํ, สมฺปยุตฺเต เอกํ, วิปฺปยุตฺเต ตีณิ, อตฺถิยา ตีณิ, นตฺถิยา เอกํ, วิคเต เอกํ, อวิคเต ตีณิ.}

\vspace{\tipitakaparskip}
\csromancenter{ปจฺจนียานุโลมํ.}
\csromancenter{(สหชาตวาโรปิ ปฏิจฺจวารสทิโส.)}
\csromanheader{2}{9. สนิทสฺสนทุก 3. ปจฺจยวาร}
\csromanheader{2}{1. ปจฺจยานุโลม}
\csromanheader{2}{\textbf{เหตุ}}
\proseitem{26}{อนิทสฺสนํ ธมฺมํ ปจฺจยา อนิทสฺสโน ธมฺโม อุปฺปชฺชติ เหตุปจฺจยา.\csromanspacer{} อนิทสฺสนํ เอกํ ขนฺธํ ปจฺจยา ตโย ขนฺธา อนิทสฺสนํ จิตฺตสมุฏฺฐานญฺจ รูปํ ฯเปฯ เทฺว ขนฺเธ ฯเปฯ.\csromanspacer{} ปฏิสนฺธิกฺขเณ ฯเปฯ.\csromanspacer{} ขนฺเธ ปจฺจยา วตฺถุ, วตฺถุํ ปจฺจยา ขนฺธา.\csromanspacer{} เอกํ มหาภูตํ ปจฺจยา ฯเปฯ มหาภูเต ปจฺจยา อนิทสฺสนํ จิตฺตสมุฏฺฐานํ รูปํ, กฏตฺตารูปํ, อุปาทารูปํ.\csromanspacer{} วตฺถุํ ปจฺจยา อนิทสฺสนา ขนฺธา. (อิตเรปิ เทฺว ปญฺหา กาตพฺพา.)}

\csromanheader{2}{\textbf{อารมฺมณ}}
\proseitem{27}{อนิทสฺสนํ ธมฺมํ ปจฺจยา อนิทสฺสโน ธมฺโม อุปฺปชฺชติ อารมฺมณปจฺจยา.\csromanspacer{} อนิทสฺสนํ เอกํ ขนฺธํ ฯเปฯ เทฺว ขนฺเธ ฯเปฯ.\csromanspacer{} ปฏิสนฺธิกฺขเณ ฯเปฯ.\csromanspacer{} วตฺถุํ ปจฺจยา ขนฺธา.\csromanspacer{} จกฺขายตนํ ปจฺจยา จกฺขุวิญฺญาณํ ฯเปฯ.\csromanspacer{} กายายตนํ ปจฺจยา กายวิญฺญาณํ.\csromanspacer{} วตฺถุํ ปจฺจยา อนิทสฺสนา ขนฺธา. (สํขิตฺตํ.)}

\vspace{\tipitakaparskip}
\csromancenter{\textbf{เหตุ}ยา ตีณิ, อารมฺมเณ เอกํ, อธิปติยา ตีณิ ฯเปฯ อวิคเต ตีณิ.}
\csromancenter{อนุโลมํ.}
\csromanheader{2}{9. สนิทสฺสนทุก \textbf{2. ปจฺจยปจฺจนีย}}
\csromanheader{2}{\textbf{นเหตุ}}
\proseitem{29}{\csromanfolio{95}อนิทสฺสนํ ธมฺมํ ปจฺจยา อนิทสฺสโน ธมฺโม อุปฺปชฺชติ นเหตุปจฺจยา.\csromanspacer{} อเหตุกํ อนิทสฺสนํ เอกํ ขนฺธํ ปจฺจยา ตโย ขนฺธา อนิทสฺสนํ จิตฺตสมุฏฺฐานญฺจ รูปํ ฯเปฯ เทฺว ขนฺเธ ฯเปฯ.\csromanspacer{} อเหตุกปฏิสนฺธิกฺขเณ ฯเปฯ. (ยาว อสญฺญสตฺตา.) จกฺขายตนํ ปจฺจยา จกฺขุวิญฺญาณํ ฯเปฯ.\csromanspacer{} กายายตนํ ปจฺจยา กายวิญฺญาณํ.\csromanspacer{} วตฺถุํ ปจฺจยา อเหตุกา อนิทสฺสนา ขนฺธา.\csromanspacer{} วิจิกิจฺฉาสหคเต อุทฺธจฺจสหคเต ขนฺเธ จ วตฺถุญฺจ ปจฺจยา วิจิกิจฺฉาสหคโต อุทฺธจฺจสหคโต โมโห. (อิตเรปิ เทฺว กาตพฺพา.\csromanspacer{} สํขิตฺตํ.)}

\vspace{\tipitakaparskip}
\csromancenter{\textbf{นเหตุ}ยา ตีณิ, น-อารมฺมเณ ตีณิ ฯเปฯ โนวิคเต ตีณิ.}
\csromancenter{ปจฺจนียํ.}
\csromanheader{2}{3. ปจฺจยานุโลมปจฺจนีย}
\proseitem{31}{เหตุปจฺจยา น-อารมฺมเณ ตีณิ ฯเปฯ นกมฺเม เอกํ ฯเปฯ นวิปฺปยุตฺเต เอกํ, โนนตฺถิยา ตีณิ, โนวิคเต ตีณิ.}

\vspace{\tipitakaparskip}
\csromancenter{อนุโลมปจฺจนียํ.}
\csromanheader{2}{4. ปจฺจยปจฺจนียานุโลม}
\vspace{\tipitakaparskip}
\csromancenter{\textbf{นเหตุ}ปจฺจยา อารมฺมเณ เอกํ ฯเปฯ มคฺเค เอกํ ฯเปฯ อวิคเตตีณิ.}
\csromancenter{ปจฺจนียานุโลมํ.}
\csromancenter{(นิสฺสยวาโรปิ เอวํ กาตพฺโพ.)}
\csromanheader{2}{9. สนิทสฺสนทุก 5. สํสฏฺฐวาร \textbf{1. ปจฺจยานุโลม}}
\proseitem{33}{อนิทสฺสนํ ธมฺมํ สํสฏฺโฐ อนิทสฺสโน ธมฺโม อุปฺปชฺชติ เหตุปจฺจยา.\csromanspacer{} อนิทสฺสนํ เอกํ ขนฺธํ สํสฏฺฐา ตโย ขนฺธา ฯเปฯ เทฺว ขนฺเธ ฯเปฯ.\csromanspacer{} ปฏิสนฺธิกฺขเณ ฯเปฯ.}

\prose{\csromanfolio{96}อนิทสฺสนํ ธมฺมํ สํสฏฺโฐ อนิทสฺสโน ธมฺโม อุปฺปชฺชติ อารมฺมณปจฺจยา. (เอวํ สพฺพํ สปฺปจฺจยคณนาหิ สทฺธิํ กาตพฺพํ.)}

\vspace{\tipitakaparskip}
\csromancenter{(สมฺปยุตฺตวาโรปิ สํสฏฺฐวารสทิโส.)}
\csromanheader{2}{9. สนิทสฺสนทุก 7. ปญฺหาวาร}
\csromanheader{2}{1. ปจฺจยานุโลม 1. วิภงฺควาร}
\csromanheader{2}{\textbf{เหตุ}}
\proseitem{34}{อนิทสฺสโน ธมฺโม อนิทสฺสนสฺส ธมฺมสฺส เหตุปจฺจเยน ปจฺจโย.\csromanspacer{} อนิทสฺสนา เหตู สมฺปยุตฺตกานํ ขนฺธานํ อนิทสฺสนานํ จิตฺตสมุฏฺฐานานญฺจ รูปานํ เหตุปจฺจเยน ปจฺจโย.\csromanspacer{} ปฏิสนฺธิกฺขเณ ฯเปฯ. (1)}

\prose{อนิทสฺสโน ธมฺโม สนิทสฺสนสฺส ธมฺมสฺส เหตุปจฺจเยน ปจฺจโย.\csromanspacer{} อนิทสฺสนา เหตู สนิทสฺสนานํ จิตฺตสมุฏฺฐานานํ รูปานํ เหตุปจฺจเยน ปจฺจโย.\csromanspacer{} ปฏิสนฺธิกฺขเณ ฯเปฯ. (2)}

\prose{อนิทสฺสโน ธมฺโม สนิทสฺสนสฺส จ อนิทสฺสนสฺส จ ธมฺมสฺส เหตุปจฺจเยน ปจฺจโย.\csromanspacer{} อนิทสฺสนา เหตู สมฺปยุตฺตกานํ ขนฺธานํ สนิทสฺสนานญฺจ อนิทสฺสนานญฺจ จิตฺตสมุฏฺฐานานญฺจ รูปานํ เหตุปจฺจเยน ปจฺจโย.\csromanspacer{} ปฏิสนฺธิกฺขเณ ฯเปฯ. (3)}

\csromanheader{2}{\textbf{อารมฺมณ}}
\proseitem{35}{ยนิทสฺสโน ธมฺโม อนิทสฺสนสฺส ธมฺมสฺส อารมฺมณปจฺจเยน ปจฺจโย.\csromanspacer{} สนิทสฺสนํ รูปํ อนิจฺจโต ฯเปฯ โทมนสฺสํ อุปฺปชฺชติ.\csromanspacer{} ทิพฺเพน จกฺขุนา รูปํ ปสฺสติ.\csromanspacer{} รูปายตนํ จกฺขุวิญฺญาณสฺส อารมฺมณปจฺจเยน ปจฺจโย.\csromanspacer{} สนิทสฺสนา ขนฺธา อิทฺธิวิธญาณสฺส, ปุพฺเพนิวาสานุสฺสติญาณสฺส, อนาคตํสญาณสฺส, อาวชฺชนาย อารมฺมณปจฺจเยน ปจฺจโย. (1)}

\prose{อนิทสฺสโน ธมฺโม อนิทสฺสนสฺส ธมฺมสฺส อารมฺมณปจฺจเยน ปจฺจโย.\csromanspacer{} ทานํ ทตฺวา, สีลํ ฯเปฯ อุโปสถกมฺมํ กตฺวา ตํ ปจฺจเวกฺขติ ปุพฺเพ สุจิณฺณานิ ปจฺจเวกฺขติ.\csromanspacer{} ฌานา วุฏฺฐหิตฺวา ฯเปฯ.\csromanspacer{} อริยา มคฺคา วุฏฺฐหิตฺวา มคฺคํ ปจฺจเวกฺขนฺติ,}

\vspace{\tipitakaparskip}
\csromancenter{สนิทสฺสนทุก ผลํ ปจฺจเวกฺขนฺติ, นิพฺพานํ ปจฺจเวกฺขนฺติ.\csromanspacer{} นิพฺพานํ โคตฺรภุสฺส, โวทานสฺส, มคฺคสฺส, ผลสฺส, อาวชฺชนาย อารมฺมณปจฺจเยน ปจฺจโย.\csromanspacer{} อริยา ปหีเน กิเลเส ปจฺจเวกฺขนฺติ, วิกฺขมฺภิเต กิเลเส ฯเปฯ.\csromanspacer{} ปุพฺเพ ฯเปฯ.\csromanspacer{} จกฺขุํ ฯเปฯ.\csromanspacer{} กายํ.\csromanspacer{} สทฺเท ฯเปฯ.\csromanspacer{} วตฺถุํ อนิทสฺสเน ขนฺเธ อนิจฺจโต ฯเปฯ โทมนสฺสํ อุปฺปชฺชติ.\csromanspacer{} ทิพฺพาย โสตธาตุยา สทฺทํ สุณาติ.\csromanspacer{} เจโตปริยญาเณน อนิทสฺสนจิตฺตสมงฺคิสฺส จิตฺตํ ชานาติ.\csromanspacer{} อากาสานญฺจายตนํ วิญฺญาณญฺจายตนสฺส ฯเปฯ.\csromanspacer{} อากิญฺจญฺญายตนํ เนวสญฺญานาสญฺญายตนสฺส ฯเปฯ.\csromanspacer{} สทฺทายตนํ โสตวิญฺญาณสฺส ฯเปฯ.\csromanspacer{} โผฏฺฐพฺพายตนํ กายวิญฺญาณสฺส ฯเปฯ.\csromanspacer{} อนิทสฺสนา ขนฺธา อิทฺธิวิธญาณสฺส, เจโตปริยญาณสฺส, ปุพฺเพนิวาสานุสฺสติญาณสฺส, ยถากมฺมูปคญาณสฺส, อนาคตํสญาณสฺส, อาวชฺชนาย อารมฺมณปจฺจเยน ปจฺจโย. (1)}
\csromanheader{2}{\textbf{อธิปติ}}
\proseitem{36}{\csromanfolio{97}สนิทสฺสโน ธมฺโม อนิทสฺสนสฺส ธมฺมสฺส อธิปติปจฺจเยน ปจฺจโย.\csromanspacer{} \textbf{อารมฺมณาธิปติ\csromandash{}}สนิทสฺสนํ รูปํ ครุํ กตฺวา อสฺสาเทติ อภินนฺทติ, ตํ ครุํ กตฺวา ราโค อุปฺปชฺชติ, ทิฏฺฐิ อุปฺปชฺชติ. (1)}

\prose{อนิทสฺสโน ธมฺโม อนิทสฺสนสฺส ธมฺมสฺส อธิปติปจฺจเยน ปจฺจโย.\csromanspacer{} \textbf{อารมฺมณาธิปติ}, สหชาตาธิปติ.\csromanspacer{}\csromanspacer{}\textbf{อารมฺมณาธิปติ\csromandash{}}ทานํ ทตฺวา, สีลํ ฯเปฯ อุโปสถกมฺมํ กตฺวา ตํ ครุํ กตฺวา ปจฺจเวกฺขติ.\csromanspacer{} ปุพฺเพ สุจิณฺณานิ ฯเปฯ.\csromanspacer{} ฌานา วุฏฺฐหิตฺวา ฯเปฯ.\csromanspacer{} อริยา มคฺคา วุฏฺฐหิตฺวา มคฺคํ ครุํ กตฺวา ปจฺจเวกฺขนฺติ, ผลํ ครุํ กตฺวา ปจฺจเวกฺขนฺติ, นิพฺพานํ ครุํ กตฺวา ปจฺจเวกฺขนฺติ.\csromanspacer{} นิพฺพานํ โคตฺรภุสฺส, โวทานสฺส, มคฺคสฺส, ผลสฺส อธิปติปจฺจเยน ปจฺจโย.\csromanspacer{} จกฺขุํ ฯเปฯ วตฺถุํ อนิทสฺสเน ขนฺเธ ครุํ กตฺวา อสฺสาเทติ อภินนฺทติ, ตํ ครุํ กตฺวา ราโค อุปฺปชฺชติ, ทิฏฺฐิ อุปฺปชฺชติ.\csromanspacer{} \textbf{สหชาตาธิปติ\csromandash{}}อนิทสฺสนาธิปติ สมฺปยุตฺตกานํ ขนฺธานํ อนิทสฺสนานํ จิตฺตสมุฏฺฐานานญฺจ รูปานํ อธิปติปจฺจเยน ปจฺจโย. (1)}

\prose{อนิทสฺสโน ธมฺโม สนิทสฺสนสฺส ธมฺมสฺส อธิปติปจฺจเยน ปจฺจโย.\csromanspacer{} \textbf{สหชาตาธิปติ\csromandash{}}อนิทสฺสนาธิปติ สนิทสฺสนานํ จิตฺตสมุฏฺฐานานํ รูปานํ อธิปติปจฺจเยน ปจฺจโย. (2)}

\prose{อนิทสฺสโน ธมฺโม สนิทสฺสนสฺส จ อนิทสฺสนสฺส จ ธมฺมสฺส อธิปติปจฺจเยน ปจฺจโย.\csromanspacer{} \textbf{สหชาตาธิปติ\csromandash{}}อนิทสฺสนาธิปติ สมฺปยุตฺตกานํ \csromanfolio{98}ขนฺธานํ สนิทสฺสนานญฺจ อนิทสฺสนานญฺจ จิตฺตสมุฏฺฐานานญฺจ รูปานํ อธิปติปจฺจเยน ปจฺจโย.(3)}

\csromanheader{2}{\textbf{อนนฺตร}}
\proseitem{37}{อนิทสฺสโน ธมฺโม อนิทสฺสนสฺส ธมฺมสฺส อนนฺตรปจฺจเยน ปจฺจโย.\csromanspacer{} ปุริมา ปุริมา อนิทสฺสนา ขนฺธา ปจฺฉิมานํ ปจฺฉิมานํ อนิทสฺสนานํ ขนฺธานํ อนนฺตรปจฺจเยน ปจฺจโย.\csromanspacer{} อนุโลมํ โคตฺรภุสฺส ฯเปฯ.\csromanspacer{} โคตฺรภุ มคฺคสฺส ฯเปฯ.\csromanspacer{} เนวสญฺญานาสญฺญายตนํ ผลสมาปตฺติยา อนนฺตรปจฺจเยน ปจฺจโย. (1)}

\csromanheader{2}{\textbf{สมนนฺตราทิ}}
\proseitem{38}{อนิทสฺสโน ธมฺโม อนิทสฺสนสฺส ธมฺมสฺส สมนนฺตรปจฺจเยน ปจฺจโย.\csromanspacer{} สหชาตปจฺจเยน ปจฺจโย.\csromanspacer{} ตีณิ.\csromanspacer{} อญฺญมญฺญปจฺจเยน ปจฺจโย.\csromanspacer{} เอกํ.\csromanspacer{} นิสฺสยปจฺจเยน ปจฺจโย.\csromanspacer{} ตีณิ.}

\csromanheader{2}{\textbf{อุปนิสฺสย}}
\proseitem{39}{สนิทสฺสโน ธมฺโม อนิทสฺสนสฺส ธมฺมสฺส อุปนิสฺสยปจฺจเยน ปจฺจโย.\csromanspacer{} \textbf{อารมฺมณูปนิสฺสโย, ปกตูปนิสฺสโย ฯเปฯ.}\csromanspacer{} ปกตูปนิสฺสโย\csromandash{}วณฺณสมฺปทํ ปตฺถยมาโน ทานํ ฯเปฯ สีลํ ฯเปฯ อุโปสถกมฺมํ กโรติ.\csromanspacer{} วณฺณสมฺปทา สทฺธาย ฯเปฯ ปตฺถนาย.\csromanspacer{} กายิกสฺส สุขสฺส, กายิกสฺส ทุกฺขสฺส, มคฺคสฺส, ผลสมาปตฺติยา อุปนิสฺสยปจฺจเยน ปจฺจโย. (1)}

\prose{อนิทสฺสโน ธมฺโม อนิทสฺสนสฺส ธมฺมสฺส อุปนิสฺสยปจฺจเยน ปจฺจโย.\csromanspacer{} \textbf{อารมฺมณูปนิสฺสโย, อนนฺตรูปนิสฺสโย, ปกตูปนิสฺสโย ฯเปฯ.}\csromanspacer{} ปกตูปนิสฺสโย\csromandash{}สทฺธํ อุปนิสฺสาย ทานํ เทติ ฯเปฯ สมาปตฺติํ อุปฺปาเทติ.\csromanspacer{} มานํ ชปฺเปติ.\csromanspacer{} ทิฏฺฐํ คณฺหาติ.\csromanspacer{} สีลํ ฯเปฯ.\csromanspacer{} เสนาสนํ อุปนิสฺสาย ทานํ เทติ ฯเปฯ สํฆํ ภินฺทติ.\csromanspacer{} สทฺธา ฯเปฯ.\csromanspacer{} เสนาสนํ สทฺธาย ฯเปฯ ผลสมาปตฺติยา อุปนิสฺสยปจฺจเยน ปจฺจโย. (1)}

\csromanheader{2}{\textbf{ปุเรชาต}}
\proseitem{40}{สนิทสฺสโน ธมฺโม อนิทสฺสนสฺส ธมฺมสฺส ปุเรชาตปจฺจเยน ปจฺจโย.\csromanspacer{} สนิทสฺสนํ รูปํ อนิจฺจโต ฯเปฯ โทมนสฺสํ อุปฺปชฺชติ.\csromanspacer{} ทิพฺเพน}

\vspace{\tipitakaparskip}
\csromancenter{สนิทสฺสนทุก จกฺขุนา รูปํ ปสฺสติ.\csromanspacer{} รูปายตนํ จกฺขุวิญฺญาณสฺส ปุเรชาตปจฺจเยน ปจฺจโย. (1)}
\prose{\csromanfolio{99}อนิทสฺสโน ธมฺโม อนิทสฺสนสฺส ธมฺมสฺส ปุเรชาตปจฺจเยน ปจฺจโย.\csromanspacer{} \textbf{อารมฺมณปุเรชาตํ, วตฺถุปุเรชาตํ}.\csromanspacer{}\csromanspacer{}\textbf{อารมฺมณปุเรชาตํ}\csromandash{} จกฺขุํ ฯเปฯ วตฺถุํ อนิจฺจโต ฯเปฯ โทมนสฺสํ อุปฺปชฺชติ.\csromanspacer{} ทิพฺพาย โสตธาตุยา สทฺทํ สุณาติ.\csromanspacer{} สทฺทายตนํ โสตวิญฺญาณสฺส ฯเปฯ.\csromanspacer{} โผฏฺฐพฺพายตนํ กายวิญฺญาณสฺส ปุเรชาตปจฺจเยน ปจฺจโย.\csromanspacer{} \textbf{วตฺถุปุเรชาตํ}\csromandash{}จกฺขายตนํ จกฺขุวิญฺญาณสฺส ฯเปฯ.\csromanspacer{} กายายตนํ กายวิญฺญาณสฺส ฯเปฯ.\csromanspacer{} วตฺถุ อนิทสฺสนานํ ขนฺธานํ ปุเรชาตปจฺจเยน ปจฺจโย. (1)}

\prose{สนิทสฺสโน จ อนิทสฺสโน จ ธมฺโม อนิทสฺสนสฺส ธมฺมสฺส ปุเรชาตปจฺจเยน ปจฺจโย.\csromanspacer{} \textbf{อารมฺมณปุเรชาตํ, วตฺถุปุเรชาตํ}\csromandash{} รูปายตนญฺจ วตฺถุ จ อนิทสฺสนานํ ขนฺธานํ ปุเรชาตปจฺจเยน ปจฺจโย.\csromanspacer{} รูปายตนญฺจ จกฺขายตนญฺจ จกฺขุวิญฺญาณสฺส ปุเรชาตปจฺจเยน ปจฺจโย. (1)}

\csromanheader{2}{\textbf{ปจฺฉาชาต}}
\proseitem{41}{อนิทสฺสโน ธมฺโม อนิทสฺสนสฺส ธมฺมสฺส ปจฺฉาชาตปจฺจเยน ปจฺจโย.\csromanspacer{} ปจฺฉาชาตา อนิทสฺสนา ขนฺธา ปุเรชาตสฺส อิมสฺส อนิทสฺสนสฺส กายสฺส ปจฺฉาชาตปจฺจเยน ปจฺจโย. (1)}

\prose{อนิทสฺสโน ธมฺโม สนิทสฺสนสฺส ธมฺมสฺส ปจฺฉาชาตปจฺจเยน ปจฺจโย.\csromanspacer{} ปจฺฉาชาตา อนิทสฺสนา ขนฺธา ปุเรชาตสฺส อิมสฺส สนิทสฺสนสฺส กายสฺส ปจฺฉาชาตปจฺจเยน ปจฺจโย. (2)}

\prose{อนิทสฺสโน ธมฺโม สนิทสฺสนสฺส จ อนิทสฺสนสฺส จ ธมฺมสฺส ปจฺฉาชาตปจฺจเยน ปจฺจโย.\csromanspacer{} ปจฺฉาชาตา อนิทสฺสนา ขนฺธา ปุเรชาตสฺส อิมสฺส สนิทสฺสนสฺส จ อนิทสฺสนสฺส จ กายสฺส ปจฺฉาชาต ปจฺจเยน ปจฺจโย. (3)}

\csromanheader{2}{\textbf{อาเสวน}}
\proseitem{42}{อนิทสฺสโน ธมฺโม อนิทสฺสนสฺส ธมฺมสฺส อาเสวนปจฺจเยน ปจฺจโย.\csromanspacer{} ปุริมา ปุริมา อนิทสฺสนา ขนฺธา ปจฺฉิมานํ ปจฺฉิมานํ อนิทสฺสนานํ ขนฺธานํ ฯเปฯ.\csromanspacer{} อนุโลมํ โคตฺรภุสฺส.\csromanspacer{} อนุโลมํ โวทานสฺส.\csromanspacer{} โคตฺรภุ มคฺคสฺส.\csromanspacer{} โวทานํ มคฺคสฺส อาเสวนปจฺจเยน ปจฺจโย. (1)}

\csromanheader{2}{\textbf{กมฺม}}
\proseitem{43}{\csromanfolio{100}อนิทสฺสโน ธมฺโม อนิทสฺสนสฺส ธมฺมสฺส กมฺมปจฺจเยน ปจฺจโย.\csromanspacer{} \textbf{สหชาตา, นานากฺขณิกา}.\csromanspacer{}\csromanspacer{}\textbf{สหชาตา} อนิทสฺสนา เจตนา สมฺปยุตฺตกานํ ขนฺธานํ อนิทสฺสนานญฺจ จิตฺตสมุฏฺฐานานํ รูปานํ กมฺมปจฺจเยน ปจฺจโย.\csromanspacer{} ปฏิสนฺธิกฺขเณ ฯเปฯ.\csromanspacer{} \textbf{นานากฺขณิกา} อนิทสฺสนา เจตนา วิปากานํ ขนฺธานํ อนิทสฺสนานญฺจ กฏตฺตารูปานํ กมฺมปจฺจเยน ปจฺจโย. (1)}

\prose{อนิทสฺสโน ธมฺโม สนิทสฺสนสฺส ธมฺมสฺส กมฺมปจฺจเยน ปจฺจโย.\csromanspacer{} \textbf{สหชาตา, นานากฺขณิกา}. (วิตฺถาเรตพฺพํ.) (2)}

\prose{อนิทสฺสโน ธมฺโม สนิทสฺสนสฺส จ อนิทสฺสนสฺส จ ธมฺมสฺส กมฺมปจฺจเยน ปจฺจโย.\csromanspacer{} \textbf{สหชาตา, นานา กฺขณิกา}. (วิตฺถาเรตพฺพํ.) (3)}

\csromanheader{2}{\textbf{วิปากาทิ}}
\proseitem{44}{อนิทสฺสโน ธมฺโม อนิทสฺสนสฺส ธมฺมสฺส วิปากปจฺจเยน ปจฺจโย.\csromanspacer{} ตีณิ.\csromanspacer{} อาหารปจฺจเยน ปจฺจโย.\csromanspacer{} ตีณิ. (ตีสุปิ กพฬีกาโร อาหาโร กาตพฺโพ.) อินฺทฺริยปจฺจเยน ปจฺจโย.\csromanspacer{} ตีณิ. (ตีสุปิ รูปชีวิตินฺทฺริยํ.) ฌานปจฺจเยน ปจฺจโย.\csromanspacer{} ตีณิ.\csromanspacer{} มคฺคปจฺจเยน ปจฺจโย.\csromanspacer{} ตีณิ.\csromanspacer{} สมฺปยุตฺตปจฺจเยน ปจฺจโย.\csromanspacer{} เอกํ.}

\csromanheader{2}{\textbf{วิปฺปยุตฺต}}
\proseitem{45}{อนิทสฺสโน ธมฺโม อนิทสฺสนสฺส ธมฺมสฺส วิปฺปยุตฺตปจฺจเยน ปจฺจโย.\csromanspacer{} สหชาตํ, ปุเรชาตํ, ปจฺฉาชาตํ.\csromanspacer{}\csromanspacer{}\textbf{สหชาตา} อนิทสฺสนา ขนฺธา อนิทสฺสนานํ จิตฺตสมุฏฺฐานานํ รูปานํ วิปฺปยุตฺตปจฺจเยน ปจฺจโย.\csromanspacer{} ปฏิสนฺธิกฺขเณ อนิทสฺสนา ขนฺธา อนิทสฺสนานํ กฏตฺตารูปานํ วิปฺปยุตฺตปจฺจเยน ปจฺจโย.\csromanspacer{} ขนฺธา วตฺถุสฺส วิปฺปยุตฺตปจฺจเยน ปจฺจโย.\csromanspacer{} วตฺถุ ขนฺธานํ วิปฺปยุตฺตปจฺจเยน ปจฺจโย.\csromanspacer{} \textbf{ปุเรชาตํ} จกฺขายตนํ จกฺขุวิญฺญาณสฺส ฯเปฯ.\csromanspacer{} กายายตนํ กายวิญฺญาณสฺส ฯเปฯ.\csromanspacer{} วตฺถุ อนิทสฺสนานํ ขนฺธานํ วิปฺปยุตฺตปจฺจเยน ปจฺจโย.\csromanspacer{} \textbf{ปจฺฉาชาตา} อนิทสฺสนา ขนฺธา ปุเรชาตสฺส อิมสฺส อนิทสฺสนสฺส กายสฺส วิปฺปยุตฺตปจฺจเยน ปจฺจโย. (1)}

\prose{อนิทสฺสโน ธมฺโม สนิทสฺสนสฺส ธมฺมสฺส วิปฺปยุตฺตปจฺจเยน ปจฺจโย.\csromanspacer{} สหชาตํ, ปจฺฉาชาตํ.\csromanspacer{}\csromanspacer{}\textbf{สหชาตา} อนิทสฺสนา ขนฺธา สนิทสฺสนานํ}

\vspace{\tipitakaparskip}
\csromancenter{สนิทสฺสนทุก จิตฺตสมุฏฺฐานานํ รูปานํ วิปฺปยุตฺตปจฺจเยน ปจฺจโย.\csromanspacer{} ปฏิสนฺธิกฺขเณ ฯเปฯ.\csromanspacer{} \textbf{ปจฺฉาชาตา} อนิทสฺสนา ขนฺธา ปุเรชาตสฺส อิมสฺส สนิทสฺสนสฺส กายสฺส วิปฺปยุตฺตปจฺจเยน ปจฺจโย. (2)}
\prose{\csromanfolio{101}อนิทสฺสโน ธมฺโม สนิทสฺสนสฺส จ อนิทสฺสนสฺส จ ธมฺมสฺส วิปฺปยุตฺตปจฺจเยน ปจฺจโย.\csromanspacer{} สหชาตํ, ปจฺฉาชาตํ.\csromanspacer{}\csromanspacer{}\textbf{สหชาตา} อนิทสฺสนา ขนฺธา สนิทสฺสนานญฺจ อนิทสฺสนานญฺจ จิตฺตสมุฏฺฐานานํ รูปานํ วิปฺปยุตฺตปจฺจเยน ปจฺจโย.\csromanspacer{} ปฏิสนฺธิกฺขเณ ฯเปฯ.\csromanspacer{} \textbf{ปจฺฉาชาตา} อนิทสฺสนา ขนฺธา ปุเรชาตสฺส อิมสฺส สนิทสฺสนสฺส จ อนิทสฺสนสฺส จ กายสฺส วิปฺปยุตฺตปจฺจเยน ปจฺจโย. (3)}

\csromanheader{2}{\textbf{อตฺถิ-อาทิ}}
\proseitem{46}{สนิทสฺสโน ธมฺโม อนิทสฺสนสฺส ธมฺมสฺส อตฺถิปจฺจเยน ปจฺจโย.\csromanspacer{} สนิทสฺสนํ รูปํ อนิจฺจโต ฯเปฯ โทมนสฺสํ อุปฺปชฺชติ.\csromanspacer{} ทิพฺเพน จกฺขุนา รูปํ ปสฺสติ.\csromanspacer{} รูปายตนํ จกฺขุวิญฺญาณสฺส อตฺถิปจฺจเยน ปจฺจโย. (1)}

\prose{อนิทสฺสโน ธมฺโม อนิทสฺสนสฺส ธมฺมสฺส อตฺถิปจฺจเยน ปจฺจโย.\csromanspacer{} สหชาตํ, ปุเรชาตํ, ปจฺฉาชาตํ, อาหารํ, อินฺทฺริยํ.\csromanspacer{}\csromanspacer{}\textbf{สหชาโต} อนิทสฺสโน เอโก ขนฺโธ ติณฺณนฺนํ ขนฺธานํ อนิทสฺสนานญฺจ จิตฺตสมุฏฺฐานานํ รูปานํ อตฺถิปจฺจเยน ปจฺจโย ฯเปฯ เทฺว ขนฺเธ ฯเปฯ. (สํขิตฺตํ.\csromanspacer{} ยาว อสญฺญสตฺตา.) \textbf{ปุเรชาตํ} จกฺขุํ ฯเปฯ วตฺถุํ อนิจฺจโต ฯเปฯ โทมนสฺสํ อุปฺปชฺชติ.\csromanspacer{} ทิพฺพาย โสตธาตุยา สทฺทํ สุณาติ.\csromanspacer{} สทฺทายตนํ โสตวิญฺญาณสฺส ฯเปฯ.\csromanspacer{} โผฏฺฐพฺพายตนํ กายวิญฺญาณสฺส ฯเปฯ.\csromanspacer{} จกฺขายตนํ จกฺขุวิญฺญาณสฺส ฯเปฯ.\csromanspacer{} กายายตนํ กายวิญฺญาณสฺส ฯเปฯ.\csromanspacer{} วตฺถุ อนิทสฺสนานํ ขนฺธานํ อตฺถิปจฺจเยน ปจฺจโย.\csromanspacer{} \textbf{ปจฺฉาชาตา} อนิทสฺสนา ขนฺธา ปุเรชาตสฺส อิมสฺส อนิทสฺสนสฺส กายสฺส อตฺถิปจฺจเยน ปจฺจโย.\csromanspacer{}\csromanspacer{}\textbf{กพฬีกาโร อาหาโร} อิมสฺส อนิทสฺสนสฺส กายสฺส อตฺถิปจฺจเยน ปจฺจโย.\csromanspacer{}\csromanspacer{}\textbf{รูปชีวิตินฺทฺริยํ} อนิทสฺสนานํ กฏตฺตารูปานํ อตฺถิปจฺจเยน ปจฺจโย. (1)}

\prose{อนิทสฺสโน ธมฺโม สนิทสฺสนสฺส ธมฺมสฺส อตฺถิปจฺจเยน ปจฺจโย.\csromanspacer{} สหชาตํ, ปจฺฉาชาตํ, อาหารํ, อินฺทฺริยํ.\csromanspacer{}\csromanspacer{}\textbf{สหชาตา} อนิทสฺสนา ขนฺธา สนิทสฺสนานํ จิตฺตสมุฏฺฐานานํ รูปานํ อตฺถิปจฺจเยน ปจฺจโย.\csromanspacer{} ปฏิสนฺธิกฺขเณ ฯเปฯ.\csromanspacer{} มหาภูตา สนิทสฺสนานํ จิตฺตสมุฏฺฐานานํ รูปานํ, \csromanfolio{102}กฏตฺตารูปานํ, อุปาทารูปานํ อตฺถิปจฺจเยน ปจฺจโย.\csromanspacer{} พาหิรา.\csromanspacer{} อาหารสมุฏฺฐานา.\csromanspacer{} อุตุสมุฏฺฐานา.\csromanspacer{} อสญฺญสตฺตานํ มหาภูตา สนิทสฺสนานํ กฏตฺตารูปานํ, อุปาทารูปานํ อตฺถิปจฺจเยน ปจฺจโย.\csromanspacer{} \textbf{ปจฺฉาชาตา} อนิทสฺสนา ขนฺธา ปุเรชาตสฺส อิมสฺส สนิทสฺสนสฺส กายสฺส อตฺถิปจฺจเยน ปจฺจโย.\csromanspacer{}\csromanspacer{}\textbf{กพฬีกาโร อาหาโร} อิมสฺส สนิทสฺสนสฺส กายสฺส อตฺถิปจฺจเยน ปจฺจโย.\csromanspacer{} \textbf{รูปชีวิตินฺทฺริยํ} สนิทสฺสนานํ กฏตฺตารูปานํ อตฺถิปจฺจเยน ปจฺจโย. (2)}

\prose{อนิทสฺสโน ธมฺโม สนิทสฺสนสฺส จ อนิทสฺสนสฺส จ ธมฺมสฺส อตฺถิปจฺจเยน ปจฺจโย.\csromanspacer{} สหชาตํ, ปจฺฉาชาตํ, อาหารํ, อินฺทฺริยํ.\csromanspacer{}\csromanspacer{}\textbf{สหชาโต} อนิทสฺสโน เอโก ขนฺโธ ติณฺณนฺนํ ขนฺธานํ สนิทสฺสนานญฺจ อนิทสฺสนานญฺจ จิตฺตสมุฏฺฐานานํ รูปานํ อตฺถิปจฺจเยน ปจฺจโย ฯเปฯ เทฺว ขนฺธา ฯเปฯ.\csromanspacer{} ปฏิสนฺธิกฺขเณ ฯเปฯ มหาภูตา สนิทสฺสนานญฺจ อนิทสฺสนานญฺจ จิตฺตสมุฏฺฐานานํ รูปานํ, กฏตฺตารูปานํ, อุปาทารูปานํ อตฺถิปจฺจเยน ปจฺจโย.\csromanspacer{} พาหิรา.\csromanspacer{} อาหารสมุฏฺฐานา.\csromanspacer{} อุตุสมุฏฺฐานา.\csromanspacer{} อสญฺญสตฺตานํ มหาภูตา สนิทสฺสนานญฺจ อนิทสฺสนานญฺจ กฏตฺตารูปานํ, อุปาทารูปานํ อตฺถิปจฺจเยน ปจฺจโย ฯเปฯ. (3)}

\proseitem{47}{สนิทสฺสโน จ อนิทสฺสโน จ ธมฺมา อนิทสฺสนสฺส ธมฺมสฺส อตฺถิปจฺจเยน ปจฺจโย.\csromanspacer{} \textbf{ปุเรชาตํ} รูปายตนญฺจ วตฺถุ จ อนิทสฺสนานํ ขนฺธานํ อตฺถิปจฺจเยน ปจฺจโย.\csromanspacer{} รูปายตนญฺจ จกฺขายตนญฺจ จกฺขุวิญฺญาณสฺส อตฺถิปจฺจเยน ปจฺจโย.}

\prose{นตฺถิปจฺจเยน ปจฺจโย.\csromanspacer{} วิคตปจฺจเยน ปจฺจโย.\csromanspacer{} อวิคตปจฺจเยน ปจฺจโย.}

\csromanheader{2}{1. ปจฺจยานุโลม 2. สงฺขฺยาวาร}
\csromanheader{2}{\textbf{สุทฺธ}}
\proseitem{48}{เหตุยา ตีณิ, อารมฺมเณ เทฺว, อธิปติยา จตฺตาริ, อนนฺตเร เอกํ, สมนนฺตเร เอกํ, สหชาเต ตีณิ, อญฺญมญฺเญ เอกํ, นิสฺสเย ตีณิ, อุปนิสฺสเย เทฺว, ปุเรชาเต ตีณิ, ปจฺฉาชาเต ตีณิ, อาเสวเน เอกํ, กมฺเม ตีณิ, วิปาเก ตีณิ, อาหาเร ตีณิ,}

\vspace{\tipitakaparskip}
\csromancenter{สนิทสฺสนทุก อินฺทฺริเย ตีณิ, ฌาเน ตีณิ, มคฺเค ตีณิ, สมฺปยุตฺเต เอกํ, วิปฺปยุตฺเต ตีณิ, อตฺถิยา ปญฺจ, นตฺถิยา เอกํ, วิคเต เอกํ, อวิคเต ปญฺจ. (เอวํ คเณตพฺพํ.)}
\csromancenter{อนุโลมํ.}
\csromanheader{2}{2. ปจฺจนียุทฺธาร}
\proseitem{49}{\csromanfolio{103}สนิทสฺสโน ธมฺโม อนิทสฺสนสฺส ธมฺมสฺส อารมฺมณปจฺจเยน ปจฺจโย.\csromanspacer{} อุปนิสฺสายปจฺจเยน ปจฺจโย. (1)}

\prose{อนิทสฺสโน ธมฺโม อนิทสฺสนสฺส ธมฺมสฺส อารมฺมณปจฺจเยน ปจฺจโย.\csromanspacer{} สหชาตปจฺจเยน ปจฺจโย.\csromanspacer{} อุปนิสฺสยปจฺจเยน ปจฺจโย.\csromanspacer{} ปุเรชาตปจฺจเยน ปจฺจโย.\csromanspacer{} ปจฺฉาชาตปจฺจเยน ปจฺจโย.\csromanspacer{} กมฺมปจฺจเยน ปจฺจโย.\csromanspacer{} อาหารปจฺจเยน ปจฺจโย.\csromanspacer{} อินฺทฺริยปจฺจเยน ปจฺจโย. (1)}

\prose{อนิทสฺสโน ธมฺโม สนิทสฺสนสฺส ธมฺมสฺส สหชาตปจฺจเยน ปจฺจโย.\csromanspacer{} ปจฺฉาชาตปจฺจเยน ปจฺจโย.\csromanspacer{} กมฺมปจฺจเยน ปจฺจโย.\csromanspacer{} อาหารปจฺจเยน ปจฺจโย.\csromanspacer{} อินฺทฺริยปจฺจเยน ปจฺจโย. (2)}

\prose{อนิทสฺสโน ธมฺโม สนิทสฺสนสฺส จ อนิทสฺสนสฺส จ ธมฺมสฺส สหชาตปจฺจเยน ปจฺจโย.\csromanspacer{} ปจฺฉาชาตปจฺจเยน ปจฺจโย.\csromanspacer{} กมฺมปจฺจเยน ปจฺจโย.\csromanspacer{} อาหารปจฺจเยน ปจฺจโย.\csromanspacer{} อินฺทฺริยปจฺจเยน ปจฺจโย. (3)}

\prose{สนิทสฺสโน จ อนิทสฺสโน จ ธมฺมา อนิทสฺสนสฺส ธมฺมสฺส ปุเรชาตปจฺจเยน ปจฺจโย. (1)}

\csromanheader{2}{2. ปจฺจยปจฺจนีย 2. สงฺขฺยาวาร}
\csromanheader{2}{\textbf{สุทฺธ}}
\proseitem{50}{นเหตุยา ปญฺจ, น-อารมฺมเณ จตฺตาริ, น-อธิปติยา ปญฺจ, น-อนนฺตเร ปญฺจ, นสมนนฺตเร ปญฺจ, นสหชาเต ปญฺจ, น-อญฺญมญฺเญ ปญฺจ, นนิสฺสเย จตฺตาริ, น-อุปนิสฺสเย ปญฺจ, นปุเรชาเต จตฺตาริ, นปจฺฉาชาเต ปญฺจ, (สพฺพตฺถ ปญฺจ.) นสมฺปยุตฺเต ปญฺจ, นวิปฺปยุตฺเต จตฺตาริ, โน-อตฺถิยา จตฺตาริ, โนนตฺถิยา ปญฺจ, โนวิคเต ปญฺจ, โน-อวิคเต จตฺตาริ.}

\vspace{\tipitakaparskip}
\csromancenter{ปจฺจนียํ.}
\csromanheader{2}{3. ปจฺจยานุโลมปจฺจนีย}
\csromanheader{2}{\textbf{เหตุทุก}}
\proseitem{51}{\csromanfolio{104}\textbf{เหตุ}ปจฺจยา น-อารมฺมเณ ตีณิ, น-อธิปติยา ตีณิ, น-อนนฺตเร ตีณิ, นสมนนฺตเร ตีณิ, น-อญฺญมญฺเญ ตีณิ, น-อุปนิสฺสเย ตีณิ ฯเปฯ นสมฺปยุตฺเต ตีณิ, นวิปฺปยุตฺเต เอกํ, โนนตฺถิยา ตีณิ, โนวิคเต ตีณิ.}

\vspace{\tipitakaparskip}
\csromancenter{อนุโลมปจฺจนียํ.}
\csromanheader{2}{4. ปจฺจปจฺจนียานุโลม}
\csromanheader{2}{\textbf{นเหตุทุก}}
\proseitem{52}{นเหตุปจฺจยา อารมฺมเณ เทฺว, อธิปติยา จตฺตาริ, อนนฺตเร เอกํ, สมนนฺตเร เอกํ, สหชาเต ตีณิ, อญฺญมญฺเญ เอกํ, นิสฺสเย ตีณิ, อุปนิสฺสเย เทฺว, ปุเรชาเต ตีณิ, ปจฺฉาชาเต ตีณิ, อาเสวเน เอกํ, กมฺเม ตีณิ ฯเปฯ มคฺเค ตีณิ, สมฺปยุตฺเต เอกํ, วิปฺปยุตฺเต ตีณิ, อตฺถิยา ปญฺจ, นตฺถิยา เอกํ, วิคเต เอกํ, อวิคเต ปญฺจ.}

\vspace{\tipitakaparskip}
\csromancenter{ปจฺจนียานุโลมํ.}
\nitthitam{สนิทสฺสนทุกํ นิฏฺฐิตํ.}
\csromanmatikaanchor{mtk.13}
\csromantocmarkref{section}{10. สปฺปฏิฆทุก}{mtk.13}
\bookmark[dest={mtk.13},level=1]{10. สปฺปฏิฆทุก}
\csromanheader{1}{10. สปฺปฏิฆทุก 1. ปฏิจฺจวาร}
\csromanheader{2}{1. ปจฺจยานุโลม 1. วิภงฺควาร}
\csromanheader{2}{\textbf{เหตุ}}
\proseitem{53}{สปฺปฏิฆํ ธมฺมํ ปฏิจฺจ สปฺปฏิโฆ ธมฺโม อุปฺปชฺชติ เหตุปจฺจยา.\csromanspacer{} สปฺปฏิฆํ เอกํ มหาภูตํ ปฏิจฺจ เทฺว มหาภูตา, เทฺว มหาภูเต ปฏิจฺจ เอกํ}

\vspace{\tipitakaparskip}
\csromancenter{สปฺปฏิฆทุก มหาภูตํ.\csromanspacer{} สปฺปฏิเฆ มหาภูเต ปฏิจฺจ สปฺปฏิฆํ จิตฺตสมุฏฺฐานํ รูปํ, กฏตฺตารูปํ, อุปาทารูปํ.\csromanspacer{} โผฏฺฐพฺพายตนํ ปฏิจฺจ จกฺขายตนํ ฯเปฯ รสายตนํ. (1)}
\prose{\csromanfolio{105}สปฺปฏิฆํ ธมฺมํ ปฏิจฺจ อปฺปฏิโฆ ธมฺโม อุปฺปชฺชติ เหตุปจฺจยา.\csromanspacer{} สปฺปฏิเฆ มหาภูเต ปฏิจฺจ อาโปธาตุ.\csromanspacer{} สปฺปฏิเฆ มหาภูเต ปฏิจฺจ อปฺปฏิฆํ จิตฺตสมุฏฺฐานํ รูปํ, กฏตฺตารูปํ, อุปาทารูปํ.\csromanspacer{} โผฏฺฐพฺพายตนํ ปฏิจฺจ อาโปธาตุ.\csromanspacer{} อิตฺถินฺทฺริยํ ฯเปฯ กพฬีกาโร อาหาโร. (2)}

\prose{สปฺปฏิฆํ ธมฺมํ ปฏิจฺจ สปฺปฏิโฆ จ อปฺปฏิโฆ จ ธมฺมา อุปฺปชฺชนฺติ เหตุปจฺจยา.\csromanspacer{} สปฺปฏิฆํ เอกํ มหาภูตํ ปฏิจฺจ เทฺว มหาภูตา อาโปธาตุ จ, เทฺว มหาภูเต ฯเปฯ สปฺปฏิเฆ มหาภูเต ปฏิจฺจ สปฺปฏิฆญฺจ อปฺปฏิฆญฺจ จิตฺตสมุฏฺฐานํ รูปํ, กฏตฺตารูปํ, อุปาทารูปํ.\csromanspacer{} โผฏฺฐพฺพายตนํ ปฏิจฺจ จกฺขายตนํ ฯเปฯ รสายตนํ อาโปธาตุ อิตฺถินฺทฺริยํ ฯเปฯ กพฬีกาโร อาหาโร. (3)}

\proseitem{54}{อปฺปฏิฆํ ธมฺมํ ปฏิจฺจ อปฺปฏิโฆ ธมฺโม อุปฺปชฺชติ เหตุปจฺจยา.\csromanspacer{} อปฺปฏิฆํ เอกํ ขนฺธํ ปฏิจฺจ ตโย ขนฺธา อปฺปฏิฆํ จิตฺตสมุฏฺฐานญฺจ รูปํ ฯเปฯ เทฺว ขนฺเธ ฯเปฯ.\csromanspacer{} ปฏิสนฺธิกฺขเณ ฯเปฯ.\csromanspacer{} ขนฺเธ ปฏิจฺจ วตฺถุ, วตฺถุํ ปฏิจฺจ ขนฺธา.\csromanspacer{} อาโปธาตุํ ปฏิจฺจ อปฺปฏิฆํ จิตฺตสมุฏฺฐานํ รูปํ, กฏตฺตารูปํ, อุปาทารูปํ.\csromanspacer{} อาโปธาตุํ ปฏิจฺจ อิตฺถินฺทฺริยํ ฯเปฯ กพฬีกาโร อาหาโร. (1)}

\prose{อปฺปฏิฆํ ธมฺมํ ปฏิจฺจ สปฺปฏิโฆ ธมฺโม อุปฺปชฺชติ เหตุปจฺจยา.\csromanspacer{} อปฺปฏิเฆ ขนฺเธ ปฏิจฺจ สปฺปฏิฆํ จิตฺตสมุฏฺฐานํ รูปํ.\csromanspacer{} ปฏิสนฺธิกฺขเณ ฯเปฯ.\csromanspacer{} อาโปธาตุํ ปฏิจฺจ สปฺปฏิฆา มหาภูตา.\csromanspacer{} อาโปธาตุํ ปฏิจฺจ สปฺปฏิฆํ จิตฺตสมุฏฺฐานํ รูปํ, กฏตฺตารูปํ, อุปาทารูปํ.\csromanspacer{} อาโปธาตุํ ปฏิจฺจ จกฺขายตนํ ฯเปฯ โผฏฺฐพฺพายตนํ.}

\prose{อปฺปฏิฆํ ธมฺมํ ปฏิจฺจ สปฺปฏิโฆ จ อปฺปฏิโฆ จ ธมฺมา อุปฺปชฺชนฺติ เหตุปจฺจยา.\csromanspacer{} อปฺปฏิฆํ เอกํ ขนฺธํ ปฏิจฺจ ตโย ขนฺธา สปฺปฏิฆญฺจ อปฺปฏิฆญฺจ จิตฺตสมุฏฺฐานํ รูปํ ฯเปฯ เทฺว ขนฺเธ ฯเปฯ.\csromanspacer{} ปฏิสนฺธิกฺขเณ ฯเปฯ.\csromanspacer{} อาโปธาตุํ ปฏิจฺจ สปฺปฏิฆญฺจ อปฺปฏิฆญฺจ จิตฺตสมุฏฺฐานํ รูปํ, กฏตฺตารูปํ, อุปาทารูปํ.\csromanspacer{} อาโปธาตุํ ปฏิจฺจ จกฺขายตนํ ฯเปฯ โผฏฺฐพฺพายตนํ, อิตฺถินฺทฺริยํ ฯเปฯ กพฬีกาโร อาหาโร. (3)}

\proseitem{55}{\csromanfolio{106}สปฺปฏิฆญฺจ อปฺปฏิฆญฺจ ธมฺมํ ปฏิจฺจ สปฺปฏิโฆ ธมฺโม อุปฺปชฺชติ เหตุปจฺจยา.\csromanspacer{} อปฺปฏิเฆ ขนฺเธ จ มหาภูเต จ ปฏิจฺจ สปฺปฏิฆํ จิตฺตสมุฏฺฐานํ รูปํ.\csromanspacer{} ปฏิสนฺธิกฺขเณ ฯเปฯ.\csromanspacer{} สปฺปฏิฆํ เอกํ มหาภูตญฺจ อาโปธาตุญฺจ ปฏิจฺจ เทฺว มหาภูตา ฯเปฯ สปฺปฏิเฆ มหาภูเต จ อาโปธาตุญฺจ ปฏิจฺจ สปฺปฏิฆํ จิตฺตสมุฏฺฐานํ รูปํ, กฏตฺตารูปํ, อุปาทารูปํ.\csromanspacer{} โผฏฺฐพฺพายตนญฺจ อาโปธาตุญฺจ ปฏิจฺจ จกฺขายตนํ ฯเปฯ รสายตนํ. (1)}

\prose{สปฺปฏิฆญฺจ อปฺปฏิฆญฺจ ธมฺมํ ปฏิจฺจ อปฺปฏิโฆ ธมฺโม อุปฺปชฺชติ เหตุปจฺจยา.\csromanspacer{} อปฺปฏิเฆ ขนฺเธ จ มหาภูเต จ ปฏิจฺจ อปฺปฏิฆํ จิตฺตสมุฏฺฐานํ รูปํ.\csromanspacer{} ปฏิสนฺธิกฺขเณ อปฺปฏิเฆ ขนฺเธ จ มหาภูเต จ ปฏิจฺจ อปฺปฏิฆํ กฏตฺตารูปํ.\csromanspacer{} โผฏฺฐพฺพายตนญฺจ อาโปธาตุญฺจ ปฏิจฺจ อปฺปฏิฆํ จิตฺตสมุฏฺฐานํ รูปํ, กฏตฺตารูปํ, อุปาทารูปํ.\csromanspacer{} โผฏฺฐพฺพายตนญฺจ อาโปธาตุญฺจ ปฏิจฺจ อิตฺถินฺทฺริยํ ฯเปฯ กพฬีกาโร อาหาโร. (2)}

\prose{สปฺปฏิฆญฺจ อปฺปฏิฆญฺจ ธมฺมํ ปฏิจฺจ สปฺปฏิโฆ จ อปฺปฏิโฆ จ ธมฺมา อุปฺปชฺชนฺติ เหตุปจฺจยา.\csromanspacer{} อปฺปฏิเฆ ขนฺเธ จ มหาภูเต จ ปฏิจฺจ สปฺปฏิฆญฺจ อปฺปฏิฆญฺจ จิตฺตสมุฏฺฐานํ รูปํ.\csromanspacer{} ปฏิสนฺธิกฺขเณ ฯเปฯ.\csromanspacer{} อปฺปฏิเฆ ขนฺเธ จ มหาภูเต จ ปฏิจฺจ สปฺปฏิฆญฺจ อปฺปฏิฆญฺจ กฏตฺตารูปํ.\csromanspacer{} โผฏฺฐพฺพายตนญฺจ อาโปธาตุญฺจ ปฏิจฺจ สปฺปฏิฆญฺจ อปฺปฏิฆญฺจ จิตฺตสมุฏฺฐานํ รูปํ, กฏตฺตารูปํ, อุปาทารูปํ.\csromanspacer{} โผฏฺฐพฺพายตนญฺจ อาโปธาตุญฺจ ปฏิจฺจ จกฺขายตนํ ฯเปฯ รสายตนํ, อิตฺถินฺทฺริยํ ฯเปฯ กพฬีกาโร อาหาโร. (3)}

\csromanheader{2}{\textbf{อารมฺมณ}}
\proseitem{56}{อปฺปฏิฆํ ธมฺมํ ปฏิจฺจ อปฺปฏิโฆ ธมฺโม อุปฺปชฺชติ อารมฺมณปจฺจยา.\csromanspacer{} อปฺปฏิฆํ เอกํ ขนฺธํ ปฏิจฺจ ตโย ขนฺธา ฯเปฯ เทฺว ขนฺเธ ฯเปฯ.\csromanspacer{} ปฏิสนฺธิกฺขเณ ฯเปฯ.\csromanspacer{} วตฺถุํ ปฏิจฺจ ขนฺธา.}

\csromanheader{2}{\textbf{อธิปตฺยาทิ}}
\proseitem{57}{สปฺปฏิฆํ ธมฺมํ ปฏิจฺจ อปฺปฏิโฆ ธมฺโม อุปฺปชฺชติ อธิปติปจฺจยา. (ปฏิสนฺธิ วชฺเชตพฺพา, กฏตฺตารูปา จ.) อนนฺตรปจฺจยา.\csromanspacer{} สมนนฺตรปจฺจยา.\csromanspacer{} สหชาตปจฺจยา. (สพฺเพ มหาภูตา กาตพฺพา.) อญฺญมญฺญปจฺจยา.\csromanspacer{} สปฺปฏิฆํ เอกํ มหาภูตํ ปฏิจฺจ เทฺว มหาภูตา.\csromanspacer{} เทฺว มหาภูเต ฯเปฯ. (1)}

\csromanheader{2}{10. สปฺปฏิฆทุก}
\prose{\csromanfolio{107}สปฺปฏิฆํ ธมฺมํ ปฏิจฺจ อปฺปฏิโฆ ธมฺโม อุปฺปชฺชติ อญฺญมญฺญปจฺจยา.\csromanspacer{} สปฺปฏิเฆ มหาภูเต ปฏิจฺจ อาโปธาตุ ฯเปฯ. (2)}

\prose{สปฺปฏิฆํ ธมฺมํ ปฏิจฺจ สปฺปฏิโฆ จ อปฺปฏิโฆ จ ธมฺมา อุปฺปชฺชนฺติ อญฺญมญฺญปจฺจยา.\csromanspacer{} สปฺปฏิฆํ เอกํ มหาภูตํ ปฏิจฺจ เทฺว มหาภูตา อาโปธาตุ จ.\csromanspacer{} เทฺว มหาภูเต ฯเปฯ. (3)}

\proseitem{58}{อปฺปฏิฆํ ธมฺมํ ปฏิจฺจ อปฺปฏิโฆ ธมฺโม อุปฺปชฺชติ อญฺญมญฺญปจฺจยา.\csromanspacer{} อปฺปฏิฆํ เอกํ ขนฺธํ ปฏิจฺจ ตโย ขนฺธา ฯเปฯ เทฺว ขนฺเธ ฯเปฯ.\csromanspacer{} ปฏิสนฺธิกฺขเณ ฯเปฯ.\csromanspacer{} ขนฺเธ ปฏิจฺจ วตฺถุ, วตฺถุํ ปฏิจฺจ ขนฺธา. (1)}

\prose{อปฺปฏิฆํ ธมฺมํ ปฏิจฺจ สปฺปฏิโฆ ธมฺโม อุปฺปชฺชติ อญฺญมญฺญปจฺจยา.\csromanspacer{} อาโปธาตุํ ปฏิจฺจ สปฺปฏิฆา มหาภูตา. (อิเม อชฺฌตฺติกพาหิรา มหาภูตา กาตพฺพา.) (2)}

\prose{สปฺปฏิฆญฺจ อปฺปฏิฆญฺจ ธมฺมํ ปฏิจฺจ สปฺปฏิโฆ ธมฺโม อุปฺปชฺชติ อญฺญมญฺญปจฺจยา.\csromanspacer{} สปฺปฏิฆํ เอกํ มหาภูตญฺจ อาโปธาตุญฺจ ปฏิจฺจ เทฺว มหาภูตา ฯเปฯ.\csromanspacer{} นิสฺสยปจฺจยา ฯเปฯ อวิคตปจฺจยา.}

\csromanheader{2}{1. ปจฺจยานุโลม 2. สงฺขฺยาวาร}
\csromanheader{2}{\textbf{สุทฺธ}}
\proseitem{59}{เหตุยา นว, อารมฺมเณ เอกํ, อธิปติยา นว, อนนฺตเร เอกํ, สมนนฺตเร เอกํ, สหชาเต นว, อญฺญมญฺเญ ฉ, นิสฺสเย นว, อุปนิสฺสเย เอกํ, ปุเรชาเต เอกํ, อาเสวเน เอกํ, กมฺเม นว, วิปาเก นว, อาหาเร นว, อินฺทฺริเย นว, ฌาเน นว, มคฺเค นว, สมฺปยุตฺเต เอกํ, วิปฺปยุตฺเต นว, อตฺถิยา นว, นตฺถิยา เอกํ, วิคเต เอกํ, อวิคเต นว.}

\vspace{\tipitakaparskip}
\csromancenter{อนุโลมํ.}
\csromanheader{2}{2. ปจฺจยปจฺจนีย 1. วิภงฺควาร}
\csromanheader{2}{\textbf{นเหตุ}}
\proseitem{60}{\csromanfolio{108}สปฺปฏิฆํ ธมฺมํ ปฏิจฺจ สปฺปฏิโฆ ธมฺโม อุปฺปชฺชติ นเหตุปจฺจยา ตีณิ.}

\prose{อปฺปฏิฆํ ธมฺมํ ปฏิจฺจ อปฺปฏิโฆ ธมฺโม อุปฺปชฺชติ นเหตุปจฺจยา.\csromanspacer{} อเหตุกํ อปฺปฏิฆํ เอกํ ขนฺธํ ปฏิจฺจ ตโย ขนฺธา อปฺปฏิฆํ จิตฺตสมุฏฺฐานญฺจ รูปํ ฯเปฯ เทฺว ขนฺเธ ฯเปฯ.\csromanspacer{} อเหตุกปฏิสนฺธิกฺขเณ ฯเปฯ.\csromanspacer{} ขนฺเธ ปฏิจฺจ วตฺถุ, วตฺถุํ ปฏิจฺจ ขนฺธา.\csromanspacer{} อาโปธาตุํ ปฏิจฺจ อปฺปฏิฆํ จิตฺตสมุฏฺฐานํ รูปํ, กฏตฺตารูปํ, อุปาทารูปํ.\csromanspacer{} อาโปธาตุํ ปฏิจฺจ อิตฺถินฺทฺริยํ ฯเปฯ กพฬีกาโร อาหาโร.\csromanspacer{} พาหิรํ.\csromanspacer{} อาหารสมุฏฺฐานํ.\csromanspacer{} อุตุสมุฏฺฐานํ.\csromanspacer{} อสญฺญสตฺตานํ.\csromanspacer{} อาโปธาตุํ ปฏิจฺจ อปฺปฏิฆํ กฏตฺตารูปํ, อุปาทารูปํ.\csromanspacer{} วิจิกิจฺฉาสหคเต อุทฺธจฺจสหคเต ขนฺเธ ปฏิจฺจ วิจิกิจฺฉาสหคโต อุทฺธจฺจสหคโต โมโห. (1)}

\prose{(อปฺปฏิฆมูลกํ อิตเรปิ เทฺว ปญฺหา กาตพฺพา.\csromanspacer{} ฆฏเนปิ ตีณิ ปญฺหา กาตพฺพา.\csromanspacer{} อชฺฌตฺติกา พาหิรา มหาภูตา สพฺเพ ชานิตฺวา กาตพฺพา.)}

\csromanheader{2}{\textbf{น-อารมฺมณาทิ}}
\proseitem{61}{สปฺปฏิฆํ ธมฺมํ ปฏิจฺจ สปฺปฏิโฆ ธมฺโม อุปฺปชฺชติ น-อารมฺมณปจฺจยา. (สพฺพํ ขิตฺตํ.) โนวิคตปจฺจยา.}

\csromanheader{2}{2. ปจฺจยปจฺจนีย 2. สงฺขฺยาวาร}
\csromanheader{2}{\textbf{สุทฺธ}}
\vspace{\tipitakaparskip}
\csromancenter{\textbf{นเหตุ}ยา นว, น-อารมฺมเณ นว, น-อธิปติยา นว, น-อนนฺตเร นว, นสมนนฺตเร นว, น-อญฺญมญฺเญ นว, น-อุปนิสฺสเย นว, นปุเรชาเต นว, นปจฺฉาชาเต นว, น-อาเสวเน นว, นกมฺเม นว, นวิปาเก นว, น-อาหาเร นว, น-อินฺทฺริเย นว, นฌาเน นว, นมคฺเค นว, นสมฺปยุตฺเต นว, นวิปฺปยุตฺเต นว, โนนตฺถิยา นว, โนวิคเต นว.}
\csromancenter{ปจฺจนียํ.}
\csromanheader{2}{10. สปฺปฏิฆทุก \textbf{3. ปจฺจยานุโลมปจฺจนีย}}
\csromanheader{2}{\textbf{เหตุทุก}}
\proseitem{63}{\csromanfolio{109}\textbf{เหตุ}ปจฺจยา น-อารมฺมเณ นว, น-อธิปติยา นว, น-อนนฺตเร นว, นสมนนฺตเร นว, น-อญฺญมญฺเญ นว, น-อุปนิสฺสเย นว, นปุเรชาเต นว, นปจฺฉาชาเต นว, น-อาเสวเน นว, นกมฺเม เอกํ, นวิปาเก นว, นสมฺปยุตฺเต นว, นวิปฺปยุตฺเต เอกํ, โนนตฺถิยา นว, โนวิคเต นว.}

\vspace{\tipitakaparskip}
\csromancenter{อนุโลมปจฺจนียํ.}
\csromanheader{2}{4. ปจฺจยปจฺจนียานุโลม}
\csromanheader{2}{\textbf{นเหตุทุก}}
\proseitem{64}{นเหตุปจฺจยา อารมฺมเณ เอกํ, อนนฺตเร เอกํ, สมนนฺตเร เอกํ, สหชาเต นว, อญฺญมญฺเญ ฉ, นิสฺสเย นว, อุปนิสฺสเย เอกํ, ปุเรชาเต เอกํ, อาเสวเน เอกํ, กมฺเม นว ฯเปฯ มคฺเค เอกํ, สมฺปยุตฺเต เอกํ, วิปฺปยุตฺเต นว, อตฺถิยา นว, นตฺถิยา เอกํ, วิคเต เอกํ, อวิคเต นว.}

\vspace{\tipitakaparskip}
\csromancenter{ปจฺจนียานุโลมํ.}
\csromancenter{(สหชาตวาโรปิ ปฏิจฺจวารสทิโส.)}
\csromanheader{2}{10. สปฺปฏิฆทุก 3. ปจฺจยวาร}
\csromanheader{2}{1. ปจฺจยานุโลม 1. วิภงฺควาร}
\csromanheader{2}{\textbf{เหตุ}}
\proseitem{65}{สปฺปฏิฆํ ธมฺมํ ปจฺจยา สปฺปฏิโฆ ธมฺโม อุปฺปชฺชติ เหตุปจฺจยา.\csromanspacer{} ตีณิ.}

\prose{\csromanfolio{110}อปฺปฏิฆํ ธมฺมํ ปจฺจยา อปฺปฏิโฆ ธมฺโม อุปฺปชฺชติ เหตุปจฺจยา.\csromanspacer{} อปฺปฏิฆํ เอกํ ขนฺธํ ปจฺจยา ตโย ขนฺธา อปฺปฏิฆํ จิตฺตสมุฏฺฐานญฺจ รูปํ ฯเปฯ เทฺว ขนฺเธ ฯเปฯ.\csromanspacer{} ปฏิสนฺธิกฺขเณ ฯเปฯ.\csromanspacer{} อาโปธาตุํ ปจฺจยา อปฺปฏิฆํ จิตฺตสมุฏฺฐานํ รูปํ กฏตฺตารูปํ.\csromanspacer{} อุปาทารูปํ.\csromanspacer{} อาโปธาตุํ ปจฺจยา อิตฺถินฺทฺริยํ ฯเปฯ กพฬีกาโร อาหาโร.\csromanspacer{} วตฺถุํ ปจฺจยา อปฺปฏิฆา ขนฺธา. (1) (อวเสสา ปญฺจ ปญฺหา ปฏิจฺจวารสทิสา.)}

\csromanheader{2}{\textbf{อารมฺมณาทิ}}
\proseitem{66}{สปฺปฏิฆํ ธมฺมํ ปจฺจยา อปฺปฏิโฆ ธมฺโม อุปฺปชฺชติ อารมฺมณปจฺจยา.\csromanspacer{} จกฺขายตนํ ปจฺจยา จกฺขุวิญฺญาณํ ฯเปฯ กายายตนํ ปจฺจยา กายวิญฺญาณํ. (1)}

\prose{อปฺปฏิฆํ ธมฺมํ ปจฺจยา อปฺปฏิโฆ ธมฺโม อุปฺปชฺชติ อารมฺมณปจฺจยา.\csromanspacer{} อปฺปฏิฆํ เอกํ ขนฺธํ ปจฺจยา ตโย ขนฺธา ฯเปฯ เทฺว ขนฺเธ ฯเปฯ.\csromanspacer{} ปฏิสนฺธิกฺขเณ ฯเปฯ.\csromanspacer{} วตฺถุํ ปจฺจยา อปฺปฏิฆา ขนฺธา. (1)}

\prose{สปฺปฏิฆญฺจ อปฺปฏิฆญฺจ ธมฺมํ ปจฺจยา อปฺปฏิโฆ ธมฺโม อุปฺปชฺชติ อารมฺมณปจฺจยา.\csromanspacer{} จกฺขุวิญฺญาณสหคตํ เอกํ ขนฺธญฺจ จกฺขายตนญฺจ ปจฺจยา ตโย ขนฺธา ฯเปฯ.\csromanspacer{} กายวิญฺญาณสหคตํ เอกํ ขนฺธญฺจ กายายตนญฺจ ปจฺจยา ตโย ขนฺธา ฯเปฯ.\csromanspacer{} อธิปติปจฺจยา. (สํขิตฺตํ.) อวิคตปจฺจยา.}

\csromanheader{2}{1. ปจฺจยานุโลม 2. สงฺขฺยาวาร}
\csromanheader{2}{\textbf{สุทฺธ}}
\proseitem{67}{เหตุยา นว, อารมฺมเณ ตีณิ, อธิปติยา นว, อนนฺตเร ตีณิ, สมนนฺตเร ตีณิ, สหชาเต นว, อญฺญมญฺเญ ฉ, นิสฺสเย นว, อุปนิสฺสเย ตีณิ, ปุเรชาเต ตีณิ, อาเสวเน เอกํ, กมฺเม นว ฯเปฯ อวิคเต นว. (เอวํ ปจฺจนียคณนาปิ กาตพฺพา.) (นิสฺสยวาโรปิ ปจฺจยวารสทิโส.)}

\prose{(สํสฏฺฐวาเรปิ สพฺพตฺถ เอกํ.\csromanspacer{} สํขิตฺตํ.\csromanspacer{} อวิคตปจฺจยา, เอกาเยว ปญฺหา.\csromanspacer{} เทฺวปิ วารา กาตพฺพา.)}

\csromanheader{2}{10. สปฺปฏิฆทุก \textbf{10. สปฺปฏิฆทุก 7. ปญฺหาวาร}}
\csromanheader{2}{1. ปจฺจยานุโลม 1. วิภงฺควาร}
\csromanheader{2}{\textbf{เหตุ}}
\proseitem{68}{\csromanfolio{111}อปฺปฏิโฆ ธมฺโม อปฺปฏิฆสฺส ธมฺมสฺส เหตุปจฺจเยน ปจฺจโย.\csromanspacer{} อปฺปฏิฆา เหตู สมฺปยุตฺตกานํ ขนฺธานํ อปฺปฏิฆานญฺจ จิตฺตสมุฏฺฐานานํ รูปานํ เหตุปจฺจเยน ปจฺจโย.\csromanspacer{} ปฏิสนฺธิกฺขเณ ฯเปฯ. (1)}

\prose{อปฺปฏิโฆ ธมฺโม สปฺปฏิฆสฺส ธมฺมสฺส เหตุปจฺจเยน ปจฺจโย.\csromanspacer{} อปฺปฏิฆา เหตู สปฺปฏิฆานํ จิตฺตสมุฏฺฐานานํ รูปานํ เหตุปจฺจเยน ปจฺจโย.\csromanspacer{} ปฏิสนฺธิกฺขเณ ฯเปฯ. (2)}

\prose{อปฺปฏิโฆ ธมฺโม สปฺปฏิฆสฺส จ อปฺปฏิฆสฺส จ ธมฺมสฺส เหตุปจฺจเยน ปจฺจโย.\csromanspacer{} อปฺปฏิฆา เหตู สมฺปยุตฺตกานํ ขนฺธานํ สปฺปฏิฆานญฺจ อปฺปฏิฆานญฺจ จิตฺตสมุฏฺฐานานํ รูปานํ เหตุปจฺจเยน ปจฺจโย.\csromanspacer{} ปฏิสนฺธิกฺขเณ ฯเปฯ. (3)}

\csromanheader{2}{\textbf{อารมฺมณ}}
\proseitem{69}{สปฺปฏิโฆ ธมฺโม อปฺปฏิฆสฺส ธมฺมสฺส อารมฺมณปจฺจเยน ปจฺจโย.\csromanspacer{} จกฺขุํ ฯเปฯ.\csromanspacer{} โผฏฺฐพฺเพ อนิจฺจโต ฯเปฯ โทมนสฺสํ อุปฺปชฺชติ.\csromanspacer{} ทิพฺเพน จกฺขุนา รูปํ ปสฺสติ.\csromanspacer{} ทิพฺพาย โสตธาตุยา สทฺทํ สุณาติ.\csromanspacer{} รูปายตนํ จกฺขุวิญฺญาณสฺส ฯเปฯ.\csromanspacer{} โผฏฺฐพฺพายตนํ กายวิญฺญาณสฺส ฯเปฯ.\csromanspacer{} สปฺปฏิฆา ขนฺธา อิทฺธิวิธญาณสฺส, ปุพฺเพนิวาสานุสฺสติญาณสฺส, อนาคตํสญาณสฺส, อาวชฺชนาย อารมฺมณปจฺจเยน ปจฺจโย. (1)}

\prose{อปฺปฏิโฆ ธมฺโม อปฺปฏิฆสฺส ธมฺมสฺส อารมฺมณปจฺจเยน ปจฺจโย.\csromanspacer{} ทานํ ฯเปฯ สีลํ ฯเปฯ อุโปสถกมฺมํ กตฺวา ตํ ปจฺจเวกฺขติ.\csromanspacer{} ปุพฺเพ สุจิณฺณานิ ปจฺจเวกฺขติ.\csromanspacer{} ฌานา วุฏฺฐหิตฺวา ฌานํ ปจฺจเวกฺขติ.\csromanspacer{} อริยา มคฺคา วุฏฺฐหิตฺวา มคฺคํ ปจฺจเวกฺขนฺติ, ผลํ ปจฺจเวกฺขนฺติ, นิพฺพานํ ปจฺจเวกฺขนฺติ.\csromanspacer{} นิพฺพานํ โคตฺรภุสฺส, โวทานสฺส, มคฺคสฺส, ผลสฺส, อาวชฺชนาย อารมฺมณปจฺจเยน ปจฺจโย.\csromanspacer{} อริยา ปหีเน กิเลเส ปจฺจเวกฺขนฺติ, วิกฺขมฺภิเต กิเลเส ปจฺจเวกฺขนฺติ.\csromanspacer{} ปุพฺเพ สมุทาจิณฺเณ กิเลเส ชานนฺติ.\csromanspacer{} วตฺถุํ ฯเปฯ.\csromanspacer{} อิตฺถินฺทฺริยํ.\csromanspacer{} ปุริสินฺทฺริยํ.\csromanspacer{} ชีวิตินฺทฺริยํ.\csromanspacer{} อาโปธาตุํ กพฬีการํ อาหารํ อนิจฺจโต ฯเปฯ โทมนสฺสํ อุปฺปชฺชติ.\csromanspacer{} เจโตปริยญาเณน อปฺปฏิฆจิตฺตสมงฺคิสฺส จิตฺตํ ชานาติ. \csromanfolio{112}อากาสานญฺจายตนํ วิญฺญาณญฺจายตนสฺส ฯเปฯ.\csromanspacer{} อากิญฺจญฺญายตนํ เนวสญฺญานาสญฺญายตนสฺส ฯเปฯ.\csromanspacer{} อปฺปฏิฆา ขนฺธา อิทฺธิวิธญาณสฺส, เจโตปริยญาณสฺส, ปุพฺเพนิวาสานุสฺสติญาณสฺส ยถากมฺมูปคญาณสฺส, อนาคตํสญาณสฺส, อาวชฺชนาย อารมฺมณปจฺจเยน ปจฺจโย. (1)}

\csromanheader{2}{\textbf{อธิปติ}}
\proseitem{70}{สปฺปฏิโฆ ธมฺโม อปฺปฏิฆสฺส ธมฺมสฺส อธิปติปจฺจเยน ปจฺจโย.\csromanspacer{} \textbf{อารมฺมณาธิปติ\csromandash{}}จกฺขุํ ฯเปฯ โผฏฺฐพฺเพ ครุํ กตฺวา อสฺสาเทติ อภินนฺทติ, ตํ ครุํ กตฺวา ราโค อุปฺปชฺชติ, ทิฏฺฐิ อุปฺปชฺชติ. (1)}

\prose{อปฺปฏิโฆ ธมฺโม อปฺปฏิฆสฺส ธมฺมสฺส อธิปติปจฺจเยน ปจฺจโย.\csromanspacer{} \textbf{อารมฺมณาธิปติ}, สหชาตาธิปติ.\csromanspacer{}\csromanspacer{}\textbf{อารมฺมณาธิปติ\csromandash{}}ทานํ ฯเปฯ สีลํ ฯเปฯ อุโปสถกมฺมํ กตฺวา ตํ ครุํ กตฺวา ปจฺจเวกฺขติ.\csromanspacer{} ปุพฺเพ สุจิณฺณานิ ครุํ กตฺวา ปจฺจเวกฺขติ.\csromanspacer{} ฌานา วุฏฺฐหิตฺวา ฌานํ ครุํ กตฺวา ปจฺจเวกฺขติ.\csromanspacer{} อริยา มคฺคา วุฏฺฐหิตฺวา มคฺคํ ครุํ กตฺวา ปจฺจเวกฺขนฺติ, ผลํ ครุํ กตฺวา ปจฺจเวกฺขนฺติ, นิพฺพานํ ครุํ กตฺวา ปจฺจเวกฺขนฺติ.\csromanspacer{} นิพฺพานํ โคตฺรภุสฺส, โวทานสฺส, มคฺคสฺส, ผลสฺส อธิปติปจฺจเยน ปจฺจโย.\csromanspacer{} วตฺถุํ ฯเปฯ.\csromanspacer{} อิตฺถินฺทฺริยํ.\csromanspacer{} ปุริสินฺทฺริยํ.\csromanspacer{} ชีวิตินฺทฺริยํ.\csromanspacer{} อาโปธาตุํ, กพฬีการํ อาหารํ ครุํ กตฺวา อสฺสาเทติ อภินนฺทติ, ตํ ครุํ กตฺวา ราโค อุปฺปชฺชติ, ทิฏฺฐิ อุปฺปชฺชติ.\csromanspacer{} \textbf{สหชาตาธิปติ\csromandash{}}อปฺปฏิฆาธิปติ สมฺปยุตฺตกานํ ขนฺธานํ อปฺปฏิฆานญฺจ จิตฺตสมุฏฺฐานานํ รูปานํ อธิปติปจฺจเยน ปจฺจโย. (1)}

\prose{อปฺปฏิโฆ ธมฺโม สปฺปฏิฆสฺส ธมฺมสฺส อธิปติปจฺจเยน ปจฺจโย.\csromanspacer{} \textbf{สหชาตาธิปติ\csromandash{}}อปฺปฏิฆาธิปติ สปฺปฏิฆานํ จิตฺตสมุฏฺฐานานํ รูปานํ อธิปติปจฺจเยน ปจฺจโย. (2)}

\prose{อปฺปฏิโฆ ธมฺโม สปฺปฏิฆสฺส จ อปฺปฏิฆสฺส จ ธมฺมสฺส อธิปติปจฺจเยน ปจฺจโย.\csromanspacer{} \textbf{สหชาตาธิปติ\csromandash{}}อปฺปฏิฆาธิปติ สมฺปยุตฺตกานํ ขนฺธานํ สปฺปฏิฆานญฺจ อปฺปฏิฆานญฺจ จิตฺตสมุฏฺฐานานํ รูปานํ อธิปติปจฺจเยน ปจฺจโย. (3)}

\csromanheader{2}{\textbf{อนนฺตร}}
\proseitem{71}{อปฺปฏิโฆ ธมฺโม อปฺปฏิฆสฺส ธมฺมสฺส อนนฺตรปจฺจเยน ปจฺจโย.\csromanspacer{} ปุริมา ปุริมา อปฺปฏิฆา ขนฺธา ฯเปฯ ผลสมาปตฺติยา อนนฺตรปจฺจเยน ปจฺจโย.}

\csromanheader{2}{10. สปฺปฏิฆทุก \textbf{สมนนฺตร}}
\vspace{\tipitakaparskip}
\csromancenter{อปฺปฏิโฆ ธมฺโม อปฺปฏิฆสฺส ธมฺมสฺส สมนนฺตรปจฺจเยน ปจฺจโย ฯเปฯ.}
\csromanheader{2}{\textbf{สหชาตาทิ}}
\proseitem{73}{\csromanfolio{113}สปฺปฏิโฆ ธมฺโม สปฺปฏิฆสฺส ธมฺมสฺส สหชาตปจฺจเยน ปจฺจโย.\csromanspacer{} นว.\csromanspacer{} อญฺญมญฺญปจฺจเยน ปจฺจโย.\csromanspacer{} ฉ.\csromanspacer{} นิสฺสยปจฺจเยน ปจฺจโย.\csromanspacer{} นว.}

\csromanheader{2}{\textbf{อุปนิสฺสย}}
\proseitem{74}{สปฺปฏิโฆ ธมฺโม อปฺปฏิฆสฺส ธมฺมสฺส อุปนิสฺสยปจฺจเยน ปจฺจโย.\csromanspacer{} \textbf{อารมฺมณูปนิสฺสโย, ปกตูปนิสฺสโย ฯเปฯ.}\csromanspacer{} ปกตูปนิสฺสโย\csromandash{} อุตุํ.\csromanspacer{} เสนาสนํ อุปนิสฺสาย ทานํ เทติ ฯเปฯ สํฆํ ภินฺทติ.\csromanspacer{} อุตุ.\csromanspacer{} เสนาสนํ สทฺธาย ฯเปฯ ผลสมาปตฺติยา อุปนิสฺสยปจฺจเยน ปจฺจโย. (1)}

\prose{อปฺปฏิโฆ ธมฺโม อปฺปฏิฆสฺส ธมฺมสฺส อุปนิสฺสยปจฺจเยน ปจฺจโย.\csromanspacer{} \textbf{อารมฺมณูปนิสฺสโย, อนนฺตรูปนิสฺสโย, ปกตูปนิสฺสโย ฯเปฯ.}\csromanspacer{} ปกตูปนิสฺสโย\csromandash{}สทฺธํ อุปนิสฺสาย ทานํ เทติ ฯเปฯ ทิฏฺฐิํ คณฺหาติ.\csromanspacer{} สีลํ ฯเปฯ.\csromanspacer{} กายิกํ สุขํ.\csromanspacer{} กายิกํ ทุกฺขํ.\csromanspacer{} โภชนํ อุปนิสฺสาย ทานํ เทติ ฯเปฯ สํฆํ ภินฺทติ.\csromanspacer{} สทฺธา ฯเปฯ.\csromanspacer{} โภชนํ สทฺธาย ฯเปฯ ผลสมาปตฺติยา อุปนิสฺสยปจฺจเยน ปจฺจโย. (1)}

\csromanheader{2}{\textbf{ปุเรชาต}}
\proseitem{75}{สปฺปฏิโฆ ธมฺโม อปฺปฏิฆสฺส ธมฺมสฺส ปุเรชาตปจฺจเยน ปจฺจโย.\csromanspacer{} \textbf{อารมฺมณปุเรชาตํ}, วตฺถุปุเรชาตํ.\csromanspacer{}\csromanspacer{}\textbf{อารมฺมณปุเรชาตํ}\csromandash{} จกฺขุํ ฯเปฯ.\csromanspacer{} โผฏฺฐพฺเพ อนิจฺจโต ฯเปฯ โทมนสฺสํ อุปฺปชฺชติ.\csromanspacer{} ทิพฺเพน จกฺขุนา รูปํ ปสฺสติ.\csromanspacer{} ทิพฺพาย โสตธาตุยา สทฺทํ สุณาติ.\csromanspacer{} รูปายตนํ จกฺขุวิญฺญาณสฺส ฯเปฯ.\csromanspacer{} โผฏฺฐพฺพายตนํ กายวิญฺญาณสฺส ปุเรชาตปจฺจเยน ปจฺจโย.\csromanspacer{} \textbf{วตฺถุปุเรชาตํ}\csromandash{}จกฺขายตนํ จกฺขุวิญฺญาณสฺส ฯเปฯ กายายตนํ กายวิญฺญาณสฺส ปุเรชาตปจฺจเยน ปจฺจโย. (1)}

\prose{อปฺปฏิโฆ ธมฺโม อปฺปฏิฆสฺส ธมฺมสฺส ปุเรชาตปจฺจเยน ปจฺจโย.\csromanspacer{} \textbf{อารมฺมณปุเรชาตํ}, วตฺถุปุเรชาตํ.\csromanspacer{}\csromanspacer{}\textbf{อารมฺมณปุเรชาตํ}\csromandash{} วตฺถุํ ฯเปฯ อิตฺถินฺทฺริยํ.\csromanspacer{} ปุริสินฺทฺริยํ.\csromanspacer{} ชีวิตินฺทฺริยํ.\csromanspacer{} อาโปธาตุํ.\csromanspacer{} กพฬีการํ อาหารํ \csromanfolio{114}อนิจฺจโต ฯเปฯ โทมนสฺสํ อุปฺปชฺชติ.\csromanspacer{}\csromanspacer{}\textbf{วตฺถุปุเรชาตํ}\csromandash{}วตฺถุ อปฺปฏิฆานํ ขนฺธานํ ปุเรชาตปจฺจเยน ปจฺจโย. (1)}

\prose{สปฺปฏิโฆ จ อปฺปฏิโฆ จ ธมฺมา อปฺปฏิฆสฺส ธมฺมสฺส ปุเรชาตปจฺจเยน ปจฺจโย.\csromanspacer{} อารมฺมณปุเรชาตํ, วตฺถุปุเรชาตํ.\csromanspacer{}\csromanspacer{}จกฺขายตนญฺจ วตฺถุ จ ฯเปฯ.\csromanspacer{} โผฏฺฐพฺพายตนญฺจ วตฺถุ จ อปฺปฏิฆานํ ขนฺธานํ ปุเรชาตปจฺจเยน ปจฺจโย. (1)}

\csromanheader{2}{\textbf{ปจฺฉาชาต}}
\proseitem{76}{อปฺปฏิโฆ ธมฺโม อปฺปฏิฆสฺส ธมฺมสฺส ปจฺฉาชาตปจฺจเยน ปจฺจโย.\csromanspacer{} ปจฺฉาชาตา อปฺปฏิฆา ขนฺธา ปุเรชาตสฺส อิมสฺส อปฺปฏิฆสฺส กายสฺส ปจฺฉาชาตปจฺจเยน ปจฺจโย. (มูลํ กาตพฺพํ.) ปจฺฉาชาตา อปฺปฏิฆา ขนฺธา ปุเรชาตสฺส อิมสฺส สปฺปฏิฆสฺส กายสฺส ปจฺฉาชาตปจฺจเยน ปจฺจโย. (มูลํ กาตพฺพํ.) ปจฺฉาชาตา อปฺปฏิฆา ขนฺธา ปุเรชาตสฺส อิมสฺส สปฺปฏิฆสฺส จ อปฺปฏิฆสฺส จ กายสฺส ปจฺฉาชาตปจฺจเยน ปจฺจโย. (ทฺวินฺนมฺปิ มูลา กาตพฺพา.) (3)}

\csromanheader{2}{\textbf{อาเสวน}}
\proseitem{77}{อปฺปฏิโฆ ธมฺโม อปฺปฏิฆสฺส ธมฺมสฺส อาเสวนปจฺจเยน ปจฺจโย.\csromanspacer{} ปุริมา ปุริมา อปฺปฏิฆา ขนฺธา ฯเปฯ.\csromanspacer{} โวทานํ มคฺคสฺส อาเสวนปจฺจเยน ปจฺจโย. (1)}

\csromanheader{2}{\textbf{กมฺม}}
\proseitem{78}{อปฺปฏิโฆ ธมฺโม อปฺปฏิฆสฺส ธมฺมสฺส กมฺมปจฺจเยน ปจฺจโย.\csromanspacer{} \textbf{สหชาตา}, นานากฺขณิกา.\csromanspacer{}\csromanspacer{}\textbf{สหชาตา} อปฺปฏิฆา เจตนา สมฺปยุตฺตกานํ ขนฺธานํ อปฺปฏิฆานญฺจ จิตฺตสมุฏฺฐานานํ รูปานํ กมฺมปจฺจเยน ปจฺจโย.\csromanspacer{} ปฏิสนฺธิกฺขเณ ฯเปฯ.\csromanspacer{} \textbf{นานากฺขณิกา} อปฺปฏิฆา เจตนา วิปากานํ ขนฺธานํ อปฺปฏิฆานญฺจ กฏตฺตารูปานํ กมฺมปจฺจเยน ปจฺจโย. (1)}

\prose{อปฺปฏิโฆ ธมฺโม สปฺปฏิฆสฺส ธมฺมสฺส กมฺมปจฺจเยน ปจฺจโย.\csromanspacer{} \textbf{สหชาตา}, นานากฺขณิกา.\csromanspacer{}\csromanspacer{}\textbf{สหชาตา} อปฺปฏิฆา เจตนา สปฺปฏิฆานํ จิตฺตสมุฏฺฐานานํ รูปานํ กมฺมปจฺจเยน ปจฺจโย.\csromanspacer{} ปฏิสนฺธิกฺขเณ ฯเปฯ.\csromanspacer{} \textbf{นานากฺขณิกา} อปฺปฏิฆา เจตนา สปฺปฏิฆานํ กฏตฺตารูปานํ กมฺมปจฺจเยน ปจฺจโย. (2)}

\csromanheader{2}{10. สปฺปฏิฆทุก}
\prose{\csromanfolio{115}อปฺปฏิโฆ ธมฺโม สปฺปฏิฆสฺส จ อปฺปฏิฆสฺส จ ธมฺมสฺส กมฺมปจฺจเยน ปจฺจโย.\csromanspacer{} \textbf{สหชาตา}, นานากฺขณิกา.\csromanspacer{}\csromanspacer{}\textbf{สหชาตา} อปฺปฏิฆา เจตนา สมฺปยุตฺตกานํ ขนฺธานํ สปฺปฏิฆานญฺจ อปฺปฏิฆานญฺจ จิตฺตสมุฏฺฐานานํ รูปานํ กมฺมปจฺจเยน ปจฺจโย.\csromanspacer{} ปฏิสนฺธิกฺขเณ ฯเปฯ.\csromanspacer{} \textbf{นานากฺขณิกา} อปฺปฏิฆา เจตนา วิปากานํ ขนฺธานํ สปฺปฏิฆานญฺจ อปฺปฏิฆานญฺจ กฏตฺตารูปานํ กมฺมปจฺจเยน ปจฺจโย. (3)}

\csromanheader{2}{\textbf{วิปาก}}
\proseitem{79}{อปฺปฏิโฆ ธมฺโม อปฺปฏิฆสฺส ธมฺมสฺส วิปากปจฺจเยน ปจฺจโย.\csromanspacer{} วิปาโก อปฺปฏิโฆ ฯเปฯ.\csromanspacer{} ตีณิ.}

\csromanheader{2}{\textbf{อาหาร}}
\proseitem{80}{อปฺปฏิโฆ ธมฺโม อปฺปฏิฆสฺส ธมฺมสฺส อาหารปจฺจเยน ปจฺจโย.\csromanspacer{} อปฺปฏิฆา อาหารา สมฺปยุตฺตกานํ ขนฺธานํ อปฺปฏิฆานญฺจ จิตฺตสมุฏฺฐานานํ รูปานํ อาหารปจฺจเยน ปจฺจโย.\csromanspacer{} ปฏิสนฺธิกฺขเณ ฯเปฯ.\csromanspacer{} กพฬีกาโร อาหาโร อิมสฺส อปฺปฏิฆสฺส กายสฺส อาหารปจฺจเยน ปจฺจโย. (อวเสสา เทฺวปิ ปญฺหา กาตพฺพา.\csromanspacer{} ปฏิสนฺธิ กพฬีกาโร อาหาโร ทฺวีสุปิ กาตพฺโพ อคฺเค.) (3)}

\csromanheader{2}{\textbf{อินฺทฺริยาทิ}}
\proseitem{81}{สปฺปฏิโฆ ธมฺโม อปฺปฏิฆสฺส ธมฺมสฺส อินฺทฺริยปจฺจเยน ปจฺจโย.\csromanspacer{} จกฺขุนฺทฺริยํ จกฺขุวิญฺญาณสฺส ฯเปฯ.\csromanspacer{} กายินฺทฺริยํ กายวิญฺญาณสฺส อินฺทฺริยปจฺจเยน ปจฺจโย. (1)}

\prose{อปฺปฏิโฆ ธมฺโม อปฺปฏิฆสฺส ธมฺมสฺส อินฺทฺริยปจฺจเยน ปจฺจโย.\csromanspacer{} ตีณิ. (ตีสุปิ ชีวิตินฺทฺริยํ อคฺเค กาตพฺพํ.) (3)}

\prose{สปฺปฏิโฆ จ อปฺปฏิโฆ จ ธมฺมา อปฺปฏิฆสฺส ธมฺมสฺส อินฺทฺริยปจฺจเยน ปจฺจโย.\csromanspacer{} จกฺขุนฺทฺริยญฺจ จกฺขุวิญฺญาณญฺจ จกฺขุวิญฺญาณสหคตานํ ขนฺธานํ อินฺทฺริยปจฺจเยน ปจฺจโย ฯเปฯ.\csromanspacer{} กายินฺทฺริยญฺจ กายวิญฺญาณญฺจ กายวิญฺญาณสหคตานํ ขนฺธานํ อินฺทฺริยปจฺจเยน ปจฺจโย.\csromanspacer{} ฌานปจฺจเยน ปจฺจโย.\csromanspacer{} ตีณิ.\csromanspacer{} มคฺคปจฺจเยน ปจฺจโย.\csromanspacer{} ตีณิ.\csromanspacer{} สมฺปยุตฺตปจฺจเยน ปจฺจโย.\csromanspacer{} เอกํ.}

\csromanheader{2}{\textbf{วิปฺปยุตฺต}}
\proseitem{82}{\csromanfolio{116}สปฺปฏิโฆ ธมฺโม อปฺปฏิฆสฺส ธมฺมสฺส วิปฺปยุตฺตปจฺจเยน ปจฺจโย.\csromanspacer{} \textbf{ปุเรชาตํ} จกฺขายตนํ จกฺขุวิญฺญาณสฺส ฯเปฯ.\csromanspacer{} กายายตนํ กายวิญฺญาณสฺส วิปฺปยุตฺตปจฺจเยน ปจฺจโย. (1)}

\prose{อปฺปฏิโฆ ธมฺโม อปฺปฏิฆสฺส ธมฺมสฺส วิปฺปยุตฺตปจฺจเยน ปจฺจโย.\csromanspacer{} สหชาตํ, ปุเรชาตํ, ปจฺฉาชาตํ.\csromanspacer{}\csromanspacer{}\textbf{สหชาตา} อปฺปฏิฆา ขนฺธา อปฺปฏิฆานํ จิตฺตสมุฏฺฐานานํ รูปานํ วิปฺปยุตฺตปจฺจเยน ปจฺจโย.\csromanspacer{} ปฏิสนฺธิกฺขเณ ฯเปฯ.\csromanspacer{} ขนฺธา วตฺถุสฺส วิปฺปยุตฺตปจฺจเยน ปจฺจโย.\csromanspacer{} วตฺถุ ขนฺธานํ วิปฺปยุตฺตปจฺจเยน ปจฺจโย.\csromanspacer{} \textbf{ปุเรชาตํ} วตฺถุ อปฺปฏิฆานํ ขนฺธานํ วิปฺปยุตฺตปจฺจเยน ปจฺจโย.\csromanspacer{} \textbf{ปจฺฉาชาตา} อปฺปฏิฆา ขนฺธา ปุเรชาตสฺส อิมสฺส อปฺปฏิฆสฺส กายสฺส วิปฺปยุตฺตปจฺจเยน ปจฺจโย. (1)}

\prose{อปฺปฏิโฆ ธมฺโม สปฺปฏิฆสฺส ธมฺมสฺส วิปฺปยุตฺตปจฺจเยน ปจฺจโย.\csromanspacer{} สหชาตํ, ปจฺฉาชาตํ.\csromanspacer{}\csromanspacer{}\textbf{สหชาตา} อปฺปฏิฆา ขนฺธา สปฺปฏิฆานํ จิตฺตสมุฏฺฐานานํ รูปานํ วิปฺปยุตฺตปจฺจเยน ปจฺจโย.\csromanspacer{} ปฏิสนฺธิกฺขเณ ฯเปฯ.\csromanspacer{} \textbf{ปจฺฉาชาตา} อปฺปฏิฆา ขนฺธา ปุเรชาตสฺส อิมสฺส สปฺปฏิฆสฺส กายสฺส วิปฺปยุตฺตปจฺจเยน ปจฺจโย. (2)}

\prose{อปฺปฏิโฆ ธมฺโม สปฺปฏิฆสฺส จ อปฺปฏิฆสฺส จ ธมฺมสฺส วิปฺปยุตฺตปจฺจเยน ปจฺจโย.\csromanspacer{} สหชาตํ, ปจฺฉาชาตํ.\csromanspacer{}\csromanspacer{}\textbf{สหชาตา} อปฺปฏิฆา ขนฺธา สปฺปฏิฆานญฺจ อปฺปฏิฆานญฺจ จิตฺตสมุฏฺฐานานํ รูปานํ วิปฺปยุตฺตปจฺจเยน ปจฺจโย.\csromanspacer{} ปฏิสนฺธิกฺขเณ ฯเปฯ.\csromanspacer{} \textbf{ปจฺฉาชาตา} อปฺปฏิฆา ขนฺธา ปุเรชาตสฺส อิมสฺส สปฺปฏิฆสฺส จ อปฺปฏิฆสฺส จ กายสฺส วิปฺปยุตฺตปจฺจเยน ปจฺจโย. (3)}

\csromanheader{2}{\textbf{อตฺถิ}}
\proseitem{83}{สปฺปฏิโฆ ธมฺโม สปฺปฏิฆสฺส ธมฺมสฺส อตฺถิปจฺจเยน ปจฺจโย.\csromanspacer{} เอกํ. (ปฏิจฺจสทิสา ปฐมปญฺหา.) (1)}

\prose{สปฺปฏิโฆ ธมฺโม อปฺปฏิฆสฺส ธมฺมสฺส อตฺถิปจฺจเยน ปจฺจโย.\csromanspacer{} สหชาตํ, ปุเรชาตํ.\csromanspacer{}\csromanspacer{}\textbf{สหชาตา} สปฺปฏิฆา มหาภูตา อาโปธาตุยา อตฺถิปจฺจเยน ปจฺจโย.\csromanspacer{} สปฺปฏิฆา มหาภูตา อปฺปฏิฆานํ จิตฺตสมุฏฺฐานานํ รูปานํ, กฏตฺตารูปานํ, อุปาทารูปานํ อตฺถิปจฺจเยน ปจฺจโย.\csromanspacer{} โผฏฺฐพฺพายตนํ อิตฺถินฺทฺริยสฺส ฯเปฯ.\csromanspacer{} กพฬีการสฺส อาหารสฺส อตฺถิปจฺจเยน}

\vspace{\tipitakaparskip}
\csromancenter{สปฺปฏิฆทุก ปจฺจโย.\csromanspacer{} พาหิรํ.\csromanspacer{} อาหารสมุฏฺฐานํ.\csromanspacer{} อุตุสมุฏฺฐานํ.\csromanspacer{} อสญฺญสตฺตานํ ฯเปฯ.\csromanspacer{} \textbf{ปุเรชาตํ} จกฺขุํ ฯเปฯ.\csromanspacer{} โผฏฺฐพฺเพ อนิจฺจโต ฯเปฯ โทมนสฺสํ อุปฺปชฺชติ.\csromanspacer{} ทิพฺเพน จกฺขุนา รูปํ ปสฺสติ.\csromanspacer{} ทิพฺพาย โสตธาตุยา สทฺทํ สุณาติ.\csromanspacer{} รูปายตนญฺจ จกฺขายตนญฺจ จกฺขุวิญฺญาณสฺส ฯเปฯ.\csromanspacer{} โผฏฺฐพฺพายตนญฺจ กายายตนญฺจ กายวิญฺญาณสฺส อตฺถิปจฺจเยน ปจฺจโย. (2)}
\prose{\csromanfolio{117}สปฺปฏิโฆ ธมฺโม สปฺปฏิฆสฺส จ อปฺปฏิฆสฺส จ ธมฺมสฺส อตฺถิปจฺจเยน ปจฺจโย.\csromanspacer{} สปฺปฏิฆํ เอกํ มหาภูตํ ทฺวินฺนํ มหาภูตานํ อาโปธาตุยา จ อตฺถิปจฺจเยน ปจฺจโย ฯเปฯ. (ปฏิจฺจสทิสํ ยาว อสญฺญสตฺตา.) (3)}

\proseitem{84}{อปฺปฏิโฆ ธมฺโม อปฺปฏิฆสฺส ธมฺมสฺส อตฺถิปจฺจเยน ปจฺจโย.\csromanspacer{} สหชาตํ, ปุเรชาตํ, ปจฺฉาชาตํ, อาหารํ, อินฺทฺริยํ.\csromanspacer{}\csromanspacer{}\textbf{สหชาโต} อปฺปฏิโฆ เอโก ขนฺโธ ติณฺณนฺนํ ขนฺธานํ ฯเปฯ. (ยาว อสญฺญสตฺตา.) \textbf{ปุเรชาตํ} วตฺถุํ ฯเปฯ อิตฺถินฺทฺริยํ.\csromanspacer{} ปุริสินฺทฺริยํ.\csromanspacer{} ชีวิตินฺทฺริยํ.\csromanspacer{} อาโปธาตุํ.\csromanspacer{} กพฬีการํ อาหารํ อนิจฺจโต ฯเปฯ โทมนสฺสํ อุปฺปชฺชติ.\csromanspacer{} วตฺถุ อปฺปฏิฆานํ ขนฺธานํ อตฺถิปจฺจเยน ปจฺจโย.\csromanspacer{} \textbf{ปจฺฉาชาตา} อปฺปฏิฆา ขนฺธา ปุเรชาตสฺส อิมสฺส อปฺปฏิฆสฺส กายสฺส อตฺถิปจฺจเยน ปจฺจโย.\csromanspacer{} \textbf{กพฬีกาโร} \textbf{อาหาโร} อิมสฺส อปฺปฏิฆสฺส กายสฺส อตฺถิปจฺจเยน ปจฺจโย.\csromanspacer{} \textbf{รูปชีวิตินฺทฺริยํ} อปฺปฏิฆานํ กฏตฺตารูปานํ อตฺถิปจฺจเยน ปจฺจโย. (1)}

\prose{อปฺปฏิโฆ ธมฺโม สปฺปฏิฆสฺส ธมฺมสฺส อตฺถิปจฺจเยน ปจฺจโย.\csromanspacer{} สหชาตํ, ปจฺฉาชาตํ, อาหารํ, อินฺทฺริยํ.\csromanspacer{}\csromanspacer{}\textbf{สหชาตา} อปฺปฏิฆา ขนฺธา สปฺปฏิฆานํ จิตฺตสมุฏฺฐานานํ รูปานํ อตฺถิปจฺจเยน ปจฺจโย.\csromanspacer{} ปฏิสนฺธิกฺขเณ ฯเปฯ.\csromanspacer{} อาโปธาตุ สปฺปฏิฆานํ มหาภูตานํ อตฺถิปจฺจเยน ปจฺจโย.\csromanspacer{} อาโปธาตุ สปฺปฏิฆานํ จิตฺตสมุฏฺฐานานํ รูปานํ, กฏตฺตารูปานํ, อุปาทารูปานํ อตฺถิปจฺจเยน ปจฺจโย.\csromanspacer{} อาโปธาตุ จกฺขายตนสฺส ฯเปฯ โผฏฺฐพฺพายตนสฺส อตฺถิปจฺจเยน ปจฺจโย.\csromanspacer{} พาหิรํ.\csromanspacer{} อาหารสมุฏฺฐานํ.\csromanspacer{} อุตุสมุฏฺฐานํ.\csromanspacer{} อสญฺญสตฺตานํ ฯเปฯ.\csromanspacer{} \textbf{ปจฺฉาชาตา} อปฺปฏิฆา ขนฺธา ปุเรชาตสฺส อิมสฺส สปฺปฏิฆสฺส กายสฺส อตฺถิปจฺจเยน ปจฺจโย.\csromanspacer{} \textbf{กพฬีกาโร} \textbf{อาหาโร} อิมสฺส สปฺปฏิฆสฺส กายสฺส อตฺถิปจฺจเยน ปจฺจโย.\csromanspacer{} \textbf{รูปชีวิตินฺทฺริยํ} สปฺปฏิฆานํ กฏตฺตารูปานํ อตฺถิปจฺจเยน ปจฺจโย. (2)}

\prose{อปฺปฏิโฆ ธมฺโม สปฺปฏิฆสฺส จ อปฺปฏิฆสฺส จ ธมฺมสฺส อตฺถิปจฺจเยน ปจฺจโย.\csromanspacer{} สหชาตํ, ปจฺฉาชาตํ, อาหารํ, อินฺทฺริยํ.\csromanspacer{}\csromanspacer{}\textbf{สหชาโต} \csromanfolio{118}อปฺปฏิโฆ เอโก ขนฺโธ ติณฺณนฺนํ ขนฺธานํ สปฺปฏิฆานญฺจ อปฺปฏิฆานญฺจ ฯเปฯ. (ปฏิจฺจสทิสํ ยาว อสญฺญสตฺตา.) \textbf{ปจฺฉาชาตา} อปฺปฏิฆา ขนฺธา ปุเรชาตสฺส อิมสฺส สปฺปฏิฆสฺส จ อปฺปฏิฆสฺส จ กายสฺส อตฺถิปจฺจเยน ปจฺจโย.\csromanspacer{} \textbf{กพฬีกาโร อาหาโร} อิมสฺส สปฺปฏิฆสฺส จ อปฺปฏิฆสฺส จ กายสฺส อตฺถิปจฺจเยน ปจฺจโย.\csromanspacer{} \textbf{รูปชีวิตินฺทฺริยํ} สปฺปฏิฆานญฺจ อปฺปฏิฆานญฺจ กฏตฺตารูปานํ อตฺถิปจฺจเยน ปจฺจโย. (3)}

\proseitem{85}{สปฺปฏิโฆ จ อปฺปฏิโฆ จ ธมฺมา สปฺปฏิฆสฺส ธมฺมสฺส อตฺถิปจฺจเยน ปจฺจโย. (ปฏิจฺจสทิสํ ยาว อสญฺญสตฺตา.) (1)}

\prose{สปฺปฏิโฆ จ อปฺปฏิโฆ จ ธมฺมา อปฺปฏิฆสฺส ธมฺมสฺส อตฺถิปจฺจเยน ปจฺจโย.\csromanspacer{} สหชาตํ, ปุเรชาตํ.\csromanspacer{}\csromanspacer{}\textbf{สหชาตา} อปฺปฏิฆา ขนฺธา จ มหาภูตา จ อปฺปฏิฆานํ จิตฺตสมุฏฺฐานานํ รูปานํ ฯเปฯ. (ปฏิจฺจสทิสํ ยาว อสญฺญสตฺตา.) \textbf{สหชาโต} จกฺขุวิญฺญาณสหคโต เอโก ขนฺโธ จ จกฺขายตนญฺจ ติณฺณนฺนํ ขนฺธานํ ฯเปฯ เทฺว ขนฺธา จ ฯเปฯ.\csromanspacer{} กายวิญฺญาณสหคโต เอโก ขนฺโธ จ กายายตนญฺจ ติณฺณนฺนํ ขนฺธานํ อตฺถิปจฺจเยน ปจฺจโย ฯเปฯ เทฺว ขนฺธา จ ฯเปฯ. (2)}

\prose{สปฺปฏิโฆ จ อปฺปฏิโฆ จ ธมฺมา สปฺปฏิฆสฺส จ อปฺปฏิฆสฺส จ ธมฺมสฺส อตฺถิปจฺจเยน ปจฺจโย. (ปฏิจฺจสทิสํ.) (3)}

\csromanheader{2}{1. ปจฺจยานุโลม 2. สงฺขฺยาวาร}
\csromanheader{2}{\textbf{สุทฺธ}}
\proseitem{86}{เหตุยา ตีณิ, อารมฺมเณ เทฺว, อธิปติยา จตฺตาริ, อนนฺตเร เอกํ, สมนนฺตเร เอกํ, สหชาเต นว, อญฺญมญฺเญ ฉ, นิสฺสเย นว, อุปนิสฺสเย เทฺว, ปุเรชาเต ตีณิ, ปจฺฉาชาเต ตีณิ, อาเสวเน เอกํ, กมฺเม ตีณิ, วิปาเก ตีณิ, อาหาเร ตีณิ, อินฺทฺริเย ปญฺจ, ฌาเน ตีณิ, มคฺเค ตีณิ, สมฺปยุตฺเต เอกํ, วิปฺปยุตฺเต จตฺตาริ, อตฺถิยา นว, นตฺถิยา เอกํ, วิคเต เอกํ, อวิคเต นว. (เอวํ คเณตพฺพํ.)}

\vspace{\tipitakaparskip}
\csromancenter{อนุโลมํ.}
\csromanheader{2}{10. สปฺปฏิฆทุก \textbf{2. ปจฺจนียุทฺธาร}}
\proseitem{87}{\csromanfolio{119}สปฺปฏิโฆ ธมฺโม สปฺปฏิฆสฺส ธมฺมสฺส สหชาตปจฺจเยน ปจฺจโย. (1)}

\prose{สปฺปฏิโฆ ธมฺโม อปฺปฏิฆสฺส ธมฺมสฺส อารมฺมณปจฺจเยน ปจฺจโย.\csromanspacer{} สหชาตปจฺจเยน ปจฺจโย.\csromanspacer{} อุปนิสฺสยปจฺจเยน ปจฺจโย.\csromanspacer{} ปุเรชาตปจฺจเยน ปจฺจโย. (2)}

\prose{สปฺปฏิโฆ ธมฺโม สปฺปฏิฆสฺส จ อปฺปฏิฆสฺส จ ธมฺมสฺส สหชาตปจฺจเยน ปจฺจโย. (3)}

\proseitem{88}{อปฺปฏิโฆ ธมฺโม อปฺปฏิฆสฺส ธมฺมสฺส อารมฺมณปจฺจเยน ปจฺจโย.\csromanspacer{} สหชาตปจฺจเยน ปจฺจโย.\csromanspacer{} อุปนิสฺสยปจฺจเยน ปจฺจโย.\csromanspacer{} ปุเรชาตปจฺจเยน ปจฺจโย.\csromanspacer{} ปจฺฉาชาตปจฺจเยน ปจฺจโย.\csromanspacer{} กมฺมปจฺจเยน ปจฺจโย.\csromanspacer{} อาหารปจฺจเยน ปจฺจโย.\csromanspacer{} อินฺทฺริยปจฺจเยน ปจฺจโย. (1)}

\prose{อปฺปฏิโฆ ธมฺโม สปฺปฏิฆสฺส ธมฺมสฺส สหชาตปจฺจเยน ปจฺจโย.\csromanspacer{} ปจฺฉาชาตปจฺจเยน ปจฺจโย.\csromanspacer{} กมฺมปจฺจเยน ปจฺจโย.\csromanspacer{} อาหารปจฺจเยน ปจฺจโย.\csromanspacer{} อินฺทฺริยปจฺจเยน ปจฺจโย. (2)}

\prose{อปฺปฏิโฆ ธมฺโม สปฺปฏิฆสฺส จ อปฺปฏิฆสฺส จ ธมฺมสฺส สหชาตปจฺจเยน ปจฺจโย.\csromanspacer{} ปจฺฉาชาตปจฺจเยน ปจฺจโย.\csromanspacer{} กมฺมปจฺจเยน ปจฺจโย.\csromanspacer{} อาหารปจฺจเยน ปจฺจโย.\csromanspacer{} อินฺทฺริยปจฺจเยน ปจฺจโย. (3)}

\proseitem{89}{สปฺปฏิโฆ จ อปฺปฏิโฆ จ ธมฺมา สปฺปฏิฆสฺส ธมฺมสฺส สหชาตปจฺจเยน ปจฺจโย. (1)}

\prose{สปฺปฏิโฆ จ อปฺปฏิโฆ จ ธมฺมา อปฺปฏิฆสฺส ธมฺมสฺส สหชาตํ.\csromanspacer{} ปุเรชาตํ. (2)}

\prose{สปฺปฏิโฆ จ อปฺปฏิโฆ จ ธมฺมา สปฺปฏิฆสฺส จ อปฺปฏิฆสฺส จ ธมฺมสฺส สหชาตปจฺจเยน ปจฺจโย. (3)}

\csromanheader{2}{2. ปจฺจยปจฺจนีย 2. สงฺขฺยาวาร}
\csromanheader{2}{\textbf{สุทฺธ}}
\proseitem{90}{\csromanfolio{120}นเหตุยา นว ฯเปฯ น-อนนฺตเร นว, นสมนนฺตเร นว, นสหชาเต จตฺตาริ, น-อญฺญมญฺเญ นว, นนิสฺสเย จตฺตาริ, น-อุปนิสฺสเย นว, นปุเรชาเต นว ฯเปฯ นสมฺปยุตฺเต นว, นวิปฺปยุตฺเต นว, โน-อตฺถิยา จตฺตาริ, โนนตฺถิยา นว, โนวิคเต นว, โน-อวิคเต จตฺตาริ.}

\vspace{\tipitakaparskip}
\csromancenter{ปจฺจนียํ.}
\csromanheader{2}{3. ปจฺจยานุโลมปจฺจนีย}
\csromanheader{2}{\textbf{เหตุทุก}}
\proseitem{91}{เหตุปจฺจยา น-อารมฺมเณ ตีณิ ฯเปฯ น-อนนฺตเร ตีณิ, นสมนนฺตเร ตีณิ, น-อญฺญมญฺเญ ตีณิ, น-อุปนิสฺสเย ตีณิ ฯเปฯ นสมฺปยุตฺเต ตีณิ, นวิปฺปยุตฺเต เอกํ, โนนตฺถิยา ตีณิ, โนวิคเต ตีณิ.}

\vspace{\tipitakaparskip}
\csromancenter{อนุโลมปจฺจนียํ.}
\csromanheader{2}{4. ปจฺจยปจฺจนียานุโลม}
\csromanheader{2}{\textbf{นเหตุทุก}}
\proseitem{92}{นเหตุปจฺจยา อารมฺมเณ เทฺว, อธิปติยา จตฺตาริ. (อนุโลมมาติกา คเณตพฺพา.) อวิคเต นว.}

\vspace{\tipitakaparskip}
\csromancenter{ปจฺจนียานุโลมํ.}
\nitthitam{สปฺปฏิฆทุกํ นิฏฺฐิตํ.}
\csromanmatikaanchor{mtk.14}
\csromantocmarkref{section}{11. รูปีทุก}{mtk.14}
\bookmark[dest={mtk.14},level=1]{11. รูปีทุก}
\csromanheader{1}{11. รูปีทุก \textbf{11. รูปีทุก 1. ปฏิจฺจวาร}}
\csromanheader{2}{1. ปจฺจยานุโลม 1. วิภงฺควาร}
\csromanheader{2}{\textbf{เหตุ}}
\proseitem{93}{\csromanfolio{121}รูปิํ ธมฺมํ ปฏิจฺจ รูปี ธมฺโม อุปฺปชฺชติ เหตุปจฺจยา.\csromanspacer{} เอกํ มหาภูตํ ปฏิจฺจ ตโย มหาภูตา ฯเปฯ เทฺว มหาภูเต ฯเปฯ มหาภูเต ปฏิจฺจ จิตฺตสมุฏฺฐานํ รูปํ, กฏตฺตารูปํ, อุปาทารูปํ. (1)}

\prose{รูปิํ ธมฺมํ ปฏิจฺจ อรูปี ธมฺโม อุปฺปชฺชติ เหตุปจฺจยา.\csromanspacer{} ปฏิสนฺธิกฺขเณ วตฺถุํ ปฏิจฺจ อรูปิโน ขนฺธา. (2)}

\prose{รูปิํ ธมฺมํ ปฏิจฺจ รูปี จ อรูปี จ ธมฺมา อุปฺปชฺชนฺติ เหตุปจฺจยา.\csromanspacer{} ปฏิสนฺธิกฺขเณ วตฺถุํ ปฏิจฺจ อรูปิโน ขนฺธา.\csromanspacer{} มหาภูเต ปฏิจฺจ กฏตฺตารูปํ. (3)}

\proseitem{94}{อรูปิํ ธมฺมํ ปฏิจฺจ อรูปี ธมฺโม อุปฺปชฺชติ เหตุปจฺจยา.\csromanspacer{} อรูปิํ เอกํ ขนฺธํ ปฏิจฺจ ตโย ขนฺธา ฯเปฯ เทฺว ขนฺเธ ฯเปฯ.\csromanspacer{} ปฏิสนฺธิกฺขเณ ฯเปฯ. (1)}

\prose{อรูปิํ ธมฺมํ ปฏิจฺจ รูปี ธมฺโม อุปฺปชฺชติ เหตุปจฺจยา.\csromanspacer{} อรูปิโน ขนฺเธ ปฏิจฺจ จิตฺตสมุฏฺฐานํ รูปํ.\csromanspacer{} ปฏิสนฺธิกฺขเณ ฯเปฯ.}

\prose{อรูปิํ ธมฺมํ ปฏิจฺจ รูปี จ อรูปี จ ธมฺมา อุปฺปชฺชนฺติ เหตุปจฺจยา.\csromanspacer{} อรูปิํ เอกํ ขนฺธํ ปฏิจฺจ ตโย ขนฺธา จิตฺตสมุฏฺฐานญฺจ รูปํ ฯเปฯ เทฺว ขนฺเธ ฯเปฯ.\csromanspacer{} ปฏิสนฺธิกฺขเณ ฯเปฯ. (3)}

\proseitem{95}{รูปิญฺจ อรูปิญฺจ ธมฺมํ ปฏิจฺจ รูปี ธมฺโม อุปฺปชฺชติ เหตุปจฺจยา.\csromanspacer{} อรูปิโน ขนฺเธ จ มหาภูเต จ ปฏิจฺจ จิตฺตสมุฏฺฐานํ รูปํ.\csromanspacer{} ปฏิสนฺธิกฺขเณ ฯเปฯ. (1)}

\prose{รูปิญฺจ อรูปิญฺจ ธมฺมํ ปฏิจฺจ อรูปี ธมฺโม อุปฺปชฺชติ เหตุปจฺจยา.\csromanspacer{} ปฏิสนฺธิกฺขเณ อรูปิํ เอกํ ขนฺธญฺจ วตฺถุญฺจ ปฏิจฺจ ตโย ขนฺธา ฯเปฯ เทฺว ขนฺเธ ฯเปฯ. (2)}

\prose{รูปิญฺจ อรูปิญฺจ ธมฺมํ ปฏิจฺจ รูปี จ อรูปี จ ธมฺมา อุปฺปชฺชนฺติ เหตุปจฺจยา.\csromanspacer{} ปฏิสนฺธิกฺขเณ อรูปิํ เอกํ ขนฺธญฺจ วตฺถุญฺจ ปฏิจฺจ ตโย ขนฺธา ฯเปฯ เทฺว ขนฺเธ ฯเปฯ.\csromanspacer{} อรูปิโน ขนฺเธ จ มหาภูเต จ ปฏิจฺจ กฏตฺตารูปํ. (สํขิตฺตํ.) (3)}

\csromanheader{2}{1. ปจฺจยานุโลม 2. สงฺขฺยาวาร}
\csromanheader{2}{\textbf{สุทฺธ}}
\proseitem{96}{\csromanfolio{122}เหตุยา นว, อารมฺมเณ ตีณิ, อธิปติยา ปญฺจ, อนนฺตเร ตีณิ, สมนนฺตเร ตีณิ, สหชาเต นว, อญฺญมญฺเญ ฉ, นิสฺสเย นว, อุปนิสฺสเย ตีณิ, ปุเรชาเต เอกํ, อาเสวเน เอกํ, กมฺเม นว, วิปาเก นว, อาหาเร นว, อินฺทฺริเย นว, ฌาเน นว, มคฺเค นว, สมฺปยุตฺเต ตีณิ, วิปฺปยุตฺเต นว, อตฺถิยา นว, นตฺถิยา ตีณิ, วิคเต ตีณิ, อวิคเต นว.}

\vspace{\tipitakaparskip}
\csromancenter{อนุโลมํ.}
\csromanheader{2}{2. ปจฺจยปจฺจนีย 1. วิภงฺควาร}
\csromanheader{2}{\textbf{นเหตุ}}
\proseitem{97}{รูปิํ ธมฺมํ ปฏิจฺจ รูปี ธมฺโม อุปฺปชฺชติ นเหตุปจฺจยา.\csromanspacer{} ตีณิ.}

\prose{อรูปิํ ธมฺมํ ปฏิจฺจ อรูปี ธมฺโม อุปฺปชฺชติ นเหตุปจฺจยา.\csromanspacer{} อเหตุกํ อรูปิํ เอกํ ขนฺธํ ปฏิจฺจ ตโย ขนฺธา ฯเปฯ เทฺว ขนฺเธ ฯเปฯ.\csromanspacer{} อเหตุกปฏิสนฺธิกฺขเณ ฯเปฯ.\csromanspacer{} วิจิกิจฺฉาสหคเต อุทฺธจฺจสหคเต ขนฺเธ ปฏิจฺจ วิจิกิจฺฉาสหคโต อุทฺธจฺจสหคโต โมโห. (\textbf{นเหตุ}ปจฺจยา นว ปญฺหา, “อเหตุกนฺ”ติ นิยาเมตพฺพํ.)}

\csromanheader{2}{2. ปจฺจยปจฺจนีย 2. สงฺขฺยาวาร}
\csromanheader{2}{\textbf{สุทฺธ}}
\vspace{\tipitakaparskip}
\csromancenter{\textbf{นเหตุ}ยา นว, น-อารมฺมเณ ตีณิ, น-อธิปติยา นว, น-อนนฺตเร ตีณิ, นสมนนฺตเร ตีณิ, น-อญฺญมญฺเญ ตีณิ, น-อุปนิสฺสเย ตีณิ, นปุเรชาเต นว, นปจฺฉาชาเต นว, น-อาเสวเน นว, นกมฺเม เทฺว, นวิปาเก ปญฺจ, น-อาหาเร เอกํ, น-อินฺทฺริเย เอกํ, นฌาเน เทฺว, นมคฺเค นว, นสมฺปยุตฺเต ตีณิ, นวิปฺปยุตฺเต เทฺว, โนนตฺถิยา ตีณิ, โนวิคเต ตีณิ.}
\csromancenter{ปจฺจนียํ.}
\csromanheader{2}{11. รูปีทุก \textbf{3. ปจฺจยานุโลมปจฺจนีย}}
\csromanheader{2}{\textbf{เหตุทุก}}
\proseitem{99}{\csromanfolio{123}\textbf{เหตุ}ปจฺจยา น-อารมฺมเณ ตีณิ, น-อธิปติยา นว, น-อนนฺตเร ตีณิ, นสมนนฺตเร ตีณิ, น-อญฺญมญฺเญ ตีณิ, น-อุปนิสฺสเย ตีณิ, นปุเรชาเต นว, นปจฺฉาชาเต นว, น-อาเสวเน นว, นกมฺเม เอกํ, นวิปาเก ปญฺจ, นสมฺปยุตฺเต ตีณิ, นวิปฺปยุตฺเต เอกํ, โนนตฺถิยา ตีณิ, โนวิคเต ตีณิ.}

\vspace{\tipitakaparskip}
\csromancenter{อนุโลมปจฺจนียํ.}
\csromanheader{2}{4. ปจฺจยปจฺจนียานุโลม}
\csromanheader{2}{\textbf{นเหตุทุก}}
\proseitem{100}{นเหตุปจฺจยา อารมฺมเณ ตีณิ, อนนฺตเร ตีณิ, สมนนฺตเร ตีณิ, สหชาเต นว, อญฺญมญฺเญ ฉ, นิสฺสเย นว, อุปนิสฺสเย ตีณิ, ปุเรชาเต เอกํ, อาเสวเน เอกํ, กมฺเม นว, วิปาเก นว, อาหาเร นว, อินฺทฺริเย นว, ฌาเน นว, มคฺเค เอกํ, สมฺปยุตฺเต ตีณิ, วิปฺปยุตฺเต นว, อตฺถิยา นว, นตฺถิยา ตีณิ, วิคเต ตีณิ, อวิคเต นว.}

\vspace{\tipitakaparskip}
\csromancenter{ปจฺจนียานุโลมํ.}
\csromancenter{(สหชาตวาโรปิ ปฏิจฺจวารสทิโส.)}
\csromanheader{2}{11. รูปีทุก 3. ปจฺจยวาร}
\csromanheader{2}{1. ปจฺจยานุโลม 1. วิภงฺควาร}
\csromanheader{2}{\textbf{เหตุ}}
\proseitem{101}{รูปิํ ธมฺมํ ปจฺจยา รูปี ธมฺโม อุปฺปชฺชติ เหตุปจฺจยา.\csromanspacer{} เอกํ มหาภูตํ ฯเปฯ. (ปฏิจฺจสทิสํ.) (1)}

\prose{\csromanfolio{124}รูปิํ ธมฺมํ ปจฺจยา อรูปี ธมฺโม อุปฺปชฺชติ เหตุปจฺจยา.\csromanspacer{} วตฺถุํ ปจฺจยา อรูปิโน ขนฺธา.\csromanspacer{} ปฏิสนฺธิกฺขเณ ฯเปฯ. (2)}

\prose{รูปิํ ธมฺมํ ปจฺจยา รูปี จ อรูปี จ ธมฺมา อุปฺปชฺชนฺติ เหตุปจฺจยา.\csromanspacer{} วตฺตุํ ปจฺจยา อรูปิโน ขนฺธา.\csromanspacer{} มหาภูเต ปจฺจยา จิตฺตสมุฏฺฐานํ รูปํ.\csromanspacer{} ปฏิสนฺธิกฺขเณ ฯเปฯ. (เอวํ อวเสสา ปญฺหา, ปวตฺติปฏิสนฺธิ วิภชิตพฺพา.) (3)}

\csromanheader{2}{\textbf{อารมฺมณ}}
\proseitem{102}{รูปิํ ธมฺมํ ปจฺจยา อรูปี ธมฺโม อุปฺปชฺชติ อารมฺมณปจฺจยา.\csromanspacer{} จกฺขายตนํ ปจฺจยา จกฺขุวิญฺญาณํ ฯเปฯ.\csromanspacer{} กายายตนํ ปจฺจยา กายวิญฺญาณํ.\csromanspacer{} วตฺถุํ ปจฺจยา อรูปิโน ขนฺธา.\csromanspacer{} ปฏิสนฺธิกฺขเณ ฯเปฯ. (1)}

\prose{อรูปิํ ธมฺมํ ปจฺจยา อรูปี ธมฺโม อุปฺปชฺชติ อารมฺมณปจฺจยา.\csromanspacer{} อรูปิํ เอกํ ขนฺธํ ฯเปฯ เทฺว ขนฺเธ ฯเปฯ. (1)}

\prose{รูปิญฺจ อรูปิญฺจ ธมฺมํ ปจฺจยา อรูปี ธมฺโม อุปฺปชฺชติ อารมฺมณปจฺจยา.\csromanspacer{} จกฺขุวิญฺญาณสหคตํ เอกํ ขนฺธญฺจ จกฺขายตนญฺจ ปจฺจยา ตโย ขนฺธา ฯเปฯ เทฺว ขนฺเธ ฯเปฯ.\csromanspacer{} กายวิญฺญาณสหคตํ ฯเปฯ.\csromanspacer{} อรูปิํ เอกํ ขนฺธญฺจ วตฺถุญฺจ ปจฺจยา ตโย ขนฺธา ฯเปฯ เทฺว ขนฺเธ ฯเปฯ. (สํขิตฺตํ.) (3)}

\csromanheader{2}{1. ปจฺจยานุโลม 2. สงฺขฺยาวาร}
\csromanheader{2}{\textbf{สุทฺธ}}
\proseitem{103}{เหตุยา นว, อารมฺมเณ ตีณิ, อธิปติยา นว, อนนฺตเร ตีณิ, สมนนฺตเร ตีณิ, สหชาเต นว, อญฺญมญฺเญ ฉ, นิสฺสเย นว, อุปนิสฺสเย ตีณิ, ปุเรชาเต ตีณิ, อาเสวเน ตีณิ, กมฺเม นว ฯเปฯ มคฺเค นว, สมฺปยุตฺเต ตีณิ, วิปฺปยุตฺเต นว, อตฺถิยา นว, นตฺถิยา ตีณิ, วิคเต ตีณิ, อวิคเต นว.}

\vspace{\tipitakaparskip}
\csromancenter{อนุโลมํ.}
\csromanheader{2}{11. รูปีทุก \textbf{2. ปจฺจยปจฺจนีย 1. วิภงฺควาร}}
\csromanheader{2}{\textbf{นเหตุ}}
\proseitem{104}{\csromanfolio{125}รูปิํ ธมฺมํ ปจฺจยา รูปี ธมฺโม อุปฺปชฺชติ นเหตุปจฺจยา.\csromanspacer{} เอกํ มหาภูตํ ฯเปฯ.\csromanspacer{} อสญฺญสตฺตานํ เอกํ มหาภูตํ ฯเปฯ. (1)}

\prose{รูปิํ ธมฺมํ ปจฺจยา อรูปี ธมฺโม อุปฺปชฺชติ นเหตุปจฺจยา.\csromanspacer{} จกฺขายตนํ ปจฺจยา จกฺขุวิญฺญาณํ ฯเปฯ.\csromanspacer{} กายายตนํ ปจฺจยา กายวิญฺญาณํ.\csromanspacer{} วตฺถุํ ปจฺจยา อเหตุกา อรูปิโน ขนฺธา.\csromanspacer{} อเหตุกปฏิสนฺธิกฺขเณ ฯเปฯ.\csromanspacer{} วตฺถุํ ปจฺจยา วิจิกิจฺฉาสหคโต อุทฺธจฺจสหคโต โมโห. (2)}

\prose{รูปิํ ธมฺมํ ปจฺจยา รูปี จ อรูปี จ ธมฺมา อุปฺปชฺชนฺติ นเหตุปจฺจยา. (ปวตฺติปฏิสนฺธิ กาตพฺพา.) (3)}

\proseitem{105}{อรูปิํ ธมฺมํ ปจฺจยา อรูปี ธมฺโม อุปฺปชฺชติ นเหตุปจฺจยา.\csromanspacer{} อเหตุกํ อรูปิํ เอกํ ขนฺธํ ฯเปฯ.\csromanspacer{} ปฏิสนฺธิกฺขเณ ฯเปฯ.\csromanspacer{} วิจิกิจฺฉาสหคเต อุทฺธจฺจสหคเต ขนฺเธ ปจฺจยา วิจิกิจฺฉาสหคโต อุทฺธจฺจสหคโต โมโห. (1)}

\prose{อรูปิํ ธมฺมํ ปจฺจยา รูปี ธมฺโม อุปฺปชฺชติ นเหตุปจฺจยา.\csromanspacer{} อเหตุเก อรูปิโน ขนฺเธ ปจฺจยา จิตฺตสมุฏฺฐานํ รูปํ.\csromanspacer{} ปฏิสนฺธิกฺขเณ ฯเปฯ. (2)}

\prose{อรูปิํ ธมฺมํ ปจฺจยา รูปี จ อรูปี จ ธมฺมา อุปฺปชฺชนฺติ นเหตุปจฺจยา.\csromanspacer{} อเหตุกํ อรูปิํ เอกํ ขนฺธํ ปจฺจยา ตโย ขนฺธา จิตฺตสมุฏฺฐานญฺจ รูปํ ฯเปฯ เทฺว ขนฺเธ ฯเปฯ.\csromanspacer{} ปฏิสนฺธิกฺขเณ ฯเปฯ. (3)}

\proseitem{106}{รูปิญฺจ อรูปิญฺจ ธมฺมํ ปจฺจยา รูปี ธมฺโม อุปฺปชฺชติ นเหตุปจฺจยา.\csromanspacer{} อเหตุเก อรูปิโน ขนฺเธ จ มหาภูเต จ ปจฺจยา จิตฺตสมุฏฺฐานํ รูปํ.\csromanspacer{} ปฏิสนฺธิกฺขเณ ฯเปฯ. (1)}

\prose{รูปิญฺจ อรูปิญฺจ ธมฺมํ ปจฺจยา อรูปี ธมฺโม อุปฺปชฺชติ นเหตุปจฺจยา.\csromanspacer{} จกฺขุวิญฺญาณสหคตํ เอกํ ขนฺธญฺจ จกฺขายตนญฺจ ปจฺจยา ตโย ขนฺธา ฯเปฯ เทฺว ขนฺเธ ฯเปฯ.\csromanspacer{} กายวิญฺญาณสหคตํ ฯเปฯ.\csromanspacer{} อรูปิํ เอกํ ขนฺธญฺจ วตฺถุญฺจ ปจฺจยา ตโย ขนฺธา ฯเปฯ เทฺว ขนฺเธ ฯเปฯ.\csromanspacer{} ปฏิสนฺธิกฺขเณ ฯเปฯ.\csromanspacer{} วิจิกิจฺฉาสหคเต อุทฺธจฺจสหคเต ขนฺเธ จ วตฺถุญฺจ ปจฺจยา วิจิกิจฺฉาสหคโต อุทฺธจฺจสหคโต โมโห. (2)}

\prose{\csromanfolio{126}รูปิญฺจ อรูปิญฺจ ธมฺมํ ปจฺจยา รูปี จ อรูปี จ ธมฺมา อุปฺปชฺชนฺติ นเหตุปจฺจยา.\csromanspacer{} อเหตุกํ อรูปิํ เอกํ ขนฺธญฺจ วตฺถุญฺจ ปจฺจยา ตโย ขนฺธา ฯเปฯ เทฺว ขนฺเธ ฯเปฯ.\csromanspacer{} อรูปิโน ขนฺเธ จ มหาภูเต จ ปจฺจยา จิตฺตสมุฏฺฐานํ รูปํ.\csromanspacer{} ปฏิสนฺธิกฺขเณ. (3)}

\csromanheader{2}{2. ปจฺจยปจฺจนีย 2. สงฺขฺยาวาร}
\csromanheader{2}{\textbf{สุทฺธ}}
\proseitem{107}{นเหตุยา นว, น-อารมฺมเณ ตีณิ, น-อธิปติยา นว, น-อนนฺตเร ตีณิ, นสมนนฺตเร ตีณิ, น-อญฺญมญฺเญ ตีณิ, น-อุปนิสฺสเย ตีณิ, นปุเรชาเต นว, นปจฺฉาชาเต นว, น-อาเสวเน นว, นกมฺเม จตฺตาริ, นวิปาเก นว, น-อาหาเร เอกํ, น-อินฺทฺริเย เอกํ, นฌาเน จตฺตาริ, นมคฺเค นว, นสมฺปยุตฺเต ตีณิ, นวิปฺปยุตฺเต เทฺว, โนนตฺถิยา ตีณิ, โนวิคเต ตีณิ.}

\vspace{\tipitakaparskip}
\csromancenter{ปจฺจนียํ.}
\csromanheader{2}{3. ปจฺจยานุโลมปจฺจนีย}
\csromanheader{2}{\textbf{เหตุทุก}}
\proseitem{108}{เหตุปจฺจยา น-อารมฺมเณ ตีณิ. (สํขิตฺตํ.\csromanspacer{} สพฺเพ กาตพฺพา.) นกมฺเม ตีณิ, นวิปาเก นว, นสมฺปยุตฺเต ตีณิ, นวิปฺปยุตฺเต เอกํ, โนนตฺถิยา ตีณิ, โนวิคเต ตีณิ.}

\vspace{\tipitakaparskip}
\csromancenter{อนุโลมปจฺจนียํ.}
\csromanheader{2}{4. ปจฺจยปจฺจนียานุโลม}
\csromanheader{2}{\textbf{นเหตุทุก}}
\proseitem{109}{นเหตุปจฺจยา อารมฺมเณ ตีณิ, (สพฺเพ กาตพฺพา.) ฯเปฯ ฌาเน นว, มคฺเค ตีณิ ฯเปฯ อวิคเต นว.}

\vspace{\tipitakaparskip}
\csromancenter{ปจฺจนียานุโลมํ.}
\csromancenter{(นิสฺสยวาโรปิ ปจฺจยวารสทิโส.)}
\csromanheader{2}{11. รูปีทุก \textbf{11. รูปีทุก 5. สํสฏฺฐวาร}}
\csromanheader{2}{1. ปจฺจยานุโลมาทิ}
\proseitem{110}{\csromanfolio{127}อรูปิํ ธมฺมํ สํสฏฺโฐ อรูปี ธมฺโม อุปฺปชฺชติ เหตุปจฺจยา.\csromanspacer{} อรูปิํ เอกํ ขนฺธํ สํสฏฺฐา ตโย ขนฺธา ฯเปฯ เทฺว ขนฺเธ ฯเปฯ.\csromanspacer{} ปฏิสนฺธิกฺขเณ ฯเปฯ.}

\prose{\textbf{เหตุ}ยา เอกํ ฯเปฯ อวิคเต เอกํ. (เอวํ ปจฺจนียาทีนิ คณนาปิ สมฺปยุตฺตวาเรปิ สพฺเพ กาตพฺพา.\csromanspacer{} เอโกเยว ปโญฺห.)}

\csromanheader{2}{11. รูปีทุก 7. ปญฺหาวาร}
\csromanheader{2}{1. ปจฺจยานุโลม 1. วิภงฺควาร}
\csromanheader{2}{\textbf{เหตุ}}
\proseitem{111}{อรูปี ธมฺโม อรูปิสฺส ธมฺมสฺส เหตุปจฺจเยน ปจฺจโย.\csromanspacer{} อรูปี เหตู สมฺปยุตฺตกานํ ขนฺธานํ เหตุปจฺจเยน ปจฺจโย.\csromanspacer{} ปฏิสนฺธิกฺขเณ ฯเปฯ. (1)}

\prose{อรูปี ธมฺโม รูปิสฺส ธมฺมสฺส เหตุปจฺจเยน ปจฺจโย.\csromanspacer{} อรูปี เหตู จิตฺตสมุฏฺฐานานํ รูปานํ เหตุปจฺจเยน ปจฺจโย.\csromanspacer{} ปฏิสนฺธิกฺขเณ ฯเปฯ. (2)}

\prose{อรูปี ธมฺโม รูปิสฺส จ อรูปิสฺส จ ธมฺมสฺส เหตุปจฺจเยน ปจฺจโย.\csromanspacer{} อรูปี เหตู สมฺปยุตฺตกานํ ขนฺธานํ จิตฺตสมุฏฺฐานานญฺจ รูปานํ เหตุปจฺจเยน ปจฺจโย.\csromanspacer{} ปฏิสนฺธิกฺขเณ ฯเปฯ. (3)}

\csromanheader{2}{\textbf{อารมฺมณ}}
\proseitem{112}{รูปี ธมฺโม อรูปิสฺส ธมฺมสฺส อารมฺมณปจฺจเยน ปจฺจโย.\csromanspacer{} จกฺขุํ ฯเปฯ วตฺถุํ.\csromanspacer{} อิตฺถินฺทฺริยํ.\csromanspacer{} ปุริสินฺทฺริยํ.\csromanspacer{} ชีวิตินฺทฺริยํ.\csromanspacer{} อาโปธาตุํ.\csromanspacer{} กพฬีการํ อาหารํ อนิจฺจโต ฯเปฯ โทมนสฺสํ อุปฺปชฺชติ.\csromanspacer{} ทิพฺเพน จกฺขุนา รูปํ ปสฺสติ.\csromanspacer{} ทิพฺพาย โสตธาตุยา สทฺทํ สุณาติ.\csromanspacer{} รูปายตนํ จกฺขุวิญฺญาณสฺส อารมฺมณปจฺจเยน ปจฺจโย ฯเปฯ.\csromanspacer{} โผฏฺฐพฺพายตนํ กายวิญฺญาณสฺส อารมฺมณปจฺจเยน ปจฺจโย.\csromanspacer{} รูปิโน ขนฺธา อิทฺธิวิธญาณสฺส, ปุพฺเพนิวาสานุสฺสติญาณสฺส, อนาคตํสญาณสฺส, อาวชฺชนาย อารมฺมณปจฺจเยน ปจฺจโย. (1)}

\proseitem{113}{\csromanfolio{128}อรูปี ธมฺโม อรูปิสฺส ธมฺมสฺส อารมฺมณปจฺจเยน ปจฺจโย.\csromanspacer{} ทานํ ฯเปฯ สีลํ ฯเปฯ อุโปสถกมฺมํ กตฺวา ตํ ปจฺจเวกฺขติ.\csromanspacer{} ปุพฺเพ สุจิณฺณานิ ปจฺจเวกฺขติ.\csromanspacer{} ฌานา ฯเปฯ.\csromanspacer{} อริยา มคฺคา วุฏฺฐหิตฺวา มคฺคํ ปจฺจเวกฺขนฺติ, ผลํ ปจฺจเวกฺขนฺติ, นิพฺพานํ ปจฺจเวกฺขนฺติ.\csromanspacer{} นิพฺพานํ โคตฺรภุสฺส, โวทานสฺส, มคฺคสฺส, ผลสฺส, อาวชฺชนาย อารมฺมณปจฺจเยน ปจฺจโย.\csromanspacer{} อริยา ปหีเน กิเลเส ฯเปฯ.\csromanspacer{} วิกฺขมฺภิเต กิเลเส ฯเปฯ.\csromanspacer{} ปุพฺเพ ฯเปฯ.\csromanspacer{} อรูปิโน ขนฺเธ อนิจฺจโต ฯเปฯ โทมนสฺสํ อุปฺปชฺชติ.\csromanspacer{} เจโตปริยญาเณน อรูปิจิตฺตสมงฺคิสฺส จิตฺตํ ชานาติ.\csromanspacer{} อากาสานญฺจายตนํ ฯเปฯ.\csromanspacer{} เนวสญฺญานาสญฺญายตนสฺส ฯเปฯ.\csromanspacer{} อรูปิโน ขนฺธา อิทฺธิวิธญาณสฺส, เจโตปริยญาณสฺส, ปุพฺเพนิวาสานุสฺสติญาณสฺส, ยถากมฺมูปคญาณสฺส, อนาคตํสญาณสฺส, อาวชฺชนาย อารมฺมณปจฺจเยน ปจฺจโย. (1)}

\csromanheader{2}{\textbf{อธิปติ}}
\proseitem{114}{รูปี ธมฺโม อรูปิสฺส ธมฺมสฺส อธิปติปจฺจเยน ปจฺจโย.\csromanspacer{} \textbf{อารมฺมณาธิปติ\csromandash{}}จกฺขุํ ฯเปฯ.\csromanspacer{} กพฬีการํ อาหารํ ครุํ กตฺวา อสฺสาเทติ อภินนฺทติ, ตํ ครุํ กตฺวา ราโค อุปฺปชฺชติ.\csromanspacer{} ทิฏฺฐิ อุปฺปชฺชติ. (1)}

\prose{อรูปี ธมฺโม อรูปิสฺส ธมฺมสฺส อธิปติปจฺจเยน ปจฺจโย.\csromanspacer{} \textbf{อารมฺมณาธิปติ}, สหชาตาธิปติ.\csromanspacer{}\csromanspacer{}\textbf{อารมฺมณาธิปติ\csromandash{}}ทานํ ฯเปฯ. (สํขิตฺตํ.) นิพฺพานํ มคฺคสฺส, ผลสฺส อธิปติปจฺจเยน ปจฺจโย.\csromanspacer{} อรูปิโน ขนฺเธ ครุํ กตฺวา อสฺสาเทติ ฯเปฯ ราโค อุปฺปชฺชติ.\csromanspacer{} ทิฏฺฐิ อุปฺปชฺชติ.\csromanspacer{} \textbf{สหชาตาธิปติ\csromandash{}}อรูปี อธิปติ สมฺปยุตฺตกานํ ขนฺธานํ อธิปติปจฺจเยน ปจฺจโย.}

\prose{อรูปี ธมฺโม รูปิสฺส ธมฺมสฺส อธิปติปจฺจเยน ปจฺจโย.\csromanspacer{} \textbf{สหชาตาธิปติ\csromandash{}}อรูปี อธิปติ จิตฺตสมุฏฺฐานานํ รูปานํ อธิปติปจฺจเยน ปจฺจโย. (2)}

\prose{อรูปี ธมฺโม รูปิสฺส จ อรูปิสฺส จ ธมฺมสฺส อธิปติปจฺจเยน ปจฺจโย.\csromanspacer{} \textbf{สหชาตาธิปติ\csromandash{}}อรูปี อธิปติ สมฺปยุตฺตกานํ ขนฺธานํ จิตฺตสมุฏฺฐานานญฺจ รูปานํ อธิปติปจฺจเยน ปจฺจโย. (3)}

\csromanheader{2}{11. รูปีทุก \textbf{อนนฺตราทิ}}
\proseitem{115}{\csromanfolio{129}อรูปี ธมฺโม อรูปิสฺส ธมฺมสฺส อนนฺตรปจฺจเยน ปจฺจโย.\csromanspacer{} ปุริมา ปุริมา อรูปิโน ขนฺโธ ปจฺฉิมานํ ปจฺฉิมานํ อรูปีนํ ขนฺธานํ ฯเปฯ ผลสมาปตฺติยา อนนฺตรปจฺจเยน ปจฺจโย.\csromanspacer{} สมนนฺตรปจฺจเยน ปจฺจโย. (สหชาตปจฺจเย สตฺต.\csromanspacer{} อิห ฆฏนา นตฺถิ.\csromanspacer{} อญฺญมญฺญปจฺจเย ฉ.\csromanspacer{} นิสฺสยปจฺจเย สตฺต ปญฺหา.\csromanspacer{} อิห ฆฏนา นตฺถิ.)}

\csromanheader{2}{\textbf{อุปนิสฺสย}}
\proseitem{116}{รูปี ธมฺโม อรูปิสฺส ธมฺมสฺส อุปนิสฺสยปจฺจเยน ปจฺจโย.\csromanspacer{} \textbf{อารมฺมณูปนิสฺสโย, ปกตูปนิสฺสโย ฯเปฯ.}\csromanspacer{} ปกตูปนิสฺสโย\csromandash{}อุตุํ.\csromanspacer{} โภชนํ.\csromanspacer{} เสนาสนํ อุปนิสฺสาย ทานํ เทติ ฯเปฯ สํฆํ ภินฺทติ.\csromanspacer{} อุตุ.\csromanspacer{} โภชนํ.\csromanspacer{} เสนาสนํ สทฺธาย ฯเปฯ ผลสมาปตฺติยา อุปนิสฺสยปจฺจเยน ปจฺจโย. (1)}

\prose{อรูปี ธมฺโม อรูปิสฺส ธมฺมสฺส อุปนิสฺสยปจฺจเยน ปจฺจโย.\csromanspacer{} \textbf{อารมฺมณูปนิสฺสโย, อนนฺตรูปนิสฺสโย, ปกตูปนิสฺสโย ฯเปฯ.}\csromanspacer{} ปกตูปนิสฺสโย\csromandash{}สทฺธํ อุปนิสฺสาย ทานํ เทติ.\csromanspacer{} สีลํ ฯเปฯ.\csromanspacer{} กายิกํ ทุกฺขํ อุปนิสฺสาย ทานํ เทติ ฯเปฯ ปํฆํ ภินฺทติ.\csromanspacer{} สทฺธา ฯเปฯ.\csromanspacer{} กายิกํ ทุกฺขํ.\csromanspacer{} สทฺธาย ฯเปฯ ผลสมาปตฺติยา อุปนิสฺสยปจฺจเยน ปจฺจโย. (1)}

\csromanheader{2}{\textbf{ปุเรชาต}}
\proseitem{117}{รูปี ธมฺโม อรูปิสฺส ธมฺมสฺส ปุเรชาตปจฺจเยน ปจฺจโย.\csromanspacer{} \textbf{อารมฺมณปุเรชาตํ}, วตฺถุปุเรชาตํ.\csromanspacer{}\csromanspacer{}\textbf{อารมฺมณปุเรชาตํ}\csromandash{}จกฺขุํ ฯเปฯ วตฺถุํ ฯเปฯ.\csromanspacer{} กพฬีการํ อาหารํ อนิจฺจโต ฯเปฯ โทมนสฺสํ อุปฺปชฺชติ.\csromanspacer{} ทิพฺเพน จกฺขุนา รูปํ ปสฺสติ.\csromanspacer{} ทิพฺพาย โสตธาตุยา สทฺทํ สุณาติ.\csromanspacer{} รูปายตนํ จกฺขุวิญฺญาณสฺส ฯเปฯ.\csromanspacer{} โผฏฺฐพฺพายตนํ กายวิญฺญาณสฺส ฯเปฯ.\csromanspacer{} \textbf{วตฺถุปุเรชาตํ}\csromandash{}จกฺขายตนํ จกฺขุวิญฺญาณสฺส ฯเปฯ.\csromanspacer{} กายายตนํ กายวิญฺญาณสฺส ฯเปฯ.\csromanspacer{} วตฺถุ อรูปีนํ ขนฺธานํ ปุเรชาตปจฺจเยน ปจฺจโย. (1)}

\csromanheader{2}{\textbf{ปจฺฉาชาต}}
\proseitem{118}{\csromanfolio{130}อรูปี ธมฺโม รูปิสฺส ธมฺมสฺส ปจฺฉาชาตปจฺจเยน ปจฺจโย.\csromanspacer{} ปจฺฉาชาตา อรูปิโน ขนฺธา ปุเรชาตสฺส อิมสฺส กายสฺส ปจฺฉาชาตปจฺจเยน ปจฺจโย. (1)}

\csromanheader{2}{\textbf{อาเสวน}}
\proseitem{119}{อรูปี ธมฺโม อรูปิสฺส ธมฺมสฺส อาเสวนปจฺจเยน ปจฺจโย.\csromanspacer{} ปุริมา ปุริมา อรูปิโน ขนฺธา ปจฺฉิมานํ ปจฺฉิมานํ อรูปีนํ ขนฺธานํ อาเสวน ปจฺจเยน ปจฺจโย. (1)}

\csromanheader{2}{\textbf{กมฺม}}
\proseitem{120}{อรูปี ธมฺโม อรูปิสฺส ธมฺมสฺส กมฺมปจฺจเยน ปจฺจโย.\csromanspacer{} \textbf{สหชาตา}, นานากฺขณิกา.\csromanspacer{}\csromanspacer{}\textbf{สหชาตา} อรูปี เจตนา สมฺปยุตฺตกานํ ขนฺธานํ กมฺมปจฺจเยน ปจฺจโย.\csromanspacer{} ปฏิสนฺธิกฺขเณ ฯเปฯ.\csromanspacer{} \textbf{นานากฺขณิกา} อรูปี เจตนา วิปากานํ ขนฺธานํ กมฺมปจฺจเยน ปจฺจโย. (1)}

\prose{อรูปี ธมฺโม รูปิสฺส ธมฺมสฺส กมฺมปจฺจเยน ปจฺจโย.\csromanspacer{} \textbf{สหชาตา}, นานากฺขณิกา.\csromanspacer{}\csromanspacer{}\textbf{สหชาตา} อรูปี เจตนา จิตฺตสมุฏฺฐานานํ รูปานํ กมฺมปจฺจเยน ปจฺจโย.\csromanspacer{} ปฏิสนฺธิกฺขเณ ฯเปฯ.\csromanspacer{} \textbf{นานากฺขณิกา} อรูปี เจตนา กฏตฺตารูปานํ กมฺมปจฺจเยน ปจฺจโย. (2)}

\prose{อรูปี ธมฺโม รูปิสฺส จ อรูปิสฺส จ ธมฺมสฺส กมฺมปจฺจเยน ปจฺจโย.\csromanspacer{} \textbf{สหชาตา}, นานากฺขณิกา.\csromanspacer{}\csromanspacer{}\textbf{สหชาตา} อรูปี เจตนา สมฺปยุตฺตกานํ ขนฺธานํ จิตฺตสมุฏฺฐานานญฺจ รูปานํ กมฺมปจฺจเยน ปจฺจโย.\csromanspacer{} ปฏิสนฺธิกฺขเณ ฯเปฯ.\csromanspacer{} \textbf{นานากฺขณิกา} อรูปี เจตนา วิปากานํ ขนฺธานํ กฏตฺตา จ รูปานํ กมฺมปจฺจเยน ปจฺจโย. (3)}

\csromanheader{2}{\textbf{วิปากาหาร}}
\proseitem{121}{อรูปี ธมฺโม อรูปิสฺส ธมฺมสฺส วิปากปจฺจเยน ปจฺจโย.\csromanspacer{} ตีณิ.}

\prose{รูปี ธมฺโม รูปิสฺส ธมฺมสฺส อาหารปจฺจเยน ปจฺจโย.\csromanspacer{} กพฬีกาโร อาหาโร อิมสฺส กายสฺส อาหารปจฺจเยน ปจฺจโย. (1)}

\prose{อรูปี ธมฺโม อรูปิสฺส ธมฺมสฺส อาหารปจฺจเยน ปจฺจโย.\csromanspacer{} ตีณิ.}

\csromanheader{2}{11. รูปีทุก \textbf{อินฺทฺริย}}
\proseitem{122}{\csromanfolio{131}รูปี ธมฺโม รูปิสฺส ธมฺมสฺส อินฺทฺริยปจฺจเยน ปจฺจโย.\csromanspacer{} รูปชีวิตินฺทฺริยํ กฏตฺตารูปานํ อินฺทฺริยปจฺจเยน ปจฺจโย. (1)}

\prose{รูปี ธมฺโม อรูปิสฺส ธมฺมสฺส อินฺทฺริยปจฺจเยน ปจฺจโย.\csromanspacer{} จกฺขุนฺทฺริยํ ฯเปฯ.\csromanspacer{} กายินฺทฺริยํ กายวิญฺญาณสฺส อินฺทฺริยปจฺจเยน ปจฺจโย. (2)}

\prose{อรูปี ธมฺโม อรูปิสฺส ธมฺมสฺส อินฺทฺริยปจฺจเยน ปจฺจโย.\csromanspacer{} ตีณิ.}

\prose{รูปี จ อรูปี จ ธมฺมา อรูปิสฺส ธมฺมสฺส อินฺทฺริยปจฺจเยน ปจฺจโย.\csromanspacer{} จกฺขุนฺทฺริยญฺจ จกฺขุวิญฺญาณญฺจ จกฺขุวิญฺญาณสหคตานํ ขนฺธานํ อินฺทฺริยปจฺจเยน ปจฺจโย ฯเปฯ.\csromanspacer{} กายินฺทฺริยญฺจ ฯเปฯ.}

\csromanheader{2}{\textbf{ฌานาทิ}}
\proseitem{123}{อรูปี ธมฺโม อรูปิสฺส ธมฺมสฺส ฌานปจฺจเยน ปจฺจโย.\csromanspacer{} ตีณิ.\csromanspacer{} มคฺคปจฺจเยน ปจฺจโย.\csromanspacer{} ตีณิ.\csromanspacer{} สมฺปยุตฺตปจฺจเยน ปจฺจโย.\csromanspacer{} เอกํ.}

\csromanheader{2}{\textbf{วิปฺปยุตฺต}}
\proseitem{124}{รูปี ธมฺโม อรูปิสฺส ธมฺมสฺส วิปฺปยุตฺตปจฺจเยน ปจฺจโย.\csromanspacer{} \textbf{สหชาตํ}, ปุเรชาตํ.\csromanspacer{}\csromanspacer{}\textbf{สหชาตํ} ปฏิสนฺธิกฺขเณ วตฺถุ อรูปีนํ ขนฺธานํ วิปฺปยุตฺตปจฺจเยน ปจฺจโย.\csromanspacer{} \textbf{ปุเรชาตํ} จกฺขายตนํ จกฺขุวิญฺญาณสฺส ฯเปฯ กายายตนํ กายวิญฺญาณสฺส ฯเปฯ.\csromanspacer{} วตฺถุ อรูปีนํ ขนฺธานํ วิปฺปยุตฺตปจฺจเยน ปจฺจโย. (1)}

\prose{อรูปี ธมฺโม รูปิสฺส ธมฺมสฺส วิปฺปยุตฺตปจฺจเยน ปจฺจโย.\csromanspacer{} \textbf{สหชาตํ} ปจฺฉาชาตํ.\csromanspacer{}\csromanspacer{}\textbf{สหชาตา} อรูปิโน ขนฺธา จิตฺตสมุฏฺฐานานํ รูปานํ วิปฺปยุตฺตปจฺจเยน ปจฺจโย.\csromanspacer{} ปฏิสนฺธิกฺขเณ อรูปิโน ขนฺธา กฏตฺตารูปานํ วิปฺปยุตฺตปจฺจเยน ปจฺจโย.\csromanspacer{} อรูปิโน ขนฺธา วตฺถุสฺส วิปฺปยุตฺตปจฺจเยน ปจฺจโย.\csromanspacer{} \textbf{ปจฺฉาชาตา} อรูปิโน ขนฺธา ปุเรชาตสฺส อิมสฺส กายสฺส วิปฺปยุตฺตปจฺจเยน ปจฺจโย. (1)}

\csromanheader{2}{\textbf{อตฺถิ-อาทิ}}
\proseitem{125}{\csromanfolio{132}รูปี ธมฺโม รูปิสฺส ธมฺมสฺส อตฺถิปจฺจเยน ปจฺจโย.\csromanspacer{} \textbf{สหชาตํ}, อาหารํ, อินฺทฺริยํ.\csromanspacer{}\csromanspacer{}\textbf{สหชาตํ} เอกํ มหาภูตํ ฯเปฯ. (ยาว อสญฺญสตฺตา.) \textbf{กพฬีกาโร อาหาโร} อิมสฺส กายสฺส อตฺถิปจฺจเยน ปจฺจโย.\csromanspacer{} \textbf{รูปชีวินฺทฺริยํ} กฏตฺตารูปานํ อตฺถิปจฺจเยน ปจฺจโย. (1)}

\prose{รูปี ธมฺโม อรูปิสฺส ธมฺมสฺส อตฺถิปจฺจเยน ปจฺจโย.\csromanspacer{} \textbf{สหชาตํ}, ปุเรชาตํ.\csromanspacer{}\csromanspacer{}\textbf{สหชาตํ} ปฏิสนฺธิกฺขเณ วตฺถุ อรูปีนํ ขนฺธานํ อตฺถิปจฺจเยน ปจฺจโย.\csromanspacer{} \textbf{ปุเรชาตํ} จกฺขุํ ฯเปฯ.\csromanspacer{} กพฬีการํ อาหารํ อนิจฺจโต ฯเปฯ โทมนสฺสํ อุปฺปชฺชติ.\csromanspacer{} ทิพฺเพน จกฺขุนา รูปํ ปสฺสติ.\csromanspacer{} ทิพฺพาย โสตธาตุยา สทฺทํ สุณาติ.\csromanspacer{} รูปายตนํ จกฺขุวิญฺญาณสฺส ฯเปฯ.\csromanspacer{} โผฏฺฐพฺพายตนํ กายวิญฺญาณสฺส ฯเปฯ.\csromanspacer{} จกฺขายตนํ จกฺขุวิญฺญาณสฺส ฯเปฯ.\csromanspacer{} กายายตนํ กายวิญฺญาณสฺส ฯเปฯ.\csromanspacer{} วตฺถุ อรูปีนํ ขนฺธานํ อตฺถิปจฺจเยน ปจฺจโย. (2)}

\proseitem{126}{อรูปี ธมฺโม อรูปิสฺส ธมฺมสฺส อตฺถิปจฺจเยน ปจฺจโย.\csromanspacer{} อรูปี เอโก ขนฺโธ ติณฺณนฺนํ ขนฺธานํ ฯเปฯ เทฺว ขนฺธา ฯเปฯ.\csromanspacer{} ปฏิสนฺธิกฺขเณ ฯเปฯ. (1)}

\prose{อรูปี ธมฺโม รูปิสฺส ธมฺมสฺส อตฺถิปจฺจเยน ปจฺจโย.\csromanspacer{} \textbf{สหชาตํ}, ปจฺฉาชาตํ.\csromanspacer{}\csromanspacer{}\textbf{สหชาตา} อรูปิโน ขนฺธา จิตฺตสมุฏฺฐานานํ รูปานํ อตฺถิปจฺจเยน ปจฺจโย.\csromanspacer{} ปฏิสนฺธิกฺขเณ ฯเปฯ.\csromanspacer{} \textbf{ปจฺฉาชาตา} อรูปิโน ขนฺธา ปุเรชาตสฺส อิมสฺส กายสฺส อตฺถิปจฺจเยน ปจฺจโย. (2)}

\prose{อรูปี ธมฺโม รูปิสฺส จ อรูปิสฺส จ ธมฺมสฺส อตฺถิปจฺจเยน ปจฺจโย.\csromanspacer{} อรูปี เอโก ขนฺโธ ติณฺณนฺนํ ขนฺธานํ จิตฺตสมุฏฺฐานานญฺจ รูปานํ อตฺถิปจฺจเยน ปจฺจโย ฯเปฯ เทฺว ขนฺธา ฯเปฯ.\csromanspacer{} ปฏิสนฺธิกฺขเณ ฯเปฯ. (3)}

\proseitem{127}{รูปี จ อรูปี จ ธมฺมา รูปิสฺส ธมฺมสฺส อตฺถิปจฺจเยน ปจฺจโย.\csromanspacer{} \textbf{สหชาตํ}, ปจฺฉาชาตํ, อาหารํ, อินฺทฺริยํ.\csromanspacer{}\csromanspacer{}\textbf{สหชาตา} อรูปี ขนฺธา จ มหาภูตา จ จิตฺตสมุฏฺฐานานํ รูปานํ อตฺถิปจฺจเยน ปจฺจโย.\csromanspacer{} ปฏิสนฺธิกฺขเณ ฯเปฯ.\csromanspacer{} \textbf{ปจฺฉาชาตา} อรูปิโน ขนฺธา จ กพฬีกาโร อาหาโร จ อิมสฺส กายสฺส อตฺถิปจฺจเยน ปจฺจโย.\csromanspacer{} \textbf{ปจฺฉาชาตา} อรูปิโน ขนฺธา จ รูปชีวิตินฺทฺริยญฺจ กฏตฺตารูปานํ อตฺถิปจฺจเยน ปจฺจโย. (1)}

\csromanheader{2}{11. รูปีทุก}
\prose{\csromanfolio{133}รูปี จ อรูปี จ ธมฺมา อรูปิสฺส ธมฺมสฺส อตฺถิปจฺจเยน ปจฺจโย.\csromanspacer{} \textbf{สหชาตํ, ปุเรชาตํ}.\csromanspacer{}\csromanspacer{}สหชาโต จกฺขุวิญฺญาณสหคโต เอโก ขนฺโธ จ จกฺขายตนญฺจ ติณฺณนฺนํ ขนฺธานํ อตฺถิปจฺจเยน ปจฺจโย ฯเปฯ เทฺว ขนฺธา จ ฯเปฯ.\csromanspacer{} กายวิญฺญาณสหคโต ฯเปฯ.\csromanspacer{} อรูปี เอโก ขนฺโธ จ วตฺถุ จ ติณฺณนฺนํ ขนฺธานํ อตฺถิปจฺจเยน ปจฺจโย ฯเปฯ เทฺว ขนฺธา จ ฯเปฯ.\csromanspacer{} ปฏิสนฺธิกฺขเณ อรูปี เอโก ขนฺโธ จ วตฺถุ จ ติณฺณนฺนํ ขนฺธานํ ฯเปฯ เทฺว ขนฺธา จ ฯเปฯ. (2)}

\prose{นตฺถิปจฺจเยน ปจฺจโย.\csromanspacer{} วิคตปจฺจเยน ปจฺจโย.\csromanspacer{} อวิคตปจฺจเยน ปจฺจโย.}

\csromanheader{2}{1. ปจฺจยานุโลม 2. สงฺขฺยาวาร}
\csromanheader{2}{\textbf{สุทฺธ}}
\proseitem{128}{เหตุยา ตีณิ, อารมฺมเณ เทฺว, อธิปติยา จตฺตาริ, อนนฺตเร เอกํ, สมนนฺตเร เอกํ, สหชาเต สตฺต, อญฺญมญฺเญ ฉ, นิสฺสเย สตฺต, อุปนิสฺสเย เทฺว, ปุเรชาเต เอกํ, ปจฺฉาชาเต เอกํ, อาเสวเน เอกํ, กมฺเม ตีณิ, วิปาเก ตีณิ, อาหาเร จตฺตาริ, อินฺทฺริเย ฉ, ฌาเน ตีณิ, มคฺเค ตีณิ, สมฺปยุตฺเต เอกํ, วิปฺปยุตฺเต เทฺว, อตฺถิยา สตฺต, นตฺถิยา เอกํ, วิคเต เอกํ, อวิคเต สตฺต.}

\vspace{\tipitakaparskip}
\csromancenter{อนุโลมํ.}
\csromanheader{2}{2. ปจฺจนียุทฺธาร}
\proseitem{129}{รูปี ธมฺโม รูปิสฺส ธมฺมสฺส สหชาตปจฺจเยน ปจฺจโย.\csromanspacer{} อาหารปจฺจเยน ปจฺจโย.\csromanspacer{} อินฺทฺริยปจฺจเยน ปจฺจโย. (1)}

\prose{รูปี ธมฺโม อรูปิสฺส ธมฺมสฺส อารมฺมณปจฺจเยน ปจฺจโย.\csromanspacer{} สหชาตปจฺจเยน ปจฺจโย.\csromanspacer{} อุปนิสฺสยปจฺจเยน ปจฺจโย.\csromanspacer{} ปุเรชาตปจฺจเยน ปจฺจโย. (2)}

\prose{อรูปี ธมฺโม อรูปิสฺส ธมฺมสฺส อารมฺมณปจฺจเยน ปจฺจโย.\csromanspacer{} สหชาตปจฺจเยน ปจฺจโย.\csromanspacer{} อุปนิสฺสยปจฺจเยน ปจฺจโย.\csromanspacer{} กมฺมปจฺจเยน ปจฺจโย. (1)}

\prose{\csromanfolio{134}อรูปี ธมฺโม รูปิสฺส ธมฺมสฺส สหชาตปจฺจเยน ปจฺจโย.\csromanspacer{} ปจฺฉาชาตปจฺจเยน ปจฺจโย.\csromanspacer{} กมฺมปจฺจเยน ปจฺจโย. (2)}

\prose{อรูปี ธมฺโม รูปิสฺส จ อรูปิสฺส จ ธมฺมสฺส สหชาตปจฺจเยน ปจฺจโย.\csromanspacer{} กมฺมปจฺจเยน ปจฺจโย. (3)}

\prose{รูปี จ อรูปี จ ธมฺมา รูปิสฺส ธมฺมสฺส สหชาตํ, ปจฺฉาชาตํ, อาหารํ, อินฺทฺริยํ. (1)}

\prose{รูปี จ อรูปี จ ธมฺมา อรูปิสฺส ธมฺมสฺส สหชาตํ, ปุเรชาตํ.}

\csromanheader{2}{2. ปจฺจยปจฺจนีย 2. สงฺขฺยาวาร}
\csromanheader{2}{\textbf{สุทฺธ}}
\proseitem{130}{นเหตุยา สตฺต, น-อารมฺมเณ สตฺต, น-อธิปติยา สตฺต น-อนนฺตเร สตฺต, นสมนนฺตเร สตฺต, นสหชาเต ฉ, น-อญฺญมญฺเญ ฉ, นนิสฺสเย ฉ, น-อุปนิสฺสเย สตฺต, นปุเรชาเต สตฺต ฯเปฯ นมคฺเค สตฺต, นสมฺปยุตฺเต ฉ, นวิปฺปยุตฺเต ปญฺจ, โน-อตฺถิยา จตฺตาริ, โนนตฺถิยา สตฺต, โนวิคเต สตฺต, โน-อวิคเต จตฺตาริ.}

\vspace{\tipitakaparskip}
\csromancenter{ปจฺจนียํ.}
\csromanheader{2}{3. ปจฺจยานุโลมปจฺจนีย}
\csromanheader{2}{\textbf{เหตุทุก}}
\proseitem{131}{เหตุปจฺจยา น-อารมฺมเณ ตีณิ, น-อธิปติยา ตีณิ, น-อนนฺตเร ตีณิ, นสมนนฺตเร ตีณิ, น-อญฺญมญฺเญ เอกํ, น-อุปนิสฺสเย ตีณิ, (สพฺพตฺถ ตีณิ.) นสมฺปยุตฺเต เอกํ, นวิปฺปยุตฺเต เอกํ, โนนตฺถิยา ตีณิ, โนวิคเต ตีณิ.}

\vspace{\tipitakaparskip}
\csromancenter{อนุโลมปจฺจนียํ.}
\csromanmatikaanchor{mtk.15}
\csromantocmarkref{section}{12. โลกิยทุก}{mtk.15}
\bookmark[dest={mtk.15},level=1]{12. โลกิยทุก}
\csromanheader{1}{12. โลกิยทุก \textbf{4. ปจฺจยปจฺจนียานุโลม}}
\csromanheader{2}{\textbf{นเหตุทุก}}
\proseitem{132}{\csromanfolio{135}นเหตุปจฺจยา อารมฺมเณ เทฺว, อธิปติยา จตฺตาริ, (อนุโลมมาติกา กาตพฺพา.) อวิคเต สตฺต.}

\vspace{\tipitakaparskip}
\csromancenter{ปจฺจนียานุโลมํ.}
\nitthitam{รูปีทุกํ นิฏฺฐิตํ.}
\csromanheader{2}{12. โลกิยทุก 1. ปฏิจฺจวาร}
\csromanheader{2}{1. ปจฺจยานุโลม 1. วิภงฺควาร}
\csromanheader{2}{\textbf{เหตุ}}
\proseitem{133}{โลกิยํ ธมฺมํ ปฏิจฺจ โลกิโย ธมฺโม อุปฺปชฺชติ เหตุปจฺจยา.\csromanspacer{} โลกิยํ เอกํ ขนฺธํ ปฏิจฺจ ตโย ขนฺธา จิตฺตสมุฏฺฐานญฺจ รูปํ ฯเปฯ เทฺว ขนฺเธ ฯเปฯ.\csromanspacer{} ปฏิสนฺธิกฺขเณ ฯเปฯ.\csromanspacer{} ขนฺเธ ปฏิจฺจ วตฺถุ, วตฺถุํ ปฏิจฺจ ขนฺธา.\csromanspacer{} เอกํ มหาภูตํ ฯเปฯ มหาภูเต ปฏิจฺจ จิตฺตสมุฏฺฐานํ รูปํ, กฏตฺตารูปํ, อุปาทารูปํ. (1)}

\proseitem{134}{โลกุตฺตรํ ธมฺมํ ปฏิจฺจ โลกุตฺตโร ธมฺโม อุปฺปชฺชติ เหตุปจฺจยา.\csromanspacer{} โลกุตฺตรํ เอกํ ขนฺธํ ปฏิจฺจ ตโย ขนฺธา ฯเปฯ เทฺว ขนฺเธ ฯเปฯ. (1)}

\prose{โลกุตฺตรํ ธมฺมํ ปฏิจฺจ โลกิโย ธมฺโม อุปฺปชฺชติ เหตุปจฺจยา.\csromanspacer{} โลกุตฺตเร ขนฺเธ ปฏิจฺจ จิตฺตสมุฏฺฐานํ รูปํ. (2)}

\prose{โลกุตฺตรํ ธมฺมํ ปฏิจฺจ โลกิโย จ โลกุตฺตโร จ ธมฺมา อุปฺปชฺชนฺติ เหตุปจฺจยา.\csromanspacer{} โลกุตฺตรํ เอกํ ขนฺธํ ปฏิจฺจ ตโย ขนฺธา จิตฺตสมุฏฺฐานญฺจ รูปํ ฯเปฯ เทฺว ขนฺเธ ฯเปฯ. (3)}

\prose{โลกิยญฺจ โลกุตฺตรญฺจ ธมฺมํ ปฏิจฺจ โลกิโย ธมฺโม อุปฺปชฺชติ เหตุปจฺจยา.\csromanspacer{} โลกุตฺตเร ขนฺเธ จ มหาภูเต จ ปฏิจฺจ จิตฺตสมุฏฺฐานํ รูปํ. (สํขิตฺตํ.) (1)}

\csromanheader{2}{1. ปจฺจยานุโลม 2. สงฺขฺยาวาร}
\csromanheader{2}{\textbf{สุทฺธ}}
\proseitem{135}{\csromanfolio{136}เหตุยา ปญฺจ, อารมฺมเณ, เทฺว, อธิปติยา ปญฺจ, อนนฺตเร เทฺว, สมนนฺตเร เทฺว, สหชาเต ปญฺจ, อญฺญมญฺเญ เทฺว, นิสฺสเย ปญฺจ, อุปนิสฺสเย เทฺว, ปุเรชาเต เทฺว, อาเสวเน เทฺว, กมฺเม ปญฺจ, วิปาเก ปญฺจ, อาหาเร ปญฺจ, อินฺทฺริเย ปญฺจ, ฌาเน ปญฺจ, มคฺเค ปญฺจ, สมฺปยุตฺเต เทฺว, วิปฺปยุตฺเต ปญฺจ, อตฺถิยา ปญฺจ, นตฺถิยา เทฺว, วิคเต เทฺว, อวิคเต ปญฺจ.}

\vspace{\tipitakaparskip}
\csromancenter{อนุโลมํ.}
\csromanheader{2}{2. ปจฺจยปจฺจนีย 1. วิภงฺควาร}
\csromanheader{2}{\textbf{นเหตุ}}
\proseitem{136}{โลกิยํ ธมฺมํ ปฏิจฺจ โลกิโย ธมฺโม อุปฺปชฺชติ นเหตุปจฺจยา.\csromanspacer{} อเหตุกํ โลกิยํ เอกํ ขนฺธํ ปฏิจฺจ ตโย ขนฺธา จิตฺตสมุฏฺฐานญฺจ รูปํ ฯเปฯ เทฺว ขนฺเธ ฯเปฯ.\csromanspacer{} อเหตุกปฏิสนฺธิกฺขเณ ฯเปฯ. (ยาว อสญฺญสตฺตา.) วิจิกิจฺฉาสหคเต อุทฺธจฺจสหคเต ขนฺเธ ปฏิจฺจ วิจิกิจฺฉาสหคโต อุทฺธจฺจสหคโต โมโห. (สํขิตฺตํ.)}

\csromanheader{2}{2. ปจฺจยปจฺจนีย 2. สงฺขฺยาวาร}
\csromanheader{2}{\textbf{สุทฺธ}}
\vspace{\tipitakaparskip}
\csromancenter{\textbf{นเหตุ}ยา เอกํ, น-อารมฺมเณ ตีณิ, น-อธิปติยา เทฺว, น-อนนฺตเร ตีณิ, นสมนนฺตเร ตีณิ, น-อญฺญมญฺเญ ตีณิ, น-อุปนิสฺสเย ตีณิ, นปุเรชาเต จตฺตาริ, นปจฺฉาชาเต ปญฺจ, น-อาเสวเน ปญฺจ, (น-อาเสวนมูลเก โลกุตฺตเร สุทฺธเก “วิปาโก”ติ นิยาเมตพฺพํ.\csromanspacer{} อวเสสา ปกติกาเยว.) นกมฺเม เทฺว, นวิปาเก ปญฺจ, น-อาหาเร เอกํ, น-อินฺทฺริเย เอกํ, นฌาเน เอกํ, นมคฺเค เอกํ, นสมฺปยุตฺเต ตีณิ, นวิปฺปยุตฺเต เทฺว, โนนตฺถิยา ตีณิ, โนวิคเต ตีณิ.}
\csromancenter{ปจฺจนียํ.}
\csromanheader{2}{12. โลกิยทุก \textbf{3. ปจฺจยานุโลมปจฺจนีย}}
\csromanheader{2}{\textbf{เหตุทุก}}
\proseitem{138}{\csromanfolio{137}\textbf{เหตุ}ปจฺจยา น-อารมฺมเณ ตีณิ, น-อธิปติยา เทฺว, (น-อนนฺตรปทาที ปจฺจนียสทิสา.) ฯเปฯ นวิปาเก ปญฺจ, นสมฺปยุตฺเต ตีณิ, นวิปฺปยุตฺเต เทฺว, โนนตฺถิยา ตีณิ, โนวิคเต ตีณิ.}

\vspace{\tipitakaparskip}
\csromancenter{อนุโลมปจฺจนียํ.}
\csromanheader{2}{4. ปจฺจยปจฺจนียานุโลม}
\csromanheader{2}{\textbf{นเหตุทุก}}
\vspace{\tipitakaparskip}
\csromancenter{นเหตุปจฺจยา อารมฺมเณ เอกํ, อนนฺตเร เอกํ ฯเปฯ อวิคเต เอกํ.}
\csromancenter{ปจฺจนียานุโลมํ.}
\csromancenter{(สหชาตวาโร ปฏิจฺจวารสทิโส.)}
\csromanheader{2}{12. โลกิยทุก 3. ปจฺจยวาร}
\csromanheader{2}{1. ปจฺจยานุโลม 1. วิภงฺควาร}
\csromanheader{2}{\textbf{เหตุ}}
\proseitem{140}{โลกิยํ ธมฺมํ ปจฺจยา โลกิโย ธมฺโม อุปฺปชฺชติ เหตุปจฺจยา.\csromanspacer{} โลกิยํ เอกํ ขนฺธํ ปจฺจยา ฯเปฯ.\csromanspacer{} เอกํ มหาภูตํ ฯเปฯ มหาภูเต ปจฺจยา จิตฺตสมุฏฺฐานํ รูปํ, กฏตฺตารูปํ, อุปาทารูปํ.\csromanspacer{} วตฺถุํ ปจฺจยา โลกิยา ขนฺธา. (1)}

\prose{โลกิยํ ธมฺมํ ปจฺจยา โลกุตฺตโร ธมฺโม อุปฺปชฺชติ เหตุปจฺจยา.\csromanspacer{} วตฺถุํ ปจฺจยา โลกุตฺตรา ขนฺธา. (2)}

\prose{\csromanfolio{138}โลกิยํ ธมฺมํ ปจฺจยา โลกิโย จ โลกุตฺตโร จ ธมฺมา อุปฺปชฺชนฺติ เหตุปจฺจยา.\csromanspacer{} วตฺถุํ ปจฺจยา โลกุตฺตรา ขนฺธา.\csromanspacer{} มหาภูเต ปจฺจยา จิตฺตสมุฏฺฐานํ รูปํ. (3)}

\proseitem{141}{โลกุตฺตรํ ธมฺมํ ปจฺจยา โลกุตฺตโร ธมฺโม อุปฺปชฺชติ เหตุปจฺจยา.\csromanspacer{} ตีณิ.}

\prose{โลกิยญฺจ โลกุตฺตรญฺจ ธมฺมํ ปจฺจยา โลกิโย ธมฺโม อุปฺปชฺชติ เหตุปจฺจยา.\csromanspacer{} โลกุตฺตเร ขนฺเธ จ มหาภูเต จ ปจฺจยา จิตฺตสมุฏฺฐานํ รูปํ. (1)}

\prose{โลกิยญฺจ โลกุตฺตรญฺจ ธมฺมํ ปจฺจยา โลกุตฺตโร ธมฺโม อุปฺปชฺชติ เหตุปจฺจยา.\csromanspacer{} โลกุตฺตรํ เอกํ ขนฺธญฺจ วตฺถุญฺจ ปจฺจยา ตโย ขนฺธา ฯเปฯ เทฺว ขนฺเธ ฯเปฯ. (2)}

\prose{โลกิยญฺจ โลกุตฺตรญฺจ ธมฺมํ ปจฺจยา โลกิโย จ โลกุตฺตโร จ ธมฺมา อุปฺปชฺชนฺติ เหตุปจฺจยา.\csromanspacer{} โลกุตฺตรํ เอกํ ขนฺธญฺจ วตฺถุญฺจ ปจฺจยา ตโย ขนฺธา ฯเปฯ เทฺว ขนฺเธ ฯเปฯ.\csromanspacer{} โลกุตฺตเร ขนฺเธ จ มหาภูเต จ ปจฺจยา จิตฺตสมุฏฺฐานํ รูปํ. (สํขิตฺตํ.) (3)}

\csromanheader{2}{1. ปจฺจยานุโลม 2. สงฺขฺยาวาร}
\csromanheader{2}{\textbf{สุทฺธ}}
\proseitem{142}{เหตุยา นว, อารมฺมเณ จตฺตาริ, อธิปติยา นว, อนนฺตเร จตฺตาริ, สมนนฺตเร จตฺตาริ, สหชาเต นว, อญฺญมญฺเญ จตฺตาริ, นิสฺสเย นว, อุปนิสฺสเย จตฺตาริ, ปุเรชาเต จตฺตาริ, อาเสวเน จตฺตาริ, กมฺเม นว, วิปาเก นว ฯเปฯ มคฺเค นว, สมฺปยุตฺเต จตฺตาริ วิปฺปยุตฺเต นว, อตฺถิยา นว, นตฺถิยา จตฺตาริ, วิคเต จตฺตาริ, อวิคเต นว.}

\vspace{\tipitakaparskip}
\csromancenter{อนุโลมํ.}
\csromanheader{2}{12. โลกิยทุก \textbf{2. ปจฺจยปจฺจนีย 1. วิภงฺควาร}}
\csromanheader{2}{\textbf{นเหตุ}}
\proseitem{143}{\csromanfolio{139}โลกิยํ ธมฺมํ ปจฺจยา โลกิโย ธมฺโม อุปฺปชฺชติ นเหตุปจฺจยา.\csromanspacer{} อเหตุกํ โลกิยํ เอกํ ขนฺธํ ฯเปฯ. (ยาว อสญฺญสตฺตา.) จกฺขายตนํ ปจฺจยา จกฺขุวิญฺญาณํ ฯเปฯ.\csromanspacer{} กายายตนํ ปจฺจยา กายวิญฺญาณํ.\csromanspacer{} วตฺถุํ ปจฺจยา อเหตุกา โลกิยา ขนฺธา.\csromanspacer{} วิจิกิจฺฉาสหคเต อุทฺธจฺจสหคเต ขนฺเธ จ วตฺถุญฺจ ปจฺจยา วิจิกิจฺฉาสหคโต อุทฺธจฺจสหคโต โมโห. (สํขิตฺตํ.)}

\csromanheader{2}{2. ปจฺจยปจฺจนีย 2. สงฺขฺยาวาร}
\csromanheader{2}{\textbf{สุทฺธ}}
\vspace{\tipitakaparskip}
\csromancenter{\textbf{นเหตุ}ยา เอกํ, น-อารมฺมเณ ตีณิ, น-อธิปติยา จตฺตาริ, น-อนนฺตเร ตีณิ ฯเปฯ น-อุปนิสฺสเย ตีณิ, นปุเรชาเต จตฺตาริ, นปจฺฉาชาเต นว, น-อาเสวเน นว, (โลกุตฺตเร อรูเป วิปากนฺติ นิยาเมตพฺพํ.) นกมฺเม จตฺตาริ, นวิปาเก นว, น-อาหาเร เอกํ, น-อินฺทฺริเย เอกํ, นฌาเน เอกํ, นมคฺเค เอกํ, นสมฺปยุตฺเต ตีณิ, นวิปฺปยุตฺเต เทฺว, โนนตฺถิยา ตีณิ, โนวิคเต ตีณิ.}
\csromancenter{ปจฺจนียํ.}
\csromanheader{2}{3. ปจฺจยานุโลมปจฺจนีย}
\csromanheader{2}{\textbf{เหตุทุก}}
\proseitem{145}{เหตุปจฺจโย น-อารมฺมเณ ตีณิ, น-อธิปติยา จตฺตาริ, น-อนนฺตเร ตีณิ. (นสมนนฺตรปทาที ปจฺจนียสทิสา.) นวิปาเก นว, นสมฺปยุตฺเต ตีณิ, นวิปฺปยุตฺเต เทฺว, โนนตฺถิยา ตีณิ, โนวิคเต ตีณิ.}

\vspace{\tipitakaparskip}
\csromancenter{อนุโลมปจฺจนียํ.}
\csromanheader{2}{4. ปจฺจยปจฺจนียานุโลม}
\csromanheader{2}{\textbf{นเหตุทุก}}
\vspace{\tipitakaparskip}
\csromancenter{นเหตุปจฺจยา อารมฺมเณ เอกํ, อนนฺตเร เอกํ ฯเปฯ อวิคเต เอกํ.}
\csromancenter{ปจฺจนียานุโลมํ.}
\csromancenter{(นิสฺสยวาโร ปจฺจยวารสทิโส)}
\csromanheader{2}{12. โลกิยทุก 5. สํสฏฺฐวาร}
\csromanheader{2}{1. ปจฺจยานุโลม 1. วิภงฺควาร}
\csromanheader{2}{\textbf{เหตุ}}
\proseitem{147}{\csromanfolio{140}โลกิยํ ธมฺมํ สํสฏฺโฐ โลกิโย ธมฺโม อุปฺปชฺชติ เหตุปจฺจยา.\csromanspacer{} โลกิยํ เอกํ ขนฺธํ สํสฏฺฐา ตโย ขนฺธา ฯเปฯ เทฺว ขนฺเธ ฯเปฯ.\csromanspacer{} ปฏิสนฺธิกฺขเณ ฯเปฯ. (1)}

\prose{โลกุตฺตรํ ธมฺมํ สํสฏฺโฐ โลกุตฺตโร ธมฺโม อุปฺปชฺชติ เหตุปจฺจยา.\csromanspacer{} โลกุตฺตรํ เอกํ ขนฺธํ สํสฏฺฐา ตโย ขนฺธา ฯเปฯ เทฺว ขนฺเธ ฯเปฯ. (1)}

\prose{(สํสฏฺฐวาโร เอวํ วิตฺถาเรตพฺโพ, สห คณนาหิ เทฺว ปญฺหา.)}

\vspace{\tipitakaparskip}
\csromancenter{(สมฺปยุตฺตวาโร สํสฏฺฐวารสทิโส.)}
\csromanheader{2}{12. โลกิยทุก 7. ปญฺหาวาร}
\csromanheader{2}{1. ปจฺจยานุโลม 1. วิภงฺควาร}
\csromanheader{2}{\textbf{เหตุ}}
\proseitem{148}{โลกิโย ธมฺโม โลกิยสฺส ธมฺมสฺส เหตุปจฺจเยน ปจฺจโย.\csromanspacer{} โลกิยา เหตู สมฺปยุตฺตกานํ ขนฺธานํ จิตฺตสมุฏฺฐานานญฺจ รูปานํ เหตุปจฺจเยน ปจฺจโย.\csromanspacer{} ปฏิสนฺธิกฺขเณ ฯเปฯ. (1)}

\csromanheader{2}{12. โลกิยทุก}
\vspace{\tipitakaparskip}
\csromancenter{โลกุตฺตโร ธมฺโม โลกุตฺตรสฺส ธมฺมสฺส เหตุปจฺจเยน ปจฺจโย.\csromanspacer{} ตีณิ.}
\csromanheader{2}{\textbf{อารมฺมณ}}
\proseitem{149}{\csromanfolio{141}โลกิโย ธมฺโม โลกิยสฺส ธมฺมสฺส อารมฺมณปจฺจเยน ปจฺจโย.\csromanspacer{} ทานํ ฯเปฯ สีลํ ฯเปฯ อุโปสถกมฺมํ กตฺวา ตํ ปจฺจเวกฺขติ.\csromanspacer{} ปุพฺเพ สุจิณฺณานิ ปจฺจเวกฺขติ.\csromanspacer{} ฌานา ฯเปฯ.\csromanspacer{} อริยา โคตฺรภุํ ปจฺจเวกฺขนฺติ.\csromanspacer{} โวทานํ ปจฺจเวกฺขนฺติ.\csromanspacer{} ปหีเน กิเลเส ปจฺจเวกฺขนฺติ.\csromanspacer{} วิกฺขมฺภิเต กิเลเส ฯเปฯ.\csromanspacer{} ปุพฺเพ สมุทาจิณฺเณ กิเลเส ชานนฺติ.\csromanspacer{} จกฺขุํ ฯเปฯ วตฺถุํ โลกิเย ขนฺเธ อนิจฺจโต ฯเปฯ โทมนสฺสํ อุปฺปชฺชติ.\csromanspacer{} ทิพฺเพน จกฺขุนา รูปํ ปสฺสติ.\csromanspacer{} ทิพฺพาย โสตธาตุยา สทฺทํ สุณาติ.\csromanspacer{} เจโตปริยญาเณน โลกิยจิตฺตสมงฺคิสฺส จิตฺตํ ชานาติ.\csromanspacer{} อากาสานญฺจายตนํ วิญฺญาณญฺจายตนสฺส ฯเปฯ.\csromanspacer{} อากิญฺจญฺญายตนํ เนวสญฺญานาสญฺญายตนสฺส ฯเปฯ.\csromanspacer{} รูปายตนํ จกฺขุวิญฺญาณสฺส ฯเปฯ.\csromanspacer{} โผฏฺฐพฺพายตนํ กายวิญฺญาณสฺส.\csromanspacer{} โลกิยา ขนฺธา อิทฺธิวิธญาณสฺส, เจโตปริยญาณสฺส, ปุพฺเพนิวาสานุสฺสติญาณสฺส, ยถากมฺมูปคญาณสฺส, อนาคตํสญาณสฺส, อาวชฺชนาย อารมฺมณปจฺจเยน ปจฺจโย. (1)}

\proseitem{150}{โลกุตฺตโร ธมฺโม โลกุตฺตรสฺส ธมฺมสฺส อารมฺมณปจฺจเยน ปจฺจโย.\csromanspacer{} นิพฺพานํ มคฺคสฺส, ผลสฺส อารมฺมณปจฺจเยน ปจฺจโย. (1)}

\prose{โลกุตฺตโร ธมฺโม โลกิยสฺส ธมฺมสฺส อารมฺมณปจฺจเยน ปจฺจโย.\csromanspacer{} อริยา มคฺคา วุฏฺฐหิตฺวา มคฺคํ ปจฺจเวกฺขนฺติ, ผลํ ปจฺจเวกฺขนฺติ, นิพฺพานํ ปจฺจเวกฺขนฺติ.\csromanspacer{} นิพฺพานํ โคตฺรภุสฺส, โวทานสฺส, อาวชฺชนาย, อารมฺมณปจฺจเยน ปจฺจโย.\csromanspacer{} อริยา เจโตปริยญาเณน โลกุตฺตรจิตฺตสมงฺคิสฺส จิตฺตํ ชานนฺติ.\csromanspacer{} โลกุตฺตรา ขนฺธา เจโตปริยญาณสฺส, ปุพฺเพนิวาสานุสฺสติญาณสฺส, อนาคตํสญาณสฺส, อาวชฺชนาย อารมฺมณปจฺจเยน ปจฺจโย. (2)}

\csromanheader{2}{\textbf{อธิปติ}}
\proseitem{151}{โลกิโย ธมฺโม โลกิยสฺส ธมฺมสฺส อธิปติปจฺจเยน ปจฺจโย.\csromanspacer{} \textbf{อารมฺมณาธิปติ}, สหชาตาธิปติ.\csromanspacer{}\csromanspacer{}\textbf{อารมฺมณาธิปติ}\csromandash{}ทานํ ทตฺวา, สีลํ ฯเปฯ อุโปสถกมฺมํ ฯเปฯ.\csromanspacer{} ปุพฺเพ ฯเปฯ.\csromanspacer{} ฌานา ฯเปฯ.\csromanspacer{} เสกฺขา โคตฺรภุํ ครุํ กตฺวา ปจฺจเวกฺขนฺติ, โวทานํ ครุํ กตฺวา ปจฺจเวกฺขนฺติ.\csromanspacer{} จกฺขุํ ฯเปฯ. \csromanfolio{142}วตฺถุํ โลกิเย ขนฺเธ ครุํ กตฺวา อสฺสาเทติ อภินนฺทติ, ตํ ครุํ กตฺวา ราโค อุปฺปชฺชติ, ทิฏฺฐิ อุปฺปชฺชติ.\csromanspacer{}\csromanspacer{}\textbf{สหชาตาธิปติ}\csromandash{}โลกิยาธิปติ สมฺปยุตฺตกานํ ขนฺธานํ จิตฺตสมุฏฺฐานานญฺจ รูปานํ อธิปติปจฺจเยน ปจฺจโย. (1)}

\proseitem{152}{โลกุตฺตโร ธมฺโม โลกุตฺตรสฺส ธมฺมสฺส อธิปติปจฺจเยน ปจฺจโย.\csromanspacer{} \textbf{อารมฺมณาธิปติ}, สหชาตาธิปติ.\csromanspacer{}\csromanspacer{}\textbf{อารมฺมณาธิปติ\csromandash{}} นิพฺพานํ มคฺคสฺส, ผลสฺส อธิปติปจฺจเยน ปจฺจโย.\csromanspacer{} \textbf{สหชาตาธิปติ}\csromandash{} โลกุตฺตราธิปติ สมฺปยุตฺตกานํ ขนฺธานํ อธิปติปจฺจเยน ปจฺจโย. (1)}

\prose{โลกุตฺตโร ธมฺโม โลกิยสฺส ธมฺมสฺส อธิปติปจฺจเยน ปจฺจโย อารมฺมณาธิปติ, สหชาตาธิปติ.\csromanspacer{}\csromanspacer{}\textbf{อารมฺมณาธิปติ\csromandash{}}อริยา มคฺคา วุฏฺฐหิตฺวา มคฺคํ ครุํ กตฺวา ปจฺจเวกฺขนฺติ, ผลํ ครุํ กตฺวา ปจฺจเวกฺขนฺติ, นิพฺพานํ ครุํ กตฺวา ปจฺจเวกฺขนฺติ.\csromanspacer{} นิพฺพานํ โคตฺรภุสฺส, โวทานสฺส อธิปติปจฺจเยน ปจฺจโย.\csromanspacer{} \textbf{สหชาตาธิปติ}\csromandash{} โลกุตฺตราธิปติ จิตฺตสมุฏฺฐานานํ รูปานํ อธิปติปจฺจเยน ปจฺจโย.}

\prose{โลกุตฺตโร ธมฺโม โลกิยสฺส จ โลกุตฺตรสฺส จ ธมฺมสฺส อธิปติปจฺจเยน ปจฺจโย.\csromanspacer{}\csromanspacer{}\textbf{สหชาตาธิปติ}\csromandash{}โลกุตฺตราธิปติ สมฺปยุตฺตกานํ ขนฺธานํ จิตฺตสมุฏฺฐานานญฺจ รูปานํ อธิปติปจฺจเยน ปจฺจโย. (3)}

\csromanheader{2}{\textbf{อนนฺตร}}
\proseitem{153}{โลกิโย ธมฺโม โลกิยสฺส ธมฺมสฺส อนนฺตรปจฺจเยน ปจฺจโย.\csromanspacer{} ปุริมา ปุริมา โลกิยา ขนฺธา ปจฺฉิมานํ ปจฺฉิมานํ โลกิยานํ ขนฺธานํ ฯเปฯ.\csromanspacer{} อนุโลมํ โคตฺรภุสฺส.\csromanspacer{} อนุโลมํ โวทานสฺส อนนฺตรปจฺจเยน ปจฺจโย. (1)}

\prose{โลกิโย ธมฺโม โลกุตฺตรสฺส ธมฺมสฺส อนนฺตรปจฺจเยน ปจฺจโย.\csromanspacer{} โคตฺรภุ มคฺคสฺส.\csromanspacer{} โวทานํ มคฺคสฺส.\csromanspacer{} อนุโลมํ ผลสมาปตฺติยา.\csromanspacer{} นิโรธา วุฏฺฐหนฺตสฺส เนวสญฺญานาสญฺญายตนํ ผลสมาปตฺติยา อนนฺตรปจฺจเยน ปจฺจโย. (2)}

\proseitem{154}{โลกุตฺตโร ธมฺโม โลกุตฺตรสฺส ธมฺมสฺส อนนฺตรปจฺจเยน ปจฺจโย.\csromanspacer{} ปุริมา ปุริมา โลกุตฺตรา ขนฺธา ปจฺฉิมานํ ปจฺฉิมานํ โลกุตฺตรานํ}

\vspace{\tipitakaparskip}
\csromancenter{โลกิยทุก ขนฺธานํ อนนฺตรปจฺจเยน ปจฺจโย.\csromanspacer{} มคฺโค ผลสฺส.\csromanspacer{} ผลํ ผลสฺส อนนฺตรปจฺจเยน ปจฺจโย. (1)}
\prose{\csromanfolio{143}โลกุตฺตโร ธมฺโม โลกิยสฺส ธมฺมสฺส อนนฺตรปจฺจเยน ปจฺจโย.\csromanspacer{} ผลํ วุฏฺฐานสฺส อนนฺตรปจฺจเยน ปจฺจโย. (2)}

\csromanheader{2}{\textbf{สมนนฺตราทิ}}
\proseitem{155}{โลกิโย ธมฺโม โลกิยสฺส ธมฺมสฺส สมนนฺตรปจฺจเยน ปจฺจโย.\csromanspacer{} สหชาตปจฺจเยน ปจฺจโย. (ปญฺจ ปญฺหา ฆฏนา นตฺถิ.) อญฺญมญฺญปจฺจเยน ปจฺจโย.\csromanspacer{} เทฺว.\csromanspacer{} นิสฺสยปจฺจเยน ปจฺจโย.}

\csromanheader{2}{\textbf{อุปนิสฺสย}}
\proseitem{156}{โลกิโย ธมฺโม โลกิยสฺส ธมฺมสฺส อุปนิสฺสยปจฺจเยน ปจฺจโย.\csromanspacer{} \textbf{อารมฺมณูปนิสฺสโย, อนนฺตรูปนิสฺสโย, ปกตูปนิสฺสโย ฯเปฯ.}\csromanspacer{} ปกตูปนิสฺสโย\csromandash{}โลกิยํ สทฺธํ อุปนิสฺสาย ทานํ เทติ ฯเปฯ วิปสฺสนํ อุปฺปาเทติ, อภิญฺญํ อุปฺปาเทติ, สมาปตฺติํ อุปฺปาเทติ.\csromanspacer{} มานํ ชปฺเปติ.\csromanspacer{} ทิฏฺฐิํ คณฺหาติ.\csromanspacer{} โลกิยํ สีลํ ฯเปฯ.\csromanspacer{} เสนาสนํ อุปนิสฺสาย ทานํ เทติ ฯเปฯ สํฆํ ภินฺทติ.\csromanspacer{} โลกิยา สทฺธา ฯเปฯ.\csromanspacer{} เสนาสนํ โลกิยาย สทฺธาย ฯเปฯ กายิกสฺส ทุกฺขสฺส อุปนิสฺสยปจฺจเยน ปจฺจโย.\csromanspacer{} กุสลากุสลํ กมฺมํ วิปากสฺส อุปนิสฺสยปจฺจเยน ปจฺจโย. (1)}

\prose{โลกิโย ธมฺโม โลกุตฺตรสฺส ธมฺมสฺส อุปนิสฺสยปจฺจเยน ปจฺจโย.\csromanspacer{} \textbf{อนนฺตรูปนิสฺสโย, ปกตูปนิสฺสโย ฯเปฯ.}\csromanspacer{} ปกตูปนิสฺสโย\csromandash{} ปฐมสฺส มคฺคสฺส ปริกมฺมํ ปฐมสฺส มคฺคสฺส อุปนิสฺสยปจฺจเยน ปจฺจโย ฯเปฯ.\csromanspacer{} จตุตฺถสฺส มคฺคสฺส ปริกมฺมํ จตุตฺถสฺส มคฺคสฺส อุปนิสฺสยปจฺจเยน ปจฺจโย. (2)}

\proseitem{157}{โลกุตฺตโร ธมฺโม โลกุตฺตรสฺส ธมฺมสฺส อุปนิสฺสยปจฺจเยน ปจฺจโย.\csromanspacer{} \textbf{อารมฺมณูปนิสฺสโย, อนนฺตรูปนิสฺสโย, ปกตูปนิสฺสโย ฯเปฯ.}\csromanspacer{} ปกตูปนิสฺสโย\csromandash{}ปฐโม มคฺโค ทุติยสฺส มคฺคสฺส อุปนิสฺสยปจฺจเยน ปจฺจโย ฯเปฯ.\csromanspacer{} ตติโย มคฺโค จตุตฺถสฺส มคฺคสฺส อุปนิสฺสยปจฺจเยน ปจฺจโย. (1)}

\prose{\csromanfolio{144}โลกุตฺตโร ธมฺโม โลกิยสฺส ธมฺมสฺส อุปนิสฺสยปจฺจเยน ปจฺจโย.\csromanspacer{} \textbf{อารมฺมณูปนิสฺสโย, อนนฺตรูปนิสฺสโย, ปกตูปนิสฺสโย ฯเปฯ.}\csromanspacer{} ปกตูปนิสฺสโย\csromandash{}อริยา มคฺคํ อุปนิสฺสาย อนุปฺปนฺนํ สมาปตฺติํ อุปฺปาเทนฺติ, อุปฺปนฺนํ สมาปชฺชนฺติ.\csromanspacer{} สงฺขาเร อนิจฺจโต ทุกฺขโต อนตฺตโต วิปสฺสนฺติ.\csromanspacer{} อริยานํ มคฺโค ฯเปฯ.\csromanspacer{} ฐานาฐานโกสลฺลสฺส อุปนิสฺสยปจฺจเยน ปจฺจโย.\csromanspacer{} ผลสมาปตฺติ กายิกสฺส สุขสฺส อุปนิสฺสยปจฺจเยน ปจฺจโย. (2)}

\csromanheader{2}{\textbf{ปุเรชาต}}
\proseitem{158}{โลกิโย ธมฺโม โลกิยสฺส ธมฺมสฺส ปุเรชาตปจฺจเยน ปจฺจโย.\csromanspacer{} \textbf{อารมฺมณปุเรชาตํ}, วตฺถุปุเรชาตํ.\csromanspacer{}\csromanspacer{}\textbf{อารมฺมณปุเรชาตํ}\csromandash{}จกฺขุํ ฯเปฯ วตฺถุํ อนิจฺจโต ทุกฺขโต อนตฺตโต ฯเปฯ โทมนสฺสํ อุปฺปชฺชติ.\csromanspacer{} ทิพฺเพน จกฺขุนา รูปํ ปสฺสติ.\csromanspacer{} ทิพฺพาย โสตธาตุยา สทฺทํ สุณาติ.\csromanspacer{} รูปายตนํ จกฺขุวิญฺญาณสฺส ฯเปฯ.\csromanspacer{} โผฏฺฐพฺพายตนํ กายวิญฺญาณสฺส ปุเรชาตปจฺจเยน ปจฺจโย.\csromanspacer{} \textbf{วตฺถุปุเรชาตํ}\csromandash{}จกฺขายตนํ จกฺขุวิญฺญาณสฺส ฯเปฯ.\csromanspacer{} กายาตนํ กายวิญฺญาณสฺส ฯเปฯ.\csromanspacer{} วตฺถุ โลกิยานํ ขนฺธานํ ปุเรชาตปจฺจเยน ปจฺจโย. (1)}

\prose{โลกิโย ธมฺโม โลกุตฺตรสฺส ธมฺมสฺส ปุเรชาตปจฺจเยน ปจฺจโย.\csromanspacer{} \textbf{วตฺถุปุเรชาตํ}\csromandash{}วตฺถุ โลกุตฺตรานํ ขนฺธานํ ปุเรชาตปจฺจเยน ปจฺจโย. (2)}

\csromanheader{2}{\textbf{ปจฺฉาชาต}}
\proseitem{159}{โลกิโย ธมฺโม โลกิยสฺส ธมฺมสฺส ปจฺฉาชาตปจฺจเยน ปจฺจโย.\csromanspacer{} ปจฺฉาชาตา โลกิยา ขนฺธา ปุเรชาตสฺส อิมสฺส กายสฺส ปจฺฉาชาตปจฺจเยน ปจฺจโย. (1)}

\prose{โลกุตฺตโร ธมฺโม โลกิยสฺส ธมฺมสฺส ปจฺฉาชาตปจฺจเยน ปจฺจโย.\csromanspacer{} ปจฺฉาชาตา โลกุตฺตรา ขนฺธา ปุเรชาตสฺส อิมสฺส กายสฺส ปจฺฉาชาตปจฺจเยน ปจฺจโย. (1)}

\csromanheader{2}{12. โลกิยทุก \textbf{อาเสวน}}
\proseitem{160}{\csromanfolio{145}โลกิโย ธมฺโม โลกิยสฺส ธมฺมสฺส อาเสวนปจฺจเยน ปจฺจโย.\csromanspacer{} ปุริมา ปุริมา โลกิยา ขนฺธา ปจฺฉิมานํ ปจฺฉิมานํ โลกิยานํ ขนฺธานํ อาเสวนปจฺจเยน ปจฺจโย.\csromanspacer{} อนุโลมํ โคตฺรภุสฺส.\csromanspacer{} อนุโลมํ โวทานสฺส อาเสวนปจฺจเยน ปจฺจโย. (1)}

\prose{โลกิโย ธมฺโม โลกุตฺตรสฺส ธมฺมสฺส อาเสวนปจฺจเยน ปจฺจโย.\csromanspacer{} โคตฺรภุ มคฺคสฺส.\csromanspacer{} โวทานํ มคฺคสฺส อาเสวนปจฺจเยน ปจฺจโย. (2)}

\csromanheader{2}{\textbf{กมฺม}}
\proseitem{161}{โลกิโย ธมฺโม โลกิยสฺส ธมฺมสฺส กมฺมปจฺจเยน ปจฺจโย.\csromanspacer{} \textbf{สหชาตา}, นานากฺขณิกา.\csromanspacer{}\csromanspacer{}\textbf{สหชาตา} โลกิยา เจตนา สมฺปยุตฺตกานํ ขนฺธานํ จิตฺตสมุฏฺฐานานญฺจ รูปานํ กมฺมปจฺจเยน ปจฺจโย.\csromanspacer{} ปฏิสนฺธิกฺขเณ ฯเปฯ.\csromanspacer{} \textbf{นานากฺขณิกา} โลกิยา เจตนา วิปากานํ ขนฺธานํ กฏตฺตา จ รูปานํ กมฺมปจฺจเยน ปจฺจโย. (1)}

\prose{โลกุตฺตโร ธมฺโม โลกุตฺตรสฺส ธมฺมสฺส กมฺมปจฺจเยน ปจฺจโย.\csromanspacer{} \textbf{สหชาตา}, นานากฺขณิกา.\csromanspacer{}\csromanspacer{}\textbf{สหชาตา} โลกุตฺตรา เจตนา สมฺปยุตฺตกานํ ขนฺธานํ กมฺมปจฺจเยน ปจฺจโย.\csromanspacer{} \textbf{นานากฺขณิกา} โลกุตฺตรา เจตนา วิปากานํ ขนฺธานํ กมฺมปจฺจเยน ปจฺจโย. (1)}

\prose{โลกุตฺตโร ธมฺโม โลกิยสฺส ธมฺมสฺส กมฺมปจฺจเยน ปจฺจโย.\csromanspacer{} โลกุตฺตรา เจตนา จิตฺตสมุฏฺฐานานํ รูปานํ กมฺมปจฺจเยน ปจฺจโย.}

\prose{โลกุตฺตโร ธมฺโม โลกิยสฺส จ โลกุตฺตรสฺส จ ธมฺมสฺส กมฺมปจฺจเยน ปจฺจโย.\csromanspacer{} โลกุตฺตรา เจตนา สมฺปยุตฺตกานํ ขนฺธานํ จิตฺตสมุฏฺฐานานญฺจ รูปานํ กมฺมปจฺจเยน ปจฺจโย. (3)}

\csromanheader{2}{\textbf{วิปาก}}
\proseitem{162}{โลกิโย ธมฺโม โลกิยสฺส ธมฺมสฺส วิปากปจฺจเยน ปจฺจโย.\csromanspacer{} วิปาโก โลกิโย เอโก ขนฺโธ ติณฺณนฺนํ ขนฺธานํ จิตฺตสมุฏฺฐานานญฺจ รูปานํ วิปากปจฺจเยน ปจฺจโย ฯเปฯ เทฺว ขนฺธา ฯเปฯ.\csromanspacer{} ปฏิสนฺธิกฺขเณ ฯเปฯ. (1)}

\prose{\csromanfolio{146}โลกุตฺตโร ธมฺโม โลกุตฺตรสฺส ธมฺมสฺส วิปากปจฺจเยน ปจฺจโย.\csromanspacer{} ตีณิ.}

\csromanheader{2}{\textbf{อาหาร}}
\proseitem{163}{โลกิโย ธมฺโม โลกิยสฺส ธมฺมสฺส อาหารปจฺจเยน ปจฺจโย.\csromanspacer{} โลกิยา อาหารา สมฺปยุตฺตกานํ ขนฺธานํ จิตฺตสมุฏฺฐานญฺจ รูปานํ อาหารปจฺจเยน ปจฺจโย.\csromanspacer{} ปฏิสนฺธิกฺขเณ ฯเปฯ.\csromanspacer{} กพฬีกาโร, อาหาโร อิมสฺส กายสฺส อาหารปจฺจเยน ปจฺจโย. (1)}

\prose{โลกุตฺตโร ธมฺโม โลกุตฺตรสฺส ธมฺมสฺส อาหารปจฺจเยน ปจฺจโย.\csromanspacer{} ตีณิ.}

\csromanheader{2}{\textbf{อินฺทฺริย}}
\proseitem{164}{โลกิโย ธมฺโม โลกิยสฺส ธมฺมสฺส อินฺทฺริยปจฺจเยน ปจฺจโย. (ปฏิสนฺธิ กาตพฺพา.) จกฺขุนฺทฺริยํ จกฺขุวิญฺญาณสฺส ฯเปฯ.\csromanspacer{} กายินฺทฺริยํ กายวิญฺญาณสฺส อินฺทฺริยปจฺจเยน ปจฺจโย.\csromanspacer{} รูปชีวิตินฺทฺริยํ กฏตฺตารูปานํ อินฺทฺริยปจฺจเยน ปจฺจโย. (1)}

\prose{โลกุตฺตโร ธมฺโม โลกุตฺตรสฺส ธมฺมสฺส อินฺทฺริยปจฺจเยน ปจฺจโย.\csromanspacer{} ตีณิ.}

\csromanheader{2}{\textbf{ฌานาทิ}}
\proseitem{165}{โลกิโย ธมฺโม โลกิยสฺส ธมฺมสฺส ฌานปจฺจเยน ปจฺจโย.\csromanspacer{} เอกํ.\csromanspacer{} โลกุตฺตโร ธมฺโม ฯเปฯ.\csromanspacer{} ตีณิ.\csromanspacer{} มคฺคปจฺจเยน ปจฺจโย.\csromanspacer{} โลกิเย เอกํ.\csromanspacer{} โลกุตฺตเร ตีณิ.}

\prose{โลกิโย ธมฺโม โลกิยสฺส ธมฺมสฺส สมฺปยุตฺตปจฺจเยน ปจฺจโย.\csromanspacer{} เอกํ.\csromanspacer{} โลกุตฺตโร ธมฺโม ฯเปฯ.\csromanspacer{} เอกํ.}

\csromanheader{2}{\textbf{วิปฺปยุตฺต}}
\proseitem{166}{โลกิโย ธมฺโม โลกิยสฺส ธมฺมสฺส วิปฺปยุตฺตปจฺจเยน ปจฺจโย.\csromanspacer{} สหชาตํ, ปุเรชาตํ, ปจฺฉาชาตํ.\csromanspacer{}\csromanspacer{}\textbf{สหชาตา} โลกิยา ขนฺธา จิตฺตสมุฏฺฐานานํ รูปานํ วิปฺปยุตฺตปจฺจเยน ปจฺจโย.\csromanspacer{} ปฏิสนฺธิกฺขเณ ฯเปฯ.\csromanspacer{} ขนฺธา วตฺถุสฺส วิปฺปยุตฺตปจฺจเยน ปจฺจโย.\csromanspacer{} วตฺถุ ขนฺธานํ วิปฺปยุตฺตปจฺจเยน ปจฺจโย.\csromanspacer{} \textbf{ปุเรชาตํ} จกฺขายตนํ จกฺขุวิญฺญาณสฺส ฯเปฯ.}

\vspace{\tipitakaparskip}
\csromancenter{โลกิยทุก กายายตนํ กายวิญฺญาณสฺส ฯเปฯ.\csromanspacer{} วตฺถุ โลกิยานํ ขนฺธานํ วิปฺปยุตฺตปจฺจเยน ปจฺจโย.\csromanspacer{} \textbf{ปจฺฉาชาตา} โลกิยา ขนฺธา ปุเรชาตสฺส อิมสฺส กายสฺส วิปฺปยุตฺตปจฺจเยน ปจฺจโย. (1)}
\prose{\csromanfolio{147}โลกิโย ธมฺโม โลกุตฺตรสฺส ธมฺมสฺส วิปฺปยุตฺตปจฺจเยน ปจฺจโย.\csromanspacer{} \textbf{ปุเรชาตํ} วตฺถุ โลกุตฺตรานํ ขนฺธานํ วิปฺปยุตฺตปจฺจเยน ปจฺจโย.}

\prose{โลกุตฺตโร ธมฺโม โลกิยสฺส ธมฺมสฺส วิปฺปยุตฺตปจฺจเยน ปจฺจโย.\csromanspacer{} สหชาตํ, ปจฺฉาชาตํ.\csromanspacer{}\csromanspacer{}\textbf{สหชาตา} โลกุตฺตรา ขนฺธา จิตฺตสมุฏฺฐานานํ รูปานํ วิปฺปยุตฺตปจฺจเยน ปจฺจโย.\csromanspacer{} \textbf{ปจฺฉาชาตา} โลกุตฺตรา ขนฺธา ปุเรชาตสฺส อิมสฺส กายสฺส วิปฺปยุตฺตปจฺจเยน ปจฺจโย. (1)}

\csromanheader{2}{\textbf{อตฺถิ-อาทิ}}
\proseitem{167}{โลกิโย ธมฺโม โลกิยสฺส ธมฺมสฺส อตฺถิปจฺจเยน ปจฺจโย.\csromanspacer{} สหชาตํ, ปุเรชาตํ, ปจฺฉาชาตํ, อาหารํ, อินฺทฺริยํ.\csromanspacer{}\csromanspacer{}\textbf{สหชาโต} โลกิโย เอโก ขนฺโธ ติณฺณนฺนํ ขนฺธานํ จิตฺตสมุฏฺฐานานญฺจ รูปานํ ฯเปฯ. (ยาว อสญฺญสตฺตา.) \textbf{ปุเรชาตํ} จกฺขุํ ฯเปฯ วตฺถุํ ฯเปฯ. (ปุเรชาตสทิสํ.) วตฺถุ โลกิยานํ ขนฺธานํ อตฺถิปจฺจเยน ปจฺจโย.\csromanspacer{} \textbf{ปจฺฉาชาตา} โลกิยา ขนฺธา ปุเรชาตสฺส อิมสฺส กายสฺส อตฺถิปจฺจเยน ปจฺจโย.\csromanspacer{} \textbf{กพฬีกาโร อาหาโร} อิมสฺส กายสฺส อตฺถิปจฺจเยน ปจฺจโย.\csromanspacer{} \textbf{รูปชีวิตินฺทฺริยํ} กฏตฺตารูปานํ อตฺถิปจฺจเยน ปจฺจโย. (1)}

\prose{โลกิโย ธมฺโม โลกุตฺตรสฺส ธมฺมสฺส อตฺถิปจฺจเยน ปจฺจโย.\csromanspacer{} \textbf{ปุเรชาตํ} วตฺถุ โลกุตฺตรานํ ขนฺธานํ อตฺถิปจฺจเยน ปจฺจโย. (2)}

\proseitem{168}{โลกุตฺตโร ธมฺโม โลกุตฺตรสฺส ธมฺมสฺส อตฺถิปจฺจเยน ปจฺจโย.\csromanspacer{} โลกุตฺตโร เอโก ขนฺโธ ติณฺณนฺนํ ขนฺธานํ อตฺถิปจฺจเยน ปจฺจโย ฯเปฯ เทฺว ขนฺธา ฯเปฯ. (1)}

\prose{โลกุตฺตโร ธมฺโม โลกิยสฺส ธมฺมสฺส อตฺถิปจฺจเยน ปจฺจโย.\csromanspacer{} สหชาตํ, ปจฺฉาชาตํ.\csromanspacer{}\csromanspacer{}\textbf{สหชาตา} โลกุตฺตรา ขนฺธา จิตฺตสมุฏฺฐานานํ รูปานํ อตฺถิปจฺจเยน ปจฺจโย.\csromanspacer{} \textbf{ปจฺฉาชาตา} โลกุตฺตรา ขนฺธา ปุเรชาตสฺส อิมสฺส กายสฺส อตฺถิปจฺจเยน ปจฺจโย. (2)}

\prose{\csromanfolio{148}โลกุตฺตโร ธมฺโม โลกิยสฺส จ โลกุตฺตรสฺส จ ธมฺมสฺส อตฺถิปจฺจเยน ปจฺจโย.\csromanspacer{} โลกุตฺตโร เอโก ขนฺโธ ติณฺณนฺนํ ขนฺธานํ จิตฺตสมุฏฺฐานานญฺจ รูปานํ อตฺถิปจฺจเยน ปจฺจโย ฯเปฯ เทฺว ขนฺธา ฯเปฯ. (3)}

\proseitem{169}{โลกิโย จ โลกุตฺตโร จ ธมฺมา โลกิยสฺส ธมฺมสฺส อตฺถิปจฺจเยน ปจฺจโย.\csromanspacer{} สหชาตํ, ปจฺฉาชาตํ, อาหารํ, อินฺทฺริยํ.\csromanspacer{} \textbf{สหชาตา} โลกุตฺตรา ขนฺธา จ มหาภูตา จ จิตฺตสมุฏฺฐานานํ รูปานํ อตฺถิปจฺจเยน ปจฺจโย.\csromanspacer{} \textbf{ปจฺฉาชาตา} โลกุตฺตรา ขนฺธา จ กพฬีกาโร อาหาโร จ อิมสฺส กายสฺส อตฺถิปจฺจเยน ปจฺจโย.\csromanspacer{} \textbf{ปจฺฉาชาตา} โลกุตฺตรา ขนฺธา จ รูปชีวิตินฺทฺริยญฺจ กฏตฺตารูปานํ อตฺถิปจฺจเยน ปจฺจโย. (1)}

\prose{โลกิโย จ โลกุตฺตโร จ ธมฺมา โลกุตฺตรสฺส ธมฺมสฺส อตฺถิปจฺจเยน ปจฺจโย.\csromanspacer{} สหชาตํ, ปุเรชาตํ.\csromanspacer{}\csromanspacer{}สหชาโต โลกุตฺตโร เอโก ขนฺโธ จ วตฺถุ จ ติณฺณนฺนํ ขนฺธานํ อตฺถิปจฺจเยน ปจฺจโย ฯเปฯ เทฺว ขนฺธา จ ฯเปฯ. (2)}

\prose{นตฺถิปจฺจเยน ปจฺจโย.\csromanspacer{} วิคตปจฺจเยน ปจฺจโย.\csromanspacer{} อวิคตปจฺจเยน ปจฺจโย. (2)}

\csromanheader{2}{1. ปจฺจยานุโลม 2. สงฺขฺยาวาร}
\csromanheader{2}{\textbf{สุทฺธ}}
\proseitem{170}{เหตุยา จตฺตาริ, อารมฺมเณ ตีณิ, อธิปติยา จตฺตาริ, อนนฺตเร จตฺตาริ, สมนนฺตเร จตฺตาริ, สหชาเต ปญฺจ, อญฺญมญฺเญ เทฺว, นิสฺสเย สตฺต, อุปนิสฺสเย จตฺตาริ, ปุเรชาเต เทฺว, ปจฺฉาชาเต เทฺว, อาเสวเน เทฺว, กมฺเม จตฺตาริ, วิปาเก จตฺตาริ, อาหาเร จตฺตาริ, อินฺทฺริเย จตฺตาริ, ฌาเน จตฺตาริ, มคฺเค จตฺตาริ, สมฺปยุตฺเต เทฺว, วิปฺปยุตฺเต ตีณิ, อตฺถิยา สตฺต, นตฺถิยา จตฺตาริ, วิคเต จตฺตาริ, อวิคเต สตฺต.}

\vspace{\tipitakaparskip}
\csromancenter{อนุโลมํ.}
\csromanheader{2}{2. ปจฺจนียุทฺธาร}
\proseitem{171}{โลกิโย ธมฺโม โลกิยสฺส ธมฺมสฺส อารมฺมณปจฺจเยน ปจฺจโย.\csromanspacer{} สหชาตปจฺจเยน ปจฺจโย.\csromanspacer{} อุปนิสฺสยปจฺจเยน ปจฺจโย.}

\vspace{\tipitakaparskip}
\csromancenter{โลกิยทุก ปุเรชาตปจฺจเยน ปจฺจโย.\csromanspacer{} ปจฺฉาชาตปจฺจเยน ปจฺจโย.\csromanspacer{} กมฺมปจฺจเยน ปจฺจโย.\csromanspacer{} อาหารปจฺจเยน ปจฺจโย.\csromanspacer{} อินฺทฺริยปจฺจเยน ปจฺจโย. (1)}
\prose{\csromanfolio{149}โลกิโย ธมฺโม โลกุตฺตรสฺส ธมฺมสฺส อุปนิสฺสยปจฺจเยน ปจฺจโย.\csromanspacer{} ปุเรชาตปจฺจเยน ปจฺจโย. (2)}

\proseitem{172}{โลกุตฺตโร ธมฺโม โลกุตฺตรสฺส ธมฺมสฺส อารมฺมณปจฺจเยน ปจฺจโย.\csromanspacer{} สหชาตปจฺจเยน ปจฺจโย.\csromanspacer{} อุปนิสฺสยปจฺจเยน ปจฺจโย. (1)}

\prose{โลกุตฺตโร ธมฺโม โลกิยสฺส ธมฺมสฺส อารมฺมณปจฺจเยน ปจฺจโย.\csromanspacer{} สหชาตปจฺจเยน ปจฺจโย.\csromanspacer{} อุปนิสฺสยปจฺจเยน ปจฺจโย.\csromanspacer{} ปจฺฉาชาตปจฺจเยน ปจฺจโย. (2)}

\prose{โลกุตฺตโร ธมฺโม โลกิยสฺส จ โลกุตฺตรสฺส จ ธมฺมสฺส สหชาตปจฺจเยน ปจฺจโย. (3)}

\prose{โลกิโย จ โลกุตฺตโร จ ธมฺมา โลกิยสฺส ธมฺมสฺส สหชาตํ.\csromanspacer{} ปจฺฉาชาตํ.\csromanspacer{} อาหารํ.\csromanspacer{} อินฺทฺริยํ. (1)}

\prose{โลกิโย จ โลกุตฺตโร จ ธมฺมา โลกุตฺตรสฺส ธมฺมสฺส สหชาตํ.\csromanspacer{} ปุเรชาตํ. (2)}

\csromanheader{2}{2. ปจฺจยปจฺจนีย 2. สงฺขฺยาวาร}
\csromanheader{2}{\textbf{สุทฺธ}}
\proseitem{173}{นเหตุยา สตฺต ฯเปฯ นสมนนฺตเร สตฺต, นสหชาเต ปญฺจ, น-อญฺญมญฺเญ ปญฺจ, นนิสฺสเย ปญฺจ, น-อุปนิสฺสเย สตฺต, นปุเรชาเต ฉ, นปจฺฉาชาเต สตฺต ฯเปฯ นมคฺเค สตฺต, นสมฺปยุตฺเต ปญฺจ, นวิปฺปยุตฺเต จตฺตาริ, โน-อตฺถิยา จตฺตาริ, โนนตฺถิยา สตฺต, โนวิคเต สตฺต, โน-อวิคเต จตฺตาริ.}

\vspace{\tipitakaparskip}
\csromancenter{ปจฺจนียํ.}
\csromanheader{2}{3. ปจฺจยานุโลมปจฺจนีย}
\csromanheader{2}{\textbf{เหตุทุก}}
\proseitem{174}{\csromanfolio{150}\textbf{เหตุ}ปจฺจยา น-อารมฺมเณ จตฺตาริ ฯเปฯ นสมนนฺตเร จตฺตาริ, น-อญฺญมญฺเญ เทฺว, น-อุปนิสฺสเย จตฺตาริ ฯเปฯ นมคฺเค จตฺตาริ, นสมฺปยุตฺเต เทฺว, นวิปฺปยุตฺเต เทฺว, โนนตฺถิยา จตฺตาริ, โนวิคเต จตฺตาริ.}

\vspace{\tipitakaparskip}
\csromancenter{อนุโลมปจฺจนียํ.}
\csromanheader{2}{4. ปจฺจยปจฺจนียานุโลม}
\csromanheader{2}{\textbf{นเหตุทุก}}
\proseitem{175}{นเหตุปจฺจยา อารมฺมเณ ตีณิ, อธิปติยา จตฺตาริ, (อนุโลมมาติกา กาตพฺพา.) อวิคเต สตฺต.}

\vspace{\tipitakaparskip}
\csromancenter{ปจฺจนียานุโลมํ.}
\nitthitam{โลกิยทุกํ นิฏฺฐิตํ.}
\csromanmatikaanchor{mtk.16}
\csromantocmarkref{section}{13. เกนจิวิญฺเญยฺยทุก}{mtk.16}
\bookmark[dest={mtk.16},level=1]{13. เกนจิวิญฺเญยฺยทุก}
\csromanheader{1}{13. เกนจิวิญฺเญยฺยทุก 1. ปฏิจฺจวาร}
\csromanheader{2}{1. ปจฺจยานุโลม 1. วิภงฺควาร}
\csromanheader{2}{\textbf{เหตุ}}
\proseitem{176}{เกนจิ วิญฺเญยฺยํ ธมฺมํ ปฏิจฺจ เกนจิ วิญฺเญยฺโย ธมฺโม อุปฺปชฺชติ เหตุปจฺจยา.\csromanspacer{} เกนจิ วิญฺเญยฺยํ เอกํ ขนฺธํ ปฏิจฺจ ตโย ขนฺธา จิตฺตสมุฏฺฐานญฺจ รูปํ ฯเปฯ เทฺว ขนฺเธ ฯเปฯ.\csromanspacer{} ปฏิสนฺธิกฺขเณ ฯเปฯ.\csromanspacer{} ขนฺเธ ปฏิจฺจ วตฺถุ, วตฺถุํ ปฏิจฺจ ขนฺธา.\csromanspacer{} เอกํ มหาภูตํ ปฏิจฺจ ตโย มหาภูตา ฯเปฯ มหาภูเต ปฏิจฺจ จิตฺตสมุฏฺฐานํ รูปํ, กฏตฺตารูปํ, อุปาทารูปํ. (1)}

\prose{เกนจิ วิญฺเญยฺยํ ธมฺมํ ปฏิจฺจ เกนจิ นวิญฺเญยฺโย ธมฺโม อุปฺปชฺชติ เหตุปจฺจยา.\csromanspacer{} เกนจิ วิญฺเญยฺยํ เอกํ ขนฺธํ ปฏิจฺจ เกนจิ นวิญฺเญยฺยา}

\vspace{\tipitakaparskip}
\csromancenter{เกนจิวิญฺเญยฺยทุก ตโย ขนฺธา จิตฺตสมุฏฺฐานญฺจ รูปํ ฯเปฯ เทฺว ขนฺเธ ฯเปฯ.\csromanspacer{} ปฏิสนฺธิกฺขเณ ฯเปฯ.\csromanspacer{} ขนฺเธ ปฏิจฺจ วตฺถุ, วตฺถุํ ปฏิจฺจ ขนฺธา.\csromanspacer{} เอกํ มหาภูตํ ฯเปฯ มหาภูเต ปฏิจฺจ จิตฺตสมุฏฺฐานํ รูปํ, กฏตฺตารูปํ, อุปาทารูปํ. (2)}
\prose{\csromanfolio{151}เกนจิ วิญฺเญยฺยํ ธมฺมํ ปฏิจฺจ เกนจิ วิญฺเญยฺโย จ เกนจิ นวิญฺเญยฺโย จ ธมฺมา อุปฺปชฺชนฺติ เหตุปจฺจยา.\csromanspacer{} เกนจิ วิญฺเญยฺยํ เอกํ ขนฺธํ ปฏิจฺจ เกนจิ วิญฺเญยฺยา จ เกนจิ นวิญฺเญยฺยา จ ตโย ขนฺธา จิตฺตสมุฏฺฐานญฺจ รูปํ ฯเปฯ เทฺว ขนฺเธ ฯเปฯ.\csromanspacer{} ปฏิสนฺธิกฺขเณ ฯเปฯ.\csromanspacer{} ขนฺเธ ปฏิจฺจ วตฺถุ, วตฺถุํ ปฏิจฺจ ขนฺธา.\csromanspacer{} เอกํ มหาภูตํ ฯเปฯ มหาภูเต ปฏิจฺจ จิตฺตสมุฏฺฐานํ รูปํ, กฏตฺตารูปํ, อุปาทารูปํ. (3)}

\proseitem{177}{เกนจิ นวิญฺเญยฺยํ ธมฺมํ ปฏิจฺจ เกนจิ นวิญฺเญยฺโย ธมฺโม อุปฺปชฺชติ เหตุปจฺจยา.\csromanspacer{} เกนจิ นวิญฺเญยฺยํ เอกํ ขนฺธํ ปฏิจฺจ เกนจิ นวิญฺเญยฺยา ตโย ขนฺธา จิตฺตสมุฏฺฐานญฺจ รูปํ ฯเปฯ เทฺว ขนฺเธ ฯเปฯ.\csromanspacer{} ปฏิสนฺธิกฺขเณ ฯเปฯ.\csromanspacer{} ขนฺเธ ปฏิจฺจ วตฺถุ, วตฺถุํ ปฏิจฺจ ขนฺธา.\csromanspacer{} เอกํ มหาภูตํ ฯเปฯ มหาภูเต ปฏิจฺจ จิตฺตสมุฏฺฐานํ รูปํ, กฏตฺตารูปํ, อุปาทารูปํ. (1)}

\prose{เกนจิ นวิญฺเญยฺยํ ธมฺมํ ปฏิจฺจ เกนจิ วิญฺเญยฺโย ธมฺโม อุปฺปชฺชติ เหตุปจฺจยา.\csromanspacer{} เกนจิ นวิญฺเญยฺยํ เอกํ ขนฺธํ ปฏิจฺจ เกนจิ วิญฺเญยฺยา ตโย ขนฺธา จิตฺตสมุฏฺฐานญฺจ รูปํ ฯเปฯ เทฺว ขนฺเธ ฯเปฯ.\csromanspacer{} ปฏิสนฺธิกฺขเณ ฯเปฯ.\csromanspacer{} ขนฺเธ ปฏิจฺจ วตฺถุ, วตฺถุํ ปฏิจฺจ ขนฺธา.\csromanspacer{} เอกํ มหาภูตํ ฯเปฯ มหาภูเต ปฏิจฺจ จิตฺตสมุฏฺฐานํ รูปํ, กฏตฺตารูปํ, อุปาทารูปํ. (2)}

\prose{เกนจิ นวิญฺเญยฺยํ ธมฺมํ ปฏิจฺจ เกนจิ วิญฺเญยฺโย จ เกนจิ นวิญฺเญยฺโย จ ธมฺมา อุปฺปชฺชนฺติ เหตุปจฺจยา.\csromanspacer{} เกนจิ นวิญฺเญยฺยํ เอกํ ขนฺธํ ปฏิจฺจ เกนจิ วิญฺเญยฺยา จ เกนจิ นวิญฺเญยฺยา จ ตโย ขนฺธา จิตฺตสมุฏฺฐานญฺจ รูปํ ฯเปฯ เทฺว ขนฺเธ ฯเปฯ.\csromanspacer{} ปฏิสนฺธิกฺขเณ ฯเปฯ.\csromanspacer{} ขนฺเธ ปฏิจฺจ วตฺถุ, วตฺถุํ ปฏิจฺจ ขนฺธา.\csromanspacer{} เอกํ มหาภูตํ ฯเปฯ มหาภูเต ปฏิจฺจ จิตฺตสมุฏฺฐานํ รูปํ, กฏตฺตารูปํ, อุปาทารูปํ. (3)}

\proseitem{178}{เกนจิ วิญฺเญยฺยญฺจ เกนจิ นวิญฺเญยฺยญฺจ ธมฺมํ ปฏิจฺจ เกนจิ วิญฺเญยฺโย ธมฺโม อุปฺปชฺชติ เหตุปจฺจยา.\csromanspacer{} เกนจิ วิญฺเญยฺยญฺจ เกนจิ นวิญฺเญยฺยญฺจ เอกํ ขนฺธํ ปฏิจฺจ เกนจิ วิญฺเญยฺยา ตโย ขนฺธา จิตฺตสมุฏฺฐานญฺจ รูปํ ฯเปฯ เทฺว ขนฺเธ ฯเปฯ.\csromanspacer{} ปฏิสนฺธิกฺขเณ ฯเปฯ.\csromanspacer{} ขนฺเธ ปฏิจฺจ วตฺถุ, วตฺถุํ ปฏิจฺจ \csromanfolio{152}ขนฺธา.\csromanspacer{} เอกํ มหาภูตํ ฯเปฯ มหาภูเต ปฏิจฺจ จิตฺตสมุฏฺฐานํ รูปํ, กฏตฺตารูปํ, อุปาทารูปํ. (1)}

\prose{เกนจิ วิญฺเญยฺยญฺจ เกนจิ นวิญฺเญยฺยญฺจ ธมฺมํ ปฏิจฺจ เกนจิ นวิญฺเญยฺโย ธมฺโม อุปฺปชฺชติ เหตุปจฺจยา.\csromanspacer{} เกนจิ วิญฺเญยฺยญฺจ เกนจิ นวิญฺเญยฺยญฺจ เอกํ ขนฺธํ ปฏิจฺจ เกนจิ นวิญฺเญยฺยา ตโย ขนฺธา จิตฺตสมุฏฺฐานญฺจ รูปํ ฯเปฯ เทฺว ขนฺเธ ฯเปฯ.\csromanspacer{} ปฏิสนฺธิกฺขเณ ฯเปฯ.\csromanspacer{} ขนฺเธ ปฏิจฺจ วตฺถุ, วตฺถุํ ปฏิจฺจ ขนฺธา.\csromanspacer{} เอกํ มหาภูตํ ฯเปฯ มหาภูเต ปฏิจฺจ จิตฺตสมุฏฺฐานํ รูปํ, กฏตฺตารูปํ, อุปาทารูปํ. (2)}

\prose{เกนจิ วิญฺเญยฺยญฺจ เกนจิ นวิญฺเญยฺยญฺจ ธมฺมํ ปฏิจฺจ เกนจิ วิญฺเญยฺโย จ เกนจิ นวิญฺเญยฺโย จ ธมฺมา อุปฺปชฺชนฺติ เหตุปจฺจยา.\csromanspacer{} เกนจิ วิญฺเญยฺยญฺจ เกนจิ นวิญฺเญยฺยญฺจ เอกํ ขนฺธํ ปฏิจฺจ เกนจิ วิญฺเญยฺยา จ เกนจิ นวิญฺเญยฺยา จ ตโย ขนฺธา จิตฺตสมุฏฺฐานญฺจ รูปํ ฯเปฯ เทฺว ขนฺเธ ฯเปฯ.\csromanspacer{} ปฏิสนฺธิกฺขเณ ฯเปฯ.\csromanspacer{} ขนฺเธ ปฏิจฺจ วตฺถุ, วตฺถุํ ปฏิจฺจ ขนฺธา.\csromanspacer{} เอกํ มหาภูตํ ฯเปฯ มหาภูเต ปฏิจฺจ จิตฺตสมุฏฺฐานํ รูปํ, กฏตฺตารูปํ, อุปาทารูปํ. (สํขิตฺตํ.) (3)}

\csromanheader{2}{1. ปจฺจยานุโลม 2. สงฺขฺยาวาร}
\vspace{\tipitakaparskip}
\csromancenter{เหตุยา นว, อารมฺมเณ นว ฯเปฯ อวิคเต นว.}
\csromancenter{อนุโลมํ.}
\csromanheader{2}{2. ปจฺจยปจฺจนีย 2. สงฺขฺยาวาร}
\proseitem{180}{นเหตุยา นว, น-อารมฺมเณ นว ฯเปฯ โนวิคเต นว. (เอวํ จตฺตาริปิ คณนา ปริปุณฺณา.) (สหชาตวาโรปิ ปจฺจยวาโรปิ นิสฺสยวาโรปิ สํสฏฺฐวาโรปิ สมฺปยุตฺตวาโรปิ เอวํ วิตฺถาเรตพฺพา.\csromanspacer{} ปจฺจยวาเร วตฺถุ จ ปญฺจายตนานิ จ ทสฺเสตพฺพานิ.\csromanspacer{} ยถา ยถา ลพฺภติ.\csromanspacer{} ตํ ตํ กาตพฺพํ.)}

\csromanheader{2}{13. เกนจิวิญฺเญยฺยทุก \textbf{13. เกนจิวิญฺเญยฺยทุก 7. ปญฺหาวาร}}
\csromanheader{2}{1. ปจฺจยานุโลมาทิ}
\proseitem{181}{\csromanfolio{153}เกนจิ วิญฺเญยฺโย ธมฺโม เกนจิ วิญฺเญยฺยสฺส ธมฺมสฺส เหตุปจฺจเยน ปจฺจโย.\csromanspacer{} เกนจิ วิญฺเญยฺยา เหตู สมฺปยุตฺตกานํ ขนฺธานํ จิตฺตสมุฏฺฐานานญฺจ รูปานํ เหตุปจฺจเยน ปจฺจโย.\csromanspacer{} ปฏิสนฺธิกฺขเณ ฯเปฯ. (สํขิตฺตํ.)}

\prose{เหตุยา นว, อารมฺมเณ นว ฯเปฯ อวิคเต นว.\csromanspacer{} อนุโลมํ.}

\prose{นเหตุยา นว ฯเปฯ โนวิคเต นว. (เอวํ จตฺตาริปิ คณนา ปริปุณฺณา.)}

\vspace{\tipitakaparskip}
\csromancenter{ปจฺจนียํ.}
\nitthitam{เกนจิวิญฺเญยฺยทุกํ นิฏฺฐิตํ.}
\nitthitam{จูฬนฺตรทุกํ นิฏฺฐิตํ.}
\csromanmatikaanchor{mtk.17}
\csromantocmarknopage{chapter}{3. อาสวโคจฺฉก}{mtk.17}
\bookmark[dest={mtk.17},level=0]{3. อาสวโคจฺฉก}
\chapterheadpageread{3. อาสวโคจฺฉก}
\csromanmatikaanchor{mtk.18}
\csromantocmarkref{section}{14. อาสวทุก}{mtk.18}
\bookmark[dest={mtk.18},level=1]{14. อาสวทุก}
\csromanheader{1}{14. อาสวทุก 1. ปฏิจฺจวาร}
\csromanheader{2}{1. ปจฺจยานุโลม 1. วิภงฺควาร}
\csromanheader{2}{\textbf{เหตุ}}
\proseitem{1}{\csromanfolio{154}อาสวํ ธมฺมํ ปฏิจฺจ อาสโว ธมฺโม อุปฺปชฺชติ เหตุปจฺจยา.\csromanspacer{} กามาสวํ ปฏิจฺจ ทิฏฺฐาสโว อวิชฺชาสโว.\csromanspacer{} ทิฏฺฐาสวํ ปฏิจฺจ กามาสโว อวิชฺชาสโว.\csromanspacer{} อวิชฺชาสวํ ปฏิจฺจ กามาสโว ทิฏฺฐาสโว.\csromanspacer{} ภวาสวํ ปฏิจฺจ อวิชฺชาสโว.\csromanspacer{} ทิฏฺฐาสวํ ปฏิจฺจ อวิชฺชาสโว. (เอเกกมฺปิ จกฺกํ กาตพฺพํ.) (1)}

\prose{อาสวํ ธมฺมํ ปฏิจฺจ โน-อาสโว ธมฺโม อุปฺปชฺชติ เหตุปจฺจยา.\csromanspacer{} อาสวํ ปฏิจฺจ อาสวสมฺปยุตฺตกา ขนฺธา จิตฺตสมุฏฺฐานญฺจ รูปํ. (2)}

\prose{อาสวํ ธมฺมํ ปฏิจฺจ อาสโว จ โน-อาสโว จ ธมฺมา อุปฺปชฺชนฺติ เหตุปจฺจยา.\csromanspacer{} กามาสวํ ปฏิจฺจ ทิฏฺฐาสโว อวิชฺชาสโว สมฺปยุตฺตกา จ ขนฺธา จิตฺตสมุฏฺฐานญฺจ รูปํ. (จกฺกํ.) (3)}

\proseitem{2}{โน-อาสวํ ธมฺมํ ปฏิจฺจ โน-อาสโว ธมฺโม อุปฺปชฺชติ เหตุปจฺจยา.\csromanspacer{} โน-อาสวํ เอกํ ขนฺธํ ปฏิจฺจ ตโย ขนฺธา จิตฺตสมุฏฺฐานํ รูปํ ฯเปฯ เทฺว ขนฺเธ ฯเปฯ.\csromanspacer{} ปฏิสนฺธิกฺขเณ ฯเปฯ.\csromanspacer{} ขนฺเธ ปฏิจฺจ วตฺถุ, วตฺถุํ ปฏิจฺจ ขนฺธา.\csromanspacer{} เอกํ มหาภูตํ ฯเปฯ มหาภูเต ปฏิจฺจ จิตฺตสมุฏฺฐานํ รูปํ, กฏตฺตารูปํ, อุปาทารูปํ. (1)}

\prose{โน-อาสวํ ธมฺมํ ปฏิจฺจ อาสโว ธมฺโม อุปฺปชฺชติ เหตุปจฺจยา.\csromanspacer{} โน-อาสเว ขนฺเธ ปฏิจฺจ อาสวา. (2)}

\prose{โน-อาสวํ ธมฺมํ ปฏิจฺจ อาสโว จ โน-อาสโว จ ธมฺมา อุปฺปชฺชนฺติ เหตุปจฺจยา.\csromanspacer{} โน-อาสวํ เอกํ ขนฺธํ ปฏิจฺจ ตโย ขนฺธา อาสวา จ จิตฺตสมุฏฺฐานญฺจ รูปํ ฯเปฯ เทฺว ขนฺเธ ฯเปฯ. (3)}

\proseitem{3}{อาสวญฺจ โน-อาสวญฺจ ธมฺมํ ปฏิจฺจ อาสโว ธมฺโม อุปฺปชฺชติ เหตุปจฺจยา.\csromanspacer{} กามาสวญฺจ สมฺปยุตฺตเก จ ขนฺเธ ปฏิจฺจ ทิฏฺฐาสโว อวิชฺชาสโว. (จกฺกํ พนฺธิตพฺพํ.) (1)}

\csromanheader{2}{14. อาสวทุก}
\prose{\csromanfolio{155}อาสวญฺจ โน-อาสวญฺจ ธมฺมํ ปฏิจฺจ โน-อาสโว ธมฺโม อุปฺปชฺชติ เหตุปจฺจยา.\csromanspacer{} โน-อาสวํ เอกํ ขนฺธญฺจ อาสเว จ ปฏิจฺจ ตโย ขนฺธา จิตฺตสมุฏฺฐานญฺจ รูปํ ฯเปฯ เทฺว ขนฺเธ ฯเปฯ. (2)}

\prose{อาสวญฺจ โน-อาสวญฺจ ธมฺมํ ปฏิจฺจ อาสโว จ โน-อาสโว จ ธมฺมา อุปฺปชฺชนฺติ เหตุปจฺจยา.\csromanspacer{} โน-อาสวํ เอกํ ขนฺธญฺจ กามาสวญฺจ ปฏิจฺจ ตโย ขนฺธา ทิฏฺฐาสโว อวิชฺชาสโว จิตฺตสมุฏฺฐานญฺจ รูปํ ฯเปฯ เทฺว ขนฺเธ จ ฯเปฯ. (จกฺกํ.\csromanspacer{} สํขิตฺตํ.) (3)}

\csromanheader{2}{1. ปจฺจยานุโลม 2. สงฺขฺยาวาร}
\proseitem{4}{เหตุยา นว, อารมฺมเณ นว, (สพฺพตฺถ นว.) วิปาเก เอกํ อาหาเร นว ฯเปฯ อวิคเต นว.}

\vspace{\tipitakaparskip}
\csromancenter{อนุโลมํ.}
\csromanheader{2}{2. ปจฺจยปจฺจนีย 1. วิภงฺควาร}
\csromanheader{2}{\textbf{นเหตุ}}
\proseitem{5}{โน-อาสวํ ธมฺมํ ปฏิจฺจ โน-อาสโว ธมฺโม อุปฺปชฺชติ นเหตุปจฺจยา.\csromanspacer{} อเหตุกํ โน-อาสวํ เอกํ ขนฺธํ ปฏิจฺจ ตโย ขนฺธา จิตฺตสมุฏฺฐานญฺจ รูปํ ฯเปฯ เทฺว ขนฺเธ ฯเปฯ.\csromanspacer{} อเหตุกปฏิสนฺธิกฺขเณ ฯเปฯ.\csromanspacer{} ขนฺเธ ปฏิจฺจ วตฺถุ, วตฺถุํ ปฏิจฺจ ขนฺธา.\csromanspacer{} เอกํ มหาภูตํ ฯเปฯ. (ยาว อสญฺญสตฺตา.) (1)}

\prose{โน-อาสวํ ธมฺมํ ปฏิจฺจ อาสโว ธมฺโม อุปฺปชฺชติ นเหตุปจฺจยา.\csromanspacer{} วิจิกิจฺฉาสหคเต อุทฺธจฺจสหคเต ขนฺเธ ปฏิจฺจ วิจิกิจฺฉาสหคโต อุทฺธจฺจสหคโต โมโห. (2)}

\csromanheader{2}{\textbf{น-อารมฺมณ}}
\proseitem{6}{อาสวํ ธมฺมํ ปฏิจฺจ โน-อาสโว ธมฺโม อุปฺปชฺชติ น-อารมฺมณปจฺจยา.\csromanspacer{} อาสเว ปฏิจฺจ จิตฺตสมุฏฺฐานํ รูปํ. (1)}

\prose{\csromanfolio{156}โน-อาสวํ ธมฺมํ ปฏิจฺจ โน-อาสโว ธมฺโม อุปฺปชฺชติ น-อารมฺมณปจฺจยา.\csromanspacer{} โน-อาสเว ขนฺเธ ปฏิจฺจ จิตฺตสมุฏฺฐานํ รูปํ.\csromanspacer{} ปฏิสนฺธิกฺขเณ ฯเปฯ.\csromanspacer{} ขนฺเธ ปฏิจฺจ วตฺถุ.\csromanspacer{} เอกํ มหาภูตํ ฯเปฯ. (ยาว อสญฺญสตฺตา.) (1)}

\prose{อาสวญฺจ โน-อาสวญฺจ ธมฺมํ ปฏิจฺจ โน-อาสโว ธมฺโม อุปฺปชฺชติ น-อารมฺมณปจฺจยา.\csromanspacer{} อาสเว จ สมฺปยุตฺตเก จ ขนฺเธ ปฏิจฺจ จิตฺตสมุฏฺฐานํ รูปํ. (สํขิตฺตํ.) (1)}

\csromanheader{2}{2. ปจฺจยปจฺจนีย 2. สงฺขฺยาวาร}
\csromanheader{2}{\textbf{สุทฺธ}}
\proseitem{7}{นเหตุยา เทฺว, น-อารมฺมเณ ตีณิ, น-อธิปติยา นว, น-อนนฺตเร ตีณิ, นสมนนฺตเร ตีณิ, น-อญฺญมญฺเญ ตีณิ, น-อุปนิสฺสเย ตีณิ, นปุเรชาเต นว, นปจฺฉาชาเต นว, น-อาเสวเน นว, นกมฺเม ตีณิ, นวิปาเก นว, น-อาหาเร เอกํ, น-อินฺทฺริเย เอกํ, นฌาเน เอกํ, นมคฺเค เอกํ, นสมฺปยุตฺเต ตีณิ, นวิปฺปยุตฺเต นว, โนนตฺถิยา ตีณิ, โนวิคเต ตีณิ.}

\vspace{\tipitakaparskip}
\csromancenter{ปจฺจนียํ.}
\csromanheader{2}{3. ปจฺจยานุโลมปจฺจนีย}
\csromanheader{2}{\textbf{เหตุทุก}}
\proseitem{8}{เหตุปจฺจยา น-อารมฺมเณ ตีณิ, น-อธิปติยา นว, น-อนนฺตเร ตีณิ, นสมนนฺตเร ตีณิ, น-อญฺญมญฺเญ ตีณิ, น-อุปนิสฺสเย ตีณิ, นปุเรชาเต นว, นปจฺฉาชาเต นว, น-อาเสวเน นว, นกมฺเม ตีณิ, นวิปาเก นว, นสมฺปยุตฺเต ตีณิ, นวิปฺปยุตฺเต นว, โนนตฺถิยา ตีณิ, โนวิคเต ตีณิ.}

\vspace{\tipitakaparskip}
\csromancenter{อนุโลมปจฺจนียํ.}
\csromanheader{2}{14. อาสวทุก \textbf{4. ปจฺจยปจฺจนียานุโลม}}
\csromanheader{2}{\textbf{นเหตุทุก}}
\proseitem{9}{\csromanfolio{157}นเหตุปจฺจยา อารมฺมเณ เทฺว, อนนฺตเร เทฺว ฯเปฯ วิปาเก เอกํ ฯเปฯ มคฺเค เอกํ ฯเปฯ อวิคเต เทฺว.}

\vspace{\tipitakaparskip}
\csromancenter{ปจฺจนียานุโลมํ.}
\csromancenter{(สหชาตวาโร ปฏิจฺจวารสทิโส.)}
\csromanheader{2}{14. อาสวทุก 3. ปจฺจยวาร}
\csromanheader{2}{1. ปจฺจยานุโลม 1. วิภงฺควาร}
\csromanheader{2}{\textbf{เหตุ}}
\proseitem{10}{อาสวํ ธมฺมํ ปจฺจยา อาสโว ธมฺโม อุปฺปชฺชติ เหตุปจฺจยา. (อาสวมูลกํ ตีณิ.\csromanspacer{} ปฏิจฺจสทิสา.)}

\prose{โน-อาสวํ ธมฺมํ ปจฺจยา โน-อาสโว ธมฺโม อุปฺปชฺชติ เหตุปจฺจยา.\csromanspacer{} โน-อาสวํ เอกํ ขนฺธํ ปจฺจยา ตโย ขนฺธา จิตฺตสมุฏฺฐานญฺจ รูปํ ฯเปฯ เทฺว ขนฺเธ ฯเปฯ.\csromanspacer{} ปฏิสนฺธิกฺขเณ ฯเปฯ.\csromanspacer{} ขนฺเธ ปจฺจยา วตฺถุ, วตฺถุํ ปจฺจยา ขนฺธา.\csromanspacer{} มหาภูเต ปจฺจยา จิตฺตสมุฏฺฐานํ รูปํ, กฏตฺตารูปํ, อุปาทารูปํ.\csromanspacer{} วตฺถุํ ปจฺจยา โน-อาสวา ขนฺธา. (1)}

\prose{โน-อาสวํ ธมฺมํ ปจฺจยา อาสโว ธมฺโม อุปฺปชฺชติ เหตุปจฺจยา.\csromanspacer{} โน-อาสเว ขนฺเธ ปจฺจยา อาสวา.\csromanspacer{} วตฺถุํ ปจฺจยา อาสวา. (2)}

\prose{โน-อาสวํ ธมฺมํ ปจฺจยา อาสโว จ โน-อาสโว จ ธมฺมา อุปฺปชฺชนฺติ เหตุปจฺจยา.\csromanspacer{} โน-อาสวํ เอกํ ขนฺธํ ปจฺจยา ตโย ขนฺธา อาสวา จ จิตฺตสมุฏฺฐานญฺจ รูปํ ฯเปฯ เทฺว ขนฺเธ ฯเปฯ.\csromanspacer{} วตฺถุํ ปจฺจยา อาสวา สมฺปยุตฺตกา จ ขนฺธา. (3)}

\proseitem{11}{อาสวญฺจ โน-อาสวญฺจ ธมฺมํ ปจฺจยา อาสโว ธมฺโม อุปฺปชฺชติ เหตุปจฺจยา.\csromanspacer{} กามาสวญฺจ สมฺปยุตฺตเก จ ขนฺเธ ปจฺจยา ทิฏฺฐาสโว \csromanfolio{158}อวิชฺชาสโว. (จกฺกํ.) กามาสวญฺจ วตฺถุญฺจ ปจฺจยา ทิฏฺฐาสโว อวิชฺชาสโว. (จกฺกํ.) (1)}

\prose{อาสวญฺจ โน-อาสวญฺจ ธมฺมํ ปจฺจยา โน-อาสโว ธมฺโม อุปฺปชฺชติ เหตุปจฺจยา.\csromanspacer{} โน-อาสวํ เอกํ ขนฺธญฺจ อาสเว จ ปจฺจยา ตโย ขนฺธา จิตฺตสมุฏฺฐานญฺจ รูปํ ฯเปฯ เทฺว ขนฺเธ ฯเปฯ.\csromanspacer{} อาสวญฺจ วตฺถุญฺจ ปจฺจยา โน-อาสวา ขนฺธา. (2)}

\prose{อาสวญฺจ โน-อาสวญฺจ ธมฺมํ ปจฺจยา อาสโว จ โน-อาสโว จ ธมฺมา อุปฺปชฺชนฺติ เหตุปจฺจยา.\csromanspacer{} โน-อาสวํ เอกํ ขนฺธญฺจ กามาสวญฺจ ปจฺจยา ตโย ขนฺธา ทิฏฺฐาสโว อวิชฺชาสโว จิตฺตสมุฏฺฐานญฺจ รูปํ ฯเปฯ เทฺว ขนฺเธ ฯเปฯ. (จกฺกํ.) กามาสวญฺจ วตฺถุญฺจ ปจฺจยา ทิฏฺฐาสโว อวิชฺชาสโว สมฺปยุตฺตกา จ ขนฺธา. (จกฺกํ.\csromanspacer{} สํขิตฺตํ.) (3)}

\csromanheader{2}{1. ปจฺจยานุโลม 2. สงฺขฺยาวาร}
\vspace{\tipitakaparskip}
\csromancenter{เหตุยา นว, อารมฺมเณ นว, อธิปติยา นว ฯเปฯ วิปาเก เอกํ ฯเปฯ อวิคเต นว.}
\csromancenter{อนุโลมํ.}
\csromanheader{2}{2. ปจฺจยปจฺจนีย 1. วิภงฺควาร}
\csromanheader{2}{\textbf{นเหตุ}}
\proseitem{13}{โน-อาสวํ ธมฺมํ ปจฺจยา โน-อาสโว ธมฺโม อุปฺปชฺชติ นเหตุปจฺจยา.\csromanspacer{} อเหตุกํ โน-อาสวํ เอกํ ขนฺธํ ปจฺจยา ตโย ขนฺธา จิตฺตสมุฏฺฐานญฺจ รูปํ ฯเปฯ เทฺว ขนฺเธ ฯเปฯ.\csromanspacer{} อเหตุกปฏิสนฺธิกฺขเณ ฯเปฯ. (ยาว อสญฺญสตฺตา.) จกฺขายตนํ ปจฺจยา จกฺขุวิญฺญาณํ ฯเปฯ.\csromanspacer{} กายายตนํ ปจฺจยา กายวิญฺญาณํ.\csromanspacer{} วตฺถุํ ปจฺจยา อเหตุกา โน-อาสวา ขนฺธา. (1)}

\prose{โน-อาสวํ ธมฺมํ ปจฺจยา อาสโว ธมฺโม อุปฺปชฺชติ นเหตุปจฺจยา.\csromanspacer{} วิจิกิจฺฉาสหคเต อุทฺธจฺจสหคเต ขนฺเธ จ วตฺถุญฺจ ปจฺจยา วิจิกิจฺฉาสหคโต อุทฺธสหคโต โมโห. (2)}

\csromanheader{2}{14. อาสวทุก \textbf{2. ปจฺจยปจฺจนีย 2. สงฺขฺยาวร}}
\proseitem{14}{\csromanfolio{159}นเหตุยา เทฺว, น-อารมฺมเณ ตีณิ, น-อธิปติยา นว ฯเปฯ นกมฺเม ตีณิ ฯเปฯ นวิปฺปยุตฺเต นว, โนนตฺถิยา ตีณิ, โนวิคเต ตีณิ.(เอวํ สพฺเพ คณนา คเณตพฺพา.)}

\vspace{\tipitakaparskip}
\csromancenter{(นิสฺสยวาโร ปจฺจยวารสทิโส.)}
\csromanheader{2}{14. อาสวทุก 5. สํสฏฺฐวาร}
\csromanheader{2}{1-4. ปจฺจยานุโลมาทิ}
\proseitem{15}{อาสวํ ธมฺมํ สํสฏฺโฐ อาสโว ธมฺโม อุปฺปชฺชติ เหตุปจฺจยา.\csromanspacer{} กามาสวํ สํสฏฺโฐ ทิฏฺฐาสโว อวิชฺชาสโว. (จกฺกํ.\csromanspacer{} สํขิตฺตํ.)}

\prose{\textbf{เหตุ}ยา นว, อารมฺมเณ นว, (สพฺพตฺถ นว.) วิปาเก เอกํ ฯเปฯ อวิคเต นว.}

\prose{นเหตุยา เทฺว, น-อธิปติยา นว, นปุเรชาเต นว, นปจฺฉาชาเต นว, น-อาเสวเน นว, นกมฺเม ตีณิ, นวิปาเก นว, นฌาเน เอกํ, นมคฺเค เอกํ, นวิปฺปยุตฺเต นว.}

\vspace{\tipitakaparskip}
\csromancenter{(คณนาปิ สมฺปยุตฺตวาโรปิ สํสฏฺฐวารสทิโส.)}
\csromanheader{2}{14. อาสวทุก 7. ปญฺหาวาร}
\csromanheader{2}{1. ปจฺจยานุโลม 1. วิภงฺควาร}
\csromanheader{2}{\textbf{เหตุ}}
\proseitem{16}{อาสโว ธมฺโม อาสวสฺส ธมฺมสฺส เหตุปจฺจเยน ปจฺจโย.\csromanspacer{} กามาสโว ทิฏฺฐาสวสฺส อวิชฺชาสวสฺส เหตุปจฺจเยน ปจฺจโย.\csromanspacer{} ภวาสโว อวิชฺชาสวสฺส เหตุปจฺจเยน ปจฺจโย. (จกฺกํ.) (1)}

\prose{อาสโว ธมฺโม โน-อาสวสฺส ธมฺมสฺส เหตุปจฺจเยน ปจฺจโย.\csromanspacer{} อาสวา เหตู สมฺปยุตฺตกานํ ขนฺธานํ จิตฺตสมุฏฺฐานานญฺจ รูปานํ เหตุปจฺจเยน ปจฺจโย. (2)}

\prose{\csromanfolio{160}อาสโว ธมฺโม อาสวสฺส จ โน-อาสวสฺส จ ธมฺมสฺส เหตุปจฺจเยน ปจฺจโย.\csromanspacer{} กามาสโว ทิฏฺฐาสวสฺส อวิชฺชาสวสฺส สมฺปยุตฺตกานญฺจ ขนฺธานํ จิตฺตสมุฏฺฐานานญฺจ รูปานํ เหตุปจฺจเยน ปจฺจโย. (3)}

\proseitem{17}{โน-อาสโว ธมฺโม โน-อาสวสฺส ธมฺมสฺส เหตุปจฺจเยน ปจฺจโย.\csromanspacer{} โน-อาสวา เหตู สมฺปยุตฺตกานํ ขนฺธานํ จิตฺตสมุฏฺฐานานญฺจ รูปานํ เหตุปจฺจเยน ปจฺจโย.\csromanspacer{} ปฏิสนฺธิกฺขเณ ฯเปฯ. (1)}

\prose{โน-อาสโว ธมฺโม อาสวสฺส ธมฺมสฺส เหตุปจฺจเยน ปจฺจโย.\csromanspacer{} โน-อาสวา เหตู สมฺปยุตฺตกานํ อาสวานํ เหตุปจฺจเยน ปจฺจโย. (2)}

\prose{โน-อาสโว ธมฺโม อาสวสฺส จ โน-อาสวสฺส จ ธมฺมสฺส เหตุปจฺจเยน ปจฺจโย.\csromanspacer{} โน-อาสวา เหตู สมฺปยุตฺตกานํ ขนฺธานํ อาสวานํ จิตฺตสมุฏฺฐานานญฺจ รูปานํ เหตุปจฺจเยน ปจฺจโย. (3)}

\prose{อาสโว จ โน-อาสโว จ ธมฺมา โน-อาสวสฺส ธมฺมสฺส เหตุปจฺจเยน ปจฺจโย.\csromanspacer{} อาสวา จ โน-อาสวา จ เหตู สมฺปยุตฺตกานํ ขนฺธานํ จิตฺตสมุฏฺฐานานญฺจ รูปานํ เหตุปจฺจเยน ปจฺจโย. (1)}

\csromanheader{2}{\textbf{อารมฺมณ}}
\proseitem{18}{อาสโว ธมฺโม อาสวสฺส ธมฺมสฺส อารมฺมณปจฺจเยน ปจฺจโย.\csromanspacer{} อาสเว อารพฺภ อาสวา อุปฺปชฺชนฺติ. (1)}

\prose{อาสโว ธมฺโม โน-อาสวสฺส ธมฺมสฺส อารมฺมณปจฺจเยน ปจฺจโย.\csromanspacer{} อาสเว อารพฺภ โน-อาสวา ขนฺธา อุปฺปชฺชนฺติ. (2)}

\prose{อาสโว ธมฺโม อาสวสฺส จ โน-อาสวสฺส จ ธมฺมสฺส อารมฺมณปจฺจเยน ปจฺจโย.\csromanspacer{} อาสเว อารพฺภ อาสวา จ สมฺปยุตฺตกา จ ขนฺธา อุปฺปชฺชนฺติ. (3)}

\proseitem{19}{โน-อาสโว ธมฺโม โน-อาสวสฺส ธมฺมสฺส อารมฺมณปจฺจเยน ปจฺจโย.\csromanspacer{} ทานํ ฯเปฯ สีลํ ฯเปฯ อุโปสถกมฺมํ ฯเปฯ.\csromanspacer{} ปุพฺเพ ฯเปฯ.\csromanspacer{} ฌานา ฯเปฯ.\csromanspacer{} อริยา มคฺคา วุฏฺฐหิตฺวา มคฺคํ ปจฺจเวกฺขนฺติ, ผลํ ปจฺจเวกฺขนฺติ, นิพฺพานํ ปจฺจเวกฺขนฺติ.\csromanspacer{} นิพฺพานํ โคตฺรภุสฺส, โวทานสฺส, มคฺคสฺส, ผลสฺส, อาวชฺชนาย อารมฺมณปจฺจเยน ปจฺจโย.\csromanspacer{} อริยา โน-อาสเว ปหีเน กิเลเส ฯเปฯ.}

\vspace{\tipitakaparskip}
\csromancenter{อาสวทุก วิกฺขมฺภิเต กิเลเส ฯเปฯ.\csromanspacer{} ปุพฺเพ สมุทาจิณฺเณ กิเลเส ชานนฺติ.\csromanspacer{} จกฺขุํ ฯเปฯ วตฺถุํ โน-อาสเว ขนฺเธ อนิจฺจโต ฯเปฯ โทมนสฺสํ อุปฺปชฺชติ.\csromanspacer{} ทิพฺเพน จกฺขุนา รูปํ ปสฺสติ.\csromanspacer{} ทิพฺพาย โสตธาตุยา สทฺทํ สุณาติ.\csromanspacer{} เจโตปริยญาเณน โน-อาสวจิตฺตสมงฺคิสฺส จิตฺตํ ชานาติ.\csromanspacer{}\csromanspacer{}อากาสานญฺจายตนํ วิญฺญาณญฺจายตนสฺส ฯเปฯ.\csromanspacer{} อากิญฺจญฺญายตนํ เนวสญฺญานาสญฺญายตนสฺส ฯเปฯ.\csromanspacer{} รูปายตนํ จกฺขุวิญฺญาณสฺส ฯเปฯ.\csromanspacer{} โผฏฺฐพฺพายตนํ กายวิญฺญาณสฺส ฯเปฯ.\csromanspacer{} โน-อาสวา ขนฺธา อิทฺธิวิธญาณสฺส, เจโตปริยญาณสฺส, ปุพฺเพนิวาสานุสฺสติญาณสฺส, ยถากมฺมูปคญาณสฺส, อนาคตํสญาณสฺส, อาวชฺชนาย อารมฺมณปจฺจเยน ปจฺจโย. (1)}
\prose{\csromanfolio{161}โน-อาสโว ธมฺโม อาสวสฺส ธมฺมสฺส อารมฺมณปจฺจเยน ปจฺจโย.\csromanspacer{} ทานํ ทตฺวา ฯเปฯ ตํ อสฺสาเทติ อภินนฺทติ, ตํ อารพฺภ อาสวา อุปฺปชฺชนฺติ.\csromanspacer{} สีลํ ฯเปฯ อุโปสถกมฺมํ ฯเปฯ.\csromanspacer{} ปุพฺเพ ฯเปฯ.\csromanspacer{} ฌานา ฯเปฯ.\csromanspacer{} จกฺขุํ ฯเปฯ วตฺถุํ โน-อาสเว ขนฺเธ อสฺสาเทติ อภินนฺทติ, ตํ อารพฺภ อาสวา อุปฺปชฺชนฺติ. (2)}

\prose{โน-อาสโว ธมฺโม อาสวสฺส จ โน-อาสวสฺส จ ธมฺมสฺส อารมฺมณปจฺจเยน ปจฺจโย.\csromanspacer{} ทานํ ทตฺวา ฯเปฯ. (ทุติยคมนํ.) โน-อาสเว ขนฺเธ อสฺสาเทติ อภินนฺทติ, ตํ อารพฺภ อาสวา จ สมฺปยุตฺตกา จ ขนฺธา อุปฺปชฺชนฺติ. (3)}

\proseitem{20}{อาสโว จ โน-อาสโว จ ธมฺมา อาสวสฺส ธมฺมสฺส อารมฺมณปจฺจเยน ปจฺจโย.\csromanspacer{} อาสเว จ สมฺปยุตฺตเก จ ขนฺเธ อารพฺภ อาสวา อุปฺปชฺชนฺติ. (1)}

\prose{อาสโว จ โน-อาสโว จ ธมฺมา โน-อาสวสฺส ธมฺมสฺส อารมฺมณปจฺจเยน ปจฺจโย.\csromanspacer{} อาสเว จ สมฺปยุตฺตเก จ ขนฺเธ อารพฺภ โน-อาสวา ขนฺธา อุปฺปชฺชนฺติ. (2)}

\prose{อาสโว จ โน-อาสโว จ ธมฺมา อาสวสฺส จ โน-อาสวสฺส จ ธมฺมสฺส อารมฺมณปจฺจเยน ปจฺจโย.\csromanspacer{} อาสเว จ สมฺปยุตฺตเก จ ขนฺเธ อารพฺภ อาสวา จ สมฺปยุตฺตกา จ ขนฺธา อุปฺปชฺชนฺติ. (3)}

\csromanheader{2}{\textbf{อธิปติ}}
\proseitem{21}{\csromanfolio{162}อาสโว ธมฺโม อาสวสฺส ธมฺมสฺส อธิปติปจฺจเยน ปจฺจโย.\csromanspacer{} \textbf{อารมฺมณาธิปติ}\csromandash{}อาสเว ครุํ กตฺวา อาสวา อุปฺปชฺชนฺติ. (ตีณิ อารมฺมณสทิสา, ครุการมฺมณา กาตพฺพา.)}

\prose{โน-อาสโว ธมฺโม โน-อาสวสฺส ธมฺมสฺส อธิปติปจฺจเยน ปจฺจโย.\csromanspacer{} \textbf{อารมฺมณาธิปติ}, สหชาตาธิปติ.\csromanspacer{}\csromanspacer{}\textbf{อารมฺมณาธิปติ}\csromandash{}ทานํ ฯเปฯ สีลํ ฯเปฯ อุโปสถกมฺมํ ฯเปฯ.\csromanspacer{} ปุพฺเพ ฯเปฯ.\csromanspacer{} ฌานา ฯเปฯ.\csromanspacer{} อริยา มคฺคา ฯเปฯ ผลํ ฯเปฯ.\csromanspacer{} นิพฺพานํ ครุํ ฯเปฯ.\csromanspacer{} นิพฺพานํ โคตฺรภุสฺส, โวทานสฺส, มคฺคสฺส, ผลสฺส อธิปติปจฺจเยน ปจฺจโย.\csromanspacer{} จกฺขุํ ฯเปฯ วตฺถุํ โน-อาสเว ขนฺเธ ครุํ กตฺวา อสฺสาเทติ อภินนฺทติ, ตํ ครุํ กตฺวา ราโค อุปฺปชฺชติ, ทิฏฺฐิ อุปฺปชฺชติ.\csromanspacer{} \textbf{สหชาตาธิปติ\csromandash{}}โน-อาสวา อธิปติ สมฺปยุตฺตกานํ ขนฺธานํ จิตฺตสมุฏฺฐานานญฺจ รูปานํ อธิปติปจฺจเยน ปจฺจโย. (1)}

\prose{โน-อาสโว ธมฺโม อาสวสฺส ธมฺมสฺส อธิปติปจฺจเยน ปจฺจโย.\csromanspacer{} \textbf{อารมฺมณาธิปติ}, สหชาตาธิปติ.\csromanspacer{}\csromanspacer{}\textbf{อารมฺมณาธิปติ}\csromandash{}ทานํ ฯเปฯ โน-อาสเว ขนฺเธ ครุํ กตฺวา อสฺสาเทติ อภินนฺทติ, ตํ ครุํ กตฺวา อาสวา อุปฺปชฺชนฺติ.\csromanspacer{} \textbf{สหชาตาธิปติ\csromandash{}}โน-อาสวา อธิปติ สมฺปยุตฺตกานํ อาสวานํ อธิปติปจฺจเยน ปจฺจโย. (2)}

\prose{โน-อาสโว ธมฺโม อาสวสฺส จ โน-อาสวสฺส จ ธมฺมสฺส อธิปติปจฺจเยน ปจฺจโย.\csromanspacer{} \textbf{อารมฺมณาธิปติ}, สหชาตาธิปติ.\csromanspacer{}\csromanspacer{}\textbf{อารมฺมณาธิปติ}\csromandash{}ทานํ ฯเปฯ โน-อาสเว ขนฺเธ ครุํ กตฺวา อสฺสาเทติ อภินนฺทติ, ตํ ครุํ กตฺวา ราโค อุปฺปชฺชติ, ทิฏฺฐิ อุปฺปชฺชติ.\csromanspacer{} \textbf{สหชาตาธิปติ\csromandash{}}โน-อาสวา อธิปติ สมฺปยุตฺตกานํ ขนฺธานํ อาสวานํ จิตฺตสมุฏฺฐานานญฺจ รูปานํ อธิปติปจฺจเยน ปจฺจโย. (3)}

\prose{อาสโว จ โน-อาสโว จ ธมฺมา อาสวสฺส ธมฺมสฺส อธิปติปจฺจเยน ปจฺจโย.\csromanspacer{} \textbf{อารมฺมณาธิปติ}\csromandash{}อาสเว จ สมฺปยุตฺตเก จ ขนฺเธ ครุํ กตฺวา อสฺสาเทติ ฯเปฯ.\csromanspacer{} อาสวา อุปฺปชฺชนฺติ. (ตีณิ.\csromanspacer{} ครุการมฺมณา.)}

\csromanheader{2}{14. อาสวทุก \textbf{อนนฺตร}}
\proseitem{22}{\csromanfolio{163}อาสโว ธมฺโม อาสวสฺส ธมฺมสฺส อนนฺตรปจฺจเยน ปจฺจโย.\csromanspacer{} ปุริมา ปุริมา อาสวา ปจฺฉิมานํ ปจฺฉิมานํ อาสวานํ อนนฺตรปจฺจเยน ปจฺจโย. (1)}

\prose{อาสโว ธมฺโม โน-อาสวสฺส ธมฺมสฺส อนนฺตรปจฺจเยน ปจฺจโย.\csromanspacer{} ปุริมา ปุริมา อาสวา ปจฺฉิมานํ ปจฺฉิมานํ โน-อาสวานํ ขนฺธานํ อนนฺตรปจฺจเยน ปจฺจโย.\csromanspacer{} อาสวา วุฏฺฐานสฺส อนนฺตรปจฺจเยน ปจฺจโย. (2)}

\prose{อาสโว ธมฺโม อาสวสฺส จ โน-อาสวสฺส จ ธมฺมสฺส อนนฺตรปจฺจเยน ปจฺจโย.\csromanspacer{} ปุริมา ปุริมา อาสวา ปจฺฉิมานํ ปจฺฉิมานํ อาสวานํ สมฺปยุตฺตกานญฺจ ขนฺธานํ อนนฺตรปจฺจเยน ปจฺจโย. (3)}

\proseitem{23}{โน-อาสโว ธมฺโม โน-อาสวสฺส ธมฺมสฺส อนนฺตรปจฺจเยน ปจฺจโย.\csromanspacer{} ปุริมา ปุริมา โน-อาสวา ขนฺธา ปจฺฉิมานํ ปจฺฉิมานํ โน-อาสวานํ ขนฺธานํ อนนฺตรปจฺจเยน ปจฺจโย.\csromanspacer{} อนุโลมํ โคตฺรภุสฺส ฯเปฯ ผลสมาปตฺติยา อนนฺตรปจฺจเยน ปจฺจโย. (1)}

\prose{โน-อาสโว ธมฺโม อาสวสฺส ธมฺมสฺส อนนฺตรปจฺจเยน ปจฺจโย.\csromanspacer{} ปุริมา ปุริมา โน-อาสวา ขนฺธา ปจฺฉิมานํ ปจฺฉิมานํ อาสวานํ อนนฺตรปจฺจเยน ปจฺจโย.\csromanspacer{} อาวชฺชนา อาสวานํ อนนฺตรปจฺจเยน ปจฺจโย. (2)}

\prose{โน-อาสโว ธมฺโม อาสวสฺส จ โน-อาสวสฺส จ ธมฺมสฺส อนนฺตรปจฺจเยน ปจฺจโย.\csromanspacer{} ปุริมา ปุริมา โน-อาสวา ขนฺธา ปจฺฉิมานํ ปจฺฉิมานํ อาสวานํ สมฺปยุตฺตกานญฺจ ขนฺธานํ อนนฺตรปจฺจเยน ปจฺจโย. (3)}

\proseitem{24}{อาสโว จ โน-อาสโว จ ธมฺมา อาสวสฺส ธมฺมสฺส อนนฺตรปจฺจเยน ปจฺจโย.\csromanspacer{} ตีณิ.}

\csromanheader{2}{\textbf{สมนนฺตราทิ}}
\proseitem{25}{อาสโว ธมฺโม อาสวสฺส ธมฺมสฺส สมนนฺตรปจฺจเยน ปจฺจโย.\csromanspacer{} สหชาตปจฺจเยน ปจฺจโย.\csromanspacer{} นว.\csromanspacer{} อญฺญมญฺญปจฺจเยน ปจฺจโย.\csromanspacer{} นว.\csromanspacer{} นิสฺสยปจฺจเยน ปจฺจโย.\csromanspacer{} นว. (วตฺถุ จ ทสฺเสตพฺพํ.)}

\csromanheader{2}{\textbf{อุปนิสฺสย}}
\proseitem{26}{\csromanfolio{164}อาสโว ธมฺโม อาสวสฺส ธมฺมสฺส อุปนิสฺสยปจฺจเยน ปจฺจโย.\csromanspacer{} อารมฺมณูปนิสฺสโย, อนนฺตรูปนิสฺสโย, ปกตูปนิสฺสโย ฯเปฯ.\csromanspacer{} ปกตูปนิสฺสโย\csromandash{}อาสวา อาสวานํ อุปนิสฺสยปจฺจเยน ปจฺจโย. (ตีณิ.)}

\proseitem{27}{โน-อาสโว ธมฺโม โน-อาสวสฺส ธมฺมสฺส อุปนิสฺสยปจฺจเยน ปจฺจโย.\csromanspacer{} อารมฺมณูปนิสฺสโย, อนนฺตรูปนิสฺสโย, ปกตูปนิสฺสโย ฯเปฯ.\csromanspacer{} ปกตูปนิสฺสโย\csromandash{}สทฺธํ อุปนิสฺสาย ทานํ เทติ ฯเปฯ สมาปตฺติํ อุปฺปาเทติ.\csromanspacer{} มานํ ชปฺเปติ.\csromanspacer{} ทิฏฺฐิํ คณฺหาติ.\csromanspacer{} สีลํ ฯเปฯ.\csromanspacer{} เสนาสนํ อุปนิสฺสาย ทานํ เทติ ฯเปฯ สํฆํ ภินฺทติ.\csromanspacer{} สทฺธา ฯเปฯ.\csromanspacer{} เสนาสนํ สทฺธาย ฯเปฯ ผลสมาปตฺติยา อุปนิสฺสยปจฺจเยน ปจฺจโย. (1)}

\prose{โน-อาสโว ธมฺโม อาสวสฺส ธมฺมสฺส อุปนิสฺสยปจฺจเยน ปจฺจโย.\csromanspacer{} อารมฺมณูปนิสฺสโย, อนนฺตรูปนิสฺสโย, ปกตูปนิสฺสโย ฯเปฯ.\csromanspacer{} ปกตูปนิสฺสโย\csromandash{}สทฺธํ อุปนิสฺสาย มานํ ชปฺเปติ.\csromanspacer{} ทิฏฺฐิํ คณฺหาติ.\csromanspacer{} สีลํ ฯเปฯ สํฆํ ภินฺทติ.\csromanspacer{} สทฺธา ฯเปฯ.\csromanspacer{} เสนาสนํ ราคสฺส ฯเปฯ ปตฺถนาย อุปนิสฺสยปจฺจเยน ปจฺจโย. (2)}

\prose{โน-อาสโว ธมฺโม อาสวสฺส จ โน-อาสวสฺส จ ธมฺมสฺส อุปนิสฺสยปจฺจเยน ปจฺจโย.\csromanspacer{} อารมฺมณูปนิสฺสโย, อนนฺตรูปนิสฺสโย, ปกตูปนิสฺสโย ฯเปฯ.\csromanspacer{} ปกตูปนิสฺสโย\csromandash{}สทฺธํ อุปนิสฺสาย มานํ ชปฺเปติ.\csromanspacer{} ทิฏฺฐิํ คณฺหาติ.\csromanspacer{} สีลํ ฯเปฯ.\csromanspacer{} เสนาสนํ อุปนิสฺสาย ฯเปฯ ปาณํ หนติ ฯเปฯ.\csromanspacer{} สํฆํ ภินฺทติ.\csromanspacer{} สทฺธา ฯเปฯ.\csromanspacer{} เสนาสนํ ราคสฺส ฯเปฯ ปตฺถนาย อุปนิสฺสยปจฺจเยน ปจฺจโย. (3)}

\prose{อาสโว จ โน-อาสโว จ ธมฺมา อาสวสฺส ธมฺมสฺส อุปนิสฺสยปจฺจเยน ปจฺจโย.\csromanspacer{} อารมฺมณูปนิสฺสโย, อนนฺตรูปนิสฺสโย, ปกตูปนิสฺสโย ฯเปฯ.\csromanspacer{} ปกตูปนิสฺสโย\csromandash{}ตีณิ.}

\csromanheader{2}{\textbf{ปุเรชาต}}
\proseitem{28}{โน-อาสโว ธมฺโม โน-อาสวสฺส ธมฺมสฺส ปุเรชาตปจฺจเยน ปจฺจโย.\csromanspacer{} อารมฺมณปุเรชาตํ, วตฺถุปุเรชาตํ.\csromanspacer{}\csromanspacer{}\textbf{อารมฺมณปุเรชาตํ\csromandash{}}}

\vspace{\tipitakaparskip}
\csromancenter{อาสวทุก จกฺขุํ ฯเปฯ วตฺถุํ. (เอวํ วิตฺถาเรตพฺพํ.) โผฏฺฐพฺพายตนํ กายวิญฺญาณสฺส ปุเรชาตปจฺจเยน ปจฺจโย.\csromanspacer{} \textbf{วตฺถุปุเรชาตํ}\csromandash{} จกฺขายตนํ จกฺขุวิญฺญาณสฺส ฯเปฯ กายายตนํ กายวิญฺญาณสฺส.\csromanspacer{} วตฺถุ โน-อาสวานํ ขนฺธานํ ปุเรชาตปจฺจเยน ปจฺจโย. (1)}
\prose{\csromanfolio{165}โน-อาสโว ธมฺโม อาสวสฺส ธมฺมสฺส ปุเรชาตปจฺจเยน ปจฺจโย.\csromanspacer{} \textbf{อารมฺมณปุเรชาตํ}, วตฺถุปุเรชาตํ.\csromanspacer{}\csromanspacer{}\textbf{อารมฺมณปุเรชาตํ}\csromandash{}จกฺขุํ ฯเปฯ วตฺถุํ อสฺสาเทติ อภินนฺทติ, ตํ อารพฺภ อาสวา อุปฺปชฺชนฺติ.\csromanspacer{} \textbf{วตฺถุปุเรชาตํ}\csromandash{}วตฺถุ อาสวานํ ปุเรชาตปจฺจเยน ปจฺจโย. (2)}

\prose{โน-อาสโว ธมฺโม อาสวสฺส จ โน-อาสวสฺส จ ธมฺมสฺส ปุเรชาตปจฺจเยน ปจฺจโย.\csromanspacer{} \textbf{อารมฺมณปุเรชาตํ}, วตฺถุปุเรชาตํ.\csromanspacer{}\csromanspacer{}\textbf{อารมฺมณปุเรชาตํ}\csromandash{}จกฺขุํ ฯเปฯ วตฺถุํ อสฺสาเทติ อภินนฺทติ, ตํ อารพฺภ อาสวา จ สมฺปยุตฺตกา จ ขนฺธา อุปฺปชฺชนฺติ.\csromanspacer{} \textbf{วตฺถุปุเรชาตํ}\csromandash{} วตฺถุ อาสวานญฺจ อาสวสมฺปยุตฺตกานญฺจ ขนฺธานํ ปุเรชาตปจฺจเยน ปจฺจโย. (3)}

\csromanheader{2}{\textbf{ปจฺฉาชาต}}
\proseitem{29}{อาสโว ธมฺโม โน-อาสวสฺส ธมฺมสฺส ปจฺฉาชาตปจฺจเยน ปจฺจโย.\csromanspacer{} ปจฺฉาชาตา อาสวา ปุเรชาตสฺส อิมสฺส กายสฺส ปจฺฉาชาตปจฺจเยน ปจฺจโย. (1)}

\prose{โน-อาสโว ธมฺโม โน-อาสวสฺส ธมฺมสฺส ปจฺฉาชาตปจฺจเยน ปจฺจโย.\csromanspacer{} ปจฺฉาชาตา โน-อาสวา ขนฺธา ปุเรชาตสฺส อิมสฺส กายสฺส ปจฺฉาชาตปจฺจเยน ปจฺจโย. (1)}

\prose{อาสโว จ โน-อาสโว จ ธมฺมา โน-อาสวสฺส ธมฺมสฺส ปจฺฉาชาตปจฺจเยน ปจฺจโย.\csromanspacer{} ปจฺฉาชาตา อาสวา จ สมฺปยุตฺตกา จ ขนฺธา ปุเรชาตสฺส อิมสฺส กายสฺส ปจฺฉาชาตปจฺจเยน ปจฺจโย. (1)}

\csromanheader{2}{\textbf{อาเสวน}}
\vspace{\tipitakaparskip}
\csromancenter{อาสโว ธมฺโม อาสวสฺส ธมฺมสฺส อาเสวนปจฺจเยน ปจฺจโย.\csromanspacer{} นว.}
\csromanheader{2}{\textbf{กมฺม}}
\proseitem{31}{\csromanfolio{166}โน-อาสโว ธมฺโม โน-อาสวสฺส ธมฺมสฺส กมฺมปจฺจเยน ปจฺจโย.\csromanspacer{} \textbf{สหชาตา}, นานากฺขณิกา.\csromanspacer{}\csromanspacer{}\textbf{สหชาตา} โน-อาสวา เจตนา สมฺปยุตฺตกานํ ขนฺธานํ จิตฺตสมุฏฺฐานานญฺจ รูปานํ กมฺมปจฺจเยน ปจฺจโย.\csromanspacer{} ปฏิสนฺธิกฺขเณ ฯเปฯ.\csromanspacer{} \textbf{นานากฺขณิกา} โน-อาสวา เจตนา วิปากานํ ขนฺธานํ กฏตฺตา จ รูปานํ กมฺมปจฺจเยน ปจฺจโย. (1)}

\prose{โน-อาสโว ธมฺโม อาสวสฺส ธมฺมสฺส กมฺมปจฺจเยน ปจฺจโย.\csromanspacer{} โน-อาสวา เจตนา สมฺปยุตฺตกานํ อาสวานํ กมฺมปจฺจเยน ปจฺจโย.}

\prose{โน-อาสโว ธมฺโม อาสวสฺส จ โน-อาสวสฺส จ ธมฺมสฺส กมฺมปจฺจเยน ปจฺจโย.\csromanspacer{} โน-อาสวา เจตนา สมฺปยุตฺตกานํ ขนฺธานํ อาสวานํ จิตฺตสมุฏฺฐานานญฺจ รูปานํ กมฺมปจฺจเยน ปจฺจโย. (3)}

\csromanheader{2}{\textbf{วิปากาหาร}}
\proseitem{32}{โน-อาสโว ธมฺโม โน-อาสวสฺส ธมฺมสฺส วิปากปจฺจเยน ปจฺจโย.\csromanspacer{} เอกํ.\csromanspacer{} อาหารปจฺจเยน ปจฺจโย.\csromanspacer{} โน-อาสวา อาหารา สมฺปยุตฺตกานํ ขนฺธานํ จิตฺตสมุฏฺฐานานญฺจ รูปานํ อาหารปจฺจเยน ปจฺจโย.\csromanspacer{} ปฏิสนฺธิกฺขเณ ฯเปฯ.\csromanspacer{} กพฬีกาโร อาหาโร อิมสฺส กายสฺส อาหารปจฺจเยน ปจฺจโย. (1)}

\prose{โน-อาสโว ธมฺโม อาสวสฺส ธมฺมสฺส อาหารปจฺจเยน ปจฺจโย.\csromanspacer{} โน-อาสวา อาหารา สมฺปยุตฺตกานํ อาสวานํ อาหารปจฺจเยน ปจฺจโย. (2)}

\prose{โน-อาสโว ธมฺโม อาสวสฺส จ โน-อาสวสฺส จ ธมฺมสฺส อาหารปจฺจเยน ปจฺจโย.\csromanspacer{} โน-อาสวา อาหารา สมฺปยุตฺตกานํ ขนฺธานํ อาสวานํ จิตฺตสมุฏฺฐานานญฺจ รูปานํ อาหารปจฺจเยน ปจฺจโย. (3)}

\csromanheader{2}{\textbf{อินฺทฺริยาทิ}}
\proseitem{33}{โน-อาสโว ธมฺโม โน-อาสวสฺส ธมฺมสฺส อินฺทฺริยปจฺจเยน ปจฺจโย.\csromanspacer{} โน-อาสวา อินฺทฺริยา สมฺปยุตฺตกานํ ขนฺธานํ จิตฺตสมุฏฺฐานานญฺจ}

\vspace{\tipitakaparskip}
\csromancenter{อาสวทุก รูปานํ อินฺทฺริยปจฺจเยน ปจฺจโย.\csromanspacer{} ปฏิสนฺธิกฺขเณ ฯเปฯ.\csromanspacer{} จกฺขุนฺทฺริยํ จกฺขุวิญฺญาณสฺส ฯเปฯ.\csromanspacer{} กายินฺทฺริยํ กายวิญฺญาณสฺส ฯเปฯ.\csromanspacer{} รูปชีวิตินฺทฺริยํ กฏตฺตารูปานํ อินฺทฺริยปจฺจเยน ปจฺจโย.\csromanspacer{} ตีณิ.\csromanspacer{} ฌานปจฺจเยน ปจฺจโย.\csromanspacer{} ตีณิ.\csromanspacer{} มคฺคปจฺจเยน ปจฺจโย.\csromanspacer{} นว.\csromanspacer{} สมฺปยุตฺตปจฺจเยน ปจฺจโย.\csromanspacer{} นว.}
\csromanheader{2}{\textbf{วิปฺปยุตฺต}}
\proseitem{34}{\csromanfolio{167}อาสโว ธมฺโม โน-อาสวสฺส ธมฺมสฺส วิปฺปยุตฺตปจฺจเยน ปจฺจโย.\csromanspacer{} สหชาตํ, ปจฺฉาชาตํ.\csromanspacer{}\csromanspacer{}\textbf{สหชาตา} อาสวา จิตฺตสมุฏฺฐานานํ รูปานํ วิปฺปยุตฺตปจฺจเยน ปจฺจโย.\csromanspacer{} \textbf{ปจฺฉาชาตา} อาสวา ปุเรชาตสฺส อิมสฺส กายสฺส วิปฺปยุตฺตปจฺจเยน ปจฺจโย. (1)}

\proseitem{35}{โน-อาสโว ธมฺโม โน-อาสวสฺส ธมฺมสฺส วิปฺปยุตฺตปจฺจเยน ปจฺจโย.\csromanspacer{} สหชาตํ, ปุเรชาตํ, ปจฺฉาชาตํ.\csromanspacer{}\csromanspacer{}\textbf{สหชาตา} โน-อาสวา ขนฺธา จิตฺตสมุฏฺฐานานํ รูปานํ วิปฺปยุตฺตปจฺจเยน ปจฺจโย.\csromanspacer{} ปฏิสนฺธิกฺขเณ ฯเปฯ.\csromanspacer{} ขนฺธา วตฺถุสฺส วิปฺปยุตฺตปจฺจเยน ปจฺจโย.\csromanspacer{} วตฺถุ ขนฺธานํ วิปฺปยุตฺตปจฺจเยน ปจฺจโย.\csromanspacer{} \textbf{ปุเรชาตํ} จกฺขายตนํ จกฺขุวิญฺญาณสฺส ฯเปฯ กายายตนํ กายวิญฺญาณสฺส ฯเปฯ วตฺถุ โน-อาสวานํ ขนฺธานํ วิปฺปยุตฺตปจฺจเยน ปจฺจโย.\csromanspacer{} \textbf{ปจฺฉาชาตา} โน-อาสวา ขนฺธา ปุเรชาตสฺส อิมสฺส กายสฺส วิปฺปยุตฺตปจฺจเยน ปจฺจโย. (1)}

\prose{โน-อาสโว ธมฺโม อาสวสฺส ธมฺมสฺส วิปฺปยุตฺตปจฺจเยน ปจฺจโย.\csromanspacer{} \textbf{ปุเรชาตํ} วตฺถุ อาสวานํ วิปฺปยุตฺตปจฺจเยน ปจฺจโย. (2)}

\prose{โน-อาสโว ธมฺโม อาสวสฺส จ โน-อาสวสฺส จ ธมฺมสฺส วิปฺปยุตฺตปจฺจเยน ปจฺจโย.\csromanspacer{} \textbf{ปุเรชาตํ} วตฺถุ อาสวานํ สมฺปยุตฺตกานญฺจ ขนฺธานํ วิปฺปยุตฺตปจฺจเยน ปจฺจโย. (3)}

\proseitem{36}{อาสโว จ โน-อาสโว จ ธมฺมา โน-อาสวสฺส ธมฺมสฺส วิปฺปยุตฺตปจฺจเยน ปจฺจโย.\csromanspacer{} สหชาตํ, ปจฺฉาชาตํ.\csromanspacer{}\csromanspacer{}\textbf{สหชาตา} อาสวา จ สมฺปยุตฺตกา จ ขนฺธา จิตฺตสมุฏฺฐานานํ รูปานํ วิปฺปยุตฺตปจฺจเยน ปจฺจโย.\csromanspacer{} \textbf{ปจฺฉาชาตา} อาสวา จ สมฺปยุตฺตกา จ ขนฺธา ปุเรชาตสฺส อิมสฺส กายสฺส วิปฺปยุตฺตปจฺจเยน ปจฺจโย. (1)}

\csromanheader{2}{\textbf{อตฺถิ}}
\proseitem{37}{\csromanfolio{168}อาสโว ธมฺโม อาสวสฺส ธมฺมสฺส อตฺถิปจฺจเยน ปจฺจโย.\csromanspacer{} กามาสโว ทิฏฺฐาสวสฺส อวิชฺชาสวสฺส อตฺถิปจฺจเยน ปจฺจโย. (จกฺกํ.) (1)}

\prose{อาสโว ธมฺโม โน-อาสวสฺส ธมฺมสฺส อตฺถิปจฺจเยน ปจฺจโย.\csromanspacer{} สหชาตํ, ปจฺฉาชาตํ.\csromanspacer{}\csromanspacer{}\textbf{สหชาตา} อาสวา สมฺปยุตฺตกานํ ขนฺธานํ จิตฺตสมุฏฺฐานานญฺจ รูปานํ อตฺถิปจฺจเยน ปจฺจโย.\csromanspacer{} \textbf{ปจฺฉาชาตา} อาสวา ปุเรชาตสฺส อิมสฺส กายสฺส อตฺถิปจฺจเยน ปจฺจโย. (2)}

\prose{อาสโว ธมฺโม อาสวสฺส จ โน-อาสวสฺส จ ธมฺมสฺส อตฺถิปจฺจเยน ปจฺจโย.\csromanspacer{} กามาสโว ทิฏฺฐาสวสฺส อวิชฺชาสวสฺส สมฺปยุตฺตกานํ ขนฺธานํ จิตฺตสมุฏฺฐานานญฺจ รูปานํ อตฺถิปจฺจเยน ปจฺจโย. (จกฺกํ.) (3)}

\proseitem{38}{โน-อาสโว ธมฺโม โน-อาสวสฺส ธมฺมสฺส อตฺถิปจฺจเยน ปจฺจโย.\csromanspacer{} สหชาตํ, ปุเรชาตํ, ปจฺฉาชาตํ, อาหารํ, อินฺทฺริยํ.\csromanspacer{}\csromanspacer{}\textbf{สหชาโต} โน-อาสโว เอโก ขนฺโธ ติณฺณนฺนํ ขนฺธานํ ฯเปฯ. (ยาว อสญฺญสตฺตา.) \textbf{ปุเรชาตํ} จกฺขุํ ฯเปฯ วตฺถุํ อนิจฺจโต ฯเปฯ โทมนสฺสํ อุปฺปชฺชติ.\csromanspacer{} ทิพฺเพน จกฺขุนา รูปํ ปสฺสติ.\csromanspacer{} ทิพฺพาย โสตธาตุยา สทฺทํ สุณาติ.\csromanspacer{} รูปายตนํ จกฺขุวิญฺญาณสฺส ฯเปฯ.\csromanspacer{} โผฏฺฐพฺพายตนํ กายวิญฺญาณสฺส ฯเปฯ.\csromanspacer{} จกฺขายตนํ จกฺขุวิญฺญาณสฺส ฯเปฯ.\csromanspacer{} กายายตนํ กายวิญฺญาณสฺส ฯเปฯ.\csromanspacer{} วตฺถุ โน-อาสวานํ ขนฺธานํ อตฺถิปจฺจเยน ปจฺจโย.\csromanspacer{} \textbf{ปจฺฉาชาตา} โน-อาสวา ขนฺธา ปุเรชาตสฺส อิมสฺส กายสฺส อตฺถิปจฺจเยน ปจฺจโย.\csromanspacer{} \textbf{กพฬีกาโร อาหาโร} อิมสฺส กายสฺส อตฺถิปจฺจเยน ปจฺจโย.\csromanspacer{} \textbf{รูปชีวิตินฺทฺริยํ} กฏตฺตารูปานํ อตฺถิปจฺจเยน ปจฺจโย. (1)}

\prose{โน-อาสโว ธมฺโม อาสวสฺส ธมฺมสฺส อตฺถิปจฺจเยน ปจฺจโย.\csromanspacer{} สหชาตํ, ปุเรชาตํ.\csromanspacer{}\csromanspacer{}\textbf{สหชาตา} โน-อาสวา ขนฺธา อาสวานํ อตฺถิปจฺจเยน ปจฺจโย.\csromanspacer{} \textbf{ปุเรชาตํ} จกฺขุํ ฯเปฯ วตฺถุํ อสฺสาเทติ อภินนฺทติ, ตํ อารพฺภ อาสวา อุปฺปชฺชนฺติ.\csromanspacer{} วตฺถุ อาสวานํ อตฺถิปจฺจเยน ปจฺจโย. (2)}

\csromanheader{2}{14. อาสวทุก}
\prose{\csromanfolio{169}โน-อาสโว ธมฺโม อาสวสฺส จ โน-อาสวสฺส จ ธมฺมสฺส อตฺถิปจฺจเยน ปจฺจโย.\csromanspacer{} สหชาตํ, ปุเรชาตํ.\csromanspacer{}\csromanspacer{}\textbf{สหชาโต} โน-อาสโว เอโก ขนฺโธ ติณฺณนฺนํ ขนฺธานํ อาสวานํ จิตฺตสมุฏฺฐานานญฺจ รูปานํ อตฺถิปจฺจเยน ปจฺจโย ฯเปฯ เทฺว ขนฺธา ฯเปฯ. (จกฺกํ.) (3)}

\proseitem{39}{อาสโว จ โน-อาสโว จ ธมฺมา อาสวสฺส ธมฺมสฺส อตฺถิปจฺจเยน ปจฺจโย.\csromanspacer{} สหชาตํ, ปุเรชาตํ.\csromanspacer{}\csromanspacer{}\textbf{สหชาโต} กามาสโว จ สมฺปยุตฺตกา จ ขนฺธา ทิฏฺฐาสวสฺส อวิชฺชาสวสฺส อตฺถิปจฺจเยน ปจฺจโย. (จกฺกํ.) กามาสโว จ วตฺถุ จ ทิฏฺฐาสวสฺส อวิชฺชาสวสฺส อตฺถิปจฺจเยน ปจฺจโย. (จกฺกํ.) (1)}

\prose{อาสโว จ โน-อาสโว จ ธมฺมา โน-อาสวสฺส ธมฺมสฺส อตฺถิปจฺจเยน ปจฺจโย.\csromanspacer{} สหชาตํ, ปุเรชาตํ, ปจฺฉาชาตํ, อาหารํ, อินฺทฺริยํ.\csromanspacer{}\csromanspacer{}\textbf{สหชาโต} โน-อาสโว เอโก ขนฺโธ จ อาสวา จ ติณฺณนฺนํ ขนฺธานํ จิตฺตสมุฏฺฐานานญฺจ รูปานํ อตฺถิปจฺจเยน ปจฺจโย.\csromanspacer{} สหชาตา อาสวา จ มหาภูตา จ จิตฺตสมุฏฺฐานานํ รูปานํ อตฺถิปจฺจเยน ปจฺจโย.\csromanspacer{} อาสวา จ วตฺถุ จ โน-อาสวานํ ขนฺธานํ อตฺถิปจฺจเยน ปจฺจโย.\csromanspacer{} \textbf{ปจฺฉาชาตา} อาสวา จ กพฬีกาโร อาหาโร จ อิมสฺส กายสฺส อตฺถิปจฺจเยน ปจฺจโย.\csromanspacer{} \textbf{ปจฺฉาชาตา} อาสวา จ รูปชีวิตินฺทฺริยญฺจ กฏตฺตารูปานํ อตฺถิปจฺจเยน ปจฺจโย. (2)}

\prose{อาสโว จ โน-อาสโว จ ธมฺมา อาสวสฺส จ โน-อาสวสฺส จ ธมฺมสฺส อตฺถิปจฺจเยน ปจฺจโย.\csromanspacer{} สหชาตํ, ปุเรชาตํ.\csromanspacer{}\csromanspacer{}\textbf{สหชาโต} โน-อาสโว เอโก ขนฺโธ จ กามาสโว จ ติณฺณนฺนํ ขนฺธานํ ทิฏฺฐาสวสฺส อวิชฺชาสวสฺส จิตฺตสมุฏฺฐานานญฺจ รูปานํ อตฺถิปจฺจเยน ปจฺจโย ฯเปฯ เทฺว ขนฺธา จ ฯเปฯ. (จกฺกํ.) \textbf{สหชาโต} กามาสโว จ วตฺถุ จ ทิฏฺฐาสวสฺส อวิชฺชาสวสฺส สมฺปยุตฺตกานญฺจ ขนฺธานํ อตฺถิปจฺจเยน ปจฺจโย. (จกฺกํ.) (3)}

\csromanheader{2}{1. ปจฺจยานุโลม 2. สงฺขฺยาวาร}
\csromanheader{2}{\textbf{สุทฺธ}}
\proseitem{40}{เหตุยา สตฺต, อารมฺมเณ นว, อธิปติยา นว, อนนฺตเร นว, สมนนฺตเร นว, สหชาเต นว, อญฺญมญฺเญ นว, นิสฺสเย นว, \csromanfolio{170}อุปนิสฺสเย นว, ปุเรชาเต ตีณิ, ปจฺฉาชาเต ตีณิ, อาเสวเน นว, กมฺเม ตีณิ, วิปาเก เอกํ, อาหาเร ตีณิ ฯเปฯ มคฺเค นว, สมฺปยุตฺเต นว, วิปฺปยุตฺเต ปญฺจ, อตฺถิยา นว, นตฺถิยา นว, วิคเต นว, อวิคเต นว.}

\vspace{\tipitakaparskip}
\csromancenter{อนุโลมํ.}
\csromanheader{2}{2. ปจฺจนียุทฺธาร}
\proseitem{41}{อาสโว ธมฺโม อาสวสฺส ธมฺมสฺส อารมฺมณปจฺจเยน ปจฺจโย.\csromanspacer{} สหชาตปจฺจเยน ปจฺจโย.\csromanspacer{} อุปนิสฺสยปจฺจเยน ปจฺจโย. (1)}

\prose{อาสโว ธมฺโม โน-อาสวสฺส ธมฺมสฺส อารมฺมณปจฺจเยน ปจฺจโย.\csromanspacer{} สหชาตปจฺจเยน ปจฺจโย.\csromanspacer{} อุปนิสฺสยปจฺจเยน ปจฺจโย.\csromanspacer{} ปจฺฉาชาตปจฺจเยน ปจฺจโย. (2)}

\prose{อาสโว ธมฺโม อาสวสฺส จ โน-อาสวสฺส จ ธมฺมสฺส อารมฺมณปจฺจเยน ปจฺจโย.\csromanspacer{} สหชาตปจฺจเยน ปจฺจโย.\csromanspacer{} อุปนิสฺสยปจฺจเยน ปจฺจโย. (3)}

\proseitem{42}{โน-อาสโว ธมฺโม โน-อาสวสฺส ธมฺมสฺส อารมฺมณปจฺจเยน ปจฺจโย.\csromanspacer{} สหชาตปจฺจเยน ปจฺจโย.\csromanspacer{} อุปนิสฺสยปจฺจเยน ปจฺจโย.\csromanspacer{} ปุเรชาตปจฺจเยน ปจฺจโย.\csromanspacer{} ปจฺฉาชาตปจฺจเยน ปจฺจโย.\csromanspacer{} กมฺมปจฺจเยน ปจฺจโย.\csromanspacer{} อาหารปจฺจเยน ปจฺจโย.\csromanspacer{} อินฺทฺริยปจฺจเยน ปจฺจโย. (1)}

\prose{โน-อาสโว ธมฺโม อาสวสฺส ธมฺมสฺส อารมฺมณปจฺจเยน ปจฺจโย.\csromanspacer{} สหชาตปจฺจเยน ปจฺจโย.\csromanspacer{} อุปนิสฺสยปจฺจเยน ปจฺจโย.\csromanspacer{} ปุเรชาตปจฺจเยน ปจฺจโย. (2)}

\prose{โน-อาสโว ธมฺโม อาสวสฺส จ โน-อาสวสฺส จ ธมฺมสฺส อารมฺมณปจฺจเยน ปจฺจโย.\csromanspacer{} สหชาตปจฺจเยน ปจฺจโย.\csromanspacer{} อุปนิสฺสยปจฺจเยน ปจฺจโย.\csromanspacer{} ปุเรชาตปจฺจเยน ปจฺจโย. (3)}

\proseitem{43}{อาสโว จ โน-อาสโว จ ธมฺมา อาสวสฺส ธมฺมสฺส อารมฺมณปจฺจเยน ปจฺจโย.\csromanspacer{} สหชาตปจฺจเยน ปจฺจโย.\csromanspacer{} อุปนิสฺสยปจฺจเยน ปจฺจโย. (1)}

\csromanheader{2}{14. อาสวทุก}
\prose{\csromanfolio{171}อาสโว จ โน-อาสโว จ ธมฺมา โน-อาสวสฺส ธมฺมสฺส อารมฺมณปจฺจเยน ปจฺจโย.\csromanspacer{} สหชาตปจฺจเยน ปจฺจโย.\csromanspacer{} อุปนิสฺสยปจฺจเยน ปจฺจโย.\csromanspacer{} ปจฺฉาชาตปจฺจเยน ปจฺจโย. (2)}

\prose{อาสโว จ โน-อาสโว จ ธมฺมา อาสวสฺส จ โน-อาสวสฺส จ ธมฺมสฺส อารมฺมณปจฺจเยน ปจฺจโย.\csromanspacer{} สหชาตปจฺจเยน ปจฺจโย.\csromanspacer{} อุปนิสฺสยปจฺจเยน ปจฺจโย. (3)}

\csromanheader{2}{2. ปจฺจยปจฺจนีย 2. สงฺขฺยาวาร}
\proseitem{44}{นเหตุยา นว, น-อารมฺมเณ นว, น-อธิปติยา นว, (สพฺพตฺถ นว,) โน-อวิคเต นว.}

\vspace{\tipitakaparskip}
\csromancenter{ปจฺจนียํ.}
\csromanheader{2}{3. ปจฺจยานุโลมปจฺจนีย}
\csromanheader{2}{\textbf{เหตุทุก}}
\proseitem{45}{เหตุปจฺจยา น-อารมฺมเณ สตฺต, น-อธิปติยา สตฺต, น-อนนฺตเร สตฺต, นสมนนฺตเร สตฺต, น-อญฺญมญฺเญ ตีณิ, น-อุปนิสฺสเย สตฺต, (สพฺพตฺถ สตฺต,) นมคฺเค สตฺต, นสมฺปยุตฺเต ตีณิ, นวิปฺปยุตฺเต สตฺต, โนนตฺถิยา สตฺต, โนวิคเต สตฺต.}

\vspace{\tipitakaparskip}
\csromancenter{อนุโลมปจฺจนียํ.}
\csromanheader{2}{4. ปจฺจยปจฺจนียานุโลม}
\csromanheader{2}{\textbf{นเหตุทุก}}
\proseitem{46}{นเหตุปจฺจยา อารมฺมเณ นว, อธิปติยา นว. (อนุโลมปทา ปริปุณฺณา) อวิคเต นว.}

\vspace{\tipitakaparskip}
\csromancenter{ปจฺจนียานุโลมํ.}
\nitthitam{อาสวทุกํ นิฏฺฐิตํ.}
\csromanmatikaanchor{mtk.19}
\csromantocmarkref{section}{15. สาสวทุก}{mtk.19}
\bookmark[dest={mtk.19},level=1]{15. สาสวทุก}
\csromanheader{1}{15. สาสวทุก 1. ปฏิจฺจวาร}
\csromanheader{2}{\textbf{เหตุ}}
\proseitem{47}{\csromanfolio{172}สาสวํ ธมฺมํ ปฏิจฺจ สาสโว ธมฺโม อุปฺปชฺชติ เหตุปจฺจยา.\csromanspacer{} สาสวํ เอกํ ขนฺธํ ปฏิจฺจ ตโย ขนฺธา จิตฺตสมุฏฺฐานญฺจ รูปํ ฯเปฯ เทฺว ขนฺเธ ฯเปฯ.\csromanspacer{} ปฏิสนฺธิกฺขเณ ฯเปฯ.\csromanspacer{} ขนฺเธ ปฏิจฺจ วตฺถุ, วตฺถุํ ปฏิจฺจ ขนฺธา.\csromanspacer{} เอกํ มหาภูตํ ฯเปฯ มหาภูเต ปฏิจฺจ จิตฺตสมุฏฺฐานํ รูปํ, กฏตฺตารูปํ, อุปาทารูปํ. (1)}

\prose{อนาสวํ ธมฺมํ ปฏิจฺจ อนาสโว ธมฺโม อุปฺปชฺชติ เหตุปจฺจยา.\csromanspacer{} อนาสวํ เอกํ ขนฺธํ ปฏิจฺจ ตโย ขนฺธา ฯเปฯ เทฺว ขนฺเธ ฯเปฯ. (1)}

\prose{อนาสวํ ธมฺมํ ปฏิจฺจ สาสโว ธมฺโม อุปฺปชฺชติ เหตุปจฺจยา.\csromanspacer{} อนาสเว ขนฺเธ ปฏิจฺจ จิตฺตสมุฏฺฐานํ รูปํ. (2)}

\prose{อนาสวํ ธมฺมํ ปฏิจฺจ สาสโว จ อนาสโว จ ธมฺมา อุปฺปชฺชนฺติ เหตุปจฺจยา.\csromanspacer{} อนาสวํ เอกํ ขนฺธํ ปฏิจฺจ ตโย ขนฺธา จิตฺตสมุฏฺฐานญฺจ รูปํ ฯเปฯ เทฺว ขนฺเธ ฯเปฯ. (3)}

\prose{สาสวญฺจ อนาสวญฺจ ธมฺมํ ปฏิจฺจ สาสโว ธมฺโม อุปฺปชฺชติ เหตุปจฺจยา.\csromanspacer{} อนาสเว ขนฺเธ จ มหาภูเต จ ปฏิจฺจ จิตฺตสมุฏฺฐานํ รูปํ. (1)}

\prose{(ยถา จูฬนฺตรทุเก โลกิยทุกํ, เอวํ กาตพฺพํ, นินฺนานากรณํ.)}

\nitthitam{สาสวทุกํ นิฏฺฐิตํ.}
\csromanmatikaanchor{mtk.20}
\csromantocmarkref{section}{16. อาสวสมฺปยุตฺตทุก}{mtk.20}
\bookmark[dest={mtk.20},level=1]{16. อาสวสมฺปยุตฺตทุก}
\csromanheader{1}{16. อาสวสมฺปยุตฺตทุก 1. ปฏิจฺจวาร}
\csromanheader{2}{1. ปจฺจยานุโลม 1. วิภงฺควาร}
\csromanheader{2}{\textbf{เหตุ}}
\proseitem{48}{อาสวสมฺปยุตฺตํ ธมฺมํ ปฏิจฺจ อาสวสมฺปยุตฺโต ธมฺโม อุปฺปชฺชติ เหตุปจฺจยา.\csromanspacer{} อาสวสมฺปยุตฺตํ เอกํ ขนฺธํ ปฏิจฺจ ตโย ขนฺธา ฯเปฯ เทฺว ขนฺเธ ฯเปฯ. (1)}

\csromanheader{2}{16. อาสวสมฺปยุตฺตทุก}
\prose{\csromanfolio{173}อาสวสมฺปยุตฺตํ ธมฺมํ ปฏิจฺจ อาสววิปฺปยุตฺโต ธมฺโม อุปฺปชฺชติ เหตุปจฺจยา.\csromanspacer{} อาสวสมฺปยุตฺเต ขนฺเธ ปฏิจฺจ จิตฺตสมุฏฺฐานํ รูปํ.\csromanspacer{} โทมนสฺสสหคเต ขนฺเธ ปฏิจฺจ โมโห จิตฺตสมุฏฺฐานญฺจ รูปํ. (2)}

\prose{อาสวสมฺปยุตฺตํ ธมฺมํ ปฏิจฺจ อาสวสมฺปยุตฺโต จ อาสววิปฺปยุตฺโต จ ธมฺมา อุปฺปชฺชนฺติ เหตุปจฺจยา.\csromanspacer{} อาสวสมฺปยุตฺตํ เอกํ ขนฺธํ ปฏิจฺจ ตโย ขนฺธา จิตฺตสมุฏฺฐานญฺจ รูปํ ฯเปฯ เทฺว ขนฺเธ ฯเปฯ.\csromanspacer{} โทมนสฺสสหคตํ เอกํ ขนฺธํ ปฏิจฺจ ตโย ขนฺธา โมโห จิตฺตสมุฏฺฐานญฺจ รูปํ ฯเปฯ เทฺว ขนฺเธ ฯเปฯ. (3)}

\proseitem{49}{อาสววิปฺปยุตฺตํ ธมฺมํ ปฏิจฺจ อาสววิปฺปยุตฺโต ธมฺโม อุปฺปชฺชติ เหตุปจฺจยา.\csromanspacer{} อาสววิปฺปยุตฺตํ เอกํ ขนฺธํ ปฏิจฺจ ตโย ขนฺธา จิตฺตสมุฏฺฐานญฺจ รูปํ ฯเปฯ เทฺว ขนฺเธ ฯเปฯ.\csromanspacer{} โทมนสฺสสหคตํ วิจิกิจฺฉาสหคตํ อุทฺธจฺจสหคตํ โมหํ ปฏิจฺจ จิตฺตสมุฏฺฐานํ รูปํ.\csromanspacer{} ปฏิสนฺธิกฺขเณ ฯเปฯ.\csromanspacer{} ขนฺเธ ปฏิจฺจ วตฺถุ, วตฺถุํ ปฏิจฺจ ขนฺธา.\csromanspacer{} เอกํ มหาภูตํ ปฏิจฺจ ตโย มหาภูตา, ตโย มหาภูเต ปฏิจฺจ เอกํ มหาภูตํ, เทฺว มหาภูเต ปฏิจฺจ เทฺว มหาภูตา.\csromanspacer{} มหาภูเต ปฏิจฺจ จิตฺตสมุฏฺฐานํ รูปํ, กฏตฺตารูปํ, อุปาทารูปํ. (1)}

\prose{อาสววิปฺปยุตฺตํ ธมฺมํ ปฏิจฺจ อาสวสมฺปยุตฺโต ธมฺโม อุปฺปชฺชติ เหตุปจฺจยา.\csromanspacer{} โทมนสฺสสหคตํ วิจิกิจฺฉาสหคตํ อุทฺธจฺจสหคตํ โมหํ ปฏิจฺจ สมฺปยุตฺตกา ขนฺธา. (2)}

\prose{อาสววิปฺปยุตฺตํ ธมฺมํ ปฏิจฺจ อาสวสมฺปยุตฺโต จ อาสววิปฺปยุตฺโต จ ธมฺมา อุปฺปชฺชนฺติ เหตุปจฺจยา.\csromanspacer{} โทมนสฺสสหคตํ วิจิกิจฺฉาสหคตํ อุทฺธจฺจสหคตํ โมหํ ปฏิจฺจ สมฺปยุตฺตกา ขนฺธา จิตฺตสมุฏฺฐานญฺจ รูปํ. (3)}

\proseitem{50}{อาสวสมฺปยุตฺตญฺจ อาสววิปฺปยุตฺตญฺจ ธมฺมํ ปฏิจฺจ อาสวสมฺปยุตฺโต ธมฺโม อุปฺปชฺชติ เหตุปจฺจยา.\csromanspacer{} โทมนสฺสสหคตํ วิจิกิจฺฉาสหคตํ อุทฺธจฺจสหคตํ เอกํ ขนฺธญฺจ โมหญฺจ ปฏิจฺจ ตโย ขนฺธา ฯเปฯ เทฺว ขนฺเธ ฯเปฯ. (1)}

\prose{อาสวสมฺปยุตฺตญฺจ อาสววิปฺปยุตฺตญฺจ ธมฺมํ ปฏิจฺจ อาสววิปฺปยุตฺโต ธมฺโม อุปฺปชฺชติ เหตุปจฺจยา.\csromanspacer{} อาสวสมฺปยุตฺเต ขนฺเธ จ มหาภูเต จ ปฏิจฺจ จิตฺตสมุฏฺฐานํ รูปํ.\csromanspacer{} โทมนสฺสสหคเต วิจิกิจฺฉาสหคเต อุทฺธจฺจสหคเต ขนฺเธ จ โมหญฺจ ปฏิจฺจ จิตฺตสมุฏฺฐานํ รูปํ. (2)}

\prose{\csromanfolio{174}อาสวสมฺปยุตฺตญฺจ อาสววิปฺปยุตฺตญฺจ ธมฺมํ ปฏิจฺจ อาสวสมฺปยุตฺโต จ อาสววิปฺปยุตฺโต จ ธมฺมา อุปฺปชฺชนฺติ เหตุปจฺจยา.\csromanspacer{} โทมนสฺสสหคตํ วิจิกิจฺฉาสหคตํ อุทฺธจฺจสหคตํ เอกํ ขนฺธญฺจ โมหญฺจ ปฏิจฺจ ตโย ขนฺธา จิตฺตสมุฏฺฐานญฺจ รูปํ.\csromanspacer{} เทฺว ขนฺเธ ฯเปฯ. (3)}

\csromanheader{2}{\textbf{อารมฺมณ}}
\proseitem{51}{อาสวสมฺปยุตฺตํ ธมฺมํ ปฏิจฺจ อาวสมฺปยุตฺโต ธมฺโม อุปฺปชฺชติ อารมฺมณปจฺจยา.\csromanspacer{} อาสวสมฺปยุตฺตํ เอกํ ขนฺธํ ปฏิจฺจ ตโย ขนฺธา ฯเปฯ เทฺว ขนฺเธ ฯเปฯ. (1)}

\prose{อาสวสมฺปยุตฺตํ ธมฺมํ ปฏิจฺจ อาสววิปฺปยุตฺโต ธมฺโม อุปฺปชฺชติ อารมฺมณปจฺจยา.\csromanspacer{} โทมนสฺสสหคเต วิจิกิจฺฉาสหคเต อุทฺธจฺจสหคเต ขนฺเธ ปฏิจฺจ โมโห. (2)}

\prose{อาสวสมฺปยุตฺตํ ธมฺมํ ปฏิจฺจ อาสวสมฺปยุตฺโต จ อาสววิปฺปยุตฺโต จ ธมฺมา อุปฺปชฺชนฺติ อารมฺมณปจฺจยา.\csromanspacer{} โทมนสฺสสหคตํ วิจิกิจฺฉาสหคตํ อุทฺธจฺจสหคตํ เอกํ ขนฺธํ ปฏิจฺจ ตโย ขนฺธา โมโห จ ฯเปฯ เทฺว ขนฺเธ ฯเปฯ. (3)}

\proseitem{52}{อาสววิปฺปยุตฺตํ ธมฺมํ ปฏิจฺจ อาสววิปฺปยุตฺโต ธมฺโม อุปฺปชฺชติ อารมฺมณปจฺจยา.\csromanspacer{} อาสววิปฺปยุตฺตํ เอกํ ขนฺธํ ปฏิจฺจ ตโย ขนฺธา ฯเปฯ เทฺว ขนฺเธ ฯเปฯ.\csromanspacer{} ปฏิสนฺธิกฺขเณ ฯเปฯ.\csromanspacer{} วตฺถุํ ปฏิจฺจ ขนฺธา. (1)}

\prose{อาสววิปฺปยุตฺตํ ธมฺมํ ปฏิจฺจ อาสวสมฺปยุตฺโต ธมฺโม อุปฺปชฺชติ อารมฺมณปจฺจยา.\csromanspacer{} โทมนสฺสสหคตํ วิจิกิจฺฉาสหคตํ อุทฺธจฺจสหคตํ โมหํ ปฏิจฺจ สมฺปยุตฺตกา ขนฺธา. (2)}

\prose{อาสวสมฺปยุตฺตญฺจ อาสววิปฺปยุตฺตญฺจ ธมฺมํ ปฏิจฺจ อาสวสมฺปยุตฺโต ธมฺโม อุปฺปชฺชติ อารมฺมณปจฺจยา.\csromanspacer{} โทมนสฺสสหคตํ วิจิกิจฺฉาสหคตํ อุทฺธจฺจสหคตํ เอกํ ขนฺธญฺจ โมหญฺจ ปฏิจฺจ ตโย ขนฺธา ฯเปฯ เทฺว ขนฺเธ ฯเปฯ. (1)}

\csromanheader{2}{\textbf{อธิปติ}}
\proseitem{53}{อาสวสมฺปยุตฺตํ ธมฺมํ ปฏิจฺจ อาสวสมฺปยุตฺโต ธมฺโม อุปฺปชฺชติ อธิปติปจฺจยา.\csromanspacer{} ตีณิ.}

\csromanheader{2}{16. อาสวสมฺปยุตฺตทุก}
\prose{\csromanfolio{175}อาสววิปฺปยุตฺตํ ธมฺมํ ปฏิจฺจ อาสววิปฺปยุตฺโต ธมฺโม อุปฺปชฺชติ อธิปติปจฺจยา.\csromanspacer{} อาสววิปฺปยุตฺตํ เอกํ ขนฺธํ ปฏิจฺจ ตโย ขนฺธา จิตฺตสมุฏฺฐานญฺจ รูปํ ฯเปฯ เทฺว ขนฺเธ ฯเปฯ.\csromanspacer{} โทมนสฺสสหคตํ โมหํ ปฏิจฺจ จิตฺตสมุฏฺฐานํ รูปํ.\csromanspacer{} เอกํ มหาภูตํ ฯเปฯ มหาภูเต ปฏิจฺจ จิตฺตสมุฏฺฐานํ รูปํ, อุปาทารูปํ. (1)}

\prose{อาสววิปฺปยุตฺตํ ธมฺมํ ปฏิจฺจ อาสวสมฺปยุตฺโต ธมฺโม อุปฺปชฺชติ อธิปติปจฺจยา.\csromanspacer{} โทมนสฺสสหคตํ โมหํ ปฏิจฺจ สมฺปยุตฺตกา ขนฺธา. (2)}

\prose{อาสววิปฺปยุตฺตํ ธมฺมํ ปฏิจฺจ อาสวสมฺปยุตฺโต จ อาสววิปฺปยุตฺโต จ ธมฺมา อุปฺปชฺชนฺติ อธิปติปจฺจยา.\csromanspacer{} โทมนสฺสสหคตํ โมหํ ปฏิจฺจ สมฺปยุตฺตกา ขนฺธา จิตฺตสมุฏฺฐานญฺจ รูปํ. (3)}

\proseitem{54}{อาสวสมฺปยุตฺตญฺจ อาสววิปฺปยุตฺตญฺจ ธมฺมํ ปฏิจฺจ อาสวสมฺปยุตฺโต ธมฺโม อุปฺปชฺชติ อธิปติปจฺจยา.\csromanspacer{} โทมนสฺสสหคตํ เอกํ ขนฺธญฺจ โมหญฺจ ปฏิจฺจ ตโย ขนฺธา ฯเปฯ เทฺว ขนฺเธ ฯเปฯ. (1)}

\prose{อาสวสมฺปยุตฺตญฺจ อาสววิปฺปยุตฺตญฺจ ธมฺมํ ปฏิจฺจ อาสววิปฺปยุตฺโต ธมฺโม อุปฺปชฺชติ อธิปติปจฺจยา.\csromanspacer{} อาสวสมฺปยุตฺเต ขนฺเธ จ มหาภูเต จ ปฏิจฺจ จิตฺตสมุฏฺฐานํ รูปํ.\csromanspacer{} โทมนสฺสสหคเต ขนฺเธ จ โมหญฺจ ปฏิจฺจ จิตฺตสมุฏฺฐานํ รูปํ.}

\prose{อาสวสมฺปยุตฺตญฺจ อาสววิปฺปยุตฺตญฺจ ธมฺมํ ปฏิจฺจ อาสวสมฺปยุตฺโต จ อาสววิปฺปยุตฺโต จ ธมฺมา อุปฺปชฺชนฺติ อธิปติปจฺจยา.\csromanspacer{} โทมนสฺสสหคตํ เอกํ ขนฺธญฺจ โมหญฺจ ปฏิจฺจ ตโย ขนฺธา จิตฺตสมุฏฺฐานญฺจ รูปํ ฯเปฯ เทฺว ขนฺเธ ฯเปฯ. (3) (เอวํ สพฺเพ ปจฺจยา วิตฺถาเรตพฺพา.\csromanspacer{} สํขิตฺตํ.)}

\csromanheader{2}{1. ปจฺจยานุโลม 2. สงฺขฺยาวาร}
\csromanheader{2}{\textbf{สุทฺธ}}
\proseitem{55}{เหตุยา นว, อารมฺมเณ ฉ, อธิปติยา นว, อนนฺตเร ฉ, สมนนฺตเร ฉ, สหชาเต นว, อญฺญมญฺเญ ฉ, นิสฺสเย นว, อุปนิสฺสเย ฉ, ปุเรชาเต ฉ, อาเสวเน ฉ, กมฺเม นว, วิปาเก \csromanfolio{176}เอกํ, อาหาเร นว, อินฺทฺริเย นว, ฌาเน นว, มคฺเค นว, สมฺปยุตฺเต ฉ, วิปฺปยุตฺเต นว, อตฺถิยา นว, นตฺถิยา ฉ, วิคเต ฉ, อวิคเต นว.}

\vspace{\tipitakaparskip}
\csromancenter{อนุโลมํ.}
\csromanheader{2}{2. ปจฺจยปจฺจนีย 1. วิภงฺควาร}
\csromanheader{2}{\textbf{นเหตุ}}
\proseitem{56}{อาสวสมฺปยุตฺตํ ธมฺมํ ปฏิจฺจ อาสววิปฺปยุตฺโต ธมฺโม อุปฺปชฺชติ นเหตุปจฺจยา.\csromanspacer{} วิจิกิจฺฉาสหคเต อุทฺธจฺจสหคเต ขนฺเธ ปฏิจฺจ วิจิกิจฺฉาสหคโต อุทฺธจฺจสหคโต โมโห. (1)}

\prose{อาสววิปฺปยุตฺตํ ธมฺมํ ปฏิจฺจ อาสววิปฺปยุตฺโต ธมฺโม อุปฺปชฺชติ นเหตุปจฺจยา.\csromanspacer{} อเหตุกํ อาสววิปฺปยุตฺตํ เอกํ ขนฺธํ ปฏิจฺจ ตโย ขนฺธา จิตฺตสมุฏฺฐานญฺจ รูปํ ฯเปฯ เทฺว ขนฺเธ ฯเปฯ.\csromanspacer{} อเหตุกปฏิสนฺธิกฺขเณ ฯเปฯ. (ยาว อสญฺญสตฺตา.) (1)}

\csromanheader{2}{\textbf{น-อารมฺมณ}}
\proseitem{57}{อาสวสมฺปยุตฺตํ ธมฺมํ ปฏิจฺจ อาสววิปฺปยุตฺโต ธมฺโม อุปฺปชฺชติ น-อารมฺมณปจฺจโย.\csromanspacer{} อาสวสมฺปยุตฺเต ขนฺเธ ปฏิจฺจ จิตฺตสมุฏฺฐานํ รูปํ. (1)}

\prose{อาสววิปฺปยุตฺตํ ธมฺมํ ปฏิจฺจ อาสววิปฺปยุตฺโต ธมฺโม อุปฺปชฺชติ น-อารมฺมณปจฺจยา.\csromanspacer{} อาสววิปยุตฺเต ขนฺเธ ปฏิจฺจ จิตฺตสมุฏฺฐานํ รูปํ.\csromanspacer{} โทมนสฺสสหคตํ วิจิกิจฺฉาสหคตํ อุทฺธจฺจสหคตํ โมหํ ปฏิจฺจ จิตฺตสมุฏฺฐานํ รูปํ.\csromanspacer{} ปฏิสนฺธิกฺขเณ ฯเปฯ.\csromanspacer{} ขนฺเธ ปฏิจฺจ วตฺถุ.\csromanspacer{} เอกํ มหาภูตํ ฯเปฯ. (ยาว อสญฺญสตฺตา.) (1)}

\prose{อาสวสมฺปยุตฺตญฺจ อาสววิปฺปยุตฺตญฺจ ธมฺมํ ปฏิจฺจ อาสววิปฺปยุตฺโต ธมฺโม อุปฺปชฺชติ น-อารมฺมณปจฺจยา.\csromanspacer{} อาสวสมฺปยุตฺเต ขนฺเธ จ มหาภูเต จ ปฏิจฺจ จิตฺตสมุฏฺฐานํ รูปํ.\csromanspacer{} โทมนสฺสสหคเต วิจิกิจฺฉาสหคเต อุทฺธจฺจสหคเต ขนฺเธ จ โมหญฺจ ปฏิจฺจ จิตฺตสมุฏฺฐานํ รูปํ. (1)}

\csromanheader{2}{16. อาสวสมฺปยุตฺตทุก \textbf{น-อธิปติ}}
\proseitem{58}{\csromanfolio{177}อาสวสมฺปยุตฺตํ ธมฺมํ ปฏิจฺจ อาสวสมฺปยุตฺโต ธมฺโม อุปฺปชฺชติ น-อธิปติปจฺจยา. (สํขิตฺตํ.)}

\csromanheader{2}{\textbf{นปุเรชาตาทิ}}
\proseitem{59}{อาสวสมฺปยุตฺตํ ธมฺมํ ปฏิจฺจ อาสวสมฺปยุตฺโต ธมฺโม อุปฺปชฺชติ นปุเรชาตปจฺจยา.\csromanspacer{} อรูเป อาสวสมฺปยุตฺตํ เอกํ ขนฺธํ ปฏิจฺจ ตโย ขนฺธา ฯเปฯ เทฺว ขนฺเธ ฯเปฯ. (1)}

\prose{อาสวสมฺปยุตฺตํ ธมฺมํ ปฏิจฺจ อาสววิปฺปยุตฺโต ธมฺโม อุปฺปชฺชติ นปุเรชาตปจฺจยา.\csromanspacer{} อรูเป วิจิกิจฺฉาสหคเต อุทฺธจฺจสหคเต ขนฺเธ ปฏิจฺจ วิจิกิจฺฉาสหคโต อุทฺธจฺจสหคโต โมโห.\csromanspacer{} อาสวสมฺปยุตฺเต ขนฺเธ ปฏิจฺจ จิตฺตสมุฏฺฐานํ รูปํ. (2)}

\prose{อาสวสมฺปยุตฺตํ ธมฺมํ ปฏิจฺจ อาสวสมฺปยุตฺโต จ อาสววิปฺปยุตฺโต จ ธมฺมา อุปฺปชฺชนฺติ นปุเรชาตปจฺจยา.\csromanspacer{} อรูเป วิจิกิจฺฉาสหคตํ อุทฺธจฺจสหคตํ เอกํ ขนฺธํ ปฏิจฺจ ตโย ขนฺธา โมโห จ ฯเปฯ เทฺว ขนฺเธ ฯเปฯ. (3)}

\proseitem{60}{อาสววิปฺปยุตฺตํ ธมฺมํ ปฏิจฺจ อาสววิปฺปยุตฺโต ธมฺโม อุปฺปชฺชติ นปุเรชาตปจฺจยา.\csromanspacer{} อรูเป อาสววิปฺปยุตฺตํ เอกํ ขนฺธํ ปฏิจฺจ ตโย ขนฺธา ฯเปฯ เทฺว ขนฺเธ ฯเปฯ.\csromanspacer{} อาสววิปฺปยุตฺเต ขนฺเธ ปฏิจฺจ จิตฺตสมุฏฺฐานํ รูปํ.\csromanspacer{} โทมนสฺสสหคตํ วิจิกิจฺฉาสหคตํ อุทฺธจฺจสหคตํ โมหํ ปฏิจฺจ จิตฺตสมุฏฺฐานํ รูปํ.\csromanspacer{} ปฏิสนฺธิกฺขเณ ฯเปฯ. (ยาว อสญฺญสตฺตา.) (1)}

\prose{อาสววิปฺปยุตฺตํ ธมฺมํ ปฏิจฺจ อาสวสมฺปยุตฺโต ธมฺโม อุปฺปชฺชติ นปุเรชาตปจฺจยา.\csromanspacer{} อรูเป วิจิกิจฺฉาสหคตํ อุทฺธจฺจสหคตํ โมหํ ปฏิจฺจ สมฺปยุตฺตกา ขนฺธา. (2)}

\proseitem{61}{อาสวสมฺปยุตฺตญฺจ อาสววิปฺปยุตฺตญฺจ ธมฺมํ ปฏิจฺจ อาสวสมฺปยุตฺโต ธมฺโม อุปฺปชฺชติ นปุเรชาตปจฺจยา.\csromanspacer{} อรูเป วิจิกิจฺฉาสหคตํ อุทฺธจฺจสหคตํ เอกํ ขนฺธญฺจ โมหญฺจ ปฏิจฺจ ตโย ขนฺธา ฯเปฯ เทฺว ขนฺเธ จ ฯเปฯ. (1)}

\prose{อาสวสมฺปยุตฺตญฺจ อาสววิปฺปยุตฺตญฺจ ธมฺมํ ปฏิจฺจ อาสววิปฺปยุตฺโต ธมฺโม อุปฺปชฺชติ นปุเรชาตปจฺจยา.\csromanspacer{} อาสวสมฺปยุตฺเต ขนฺเธ จ มหาภูเต \csromanfolio{178}จ ปฏิจฺจ จิตฺตสมุฏฺฐานํ รูปํ.\csromanspacer{} โทมนสฺสสหคเต วิจิกิจฺฉาสหคเต อุทฺธจฺจสหคเต ขนฺเธ จ โมหญฺจ ปฏิจฺจ จิตฺตสมุฏฺฐานํ รูปํ. (2) (นปจฺฉาชาตปจฺจยา นว, น-อาเสวนปจฺจยา นว.)}

\csromanheader{2}{\textbf{นกมฺมาทิ}}
\proseitem{62}{อาสวสมฺปยุตฺตํ ธมฺมํ ปฏิจฺจ อาสวสมฺปยุตฺโต ธมฺโม อุปฺปชฺชติ นกมฺมปจฺจยา.\csromanspacer{} อาสวสมฺปยุตฺเต ขนฺเธ ปฏิจฺจ สมฺปยุตฺตกา เจตนา. (1)}

\prose{อาสววิปฺปยุตฺตํ ธมฺมํ ปฏิจฺจ อาสววิปฺปยุตฺโต ธมฺโม อุปฺปชฺชติ นกมฺมปจฺจยา.\csromanspacer{} อาสววิปฺปยุตฺเต ขนฺเธ ปฏิจฺจ วิปฺปยุตฺตกา เจตนา. (1)}

\prose{อาสววิปฺปยุตฺตํ ธมฺมํ ปฏิจฺจ อาสวสมฺปยุตฺโต ธมฺโม อุปฺปชฺชติ นกมฺมปจฺจยา.\csromanspacer{} โทมนสฺสสหคตํ วิจิกิจฺฉาสหคตํ อุทฺธจฺจสหคตํ โมหํ ปฏิจฺจ สมฺปยุตฺตกา เจตนา. (2)}

\prose{อาสวสมฺปยุตฺตญฺจ อาสววิปฺปยุตฺตญฺจ ธมฺมํ ปฏิจฺจ อาสวสมฺปยุตฺโต ธมฺโม อุปฺปชฺชติ นกมฺมปจฺจยา.\csromanspacer{} โทมนสฺสสหคเต วิจิกิจฺฉาสหคเต อุทฺธจฺจสหคเต ขนฺเธ จ โมหญฺจ ปฏิจฺจ สมฺปยุตฺตกา เจตนา.\csromanspacer{} นวิปากปจฺจยา.\csromanspacer{} น-อาหารปจฺจยา.\csromanspacer{} น-อินฺทฺริยปจฺจยา.\csromanspacer{} นฌานปจฺจยา.\csromanspacer{} นมคฺคปจฺจยา.\csromanspacer{} นสมฺปยุตฺตปจฺจยา.\csromanspacer{} นวิปฺปยุตฺตปจฺจยา.\csromanspacer{} โนนตฺถิปจฺจยา.\csromanspacer{} โนวิคตปจฺจยา.}

\csromanheader{2}{2. ปจฺจยปจฺจนีย 2. สงฺขฺยาวาร}
\csromanheader{2}{\textbf{สุทฺธ}}
\proseitem{63}{นเหตุยา เทฺว, น-อารมฺมเณ ตีณิ, น-อธิปติยา นว, น-อนนฺตเร ตีณิ, นสมนนฺตเร ตีณิ, น-อญฺญมญฺเญ ตีณิ, น-อุปนิสฺสเย ตีณิ, นปุเรชาเต สตฺต, นปจฺฉาชาเต นว, น-อาเสวเน นว, นกมฺเม จตฺตาริ, นวิปาเก นว, น-อาหาเร เอกํ, น-อินฺทฺริเย เอกํ, นฌาเน เอกํ, นมคฺเค เอกํ, นสมฺปยุตฺเต ตีณิ, นวิปฺปยุตฺเต ฉ, โนนตฺถิยา ตีณิ, โนวิคเต ตีณิ.}

\vspace{\tipitakaparskip}
\csromancenter{ปจฺจนียํ.}
\csromanheader{2}{16. อาสวสมฺปยุตฺตทุก \textbf{3. ปจฺจยานุโลมปจฺจนีย}}
\csromanheader{2}{\textbf{เหตุทุก}}
\proseitem{64}{\csromanfolio{179}\textbf{เหตุ}ปจฺจยา น-อารมฺมเณ ตีณิ ฯเปฯ นปุเรชาเต ฉ ฯเปฯ นวิปาเก นว ฯเปฯ นสมฺปยุตฺเต ตีณิ, นวิปฺปยุตฺเต จตฺตาริ, โนนตฺถิยา ตีณิ, โนวิคเต ตีณิ.}

\vspace{\tipitakaparskip}
\csromancenter{อนุโลมปจฺจนียํ.}
\csromanheader{2}{4. ปจฺจยปจฺจนียานุโลม}
\csromanheader{2}{\textbf{นเหตุทุก}}
\proseitem{65}{นเหตุปจฺจยา อารมฺมเณ เทฺว, อนนฺตเร เทฺว, สมนนฺตเร เทฺว ฯเปฯ กมฺเม เทฺว, วิปาเก เอกํ, อาหาเร เทฺว ฯเปฯ มคฺเค เอกํ, สมฺปยุตฺเต เทฺว, วิปฺปยุตฺเต เทฺว ฯเปฯ อวิคเต เทฺว.}

\vspace{\tipitakaparskip}
\csromancenter{(สหชาตวาโร ปฏิจฺจวารสทิโส.)}
\csromanheader{2}{16. อาสวสมฺปยุตฺตทุก 3. ปจฺจยวาร}
\csromanheader{2}{1. ปจฺจยานุโลม 1. วิภงฺควาร}
\csromanheader{2}{\textbf{เหตุ}}
\proseitem{66}{อาสวสมฺปยุตฺตํ ธมฺมํ ปจฺจยา อาสวสมฺปยุตฺโต ธมฺโม อุปฺปชฺชติ เหตุปจฺจยา.\csromanspacer{} ตีณิ. (ปฏิจฺจวารสทิโส.)}

\prose{อาสววิปฺปยุตฺตํ ธมฺมํ ปจฺจยา อาสววิปฺปยุตฺโต ธมฺโม อุปฺปชฺชติ เหตุปจฺจยา.\csromanspacer{} อาสววิปฺปยุตฺตํ เอกํ ขนฺธํ ปจฺจยา ตโย ขนฺธา จิตฺตสมุฏฺฐานญฺจ รูปํ ฯเปฯ เทฺว ขนฺเธ ฯเปฯ.\csromanspacer{} โทมนสฺสสหคตํ วิจิกิจฺฉาสหคตํ อุทฺธจฺจสหคตํ โมหํ ปจฺจยา จิตฺตสมุฏฺฐานํ รูปํ.\csromanspacer{} ปฏิสนฺธิกฺขเณ ฯเปฯ.\csromanspacer{} ขนฺเธ ปจฺจยา วตฺถุ, วตฺถุํ ปจฺจยา ขนฺธา.\csromanspacer{} เอกํ มหาภูตํ ฯเปฯ มหาภูเต ปจฺจยา จิตฺตสมุฏฺฐานํ \csromanfolio{180}รูปํ, กฏตฺตารูปํ, อุปาทารูปํ.\csromanspacer{} วตฺถุํ ปจฺจยา อาสววิปฺปยุตฺตา ขนฺธา.\csromanspacer{} วตฺถุํ ปจฺจยา โทมนสฺสสหคโต โมโห. (1)}

\prose{อาสววิปฺปยุตฺตํ ธมฺมํ ปจฺจยา อาสวสมฺปยุตฺโต ธมฺโม อุปฺปชฺชติ เหตุปจฺจยา.\csromanspacer{} วตฺถุํ ปจฺจยา อาสวสมฺปยุตฺตกา ขนฺธา.\csromanspacer{} โทมนสฺสสหคตํ วิจิกิจฺฉาสหคตํ อุทฺธจฺจสหคตํ โมหํ ปจฺจยา สมฺปยุตฺตกา ขนฺธา. (2)}

\prose{อาสววิปฺปยุตฺตํ ธมฺมํ ปจฺจยา อาสวสมฺปยุตฺโต จ อาสววิปฺปยุตฺโต จ ธมฺมา อุปฺปชฺชนฺติ เหตุปจฺจยา.\csromanspacer{} วตฺถุํ ปจฺจยา อาสวสมฺปยุตฺตกา ขนฺธา.\csromanspacer{} มหาภูเต ปจฺจยา จิตฺตสมุฏฺฐานํ รูปํ.\csromanspacer{} โทมนสฺสสหคตํ วิจิกิจฺฉาสหคตํ อุทฺธจฺจสหคตํ โมหํ ปจฺจยา สมฺปยุตฺตกา ขนฺธา จิตฺตสมุฏฺฐานญฺจ รูปํ.\csromanspacer{} วตฺถุํ ปจฺจยา โทมนสฺสสหคตา ขนฺธา จ โมโห จ. (3)}

\proseitem{67}{อาสวสมฺปยุตฺตญฺจ อาสววิปฺปยุตฺตญฺจ ธมฺมํ ปจฺจยา อาสวสมฺปยุตฺโต ธมฺโม อุปฺปชฺชติ เหตุปจฺจยา.\csromanspacer{} อาสวสมฺปยุตฺตํ เอกํ ขนฺธญฺจ วตฺถุญฺจ ปจฺจยา ตโย ขนฺธา ฯเปฯ เทฺว ขนฺเธ จ ฯเปฯ.\csromanspacer{} โทมนสฺสสหคตํ วิจิกิจฺฉาสหคตํ อุทฺธจฺจสหคตํ เอกํ ขนฺธญฺจ โมหญฺจ ปจฺจยา ตโย ขนฺธา ฯเปฯ เทฺว ขนฺเธ จ ฯเปฯ. (1)}

\prose{อาสวสมฺปยุตฺตญฺจ อาสววิปฺปยุตฺตญฺจ ธมฺมํ ปจฺจยา อาสววิปฺปยุตฺโต ธมฺโม อุปฺปชฺชติ เหตุปจฺจยา.\csromanspacer{} อาสวสมฺปยุตฺเต ขนฺเธ จ มหาภูเต จ ปจฺจยา จิตฺตสมุฏฺฐานํ รูปํ.\csromanspacer{} โทมนสฺสสหคเต วิจิกิจฺฉาสหคเต อุทฺธจฺจสหคเต ขนฺเธ จ โมหญฺจ ปจฺจยา จิตฺตสมุฏฺฐานํ รูปํ.\csromanspacer{} โทมนสฺสสหคเต ขนฺเธ จ วตฺถุญฺจ ปจฺจยา โทมนสฺสสหคโต โมโห. (2)}

\prose{อาสวสมฺปยุตฺตญฺจ อาสววิปฺปยุตฺตญฺจ ธมฺมํ ปจฺจยา อาสวสมฺปยุตฺโต จ อาสววิปฺปยุตฺโต จ ธมฺมา อุปฺปชฺชนฺติ เหตุปจฺจยา.\csromanspacer{} อาสวสมฺปยุตฺตํ เอกํ ขนฺธญฺจ วตฺถุญฺจ ปจฺจยา ตโย ขนฺธา ฯเปฯ เทฺว ขนฺเธ ฯเปฯ.\csromanspacer{} อาสวสมฺปยุตฺเต ขนฺเธ จ มหาภูเต จ ปจฺจยา จิตฺตสมุฏฺฐานํ รูปํ.\csromanspacer{} โทมนสฺสสหคตํ วิจิกิจฺฉาสหคตํ อุทฺธจฺจสหคตํ เอกํ ขนฺธญฺจ โมหญฺจ ปจฺจยา ตโย ขนฺธา จิตฺตสมุฏฺฐานญฺจ รูปํ ฯเปฯ เทฺว ขนฺเธ ฯเปฯ.\csromanspacer{} โทมนสฺสสหคตํ เอกํ ขนฺธญฺจ วตฺถุญฺจ ปจฺจยา ตโย ขนฺธา โมโห จ ฯเปฯ เทฺว ขนฺเธ ฯเปฯ. (3)}

\csromanheader{2}{16. อาสวสมฺปยุตฺตทุก \textbf{อารมฺมณาทิ}}
\proseitem{68}{\csromanfolio{181}อาสวสมฺปยุตฺตํ ธมฺมํ ปจฺจยา อาสวสมฺปยุตฺโต ธมฺโม อุปฺปชฺชติ อารมฺมณปจฺจยา.\csromanspacer{} ตีณิ. (ปฏิจฺจวารสทิสา.)}

\prose{อาสววิปฺปยุตฺตํ ธมฺมํ ปจฺจยา อาสววิปฺปยุตฺโต ธมฺโม อุปฺปชฺชติ อารมฺมณปจฺจยา.\csromanspacer{} อาสววิปฺปยุตฺตํ เอกํ ขนฺธํ ปจฺจยา ตโย ขนฺธา ฯเปฯ เทฺว ขนฺเธ ฯเปฯ.\csromanspacer{} ปฏิสนฺธิกฺขเณ ฯเปฯ.\csromanspacer{} วตฺถุํ ปจฺจยา ขนฺธา.\csromanspacer{} จกฺขายตนํ ปจฺจยา จกฺขุวิญฺญาณํ ฯเปฯ.\csromanspacer{} กายายตนํ ปจฺจยา กายวิญฺญาณํ.\csromanspacer{} วตฺถุํ ปจฺจยา อาสววิปฺปยุตฺตา ขนฺธา.\csromanspacer{} วตฺถุํ ปจฺจยา โทมนสฺสสหคโต วิจิกิจฺฉาสหคโต อุทฺธจฺจสหคโต โมโห. (1)}

\prose{อาสววิปฺปยุตฺตํ ธมฺมํ ปจฺจยา อาสวสมฺปยุตฺโต ธมฺโม อุปฺปชฺชติ อารมฺมณปจฺจยา.\csromanspacer{} วตฺถุํ ปจฺจยา อาสวสมฺปยุตฺตกา ขนฺธา.\csromanspacer{} โทมนสฺสสหคตํ วิจิกิจฺฉาสหคตํ อุทฺธจฺจสหคตํ โมหํ ปจฺจยา สมฺปยุตฺตกา ขนฺธา. (2)}

\prose{อาสววิปฺปยุตฺตํ ธมฺมํ ปจฺจยา อาสวสมฺปยุตฺโต จ อาสววิปฺปยุตฺโต จ ธมฺมา อุปฺปชฺชนฺติ อารมฺมณปจฺจยา.\csromanspacer{} วตฺถุํ ปจฺจยา โทมนสฺสสหคตา วิจิกิจฺฉาสหคตา อุทฺธจฺจสหคตา ขนฺธา โมโห จ. (3)}

\proseitem{69}{อาสวสมฺปยุตฺตญฺจ อาสววิปฺปยุตฺตญฺจ ธมฺมํ ปจฺจยา อาสวสมฺปยุตฺโต ธมฺโม อุปฺปชฺชติ อารมฺมณปจฺจยา.\csromanspacer{} อาสวสมฺปยุตฺตํ เอกํ ขนฺธญฺจ วตฺถุญฺจ ปจฺจยา ตโย ขนฺธา ฯเปฯ เทฺว ขนฺเธ ฯเปฯ. (1)}

\prose{อาสวสมฺปยุตฺตญฺจ อาสววิปฺปยุตฺตญฺจ ธมฺมํ ปจฺจยา อาสววิปฺปยุตฺโต ธมฺโม อุปฺปชฺชติ อารมฺมณปจฺจยา.\csromanspacer{} โทมนสฺสสหคเต วิจิกิจฺฉาสหคเต อุทฺธจฺจสหคเต ขนฺเธ จ วตฺถุญฺจ ปจฺจยา โทมนสฺสสหคโต วิจิกิจฺฉาสหคโต อุทฺธจฺจสหคโต โมโห. (2)}

\prose{อาสวสมฺปยุตฺตญฺจ อาสววิปฺปยุตฺตญฺจ ธมฺมํ ปจฺจยา อาสวสมฺปยุตฺโต จ อาสววิปฺปยุตฺโต จ ธมฺมา อุปฺปชฺชนฺติ อารมฺมณปจฺจยา.\csromanspacer{} โทมนสฺสสหคตํ วิจิกิจฺฉาสหคตํ อุทฺธจฺจสหคตํ เอกํ ขนฺธญฺจ วตฺถุญฺจ ปจฺจยา ตโย ขนฺธา โมโห จ ฯเปฯ เทฺว ขนฺเธ ฯเปฯ.\csromanspacer{} อธิปติปจฺจยา.\csromanspacer{} อนนฺตรปจฺจยา ฯเปฯ อวิคตปจฺจยา. (3)}

\csromanheader{2}{1. ปจฺจยานุโลม 2. สงฺขฺยาวาร}
\proseitem{70}{\csromanfolio{182}เหตุยา นว, อารมฺมเณ นว, (สพฺพตฺถ นว.) กมฺเม นว, วิปาเก เอกํ ฯเปฯ อวิคเต นว.}

\vspace{\tipitakaparskip}
\csromancenter{อนุโลมํ.}
\csromanheader{2}{2. ปจฺจยปจฺจนีย 1. วิภงฺควาร}
\csromanheader{2}{\textbf{นเหตุ}}
\proseitem{71}{อาสวสมฺปยุตฺตํ ธมฺมํ ปจฺจยา อาสววิปฺปยุตฺโต ธมฺโม อุปฺปชฺชติ นเหตุปจฺจยา.\csromanspacer{} วิจิกิจฺฉาสหคเต อุทฺธจฺจสหคเต ขนฺเธ ปจฺจยา วิจิกิจฺฉาสหคโต อุทฺธจฺจสหคโต โมโห. (1)}

\prose{อาสววิปฺปยุตฺตํ ธมฺมํ ปจฺจยา อาสววิปฺปยุตฺโต ธมฺโม อุปฺปชฺชติ นเหตุปจฺจยา.\csromanspacer{} อเหตุกํ อาสววิปฺปยุตฺตํ เอกํ ขนฺธํ ปจฺจยา ตโย ขนฺธา จิตฺตสมุฏฺฐานญฺจ รูปํ ฯเปฯ เทฺว ขนฺเธ ฯเปฯ.\csromanspacer{} อเหตุกปฏิสนฺธิกฺขเณ ฯเปฯ. (ยาว อสญฺญสตฺตา.) จกฺขายตนํ ปจฺจยา จกฺขุวิญฺญาณํ ฯเปฯ.\csromanspacer{} กายายตนํ ปจฺจยา กายวิญฺญาณํ.\csromanspacer{} วตฺถุํ ปจฺจยา อเหตุกา อาสววิปฺปยุตฺตา ขนฺธา.\csromanspacer{} วตฺถุํ ปจฺจยา วิจิกิจฺฉาสหคโต อุทฺธจฺจสหคโต โมโห. (1)}

\prose{อาสวสมฺปยุตฺตญฺจ อาสววิปฺปยุตฺตญฺจ ธมฺมํ ปจฺจยา อาสววิปฺปยุตฺโต ธมฺโม อุปฺปชฺชติ นเหตุปจฺจยา.\csromanspacer{} วิจิกิจฺฉาสหคโต อุทฺธจฺจสหคเต ขนฺเธ จ วตฺถุญฺจ ปจฺจยา วิจิกิจฺฉาสหคโต อุทฺธจฺจสหคโต โมโห. (1) (สํขิตฺตํ.)}

\csromanheader{2}{2. ปจฺจยปจฺจนีย 2. สงฺขฺยาวาร}
\csromanheader{2}{\textbf{สุทฺธ}}
\vspace{\tipitakaparskip}
\csromancenter{\textbf{นเหตุ}ยา ตีณิ, น-อารมฺมเณ ตีณิ, น-อธิปติยา นว, น-อนนฺตเร ตีณิ, นสมนนฺตเร ตีณิ, น-อญฺญมญฺเญ ตีณิ, น-อุปนิสฺสเย ตีณิ, นปุเรชาเต สตฺต, นปจฺฉาชาเต นว, น-อาเสวเน นว, นกมฺเม}
\csromancenter{อาสวสมฺปยุตฺตทุก จตฺตาริ, นวิปาเก นว, น-อาหาเร เอกํ, น-อินฺทฺริเย เอกํ, นฌาเน เอกํ, นมคฺเค เอกํ, นสมฺปยุตฺเต ตีณิ, นวิปฺปยุตฺเต ฉ, โนนตฺถิยา ตีณิ, โนวิคเต ตีณิ. (เอวํ อิตเรปิ เทฺว คณนา กาตพฺพา.)}
\csromancenter{(นิสฺสยวาโรปิ ปจฺจยวารสทิโส.)}
\csromanheader{2}{16. อาสวสมฺปยุตฺตทุก 5. สํสฏฺฐวาร}
\csromanheader{2}{1-4. ปจฺจยานุโลมาทิ}
\proseitem{73}{\csromanfolio{183}อาสวสมฺปยุตฺตํ ธมฺมํ สํสฏฺโฐ อาสวสมฺปยุตฺโต ธมฺโม อุปฺปชฺชติ เหตุปจฺจยา. (สํขิตฺตํ.)}

\prose{\textbf{เหตุ}ยา ฉ, อารมฺมเณ ฉ, อธิปติยา ฉ, (สพฺพตฺถ ฉ.) วิปาเก เอกํ ฯเปฯ อวิคเต ฉ.}

\prose{อาสวสมฺปยุตฺตํ ธมฺมํ สํสฏฺโฐ อาสววิปฺปยุตฺโต ธมฺโม อุปฺปชฺชติ นเหตุปจฺจยา.\csromanspacer{} วิจิกิจฺฉาสหคเต อุทฺธจฺจสหคเต ขนฺเธ สํสฏฺโฐ วิจิกิจฺฉาสหคโต อุทฺธจฺจสหคโต โมโห. (1)}

\prose{อาสววิปฺปยุตฺตํ ธมฺมํ สํสฏฺโฐ อาสววิปฺปยุตฺโต ธมฺโม อุปฺปชฺชติ นเหตุปจฺจยา ฯเปฯ.}

\prose{นเหตุยา เทฺว, น-อธิปติยา ฉ, นปุเรชาเต ฉ, นปจฺฉาชาเต ฉ, น-อาเสวเน ฉ, นกมฺเม จตฺตาริ, นวิปาเก ฉ, นฌาเน เอกํ, นมคฺเค เอกํ, นวิปฺปยุตฺเต ฉ. (เอวํ อิตเรปิ เทฺว คณนา กาตพฺพา.)}

\vspace{\tipitakaparskip}
\csromancenter{(สมฺปยุตฺตวาโรปิ สํสฏฺฐวารสทิโส.)}
\csromanheader{2}{16. อาสวสมฺปยุตฺตทุก 7. ปญฺหาวาร}
\csromanheader{2}{1. ปจฺจยานุโลม 1. วิภงฺควาร}
\csromanheader{2}{\textbf{เหตุ}}
\proseitem{74}{อาสวสมฺปยุตฺโต ธมฺโม อาสวสมฺปยุตฺตสฺส ธมฺมสฺส เหตุปจฺจเยน ปจฺจโย.\csromanspacer{} อาสวสมฺปยุตฺตา เหตู สมฺปยุตฺตกานํ ขนฺธานํ เหตุปจฺจเยน ปจฺจโย. (1)}

\prose{\csromanfolio{184}อาสวสมฺปยุตฺโต ธมฺโม อาสววิปฺปยุตฺตสฺส ธมฺมสฺส เหตุปจฺจเยน ปจฺจโย.\csromanspacer{} อาสวสมฺปยุตฺตา เหตู จิตฺตสมุฏฺฐานานํ รูปานํ เหตุปจฺจเยน ปจฺจโย.\csromanspacer{} โทโส โมหสฺส จิตฺตสมุฏฺฐานานญฺจ รูปานํ เหตุปจฺจเยน ปจฺจโย. (2)}

\prose{อาสวสมฺปยุตฺโต ธมฺโม อาสวสมฺปยุตฺตสฺส จ อาสววิปฺปยุตฺตสฺส จ ธมฺมสฺส เหตุปจฺจเยน ปจฺจโย.\csromanspacer{} อาสวสมฺปยุตฺตา เหตู สมฺปยุตฺตกานํ ขนฺธานํ จิตฺตสมุฏฺฐานานญฺจ รูปานํ เหตุปจฺจเยน ปจฺจโย.\csromanspacer{} โทโส สมฺปยุตฺตกานํ ขนฺธานํ โมหสฺส จิตฺตสมุฏฺฐานานญฺจ รูปานํ เหตุปจฺจเยน ปจฺจโย. (3)}

\proseitem{75}{อาสววิปฺปยุตฺโต ธมฺโม อาสววิปฺปยุตฺตสฺส ธมฺมสฺส เหตุปจฺจเยน ปจฺจโย.\csromanspacer{} อาสววิปฺปยุตฺตา เหตู สมฺปยุตฺตกานํ ขนฺธานํ จิตฺตสมุฏฺฐานานญฺจ รูปานํ เหตุปจฺจเยน ปจฺจโย.\csromanspacer{} โทมนสฺสสหคโต วิจิกิจฺฉาสหคโต อุทฺธจฺจสหคโต โมโห จิตฺตสมุฏฺฐานานํ รูปานํ เหตุปจฺจเยน ปจฺจโย.\csromanspacer{} ปฏิสนฺธิกฺขเณ ฯเปฯ. (1)}

\prose{อาสววิปฺปยุตฺโต ธมฺโม อาสวสมฺปยุตฺตสฺส ธมฺมสฺส เหตุปจฺจเยน ปจฺจโย.\csromanspacer{} โทมนสฺสสหคโต วิจิกิจฺฉาสหคโต อุทฺธจฺจสหคโต โมโห สมฺปยุตฺตกานํ ขนฺธานํ เหตุปจฺจเยน ปจฺจโย. (2)}

\prose{อาสววิปฺปยุตฺโต ธมฺโม อาสวสมฺปยุตฺตสฺส จ อาสววิปฺปยุตฺตสฺส จ ธมฺมสฺส เหตุปจฺจเยน ปจฺจโย.\csromanspacer{} โทมนสฺสสหคโต วิจิกิจฺฉาสหคโต อุทฺธจฺจสหคโต โมโห สมฺปยุตฺตกานํ ขนฺธานํ จิตฺตสมุฏฺฐานานญฺจ รูปานํ เหตุปจฺจเยน ปจฺจโย. (3)}

\proseitem{76}{อาสวสมฺปยุตฺโต จ อาสววิปฺปยุตฺโต จ ธมฺมา อาสวสมฺปยุตฺตสฺส ธมฺมสฺส เหตุปจฺจเยน ปจฺจโย.\csromanspacer{} โทโส จ โมโห จ อาสวสมฺปยุตฺตกานํ ขนฺธานํ เหตุปจฺจเยน ปจฺจโย. (1)}

\prose{อาสวสมฺปยุตฺโต จ อาสววิปฺปยุตฺโต จ ธมฺมา อาสววิปฺปยุตฺตสฺส ธมฺมสฺส เหตุปจฺจเยน ปจฺจโย.\csromanspacer{} โทโส จ โมโห จ จิตฺตสมุฏฺฐานานํ รูปานํ เหตุปจฺจเยน ปจฺจโย. (2)}

\prose{อาสวสมฺปยุตฺโต จ อาสววิปฺปยุตฺโต จ ธมฺมา อาสวสมฺปยุตฺตสฺส จ อาสววิปฺปยุตฺตสฺส จ ธมฺมสฺส เหตุปจฺจเยน ปจฺจโย.\csromanspacer{} โทโส จ}

\vspace{\tipitakaparskip}
\csromancenter{อาสวสมฺปยุตฺตทุก โมโห จ สมฺปยุตฺตกานํ ขนฺธานํ จิตฺตสมุฏฺฐานานญฺจ รูปานํ เหตุปจฺจเยน ปจฺจโย. (3)}
\csromanheader{2}{\textbf{อารมฺมณ}}
\proseitem{77}{\csromanfolio{185}อาสวสมฺปยุตฺโต ธมฺโม อาสวสมฺปยุตฺตสฺส ธมฺมสฺส อารมฺมณปจฺจเยน ปจฺจโย.\csromanspacer{} อาสวสมฺปยุตฺเต ขนฺเธ อารพฺภ อาสวสมฺปยุตฺตกา ขนฺธา อุปฺปชฺชนฺติ. (1)}

\prose{อาสวสมฺปยุตฺโต ธมฺโม อาสววิปฺปยุตฺตสฺส ธมฺมสฺส อารมฺมณปจฺจเยน ปจฺจโย.\csromanspacer{} อาสวสมฺปยุตฺเต ขนฺเธ อารพฺภ อาสววิปฺปยุตฺตา ขนฺธา จ โมโห จ อุปฺปชฺชนฺติ. (2)}

\prose{อาสวสมฺปยุตฺโต ธมฺโม อาสวสมฺปยุตฺตสฺส จ อาสววิปฺปยุตฺตสฺส จ ธมฺมสฺส อารมฺมณปจฺจเยน ปจฺจโย.\csromanspacer{} อาสวสมฺปยุตฺเต ขนฺเธ อารพฺภ โทมนสฺสสหคตา วิจิกิจฺฉาสหคตา อุทฺธจฺจสหคตา ขนฺธา จ โมโห จ อุปฺปชฺชนฺติ. (3)}

\proseitem{78}{อาสววิปฺปยุตฺโต ธมฺโม อาสววิปฺปยุตฺตสฺส ธมฺมสฺส อารมฺมณปจฺจเยน ปจฺจโย.\csromanspacer{} ทานํ ทตฺวา, สีลํ ฯเปฯ อุโปสถกมฺมํ กตฺวา ตํ ปจฺจเวกฺขติ.\csromanspacer{} ปุพฺเพ สุจิณฺณานิ ฯเปฯ.\csromanspacer{} ฌานา ฯเปฯ.\csromanspacer{} อริยา มคฺคํ ฯเปฯ.\csromanspacer{} ผลํ ฯเปฯ.\csromanspacer{} นิพฺพานํ ฯเปฯ.\csromanspacer{} นิพฺพานํ โคตฺรภุสฺส, โวทานสฺส, มคฺคสฺส, ผลสฺส, อาวชฺชนาย อารมฺมณปจฺจเยน ปจฺจโย.\csromanspacer{} อริยา อาสววิปฺปยุตฺเต ปหีเน กิเลเส ปจฺจเวกฺขนฺติ.\csromanspacer{} วิกฺขมฺภิเต ฯเปฯ.\csromanspacer{} ปุพฺเพ ฯเปฯ.\csromanspacer{} จกฺขุํ ฯเปฯ วตฺถุํ อาสววิปฺปยุตฺเต ขนฺเธ อนิจฺจโต ทุกฺขโต อนตฺตโต วิปสฺสติ. (อิธ อสฺสาทนา นตฺถิ.) ทิพฺเพน จกฺขุนา รูปํ ปสฺสติ.\csromanspacer{} ทิพฺพาย โสตธาตุยา สทฺทํ สุณาติ.\csromanspacer{} เจโตปริยญาเณน อาสววิปฺปยุตฺตจิตฺตสมงฺคิสฺส จิตฺตํ ชานาติ.\csromanspacer{} อากาสานญฺจายตนํ วิญฺญาณญฺจายตนสฺส ฯเปฯ.\csromanspacer{} อากิญฺจญฺญายตนํ เนวสญฺญานาสญฺญายตนสฺส ฯเปฯ.\csromanspacer{} รูปายตนํ จกฺขุวิญฺญาณสฺส ฯเปฯ.\csromanspacer{} โผฏฺฐพฺพายตนํ กายวิญฺญาณสฺส ฯเปฯ.\csromanspacer{} อาสววิปฺปยุตฺตา ขนฺธา อิทฺธิวิธญาณสฺส, เจโตปริยญาณสฺส, ปุพฺเพนิวาสานุสฺสติญาณสฺส, ยถากมฺมูปคญาณสฺส, อนาคตํสญาณสฺส, อาวชฺชนาย โมหสฺส อารมฺมณปจฺจเยน ปจฺจโย. (1) \csromanfolio{186}อาสววิปฺปยุตฺโต ธมฺโม อาสวสมฺปยุตฺตสฺส ธมฺมสฺส อารมฺมณปจฺจเยน ปจฺจโย.\csromanspacer{} ทานํ ทตฺวา, สีลํ ฯเปฯ อุโปสถกมฺมํ กตฺวา ตํ อสฺสาเทติ อภินนฺทติ, ตํ อารพฺภ ราโค อุปฺปชฺชติ.\csromanspacer{} ทิฏฺฐิ ฯเปฯ วิจิกิจฺฉา ฯเปฯ อุทฺธจฺจํ ฯเปฯ โทมนสฺสํ อุปฺปชฺชติ.\csromanspacer{} ปุพฺเพ สุจิณฺณานิ ฯเปฯ.\csromanspacer{} ฌานา วุฏฺฐหิตฺวา ฌานํ ฯเปฯ.\csromanspacer{} จกฺขุํ ฯเปฯ วตฺถุํ อาสววิปฺปยุตฺเต ขนฺเธ อสฺสาเทติ อภินนฺทติ, ตํ อารพฺภ ราโค ฯเปฯ.\csromanspacer{} ทิฏฺฐิ.\csromanspacer{} โทมนสฺสํ.\csromanspacer{} วิจิกิจฺฉา.\csromanspacer{} อุทฺธจฺจํ อุปฺปชฺชติ. (2) อาสววิปฺปยุตฺโต ธมฺโม อาสวสมฺปยุตฺตสฺส จ อาสววิปฺปยุตฺตสฺส จ ธมฺมสฺส อารมฺมณปจฺจเยน ปจฺจโย.\csromanspacer{} จกฺขุํ ฯเปฯ วตฺถุํ อาสววิปฺปยุตฺเต ขนฺเธ อารพฺภ โทมนสฺสสหคตา วิจิกิจฺฉาสหคตา อุทฺธจฺจสหคตา ขนฺธา จ โมโห จ อุปฺปชฺชนฺติ. (3)}

\proseitem{79}{อาสวสมฺปยุตฺโต จ อาสววิปฺปยุตฺโต จ ธมฺมา อาสวสมฺปยุตฺตสฺส ธมฺมสฺส อารมฺมณปจฺจเยน ปจฺจโย.\csromanspacer{} โทมนสฺสสหคเต วิจิกิจฺฉาสหคเต อุทฺธจฺจสหคเต ขนฺเธ จ โมหญฺจ อารพฺภ อาสวสมฺปยุตฺตา ขนฺธา อุปฺปชฺชนฺติ. (1)}

\prose{อาสวสมฺปยุตฺโต จ อาสววิปฺปยุตฺโต จ ธมฺมา อาสววิปฺปยุตฺตสฺส ธมฺมสฺส อารมฺมณปจฺจเยน ปจฺจโย.\csromanspacer{} โทมนสฺสสหคเต วิจิกิจฺฉาสหคเต อุทฺธจฺจสหคเต ขนฺเธ จ โมหญฺจ อารพฺภ อาสววิปฺปยุตฺตา ขนฺธา จ โมโห จ อุปฺปชฺชนฺติ. (2)}

\prose{อาสวสมฺปยุตฺโต จ อาสววิปฺปยุตฺโต จ ธมฺมา อาสวสมฺปยุตฺตสฺส จ อาสววิปฺปยุตฺตสฺส จ ธมฺมสฺส อารมฺมณปจฺจเยน ปจฺจโย.\csromanspacer{} โทมนสฺสสหคเต วิจิกิจฺฉาสหคเต อุทฺธจฺจสหคเต ขนฺเธ จ โมหญฺจ อารพฺภ โทมนสฺสสหคตา วิจิกิจฺฉาสหคตา อุทฺธจฺจสหคตา ขนฺธา จ โมโห จ อุปฺปชฺชนฺติ. (3)}

\csromanheader{2}{\textbf{อธิปติ}}
\proseitem{80}{อาสวสมฺปยุตฺโต ธมฺโม อาสวสมฺปยุตฺตสฺส ธมฺมสฺส อธิปติปจฺจเยน ปจฺจโย.\csromanspacer{} \textbf{อารมฺมณาธิปติ}, สหชาตาธิปติ.\csromanspacer{}\csromanspacer{}\textbf{อารมฺมณาธิปติ}\csromandash{}อาสวสมฺปยุตฺเต ขนฺเธ ครุํ กตฺวา อาสวสมฺปยุตฺตกา ขนฺธา อุปฺปชฺชนฺติ.\csromanspacer{} \textbf{สหชาตาธิปติ}\csromandash{} อาสวสมฺปยุตฺตาธิปติ สมฺปยุตฺตกานํ ขนฺธานํ อธิปติปจฺจเยน ปจฺจโย. (1)}

\csromanheader{2}{16. อาสวสมฺปยุตฺตทุก}
\prose{\csromanfolio{187}อาสวสมฺปยุตฺโต ธมฺโม อาสววิปฺปยุตฺตสฺส ธมฺมสฺส อธิปติปจฺจเยน ปจฺจโย.\csromanspacer{} \textbf{สหชาตาธิปติ}\csromandash{}อาสวสมฺปยุตฺตาธิปติ จิตฺตสมุฏฺฐานานํ รูปานํ อธิปติปจฺจเยน ปจฺจโย.\csromanspacer{} โทมนสฺสสหคตาธิปติ โมหสฺส จิตฺตสมุฏฺฐานานญฺจ รูปานํ อธิปติปจฺจเยน ปจฺจโย. (2)}

\prose{อาสวสมฺปยุตฺโต ธมฺโม อาสวสมฺปยุตฺตสฺส จ อาสววิปฺปยุตฺตสฺส จ ธมฺมสฺส อธิปติปจฺจเยน ปจฺจโย.\csromanspacer{} \textbf{สหชาตาธิปติ}\csromandash{} อาสวสมฺปยุตฺตาธิปติ สมฺปยุตฺตกานํ ขนฺธานํ จิตฺตสมุฏฺฐานานญฺจ รูปานํ อธิปติปจฺจเยน ปจฺจโย.\csromanspacer{} โทมนสฺสสหคตาธิปติ สมฺปยุตฺตกานํ ขนฺธานํ โมหสฺส จ จิตฺตสมุฏฺฐานานญฺจ รูปานํ อธิปติปจฺจเยน ปจฺจโย. (3)}

\proseitem{81}{อาสววิปฺปยุตฺโต ธมฺโม อาสววิปฺปยุตฺตสฺส ธมฺมสฺส อธิปติปจฺจเยน ปจฺจโย.\csromanspacer{} \textbf{อารมฺมณาธิปติ}, สหชาตาธิปติ.\csromanspacer{}\csromanspacer{}\textbf{อารมฺมณาธิปติ}\csromandash{}ทานํ ทตฺวา, สีลํ ฯเปฯ อุโปสถกมฺมํ กตฺวา ตํ ครุํ กตฺวา ปจฺจเวกฺขติ.\csromanspacer{} ปุพฺเพ สุจิณฺณานิ ฯเปฯ.\csromanspacer{} ฌานา ฯเปฯ.\csromanspacer{} อริยา มคฺคา วุฏฺฐหิตฺวา มคฺคํ ฯเปฯ.\csromanspacer{} ผลํ ฯเปฯ.\csromanspacer{} นิพฺพานํ ครุํ กตฺวา ปจฺจเวกฺขนฺติ.\csromanspacer{} นิพฺพานํ โคตฺรภุสฺส, โวทานสฺส, มคฺคสฺส, ผลสฺส อธิปติปจฺจเยน ปจฺจโย.\csromanspacer{} \textbf{สหชาตาธิปติ}\csromandash{}อาสววิปฺปยุตฺตาธิปติ สมฺปยุตฺตกานํ ขนฺธานํ จิตฺตสมุฏฺฐานานญฺจ รูปานํ อธิปติปจฺจเยน ปจฺจโย. (1)}

\prose{อาสววิปฺปยุตฺโต ธมฺโม อาสวสมฺปยุตฺตสฺส ธมฺมสฺส อธิปติปจฺจเยน ปจฺจโย.\csromanspacer{} \textbf{อารมฺมณาธิปติ}\csromandash{}ทานํ ทตฺวา, สีลํ ฯเปฯ อุโปสถกมฺมํ กตฺวา ตํ ครุํ กตฺวา อสฺสาเทติ อภินนฺทติ, ตํ ครุํ กตฺวา ราโค อุปฺปชฺชติ.\csromanspacer{} ทิฏฺฐิ อุปฺปชฺชติ.\csromanspacer{} ปุพฺเพ สุจิณฺณานิ ฯเปฯ.\csromanspacer{} ฌานา ฯเปฯ.\csromanspacer{} จกฺขุํ ฯเปฯ วตฺถุํ อาสววิปฺปยุตฺเต ขนฺเธ ครุํ กตฺวา อสฺสาเทติ อภินนฺทติ, ตํ ครุํ กตฺวา ราโค อุปฺปชฺชติ.\csromanspacer{} ทิฏฺฐิ อุปฺปชฺชติ. (2)}

\csromanheader{2}{\textbf{อนนฺตร}}
\proseitem{82}{อาสวสมฺปยุตฺโต ธมฺโม อาสวสมฺปยุตฺตสฺส ธมฺมสฺส อนนฺตรปจฺจเยน ปจฺจโย.\csromanspacer{} ปุริมา ปุริมา อาสวสมฺปยุตฺตกา ขนฺธา ปจฺฉิมานํ ปจฺฉิมานํ อาสวสมฺปยุตฺตกานํ ขนฺธานํ อนนฺตรปจฺจเยน ปจฺจโย. (1)}

\prose{อาสวสมฺปยุตฺโต ธมฺโม อาสววิปฺปยุตฺตสฺส ธมฺมสฺส อนนฺตรปจฺจเยน ปจฺจโย.\csromanspacer{} ปุริมา ปุริมา โทมนสฺสสหคตา วิจิกิจฺฉาสหคตา \csromanfolio{188}อุทฺธจฺจสหคตา ขนฺธา ปจฺฉิมสฺส ปจฺฉิมสฺส โทมนสฺสสหคตสฺส วิจิกิจฺฉาสหคตสฺส อุทฺธจฺจสหคตสฺส โมหสฺส อนนฺตรปจฺจเยน ปจฺจโย.\csromanspacer{} อาสวสมฺปยุตฺตา ขนฺธา วุฏฺฐานสฺส อนนฺตรปจฺจเยน ปจฺจโย. (2)}

\prose{อาสวสมฺปยุตฺโต ธมฺโม อาสวสมฺปยุตฺตสฺส จ อาสววิปฺปยุตฺตสฺส จ ธมฺมสฺส อนนฺตรปจฺจเยน ปจฺจโย.\csromanspacer{} ปุริมา ปุริมา โทมนสฺสสหคตา วิจิกิจฺฉาสหคตา อุทฺธจฺจสหคตา ขนฺธา ปจฺฉิมานํ ปจฺฉิมานํ โทมนสฺสสหคตานํ วิจิกิจฺฉาสหคตานํ อุทฺธจฺจสหคตานํ ขนฺธานํ โมหสฺส จ อนนฺตรปจฺจเยน ปจฺจโย. (3)}

\proseitem{83}{อาสววิปฺปยุตฺโต ธมฺโม อาสววิปฺปยุตฺตสฺส ธมฺมสฺส อนนฺตรปจฺจเยน ปจฺจโย.\csromanspacer{} ปุริโม ปุริโม โทมนสฺสสหคโต วิจิกิจฺฉาสหคโต อุทฺธจฺจสหคโต โมโห ปจฺฉิมสฺส ปจฺฉิมสฺส โทมนสฺสสหคตสฺส วิจิกิจฺฉาสหคตสฺส อุทฺธจฺจสหคตสฺส โมหสฺส อนนฺตรปจฺจเยน ปจฺจโย.\csromanspacer{} ปุริมา ปุริมา อาสววิปฺปยุตฺตา ขนฺธา ปจฺฉิมานํ ปจฺฉิมานํ อาสววิปฺปยุตฺตานํ ขนฺธานํ อนนฺตรปจฺจเยน ปจฺจโย.\csromanspacer{} อนุโลมํ โคตฺรภุสฺส ฯเปฯ ผลสมาปตฺติยา อนนฺตรปจฺจเยน ปจฺจโย. (1)}

\prose{อาสววิปฺปยุตฺโต ธมฺโม อาสวสมฺปยุตฺตสฺส ธมฺมสฺส อนนฺตรปจฺจเยน ปจฺจโย.\csromanspacer{} ปุริโม ปุริโม โทมนสฺสสหคโต วิจิกิจฺฉาสหคโต อุทฺธจฺจสหคโต โมโห ปจฺฉิมานํ ปจฺฉิมานํ โทมนสฺสสหคตานํ วิจิกิจฺฉาสหคตานํ อุทฺธจฺจสหคตานํ ขนฺธานํ อนนฺตรปจฺจเยน ปจฺจโย.\csromanspacer{} อาวชฺชนา อาสวสมฺปยุตฺตกานํ ขนฺธานํ อนนฺตรปจฺจเยน ปจฺจโย. (2)}

\prose{อาสววิปฺปยุตฺโต ธมฺโม อาสวสมฺปยุตฺตสฺส จ อาสววิปฺปยุตฺตสฺส จ ธมฺมสฺส อนนฺตรปจฺจเยน ปจฺจโย.\csromanspacer{} ปุริโม ปุริโม โทมนสฺสสหคโต วิจิกิจฺฉาสหคโต อุทฺธจฺจสหคโต โมโห ปจฺฉิมานํ ปจฺฉิมานํ โทมนสฺสสหคตานํ วิจิกิจฺฉาสหคตานํ อุทฺธจฺจสหคตานํ ขนฺธานํ โมหสฺส จ อนนฺตรปจฺจเยน ปจฺจโย.\csromanspacer{} อาวชฺชนา โทมนสฺสสหคตานํ วิจิกิจฺฉาสหคตานํ อุทฺธจฺจสหคตานํ ขนฺธานํ โมหสฺส จ อนนฺตรปจฺจเยน ปจฺจโย. (3)}

\csromanheader{2}{16. อาสวสมฺปยุตฺตทุก}
\proseitem{84}{\csromanfolio{189}อาสวสมฺปยุตฺโต จ อาสววิปฺปยุตฺโต จ ธมฺมา อาสวสมฺปยุตฺตสฺส ธมฺมสฺส อนนฺตรปจฺจเยน ปจฺจโย.\csromanspacer{} ปุริมา ปุริมา โทมนสฺสสหคตา วิจิกิจฺฉาสหคตา อุทฺธจฺจสหคตา ขนฺธา จ โมโห จ ปจฺฉิมานํ ปจฺฉิมานํ โทมนสฺสสหคตานํ วิจิกิจฺฉาสหคตานํ อุทฺธจฺจสหคตานํ ขนฺธานํ อนนฺตรปจฺจเยน ปจฺจโย. (1)}

\prose{อาสวสมฺปยุตฺโต จ อาสววิปฺปยุตฺโต จ ธมฺมา อาสววิปฺปยุตฺตสฺส ธมฺมสฺส อนนฺตรปจฺจเยน ปจฺจโย.\csromanspacer{} ปุริมา ปุริมา โทมนสฺสสหคตา วิจิกิจฺฉาสหคตา อุทฺธจฺจสหคตา ขนฺธา จ โมโห จ ปจฺฉิมสฺส ปจฺฉิมสฺส โทมนสฺสสหคตสฺส วิจิกิจฺฉาสหคตสฺส อุทฺธจฺจสหคตสฺส โมหสฺส อนนฺตรปจฺจเยน ปจฺจโย.\csromanspacer{} โทมนสฺสสหคตา วิจิกิจฺฉาสหคตา อุทฺธจฺจสหคตา ขนฺธา จ โมโห จ วุฏฺฐานสฺส อนนฺตรปจฺจเยน ปจฺจโย. (2)}

\prose{อาสวสมฺปยุตฺโต จ อาสววิปฺปยุตฺโต จ ธมฺมา อาสวสมฺปยุตฺตสฺส จ อาสววิปฺปยุตฺตสฺส จ ธมฺมสฺส อนนฺตรปจฺจเยน ปจฺจโย.\csromanspacer{} ปุริมา ปุริมา โทมนสฺสสหคตา วิจิกิจฺฉาสหคตา อุทฺธจฺจสหคตา ขนฺธา จ โมโห จ ปจฺฉิมานํ ปจฺฉิมานํ โทมนสฺสสหคตานํ วิจิกิจฺฉาสหคตานํ อุทฺธจฺจสหคตานํ ขนฺธานํ โมหสฺส จ อนนฺตรปจฺจเยน ปจฺจโย. (3)}

\csromanheader{2}{\textbf{สมนนฺตราทิ}}
\proseitem{85}{อาสวสมฺปยุตฺโต ธมฺโม อาสวสมฺปยุตฺตสฺส ธมฺมสฺส สมนนฺตรปจฺจเยน ปจฺจโย.\csromanspacer{} สหชาตปจฺจเยน ปจฺจโย.\csromanspacer{} นว.\csromanspacer{} อญฺญมญฺญปจฺจเยน ปจฺจโย.\csromanspacer{} ฉ.\csromanspacer{} นิสฺสยปจฺจเยน ปจฺจโย.\csromanspacer{} นว.}

\csromanheader{2}{\textbf{อุปนิสฺสย}}
\proseitem{86}{อาสวสมฺปยุตฺโต ธมฺโม อาสวสมฺปยุตฺตสฺส ธมฺมสฺส อุปนิสฺสยปจฺจเยน ปจฺจโย.\csromanspacer{} \textbf{อารมฺมณูปนิสฺสโย, อนนฺตรูปนิสฺสโย, ปกตูปนิสฺสโย ฯเปฯ.}\csromanspacer{} ปกตูปนิสฺสโย\csromandash{}อาสวสมฺปยุตฺตา ขนฺธา อาสวสมฺปยุตฺตานํ ขนฺธานํ อุปนิสฺสยปจฺจเยน ปจฺจโย. (1)}

\prose{อาสวสมฺปยุตฺโต ธมฺโม อาสววิปฺปยุตฺตสฺส ธมฺมสฺส อุปนิสฺสยปจฺจเยน ปจฺจโย.\csromanspacer{} \textbf{อนนฺตรูปนิสฺสโย, ปกตูปนิสฺสโย ฯเปฯ.} \csromanfolio{190}ปกตูปนิสฺสโย\csromandash{}อาสวสมฺปยุตฺตา ขนฺธา อาสววิปฺปยุตฺตานํ ขนฺธานํ โมหสฺส จ อุปนิสฺสยปจฺจเยน ปจฺจโย. (2)}

\prose{อาสวสมฺปยุตฺโต ธมฺโม อาสวสมฺปยุตฺตสฺส จ อาสววิปฺปยุตฺตสฺส จ ธมฺมสฺส อุปนิสฺสยปจฺจเยน ปจฺจโย.\csromanspacer{} \textbf{อนนฺตรูปนิสฺสโย, ปกตูปนิสฺสโย ฯเปฯ.}\csromanspacer{} ปกตูปนิสฺสโย\csromandash{}อาสวสมฺปยุตฺตกา ขนฺธา โทมนสฺสสหคตานํ วิจิกิจฺฉาสหคตานํ อุทฺธจฺจสหคตานํ ขนฺธานํ โมหสฺส จ อุปนิสฺสยปจฺจเยน ปจฺจโย. (3)}

\proseitem{87}{อาสววิปฺปยุตฺโต ธมฺโม อาสววิปฺปยุตฺตสฺส ธมฺมสฺส อุปนิสฺสยปจฺจเยน ปจฺจโย.\csromanspacer{} \textbf{อารมฺมณูปนิสฺสโย, อนนฺตรูปนิสฺสโย, ปกตูปนิสฺสโย ฯเปฯ.}\csromanspacer{} ปกตูปนิสฺสโย\csromandash{}สทฺธํ อุปนิสฺสาย ทานํ เทติ ฯเปฯ สีลํ ฯเปฯ ปญฺญํ.\csromanspacer{} กายิกํ สุขํ ฯเปฯ.\csromanspacer{} เสนาสนํ.\csromanspacer{} โมหํ อุปนิสฺสาย ทานํ เทติ ฯเปฯ สมาปตฺติํ อุปฺปาเทติ, สทฺธา ฯเปฯ ปญฺญา.\csromanspacer{} กายิกํ สุขํ.\csromanspacer{} กายิกํ ทุกฺขํ.\csromanspacer{} โมโห จ สทฺธาย ฯเปฯ โมหสฺส จ อุปนิสฺสยปจฺจเยน ปจฺจโย. (1)}

\prose{อาสววิปฺปยุตฺโต ธมฺโม อาสวสมฺปยุตฺตสฺส ธมฺมสฺส อุปนิสฺสยปจฺจเยน ปจฺจโย.\csromanspacer{} \textbf{อารมฺมณูปนิสฺสโย, อนนฺตรูปนิสฺสโย, ปกตูปนิสฺสโย ฯเปฯ.}\csromanspacer{} ปกตูปนิสฺสโย\csromandash{}สทฺธํ อุปนิสฺสาย มานํ ชปฺเปติ, ทิฏฺฐํ คณฺหาติ, สีลํ ฯเปฯ ปญฺญํ.\csromanspacer{} กายิกํ สุขํ.\csromanspacer{} กายิกํ ทุกฺขํ.\csromanspacer{} อุตุํ.\csromanspacer{} โภชนํ.\csromanspacer{} เสนาสนํ.\csromanspacer{} โมหํ อุปนิสฺสาย ปาณํ หนติ ฯเปฯ สํฆํ ภินฺทติ.\csromanspacer{} สทฺธา ฯเปฯ โมโห จ ราคสฺส ฯเปฯ ปตฺถนาย อุปนิสฺสยปจฺจเยน ปจฺจโย. (2)}

\prose{อาสววิปฺปยุตฺโต ธมฺโม อาสวสมฺปยุตฺตสฺส จ อาสววิปฺปยุตฺตสฺส จ ธมฺมสฺส อุปนิสฺสยปจฺจเยน ปจฺจโย.\csromanspacer{} \textbf{อนนฺตรูปนิสฺสโย, ปกตูปนิสฺสโย ฯเปฯ.}\csromanspacer{} ปกตูปนิสฺสโย\csromandash{}สทฺธา.\csromanspacer{} สีลํ ฯเปฯ โมโห โทมนสฺสสหคตานํ วิจิกิจฺฉาสหคตานํ อุทฺธจฺจสหคตานํ ขนฺธานํ โมหสฺส จ อุปนิสฺสยปจฺจเยน ปจฺจโย. (3)}

\proseitem{88}{อาสวสมฺปยุตฺโต จ อาสววิปฺปยุตฺโต จ ธมฺมา อาสวสมฺปยุตฺตสฺส ธมฺมสฺส อุปนิสฺสยปจฺจเยน ปจฺจโย.\csromanspacer{} \textbf{อนนฺตรูปนิสฺสโย,}}

\vspace{\tipitakaparskip}
\csromancenter{อาสวสมฺปยุตฺตทุก \textbf{ปกตูปนิสฺสโย ฯเปฯ.}\csromanspacer{} ปกตูปนิสฺสโย\csromandash{}โทมนสฺสสหคตา วิจิกิจฺฉาสหคตา อุทฺธจฺจสหคตา ขนฺธา จ โมโห จ อาสวสมฺปยุตฺตกานํ ขนฺธานํ อุปนิสฺสยปจฺจเยน ปจฺจโย. (ปุจฺฉิตพฺพํ มูลํ.) โทมนสฺสสหคตา วิจิกิจฺฉาสหคตา อุทฺธจฺจสหคตา ขนฺธา จ โมโห จ อาสวสมฺปยุตฺตกานํ ขนฺธานํ อุปนิสฺสยปจฺจเยน ปจฺจโย. (1)}
\prose{\csromanfolio{191}อาสวสมฺปยุตฺโต จ อาสววิปฺปยุตฺโต จ ธมฺมา อาสวสมฺปยุตฺตสฺส จ อาสววิปฺปยุสฺส จ ธมฺมสฺส อุปนิสฺสยปจฺจเยน ปจฺจโย.\csromanspacer{} \textbf{อนนฺตรูปนิสฺสโย, ปกตูปนิสฺสโย ฯเปฯ.}\csromanspacer{} ปกตูปนิสฺสโย\csromandash{} โทมนสฺสสหคตา วิจิกิจฺฉาสหคตา อุทฺธจฺจสหคตา ขนฺธา จ โมโห จ โทมนสฺสสหคตานํ วิจิกิจฺฉาสหคตานํ อุทฺธจฺจสหคตานํ ขนฺธานํ โมหสฺส จ อุปนิสฺสยปจฺจเยน ปจฺจโย. (2)}

\csromanheader{2}{\textbf{ปุเรชาต}}
\proseitem{89}{อาสววิปฺปยุตฺโต ธมฺโม อาสววิปฺปยุตฺตสฺส ธมฺมสฺส ปุเรชาตปจฺจเยน ปจฺจโย.\csromanspacer{} \textbf{อารมฺมณปุเรชาตํ}, วตฺถุปุเรชาตํ.\csromanspacer{}\csromanspacer{}\textbf{อารมฺมณปุเรชาตํ}\csromandash{}จกฺขุํ ฯเปฯ วตฺถุํ อนิจฺจโต ทุกฺขโต อนตฺตโต วิปสฺสติ.\csromanspacer{} ทิพฺเพน จกฺขุนา รูปํ ปสฺสติ.\csromanspacer{} ทิพฺพาย โสตธาตุยา สทฺทํ สุณาติ.\csromanspacer{} รูปายตนํ จกฺขุวิญฺญาณสฺส ฯเปฯ.\csromanspacer{} โผฏฺฐพฺพายตนํ กายวิญฺญาณสฺส ปุเรชาตปจฺจเยน ปจฺจโย.\csromanspacer{} \textbf{วตฺถุปุเรชาตํ}\csromandash{} จกฺขายตนํ จกฺขุวิญฺญาณสฺส ฯเปฯ.\csromanspacer{} กายายตนํ กายวิญฺญาณสฺส ฯเปฯ.\csromanspacer{} วตฺถุ อาสววิปฺปยุตฺตานํ ขนฺธานํ โมหสฺส จ ปุเรชาตปจฺจเยน ปจฺจโย. (1)}

\prose{อาสววิปฺปยุตฺโต ธมฺโม อาสวสมฺปยุตฺตสฺส ธมฺมสฺส ปุเรชาตปจฺจเยน ปจฺจโย.\csromanspacer{} \textbf{อารมฺมณปุเรชาตํ}, วตฺถุปุเรชาตํ.\csromanspacer{}\csromanspacer{}\textbf{อารมฺมณปุเรชาตํ}\csromandash{}จกฺขุํ ฯเปฯ วตฺถุํ อสฺสาเทติ อภินนฺทติ, ตํ อารพฺภ ราโค อุปฺปชฺชติ ฯเปฯ โทมนสฺสํ อุปฺปชฺชติ.\csromanspacer{} \textbf{วตฺถุปุเรชาตํ}\csromandash{}วตฺถุ อาสวสมฺปยุตฺตกานํ ขนฺธานํ ปุเรชาตปจฺจเยน ปจฺจโย. (2)}

\prose{อาสววิปฺปยุตฺโต ธมฺโม อาสวสมฺปยุตฺตสฺส จ อาสววิปฺปยุตฺตสฺส จ ธมฺมสฺส ปุเรชาตปจฺจเยน ปจฺจโย.\csromanspacer{} \textbf{อารมฺมณปุเรชาตํ}, วตฺถุปุเรชาตํ.\csromanspacer{}\csromanspacer{}\textbf{อารมฺมณปุเรชาตํ}\csromandash{}จกฺขุํ ฯเปฯ วตฺถุํ อารพฺภ โทมนสฺสสหคตา วิจิกิจฺฉาสหคตา อุทฺธจฺจสหคตา ขนฺธา จ โมโห จ \csromanfolio{192}อุปฺปชฺชนฺติ.\csromanspacer{} \textbf{วตฺถุปุเรชาตํ}\csromandash{}วตฺถุ โทมนสฺสสหคตานํ วิจิกิจฺฉาสหคตานํ อุทฺธจฺจสหคตานํ ขนฺธานํ โมหสฺส จ ปุเรชาตปจฺจเยน ปจฺจโย. (3)}

\csromanheader{2}{\textbf{ปจฺฉาชาต}}
\proseitem{90}{อาสวสมฺปยุตฺโต ธมฺโม อาสววิปฺปยุตฺตสฺส ธมฺมสฺส ปจฺฉาชาตปจฺจเยน ปจฺจโย.\csromanspacer{} ปจฺฉาชาตา อาสวสมฺปยุตฺตกา ขนฺธา ปุเรชาตสฺส อิมสฺส กายสฺส ปจฺฉาชาตปจฺจเยน ปจฺจโย. (1)}

\prose{อาสววิปฺปยุตฺโต ธมฺโม อาสววิปฺปยุตฺตสฺส ธมฺมสฺส ปจฺฉาชาตปจฺจเยน ปจฺจโย.\csromanspacer{} ปจฺฉาชาตา อาสววิปฺปยุตฺโต ขนฺธา จ โมโห จ ปุเรชาตสฺส อิมสฺส กายสฺส ปจฺฉาชาตปจฺจเยน ปจฺจโย. (1)}

\prose{อาสวสมฺปยุตฺโต จ อาสววิปฺปยุตฺโต จ ธมฺมา อาสววิปฺปยุตฺตสฺส ธมฺมสฺส ปจฺฉาชาตปจฺจเยน ปจฺจโย.\csromanspacer{} ปจฺฉาชาตา โทมนสฺสสหคตา วิจิกิจฺฉาสหคตา อุทฺธจฺจสหคตา ขนฺธา จ โมโห จ ปุเรชาตสฺส อิมสฺส กายสฺส ปจฺฉาชาตปจฺจเยน ปจฺจโย. (1)}

\csromanheader{2}{\textbf{อาเสวน}}
\proseitem{91}{อาสวสมฺปยุตฺโต ธมฺโม อาสวสมฺปยุตฺตสฺส ธมฺมสฺส อาเสวนปจฺจเยน ปจฺจโย.\csromanspacer{} นว. (อาวชฺชนาปิ วุฏฺฐานมฺปิ นตฺถิ.)}

\csromanheader{2}{\textbf{กมฺม}}
\proseitem{92}{อาสวสมฺปยุตฺโต ธมฺโม อาสวสมฺปยุตฺตสฺส ธมฺมสฺส กมฺมปจฺจเยน ปจฺจโย.\csromanspacer{} อาสวสมฺปยุตฺตกา เจตนา สมฺปยุตฺตกานํ ขนฺธานํ กมฺมปจฺจเยน ปจฺจโย. (1)}

\prose{อาสวสมฺปยุตฺโต ธมฺโม อาสววิปฺปยุตฺตสฺส ธมฺมสฺส กมฺมปจฺจเยน ปจฺจโย.\csromanspacer{} \textbf{สหชาตา}, นานากฺขณิกา.\csromanspacer{}\csromanspacer{}\textbf{สหชาตา} อาสวสมฺปยุตฺตา เจตนา จิตฺตสมุฏฺฐานานํ รูปานํ กมฺมปจฺจเยน ปจฺจโย.\csromanspacer{} โทมนสฺสสหคตา วิจิกิจฺฉาสหคตา อุทฺธจฺจสหคตา เจตนา โมหสฺส จิตฺตสมุฏฺฐานานญฺจ รูปานํ กมฺมปจฺจเยน ปจฺจโย.\csromanspacer{} \textbf{นานากฺขณิกา} อาสวสมฺปยุตฺตกา เจตนา วิปากานํ ขนฺธานํ กฏตฺตา จ รูปานํ กมฺมปจฺจเยน ปจฺจโย. (2)}

\csromanheader{2}{16. อาสวสมฺปยุตฺตทุก}
\prose{\csromanfolio{193}อาสวสมฺปยุตฺโต ธมฺโม อาสวสมฺปยุตฺตสฺส จ อาสววิปฺปยุตฺตสฺส จ ธมฺมสฺส กมฺมปจฺจเยน ปจฺจโย.\csromanspacer{} อาสวสมฺปยุตฺตา เจตนา สมฺปยุตฺตกานํ ขนฺธานํ จิตฺตสมุฏฺฐานานญฺจ รูปานํ กมฺมปจฺจเยน ปจฺจโย.\csromanspacer{} โทมนสฺสสหคตา วิจิกิจฺฉาสหคตา อุทฺธจฺจสหคตา เจตนา สมฺปยุตฺตกานํ ขนฺธานํ โมหสฺส จ จิตฺตสมุฏฺฐานานญฺจ รูปานํ กมฺมปจฺจเยน ปจฺจโย. (3)}

\prose{อาสววิปฺปยุตฺโต ธมฺโม อาสววิปฺปยุตฺตสฺส ธมฺมสฺส กมฺมปจฺจเยน ปจฺจโย.\csromanspacer{} \textbf{สหชาตา}, นานากฺขณิกา.\csromanspacer{}\csromanspacer{}\textbf{สหชาตา} อาสววิปฺปยุตฺตา เจตนา สมฺปยุตฺตกานํ ขนฺธานํ จิตฺตสมุฏฺฐานานญฺจ รูปานํ กมฺมปจฺจเยน ปจฺจโย.\csromanspacer{}\csromanspacer{}ปฏิสนฺธิกฺขเณ ฯเปฯ.\csromanspacer{} \textbf{นานากฺขณิกา} อาสววิปฺปยุตฺตา เจตนา วิปากานํ ขนฺธานํ กฏตฺตา จ รูปานํ กมฺมปจฺจเยน ปจฺจโย. (1)}

\csromanheader{2}{\textbf{วิปาก}}
\proseitem{93}{อาสววิปฺปยุตฺโต ธมฺโม อาสววิปฺปยุตฺตสฺส ธมฺมสฺส วิปากปจฺจเยน ปจฺจโย.\csromanspacer{} เอกํ.}

\csromanheader{2}{\textbf{อาหาร}}
\proseitem{94}{อาสวสมฺปยุตฺโต ธมฺโม อาสวสมฺปยุตฺตสฺส ธมฺมสฺส อาหารปจฺจเยน ปจฺจโย.\csromanspacer{} อาสวสมฺปยุตฺตา อาหารา สมฺปยุตฺตกานํ ขนฺธานํ อาหารปจฺจเยน ปจฺจโย. (1)}

\prose{อาสวสมฺปยุตฺโต ธมฺโม อาสววิปฺปยุตฺตสฺส ธมฺมสฺส อาหารปจฺจเยน ปจฺจโย.\csromanspacer{} อาสวสมฺปยุตฺตา อาหารา จิตฺตสมุฏฺฐานานํ รูปานํ อาหารปจฺจเยน ปจฺจโย.\csromanspacer{} โทมนสฺสสหคตา วิจิกิจฺฉาสหคตา อุทฺธจฺจสหคตา อาหารา โมหสฺส จิตฺตสมุฏฺฐานานญฺจ รูปานํ อาหารปจฺจเยน ปจฺจโย. (2)}

\prose{อาสวสมฺปยุตฺโต ธมฺโม อาสวสมฺปยุตฺตสฺส จ อาสววิปฺปยุตฺตสฺส จ ธมฺมสฺส อาหารปจฺจเยน ปจฺจโย.\csromanspacer{} อาสวสมฺปยุตฺตา อาหารา สมฺปยุตฺตกานํ ขนฺธานํ จิตฺตสมุฏฺฐานานญฺจ รูปานํ อาหารปจฺจเยน ปจฺจโย.\csromanspacer{} โทมนสฺสสหคตา วิจิกิจฺฉาสหคตา อุทฺธจฺจสหคตา อาหารา สมฺปยุตฺตกานํ ขนฺธานํ โมหสฺส จ จิตฺตสมุฏฺฐานานญฺจ รูปานํ อาหารปจฺจเยน ปจฺจโย. (3)}

\proseitem{95}{\csromanfolio{194}อาสววิปฺปยุตฺโต ธมฺโม อาสววิปฺปยุตฺตสฺส ธมฺมสฺส อาหารปจฺจเยน ปจฺจโย.\csromanspacer{} อาสววิปฺปยุตฺตา อาหารา สมฺปยุตฺตกานํ ขนฺธานํ จิตฺตสมุฏฺฐานานญฺจ รูปานํ อาหารปจฺจเยน ปจฺจโย.\csromanspacer{} ปฏิสนฺธิกฺขเณ ฯเปฯ.\csromanspacer{} กพฬีกาโร อาหาโร อิมสฺส กายสฺส อาหารปจฺจเยน ปจฺจโย. (1)}

\csromanheader{2}{\textbf{อินฺทฺริยาทิ}}
\proseitem{96}{อาสวสมฺปยุตฺโต ธมฺโม อาสวสมฺปยุตฺตสฺส ธมฺมสฺส อินฺทฺริยปจฺจเยน ปจฺจโย.\csromanspacer{} จตฺตาริ.\csromanspacer{} ฌานปจฺจเยน ปจฺจโย.\csromanspacer{} จตฺตาริ.\csromanspacer{} มคฺคปจฺจเยน ปจฺจโย.\csromanspacer{} จตฺตาริ.\csromanspacer{} สมฺปยุตฺตปจฺจเยน ปจฺจโย.\csromanspacer{} ฉ.}

\csromanheader{2}{\textbf{วิปฺปยุตฺต}}
\proseitem{97}{อาสวสมฺปยุตฺโต ธมฺโม อาสววิปฺปยุตฺตสฺส ธมฺมสฺส วิปฺปยุตฺตปจฺจเยน ปจฺจโย.\csromanspacer{} สหชาตํ, ปจฺฉาชาตํ.\csromanspacer{}\csromanspacer{}\textbf{สหชาตา} อาสวสมฺปยุตฺโต ขนฺธา จิตฺตสมุฏฺฐานานํ รูปานํ วิปฺปยุตฺตปจฺจเยน ปจฺจโย.\csromanspacer{} \textbf{ปจฺฉาชาตา} อาสวสมฺปยุตฺตา ขนฺธา ปุเรชาตสฺส อิมสฺส กายสฺส วิปฺปยุตฺตปจฺจเยน ปจฺจโย. (1)}

\prose{อาสววิปฺปยุตฺโต ธมฺโม อาสววิปฺปยุตฺตสฺส ธมฺมสฺส วิปฺปยุตฺตปจฺจเยน ปจฺจโย.\csromanspacer{} \textbf{สหชาตํ.}\csromanspacer{}\textbf{ ปุเรชาตํ.}\csromanspacer{}\textbf{ ปจฺฉาชาตํ.} (สํขิตฺตํ.\csromanspacer{} วิตฺถาเรตพฺพํ.) (1)}

\proseitem{98}{อาสววิปฺปยุตฺโต ธมฺโม อาสวสมฺปยุตฺตสฺส ธมฺมสฺส วิปฺปยุตฺตปจฺจเยน ปจฺจโย.\csromanspacer{} \textbf{ปุเรชาตํ} วตฺถุ อาสวสมฺปยุตฺตกานํ ขนฺธานํ วิปฺปยุตฺตปจฺจเยน ปจฺจโย. (1)}

\prose{อาสววิปฺปยุตฺโต ธมฺโม อาสวสมฺปยุตฺตสฺส จ อาสววิปฺปยุตฺตสฺส จ ธมฺมสฺส วิปฺปยุตฺตปจฺจเยน ปจฺจโย.\csromanspacer{} \textbf{ปุเรชาตํ} วตฺถุ โทมนสฺสสหคตานํ วิจิกิจฺฉาสหคตานํ อุทฺธจฺจสหคตานํ ขนฺธานํ โมหสฺส จ วิปฺปยุตฺตปจฺจเยน ปจฺจโย. (2)}

\prose{อาสวสมฺปยุตฺโต จ อาสววิยุตฺโต จ ธมฺมา อาสววิปฺปยุตฺตสฺส ธมฺมสฺส วิปฺปยุตฺตปจฺจเยน ปจฺจโย.\csromanspacer{} สหชาตํ, ปจฺฉาชาตํ.\csromanspacer{}\csromanspacer{}\textbf{สหชาตา} โทมนสฺสสหคตา วิจิกิจฺฉาสหคตา อุทฺธจฺจสหคตา}

\vspace{\tipitakaparskip}
\csromancenter{อาสวสมฺปยุตฺตทุก ขนฺธา จ โมโห จ จิตฺตสมุฏฺฐานํ รูปานํ วิปฺปยุตฺตปจฺจเยน ปจฺจโย.\csromanspacer{} \textbf{ปจฺฉาชาตา} โทมนสฺสสหคตา วิจิกิจฺฉาสหคตา อุทฺธจฺจสหคตา ขนฺธา จ โมโห จ ปุเรชาตสฺส อิมสฺส กายสฺส วิปฺปยุตฺตปจฺจเยน ปจฺจโย. (1)}
\csromanheader{2}{\textbf{อตฺถิ}}
\proseitem{99}{\csromanfolio{195}อาสวสมฺปยุตฺโต ธมฺโม อาสวสมฺปยุตฺตสฺส ธมฺมสฺส อตฺถิปจฺจเยน ปจฺจโย.\csromanspacer{} เอกํ. (1)}

\prose{อาสวสมฺปยุตฺโต ธมฺโม อาสววิปฺปยุตฺตสฺส ธมฺมสฺส อตฺถิปจฺจเยน ปจฺจโย.\csromanspacer{} สหชาตํ, ปจฺฉาชาตํ.\csromanspacer{}\csromanspacer{}\textbf{สหชาตา} อาสวสมฺปยุตฺตา ขนฺธา จิตฺตสมุฏฺฐานานํ รูปานํ อตฺถิปจฺจเยน ปจฺจโย.\csromanspacer{} โทมนสฺสสหคตา วิจิกิจฺฉาสหคตา อุทฺธจฺจสหคตา ขนฺธา โมหสฺส จ จิตฺตสมุฏฺฐานานญฺจ รูปานํ อตฺถิปจฺจเยน ปจฺจโย.\csromanspacer{} \textbf{ปจฺฉาชาตา} อาสวสมฺปยุตฺตา ขนฺธา ปุเรชาตสฺส อิมสฺส กายสฺส อตฺถิปจฺจเยน ปจฺจโย. (2)}

\prose{อาสวสมฺปยุตฺโต ธมฺโม อาสวสมฺปยุตฺตสฺส จ อาสววิปฺปยุตฺตสฺส จ ธมฺมสฺส อตฺถิปจฺจเยน ปจฺจโย. (สหชาตสทิสํ.) (3)}

\proseitem{100}{อาสววิปฺปยุตฺโต ธมฺโม อาสววิปฺปยุตฺตสฺส ธมฺมสฺส อตฺถิปจฺจเยน ปจฺจโย.\csromanspacer{} \textbf{สหชาตํ, ปุเรชาตํ, ปจฺฉาชาตํ, อาหารํ,} อินฺทฺริยํ. (สํขิตฺตํ.\csromanspacer{} วิตฺถาเรตพฺพํ.) (1)}

\prose{อาสววิปฺปยุตฺโต ธมฺโม อาสวสมฺปยุตฺตสฺส ธมฺมสฺส อตฺถิปจฺจเยน ปจฺจโย.\csromanspacer{} สหชาตํ, ปุเรชาตํ.\csromanspacer{}\csromanspacer{}\textbf{สหชาโต} โทนมสฺสสหคโต วิจิกิจฺฉาสหคโต อุทฺธจฺจสหคโต โมโห สมฺปยุตฺตกานํ ขนฺธานํ อตฺถิปจฺจเยน ปจฺจโย.\csromanspacer{} \textbf{ปุเรชาตํ} จกฺขุํ ฯเปฯ วตฺถุํ อสฺสาเทติ อภินนฺทติ, ตํ อารพฺภ ราโค อุปฺปชฺชติ, ทิฏฺฐิ อุปฺปชฺชติ, วิจิกิจฺฉา อุปฺปชฺชติ, อุทฺธจฺจํ อุปฺปชฺชติ, โทมนสฺสํ อุปฺปชฺชติ.\csromanspacer{} วตฺถุ อาสวสมฺปยุตฺตกานํ ขนฺธานํ อตฺถิปจฺจเยน ปจฺจโย. (2)}

\prose{อาสววิปฺปยุตฺโต ธมฺโม อาสวสมฺปยุตฺตสฺส จ อาสววิปฺปยุตฺตสฺส จ ธมฺมสฺส อตฺถิปจฺจเยน ปจฺจโย.\csromanspacer{} สหชาตํ, ปุเรชาตํ.\csromanspacer{}\csromanspacer{}\textbf{สหชาโต} \csromanfolio{196}โทมนสฺสสหคโต วิจิกิจฺฉาสหคโต อุทฺธจฺจสหคโต โมโห สมฺปยุตฺตกานํ ขนฺธานํ จิตฺตสมุฏฺฐานานญฺจ รูปานํ อตฺถิปจฺจเยน ปจฺจโย.\csromanspacer{} \textbf{ปุเรชาตํ} จกฺขุํ ฯเปฯ วตฺถุํ อารพฺภ โทมนสฺสสหคตา วิจิกิจฺฉาสหคตา อุทฺธจฺจสหคตา ขนฺธา จ โมโห จ อุปฺปชฺชนฺติ. (3)}

\proseitem{101}{อาสวสมฺปยุตฺโต จ อาสววิปฺปยุตฺโต จ ธมฺมา อาสวสมฺปยุตฺตสฺส ธมฺมสฺส อตฺถิปจฺจเยน ปจฺจโย.\csromanspacer{} สหชาตํ, ปุเรชาตํ.\csromanspacer{}\csromanspacer{}\textbf{สหชาโต} อาสวสมฺปยุตฺโต เอโก ขนฺโธ จ วตฺถุ จ ติณฺณนฺนํ ขนฺธานํ อตฺถิปจฺจเยน ปจฺจโย ฯเปฯ เทฺว ขนฺธา จ ฯเปฯ.\csromanspacer{} โทมนสฺสสหคโต วิจิกิจฺฉาสหคโต อุทฺธจฺจสหคโต เอโก ขนฺโธ จ โมโห จ ติณฺณนฺนํ ขนฺธานํ อตฺถิปจฺจเยน ปจฺจโย.\csromanspacer{} เทฺว ขนฺธา จ ฯเปฯ. (1)}

\prose{อาสวสมฺปยุตฺโต จ อาสววิปฺปยุตฺโต จ ธมฺมา อาสววิปฺปยุตฺตสฺส ธมฺมสฺส อตฺถิปจฺจเยน ปจฺจโย.\csromanspacer{} สหชาตํ, ปุเรชาตํ, ปจฺฉาชาตํ, อาหารํ, อินฺทฺริยํ.\csromanspacer{}\csromanspacer{}\textbf{สหชาตา} อาสวสมฺปยุตฺตา ขนฺธา จ มหาภูตา จ จิตฺตสมุฏฺฐานานํ รูปานํ อตฺถิปจฺจเยน ปจฺจโย.\csromanspacer{} โทมนสฺสสหคตา วิจิกิจฺฉาสหคตา อุทฺธจฺจสหคตา ขนฺธา จ โมโห จ จิตฺตสมุฏฺฐานานํ รูปานํ อตฺถิปจฺจเยน ปจฺจโย.\csromanspacer{} โทมนสฺสสหคตา วิจิกิจฺฉาสหคตา อุทฺธจฺจสหคตา ขนฺธา จ วตฺถุ จ โมหสฺส อตฺถิปจฺจเยน ปจฺจโย.\csromanspacer{} \textbf{ปจฺฉาชาตา} อาสวสมฺปยุตฺตา ขนฺธา จ กพฬีกาโร อาหาโร จ อิมสฺส กายสฺส อตฺถิปจฺจเยน ปจฺจโย.\csromanspacer{} \textbf{ปจฺฉาชาตา} อาสวสมฺปยุตฺตา ขนฺธา จ รูปชีวิตินฺทฺริยญฺจ กฏตฺตารูปานํ อตฺถิปจฺจเยน ปจฺจโย. (2)}

\prose{อาสวสมฺปยุตฺโต จ อาสววิปฺปยุตฺโต จ ธมฺมา อาสวสมฺปยุตฺตสฺส จ อาสววิปฺปยุตฺตสฺส จ ธมฺมสฺส อตฺถิปจฺจเยน ปจฺจโย.\csromanspacer{} สหชาตํ, ปุเรชาตํ.\csromanspacer{}\csromanspacer{}\textbf{สหชาโต} โทมนสฺสสหคโต วิจิกิจฺฉาสหคโต อุทฺธจฺจสหคโต เอโก ขนฺโธ จ โมโห จ ติณฺณนฺนํ ขนฺธานํ จิตฺตสมุฏฺฐานานญฺจ รูปานํ อตฺถิปจฺจเยน ปจฺจโย ฯเปฯ เทฺว ขนฺธา จ ฯเปฯ.\csromanspacer{} \textbf{สหชาโต} โทมนสฺสสหคโต วิจิกิจฺฉาสหคโต อุทฺธจฺจสหคโต เอโก ขนฺโธ จ วตฺถุ จ ติณฺณนฺนํ ขนฺธานํ โมหสฺส จ ฯเปฯ เทฺว ขนฺธา จ ฯเปฯ. (3)}

\csromanheader{2}{16. อาสวสมฺปยุตฺตทุก \textbf{1. ปจฺจยานุโลม 2. สงฺขฺยาวาร}}
\csromanheader{2}{\textbf{สุทฺธ}}
\proseitem{102}{\csromanfolio{197}เหตุยา นว, อารมฺมเณ นว, อธิปติยา ปญฺจ, อนนฺตเร นว, สมนนฺตเร นว, สหชาเต นว, อญฺญมญฺเญ ฉ, นิสฺสเย นว, อุปนิสฺสเย นว, ปุเรชาเต ตีณิ, ปจฺฉาชาเต ตีณิ, อาเสวเน นว, กมฺเม จตฺตาริ, วิปาเก เอกํ, อาหาเร จตฺตาริ, อินฺทฺริเย จตฺตาริ, ฌาเน จตฺตาริ, มคฺเค จตฺตาริ, สมฺปยุตฺเต ฉ, วิปฺปยุตฺเต ปญฺจ, อตฺถิยา นว, นตฺถิยา นว, วิคเต นว, อวิคเต นว.}

\csromanheader{2}{2. ปจฺจนียุทฺธาร}
\proseitem{103}{อาสวสมฺปยุตฺโต ธมฺโม อาสวสมฺปยุตฺตสฺส ธมฺมสฺส อารมฺมณปจฺจเยน ปจฺจโย.\csromanspacer{} สหชาตปจฺจเยน ปจฺจโย.\csromanspacer{} อุปนิสฺสยปจฺจเยน ปจฺจโย. (1)}

\prose{อาสวสมฺปยุตฺโต ธมฺโม อาสววิปฺปยุตฺตสฺส ธมฺมสฺส อารมฺมณปจฺจเยน ปจฺจโย.\csromanspacer{} สหชาตปจฺจเยน ปจฺจโย.\csromanspacer{} อุปนิสฺสยปจฺจเยน ปจฺจโย.\csromanspacer{} ปจฺฉาชาตปจฺจเยน ปจฺจโย.\csromanspacer{} กมฺมปจฺจเยน ปจฺจโย. (2)}

\prose{อาสวสมฺปยุตฺโต ธมฺโม อาสวสมฺปยุตฺตสฺส จ อาสววิปฺปยุตฺตสฺส จ ธมฺมสฺส อารมฺมณปจฺจเยน ปจฺจโย.\csromanspacer{} สหชาตปจฺจเยน ปจฺจโย.\csromanspacer{} อุปนิสฺสยปจฺจเยน ปจฺจโย. (3)}

\proseitem{104}{อาสววิปฺปยุตฺโต ธมฺโม อาสววิปฺปยุตฺตสฺส ธมฺมสฺส อารมฺมณปจฺจเยน ปจฺจโย.\csromanspacer{} สหชาตปจฺจเยน ปจฺจโย.\csromanspacer{} อุปนิสฺสยปจฺจเยน ปจฺจโย.\csromanspacer{} ปุเรชาตปจฺจเยน ปจฺจโย.\csromanspacer{} ปจฺฉาชาตปจฺจเยน ปจฺจโย.\csromanspacer{} กมฺมปจฺจเยน ปจฺจโย.\csromanspacer{} อาหารปจฺจเยน ปจฺจโย.\csromanspacer{} อินฺทฺริยปจฺจเยน ปจฺจโย. (1)}

\prose{อาสววิปฺปยุตฺโต ธมฺโม อาสวสมฺปยุตฺตสฺส ธมฺมสฺส อารมฺมณปจฺจเยน ปจฺจโย.\csromanspacer{} สหชาตปจฺจเยน ปจฺจโย.\csromanspacer{} อุปนิสฺสยปจฺจเยน ปจฺจโย.\csromanspacer{} ปุเรชาตปจฺจเยน ปจฺจโย. (2)}

\prose{\csromanfolio{198}อาสววิปฺปยุตฺโต ธมฺโม อาสวสมฺปยุตฺตสฺส จ อาสววิปฺปยุตฺตสฺส จ ธมฺมสฺส อารมฺมณปจฺจเยน ปจฺจโย.\csromanspacer{} สหชาตปจฺจเยน ปจฺจโย.\csromanspacer{} อุปนิสฺสยปจฺจเยน ปจฺจโย.\csromanspacer{} ปุเรชาตปจฺจเยน ปจฺจโย. (3)}

\proseitem{105}{อาสวสมฺปยุตฺโต จ อาสววิปฺปยุตฺโต จ ธมฺมา อาสวสมฺปยุตฺตสฺส ธมฺมสฺส อารมฺมณปจฺจเยน ปจฺจโย.\csromanspacer{} สหชาตปจฺจเยน ปจฺจโย.\csromanspacer{} อุปนิสฺสยปจฺจเยน ปจฺจโย. (1)}

\prose{อาสวสมฺปยุตฺโต จ อาสววิปฺปยุตฺโต จ ธมฺมา อาสววิปฺปยุตฺตสฺส ธมฺมสฺส อารมฺมณปจฺจเยน ปจฺจโย.\csromanspacer{} สหชาตปจฺจเยน ปจฺจโย.\csromanspacer{} อุปนิสฺสยปจฺจเยน ปจฺจโย.\csromanspacer{} ปจฺฉาชาตปจฺจเยน ปจฺจโย. (2)}

\prose{อาสวสมฺปยุตฺโต จ อาสววิปฺปยุตฺโต จ ธมฺมา อาสวสมฺปยุตฺตสฺส จ อาสววิปฺปยุตฺตสฺส จ ธมฺมสฺส อารมฺมณปจฺจเยน ปจฺจโย.\csromanspacer{} สหชาตปจฺจเยน ปจฺจโย.\csromanspacer{} อุปนิสฺสยปจฺจเยน ปจฺจโย. (3)}

\csromanheader{2}{2. ปจฺจยปจฺจนีย 2. สงฺขฺยาวาร}
\csromanheader{2}{\textbf{สุทฺธ}}
\proseitem{106}{นเหตุยา นว, น-อารมฺมเณ นว, น-อธิปติยา นว, (สพฺพตฺถ นว,) โน-อวิคเต นว.}

\csromanheader{2}{3. ปจฺจยานุโลมปจฺจนีย}
\csromanheader{2}{\textbf{เหตุทุก}}
\proseitem{107}{เหตุปจฺจยา น-อารมฺมเณ นว, น-อธิปติยา นว ฯเปฯ นสมนนฺตเร นว, น-อญฺญมญฺเญ ตีณิ, น-อุปนิสฺสเย นว ฯเปฯ นมคฺเค นว, นสมฺปยุตฺเต ตีณิ, นวิปฺปยุตฺเต ฉ, โนนตฺถิยา นว, โนวิคเต นว.}

\csromanmatikaanchor{mtk.21}
\csromantocmarkref{section}{17. อาสวสาสวทุก}{mtk.21}
\bookmark[dest={mtk.21},level=1]{17. อาสวสาสวทุก}
\csromanheader{1}{17. อาสวสาสวทุก \textbf{4. ปจฺจยปจฺจนียานุโลม}}
\csromanheader{2}{\textbf{นเหตุทุก}}
\vspace{\tipitakaparskip}
\csromancenter{นเหตุปจฺจยา อารมฺมเณ นว, อธิปติยา ปญฺจ ฯเปฯ อวิคเต นว.}
\csromancenter{ปจฺจนียานุโลมํ.}
\nitthitam{อาสวสมฺปยุตฺตทุกํ นิฏฺฐิตํ.}
\csromanheader{2}{17. อาสวสาสวทุก 1. ปฏิจฺจวาร}
\csromanheader{2}{1. ปจฺจยานุโลม 1. วิภงฺควาร}
\csromanheader{2}{\textbf{เหตุ}}
\proseitem{109}{\csromanfolio{199}อาสวญฺเจว สาสวญฺจ ธมฺมํ ปฏิจฺจ อาสโว เจว สาสโว จ ธมฺโม อุปฺปชฺชติ เหตุปจฺจยา.\csromanspacer{} กามาสวํ ปฏิจฺจ ทิฏฺฐาสโว อวิชฺชาสโว. (จกฺกํ พนฺธิตพฺพํ.) ภวาสวํ ปฏิจฺจ อวิชฺชาสโว. (จกฺกํ พนฺธิตพฺพํ.) ทิฏฺฐาสวํ ปฏิจฺจ อวิชฺชาสโว. (1)}

\prose{อาสวญฺเจว สาสวญฺจ ธมฺมํ ปฏิจฺจ สาสโว เจว โน จ อาสโว ธมฺโม อุปฺปชฺชติ เหตุปจฺจยา.\csromanspacer{} อาสเว ปฏิจฺจ สมฺปยุตฺตกา ขนฺธา จิตฺตสมุฏฺฐานญฺจ รูปํ. (2)}

\prose{อาสวญฺเจว สาสวญฺจ ธมฺมํ ปฏิจฺจ อาสโว เจว สาสโว จ สาสโว เจว โน จ อาสโว จ ธมฺมา อุปฺปชฺชนฺติ เหตุปจฺจยา.\csromanspacer{} กามาสวํ ปฏิจฺจ ทิฏฺฐาสโว อวิชฺชาสโว สมฺปยุตฺตกา ขนฺธา จิตฺตสมุฏฺฐานญฺจ รูปํ.\csromanspacer{} ภวาสวํ (จกฺกํ พนฺธิตพฺพํ.) (3)}

\proseitem{110}{สาสวญฺเจว โน จ อาสวํ ธมฺมํ ปฏิจฺจ สาสโว เจว โน จ อาสโว ธมฺโม อุปฺปชฺชติ เหตุปจฺจยา.\csromanspacer{} สาสวญฺเจว โน จ อาสวํ เอกํ ขนฺธํ ปฏิจฺจ ตโย ขนฺธา จิตฺตสมุฏฺฐานญฺจ รูปํ ฯเปฯ เทฺว ขนฺเธ ฯเปฯ.\csromanspacer{} ปฏิสนฺธิกฺขเณ ฯเปฯ.\csromanspacer{} ขนฺเธ ปฏิจฺจ วตฺถุ, วตฺถุํ ปฏิจฺจ ขนฺธา.\csromanspacer{} เอกํ มหาภูตํ ฯเปฯ. (1)}

\prose{\csromanfolio{200}สาสวญฺเจว โน จ อาสวํ ธมฺมํ ปฏิจฺจ อาสโว เจว สาสโว จ ธมฺโม อุปฺปชฺชติ เหตุปจฺจยา.\csromanspacer{} สาสเว เจว โน จ อาสเว ขนฺเธ ปฏิจฺจ อาสวา. (2)}

\prose{สาสวญฺเจว โน จ อาสวํ ธมฺมํ ปฏิจฺจ อาสโว เจว สาสโว จ สาสโว เจว โน จ อาสโว จ ธมฺมา อุปฺปชฺชนฺติ เหตุปจฺจยา.\csromanspacer{} สาสวญฺเจว โน จ อาสวํ เอกํ ขนฺธํ ปฏิจฺจ ตโย ขนฺธา อาสวา จ จิตฺตสมุฏฺฐานญฺจ รูปํ.\csromanspacer{} เทฺว ขนฺเธ ฯเปฯ. (3)}

\proseitem{111}{อาสวญฺเจว สาสวญฺจ สาสวญฺเจว โน จ อาสวญฺจ ธมฺมํ ปฏิจฺจ อาสโว เจว สาสโว จ ธมฺโม อุปฺปชฺชติ เหตุปจฺจยา.\csromanspacer{} กามาสวญฺจ สมฺปยุตฺตเก จ ขนฺเธ ปฏิจฺจ ทิฏฺฐาสโว อวิชฺชาสโว. (เอวํ จกฺกํ พนฺธิตพฺพํ.) (1)}

\prose{อาสวญฺเจว สาสวญฺจ สาสวญฺเจว โน จ อาสวญฺจ ธมฺมํ ปฏิจฺจ สาสโว เจว โน จ อาสโว ธมฺโม อุปฺปชฺชติ เหตุปจฺจยา.\csromanspacer{} สาสวญฺเจว โน จ อาสวํ เอกํ ขนฺธญฺจ อาสเว จ ปฏิจฺจ ตโย ขนฺธา จิตฺตสมุฏฺฐานญฺจ รูปํ.\csromanspacer{} เทฺว ขนฺเธ จ ฯเปฯ. (2)}

\prose{อาสวญฺเจว สาสวญฺจ สาสวญฺเจว โน จ อาสวญฺจ ธมฺมํ ปฏิจฺจ อาสโว เจว สาสโว จ สาสโว เจว โน จ อาสโว จ ธมฺมา อุปฺปชฺชนฺติ เหตุปจฺจยา.\csromanspacer{} สาสวญฺเจว โน จ อาสวํ เอกํ ขนฺธญฺจ กามาสวญฺจ ปฏิจฺจ ตโย ขนฺธา ทิฏฺฐาสโว อวิชฺชาสโว จิตฺตสมุฏฺฐานญฺจ รูปํ.\csromanspacer{} เทฺว ขนฺเธ จ ฯเปฯ. (จกฺกํ สํขิตฺตํ.) (3)}

\prose{(เอวํ ปฏิจฺจวาโรปิ สหชาตวาโรปิ ปจฺจยวาโรปิ นิสฺสายวาโรปิ สํสฏฺฐวาโรปิ สมฺปยุตฺตวาโรปิ ยถา อาสวทุกํ, เอวํ กาตพฺพํ.\csromanspacer{} นินฺนานํ.)}

\csromanheader{2}{17. อาสวสาสวทุก 7. ปญฺหาวาร}
\vspace{\tipitakaparskip}
\csromancenter{ปจฺจยานุโลม 1.\csromanspacer{} วิภงฺควาร (ปญฺหาวาเร เหตุปจฺจเยปิ อารมฺมณปจฺจเยปิ โลกุตฺตรํ นกาตพฺพํ.\csromanspacer{} เสกฺขา โคตฺรภุํ ปจฺจเวกฺขนฺติ.\csromanspacer{} โวทานํ ปจฺจเวกฺขนฺตีติ กาตพฺพา.\csromanspacer{}\csromanspacer{}อธิปติปจฺจยมฺปิ สพฺพํ ชานิตฺวา กาตพฺพํ.)}
\csromanheader{2}{17. อาสวสาสวทุก \textbf{อนนฺตร}}
\proseitem{112}{\csromanfolio{201}อาสโว เจว สาสโว จ ธมฺโม อาสวสฺส เจว สาสวสฺส จ ธมฺมสฺส อนนฺตรปจฺจเยน ปจฺจโย.\csromanspacer{} ปุริมา ปุริมา อาสวา ปจฺฉิมานํ ปจฺฉิมานํ อาสวานํ อนนฺตรปจฺจเยน ปจฺจโย. (1)}

\prose{อาสโว เจว สาสโว จ ธมฺโม สาสวสฺส เจว โน จ อาสวสฺส ธมฺมสฺส อนนฺตรปจฺจเยน ปจฺจโย.\csromanspacer{} ปุริมา ปุริมา อาสวา ปจฺฉิมานํ ปจฺฉิมานํ สาสวานญฺเจว โน จ อาสวานํ ขนฺธานํ อนนฺตรปจฺจเยน ปจฺจโย.\csromanspacer{} อาสวา วุฏฺฐานสฺส อนนฺตรปจฺจเยน ปจฺจโย.}

\prose{อาสโว เจว สาสโว จ ธมฺโม อาสวสฺส เจว สาสวสฺส จ สาสวสฺส เจว โน จ อาสวสฺส จ ธมฺมสฺส อนนฺตรปจฺจเยน ปจฺจโย.\csromanspacer{} ปุริมา ปุริมา อาสวา ปจฺฉิมานํ ปจฺฉิมานํ อาสวานํ สมฺปยุตฺตกานญฺจ ขนฺธานํ อนนฺตรปจฺจเยน ปจฺจโย. (3)}

\proseitem{113}{สาสโว เจว โน จ อาสโว ธมฺโม สาสวสฺส เจว โน จ อาสวสฺส ธมฺมสฺส อนนฺตรปจฺจเยน ปจฺจโย.\csromanspacer{} ปุริมา ปุริมา สาสวา เจว โน จ อาสวา ขนฺธา ปจฺฉิมานํ ปจฺฉิมานํ สาสวานญฺเจว โน จ อาสวานํ ขนฺธานํ อนนฺตรปจฺจเยน ปจฺจโย.\csromanspacer{} อนุโลมํ โคตฺรภุสฺส.\csromanspacer{} อนุโลมํ โวทานสฺส.\csromanspacer{} อาวชฺชนา สาสวานญฺเจว โน จ อาสวานํ ขนฺธานํ อนนฺตรปจฺจเยน ปจฺจโย. (1)}

\prose{สาสโว เจว โน จ อาสโว ธมฺโม อาสวสฺส เจว สาสวสฺส จ ธมฺมสฺส อนนฺตรปจฺจเยน ปจฺจโย.\csromanspacer{} ปุริมา ปุริมา สาสวา เจว โน จ อาสวา ขนฺธา ปจฺฉิมานํ ปจฺฉิมานํ อาสวานํ อนนฺตรปจฺจเยน ปจฺจโย.\csromanspacer{} อาวชฺชนา อาสวานํ อนนฺตรปจฺจเยน ปจฺจโย. (2)}

\prose{สาสโว เจว โน จ อาสโว ธมฺโม อาสวสฺส เจว สาสวสฺส จ สาสวสฺส เจว โน จ อาสวสฺส จ ธมฺมสฺส อนนฺตรปจฺจเยน ปจฺจโย.\csromanspacer{} ปุริมา ปุริมา สาสวา เจว โน จ อาสวา ขนฺธา ปจฺฉิมานํ ปจฺฉิมานํ อาสวานํ สมฺปยุตฺตกานญฺจ ขนฺธานํ อนนฺตรปจฺจเยน ปจฺจโย.\csromanspacer{} อาวชฺชนา อาสวานํ สมฺปยุตฺตกานญฺจ ขนฺธานํ อนนฺตรปจฺจเยน ปจฺจโย. (3)}

\proseitem{114}{อาสโว เจว สาสโว จ สาสโว เจว โน จ อาสโว จ ธมฺมา อาสวสฺส เจว สาสวสฺส จ ธมฺมสฺส อนนฺตรปจฺจเยน \csromanfolio{202}ปจฺจโย.\csromanspacer{} ปุริมา ปุริมา อาสวา จ สมฺปยุตฺตกา จ ขนฺธา ปจฺฉิมานํ ปจฺฉิมานํ อาสวานํ อนนฺตรปจฺจเยน ปจฺจโย. (1)}

\prose{อาสโว เจว สาสโว จ สาสโว เจว โน จ อาสโว จ ธมฺมา สาสวสฺส เจว โน จ อาสวสฺส ธมฺมสฺส อนนฺตรปจฺจเยน ปจฺจโย.\csromanspacer{} ปุริมา ปุริมา อาสวา จ สมฺปยุตฺตกา จ ขนฺธา ปจฺฉิมานํ ปจฺฉิมานํ สาสวานญฺเจว โน จ อาสวานํ ขนฺธานํ อนนฺตรปจฺจเยน ปจฺจโย.\csromanspacer{} อาสวา จ สมฺปยุตฺตกา จ ขนฺธา วุฏฺฐานสฺส อนนฺตรปจฺจเยน ปจฺจโย.}

\prose{อาสโว เจว สาสโว จ สาสโว เจว โน จ อาสโว จ ธมฺมา อาสวสฺส เจว สาสวสฺส จ สาสวสฺส เจว โน จ อาสวสฺส จ ธมฺมสฺส อนนฺตรปจฺจเยน ปจฺจโย.\csromanspacer{} ปุริมา ปุริมา อาสวา จ สมฺปยุตฺตกา จ ขนฺธา ปจฺฉิมานํ ปจฺฉิมานํ อาสวานํ สมฺปยุตฺตกานญฺจ ขนฺธานํ อนนฺตรปจฺจเยน ปจฺจโย. (เอวํ สพฺพํ วิตฺถาเรตพฺพํ.) (3)}

\prose{(อาสวทุเกปิ อนนฺตรํ อิมินา สทิสํ กาตพฺพํ.\csromanspacer{} อาวชฺชนาปิ วุฏฺฐานมฺปิ เอวํ สมุทฺทิฏฺฐํ สํขิตฺตํ.\csromanspacer{} สพฺพํ ปริปุณฺณํ.\csromanspacer{} อาสวทุกสทิสํ กาตพฺพํ.\csromanspacer{} นินฺนานํ.)}

\nitthitam{อาสวสาสวทุกํ นิฏฺฐิตํ.}
\csromanmatikaanchor{mtk.22}
\csromantocmarkref{section}{18. อาสว-อาสวสมฺปยุตฺตทุก}{mtk.22}
\bookmark[dest={mtk.22},level=1]{18. อาสว-อาสวสมฺปยุตฺตทุก}
\csromanheader{1}{18. อาสว-อาสวสมฺปยุตฺตทุก 1. ปฏิจฺจวาร}
\csromanheader{2}{1. ปจฺจยานุโลม 1. วิภงฺควาร}
\csromanheader{2}{\textbf{เหตุ}}
\proseitem{115}{อาสวญฺเจว อาสวสมฺปยุตฺตญฺจ ธมฺมํ ปฏิจฺจ อาสโว เจว อาสวสมฺปยุตฺโต จ ธมฺโม อุปฺปชฺชติ เหตุปจฺจยา.\csromanspacer{} กามาสวํ ปฏิจฺจ ทิฏฺฐาสโว อวิชฺชาสโว. (จกฺกํ พนฺธิตพฺพํ.) ภวาสวํ ปฏิจฺจ อวิชฺชาสโว. (จกฺกํ พนฺธิตพฺพํ.) ทิฏฺฐาสวํ ปฏิจฺจ อวิชฺชาสโว. (1)}

\prose{อาสวญฺเจว อาสวสมฺปยุตฺตญฺจ ธมฺมํ ปฏิจฺจ อาสวสมฺปยุตฺโต เจว โน จ อาสโว ธมฺโม อุปฺปชฺชติ เหตุปจฺจยา.\csromanspacer{} อาสเว ปฏิจฺจ สมฺปยุตฺตกา ขนฺธา. (2)}

\csromanheader{2}{18. อาสว-อาสวสมฺปยุตฺตทุก}
\prose{\csromanfolio{203}อาสวญฺเจว อาสวสมฺปยุตฺตญฺจ ธมฺมํ ปฏิจฺจ อาสโว เจว อาสวสมฺปยุตฺโต จ อาสวสมฺปยุตฺโต เจว โน จ อาสโว จ ธมฺมา อุปฺปชฺชนฺติ เหตุปจฺจยา.\csromanspacer{} อาสวสมฺปยุตฺตํ กามาสวํ ปฏิจฺจ ทิฏฺฐาสโว อวิชฺชาสโว สมฺปยุตฺตกา จ ขนฺธา. (สพฺพํ จกฺกํ.) (3)}

\proseitem{116}{อาสวสมฺปยุตฺตญฺเจว โน จ อาสวํ ธมฺมํ ปฏิจฺจ อาสวสมฺปยุตฺโต เจว โน จ อาสโว ธมฺโม อุปฺปชฺชติ เหตุปจฺจยา.\csromanspacer{} อาสวสมฺปยุตฺตญฺเจว โน จ อาสวํ เอกํ ขนฺธํ ปฏิจฺจ ตโย ขนฺธา ฯเปฯ เทฺว ขนฺเธ ฯเปฯ. (1)}

\prose{อาสวสมฺปยุตฺตญฺเจว โน จ อาสวํ ธมฺมํ ปฏิจฺจ อาสโว เจว อาสวสมฺปยุตฺโต จ ธมฺโม อุปฺปชฺชติ เหตุปจฺจยา.\csromanspacer{} อาสวสมฺปยุตฺตญฺเจว โน จ อาสเว ขนฺเธ ปฏิจฺจ อาสวา. (2)}

\prose{อาสวสมฺปยุตฺตญฺเจว โน จ อาสวํ ธมฺมํ ปฏิจฺจ อาสโว เจว อาสวสมฺปยุตฺโต จ อาสวสมฺปยุตฺโต เจว โน จ อาสโว จ ธมฺมา อุปฺปชฺชนฺติ เหตุปจฺจยา.\csromanspacer{} อาสวสมฺปยุตฺตญฺเจว โน จ อาสวํ เอกํ ขนฺธํ ปฏิจฺจ ตโย ขนฺธา อาสวา จ ฯเปฯ เทฺว ขนฺเธ ฯเปฯ. (3)}

\proseitem{117}{อาสวญฺเจว อาสวสมฺปยุตฺตญฺจ อาสวสมฺปยุตฺตญฺเจว โน จ อาสวญฺจ ธมฺมํ ปฏิจฺจ อาสโว เจว อาสวสมฺปยุตฺโต จ ธมฺโม อุปฺปชฺชติ เหตุปจฺจยา.\csromanspacer{} กามาสวญฺจ สมฺปยุตฺตเก จ ขนฺเธ ปฏิจฺจ ทิฏฺฐาสโว อวิชฺชาสโว. (สพฺพํ จกฺกํ.) (1)}

\prose{อาสวญฺเจว อาสวสมฺปยุตฺตญฺจ อาสวสมฺปยุตฺตญฺเจว โน จ อาสวญฺจ ธมฺมํ ปฏิจฺจ อาสวสมฺปยุตฺโต เจว โน จ อาสโว ธมฺโม อุปฺปชฺชติ เหตุปจฺจยา.\csromanspacer{} อาสวสมฺปยุตฺตญฺเจว โน จ อาสวํ เอกํ ขนฺธญฺจ อาสเว จ ปฏิจฺจ ตโย ขนฺธา ฯเปฯ เทฺว ขนฺเธ ฯเปฯ. (2)}

\prose{อาสวญฺเจว อาสวสมฺปยุตฺตญฺจ อาสวสมฺปยุตฺตญฺเจว โน จ อาสวญฺจ ธมฺมํ ปฏิจฺจ อาสโว เจว อาสวสมฺปยุตฺโต จ อาสวสมฺปยุตฺโต เจว โน จ อาสโว จ ธมฺมา อุปฺปชฺชนฺติ เหตุปจฺจยา.\csromanspacer{} อาสวสมฺปยุตฺตญฺเจว โน จ อาสวํ เอกํ ขนฺธญฺจ กามาสวญฺจ ปฏิจฺจ ตโย ขนฺธา ทิฏฺฐาสโว อวิชฺชาสโว ฯเปฯ เทฺว ขนฺเธ ฯเปฯ. (3)}

\vspace{\tipitakaparskip}
\csromancenter{(จกฺกํ.\csromanspacer{} เอวํ สพฺเพ ปจฺจยา กาตพฺพา.)}
\csromanheader{2}{1. ปจฺจยานุโลม 2. สงฺขฺยาวาร}
\csromanheader{2}{\textbf{สุทฺธ}}
\proseitem{118}{\csromanfolio{204}\textbf{เหตุ}ยา นว, อารมฺมเณ นว, อธิปติยา นว, (สพฺพตฺถ นว, สํขิตฺตํ.) กมฺเม นว, (วิปากํ นตฺถิ.) อาหาเร นว ฯเปฯ อวิคเต นว.}

\vspace{\tipitakaparskip}
\csromancenter{อนุโลมํ.}
\csromanheader{2}{2. ปจฺจยปจฺจนีย 1. วิภงฺควาร}
\csromanheader{2}{\textbf{น-อธิปตฺยาทิ}}
\proseitem{119}{อาสวญฺเจว อาสวสมฺปยุตฺตํ ธมฺมํ ปฏิจฺจ อาสโว เจว อาสวสมฺปยุตฺโต จ ธมฺโม อุปฺปชฺชติ น-อธิปติปจฺจยา. (นเหตุมูลกํ นตฺถิ.) นปุเรชาตปจฺจยา.\csromanspacer{} นปจฺฉาชาตปจฺจยา. (สํขิตฺตํ.)}

\csromanheader{2}{2. ปจฺจยปจฺจนีย 2. สงฺขฺยาวาร}
\csromanheader{2}{\textbf{สุทฺธ}}
\proseitem{120}{น-อธิปติยา นว, นปุเรชาเต นว, นปจฺฉาชาเต นว, น-อาเสวเน นว, นกมฺเม ตีณิ, นวิปาเก นว, นวิปฺปยุตฺเต นว.}

\prose{(เอวํ อิตเร เทฺว คณนาปิ สหชาตวารมฺปิ ปจฺจยวารมฺปิ นิสฺสยวารมฺปิ สํสฏฺฐวารมฺปิ สมฺปยุตฺตวารมฺปิ ปริปุณฺณํ ปฏิจฺจสทิสา.)}

\csromanheader{2}{18. อาสว-อาสวสมฺปยุตฺตทุก 7. ปญฺหาวาร}
\csromanheader{2}{1. ปจฺจยานุโลม 1. วิภงฺควาร}
\csromanheader{2}{\textbf{เหตุ}}
\proseitem{121}{อาสโว เจว อาสวสมฺปยุตฺโต จ ธมฺโม อาสวสฺส เจว อาสวสมฺปยุตฺตสฺส จ ธมฺมสฺส เหตุปจฺจเยน ปจฺจโย.\csromanspacer{} ตีณิ.}

\csromanheader{2}{18. อาสว-อาสวสมฺปยุตฺตทุก}
\prose{\csromanfolio{205}อาสวสมฺปยุตฺโต เจว โน จ อาสโว ธมฺโม อาสวสมฺปยุตฺตสฺส เจว โน จ อาสวสฺส ธมฺมสฺส เหตุปจฺจเยน ปจฺจโย.\csromanspacer{} อาสวสมฺปยุตฺตา เจว โน จ อาสวา เหตู สมฺปยุตฺตกานํ ขนฺธานํ เหตุปจฺจเยน ปจฺจโย. (1)}

\csromanheader{2}{\textbf{อารมฺมณาทิ}}
\proseitem{122}{อาสโว เจว อาสวสมฺปยุตฺโต จ ธมฺโม อาสวสฺส เจว อาสวสมฺปยุตฺตสฺส จ ธมฺมสฺส อารมฺมณปจฺจเยน ปจฺจโย.\csromanspacer{} ตีณิ.}

\prose{อาสวสมฺปยุตฺโต เจว โน จ อาสโว ธมฺโม อาสวสมฺปยุตฺตสฺส เจว โน จ อาสวสฺส ธมฺมสฺส อารมฺมณปจฺจเยน ปจฺจโย.\csromanspacer{} อาสวสมฺปยุตฺเต เจว โน จ อาสเว ขนฺเธ อารพฺภ อาสวสมฺปยุตฺตา เจว โน จ อาสวา ขนฺธา อุปฺปชฺชนฺติ. (1)}

\prose{อาสวสมฺปยุตฺโต เจว โน จ อาสโว ธมฺโม อาสวสฺส เจว อาสวสมฺปยุตฺตสฺส จ ธมฺมสฺส อารมฺมณปจฺจเยน ปจฺจโย.\csromanspacer{} อาสวสมฺปยุตฺเต เจว โน จ อาสเว ขนฺเธ อารพฺภ อาสวา อุปฺปชฺชนฺติ. (2)}

\prose{อาสวสมฺปยุตฺโต เจว โน จ อาสโว ธมฺโม อาสวสฺส เจว อาสวสมฺปยุตฺตสฺส จ อาสวสมฺปยุตฺตสฺส เจว โน จ อาสวสฺส จ ธมฺมสฺส อารมฺมณปจฺจเยน ปจฺจโย.\csromanspacer{} อาสวสมฺปยุตฺเต เจว โน จ อาสเว ขนฺเธ อารพฺภ อาสวา จ อาสวสมฺปยุตฺตกา จ ขนฺธา อุปฺปชฺชนฺติ. (3)}

\prose{อาสโว เจว อาสวสมฺปยุตฺโต จ อาสวสมฺปยุตฺโต เจว โน จ อาสโว จ ธมฺมา อาสวสฺส เจว อาสวสมฺปยุตฺตสฺส จ ธมฺมสฺส อารมฺมณปจฺจเยน ปจฺจโย.\csromanspacer{} ตีณิ.}

\prose{อธิปติปจฺจยา. (อารมฺมณสทิสา, ครุการมฺมณา.) อนนฺตรปจฺจยา. (อารมฺมณสทิสาเยว, “ปุริมา ปุริมา”ติ กาตพฺพา.) สมนนฺตรปจฺจยา.สหชาตปจฺจยา.\csromanspacer{} อญฺญมญฺญปจฺจยา.\csromanspacer{} นิสฺสยปจฺจยา.\csromanspacer{} อุปนิสฺสยปจฺจยา. (อารมฺมณสทิสํเยว.\csromanspacer{} วิภชนา นตฺถิ.\csromanspacer{} ตีณิ.\csromanspacer{} อุปนิสฺสยํ สพฺพํ กาตพฺพํ.)}

\csromanheader{2}{\textbf{กมฺมาทิ}}
\proseitem{123}{\csromanfolio{206}อาสวสมฺปยุตฺโต เจว โน จ อาสโว ธมฺโม อาสวสมฺปยุตฺตสฺส เจว โน จ อาสวสฺส ธมฺมสฺส กมฺมปจฺจเยน ปจฺจโย.\csromanspacer{} ตีณิ.\csromanspacer{} อาหารปจฺจเยน ปจฺจโย.\csromanspacer{} ตีณิ.\csromanspacer{} อินฺทฺริยปจฺจเยน ปจฺจโย.\csromanspacer{} ตีณิ.\csromanspacer{} ฌานปจฺจเยน ปจฺจโย.\csromanspacer{} ตีณิ.\csromanspacer{} มคฺคปจฺจเยน ปจฺจโย.\csromanspacer{} นว.\csromanspacer{} สมฺปยุตฺตปจฺจเยน ปจฺจโย.\csromanspacer{} นว.\csromanspacer{} อตฺถิปจฺจเยน ปจฺจโย.\csromanspacer{} นตฺถิปจฺจเยน ปจฺจโย.\csromanspacer{} วิคตปจฺจเยน ปจฺจโย.\csromanspacer{} อวิคตปจฺจเยน ปจฺจโย.\csromanspacer{} นว.}

\csromanheader{2}{1. ปจฺจยานุโลม 2. สงฺขฺยาวาร}
\csromanheader{2}{\textbf{สุทฺธ}}
\proseitem{124}{เหตุยา จตฺตาริ, อารมฺมเณ นว, อธิปติยา นว, อนนฺตเร นว, สมนนฺตเร นว, สหชาเต นว, อญฺญมญฺเญ นว, นิสฺสเย นว, อุปนิสฺสเย นว, อาเสวเน นว, กมฺเม ตีณิ, อาหาเร ตีณิ, อินฺทฺริเย ตีณิ, ฌาเน ตีณิ, มคฺเค นว, สมฺปยุตฺเต นว, อตฺถิยา นว, นตฺถิยา นว, วิคเต นว, อวิคเต นว.}

\csromanheader{2}{2. ปจฺจนียุทฺธาร}
\proseitem{125}{อาสโว เจว อาสวสมฺปยุตฺโต จ ธมฺโม อาสวสฺส เจว อาสวสมฺปยุตฺตสฺส จ ธมฺมสฺส อารมฺมณปจฺจเยน ปจฺจโย.\csromanspacer{} สหชาตปจฺจเยน ปจฺจโย.\csromanspacer{} อุปนิสฺสยปจฺจเยน ปจฺจโย. (1)}

\prose{อาสโว เจว อาสวสมฺปยุตฺโต จ ธมฺโม อาสวสมฺปยุตฺตสฺส เจว โน จ อาสวสฺส ธมฺมสฺส อารมฺมณปจฺจเยน ปจฺจโย.\csromanspacer{} สหชาตปจฺจเยน ปจฺจโย.\csromanspacer{} อุปนิสฺสยปจฺจเยน ปจฺจโย. (2)}

\prose{อาสโว เจว อาสวสมฺปยุตฺโต จ ธมฺโม อาสวสฺส เจว อาสวสมฺปยุตฺตสฺส จ อาสวสมฺปยุตฺตสฺส เจว โน จ อาสวสฺส จ ธมฺมสฺส อารมฺมณปจฺจเยน ปจฺจโย.\csromanspacer{} สหชาตปจฺจเยน ปจฺจโย.\csromanspacer{} อุปนิสฺสยปจฺจเยน ปจฺจโย. (3)}

\csromanheader{2}{18. อาสว-อาสวสมฺปยุตฺตทุก}
\proseitem{126}{\csromanfolio{207}อาสวสมฺปยุตฺโต เจว โน จ อาสโว ธมฺโม อาสวสมฺปยุตฺตสฺส เจว โน จ อาสวสฺส ธมฺมสฺส อารมฺมณปจฺจเยน ปจฺจโย.\csromanspacer{} สหชาตปจฺจเยน ปจฺจโย.\csromanspacer{} อุปนิสฺสยปจฺจเยน ปจฺจโย. (1)}

\prose{อาสวสมฺปยุตฺโต เจว โน จ อาสโว ธมฺโม อาสวสฺส เจว อาสวสมฺปยุตฺตสฺส จ ธมฺมสฺส อารมฺมณปจฺจเยน ปจฺจโย.\csromanspacer{} สหชาตปจฺจเยน ปจฺจโย.\csromanspacer{} อุปนิสฺสยปจฺจเยน ปจฺจโย. (2)}

\prose{อาสวสมฺปยุตฺโต เจว โน จ อาสโว ธมฺโม อาสวสฺส เจว อาสวสมฺปยุตฺตสฺส จ อาสวสมฺปยุตฺตสฺส เจว โน จ อาสวสฺส จ ธมฺมสฺส อารมฺมณปจฺจเยน ปจฺจโย.\csromanspacer{} สหชาตปจฺจเยน ปจฺจโย.\csromanspacer{} อุปนิสฺสยปจฺจเยน ปจฺจโย. (3)}

\proseitem{127}{อาสโว เจว อาสวสมฺปยุตฺโต จ อาสวสมฺปยุตฺโต เจว โน จ อาสโว จ ธมฺมา อาสวสฺส เจว อาสวสมฺปยุตฺตสฺส จ ธมฺมสฺส อารมฺมณปจฺจเยน ปจฺจโย.\csromanspacer{} สหชาตปจฺจเยน ปจฺจโย.\csromanspacer{} อุปนิสฺสยปจฺจเยน ปจฺจโย. (1)}

\prose{อาสโว เจว อาสวสมฺปยุตฺโต จ อาสวสมฺปยุตฺโต เจว โน จ อาสโว จ ธมฺมา อาสวสมฺปยุตฺตสฺส เจว โน จ อาสวสฺส ธมฺมสฺส อารมฺมณปจฺจเยน ปจฺจโย.\csromanspacer{} สหชาตปจฺจเยน ปจฺจโย.\csromanspacer{} อุปนิสฺสยปจฺจเยน ปจฺจโย. (2)}

\prose{อาสโว เจว อาสวสมฺปยุตฺโต จ อาสวสมฺปยุตฺโต เจว โน จ อาสโว จ ธมฺมา อาสวสฺส เจว อาสวสมฺปยุตฺตสฺส จ อาสวสมฺปยุตฺตสฺส เจว โน จ อาสวสฺส จ ธมฺมสฺส อารมฺมณปจฺจเยน ปจฺจโย.\csromanspacer{} สหชาตปจฺจเยน ปจฺจโย.\csromanspacer{} อุปนิสฺสยปจฺจเยน ปจฺจโย. (3)}

\csromanheader{2}{2. ปจฺจยปจฺจนีย 2. สงฺขฺยาวาร}
\csromanheader{2}{\textbf{สุทฺธ}}
\vspace{\tipitakaparskip}
\csromancenter{นเหตุยา นว, น-อารมฺมเณ นว, (สพฺพตฺถ นว.) โน-อวิคเต นว.}
\csromanheader{2}{3. ปจฺจยานุโลมปจฺจนีย}
\csromanheader{2}{\textbf{เหตุทุก}}
\proseitem{129}{\csromanfolio{208}\textbf{เหตุ}ปจฺจยา น-อารมฺมเณ จตฺตาริ ฯเปฯ นสมนนฺตเร จตฺตาริ, น-อุปนิสฺสเย จตฺตาริ ฯเปฯ นมคฺเค จตฺตาริ ฯเปฯ นวิปฺปยุตฺเต จตฺตาริ, โนนตฺถิยา จตฺตาริ, โนวิคเต จตฺตาริ.}

\csromanheader{2}{4. ปจฺจยปจฺจนียานุโลม}
\proseitem{130}{นเหตุปจฺจยา อารมฺมเณ นว, อธิปติยา นว, (อนุโลมปทานิ คณิตพฺพานิ.) ฯเปฯ อวิคเต นว.}

\nitthitam{อาสว-อาสวสมฺปยุตฺตทุกํ นิฏฺฐิตํ.}
\csromanmatikaanchor{mtk.23}
\csromantocmarkref{section}{19. อาสววิปฺปยุตฺตสาสวทุก}{mtk.23}
\bookmark[dest={mtk.23},level=1]{19. อาสววิปฺปยุตฺตสาสวทุก}
\csromanheader{1}{19. อาสววิปฺปยุตฺตสาสวทุก 1. ปฏิจฺจวาร}
\csromanheader{2}{\textbf{เหตุ}}
\proseitem{131}{อาสววิปฺปยุตฺตํ สาสวํ ธมฺมํ ปฏิจฺจ อาสววิปฺปยุตฺโต สาสโว ธมฺโม อุปฺปชฺชติ เหตุปจฺจยา.\csromanspacer{} อาสววิปฺปยุตฺตํ สาสวํ เอกํ ขนฺธํ ปฏิจฺจ ตโย ขนฺธา จิตฺตสมุฏฺฐานญฺจ รูปํ ฯเปฯ เทฺว ขนฺเธ ฯเปฯ.\csromanspacer{} ปฏิสนฺธิกฺขเณ ฯเปฯ.\csromanspacer{} ขนฺเธ ปฏิจฺจ วตฺถุ, วตฺถุํ ปฏิจฺจ ขนฺธา.\csromanspacer{} เอกํ มหาภูตํ ฯเปฯ. (1)}

\prose{อาสววิปฺปยุตฺตํ อนาสวํ ธมฺมํ ปฏิจฺจ อาสววิปฺปยุตฺโต อนาสโว ธมฺโม อุปฺปชฺชติ เหตุปจฺจยา.\csromanspacer{} ตีณิ.}

\prose{(ยถา จูฬนฺตรทุเก โลกิยทุกํ, เอวํ วิตฺถาเรตพฺพํ.\csromanspacer{} นินฺนานากรณํ.\csromanspacer{} สํขิตฺตํ.)}

\nitthitam{อาสววิปฺปยุตฺตสาสวทุกํ นิฏฺฐิตํ.}
\nitthitam{อาสวโคจฺฉกํ นิฏฺฐิตํ.}
\csromanmatikaanchor{mtk.24}
\csromanmatikaanchor{mtk.25}
\csromantocmarknopage{chapter}{4. สญฺโญชนโคจฺฉก}{mtk.24}
\bookmark[dest={mtk.24},level=0]{4. สญฺโญชนโคจฺฉก}
\csromantocmarkref{section}{20. สญฺโญชนทุก}{mtk.25}
\bookmark[dest={mtk.25},level=1]{20. สญฺโญชนทุก}
\csromanheader{1}{20. สญฺโญชนทุก \textbf{4. สญฺโญชนโคจฺฉก}}
\csromanheader{2}{20. สญฺโญชนทุก 1. ปฏิจฺจวาร}
\csromanheader{2}{1. ปจฺจยานุโลม 1. วิภงฺควาร}
\csromanheader{2}{\textbf{เหตุ}}
\proseitem{1}{\csromanfolio{209}สญฺโญชนํ ธมฺมํ ปฏิจฺจ สญฺโญชโน ธมฺโม อุปฺปชฺชติ เหตุปจฺจยา.\csromanspacer{} กามราคสญฺโญชนํ ปฏิจฺจ ทิฏฺฐิสญฺโญชนํ อวิชฺชาสญฺโญชนํ.\csromanspacer{} กามราคสญฺโญชนํ ปฏิจฺจ สีลพฺพตปรามาสสญฺโญชนํ อวิชฺชาสญฺโญชนํ.\csromanspacer{} กามราคสญฺโญชนํ ปฏิจฺจ มานสญฺโญชนํ อวิชฺชาสญฺโญชนํ.\csromanspacer{} กามราคสญฺโญชนํ ปฏิจฺจ อวิชฺชาสญฺโญชนํ.\csromanspacer{} ปฏิฆสญฺโญชนํ ปฏิจฺจ อิสฺสาสญฺโญชนํ อวิชฺชาสญฺโญชนํ.\csromanspacer{} ปฏิฆสญฺโญชนํ ปฏิจฺจ มจฺฉริยสญฺโญชนํ อวิชฺชาสญฺโญชนํ.\csromanspacer{} ปฏิฆสญฺโญชนํ ปฏิจฺจ อวิชฺชาสญฺโญชนํ.\csromanspacer{} มานสญฺโญชนํ ปฏิจฺจ ภวราคสญฺโญชนํ อวิชฺชาสญฺโญชนํ.\csromanspacer{} ภวราคสญฺโญชนํ ปฏิจฺจ อวิชฺชาสญฺโญชนํ.\csromanspacer{} วิจิกิจฺฉาสญฺโญชนํ ปฏิจฺจ อวิชฺชาสญฺโญชนํ. (1)}

\prose{สญฺโญชนํ ธมฺมํ ปฏิจฺจ โนสญฺโญชโน ธมฺโม อุปฺปชฺชติ เหตุปจฺจยา.\csromanspacer{} สญฺโญชเน ปฏิจฺจ สมฺปยุตฺตกา ขนฺธา จิตฺตสมุฏฺฐานญฺจ รูปํ. (2)}

\prose{สญฺโญชนํ ธมฺมํ ปฏิจฺจ สญฺโญชโน จ โนสญฺโญชโน จ ธมฺมา อุปฺปชฺชนฺติ เหตุปจฺจยา.\csromanspacer{} กามราคสญฺโญชนํ ปฏิจฺจ ทิฏฺฐิสญฺโญชนํ อวิชฺชาสญฺโญชนํ สมฺปยุตฺตกา จ ขนฺธา จ จิตฺตสมุฏฺฐานญฺจ รูปํ. (จกฺกํ พนฺธิตพฺพํ.) (3)}

\proseitem{2}{โนสญฺโญชนํ ธมฺมํ ปฏิจฺจ โนสญฺโญชโน ธมฺโม อุปฺปชฺชติ เหตุปจฺจยา.\csromanspacer{} โนสญฺโญชนํ เอกํ ขนฺธํ ปฏิจฺจ ตโย ขนฺธา จิตฺตสมุฏฺฐานญฺจ รูปํ ฯเปฯ เทฺว ขนฺเธ ฯเปฯ.\csromanspacer{} ปฏิสนฺธิกฺขเณ ฯเปฯ.\csromanspacer{} ขนฺเธ ปฏิจฺจ วตฺถุ, วตฺถุํ ปฏิจฺจ ขนฺธา.\csromanspacer{} เอกํ มหาภูตํ ฯเปฯ มหาภูเต ปฏิจฺจ จิตฺตสมุฏฺฐานํ รูปํ, กฏตฺตารูปํ, อุปาทารูปํ. (1)}

\prose{โนสญฺโญชนํ ธมฺมํ ปฏิจฺจ สญฺโญชโน ธมฺโม อุปฺปชฺชติ เหตุปจฺจยา.\csromanspacer{} โนสญฺโญชเน ขนฺเธ ปฏิจฺจ สญฺโญชนา. (2)}

\prose{\csromanfolio{210}โนสญฺโญชนํ ธมฺมํ ปฏิจฺจ สญฺโญชโน จ โนสญฺโญชโน จ ธมฺมา อุปฺปชฺชนฺติ เหตุปจฺจยา.\csromanspacer{} โนสญฺโญชนํ เอกํ ขนฺธํ ปฏิจฺจ ตโย ขนฺธา สญฺโญชนา จ จิตฺตสมุฏฺฐานญฺจ รูปํ.\csromanspacer{} เทฺว ขนฺเธ ฯเปฯ. (3)}

\proseitem{3}{สญฺโญชนญฺจ โนสญฺโญชนญฺจ ธมฺมํ ปฏิจฺจ สญฺโญชโน ธมฺโม อุปฺปชฺชติ เหตุปจฺจยา.\csromanspacer{} กามราคสญฺโญชนญฺจ สมฺปยุตฺตเก จ ขนฺเธ ปฏิจฺจ ทิฏฺฐิสญฺโญชนํ อวิชฺชาสญฺโญชนํ. (จกฺกํ พนฺธิตพฺพํ.) (1)}

\prose{สญฺโญชนญฺจ โนสญฺโญชนญฺจ ธมฺมํ ปฏิจฺจ โนสญฺโญชโน ธมฺโม อุปฺปชฺชติ เหตุปจฺจยา.\csromanspacer{} โนสญฺโญชนํ เอกํ ขนฺธญฺจ สญฺโญชเน จ ปฏิจฺจ ตโย ขนฺธา จิตฺตสมุฏฺฐานญฺจ รูปํ ฯเปฯ เทฺว ขนฺเธ จ ฯเปฯ.}

\prose{สญฺโญชนญฺจ โนสญฺโญชนญฺจ ธมฺมํ ปฏิจฺจ สญฺโญชโน จ โนสญฺโญชโน จ ธมฺมา อุปฺปชฺชนฺติ เหตุปจฺจยา.\csromanspacer{} โนสญฺโญชนํ เอกํ ขนฺธญฺจ กามราคสญฺโญชนญฺจ ปฏิจฺจ ตโย ขนฺธา ทิฏฺฐิสญฺโญชนํ อวิชฺชาสญฺโญชนํ จิตฺตสมุฏฺฐานญฺจ รูปํ ฯเปฯ. (จกฺกํ พนฺธิตพฺพํ.) (3)}

\prose{(อารมฺมณปจฺจเย รูปํ นตฺถิ.\csromanspacer{} อธิปติปจฺจโย เหตุสทิโส.\csromanspacer{} วิจิกิจฺฉาสญฺโญชนํ นตฺถิ.) อนนฺตรปจฺจยา ฯเปฯ อวิคตปจฺจยา.}

\csromanheader{2}{1. ปจฺจยานุโลม 2. สงฺขฺยาวาร}
\csromanheader{2}{\textbf{สุทฺธ}}
\proseitem{4}{เหตุยา นว, อารมฺมเณ นว, อธิปติยา, นว, อนนฺตเร นว, (สพฺพตฺถ นว.) วิปาเก เอกํ, อาหาเร นว ฯเปฯ อวิคเต นว.}

\vspace{\tipitakaparskip}
\csromancenter{อนุโลมํ.}
\csromanheader{2}{2. ปจฺจยปจฺจนีย 1. วิภงฺควาร}
\csromanheader{2}{\textbf{นเหตุ}}
\proseitem{5}{สญฺโญชนํ ธมฺมํ ปฏิจฺจ สญฺโญชโน ธมฺโม อุปฺปชฺชติ นเหตุปจฺจยา.\csromanspacer{} วิจิกิจฺฉาสญฺโญชนํ ปฏิจฺจ อวิชฺชาสญฺโญชนํ. (1)}

\csromanheader{2}{20. สญฺโญชนทุก}
\prose{\csromanfolio{211}โนสญฺโญชนํ ธมฺมํ ปฏิจฺจ โนสญฺโญชโน ธมฺโม อุปฺปชฺชติ นเหตุปจฺจยา.\csromanspacer{} อเหตุกํ โนสญฺโญชนํ เอกํ ขนฺธํ ปฏิจฺจ ตโย ขนฺธา จิตฺตสมุฏฺฐานญฺจ รูปํ ฯเปฯ เทฺว ขนฺเธ ฯเปฯ. (ยาว อสญฺญสตฺตา.) (1)}

\prose{โนสญฺโญชนํ ธมฺมํ ปฏิจฺจ สญฺโญชโน ธมฺโม อุปฺปชฺชติ นเหตุปจฺจยา.\csromanspacer{} วิจิกิจฺฉาสหคเต อุทฺธจฺจสหคเต ขนฺเธ ปฏิจฺจ อวิชฺชาสญฺโญชนํ. (2)}

\prose{สญฺโญชนญฺจ โนสญฺโญชนญฺจ ธมฺมํ ปฏิจฺจ สญฺโญชโน ธมฺโม อุปฺปชฺชติ นเหตุปจฺจยา.\csromanspacer{} วิจิกิจฺฉาสญฺโญชนญฺจ สมฺปยุตฺตเก จ ขนฺเธ ปฏิจฺจ อวิชฺชาสญฺโญชนํ. (3) (สํขิตฺตํ.\csromanspacer{} อาสวโคจฺฉกสทิสํ.\csromanspacer{} น-อารมฺมณาปิ สพฺเพ อุทฺธริตพฺพา.)}

\csromanheader{2}{2. ปจฺจยปจฺจนีย 2. สงฺขฺยาวาร}
\csromanheader{2}{\textbf{สุทฺธ}}
\proseitem{6}{นเหตุยา จตฺตาริ, น-อารมฺมเณ ตีณิ, น-อธิปติยา นว, น-อนนฺตเร ตีณิ, นสมนนฺตเร ตีณิ, น-อญฺญมญฺเญ ตีณิ, น-อุปนิสฺสเย ตีณิ, นปุเรชาเต นว, นปจฺฉาชาเต นว, น-อาเสวเน นว, นกมฺเม ตีณิ, นวิปาเก นว, น-อาหาเร เอกํ, น-อินฺทฺริเย เอกํ, นฌาเน เอกํ, นมคฺเค เอกํ, นสมฺปยุตฺเต ตีณิ, นวิปฺปยุตฺเต นว, โนนตฺถิยา ตีณิ, โนวิคเต ตีณิ.}

\vspace{\tipitakaparskip}
\csromancenter{ปจฺจนียํ.}
\csromanheader{2}{3. ปจฺจยานุโลมปจฺจนีย}
\csromanheader{2}{\textbf{เหตุทุก}}
\proseitem{7}{เหตุปจฺจยา น-อารมฺมเณ ตีณิ, น-อธิปติยา นว. (เอวํ สพฺพํ คเณตพฺพํ.)}

\csromanheader{2}{4. ปจฺจยปจฺจนียานุโลม}
\csromanheader{2}{\textbf{นเหตุทุก}}
\proseitem{8}{\csromanfolio{212}นเหตุปจฺจยา อารมฺมเณ จตฺตาริ, (สพฺพตฺถ จตฺตาริ.) วิปาเก เอกํ, อาหาเร จตฺตาริ ฯเปฯ มคฺเค ตีณิ, สมฺปยุตฺเต จตฺตาริ ฯเปฯ อวิคเต จตฺตาริ.}

\vspace{\tipitakaparskip}
\csromancenter{ปจฺจนียานุโลมํ.}
\csromanheader{2}{20. สญฺโญชนทุก 2. สหชาตวาร}
\proseitem{9}{สญฺโญชนํ ธมฺมํ สหชาโต สญฺโญชโน ธมฺโม อุปฺปชฺชติ เหตุปจฺจยา. (ปฏิจฺจวารสทิสํ.)}

\csromanheader{2}{20. สญฺโญชนทุก 3. ปจฺจยวาร}
\csromanheader{2}{1. ปจฺจยานุโลมาทิ}
\csromanheader{2}{\textbf{เหตุ}}
\proseitem{10}{สญฺโญชนํ ธมฺมํ ปจฺจยา สญฺโญชโน ธมฺโม อุปฺปชฺชติ เหตุปจฺจยา.\csromanspacer{} ตีณิ. (ปฏิจฺจสทิสํ.)}

\prose{โนสญฺโญชนํ ธมฺมํ ปจฺจยา โนสญฺโญชโน ธมฺโม อุปฺปชฺชติ เหตุปจฺจยา.\csromanspacer{} โนสญฺโญชนํ เอกํ ขนฺธํ ปจฺจยา ตโย ขนฺธา จิตฺตสมุฏฺฐานญฺจ รูปํ ฯเปฯ เทฺว ขนฺเธ ฯเปฯ.\csromanspacer{} ปฏิสนฺธิกฺขเณ ฯเปฯ.\csromanspacer{} ขนฺเธ ปจฺจยา วตฺถุ, วตฺถุํ ปจฺจยา ขนฺธา.\csromanspacer{} เอกํ มหาภูตํ ฯเปฯ มหาภูเต ปจฺจยา จิตฺตสมุฏฺฐานํ รูปํ, กฏตฺตารูปํ, อุปาทารูปํ.\csromanspacer{} วตฺถุํ ปจฺจยา โนสญฺโญชนา ขนฺธา. (1)}

\prose{โนสญฺโญชนํ ธมฺมํ ปจฺจยา สญฺโญชโน ธมฺโม อุปฺปชฺชติ เหตุปจฺจยา.\csromanspacer{} โนสญฺโญชเน ขนฺเธ ปจฺจยา สญฺโญชนา.\csromanspacer{} วตฺถุํ ปจฺจยา สญฺโญชนา. (2)}

\prose{โนสญฺโญชนํ ธมฺมํ ปจฺจยา สญฺโญชโน จ โนสญฺโญชโน จ ธมฺมา อุปฺปชฺชนฺติ เหตุปจฺจยา.\csromanspacer{} โนสญฺโญชนํ เอกํ ขนฺธํ ปจฺจยา ตโย ขนฺธา สญฺโญชนญฺจ จิตฺตสมุฏฺฐานญฺจ รูปํ ฯเปฯ เทฺว ขนฺเธ ฯเปฯ.\csromanspacer{} วตฺถุํ ปจฺจยา สญฺโญชนา.\csromanspacer{} มหาภูเต ปจฺจยา จิตฺตสมุฏฺฐานํ รูปํ.\csromanspacer{} วตฺถุํ ปจฺจยา สญฺโญชนา สมฺปยุตฺตกา จ ขนฺธา. (3)}

\csromanheader{2}{20. สญฺโญชนทุก}
\proseitem{11}{\csromanfolio{213}สญฺโญชนญฺจ โนสญฺโญชนญฺจ ธมฺมํ ปจฺจยา สญฺโญชโน ธมฺโม อุปฺปชฺชติ เหตุปจฺจยา.\csromanspacer{} กามราคสญฺโญชนญฺจ สมฺปยุตฺตเก จ ขนฺเธ ปจฺจยา ทิฏฺฐิสญฺโญชนํ อวิชฺชาสญฺโญชนํ.\csromanspacer{} กามราคสญฺโญชนญฺจ วตฺถุญฺจ ปจฺจยา ทิฏฺฐิสญฺโญชนํ อวิชฺชาสญฺโญชนํ. (จกฺกํ.) (1)}

\prose{สญฺโญชนญฺจ โนสญฺโญชนญฺจ ธมฺมํ ปจฺจยา โนสญฺโญชโน ธมฺโม อุปฺปชฺชติ เหตุปจฺจยา.\csromanspacer{} โนสญฺโญชนํ เอกํ ขนฺธญฺจ สญฺโญชเน จ ปจฺจยา ตโย ขนฺธา จิตฺตสมุฏฺฐานญฺจ รูปํ ฯเปฯ เทฺว ขนฺเธ ฯเปฯ. (จกฺกํ.) สญฺโญชเน จ วตฺถุญฺจ ปจฺจยา โนสญฺโญชนา ขนฺธา. (2)}

\prose{สญฺโญชนญฺจ โนสญฺโญชนญฺจ ธมฺมํ ปจฺจยา สญฺโญชโน จ โนสญฺโญชโน จ ธมฺมา อุปฺปชฺชนฺติ เหตุปจฺจยา.\csromanspacer{} โนสญฺโญชนํ เอกํ ขนฺธญฺจ กามราคสญฺโญชนญฺจ ปจฺจยา ตโย ขนฺธา ทิฏฺฐิสญฺโญชนํ อวิชฺชาสญฺโญชนํ จิตฺตสมุฏฺฐานญฺจ รูปํ.\csromanspacer{} เทฺว ขนฺเธ ฯเปฯ. (จกฺกํ.) กามราคสญฺโญชนญฺจ วตฺถุญฺจ ปจฺจยา ทิฏฺฐิสญฺโญชนํ อวิชฺชาสญฺโญชนํ สมฺปยุตฺตกา จ ขนฺธา. (จกฺกํ.\csromanspacer{} สํขิตฺตํ.) (3)}

\proseitem{12}{\textbf{เหตุ}ยา นว, อารมฺมเณ นว, (สพฺพตฺถ นว.) วิปาเก เอกํ ฯเปฯ อวิคเต นว.}

\proseitem{13}{นเหตุยา จตฺตาริ, (ยตฺถ ยตฺถ วตฺถุ ลพฺภติ, ตตฺถ ตตฺถ นินฺเนตพฺพํ.) น-อารมฺมเณ ตีณิ ฯเปฯ โนวิคเต ตีณิ. (เอวํ อิตเรปิ เทฺว คณนา จ นิสฺสยวาโร จ กาตพฺโพ.)}

\csromanheader{2}{20. สญฺโญชนทุก 5. สํสฏฺฐวาร}
\csromanheader{2}{1. ปจฺจยานุโลม}
\csromanheader{2}{\textbf{เหตุ}}
\proseitem{14}{สญฺโญชนํ ธมฺมํ สํสฏฺโฐ สญฺโญชโน ธมฺโม อุปฺปชฺชติ เหตุปจฺจยา.\csromanspacer{} กามราคสญฺโญชนํ สํสฏฺฐํ ทิฏฺฐิสญฺโญชนํ อวิชฺชาสญฺโญชนํ. (เอวํ นว ปญฺหา.\csromanspacer{} อรูปาเยว กาตพฺพา.)}

\vspace{\tipitakaparskip}
\csromancenter{(สํสฏฺฐวาโรปิ สมฺปยุตฺตวาโรปิ เอวํ กาตพฺพา.)}
\csromanheader{2}{20. สญฺโญชนทุก 7. ปญฺหาวาร}
\csromanheader{2}{1. ปจฺจยานุโลม 1. วิภงฺควาร}
\csromanheader{2}{\textbf{เหตุ}}
\proseitem{15}{\csromanfolio{214}สญฺโญชโน ธมฺโม สญฺโญชนสฺส ธมฺมสฺส เหตุปจฺจเยน ปจฺจโย.\csromanspacer{} สญฺโญชนา เหตู สมฺปยุตฺตกานํ สญฺโญชนานํ เหตุปจฺจเยน ปจฺจโย. (1)}

\prose{สญฺโญชโน ธมฺโม โนสญฺโญชนสฺส ธมฺมสฺส เหตุปจฺจเยน ปจฺจโย.\csromanspacer{} สญฺโญชนา เหตู สมฺปยุตฺตกานํ ขนฺธานํ จิตฺตสมุฏฺฐานานญฺจ รูปานํ เหตุปจฺจเยน ปจฺจโย. (2)}

\prose{สญฺโญชนา ธมฺโม สญฺโญชนสฺส จ โนสญฺโญชนสฺส จ ธมฺมสฺส เหตุปจฺจเยน ปจฺจโย.\csromanspacer{} สญฺโญชนา เหตู สมฺปยุตฺตกานํ ขนฺธานํ สญฺโญชนานํ จิตฺตสมุฏฺฐานานญฺจ รูปานํ เหตุปจฺจเยน ปจฺจโย. (3)}

\proseitem{16}{โนสญฺโญชโน ธมฺโม โนสญฺโญชนสฺส ธมฺมสฺส เหตุปจฺจเยน ปจฺจโย.\csromanspacer{} โนสญฺโญชนา เหตู สมฺปยุตฺตกานํ ขนฺธานํ จิตฺตสมุฏฺฐานานญฺจ รูปานํ เหตุปจฺจเยน ปจฺจโย.\csromanspacer{} ปฏิสนฺธิกฺขเณ ฯเปฯ. (1)}

\csromanheader{2}{\textbf{อารมฺมณ}}
\proseitem{17}{สญฺโญชโน ธมฺโม สญฺโญชนสฺส ธมฺมสฺส อารมฺมณปจฺจเยน ปจฺจโย.\csromanspacer{} สญฺโญชเน อารพฺภ สญฺโญชนา อุปฺปชฺชนฺติ. (มูลํ กาตพฺพํ.) สญฺโญชเน อารพฺภ โนสญฺโญชนา ขนฺธา อุปฺปชฺชนฺติ. (มูลํ กาตพฺพํ.) สญฺโญชเน อารพฺภ สญฺโญชนา จ สมฺปยุตฺตกา จ ขนฺธา อุปฺปชฺชนฺติ. (3)}

\proseitem{18}{โนสญฺโญชโน ธมฺโม โนสญฺโญชนสฺส ธมฺมสฺส อารมฺมณปจฺจเยน ปจฺจโย.\csromanspacer{} ทานํ, ทตฺวา, สีลํ ฯเปฯ อุโปสถกมฺมํ กตฺวา ตํ ปจฺจเวกฺขติ.\csromanspacer{} ปุพฺเพ สุจิณฺณานิ ฯเปฯ.\csromanspacer{} ฌานา ฯเปฯ.\csromanspacer{} อริยา มคฺคา วุฏฺฐหิตฺวา มคฺคํ ปจฺจเวกฺขนฺติ, ผลํ ฯเปฯ.\csromanspacer{} นิพฺพานํ ปจฺจเวกฺขนฺติ.\csromanspacer{} นิพฺพานํ โคตฺรภุสฺส, โวทานสฺส, มคฺคสฺส, ผลสฺส, อาวชฺชนาย อารมฺมณปจฺจเยน ปจฺจโย.\csromanspacer{} อริยา โนสญฺโญชเน ปหีเน กิเลเส ฯเปฯ.\csromanspacer{} วิกฺขมฺภิเต กิเลเส ฯเปฯ.\csromanspacer{} ปุพฺเพ ฯเปฯ.}

\vspace{\tipitakaparskip}
\csromancenter{สญฺโญชนทุก จกฺขุํ ฯเปฯ วตฺถุํ โนสญฺโญชเน ขนฺเธ อนิจฺจโต ฯเปฯ โทมนสฺสํ อุปฺปชฺชติ.\csromanspacer{} ทิพฺเพน จกฺขุนา รูปํ ปสฺสติ.\csromanspacer{} ทิพฺพาย โสตธาตุยา สทฺทํ สุณาติ.\csromanspacer{} เจโตปริยญาเณน โนสญฺโญชนจิตฺตสมงฺคิสฺส จิตฺตํ ชานาติ.\csromanspacer{} อากาสานญฺจายตนํ วิญฺญาณญฺจายตนสฺส ฯเปฯ.\csromanspacer{} อากิญฺจญฺญายตนํ เนวสญฺญานาสญฺญายตนสฺส ฯเปฯ.\csromanspacer{} รูปายตนํ จกฺขุวิญฺญาณสฺส ฯเปฯ.\csromanspacer{} โผฏฺฐพฺพายตนํ กายวิญฺญาณสฺส ฯเปฯ.\csromanspacer{} โนสญฺโญชนา ขนฺธา อิทฺธิวิธญาณสฺส, เจโตปริยญาณสฺส, ปุพฺเพนิวาสานุสฺสติญาณสฺส, ยถากมฺมูปคญาณสฺส, อนาคตํสญาณสฺส, อาวชฺชนาย อารมฺมณปจฺจเยน ปจฺจโย. (1)}
\prose{\csromanfolio{215}โนสญฺโญชโน ธมฺโม สญฺโญชนสฺส ธมฺมสฺส อารมฺมณปจฺจเยน ปจฺจโย.\csromanspacer{} ทานํ ฯเปฯ สีลํ ฯเปฯ อุโปสถกมฺมํ ฯเปฯ.\csromanspacer{} ปุพฺเพ ฯเปฯ.\csromanspacer{} ฌานา ฯเปฯ.\csromanspacer{} จกฺขุํ ฯเปฯ วตฺถุํ โนสญฺโญชเน ขนฺเธ อสฺสาเทติ อภินนฺทติ, ตํ อารพฺภ ราโค อุปฺปชฺชติ ฯเปฯ โทมนสฺสํ อุปฺปชฺชติ. (2)}

\prose{โนสญฺโญชโน ธมฺโม สญฺโญชนสฺส จ โนสญฺโญชนสฺส จ ธมฺมสฺส อารมฺมณปจฺจเยน ปจฺจโย.\csromanspacer{} ทานํ ทตฺวา, สีลํ ฯเปฯ อุโปสถกมฺมํ ฯเปฯ.\csromanspacer{} ปุพฺเพ ฯเปฯ.\csromanspacer{} ฌานา ฯเปฯ.\csromanspacer{} จกฺขุํ ฯเปฯ วตฺถุํ โนสญฺโญชเน ขนฺเธ อสฺสาเทติ อภินนฺทติ, ตํ อารพฺภ สญฺโญชนา จ สญฺโญชนสมฺปยุตฺตกา จ ขนฺธา อุปฺปชฺชนฺติ. (3)}

\prose{สญฺโญชโน จ โนสญฺโญชโน จ ธมฺมา สญฺโญชนสฺส ธมฺมสฺส อารมฺมณปจฺจเยน ปจฺจโย.\csromanspacer{} ตีณิ. (อารพฺภเยว กาตพฺพา.)}

\csromanheader{2}{\textbf{อธิปติ}}
\proseitem{19}{สญฺโญชโน ธมฺโม สญฺโญชนสฺส ธมฺมสฺส อธิปติปจฺจเยน ปจฺจโย.\csromanspacer{} \textbf{อารมฺมณาธิปติ\csromandash{}}สญฺโญชนํ ครุํ กตฺวา ฯเปฯ.\csromanspacer{} ตีณิ. (ครุการมฺมณา)}

\prose{โนสญฺโญชโน ธมฺโม โนสญฺโญชนสฺส ธมฺมสฺส อธิปติปจฺจเยน ปจฺจโย.\csromanspacer{} \textbf{อารมฺมณาธิปติ}, สหชาตาธิปติ.\csromanspacer{}\csromanspacer{}\textbf{อารมฺมณาธิปติ\csromandash{}}ทานํ ทตฺวา, สีลํ ฯเปฯ.\csromanspacer{} ตีณิ. (ติณฺณมฺปิ อารมฺมณาธิปติ.\csromanspacer{} สหชาตาธิปติปิ กาตพฺพา.\csromanspacer{} วิภชิตพฺพา.\csromanspacer{} ตีณิปิ.)}

\prose{\csromanfolio{216}สญฺโญชโน จ โนสญฺโญชโน จ ธมฺมา สญฺโญชนสฺส ธมฺมสฺส อธิปติปจฺจเยน ปจฺจโย.\csromanspacer{} \textbf{อารมฺมณาธิปติ}\csromandash{}สญฺโญชเน จ สมฺปยุตฺตเก จ ขนฺเธ ครุํ กตฺวา ฯเปฯ.\csromanspacer{} ตีณิ.}

\csromanheader{2}{\textbf{อนนฺตร}}
\proseitem{20}{สญฺโญชโน ธมฺโม สญฺโญชนสฺส ธมฺมสฺส อนนฺตรปจฺจเยน ปจฺจโย.\csromanspacer{} ปุริมา ปุริมา สญฺโญชนา ปจฺฉิมานํ ปจฺฉิมานํ สญฺโญชนานํ อนนฺตรปจฺจเยน ปจฺจโย.\csromanspacer{} ตีณิ.}

\prose{โนสญฺโญชโน ธมฺโม โนสญฺโญชนสฺส ธมฺมสฺส อนนฺตรปจฺจเยน ปจฺจโย.\csromanspacer{} ปุริมา ปุริมา โนสญฺโญชนา ขนฺธา ปจฺฉิมานํ ฯเปฯ ผลสมาปตฺติยา อนนฺตรปจฺจเยน ปจฺจโย. (1)}

\prose{โนสญฺโญชโน ธมฺโม สญฺโญชนสฺส ธมฺมสฺส อนนฺตรปจฺจเยน ปจฺจโย.\csromanspacer{} ปุริมา ปุริมา โนสญฺโญชนา ขนฺธา ปจฺฉิมานํ ปจฺฉิมานํ สญฺโญชนานํ อนนฺตรปจฺจเยน ปจฺจโย.\csromanspacer{} อาวชฺชนา สญฺโญชนานํ อนนฺตรปจฺจเยน ปจฺจโย. (เอวํ เทฺวปิ กาตพฺพา.) (3)}

\prose{สญฺโญชโน จ โนสญฺโญชโน จ ธมฺมา สญฺโญชนสฺส ธมฺมสฺส อนนฺตรปจฺจเยน ปจฺจโย.\csromanspacer{} ตีณิ.}

\csromanheader{2}{\textbf{สมนนฺตราทิ}}
\proseitem{21}{สญฺโญชโน ธมฺโม สญฺโญชนสฺส ธมฺมสฺส สมนนฺตรปจฺจเยน ปจฺจโย.\csromanspacer{} นว.\csromanspacer{} สหชาตปจฺจเยน ปจฺจโย.\csromanspacer{} นว.\csromanspacer{} อญฺญมญฺญปจฺจเยน ปจฺจโย.\csromanspacer{} นว.\csromanspacer{} นิสฺสยปจฺจเยน ปจฺจโย.\csromanspacer{} นว.}

\csromanheader{2}{\textbf{อุปนิสฺสย}}
\proseitem{22}{สญฺโญชโน ธมฺโม สญฺโญชนสฺส ธมฺมสฺส อุปนิสฺสยปจฺจเยน ปจฺจโย.\csromanspacer{} อารมฺมณูปนิสฺสโย, อนนฺตรูปนิสฺสโย, ปกตูปนิสฺสโย ฯเปฯ.\csromanspacer{} ปกตูปนิสฺสโย\csromandash{}สญฺโญชนา สญฺโญชนานํ อุปนิสฺสยปจฺจเยน ปจฺจโย. (เอวํ ตีณิปิ.)}

\csromanheader{2}{20. สญฺโญชนทุก}
\proseitem{23}{\csromanfolio{217}โนสญฺโญชโน ธมฺโม โนสญฺโญชนสฺส ธมฺมสฺส อุปนิสฺสยปจฺจเยน ปจฺจโย.\csromanspacer{} อารมฺมณูปนิสฺสโย, อนนฺตรูปนิสฺสโย, ปกตูปนิสฺสโย ฯเปฯ.\csromanspacer{} ปกตูปนิสฺสโย\csromandash{}สทฺธํ อุปนิสฺสาย ทานํ เทติ ฯเปฯ สมาปตฺติํ อุปฺปาเทติ.\csromanspacer{} มานํ ชปฺเปติ.\csromanspacer{} ทิฏฺฐิํ คณฺหาติ.\csromanspacer{} สีลํ ฯเปฯ.\csromanspacer{} ปญฺญํ.\csromanspacer{} ราคํ ฯเปฯ ปตฺถนํ ฯเปฯ.\csromanspacer{} เสนาสนํ อุปนิสฺสาย ทานํ เทติ ฯเปฯ สํฆํ ภินฺทติ.\csromanspacer{} สทฺธา ฯเปฯ.\csromanspacer{} เสนาสนํ สทฺธาย ฯเปฯ ผลสมาปตฺติยา อุปนิสฺสยปจฺจเยน ปจฺจโย. (1)}

\prose{โนสญฺโญชโน ธมฺโม สญฺโญชนสฺส ธมฺมสฺส อุปนิสฺสยปจฺจเยน ปจฺจโย.\csromanspacer{} อารมฺมณูปนิสฺสโย, อนนฺตรูปนิสฺสโย, ปกตูปนิสฺสโย ฯเปฯ.\csromanspacer{} ปกตูปนิสฺสโย\csromandash{}สทฺธํ อุปนิสฺสาย มานํ ชปฺเปติ.\csromanspacer{} ทิฏฺฐิํ คณฺหาติ.\csromanspacer{} สีลํ ฯเปฯ.\csromanspacer{} เสนาสนํ อุปนิสฺสาย ปาณํ หนติ ฯเปฯ สํฆํ ภินฺทติ.\csromanspacer{} สทฺธา ฯเปฯ.\csromanspacer{} เสนาสนํ ราคสฺส ฯเปฯ ปตฺถนาย อุปนิสฺสยปจฺจเยน ปจฺจโย. (2)}

\prose{โนสญฺโญชโน ธมฺโม สญฺโญชนสฺส จ โนสญฺโญชนสฺส จ ธมฺมสฺส อุปนิสฺสยปจฺจเยน ปจฺจโย.\csromanspacer{} อารมฺมณูปนิสฺสโย, อนนฺตรูปนิสฺสโย, ปกตูปนิสฺสโย ฯเปฯ.\csromanspacer{} ปกตูปนิสฺสโย\csromandash{}สทฺธํ อุปนิสฺสาย มานํ ชปฺเปติ.\csromanspacer{} ทิฏฺฐิํ คณฺหาติ.\csromanspacer{} สีลํ ฯเปฯ.\csromanspacer{} เสนาสนํ อุปนิสฺสาย ปาณํ หนติ ฯเปฯ สํฆํ ภินฺทติ.\csromanspacer{} สทฺธา ฯเปฯ.\csromanspacer{} เสนาสนํ สญฺโญชนานํ สมฺปยุตฺตกานญฺจ ขนฺธานํ อุปนิสฺสยปจฺจเยน ปจฺจโย. (3)}

\prose{สญฺโญชโน จ โนสญฺโญชโน จ ธมฺมา สญฺโญชนสฺส ธมฺมสฺส อุปนิสฺสยปจฺจเยน ปจฺจโย.\csromanspacer{} ตีณิ.}

\csromanheader{2}{\textbf{ปุเรชาต}}
\proseitem{24}{โนสญฺโญชโน ธมฺโม โนสญฺโญชนสฺส ธมฺมสฺส ปุเรชาตปจฺจเยน ปจฺจโย.\csromanspacer{} \textbf{อารมฺมณปุเรชาตํ}, วตฺถุปุเรชาตํ.\csromanspacer{}\csromanspacer{}\textbf{อารมฺมณปุเรชาตํ}\csromandash{}จกฺขุํ ฯเปฯ วตฺถุํ อนิจฺจโต ฯเปฯ โทมนสฺสํ อุปฺปชฺชติ.\csromanspacer{} ทิพฺเพน จกฺขุนา รูปํ ปสฺสติ.\csromanspacer{} ทิพฺพาย โสตธาตุยา สทฺทํ สุณาติ.\csromanspacer{} รูปายตนํ จกฺขุวิญฺญาณสฺส ฯเปฯ.\csromanspacer{} โผฏฺฐพฺพายตนํ กายวิญฺญาณสฺส ฯเปฯ.\csromanspacer{} \textbf{วตฺถุปุเรชาตํ}\csromandash{}จกฺขายตนํ จกฺขุวิญฺญาณสฺส ฯเปฯ กายายตนํ กายวิญฺญาณสฺส ฯเปฯ.\csromanspacer{} วตฺถุ โนสญฺโญชนานํ ขนฺธานํ ปุเรชาตปจฺจเยน ปจฺจโย. (1)}

\prose{\csromanfolio{218}โนสญฺโญชโน ธมฺโม สญฺโญชนสฺส ธมฺมสฺส ปุเรชาตปจฺจเยน ปจฺจโย.\csromanspacer{} \textbf{อารมฺมณปุเรชาตํ}, วตฺถุปุเรชาตํ.\csromanspacer{}\csromanspacer{}\textbf{อารมฺมณปุเรชาตํ}\csromandash{} จกฺขุํ ฯเปฯ วตฺถุํ อสฺสาเทติ อภินนฺทติ, ตํ อารพฺภ ราโค อุปฺปชฺชติ ฯเปฯ โทมนสฺสํ อุปฺปชฺชติ.\csromanspacer{} \textbf{วตฺถุปุเรชาตํ}\csromandash{}วตฺถุ สญฺโญชนานํ ปุเรชาตปจฺจเยน ปจฺจโย. (2)}

\prose{โนสญฺโญชโน ธมฺโม สญฺโญชนสฺส จ โนสญฺโญชนสฺส จ ธมฺมสฺส ปุเรชาตปจฺจเยน ปจฺจโย.\csromanspacer{} \textbf{อารมฺมณปุเรชาตํ}, วตฺถุปุเรชาตํ.\csromanspacer{}\csromanspacer{}\textbf{อารมฺมณปุเรชาตํ}\csromandash{}จกฺขุํ ฯเปฯ วตฺถุํ อสฺสาเทติ อภินนฺทติ, ตํ อารพฺภ สญฺโญชนา จ สมฺปยุตฺตกา จ ขนฺธา อุปฺปชฺชนฺติ.\csromanspacer{} \textbf{วตฺถุปุเรชาตํ}\csromandash{}วตฺถุ สญฺโญชนานํ สมฺปยุตฺตกานญฺจ ขนฺธานํ ปุเรชาตปจฺจเยน ปจฺจโย. (3)}

\csromanheader{2}{\textbf{ปจฺฉาชาต}}
\proseitem{25}{สญฺโญชโน ธมฺโม โนสญฺโญชนสฺส ธมฺมสฺส ปจฺฉาชาตปจฺจเยน ปจฺจโย.\csromanspacer{} เอกํ.}

\prose{โนสญฺโญชโน ธมฺโม โนสญฺโญชนสฺส ธมฺมสฺส ปจฺฉาชาตปจฺจเยน ปจฺจโย ฯเปฯ. (1)}

\prose{สญฺโญชโน จ โนสญฺโญชโน จ ธมฺมา โนสญฺโญชนสฺส ธมฺมสฺส ปจฺฉาชาตปจฺจเยน ปจฺจโย ฯเปฯ. (1)}

\csromanheader{2}{\textbf{อาเสวน}}
\vspace{\tipitakaparskip}
\csromancenter{สญฺโญชโน ธมฺโม สญฺโญชนสฺส ธมฺมสฺส อาเสวนปจฺจเยน ปจฺจโย.\csromanspacer{} นว.}
\csromanheader{2}{\textbf{กมฺมาทิ}}
\proseitem{27}{โนสญฺโญชโน ธมฺโม โนสญฺโญชนสฺส ธมฺมสฺส กมฺมปจฺจเยน ปจฺจโย.\csromanspacer{} ตีณิ.\csromanspacer{} วิปากปจฺจเยน ปจฺจโย.\csromanspacer{} เอกํ.\csromanspacer{} อาหารปจฺจเยน ปจฺจโย.\csromanspacer{} ตีณิ.\csromanspacer{} อินฺทฺริยปจฺจเยน ปจฺจโย.\csromanspacer{} ตีณิ.\csromanspacer{} ฌานปจฺจเยน ปจฺจโย.\csromanspacer{} ตีณิ.\csromanspacer{} มคฺคปจฺจเยน ปจฺจโย.\csromanspacer{} นว.\csromanspacer{} สมฺปยุตฺตปจฺจเยน ปจฺจโย.\csromanspacer{} นว.}

\csromanheader{2}{20. สญฺโญชนทุก \textbf{วิปฺปยุตฺต}}
\proseitem{28}{\csromanfolio{219}สญฺโญชโน ธมฺโม โนสญฺโญชนสฺส ธมฺมสฺส วิปฺปยุตฺตปจฺจเยน ปจฺจโย.\csromanspacer{} สหชาตํ, ปจฺฉาชาตํ. (วิภชิตพฺพํ.) (1)}

\prose{โนสญฺโญชโน ธมฺโม โนสญฺโญชนสฺส ธมฺมสฺส วิปฺปยุตฺตปจฺจเยน ปจฺจโย.\csromanspacer{} สหชาตํ, ปุเรชาตํ, ปจฺฉาชาตํ. (วิภชิตพฺพํ.) (1)}

\prose{โนสญฺโญชโน ธมฺโม สญฺโญชนสฺส ธมฺมสฺส วิปฺปยุตฺตปจฺจเยน ปจฺจโย.\csromanspacer{} \textbf{ปุเรชาตํ} วตฺถุ สญฺโญชนานํ วิปฺปยุตฺตปจฺจเยน ปจฺจโย. (2)}

\prose{โนสญฺโญชโน ธมฺโม สญฺโญชนสฺส จ โนสญฺโญชนสฺส จ ธมฺมสฺส วิปฺปยุตฺตปจฺจเยน ปจฺจโย.\csromanspacer{} \textbf{ปุเรชาตํ} วตฺถุ สญฺโญชนานํ สมฺปยุตฺตกานญฺจ ขนฺธานํ วิปฺปยุตฺตปจฺจเยน ปจฺจโย. (3)}

\prose{สญฺโญชโน จ โนสญฺโญชโน จ ธมฺมา โนสญฺโญชนสฺส ธมฺมสฺส วิปฺปยุตฺตปจฺจเยน ปจฺจโย.\csromanspacer{} สหชาตํ, ปจฺฉาชาตํ. (วิภชิตพฺพํ.) (1)}

\csromanheader{2}{\textbf{อตฺถิ}}
\proseitem{29}{สญฺโญชโน ธมฺโม สญฺโญชนสฺส ธมฺมสฺส อตฺถิปจฺจเยน ปจฺจโย.\csromanspacer{} เอกํ. (ปฏิจฺจสทิสํ.) (1)}

\prose{สญฺโญชโน ธมฺโม โนสญฺโญชนสฺส ธมฺมสฺส อตฺถิปจฺจเยน ปจฺจโย.\csromanspacer{} สหชาตํ, ปจฺฉาชาตํ. (สํขิตฺตํ.) (2)}

\prose{สญฺโญชโน ธมฺโม สญฺโญชนสฺส จ โนสญฺโญชนสฺส จ ธมฺมสฺส อตฺถิปจฺจเยน ปจฺจโย. (ปฏิจฺจสทิสํ.) (3)}

\proseitem{30}{โนสญฺโญชโน ธมฺโม โนสญฺโญชนสฺส ธมฺมสฺส อตฺถิปจฺจเยน ปจฺจโย.\csromanspacer{} สหชาตํ, ปุเรชาตํ, ปจฺฉาชาตํ, อาหารํ, อินฺทฺริยํ (สํขิตฺตํ.) (1)}

\prose{โนสญฺโญชโน ธมฺโม สญฺโญชนสฺส ธมฺมสฺส อตฺถิปจฺจเยน ปจฺจโย.\csromanspacer{} สหชาตํ, ปุเรชาตํ.\csromanspacer{}\csromanspacer{}\textbf{สหชาตา} โนสญฺโญชนา ขนฺธา สมฺปยุตฺตกานํ สญฺโญชนานํ อตฺถิปจฺจเยน ปจฺจโย.\csromanspacer{} \textbf{ปุเรชาตํ} จกฺขุํ ฯเปฯ วตฺถุํ อสฺสาเทติ อภินนฺทติ, ตํ อารพฺภ ราโค อุปฺปชฺชติ ฯเปฯ โทมนสฺสํ อุปฺปชฺชติ.\csromanspacer{} วตฺถุ อญฺโญชนานํ อตฺถิปจฺจเยน ปจฺจโย. (2)}

\prose{\csromanfolio{220}โนสญฺโญชโน ธมฺโม สญฺโญชนสฺส จ โนสญฺโญชนสฺส จ ธมฺมสฺส อตฺถิปจฺจเยน ปจฺจโย.\csromanspacer{} สหชาตํ, ปุเรชาตํ.\csromanspacer{}\csromanspacer{}\textbf{สหชาโต} โนสญฺโญชโน ฯเปฯ. (สํขิตฺตํ.\csromanspacer{} อาสวสทิสํ.) (3)}

\proseitem{31}{สญฺโญชโน จ โนสญฺโญชโน จ ธมฺมา สญฺโญชนสฺส ธมฺมสฺส อตฺถิปจฺจเยน ปจฺจโย.\csromanspacer{} สหชาตํ, ปุเรชาตํ. (อาสวสทิสํ.) (1)}

\prose{สญฺโญชโน จ โนสญฺโญชโน จ ธมฺมา โนสญฺโญชนสฺส ธมฺมสฺส อตฺถิปจฺจเยน ปจฺจโย.\csromanspacer{} สหชาตํ, ปุเรชาตํ, ปจฺฉาชาตํ, อาหารํ, อินฺทฺริยํ. (วิภชิตพฺพํ.\csromanspacer{} อาสวสทิสํ.) (2)}

\prose{สญฺโญชโน จ โนสญฺโญชโน จ ธมฺมา สญฺโญชนสฺส จ โนสญฺโญชนสฺส จ ธมฺมสฺส อตฺถิปจฺจเยน ปจฺจโย.\csromanspacer{} สหชาตํ, ปุเรชาตํ. (วิภชิตพฺพํ.\csromanspacer{} อาสวสทิสํ.) (3)}

\csromanheader{2}{1. ปจฺจยานุโลม 2. สงฺขฺยาวาร}
\csromanheader{2}{\textbf{สุทฺธ}}
\proseitem{32}{เหตุยา จตฺตาริ, อารมฺมเณ นว, อธิปติยา นว, อนนฺตเร นว, สมนนฺตเร นว, สหชาเต นว, อญฺญมญฺเญ นว, นิสฺสเย นว, อุปนิสฺสเย นว, ปุเรชาเต ตีณิ, ปจฺฉาชาเต ตีณิ, อาเสวเน นว, กมฺเม ตีณิ, วิปาเก เอกํ, อาหาเร ตีณิ, อินฺทฺริเย ตีณิ, ฌาเน ตีณิ, มคฺเค นว, สมฺปยุตฺเต นว, วิปฺปยุตฺเต ปญฺจ, อตฺถิยา นว, นตฺถิยา นว, วิคเต นว, อวิคเต นว.}

\vspace{\tipitakaparskip}
\csromancenter{อนุโลมํ.}
\csromanheader{2}{2. ปจฺจนียุทฺธาร}
\proseitem{33}{สญฺโญชโน ธมฺโม สญฺโญชนสฺส ธมฺมสฺส อารมฺมณปจฺจเยน ปจฺจโย.\csromanspacer{} สหชาตปจฺจเยน ปจฺจโย.\csromanspacer{} อุปนิสฺสยปจฺจเยน ปจฺจโย. (1)}

\prose{สญฺโญชโน ธมฺโม โนสญฺโญชนสฺส ธมฺมสฺส อารมฺมณปจฺจเยน ปจฺจโย.\csromanspacer{} สหชาตปจฺจเยน ปจฺจโย.\csromanspacer{} อุปนิสฺสยปจฺจเยน ปจฺจโย.\csromanspacer{} ปจฺฉาชาตปจฺจเยน ปจฺจโย. (2)}

\csromanheader{2}{20. สญฺโญชนทุก}
\prose{\csromanfolio{221}สญฺโญชโน ธมฺโม สญฺโญชนสฺส จ โนสญฺโญชนสฺส จ ธมฺมสฺส อารมฺมณปจฺจเยน ปจฺจโย.\csromanspacer{} สหชาตปจฺจเยน ปจฺจโย.\csromanspacer{} อุปนิสฺสยปจฺจเยน ปจฺจโย. (3)}

\proseitem{34}{โนสญฺโญชโน ธมฺโม โนสญฺโญชนสฺส ธมฺมสฺส อารมฺมณปจฺจเยน ปจฺจโย.\csromanspacer{} สหชาตปจฺจเยน ปจฺจโย.\csromanspacer{} อุปนิสฺสยปจฺจเยน ปจฺจโย.\csromanspacer{} ปุเรชาตปจฺจเยน ปจฺจโย.\csromanspacer{} ปจฺฉาชาตปจฺจเยน ปจฺจโย.\csromanspacer{} กมฺมปจฺจเยน ปจฺจโย.\csromanspacer{} อาหารปจฺจเยน ปจฺจโย.\csromanspacer{} อินฺทฺริยปจฺจเยน ปจฺจโย. (1)}

\prose{โนสญฺโญชโน ธมฺโม สญฺโญชนสฺส ธมฺมสฺส อารมฺมณปจฺจเยน ปจฺจโย.\csromanspacer{} สหชาตปจฺจเยน ปจฺจโย.\csromanspacer{} อุปนิสฺสยปจฺจเยน ปจฺจโย.\csromanspacer{} ปุเรชาตปจฺจเยน ปจฺจโย. (2)}

\prose{โนสญฺโญชโน ธมฺโม สญฺโญชนสฺส จ โนสญฺโญชนสฺส จ ธมฺมสฺส อารมฺมณปจฺจเยน ปจฺจโย.\csromanspacer{} สหชาตปจฺจเยน ปจฺจโย.\csromanspacer{} อุปนิสฺสยปจฺจเยน ปจฺจโย.\csromanspacer{} ปุเรชาตปจฺจเยน ปจฺจโย. (3)}

\proseitem{35}{สญฺโญชโน จ โนสญฺโญชโน จ ธมฺมา สญฺโญชนสฺส ธมฺมสฺส อารมฺมณปจฺจเยน ปจฺจโย.\csromanspacer{} สหชาตปจฺจเยน ปจฺจโย.\csromanspacer{} อุปนิสฺสยปจฺจเยน ปจฺจโย. (1)}

\prose{สญฺโญชโน จ โนสญฺโญชโน จ ธมฺมา โนสญฺโญชนสฺส ธมฺมสฺส อารมฺมณปจฺจเยน ปจฺจโย.\csromanspacer{} สหชาตปจฺจเยน ปจฺจโย.\csromanspacer{} อุปนิสฺสยปจฺจเยน ปจฺจโย.\csromanspacer{} ปจฺฉาชาตปจฺจเยน ปจฺจโย. (2)}

\prose{สญฺโญชโน จ โนสญฺโญชโน จ ธมฺมา สญฺโญชนสฺส จ โนสญฺโญชนสฺส จ ธมฺมสฺส อารมฺมณปจฺจเยน ปจฺจโย.\csromanspacer{} สหชาตปจฺจเยน ปจฺจโย.\csromanspacer{} อุปนิสฺสยปจฺจเยน ปจฺจโย. (3)}

\csromanheader{2}{2. ปจฺจยปจฺจนีย 2. สงฺขฺยาวาร}
\vspace{\tipitakaparskip}
\csromancenter{นเหตุยา นว, น-อารมฺมเณ นว, (สพฺพตฺถ นว.) โน-อวิคเต นว.}
\csromanheader{2}{3. ปจฺจยานุโลมปจฺจนีย}
\csromanheader{2}{\textbf{เหตุทุก}}
\proseitem{37}{\csromanfolio{222}\textbf{เหตุ}ปจฺจยา น-อารมฺมเณ จตฺตาริ ฯเปฯ นสมนนฺตเร จตฺตาริ, น-อญฺญมญฺเญ เทฺว, น-อุปนิสฺสเย จตฺตาริ ฯเปฯ นมคฺเค จตฺตาริ, นสมฺปยุตฺเต เทฺว, นวิปฺปยุตฺเต จตฺตาริ, โนนตฺถิยา จตฺตาริ, โนวิคเต จตฺตาริ.}

\csromanheader{2}{4. ปจฺจยปจฺจนียานุโลม}
\csromanheader{2}{\textbf{นเหตุทุก}}
\csromanheader{2}{38. นเหตุปจฺจยา อารมฺมเณ นว, อธิปติยา นว, (อนุโลมมาติกา) อวิคเต นว.}
\nitthitam{สญฺโญชนทุกํ นิฏฺฐิตํ.}
\csromanmatikaanchor{mtk.26}
\csromantocmarkref{section}{21. สญฺโญชนิยทุก}{mtk.26}
\bookmark[dest={mtk.26},level=1]{21. สญฺโญชนิยทุก}
\csromanheader{1}{21. สญฺโญชนิยทุก 1-7. ปฏิจฺจาทิวาร}
\proseitem{39}{สญฺโญชนิยํ ธมฺมํ ปฏิจฺจ สญฺโญชนิโย ธมฺโม อุปฺปชฺชติ เหตุปจฺจยา.\csromanspacer{} สญฺโญชนิยํ เอกํ ขนฺธํ ปฏิจฺจ ตโย ขนฺธา จิตฺตสมุฏฺฐานญฺจ รูปํ ฯเปฯ เทฺว ขนฺเธ ฯเปฯ.\csromanspacer{} ปฏิสนฺธิกฺขเณ ฯเปฯ.\csromanspacer{} ขนฺเธ ปฏิจฺจ วตฺถุ, วตฺถุํ ปฏิจฺจ ขนฺธา.\csromanspacer{} เอกํ มหาภูตํ ฯเปฯ. (จูฬนฺตรทุเก โลกิยทุกสทิสํ, นินฺนานากรณํ.)}

\nitthitam{สญฺโญชนิยทุกํ นิฏฺฐิตํ.}
\csromanmatikaanchor{mtk.27}
\csromantocmarkref{section}{22. สญฺโญชนสมฺปยุตฺตทุก}{mtk.27}
\bookmark[dest={mtk.27},level=1]{22. สญฺโญชนสมฺปยุตฺตทุก}
\csromanheader{1}{22. สญฺโญชนสมฺปยุตฺตทุก 1. ปฏิจฺจวาร}
\csromanheader{2}{1. ปจฺจยานุโลม 1. วิภงฺควาร}
\csromanheader{2}{\textbf{เหตุ}}
\proseitem{40}{สญฺโญชนสมฺปยุตฺตํ ธมฺมํ ปฏิจฺจ สญฺโญชนสมฺปยุตฺโต ธมฺโม อุปฺปชฺชติ เหตุปจฺจยา.\csromanspacer{} สญฺโญชนสมฺปยุตฺตํ เอกํ ขนฺธํ ปฏิจฺจ ตโย ขนฺธา ฯเปฯ เทฺว ขนฺเธ ฯเปฯ. (1)}

\csromanheader{2}{22. สญฺโญชนสมฺปยุตฺตทุก}
\prose{\csromanfolio{223}สญฺโญชนสมฺปยุตฺตํ ธมฺมํ ปฏิจฺจ สญฺโญชนวิปฺปยุตฺโต ธมฺโม อุปฺปชฺชติ เหตุปจฺจยา.\csromanspacer{} สญฺโญชนสมฺปยุตฺเต ขนฺเธ ปฏิจฺจ จิตฺตสมุฏฺฐานํ รูปํ. (2)}

\prose{สญฺโญชนสมฺปยุตฺตํ ธมฺมํ ปฏิจฺจ สญฺโญชนสมฺปยุตฺโต จ สญฺโญชนวิปฺปยุตฺโต จ ธมฺมา อุปฺปชฺชนฺติ เหตุปจฺจยา.\csromanspacer{} สญฺโญชนสมฺปยุตฺตํ เอกํ ขนฺธํ ปฏิจฺจ ตโย ขนฺธา จิตฺตสมุฏฺฐานญฺจ รูปํ ฯเปฯ เทฺว ขนฺเธ ฯเปฯ. (3)}

\proseitem{41}{สญฺโญชนวิปฺปยุตฺตํ ธมฺมํ ปฏิจฺจ สญฺโญชนวิปฺปยุตฺโต ธมฺโม อุปฺปชฺชติ เหตุปจฺจยา.\csromanspacer{} สญฺโญชนวิปฺปยุตฺตํ เอกํ ขนฺธํ ปฏิจฺจ ตโย ขนฺธา จิตฺตสมุฏฺฐานญฺจ รูปํ ฯเปฯ เทฺว ขนฺเธ ฯเปฯ.\csromanspacer{} อุทฺธจฺจสหคตํ โมหํ ปฏิจฺจ จิตฺตสมุฏฺฐานํ รูปํ.\csromanspacer{} ปฏิสนฺธิกฺขเณ ฯเปฯ.\csromanspacer{} เอกํ มหาภูตํ ฯเปฯ. (1)}

\prose{สญฺโญชนวิปฺปยุตฺตํ ธมฺมํ ปฏิจฺจ สญฺโญชนสมฺปยุตฺโต ธมฺโม อุปฺปชฺชติ เหตุปจฺจยา.\csromanspacer{} อุทฺธจฺจสหคตํ โมหํ ปฏิจฺจ สมฺปยุตฺตกา ขนฺธา. (2)}

\prose{สญฺโญชนวิปฺปยุตฺตํ ธมฺมํ ปฏิจฺจ สญฺโญชนสมฺปยุตฺโต จ สญฺโญชนวิปฺปยุตฺโต จ ธมฺมา อุปฺปชฺชนฺติ เหตุปจฺจยา.\csromanspacer{} อุทฺธจฺจสหคตํ โมหํ ปฏิจฺจ สมฺปยุตฺตกา ขนฺธา จิตฺตสมุฏฺฐานญฺจ รูปํ. (3)}

\proseitem{42}{สญฺโญชนสมฺปยุตฺตญฺจ สญฺโญชนวิปฺปยุตฺตญฺจ ธมฺมํ ปฏิจฺจ สญฺโญชนสมฺปยุตฺโต ธมฺโม อุปฺปชฺชติ เหตุปจฺจยา.\csromanspacer{} อุทฺธจฺจสหคตํ เอกํ ขนฺธญฺจ โมหญฺจ ปฏิจฺจ ตโย ขนฺธา ฯเปฯ เทฺว ขนฺเธ ฯเปฯ. (1)}

\prose{สญฺโญชนสมฺปยุตฺตญฺจ สญฺโญชนวิปฺปยุตฺตญฺจ ธมฺมํ ปฏิจฺจ สญฺโญชนวิปฺปยุตฺโต ธมฺโม อุปฺปชฺชติ เหตุปจฺจยา.\csromanspacer{} สญฺโญชนสมฺปยุตฺเต ขนฺเธ จ มหาภูเต จ ปฏิจฺจ จิตฺตสมุฏฺฐานํ รูปํ.\csromanspacer{} อุทฺธจฺจสหคเต ขนฺเธ จ โมหญฺจ ปฏิจฺจ จิตฺตสมุฏฺฐานํ รูปํ.}

\prose{สญฺโญชนสมฺปยุตฺตญฺจ สญฺโญชนวิปฺปยุตฺตญฺจ ธมฺมํ ปฏิจฺจ สญฺโญชนสมฺปยุตฺโต จ สญฺโญชนวิปฺปยุตฺโต จ ธมฺมา อุปฺปชฺชนฺติ เหตุปจฺจยา.\csromanspacer{} อุทฺธจฺจสหคตํ เอกํ ขนฺธญฺจ โมหญฺจ ปฏิจฺจ ตโย ขนฺธา จิตฺตสมุฏฺฐานญฺจ รูปํ ฯเปฯ เทฺว ขนฺเธ ฯเปฯ. (3)}

\csromanheader{2}{\textbf{อารมฺมณ}}
\proseitem{43}{\csromanfolio{224}สญฺโญชนสมฺปยุตฺตํ ธมฺมํ ปฏิจฺจ สญฺโญชนสมฺปยุตฺโต ธมฺโม อุปฺปชฺชติ อารมฺมณปจฺจยา.\csromanspacer{} สญฺโญชนสมฺปยุตฺตํ เอกํ ขนฺธํ ฯเปฯ เทฺว ขนฺเธ ฯเปฯ. (1)}

\prose{สญฺโญชนสมฺปยุตฺตํ ธมฺมํ ปฏิจฺจ สญฺโญชนวิปฺปยุตฺโต ธมฺโม อุปฺปชฺชติ อารมฺมณปจฺจยา.\csromanspacer{} อุทฺธจฺจสหคเต ขนฺเธ ปฏิจฺจ อุทฺธจฺจสหคโต โมโห. (2)}

\prose{สญฺโญชนสมฺปยุตฺตํ ธมฺมํ ปฏิจฺจ สญฺโญชนสมฺปยุตฺโต จ สญฺโญชนวิปฺปยุตฺโต จ ธมฺมา อุปฺปชฺชนฺติ อารมฺมณปจฺจยา.\csromanspacer{} อุทฺธจฺจสหคตํ เอกํ ขนฺธํ ปฏิจฺจ ตโย ขนฺธา โมโห จ ฯเปฯ เทฺว ขนฺเธ ฯเปฯ. (3)}

\proseitem{44}{สญฺโญชนวิปฺปยุตฺตํ ธมฺมํ ปฏิจฺจ สญฺโญชนวิปฺปยุตฺโต ธมฺโม อุปฺปชฺชติ อารมฺมณปจฺจยา.\csromanspacer{} สญฺโญชนวิปฺปยุตฺตํ เอกํ ขนฺธํ ปฏิจฺจ ตโย ขนฺธา ฯเปฯ เทฺว ขนฺเธ ฯเปฯ.\csromanspacer{} ปฏิสนฺธิกฺขเณ ฯเปฯ.\csromanspacer{} วตฺถุํ ปฏิจฺจ ขนฺธา. (1)}

\prose{สญฺโญชนวิปฺปยุตฺตํ ธมฺมํ ปฏิจฺจ สญฺโญชนสมฺปยุตฺโต ธมฺโม อุปฺปชฺชติ อารมฺมณปจฺจยา.\csromanspacer{} อุทฺธจฺจสหคตํ โมหํ ปฏิจฺจ สมฺปยุตฺตกา ขนฺธา. (2)}

\prose{สญฺโญชนสมฺปยุตฺตญฺจ สญฺโญชนวิปฺปยุตฺตญฺจ ธมฺมํ ปฏิจฺจ สญฺโญชนสมฺปยุตฺโต ธมฺโม อุปฺปชฺชติ อารมฺมณปจฺจยา.\csromanspacer{} อุทฺธจฺจสหคตํ เอกํ ขนฺธญฺจ โมหญฺจ ปฏิจฺจ ตโย ขนฺธา ฯเปฯ เทฺว ขนฺเธ ฯเปฯ. (1)}

\csromanheader{2}{\textbf{อธิปติ}}
\proseitem{45}{สญฺโญชนสมฺปยุตฺตํ ธมฺมํ ปฏิจฺจ สญฺโญชนสมฺปยุตฺโต ธมฺโม อุปฺปชฺชติ อธิปติปจฺจยา.\csromanspacer{} ตีณิ.}

\prose{สญฺโญชนวิปฺปยุตฺตํ ธมฺมํ ปฏิจฺจ สญฺโญชนวิปฺปยุตฺโต ธมฺโม อุปฺปชฺชติ อธิปติปจฺจยา.\csromanspacer{} เอกํ.}

\prose{สญฺโญชนสมฺปยุตฺตญฺจ สญฺโญชนวิปฺปยุตฺตญฺจ ธมฺมํ ปฏิจฺจ สญฺโญชนวิปฺปยุตฺโต ธมฺโม อุปฺปชฺชติ อธิปติปจฺจยา.\csromanspacer{} สญฺโญชนสมฺปยุตฺตเก ขนฺเธ จ มหาภูเต จ ปฏิจฺจ จิตฺตสมุฏฺฐานํ รูปํ. (1) (สํขิตฺตํ.)}

\csromanheader{2}{22. สญฺโญชนสมฺปยุตฺตทุก \textbf{1. ปจฺจยานุโลม 2. สงฺขฺยาวาร}}
\csromanheader{2}{\textbf{สุทฺธ}}
\proseitem{46}{\csromanfolio{225}เหตุยา นว, อารมฺมเณ ฉ, อธิปติยา ปญฺจ, อนนฺตเร ฉ, สมนนฺตเร ฉ, สหชาเต นว, อญฺญมญฺเญ ฉ, นิสฺสเย นว, อุปนิสฺสเย ฉ, ปุเรชาเต ฉ, อาเสวเน ฉ, กมฺเม นว, วิปาเก เอกํ, อาหาเร นว, อินฺทฺริเย นว, ฌาเน นว, มคฺเค นว, สมฺปยุตฺเต ฉ, วิปฺปยุตฺเต นว, อตฺถิยา นว, นตฺถิยา ฉ, วิคเต ฉ, อวิคเต นว.}

\vspace{\tipitakaparskip}
\csromancenter{อนุโลมํ.}
\csromanheader{2}{2. ปจฺจยปจฺจนีย 1. วิภงฺควาร}
\csromanheader{2}{\textbf{นเหตุ}}
\proseitem{47}{สญฺโญชนสมฺปยุตฺตํ ธมฺมํ ปฏิจฺจ สญฺโญชนสมฺปยุตฺโต ธมฺโม อุปฺปชฺชติ นเหตุปจฺจยา.\csromanspacer{} วิจิกิจฺฉาสหคเต ขนฺเธ ปฏิจฺจ วิจิกิจฺฉาสหคโต โมโห. (1)}

\prose{สญฺโญชนสมฺปยุตฺตํ ธมฺมํ ปฏิจฺจ สญฺโญชนวิปฺปยุตฺโต ธมฺโม อุปฺปชฺชติ นเหตุปจฺจยา.\csromanspacer{} อุทฺธจฺจสหคเต ขนฺเธ ปฏิจฺจ อุทฺธจฺจสหคโต โมโห. (2)}

\prose{สญฺโญชนวิปฺปยุตฺตํ ธมฺมํ ปฏิจฺจ สญฺโญชนวิปฺปยุตฺโต ธมฺโม อุปฺปชฺชติ นเหตุปจฺจยา.\csromanspacer{} อเหตุกํ สญฺโญชนวิปฺปยุตฺตํ เอกํ ขนฺธํ ฯเปฯ. (ยาว อสญฺญสตฺตา.\csromanspacer{} สํขิตฺตํ.) (1)}

\csromanheader{2}{2. ปจฺจยปจฺจนีย 2. สงฺขฺยาวาร}
\csromanheader{2}{\textbf{สุทฺธ}}
\vspace{\tipitakaparskip}
\csromancenter{\textbf{นเหตุ}ยา ตีณิ, น-อารมฺมเณ ตีณิ, น-อธิปติยา นว, น-อนนฺตเร ตีณิ, นสมนนฺตเร ตีณิ, น-อญฺญมญฺเญ ตีณิ, น-อุปนิสฺสเย ตีณิ, นปุเรชาเต สตฺต, นปจฺฉาชาเต นว, น-อาเสวเน นว, นกมฺเม}
\prosecont{\csromanfolio{226}จตฺตาริ, นวิปาเก นว, น-อาหาเร เอกํ, น-อินฺทฺริเย เอกํ, นฌาเน เอกํ, นมคฺเค เอกํ, นสมฺปยุตฺเต ตีณิ, นวิปฺปยุตฺเต ฉ, โนนตฺถิยา ตีณิ, โนวิคเต ตีณิ.}

\vspace{\tipitakaparskip}
\csromancenter{ปจฺจนียํ.}
\csromanheader{2}{3. ปจฺจยานุโลมปจฺจนีย}
\csromanheader{2}{\textbf{เหตุทุก}}
\proseitem{49}{\textbf{เหตุ}ปจฺจยา น-อารมฺมเณ ตีณิ, น-อธิปติยา นว ฯเปฯ น-อุปนิสฺสเย ตีณิ, นปุเรชาเต ฉ, นปจฺฉาชาเต นว, น-อาเสวเน นว, นกมฺเม จตฺตาริ, นวิปาเก นว, นสมฺปยุตฺเต ตีณิ, นวิปฺปยุตฺเต จตฺตาริ, โนนตฺถิยา ตีณิ, โนวิคเต ตีณิ.}

\csromanheader{2}{4. ปจฺจยปจฺจนียานุโลม}
\csromanheader{2}{\textbf{นเหตุทุก}}
\proseitem{50}{นเหตุปจฺจยา อารมฺมเณ ตีณิ ฯเปฯ วิปาเก เอกํ, อาหาเร ตีณิ ฯเปฯ มคฺเค เทฺว ฯเปฯ อวิคเต ตีณิ.}

\vspace{\tipitakaparskip}
\csromancenter{(สหชาตวาโร ปฏิจฺจวารสทิโส.)}
\csromanheader{2}{22. สญฺโญชนสมฺปยุตฺตทุก 3. ปจฺจยวาร}
\csromanheader{2}{1. ปจฺจยานุโลม 1. วิภงฺควาร}
\csromanheader{2}{\textbf{เหตุ}}
\proseitem{51}{สญฺโญชนสมฺปยุตฺตํ ธมฺมํ ปจฺจยา สญฺโญชนสมฺปยุตฺโต ธมฺโม อุปฺปชฺชติ เหตุปจฺจยา.\csromanspacer{} ตีณิ. (ปฏิจฺจสทิสํ.)}

\csromanheader{2}{22. สญฺโญชนสมฺปยุตฺตทุก}
\prose{\csromanfolio{227}สญฺโญชนวิปฺปยุตฺตํ ธมฺมํ ปจฺจยา สญฺโญชนวิปฺปยุตฺโต ธมฺโม อุปฺปชฺชติ เหตุปจฺจยา. (ยาว ปฏิสนฺธิ.) เอกํ มหาภูตํ ฯเปฯ วตฺถุํ ปจฺจยา สญฺโญชนวิปฺปยุตฺตา ขนฺธา.\csromanspacer{} วตฺถุํ ปจฺจยา อุทฺธจฺจสหคโต โมโห. (1)}

\prose{สญฺโญชนวิปฺปยุตฺตํ ธมฺมํ ปจฺจยา สญฺโญชนสมฺปยุตฺโต ธมฺโม อุปฺปชฺชติ เหตุปจฺจยา.\csromanspacer{} วตฺถุํ ปจฺจยา สญฺโญชนสมฺปยุตฺตกา ขนฺธา.\csromanspacer{} อุทฺธจฺจสหคตํ โมหํ ปจฺจยา สมฺปยุตฺตกา ขนฺธา. (2)}

\prose{สญฺโญชนวิปฺปยุตฺตํ ธมฺมํ ปจฺจยา สญฺโญชนสมฺปยุตฺโต จ สญฺโญชนวิปฺปยุตฺโต จ ธมฺมา อุปฺปชฺชนฺติ เหตุปจฺจยา.\csromanspacer{} วตฺถุํ ปจฺจยา สญฺโญชนสมฺปยุตฺตกา ขนฺธา.\csromanspacer{} มหาภูเต ปจฺจยา จิตฺตสมุฏฺฐานํ รูปํ.\csromanspacer{} อุทฺธจฺจสหคตํ โมหํ ปจฺจยา สมฺปยุตฺตกา ขนฺธา จิตฺตสมุฏฺฐานญฺจ รูปํ. (3)}

\proseitem{52}{สญฺโญชนสมฺปยุตฺตญฺจ สญฺโญชนวิปฺปยุตฺตญฺจ ธมฺมํ ปจฺจยา สญฺโญชนสมฺปยุตฺโต ธมฺโม อุปฺปชฺชติ เหตุปจฺจยา.\csromanspacer{} สญฺโญชนสมฺปยุตฺตํ เอกํ ขนฺธญฺจ วตฺถุญฺจ ปจฺจยา ตโย ขนฺธา ฯเปฯ เทฺว ขนฺเธ ฯเปฯ.\csromanspacer{} อุทฺธจฺจสหคตํ เอกํ ขนฺธญฺจ โมหญฺจ ปจฺจยา ตโย ขนฺธา ฯเปฯ เทฺว ขนฺเธ ฯเปฯ. (1)}

\prose{สญฺโญชนสมฺปยุตฺตญฺจ สญฺโญชนวิปฺปยุตฺตญฺจ ธมฺมํ ปจฺจยา สญฺโญชนวิปฺปยุตฺโต ธมฺโม อุปฺปชฺชติ เหตุปจฺจยา.\csromanspacer{} สญฺโญชนสมฺปยุตฺเต ขนฺเธ จ มหาภูเต จ ปจฺจยา จิตฺตสมุฏฺฐานํ รูปํ.\csromanspacer{} อุทฺธจฺจสหคเต ขนฺเธ จ โมหญฺจ ปจฺจยา จิตฺตสมุฏฺฐานํ รูปํ.}

\prose{สญฺโญชนสมฺปยุตฺตญฺจ สญฺโญชนวิปฺปยุตฺตญฺจ ธมฺมํ ปจฺจยา สญฺโญชนสมฺปยุตฺโต จ สญฺโญชนวิปฺปยุตฺโต จ ธมฺมา อุปฺปชฺชนฺติ เหตุปจฺจยา.\csromanspacer{} สญฺโญชนสมฺปยุตฺตํ เอกํ ขนฺธญฺจ วตฺถุญฺจ ปจฺจยา ตโย ขนฺธา ฯเปฯ เทฺว ขนฺเธ ฯเปฯ.\csromanspacer{} สญฺโญชนสมฺปยุตฺเต ขนฺเธ จ มหาภูเต จ ปจฺจยา จิตฺตสมุฏฺฐานํ รูปํ.\csromanspacer{} อุทฺธจฺจสหคตํ เอกํ ขนฺธญฺจ โมหญฺจ ปจฺจยา ตโย ขนฺธา จิตฺตสมุฏฺฐานญฺจ รูปํ ฯเปฯ เทฺว ขนฺเธ ฯเปฯ. (3)}

\csromanheader{2}{\textbf{อารมฺมณ}}
\proseitem{53}{สญฺโญชนสมฺปยุตฺตํ ธมฺมํ ปจฺจยา สญฺโญชนสมฺปยุตฺโต ธมฺโม อุปฺปชฺชติ อารมฺมณปจฺจยา.\csromanspacer{} ตีณิ. (ปฏิจฺจสทิสา.)}

\prose{สญฺโญชนวิปฺปยุตฺตํ ธมฺมํ ปจฺจยา สญฺโญชนวิปฺปยุตฺโต ธมฺโม อุปฺปชฺชติ อารมฺมณปจฺจยา. (ยาว ปฏิสนฺธิ.) วตฺถุํ ปจฺจยา สญฺโญชนวิปฺปยุตฺตา \csromanfolio{228}ขนฺธา.\csromanspacer{} จกฺขายตนํ ปจฺจยา จกฺขุวิญฺญาณํ ฯเปฯ.\csromanspacer{} กายายตนํ ปจฺจยา กายวิญฺญาณํ.\csromanspacer{} วตฺถุํ ปจฺจยา สญฺโญชนวิปฺปยุตฺตา ขนฺธา.\csromanspacer{} วตฺถุํ ปจฺจยา อุทฺธจฺจสหคโต โมโห. (1)}

\prose{สญฺโญชนวิปฺปยุตฺตํ ธมฺมํ ปจฺจยา สญฺโญชนสมฺปยุตฺโต ธมฺโม อุปฺปชฺชติ อารมฺมณปจฺจยา.\csromanspacer{} วตฺถุํ ปจฺจยา สญฺโญชนสมฺปยุตฺตกา ขนฺธา.\csromanspacer{} อุทฺธจฺจสหคตํ โมหํ ปจฺจยา สมฺปยุตฺตกา ขนฺธา. (2)}

\prose{สญฺโญชนวิปฺปยุตฺตํ ธมฺมํ ปจฺจยา สญฺโญชนสมฺปยุตฺโต จ สญฺโญชนวิปฺปยุตฺโต จ ธมฺมา อุปฺปชฺชนฺติ อารมฺมณปจฺจยา.\csromanspacer{} วตฺถุํ ปจฺจยา อุทฺธจฺจสหคตา ขนฺธา จ โมโห จ. (3)}

\proseitem{54}{สญฺโญชนสมฺปยุตฺตญฺจ สญฺโญชนวิปฺปยุตฺตญฺจ ธมฺมํ ปจฺจยา สญฺโญชนสมฺปยุตฺโต ธมฺโม อุปฺปชฺชติ อารมฺมณปจฺจยา.\csromanspacer{} สญฺโญชนสมฺปยุตฺตํ เอกํ ขนฺธญฺจ วตฺถุญฺจ ปจฺจยา ตโย ขนฺธา ฯเปฯ เทฺว ขนฺเธ ฯเปฯ.\csromanspacer{} อุทฺธจฺจสหคตํ เอกํ ขนฺธญฺจ โมหญฺจ ปจฺจยา ตโย ขนฺธา ฯเปฯ เทฺว ขนฺเธ ฯเปฯ. (1)}

\prose{สญฺโญชนสมฺปยุตฺตญฺจ สญฺโญชนวิปฺปยุตฺตญฺจ ธมฺมํ ปจฺจยา สญฺโญชนวิปฺปยุตฺโต ธมฺโม อุปฺปชฺชติ อารมฺมณปจฺจยา.\csromanspacer{} อุทฺธจฺจสหคเต ขนฺเธ จ วตฺถุญฺจ ปจฺจยา อุทฺธจฺจสหคโต โมโห.}

\prose{สญฺโญชนสมฺปยุตฺตญฺจ สญฺโญชนวิปฺปยุตฺตญฺจ ธมฺมํ ปจฺจยา สญฺโญชนสมฺปยุตฺโต จ สญฺโญชนวิปฺปยุตฺโต จ ธมฺมา อุปฺปชฺชนฺติ อารมฺมณปจฺจยา.\csromanspacer{} อุทฺทจฺจสหคตํ เอกํ ขนฺธญฺจ วตฺถุญฺจ ปจฺจยา ตโย ขนฺธา โมโห จ ฯเปฯ เทฺว ขนฺเธ ฯเปฯ. (3)}

\csromanheader{2}{\textbf{อธิปตฺยาทิ}}
\proseitem{55}{สญฺโญชนสมฺปยุตฺตํ ธมฺมํ ปจฺจยา สญฺโญชนสมฺปยุตฺโต ธมฺโม อุปฺปชฺชติ อธิปติปจฺจยา ฯเปฯ อวิคตปจฺจยา ฯเปฯ.}

\csromanheader{2}{1. ปจฺจยานุโลม 2. สงฺขฺยาวาร}
\csromanheader{2}{\textbf{สุทฺธ}}
\proseitem{56}{เหตุยา นว, อารมฺมเณ นว, อธิปติยา นว, (สพฺพตฺถ นว.) วิปาเก เอกํ, อาหาเร นว ฯเปฯ อวิคเต นว.}

\vspace{\tipitakaparskip}
\csromancenter{อนุโลมํ.}
\csromanheader{2}{22. สญฺโญชนสมฺปยุตฺตทุก \textbf{2. ปจฺจยปจฺจนีย 1. วิภงฺควาร}}
\csromanheader{2}{\textbf{นเหตุ}}
\proseitem{57}{\csromanfolio{229}สญฺโญชนสมฺปยุตฺตํ ธมฺมํ ปจฺจยา สญฺโญชนสมฺปยุตฺโต ธมฺโม อุปฺปชฺชติ นเหตุปจฺจยา.\csromanspacer{} วิจิกิจฺฉาสหคเต ขนฺเธ ปจฺจยา วิจิกิจฺฉาสหคโต โมโห. (1)}

\prose{สญฺโญชนสมฺปยุตฺตํ ธมฺมํ ปจฺจยา สญฺโญชนวิปฺปยุตฺโต ธมฺโม อุปฺปชฺชติ นเหตุปจฺจยา.\csromanspacer{} อุทฺธจฺจสหคเต ขนฺเธ ปจฺจยา อุทฺธจฺจสหคโต โมโห. (2)}

\proseitem{58}{สญฺโญชนวิปฺปยุตฺตํ ธมฺมํ ปจฺจยา สญฺโญชนวิปฺปยุตฺโต ธมฺโม อุปฺปชฺชติ นเหตุปจฺจยา.\csromanspacer{} อเหตุกํ สญฺโญชนวิปฺปยุตฺตํ เอกํ ขนฺธํ ฯเปฯ. (ยาว อสญฺญสตฺตา.) จกฺขายตนํ ปจฺจยา จกฺขุวิญฺญาณํ ฯเปฯ.\csromanspacer{} กายายตนํ ปจฺจยา กายวิญฺญาณํ.\csromanspacer{} วตฺถุํ ปจฺจยา อเหตุกา สญฺโญชนวิปฺปยุตฺตา ขนฺธา จ อุทฺธจฺจสหคโต โมโห จ. (1)}

\prose{สญฺโญชนวิปฺปยุตฺตํ ธมฺมํ ปจฺจยา สญฺโญชนสมฺปยุตฺโต ธมฺโม อุปฺปชฺชติ นเหตุปจฺจยา.\csromanspacer{} วตฺถุํ ปจฺจยา วิจิกิจฺฉาสหคโต โมโห. (2)}

\proseitem{59}{สญฺโญชนสมฺปยุตฺตญฺจ สญฺโญชนวิปฺปยุตฺตญฺจ ธมฺมํ ปจฺจยา สญฺโญชนสมฺปยุตฺโต ธมฺโม อุปฺปชฺชติ นเหตุปจฺจยา.\csromanspacer{} วิจิกิจฺฉาสหคเต ขนฺเธ จ วตฺถุญฺจ ปจฺจยา วิจิกิจฺฉาสหคโต โมโห. (1)}

\prose{สญฺโญชนสมฺปยุตฺตญฺจ สญฺโญชนวิปฺปยุตฺตญฺจ ธมฺมํ ปจฺจยา สญฺโญชนวิปฺปยุตฺโต ธมฺโม อุปฺปชฺชติ นเหตุปจฺจยา.\csromanspacer{} อุทฺธจฺจสหคเต ขนฺเธ จ วตฺถุญฺจ ปจฺจยา อุทฺธจฺจสหคโต โมโห. (สํขิตฺตํ.) (2)}

\csromanheader{2}{2. ปจฺจยปจฺจนีย 2. สงฺขฺยาวาร}
\csromanheader{2}{\textbf{สุทฺธ}}
\vspace{\tipitakaparskip}
\csromancenter{\textbf{นเหตุ}ยา ฉ, น-อารมฺมเณ ตีณิ, น-อธิปติยา นว, น-อนนฺตเร ตีณิ, นสมนนฺตเร ตีณิ, น-อญฺญมญฺเญ ตีณิ, น-อุปนิสฺสเย ตีณิ, นปุเรชาเต สตฺต, นปจฺฉาชาเต นว, น-อาเสวเน นว, นกมฺเม จตฺตาริ,}
\prosecont{\csromanfolio{230}นวิปาเก นว, น-อาหาเร เอกํ, น-อินฺทฺริเย เอกํ, นฌาเน เอกํ, นมคฺเค เอกํ, นสมฺปยุตฺเต ตีณิ, นวิปฺปยุตฺเต ฉ, โนนตฺถิยา ตีณิ, โนวิคเต ตีณิ.}

\csromanheader{2}{3. ปจฺจยานุโลมปจฺจนีย}
\csromanheader{2}{\textbf{เหตุทุก}}
\proseitem{61}{เหตุปจฺจยา น-อารมฺมเณ ตีณิ, น-อธิปติยา นว. (สํขิตฺตํ.\csromanspacer{} เอวํ กาตพฺพํ.)}

\csromanheader{2}{4. ปจฺจยปจฺจนียานุโลม}
\csromanheader{2}{\textbf{นเหตุทุก}}
\proseitem{62}{นเหตุปจฺจยา อารมฺมเณ ฉ, (สพฺพตฺถ ฉ,) วิปาเก เอกํ, อาหาเร ฉ ฯเปฯ มคฺเค ฉ ฯเปฯ อวิคเต ฉ.}

\vspace{\tipitakaparskip}
\csromancenter{(นิสฺสยวาโรปิ ปจฺจยวารสทิโส.)}
\csromanheader{2}{22. สญฺโญชนสมฺปยุตฺตทุก 5. สํสฏฺฐวาร}
\csromanheader{2}{1. ปจฺจยานุโลมาทิ}
\proseitem{63}{สญฺโญชนสมฺปยุตฺตํ ธมฺมํ สํสฏฺโฐ สญฺโญชนสมฺปยุตฺโต ธมฺโม อุปฺปชฺชติ เหตุปจฺจยา.\csromanspacer{} สญฺโญชนสมฺปยุตฺตํ เอกํ ขนฺธํ สํสฏฺฐา ตโย ขนฺธา ฯเปฯ เทฺว ขนฺเธ ฯเปฯ. (1)}

\prose{สญฺโญชนวิปฺปยุตฺตํ ธมฺมํ สํสฏฺโฐ สญฺโญชนวิปฺปยุตฺโต ธมฺโม อุปฺปชฺชติ เหตุปจฺจยา.\csromanspacer{} สญฺโญชนวิปฺปยุตฺตํ เอกํ ขนฺธํ สํสฏฺฐา ตโย ขนฺธา ฯเปฯ เทฺว ขนฺเธ ฯเปฯ.\csromanspacer{} ปฏิสนฺธิกฺขเณ ฯเปฯ. (1)}

\prose{สญฺโญชนวิปฺปยุตฺตํ ธมฺมํ สํสฏฺโฐ สญฺโญชนสมฺปยุตฺโต ธมฺโม อุปฺปชฺชติ เหตุปจฺจยา.\csromanspacer{} อุทฺธจฺจสหคตํ โมหํ สํสฏฺฐา สมฺปยุตฺตกา ขนฺธา. (2)}

\csromanheader{2}{22. สญฺโญชนสมฺปยุตฺตทุก}
\prose{\csromanfolio{231}สญฺโญชนสมฺปยุตฺตญฺจ สญฺโญชนวิปฺปยุตฺตญฺจ ธมฺมํ สํสฏฺโฐ สญฺโญชนสมฺปยุตฺโต ธมฺโม อุปฺปชฺชติ เหตุปจฺจยา.\csromanspacer{} อุทฺธจฺจสหคตํ เอกํ ขนฺธญฺจ โมหญฺจ สํสฏฺฐา ตโย ขนฺธา ฯเปฯ เทฺว ขนฺเธ ฯเปฯ. (1) (สํขิตฺตํ.)}

\proseitem{64}{\textbf{เหตุ}ยา จตฺตาริ, อารมฺมเณ ฉ, อธิปติยา เทฺว, อนนฺตเร ฉ, สมนนฺตเร ฉ ฯเปฯ อวิคเต ฉ.}

\proseitem{65}{นเหตุยา ตีณิ, น-อธิปติยา ฉ, นปุเรชาเต ฉ, นปจฺฉาชาเต ฉ, น-อาเสวเน ฉ, นกมฺเม จตฺตาริ, นวิปาเก ฉ, นฌาเน เอกํ, นมคฺเค เอกํ, นวิปฺปยุตฺเต ฉ. (อิตเร เทฺว คณนาปิ สมฺปยุตฺตวาโรปิ กาตพฺโพ.)}

\csromanheader{2}{22. สญฺโญชนสมฺปยุตฺตทุก 7. ปญฺหาวาร}
\csromanheader{2}{1. ปจฺจยานุโลม 1. วิภงฺควาร}
\csromanheader{2}{\textbf{เหตุ}}
\proseitem{66}{สญฺโญชนสมฺปยุตฺโต ธมฺโม สญฺโญชนสมฺปยุตฺตสฺส ธมฺมสฺส เหตุปจฺจเยน ปจฺจโย.\csromanspacer{} สญฺโญชนสมฺปยุตฺตา เหตู สมฺปยุตฺตกานํ ขนฺธานํ เหตุปจฺจเยน ปจฺจโย. (1)}

\prose{สญฺโญชนสมฺปยุตฺโต ธมฺโม สญฺโญชนวิปฺปยุตฺตสฺส ธมฺมสฺส เหตุปจฺจเยน ปจฺจโย.\csromanspacer{} สญฺโญชนสมฺปยุตฺตา เหตู จิตฺตสมุฏฺฐานานํ รูปานํ เหตุปจฺจเยน ปจฺจโย. (2)}

\prose{สญฺโญชนสมฺปยุตฺโต ธมฺโม สญฺโญชนสมฺปยุตฺตสฺส จ สญฺโญชนวิปฺปยุตฺตสฺส จ ธมฺมสฺส เหตุปจฺจเยน ปจฺจโย.\csromanspacer{} สญฺโญชนสมฺปยุตฺตา เหตู สมฺปยุตฺตกานํ ขนฺธานํ จิตฺตสมุฏฺฐานานญฺจ รูปานํ เหตุปจฺจเยน ปจฺจโย. (3)}

\proseitem{67}{สญฺโญชนวิปฺปยุตฺโต ธมฺโม สญฺโญชนวิปฺปยุตฺตสฺส ธมฺมสฺส เหตุปจฺจเยน ปจฺจโย.\csromanspacer{} สญฺโญชนวิปฺปยุตฺโต เหตู สมฺปยุตฺตกานํ ขนฺธานํ จิตฺตสมุฏฺฐานานญฺจ รูปานํ เหตุปจฺจเยน ปจฺจโย.\csromanspacer{} อุทฺธจฺจสหคโต \csromanfolio{232}โมโห จิตฺตสมุฏฺฐานานํ รูปานํ เหตุปจฺจเยน ปจฺจโย.\csromanspacer{} ปฏิสนฺธิกฺขเณ ฯเปฯ. (1)}

\prose{สญฺโญชนวิปฺปยุตฺโต ธมฺโม สญฺโญชนสมฺปยุตฺตสฺส ธมฺมสฺส เหตุปจฺจเยน ปจฺจโย.\csromanspacer{} อุทฺธจฺจสหคโต โมโห สมฺปยุตฺตกานํ ขนฺธานํ เหตุปจฺจเยน ปจฺจโย. (2)}

\prose{สญฺโญชนวิปฺปยุตฺโต ธมฺโม สญฺโญชนสมฺปยุตฺตสฺส จ สญฺโญชนวิปฺปยุตฺตสฺส จ ธมฺมสฺส เหตุปจฺจเยน ปจฺจโย.\csromanspacer{} อุทฺธจฺจสหคโต โมโห สมฺปยุตฺตกานํ ขนฺธาน จิตฺตสมุฏฺฐานาญฺจ รูปานํ เหตุปจฺจเยน ปจฺจโย. (3)}

\csromanheader{2}{\textbf{อารมฺมณ}}
\proseitem{68}{สญฺโญชนสมฺปยุตฺโต ธมฺโม สญฺโญชนสมฺปยุตฺตสฺส ธมฺมสฺส อารมฺมณปจฺจเยน ปจฺจโย.\csromanspacer{} สญฺโญชนสมฺปยุตฺเต ขนฺเธ อารพฺภ สญฺโญชนสมฺปยุตฺตกา ขนฺธา อุปฺปชฺชนฺติ. (มูลํ กาตพฺพํ.) สญฺโญชนสมฺปยุตฺเต ขนฺเธ อารพฺภ สญฺโญชนวิปฺปยุตฺตา ขนฺธา จ โมโห จ อุปฺปชฺชนฺติ. (มูลํ กาตพฺพํ.) สญฺโญชนสมฺปยุตฺเต ขนฺเธ อารพฺภ อุทฺธจฺจสหคตา ขนฺธา จ โมโห จ อุปฺปชฺชนฺติ. (3)}

\proseitem{69}{สญฺโญชนวิปฺปยุตฺโต ธมฺโม สญฺโญชนวิปฺปยุตฺตสฺส ธมฺมสฺส อารมฺมณปจฺจเยน ปจฺจโย.\csromanspacer{} ทานํ ทตฺวา, สีลํ ฯเปฯ อุโปสถกมฺมํ กตฺวา ตํ ปจฺจเวกฺขติ.\csromanspacer{} ปุพฺเพ สุจิณฺณานิ ฯเปฯ.\csromanspacer{} ฌานา ฯเปฯ.\csromanspacer{} อริยา มคฺคา วุฏฺฐหิตฺวา มคฺคํ ปจฺจเวกฺขนฺติ, ผลํ ปจฺจเวกฺขนฺติ, นิพฺพานํ ปจฺจเวกฺขนฺติ.\csromanspacer{} นิพฺพานํ โคตฺรภุสฺส, โวทานสฺส, มคฺคสฺส, ผลสฺส, อาวชฺชนาย อารมฺมณปจฺจเยน ปจฺจโย.\csromanspacer{} อริยา สญฺโญชนวิปฺปยุตฺเต ปหีเน กิเลเส ปจฺจเวกฺขนฺติ, วิกฺขมฺภิเต กิเลเส ปจฺจเวกฺขนฺติ.\csromanspacer{} ปุพฺเพ สมุทาจิณฺเณ กิเลเส ชานนฺติ.\csromanspacer{} จกฺขุํ ฯเปฯ วตฺถุํ สญฺโญชนวิปฺปยุตฺเต ขนฺเธ จ โมหญฺจ อนิจฺจโต ฯเปฯ วิปสฺสติ ฯเปฯ.\csromanspacer{} ทิพฺเพน จกฺขุนา รูปํ ปสฺสติ.\csromanspacer{} ทิพฺพาย โสตธาตุยา สทฺทํ สุณาติ.\csromanspacer{} เจโตปริยญาเณน สญฺโญชนวิปฺปยุตฺตจิตฺตสมงฺคิสฺส จิตฺตํ ชานาติ.\csromanspacer{} อากาสานญฺจายตนํ วิญฺญาณญฺจายตนสฺส ฯเปฯ.\csromanspacer{} รูปายตนํ ฯเปฯ.\csromanspacer{} โผฏฺฐพฺพายตนํ ฯเปฯ.\csromanspacer{} สญฺโญชนวิปฺปยุตฺตา ขนฺธา อิทฺธิวิธญาณสฺส, เจโตปริยญาณสฺส, ปุพฺเพนิวาสานุสฺสติญาณสฺส, ยถากมฺมูปคญาณสฺส,}

\vspace{\tipitakaparskip}
\csromancenter{สญฺโญชนสมฺปยุตฺตทุก อนาคตํสญาณสฺส, อาวชฺชนาย โมหสฺส จ อารมฺมณปจฺจเยน ปจฺจโย. (1)}
\prose{\csromanfolio{233}สญฺโญชนวิปฺปยุตฺโต ธมฺโม สญฺโญชนสมฺปยุตฺตสฺส ธมฺมสฺส อารมฺมณปจฺจเยน ปจฺจโย.\csromanspacer{} ทานํ ทตฺวา, สีลํ ฯเปฯ อุโปสถกมฺมํ ฯเปฯ.\csromanspacer{} ปุพฺเพ สุจิณฺณานิ ฯเปฯ.\csromanspacer{} ฌานา วุฏฺฐหิตฺวา ฌานํ ฯเปฯ.\csromanspacer{} จกฺขุํ ฯเปฯ วตฺถุํ สญฺโญชนวิปฺปยุตฺเต ขนฺเธ จ โมหญฺจ อสฺสาเทติ อภินนฺทติ, ตํ อารพฺภ ราโค อุปฺปชฺชติ ฯเปฯ โทมนสฺสํ อุปฺปชฺชติ. (2)}

\prose{สญฺโญชนวิปฺปยุตฺโต ธมฺโม สญฺโญชนสมฺปยุตฺตสฺส จ สญฺโญชนวิปฺปยุตฺตสฺส จ ธมฺมสฺส อารมฺมณปจฺจเยน ปจฺจโย.\csromanspacer{} จกฺขุํ ฯเปฯ วตฺถุํ สญฺโญชนวิปฺปยุตฺเต ขนฺเธ จ โมหญฺจ อารพฺภ อุทฺธจฺจสหคตา ขนฺธา จ โมโห จ อุปฺปชฺชนฺติ. (3)}

\proseitem{70}{สญฺโญชนสมฺปยุตฺโต จ สญฺโญชนวิปฺปยุตฺโต จ ธมฺมา สญฺโญชนสมฺปยุตฺตสฺส ธมฺมสฺส อารมฺมณปจฺจเยน ปจฺจโย.\csromanspacer{} อุทฺธจฺจสหคเต ขนฺเธ จ โมหญฺจ อารพฺภ สญฺโญชนสมฺปยุตฺตกา ขนฺธา อุปฺปชฺชนฺติ. (มูลํ กาตพฺพํ.) อุทฺธจฺจสหคเต ขนฺเธ จ โมหญฺจ อารพฺภ สญฺโญชนวิปฺปยุตฺตา ขนฺธา จ โมโห จ อุปฺปชฺชนฺติ. (มูลํ กาตพฺพํ.) อุทฺธจฺจสหคเต ขนฺเธ จ โมหญฺจ อารพฺภ อุทฺธจฺจสหคตา ขนฺธา จ โมโห จ อุปฺปชฺชนฺติ. (3)}

\csromanheader{2}{\textbf{อธิปติ}}
\proseitem{71}{สญฺโญชนสมฺปยุตฺโต ธมฺโม สญฺโญชนสมฺปยุตฺตสฺส ธมฺมสฺส อธิปติปจฺจเยน ปจฺจโย.\csromanspacer{} \textbf{อารมฺมณาธิปติ}, สหชาตาธิปติ.\csromanspacer{}\csromanspacer{}\textbf{อารมฺมณาธิปติ}\csromandash{}ราคํ ฯเปฯ.\csromanspacer{} ทิฏฺฐิํ ครุํ กตฺวา อสฺสาเทติ อภินนฺทติ, ตํ ครุํ กตฺวา ราโค อุปฺปชฺชติ, ทิฏฺฐิ อุปฺปชฺชติ.\csromanspacer{} \textbf{สหชาตาธิปติ}\csromandash{} สญฺโญชนสมฺปยุตฺตาธิปติ สมฺปยุตฺตกานํ ขนฺธานํ อธิปติปจฺจเยน ปจฺจโย. (1)}

\prose{สญฺโญชนสมฺปยุตฺโต ธมฺโม สญฺโญชนวิปฺปยุตฺตสฺส ธมฺมสฺส อธิปติปจฺจเยน ปจฺจโย.\csromanspacer{} \textbf{สหชาตาธิปติ}\csromandash{}สญฺโญชนสมฺปยุตฺตาธิปติ จิตฺตสมุฏฺฐานานํ รูปานํ อธิปติปจฺจเยน ปจฺจโย. (2)}

\prose{\csromanfolio{234}สญฺโญชนสมฺปยุตฺโต ธมฺโม สญฺโญชนสมฺปยุตฺตสฺส จ สญฺโญชนวิปฺปยุตฺตสฺส จ ธมฺมสฺส อธิปติปจฺจเยน ปจฺจโย.\csromanspacer{} \textbf{สหชาตาธิปติ\csromandash{}}สญฺโญชนสมฺปยุตฺตาธิปติ สมฺปยุตฺตกานํ ขนฺธานํ จิตฺตสมุฏฺฐานานญฺจ รูปานํ อธิปติปจฺจเยน ปจฺจโย. (3)}

\proseitem{72}{สญฺโญชนวิปฺปยุตฺโต ธมฺโม สญฺโญชนวิปฺปยุตฺตสฺส ธมฺมสฺส อธิปติปจฺจเยน ปจฺจโย.\csromanspacer{} \textbf{อารมฺมณาธิปติ}, สหชาตาธิปติ.\csromanspacer{}\csromanspacer{}\textbf{อารมฺมณาธิปติ\csromandash{}}ทานํ ทตฺวา, สีลํ ฯเปฯ อุโปสถกมฺมํ กตฺวา ตํ ครุํ กตฺวา ปจฺจเวกฺขติ.\csromanspacer{} ปุพฺเพ สุจิณฺณานิ ฯเปฯ.\csromanspacer{} ฌานา ฯเปฯ.\csromanspacer{} อริยา มคฺคา ฯเปฯ.\csromanspacer{} ผลํ ฯเปฯ.\csromanspacer{} นิพฺพานํ ฯเปฯ.\csromanspacer{} นิพฺพานํ โคตฺรภุสฺส, โวทานสฺส, มคฺคสฺส, ผลสฺส อธิปติปจฺจเยน ปจฺจโย.\csromanspacer{} \textbf{สหชาตาธิปติ\csromandash{}} สญฺโญชนวิปฺปยุตฺตาธิปติ สมฺปยุตฺตกานํ ขนฺธานํ จิตฺตสมุฏฺฐานานญฺจ รูปานํ อธิปติปจฺจเยน ปจฺจโย. (1)}

\prose{สญฺโญชนวิปฺปยุตฺโต ธมฺโม สญฺโญชนสมฺปยุตฺตสฺส ธมฺมสฺส อธิปติปจฺจเยน ปจฺจโย.\csromanspacer{} \textbf{อารมฺมณาธิปติ\csromandash{}}ทานํ ทตฺวา, สีลํ ฯเปฯ อุโปสถกมฺมํ กตฺวา ตํ ครุํ กตฺวา อสฺสาเทติ อภินนฺทติ, ตํ ครุํ กตฺวา ราโค อุปฺปชฺชติ, ทิฏฺฐิ อุปฺปชฺชติ.\csromanspacer{} ปุพฺเพ สุจิณฺณานิ ฯเปฯ.\csromanspacer{} ฌานา ฯเปฯ.\csromanspacer{} จกฺขุํ ฯเปฯ วตฺถุํ สญฺโญชนวิปฺปยุตฺเต ขนฺเธ ครุํ กตฺวา อสฺสาเทติ อภินนฺทติ, ตํ ครุํ กตฺวา ราโค อุปฺปชฺชติ, ทิฏฺฐิ อุปฺปชฺชติ. (2)}

\csromanheader{2}{\textbf{อนนฺตร}}
\proseitem{73}{สญฺโญชนสมฺปยุตฺโต ธมฺโม สญฺโญชนสมฺปยุตฺตสฺส ธมฺมสฺส อนนฺตรปจฺจเยน ปจฺจโย.\csromanspacer{} ปุริมา ปุริมา สญฺโญชนสมฺปยุตฺตา ขนฺธา ปจฺฉิมานํ ปจฺฉิมานํ สญฺโญชนสมฺปยุตฺตกานํ ขนฺธานํ อนนฺตรปจฺจเยน ปจฺจโย. (1)}

\prose{สญฺโญชนสมฺปยุตฺโต ธมฺโม สญฺโญชนวิปฺปยุตฺตสฺส ธมฺมสฺส อนนฺตรปจฺจเยน ปจฺจโย.\csromanspacer{} ปุริมา ปุริมา อุทฺธจฺจสหคตา ขนฺธา ปจฺฉิมสฺส ปจฺฉิมสฺส อุทฺธจฺจสหคตสฺส โมหสฺส อนนฺตรปจฺจเยน ปจฺจโย.\csromanspacer{} สญฺโญชนสมฺปยุตฺตา ขนฺธา วุฏฺฐานสฺส อนนฺตรปจฺจเยน ปจฺจโย. (2)}

\prose{สญฺโญชนสมฺปยุตฺโต ธมฺโม สญฺโญชนสมฺปยุตฺตสฺส จ สญฺโญชนวิปฺปยุตฺตสฺส จ ธมฺมสฺส อนนฺตรปจฺจเยน ปจฺจโย.\csromanspacer{} ปุริมา ปุริมา}

\vspace{\tipitakaparskip}
\csromancenter{สญฺโญชนสมฺปยุตฺตทุก อุทฺธจฺจสหคตา ขนฺธา ปจฺฉิมานํ ปจฺฉิมานํ อุทฺธจฺจสหคตานํ ขนฺธานํ โมหสฺส จ อนนฺตรปจฺจเยน ปจฺจโย. (3)}
\proseitem{74}{\csromanfolio{235}สญฺโญชนวิปฺปยุตฺโต ธมฺโม สญฺโญชนวิปฺปยุตฺตสฺส ธมฺมสฺส อนนฺตรปจฺจเยน ปจฺจโย.\csromanspacer{} ปุริโม ปุริโม อุทฺธจฺจสหคโต โมโห ปจฺฉิมสฺส ปจฺฉิมสฺส อุทฺธจฺจสหคตสฺส โมหสฺส อนนฺตรปจฺจเยน ปจฺจโย.\csromanspacer{} ปุริมา ปุริมา สญฺโญชนวิปฺปยุตฺตา ขนฺธา ปจฺฉิมานํ ฯเปฯ ผลสมาปตฺติยา อนนฺตรปจฺจเยน ปจฺจโย. (1)}

\prose{สญฺโญชนวิปฺปยุตฺโต ธมฺโม สญฺโญชนสมฺปยุตฺตสฺส ธมฺมสฺส อนนฺตรปจฺจเยน ปจฺจโย.\csromanspacer{} ปุริโม ปุริโม อุทฺธจฺจสหคโต โมโห ปจฺฉิมานํ ปจฺฉิมานํ อุทฺธจฺจสหคตานํ ขนฺธานํ อนนฺตรปจฺจเยน ปจฺจโย.\csromanspacer{} อาวชฺชนา สญฺโญชนสมฺปยุตฺตกานํ ขนฺธานํ อนนฺตรปจฺจเยน ปจฺจโย. (2)}

\prose{สญฺโญชนวิปฺปยุตฺโต ธมฺโม สญฺโญชนสมฺปยุตฺตสฺส จ สญฺโญชนวิปฺปยุตฺตสฺส จ ธมฺมสฺส อนนฺตรปจฺจเยน ปจฺจโย.\csromanspacer{} ปุริโม ปุริโม อุทฺธจฺจสหคโต โมโห ปจฺฉิมานํ ปจฺฉิมานํ อุทฺธจฺจสหคตานํ ขนฺธานํ โมหสฺส จ อนนฺตรปจฺจเยน ปจฺจโย.\csromanspacer{} อาวชฺชนา อุทฺธจฺจสหคตานํ ขนฺธานํ โมหสฺส จ อนนฺตรปจฺจเยน ปจฺจโย. (3)}

\proseitem{75}{สญฺโญชนสมฺปยุตฺโต จ สญฺโญชนวิปฺปยุตฺโต จ ธมฺมา สญฺโญชนสมฺปยุตฺตสฺส ธมฺมสฺส อนนฺตรปจฺจเยน ปจฺจโย.\csromanspacer{} ปุริมา ปุริมา อุทฺธจฺจสหคตา ขนฺธา จ โมโห จ ปจฺฉิมานํ ปจฺฉิมานํ อุทฺธจฺจสหคตานํ ขนฺธานํ อนนฺตรปจฺจเยน ปจฺจโย. (1)}

\prose{สญฺโญชนสมฺปยุตฺโต จ สญฺโญชนวิปฺปยุตฺโต จ ธมฺมา สญฺโญชนวิปฺปยุตฺตสฺส ธมฺมสฺส อนนฺตรปจฺจเยน ปจฺจโย.\csromanspacer{} ปุริมา ปุริมา อุทฺธจฺจสหคตา ขนฺธา จ โมโห จ ปจฺฉิมสฺส ปจฺฉิมสฺส อุทฺธจฺจสหคตสฺส โมหสฺส อนนฺตรปจฺจเยน ปจฺจโย.\csromanspacer{} อุทฺธจฺจสหคตา ขนฺธา จ โมโห จ วุฏฺฐานสฺส อนนฺตรปจฺจเยน ปจฺจโย. (2)}

\prose{สญฺโญชนสมฺปยุตฺโต จ สญฺโญชนวิปฺปยุตฺโต จ ธมฺมา สญฺโญชนสมฺปยุตฺตสฺส จ สญฺโญชนวิปฺปยุตฺตสฺส จ ธมฺมสฺส อนนฺตรปจฺจเยน ปจฺจโย.\csromanspacer{} ปุริมา ปุริมา อุทฺธจฺจสหคตา ขนฺธา จ โมโห จ ปจฺฉิมานํ \csromanfolio{236}ปจฺฉิมานํ อุทฺธจฺจสหคตานํ ขนฺธานํ โมหสฺส จ อนนฺตรปจฺจเยน ปจฺจโย. (3)}

\csromanheader{2}{\textbf{สมนนฺตราทิ}}
\proseitem{76}{สญฺโญชนสมฺปยุตฺโต ธมฺโม สญฺโญชนสมฺปยุตฺตสฺส ธมฺมสฺส สมนนฺตรปจฺจเยน ปจฺจโย.\csromanspacer{} นว.\csromanspacer{} สหชาตปจฺจเยน ปจฺจโย.\csromanspacer{} นว.\csromanspacer{} อญฺญมญฺญปจฺจเยน ปจฺจโย.\csromanspacer{} ฉ.\csromanspacer{} นิสฺสยปจฺจเยน ปจฺจโย.\csromanspacer{} นว.}

\csromanheader{2}{\textbf{อุปนิสฺสย}}
\proseitem{77}{สญฺโญชนสมฺปยุตฺโต ธมฺโม สญฺโญชนสมฺปยุตฺตสฺส ธมฺมสฺส อุปนิสฺสยปจฺจเยน ปจฺจโย.\csromanspacer{} \textbf{อารมฺมณูปนิสฺสโย, อนนฺตรูปนิสฺสโย,} \textbf{ปกตูปนิสฺสโย ฯเปฯ.}\csromanspacer{} ปกตูปนิสฺสโย\csromandash{}สญฺโญชนสมฺปยุตฺตา ขนฺธา สญฺโญชนสมฺปยุตฺตกานํ ขนฺธานํ อุปนิสฺสยปจฺจเยน ปจฺจโย. (1)}

\prose{สญฺโญชนสมฺปยุตฺโต ธมฺโม สญฺโญชนวิปฺปยุตฺตสฺส ธมฺมสฺส อุปนิสฺสยปจฺจเยน ปจฺจโย.\csromanspacer{} \textbf{อนนฺตรูปนิสฺสโย, ปกตูปนิสฺสโย ฯเปฯ.}\csromanspacer{} ปกตูปนิสฺสโย\csromandash{}สญฺโญชนสมฺปยุตฺตา ขนฺธา สญฺโญชนวิปฺปยุตฺตานํ ขนฺธานํ โมหสฺส จ อุปนิสฺสยปจฺจเยน ปจฺจโย. (2)}

\prose{สญฺโญชนสมฺปยุตฺโต ธมฺโม สญฺโญชนสมฺปยุตฺตสฺส จ สญฺโญชนวิปฺปยุตฺตสฺส จ ธมฺมสฺส อุปนิสฺสยปจฺจเยน ปจฺจโย.\csromanspacer{} \textbf{อนนฺตรูปนิสฺสโย.}\csromanspacer{}\textbf{ ปกตูปนิสฺสโย ฯเปฯ.}\csromanspacer{} ปกตูปนิสฺสโย\csromandash{} สญฺโญชนสมฺปยุตฺตา ขนฺธา อุทฺธจฺจสหคตานํ ขนฺธานํ โมหสฺส จ อุปนิสฺสยปจฺจเยน ปจฺจโย. (3)}

\proseitem{78}{สญฺโญชนวิปฺปยุตฺโต ธมฺโม สญฺโญชนวิปฺปยุตฺตสฺส ธมฺมสฺส อุปนิสฺสยปจฺจเยน ปจฺจโย.\csromanspacer{} \textbf{อารมฺมณูปนิสฺสโย, อนนฺตรูปนิสฺสโย,} \textbf{ปกตูปนิสฺสโย ฯเปฯ.}\csromanspacer{} ปกตูปนิสฺสโย\csromandash{}สทฺธํ อุปนิสฺสาย ทานํ เทติ ฯเปฯ สมาปตฺติํ อุปฺปาเทติ.\csromanspacer{} สีลํ ฯเปฯ.\csromanspacer{} ปญฺญํ.\csromanspacer{} กายิกํ สุขํ.\csromanspacer{} กายิกํ ทุกฺขํ.\csromanspacer{} อุตุํ.\csromanspacer{} โภชนํ.\csromanspacer{} เสนาสนํ โมหํ อุปนิสฺสาย ทานํ เทติ ฯเปฯ สมาปตฺติํ อุปฺปาเทติ.\csromanspacer{} สทฺธา ฯเปฯ.\csromanspacer{} เสนาสนํ โมโห จ สทฺธาย ฯเปฯ ผลสมาปตฺติยา โมหสฺส จ อุปนิสฺสยปจฺจเยน ปจฺจโย. (1)}

\csromanheader{2}{22. สญฺโญชนสมฺปยุตฺตทุก}
\prose{\csromanfolio{237}สญฺโญชนวิปฺปยุตฺโต ธมฺโม สญฺโญชนสมฺปยุตฺตสฺส ธมฺมสฺส อุปนิสฺสยปจฺจเยน ปจฺจโย.\csromanspacer{} \textbf{อารมฺมณูปนิสฺสโย, อนนฺตรูปนิสฺสโย, ปกตูปนิสฺสโย ฯเปฯ.}\csromanspacer{} ปกตูปนิสฺสโย\csromandash{}สทฺธํ อุปนิสฺสาย มานํ ชปฺเปติ.\csromanspacer{} ทิฏฺฐิํ คณฺหาติ.\csromanspacer{} สีลํ ฯเปฯ.\csromanspacer{} ปญฺญํ.\csromanspacer{} กายิกํ สุขํ.\csromanspacer{} กายิกํ ทุกฺขํ.\csromanspacer{} อุตุํ.\csromanspacer{} โภชนํ.\csromanspacer{} เสนาสนํ โมหํ อุปนิสฺสาย ปาณํ หนติ ฯเปฯ สํฆํ ภินฺทติ.\csromanspacer{} สทฺธา ฯเปฯ.\csromanspacer{} เสนาสนํ โมโห ราคสฺส ฯเปฯ ปตฺถนาย อุปนิสฺสยปจฺจเยน ปจฺจโย. (2)}

\prose{สญฺโญชนวิปฺปยุตฺโต ธมฺโม สญฺโญชนสมฺปยุตฺตสฺส จ สญฺโญชนวิปฺปยุตฺตสฺส จ ธมฺมสฺส อุปนิสฺสยปจฺจเยน ปจฺจโย.\csromanspacer{} \textbf{อนนฺตรูปนิสฺสโย, ปกตูปนิสฺสโย ฯเปฯ} ปกตูปนิสฺสโย\csromandash{}สทฺธา ฯเปฯ ปญฺญา.\csromanspacer{} กายิกํ สุขํ ฯเปฯ.\csromanspacer{} เสนาสนํ โมโห จ อุทฺธจฺจสหคตานํ ขนฺธานํ โมหสฺส จ อุปนิสฺสยปจฺจเยน ปจฺจโย. (3)}

\proseitem{79}{สญฺโญชนสมฺปยุตฺโต จ สญฺโญชนวิปฺปยุตฺโต จ ธมฺมา สญฺโญชนสมฺปยุตฺตสฺส ธมฺมสฺส อุปนิสฺสยปจฺจเยน ปจฺจโย.\csromanspacer{} \textbf{อนนฺตรูปนิสฺสโย, ปกตูปนิสฺสโย ฯเปฯ.}\csromanspacer{} ปกตูปนิสฺสโย\csromandash{} อุทฺธจฺจสหคตา ขนฺธา จ โมโห จ สญฺโญชนสมฺปยุตฺตกานํ ขนฺธานํ อุปนิสฺสยปจฺจเยน ปจฺจโย. (1)}

\prose{สญฺโญชนสมฺปยุตฺโต จ สญฺโญชนวิปฺปยุตฺโต จ ธมฺมา สญฺโญชนวิปฺปยุตฺตสฺส ธมฺมสฺส อุปนิสฺสยปจฺจเยน ปจฺจโย.\csromanspacer{} \textbf{อนนฺตรูปนิสฺสโย, ปกตูปนิสฺสโย ฯเปฯ.}\csromanspacer{} ปกตูปนิสฺสโย\csromandash{} อุทฺธจฺจสหคตา ขนฺธา จ โมโห จ สญฺโญชนวิปฺปยุตฺตานํ ขนฺธานํ โมหสฺส จ อุปนิสฺสยปจฺจเยน ปจฺจโย. (2)}

\prose{สญฺโญชนสมฺปยุตฺโต จ สญฺโญชนวิปฺปยุตฺโต จ ธมฺมา สญฺโญชนสมฺปยุตฺตสฺส จ สญฺโญชนวิปฺปยุตฺตสฺส จ ธมฺมสฺส อุปนิสฺสยปจฺจเยน ปจฺจโย.\csromanspacer{} \textbf{อนนฺตรูปนิสฺสโย, ปกตูปนิสฺสโย ฯเปฯ.}\csromanspacer{} ปกตูปนิสฺสโย\csromandash{}อุทฺธจฺจสหคตา ขนฺธา จ โมโห จ อุทฺธจฺจสหคตานํ ขนฺธานํ โมหสฺส จ อุปนิสฺสยปจฺจเยน ปจฺจโย. (3)}

\csromanheader{2}{\textbf{ปุเรชาต}}
\proseitem{80}{\csromanfolio{238}สญฺโญชนวิปฺปยุตฺโต ธมฺโม สญฺโญชนวิปฺปยุตฺตสฺส ธมฺมสฺส ปุเรชาตปจฺจเยน ปจฺจโย.\csromanspacer{} \textbf{อารมฺมณปุเรชาตํ}, วตฺถุปุเรชาตํ.\csromanspacer{}\csromanspacer{}\textbf{อารมฺมณปุเรชาตํ}\csromandash{}จกฺขุํ ฯเปฯ วตฺถุํ อนิจฺจโต ฯเปฯ วิปสฺสติ.\csromanspacer{} ทิพฺเพน จกฺขุนา รูปํ ปสฺสติ.\csromanspacer{} ทิพฺพาย โสตธาตุยา สทฺทํ สุณาติ.\csromanspacer{} รูปายตนํ จกฺขุวิญฺญาณสฺส ฯเปฯ.\csromanspacer{} โผฏฺฐพฺพายตนํ กายวิญฺญาณสฺส ฯเปฯ.\csromanspacer{} \textbf{วตฺถุปุเรชาตํ}\csromandash{}จกฺขายตนํ จกฺขุวิญฺญาณสฺส ฯเปฯ กายายตนํ กายวิญฺญาณสฺส ฯเปฯ.\csromanspacer{} วตฺถุ สญฺโญชนวิปฺปยุตฺตานํ ขนฺธานํ โมหสฺส จ ปุเรชาตปจฺจเยน ปจฺจโย. (1)}

\prose{สญฺโญชนวิปฺปยุตฺโต ธมฺโม สญฺโญชนสมฺปยุตฺตสฺส ธมฺมสฺส ปุเรชาตปจฺจเยน ปจฺจโย.\csromanspacer{} \textbf{อารมฺมณปุเรชาตํ}, วตฺถุปุเรชาตํ.\csromanspacer{}\csromanspacer{}\textbf{อารมฺมณปุเรชาตํ}\csromandash{}จกฺขุํ ฯเปฯ วตฺถุํ อสฺสาเทติ อภินนฺทติ, ตํ อารพฺภ ราโค อุปฺปชฺชติ ฯเปฯ โทมนสฺสํ อุปฺปชฺชติ.\csromanspacer{} \textbf{วตฺถุปุเรชาตํ}\csromandash{}วตฺถุ สญฺโญชนสมฺปยุตฺตกานํ ขนฺธานํ ปุเรชาตปจฺจเยน ปจฺจโย. (2)}

\prose{สญฺโญชนวิปฺปยุตฺโต ธมฺโม สญฺโญชนสมฺปยุตฺตสฺส จ สญฺโญชนวิปฺปยุตฺตสฺส จ ธมฺมสฺส ปุเรชาตปจฺจเยน ปจฺจโย.\csromanspacer{} \textbf{อารมฺมณปุเรชาตํ}, วตฺถุปุเรชาตํ.\csromanspacer{}\csromanspacer{}\textbf{อารมฺมณปุเรชาตํ}\csromandash{}จกฺขุํ ฯเปฯ วตฺถุํ อารพฺภ อุทฺธจฺจสหคตา ขนฺธา จ โมโห จ อุปฺปชฺชนฺติ.\csromanspacer{} \textbf{วตฺถุปุเรชาตํ}\csromandash{}วตฺถุ อุทฺธจฺจสหคตานํ ขนฺธานํ โมหสฺส จ ปุเรชาตปจฺจเยน ปจฺจโย. (3)}

\csromanheader{2}{\textbf{ปจฺฉาชาต}}
\proseitem{81}{สญฺโญชนสมฺปยุตฺโต ธมฺโม สญฺโญชนวิปฺปยุตฺตสฺส ธมฺมสฺส ปจฺฉาชาตปจฺจเยน ปจฺจโย.\csromanspacer{} ปจฺฉาชาตา สญฺโญชนสมฺปยุตฺตา ขนฺธา ปุเรชาตสฺส อิมสฺส กายสฺส ปจฺฉาชาตปจฺจเยน ปจฺจโย. (1)}

\prose{สญฺโญชนวิปฺปยุตฺโต ธมฺโม สญฺโญชนวิปฺปยุตฺตสฺส ธมฺมสฺส ปจฺฉาชาตปจฺจเยน ปจฺจโย.\csromanspacer{} ปจฺฉาชาตา สญฺโญชนวิปฺปยุตฺตา ขนฺธา จ โมโห จ ปุเรชาตสฺส อิมสฺส กายสฺส ปจฺฉาชาตปจฺจเยน ปจฺจโย. (1)}

\csromanheader{2}{22. สญฺโญชนสมฺปยุตฺตทุก}
\prose{\csromanfolio{239}สญฺโญชนสมฺปยุตฺโต จ สญฺโญชนวิปฺปยุตฺโต จ ธมฺมา สญฺโญชนวิปฺปยุตฺตสฺส ธมฺมสฺส ปจฺฉาชาตปจฺจเยน ปจฺจโย.\csromanspacer{} ปจฺฉาชาตา อุทฺธจฺจสหคตา ขนฺธา จ โมโห จ ปุเรชาตสฺส อิมสฺส กายสฺส ปจฺฉาชาตปจฺจเยน ปจฺจโย. (1)}

\csromanheader{2}{\textbf{อาเสวน}}
\proseitem{82}{สญฺโญชนสมฺปยุตฺโต ธมฺโม สญฺโญชนสมฺปยุตฺตสฺส ธมฺมสฺส อาเสวนปจฺจเยน ปจฺจโย.\csromanspacer{} นว. (อาวชฺชนาปิ วุฏฺฐานมฺปิ นตฺถิ.)}

\csromanheader{2}{\textbf{กมฺม}}
\proseitem{83}{สญฺโญชนสมฺปยุตฺโต ธมฺโม สญฺโญชนสมฺปยุตฺตสฺส ธมฺมสฺส กมฺมปจฺจเยน ปจฺจโย.\csromanspacer{} สญฺโญชนสมฺปยุตฺตา เจตนา สญฺโญชนสมฺปยุตฺตกานํ ขนฺธานํ กมฺมปจฺจเยน ปจฺจโย. (1)}

\prose{สญฺโญชนสมฺปยุตฺโต ธมฺโม สญฺโญชนวิปฺปยุตฺตสฺส ธมฺมสฺส กมฺมปจฺจเยน ปจฺจโย.\csromanspacer{} \textbf{สหชาตา}, นานากฺขณิกา.\csromanspacer{}\csromanspacer{}\textbf{สหชาตา} สญฺโญชนสมฺปยุตฺตา เจตนา จิตฺตสมุฏฺฐานานํ รูปานํ กมฺมปจฺจเยน ปจฺจโย.\csromanspacer{} อุทฺธจฺจสหคตา เจตนา โมหสฺส จิตฺตสมุฏฺฐานานญฺจ รูปานํ กมฺมปจฺจเยน ปจฺจโย.\csromanspacer{} \textbf{นานากฺขณิกา} สญฺโญชนสมฺปยุตฺตกา เจตนา วิปากานํ ขนฺธานํ กฏตฺตา จ รูปานํ กมฺมปจฺจเยน ปจฺจโย. (2)}

\prose{สญฺโญชนสมฺปยุตฺโต ธมฺโม สญฺโญชนสมฺปยุตฺตสฺส จ สญฺโญชนวิปฺปยุตฺตสฺส จ ธมฺมสฺส กมฺมปจฺจเยน ปจฺจโย.\csromanspacer{} สญฺโญชนสมฺปยุตฺตา เจตนา สมฺปยุตฺตกานํ ขนฺธานํ จิตฺตสมุฏฺฐานานญฺจ รูปานํ กมฺมปจฺจเยน ปจฺจโย.\csromanspacer{} อุทฺธจฺจสหคตา เจตนา สมฺปยุตฺตกานํ ขนฺธานํ โมหสฺส จ จิตฺตสมุฏฺฐานานญฺจ รูปานํ กมฺมปจฺจเยน ปจฺจโย. (3)}

\prose{สญฺโญชนวิปฺปยุตฺโต ธมฺโม สญฺโญชนวิปฺปยุตฺตสฺส ธมฺมสฺส กมฺมปจฺจเยน ปจฺจโย.\csromanspacer{} \textbf{สหชาตา}, นานากฺขณิกา.\csromanspacer{}\csromanspacer{}\textbf{สหชาตา} สญฺโญชนวิปฺปยุตฺตา เจตนา สมฺปยุตฺตกานํ ขนฺธานํ จิตฺตสมุฏฺฐานานญฺจ รูปานํ กมฺมปจฺจเยน ปจฺจโย.\csromanspacer{} ปฏิสนฺธิกฺขเณ ฯเปฯ.\csromanspacer{} \textbf{นานากฺขณิกา} สญฺโญชนวิปฺปยุตฺตา เจตนา วิปากานํ ขนฺธานํ กฏตฺตา จ รูปานํ กมฺมปจฺจเยน ปจฺจโย.}

\csromanheader{2}{\textbf{วิปาก}}
\proseitem{84}{\csromanfolio{240}สญฺโญชนวิปฺปยุตฺโต ธมฺโม สญฺโญชนวิปฺปยุตฺตสฺส ธมฺมสฺส วิปากปจฺจเยน ปจฺจโย.\csromanspacer{} เอกํ.}

\csromanheader{2}{\textbf{อาหาราทิ}}
\proseitem{85}{สญฺโญชนสมฺปยุตฺโต ธมฺโม สญฺโญชนสมฺปยุตฺตสฺส ธมฺมสฺส อาหารปจฺจเยน ปจฺจโย.\csromanspacer{} จตฺตาริ.\csromanspacer{} อินฺทฺริยปจฺจเยน ปจฺจโย.\csromanspacer{} จตฺตาริ.\csromanspacer{} ฌานปจฺจเยน ปจฺจโย.\csromanspacer{} จตฺตาริ.\csromanspacer{} มคฺคปจฺจเยน ปจฺจโย.\csromanspacer{} จตฺตาริ.\csromanspacer{} สมฺปยุตฺตปจฺจเยน ปจฺจโย.\csromanspacer{} ฉ.}

\csromanheader{2}{\textbf{วิปฺปยุตฺต}}
\proseitem{86}{สญฺโญชนสมฺปยุตฺโต ธมฺโม สญฺโญชนวิปฺปยุตฺตสฺส ธมฺมสฺส วิปฺปยุตฺตปจฺจเยน ปจฺจโย.\csromanspacer{} สหชาตํ, ปจฺฉาชาตํ. (สํขิตฺตํ.) (1)}

\prose{สญฺโญชนวิปฺปยุตฺโต ธมฺโม สญฺโญชนวิปฺปยุตฺตสฺส ธมฺมสฺส วิปฺปยุตฺตปจฺจเยน ปจฺจโย.\csromanspacer{} สหชาตํ, ปุเรชาตํ, ปจฺฉาชาตํ. (สํขิตฺตํ.) (1)}

\prose{สญฺโญชนวิปฺปยุตฺโต ธมฺโม สญฺโญชนสมฺปยุตฺตสฺส ธมฺมสฺส วิปฺปยุตฺตปจฺจเยน ปจฺจโย.\csromanspacer{} \textbf{ปุเรชาตํ} วตฺถุ สญฺโญชนสมฺปยุตฺตกานํ ขนฺธานํ วิปฺปยุตฺตปจฺจเยน ปจฺจโย. (2)}

\prose{สญฺโญชนวิปฺปยุตฺโต ธมฺโม สญฺโญชนสมฺปยุตฺตสฺส จ สญฺโญชนวิปฺปยุตฺตสฺส จ ธมฺมสฺส วิปฺปยุตฺตปจฺจเยน ปจฺจโย.\csromanspacer{} \textbf{ปุเรชาตํ} วตฺถุ อุทฺธจฺจสหคตานํ ขนฺธานํ โมหสฺส จ วิปฺปยุตฺตปจฺจเยน ปจฺจโย. (3)}

\prose{สญฺโญชนสมฺปยุตฺโต จ สญฺโญชนวิปฺปยุตฺโต จ ธมฺมา สญฺโญชนวิปฺปยุตฺตสฺส ธมฺมสฺส วิปฺปยุตฺตปจฺจเยน ปจฺจโย.\csromanspacer{} สหชาตํ, ปจฺฉาชาตํ. (สํขิตฺตํ.) (1)}

\csromanheader{2}{\textbf{อตฺถิ-อาทิ}}
\proseitem{87}{สญฺโญชนสมฺปยุตฺโต ธมฺโม สญฺโญชนสมฺปยุตฺตสฺส ธมฺมสฺส อตฺถิปจฺจเยน ปจฺจโย.\csromanspacer{} เอกํ. (ปฏิจฺจวารสทิสํ.) (1)}

\csromanheader{2}{22. สญฺโญชนสมฺปยุตฺตทุก}
\prose{\csromanfolio{241}สญฺโญชนสมฺปยุตฺโต ธมฺโม สญฺโญชนวิปฺปยุตฺตสฺส ธมฺมสฺส อตฺถิปจฺจเยน ปจฺจโย.\csromanspacer{} สหชาตํ, ปจฺฉาชาตํ.\csromanspacer{}\csromanspacer{}\textbf{สหชาตา} สญฺโญชนสมฺปยุตฺตา ขนฺธา จิตฺตสมุฏฺฐานานํ รูปานํ อตฺถิปจฺจเยน ปจฺจโย.\csromanspacer{} อุทฺธจฺจสหคตา ขนฺธา โมหสฺส จิตฺตสมุฏฺฐานานญฺจ รูปานํ อตฺถิปจฺจเยน ปจฺจโย.\csromanspacer{} \textbf{ปจฺฉาชาตา} สญฺโญชนสมฺปยุตฺตา ขนฺธา ปุเรชาตสฺส อิมสฺส กายสฺส อตฺถิปจฺจเยน ปจฺจโย. (2)}

\prose{สญฺโญชนสมฺปยุตฺโต ธมฺโม สญฺโญชนสมฺปยุตฺตสฺส จ สญฺโญชนวิปฺปยุตฺตสฺส จ ธมฺมสฺส อตฺถิปจฺจเยน ปจฺจโย.\csromanspacer{} สญฺโญชนสมฺปยุตฺโต เอโก ขนฺโธ ติณฺณินฺนํ ขนฺธานํ จิตฺตสมุฏฺฐานานญฺจ รูปานํ อตฺถิปจฺจเยน ปจฺจโย ฯเปฯ เทฺว ขนฺธา ฯเปฯ.\csromanspacer{} อุทฺธจฺจสหคโต เอโก ขนฺโธ ติณฺณนฺนํ ขนฺธานํ โมหสฺส จ จิตฺตสมุฏฺฐานานญฺจ รูปานํ อตฺถิปจฺจเยน ปจฺจโย ฯเปฯ. (3)}

\proseitem{88}{สญฺโญชนวิปฺปยุตฺโต ธมฺโม สญฺโญชนวิปฺปยุตฺตสฺส ธมฺมสฺส อตฺถิปจฺจเยน ปจฺจโย.\csromanspacer{} สหชาตํ, ปุเรชาตํ, ปจฺฉาชาตํ, อาหารํ อินฺทฺริยํ. (สํขิตฺตํ.\csromanspacer{} วิตฺถาเรตพฺพํ.) (1)}

\prose{สญฺโญชนวิปฺปยุตฺโต ธมฺโม สญฺโญชนสมฺปยุตฺตสฺส ธมฺมสฺส อตฺถิปจฺจเยน ปจฺจโย.\csromanspacer{} สหชาตํ, ปุเรชาตํ.\csromanspacer{}\csromanspacer{}\textbf{สหชาโต} อุทฺธจฺจสหคโต โมโห สมฺปยุตฺตกานํ ขนฺธานํ อตฺถิปจฺจเยน ปจฺจโย.\csromanspacer{} \textbf{ปุเรชาตํ} จกฺขุํ ฯเปฯ วตฺถุํ อสฺสาเทติ อภินนฺทติ, ตํ อารพฺภ ราโค อุปฺปชฺชติ ฯเปฯ โทมนสฺสํ อุปฺปชฺชติ.\csromanspacer{} วตฺถุ สญฺโญชนสมฺปยุตฺตกานํ ขนฺธานํ อตฺถิปจฺจเยน ปจฺจโย. (2)}

\prose{สญฺโญชนวิปฺปยุตฺโต ธมฺโม สญฺโญชนสมฺปยุตฺตสฺส จ สญฺโญชนวิปฺปยุตฺตสฺส จ ธมฺมสฺส อตฺถิปจฺจเยน ปจฺจโย.\csromanspacer{} สหชาตํ, ปุเรชาตํ.\csromanspacer{}\csromanspacer{}\textbf{สหชาโต} อุทฺธจฺจสหคโต โมโห สมฺปยุตฺตกานํ ขนฺธานํ จิตฺตสมุฏฺฐานานญฺจ รูปานํ อตฺถิปจฺจเยน ปจฺจโย.\csromanspacer{} \textbf{ปุเรชาตํ} จกฺขุํ ฯเปฯ วตฺถุํ อารพฺภ อุทฺธจฺจสหคตา ขนฺธา จ โมโห จ อุปฺปชฺชนฺติ.\csromanspacer{} วตฺถุ อุทฺธจฺจสหคตานํ ขนฺธานํ โมหสฺส จ อตฺถิปจฺจเยน ปจฺจโย. (3)}

\proseitem{89}{สญฺโญชนสมฺปยุตฺโต จ สญฺโญชนวิปฺปยุตฺโต จ ธมฺมา สญฺโญชนสมฺปยุตฺตสฺส ธมฺมสฺส อตฺถิปจฺจเยน ปจฺจโย.\csromanspacer{} สหชาตํ, ปุเรชาตํ.\csromanspacer{}\csromanspacer{}\textbf{สหชาโต} สญฺโญชนสมฺปยุตฺโต เอโก ขนฺโธ จ \csromanfolio{242}วตฺถุ จ ติณฺณนฺนํ ขนฺธานํ อตฺถิปจฺจเยน ปจฺจโย ฯเปฯ เทฺว ขนฺธา จ ฯเปฯ.\csromanspacer{} อุทฺธจฺจสหคโต เอโก ขนฺโธ จ โมโห จ ติณฺณนฺนํ ขนฺธานํ อตฺถิปจฺจเยน ปจฺจโย ฯเปฯ เทฺว ขนฺธา จ ฯเปฯ. (1)}

\prose{สญฺโญชนสมฺปยุตฺโต จ สญฺโญชนวิปฺปยุตฺโต จ ธมฺมา สญฺโญชนวิปฺปยุตฺตสฺส ธมฺมสฺส อตฺถิปจฺจเยน ปจฺจโย.\csromanspacer{} \textbf{สหชาตํ, ปุเรชาตํ,} ปจฺฉาชาตํ, อาหารํ อินฺทฺริยํ.\csromanspacer{}\csromanspacer{}\textbf{สหชาตา} สญฺโญชนสมฺปยุตฺตา ขนฺธา จ โมโห จ มหาภูตา จ จิตฺตสมุฏฺฐานานํ รูปานํ อตฺถิปจฺจเยน ปจฺจโย.\csromanspacer{} อุทฺธจฺจสหคตา ขนฺธา จ โมโห จ จิตฺตสมุฏฺฐานานํ รูปานํ อตฺถิปจฺจเยน ปจฺจโย.\csromanspacer{} อุทฺธจฺจสหคตา ขนฺธา จ วตฺถุ จ โมหสฺส อตฺถิปจฺจเยน ปจฺจโย.\csromanspacer{} \textbf{ปจฺฉาชาตา} อุทฺธจฺจสหคตา ขนฺธา จ โมโห จ ปุเรชาตสฺส อิมสฺส กายสฺส อตฺถิปจฺจเยน ปจฺจโย.\csromanspacer{} \textbf{ปจฺฉาชาตา} สญฺโญชนสมฺปยุตฺตา ขนฺธา จ กพฬีกาโร อาหาโร จ อิมสฺส กายสฺส อตฺถิปจฺจเยน ปจฺจโย.\csromanspacer{} \textbf{ปจฺฉาชาตา} สญฺโญชนสมฺปยุตฺตา ขนฺธา จ รูปชีวิตินฺทฺริยญฺจ กฏตฺตารูปานํ อตฺถิปจฺจเยน ปจฺจโย. (2)}

\prose{สญฺโญชนสมฺปยุตฺโต จ สญฺโญชนวิปฺปยุตฺโต จ ธมฺมา สญฺโญชนสมฺปยุตฺตสฺส จ สญฺโญชนวิปฺปยุตฺตสฺส จ ธมฺมสฺส อตฺถิปจฺจเยน ปจฺจโย.\csromanspacer{} สหชาตํ, ปุเรชาตํ.\csromanspacer{}\csromanspacer{}\textbf{สหชาโต} อุทฺธจฺจสหคโต เอโก ขนฺโธ จ โมโห จ ติณฺณนฺนํ ขนฺธานํ จิตฺตสมุฏฺฐานานญฺจ รูปานํ อตฺถิปจฺจเยน ปจฺจโย ฯเปฯ เทฺว ขนฺธา จ ฯเปฯ.\csromanspacer{} \textbf{สหชาโต} อุทฺธจฺจสหคโต เอโก ขนฺโธ จ วตฺถุ จ ติณฺณนฺนํ ขนฺธานํ โมหสฺส จ อตฺถิปจฺจเยน ปจฺจโย ฯเปฯ เทฺว ขนฺธา ฯเปฯ. (3)}

\prose{นตฺถิปจฺจเยน ปจฺจโย.\csromanspacer{} วิคตปจฺจเยน ปจฺจโย.\csromanspacer{} อวิคตปจฺจเยน ปจฺจโย.}

\csromanheader{2}{1. ปจฺจยานุโลม 2. สงฺขฺยาวาร}
\csromanheader{2}{\textbf{สุทฺธ}}
\proseitem{90}{เหตุยา ฉ, อารมฺมเณ นว, อธิปติยา ปญฺจ, อนนฺตเร นว, สมนนฺตเร นว, สหชาเต นว, อญฺญมญฺเญ ฉ, นิสฺสเย นว, อุปนิสฺสเย นว, ปุเรชาเต ตีณิ, ปจฺฉาชาเต ตีณิ, อาเสวเน นว, กมฺเม จตฺตาริ, วิปาเก เอกํ, อาหาเร จตฺตาริ, อินฺทฺริเย จตฺตาริ, ฌาเน จตฺตาริ,}

\vspace{\tipitakaparskip}
\csromancenter{สญฺโญชนสมฺปยุตฺตทุก มคฺเค จตฺตาริ, สมฺปยุตฺเต ฉ, วิปฺปยุตฺเต ปญฺจ, อตฺถิยา นว, นตฺถิยา นว, วิคเต นว, อวิคเต นว.}
\csromancenter{อนุโลมํ.}
\csromanheader{2}{2. ปจฺจนียุทฺธาร}
\proseitem{91}{\csromanfolio{243}สญฺโญชนสมฺปยุตฺโต ธมฺโม สญฺโญชนสมฺปยุตฺตสฺส ธมฺมสฺส อารมฺมณปจฺจเยน ปจฺจโย.\csromanspacer{} สหชาตปจฺจเยน ปจฺจโย.\csromanspacer{} อุปนิสฺสยปจฺจเยน ปจฺจโย. (1)}

\prose{สญฺโญชนสมฺปยุตฺโต ธมฺโม สญฺโญชนวิปฺปยุตฺตสฺส ธมฺมสฺส อารมฺมณปจฺจเยน ปจฺจโย.\csromanspacer{} สหชาตปจฺจเยน ปจฺจโย.\csromanspacer{} อุปนิสฺสยปจฺจเยน ปจฺจโย.\csromanspacer{} ปจฺฉาชาตปจฺจเยน ปจฺจโย.\csromanspacer{} กมฺมปจฺจเยน ปจฺจโย. (2)}

\prose{สญฺโญชนสมฺปยุตฺโต ธมฺโม สญฺโญชนสมฺปยุตฺตสฺส จ สญฺโญชนวิปฺปยุตฺตสฺส จ ธมฺมสฺส อารมฺมณปจฺจเยน ปจฺจโย.\csromanspacer{} สหชาตปจฺจเยน ปจฺจโย.\csromanspacer{} อุปนิสฺสยปจฺจเยน ปจฺจโย. (3)}

\proseitem{92}{สญฺโญชนวิปฺปยุตฺโต ธมฺโม สญฺโญชนวิปฺปยุตฺตสฺส ธมฺมสฺส อารมฺมณปจฺจเยน ปจฺจโย.\csromanspacer{} สหชาตปจฺจเยน ปจฺจโย.\csromanspacer{} อุปนิสฺสยปจฺจเยน ปจฺจโย.\csromanspacer{} ปุเรชาตปจฺจเยน ปจฺจโย.\csromanspacer{} ปจฺฉาชาตปจฺจเยน ปจฺจโย.\csromanspacer{} กมฺมปจฺจเยน ปจฺจโย.\csromanspacer{} อาหารปจฺจเยน ปจฺจโย.\csromanspacer{} อินฺทฺริยปจฺจเยน ปจฺจโย. (1)}

\prose{สญฺโญชนวิปฺปยุตฺโต ธมฺโม สญฺโญชนสมฺปยุตฺตสฺส ธมฺมสฺส อารมฺมณปจฺจเยน ปจฺจโย.\csromanspacer{} สหชาตปจฺจเยน ปจฺจโย.\csromanspacer{} อุปนิสฺสยปจฺจเยน ปจฺจโย.\csromanspacer{} ปุเรชาตปจฺจเยน ปจฺจโย. (2)}

\prose{สญฺโญชนวิปฺปยุตฺโต ธมฺโม สญฺโญชนสมฺปยุตฺตสฺส จ สญฺโญชนวิปฺปยุตฺตสฺส จ ธมฺมสฺส อารมฺมณปจฺจเยน ปจฺจโย.\csromanspacer{} สหชาตปจฺจเยน ปจฺจโย.\csromanspacer{} อุปนิสฺสยปจฺจเยน ปจฺจโย.\csromanspacer{} ปุเรชาตปจฺจเยน ปจฺจโย. (3)}

\proseitem{93}{สญฺโญชนสมฺปยุตฺโต จ สญฺโญชนวิปฺปยุตฺโต จ ธมฺมา สญฺโญชนสมฺปยุตฺตสฺส ธมฺมสฺส อารมฺมณปจฺจเยน ปจฺจโย.\csromanspacer{} สหชาตปจฺจเยน ปจฺจโย.\csromanspacer{} อุปนิสฺสยปจฺจเยน ปจฺจโย. (1)}

\prose{\csromanfolio{244}สญฺโญชนสมฺปยุตฺโต จ สญฺโญชนวิปฺปยุตฺโต จ ธมฺมา สญฺโญชนวิปฺปยุตฺตสฺส ธมฺมสฺส อารมฺมณปจฺจเยน ปจฺจโย.\csromanspacer{} สหชาตปจฺจเยน ปจฺจโย.\csromanspacer{} อุปนิสฺสยปจฺจเยน ปจฺจโย.\csromanspacer{} ปจฺฉาชาตปจฺจเยน ปจฺจโย. (2)}

\prose{สญฺโญชนสมฺปยุตฺโต จ สญฺโญชนวิปฺปยุตฺโต จ ธมฺมา สญฺโญชนสมฺปยุตฺตสฺส จ สญฺโญชนวิปฺปยุตฺตสฺส จ ธมฺมสฺส อารมฺมณปจฺจเยน ปจฺจโย.\csromanspacer{} สหชาตปจฺจเยน ปจฺจโย.\csromanspacer{} อุปนิสฺสยปจฺจเยน ปจฺจโย. (3)}

\csromanheader{2}{2. ปจฺจยปจฺจนีย 2. สงฺขฺยาวาร}
\proseitem{94}{นเหตุยา นว, น-อารมฺมเณ นว, (สพฺพตฺถ นว.) โนวิคเต นว, โน-อวิคเต นว.}

\csromanheader{2}{3. ปจฺจยานุโลมปจฺจนีย}
\csromanheader{2}{\textbf{เหตุทุก}}
\proseitem{95}{\textbf{เหตุ}ปจฺจยา น-อารมฺมเณ ฉ ฯเปฯ นสมนนฺตเร ฉ.\csromanspacer{} น-อญฺญมญฺเญ เทฺว, น-อุปนิสฺสเย ฉ ฯเปฯ นมคฺเค ฉ, นสมฺปยุตฺเต เทฺว, นวิปฺปยุตฺเต ตีณิ, โนนตฺถิยา ฉ, โนวิคเต ฉ.}

\csromanheader{2}{4. ปจฺจยปจฺจนียานุโลม}
\proseitem{96}{นเหตุปจฺจยา อารมฺมเณ นว, อธิปติยา ปญฺจ, (อนุโลมปทานิ คณิตพฺพานิ.) อวิคเต นว.}

\nitthitam{สญฺโญชนสมฺปยุตฺตทุกํ นิฏฺฐิตํ.}
\csromanmatikaanchor{mtk.28}
\csromantocmarkref{section}{23. สญฺโญชนสญฺโญชนิยทุก}{mtk.28}
\bookmark[dest={mtk.28},level=1]{23. สญฺโญชนสญฺโญชนิยทุก}
\csromanheader{1}{23. สญฺโญชนสญฺโญชนิยทุก 1. ปฏิจฺจวาร}
\csromanheader{2}{1. ปจฺจยานุโลม 1. วิภงฺควาร}
\csromanheader{2}{\textbf{เหตุ}}
\proseitem{97}{สญฺโญชนญฺเจว สญฺโญชนิยญฺจ ธมฺมํ ปฏิจฺจ สญฺโญชโน เจว สญฺโญชนิโย จ ธมฺโม อุปฺปชฺชติ เหตุปจฺจยา.\csromanspacer{} กามราคสญฺโญชนํ ปฏิจฺจ ทิฏฺฐิสญฺโญชนํ อวิชฺชาสญฺโญชนํ. (จกฺกํ.) (1)}

\csromanheader{2}{23. สญฺโญชนสญฺโญชนิยทุก}
\prose{\csromanfolio{245}สญฺโญชนญฺเจว สญฺโญชนิยญฺจ ธมฺมํ ปฏิจฺจ สญฺโญชนิโย เจว โน จ สญฺโญชโน ธมฺโม อุปฺปชฺชติ เหตุปจฺจยา.\csromanspacer{} สญฺโญชเน ปฏิจฺจ สมฺปยุตฺตกา ขนฺธา จิตฺตสมุฏฺฐานญฺจ รูปํ. (2)}

\prose{สญฺโญชนญฺเจว สญฺโญชนิยญฺจ ธมฺมํ ปฏิจฺจ สญฺโญชโน เจว สญฺโญชนิโย จ สญฺโญชนิโย เจว โน จ สญฺโญชโน จ ธมฺมา อุปฺปชฺชนฺติ เหตุปจฺจยา.\csromanspacer{} กามราคสญฺโญชนํ ปฏิจฺจ ทิฏฺฐิสญฺโญชนํ อวิชฺชาสญฺโญชนํ สมฺปยุตฺตกา ขนฺธา จิตฺตสมุฏฺฐานญฺจ รูปํ. (จกฺกํ.) (3)}

\proseitem{98}{สญฺโญชนิยญฺเจว โน จ สญฺโญชนํ ธมฺมํ ปฏิจฺจ สญฺโญชนิโย เจว โน จ สญฺโญชโน ธมฺโม อุปฺปชฺชติ เหตุปจฺจยา.\csromanspacer{} สญฺโญชนิยญฺเจว โน จ สญฺโญชนํ เอกํ ขนฺธํ ปฏิจฺจ ตโย ขนฺธา จิตฺตสมุฏฺฐานญฺจ รูปํ ฯเปฯ เทฺว ขนฺเธ ฯเปฯ.\csromanspacer{} ปฏิสนฺธิกฺขเณ ฯเปฯ. (ยาว มหาภูตา.) (1)}

\prose{สญฺโญชนิยญฺเจว โน จ สญฺโญชนํ ธมฺมํ ปฏิจฺจ สญฺโญชโน เจว สญฺโญชนิโย จ ธมฺโม อุปฺปชฺชติ เหตุปจฺจยา.\csromanspacer{} สญฺโญชนิเย เจว โน จ สญฺโญชเน ขนฺเธ ปฏิจฺจ สญฺโญชนา. (2)}

\prose{สญฺโญชนิยญฺเจว โน จ สญฺโญชนํ ธมฺมํ ปฏิจฺจ สญฺโญชโน เจว สญฺโญชนิโย จ สญฺโญชนิโย เจว โน จ สญฺโญชโน จ ธมฺมา อุปฺปชฺชนฺติ เหตุปจฺจยา.\csromanspacer{} สญฺโญชนิยญฺเจว โน จ สญฺโญชนํ เอกํ ขนฺธํ ปฏิจฺจ ตโย ขนฺธา สญฺโญชนา จ จิตฺตสมุฏฺฐานญฺจ รูปํ ฯเปฯ เทฺว ขนฺเธ ฯเปฯ. (3)}

\proseitem{99}{สญฺโญชนญฺเจว สญฺโญชนิยญฺจ สญฺโญชนิยญฺเจว โน จ สญฺโญชนํ ธมฺมํ ปฏิจฺจ สญฺโญชโน เจว สญฺโญชนิโย จ ธมฺโม อุปฺปชฺชติ เหตุปจฺจยา.\csromanspacer{} กามราคสญฺโญชนญฺจ สมฺปยุตฺตเก จ ขนฺเธ ปฏิจฺจ ทิฏฺฐิสญฺโญชนํ อวิชฺชาสญฺโญชนํ. (จกฺกํ.) (1)}

\prose{สญฺโญชนญฺเจว สญฺโญชนิยญฺจ สญฺโญชนิยญฺเจว โน จ สญฺโญชนญฺจ ธมฺมํ ปฏิจฺจ สญฺโญชนิโย เจว โน จ สญฺโญชโน ธมฺโม อุปฺปชฺชติ เหตุปจฺจยา.\csromanspacer{} สญฺโญชนิยญฺเจว โน จ สญฺโญชนํ เอกํ ขนฺธญฺจ สญฺโญชเน จ ปฏิจฺจ ตโย ขนฺธา จิตฺตสมุฏฺฐานญฺจ รูปํ ฯเปฯ เทฺว ขนฺเธ ฯเปฯ. (2)}

\prose{สญฺโญชนญฺเจว สญฺโญชนิยญฺจ สญฺโญชนิยญฺเจว โน จ สญฺโญชนญฺจ ธมฺมํ ปฏิจฺจ สญฺโญชโน เจว สญฺโญชนิโย จ สญฺโญชนิโย \csromanfolio{246}เจว โน จ สญฺโญชโน จ ธมฺมา อุปฺปชฺชนฺติ เหตุปจฺจยา.\csromanspacer{} สญฺโญชนิยญฺเจว โน จ สญฺโญชนํ เอกํ ขนฺธญฺจ กามราคสญฺโญชนญฺจ ปฏิจฺจ ตโย ขนฺธา ทิฏฺฐิสญฺโญชนํ อวิชฺชาสญฺโญชนํ จิตฺตสมุฏฺฐานญฺจ รูปํ ฯเปฯ เทฺว ขนฺเธ จ ฯเปฯ. (จกฺกํ พนฺธิตพฺพํ.) (3) (สญฺโญชนโคจฺฉเก ปฐมทุกสทิสํ.)}

\prose{(เอวํ อิมมฺปิ ทุกํ วิตฺถาเรตพฺพํ.\csromanspacer{} นินฺนานากรณํ ฐเปตฺวา โลกุตฺตรํ.) สญฺโญชนสญฺโญชนิยทุกํ นิฏฺฐิตํ.}

\csromanmatikaanchor{mtk.29}
\csromantocmarkref{section}{24. สญฺโญชนสญฺโญชนสมฺปยุตฺตทุก}{mtk.29}
\bookmark[dest={mtk.29},level=1]{24. สญฺโญชนสญฺโญชนสมฺปยุตฺตทุก}
\csromanheader{1}{24. สญฺโญชนสญฺโญชนสมฺปยุตฺตทุก 1. ปฏิจฺจวาร}
\csromanheader{2}{1. ปจฺจยานุโลม 1. วิภงฺควาร}
\csromanheader{2}{\textbf{เหตุ}}
\proseitem{100}{สญฺโญชนญฺเจว สญฺโญชนสมฺปยุตฺตญฺจ ธมฺมํ ปฏิจฺจ สญฺโญชโน เจว สญฺโญชนสมฺปยุตฺโต จ ธมฺโม อุปฺปชฺชติ เหตุปจฺจยา.\csromanspacer{} กามราคสญฺโญชนํ ปฏิจฺจ ทิฏฺฐิสญฺโญชนํ อวิชฺชาสญฺโญชนํ. (จกฺกํ.) (1)}

\prose{สญฺโญชนญฺเจว สญฺโญชนสมฺปยุตฺตญฺจ ธมฺมํ ปฏิจฺจ สญฺโญชนสมฺปยุตฺโต เจว โน จ สญฺโญชโน ธมฺโม อุปฺปชฺชติ เหตุปจฺจยา.\csromanspacer{} สญฺโญชเน ปฏิจฺจ สมฺปยุตฺตกา ขนฺธา. (2)}

\prose{สญฺโญชนญฺเจว สญฺโญชนสมฺปยุตฺตญฺจ ธมฺมํ ปฏิจฺจ สญฺโญชโน เจว สญฺโญชนสมฺปยุตฺโต จ สญฺโญชนสมฺปยุตฺโต เจว โน จ สญฺโญชโน จ ธมฺมา อุปฺปชฺชนฺติ เหตุปจฺจยา.\csromanspacer{} กามราคสญฺโญชนํ ปฏิจฺจ ทิฏฺฐิสญฺโญชนํ อวิชฺชาสญฺโญชนํ สมฺปยุตฺตกา จ ขนฺธา. (จกฺกํ.) (3)}

\proseitem{101}{สญฺโญชนสมฺปยุตฺตญฺเจว โน จ สญฺโญชนํ ธมฺมํ ปฏิจฺจ สญฺโญชนสมฺปยุตฺโต เจว โน จ สญฺโญชโน ธมฺโม อุปฺปชฺชติ เหตุปจฺจยา.\csromanspacer{} สญฺโญชนสมฺปยุตฺตญฺเจว โน จ สญฺโญชนํ เอกํ ขนฺธํ ปฏิจฺจ ตโย ขนฺธา ฯเปฯ เทฺว ขนฺเธ ฯเปฯ. (1)}

\csromanheader{2}{24. สญฺโญชนสญฺโญชนสมฺปยุตฺตทุก}
\prose{\csromanfolio{247}สญฺโญชนสมฺปยุตฺตญฺเจว โน จ สญฺโญชนํ ธมฺมํ ปฏิจฺจ สญฺโญชโน เจว สญฺโญชนสมฺปยุตฺโต จ ธมฺโม อุปฺปชฺชติ เหตุปจฺจยา.\csromanspacer{} สญฺโญชนสมฺปยุตฺเต เจว โน จ สญฺโญชเน ขนฺเธ ปฏิจฺจ สญฺโญชนา. (2)}

\prose{สญฺโญชนสมฺปยุตฺตญฺเจว โน จ สญฺโญชนํ ธมฺมํ ปฏิจฺจ สญฺโญชโน เจว สญฺโญชนสมฺปยุตฺโต จ สญฺโญชนสมฺปยุตฺโต เจว โน จ สญฺโญชโน จ ธมฺมา อุปฺปชฺชนฺติ เหตุปจฺจยา.\csromanspacer{} สญฺโญชนสมฺปยุตฺตญฺเจว โน จ สญฺโญชนํ เอกํ ขนฺธํ ปฏิจฺจ ตโย ขนฺธา สญฺโญชนา จ ฯเปฯ เทฺว ขนฺเธ. (3)}

\proseitem{102}{สญฺโญชนญฺเจว สญฺโญชนสมฺปยุตฺตญฺจ สญฺโญชนสมฺปยุตฺตญฺเจว โน จ สญฺโญชนญฺจ ธมฺมํ ปฏิจฺจ สญฺโญชโน เจว สญฺโญชนสมฺปยุตฺโต จ ธมฺโม อุปฺปชฺชติ เหตุปจฺจยา.\csromanspacer{} กามราคสญฺโญชนญฺจ สมฺปยุตฺตเก จ ขนฺเธ ปฏิจฺจ ทิฏฺฐิสญฺโญชนํ อวิชฺชาสญฺโญชนํ. (จกฺกํ.) (1)}

\prose{สญฺโญชนญฺเจว สญฺโญชนสมฺปยุตฺตญฺจ สญฺโญชนสมฺปยุตฺตญฺเจว โน จ สญฺโญชนญฺจ ธมฺมํ ปฏิจฺจ สญฺโญชนสมฺปยุตฺโต เจว โน จ สญฺโญชโน ธมฺโม อุปฺปชฺชติ เหตุปจฺจยา.\csromanspacer{} สญฺโญชนสมฺปยุตฺตญฺเจว โน จ สญฺโญชนํ เอกํ ขนฺธญฺจ สญฺโญชเน จ ปฏิจฺจ ตโย ขนฺธา ฯเปฯ เทฺว ขนฺเธ จ ฯเปฯ. (2)}

\prose{สญฺโญชนญฺเจว สญฺโญชนสมฺปยุตฺตญฺจ สญฺโญชนสมฺปยุตฺตญฺเจว โน จ สญฺโญชนญฺจ ธมฺมํ ปฏิจฺจ สญฺโญชโน เจว สญฺโญชนสมฺปยุตฺโต จ สญฺโญชนสมฺปยุตฺโต เจว โน จ สญฺโญชโน จ ธมฺมา อุปฺปชฺชนฺติ เหตุปจฺจยา.\csromanspacer{} สญฺโญชนสมฺปยุตฺตญฺเจว โน จ สญฺโญชนํ เอกํ ขนฺธญฺจ กามราคสญฺโญชนญฺจ ปฏิจฺจ ตโย ขนฺธา ทิฏฺฐิสญฺโญชนํ อวิชฺชาสญฺโญชนํ ฯเปฯ เทฺว ขนฺเธ จ ฯเปฯ. (จกฺกํ.) (3)}

\csromanheader{2}{1. ปจฺจยานุโลม 2. สงฺขฺยาวาร}
\csromanheader{2}{\textbf{สุทฺธ}}
\proseitem{103}{เหตุยา นว, อารมฺมเณ นว, อธิปติยา นว, (สพฺพตฺถ นว.) กมฺเม นว, อาหาเร นว ฯเปฯ อวิคเต นว.}

\vspace{\tipitakaparskip}
\csromancenter{อนุโลมํ.}
\csromanheader{2}{2. ปจฺจยปจฺจนีย 1. วิภงฺควาร}
\csromanheader{2}{\textbf{นเหตุ}}
\proseitem{104}{\csromanfolio{248}สญฺโญชนญฺเจว สญฺโญชนสมฺปยุตฺตญฺจ ธมฺมํ ปฏิจฺจ สญฺโญชโน เจว สญฺโญชนสมฺปยุตฺโต จ ธมฺโม อุปฺปชฺชติ นเหตุปจฺจยา.\csromanspacer{} วิจิกิจฺฉาสญฺโญชนํ ปฏิจฺจ อวิชฺชาสญฺโญชนํ. (1)}

\prose{สญฺโญชนสมฺปยุตฺตญฺเจว โน จ สญฺโญชนํ ธมฺมํ ปฏิจฺจ สญฺโญชโน เจว สญฺโญชนสมฺปยุตฺโต จ ธมฺโม อุปฺปชฺชติ นเหตุปจฺจยา.\csromanspacer{} วิจิกิจฺฉาสหคเต ขนฺเธ ปฏิจฺจ วิจิกิจฺฉาสหคโต โมโห. (1)}

\prose{สญฺโญชนญฺเจว สญฺโญชนสมฺปยุตฺตญฺจ สญฺโญชนสมฺปยุตฺตญฺเจว โน จ สญฺโญชนญฺจ ธมฺมํ ปฏิจฺจ สญฺโญชโน เจว สญฺโญชนสมฺปยุตฺโต จ ธมฺโม อุปฺปชฺชติ นเหตุปจฺจยา.\csromanspacer{} วิจิกิจฺฉาสญฺโญชนญฺจ สมฺปยุตฺตเก จ ขนฺเธ ปฏิจฺจ อวิชฺชาสญฺโญชนํ. (1)}

\csromanheader{2}{2. ปจฺจยปจฺจนีย 2. สงฺขฺยาวาร}
\csromanheader{2}{\textbf{สุทฺธ}}
\vspace{\tipitakaparskip}
\csromancenter{\textbf{นเหตุ}ยา ตีณิ, น-อธิปติยา นว, นปุเรชาเต นว, นปจฺฉาชาเต นว, น-อาเสวเน นว, นกมฺเม ตีณิ, นวิปาเก นว, นวิปฺปยุตฺเต นว.}
\prose{(เอวํ อิตเร เทฺว คณนาปิ สหชาตวาโรปิ กาตพฺโพ.\csromanspacer{} ปจฺจยวาโรปิ นิสฺสยวาโรปิ สํสฏฺฐวาโรปิ สมฺปยุตฺตวาโรปิ ปฏิจฺจวารสทิสา.)}

\csromanheader{2}{24. สญฺโญชนสญฺโญชนสมฺปยุตฺตทุก 7. ปญฺหาวาร}
\csromanheader{2}{1. ปจฺจยานุโลม 1. วิภงฺควาร}
\csromanheader{2}{\textbf{เหตุ}}
\proseitem{106}{สญฺโญชโน เจว สญฺโญชนสมฺปยุตฺโต จ ธมฺโม สญฺโญชนสฺส เจว สญฺโญชนสมฺปยุตฺตสฺส จ ธมฺมสฺส เหตุปจฺจเยน ปจฺจโย}

\vspace{\tipitakaparskip}
\csromancenter{สญฺโญชนสญฺโญชนสมฺปยุตฺตทุก กามราคสญฺโญชโน ทิฏฺฐิสญฺโญชนสฺส อวิชฺชาสญฺโญชนสฺส เหตุปจฺจเยน ปจฺจโย. (จกฺกํ.) (1)}
\prose{\csromanfolio{249}สญฺโญชโน เจว สญฺโญชนสมฺปยุตฺโต จ ธมฺโม สญฺโญชนสมฺปยุตฺตสฺส เจว โน จ สญฺโญชนสฺส ธมฺมสฺส เหตุปจฺจเยน ปจฺจโย.\csromanspacer{} สญฺโญชนา เจว สญฺโญชนสมฺปยุตฺตา จ เหตู สมฺปยุตฺตกานํ ขนฺธานํ เหตุปจฺจเยน ปจฺจโย. (2)}

\prose{สญฺโญชโน เจว สญฺโญชนสมฺปยุตฺโต จ ธมฺโม สญฺโญชนสฺส เจว สญฺโญชนสมฺปยุตฺตสฺส จ สญฺโญชนสมฺปยุตฺตสฺส เจว โน จ สญฺโญชนสฺส จ ธมฺมสฺส เหตุปจฺจเยน ปจฺจโย.\csromanspacer{} กามราคสญฺโญชโน ทิฏฺฐิสญฺโญชนสฺส อวิชฺชาสญฺโญชนสฺส สมฺปยุตฺตกานญฺจ ขนฺธานํ เหตุปจฺจเยน ปจฺจโย. (จกฺกํ.) (3)}

\csromanheader{2}{\textbf{อารมฺมณ}}
\proseitem{107}{สญฺโญชโน เจว สญฺโญชนสมฺปยุตฺโต จ ธมฺโม สญฺโญชนสฺส เจว สญฺโญชนสมฺปยุตฺตสฺส จ ธมฺมสฺส อารมฺมณปจฺจเยน ปจฺจโย.\csromanspacer{} สญฺโญชเน อารพฺภ สญฺโญชนา อุปฺปชฺชนฺติ. (มูลํ กาตพฺพํ.) สญฺโญชเน อารพฺภ สญฺโญชนสมฺปยุตฺตา เจว โน จ สญฺโญชนา ขนฺธา อุปฺปชฺชนฺติ. (มูลํ กาตพฺพํ.) สญฺโญชเน อารพฺภ สญฺโญชนา จ สญฺโญชนสมฺปยุตฺตกา จ ขนฺธา อุปฺปชฺชนฺติ. (1)}

\prose{สญฺโญชนสมฺปยุตฺโต เจว โน จ สญฺโญชโน ธมฺโม สญฺโญชนสมฺปยุตฺตสฺส เจว โน จ สญฺโญชนสฺส ธมฺมสฺส อารมฺมณปจฺจเยน ปจฺจโย.\csromanspacer{} สญฺโญชนสมฺปยุตฺเต เจว โน จ สญฺโญชเน ขนฺเธ อารพฺภ สญฺโญชนสมฺปยุตฺตา เจว โน จ สญฺโญชนา ขนฺธา อุปฺปชฺชนฺติ. (1)}

\prose{สญฺโญชนสมฺปยุตฺโต เจว โน จ สญฺโญชโน ธมฺโม สญฺโญชนสฺส เจว สญฺโญชนสมฺปยุตฺตสฺส จ ธมฺมสฺส อารมฺมณปจฺจเยน ปจฺจโย.\csromanspacer{} สญฺโญชนสมฺปยุตฺเต เจว โน จ สญฺโญชเน ขนฺเธ อารพฺภ สญฺโญชนา อุปฺปชฺชนฺติ. (2)}

\prose{สญฺโญชนสมฺปยุตฺโต เจว โน จ สญฺโญชโน ธมฺโม สญฺโญชนสฺส เจว สญฺโญชนสมฺปยุตฺตสฺส จ สญฺโญชนสมฺปยุตฺตสฺส เจว โน จ สญฺโญชนสฺส จ ธมฺมสฺส อารมฺมณปจฺจเยน ปจฺจโย.\csromanspacer{} สญฺโญชนสมฺปยุตฺเต \csromanfolio{250}เจว โน จ สญฺโญชเน ขนฺเธ อารพฺภ สญฺโญชนา จ สญฺโญชนสมฺปยุตฺตกา จ ขนฺธา อุปฺปชฺชนฺติ. (3)}

\prose{สญฺโญชโน เจว สญฺโญชนสมฺปยุตฺโต จ สญฺโญชนสมฺปยุตฺโต เจว โน จ สญฺโญชโน จ ธมฺมา สญฺโญชนสฺส เจว สญฺโญชนสมฺปยุตฺตสฺส จ ธมฺมสฺส อารมฺมณปจฺจเยน ปจฺจโย.\csromanspacer{} ตีณิ.}

\csromanheader{2}{\textbf{อธิปติ}}
\proseitem{108}{สญฺโญชโน เจว สญฺโญชนสมฺปยุตฺโต จ ธมฺโม สญฺโญชนสฺส เจว สญฺโญชนสมฺปยุตฺตสฺส จ ธมฺมสฺส อธิปติปจฺจเยน ปจฺจโย.\csromanspacer{} อารมฺมณาธิปติ.\csromanspacer{} ตีณิ.}

\prose{สญฺโญชนสมฺปยุตฺโต เจว โน จ สญฺโญชโน ธมฺโม สญฺโญชนสมฺปยุตฺตสฺส เจว โน จ สญฺโญชนสฺส ธมฺมสฺส อธิปติปจฺจเยน ปจฺจโย.\csromanspacer{} อารมฺมณาธิปติ, สหชาตาธิปติ.\csromanspacer{} ตีณิ. (อิมาสุ ตีสุปิ ปญฺหาสุ อารมฺมณาธิปติปิ สหชาตาธิปติปิ กาตพฺพา.)}

\prose{สญฺโญชโน เจว สญฺโญชนสมฺปยุตฺโต จ สญฺโญชนสมฺปยุตฺโต เจว โน จ สญฺโญชโน จ ธมฺมา สญฺโญชนสฺส เจว สญฺโญชนสมฺปยุตฺตสฺส จ ธมฺมสฺส อธิปติปจฺจเยน ปจฺจโย.\csromanspacer{} อารมฺมณาธิปติ.\csromanspacer{} ตีณิ.}

\csromanheader{2}{\textbf{อนนฺตราทิ}}
\proseitem{109}{สญฺโญชโน เจว สญฺโญชนสมฺปยุตฺโต จ ธมฺโม สญฺโญชนสฺส เจว สญฺโญชนสมฺปยุตฺตสฺส จ ธมฺมสฺส อนนฺตรปจฺจเยน ปจฺจโย.\csromanspacer{} นว. (นินฺนานากรณํ, วิภชนา นตฺถิ.\csromanspacer{} อารมฺมณสทิสา.) สมนนฺตรปจฺจเยน ปจฺจโย.\csromanspacer{} นว.\csromanspacer{} สหชาตปจฺจเยน ปจฺจโย.\csromanspacer{} นว.\csromanspacer{} อญฺญมญฺญปจฺจเยน ปจฺจโย.\csromanspacer{} นว.\csromanspacer{} นิสฺสยปจฺจเยน ปจฺจโย.\csromanspacer{} นว.\csromanspacer{} อุปนิสฺสยปจฺจเยน ปจฺจโย.\csromanspacer{} นว. (อารมฺมณนเยน กาตพฺพา.) อาเสวนปจฺจเยน ปจฺจโย.\csromanspacer{} นว.}

\csromanheader{2}{\textbf{กมฺมาทิ}}
\proseitem{110}{สญฺโญชนสมฺปยุตฺโต เจว โน จ สญฺโญชโน ธมฺโม สญฺโญชนสมฺปยุตฺตสฺส เจว โน จ สญฺโญชนสฺส ธมฺมสฺส กมฺมปจฺจเยน ปจฺจโย.\csromanspacer{} ตีณิ.\csromanspacer{} อาหารปจฺจเยน ปจฺจโย.\csromanspacer{} ตีณิ.\csromanspacer{} อินฺทฺริยปจฺจเยน ปจฺจโย.\csromanspacer{} ตีณิ.\csromanspacer{} ฌานปจฺจเยน ปจฺจโย.\csromanspacer{} ตีณิ.\csromanspacer{} มคฺคปจฺจเยน ปจฺจโย.}

\vspace{\tipitakaparskip}
\csromancenter{สญฺโญชนสญฺโญชนสมฺปยุตฺตทุก นว.\csromanspacer{} สมฺปยุตฺตปจฺจเยน ปจฺจโย.\csromanspacer{} นว.\csromanspacer{} อตฺถิปจฺจเยน ปจฺจโย.\csromanspacer{} นว.\csromanspacer{} นตฺถิปจฺจเยน ปจฺจโย.\csromanspacer{} นว.\csromanspacer{} วิคตปจฺจเยน ปจฺจโย.\csromanspacer{} นว.\csromanspacer{} อวิคตปจฺจเยน ปจฺจโย.\csromanspacer{} นว.}
\csromanheader{2}{1. ปจฺจยานุโลม 2. สงฺขฺยาวาร}
\csromanheader{2}{\textbf{สุทฺธ}}
\proseitem{111}{\csromanfolio{251}เหตุยา ตีณิ, อารมฺมเณ นว, อธิปติยา นว, อนนฺตเร นว, สมนนฺตเร นว, สหชาเต นว, อญฺญมญฺเญ นว, นิสฺสเย นว, อุปนิสฺสเย นว, อาเสวเน นว, กมฺเม ตีณิ, อาหาเร ตีณิ, อินฺทฺริเย ตีณิ, ฌาเน ตีณิ, มคฺเค นว, สมฺปยุตฺเต นว, อตฺถิยา นว, นตฺถิยา นว, วิคเต นว, อวิคเต นว.}

\vspace{\tipitakaparskip}
\csromancenter{อนุโลมํ.}
\csromanheader{2}{2. ปจฺจนียุทฺธาร}
\proseitem{112}{สญฺโญชโน เจว สญฺโญชนสมฺปยุตฺโต จ ธมฺโม สญฺโญชนสฺส เจว สญฺโญชนสมฺปยุตฺตสฺส จ ธมฺมสฺส อารมฺมณปจฺจเยน ปจฺจโย.\csromanspacer{} สหชาตปจฺจเยน ปจฺจโย.\csromanspacer{} อุปนิสฺสยปจฺจเยน ปจฺจโย. (สํขิตฺตํ.\csromanspacer{} เอวํ นว ปญฺหา กาตพฺพา.\csromanspacer{} ตีสุเยว ปเทสุ ปริวตฺเตตพฺพา.\csromanspacer{} นานากฺขณิกา นตฺถิ.)}

\csromanheader{2}{2. ปจฺจยปจฺจนีย 2. สงฺขฺยาวาร}
\vspace{\tipitakaparskip}
\csromancenter{นเหตุยา นว, น-อารมฺมเณ นว, (สพฺพตฺถ นว.) โน-อวิคเต นว.}
\csromanheader{2}{3. ปจฺจยานุโลมปจฺจนีย}
\csromanheader{2}{\textbf{เหตุทุก}}
\proseitem{114}{เหตุปจฺจยา น-อารมฺมเณ ตีณิ ฯเปฯ นสมนนฺตเร ตีณิ, น-อุปนิสฺสเย ตีณิ, นปุเรชาเต ตีณิ ฯเปฯ นมคฺเค ตีณิ, นสมฺปยุตฺเต ตีณิ, นวิปฺปยุตฺเต ตีณิ, โนนตฺถิยา ตีณิ, โนวิคเต ตีณิ.}

\csromanheader{2}{4. ปจฺจยปจฺจนียานุโลม}
\proseitem{115}{\csromanfolio{252}นเหตุปจฺจยา อารมฺมเณ นว, อธิปติยา นว, (อนุโลมปทานิ กาตพฺพานิ.) ฯเปฯ อวิคเต นว.}

\nitthitam{สญฺโญชนสญฺโญชนสมฺปยุตฺตทุกํ นิฏฺฐิตํ.}
\csromanmatikaanchor{mtk.30}
\csromantocmarkref{section}{25. สญฺโญชนวิปฺปยุตฺตสญฺโญชนิยทุก}{mtk.30}
\bookmark[dest={mtk.30},level=1]{25. สญฺโญชนวิปฺปยุตฺตสญฺโญชนิยทุก}
\csromanheader{1}{25. สญฺโญชนวิปฺปยุตฺตสญฺโญชนิยทุก 1. ปฏิจฺจวาร}
\csromanheader{2}{1. ปจฺจยานุโลม 1. วิภงฺควาร}
\proseitem{116}{สญฺโญชนวิปฺปยุตฺตํ สญฺโญชนิยํ ธมฺมํ ปฏิจฺจ สญฺโญชนวิปฺปยุตฺโต สญฺโญชนิโย ธมฺโม อุปฺปชฺชติ เหตุปจฺจยา.\csromanspacer{} สญฺโญชนวิปฺปยุตฺตํ สญฺโญชนิยํ เอกํ ขนฺธํ ปฏิจฺจ ตโย ขนฺธา จิตฺตสมุฏฺฐานญฺจ รูปํ ฯเปฯ เทฺว ขนฺเธ ฯเปฯ.\csromanspacer{} ปฏิสนฺธิกฺขเณ ฯเปฯ. (ยาว อสญฺญสตฺตา.\csromanspacer{} มหาภูตา.)}

\prose{สญฺโญชนวิปฺปยุตฺตํ อสญฺโญชนิยํ ธมฺมํ ปฏิจฺจ สญฺโญชนวิปฺปยุตฺโต อสญฺโญชนิโย ธมฺโม อุปฺปชฺชติ เหตุปจฺจยา.\csromanspacer{} สญฺโญชนวิปฺปยุตฺตํ อสญฺโญชนิยํ เอกํ ขนฺธํ ปฏิจฺจ ตโย ขนฺธา ฯเปฯ เทฺว ขนฺเธ ฯเปฯ.}

\prose{สญฺโญชนวิปฺปยุตฺตํ อสญฺโญชนิยํ ธมฺมํ ฯเปฯ. (เทฺว ปญฺหา กาตพฺพา.)}

\prose{(อิมํ ทุกํ จูฬนฺตรทุเก โลกิยทุกสทิสํ, นินฺนานากรณํ.)}

\nitthitam{สญฺโญชนวิปฺปยุตฺตสญฺโญชนิยทุกํ นิฏฺฐิตํ.}
\nitthitam{สญฺโญชนโคจฺฉกํ นิฏฺฐิตํ.}
\csromanmatikaanchor{mtk.31}
\csromanmatikaanchor{mtk.32}
\csromantocmarknopage{chapter}{5. คนฺถโคจฺฉก}{mtk.31}
\bookmark[dest={mtk.31},level=0]{5. คนฺถโคจฺฉก}
\csromantocmarkref{section}{26. คนฺถทุก}{mtk.32}
\bookmark[dest={mtk.32},level=1]{26. คนฺถทุก}
\csromanheader{1}{26. คนฺถทุก \textbf{5. คนฺถโคจฺฉก}}
\csromanheader{2}{26. คนฺถทุก 1. ปฏิจฺจวาร}
\csromanheader{2}{1. ปจฺจยานุโลม 1. วิภงฺควาร}
\csromanheader{2}{\textbf{เหตุ}}
\proseitem{1}{\csromanfolio{253}คนฺถํ ธมฺมํ ปฏิจฺจ คนฺโถ ธมฺโม อุปฺปชฺชติ เหตุปจฺจยา.\csromanspacer{} สีลพฺพตปรามาสํ กายคนฺถํ ปฏิจฺจ อภิชฺฌากายคนฺโถ.\csromanspacer{} อภิชฺฌากายคนฺถํ ปฏิจฺจ สีลพฺพตปรามาโส กายคนฺโถ.\csromanspacer{} อิทํสจฺจาภินิเวสํ กายคนฺถํ ปฏิจฺจ อภิชฺฌากายคนฺโถ.\csromanspacer{} อภิชฺฌากายคนฺถํ ปฏิจฺจ อิทํสจฺจาภินิเวโส กายคนฺโถ. (1)}

\prose{คนฺถํ ธมฺมํ ปฏิจฺจ โนคนฺโถ ธมฺโม อุปฺปชฺชติ เหตุปจฺจยา.\csromanspacer{} คนฺเถ ปฏิจฺจ สมฺปยุตฺตกา ขนฺธา จิตฺตสมุฏฺฐานญฺจ รูปํ. (2)}

\prose{คนฺถํ ธมฺมํ ปฏิจฺจ คนฺโถ จ โนคนฺโถ จ ธมฺมา อุปฺปชฺชนฺติ เหตุปจฺจยา.\csromanspacer{} สีลพฺพตปรามาสํ กายคนฺถํ ปฏิจฺจ อภิชฺฌากายคนฺโถ สมฺปยุตฺตกา จ ขนฺธา จิตฺตสมุฏฺฐานญฺจ รูปํ. (จกฺกํ.) (3)}

\proseitem{2}{โนคนฺถํ ธมฺมํ ปฏิจฺจ โนคนฺโถ ธมฺโม อุปฺปชฺชติ เหตุปจฺจยา.\csromanspacer{} โนคนฺถํ เอกํ ขนฺธํ ปฏิจฺจ ตโย ขนฺธา จิตฺตสมุฏฺฐานญฺจ รูปํ ฯเปฯ เทฺว ขนฺเธ ฯเปฯ.\csromanspacer{} ปฏิสนฺธิกฺขเณ ฯเปฯ.\csromanspacer{} ขนฺเธ ปฏิจฺจ วตฺถุ, วตฺถุํ ปฏิจฺจ ขนฺธา.\csromanspacer{} เอกํ มหาภูตํ ฯเปฯ. (1)}

\prose{โนคนฺถํ ธมฺมํ ปฏิจฺจ คนฺโถ ธมฺโม อุปฺปชฺชติ เหตุปจฺจยา.\csromanspacer{} โนคนฺเถ ขนฺเธ ปฏิจฺจ คนฺถา. (2)}

\prose{โนคนฺถํ ธมฺมํ ปฏิจฺจ คนฺโถ จ โนคนฺโถ จ ธมฺมา อุปฺปชฺชนฺติ เหตุปจฺจยา.\csromanspacer{} โนคนฺถํ เอกํ ขนฺธํ ปฏิจฺจ ตโย ขนฺธา คนฺถา จ จิตฺตสมุฏฺฐานญฺจ รูปํ ฯเปฯ เทฺว ขนฺเธ ฯเปฯ. (3)}

\proseitem{3}{คนฺถญฺจ โนคนฺถญฺจ ธมฺมํ ปฏิจฺจ คนฺโถ ธมฺโม อุปฺปชฺชติ เหตุปจฺจยา.\csromanspacer{} สีลพฺพตปรามาสํ กายคนฺถญฺจ สมฺปยุตฺตเก จ ขนฺเธ ปฏิจฺจ อภิชฺฌากายคนฺโถ. (จกฺกํ.) (1)}

\prose{\csromanfolio{254}คนฺถญฺจ โนคนฺถญฺจ ธมฺมํ ปฏิจฺจ โนคนฺโถ ธมฺโม อุปฺปชฺชติ เหตุปจฺจยา.\csromanspacer{} โนคนฺถํ เอกํ ขนฺธญฺจ คนฺเถ จ ปฏิจฺจ ตโย ขนฺธา จิตฺตสมุฏฺฐานญฺจ รูปํ ฯเปฯ เทฺว ขนฺเธ จ ฯเปฯ.\csromanspacer{} คนฺเถ จ สมฺปยุตฺตเก ขนฺเธ จ ปฏิจฺจ จิตฺตสมุฏฺฐานํ รูปํ. (2)}

\prose{คนฺถญฺจ โนคนฺถญฺจ ธมฺมํ ปฏิจฺจ คนฺโถ จ โนคนฺโถ จ ธมฺมา อุปฺปชฺชนฺติ เหตุปจฺจยา.\csromanspacer{} โนคนฺถํ เอกํ ขนฺธญฺจ สีลพฺพตปรามาสกายคนฺถญฺจ ปฏิจฺจ ตโย ขนฺธา อภิชฺฌากายคนฺโถ จิตฺตสมุฏฺฐานญฺจ รูปํ ฯเปฯ เทฺว ขนฺเธ จ ฯเปฯ. (3) (จกฺกํ.\csromanspacer{} สํขิตฺตํ.) อารมฺมณปจฺจยา ฯเปฯ อวิคตปจฺจยา.}

\csromanheader{2}{1. ปจฺจยานุโลม 2. สงฺขฺยาวาร}
\csromanheader{2}{\textbf{สุทฺธ}}
\proseitem{4}{เหตุยา นว, อารมฺมเณ นว, อธิปติยา นว, (สพฺพตฺถ นว.) วิปาเก เอกํ, อาหาเร นว ฯเปฯ อวิคเต นว.}

\vspace{\tipitakaparskip}
\csromancenter{อนุโลมํ.}
\csromanheader{2}{2. ปจฺจยปจฺจนีย 1. วิภงฺควาร}
\csromanheader{2}{\textbf{นเหตุ}}
\proseitem{5}{โนคนฺถํ ธมฺมํ ปฏิจฺจ โนคนฺโถ ธมฺโม อุปฺปชฺชติ นเหตุปจฺจยา.\csromanspacer{} อเหตุกํ โนคนฺถํ เอกํ ขนฺธํ ปฏิจฺจ ตโย ขนฺธา จิตฺตสมุฏฺฐานญฺจ รูปํ ฯเปฯ เทฺว ขนฺเธ ฯเปฯ.\csromanspacer{} อเหตุกปฏิสนฺธิกฺขเณ. (ยาว อสญฺญสตฺตา.) วิจิกิจฺฉาสหคเต อุทฺธจฺจสหคเต ขนฺเธ ปฏิจฺจ วิจิกิจฺฉาสหคโต อุทฺธจฺจสหคโต โมโห. (1)}

\csromanheader{2}{\textbf{น-อารมฺมณาทิ}}
\proseitem{6}{คนฺถํ ธมฺมํ ปฏิจฺจ โนคนฺโถ ธมฺโม อุปฺปชฺชติ น-อารมฺมณปจฺจยา.\csromanspacer{} คนฺเถ ปฏิจฺจ จิตฺตสมุฏฺฐานํ รูปํ. (1)}

\prose{โนคนฺถํ ธมฺมํ ปฏิจฺจ โนคนฺโถ ธมฺโม อุปฺปชฺชติ น-อารมฺมณปจฺจยา.\csromanspacer{} โนคนฺเถ ขนฺเธ ปฏิจฺจ จิตฺตสมุฏฺฐานํ รูปํ.\csromanspacer{} ปฏิสนฺธิกฺขเณ ฯเปฯ. (ยาว อสญฺญสตฺตา.) (1)}

\csromanheader{2}{26. คนฺถทุก}
\prose{\csromanfolio{255}คนฺถญฺจ โนคนฺถญฺจ ธมฺมํ ปฏิจฺจ โนคนฺโถ ธมฺโม อุปฺปชฺชติ น-อารมฺมณปจฺจยา.\csromanspacer{} คนฺเถ จ สมฺปยุตฺตเก จ ขนฺเธ ปฏิจฺจ จิตฺตสมุฏฺฐานํ รูปํ. (สํขิตฺตํ.)}

\prose{น-อธิปติปจฺจยา.\csromanspacer{} นว.\csromanspacer{} น-อนนฺตรปจฺจยา.\csromanspacer{} ตีณิ.\csromanspacer{} นสมนนฺตรปจฺจยา.\csromanspacer{} ตีณิ.\csromanspacer{} น-อญฺญมญฺญปจฺจยา.\csromanspacer{} ตีณิ.\csromanspacer{} น-อุปนิสฺสยปจฺจยา.\csromanspacer{} ตีณิ.}

\csromanheader{2}{\textbf{นปุเรชาตาทิ}}
\proseitem{7}{คนฺถํ ธมฺมํ ปฏิจฺจ คนฺโถ ธมฺโม อุปฺปชฺชติ นปุเรชาตปจฺจยา.\csromanspacer{} อรูเป อิทํสจฺจาภินิเวสํ กายคนฺถํ ปฏิจฺจ อภิชฺฌากายคนฺโถ.\csromanspacer{} อภิชฺฌากายคนฺถํ ปฏิจฺจ อิทํสจฺจาภินิเวโส กายคนฺโถ. (อรูเป สีลพฺพตปรามาโส นตฺถิ, เอวํ นว ปญฺหา กาตพฺพา.) นปจฺฉาชาเต นว, น-อาเสวเน นว, นกมฺเม ตีณิ, นวิปาเก นว, น-อาหาเร เอกํ, น-อินฺทฺริเย เอกํ, นฌาเน เอกํ, นมคฺเค เอกํ, นสมฺปยุตฺเต ตีณิ, นวิปฺปยุตฺเต นว, โนนตฺถิยา ตีณิ, โนวิคเต ตีณิ.}

\csromanheader{2}{2. ปจฺจยปจฺจนีย 2. สงฺขฺยาวาร}
\csromanheader{2}{\textbf{สุทฺธ}}
\proseitem{8}{นเหตุยา เอกํ, น-อารมฺมเณ ตีณิ, น-อธิปติยา นว, น-อนนฺตเร ตีณิ, นสมนนฺตเร ตีณิ, น-อญฺญมญฺเญ ตีณิ, น-อุปนิสฺสเย ตีณิ, นปุเรชาเต นว, นปจฺฉาชาเต นว, น-อาเสวเน นว, นกมฺเม ตีณิ, นวิปาเก นว, น-อาหาเร เอกํ, น-อินฺทฺริเย เอกํ, นฌาเน เอกํ, นมคฺเค เอกํ, นสมฺปยุตฺเต ตีณิ, นวิปฺปยุตฺเต นว, โนนตฺถิยา ตีณิ, โนวิคเต ตีณิ.}

\csromanheader{2}{3. ปจฺจยานุโลมปจฺจนีย}
\proseitem{9}{เหตุปจฺจยา น-อารมฺมเณ ตีณิ, น-อธิปติยา นว. (สํขิตฺตํ.\csromanspacer{} เอวํ คเณตพฺพํ.)}

\csromanheader{2}{4. ปจฺจยปจฺจนียานุโลม}
\vspace{\tipitakaparskip}
\csromancenter{นเหตุปจฺจยา อารมฺมเณ เอกํ ฯเปฯ อนนฺตเร เอกํ ฯเปฯ อวิคเต เอกํ.}
\csromancenter{(สหชาตวาโร ปฏิจฺจวารสทิโส.)}
\csromanheader{2}{26. คนฺถทุก 3. ปจฺจยวาร}
\csromanheader{2}{1. ปจฺจยานุโลม 1. วิภงฺควาร}
\csromanheader{2}{\textbf{เหตุ}}
\proseitem{11}{\csromanfolio{256}คนฺถํ ธมฺมํ ปจฺจยา คนฺโถ ธมฺโม อุปฺปชฺชติ เหตุปจฺจยา.\csromanspacer{} ตีณิ. (ปฏิจฺจวารสทิโส.)}

\prose{โนคนฺถํ ธมฺมํ ปจฺจยา โนคนฺโถ ธมฺโม อุปฺปชฺชติ เหตุปจฺจยา.\csromanspacer{} โนคนฺถํ เอกํ ขนฺธํ ปจฺจยา ตโย ขนฺธา จิตฺตสมุฏฺฐานญฺจ รูปํ ฯเปฯ เทฺว ขนฺเธ ฯเปฯ.\csromanspacer{} ปฏิสนฺธิกฺขเณ ฯเปฯ.\csromanspacer{} ขนฺเธ ปจฺจยา วตฺถุ, วตฺถุํ ปจฺจยา ขนฺธา.\csromanspacer{} เอกํ มหาภูตํ ฯเปฯ.\csromanspacer{} วตฺถุํ ปจฺจยา โนคนฺถา ขนฺธา. (1)}

\prose{โนคนฺถํ ธมฺมํ ปจฺจยา คนฺโถ ธมฺโม อุปฺปชฺชติ เหตุปจฺจยา.\csromanspacer{} โนคนฺเถ ขนฺเธ ปจฺจยา คนฺถา.\csromanspacer{} วตฺถุํ ปจฺจยา คนฺถา. (2)}

\prose{โนคนฺถํ ธมฺมํ ปจฺจยา คนฺโถ จ โนคนฺโถ จ ธมฺมา อุปฺปชฺชนฺติ เหตุปจฺจยา.\csromanspacer{} โนคนฺถํ เอกํ ขนฺธํ ปจฺจยา ตโย ขนฺธา คนฺถา จ จิตฺตสมุฏฺฐานญฺจ รูปํ ฯเปฯ เทฺว ขนฺเธ ฯเปฯ.\csromanspacer{} วตฺถุํ ปจฺจยา คนฺถา สมฺปยุตฺตกา จ ขนฺธา. (3)}

\proseitem{12}{คนฺถญฺจ โนคนฺถญฺจ ธมฺมํ ปจฺจยา คนฺโถ ธมฺโม อุปฺปชฺชติ เหตุปจฺจยา.\csromanspacer{} สีลพฺพตปรามาสํ กายคนฺถญฺจ สมฺปยุตฺตเก จ ขนฺเธ ปจฺจยา อภิชฺฌากายคนฺโถ. (จกฺกํ.) สีลพฺพตปรามาสํ กายคนฺถญฺจ วตฺถุญฺจ ปจฺจยา อภิชฺฌากายคนฺโถ. (จกฺกํ.) (1)}

\prose{คนฺถญฺจ โนคนฺถญฺจ ธมฺมํ ปจฺจยา โนคนฺโถ ธมฺโม อุปฺปชฺชติ เหตุปจฺจยา.\csromanspacer{} โนคนฺถํ เอกํ ขนฺธญฺจ คนฺเถ จ ปจฺจยา ตโย ขนฺธา จิตฺตสมุฏฺฐานญฺจ รูปํ ฯเปฯ เทฺว ขนฺเธ ฯเปฯ.\csromanspacer{} คนฺเถ จ วตฺถุญฺจ ปจฺจยา โนคนฺถา ขนฺธา. (2)}

\prose{คนฺถญฺจ โนคนฺถญฺจ ธมฺมํ ปจฺจยา คนฺโถ จ โนคนฺโถ จ ธมฺมา อุปฺปชฺชนฺติ เหตุปจฺจยา.\csromanspacer{} โนคนฺถํ เอกํ ขนฺธญฺจ สีลพฺพตรามาสํ กายคนฺถญฺจ ปจฺจยา ตโย ขนฺธา อภิชฺฌากายคนฺโถ จ จิตฺตสมุฏฺฐานญฺจ รูปํ ฯเปฯ เทฺว ขนฺเธ ฯเปฯ. (จกฺกํ.) สีลพฺพตปรามาสํ กายคนฺถญฺจ วตฺถุญฺจ ปจฺจยา อภิชฺฌากายคนฺโถ จ สมฺปยุตฺตกา จ ขนฺธา. (จกฺกํ.) (3)}

\csromanheader{2}{26. คนฺถทุก \textbf{1. ปจฺจยานุโลม 2. สงฺขฺยาวาร}}
\vspace{\tipitakaparskip}
\csromancenter{เหตุยา นว, อารมฺมเณ นว, อธิปติยา นว ฯเปฯ อวิคเต นว.}
\csromancenter{อนุโลมํ.}
\csromanheader{2}{2. ปจฺจยปจฺจนีย 1. วิภงฺควาร}
\csromanheader{2}{\textbf{นเหตุ}}
\proseitem{14}{\csromanfolio{257}โนคนฺถํ ธมฺมํ ปจฺจยา โนคนฺโถ ธมฺโม อุปฺปชฺชติ นเหตุปจฺจยา.\csromanspacer{} อเหตุกํ โนคนฺถํ เอกํ ขนฺธํ ปจฺจยา ตโย ขนฺธา จิตฺตสมุฏฺฐานญฺจ รูปํ ฯเปฯ เทฺว ขนฺเธ ฯเปฯ.\csromanspacer{} ปฏิสนฺธิกฺขเณ ฯเปฯ. (ยาว อสญฺญสตฺตา.) จกฺขายตนํ ปจฺจยา จกฺขุวิญฺญาณํ ฯเปฯ.\csromanspacer{} กายายตนํ ปจฺจยา กายวิญฺญาณํ.\csromanspacer{} วตฺถุํ ปจฺจยา อเหตุกา โนคนฺถา ขนฺธา.\csromanspacer{} วิจิกิจฺฉาสหคเต อุทฺธจฺจสหคเต ขนฺเธ จ วตฺถุญฺจ ปจฺจยา วิจิกิจฺฉาสหคโต อุทฺธจฺจสหคโต โมโห (สํขิตฺตํ.)}

\csromanheader{2}{2. ปจฺจยปจฺจนีย 2. สงฺขฺยาวาร}
\vspace{\tipitakaparskip}
\csromancenter{\textbf{นเหตุ}ยา เอกํ, น-อารมฺมเณ ตีณิ, น-อธิปติยา นว. (เอวํ คเณตพฺพํ.)}
\csromanheader{2}{3. ปจฺจยานุโลมปจฺจนิย}
\vspace{\tipitakaparskip}
\csromancenter{เหตุปจฺจยา น-อารมฺมเณ ตีณิ, น-อธิปติยา นว ฯเปฯ โนวิคเต ตีณิ.}
\csromanheader{2}{4. ปจฺจยปจฺจนียานุโลม}
\vspace{\tipitakaparskip}
\csromancenter{\textbf{นเหตุ}ปจฺจยา อารมฺมเณ เอกํ ฯเปฯ อวิคเต เอกํ.}
\prose{(นิสฺสยวาโร ปจฺจยวารสทิโสว.\csromanspacer{} สํสฏฺฐวาโรปิ สมฺปยุตฺตวาโรปิ นว ปญฺหา กาตพฺพา.\csromanspacer{} รูปํ นตฺถิ.)}

\csromanheader{2}{26. คนฺถทุก 7. ปญฺหาวาร}
\csromanheader{2}{1. ปจฺจยานุโลม 1. วิภงฺควาร}
\csromanheader{2}{\textbf{เหตุ}}
\proseitem{18}{\csromanfolio{258}คนฺโถ ธมฺโม คนฺถสฺส ธมฺมสฺส เหตุปจฺจเยน ปจฺจโย.\csromanspacer{} คนฺถา เหตู สมฺปยุตฺตกานํ คนฺถานํ เหตุปจฺจเยน ปจฺจโย. (1)}

\prose{คนฺโถ ธมฺโม โนคนฺถสฺส ธมฺมสฺส เหตุปจฺจเยน ปจฺจโย.\csromanspacer{} คนฺถา เหตู สมฺปยุตฺตกานํ ขนฺธานํ จิตฺตสมุฏฺฐานานญฺจ รูปานํ เหตุปจฺจเยน ปจฺจโย. (2)}

\prose{คนฺโถ ธมฺโม คนฺถสฺส จ โนคนฺถสฺส จ ธมฺมสฺส เหตุปจฺจเยน ปจฺจโย.\csromanspacer{} คนฺถา เหตู สมฺปยุตฺตกานํ ขนฺธานํ คนฺถานญฺจ จิตฺตสมุฏฺฐานญฺจ รูปานํ เหตุปจฺจเยน ปจฺจโย. (3)}

\proseitem{19}{โนคนฺโถ ธมฺโม โนคนฺถสฺส ธมฺมสฺส เหตุปจฺจเยน ปจฺจโย.\csromanspacer{} โนคนฺถา เหตู สมฺปยุตฺตกานํ ขนฺธานํ จิตฺตสมุฏฺฐานานญฺจ รูปานํ เหตุปจฺจเยน ปจฺจโย.\csromanspacer{} ปฏิสนฺธิกฺขเณ ฯเปฯ. (1)}

\prose{โนคนฺโถ ธมฺโม คนฺถสฺส ธมฺมสฺส เหตุปจฺจเยน ปจฺจโย.\csromanspacer{} โนคนฺถา เหตู สมฺปยุตฺตกานํ คนฺถานํ เหตุปจฺจเยน ปจฺจโย. (2)}

\prose{โนคนฺโถ ธมฺโม คนฺถสฺส จ โนคนฺถสฺส จ ธมฺมสฺส เหตุปจฺจเยน ปจฺจโย.\csromanspacer{} โนคนฺถา เหตู สมฺปยุตฺตกานํ ขนฺธานํ คนฺถานญฺจ จิตฺตสมุฏฺฐานานญฺจ รูปานํ เหตุปจฺจเยน ปจฺจโย. (3)}

\proseitem{20}{คนฺโถ จ โนคนฺโถ จ ธมฺมา คนฺถสฺส ธมฺมสฺส เหตุปจฺจเยน ปจฺจโย.\csromanspacer{} คนฺถา จ โนคนฺถา จ เหตู สมฺปยุตฺตกานํ คนฺถานํ เหตุปจฺจเยน ปจฺจโย. (1)}

\prose{คนฺโถ จ โนคนฺโถ จ ธมฺมา โนคนฺถสฺส ธมฺมสฺส เหตุปจฺจเยน ปจฺจโย.\csromanspacer{} คนฺถา จ โนคนฺถา จ เหตู สมฺปยุตฺตกานํ ขนฺธานํ จิตฺตสมุฏฺฐานานญฺจ รูปานํ เหตุปจฺจเยน ปจฺจโย. (2)}

\csromanheader{2}{26. คนฺถทุก}
\prose{\csromanfolio{259}คนฺโถ จ โนคนฺโถ จ ธมฺมา คนฺถสฺส จ โนคนฺถสฺส จ ธมฺมสฺส เหตุปจฺจเยน ปจฺจโย.\csromanspacer{} คนฺถา จ โนคนฺถา จ เหตู สมฺปยุตฺตกานํ ขนฺธานํ คนฺถานญฺจ จิตฺตสมุฏฺฐานญฺจ รูปานํ เหตุปจฺจเยน ปจฺจโย. (3)}

\csromanheader{2}{\textbf{อารมฺมณ}}
\proseitem{21}{คนฺโถ ธมฺโม คนฺถสฺส ธมฺมสฺส อารมฺมณปจฺจเยน ปจฺจโย.\csromanspacer{} คนฺเถ อารพฺภ คนฺถา อุปฺปชฺชนฺติ. (มูลํ กาตพฺพํ.) คนฺเถ อารพฺภ โนคนฺถา ขนฺธา อุปฺปชฺชนฺติ. (มูลํ กาตพฺพํ.) คนฺเถ อารพฺภ คนฺถา จ สมฺปยุตฺตกา จ ขนฺธา อุปฺปชฺชนฺติ. (3)}

\proseitem{22}{โนคนฺโถ ธมฺโม โนคนฺถสฺส ธมฺมสฺส อารมฺมณปจฺจเยน ปจฺจโย.\csromanspacer{} ทานํ ทตฺวา, สีลํ ฯเปฯ อุโปสถกมฺมํ กตฺวา ตํ ปจฺจเวกฺขติ.\csromanspacer{} ปุพฺเพ สุจิณฺณานิ ฯเปฯ.\csromanspacer{} ฌานา ฯเปฯ.\csromanspacer{} อริยา มคฺคา วุฏฺฐหิตฺวา มคฺคํ ปจฺจเวกฺขนฺติ, ผลํ, ปจฺจเวกฺขนฺติ, นิพฺพานํ ปจฺจเวกฺขนฺติ.\csromanspacer{} นิพฺพานํ โคตฺรภุสฺส, โวทานสฺส, มคฺคสฺส, ผลสฺส, อาวชฺชนาย อารมฺมณปจฺจเยน ปจฺจโย.\csromanspacer{} อริยา โนคนฺเถ ปหีเน กิเลเส ฯเปฯ.\csromanspacer{} วิกฺขมฺภิเต กิเลเส ปจฺจเวกฺขนฺติ.\csromanspacer{} ปุพฺเพ ฯเปฯ.\csromanspacer{} จกฺขุํ ฯเปฯ วตฺถุํ โนคนฺเถ ขนฺเธ อนิจฺจโต ฯเปฯ โทมนสฺสํ อุปฺปชฺชติ.\csromanspacer{} ทิพฺเพน จกฺขุนา รูปํ ปสฺสติ.\csromanspacer{} ทิพฺพาย โสตธาตุยา สทฺทํ สุณาติ.\csromanspacer{} เจโตปริยญาเณน โนคนฺถจิตฺตสมงฺคิสฺส จิตฺตํ ชานาติ.\csromanspacer{} อากาสานญฺจายตนํ วิญฺญาณญฺจายตนสฺส ฯเปฯ.\csromanspacer{} อากิญฺจญฺญายตนํ เนวสญฺญานาสญฺญายตนสฺส ฯเปฯ.\csromanspacer{} รูปายตนํ จกฺขุวิญฺญาณสฺส ฯเปฯ.\csromanspacer{} โผฏฺฐพฺพายตนํ กายวิญฺญาณสฺส ฯเปฯ.\csromanspacer{} โนคนฺถา ขนฺธา อิทฺธิวิธญาณสฺส, เจโตปริยญาณสฺส, ปุพฺเพนิวาสานุสฺสติญาณสฺส, ยถากมฺมูปคญาณสฺส, อนาคตํสญาณสฺส, อาวชฺชนาย อารมฺมณปจฺจเยน ปจฺจโย. (1)}

\prose{โนคนฺโถ ธมฺโม คนฺถสฺส ธมฺมสฺส อารมฺมณปจฺจเยน ปจฺจโย.\csromanspacer{} ทานํ ทตฺวา, สีลํ ฯเปฯ อุโปสถกมฺมํ กตฺวา ตํ อสฺสาเทติ อภินนฺทติ, ตํ อารพฺภ ราโค อุปฺปชฺชติ, ทิฏฺฐิ อุปฺปชฺชติ, โทมนสฺสํ อุปฺปชฺชติ.\csromanspacer{} ปุพฺเพ สุจิณฺณานิ ฯเปฯ.\csromanspacer{} ฌานา ฯเปฯ.\csromanspacer{} จกฺขุํ ฯเปฯ วตฺถุํ โนคนฺเถ ขนฺเธ อสฺสาเทติ อภินนฺทติ, ตํ อารพฺภ ราโค อุปฺปชฺชติ ฯเปฯ โทมนสฺสํ อุปฺปชฺชติ. (2)}

\prose{โนคนฺโถ ธมฺโม คนฺถสฺส จ โนคนฺถสฺส จ ธมฺมสฺส อารมฺมณปจฺจเยน ปจฺจโย.\csromanspacer{} ทานํ ทตฺวา, สีลํ ฯเปฯ อุโปสถกมฺมํ กตฺวา ตํ \csromanfolio{260}อสฺสาเทติ อภินนฺทติ, ตํ อารพฺภ คนฺถา จ สมฺปยุตฺตกา จ ขนฺธา อุปฺปชฺชนฺติ.\csromanspacer{} ปุพฺเพ สุจิณฺณานิ ฯเปฯ.\csromanspacer{} ฌานา ฯเปฯ.\csromanspacer{} จกฺขุํ ฯเปฯ วตฺถุํ โนคนฺเถ ขนฺเธ อสฺสาเทติ อภินนฺทติ, ตํ อารพฺภ คนฺถา จ สมฺปยุตฺตกา จ ขนฺธา อุปฺปชฺชนฺติ. (3)}

\prose{คนฺโถ จ โนคนฺโถ จ ธมฺมา คนฺถสฺส ธมฺมสฺส อารมฺมณปจฺจเยน ปจฺจโย.\csromanspacer{} ตีณิ. (อารพฺภ กาตพฺพา.)}

\csromanheader{2}{\textbf{อธิปติ}}
\proseitem{23}{คนฺโถ ธมฺโม คนฺถสฺส ธมฺมสฺส อธิปติปจฺจเยน ปจฺจโย.\csromanspacer{} ตีณิ. (อารมฺมณสทิสา.\csromanspacer{} ครุการมฺมณา กาตพฺพา.)}

\prose{โนคนฺโถ ธมฺโม โนคนฺถสฺส ธมฺมสฺส อธิปติปจฺจเยน ปจฺจโย.\csromanspacer{} \textbf{อารมฺมณาธิปติ}, สหชาตาธิปติ.\csromanspacer{}\csromanspacer{}\textbf{อารมฺมณาธิปติ\csromandash{}}ทานํ ฯเปฯ สีลํ ฯเปฯ อุโปสถกมฺมํ ฯเปฯ.\csromanspacer{} ปุพฺเพ ฯเปฯ.\csromanspacer{} ฌานา ฯเปฯ.\csromanspacer{} อริยา มคฺคา ฯเปฯ.\csromanspacer{} ผลํ ฯเปฯ.\csromanspacer{} นิพฺพานํ ฯเปฯ.\csromanspacer{} นิพฺพานํ โคตฺรภุสฺส, โวทานสฺส, มคฺคสฺส, ผลสฺส, อธิปติปจฺจเยน ปจฺจโย.\csromanspacer{} จกฺขุํ ฯเปฯ วตฺถุํ โนคนฺเถ ขนฺเธ ครุํ กตฺวา อสฺสาเทติ อภินนฺทติ, ตํ ครุํ กตฺวา ราโค อุปฺปชฺชติ, ทิฏฺฐิ อุปฺปชฺชติ.\csromanspacer{} \textbf{สหชาตาธิปติ}\csromandash{}โนคนฺถาธิปติ สมฺปยุตฺตกานํ ขนฺธานํ จิตฺตสมุฏฺฐานานญฺจ รูปานํ อธิปติปจฺจเยน ปจฺจโย. (1)}

\prose{โนคนฺโถ ธมฺโม คนฺถสฺส ธมฺมสฺส อธิปติปจฺจเยน ปจฺจโย.\csromanspacer{} \textbf{อารมฺมณาธิปติ}, สหชาตาธิปติ.\csromanspacer{}\csromanspacer{}\textbf{อารมฺมณาธิปติ\csromandash{}}ทานํ ฯเปฯ สีลํ ฯเปฯ อุโปสถกมฺมํ ฯเปฯ.\csromanspacer{} ปุพฺเพ สุจิณฺณานิ ฯเปฯ.\csromanspacer{} ฌานา ฯเปฯ.\csromanspacer{} จกฺขุํ ฯเปฯ วตฺถุํ โนคนฺเถ ขนฺเธ ครุํ กตฺวา อสฺสาเทติ อภินนฺทติ, ตํ ครุํ กตฺวา ราโค อุปฺปชฺชติ, ทิฏฺฐิ อุปฺปชฺชติ.\csromanspacer{} \textbf{สหชาตาธิปติ}\csromandash{} โนคนฺถาธิปติ คนฺถานํ อธิปติปจฺจเยน ปจฺจโย. (2)}

\prose{โนคนฺโถ ธมฺโม คนฺถสฺส จ โนคนฺถสฺส จ ธมฺมสฺส อธิปติปจฺจเยน ปจฺจโย.\csromanspacer{} \textbf{อารมฺมณาธิปติ}, สหชาตาธิปติ.\csromanspacer{}\csromanspacer{}\textbf{อารมฺมณาธิปติ\csromandash{}}ทานํ ฯเปฯ โนคนฺเถ ขนฺเธ ครุํ กตฺวา อสฺสาเทติ อภินนฺทติ, ตํ ครุํ กตฺวา คนฺถา จ สมฺปยุตฺตกา จ ขนฺธา อุปฺปชฺชนฺติ.\csromanspacer{} \textbf{สหชาตาธิปติ}\csromandash{}โนคนฺถาธิปติ สมฺปยุตฺตกานํ ขนฺธานํ คนฺถานญฺจ จิตฺตสมุฏฺฐานานญฺจ รูปานํ อธิปติปจฺจเยน ปจฺจโย. (3)}

\csromanheader{2}{26. คนฺถทุก}
\proseitem{24}{\csromanfolio{261}คนฺโถ จ โนคนฺโถ จ ธมฺมา คนฺถสฺส ธมฺมสฺส อธิปติปจฺจเยน ปจฺจโย.\csromanspacer{} อารมฺมณาธิปติ.\csromanspacer{} ตีณิ.}

\csromanheader{2}{\textbf{อนนฺตร}}
\proseitem{25}{คนฺโถ ธมฺโม คนฺถสฺส ธมฺมสฺส อนนฺตรปจฺจเยน ปจฺจโย.\csromanspacer{} ปุริมา ปุริมา คนฺถา ปจฺฉิมานํ ปจฺฉิมานํ คนฺถานํ อนนฺตรปจฺจเยน ปจฺจโย. (1)}

\prose{คนฺโถ ธมฺโม โนคนฺถสฺส ธมฺมสฺส อนนฺตรปจฺจเยน ปจฺจโย.\csromanspacer{} ปุริมา ปุริมา คนฺถา ปจฺฉิมานํ ปจฺฉิมานํ โนคนฺถานํ ขนฺธานํ อนนฺตรปจฺจเยน ปจฺจโย.\csromanspacer{} คนฺถา วุฏฺฐานสฺส อนนฺตรปจฺจเยน ปจฺจโย. (2)}

\prose{คนฺโถ ธมฺโม คนฺถสฺส จ โนคนฺถสฺส จ ธมฺมสฺส อนนฺตรปจฺจเยน ปจฺจโย.\csromanspacer{} ปุริมา ปุริมา คนฺถา ปจฺฉิมานํ ปจฺฉิมานํ คนฺถา นํ สมฺปยุตฺตกานญฺจ ขนฺธานํ อนนฺตรปจฺจเยน ปจฺจโย. (3)}

\proseitem{26}{โนคนฺโถ ธมฺโม โนคนฺถสฺส ธมฺมสฺส อนนฺตรปจฺจเยน ปจฺจโย.\csromanspacer{} ตีณิ. (เทฺว อาวชฺชนา กาตพฺพา.\csromanspacer{} ปฐโม นตฺถิ.)}

\prose{คนฺโถ จ โนคนฺโถ จ ธมฺมา คนฺถสฺส ธมฺมสฺส อนนฺตรปจฺจเยน ปจฺจโย.\csromanspacer{} ตีณิ. (เอกมฺปิ วุฏฺฐานํ กาตพฺพํ.\csromanspacer{} มชฺเฌ.)}

\csromanheader{2}{\textbf{สมนนฺตราทิ}}
\proseitem{27}{คนฺโถ ธมฺโม คนฺถสฺส ธมฺมสฺส สมนนฺตรปจฺจเยน ปจฺจโย.\csromanspacer{} นว.\csromanspacer{} สหชาตปจฺจเยน ปจฺจโย.\csromanspacer{} นว.\csromanspacer{} อญฺญมญฺญปจฺจเยน ปจฺจโย.\csromanspacer{} นว.\csromanspacer{} นิสฺสยปจฺจเยน ปจฺจโย.\csromanspacer{} นว.}

\csromanheader{2}{\textbf{อุปนิสฺสย}}
\proseitem{28}{คนฺโถ ธมฺโม คนฺถสฺส ธมฺมสฺส อุปนิสฺสยปจฺจเยน ปจฺจโย.\csromanspacer{} อารมฺมณูปนิสฺสโย, อนนฺตรูปนิสฺสโย, ปกตูปนิสฺสโย ฯเปฯ.\csromanspacer{} ปกตูปนิสฺสโย\csromandash{}คนฺถา คนฺถานํ อุปนิสฺสยปจฺจเยน ปจฺจโย.\csromanspacer{} ตีณิ.}

\proseitem{29}{โนคนฺโถ ธมฺโม โนคนฺถสฺส ธมฺมสฺส อุปนิสฺสยปจฺจเยน ปจฺจโย.\csromanspacer{} อารมฺมณูปนิสฺสโย, อนนฺตรูปนิสฺสโย, ปกตูปนิสฺสโย \csromanfolio{262}ฯเปฯ.\csromanspacer{} ปกตูปนิสฺสโย\csromandash{}สทฺธํ อุปนิสฺสาย ทานํ เทติ ฯเปฯ มานํ ชปฺเปติ.\csromanspacer{} ทิฏฺฐิํ คณฺหาติ.\csromanspacer{} สีลํ ฯเปฯ.\csromanspacer{} เสนาสนํ อุปนิสฺสาย ทานํ เทติ ฯเปฯ สํฆํ ภินฺทติ.\csromanspacer{} สทฺธา ฯเปฯ.\csromanspacer{} เสนาสนํ สทฺธาย ฯเปฯ ผลสมาปตฺติยา อุปนิสฺสยปจฺจเยน ปจฺจโย. (1)}

\prose{โนคนฺโถ ธมฺโม คนฺถสฺส ธมฺมสฺส อุปนิสฺสยปจฺจเยน ปจฺจโย.\csromanspacer{} อารมฺมณูปนิสฺสโย, อนนฺตรูปนิสฺสโย, ปกตูปนิสฺสโย ฯเปฯ.\csromanspacer{} ปกตูปนิสฺสโย\csromandash{}สทฺธํ อุปนิสฺสาย ทิฏฺฐิํ คณฺหาติ.\csromanspacer{} สีลํ ฯเปฯ.\csromanspacer{} เสนาสนํ อุปนิสฺสาย ปาณํ หนติ ฯเปฯ สํฆํ ภินฺทติ.\csromanspacer{} สทฺธา ฯเปฯ.\csromanspacer{} เสนาสนํ ราคสฺส ฯเปฯ ปตฺถนาย อุปนิสฺสยปจฺจเยน ปจฺจโย. (2)}

\prose{โนคนฺโถ ธมฺโม คนฺถสฺส จ โนคนฺถสฺส จ ธมฺมสฺส อุปนิสฺสยปจฺจเยน ปจฺจโย.\csromanspacer{} อารมฺมณูปนิสฺสโย, อนนฺตรูปนิสฺสโย, ปกตูปนิสฺสโย ฯเปฯ.\csromanspacer{} ปกตูปนิสฺสโย\csromandash{}สทฺธํ อุปนิสฺสาย มานํ ชปฺเปติ.\csromanspacer{} ทิฏฺฐิํ คณฺหาติ.\csromanspacer{} สีลํ ฯเปฯ.\csromanspacer{} เสนาสนํ อุปนิสฺสาย ปาณํ หนติ ฯเปฯ สํฆํ ภินฺทติ.\csromanspacer{} สทฺธา ฯเปฯ.\csromanspacer{} เสนาสนํ คนฺถานํ สมฺปยุตฺตกานญฺจ ขนฺธานํ อุปนิสฺสยปจฺจเยน ปจฺจโย. (3)}

\prose{คนฺโถ จ โนคนฺโถ จ ธมฺมา คนฺถสฺส ธมฺมสฺส อุปนิสฺสยปจฺจเยน ปจฺจโย.\csromanspacer{} อารมฺมณูปนิสฺสโย, อนนฺตรูปนิสฺสโย, ปกตูปนิสฺสโย ฯเปฯ.\csromanspacer{} ปกตูปนิสฺสโย.\csromanspacer{} ตีณิ. (อารมฺมณนเยน กาตพฺพา.)}

\csromanheader{2}{\textbf{ปุเรชาตาทิ}}
\proseitem{30}{โนคนฺโถ ธมฺโม โนคนฺถสฺส ธมฺมสฺส ปุเรชาตปจฺจเยน ปจฺจโย.\csromanspacer{} \textbf{อารมฺมณปุเรชาตํ}, วตฺถุปุเรชาตํ.\csromanspacer{}\csromanspacer{}\textbf{อารมฺมณปุเรชาตํ}\csromandash{} จกฺขุํ ฯเปฯ วตฺถุํ อนิจฺจโต ฯเปฯ โทมนสฺสํ อุปฺปชฺชติ.\csromanspacer{} ทิพฺเพน จกฺขุนา รูปํ ปสฺสติ.\csromanspacer{} ทิพฺพาย โสตธาตุยา สทฺทํ สุณาติ.\csromanspacer{} รูปายตนํ จกฺขุวิญฺญาณสฺส ฯเปฯ.\csromanspacer{} โผฏฺฐพฺพายตนํ กายวิญฺญาณสฺส ฯเปฯ.\csromanspacer{} \textbf{วตฺถุปุเรชาตํ}\csromandash{}จกฺขายตนํ จกฺขุวิญฺญาณสฺส ฯเปฯ กายายตนํ กายวิญฺญาณสฺส ฯเปฯ.\csromanspacer{} วตฺถุ โนคนฺถานํ ขนฺธานํ ปุเรชาตปจฺจเยน ปจฺจโย. (1)}

\prose{โนคนฺโถ ธมฺโม คนฺถสฺส ธมฺมสฺส ปุเรชาตปจฺจเยน ปจฺจโย.\csromanspacer{} \textbf{อารมฺมณปุเรชาตํ}, วตฺถุปุเรชาตํ.\csromanspacer{}\csromanspacer{}\textbf{อารมฺมณปุเรชาตํ}\csromandash{}จกฺขุํ ฯเปฯ วตฺถุํ อสฺสาเทติ อภินนฺทติ, ตํ อารพฺภ ราโค อุปฺปชฺชติ, ทิฏฺฐิ อุปฺปชฺชติ,}

\vspace{\tipitakaparskip}
\csromancenter{คนฺถทุก โทมนสฺสํ อุปฺปชฺชติ.\csromanspacer{} \textbf{วตฺถุปุเรชาตํ}\csromandash{}วตฺถุ คนฺถานํ ปุเรชาตปจฺจเยน ปจฺจโย. (2)}
\prose{\csromanfolio{263}โนคนฺโถ ธมฺโม คนฺถสฺส จ โนคนฺถสฺส จ ธมฺมสฺส ปุเรชาตปจฺจเยน ปจฺจโย.\csromanspacer{} \textbf{อารมฺมณปุเรชาตํ}, วตฺถุปุเรชาตํ.\csromanspacer{}\csromanspacer{}\textbf{อารมฺมณปุเรชาตํ}\csromandash{}จกฺขุํ ฯเปฯ วตฺถุํ อสฺสาเทติ อภินนฺทติ, ตํ อารพฺภ คนฺถา จ สมฺปยุตฺตกา จ ขนฺธา อุปฺปชฺชนฺติ.\csromanspacer{} \textbf{วตฺถุปุเรชาตํ}\csromandash{} วตฺถุ คนฺถานํ สมฺปยุตฺตกานญฺจ ขนฺธานํ ปุเรชาตปจฺจเยน ปจฺจโย. (3) ปจฺฉาชาตปจฺจเยน ปจฺจโย.\csromanspacer{} ตีณิ.\csromanspacer{} อาเสวนปจฺจเยน ปจฺจโย.\csromanspacer{} นว.}

\csromanheader{2}{\textbf{กมฺม}}
\proseitem{31}{โนคนฺโถ ธมฺโม โนคนฺถสฺส ธมฺมสฺส กมฺมปจฺจเยน ปจฺจโย.\csromanspacer{} \textbf{สหชาตา}, นานากฺขณิกา.\csromanspacer{}\csromanspacer{}\textbf{สหชาตา} โนคนฺถา เจตนา สมฺปยุตฺตกานํ ขนฺธานํ จิตฺตสมุฏฺฐานานญฺจ รูปานํ กมฺมปจฺจเยน ปจฺจโย.\csromanspacer{} ปฏิสนฺธิกฺขเณ ฯเปฯ.\csromanspacer{} \textbf{นานากฺขณิกา} โนคนฺถา เจตนา วิปากานํ ขนฺธานํ กฏตฺตา จ รูปานํ กมฺมปจฺจเยน ปจฺจโย. (1)}

\prose{โนคนฺโถ ธมฺโม คนฺถสฺส ธมฺมสฺส กมฺมปจฺจเยน ปจฺจโย.\csromanspacer{} โนคนฺถา เจตนา สมฺปยุตฺตกานํ คนฺถานํ กมฺมปจฺจเยน ปจฺจโย.}

\prose{โนคนฺโถ ธมฺโม คนฺถสฺส จ โนคนฺถสฺส จ ธมฺมสฺส กมฺมปจฺจเยน ปจฺจโย.\csromanspacer{} โนคนฺถา เจตนา สมฺปยุตฺตกานํ ขนฺธานํ คนฺถานญฺจ จิตฺตสมุฏฺฐานานญฺจ รูปานํ กมฺมปจฺจเยน ปจฺจโย. (3)}

\csromanheader{2}{\textbf{วิปากาทิ}}
\proseitem{32}{โนคนฺโถ ธมฺโม โนคนฺถสฺส ธมฺมสฺส วิปากปจฺจเยน ปจฺจโย.\csromanspacer{} เอกํ.\csromanspacer{} อาหารปจฺจเยน ปจฺจโย.\csromanspacer{} ตีณิ.\csromanspacer{} อินฺทฺริยปจฺจเยน ปจฺจโย.\csromanspacer{} ตีณิ.\csromanspacer{} ฌานปจฺจเยน ปจฺจโย.\csromanspacer{} ตีณิ.\csromanspacer{} มคฺคปจฺจเยน ปจฺจโย.\csromanspacer{} นว.\csromanspacer{} สมฺปยุตฺตปจฺจเยน ปจฺจโย.\csromanspacer{} นว.}

\csromanheader{2}{\textbf{วิปฺปยุตฺต}}
\proseitem{33}{คนฺโถ ธมฺโม โนคนฺถสฺส ธมฺมสฺส วิปฺปยุตฺตปจฺจเยน ปจฺจโย.\csromanspacer{} สหชาตํ, ปจฺฉาชาตํ. (สํขิตฺตํ.) (1)}

\prose{\csromanfolio{264}โนคนฺโถ ธมฺโม โนคนฺถสฺส ธมฺมสฺส วิปฺปยุตฺตปจฺจเยน ปจฺจโย.\csromanspacer{} สหชาตํ, ปุเรชาตํ, ปจฺฉาชาตํ. (สํขิตฺตํ.) (1)}

\prose{โนคนฺโถ ธมฺโม คนฺถสฺส ธมฺมสฺส วิปฺปยุตฺตปจฺจเยน ปจฺจโย.\csromanspacer{} \textbf{ปุเรชาตํ} วตฺถุ คนฺถานํ วิปฺปยุตฺตปจฺจเยน ปจฺจโย. (2)}

\prose{โนคนฺโถ ธมฺโม คนฺถสฺส จ โนคนฺถสฺส จ ธมฺมสฺส วิปฺปยุตฺตปจฺจเยน ปจฺจโย.\csromanspacer{} \textbf{ปุเรชาตํ} วตฺถุ คนฺถานํ สมฺปยุตฺตกานญฺจ ขนฺธานํ วิปฺปยุตฺตปจฺจเยน ปจฺจโย. (3)}

\prose{คนฺโธ จ โนคนฺโถ จ ธมฺมา โนคนฺถสฺส ธมฺมสฺส วิปฺปยุตฺตปจฺจเยน ปจฺจโย.\csromanspacer{} สหชาตํ, ปจฺฉาชาตํ. (สํขิตฺตํ.) (1)}

\csromanheader{2}{\textbf{อตฺถิ}}
\proseitem{34}{คนฺโถ ธมฺโม คนฺถสฺส ธมฺมสฺส อตฺถิปจฺจเยน ปจฺจโย.\csromanspacer{} เอกํ. (ปฏิจฺจสทิสํ.)}

\prose{คนฺโถ ธมฺโม โนคนฺถสฺส ธมฺมสฺส อตฺถิปจฺจเยน ปจฺจโย.\csromanspacer{} สหชาตํ, ปจฺฉาชาตํ.\csromanspacer{}\csromanspacer{}\textbf{สหชาตา} คนฺถา สมฺปยุตฺตกานํ ขนฺธานํ จิตฺตสมุฏฺฐานานญฺจ รูปานํ อตฺถิปจฺจเยน ปจฺจโย.\csromanspacer{} \textbf{ปจฺฉาชาตา} คนฺถา ปุเรชาตสฺส อิมสฺส กายสฺส อตฺถิปจฺจเยน ปจฺจโย. (2)}

\prose{คนฺโถ ธมฺโม คนฺถสฺส จ โนคนฺถสฺส จ ธมฺมสฺส อตฺถิปจฺจเยน ปจฺจโย.\csromanspacer{} เอกํ. (ปฏิจฺจสทิสํ.) (3)}

\proseitem{35}{โนคนฺโถ ธมฺโม โนคนฺถสฺส ธมฺมสฺส อตฺถิปจฺจเยน ปจฺจโย.\csromanspacer{} สหชาตํ, ปุเรชาตํ, ปจฺฉาชาตํ, อาหารํ, อินฺทฺริยํ. (สํขิตฺตํ.) (1)}

\prose{โนคนฺโถ ธมฺโม คนฺถสฺส ธมฺมสฺส อตฺถิปจฺจเยน ปจฺจโย.\csromanspacer{} สหชาตํ, ปุเรชาตํ.\csromanspacer{}\csromanspacer{}\textbf{สหชาตา} โนคนฺถา ขนฺธา คนฺถานํ อตฺถิปจฺจเยน ปจฺจโย.\csromanspacer{} \textbf{ปุเรชาตํ} จกฺขุํ ฯเปฯ วตฺถุํ อสฺสาเทติ อภินนฺทติ, ตํ อารพฺภ ราโค อุปฺปชฺชติ, ทิฏฺฐิ อุปฺปชฺชติ, โทมนสฺสํ อุปฺปชฺชติ.\csromanspacer{} วตฺถุ คนฺถานํ อตฺถิปจฺจเยน ปจฺจโย. (2)}

\csromanheader{2}{26. คนฺถทุก}
\prose{\csromanfolio{265}โนคนฺโถ ธมฺโม คนฺถสฺส จ โนคนฺถสฺส จ ธมฺมสฺส อตฺถิปจฺจเยน ปจฺจโย.\csromanspacer{} สหชาตํ, ปุเรชาตํ.\csromanspacer{}\csromanspacer{}\textbf{สหชาโต} โนคนฺโถ เอโก ขนฺโธ ติณฺณนฺนํ ขนฺธานํ คนฺถานญฺจ จิตฺตสมุฏฺฐานานญฺจ รูปานํ อตฺถิปจฺจเยน ปจฺจโย ฯเปฯ.\csromanspacer{} \textbf{ปุเรชาตํ} จกฺขุํ ฯเปฯ วตฺถุํ อสฺสาเทติ อภินนฺทติ, ตํ อารพฺภ คนฺถา จ สมฺปยุตฺตกา จ ขนฺธา อุปฺปชฺชนฺติ.\csromanspacer{} วตฺถุ คนฺถานํ สมฺปยุตฺตกานญฺจ ขนฺธานํ อตฺถิปจฺจเยน ปจฺจโย. (3)}

\proseitem{36}{คนฺโถ จ โนคนฺโถ จ ธมฺมา คนฺถสฺส ธมฺมสฺส อตฺถิปจฺจเยน ปจฺจโย.\csromanspacer{} สหชาตํ, ปุเรชาตํ.\csromanspacer{}\csromanspacer{}\textbf{สหชาโต} สีลพฺพตปรามาโส กายคนฺโถ จ สมฺปยุตฺตกา จ ขนฺธา อภิชฺโฌกายคนฺถสฺส อตฺถิปจฺจเยน ปจฺจโย. (จกฺกํ.) \textbf{สหชาโต} สีลพฺพตปรามาโส กายคนฺโถ จ วตฺถุ จ อภิชฺฌากายคนฺถสฺส อตฺถิปจฺจเยน ปจฺจโย. (จกฺกํ.) (1)}

\prose{คนฺโถ จ โนคนฺโถ จ ธมฺมา โนคนฺถสฺส ธมฺมสฺส อตฺถิปจฺจเยน ปจฺจโย.\csromanspacer{} สหชาตํ, ปุเรชาตํ, ปจฺฉาชาตํ, อาหารํ, อินฺทฺริยํ.\csromanspacer{}\csromanspacer{}\textbf{สหชาโต} โนคนฺโถ เอโก ขนฺโธ จ คนฺโถ จ ติณฺณนฺนํ ขนฺธานํ จิตฺตสมุฏฺฐานานญฺจ รูปานํ อตฺถิปจฺจเยน ปจฺจโย ฯเปฯ เทฺว ขนฺธา จ ฯเปฯ.\csromanspacer{} \textbf{สหชาตา} คนฺถา จ วตฺถุ จ โนคนฺถานํ ขนฺธานํ อตฺถิปจฺจเยน ปจฺจโย.\csromanspacer{} \textbf{สหชาตา} คนฺถา จ มหาภูตา จ จิตฺตสมุฏฺฐานานญฺจ รูปานํ อตฺถิปจฺจเยน ปจฺจโย.\csromanspacer{} \textbf{ปจฺฉาชาตา} คนฺถา จ สมฺปยุตฺตกา จ ขนฺธา ปุเรชาตสฺส อิมสฺส กายสฺส อตฺถิปจฺจเยน ปจฺจโย.\csromanspacer{} \textbf{ปจฺฉาชาตา} คนฺถา จ กพฬีกาโร อาหาโร จ อิมสฺส กายสฺส อตฺถิปจฺจเยน ปจฺจโย.\csromanspacer{} \textbf{ปจฺฉาชาตา} คนฺถา จ รูปชีวิตินฺทฺริยญฺจ กฏตฺตารูปานํ อตฺถิปจฺจเยน ปจฺจโย. (2)}

\prose{คนฺโถ จ โนคนฺโถ จ ธมฺมา คนฺถสฺส จ โนคนฺถสฺส จ ธมฺมสฺส อตฺถิปจฺจเยน ปจฺจโย.\csromanspacer{} สหชาตํ, ปุเรชาตํ.\csromanspacer{}\csromanspacer{}\textbf{สหชาโต} โนคนฺโถ เอโก ขนฺโธ จ สีลพฺพตปรามาโส กายคนฺโถ จ ติณฺณนฺนํ ขนฺธานํ อภิชฺฌากายคนฺถสฺส จ จิตฺตสมุฏฺฐานานญฺจ รูปานํ อตฺถิปจฺจเยน ปจฺจโย ฯเปฯ.\csromanspacer{} \textbf{สหชาโต} สีลพฺพตปรามาโส กายคนฺโถ จ วตฺถุ จ อภิชฺฌากายคนฺถสฺส จ สมฺปยุตฺตกานญฺจ ขนฺธานํ อตฺถิปจฺจเยน ปจฺจโย. (จกฺกํ.) (3)}

\csromanheader{2}{1. ปจฺจยานุโลม 2. สงฺขฺยาวาร}
\csromanheader{2}{\textbf{สุทฺธ}}
\proseitem{37}{\csromanfolio{266}เหตุยา นว, อารมฺมเณ นว, (สพฺพตฺถ นว,) อุปนิสฺสเย นว, ปุเรชาเต ตีณิ, ปจฺฉาชาเต ตีณิ, อาเสวเน นว, กมฺเม ตีณิ, วิปาเก เอกํ, อาหาเร ตีณิ, อินฺทฺริเย ตีณิ, ฌาเน ตีณิ, มคฺเค นว, สมฺปยุตฺเต นว, วิปฺปยุตฺเต ปญฺจ, อตฺถิยา นว, นตฺถิยา นว, วิคเต นว, อวิคเต นว.}

\vspace{\tipitakaparskip}
\csromancenter{อนุโลมํ.}
\csromanheader{2}{2. ปจฺจนียุทฺธาร}
\proseitem{38}{คนฺโถ ธมฺโม คนฺถสฺส ธมฺมสฺส อารมฺมณปจฺจเยน ปจฺจโย.\csromanspacer{} สหชาตปจฺจเยน ปจฺจโย.\csromanspacer{} อุปนิสฺสยปจฺจเยน ปจฺจโย. (1)}

\prose{คนฺโถ ธมฺโม โนคนฺถสฺส ธมฺมสฺส อารมฺมณปจฺจเยน ปจฺจโย.\csromanspacer{} สหชาตปจฺจเยน ปจฺจโย.\csromanspacer{} อุปนิสฺสยปจฺจเยน ปจฺจโย.\csromanspacer{} ปจฺฉาชาตปจฺจเยน ปจฺจโย. (2)}

\prose{คนฺโถ ธมฺโม คนฺถสฺส จ โนคนฺถสฺส จ ธมฺมสฺส อารมฺมณปจฺจเยน ปจฺจโย.\csromanspacer{} สหชาตปจฺจเยน ปจฺจโย.\csromanspacer{} อุปนิสฺสยปจฺจเยน ปจฺจโย. (3)}

\proseitem{39}{โนคนฺโถ ธมฺโม โนคนฺถสฺส ธมฺมสฺส อารมฺมณปจฺจเยน ปจฺจโย.\csromanspacer{} สหชาตปจฺจเยน ปจฺจโย.\csromanspacer{} อุปนิสฺสยปจฺจเยน ปจฺจโย.\csromanspacer{} ปุเรชาตปจฺจเยน ปจฺจโย.\csromanspacer{} ปจฺฉาชาตปจฺจเยน ปจฺจโย.\csromanspacer{} กมฺมปจฺจเยน ปจฺจโย.\csromanspacer{} อาหารปจฺจเยน ปจฺจโย.\csromanspacer{} อินฺทฺริยปจฺจเยน ปจฺจโย. (1)}

\prose{โนคนฺโถ ธมฺโม คนฺถสฺส ธมฺมสฺส อารมฺมณปจฺจเยน ปจฺจโย.\csromanspacer{} สหชาตปจฺจเยน ปจฺจโย.\csromanspacer{} อุปนิสฺสยปจฺจเยน ปจฺจโย.\csromanspacer{} ปุเรชาตปจฺจเยน ปจฺจโย. (2)}

\prose{โนคนฺโถ ธมฺโม คนฺถสฺส จ โนคนฺถสฺส จ ธมฺมสฺส อารมฺมณปจฺจเยน ปจฺจโย.\csromanspacer{} สหชาตปจฺจเยน ปจฺจโย.\csromanspacer{} อุปนิสฺสยปจฺจเยน ปจฺจโย.\csromanspacer{} ปุเรชาตปจฺจเยน ปจฺจโย. (3)}

\csromanmatikaanchor{mtk.33}
\csromantocmarkref{section}{27. คนฺถนิยทุก}{mtk.33}
\bookmark[dest={mtk.33},level=1]{27. คนฺถนิยทุก}
\csromanheader{1}{27. คนฺถนิยทุก}
\proseitem{40}{\csromanfolio{267}คนฺโถ จ โนคนฺโถ จ ธมฺมา คนฺถสฺส ธมฺมสฺส อารมฺมณปจฺจเยน ปจฺจโย.\csromanspacer{} สหชาตปจฺจเยน ปจฺจโย.\csromanspacer{} อุปนิสฺสยปจฺจเยน ปจฺจโย. (1)}

\prose{คนฺโถ จ โนคนฺโถ จ ธมฺมา โนคนฺถสฺส ธมฺมสฺส อารมฺมณปจฺจเยน ปจฺจโย.\csromanspacer{} สหชาตปจฺจเยน ปจฺจโย.\csromanspacer{} อุปนิสฺสยปจฺจเยน ปจฺจโย.\csromanspacer{} ปจฺฉาชาตปจฺจเยน ปจฺจโย. (2)}

\prose{คนฺโถ จ โนคนฺโถ จ ธมฺมา คนฺถสฺส จ โนคนฺถสฺส จ ธมฺมสฺส อารมฺมณปจฺจเยน ปจฺจโย.\csromanspacer{} สหชาตปจฺจเยน ปจฺจโย.\csromanspacer{} อุปนิสฺสยปจฺจเยน ปจฺจโย. (3)}

\csromanheader{2}{2. ปจฺจยปจฺจนีย 2. สงฺขฺยาวาร}
\vspace{\tipitakaparskip}
\csromancenter{นเหตุยา นว, น-อารมฺมเณ นว, (สพฺพตฺถ นว,) โน-อวิคเต นว.}
\csromanheader{2}{3. ปจฺจยานุโลมปจฺจนีย}
\csromanheader{2}{\textbf{เหตุทุก}}
\proseitem{42}{เหตุปจฺจยา น-อารมฺมเณ นว ฯเปฯ นสมนนฺตเร นว, น-อญฺญมญฺเญ ตีณิ, น-อุปนิสฺสเย นว, (สพฺพตฺถ นว,) นมคฺเค นว, นสมฺปยุตฺเต ตีณิ, นวิปฺปยุตฺเต นว, โนนตฺถิยา นว, โนวิคเต นว.}

\csromanheader{2}{4. ปจฺจยปจฺจนียานุโลม}
\proseitem{43}{นเหตุปจฺจยา อารมฺมเณ นว, อธิปติยา นว, (อนุโลมปทานิ ปริปุณฺณานิ กาตพฺพานิ.) อวิคเต นว.}

\nitthitam{คนฺถทุกํ นิฏฺฐิตํ.}
\csromanheader{2}{27. คนฺถนิยทุก 1-7. วารสตฺตก}
\proseitem{44}{คนฺถนิยํ ธมฺมํ ปฏิจฺจ คนฺถนิโย ธมฺโม อุปฺปชฺชติ เหตุปจฺจยา.\csromanspacer{} คนฺถนิยํ เอกํ ขนฺธํ. (สํขิตฺตํ.)}

\prose{(ยถา จูฬนฺตรทุเก โลกิยทุกํ, เอวํ วิภชิตพฺพํ.\csromanspacer{} นินฺนานากรณํ.)}

\nitthitam{คนฺถนิยทุกํ นิฏฺฐิตํ.}
\csromanmatikaanchor{mtk.34}
\csromantocmarkref{section}{28. คนฺถสมฺปยุตฺตทุก}{mtk.34}
\bookmark[dest={mtk.34},level=1]{28. คนฺถสมฺปยุตฺตทุก}
\csromanheader{1}{28. คนฺถสมฺปยุตฺตทุก 1. ปฏิจฺจวาร}
\csromanheader{2}{1. ปจฺจยานุโลม 1. วิภงฺควาร}
\csromanheader{2}{\textbf{เหตุ}}
\proseitem{45}{\csromanfolio{268}คนฺถสมฺปยุตฺตํ ธมฺมํ ปฏิจฺจ คนฺถสมฺปยุตฺโต ธมฺโม อุปฺปชฺชติ เหตุปจฺจยา คนฺถสมฺปยุตฺตํ เอกํ ขนฺธํ ปฏิจฺจ ตโย ขนฺธา ฯเปฯ เทฺว ขนฺเธ ฯเปฯ. (1)}

\prose{คนฺถสมฺปยุตฺตํ ธมฺมํ ปฏิจฺจ คนฺถวิปฺปยุตฺโต ธมฺโม อุปฺปชฺชติ เหตุปจฺจยา.\csromanspacer{} คนฺถสมฺปยุตฺเต ขนฺเธ ปฏิจฺจ จิตฺตสมุฏฺฐานํ รูปํ.\csromanspacer{} ทิฏฺฐิคตวิปฺปยุตฺตโลภสหคเต ขนฺเธ ปฏิจฺจ โลโภ จิตฺตสมุฏฺฐานญฺจ รูปํ.\csromanspacer{} โทมนสฺสสหคเต ขนฺเธ ปฏิจฺจ ปฏิฆํ จิตฺตสมุฏฺฐานญฺจ รูปํ. (2)}

\prose{คนฺถสมฺปยุตฺตํ ธมฺมํ ปฏิจฺจ คนฺถสมฺปยุตฺโต จ คนฺถวิปฺปยุตฺโต จ ธมฺมา อุปฺปชฺชนฺติ เหตุปจฺจยา.\csromanspacer{} คนฺถสมฺปยุตฺตํ เอกํ ขนฺธํ ปฏิจฺจ ตโย ขนฺธา จิตฺตสมุฏฺฐานญฺจ รูปํ ฯเปฯ เทฺว ขนฺเธ ฯเปฯ.\csromanspacer{} ทิฏฺฐิคตวิปฺปยุตฺตโลภสหคตํ เอกํ ขนฺธํ ปฏิจฺจ ตโย ขนฺธา โลโภ จ จิตฺตสมุฏฺฐานญฺจ รูปํ ฯเปฯ เทฺว ขนฺเธ ฯเปฯ.\csromanspacer{} โทมนสฺสสหคตํ เอกํ ขนฺธํ ปฏิจฺจ ตโย ขนฺธา ปฏิฆญฺจ จิตฺตสมุฏฺฐานญฺจ รูปํ ฯเปฯ เทฺว ขนฺเธ ฯเปฯ. (3)}

\proseitem{46}{คนฺถวิปฺปยุตฺตํ ธมฺมํ ปฏิจฺจ คนฺถวิปฺปยุตฺโต ธมฺโม อุปฺปชฺชติ เหตุปจฺจยา.\csromanspacer{} คนฺถวิปฺปยุตฺตํ เอกํ ขนฺธํ ปฏิจฺจ ตโย ขนฺธา จิตฺตสมุฏฺฐานญฺจ รูปํ ฯเปฯ เทฺว ขนฺเธ ฯเปฯ.\csromanspacer{} ทิฏฺฐิคตวิปฺปยุตฺตํ โลภํ ปฏิจฺจ จิตฺตสมุฏฺฐานํ รูปํ.\csromanspacer{} ปฏิฆํ ปฏิจฺจ จิตฺตสมุฏฺฐานํ รูปํ.\csromanspacer{} ปฏิสนฺธิกฺขเณ ฯเปฯ ขนฺเธ ปฏิจฺจ. (สํขิตฺตํ.) (1)}

\prose{คนฺถวิปฺปยุตฺตํ ธมฺมํ ปฏิจฺจ คนฺถสมฺปยุตฺโต ธมฺโม อุปฺปชฺชติ เหตุปจฺจยา.\csromanspacer{} ทิฏฺฐิคตวิปฺปยุตฺตํ โลภํ ปฏิจฺจ สมฺปยุตฺตกา ขนฺธา.\csromanspacer{} ปฏิฆํ ปฏิจฺจ สมฺปยุตฺตกา ขนฺธา. (2)}

\prose{คนฺถวิปฺปยุตฺตํ ธมฺมํ ปฏิจฺจ คนฺถสมฺปยุตฺโต จ คนฺถวิปฺปยุตฺโต จ ธมฺมา อุปฺปชฺชนฺติ เหตุปจฺจยา.\csromanspacer{} ทิฏฺฐคตวิปฺปยุตฺตํ โลภํ ปฏิจฺจ สมฺปยุตฺตกา ขนฺธา จิตฺตสมุฏฺฐานญฺจ รูปํ.\csromanspacer{} ปฏิฆํ ปฏิจฺจ สมฺปยุตฺตกา ขนฺธา จิตฺตสมุฏฺฐานญฺจ รูปํ. (3)}

\proseitem{47}{คนฺถสมฺปยุตฺตญฺจ คนฺถวิปฺปยุตฺตญฺจ ธมฺมํ ปฏิจฺจ คนฺถสมฺปยุตฺโต ธมฺโม อุปฺปชฺชติ เหตุปจฺจยา.\csromanspacer{} ทิฏฺฐิคตวิปฺปยุตฺตโลภสหคตํ เอกํ ขนฺธญฺจ}

\vspace{\tipitakaparskip}
\csromancenter{คนฺถสมฺปยุตฺตทุก โลภญฺจ ปฏิจฺจ ตโย ขนฺธา ฯเปฯ เทฺว ขนฺเธ ฯเปฯ.\csromanspacer{} โทมนสฺสสหคตํ เอกํ ขนฺธญฺจ ปฏิฆญฺจ ปฏิจฺจ ตโย ขนฺธา ฯเปฯ เทฺว ขนฺเธ ฯเปฯ. (1)}
\prose{\csromanfolio{269}คนฺถสมฺปยุตฺตญฺจ คนฺถวิปฺปยุตฺตญฺจ ธมฺมํ ปฏิจฺจ คนฺถวิปฺปยุตฺโต ธมฺโม อุปฺปชฺชติ เหตุปจฺจยา.\csromanspacer{} คนฺถสมฺปยุตฺเต ขนฺเธ จ มหาภูเต จ ปฏิจฺจ จิตฺตสมุฏฺฐานํ รูปํ.\csromanspacer{} ทิฏฺฐิคตวิปฺปยุตฺตโลภสหคเต ขนฺเธ จ โลภญฺจ ปฏิจฺจ จิตฺตสมุฏฺฐานํ รูปํ.\csromanspacer{} โทมนสฺสสหคเต ขนฺเธ จ ปฏิฆญฺจ ปฏิจฺจ จิตฺตสมุฏฺฐานํ รูปํ. (2)}

\prose{คนฺถสมฺปยุตฺตญฺจ คนฺถวิปฺปยุตฺตญฺจ ธมฺมํ ปฏิจฺจ คนฺถสมฺปยุตฺโต จ คนฺถวิปฺปยุตฺโต จ ธมฺมา อุปฺปชฺชนฺติ เหตุปจฺจยา.\csromanspacer{} ทิฏฺฐิคตวิปฺปยุตฺตโลภสหคตํ เอกํ ขนฺธญฺจ โลภญฺจ ปฏิจฺจ ตโย ขนฺธา จิตฺตสมุฏฺฐานญฺจ รูปํ ฯเปฯ เทฺว ขนฺเธ ฯเปฯ.\csromanspacer{} โทมนสฺสสหคตํ เอกํ ขนฺธญฺจ ปฏิฆญฺจ ปฏิจฺจ ตโย ขนฺธา จิตฺตสมุฏฺฐานญฺจ รูปํ ฯเปฯ เทฺว ขนฺเธ ฯเปฯ. (3)}

\csromanheader{2}{\textbf{อารมฺมณ}}
\proseitem{48}{คนฺถสมฺปยุตฺตํ ธมฺมํ ปฏิจฺจ คนฺถสมฺปยุตฺโต ธมฺโม อุปฺปชฺชติ อารมฺมณปจฺจยา.\csromanspacer{} คนฺถสมฺปยุตฺตํ เอกํ ขนฺธํ ปฏิจฺจ ตโย ขนฺธา ฯเปฯ เทฺว ขนฺเธ ฯเปฯ. (1)}

\prose{คนฺถสมฺปยุตฺตํ ธมฺมํ ปฏิจฺจ คนฺถวิปฺปยุตฺโต ธมฺโม อุปฺปชฺชติ อารมฺมณปจฺจยา.\csromanspacer{} ทิฏฺฐิคตวิปฺปยุตฺตโลภสหคเต ขนฺเธ ปฏิจฺจ โลโภ.\csromanspacer{} โทมนสฺสสหคเต ขนฺเธ ปฏิจฺจ ปฏิฆํ. (2)}

\prose{คนฺถสมฺปยุตฺตํ ธมฺมํ ปฏิจฺจ คนฺถสมฺปยุตฺโต จ คนฺถวิปฺปยุตฺโต จ ธมฺมา อุปฺปชฺชนฺติ อารมฺมณปจฺจยา.\csromanspacer{} ทิฏฺฐิคตวิปฺปยุตฺตโลภสหคตํ เอกํ ขนฺธํ ปฏิจฺจ ตโย ขนฺธา โลโภ จ ฯเปฯ เทฺว ขนฺเธ ฯเปฯ.\csromanspacer{} โทมนสฺสสหคตํ เอกํ ขนฺธํ ปฏิจฺจ ตโย ขนฺธา ปฏิฆญฺจ ฯเปฯ เทฺว ขนฺเธ ฯเปฯ. (3)}

\proseitem{49}{คนฺถวิปฺปยุตฺตํ ธมฺมํ ปฏิจฺจ คนฺถวิปฺปยุตฺโต ธมฺโม อุปฺปชฺชติ อารมฺมณปจฺจยา.\csromanspacer{} คนฺถวิปฺปยุตฺตํ เอกํ ขนฺธํ ปฏิจฺจ ตโย ขนฺธา ฯเปฯ เทฺว ขนฺเธ ฯเปฯ.\csromanspacer{} ปฏิสนฺธิกฺขเณ ฯเปฯ วตฺถุํ ปฏิจฺจ ขนฺธา. (1)}

\prose{คนฺถวิปฺปยุตฺตํ ธมฺมํ ปฏิจฺจ คนฺถสมฺปยุตฺโต ธมฺโม อุปฺปชฺชติ อารมฺมณปจฺจยา.\csromanspacer{} ทิฏฺฐิคตวิปฺปยุตฺตํ โลภํ ปฏิจฺจ สมฺปยุตฺตกา ขนฺธา.\csromanspacer{} ปฏิฆํ ปฏิจฺจ สมฺปยุตฺตกา ขนฺธา. (2)}

\prose{\csromanfolio{270}คนฺถสมฺปยุตฺตญฺจ คนฺถวิปฺปยุตฺตญฺจ ธมฺมํ ปฏิจฺจ คนฺถสมฺปยุตฺโต ธมฺโม อุปฺปชฺชติ อารมฺมณปจฺจยา.\csromanspacer{} ทิฏฺฐิคตวิปฺปยุตฺตโลภสหคตํ เอกํ ขนฺธญฺจ โลภญฺจ ปฏิจฺจ ตโย ขนฺธา ฯเปฯ เทฺว ขนฺเธ ฯเปฯ.\csromanspacer{} โทมนสฺสสหคตํ เอกํ ขนฺธญฺจ ปฏิฆญฺจ ปฏิจฺจ ตโย ขนฺธา ฯเปฯ เทฺว ขนฺเธ ฯเปฯ. (สํขิตฺตํ.) (1)}

\csromanheader{2}{1. ปจฺจยานุโลม 2. สงฺขฺยาวาร}
\csromanheader{2}{\textbf{สุทฺธ}}
\proseitem{50}{เหตุยา นว, อารมฺมเณ ฉ, อธิปติยา นว, อนนฺตเร ฉ, สมนนฺตเร ฉ, สหชาเต นว, อญฺญมญฺเญ ฉ, นิสฺสเย นว, อุปนิสฺสเย ฉ, ปุเรชาเต ฉ, อาเสวเน ฉ, กมฺเม นว, วิปาเก เอกํ, อาหาเร นว, อินฺทฺริเย นว, ฌาเน นว, มคฺเค นว, สมฺปยุตฺเต ฉ, วิปฺปยุตฺเต นว, อตฺถิยา นว, นตฺถิยา ฉ, วิคเต ฉ, อวิคเต นว.}

\vspace{\tipitakaparskip}
\csromancenter{อนุโลมํ.}
\csromanheader{2}{2. ปจฺจยปจฺจนีย 1. วิภงฺควาร}
\csromanheader{2}{\textbf{นเหตุ}}
\proseitem{51}{คนฺถวิปฺปยุตฺตํ ธมฺมํ ปฏิจฺจ คนฺถวิปฺปยุตฺโต ธมฺโม อุปฺปชฺชติ นเหตุปจฺจยา อเหตุกํ คนฺถวิปฺปยุตฺตํ เอกํ ขนฺธํ ปฏิจฺจ ตโย ขนฺธา จิตฺตสมุฏฺฐานญฺจ รูปํ ฯเปฯ เทฺว ขนฺเธ ฯเปฯ.\csromanspacer{} อเหตุกปฏิสนฺธิกฺขเณ ฯเปฯ. (ยาว อสญฺญสตฺตา.) วิจิกิจฺฉาสหคเต อุทฺธจฺจสหคเต ขนฺเธ ปฏิจฺจ วิจิกิจฺฉาสหคโต อุทฺธจฺจสหคโต โมโห. (สํขิตฺตํ.) (1)}

\csromanheader{2}{2. ปจฺจยปจฺจนีย 2. สงฺขฺยาวาร}
\csromanheader{2}{\textbf{สุทฺธ}}
\vspace{\tipitakaparskip}
\csromancenter{\textbf{นเหตุ}ยา เอกํ, น-อารมฺมเณ ตีณิ, น-อธิปติยา นว, น-อนนฺตเร ตีณิ, นสมนนฺตเร ตีณิ, น-อญฺญมญฺเญ ตีณิ, น-อุปนิสฺสเย ตีณิ, นปุเรชาเต สตฺต. (นปุเรชาเต วิภชนฺเตน อรูปํ ปฐมํ}
\csromancenter{คนฺถสมฺปยุตฺตทุก กาตพฺพํ.\csromanspacer{} รูปํ ยตฺถ ลพฺภติ, ปจฺฉา กาตพฺพํ.\csromanspacer{} ปฏิฆญฺจ อรูเป นตฺถิ.) นปจฺฉาชาเต นว, น-อาเสวเน นว, นกมฺเม จตฺตาริ, นวิปาเก นว, น-อาหาเร เอกํ, น-อินฺทฺริเย เอกํ, นฌาเน เอกํ, นมคฺเค เอกํ, นสมฺปยุตฺเต ตีณิ, นวิปฺปยุตฺเต ฉ, โนนตฺถิยา ตีณิ, โนวิคเต ตีณิ.}
\csromanheader{2}{3. ปจฺจยานุโลมปจฺจนีย}
\proseitem{53}{\csromanfolio{271}\textbf{เหตุ}ปจฺจยา น-อารมฺมเณ ตีณิ, น-อธิปติยา นว, (เอวํ คเณตพฺพํ.\csromanspacer{} สํขิตฺตํ.) โนวิคเต ตีณิ.}

\csromanheader{2}{4. ปจฺจยปจฺจนียานุโลม}
\vspace{\tipitakaparskip}
\csromancenter{นเหตุปจฺจยา อารมฺมเณ เอกํ ฯเปฯ อวิคเต เอกํ.}
\csromancenter{(สหชาตวาโรปิ เอวํ กาตพฺโพ.)}
\csromanheader{2}{28. คนฺถสมฺปยุตฺตทุก 3. ปจฺจยวาร}
\csromanheader{2}{1. ปจฺจยานุโลม 1. วิภงฺควาร}
\csromanheader{2}{\textbf{เหตุ}}
\proseitem{55}{คนฺถสมฺปยุตฺตํ ธมฺมํ ปจฺจยา คนฺถสมฺปยุตฺโต ธมฺโม อุปฺปชฺชติ เหตุปจฺจยา.\csromanspacer{} ตีณิ. (ปฏิจฺจสทิสา.)}

\prose{คนฺถวิปฺปยุตฺตํ ธมฺมํ ปจฺจยา คนฺถวิปฺปยุตฺโต ธมฺโม อุปฺปชฺชติ เหตุปจฺจยา.\csromanspacer{} คนฺถวิปฺปยุตฺตํ เอกํ ขนฺธํ ปจฺจยา ฯเปฯ.\csromanspacer{} ปฏิสนฺธิกฺขเณ ฯเปฯ.\csromanspacer{} ขนฺเธ ปจฺจยา วตฺถุ, วตฺถุํ ปจฺจยา ขนฺธา.\csromanspacer{} เอกํ มหาภูตํ ฯเปฯ วตฺถุํ ปจฺจยา คนฺถวิปฺปยุตฺตา ขนฺธา. (1)}

\prose{\csromanfolio{272}คนฺถวิปฺปยุตฺตํ ธมฺมํ ปจฺจยา คนฺถสมฺปยุตฺโต ธมฺโม อุปฺปชฺชติ เหตุปจฺจยา.\csromanspacer{} วตฺถุํ ปจฺจยา คนฺถสมฺปยุตฺตกา ขนฺธา.\csromanspacer{} ทิฏฺฐิคตวิปฺปยุตฺตํ โลภํ ปจฺจยา สมฺปยุตฺตกา ขนฺธา.\csromanspacer{} ปฏิฆํ ปจฺจยา สมฺปยุตฺตกา ขนฺธา. (2)}

\prose{คนฺถวิปฺปยุตฺตํ ธมฺมํ ปจฺจยา คนฺถสมฺปยุตฺโต จ คนฺถวิปฺปยุตฺโต จ ธมฺมา อุปฺปชฺชนฺติ เหตุปจฺจยา.\csromanspacer{} วตฺถุํ ปจฺจยา คนฺถสมฺปยุตฺตกา ขนฺธา.\csromanspacer{} มหาภูเต ปจฺจยา จิตฺตสมุฏฺฐานํ รูปํ.\csromanspacer{} ทิฏฺฐิคตวิปฺปยุตฺตํ โลภํ ปจฺจยา สมฺปยุตฺตกา ขนฺธา จิตฺตสมุฏฺฐานญฺจ รูปํ.\csromanspacer{} ปฏิฆํ ปจฺจยา สมฺปยุตฺตกา ขนฺธา จิตฺตสมุฏฺฐานญฺจ รูปํ.\csromanspacer{} วตฺถุํ ปจฺจยา ทิฏฺฐิคตวิปฺปยุตฺตโลภสหคตา ขนฺธา จ โลโภ จ.\csromanspacer{} วตฺถุํ ปจฺจยา โทมนสฺสสหคตา ขนฺธา จ ปฏิฆญฺจ. (3)}

\proseitem{56}{คนฺถสมฺปยุตฺตญฺจ คนฺถวิปฺปยุตฺตญฺจ ธมฺมํ ปจฺจยา คนฺถสมฺปยุตฺโต ธมฺโม อุปฺปชฺชติ เหตุปจฺจยา.\csromanspacer{} คนฺถสมฺปยุตฺตํ เอกํ ขนฺธญฺจ วตฺถุญฺจ ปจฺจยา ตโย ขนฺธา ฯเปฯ เทฺว ขนฺเธ ฯเปฯ.\csromanspacer{} ทิฏฺฐิคตวิปฺปยุตฺตโลภสหคตํ เอกํ ขนฺธญฺจ วตฺถุญฺจ โลภญฺจ ปจฺจยา ตโย ขนฺธา ฯเปฯ เทฺว ขนฺเธ ฯเปฯ.\csromanspacer{} โทมนสฺสสหคตํ เอกํ ขนฺธญฺจ วตฺถุญฺจ ปฏิฆญฺจ ปจฺจยา ตโย ขนฺธา ฯเปฯ เทฺว ขนฺเธ ฯเปฯ. (1)}

\prose{คนฺถสมฺปยุตฺตญฺจ คนฺถวิปฺปยุตฺตญฺจ ธมฺมํ ปจฺจยา คนฺถวิปฺปยุตฺโต ธมฺโม อุปฺปชฺชติ เหตุปจฺจยา.\csromanspacer{} คนฺถสมฺปยุตฺเต ขนฺเธ จ มหาภูเต จ ปจฺจยา จิตฺตสมุฏฺฐานํ รูปํ.\csromanspacer{} ทิฏฺฐิคตวิปฺปยุตฺตโลภสหคเต ขนฺเธ จ โลภญฺจ ปจฺจยา จิตฺตสมุฏฺฐานํ รูปํ.\csromanspacer{} โทมนสฺสสหคเต ขนฺเธ จ ปฏิฆญฺจ ปจฺจยา จิตฺตสมุฏฺฐานํ รูปํ.\csromanspacer{} ทิฏฺฐิคตวิปฺปยุตฺตโลภสหคเต ขนฺเธ จ วตฺถุญฺจ ปจฺจยา โลโภ.\csromanspacer{} โทมนสฺสสหคเต ขนฺเธ จ วตฺถุญฺจ ปจฺจยา ปฏิฆํ. (2)}

\prose{คนฺถสมฺปยุตฺตญฺจ คนฺถวิปฺปยุตฺตญฺจ ธมฺมํ ปจฺจยา คนฺถสมฺปยุตฺโต จ คนฺถวิปฺปยุตฺโต จ ธมฺมา อุปฺปชฺชนฺติ เหตุปจฺจยา.\csromanspacer{} คนฺถสมฺปยุตฺตํ เอกํ ขนฺธญฺจ วตฺถุญฺจ ปจฺจยา ตโย ขนฺธา ฯเปฯ เทฺว ขนฺเธ ฯเปฯ.\csromanspacer{} คนฺถสมฺปยุตฺเต ขนฺเธ จ มหาภูเต จ ปจฺจยา จิตฺตสมุฏฺฐานํ รูปํ.\csromanspacer{} ทิฏฺฐิคตวิปฺปยุตฺตโลภสหคตํ เอกํ ขนฺธญฺจ โลภญฺจ ปจฺจยา ตโย ขนฺธา จิตฺตสมุฏฺฐานญฺจ รูปํ ฯเปฯ เทฺว ขนฺเธ ฯเปฯ.\csromanspacer{} โทมนสฺสสหคตํ เอกํ ขนฺธญฺจ ปฏิฆญฺจ ปจฺจยา ตโย ขนฺธา จิตฺตสมุฏฺฐานญฺจ รูปํ ฯเปฯ เทฺว ขนฺเธ ฯเปฯ.\csromanspacer{} ทิฏฺฐิคตวิปฺปยุตฺตโลภสหคตํ เอกํ ขนฺธญฺจ วตฺถุญฺจ ปจฺจยา ตโย ขนฺธา โลโภ จ ฯเปฯ เทฺว ขนฺเธ ฯเปฯ.\csromanspacer{} โทมนสฺสสหคตํ เอกํ ขนฺธญฺจ วตฺถุญฺจ ปจฺจยา ตโย ขนฺธา ปฏิฆญฺจ ฯเปฯ เทฺว ขนฺเธ ฯเปฯ. (สํขิตฺตํ.) (3)}

\csromanheader{2}{28. คนฺถสมฺปยุตฺตทุก \textbf{1. ปจฺจยานุโลม 2. สงฺขฺยาวาร}}
\csromanheader{2}{\textbf{สุทฺธ}}
\proseitem{57}{\csromanfolio{273}เหตุยา นว, อารมฺมเณ นว, อธิปติยา นว, อนนฺตเร นว, สมนนฺตเร นว, (สพฺพตฺถ นว.) วิปาเก เอกํ, อาหาเร นว ฯเปฯ อวิคเต นว.}

\vspace{\tipitakaparskip}
\csromancenter{อนุโลมํ.}
\csromanheader{2}{2. ปจฺจยปจฺจนีย 1. วิภงฺควาร}
\csromanheader{2}{\textbf{นเหตุ}}
\proseitem{58}{คนฺถวิปฺปยุตฺตํ ธมฺมํ ปจฺจยา คนฺถวิปฺปยุตฺโต ธมฺโม อุปฺปชฺชติ นเหตุปจฺจยา.\csromanspacer{} อเหตุกํ คนฺถวิปฺปยุตฺตํ ฯเปฯ.\csromanspacer{} อเหตุกปฏิสนฺธิกฺขเณ. (ยาว อสญฺญสตฺตา.) จกฺขายตนํ ปจฺจยา จกฺขุวิญฺญาณํ ฯเปฯ.\csromanspacer{} กายายตนํ ปจฺจยา กายวิญฺญาณํ.\csromanspacer{} วตฺถุํ ปจฺจยา อเหตุกา คนฺถวิปฺปยุตฺตา ขนฺธา.\csromanspacer{} วิจิกิจฺฉาสหคเต อุทฺธจฺจสหคเต ขนฺเธ จ วตฺถุญฺจ ปจฺจยา วิจิกิจฺฉาสหคโต อุทฺธจฺจสหคโต โมโห. (สํขิตฺตํ.)}

\csromanheader{2}{2. ปจฺจยปจฺจนีย 2. สงฺขฺยาวาร}
\csromanheader{2}{\textbf{สุทฺธ}}
\vspace{\tipitakaparskip}
\csromancenter{\textbf{นเหตุ}ยา เอกํ, น-อารมฺมเณ ตีณิ, น-อธิปติยา นว, น-อนนฺตเร ตีณิ ฯเปฯ น-อุปนิสฺสเย ตีณิ, นปุเรชาเต สตฺต, นปจฺฉาชาเต นว, น-อาเสวเน นว, นกมฺเม จตฺตาริ, นวิปาเก นว, น-อาหาเร เอกํ, น-อินฺทฺริเย เอกํ, นฌาเน เอกํ, นมคฺเค เอกํ, นสมฺปยุตฺเต ตีณิ, นวิปฺปยุตฺเต ฉ, โนนตฺถิยา ตีณิ, โนวิคเต ตีณิ.}
\csromanheader{2}{3. ปจฺจยานุโลมปจฺจนีย}
\proseitem{60}{เหตุปจฺจยา น-อารมฺมเณ ตีณิ, น-อธิปติยา นว. (เอวํ คเณตพฺพํ.)}

\csromanheader{2}{4. ปจฺจยปจฺจนียานุโลม}
\vspace{\tipitakaparskip}
\csromancenter{นเหตุปจฺจยา อารมฺมเณ เอกํ ฯเปฯ อวิคเต เอกํ.}
\csromancenter{(นิสฺสยวาโร ปจฺจยวารสทิโส)}
\csromanheader{2}{28. คนฺถสมฺปยุตฺตทุก 5. สํสฏฺฐวาร}
\csromanheader{2}{1. ปจฺจยานุโลม 1. วิภงฺควาร}
\csromanheader{2}{\textbf{เหตุ}}
\proseitem{62}{\csromanfolio{274}คนฺถสมฺปยุตฺตํ ธมฺมํ สํสฏฺโฐ คนฺถสมฺปยุตฺโต ธมฺโม อุปฺปชฺชติ เหตุปจฺจยา.\csromanspacer{} คนฺถสมฺปยุตฺตํ เอกํ ขนฺธํ สํสฏฺฐา ตโย ขนฺธา ฯเปฯ เทฺว ขนฺเธ ฯเปฯ. (1)}

\prose{คนฺถสมฺปยุตฺตํ ธมฺมํ สํสฏฺโฐ คนฺถวิปฺปยุตฺโต ธมฺโม อุปฺปชฺชติ เหตุปจฺจยา.\csromanspacer{} ทิฏฺฐิคตวิปฺปยุตฺตโลภสหคเต ขนฺเธ สํสฏฺโฐ โลโภ.\csromanspacer{} โทมนสฺสสหคเต ขนฺเธ สํสฏฺฐํ ปฏิฆํ. (2)}

\prose{คนฺถสมฺปยุตฺตํ ธมฺมํ สํสฏฺโฐ คนฺถสมฺปยุตฺโต จ คนฺถวิปฺปยุตฺโต จ ธมฺมา อุปฺปชฺชนฺติ เหตุปจฺจยา.\csromanspacer{} ทิฏฺฐิคตวิปฺปยุตฺตโลภสหคตํ เอกํ ขนฺธํ สํสฏฺฐา ตโย ขนฺธา โลโภ จ ฯเปฯ เทฺว ขนฺเธ ฯเปฯ.\csromanspacer{} โทมนสฺสสหคตํ เอกํ ขนฺธํ สํสฏฺฐา ตโย ขนฺธา ปฏิฆญฺจ.\csromanspacer{} เทฺว ขนฺเธ ฯเปฯ. (3)}

\proseitem{63}{คนฺถวิปฺปยุตฺตํ ธมฺมํ สํสฏฺโฐ คนฺถวิปฺปยุตฺโต ธมฺโม อุปฺปชฺชติ เหตุปจฺจยา.\csromanspacer{} คนฺถวิปฺปยุตฺตํ เอกํ ขนฺธํ สํสฏฺฐา ตโย ขนฺธา ฯเปฯ เทฺว ขนฺเธ ฯเปฯ.\csromanspacer{} ปฏิสนฺธิกฺขเณ ฯเปฯ. (1)}

\prose{คนฺถวิปฺปยุตฺตํ ธมฺมํ สํสฏฺโฐ คนฺถสมฺปยุตฺโต ธมฺโม อุปฺปชฺชติ เหตุปจฺจยา.\csromanspacer{} ทิฏฺฐิคตวิปฺปยุตฺตํ โลภํ สํสฏฺฐา สมฺปยุตฺตกา ขนฺธา.\csromanspacer{} ปฏิฆํ สํสฏฺฐา สมฺปยุตฺตกา ขนฺธา. (2)}

\prose{คนฺถสมฺปยุตญฺจ คนฺถวิปฺปยุตฺตญฺจ ธมฺมํ สํสฏฺโฐ คนฺถสมฺปยุตฺโต ธมฺโม อุปฺปชฺชติ เหตุปจฺจยา.\csromanspacer{} ทิฏฺฐิคตวิปฺปยุตฺตํ โลภสหคตํ เอกํ ขนฺธญฺจ}

\vspace{\tipitakaparskip}
\csromancenter{คนฺถสมฺปยุตฺตทุก โลภญฺจ สํสฏฺฐา ตโย ขนฺธา ฯเปฯ เทฺว ขนฺเธ ฯเปฯ.\csromanspacer{} โทมนสฺสสหคตํ เอกํ ขนฺธญฺจ ปฏิฆญฺจ สํสฏฺฐา ตโย ขนฺธา ฯเปฯ เทฺว ขนฺเธ ฯเปฯ. (สํขิตฺตํ.)}
\csromanheader{2}{1. ปจฺจยานุโลม 2. สงฺขฺยาวาร}
\proseitem{64}{\csromanfolio{275}เหตุยา ฉ, อารมฺมเณ ฉ, อธิปติยา ฉ, (สพฺพตฺถ ฉ.) วิปาเก เอกํ, อาหาเร ฉ ฯเปฯ อวิคเต ฉ.}

\vspace{\tipitakaparskip}
\csromancenter{อนุโลมํ.}
\csromanheader{2}{2. ปจฺจยปจฺจนีย 1. วิภงฺควาร}
\csromanheader{2}{\textbf{นเหตุ}}
\proseitem{65}{คนฺถวิปฺปยุตฺตํ ธมฺมํ สํสฏฺโฐ คนฺถวิปฺปยุตฺโต ธมฺโม อุปฺปชฺชติ นเหตุปจฺจยา.\csromanspacer{} อเหตุกํ คนฺถวิปฺปยุตฺตํ ฯเปฯ.\csromanspacer{} อเหตุกปฏิสนฺธิกฺขเณ ฯเปฯ.\csromanspacer{} วิจิกิจฺฉาสหคเต อุทฺธจฺจสหคเต ขนฺเธ สํสฏฺโฐ วิจิกิจฺฉาสหคโต อุทฺธจฺจสหคโต โมโห. (สํขิตฺตํ.)}

\csromanheader{2}{2. ปจฺจยปจฺจนีย 2. สงฺขฺยาวาร}
\vspace{\tipitakaparskip}
\csromancenter{\textbf{นเหตุ}ยา เอกํ, น-อธิปติยา ฉ, นปุเรชาเต ฉ, นปจฺฉาชาเต ฉ, น-อาเสวเน ฉ, นกมฺเม จตฺตาริ, นวิปาเก ฉ, นฌาเน เอกํ, นมคฺเค เอกํ, นวิปฺปยุตฺเต ฉ.}
\csromanheader{2}{3. ปจฺจยานุโลมปจฺจนีย}
\proseitem{67}{เหตุปจฺจยา น-อธิปติยา ฉ, นปุเรชาเต ฉ, นปจฺฉาชาเต ฉ, น-อาเสวเน ฉ, นกมฺเม จตฺตาริ, นวิปาเก ฉ, นวิปฺปยุตฺเต ฉ.}

\csromanheader{2}{4. ปจฺจยปจฺจนียานุโลม}
\vspace{\tipitakaparskip}
\csromancenter{นเหตุปจฺจยา อารมฺมเณ เอกํ, อนนฺตเร เอกํ ฯเปฯ อวิคเต เอกํ.}
\csromancenter{(สมฺปยุตฺตวาโร สํสฏฺฐวารสทิโส)}
\csromanheader{2}{28. คนฺถสมฺปยุตฺตทุก 7. ปญฺหาวาร}
\csromanheader{2}{1. ปจฺจยานุโลม 1. วิภงฺควาร}
\csromanheader{2}{\textbf{เหตุ}}
\proseitem{69}{\csromanfolio{276}คนฺถสมฺปยุตฺโต ธมฺโม คนฺถสมฺปยุตฺตสฺส ธมฺมสฺส เหตุปจฺจเยน ปจฺจโย.\csromanspacer{} คนฺถสมฺปยุตฺตา เหตู สมฺปยุตฺตกานํ ขนฺธานํ เหตุปจฺจเยน ปจฺจโย. (1)}

\prose{คนฺถสมฺปยุตฺโต ธมฺโม คนฺถวิปฺปยุตฺตสฺส ธมฺมสฺส เหตุปจฺจเยน ปจฺจโย.\csromanspacer{} คนฺถสมฺปยุตฺตา เหตู จิตฺตสมุฏฺฐานานํ รูปานํ เหตุปจฺจเยน ปจฺจโย.\csromanspacer{} ทิฏฺฐิคตวิปฺปยุตฺตโลภสหคโต เหตุ โลภสฺส จิตฺตสมุฏฺฐานานญฺจ รูปานํ เหตุปจฺจเยน ปจฺจโย.\csromanspacer{} โทมนสฺสสหคโต เหตุ ปฏิฆสฺส จิตฺตสมุฏฺฐานานญฺจ รูปานํ เหตุปจฺจเยน ปจฺจโย. (2)}

\prose{คนฺถสมฺปยุตฺโต ธมฺโม คนฺถสมฺปยุตฺตสฺส จ คนฺถวิปฺปยุตฺตสฺส จ ธมฺมสฺส เหตุปจฺจเยน ปจฺจโย.\csromanspacer{} คนฺถสมฺปยุตฺตา เหตู สมฺปยุตฺตกานํ ขนฺธานํ จิตฺตสมุฏฺฐานานญฺจ รูปานํ เหตุปจฺจเยน ปจฺจโย.\csromanspacer{} ทิฏฺฐิคตวิปฺปยุตฺตโลภสหคโต เหตุ สมฺปยุตฺตกานํ ขนฺธานํ โลภสฺส จ จิตฺตสมุฏฺฐานานญฺจ รูปานํ เหตุปจฺจเยน ปจฺจโย.\csromanspacer{} โทมนสฺสสหคโต เหตุ สมฺปยุตฺตกานํ ขนฺธานํ ปฏิฆสฺส จ จิตฺตสมุฏฺฐานานญฺจ รูปานํ เหตุปจฺจเยน ปจฺจโย. (3)}

\proseitem{70}{คนฺถวิปฺปยุตฺโต ธมฺโม คนฺถวิปฺปยุตฺตสฺส ธมฺมสฺส เหตุปจฺจเยน ปจฺจโย.\csromanspacer{} คนฺถวิปฺปยุตฺตา เหตู สมฺปยุตฺตกานํ ขนฺธานํ จิตฺตสมุฏฺฐานานญฺจ}

\vspace{\tipitakaparskip}
\csromancenter{คนฺถสมฺปยุตฺตทุก รูปานํ เหตุปจฺจเยน ปจฺจโย.\csromanspacer{} ทิฏฺฐิคตวิปฺปยุตฺโต โลโภ จิตฺตสมุฏฺฐานานํ รูปานํ เหตุปจฺจเยน ปจฺจโย.\csromanspacer{} ปฏิฆํ จิตฺตสมุฏฺฐานานํ รูปานํ เหตุปจฺจเยน ปจฺจโย.\csromanspacer{} ปฏิสนฺธิกฺขเณ ฯเปฯ. (1)}
\prose{\csromanfolio{277}คนฺถวิปฺปยุตฺโต ธมฺโม คนฺถสมฺปยุตฺตสฺส ธมฺมสฺส เหตุปจฺจเยน ปจฺจโย.\csromanspacer{} ทิฏฺฐิคตวิปฺปยุตฺโต โลโภ สมฺปยุตฺตกานํ ขนฺธานํ เหตุปจฺจเยน ปจฺจโย.\csromanspacer{} ปฏิฆํ สมฺปยุตฺตกานํ ขนฺธานํ เหตุปจฺจเยน ปจฺจโย. (2)}

\prose{คนฺถวิปฺปยุตฺโต ธมฺโม คนฺถสมฺปยุตฺตสฺส จ คนฺถวิปฺปยุตฺตสฺส จ ธมฺมสฺส เหตุปจฺจเยน ปจฺจโย.\csromanspacer{} ทิฏฺฐิคตวิปฺปยุตฺโต โลโภ สมฺปยุตฺตกานํ ขนฺธานํ จิตฺตสมุฏฺฐานานญฺจ รูปานํ เหตุปจฺจเยน ปจฺจโย.\csromanspacer{} ปฏิฆํ สมฺปยุตฺตกานํ ขนฺธานํ จิตฺตสมุฏฺฐานานญฺจ รูปานํ เหตุปจฺจเยน ปจฺจโย. (3)}

\proseitem{71}{คนฺถสมฺปยุตฺโต จ คนฺถวิปฺปยุตฺโต จ ธมฺมา คนฺถสมฺปยุตฺตสฺส ธมฺมสฺส เหตุปจฺจเยน ปจฺจโย.\csromanspacer{} ทิฏฺฐิคตวิปฺปยุตฺตโลภสหคโต เหตุ จ โลโภ จ สมฺปยุตฺตกานํ ขนฺธานํ เหตุปจฺจเยน ปจฺจโย.\csromanspacer{} โทมนสฺสสหคโต เหตุ จ ปฏิฆญฺจ สมฺปยุตฺตกานํ ขนฺธานํ เหตุปจฺจเยน ปจฺจโย. (1)}

\prose{คนฺถสมฺปยุตฺโต จ คนฺถวิปฺปยุตฺโต จ ธมฺมา คนฺถวิปฺปยุตฺตสฺส ธมฺมสฺส เหตุปจฺจเยน ปจฺจโย.\csromanspacer{} ทิฏฺฐิคตวิปฺปยุตฺตโลภสหคโต เหตุ จ โลโภ จ จิตฺตสมุฏฺฐานานํ รูปานํ เหตุปจฺจเยน ปจฺจโย.\csromanspacer{} โทมนสฺสสหคโต เหตุ จ ปฏิฆญฺจ จิตฺตสมุฏฺฐานานํ รูปานํ เหตุปจฺจเยน ปจฺจโย. (2)}

\prose{คนฺถสมฺปยุตฺโต จ คนฺถวิปฺปยุตฺโต จ ธมฺมา คนฺถสมฺปยุตฺตสฺส จ คนฺถวิปฺปยุตฺตสฺส จ ธมฺมสฺส เหตุปจฺจเยน ปจฺจโย.\csromanspacer{} ทิฏฺฐิคตวิปฺปยุตฺตโลภสหคโต เหตุ จ โลโภ จ สมฺปยุตฺตกานํ ขนฺธานํ จิตฺตสมุฏฺฐานานญฺจ รูปานํ เหตุปจฺจเยน ปจฺจโย.\csromanspacer{} โทมนสฺสสหคโต เหตุ จ ปฏิฆญฺจ สมฺปยุตฺตกานํ ขนฺธานํ จิตฺตสมุฏฺฐานานญฺจ รูปานํ เหตุปจฺจเยน ปจฺจโย. (3)}

\csromanheader{2}{\textbf{อารมฺมณ}}
\proseitem{72}{คนฺถสมฺปยุตฺโต ธมฺโม คนฺถสมฺปยุตฺตสฺส ธมฺมสฺส อารมฺมณปจฺจเยน ปจฺจโย.\csromanspacer{} คนฺถสมฺปยุตฺเต ขนฺเธ อารพฺภ คนฺถสมฺปยุตฺตกา ขนฺธา \csromanfolio{278}อุปฺปชฺชนฺติ. (ตีสุปิ มูลา ปุจฺฉิตพฺพา.) คนฺถสมฺปยุตฺเต ขนฺเธ อารพฺภ คนฺถวิปฺปยุตฺตา ขนฺธา อุปฺปชฺชนฺติ.\csromanspacer{} คนฺถสมฺปยุตฺเต ขนฺเธ อารพฺภ ทิฏฺฐิคตวิปฺปยุตฺตโลภสหคตา ขนฺธา จ โลโภ จ อุปฺปชฺชนฺติ.\csromanspacer{} โทมนสฺสสหคตา ขนฺธา จ ปฏิฆญฺจ อุปฺปชฺชนฺติ. (3)}

\proseitem{73}{คนฺถวิปฺปยุตฺโต ธมฺโม คนฺถวิปฺปยุตฺตสฺส ธมฺมสฺส อารมฺมณปจฺจเยน ปจฺจโย.\csromanspacer{} ทานํ ทตฺวา, สีลํ ฯเปฯ อุโปสถกมฺมํ กตฺวา ตํ ปจฺจเวกฺขติ, ปุพฺเพ สุจิณฺณานิ ปจฺจเวกฺขติ, ฌานา วุฏฺฐหิตฺวา ฌานํ ปจฺจเวกฺขติ.\csromanspacer{} อริยา มคฺคา วุฏฺฐหิตฺวา มคฺคํ ปจฺจเวกฺขนฺติ, ผลํ ปจฺจเวกฺขนฺติ, นิพฺพานํ ปจฺจเวกฺขนฺติ.\csromanspacer{} นิพฺพานํ โคตฺรภุสฺส, โวทานสฺส, มคฺคสฺส, ผลสฺส, อาวชฺชนาย อารมฺมณปจฺจเยน ปจฺจโย.\csromanspacer{} อริยา คนฺถวิปฺปยุตฺเต ปหีเน กิเลเส ปจฺจเวกฺขนฺติ, วิกฺขมฺภิเต กิเลเส ปจฺจเวกฺขนฺติ.\csromanspacer{} ปุพฺเพ สมุทาจิณฺเณ กิเลเส ชานนฺติ.\csromanspacer{} จกฺขุํ ฯเปฯ วตฺถุํ คนฺถวิปฺปยุตฺเต ขนฺเธ จ โลภญฺจ ปฏิฆญฺจ อนิจฺจโต ฯเปฯ วิปสฺสติ, อสฺสาเทติ อภินนฺทติ, ตํ อารพฺภ คนฺถวิปฺปยุตฺโต ราโค อุปฺปชฺชติ.\csromanspacer{} วิจิกิจฺฉา ฯเปฯ อุทฺธจฺจํ ฯเปฯ โทมนสฺสํ อุปฺปชฺชติ.\csromanspacer{} ทิพฺเพน จกฺขุนา รูปํ ปสฺสติ.\csromanspacer{} ทิพฺพาย โสตธาตุยา สทฺทํ สุณาติ.\csromanspacer{} เจโตปริยญาเณน คนฺถวิปฺปยุตฺตจิตฺตสมงฺคิสฺส จิตฺตํ ชานาติ.\csromanspacer{} อากาสานญฺจายตนํ วิญฺญาณญฺจายตนสฺส ฯเปฯ.\csromanspacer{} อากิญฺจญฺญายตนํ เนวสญฺญานาสญฺญายตนสฺส ฯเปฯ.\csromanspacer{} รูปายตนํ จกฺขุวิญฺญาณสฺส ฯเปฯ.\csromanspacer{} โผฏฺฐพฺพายตนํ กายวิญฺญาณสฺส ฯเปฯ.\csromanspacer{} คนฺถวิปฺปยุตฺตา ขนฺธา อิทฺธิวิธญาณสฺส, เจโตปริยญาณสฺส, ปุพฺเพนิวาสานุสฺสติญาณสฺส, ยถากมฺมูปคญาณสฺส, อนาคตํสญาณสฺส, อาวชฺชนาย อารมฺมณปจฺจเยน ปจฺจโย. (1)}

\prose{คนฺถวิปฺปยุตฺโต ธมฺโม คนฺถสมฺปยุตฺตสฺส ธมฺมสฺส อารมฺมณปจฺจเยน ปจฺจโย.\csromanspacer{} ทานํ ทตฺวา, สีลํ ฯเปฯ อุโปสถกมฺมํ ฯเปฯ.\csromanspacer{} ปุพฺเพ ฯเปฯ.\csromanspacer{} ฌานา ฯเปฯ ตํ อสฺสาเทติ อภินนฺทติ, ตํ อารพฺภ คนฺถสมฺปยุตฺโต ราโค อุปฺปชฺชติ.\csromanspacer{} ทิฏฺฐิ อุปฺปชฺชติ.\csromanspacer{} โทมนสฺสํ อุปฺปชฺชติ. (สํขิตฺตํ.) (2)}

\prose{คนฺถวิปฺปยุตฺโต ธมฺโม คนฺถสมฺปยุตฺตสฺส จ คนฺถวิปฺปยุตฺตสฺส จ ธมฺมสฺส อารมฺมณปจฺจเยน ปจฺจโย.\csromanspacer{} จกฺขุํ ฯเปฯ วตฺถุํ คนฺถวิปฺปยุตฺเต ขนฺเธ จ โลภญฺจ ปฏิฆญฺจ อารพฺภ ทิฏฺฐิคตวิปฺปยุตฺตโลภสหคตา ขนฺธา จ โลโภ จ โทมนสฺสสหคตา ขนฺธา จ ปฏิฆญฺจ อุปฺปชฺชนฺติ. (3)}

\csromanheader{2}{28. คนฺถสมฺปยุตฺตทุก}
\proseitem{74}{\csromanfolio{279}คนฺถสมฺปยุตฺโต จ คนฺถวิปฺปยุตฺโต จ ธมฺมา คนฺถสมฺปยุตฺตสฺส ธมฺมสฺส อารมฺมณปจฺจเยน ปจฺจโย.\csromanspacer{} ทิฏฺฐิคตวิปฺปยุตฺตโลภสหคเต ขนฺเธ จ โลภญฺจ โทมนสฺสสหคเต ขนฺเธ จ ปฏิฆญฺจ อารพฺภ คนฺถสมฺปยุตฺตกา ขนฺธา อุปฺปชฺชนฺติ. (มูลํ ปุจฺฉิตพฺพํ.) ทิฏฺฐิคตวิปฺปยุตฺตโลภสหคเต ขนฺเธ จ โลภญฺจ โทมนสฺสสหคเต ขนฺเธ จ ปฏิฆญฺจ อารพฺภ คนฺถวิปฺปยุตฺตา ขนฺธา อุปฺปชฺชนฺติ.\csromanspacer{} ทิฏฺฐิคตวิปฺปยุตฺตโลภสหคเต ขนฺเธ จ โลภญฺจ โทมนสฺสสหคเต ขนฺเธ จ ปฏิฆญฺจ อารพฺภ ทิฏฺฐิคตวิปฺปยุตฺตโลภสหคตา ขนฺธา จ โลโภ จ โทมนสฺสสหคตา ขนฺธา จ ปฏิฆญฺจ อุปฺปชฺชนฺติ. (3)}

\csromanheader{2}{\textbf{อธิปติ}}
\proseitem{75}{คนฺถสมฺปยุตฺโต ธมฺโม คนฺถสมฺปยุตฺตสฺส ธมฺมสฺส อธิปติปจฺจเยน ปจฺจโย.\csromanspacer{} \textbf{อารมฺมณาธิปติ}, สหชาตาธิปติ.\csromanspacer{}\csromanspacer{}\textbf{อารมฺมณาธิปติ}\csromandash{}คนฺถสมฺปยุตฺเต ขนฺเธ ครุํ กตฺวา คนฺถสมฺปยุตฺตกา ขนฺธา อุปฺปชฺชนฺติ.\csromanspacer{}\csromanspacer{}\textbf{สหชาตาธิปติ\csromandash{}} คนฺถสมฺปยุตฺตาธิปติ สมฺปยุตฺตกานํ ขนฺธานํ อธิปติปจฺจเยน ปจฺจโย. (1)}

\prose{คนฺถสมฺปยุตฺโต ธมฺโม คนฺถวิปฺปยุตฺตสฺส ธมฺมสฺส อธิปติปจฺจเยน ปจฺจโย.\csromanspacer{} \textbf{อารมฺมณาธิปติ}, สหชาตาธิปติ.\csromanspacer{}\csromanspacer{}\textbf{อารมฺมณาธิปติ}\csromandash{}คนฺถสมฺปยุตฺเต ขนฺเธ ครุํ กตฺวา ทิฏฺฐิคตวิปฺปยุตฺโต โลโภ อุปฺปชฺชติ.\csromanspacer{}\csromanspacer{}\textbf{สหชาตาธิปติ\csromandash{}} คนฺถสมฺปยุตฺตาธิปติ จิตฺตสมุฏฺฐานานํ รูปานํ อธิปติปจฺจเยน ปจฺจโย.\csromanspacer{} ทิฏฺฐิคตวิปฺปยุตฺตโลภสหคตาธิปติ โลภสฺส จ จิตฺตสมุฏฺฐานานญฺจ รูปานํ อธิปติปจฺจเยน ปจฺจโย.\csromanspacer{} โทมนสฺสสหคตาธิปติ ปฏิฆสฺส จ จิตฺตสมุฏฺฐานานญฺจ รูปานํ อธิปติปจฺจเยน ปจฺจโย. (2)}

\prose{คนฺถสมฺปยุตฺโต ธมฺโม คนฺถสมฺปยุตฺตสฺส จ คนฺถวิปฺปยุตฺตสฺส จ ธมฺมสฺส อธิปติปจฺจเยน ปจฺจโย.\csromanspacer{} \textbf{อารมฺมณาธิปติ}, สหชาตาธิปติ.\csromanspacer{}\csromanspacer{}\textbf{อารมฺมณาธิปติ}\csromandash{}คนฺถสมฺปยุตฺเต ขนฺเธ ครุํ กตฺวา ทิฏฺฐิคตวิปฺปยุตฺตโลภสหคตา ขนฺธา จ โลโภ จ อุปฺปชฺชนฺติ.\csromanspacer{}\csromanspacer{}\textbf{สหชาตาธิปติ\csromandash{}}ทิฏฺฐิคตวิปฺปยุตฺตโลภสหคตาธิปติ สมฺปยุตฺตกานํ ขนฺธานํ โลภสฺส จ จิตฺตสมุฏฺฐานานญฺจ รูปานํ อธิปติปจฺจเยน ปจฺจโย.\csromanspacer{} โทมนสฺสสหคตาธิปติ สมฺปยุตฺตกานํ ขนฺธานํ ปฏิฆสฺส จ จิตฺตสมุฏฺฐานานญฺจ รูปานํ อธิปติปจฺจเยน ปจฺจโย. (3)}

\proseitem{76}{\csromanfolio{280}คนฺถวิปฺปยุตฺโต ธมฺโม คนฺถวิปฺปยุตฺตสฺส ธมฺมสฺส อธิปติปจฺจเยน ปจฺจโย.\csromanspacer{} \textbf{อารมฺมณาธิปติ}, สหชาตาธิปติ.\csromanspacer{}\csromanspacer{}\textbf{อารมฺมณาธิปติ}\csromandash{}ทานํ ฯเปฯ สีลํ ฯเปฯ อุโปสถกมฺมํ กตฺวา ตํ ครุํ กตฺวา ปจฺจเวกฺขติ, ปุพฺเพ สุจิณฺณานิ ครุํ กตฺวา ปจฺจเวกฺขติ.\csromanspacer{} ฌานา วุฏฺฐหิตฺวา ฌานํ ครุํ กตฺวา ปจฺจเวกฺขติ.\csromanspacer{} อริยา มคฺคา วุฏฺฐหิตฺวา มคฺคํ ครุํ กตฺวา ปจฺจเวกฺขนฺติ, ผลา วุฏฺฐหิตฺวา ผลํ ครุํ กตฺวา ปจฺจเวกฺขนฺติ, นิพฺพานํ ครุํ กตฺวา ปจฺจเวกฺขนฺติ.\csromanspacer{} นิพฺพานํ โคตฺรภุสฺส, โวทานสฺส, มคฺคสฺส, ผลสฺส อธิปติปจฺจเยน ปจฺจโย.\csromanspacer{} จกฺขุํ ฯเปฯ วตฺถุํ คนฺถวิปฺปยุตฺเต ขนฺเธ จ โลภญฺจ ครุํ กตฺวา อสฺสาเทติ อภินนฺทติ, ตํ ครุํ กตฺวา คนฺถวิปฺปยุตฺโต ราโค อุปฺปชฺชติ.\csromanspacer{}\csromanspacer{}\textbf{สหชาตาธิปติ}\csromandash{}คนฺถวิปฺปยุตฺตาธิปติ สมฺปยุตฺตกานํ ขนฺธานํ จิตฺตสมุฏฺฐานานญฺจ รูปานํ อธิปติปจฺจเยน ปจฺจโย. (1)}

\prose{คนฺถวิปฺปยุตฺโต ธมฺโม คนฺถสมฺปยุตฺตสฺส ธมฺมสฺส อธิปติปจฺจเยน ปจฺจโย.\csromanspacer{} \textbf{อารมฺมณาธิปติ}\csromandash{}ทานํ ฯเปฯ สีลํ ฯเปฯ อุโปสถกมฺมํ ฯเปฯ.\csromanspacer{} ปุพฺเพ สุจิณฺณานิ ฯเปฯ.\csromanspacer{} ฌานา วุฏฺฐหิตฺวา ฌานํ ฯเปฯ.\csromanspacer{} จกฺขุํ ฯเปฯ วตฺถุํ คนฺถวิปฺปยุตฺเต ขนฺเธ จ โลภญฺจ ครุํ กตฺวา อสฺสาเทติ อภินนฺทติ, ตํ ครุํ กตฺวา คนฺถสมฺปยุตฺโต ราโค อุปฺปชฺชติ.\csromanspacer{} ทิฏฺฐิ อุปฺปชฺชติ. (2)}

\prose{คนฺถวิปฺปยุตฺโต ธมฺโม คนฺถสมฺปยุตฺตสฺส จ คนฺถวิปฺปยุตฺตสฺส จ ธมฺมสฺส อธิปติปจฺจเยน ปจฺจโย.\csromanspacer{} \textbf{อารมฺมณาธิปติ}\csromandash{}จกฺขุํ ฯเปฯ วตฺถุํ คนฺถวิปฺปยุตฺเต ขนฺเธ จ โลภญฺจ ครุํ กตฺวา ทิฏฺฐิคตวิปฺปยุตฺตโลภสหคตา ขนฺธา จ โลโภ จ อุปฺปชฺชนฺติ. (3)}

\proseitem{77}{คนฺถสมฺปยุตฺโต จ คนฺถวิปฺปยุตฺโต จ ธมฺมา คนฺถสมฺปยุตฺตสฺส ธมฺมสฺส อธิปติปจฺจเยน ปจฺจโย.\csromanspacer{} \textbf{อารมฺมณาธิปติ}\csromandash{}ทิฏฺฐิคตวิปฺปยุตฺตโลภสหคเต ขนฺเธ จ โลภญฺจ ครุํ กตฺวา คนฺถสมฺปยุตฺตกา ขนฺธา อุปฺปชฺชนฺติ. (มูลา ปุจฺฉิตพฺพา.) ทิฏฺฐิคตวิปฺปยุตฺตโลภสหคเต ขนฺเธ จ โลภญฺจ ครุํ กตฺวา ทิฏฺฐิคตวิปฺปยุตฺโต โลโภ อุปฺปชฺชติ. (มูลา ปุจฺฉิตพฺพา.) ทิฏฺฐิคตวิปฺปยุตฺตโลภสหคเต ขนฺเธ จ โลภญฺจ ครุํ กตฺวา ทิฏฺฐิคตวิปฺปยุตฺตโลภสหคตา ขนฺธา จ โลโภ จ อุปฺปชฺชนฺติ. (3)}

\csromanheader{2}{28. คนฺถสมฺปยุตฺตทุก \textbf{อนนฺตร}}
\proseitem{78}{\csromanfolio{281}คนฺถสมฺปยุตฺโต ธมฺโม คนฺถสมฺปยุตฺตสฺส ธมฺมสฺส อนนฺตรปจฺจเยน ปจฺจโย.\csromanspacer{} ปุริมา ปุริมา คนฺถสมฺปยุตฺตา ขนฺธา ปจฺฉิมานํ ปจฺฉิมานํ คนฺถสมฺปยุตฺตกานํ ขนฺธานํ อนนฺตรปจฺจเยน ปจฺจโย. (มูลา ปุจฺฉิตพฺพา.) ปุริมา ปุริมา ทิฏฺฐิคตวิปฺปยุตฺตโลภสหคตา ขนฺธา ปจฺฉิมสฺส ปจฺฉิมสฺส ทิฏฺฐิคกวิปฺปยุตฺตสฺส โลภสฺส อนนฺตรปจฺจเยน ปจฺจโย.\csromanspacer{} ปุริมา ปุริมา โทมนสฺสสหคตา ขนฺธา ปจฺฉิมสฺส ปจฺฉิมสฺส ปฏิฆสฺส อนนฺตรปจฺจเยน ปจฺจโย.\csromanspacer{} คนฺถสมฺปยุตฺตา ขนฺธา วุฏฺฐานสฺส อนนฺตรปจฺจเยน ปจฺจโย. (มูลา ปุจฺฉิตพฺพา.) ปุริมา ปุริมา ทิฏฺฐิคตวิปฺปยุตฺตโลภสหคตา ขนฺธา ปจฺฉิมานํ ปจฺฉิมานํ ทิฏฺฐิคตวิปฺปยุตฺตโลภสหคตานํ ขนฺธานํ โลภสฺส จ อนนฺตรปจฺจเยน ปจฺจโย.\csromanspacer{} ปุริมา ปุริมา โทมนสฺสสหคตา ขนฺธา ปจฺฉิมานํ ปจฺฉิมานํ โทมนสฺสสหคตานํ ขนฺธานํ ปฏิฆสฺส จ อนนฺตรปจฺจเยน ปจฺจโย. (3)}

\proseitem{79}{คนฺถวิปฺปยุตฺโต ธมฺโม คนฺถวิปฺปยุตฺตสฺส ธมฺมสฺส อนนฺตรปจฺจเยน ปจฺจโย.\csromanspacer{} ปุริโม ปุริโม ทิฏฺฐิคตวิปฺปยุตฺโต โลโภ ปจฺฉิมสฺส ปจฺฉิมสฺส ทิฏฺฐิคตวิปฺปยุตฺตสฺส โลภสฺส อนนฺตรปจฺจเยน ปจฺจโย.\csromanspacer{} ปุริมํ ปุริมํ ปฏิฆํ ปจฺฉิมสฺส ปจฺฉิมสฺส ปฏิฆสฺส อนนฺตรปจฺจเยน ปจฺจโย.\csromanspacer{} ปุริมา ปุริมา คนฺถวิปฺปยุตฺตา ขนฺธา ปจฺฉิมานํ ปจฺฉิมานํ คนฺถวิปฺปยุตฺตานํ ขนฺธานํ อนนฺตรปจฺจเยน ปจฺจโย.\csromanspacer{} อนุโลมํ โคตฺรภุสฺส ฯเปฯ ผลสมาปตฺติยา อนนฺตรปจฺจเยน ปจฺจโย. (1)}

\prose{คนฺถวิปฺปยุตฺโต ธมฺโม คนฺถสมฺปยุตฺตสฺส ธมฺมสฺส อนนฺตรปจฺจเยน ปจฺจโย.\csromanspacer{} ปุริโม ปุริโม ทิฏฺฐิคตวิปฺปยุตฺโต โลโภ ปจฺฉิมานํ ปจฺฉิมานํ ทิฏฺฐิคตวิปฺปยุตฺตโลภสหคตานํ ขนฺธานํ อนนฺตรปจฺจเยน ปจฺจโย.\csromanspacer{} ปุริมํ ปุริมํ ปฏิฆํ ปจฺฉิมานํ ปจฺฉิมานํ โทมนสฺสสหคตานํ ขนฺธานํ อนนฺตรปจฺจเยน ปจฺจโย.\csromanspacer{} อาวชฺชนา คนฺถสมฺปยุตฺตกานํ ขนฺธานํ อนนฺตรปจฺจเยน ปจฺจโย. (2)}

\prose{คนฺถวิปฺปยุตฺโต ธมฺโม คนฺถสมฺปยุตฺตสฺส จ คนฺถวิปฺปยุตฺตสฺส จ ธมฺมสฺส อนนฺตรปจฺจเยน ปจฺจโย.\csromanspacer{} ปุริโม ปุริโม ทิฏฺฐิคตวิปฺปยุตฺโต โลโภ ปจฺฉิมานํ ปจฺฉิมานํ ทิฏฺฐิคตวิปฺปยุตฺตโลภสหคตานํ ขนฺธานํ โลภสฺส จ \csromanfolio{282}อนนฺตรปจฺจเยน ปจฺจโย.\csromanspacer{} ปุริมํ ปุริมํ ปฏิฆํ ปจฺฉิมานํ ปจฺฉิมานํ โทมนสฺสสหคตานํ ขนฺธานํ ปฏิฆสฺส จ อนนฺตรปจฺจเยน ปจฺจโย.\csromanspacer{} อาวชฺชนา ทิฏฺฐิคตวิปฺปยุตฺตโลภสหคตานํ ขนฺธานํ โลภสฺส จ, โทมนสฺสสหคตานํ ขนฺธานํ ปฏิฆสฺส จ อนนฺตรปจฺจเยน ปจฺจโย. (3)}

\proseitem{80}{คนฺถสมฺปยุตฺโต จ คนฺถวิปฺปยุตฺโต จ ธมฺมา คนฺถสมฺปยุตฺตสฺส ธมฺมสฺส อนนฺตรปจฺจเยน ปจฺจโย.\csromanspacer{} ปุริมา ปุริมา ทิฏฺฐิคตวิปฺปยุตฺตโลภสหคตา ขนฺธา จ โลโภ จ ปจฺฉิมานํ ปจฺฉิมานํ ทิฏฺฐิคตวิปฺปยุตฺตโลภสหคตานํ ขนฺธานํ อนนฺตรปจฺจเยน ปจฺจโย.\csromanspacer{} ปุริมา ปุริมา โทมนสฺสสหคตา ขนฺธา จ ปฏิฆญฺจ ปจฺฉิมานํ ปจฺฉิมานํ โทมนสฺสสหคตานํ ขนฺธานํ อนนฺตรปจฺจเยน ปจฺจโย. (มูลา ปุจฺฉิตพฺพา.) ปุริมา ปุริมา ทิฏฺฐิคตวิปฺปยุตฺตโลภสหคตา ขนฺธา จ โลโภ จ ปจฺฉิมสฺส ปจฺฉิมสฺส ทิฏฺฐิคตวิปฺปยุตฺตสฺส โลภสฺส อนนฺตรปจฺจเยน ปจฺจโย.\csromanspacer{} ปุริมา ปุริมา โทมนสฺสสหคตา ขนฺธา จ ปฏิฆญฺจ ปจฺฉิมสฺส ปจฺฉิมสฺส ปฏิฆสฺส อนนฺตรปจฺจเยน ปจฺจโย.\csromanspacer{} ทิฏฺฐิคตวิปฺปยุตฺตโลภสหคตา ขนฺธา จ โลโภ จ, โทมนสฺสสหคตา ขนฺธา จ ปฏิฆญฺจ วุฏฺฐานสฺส อนนฺตรปจฺจเยน ปจฺจโย. (มูลา ปุจฺฉิตพฺพา.) ปุริมา ปุริมา ทิฏฺฐิคตวิปฺปยุตฺตโลภสหคตา ขนฺธา จ โลโภ จ ปจฺฉิมานํ ปจฺฉิมานํ ทิฏฺฐิคตวิปฺปยุตฺตโลภสหคตานํ ขนฺธานํ โลภสฺส จ อนนฺตรปจฺจเยน ปจฺจโย.\csromanspacer{} ปุริมา ปุริมา โทมนสฺสสหคตา ขนฺธา จ ปฏิฆญฺจ ปจฺฉิมานํ ปจฺฉิมานํ โทมนสฺสสหคตานํ ขนฺธานํ ปฏิฆสฺส จ อนนฺตรปจฺจเยน ปจฺจโย. (3)}

\csromanheader{2}{\textbf{สมนนฺตราทิ}}
\proseitem{81}{คนฺถสมฺปยุตฺโต ธมฺโม คนฺถสมฺปยุตฺตสฺส ธมฺมสฺส สมนนฺตรปจฺจเยน ปจฺจโย.\csromanspacer{} สหชาตปจฺจเยน ปจฺจโย.\csromanspacer{} อญฺญมญฺญปจฺจเยน ปจฺจโย.\csromanspacer{} นิสฺสยปจฺจเยน ปจฺจโย.}

\csromanheader{2}{\textbf{อุปนิสฺสย}}
\proseitem{82}{คนฺถสมฺปยุตฺโต ธมฺโม คนฺถสมฺปยุตฺตสฺส ธมฺมสฺส อุปนิสฺสยปจฺจเยน ปจฺจโย.\csromanspacer{} อารมฺมณูปนิสฺสโย, อนนฺตรูปนิสฺสโย, ปกตูปนิสฺสโย ฯเปฯ.\csromanspacer{} ปกตูปนิสฺสโย\csromandash{}คนฺถสมฺปยุตฺตา ขนฺธา คนฺถสมฺปยุตฺตกานํ ขนฺธานํ อุปนิสฺสยปจฺจเยน ปจฺจโย. (มูลํ ปุจฺฉิตพฺพํ.}

\vspace{\tipitakaparskip}
\csromancenter{คนฺถสมฺปยุตฺตทุก ตีณิปิ อุปนิสฺสยา.) คนฺถสมฺปยุตฺตา ขนฺธา คนฺถวิปฺปยุตฺตานํ ขนฺธานํ อุปนิสฺสยปจฺจเยน ปจฺจโย. (มูลํ.\csromanspacer{} ตีณิปิ อุปนิสฺสยา.) คนฺถสมฺปยุตฺตา ขนฺธา ทิฏฺฐิคตวิปฺปยุตฺตโลภสหคตานํ ขนฺธานํ โลภสฺส จ อุปนิสฺสยปจฺจเยน ปจฺจโย.\csromanspacer{} โทมนสฺสสหคตานํ ขนฺธานํ ปฏิฆสฺส จ อุปนิสฺสยปจฺจเยน ปจฺจโย. (3)}
\proseitem{83}{\csromanfolio{283}คนฺถวิปฺปยุตฺโต ธมฺโม คนฺถวิปฺปยุตฺตสฺส ธมฺมสฺส อุปนิสฺสยปจฺจเยน ปจฺจโย.\csromanspacer{} อารมฺมณูปนิสฺสโย, อนนฺตรูปนิสฺสโย, ปกตูปนิสฺสโย ฯเปฯ.\csromanspacer{} ปกตูปนิสฺสโย\csromandash{}สทฺธํ อุปนิสฺสาย ทานํ เทติ ฯเปฯ สมาปตฺติํ อุปฺปาเทติ.\csromanspacer{} มานํ ชปฺเปติ.\csromanspacer{} สีลํ ฯเปฯ.\csromanspacer{} ปญฺญํ.\csromanspacer{} ราคํ.\csromanspacer{} โทสํ.\csromanspacer{} โมหํ.\csromanspacer{} มานํ.\csromanspacer{} ปตฺถนํ.\csromanspacer{} กายิกํ สุขํ.\csromanspacer{} กายิกํ ทุกฺขํ.\csromanspacer{} อุตุํ.\csromanspacer{} โภชนํ.\csromanspacer{} เสนาสนํ อุปนิสฺสาย ทานํ เทติ ฯเปฯ สมาปตฺติํ อุปฺปาเทติ.\csromanspacer{} ปาณํ หนติ ฯเปฯ สํฆํ ภินฺทติ.\csromanspacer{} สทฺธา ฯเปฯ.\csromanspacer{} ปญฺญา.\csromanspacer{} ราโค ฯเปฯ.\csromanspacer{} ปตฺถนา ฯเปฯ.\csromanspacer{} เสนาสนํ สทฺธาย ฯเปฯ ปญฺญาย.\csromanspacer{} ราคสฺส.\csromanspacer{} โทสสฺส.\csromanspacer{} โมหสฺส.\csromanspacer{} มานสฺส.\csromanspacer{} ปตฺถนาย ฯเปฯ ผลสมาปตฺติยา อุปนิสฺสยปจฺจเยน ปจฺจโย. (1)}

\prose{คนฺถวิปฺปยุตฺโต ธมฺโม คนฺถสมฺปยุตฺตสฺส ธมฺมสฺส อุปนิสฺสยปจฺจเยน ปจฺจโย.\csromanspacer{} อารมฺมณูปนิสฺสโย, อนนฺตรูปนิสฺสโย, ปกตูปนิสฺสโย ฯเปฯ.\csromanspacer{} ปกตูปนิสฺสโย\csromandash{}สทฺธํ อุปนิสฺสาย มานํ ชปฺเปติ.\csromanspacer{} ทิฏฺฐิํ คณฺหาติ.\csromanspacer{} สีลํ ฯเปฯ.\csromanspacer{} ปญฺญํ.\csromanspacer{} ราคํ ฯเปฯ.\csromanspacer{} มานํ.\csromanspacer{} ปตฺถนํ ฯเปฯ.\csromanspacer{} เสนาสนํ อุปนิสฺสาย ปาณํ หนติ ฯเปฯ สํฆํ ภินฺทติ.\csromanspacer{} สทฺธา ฯเปฯ.\csromanspacer{} เสนาสนํ ราคสฺส.\csromanspacer{} โทสสฺส.\csromanspacer{} โมหสฺส.\csromanspacer{} มานสฺส.\csromanspacer{} ทิฏฺฐิยา.\csromanspacer{} ปตฺถนาย อุปนิสฺสยปจฺจเยน ปจฺจโย. (2)}

\prose{คนฺถวิปฺปยุตฺโต ธมฺโม คนฺถสมฺปยุตฺตสฺส จ คนฺถวิปฺปยุตฺตสฺส จ ธมฺมสฺส อุปนิสฺสยปจฺจเยน ปจฺจโย. (ตีณิ อุปนิสฺสยา.) สทฺธํ อุปนิสฺสาย มานํ ชปฺเปติ.\csromanspacer{} สีลํ ฯเปฯ.\csromanspacer{} ปญฺญํ.\csromanspacer{} ราคํ.\csromanspacer{} โทสํ.\csromanspacer{} โมหํ.\csromanspacer{} มานํ.\csromanspacer{} ปตฺถนํ.\csromanspacer{} กายิกํ สุขํ ฯเปฯ.\csromanspacer{} เสนาสนํ อุปนิสฺสาย ปาณํ หนติ.\csromanspacer{} สํฆํ ภินฺทติ.\csromanspacer{} สทฺธา ฯเปฯ.\csromanspacer{} ปญฺญา.\csromanspacer{} ราโค.\csromanspacer{} โทโส.\csromanspacer{} โมโห.\csromanspacer{} มาโน.\csromanspacer{} ปตฺถนา ฯเปฯ.\csromanspacer{} เสนาสนํ ทิฏฺฐิคตวิปฺปยุตฺตโลภสหคตานํ ขนฺธานํ โลภสฺส จ, โทมนสฺสสหคตานํ ขนฺธานํ ปฏิฆสฺส จ อุปนิสฺสยปจฺจเยน ปจฺจโย. (3)}

\proseitem{84}{คนฺถสมฺปยุตฺโต จ คนฺถวิปฺปยุตฺโต จ ธมฺมา คนฺถ สมฺปยุตฺตสฺส ธมฺมสฺส อุปนิสฺสยปจฺจเยน ปจฺจโย.\csromanspacer{} อารมฺมณูปนิสฺสโย, อนนฺตรูปนิสฺสโย, \csromanfolio{284}ปกตูปนิสฺสโย ฯเปฯ.\csromanspacer{} ปกตูปนิสฺสโย\csromandash{}ทิฏฺฐิคตวิปฺปยุตฺตโลภสหคตา ขนฺธา จ โลโภ จ, โทมนสฺสสหคตา ขนฺธา จ ปฏิฆญฺจ คนฺถสมฺปยุตฺตกานํ ขนฺธานํ อุปนิสฺสยปจฺจเยน ปจฺจโย. (มูลํ กาตพฺพํ.) ทิฏฺฐิคตวิปฺปยุตฺตโลภสหคตา ขนฺธา จ โลโภ จ, โทมนสฺสสหคตา ขนฺธา จ ปฏิฆญฺจ คนฺถวิปฺปยุตฺตานํ ขนฺธานํ ทิฏฺฐิคตวิปฺปยุตฺตโลภสฺส จ ปฏิฆสฺส จ อุปนิสฺสยปจฺจเยน ปจฺจโย. (มูลํ ปุจฺฉิตพฺพํ.) ทิฏฺฐิคตวิปฺปยุตฺตโลภสหคตา ขนฺธา จ โลโภ จ, โทมนสฺสสหคตา ขนฺธา จ ปฏิฆญฺจ ทิฏฺฐิคตวิปฺปยุตฺตโลภสหคตานํ ขนฺธานํ โลภสฺส จ, โทมนสฺสสหคตานํ ขนฺธานํ ปฏิฆสฺส จ อุปนิสฺสยปจฺจเยน ปจฺจโย. (3)}

\csromanheader{2}{\textbf{ปุเรชาต}}
\proseitem{85}{คนฺถวิปฺปยุตฺโต ธมฺโม คนฺถวิปฺปยุตฺตสฺส ธมฺมสฺส ปุเรชาตปจฺจเยน ปจฺจโย.\csromanspacer{} \textbf{อารมฺมณปุเรชาตํ}, วตฺถุปุเรชาตํ.\csromanspacer{}\csromanspacer{}\textbf{อารมฺมณปุเรชาตํ}\csromandash{}จกฺขุํ ฯเปฯ วตฺถุํ อนิจฺจโต วิปสฺสติ, อสฺสาเทติ อภินนฺทติ, ตํ อารพฺภ คนฺถวิปฺปยุตฺโต ราโค อุปฺปชฺชติ.\csromanspacer{} วิจิกิจฺฉา ฯเปฯ อุทฺธจฺจํ ฯเปฯ โทมนสฺสํ อุปฺปชฺชติ.\csromanspacer{} ทิพฺเพน จกฺขุนา รูปํ ปสฺสติ.\csromanspacer{} ทิพฺพาย โสตธาตุยา สทฺทํ สุณาติ.\csromanspacer{} รูปายตนํ จกฺขุวิญฺญาณสฺส ฯเปฯ.\csromanspacer{} โผฏฺฐพฺพายตนํ กายวิญฺญาณสฺส ฯเปฯ.\csromanspacer{} \textbf{วตฺถุปุเรชาตํ}\csromandash{}จกฺขายตนํ จกฺขุวิญฺญาณสฺส ฯเปฯ กายายตนํ กายวิญฺญาณสฺส ฯเปฯ.\csromanspacer{} วตฺถุ คนฺถวิปฺปยุตฺตานํ ขนฺธานํ ทิฏฺฐิคตวิปฺปยุตฺตสฺส โลภสฺส จ ปฏิฆสฺส จ ปุเรชาตปจฺจเยน ปจฺจโย. (1)}

\prose{คนฺถวิปฺปยุตฺโต ธมฺโม คนฺถสมฺปยุตฺตสฺส ธมฺมสฺส ปุเรชาตปจฺจเยน ปจฺจโย.\csromanspacer{} \textbf{อารมฺมณปุเรชาตํ}, วตฺถุปุเรชาตํ.\csromanspacer{}\csromanspacer{}\textbf{อารมฺมณปุเรชาตํ}\csromandash{}จกฺขุํ ฯเปฯ วตฺถุํ อสฺสาเทติ อภินนฺทติ, ตํ อารพฺภ คนฺถสมฺปยุตฺโต ราโค อุปฺปชฺชติ.\csromanspacer{} ทิฏฺฐิ ฯเปฯ โทมนสฺสํ อุปฺปชฺชติ.\csromanspacer{}\csromanspacer{}\textbf{วตฺถุปุเรชาตํ}\csromandash{}วตฺถุ คนฺถสมฺปยุตฺตกานํ ขนฺธานํ ปุเรชาตปจฺจเยน ปจฺจโย. (2)}

\prose{คนฺถวิปฺปยุตฺโต ธมฺโม คนฺถสมฺปยุตฺตสฺส จ คนฺถวิปฺปยุตฺตสฺส จ ธมฺมสฺส ปุเรชาตปจฺจเยน ปจฺจโย.\csromanspacer{} \textbf{อารมฺมณปุเรชาตํ}, วตฺถุปุเรชาตํ.\csromanspacer{}\csromanspacer{}\textbf{อารมฺมณปุเรชาตํ}\csromandash{}จกฺขุํ ฯเปฯ วตฺถุํ อารพฺภ ทิฏฺฐิคตวิปฺปยุตฺตโลภสหคตา ขนฺธา จ โลโภ จ, โทมนสฺสสหคตา ขนฺธา จ ปฏิฆญฺจ อุปฺปชฺชนฺติ.\csromanspacer{}\csromanspacer{}\textbf{วตฺถุปุเรชาตํ}\csromandash{}วตฺถุ ทิฏฺฐิคตวิปฺปยุตฺตโลภสหคตานํ}

\vspace{\tipitakaparskip}
\csromancenter{คนฺถสมฺปยุตฺตทุก ขนฺธานํ โลภสฺส จ, โทมนสฺสสหคตานํ ขนฺธานํ ปฏิฆสฺส จ ปุเรชาตปจฺจเยน ปจฺจโย. (3)}
\csromanheader{2}{\textbf{ปจฺฉาชาต}}
\proseitem{86}{\csromanfolio{285}คนฺถสมฺปยุตฺโต ธมฺโม คนฺถวิปฺปยุตฺตสฺส ธมฺมสฺส ปจฺฉาชาตปจฺจเยน ปจฺจโย.\csromanspacer{} เอกํ.}

\prose{คนฺถวิปฺปยุตฺโต ธมฺโม คนฺถวิปฺปยุตฺตสฺส ธมฺมสฺส ปจฺฉาชาตปจฺจเยน ปจฺจโย.\csromanspacer{} เอกํ.}

\prose{คนฺถสมฺปยุตฺโต จ คนฺถวิปฺปยุตฺโต จ ธมฺมา คนฺถวิปฺปยุตฺตสฺส ธมฺมสฺส ปจฺฉาชาตปจฺจเยน ปจฺจโย.\csromanspacer{} ปจฺฉาชาตา ทิฏฺฐิคตวิปฺปยุตฺตโลภสหคตา ขนฺธา จ โลโภ จ, โทมนสฺสสหคตา ขนฺธา จ ปฏิฆญฺจ ปุเรชาตสฺส อิมสฺส กายสฺส ปจฺฉาชาตปจฺจเยน ปจฺจโย. (1)}

\csromanheader{2}{\textbf{อาเสวน}}
\proseitem{87}{คนฺถสมฺปยุตฺโต ธมฺโม คนฺถสมฺปยุตฺตสฺส ธมฺมสฺส อาเสวนปจฺจเยน ปจฺจโย. (อนนฺตรสทิสํ.\csromanspacer{} อาวชฺชนาปิ วุฏฺฐานมฺปิ นตฺถิ.)}

\csromanheader{2}{\textbf{กมฺม}}
\proseitem{88}{คนฺถสมฺปยุตฺโต ธมฺโม คนฺถสมฺปยุตฺตสฺส ธมฺมสฺส กมฺมปจฺจเยน ปจฺจโย.\csromanspacer{} คนฺถสมฺปยุตฺตา เจตนา สมฺปยุตฺตกานํ ขนฺธานํ กมฺมปจฺจเยน ปจฺจโย. (1)}

\prose{คนฺถสมฺปยุตฺโต ธมฺโม คนฺถวิปฺปยุตฺตสฺส ธมฺมสฺส กมฺมปจฺจเยน ปจฺจโย.\csromanspacer{} \textbf{สหชาตา}, นานากฺขณิกา.\csromanspacer{}\csromanspacer{}\textbf{สหชาตา} คนฺถสมฺปยุตฺตา เจตนา จิตฺตสมุฏฺฐานานํ รูปานํ กมฺมปจฺจเยน ปจฺจโย.\csromanspacer{} ทิฏฺฐิคตวิปฺปยุตฺตโลภสหคตา เจตนา โลภสฺส จ จิตฺตสมุฏฺฐานานญฺจ รูปานํ กมฺมปจฺจเยน ปจฺจโย.\csromanspacer{} โทมนสฺสสหคตา เจตนา ปฏิฆสฺส จ จิตฺตสมุฏฺฐานานญฺจ รูปานํ กมฺมปจฺจเยน ปจฺจโย.\csromanspacer{}\csromanspacer{}\textbf{นานากฺขณิกา} คนฺถสมฺปยุตฺตา เจตนา วิปากานํ ขนฺธานํ กฏตฺตา จ รูปานํ กมฺมปจฺจเยน ปจฺจโย. (2)}

\prose{คนฺถสมฺปยุตฺโต ธมฺโม คนฺถสมฺปยุตฺตสฺส จ คนฺถวิปฺปยุตฺตสฺส จ ธมฺมสฺส กมฺมปจฺจเยน ปจฺจโย.\csromanspacer{} คนฺถสมฺปยุตฺตา เจตนา สมฺปยุตฺตกานํ ขนฺธานํ \csromanfolio{286}จิตฺตสมุฏฺฐานานญฺจ รูปานํ กมฺมปจฺจเยน ปจฺจโย.\csromanspacer{} ทิฏฺฐิคตวิปฺปยุตฺตโลภสหคตา เจตนา สมฺปยุตฺตกานํ ขนฺธานํ โลภสฺส จ จิตฺตสมุฏฺฐานานญฺจ รูปานํ กมฺมปจฺจเยน ปจฺจโย.\csromanspacer{} โทมนสฺสสหคตา เจตนา สมฺปยุตฺตกานํ ขนฺธานํ ปฏิฆสฺส จ จิตฺตสมุฏฺฐานานญฺจ รูปานํ กมฺมปจฺจเยน ปจฺจโย. (3)}

\proseitem{89}{คนฺถวิปฺปยุตฺโต ธมฺโม คนฺถวิปฺปยุตฺตสฺส ธมฺมสฺส กมฺมปจฺจเยน ปจฺจโย.\csromanspacer{} \textbf{สหชาตา}, นานากฺขณิกา.\csromanspacer{}\csromanspacer{}\textbf{สหชาตา} คนฺถวิปฺปยุตฺตา เจตนา สมฺปยุตฺตกานํ ขนฺธานํ จิตฺตสมุฏฺฐานานญฺจ รูปานํ กมฺมปจฺจเยน ปจฺจโย.\csromanspacer{} ปฏิสนฺธิกฺขเณ ฯเปฯ.\csromanspacer{} \textbf{นานากฺขณิกา} คนฺถวิปฺปยุตฺตา เจตนา วิปากานํ ขนฺธานํ กฏตฺตา จ รูปานํ กมฺมปจฺจเยน ปจฺจโย. (1)}

\csromanheader{2}{\textbf{วิปากาทิ}}
\proseitem{90}{คนฺถวิปฺปยุตฺโต ธมฺโม คนฺถวิปฺปยุตฺตสฺส ธมฺมสฺส วิปากปจฺจเยน ปจฺจโย.\csromanspacer{} เอกํ.}

\prose{คนฺถสมฺปยุตฺโต ธมฺโม คนฺถสมฺปยุตฺตสฺส ธมฺมสฺส อาหารปจฺจเยน ปจฺจโย.\csromanspacer{} จตฺตาริ.\csromanspacer{} อินฺทฺริยปจฺจเยน ปจฺจโย.\csromanspacer{} จตฺตาริ.\csromanspacer{} ฌานปจฺจเยน ปจฺจโย.\csromanspacer{} จตฺตาริ.\csromanspacer{} มคฺคปจฺจเยน ปจฺจโย.\csromanspacer{} จตฺตาริ.\csromanspacer{} สมฺปยุตฺตปจฺจเยน ปจฺจโย.\csromanspacer{} ฉ.}

\csromanheader{2}{\textbf{วิปฺปยุตฺต}}
\proseitem{91}{คนฺถสมฺปยุตฺโต ธมฺโม คนฺถวิปฺปยุตฺตสฺส ธมฺมสฺส วิปฺปยุตฺตปจฺจเยน ปจฺจโย.\csromanspacer{} สหชาตํ, ปจฺฉาชาตํ. (สํขิตฺตํ.\csromanspacer{} วิภชิตพฺพํ.) (1)}

\prose{คนฺถวิปฺปยุตฺโต ธมฺโม คนฺถวิปฺปยุตฺตสฺส ธมฺมสฺส วิปฺปยุตฺตปจฺจเยน ปจฺจโย.\csromanspacer{} สหชาตํ, ปุเรชาตํ, ปจฺฉาชาตํ. (สํขิตฺตํ.) (1)}

\prose{คนฺถวิปฺปยุตฺโต ธมฺโม คนฺถสมฺปยุตฺตสฺส ธมฺมสฺส วิปฺปยุตฺตปจฺจเยน ปจฺจโย.\csromanspacer{} \textbf{ปุเรชาตํ} วตฺถุ คนฺถสมฺปยุตฺตกานํ ขนฺธานํ วิปฺปยุตฺตปจฺจเยน ปจฺจโย.\csromanspacer{}\csromanspacer{}(2)}

\prose{คนฺถวิปฺปยุตฺโต ธมฺโม คนฺถสมฺปยุตฺตสฺส จ คนฺถวิปฺปยุตฺตสฺส จ ธมฺมสฺส วิปฺปยุตฺตปจฺจเยน ปจฺจโย.\csromanspacer{} \textbf{ปุเรชาตํ} วตฺถุ ทิฏฺฐิคตวิปฺปยุตฺตโลภสหคตานํ ขนฺธานํ โลภสฺส จ, โทมนสฺสสหคตานํ ขนฺธานํ ปฏิฆสฺส จ วิปฺปยุตฺตปจฺจเยน ปจฺจโย. (3)}

\csromanheader{2}{28. คนฺถสมฺปยุตฺตทุก}
\proseitem{92}{\csromanfolio{287}คนฺถสมฺปยุตฺโต จ คนฺถวิปฺปยุตฺโต จ ธมฺมา คนฺถวิปฺปยุตฺตสฺส ธมฺมสฺส วิปฺปยุตฺตปจฺจเยน ปจฺจโย.\csromanspacer{} สหชาตํ, ปจฺฉาชาตํ.\csromanspacer{}\csromanspacer{}\textbf{สหชาตา} ทิฏฺฐิคตวิปฺปยุตฺตโลภสหคตา ขนฺธา จ โลโภ จ จิตฺตสมุฏฺฐานานํ รูปานํ วิปฺปยุตฺตปจฺจเยน ปจฺจโย.\csromanspacer{} โทมนสฺสสหคตา ขนฺธา จ ปฏิฆญฺจ จิตฺตสมุฏฺฐานานํ รูปานํ วิปฺปยุตฺตปจฺจเยน ปจฺจโย.\csromanspacer{}\csromanspacer{}\textbf{ปจฺฉาชาตา} ทิฏฺฐิคตวิปฺปยุตฺตโลภสหคตา ขนฺธา จ โลโภ จ, โทมนสฺสสหคตา ขนฺธา จ ปฏิฆญฺจ ปุเรชาตสฺส อิมสฺส กายสฺส วิปฺปยุตฺตปจฺจเยน ปจฺจโย.}

\csromanheader{2}{\textbf{อตฺถิ}}
\proseitem{93}{คนฺถสมฺปยุตฺโต ธมฺโม คนฺถสมฺปยุตฺตสฺส ธมฺมสฺส อตฺถิปจฺจเยน ปจฺจโย.\csromanspacer{} เอกํ. (ปฏิจฺจสทิสํ.) (1)}

\prose{คนฺถสมฺปยุตฺโต ธมฺโม คนฺถวิปฺปยุตฺตสฺส ธมฺมสฺส อตฺถิปจฺจเยน ปจฺจโย.\csromanspacer{} สหชาตํ, ปจฺฉาชาตํ. (สํขิตฺตํ.\csromanspacer{} วิภชิตพฺพํ.) (2)}

\prose{คนฺถสมฺปยุตฺโต ธมฺโม คนฺถสมฺปยุตฺตสฺส จ คนฺถวิปฺปยุตฺตสฺส จ ธมฺมสฺส อตฺถิปจฺจเยน ปจฺจโย.\csromanspacer{} เอกํ. (ปฏิจฺจสทิสํ.) (3)}

\proseitem{94}{คนฺถวิปฺปยุตฺโต ธมฺโม คนฺถวิปฺปยุตฺตสฺส ธมฺมสฺส อตฺถิปจฺจเยน ปจฺจโย.\csromanspacer{} สหชาตํ, ปุเรชาตํ, ปจฺฉาชาตํ, อาหารํ, อินฺทฺริยํ. (สํขิตฺตํ.\csromanspacer{} วิภชิตพฺพํ.) (1)}

\prose{คนฺถวิปฺปยุตฺโต ธมฺโม คนฺถสมฺปยุตฺตสฺส ธมฺมสฺส อตฺถิปจฺจเยน ปจฺจโย.\csromanspacer{} สหชาตํ, ปุเรชาตํ.\csromanspacer{}\csromanspacer{}\textbf{สหชาโต} ทิฏฺฐิคตวิปฺปยุตฺโต โลโภ สมฺปยุตฺตกานํ ขนฺธานํ อตฺถิปจฺจเยน ปจฺจโย.\csromanspacer{} โทมนสฺสสหคตํ ปฏิฆํ สมฺปยุตฺตกานํ ขนฺธานํ อตฺถิปจฺจเยน ปจฺจโย.\csromanspacer{}\csromanspacer{}\textbf{ปุเรชาตํ} จกฺขุํ ฯเปฯ วตฺถุํ อสฺสาเทติ อภินนฺทติ, ตํ อารพฺภ คนฺถสมฺปยุตฺโต ราโค ฯเปฯ ทิฏฺฐิ ฯเปฯ โทมนสฺสํ อุปฺปชฺชติ.\csromanspacer{} วตฺถุ คนฺถสมฺปยุตฺตกานํ ขนฺธานํ อตฺถิปจฺจเยน ปจฺจโย. (2)}

\prose{คนฺถวิปฺปยุตฺโต ธมฺโม คนฺถสมฺปยุตฺตสฺส จ คนฺถวิปฺปยุตฺตสฺส จ ธมฺมสฺส อตฺถิปจฺจเยน ปจฺจโย.\csromanspacer{} สหชาตํ, ปุเรชาตํ.\csromanspacer{}\csromanspacer{}\textbf{สหชาโต} ทิฏฺฐิคตวิปฺปยุตฺโต โลโภ สมฺปยุตฺตกานํ ขนฺธานํ จิตฺตสมุฏฺฐานานญฺจ รูปานํ อตฺถิปจฺจเยน ปจฺจโย.\csromanspacer{} ปฏิฆํ สมฺปยุตฺตกานํ ขนฺธานํ จิตฺตสมุฏฺฐานานญฺจ \csromanfolio{288}รูปานํ อตฺถิปจฺจเยน ปจฺจโย.\csromanspacer{} \textbf{ปุเรชาตํ} จกฺขุํ ฯเปฯ วตฺถุํ อารพฺภ ทิฏฺฐิคตวิปฺปยุตฺตโลภสหคตา ขนฺธา จ โลโภ จ, โทมนสฺสสหคตา ขนฺธา จ ปฏิฆญฺจ อุปฺปชฺชนฺติ.\csromanspacer{} วตฺถุ ทิฏฺฐิคตวิปฺปยุตฺตโลภสหคตานํ ขนฺธานํ โลภสฺส จ, โทมนสฺสสหคตานํ ขนฺธานํ ปฏิฆสฺส จ อตฺถิปจฺจเยน ปจฺจโย. (3)}

\proseitem{95}{คนฺถสมฺปยุตฺโต จ คนฺถวิปฺปยุตฺโต จ ธมฺมา คนฺถสมฺปยุตฺตสฺส ธมฺมสฺส อตฺถิปจฺจเยน ปจฺจโย.\csromanspacer{} สหชาตํ, ปุเรชาตํ.\csromanspacer{}\csromanspacer{}\textbf{สหชาโต} คนฺถสมฺปยุตฺโต เอโก ขนฺโธ จ วตฺถุ จ ติณฺณนฺนํ ขนฺธานํ อตฺถิปจฺจเยน ปจฺจโย ฯเปฯ.\csromanspacer{} เทฺว ขนฺธา จ ฯเปฯ.\csromanspacer{} \textbf{สหชาโต} ทิฏฺฐิคตวิปฺปยุตฺตโลภสหคโต เอโก ขนฺโธ จ โลโภ จ ติณฺณนฺนํ ขนฺธานํ อตฺถิปจฺจเยน ปจฺจโย ฯเปฯ เทฺว ขนฺธา จ ฯเปฯ.\csromanspacer{} โทมนสฺสสหคโต เอโก ขนฺโธ จ ปฏิฆญฺจ ติณฺณนฺนํ ขนฺธานํ อตฺถิปจฺจเยน ปจฺจโย ฯเปฯ เทฺว ขนฺธา จ ฯเปฯ. (1)}

\prose{คนฺถสมฺปยุตฺโต จ คนฺถวิปฺปยุตฺโต จ ธมฺมา คนฺถวิปฺปยุตฺตสฺส ธมฺมสฺส อตฺถิปจฺจเยน ปจฺจโย.\csromanspacer{} สหชาตํ, ปุเรชาตํ, ปจฺฉาชาตํ, อาหารํ, อินฺทฺริยํ.\csromanspacer{}\csromanspacer{}\textbf{สหชาตา} คนฺถสมฺปยุตฺตา ขนฺธา จ มหาภูตา จ จิตฺตสมุฏฺฐานานํ รูปานํ อตฺถิปจฺจเยน ปจฺจโย.\csromanspacer{} \textbf{สหชาตา} ทิฏฺฐิคตวิปฺปยุตฺตโลภสหคตา ขนฺธา จ โลโภ จ จิตฺตสมุฏฺฐานานํ รูปานํ อตฺถิปจฺจเยน ปจฺจโย.\csromanspacer{} \textbf{สหชาตา} โทมนสฺสสหคตา ขนฺธา จ ปฏิฆญฺจ จิตฺตสมุฏฺฐานานํ รูปานํ อตฺถิปจฺจเยน ปจฺจโย.\csromanspacer{} ทิฏฺฐิคตวิปฺปยุตฺตโลภสหคตา ขนฺธา จ วตฺถุ จ โลภสฺส อตฺถิปจฺจเยน ปจฺจโย.\csromanspacer{} โทมนสฺสสหคตา ขนฺธา จ วตฺถุ จ ปฏิฆสฺส อตฺถิปจฺจเยน ปจฺจโย.\csromanspacer{} \textbf{ปจฺฉาชาตา} ทิฏฺฐิคตวิปฺปยุตฺตโลภสหคตา ขนฺธา จ โลโภ จ โทมนสฺสสหคตา ขนฺธา จ ปฏิฆญฺจ ปุเรชาตสฺส อิมสฺส กายสฺส อตฺถิปจฺจเยน ปจฺจโย.\csromanspacer{} \textbf{ปจฺฉาชาตา} คนฺถสมฺปยุตฺตา ขนฺธา จ กพฬีกาโร อาหาโร จ อิมสฺส กายสฺส อตฺถิปจฺจเยน ปจฺจโย.\csromanspacer{} \textbf{ปจฺฉาชาตา} คนฺถสมฺปยุตฺตา ขนฺธา จ รูปชีวิตินฺทฺริยญฺจ กฏตฺตารูปานํ อตฺถิปจฺจเยน ปจฺจโย. (2)}

\prose{คนฺถสมฺปยุตฺโต จ คนฺถวิปฺปยุตฺโต จ ธมฺมา คนฺถสมฺปยุตฺตสฺส จ คนฺถวิปฺปยุตฺตสฺส จ ธมฺมสฺส อตฺถิปจฺจเยน ปจฺจโย.\csromanspacer{} สหชาตํ, ปุเรชาตํ.\csromanspacer{}\csromanspacer{}\textbf{สหชาโต} ทิฏฺฐิคตวิปฺปยุตฺตโลภสหคโต เอโก ขนฺโธ จ}

\vspace{\tipitakaparskip}
\csromancenter{คนฺถสมฺปยุตฺตทุก โลโภ จ ติณฺณนฺนํ ขนฺธานํ จิตฺตสมุฏฺฐานานญฺจ รูปานํ อตฺถิปจฺจเยน ปจฺจโย ฯเปฯ เทฺว ขนฺธา จ ฯเปฯ.\csromanspacer{} \textbf{สหชาโต} โทมนสฺสสหคโต เอโก ขนฺโธ จ ปฏิฆญฺจ ติณฺณนฺนํ ขนฺธานํ จิตฺตสมุฏฺฐานานญฺจ รูปานํ อตฺถิปจฺจเยน ปจฺจโย ฯเปฯ เทฺว ขนฺธา จ ฯเปฯ.\csromanspacer{} \textbf{สหชาโต} ทิฏฺฐิคตวิปฺปยุตฺตโลภสหคโต เอโก ขนฺโธ จ วตฺถุ จ ติณฺณนฺนํ ขนฺธานํ โลภสฺส จ อตฺถิปจฺจเยน ปจฺจโย ฯเปฯ เทฺว ขนฺธา จ ฯเปฯ.\csromanspacer{} \textbf{สหชาโต} โทมนสฺสสหคโต เอโก ขนฺโธ จ วตฺถุ จ ติณฺณนฺนํ ขนฺธานํ ปฏิฆสฺส จ อตฺถิปจฺจเยน ปจฺจโย ฯเปฯ เทฺว ขนฺธา จ ฯเปฯ. (3)}
\csromanheader{2}{1. ปจฺจยานุโลม 2. สงฺขฺยาวาร}
\csromanheader{2}{\textbf{สุทฺธ}}
\proseitem{96}{\csromanfolio{289}เหตุยา นว, อารมฺมเณ นว, อธิปติยา นว, อนนฺตเร นว, สมนนฺตเร นว, สหชาเต นว, อญฺญมญฺเญ ฉ, นิสฺสเย นว, อุปนิสฺสเย นว, ปุเรชาเต ตีณิ, ปจฺฉาชาเต ตีณิ, อาเสวเน นว, กมฺเม จตฺตาริ, วิปาเก เอกํ, อาหาเร จตฺตาริ, อินฺทฺริเย จตฺตาริ, ฌาเน จตฺตาริ, มคฺเค จตฺตาริ, สมฺปยุตฺเต ฉ, วิปฺปยุตฺเต ปญฺจ, อตฺถิยา นว, นตฺถิยา นว, วิคเต นว, อวิคเต นว.}

\vspace{\tipitakaparskip}
\csromancenter{อนุโลมํ.}
\csromanheader{2}{2. ปจฺจนียุทฺธาร}
\proseitem{97}{คนฺถสมฺปยุตฺโต ธมฺโม คนฺถสมฺปยุตฺตสฺส ธมฺมสฺส อารมฺมณปจฺจเยน ปจฺจโย.\csromanspacer{} สหชาตปจฺจเยน ปจฺจโย.\csromanspacer{} อุปนิสฺสยปจฺจเยน ปจฺจโย. (1)}

\prose{คนฺถสมฺปยุตฺโต ธมฺโม คนฺถวิปฺปยุตฺตสฺส ธมฺมสฺส อารมฺมณปจฺจเยน ปจฺจโย.\csromanspacer{} สหชาตปจฺจเยน ปจฺจโย.\csromanspacer{} อุปนิสฺสยปจฺจเยน ปจฺจโย.\csromanspacer{} ปจฺฉาชาตปจฺจเยน ปจฺจโย.\csromanspacer{} กมฺมปจฺจเยน ปจฺจโย. (2)}

\prose{คนฺถสมฺปยุตฺโต ธมฺโม คนฺถสมฺปยุตฺตสฺส จ คนฺถวิปฺปยุตฺตสฺส จ ธมฺมสฺส อารมฺมณปจฺจเยน ปจฺจโย.\csromanspacer{} สหชาตปจฺจเยน ปจฺจโย.\csromanspacer{} อุปนิสฺสยปจฺจเยน ปจฺจโย. (93)}

\proseitem{98}{\csromanfolio{290}คนฺถวิปฺปยุตฺโต ธมฺโม คนฺถวิปฺปยุตฺตสฺส ธมฺมสฺส อารมฺมณปจฺจเยน ปจฺจโย.\csromanspacer{} สหชาตปจฺจเยน ปจฺจโย.\csromanspacer{} อุปนิสฺสยปจฺจเยน ปจฺจโย.\csromanspacer{} ปุเรชาตปจฺจเยน ปจฺจโย.\csromanspacer{} ปจฺฉาชาตปจฺจเยน ปจฺจโย.\csromanspacer{} กมฺมปจฺจเยน ปจฺจโย.\csromanspacer{} อาหารปจฺจเยน ปจฺจโย.\csromanspacer{} อินฺทฺริยปจฺจเยน ปจฺจโย. (1)}

\prose{คนฺถวิปฺปยุตฺโต ธมฺโม คนฺถสมฺปยุตฺตสฺส ธมฺมสฺส อารมฺมณปจฺจเยน ปจฺจโย.\csromanspacer{} สหชาตปจฺจเยน ปจฺจโย.\csromanspacer{} อุปนิสฺสยปจฺจเยน ปจฺจโย.\csromanspacer{} ปุเรชาตปจฺจเยน ปจฺจโย. (2)}

\prose{คนฺถวิปฺปยุตฺโต ธมฺโม คนฺถสมฺปยุตฺตสฺส จ คนฺถวิปฺปยุตฺตสฺส จ ธมฺมสฺส อารมฺมณปจฺจเยน ปจฺจโย.\csromanspacer{} สหชาตปจฺจเยน ปจฺจโย.\csromanspacer{} อุปนิสฺสยปจฺจเยน ปจฺจโย.\csromanspacer{} ปุเรชาตปจฺจเยน ปจฺจโย. (3)}

\proseitem{99}{คนฺถสมฺปยุตฺโต จ คนฺถวิปฺปยุตฺโต จ ธมฺมา คนฺถสมฺปยุตฺตสฺส ธมฺมสฺส อารมฺมณปจฺจเยน ปจฺจโย.\csromanspacer{} สหชาตปจฺจเยน ปจฺจโย.\csromanspacer{} อุปนิสฺสยปจฺจเยน ปจฺจโย. (1)}

\prose{คนฺถสมฺปยุตฺโต จ คนฺถวิปฺปยุตฺโต จ ธมฺมา คนฺถวิปฺปยุตฺตสฺส ธมฺมสฺส อารมฺมณปจฺจเยน ปจฺจโย.\csromanspacer{} สหชาตปจฺจเยน ปจฺจโย.\csromanspacer{} อุปนิสฺสยปจฺจเยน ปจฺจโย.\csromanspacer{} ปจฺฉาชาตปจฺจเยน ปจฺจโย. (2)}

\prose{คนฺถสมฺปยุตฺโต จ คนฺถวิปฺปยุตฺโต จ ธมฺมา คนฺถสมฺปยุตฺตสฺส จ คนฺถวิปฺปยุตฺตสฺส จ ธมฺมสฺส อารมฺมณปจฺจเยน ปจฺจโย.\csromanspacer{} สหชาตปจฺจเยน ปจฺจโย.\csromanspacer{} อุปนิสฺสยปจฺจเยน ปจฺจโย. (3)}

\csromanheader{2}{2. ปจฺจยปจฺจนีย 2. สงฺขฺยาวาร}
\vspace{\tipitakaparskip}
\csromancenter{นเหตุยา นว, น-อารมฺมเณ นว, (สพฺพตฺถ นว,) โน-อวิคเต นว.}
\csromanheader{2}{3. ปจฺจยานุโลมปจฺจนีย}
\csromanheader{2}{\textbf{เหตุทุก}}
\proseitem{101}{เหตุปจฺจยา น-อารมฺมเณ นว, น-อธิปติยา นว ฯเปฯ นสมนนฺตเร นว, น-อญฺญมญฺเญ ตีณิ, น-อุปนิสฺสเย นว ฯเปฯ นมคฺเค}

\vspace{\tipitakaparskip}
\csromancenter{คนฺถคนฺถนิยทุก นว, นสมฺปยุตฺเต ตีณิ, นวิปฺปยุตฺเต ฉ, โนนตฺถิยา นว, โนวิคเต นว.}
\csromanheader{2}{4. ปจฺจยปจฺจนียานุโลม}
\csromanheader{2}{\textbf{นเหตุทุก}}
\proseitem{102}{\csromanfolio{291}นเหตุปจฺจยา อารมฺมเณ นว, อธิปติยา นว. (อนุโลมมาติกา วิตฺถาเรตพฺพา.) ฯเปฯ อวิคเต นว.}

\nitthitam{คนฺถสมฺปยุตฺตทุกํ นิฏฺฐิตํ.}
\csromanmatikaanchor{mtk.35}
\csromantocmarkref{section}{29. คนฺถคนฺถนิยทุก}{mtk.35}
\bookmark[dest={mtk.35},level=1]{29. คนฺถคนฺถนิยทุก}
\csromanheader{1}{29. คนฺถคนฺถนิยทุก 1. ปฏิจฺจวาร}
\csromanheader{2}{\textbf{เหตุ}}
\proseitem{103}{คนฺถญฺเจว คนฺถนิยญฺจ ธมฺมํ ปฏิจฺจ คนฺโถ เจว คนฺถนิโย จ ธมฺโม อุปฺปชฺชติ เหตุปจฺจยา.\csromanspacer{} สีลพฺพตปรามาสํ กายคนฺถํ ปฏิจฺจ อภิชฺฌากายคนฺโถ, อภิชฺฌากายคนฺถํ ปฏิจฺจ สีลพฺพตปรามาโส กายคนฺโถ, อิทํสจฺจาภินิเวสกายคนฺถํ ปฏิจฺจ อภิชฺฌากายคนฺโถ, อภิชฺฌากายคนฺถํ ปฏิจฺจ อิทํสจฺจาภินิเวสกายคนฺโถ. (1)}

\prose{คนฺถญฺเจว คนฺถนิยญฺจ ธมฺมํ ปฏิจฺจ คนฺถนิโย เจว โน จ คนฺโถ ธมฺโม อุปฺปชฺชติ เหตุปจฺจยา.\csromanspacer{} คนฺเถ ปฏิจฺจ สมฺปยุตฺตกา ขนฺธา จิตฺตสมุฏฺฐานญฺจ รูปํ. (2)}

\prose{คนฺถญฺเจว คนฺถนิยญฺจ ธมฺมํ ปฏิจฺจ คนฺโถ เจว คนฺถนิโย จ คนฺถนิโย เจว โน จ คนฺโถ จ ธมฺมา อุปฺปชฺชนฺติ เหตุปจฺจยา. (1)}

\prose{(ปฏิจฺจวารมฺปิ สหชาตวารมฺปิ ปจฺจยวารมฺปิ นิสฺสยวารมฺปิ สํสฏฺฐวารมฺปิ สมฺปยุตฺตวารมฺปิ คนฺถทุกสทิสํ, นินฺนานากรณํ.)}

\csromanheader{2}{29. คนฺถคนฺถนิยทุก 7. ปญฺหาวาร}
\csromanheader{2}{1. ปจฺจยานุโลม 1. วิภงฺควาร}
\csromanheader{2}{\textbf{เหตุ}}
\proseitem{104}{\csromanfolio{292}คนฺโถ เจว คนฺถนิโย จ ธมฺโม คนฺถสฺส เจว คนฺถนิยสฺส จ ธมฺมสฺส เหตุปจฺจเยน ปจฺจโย.\csromanspacer{} คนฺถา เหตู สมฺปยุตฺตกานํ คนฺถานํ เหตุปจฺจเยน ปจฺจโย. (เอวํ นว ปญฺหา วิตฺถาเรตพฺพา.)}

\csromanheader{2}{\textbf{อารมฺมณ}}
\proseitem{105}{คนฺโถ เจว คนฺถนิโย จ ธมฺโม คนฺถสฺส เจว คนฺถนิยสฺส จ ธมฺมสฺส อารมฺมณปจฺจเยน ปจฺจโย.\csromanspacer{} คนฺเถ อารพฺภ คนฺถา อุปฺปชฺชนฺติ. (มูลํ ปุจฺฉิตพฺพํ.) คนฺเถ อารพฺภ คนฺถนิยา เจว โน จ คนฺถา ขนฺธา อุปฺปชฺชนฺติ. (มูลํ ปุจฺฉิตพฺพํ.) คนฺเถ อารพฺภ คนฺถา จ สมฺปยุตฺตกา จ ขนฺธา อุปฺปชฺชนฺติ. (3)}

\proseitem{106}{คนฺถนิโย เจว โน จ คนฺโถ ธมฺโม คนฺถนิยสฺส เจว โน จ คนฺถสฺส จ ธมฺมสฺส อารมฺมณปจฺจเยน ปจฺจโย.\csromanspacer{} ทานํ ฯเปฯ สีลํ ฯเปฯ อุโปสถกมฺมํ กตฺวา ตํ ปจฺจเวกฺขติ.\csromanspacer{} ปุพฺเพ สุจิณฺณานิ ฯเปฯ.\csromanspacer{} ฌานา ฯเปฯ.\csromanspacer{} อริยา โคตฺรภุํ ปจฺจเวกฺขนฺติ, โวทานํ ปจฺจเวกฺขนฺติ.\csromanspacer{} ปหีเน กิเลเส ฯเปฯ.\csromanspacer{} วิกฺขมฺภิเต กิเลเส ปจฺจเวกฺขนฺติ.\csromanspacer{} ปุพฺเพ สมุทาจิณฺเณ กิเลเส ชานนฺติ.\csromanspacer{} จกฺขุํ ฯเปฯ วตฺถุํ คนฺถนิเย เจว โน จ คนฺเถ ขนฺเธ อนิจฺจโต ฯเปฯ โทมนสฺสํ อุปฺปชฺชติ.\csromanspacer{} ทิพฺเพน จกฺขุนา รูปํ ปสฺสติ.\csromanspacer{} ทิพฺพาย โสตธาตุยา สทฺทํ สุณาติ. (สพฺพํ วิตฺถาเรตพฺพํ.) อาวชฺชนาย อารมฺมณปจฺจเยน ปจฺจโย. (1)}

\prose{คนฺถนิโย เจว โน จ คนฺโถ ธมฺโม คนฺถสฺส เจว คนฺถนิยสฺส จ ธมฺมสฺส อารมฺมณปจฺจเยน ปจฺจโย.\csromanspacer{} ทานํ ฯเปฯ สีลํ ฯเปฯ อุโปสถกมฺมํ กตฺวา ตํ อสฺสาเทติ อภินนฺทติ, ตํ อารพฺภ ราโค อุปฺปชฺชติ, ทิฏฺฐิ ฯเปฯ วิจิกิจฺฉา ฯเปฯ โทมนสฺสํ อุปฺปชฺชติ.\csromanspacer{} ปุพฺเพ สุจิณฺณานิ ฯเปฯ.\csromanspacer{} ฌานา วุฏฺฐหิตฺวา ฌานํ ฯเปฯ.\csromanspacer{} จกฺขุํ ฯเปฯ วตฺถุํ คนฺถนิเย เจว โน จ คนฺเถ ขนฺเธ อสฺสาเทติ อภินนฺทติ, ตํ อารพฺภ ราโค อุปฺปชฺชติ, ทิฏฺฐิ ฯเปฯ โทมนสฺสํ อุปฺปชฺชติ. (2)}

\csromanmatikaanchor{mtk.36}
\csromantocmarkref{section}{30. คนฺถคนฺถสมฺปยุตฺตทุก}{mtk.36}
\bookmark[dest={mtk.36},level=1]{30. คนฺถคนฺถสมฺปยุตฺตทุก}
\csromanheader{1}{30. คนฺถคนฺถสมฺปยุตฺตทุก}
\prose{\csromanfolio{293}คนฺถนิโย เจว โน จ คนฺโถ ธมฺโม คนฺถสฺส เจว คนฺถนิยสฺส จ คนฺถนิยสฺส เจว โน จ คนฺถสฺส จ ธมฺมสฺส อารมฺมณปจฺจเยน ปจฺจโย.\csromanspacer{} ทานํ ฯเปฯ สีลํ ฯเปฯ อุโปสถกมฺมํ ฯเปฯ.\csromanspacer{} ปุพฺเพ สุจิณฺณานิ ฯเปฯ.\csromanspacer{} ฌานา วุฏฺฐหิตฺวา ฌานํ ฯเปฯ.\csromanspacer{} จกฺขุํ ฯเปฯ วตฺถุํ คนฺถนิเย เจว โน จ คนฺเถ ขนฺเธ อสฺสาเทติ อภินนฺทติ, ตํ อารพฺภ คนฺถา จ สมฺปยุตฺตกา จ ขนฺธา อุปฺปชฺชนฺติ. (เอวํ อิตเรปิ ตีณิ วิตฺถาเรตพฺพา.) (3)}

\prose{(อารพฺภ กาตพฺพา.\csromanspacer{} อิมสฺมิํ ทุเก โลกุตฺตรํ นตฺถิ, คนฺถทุกสทิสํ, นินฺนานากรณํ. “คนฺถนิยนฺ”ติ นิยาเมตพฺพํ.\csromanspacer{}\csromanspacer{}มคฺเค นว ปญฺหา กาตพฺพา.)}

\nitthitam{คนฺถคนฺถนิยทุกํ นิฏฺฐิตํ.}
\csromanheader{2}{30. คนฺถคนฺถสมฺปยุตฺตทุก 1. ปฏิจฺจวาร}
\csromanheader{2}{1-4. ปจฺจยจตุกฺก}
\csromanheader{2}{\textbf{เหตุ}}
\proseitem{107}{คนฺถญฺเจว คนฺถสมฺปยุตฺตญฺจ ธมฺมํ ปฏิจฺจ คนฺโถ เจว คนฺถสมฺปยุตฺโต จ ธมฺโม อุปฺปชฺชติ เหตุปจฺจยา.\csromanspacer{} สีลพฺพตปรามาสํ กายคนฺถํ ปฏิจฺจ อภิชฺฌากายคนฺโถ, อภิชฺฌากายคนฺถํ ปฏิจฺจ สีลพฺพตปรามาโส กายคนฺโถ, อิทํสจฺจาภินิเวสกายคนฺถํ ปฏิจฺจ อภิชฺฌากายคนฺโถ, อภิชฺฌากายคนฺถํ ปฏิจฺจ อิทํสจฺจาภินิเวโส กายคนฺโถ. (1)}

\prose{คนฺถญฺเจว คนฺถสมฺปยุตฺตญฺจ ธมฺมํ ปฏิจฺจ คนฺถสมฺปยุตฺโต เจว โน จ คนฺโถ ธมฺโม อุปฺปชฺชติ เหตุปจฺจยา.\csromanspacer{} คนฺเถ ปฏิจฺจ สมฺปยุตฺตกา ขนฺธา. (2)}

\prose{คนฺถญฺเจว คนฺถสมฺปยุตฺตญฺจ ธมฺมํ ปฏิจฺจ คนฺโถ เจว คนฺถสมฺปยุตฺโต จ คนฺถสมฺปยุตฺโต เจว โน จ คนฺโถ จ ธมฺมา อุปฺปชฺชนฺติ เหตุปจฺจยา.\csromanspacer{} สีลพฺพตปรามาสํ กายคนฺถํ ปฏิจฺจ อภิชฺฌากายคนฺโถ สมฺปยุตฺตกา จ ขนฺธา. (จกฺกํ.) (3)}

\proseitem{108}{คนฺถสมฺปยุตฺตญฺเจว โน จ คนฺถํ ธมฺมํ ปฏิจฺจ คนฺถสมฺปยุตฺโต เจว โน จ คนฺโถ ธมฺโม อุปฺปชฺชติ เหตุปจฺจยา.\csromanspacer{} คนฺถสมฺปยุตฺตญฺเจว โน จ คนฺถํ เอกํ ขนฺธํ ปฏิจฺจ ตโย ขนฺธา ฯเปฯ เทฺว ขนฺเธ ฯเปฯ. (1)}

\prose{\csromanfolio{294}คนฺถสมฺปยุตฺตญฺเจว โน จ คนฺถํ ธมฺมํ ปฏิจฺจ คนฺโถ เจว คนฺถสมฺปยุตฺโต จ ธมฺโม อุปฺปชฺชติ เหตุปจฺจยา.\csromanspacer{} คนฺถสมฺปยุตฺเต เจว โน จ คนฺเถ ขนฺเธ ปฏิจฺจ คนฺถา. (2)}

\prose{คนฺถสมฺปยุตฺตญฺเจว โน จ คนฺถํ ธมฺมํ ปฏิจฺจ คนฺโถ เจว คนฺถสมฺปยุตฺโต จ คนฺถสมฺปยุตฺโต เจว โน จ คนฺโถ จ ธมฺมา อุปฺปชฺชนฺติ เหตุปจฺจยา.\csromanspacer{} คนฺถสมฺปยุตฺตญฺเจว โน จ คนฺถํ เอกํ ขนฺธํ ปฏิจฺจ ตโย ขนฺธา คนฺถา จ ฯเปฯ เทฺว ขนฺเธ ฯเปฯ. (3)}

\proseitem{109}{คนฺถญฺเจว คนฺถสมฺปยุตฺตญฺจ คนฺถสมฺปยุตฺตญฺเจว โน จ คนฺถญฺจ ธมฺมํ ปฏิจฺจ คนฺโถ เจว คนฺถสมฺปยุตฺโต จ ธมฺโม อุปฺปชฺชติ เหตุปจฺจยา.\csromanspacer{} คนฺเถ จ สมฺปยุตฺตเก จ ขนฺเธ ปฏิจฺจ คนฺถา. (1)}

\prose{คนฺถญฺเจว คนฺถสมฺปยุตฺตญฺจ คนฺถสมฺปยุตฺตญฺเจว โน จ คนฺถญฺจ ธมฺมํ ปฏิจฺจ คนฺถสมฺปยุตฺโต เจว โน จ คนฺโถ ธมฺโม อุปฺปชฺชติ เหตุปจฺจยา.\csromanspacer{} คนฺถสมฺปยุตฺตญฺเจว โน จ คนฺถํ เอกํ ขนฺธญฺจ คนฺเถ จ ปฏิจฺจ ตโย ขนฺธา ฯเปฯ เทฺว ขนฺเธ จ ฯเปฯ. (2)}

\prose{คนฺถญฺเจว คนฺถสมฺปยุตฺตญฺจ คนฺถสมฺปยุตฺตญฺเจว โน จ คนฺถญฺจ ธมฺมํ ปฏิจฺจ คนฺโถ เจว คนฺถสมฺปยุตฺโต จ คนฺถสมฺปยุตฺโต เจว โน จ คนฺโถ จ ธมฺมา อุปฺปชฺชนฺติ เหตุปจฺจยา.\csromanspacer{} คนฺถสมฺปยุตฺตญฺเจว โน จ คนฺถํ เอกํ ขนฺธญฺจ สีลพฺพตปรามาสํ กายคนฺถญฺจ ปฏิจฺจ ตโย ขนฺธา อภิชฺฌากายคนฺโถ จ ฯเปฯ เทฺว ขนฺเธ จ ฯเปฯ. (จกฺกํ พนฺธิตพฺพํ.\csromanspacer{} สํขิตฺตํ.) (3)}

\proseitem{110}{เหตุยา นว, อารมฺมเณ นว, (สพฺพตฺถ นว,) กมฺเม นว, อาหาเร นว ฯเปฯ อวิคเต นว.}

\vspace{\tipitakaparskip}
\csromancenter{อนุโลมํ.}
\csromanheader{2}{2. ปจฺจนีย}
\proseitem{111}{คนฺถญฺเจว คนฺถสมฺปยุตฺตญฺจ ธมฺมํ ปฏิจฺจ คนฺโถ เจว คนฺถสมฺปยุตฺโต จ ธมฺโม อุปฺปชฺชติ น-อธิปติปจฺจยา. (สํขิตฺตํ.)}

\csromanheader{2}{30. คนฺถคนฺถสมฺปยุตฺตทุก}
\prose{\csromanfolio{295}(อิธ นเหตุปจฺจโย นตฺถิ.) น-อธิปติยา นว, นปุเรชาเต นว, นปจฺฉาชาเต นว, น-อาเสวเน นว, นกมฺเม ตีณิ, นวิปาเก นว, นวิปฺปยุตฺเต นว.}

\prose{(เอวํ อิตเร เทฺว คณนาปิ สหชาตวาโรปิ ปจฺจยวาโรปิ นิสฺสยวาโรปิ สํสฏฺฐวาโรปิ สมฺปยุตฺตวาโรปิ ปฏิจฺจวารสทิสา.)}

\csromanheader{2}{30. คนฺถคนฺถสมฺปยุตฺตทุก 7. ปญฺหาวาร}
\csromanheader{2}{1. ปจฺจยานุโลม 1. วิภงฺควาร}
\csromanheader{2}{\textbf{เหตุ}}
\proseitem{112}{คนฺโถ เจว คนฺถสมฺปยุตฺโต จ ธมฺโม คนฺถสฺส เจว คนฺถสมฺปยุตฺตสฺส จ ธมฺมสฺส เหตุปจฺจเยน ปจฺจโย.\csromanspacer{} คนฺถา เจว คนฺถสมฺปยุตฺตา จ เหตู สมฺปยุตฺตกานํ คนฺถานํ เหตุปจฺจเยน ปจฺจโย. (1)}

\prose{คนฺโถ เจว คนฺถสมฺปยุตฺโต จ ธมฺโม คนฺถสมฺปยุตฺตสฺส เจว โน จ คนฺถสฺส ธมฺมสฺส เหตุปจฺจเยน ปจฺจโย.\csromanspacer{} คนฺถา เจว คนฺถสมฺปยุตฺตา จ เหตู สมฺปยุตฺตกานํ ขนฺธานํ เหตุปจฺจเยน ปจฺจโย. (2)}

\prose{คนฺโถ เจว คนฺถสมฺปยุตฺโต จ ธมฺโม คนฺถสฺส เจว คนฺถสมฺปยุตฺตสฺส จ คนฺถสมฺปยุตฺตสฺส เจว โน จ คนฺถสฺส จ ธมฺมสฺส เหตุปจฺจเยน ปจฺจโย.\csromanspacer{} คนฺถา เจว คนฺถสมฺปยุตฺตา จ เหตู สมฺปยุตฺตกานํ ขนฺธานํ คนฺถานญฺจ เหตุปจฺจเยน ปจฺจโย. (3)}

\proseitem{113}{คนฺถสมฺปยุตฺโต เจว โน จ คนฺโถ ธมฺโม คนฺถสมฺปยุตฺตสฺส เจว โน จ คนฺถสฺส ธมฺมสฺส เหตุปจฺจเยน ปจฺจโย.\csromanspacer{} คนฺถสมฺปยุตฺตา เจว โน จ คนฺถา เหตู สมฺปยุตฺตกานํ คนฺถานํ เหตุปจฺจเยน ปจฺจโย. (1)}

\prose{คนฺถสมฺปยุตฺโต เจว โน จ คนฺโถ ธมฺโม คนฺถสฺส เจว คนฺถสมฺปยุตฺตสฺส จ ธมฺมสฺส เหตุปจฺจเยน ปจฺจโย.\csromanspacer{} คนฺถสมฺปยุตฺตา เจว โนจ คนฺถา เหตู สมฺปยุตฺตกานํ คนฺถานํ เหตุปจฺจเยน ปจฺจโย. (2)}

\prose{คนฺถสมฺปยุตฺโต เจว โน จ คนฺโถ ธมฺโม คนฺถสฺส เจว คนฺถสมฺปยุตฺตสฺส จ คนฺถสมฺปยุตฺตสฺส เจว โน จ คนฺถสฺส จ ธมฺมสฺส \csromanfolio{296}เหตุปจฺจเยน ปจฺจโย.\csromanspacer{} คนฺถสมฺปยุตฺตา เจว โน จ คนฺถา เหตู สมฺปยุตฺตกานํ ขนฺธานํ คนฺถานญฺจ เหตุปจฺจเยน ปจฺจโย. (3)}

\proseitem{114}{คนฺโถ เจว คนฺถสมฺปยุตฺโต จ คนฺถสมฺปยุตฺโต เจว โน จ คนฺโถ จ ธมฺมา คนฺถสฺส เจว คนฺถสมฺปยุตฺตสฺส จ ธมฺมสฺส เหตุปจฺจเยน ปจฺจโย.\csromanspacer{} คนฺถา เจว คนฺถสมฺปยุตฺตา จ คนฺถสมฺปยุตฺตา เจว โน จ คนฺถา จ เหตู สมฺปยุตฺตกานํ ขนฺธานํ เหตุปจฺจเยน ปจฺจโย. (1)}

\prose{คนฺโถ เจว คนฺถสมฺปยุตฺโต จ คนฺถสมฺปยุตฺโต เจว โน จ คนฺโถ จ ธมฺมา คนฺถสมฺปยุตฺตสฺส เจว โน จ คนฺถสฺส ธมฺมสฺส เหตุปจฺจเยน ปจฺจโย.\csromanspacer{} คนฺถา เจว คนฺถสมฺปยุตฺตา จ คนฺถสมฺปยุตฺตา เจว โน จ คนฺถา จ เหตู สมฺปยุตฺตกานํ ขนฺธานํ เหตุปจฺจเยน ปจฺจโย. (2)}

\prose{คนฺโถ เจว คนฺถสมฺปยุตฺโต จ คนฺถสมฺปยุตฺโต เจว โน จ คนฺโถ จ ธมฺมา คนฺถสฺส เจว คนฺถสมฺปยุตฺตสฺส จ คนฺถสมฺปยุตฺตสฺส เจว โน จ คนฺถสฺส จ ธมฺมสฺส เหตุปจฺจเยน ปจฺจโย.\csromanspacer{} คนฺถา เจว คนฺถสมฺปยุตฺตา จ คนฺถสมฺปยุตฺตา เจว โน จ คนฺถา จ เหตู สมฺปยุตฺตกานํ ขนฺธานํ คนฺถานญฺจ เหตุปจฺจเยน ปจฺจโย. (3)}

\csromanheader{2}{\textbf{อารมฺมณาทิ}}
\proseitem{115}{คนฺโถ เจว คนฺถสมฺปยุตฺโต จ ธมฺโม คนฺถสฺส เจว คนฺถสมฺปยุตฺตสฺส จ ธมฺมสฺส อารมฺมณปจฺจเยน ปจฺจโย.\csromanspacer{} คนฺเถ อารพฺภ คนฺถา อุปฺปชฺชนฺติ. (มูลํ กาตพฺพํ.) คนฺเถ อารพฺภ คนฺถสมฺปยุตฺตา เจว โน จ คนฺถา ขนฺธา อุปฺปชฺชนฺติ. (มูลํ กาตพฺพํ.).\csromanspacer{} คนฺเถ อารพฺภ คนฺถา จ คนฺถสมฺปยุตฺตกา จ ขนฺธา อุปฺปชฺชนฺติ. (3)}

\prose{คนฺถสมฺปยุตฺโต เจว โน จ คนฺโถ ธมฺโม คนฺถสมฺปยุตฺตสฺส เจว โน จ คนฺถสฺส ธมฺมสฺส อารมฺมณปจฺจเยน ปจฺจโย.\csromanspacer{} คนฺถสมฺปยุตฺเต เจว โน จ คนฺเถ ขนฺเธ อารพฺภ คนฺถสมฺปยุตฺตา เจว โน จ คนฺถา ขนฺธา อุปฺปชฺชนฺติ. (มูลํ กาตพฺพํ.) คนฺถสมฺปยุตฺเต เจว โน จ คนฺเถ ขนฺเธ อารพฺภ คนฺถา อุปฺปชฺชนฺติ. (มูลํ กาตพฺพํ.) คนฺถสมฺปยุตฺเต เจว โน จ คนฺเถ ขนฺเธ อารพฺภ คนฺถา จ คนฺธสมฺปยุตฺตกา จ ขนฺธา อุปฺปชฺชนฺติ. (3)}

\prose{(เอวํ อิตเรปิ ตีณิ ปญฺหา กาตพฺพา.\csromanspacer{} อารมฺมณสทิสํ เยว.\csromanspacer{} อธิปติยาปิ อนนฺตเรปิ อุปนิสฺสเยปิ วิภาโค นตฺถิ.)}

\csromanmatikaanchor{mtk.37}
\csromanmatikaanchor{mtk.39}
\csromantocmarkref{section}{31. คนฺถวิปฺปยุตฺตคนฺถนิยทุก}{mtk.37}
\bookmark[dest={mtk.37},level=1]{31. คนฺถวิปฺปยุตฺตคนฺถนิยทุก}
\csromanheader{2}{32-43. โอฆาทิทุก \textbf{1. ปจฺจยานุโลม 2. สงฺขฺยาวาร}}
\csromanheader{2}{\textbf{สุทฺธ}}
\proseitem{116}{\csromanfolio{297}\textbf{เหตุ}ยา นว, อารมฺมเณ นว, อธิปติยา นว, อนนฺตเร นว, สมนนฺตเร นว, สหชาเต นว, อญฺญมญฺเญ นว, นิสฺสเย นว, อุปนิสฺสเย นว, อาเสวเน นว, กมฺเม ตีณิ, อาหาเร ตีณิ, อินฺทฺริเย ตีณิ, ฌาเน ตีณิ, มคฺเค นว, สมฺปยุตฺเต นว, อตฺถิยา นว, นตฺถิยา นว, วิคเต นว, อวิคเต นว. (อรูปํ เยว ปจฺจยํ, เอเกกสฺส ตีณิ ตีณิ กาตพฺพา.\csromanspacer{}\csromanspacer{}อารมฺมณญฺจ สหชาตญฺจ อุปนิสฺสยญฺจ นวสุปิ ปริวตฺเตตพฺพํ.\csromanspacer{} เอวํ ปญฺหาวาเรปิ สพฺพํ กาตพฺพํ.)}

\nitthitam{คนฺถคนฺถสมฺปยุตฺตทุกํ นิฏฺฐิตํ.}
\csromanheader{1}{31. คนฺถวิปฺปยุตฺตคนฺถนิยทุก 1. ปฏิจฺจวาร}
\csromanheader{2}{\textbf{เหตุ}}
\proseitem{117}{คนฺถวิปฺปยุตฺตคนฺถนิยํ ธมฺมํ ปฏิจฺจ คนฺถวิปฺปยุตฺตคนฺถนิโย ธมฺโม อุปฺปชฺชติ เหตุปจฺจยา.\csromanspacer{} คนฺถวิปฺปยุตฺตํ คนฺถนิยํ เอกํ ขนฺธํ ปฏิจฺจ ตโย ขนฺธา จิตฺตสมุฏฺฐานญฺจ รูปํ ฯเปฯ เทฺว ขนฺเธ ฯเปฯ.\csromanspacer{} ปฏิสนฺธิกฺขเณ ฯเปฯ.\csromanspacer{} เอกํ มหาภูตํ ฯเปฯ. (ยถา จูฬนฺตรทุเก โลกิยทุกํ, เอวํ วิตฺถาเรตพฺพํ, นินฺนานากรณํ.) คนฺถวิปฺปยุตฺตคนฺถนิยทุกํ นิฏฺฐิตํ.\csromanspacer{} คนฺถโคจฺฉกํ นิฏฺฐิตํ.}

\csromanmatikaanchor{mtk.38}
\csromantocmarknopage{section}{6-7. โอฆโยคโคจฺฉก}{mtk.38}
\bookmark[dest={mtk.38},level=1]{6-7. โอฆโยคโคจฺฉก}
\csromantocmarkref{subsection}{32-43. โอฆาทิทุก}{mtk.39}
\bookmark[dest={mtk.39},level=2]{32-43. โอฆาทิทุก}
\csromanheader{1}{6-7. โอฆโยคโคจฺฉก}
\csromanheader{2}{32-43. โอฆาทิทุก}
\proseitem{1}{โอฆํ ธมฺมํ ปฏิจฺจ โอโฆ ธมฺโม อุปฺปชฺชติ เหตุปจฺจยา ฯเปฯ.}

\proseitem{2}{โยคํ ธมฺมํ ปฏิจฺจ โยโค ธมฺโม อุปฺปชฺชติ เหตุปจฺจยา ฯเปฯ. (เทฺวปิ โคจฺฉกา อาสวโคจฺฉกสทิสา, นินฺนานากรณา.)}

\nitthitam{โอฆโยคโคจฺฉกํ นิฏฺฐิตํ.}
\csromanmatikaanchor{mtk.40}
\csromantocmarknopage{section}{8. นีวรณโคจฺฉก}{mtk.40}
\bookmark[dest={mtk.40},level=1]{8. นีวรณโคจฺฉก}
\csromanheader{1}{8. นีวรณโคจฺฉก}
\csromanmatikaanchor{mtk.41}
\csromantocmarkref{subsection}{44. นีวรณทุก}{mtk.41}
\bookmark[dest={mtk.41},level=2]{44. นีวรณทุก}
\csromanheader{2}{44. นีวรณทุก 1. ปฏิจฺจวาร}
\csromanheader{2}{1. ปจฺจยานุโลม 1. วิภงฺควาร}
\csromanheader{2}{\textbf{เหตุ}}
\proseitem{1}{\csromanfolio{298}นีวรณํ ธมฺมํ ปฏิจฺจ นีวรโณ ธมฺโม อุปฺปชฺชติ เหตุปจฺจยา.\csromanspacer{} กามจฺฉนฺทนีวรณํ ปฏิจฺจ ถินมิทฺธนีวรณํ\footnote{ถีนมิทฺธนีวรณํ (สฺยา)}, อุทฺธจฺจนีวรณํ, อวิชฺชานีวรณํ.\csromanspacer{}\csromanspacer{}กามจฺฉนฺทนีวรณํ ปฏิจฺจ อุทฺธจฺจนีวรณํ, อวิชฺชานีวรณํ.\csromanspacer{}\csromanspacer{}พฺยาปาทนีวรณํ ปฏิจฺจ ถินมิทฺธนีวรณํ, อุทฺธจฺจนีวรณํ, อวิชฺชานีวรณํ.\csromanspacer{}\csromanspacer{}พฺยาปาทนีวรณํ ปฏิจฺจ อุทฺธจฺจนีวรณํ, อวิชฺชานีวรณํ.\csromanspacer{}\csromanspacer{}พฺยาปาทนีวรณํ ปฏิจฺจ ถินมิทฺธนีวรณํ, อุทฺธจฺจนีวรณํ, กุกฺกุจฺจนีวรณํ, อวิชฺชานีวรณํ.\csromanspacer{}\csromanspacer{}พฺยาปาทนีวรณํ ปฏิจฺจ อุทฺธจฺจนีวรณํ, กุกฺกุจฺจนีวรณํ, อวิชฺชานีวรณํ.\csromanspacer{}\csromanspacer{}วิจิกิจฺฉานีวรณํ ปฏิจฺจ อุทฺธจฺจนีวรณํ.\csromanspacer{}\csromanspacer{}อุทฺธจฺจนีวรณํ ปฏิจฺจ อวิชฺชานีวรณํ. (1)}

\prose{นีวรณํ ธมฺมํ ปฏิจฺจ โนนีวรโณ ธมฺโม อุปฺปชฺชติ เหตุปจฺจยา.\csromanspacer{} นีวรเณ ปฏิจฺจ สมฺปยุตฺตกา ขนฺธา จิตฺตสมุฏฺฐานญฺจ รูปํ. (2)}

\prose{นีวรณํ ธมฺมํ ปฏิจฺจ นีวรโณ จ โนนีวรโณ จ ธมฺมา อุปฺปชฺชนฺติ เหตุปจฺจยา.\csromanspacer{} กามจฺฉนฺทนีวรณํ ปฏิจฺจ ถินมิทฺธนีวรณํ, อุทฺธจฺจนีวรณํ, อวิชฺชานีวรณํ สมฺปยุตฺตกา จ ขนฺธา จิตฺตสมุฏฺฐานญฺจ รูปํ. (จกฺกํ.) (3)}

\proseitem{2}{โนนีวรณํ ธมฺมํ ปฏิจฺจ โนนีวรโณ ธมฺโม อุปฺปชฺชติ เหตุปจฺจยา.\csromanspacer{} โนนีวรณํ เอกํ ขนฺธํ ปฏิจฺจ ตโย ขนฺธา จิตฺตสมุฏฺฐานญฺจ รูปํ ฯเปฯ เทฺว ขนฺเธ ฯเปฯ.\csromanspacer{} ปฏิสนฺธิกฺขเณ ฯเปฯ. (1)}

\prose{โนนีวรณํ ธมฺมํ ปฏิจฺจ นีวรโณ ธมฺโม อุปฺปชฺชติ เหตุปจฺจยา.\csromanspacer{} โนนีวรเณ ขนฺเธ ปฏิจฺจ นีวรณา. (2)}

\prose{โนนีวรณํ ธมฺมํ ปฏิจฺจ นีวรโณ จ โนนีวรโณ จ ธมฺมา อุปฺปชฺชนฺติ เหตุปจฺจยา.\csromanspacer{} โนนีวรณํ เอกํ ขนฺธํ ปฏิจฺจ ตโย ขนฺธา นีวรณา จ จิตฺตสมุฏฺฐานญฺจ รูปํ ฯเปฯ เทฺว ขนฺเธ ฯเปฯ. (3)}

\csromanheader{2}{44. นีวรณทุก}
\proseitem{3}{\csromanfolio{299}นีวรณญฺจ โนนีวรณญฺจ ธมฺมํ ปฏิจฺจ นีวรโณ ธมฺโม อุปฺปชฺชติ เหตุปจฺจยา.\csromanspacer{} กามจฺฉนฺทนีวรณญฺจ สมฺปยุตฺตเก จ ขนฺเธ ปฏิจฺจ ถินมิทฺธนีวรณํ, อุทฺธจฺจนีวรณํ, อวิชฺชานีวรณํ. (จกฺกํ.) (1)}

\prose{นีวรณญฺจ โนนีวรณญฺจ ธมฺมํ ปฏิจฺจ โนนีวรโณ ธมฺโม อุปฺปชฺชติ เหตุปจฺจยา.\csromanspacer{} โนนีวรณํ เอกํ ขนฺธญฺจ นีวรณญฺจ ปฏิจฺจ ตโย ขนฺธา จิตฺตสมุฏฺฐานญฺจ รูปํ ฯเปฯ เทฺว ขนฺเธ จ ฯเปฯ.\csromanspacer{} นีวรเณ จ มหาภูเต จ ปฏิจฺจ จิตฺตสมุฏฺฐานํ รูปํ. (2)}

\prose{นีวรณญฺจ โนนีวรณญฺจ ธมฺมํ ปฏิจฺจ นีวรโณ จ โนนีวรโณ จ ธมฺมา อุปฺปชฺชนฺติ เหตุปจฺจยา.\csromanspacer{} โนนีวรณํ เอกํ ขนฺธญฺจ กามจฺฉนฺทนีวรณญฺจ ปฏิจฺจ ตโย ขนฺธา ถินมิทฺธนีวรณํ, อุทฺธจฺจนีวรณํ, อวิชฺชานีวรณํ ฯเปฯ เทฺว ขนฺเธ จ ฯเปฯ. (จกฺกํ.\csromanspacer{} สํขิตฺตํ.) (3)}

\csromanheader{2}{1. ปจฺจยานุโลม 2. สงฺขฺยาวาร}
\csromanheader{2}{\textbf{สุทฺธ}}
\proseitem{4}{เหตุยา นว, อารมฺมเณ นว, อธิปติยา นว, อนนฺตเร นว, สมนนฺตเร นว, (สพฺพตฺถ นว,) วิปาเก เอกํ, อาหาเร นว ฯเปฯ อวิคเต นว.}

\vspace{\tipitakaparskip}
\csromancenter{อนุโลมํ.}
\csromanheader{2}{2. ปจฺจยปจฺจนีย 1. วิภงฺควาร}
\csromanheader{2}{\textbf{นเหตุ}}
\proseitem{5}{นีวรณํ ธมฺมํ ปฏิจฺจ นีวรโณ ธมฺโม อุปฺปชฺชติ นเหตุปจฺจยา.\csromanspacer{} วิจิกิจฺฉานีวรณํ ปฏิจฺจ อวิชฺชานีวรณํ.\csromanspacer{} อุทฺธจฺจนีวรณํ ปฏิจฺจ อวิชฺชานีวรณํ. (1)}

\prose{โนนีวรณํ ธมฺมํ ปฏิจฺจ โนนีวรโณ ธมฺโม อุปฺปชฺชติ นเหตุปจฺจยา.\csromanspacer{} อเหตุกํ โนนีวรณํ เอกํ ขนฺธํ ปฏิจฺจ ตโย ขนฺธา จิตฺตสมุฏฺฐานญฺจ รูปํ ฯเปฯ เทฺว ฯเปฯ.\csromanspacer{} อเหตุกปฏิสนฺธิกฺขเณ ฯเปฯ. (ยาว อสญฺญสตฺตา.) (1)}

\prose{\csromanfolio{300}โนนีวรณํ ธมฺมํ ปฏิจฺจ นีวรโณ ธมฺโม อุปฺปชฺชติ นเหตุปจฺจยา.\csromanspacer{} วิจิกิจฺฉาสหคเต อุทฺธจฺจสหคเต ขนฺเธ ปฏิจฺจ อวิชฺชานีวรณํ. (2)}

\proseitem{6}{นีวรณญฺจ โนนีวรณญฺจ ธมฺมํ ปฏิจฺจ นีวรโณ ธมฺโม อุปฺปชฺชติ นเหตุปจฺจยา.\csromanspacer{} วิจิกิจฺฉานีวรณญฺจ สมฺปยุตฺตเก จ ขนฺเธ ปฏิจฺจ อวิชฺชานีวรณํ.\csromanspacer{} อุทฺธจฺจนีวรณญฺจ สมฺปยุตฺตเก จ ขนฺเธ ปฏิจฺจ อวิชฺชานีวรณํ. (1)}

\csromanheader{2}{\textbf{น-อารมฺมณาทิ}}
\proseitem{7}{นีวรณํ ธมฺมํ ปฏิจฺจ โนนีวรโณ ธมฺโม อุปฺปชฺชติ น-อารมฺมณปจฺจยา.\csromanspacer{} นีวรเณ ปฏิจฺจ จิตฺตสมุฏฺฐานํ รูปํ. (1)}

\prose{โนนีวรณํ ธมฺมํ ปฏิจฺจ โนนีวรโณ ธมฺโม อุปฺปชฺชติ น-อารมฺมณปจฺจยา.\csromanspacer{} โนนีวรเณ ขนฺเธ ปฏิจฺจ จิตฺตสมุฏฺฐานํ รูปํ.\csromanspacer{} ปฏิสนฺธิกฺขเณ ฯเปฯ. (ยาว อสญฺญสตฺตา.) (1)}

\prose{นีวรณญฺจ โนนีวรณญฺจ ธมฺมํ ปฏิจฺจ โนนีวรโณ ธมฺโม อุปฺปชฺชติ น-อารมฺมณปจฺจยา.\csromanspacer{} นีวรเณ จ สมฺปยุตฺตเก จ ขนฺเธ ปฏิจฺจ จิตฺตสมุฏฺฐานํ รูปํ. (สํขิตฺตํ.) น-อธิปติปจฺจยา.\csromanspacer{} น-อนนฺตรปจฺจยา.\csromanspacer{} นสมนนฺตรปจฺจยา.\csromanspacer{} น-อญฺญมญฺญปจฺจยา.\csromanspacer{} น-อุปนิสฺสยปจฺจยา.}

\csromanheader{2}{\textbf{นปุเรชาต}}
\proseitem{8}{นีวรณํ ธมฺมํ ปฏิจฺจ นีวรโณ ธมฺโม อุปฺปชฺชติ นปุเรชาตปจฺจยา.\csromanspacer{} อรูเป กามจฺฉนฺทนีวรณํ ปฏิจฺจ ถินมิทฺธนีวรณํ, อุทฺธจฺจนีวรณํ, อวิชฺชานีวรณํ.\csromanspacer{} อรูเป กามจฺฉนฺทนีวรณํ ปฏิจฺจ อุทฺธจฺจนีวรณํ, อวิชฺชานีวรณํ.\csromanspacer{} อรูเป วิจิกิจฺฉานีวรณํ ปฏิจฺจ อุทฺธจฺจนีวรณํ, อวิชฺชานีวรณํ.\csromanspacer{} อรูเป อุทฺธจฺจนีวรณํ ปฏิจฺจ อวิชฺชานีวรณํ. (1)}

\prose{นีวรณํ ธมฺมํ ปฏิจฺจ โนนีวรโณ ธมฺโม อุปฺปชฺชติ นปุเรชาตปจฺจยา.\csromanspacer{} อรูเป นีวรเณ ปฏิจฺจ สมฺปยุตฺตกา ขนฺธา.\csromanspacer{} นีวรเณ ปฏิจฺจ จิตฺตสมุฏฺฐานํ รูปํ. (อวเสสา ปญฺหา สพฺเพ วิตฺถาเรตพฺพา.\csromanspacer{} อรูปํ ปฐมํ กาตพฺพํ.\csromanspacer{} รูปํ ปจฺฉา ยถา ลภติ.)}

\csromanheader{2}{44. นีวรณทุก}
\proseitem{9}{\csromanfolio{301}นีวรณญฺจ โนนีวรณญฺจ ธมฺมํ ปฏิจฺจ นีวรโณ ธมฺโม อุปฺปชฺชติ นปุเรชาตปจฺจยา.\csromanspacer{} อรูเป โนนีวรเณ ขนฺเธ จ กามจฺฉนฺทนีวรณญฺจ ปฏิจฺจ ถินมิทฺธนีวรณํ, อุทฺธจฺจนีวรณํ. (จกฺกํ.) (1)}

\prose{นีวรณญฺจ โนนีวรณญฺจ ธมฺมํ ปฏิจฺจ โนนีวรโณ ธมฺโม อุปฺปชฺชติ นปุเรชาตปจฺจยา.\csromanspacer{} อรูเป โนนีวรณํ เอกํ ขนฺธญฺจ นีวรเณ จ ปฏิจฺจ ตโย ขนฺธา ฯเปฯ เทฺว ขนฺเธ จ ฯเปฯ.\csromanspacer{} นีวรเณ จ สมฺปยุตฺตเก จ ขนฺเธ ปฏิจฺจ จิตฺตสมุฏฺฐานํ รูปํ.\csromanspacer{} นีวรเณ จ มหาภูเต จ ปฏิจฺจ จิตฺตสมุฏฺฐานํ รูปํ. (2)}

\prose{นีวรณญฺจ โนนีวรณญฺจ ธมฺมํ ปฏิจฺจ นีวรโณ จ โนนีวรโณ จ ธมฺมา อุปฺปชฺชนฺติ นปุเรชาตปจฺจยา.\csromanspacer{} อรูเป โนนีวรณํ เอกํ ขนฺธญฺจ กมจฺฉนฺทนีวรณญฺจ ปฏิจฺจ ตโย ขนฺธา ถินมิทฺธนีวรณํ, อุทฺธจฺจนีวรณํ, อวิชฺชานีวรณํ. (จกฺกํ.\csromanspacer{} สํขิตฺตํ.) (3)}

\csromanheader{2}{2. ปจฺจยปจฺจนีย 2. สงฺขฺยาวาร}
\csromanheader{2}{\textbf{สุทฺธ}}
\proseitem{10}{นเหตุยา จตฺตาริ, น-อารมฺมเณ ตีณิ, น-อธิปติยา นว, น-อนนฺตเร ตีณิ, นสมนนฺตเร ตีณิ, น-อญฺญมญฺเญ ตีณิ, น-อุปนิสฺสเย ตีณิ, นปุเรชาเต นว, นปจฺฉาชาเต นว, น-อาเสวเน นว, นกมฺเม ตีณิ, นวิปาเก นว, น-อาหาเร เอกํ, น-อินฺทฺริเย เอกํ, นฌาเน เอกํ, นมคฺเค เอกํ, นสมฺปยุตฺเต ตีณิ, นวิปฺปยุตฺเต นว, โนนตฺถิยา ตีณิ, โนวิคเต ตีณิ.}

\csromanheader{2}{3. ปจฺจยานุโลมปจฺจนีย}
\proseitem{11}{เหตุปจฺจยา น-อารมฺมเณ ตีณิ, น-อธิปติยา นว. (สํขิตฺตํ.)}

\csromanheader{2}{4. ปจฺจยปจฺจนียานุโลม}
\vspace{\tipitakaparskip}
\csromancenter{นเหตุปจฺจยา อารมฺมเณ จตฺตาริ ฯเปฯ มคฺเค ตีณิ ฯเปฯ อวิคเต จตฺตาริ.}
\csromancenter{(สหชาตวาโรปิ เอวํ วิตฺถาเรตพฺโพ.)}
\csromanheader{2}{44. นีวรณทุก 3. ปจฺจยวาร}
\csromanheader{2}{1. ปจฺจยานุโลม 1. วิภงฺควาร}
\csromanheader{2}{\textbf{เหตุ}}
\proseitem{13}{\csromanfolio{302}นีวรณํ ธมฺมํ ปจฺจยา นีวรโณ ธมฺโม อุปฺปชฺชติ เหตุปจฺจยา.\csromanspacer{} ตีณิ.}

\prose{โนนีวรณํ ธมฺมํ ปจฺจยา โนนีวรโณ ธมฺโม อุปฺปชฺชติ เหตุปจฺจยา.\csromanspacer{} โนนีวรณํ เอกํ ขนฺธํ ปจฺจยา ตโย ขนฺธา จิตฺตสมุฏฺฐานญฺจ รูปํ ฯเปฯ. (ยาว อชฺฌตฺติกา มหาภูตา.) วตฺถุํ ปจฺจยา โนนีวรณา ขนฺธา. (1)}

\prose{โนนีวรณํ ธมฺมํ ปจฺจยา นีวรโณ ธมฺโม อุปฺปชฺชติ เหตุปจฺจยา.\csromanspacer{} โนนีวรเณ ขนฺเธ ปจฺจยา นีวรณา.\csromanspacer{} วตฺถุํ ปจฺจยา นีวรณา. (2)}

\prose{โนนีวรณํ ธมฺมํ ปจฺจยา นีวรโณ จ โนนีวรโณ จ ธมฺมา อุปฺปชฺชนฺติ เหตุปจฺจยา.\csromanspacer{} โนนีวรณํ เอกํ ขนฺธํ ปจฺจยา ตโย ขนฺธา นีวรณา จ จิตฺตสมุฏฺฐานญฺจ รูปํ ฯเปฯ เทฺว ขนฺเธ ฯเปฯ.\csromanspacer{} วตฺถุํ ปจฺจยา นีวรณา.\csromanspacer{} มหาภูเต ปจฺจยา จิตฺตสมุฏฺฐานํ รูปํ.\csromanspacer{} วตฺถุํ ปจฺจยา นีวรณา จ สมฺปยุตฺตกา จ ขนฺธา. (3)}

\proseitem{14}{นีวรณญฺจ โนนีวรณญฺจ ธมฺมํ ปจฺจยา นีวรโณ ธมฺโม อุปฺปชฺชติ เหตุปจฺจยา.\csromanspacer{} กามจฺฉนฺทนีวรณญฺจ สมฺปยุตฺตเก จ ขนฺเธ ปจฺจยา ถินมิทฺธนีวรณํ อุทฺธจฺจนีวรณํ อวิชฺชานีวรณํ. (จกฺกํ.) กามจฺฉนฺทนีวรณญฺจ วตฺถุญฺจ ปจฺจยา ถินมิทฺธนีวรณํ อุทฺธจฺจนีวรณํ อวิชฺชานีวรณํ. (จกฺกํ.) (1)}

\prose{นีวรณญฺจ โนนีวรณญฺจ ธมฺมํ ปจฺจยา โนนีวรโณ ธมฺโม อุปฺปชฺชติ เหตุปจฺจยา.\csromanspacer{} โนนีวรณํ เอกํ ขนฺธญฺจ นีวรณญฺจ ปจฺจยา ตโย ขนฺธา จิตฺตสมุฏฺฐานญฺจ รูปํ ฯเปฯ เทฺว ขนฺเธ ฯเปฯ.\csromanspacer{} นีวรณญฺจ วตฺถุญฺจ ปจฺจยา สมฺปยุตฺตกา ขนฺธา.\csromanspacer{} นีวรณญฺจ สมฺปยุตฺตเก จ ขนฺเธ ปจฺจยา จิตฺตสมุฏฺฐานํ รูปํ.\csromanspacer{} นีวรเณ จ มหาภูเต จ ปจฺจยา จิตฺตสมุฏฺฐานํ รูปํ. (2)}

\prose{นีวรณญฺจ โนนีวรณญฺจ ธมฺมํ ปจฺจยา นีวรโณ จ โนนีวรโณ จ ธมฺมา อุปฺปชฺชนฺติ เหตุปจฺจยา.\csromanspacer{} โนนีวรณํ เอกํ ขนฺธญฺจ กามจฺฉนฺทนีวรณญฺจ ปจฺจยา ตโย ขนฺธา ถินมิทฺธนีวรณํ อุทฺธจฺจนีวรณํ อวิชฺชานีวรณํ ฯเปฯ เทฺว}

\vspace{\tipitakaparskip}
\csromancenter{นีวรณทุก ขนฺเธ ฯเปฯ. (จกฺกํ.) กามจฺฉนฺทนีวรณญฺจ วตฺถุญฺจ ปจฺจยา ถินมิทฺธนีวรณํ อุทฺธจฺจนีวรณํ อวิชฺชานีวรณํ สมฺปยุตฺตกา จ ขนฺธา. (จกฺกํ.\csromanspacer{} สํขิตฺตํ.) (3)}
\csromanheader{2}{1. ปจฺจยานุโลม 2. สงฺขฺยาวาร}
\vspace{\tipitakaparskip}
\csromancenter{เหตุยา นว, อารมฺมเณ นว, (สพฺพตฺถ นว,) วิปาเก เอกํ ฯเปฯ อวิคเต นว.}
\csromancenter{อนุโลมํ.}
\csromanheader{2}{2. ปจฺจยปจฺจนีย 1. วิภงฺควาร}
\csromanheader{2}{\textbf{นเหตุ}}
\proseitem{16}{\csromanfolio{303}นีวรณํ ธมฺมํ ปจฺจยา นีวรโณ ธมฺโม อุปฺปชฺชติ นเหตุปจฺจยา.\csromanspacer{} วิจิกิจฺฉานีวรณํ ปจฺจยา อวิชฺชานีวรณํ.\csromanspacer{} อุทฺธจฺจนีวรณํ ปจฺจยา อวิชฺชานีวรณํ. (1)}

\prose{โนนีวรณํ ธมฺมํ ปจฺจยา โนนีวรโณ ธมฺโม อุปฺปชฺชติ นเหตุปจฺจยา.\csromanspacer{} อเหตุกํ โนนีวรณํ เอกํ ขนฺธํ ปจฺจยา ตโย ขนฺธา จิตฺตสมุฏฺฐานญฺจ รูปํ ฯเปฯ เทฺว ขนฺเธ ฯเปฯ. (ยาว อสญฺญสตฺตา.) จกฺขายตนํ ปจฺจยา จกฺขุวิญฺญาณํ ฯเปฯ.\csromanspacer{} กายายตนํ ปจฺจยา กายวิญฺญาณํ.\csromanspacer{} วตฺถุํ ปจฺจยา อเหตุกา โนนีวรณา ขนฺธา. (1)}

\prose{โนนีวรณํ ธมฺมํ ปจฺจยา นีวรโณ ธมฺโม อุปฺปชฺชติ นเหตุปจฺจยา.\csromanspacer{} วิจิกิจฺฉาสหคเต อุทฺธจฺจสหคเต ขนฺเธ ปจฺจยา อวิชฺชานีวรณํ.\csromanspacer{} วตฺถุํ ปจฺจยา อวิชฺชานีวรณํ. (2)}

\proseitem{17}{นีวรณญฺจ โนนีวรณญฺจ ธมฺมํ ปจฺจยา นีวรโณ ธมฺโม อุปฺปชฺชติ นเหตุปจฺจยา.\csromanspacer{} วิจิกิจฺฉานีวรณญฺจ สมฺปยุตฺตเก จ ขนฺเธ ปจฺจยา อวิชฺชานีวรณํ.\csromanspacer{} อุทฺธจฺจนีวรณญฺจ สมฺปยุตฺตเก จ ขนฺเธ ปจฺจยา อวิชฺชานีวรณํ.\csromanspacer{} วิจิกิจฺฉานีวรณญฺจ วตฺถุญฺจ ปจฺจยา อวิชฺชานีวรณํ.\csromanspacer{} อุทฺธจฺจนีวรณญฺจ วตฺถุญฺจ ปจฺจยา อวิชฺชานีวรณํ. (สํขิตฺตํ.) (1)}

\csromanheader{2}{2. ปจฺจยปจฺจนีย 2. สงฺขฺยาวาร}
\csromanheader{2}{\textbf{สุทฺธ}}
\proseitem{18}{\csromanfolio{304}นเหตุยา จตฺตาริ, น-อารมฺมเณ ตีณิ, น-อธิปติยา นว, น-อนนฺตเร ตีณิ, นสมนนฺตเร ตีณิ, น-อญฺญมญฺเญ ตีณิ, น-อุปนิสฺสเย ตีณิ, นปุเรชาเต นว, นปจฺฉาชาเต นว, น-อาเสวเน นว, นกมฺเม ตีณิ, นวิปาเก นว, น-อาหาเร เอกํ, น-อินฺทฺริเย เอกํ, นฌาเน เอกํ, นมคฺเค เอกํ, นสมฺปยุตฺเต ตีณิ, นวิปฺปยุตฺเต นว, โนนตฺถิยา ตีณิ, โนวิคเต ตีณิ.}

\csromanheader{2}{3. ปจฺจยานุโลมปจฺจนีย}
\proseitem{19}{เหตุปจฺจยา น-อารมฺมเณ ตีณิ, น-อธิปติยา นว. (สํขิตฺตํ.)}

\csromanheader{2}{4. ปจฺจยปจฺจนียานุโลม}
\proseitem{20}{นเหตุปจฺจยา อารมฺมเณ จตฺตาริ, อนนฺตเร จตฺตาริ, สมนนฺตเร จตฺตาริ ฯเปฯ มคฺเค ตีณิ, อวิคเต จตฺตาริ.}

\csromanheader{2}{44. นีวรณทุก 5. สํสฏฺฐวาร}
\csromanheader{2}{1-4. ปจฺจยจตุกฺก}
\proseitem{21}{นีวรณํ ธมฺมํ สํสฏฺโฐ นีวรโณ ธมฺโม อุปฺปชฺชติ เหตุปจฺจยา.\csromanspacer{} กามจฺฉนฺทนีวรณํ สํสฏฺฐํ ถินมิทฺธนีวรณํ อุทฺธจฺจนีวรณํ อวิชฺชานีวรณํ. (จกฺกํ.\csromanspacer{} สพฺพํ นีวรณํ วิตฺถาเรตพฺพํ.)}

\proseitem{22}{เหตุยา นว, อารมฺมเณ นว, (สพฺพตฺถ นว,) วิปาเก เอกํ ฯเปฯ อวิคเต นว.}

\vspace{\tipitakaparskip}
\csromancenter{อนุโลมํ.}
\csromanheader{2}{44. นีวรณทุก}
\proseitem{23}{\csromanfolio{305}นีวรณํ ธมฺมํ สํสฏฺโฐ นีวรโณ ธมฺโม อุปฺปชฺชติ นเหตุปจฺจยา.\csromanspacer{} วิจิกิจฺฉานีวรณํ สํสฏฺฐํ อวิชฺชานีวรณํ.\csromanspacer{} อุทฺธจฺจนีวรณํ สํสฏฺฐํ อวิชฺชานีวรณํ. (สํขิตฺตํ.)}

\proseitem{24}{นเหตุยา จตฺตาริ, น-อธิปติยา นว, นปุเรชาเต นว, นปจฺฉาชาเต นว, น-อาเสวเน นว, นกมฺเม ตีณิ, นวิปาเก นว, นฌาเน เอกํ, นมคฺเค เอกํ, นวิปฺปยุตฺเต นว.}

\vspace{\tipitakaparskip}
\csromancenter{ปจฺจนียํ.}
\csromancenter{(เอวํ อิตเร เทฺว คณนาปิ สมฺปยุตฺตวาโรปิ กาตพฺโพ.)}
\csromanheader{2}{44. นีวรณทุก 7. ปญฺหาวาร}
\csromanheader{2}{1. ปจฺจยานุโลม 1. วิภงฺควาร}
\csromanheader{2}{\textbf{เหตุ}}
\proseitem{25}{นีวรโณ ธมฺโม นีวรณสฺส ธมฺมสฺส เหตุปจฺจเยน ปจฺจโย.\csromanspacer{} นีวรณา เหตู สมฺปยุตฺตกานํ นีวรณานํ เหตุปจฺจเยน ปจฺจโย. (1)}

\prose{นีวรโณ ธมฺโม โนนีวรณสฺส ธมฺมสฺส เหตุปจฺจเยน ปจฺจโย.\csromanspacer{} นีวรณา เหตู สมฺปยุตฺตกานํ ขนฺธานํ จิตฺตสมุฏฺฐานานญฺจ รูปานํ เหตุปจฺจเยน ปจฺจโย. (2)}

\prose{นีวรโณ ธมฺโม นีวรณสฺส จ โนนีวรณสฺส จ ธมฺมสฺส เหตุปจฺจเยน ปจฺจโย.\csromanspacer{} นีวรณา เหตู สมฺปยุตฺตกานํ ขนฺธานํ นีวรณานญฺจ จิตฺตสมุฏฺฐานานญฺจ รูปานํ เหตุปจฺจเยน ปจฺจโย. (3)}

\prose{โนนีวรโณ ธมฺโม โนนีวรณสฺส ธมฺมสฺส เหตุปจฺจเยน ปจฺจโย.\csromanspacer{} โนนีวรณา เหตู สมฺปยุตฺตกานํ ขนฺธานํ จิตฺตสมุฏฺฐานานญฺจ รูปานํ เหตุปจฺจเยน ปจฺจโย.\csromanspacer{} ปฏิสนฺธิกฺขเณ ฯเปฯ. (1)}

\csromanheader{2}{\textbf{อารมฺมณ}}
\proseitem{26}{\csromanfolio{306}นีวรโณ ธมฺโม นีวรณสฺส ธมฺมสฺส อารมฺมณปจฺจเยน ปจฺจโย.\csromanspacer{} นีวรเณ อารพฺภ นีวรณา อุปฺปชฺชนฺติ. (มูลํ ปุจฺฉิตพฺพํ.) นีวรเณ อารพฺภ โนนีวรณา ขนฺธา อุปฺปชฺชนฺติ. (มูลํ กาตพฺพํ.) นีวรเณ อารพฺภ นีวรณา จ สมฺปยุตฺตกา จ ขนฺธา อุปฺปชฺชนฺติ. (3)}

\proseitem{27}{โนนีวรโณ ธมฺโม โนนีวรณสฺส ธมฺมสฺส อารมฺมณปจฺจเยน ปจฺจโย.\csromanspacer{} ทานํ ฯเปฯ สีลํ ฯเปฯ อุโปสถกมฺมํ กตฺวา ตํ ปจฺจเวกฺขติ, ปุพฺเพ สุจิณฺณานิ ฯเปฯ ฌานา วุฏฺฐหิตฺวา ฌานํ ฯเปฯ.\csromanspacer{} อริยา มคฺคา วุฏฺฐหิตฺวา มคฺคํ ฯเปฯ.\csromanspacer{} ผลํ ฯเปฯ.\csromanspacer{} นิพฺพานํ ฯเปฯ.\csromanspacer{} นิพฺพานํ โคตฺรภุสฺส, โวทานสฺส, มคฺคสฺส, ผลสฺส, อาวชฺชนาย อารมฺมณปจฺจเยน ปจฺจโย.\csromanspacer{} อริยา โนนีวรเณ ปหีเน กิเลเส ปจฺจเวกฺขนฺติ.\csromanspacer{} วิกฺขมฺภิเต กิเลเส ฯเปฯ.\csromanspacer{} ปุพฺเพ สมุทาจิณฺเณ ฯเปฯ.\csromanspacer{} จกฺขุํ ฯเปฯ วตฺถุํ โนนีวรเณ ขนฺเธ อนิจฺจโต ฯเปฯ โทมนสฺสํ อุปฺปชฺชติ.\csromanspacer{} ทิพฺเพน จกฺขุนา รูปํ ปสฺสติ.\csromanspacer{} ทิพฺพาย โสตธาตุยา สทฺทํ สุณาติ.\csromanspacer{} เจโตปริยญาเณน โนนีวรณจิตฺตสมงฺคิสฺส จิตฺตํ ชานาติ.\csromanspacer{} อากาสานญฺจายตนํ วิญฺญาณญฺจายตนสฺส ฯเปฯ.\csromanspacer{} อากิญฺจญฺญายตนํ เนวสญฺญานาสญฺญายตนสฺส ฯเปฯ.\csromanspacer{} รูปายตนํ จกฺขุวิญฺญาณสฺส ฯเปฯ.\csromanspacer{} โผฏฺฐพฺพายตนํ กายวิญฺญาณสฺส ฯเปฯ.\csromanspacer{} โนนีวรณา ขนฺธา อิทฺธิวิธญาณสฺส, เจโตปริยญาณสฺส, ปุพฺเพนิวาสานุสฺสติญาณสฺส, ยถากมฺมูปคญาณสฺส, อนาคตํสญาณสฺส, อาวชฺชนาย อารมฺมณปจฺจเยน ปจฺจโย. (1)}

\prose{โนนีวรโณ ธมฺโม นีวรณสฺส ธมฺมสฺส อารมฺมณปจฺจเยน ปจฺจโย.\csromanspacer{} ทานํ ฯเปฯ สีลํ ฯเปฯ อุโปสถกมฺมํ ฯเปฯ.\csromanspacer{} ปุพฺเพ สุจิณฺณานิ ฯเปฯ.\csromanspacer{} ฌานา ฯเปฯ.\csromanspacer{} จกฺขุํ ฯเปฯ วตฺถุํ โนนีวรเณ ขนฺเธ อสฺสาเทติ อภินนฺทติ, ตํ อารพฺภ ราโค อุปฺปชฺชติ.\csromanspacer{} ทิฏฺฐิ ฯเปฯ.\csromanspacer{} วิจิกิจฺฉา ฯเปฯ.\csromanspacer{} อุทฺธจฺจํ ฯเปฯ.\csromanspacer{} โทมนสฺสํ อุปฺปชฺชติ. (2)}

\prose{โนนีวรโณ ธมฺโม นีวรณสฺส จ โนนีวรณสฺส จ ธมฺมสฺส อารมฺมณปจฺจเยน ปจฺจโย.\csromanspacer{} ทานํ ฯเปฯ สีลํ ฯเปฯ อุโปสถกมฺมํ ฯเปฯ.\csromanspacer{} ปุพฺเพ สุจิณฺณานิ ฯเปฯ.\csromanspacer{} ฌานา ฯเปฯ.\csromanspacer{} จกฺขุํ ฯเปฯ วตฺถุํ โนนีวรเณ ขนฺเธ อสฺสาเทติ อภินนฺทติ ตํ อารพฺภ นีวรณา จ สมฺปยุตฺตกา จ ขนฺธา อุปฺปชฺชนฺติ. (3)}

\csromanheader{2}{44. นีวรณทุก}
\prose{\csromanfolio{307}นีวรโณ จ โนนีวรโณ จ ธมฺมา นีวรณสฺส ธมฺมสฺส อารมฺมณปจฺจเยน ปจฺจโย. (อารพฺภ กาตพฺพา.)}

\csromanheader{2}{\textbf{อธิปติ}}
\proseitem{28}{นีวรโณ ธมฺโม นีวรณสฺส ธมฺมสฺส อธิปติปจฺจเยน ปจฺจโย.\csromanspacer{} \textbf{อารมฺมณาธิปติ}\csromandash{}นีวรเณ ครุํ กตฺวา นีวรณา อุปฺปชฺชนฺติ.\csromanspacer{} ตีณิ. (อารมฺมณสทิสํ.) (3)}

\proseitem{29}{โนนีวรโณ ธมฺโม โนนีวรณสฺส ธมฺมสฺส อธิปติปจฺจเยน ปจฺจโย.\csromanspacer{} \textbf{อารมฺมณาธิปติ}, สหชาตาธิปติ.\csromanspacer{}\csromanspacer{}\textbf{อารมฺมณาธิปติ}\csromandash{}ทานํ ฯเปฯ สีลํ ฯเปฯ อุโปสถกมฺมํ กตฺวา ตํ ครุํ กตฺวา ปจฺจเวกฺขติ, ปุพฺเพ สุจิณฺณานิ ฯเปฯ.\csromanspacer{} ฌานา ฯเปฯ.\csromanspacer{} อริยา มคฺคา วุฏฺฐหิตฺวา มคฺคํ ฯเปฯ.\csromanspacer{} ผลํ ฯเปฯ.\csromanspacer{} นิพฺพานํ ฯเปฯ.\csromanspacer{} นิพฺพานํ โคตฺรภุสฺส, โวทานสฺส, มคฺคสฺส, ผลสฺส อธิปติปจฺจเยน ปจฺจโย.\csromanspacer{} จกฺขุํ ฯเปฯ วตฺถุํ โนนีวรเณ ขนฺเธ ครุํ กตฺวา อสฺสาเทติ อภินนฺทติ, ตํ ครุํ กตฺวา ราโค อุปฺปชฺชติ, ทิฏฺฐิ อุปฺปชฺชติ.\csromanspacer{} \textbf{สหชาตาธิปติ}\csromandash{}โนนีวรณาธิปติ สมฺปยุตฺตกานํ ขนฺธานํ จิตฺตสมุฏฺฐานานญฺจ รูปานํ อธิปติปจฺจเยน ปจฺจโย. (1)}

\prose{โนนีวรโณ ธมฺโม นีวรณสฺส ธมฺมสฺส อธิปติปจฺจเยน ปจฺจโย.\csromanspacer{} \textbf{อารมฺมณาธิปติ}, สหชาตาธิปติ.\csromanspacer{}\csromanspacer{}\textbf{อารมฺมณาธิปติ}\csromandash{}ทานํ ฯเปฯ สีลํ ฯเปฯ อุโปสถกมฺมํ ฯเปฯ.\csromanspacer{} ปุพฺเพ ฯเปฯ.\csromanspacer{} ฌานา ฯเปฯ.\csromanspacer{} จกฺขุํ ฯเปฯ วตฺถุํ โนนีวรเณ ขนฺเธ ครุํ กตฺวา อสฺสาเทติ อภินนฺทติ, ตํ ครุํ กตฺวา ราโค อุปฺปชฺชติ, ทิฏฺฐิ อุปฺปชฺชติ.\csromanspacer{} \textbf{สหชาตาธิปติ}\csromandash{} โนนีวรณาธิปติ สมฺปยุตฺตกานํ นีวรณานํ อธิปติปจฺจเยน ปจฺจโย. (2)}

\prose{โนนีวรโณ ธมฺโม นีวรณสฺส จ โนนีวรณสฺส จ ธมฺมสฺส อธิปติปจฺจเยน ปจฺจโย.\csromanspacer{} \textbf{อารมฺมณาธิปติ}, สหชาตาธิปติ.\csromanspacer{}\csromanspacer{}\textbf{อารมฺมณาธิปติ}\csromandash{}ทานํ ฯเปฯ สีลํ ฯเปฯ อุโปสถกมฺมํ ฯเปฯ.\csromanspacer{} ปุพฺเพ ฯเปฯ.\csromanspacer{} ฌานา ฯเปฯ.\csromanspacer{} จกฺขุํ ฯเปฯ วตฺถุํ โนนีวรโณ ขนฺเธ ครุํ กตฺวา อสฺสาเทติ อภินนฺทติ, ตํ ครุํ กตฺวา นีวรณา จ สมฺปยุตฺตกา จ ขนฺธา อุปฺปชฺชนฺติ.\csromanspacer{} \textbf{สหชาตาธิปติ}\csromandash{}โนนีวรณาธิปติ สมฺปยุตฺตกานํ ขนฺธานํ นีวรณานญฺจ จิตฺตสมุฏฺฐานานญฺจ รูปานํ อธิปติปจฺจเยน ปจฺจโย. (3)}

\prose{\csromanfolio{308}นีวรโณ จ โนนีวรโณ จ ธมฺมา นีวรณสฺส ธมฺมสฺส อธิปติปจฺจเยน ปจฺจโย.\csromanspacer{} อารมฺมณาธิปติ.\csromanspacer{} ตีณิ. (อารมฺมณาธิปติเยว.)}

\csromanheader{2}{\textbf{อนนฺตร}}
\proseitem{30}{นีวรโณ ธมฺโม นีวรณสฺส ธมฺมสฺส อนนฺตรปจฺจเยน ปจฺจโย.\csromanspacer{} ปุริมา ปุริมา นีวรณา ปจฺฉิมานํ ปจฺฉิมานํ นีวรณานํ อนนฺตรปจฺจเยน ปจฺจโย. (มูลํ ปุจฺฉิตพฺพํ.) ปุริมา ปุริมา นีวรณา ปจฺฉิมานํ ปจฺฉิมานํ โนนีวรณานํ ขนฺธานํ อนนฺตรปจฺจเยน ปจฺจโย.\csromanspacer{} นีวรณา วุฏฺฐานสฺส อนนฺตรปจฺจเยน ปจฺจโย. (มูลํ ปุจฺฉิตพฺพํ.) ปุริมา ปุริมา นีวรณา ปจฺฉิมานํ ปจฺฉิมานํ นีวรณานํ สมฺปยุตฺตกานญฺจ ขนฺธานํ อนนฺตรปจฺจเยน ปจฺจโย. (3)}

\proseitem{31}{โนนีวรโณ ธมฺโม โนนีวรณสฺส ธมฺมสฺส อนนฺตรปจฺจเยน ปจฺจโย.\csromanspacer{} ปุริมา ปุริมา โนนีวรณา ขนฺธา ปจฺฉิมานํ ปจฺฉิมานํ โนนีวรณานํ ขนฺธานํ อนนฺตรปจฺจเยน ปจฺจโย ฯเปฯ.\csromanspacer{} นิโรธา วุฏฺฐหนฺตสฺส เนวสญฺญานาสญฺญายตนํ ผลสมาปตฺติยา อนนฺตรปจฺจเยน ปจฺจโย. (มูลํ ปุจฺฉิตพฺพํ.) ปุริมา ปุริมา โนนีวรณา ขนฺธา ปจฺฉิมานํ ปจฺฉิมานํ นีวรณานํ อนนฺตรปจฺจเยน ปจฺจโย.\csromanspacer{} อาวชฺชนา นีวรณานํ อนนฺตรปจฺจเยน ปจฺจโย. (มูลํ ปุจฺฉิตพฺพํ.) ปุริม ปุริมา โนนีวรณา ขนฺธา ปจฺฉิมานํ ปจฺฉิมานํ นีวรณานํ สมฺปยุตฺตกานญฺจ ขนฺธานํ อนนฺตรปจฺจเยน ปจฺจโย.\csromanspacer{} อาวชฺชนา นีวรณานํ สมฺปยุตฺตกานญฺจ ขนฺธานํ อนนฺตรปจฺจเยน ปจฺจโย. (3)}

\proseitem{32}{นีวรโณ จ โนนีวรโณ จ ธมฺมา นีวรณสฺส ธมฺมสฺส อนนฺตรปจฺจเยน ปจฺจโย.\csromanspacer{} ปุริมา ปุริมา นีวรณา จ สมฺปยุตฺตกา จ ขนฺธา ปจฺฉิมานํ ปจฺฉิมานํ นีวรณานํ อนนฺตรปจฺจเยน ปจฺจโย. (มูลํ ปุจฺฉิตพฺพํ.) ปุริมา ปุริมา นีวรณา จ สมฺปยุตฺตกา จ ขนฺธา ปจฺฉิมานํ ปจฺฉิมานํ โนนีวรณานํ ขนฺธานํ อนนฺตรปจฺจเยน ปจฺจโย.\csromanspacer{} นีวรณา จ สมฺปยุตฺตกา จ ขนฺธา วุฏฺฐานสฺส อนนฺตรปจฺจเยน ปจฺจโย. (มูลํ ปุจฺฉิตพฺพํ.) ปุริมา ปุริมา นีวรณา จ สมฺปยุตฺตกา จ ขนฺธา ปจฺฉิมานํ ปจฺฉิมานํ นีวรณานญฺจ สมฺปยุตฺตกานญฺจ ขนฺธานํ อนนฺตรปจฺจเยน ปจฺจโย. (3)}

\csromanheader{2}{44. นีวรณทุก \textbf{สมนนฺตราทิ}}
\proseitem{33}{\csromanfolio{309}นีวรโณ ธมฺโม นีวรณสฺส ธมฺมสฺส สมนนฺตรปจฺจเยน ปจฺจโย.\csromanspacer{} สหชาตปจฺจเยน ปจฺจโย.\csromanspacer{} อญฺญมญฺญปจฺจเยน ปจฺจโย.\csromanspacer{} นิสฺสยปจฺจเยน ปจฺจโย.}

\csromanheader{2}{\textbf{อุปนิสฺสย}}
\proseitem{34}{นีวรโณ ธมฺโม นีวรณสฺส ธมฺมสฺส อุปนิสฺสยปจฺจเยน ปจฺจโย.\csromanspacer{} อารมฺมณูปนิสฺสโย, อนนฺตรูปนิสฺสโย, ปกตูปนิสฺสโย ฯเปฯ.\csromanspacer{} ปกตูปนิสฺสโย\csromandash{}นีวรณานิ นีวรณานํ อุปนิสฺสยปจฺจเยน ปจฺจโย.\csromanspacer{} ตีณิ.}

\proseitem{35}{โนนีวรโณ ธมฺโม โนนีวรณสฺส ธมฺมสฺส อุปนิสฺสยปจฺจเยน ปจฺจโย.\csromanspacer{}\csromanspacer{}อารมฺมณูปนิสฺสโย, อนนฺตรูปนิสฺสโย, ปกตูปนิสฺสโย ฯเปฯ.\csromanspacer{} ปกตูปนิสฺสโย\csromandash{}สทฺธํ อุปนิสฺสาย ทานํ เทติ ฯเปฯ สีลํ ฯเปฯ อุโปสถกมฺมํ ฯเปฯ.\csromanspacer{} ฌานํ อุปฺปาเทติ, วิปสฺสนํ ฯเปฯ มคฺคํ ฯเปฯ อภิญฺญํ ฯเปฯ สมาปตฺติํ อุปฺปาเทติ, มานํ ชปฺเปติ, ทิฏฺฐิํ คณฺหาติ.\csromanspacer{} สีลํ ฯเปฯ.\csromanspacer{} ปญฺญํ ฯเปฯ มานํ.\csromanspacer{} ทิฏฺฐิํ.\csromanspacer{} ปตฺถนํ.\csromanspacer{} กายิกํ สุขํ.\csromanspacer{} กายิกํ ทุกฺขํ.\csromanspacer{} อุตุํ.\csromanspacer{} โภชนํ.\csromanspacer{} เสนาสนํ อุปนิสฺสาย ทานํ เทติ ฯเปฯ สํฆํ ภินฺทติ.\csromanspacer{} สทฺธา ฯเปฯ.\csromanspacer{} เสนาสนํ สทฺธาย ฯเปฯ ปญฺญาย.\csromanspacer{} ราคสฺส ฯเปฯ ปตฺถนาย.\csromanspacer{} กายิกสฺส สุขสฺส, กายิกสฺส ทุกฺขสฺส, มคฺคสฺส, ผลสมาปตฺติยา อุปนิสฺสยปจฺจเยน ปจฺจโย. (1)}

\prose{โนนีวรโณ ธมฺโม นีวรณสฺส ธมฺมสฺส อุปนิสฺสยปจฺจเยน ปจฺจโย. (ตีณิ อุปนิสฺสยา.) ปกตูปนิสฺสโย\csromandash{}สทฺธํ อุปนิสฺสาย มานํ ชปฺเปติ, ทิฏฺฐิํ คณฺหาติ.\csromanspacer{} สีลํ ฯเปฯ.\csromanspacer{} เสนาสนํ อุปนิสฺสาย ปาณํ หนติ ฯเปฯ สํฆํ ภินฺทติ.\csromanspacer{} สทฺธา ฯเปฯ.\csromanspacer{} เสนาสนํ ราคสฺส ฯเปฯ ปตฺถนาย อุปนิสฺสยปจฺจเยน ปจฺจโย. (2)}

\prose{โนนีวรโณ ธมฺโม นีวรณสฺส จ โนนีวรณสฺส จ ธมฺมสฺส อุปนิสฺสยปจฺจเยน ปจฺจโย.\csromanspacer{} อารมฺมณูปนิสฺสโย, อนนฺตรูปนิสฺสโย, ปกตูปนิสฺสโย ฯเปฯ.\csromanspacer{} ปกตูปนิสฺสโย\csromandash{}สทฺธํ อุปนิสฺสาย มานํ ชปฺเปติ, ทิฏฺฐิํ คณฺหาติ.\csromanspacer{} สีลํ ฯเปฯ.\csromanspacer{} เสนาสนํ อุปนิสฺสาย ปาณํ หนติ ฯเปฯ สํฆํ ภินฺทติ.\csromanspacer{} สทฺธา ฯเปฯ.\csromanspacer{} เสนาสนํ นีวรณานํ สมฺปยุตฺตกานญฺจ ขนฺธานํ อุปนิสฺสยปจฺจเยน ปจฺจโย. (3)}

\prose{\csromanfolio{310}นีวรโณ จ โนนีวรโณ จ ธมฺมา นีวรณสฺส ธมฺมสฺส อุปนิสฺสยปจฺจเยน ปจฺจโย.\csromanspacer{} อารมฺมณูปนิสฺสโย, อนนฺตรูปนิสฺสโย, ปกตูปนิสฺสโย ฯเปฯ.\csromanspacer{} ปกตูปนิสฺสโย\csromandash{}นีวรณา จ สมฺปยุตฺตกา จ ขนฺธา นีวรณานํ อุปนิสฺสยปจฺจเยน ปจฺจโย.\csromanspacer{} ตีณิ.}

\csromanheader{2}{\textbf{ปุเรชาต}}
\proseitem{36}{โนนีวรโณ ธมฺโม โนนีวรณสฺส ธมฺมสฺส ปุเรชาตปจฺจเยน ปจฺจโย.\csromanspacer{} \textbf{อารมฺมณปุเรชาตํ}, วตฺถุปุเรชาตํ.\csromanspacer{}\csromanspacer{}\textbf{อารมฺมณปุเรชาตํ}\csromandash{} จกฺขุํ ฯเปฯ วตฺถุํ อนิจฺจโต ฯเปฯ โทมนสฺสํ อุปฺปชฺชติ ฯเปฯ.\csromanspacer{} ตีณิ. (\textbf{ปุเรชาต}ํ อารมฺมณสทิสํ, กุสลากุสลสฺส วิภชิตพฺพํ.)}

\csromanheader{2}{\textbf{ปจฺฉาชาต อาเสวน}}
\proseitem{37}{นีวรโณ ธมฺโม โนนีวรณสฺส ธมฺมสฺส ปจฺฉาชาตปจฺจเยน ปจฺจโย.\csromanspacer{} ตีณิ.\csromanspacer{} อาเสวนปจฺจเยน ปจฺจโย.\csromanspacer{} นว.}

\csromanheader{2}{\textbf{กมฺม}}
\proseitem{38}{โนนีวรโณ ธมฺโม โนนีวรณสฺส ธมฺมสฺส กมฺมปจฺจเยน ปจฺจโย.\csromanspacer{} \textbf{สหชาตา}, นานากฺขณิกา.\csromanspacer{}\csromanspacer{}\textbf{สหชาตา} โนนีวรณา เจตนา สมฺปยุตฺตกานํ ขนฺธานํ จิตฺตสมุฏฺฐานานญฺจ รูปานํ กมฺมปจฺจเยน ปจฺจโย.\csromanspacer{} ปฏิสนฺธิกฺขเณ ฯเปฯ.\csromanspacer{} \textbf{นานากฺขณิกา} โนนีวรณา เจตนา วิปากานํ ขนฺธานํ กฏตฺตา จ รูปานํ กมฺมปจฺจเยน ปจฺจโย. (มูลํ ปุจฺฉิตพฺพํ.) โนนีวรณา เจตนา สมฺปยุตฺตกานํ นีวรณานํ กมฺมปจฺจเยน ปจฺจโย. (มูลํ ปุจฺฉิตพฺพํ.) โนนีวรณา เจตนา สมฺปยุตฺตกานํ ขนฺธานํ นีวรณานญฺจ จิตฺตสมุฏฺฐานานญฺจ รูปานํ กมฺมปจฺจเยน ปจฺจโย. (3)}

\csromanheader{2}{\textbf{วิปากาทิ}}
\proseitem{39}{โนนีวรโณ ธมฺโม โนนีวรณสฺส ธมฺมสฺส วิปากปจฺจเยน ปจฺจโย.\csromanspacer{} เอกํ.\csromanspacer{} อาหารปจฺจเยน ปจฺจโย.\csromanspacer{} อินฺทฺริยปจฺจเยน ปจฺจโย.\csromanspacer{} ฌานปจฺจเยน ปจฺจโย.\csromanspacer{} มคฺคปจฺจเยน ปจฺจโย.\csromanspacer{} สมฺปยุตฺตปจฺจเยน ปจฺจโย.}

\csromanheader{2}{44. นีวรณทุก \textbf{วิปฺปยุตฺต}}
\proseitem{40}{\csromanfolio{311}นีวรโณ ธมฺโม โนนีวรณสฺส ธมฺมสฺส วิปฺปยุตฺตปจฺจเยน ปจฺจโย.\csromanspacer{} สหชาตํ, ปจฺฉาชาตํ. (เอวํ อวเสสา จตฺตาริ ปญฺหา กาตพฺพา.)}

\csromanheader{2}{\textbf{อตฺถิ}}
\proseitem{41}{นีวรโณ ธมฺโม นีวรณสฺส ธมฺมสฺส อตฺถิปจฺจเยน ปจฺจโย.\csromanspacer{} กามจฺฉนฺทนีวรณํ ถินมิทฺธนีวรณสฺส, อุทฺธจฺจนีวรณสฺส, อวิชฺชานีวรณสฺส อตฺถิปจฺจเยน ปจฺจโย. (จกฺกํ.) (1)}

\prose{นีวรโณ ธมฺโม โนนีวรณสฺส ธมฺมสฺส อตฺถิปจฺจเยน ปจฺจโย.\csromanspacer{} สหชาตํ, ปจฺฉาชาตํ. (เอวํ นีวรณมูเล ตีณิ.) (3)}

\proseitem{42}{โนนีวรโณ ธมฺโม โนนีวรณสฺส ธมฺมสฺส อตฺถิปจฺจเยน ปจฺจโย.\csromanspacer{} สหชาตํ, ปุเรชาตํ, ปจฺฉาชาตํ, อาหารํ, อินฺทฺริยํ. (สํขิตฺตํ.) (1)}

\prose{โนนีวรโณ ธมฺโม นีวรณสฺส ธมฺมสฺส อตฺถิปจฺจเยน ปจฺจโย.\csromanspacer{} สหชาตํ ปุเรชาตํ. (สํขิตฺตํ.) (2)}

\prose{โนนีวรโณ ธมฺโม นีวรณสฺส จ โนนีวรณสฺส จ ธมฺมสฺส อตฺถิปจฺจเยน ปจฺจโย.\csromanspacer{} สหชาตํ, ปุเรชาตํ. (สํขิตฺตํ.) (3)}

\proseitem{43}{นีวรโณ จ โนนีวรโณ จ ธมฺมา นีวรณสฺส ธมฺมสฺส อตฺถิปจฺจเยน ปจฺจโย.\csromanspacer{} สหชาตํ, ปุเรชาตํ.\csromanspacer{}\csromanspacer{}สหชาตํ กามจฺฉนฺทนีวรณญฺจ สมฺปยุตฺตกา จ ขนฺธา ถินมิทฺธนีวรณสฺส, อุทฺธจฺจนีวรณสฺส, อวิชฺชานีวรณสฺส อตฺถิปจฺจเยน ปจฺจโย.\csromanspacer{} กามจฺฉนฺทนีวรณญฺจ วตฺถุ จ ถินมิทฺธนีวรณสฺส, อุทฺธจฺจนีวรณสฺส, อวิชฺชานีวรณสฺส อตฺถิปจฺจเยน ปจฺจโย. (1)}

\prose{นีวรโณ จ โนนีวรโณ จ ธมฺมา โนนีวรณสฺส ธมฺมสฺส อตฺถิปจฺจเยน ปจฺจโย.\csromanspacer{} สหชาตํ, ปุเรชาตํ, ปจฺฉาชาตํ, อาหารํ, อินฺทฺริยํ.\csromanspacer{}\csromanspacer{}\textbf{สหชาโต} โนนีวรโณ เอโก ขนฺโธ จ นีวรณา จ ติณฺณนฺนํ ขนฺธานํ จิตฺตสมุฏฺฐานานญฺจ รูปานํ อตฺถิปจฺจเยน ปจฺจโย ฯเปฯ เทฺว ขนฺธา ฯเปฯ.\csromanspacer{} นีวรณา จ วตฺถุ จ โนนีวรณานํ ขนฺธานํ อตฺถิปจฺจเยน ปจฺจโย.\csromanspacer{} นีวรณา จ สมฺปยุตฺตกา จ ขนฺธา จิตฺตสมุฏฺฐานานํ รูปานํ อตฺถิปจฺจเยน \csromanfolio{312}ปจฺจโย.\csromanspacer{} นีวรณา จ มหาภูตา จ จิตฺตสมุฏฺฐานานํ รูปานํ อตฺถิปจฺจเยน ปจฺจโย.\csromanspacer{} \textbf{ปจฺฉาชาตา} นีวรณา จ สมฺปยุตฺตกา จ ขนฺธา ปุเรชาตสฺส อิมสฺส กายสฺส อตฺถิปจฺจเยน ปจฺจโย.\csromanspacer{} \textbf{ปจฺฉาชาตา} นีวรณา จ สมฺปยุตฺตกา ขนฺธา จ กพฬีกาโร อาหาโร จ อิมสฺส กายสฺส อตฺถิปจฺจเยน ปจฺจโย.\csromanspacer{} \textbf{ปจฺฉาชาตา} นีวรณา จ สมฺปยุตฺตกา ขนฺธา จ รูปชีวิตินฺทฺริยญฺจ กฏตฺตารูปานํ อตฺถิปจฺจเยน ปจฺจโย. (2)}

\prose{นีวรโณ จ โนนีวรโณ จ ธมฺมา นีวรณสฺส จ โนนีวรณสฺส จ ธมฺมสฺส อตฺถิปจฺจเยน ปจฺจโย.\csromanspacer{} สหชาตํ, ปุเรชาตํ.\csromanspacer{}\csromanspacer{}สหชาโต โนนีวรโณ เอโก ขนฺโธ จ กามจฺฉนฺทนีวรณญฺจ ติณฺณนฺนํ ขนฺธานํ ถินมิทฺธนีวรณสฺส, อุทฺธจฺจนีวรณสฺส, อวิชฺชานีวรณสฺส จ จิตฺตสมุฏฺฐานานญฺจ รูปานํ อตฺถิปจฺจเยน ปจฺจโย.\csromanspacer{} เทฺว ขนฺธา จ ฯเปฯ.\csromanspacer{} กามจฺฉนฺทนีวรณญฺจ วตฺถุ จ ถินมิทฺธนีวรณสฺส, อุทฺธจฺจนีวรณสฺส, อวิชฺชานีวรณสฺส จ สมฺปยุตฺตกานญฺจ ขนฺธานํ อตฺถิปจฺจเยน ปจฺจโย. (จกฺกํ.) (3)}

\csromanheader{2}{1. ปจฺจยานุโลม 2. สงฺขฺยาวาร}
\csromanheader{2}{\textbf{สุทฺธ}}
\proseitem{44}{เหตุยา จตฺตาริ, อารมฺมเณ นว, อธิปติยา นว, อนนฺตเร นว, สมนนฺตเร นว, สหชาเต นว, อญฺญมญฺเญ นว, นิสฺสเย นว, อุปนิสฺสเย นว, ปุเรชาเต ตีณิ, ปจฺฉาชาเต ตีณิ, อาเสวเน นว, กมฺเม ตีณิ, วิปาเก เอกํ, อาหาเร ตีณิ, อินฺทฺริเย ตีณิ, ฌาเน ตีณิ, มคฺเค ตีณิ, สมฺปยุตฺเต นว, วิปฺปยุตฺเต ปญฺจ, อตฺถิยา นว, นตฺถิยา นว, วิคเต นว, อวิคเต นว.}

\vspace{\tipitakaparskip}
\csromancenter{อนุโลมํ.}
\csromanheader{2}{2. ปจฺจนียุทฺธาร}
\proseitem{45}{นีวรโณ ธมฺโม นีวรณสฺส ธมฺมสฺส อารมฺมณปจฺจเยน ปจฺจโย.\csromanspacer{} สหชาตปจฺจเยน ปจฺจโย.\csromanspacer{} อุปนิสฺสยปจฺจเยน ปจฺจโย. (1)}

\csromanheader{2}{44. นีวรณทุก}
\prose{\csromanfolio{313}นีวรโณ ธมฺโม โนนีวรณสฺส ธมฺมสฺส อารมฺมณปจฺจเยน ปจฺจโย.\csromanspacer{} สหชาตปจฺจเยน ปจฺจโย.\csromanspacer{} อุปนิสฺสยปจฺจเยน ปจฺจโย.\csromanspacer{} ปจฺฉาชาตปจฺจเยน ปจฺจโย. (2)}

\prose{นีวรโณ ธมฺโม นีวรณสฺส จ โนนีวรณสฺส จ ธมฺมสฺส อารมฺมณปจฺจเยน ปจฺจโย.\csromanspacer{} สหชาตปจฺจเยน ปจฺจโย.\csromanspacer{} อุปนิสฺสยปจฺจเยน ปจฺจโย. (3)}

\proseitem{46}{โนนีวรโณ ธมฺโม โนนีวรณสฺส ธมฺมสฺส อารมฺมณปจฺจเยน ปจฺจโย.\csromanspacer{} สหชาตปจฺจเยน ปจฺจโย.\csromanspacer{} อุปนิสฺสยปจฺจเยน ปจฺจโย.\csromanspacer{} ปุเรชาตปจฺจเยน ปจฺจโย.\csromanspacer{} ปจฺฉาชาตปจฺจเยน ปจฺจโย.\csromanspacer{} กมฺมปจฺจเยน ปจฺจโย.\csromanspacer{} อาหารปจฺจเยน ปจฺจโย.\csromanspacer{} อินฺทฺริยปจฺจเยน ปจฺจโย. (1)}

\prose{โนนีวรโณ ธมฺโม นีวรณสฺส ธมฺมสฺส อารมฺมณปจฺจเยน ปจฺจโย.\csromanspacer{} สหชาตปจฺจเยน ปจฺจโย.\csromanspacer{} อุปนิสฺสยปจฺจเยน ปจฺจโย.\csromanspacer{} ปุเรชาตปจฺจเยน ปจฺจโย. (2)}

\prose{โนนีวรโณ ธมฺโม นีวรณสฺส จ โนนีวรณสฺส จ ธมฺมสฺส อารมฺมณปจฺจเยน ปจฺจโย.\csromanspacer{} สหชาตปจฺจเยน ปจฺจโย.\csromanspacer{} อุปนิสฺสยปจฺจเยน ปจฺจโย.\csromanspacer{} ปุเรชาตปจฺจเยน ปจฺจโย. (3)}

\proseitem{47}{นีวรโณ จ โนนีวรโณ จ ธมฺมา นีวรณสฺส ธมฺมสฺส อารมฺมณปจฺจเยน ปจฺจโย.\csromanspacer{} สหชาตปจฺจเยน ปจฺจโย.\csromanspacer{} อุปนิสฺสยปจฺจเยน ปจฺจโย. (1)}

\prose{นีวรโณ จ โนนีวรโณ จ ธมฺมา โนนีวรณสฺส ธมฺมสฺส อารมฺมณปจฺจเยน ปจฺจโย.\csromanspacer{} สหชาตปจฺจเยน ปจฺจโย.\csromanspacer{} อุปนิสฺสยปจฺจเยน ปจฺจโย.\csromanspacer{} ปจฺฉาชาตปจฺจเยน ปจฺจโย. (2)}

\prose{นีวรโณ จ โนนีวรโณ จ ธมฺมา นีวรณสฺส จ โนนีวรณสฺส จ ธมฺมสฺส อารมฺมณปจฺจเยน ปจฺจโย.\csromanspacer{} สหชาตปจฺจเยน ปจฺจโย.\csromanspacer{} อุปนิสฺสยปจฺจเยน ปจฺจโย. (3)}

\csromanheader{2}{2. ปจฺจยปจฺจนีย 2. สงฺขฺยาวาร}
\proseitem{48}{นเหตุยา นว, น-อารมฺมเณ นว, น-อธิปติยา นว, (สพฺพตฺถ นว.) โนวิคเต นว, โน-อวิคเต นว.}

\csromanheader{2}{3. ปจฺจยานุโลมปจฺจนีย}
\csromanheader{2}{\textbf{เหตุทุก}}
\proseitem{49}{\csromanfolio{314}\textbf{เหตุ}ปจฺจยา น-อารมฺมเณ จตฺตาริ ฯเปฯ นสมนนฺตเร จตฺตาริ, น-อญฺญมญฺเญ เทฺว, น-อุปนิสฺสเย จตฺตาริ ฯเปฯ นมคฺเค จตฺตาริ, นสมฺปยุตฺเต เทฺว, นวิปฺปยุตฺเต จตฺตาริ, โนนตฺถิยา จตฺตาริ, โนวิคเต จตฺตาริ.}

\csromanheader{2}{4. ปจฺจยปจฺจนียานุโลม}
\proseitem{50}{นเหตุปจฺจยา อารมฺมเณ นว, อธิปติยา นว, (อนุโลมมาติกา.) ฯเปฯ อวิคเต นว.}

\nitthitam{นีวรณทุกํ นิฏฺฐิตํ.}
\csromanmatikaanchor{mtk.42}
\csromantocmarkref{subsection}{45. นีวรณิยทุก}{mtk.42}
\bookmark[dest={mtk.42},level=2]{45. นีวรณิยทุก}
\csromanheader{2}{45. นีวรณิยทุก 1. ปฏิจฺจวาร}
\proseitem{51}{นีวรณิยํ ธมฺมํ ปฏิจฺจ นีวรณิโย ธมฺโม อุปฺปชฺชติ เหตุปจฺจยา ฯเปฯ. (นีวรณิยทุกํ ยถา โลกิยทุกํ, เอวํ กาตพฺพํ.\csromanspacer{} นินฺนานากรณํ.)}

\csromangathameasure{ทฺวิกฺขตฺตุํ กามจฺฉนฺเทน, จตุกฺขตฺตุํ ปฏิเฆน จ. \\ อุทฺธจฺจํ วิจิกิจฺฉา จ, อุโภเปเต สกิํ สกิํ.}
\csromangathagroup{\csromangathastanza{ทฺวิกฺขตฺตุํ กามจฺฉนฺเทน, จตุกฺขตฺตุํ ปฏิเฆน จ. \\ อุทฺธจฺจํ วิจิกิจฺฉา จ, อุโภเปเต สกิํ สกิํ.}}

\prose{นีวรณานํ นีวรเณหิ, อฏฺฐวิธํ ปโยชนํ. (นีวรณทุกสฺส มาติกา อิธ กตา.)}

\nitthitam{นีวรณิยทุกํ นิฏฺฐิตํ.}
\csromanmatikaanchor{mtk.43}
\csromantocmarkref{subsection}{46. นีวรณสมฺปยุตฺตทุก}{mtk.43}
\bookmark[dest={mtk.43},level=2]{46. นีวรณสมฺปยุตฺตทุก}
\csromanheader{2}{46. นีวรณสมฺปยุตฺตทุก 1. ปฏิจฺจวาร}
\csromanheader{2}{1. ปจฺจยานุโลม 1. วิภงฺควาร}
\csromanheader{2}{\textbf{เหตุ}}
\proseitem{52}{นีวรณสมฺปยุตฺตํ ธมฺมํ ปฏิจฺจ นีวรณสมฺปยุตฺโต ธมฺโม อุปฺปชฺชติ เหตุปจฺจยา.\csromanspacer{} นีวรณสมฺปยุตฺตํ เอกํ ขนฺธํ ปฏิจฺจ ตโย ขนฺธา ฯเปฯ เทฺว ขนฺเธ ฯเปฯ. (1)}

\csromanheader{2}{46. นีวรณสมฺปยุตฺตทุก}
\prose{\csromanfolio{315}นีวรณสมฺปยุตฺตํ ธมฺมํ ปฏิจฺจ นีวรณวิปฺปยุตฺโต ธมฺโม อุปฺปชฺชติ เหตุปจฺจยา.\csromanspacer{} นีวรณสมฺปยุตฺเต ขนฺเธ ปฏิจฺจ จิตฺตสมุฏฺฐานํ รูปํ. (2)}

\prose{นีวรณสมฺปยุตฺตํ ธมฺมํ ปฏิจฺจ นีวรณสมฺปยุตฺโต จ นีวรณวิปฺปยุตฺโต จ ธมฺมา อุปฺปชฺชนฺติ เหตุปจฺจยา.\csromanspacer{} นีวรณสมฺปยุตฺตํ เอกํ ขนฺธํ ปฏิจฺจ ตโย ขนฺธา จิตฺตสมุฏฺฐานญฺจ รูปํ ฯเปฯ เทฺว ขนฺเธ ฯเปฯ. (3)}

\proseitem{53}{นีวรณวิปฺปยุตฺตํ ธมฺมํ ปฏิจฺจ นีวรณวิปฺปยุตฺโต ธมฺโม อุปฺปชฺชติ เหตุปจฺจยา.\csromanspacer{} นีวรณวิปฺปยุตฺตํ เอกํ ขนฺธํ ปฏิจฺจ ตโย ขนฺธา จิตฺตสมุฏฺฐานญฺจ รูปํ ฯเปฯ เทฺว ขนฺเธ ฯเปฯ.\csromanspacer{} ปฏิสนฺธิกฺขเณ ฯเปฯ. (ยาว อชฺฌตฺติกา มหาภูตา.) (1)}

\prose{นีวรณสมฺปยุตฺตญฺจ นีวรณวิปฺปยุตฺตญฺจ ธมฺมํ ปฏิจฺจ นีวรณวิปฺปยุตฺโต ธมฺโม อุปฺปชฺชติ เหตุปจฺจยา.\csromanspacer{} นีวรณสมฺปยุตฺเต ขนฺเธ จ มหาภูเต จ ปฏิจฺจ จิตฺตสมุฏฺฐานํ รูปํ. (สํขิตฺตํ.)}

\csromanheader{2}{1. ปจฺจยานุโลม 2. สงฺขฺยาวาร}
\proseitem{54}{เหตุยา ปญฺจ, อารมฺมเณ เทฺว, อธิปติยา ปญฺจ ฯเปฯ นตฺถิยา เทฺว ฯเปฯ อวิคเต ปญฺจ.}

\vspace{\tipitakaparskip}
\csromancenter{อนุโลมํ.}
\csromanheader{2}{2. ปจฺจยปจฺจนีย 1. วิภงฺควาร}
\csromanheader{2}{\textbf{นเหตุ}}
\proseitem{55}{นีวรณสมฺปยุตฺตํ ธมฺมํ ปฏิจฺจ นีวรณสมฺปยุตฺโต ธมฺโม อุปฺปชฺชติ นเหตุปจฺจยา.\csromanspacer{} วิจิกิจฺฉาสหคเต อุทฺธจฺจสหคเต ขนฺเธ ปฏิจฺจ อวิชฺชานีวรณํ. (1)}

\prose{นีวรณวิปฺปยุตฺตํ ธมฺมํ ปฏิจฺจ นีวรณวิปฺปยุตฺโต ธมฺโม อุปฺปชฺชติ นเหตุปจฺจยา.\csromanspacer{} อเหตุกํ นีวรณวิปฺปยุตฺตํ ฯเปฯ. (ยาว อสญฺญสตฺตา.\csromanspacer{} สํขิตฺตํ.)}

\csromanheader{2}{2. ปจฺจยปจฺจนีย 2. สงฺขฺยาวาร}
\csromanheader{2}{\textbf{สุทฺธ}}
\proseitem{56}{\csromanfolio{316}นเหตุยา เทฺว, น-อารมฺมเณ ตีณิ, น-อธิปติยา ปญฺจ, น-อนนฺตเร, นสมนนฺตเร, น-อญฺญมญฺเญ, น-อุปนิสฺสเย ตีณิ, นปุเรชาเต จตฺตาริ, นปจฺฉาชาเต ปญฺจ, น-อาเสวเน ปญฺจ, นกมฺเม เทฺว, นวิปาเก ปญฺจ, น-อาหาเร เอกํ, น-อินฺทฺริเย เอกํ, นฌาเน เอกํ, นมคฺเค เอกํ, นสมฺปยุตฺเต ตีณิ, นวิปฺปยุตฺเต เทฺว, โนนตฺถิยา ตีณิ, โนวิคเต ตีณิ.}

\csromanheader{2}{3. ปจฺจยานุโลมปจฺจนีย}
\csromanheader{2}{\textbf{เหตุทุก}}
\proseitem{57}{\textbf{เหตุ}ปจฺจยา น-อารมฺมเณ ตีณิ, น-อธิปติยา ปญฺจ, นปุเรชาเต จตฺตาริ ฯเปฯ นกมฺเม เทฺว, นวิปาเก ปญฺจ, นสมฺปยุตฺเต ตีณิ, นวิปฺปยุตฺเต เทฺว, โนนตฺถิยา ตีณิ, โนวิคเต ตีณิ.}

\csromanheader{2}{4. ปจฺจยปจฺจนียานุโลม}
\vspace{\tipitakaparskip}
\csromancenter{นเหตุปจฺจยา อารมฺมเณ เทฺว ฯเปฯ มคฺเค เอกํ ฯเปฯ อวิคเต เทฺว.}
\csromancenter{(สหชาตวาโรปิ เอวํ กาตพฺโพ.)}
\csromanheader{2}{46. นีวรณสมฺปยุตฺตทุก 3. ปจฺจยวาร}
\csromanheader{2}{1. ปจฺจยานุโลม 1. วิภงฺควาร}
\csromanheader{2}{\textbf{เหตุ}}
\proseitem{59}{นีวรณสมฺปยุตฺตํ ธมฺมํ ปจฺจยา นีวรณสมฺปยุตฺโต ธมฺโม อุปฺปชฺชติ เหตุปจฺจยา.\csromanspacer{} ตีณิ.}

\prose{นีวรณวิปฺปยุตฺตํ ธมฺมํ ปจฺจยา นีวรณวิปฺปยุตฺโต ธมฺโม อุปฺปชฺชติ เหตุปจฺจยา.\csromanspacer{} นีวรณวิปฺปยุตฺตํ เอกํ ขนฺธํ ปจฺจยา ตโย ขนฺธา จิตฺตสมุฏฺฐานญฺจ}

\vspace{\tipitakaparskip}
\csromancenter{นีวรณสมฺปยุตฺตทุก รูปํ ฯเปฯ เทฺว ขนฺเธ ฯเปฯ.\csromanspacer{} ปฏิสนฺธิกฺขเณ ฯเปฯ.\csromanspacer{} เอกํ มหาภูตํ ฯเปฯ.\csromanspacer{} วตฺถุํ ปจฺจยา นีวรณวิปฺปยุตฺตา ขนฺธา. (1)}
\prose{\csromanfolio{317}นีวรณวิปฺปยุตฺตํ ธมฺมํ ปจฺจยา นีวรณสมฺปยุตฺโต ธมฺโม อุปฺปชฺชติ เหตุปจฺจยา.\csromanspacer{} วตฺถุํ ปจฺจยา นีวรณสมฺปยุตฺตา ขนฺธา. (2)}

\prose{นีวรณวิปฺปยุตฺตํ ธมฺมํ ปจฺจยา นีวรณสมฺปยุตฺโต จ นีวรณวิปฺปยุตฺโต จ ธมฺมา อุปฺปชฺชนฺติ เหตุปจฺจยา.\csromanspacer{} วตฺถุํ ปจฺจยา นีวรณสมฺปยุตฺตา ขนฺธา.\csromanspacer{} มหาภูเต ปจฺจยา จิตฺตสมุฏฺฐานํ รูปํ. (3)}

\proseitem{60}{นีวรณสมฺปยุตฺตญฺจ นีวรณวิปฺปยุตฺตญฺจ ธมฺมํ ปจฺจยา นีวรณสมฺปยุตฺโต ธมฺโม อุปฺปชฺชติ เหตุปจฺจยา.\csromanspacer{} นีวรณสมฺปยุตฺตํ เอกํ ขนฺธญฺจ วตฺถุญฺจ ปจฺจยา ตโย ขนฺธา ฯเปฯ เทฺว ขนฺเธ ฯเปฯ. (1)}

\prose{นีวรณสมฺปยุตฺตญฺจ นีวรณวิปฺปยุตฺตญฺจ ธมฺมํ ปจฺจยา นีวรณวิปฺปยุตฺโต ธมฺโม อุปฺปชฺชติ เหตุปจฺจยา.\csromanspacer{} นีวรณสมฺปยุตฺเต ขนฺเธ จ มหาภูเต จ ปจฺจยา จิตฺตสมุฏฺฐานํ รูปํ. (2)}

\prose{นีวรณสมฺปยุตฺตญฺจ นีวรณวิปฺปยุตฺตญฺจ ธมฺมํ ปจฺจยา นีวรณสมฺปยุตฺโต จ นีวรณวิปฺปยุตฺโต จ ธมฺมา อุปฺปชฺชนฺติ เหตุปจฺจยา.\csromanspacer{} นีวรณสมฺปยุตฺตํ เอกํ ขนฺธญฺจ วตฺถุญฺจ ปจฺจยา ตโย ขนฺธา ฯเปฯ เทฺว ขนฺเธ ฯเปฯ.\csromanspacer{} นีวรณสมฺปยุตฺเต ขนฺเธ จ มหาภูเต จ ปจฺจยา จิตฺตสมุฏฺฐานํ รูปํ. (สํขิตฺตํ.) (3)}

\csromanheader{2}{1. ปจฺจยานุโลม 2. สงฺขฺยาวาร}
\csromanheader{2}{\textbf{สุทฺธ}}
\proseitem{61}{เหตุยา นว, อารมฺมเณ จตฺตาริ, อธิปติยา นว, อนนฺตเร จตฺตาริ ฯเปฯ วิปาเก เอกํ ฯเปฯ อวิคเต นว.}

\vspace{\tipitakaparskip}
\csromancenter{อนุโลมํ.}
\csromanheader{2}{2. ปจฺจยปจฺจนีย 1. วิภงฺควาร}
\csromanheader{2}{\textbf{นเหตุ}}
\proseitem{62}{\csromanfolio{318}นีวรณสมฺปยุตฺตํ ธมฺมํ ปจฺจยา นีวรณสมฺปยุตฺโต ธมฺโม อุปฺปชฺชติ นเหตุปจฺจยา.\csromanspacer{} วิจิกิจฺฉาสหคเต อุทฺธจฺจสหคเต ขนฺเธ ปจฺจยา อวิชฺชานีวรณํ. (1)}

\prose{นีวรณวิปฺปยุตฺตํ ธมฺมํ ปจฺจยา นีวรณวิปฺปยุตฺโต ธมฺโม อุปฺปชฺชติ นเหตุปจฺจยา.\csromanspacer{} อเหตุกํ นีวรณวิปฺปยุตฺตํ เอกํ ขนฺธํ ปจฺจยา ตโย ขนฺธา จิตฺตสมุฏฺฐานญฺจ รูปํ ฯเปฯ. (ยาว อสญฺญสตฺตา.) จกฺขายตนํ ปจฺจยา จกฺขุวิญฺญาณํ ฯเปฯ.\csromanspacer{} กายายตนํ ปจฺจยา กายวิญฺญาณํ.\csromanspacer{} วตฺถุํ ปจฺจยา อเหตุกา นีวรณวิปฺปยุตฺตา ขนฺธา. (1)}

\prose{นีวรณวิปฺปยุตฺตํ ธมฺมํ ปจฺจยา นีวรณสมฺปยุตฺโต ธมฺโม อุปฺปชฺชติ นเหตุปจฺจยา.\csromanspacer{} วตฺถุํ ปจฺจยา วิจิกิจฺฉาสหคโต อุทฺธจฺจสหคโต โมโห. (2)}

\prose{นีวรณสมฺปยุตฺตญฺจ นีวรณวิปฺปยุตฺตญฺจ ธมฺมํ ปจฺจยา นีวรณสมฺปยุตฺโต ธมฺโม อุปฺปชฺชติ นเหตุปจฺจยา.\csromanspacer{} วิจิกิจฺฉาสหคเต อุทฺธจฺจสหคเต ขนฺเธ จ วตฺถุญฺจ ปจฺจยา วิจิกิจฺฉาสหคโต อุทฺธจฺจสหคโต โมโห. (สํขิตฺตํ.) (1)}

\csromanheader{2}{2. ปจฺจยปจฺจนีย 2. สงฺขฺยาวาร}
\csromanheader{2}{\textbf{สุทฺธ}}
\vspace{\tipitakaparskip}
\csromancenter{\textbf{นเหตุ}ยา จตฺตาริ, น-อารมฺมเณ ตีณิ ฯเปฯ นปุเรชาเต จตฺตาริ, นปจฺฉาชาเต นว, น-อาเสวเน นว, นกมฺเม จตฺตาริ, นวิปาเก นว ฯเปฯ นสมฺปยุตฺเต ตีณิ, นวิปฺปยุตฺเต เทฺว, โนนตฺถิยา ตีณิ, โนวิคเต ตีณิ.}
\csromancenter{(เอวํ อิตเร เทฺว คณนาปิ นิสฺสยวาโรปิ กาตพฺโพ.)}
\csromancenter{นีวรณสมฺปยุตฺตทุก \textbf{46.}\csromanspacer{}\textbf{ นีวรณสมฺปยุตฺตทุก 5.}\csromanspacer{}\textbf{ สํสฏฺฐวาร}}
\csromanheader{2}{1-4. ปจฺจยจตุกฺก}
\proseitem{64}{\csromanfolio{319}นีวรณสมฺปยุตฺตํ ธมฺมํ สํสฏฺโฐ นีวรณสมฺปยุตฺโต ธมฺโม อุปฺปชฺชติ เหตุปจฺจยา ฯเปฯ.}

\prose{\textbf{เหตุ}ยา เทฺว, อารมฺมเณ เทฺว, (สพฺพตฺถ เทฺว.) วิปาเก เอกํ ฯเปฯ อวิคเต เทฺว.}

\vspace{\tipitakaparskip}
\csromancenter{อนุโลมํ.}
\proseitem{65}{นีวรณสมฺปยุตฺตํ ธมฺมํ สํสฏฺโฐ นีวรณสมฺปยุตฺโต ธมฺโม อุปฺปชฺชติ นเหตุปจฺจยา.\csromanspacer{} วิจิกิจฺฉาสหคเต อุทฺธจฺจสหคเต ขนฺเธ สํสฏฺโฐ วิจิกิจฺฉาสหคโต อุทฺธจฺจสหคโต โมโห.}

\csromanheader{2}{นีวรณวิปฺปยุตฺตํ ธมฺมํ สํสฏฺโฐ. (สํขิตฺตํ.)}
\prose{นเหตุยา เทฺว, น-อธิปติยา เทฺว, นปุเรชาเต เทฺว, นกมฺเม เทฺว, นวิปาเก เทฺว, นฌาเน เอกํ, นมคฺเค เอกํ, นวิปฺปยุตฺเต เทฺว.}

\vspace{\tipitakaparskip}
\csromancenter{ปจฺจนียํ.}
\csromancenter{(เอวํ อิตเร เทฺว คณนาปิ สมฺปยุตฺตวาโรปิ กาตพฺโพ.)}
\csromanheader{2}{46. นีวรณสมฺปยุตฺตทุก 7. ปญฺหาวาร}
\csromanheader{2}{1. ปจฺจยานุโลม 1. วิภงฺควาร}
\csromanheader{2}{\textbf{เหตุ}}
\proseitem{66}{นีวรณสมฺปยุตฺโต ธมฺโม นีวรณสมฺปยุตฺตสฺส ธมฺมสฺส เหตุปจฺจเยน ปจฺจโย.\csromanspacer{} นีวรณสมฺปยุตฺตา เหตู สมฺปยุตฺตกานํ ขนฺธานํ เหตุปจฺจเยน ปจฺจโย. (มูลํ ปุจฺฉิตพฺพํ.) นีวรณสมฺปยุตฺตา เหตู จิตฺตสมุฏฺฐานานํ รูปานํ เหตุปจฺจเยน ปจฺจโย. (มูลํ ปุจฺฉิตพฺพํ.) \csromanfolio{320}นีวรณสมฺปยุตฺตา เหตู สมฺปยุตฺตกานํ ขนฺธานํ จิตฺตสมุฏฺฐานานญฺจ รูปานํ เหตุปจฺจเยน ปจฺจโย. (3)}

\prose{นีวรณวิปฺปยุตฺโต ธมฺโม นีวรณวิปฺปยุตฺตสฺส ธมฺมสฺส เหตุปจฺจเยน ปจฺจโย.\csromanspacer{} นีวรณวิปฺปยุตฺตา เหตู สมฺปยุตฺตกานํ ขนฺธานํ จิตฺตสมุฏฺฐานานญฺจ รูปานํ เหตุปจฺจเยน ปจฺจโย.\csromanspacer{} ปฏิสนฺธิกฺขเณ ฯเปฯ. (1)}

\csromanheader{2}{\textbf{อารมฺมณ}}
\proseitem{67}{นีวรณสมฺปยุตฺโต ธมฺโม นีวรณสมฺปยุตฺตสฺส ธมฺมสฺส อารมฺมณปจฺจเยน ปจฺจโย.\csromanspacer{} ราคํ อสฺสาเทติ อภินนฺทติ, ตํ อารพฺภ ราโค อุปฺปชฺชติ, ทิฏฺฐิ ฯเปฯ วิจิกิจฺฉา ฯเปฯ อุทฺธจฺจํ ฯเปฯ โทมนสฺสํ อุปฺปชฺชติ.\csromanspacer{} ทิฏฺฐิํ อสฺสาเทติ อภินนฺทติ, ตํ อารพฺภ ราโค อุปฺปชฺชติ.\csromanspacer{} ทิฏฺฐิ ฯเปฯ วิจิกิจฺฉา ฯเปฯ อุทฺธจฺจํ ฯเปฯ โทมนสฺสํ อุปฺปชฺชติ.\csromanspacer{} วิจิกิจฺฉํ อารพฺภ วิจิกิจฺฉา ฯเปฯ ทิฏฺฐิ ฯเปฯ อุทฺธจฺจํ ฯเปฯ โทมนสฺสํ อุปฺปชฺชติ.\csromanspacer{} อุทฺธจฺจํ อารพฺภ อุทฺธจฺจํ อุปฺปชฺชติ.\csromanspacer{} ทิฏฺฐิ ฯเปฯ วิจิกิจฺฉา.\csromanspacer{} โทมนสฺสํ อุปฺปชฺชติ.\csromanspacer{} โทมนสฺสํ อารพฺภ โทมนสฺสํ อุปฺปชฺชติ.\csromanspacer{} ทิฏฺฐิ ฯเปฯ วิจิกิจฺฉา ฯเปฯ อุทฺธจฺจํ อุปฺปชฺชติ. (1)}

\prose{นีวรณสมฺปยุตฺโต ธมฺโม นีวรณวิปฺปยุตฺตสฺส ธมฺมสฺส อารมฺมณปจฺจเยน ปจฺจโย.\csromanspacer{} อริยา นีวรณสมฺปยุตฺเต ปหีเน กิเลเส ปจฺจเวกฺขนฺติ, วิกฺขมฺภิเต กิเลเส ปจฺจเวกฺขนฺติ.\csromanspacer{} ปุพฺเพ สมุทาจิณฺเณ กิเลเส ชานนฺติ.\csromanspacer{} นีวรณสมฺปยุตฺเต ขนฺเธ อนิจฺจโต ทุกฺขโต อนตฺตโต วิปสฺสติ.\csromanspacer{} เจโตปริยญาเณน นีวรณสมฺปยุตฺตจิตฺตสมงฺคิสฺส จิตฺตํ ชานาติ.\csromanspacer{} นีวรณสมฺปยุตฺตา ขนฺธา เจโตปริยญาณสฺส, ปุพฺเพนิวาสานุสฺสติญาณสฺส, ยถากมฺมูปคญาณสฺส, อนาคตํสญาณสฺส, อาวชฺชนาย อารมฺมณปจฺจเยน ปจฺจโย. (2)}

\proseitem{68}{นีวรณวิปฺปยุตฺโต ธมฺโม นีวรณวิปฺปยุตฺตสฺส ธมฺมสฺส อารมฺมณปจฺจเยน ปจฺจโย.\csromanspacer{} ทานํ ฯเปฯ สีลํ ฯเปฯ อุโปสถกมฺมํ กตฺวา ตํ ปจฺจเวกฺขติ.\csromanspacer{} ปุพฺเพ สุจิณฺณานิ ฯเปฯ.\csromanspacer{} ฌานา วุฏฺฐหิตฺวา ฌานํ ฯเปฯ.\csromanspacer{} อริยา มคฺคา วุฏฺฐหิตฺวา มคฺคํ ฯเปฯ.\csromanspacer{} ผลํ ฯเปฯ.\csromanspacer{} นิพฺพานํ โคตฺรภุสฺส, โวทานสฺส, มคฺคสฺส, ผลสฺส, อาวชฺชนาย อารมฺมณปจฺจเยน ปจฺจโย.\csromanspacer{} จกฺขุํ ฯเปฯ วตฺถุํ นีวรณวิปฺปยุตฺเต ขนฺเธ อนิจฺจโต ทุกฺขโต อนตฺตโต วิปสฺสติ.\csromanspacer{} ทิพฺเพน จกฺขุนา รูปํ ปสฺสติ. (ยาว อาวชฺชนาย.) (1)}

\csromanheader{2}{46. นีวรณสมฺปยุตฺตทุก}
\prose{\csromanfolio{321}นีวรณาวิปฺปยุตฺโต ธมฺโม นีวรณสมฺปยุตฺตสฺส ธมฺมสฺส อารมฺมณปจฺจเยน ปจฺจโย.\csromanspacer{} ทานํ ฯเปฯ สีลํ ฯเปฯ อุโปสถกมฺมํ ฯเปฯ.\csromanspacer{} ปุพฺเพ ฯเปฯ.\csromanspacer{} ฌานา ฯเปฯ.\csromanspacer{} จกฺขุํ ฯเปฯ วตฺถุํ นีวรณวิปฺปยุตฺเต ขนฺเธ อสฺสาเทติ อภินนฺทติ, ตํ อารพฺภ ราโค อุปฺปชฺชติ, ทิฏฺฐิ ฯเปฯ วิจิกิจฺฉา ฯเปฯ อุทฺธจฺจํ ฯเปฯ โทมนสฺสํ อุปฺปชฺชติ. (2)}

\csromanheader{2}{\textbf{อธิปติ}}
\proseitem{69}{นีวรณสมฺปยุตฺโต ธมฺโม นีวรณสมฺปยุตฺตสฺส ธมฺมสฺส อธิปติปจฺจเยน ปจฺจโย.\csromanspacer{} \textbf{อารมฺมณาธิปติ}, สหชาตาธิปติ.\csromanspacer{}\csromanspacer{}\textbf{อารมฺมณาธิปติ}\csromandash{}ราคํ ครุํ กตฺวา อสฺสาเทติ อภินนฺทติ, ตํ ครุํ กตฺวา ราโค อุปฺปชฺชติ, ทิฏฺฐิ อุปฺปชฺชติ.\csromanspacer{} ทิฏฺฐิํ ครุํ กตฺวา อสฺสาเทติ ฯเปฯ.\csromanspacer{} \textbf{สหชาตาธิปติ}\csromandash{}นีวรณสมฺปยุตฺตาธิปติ นีวรณสมฺปยุตฺตกานํ ขนฺธานํ อธิปติปจฺจเยน ปจฺจโย. (1)}

\prose{นีวรณสมฺปยุตฺโต ธมฺโม นีวรณวิปฺปยุตฺตสฺส ธมฺมสฺส อธิปติปจฺจเยน ปจฺจโย.\csromanspacer{} \textbf{สหชาตาธิปติ}\csromandash{}นีวรณสมฺปยุตฺตาธิปติ จิตฺตสมุฏฺฐานานํ รูปานํ อธิปติปจฺจเยน ปจฺจโย. (มูลํ ปุจฺฉิตพฺพํ.) นีวรณสมฺปยุตฺตาธิปติ สมฺปยุตฺตกานํ ขนฺธานํ จิตฺตสมุฏฺฐานานญฺจ รูปานํ อธิปติปจฺจเยน ปจฺจโย. (3)}

\proseitem{70}{นีวรณวิปฺปยุตฺโต ธมฺโม นีวรณวิปฺปยุตฺตสฺส ธมฺมสฺส อธิปติปจฺจเยน ปจฺจโย.\csromanspacer{} \textbf{อารมฺมณาธิปติ}, สหชาตาธิปติ.\csromanspacer{}\csromanspacer{}\textbf{อารมฺมณาธิปติ}\csromandash{}ทานํ ฯเปฯ สีลํ ฯเปฯ อุโปสถกมฺมํ กตฺวา ตํ ครุํ กตฺวา ปจฺจเวกฺขติ.\csromanspacer{} ปุพฺเพ สุจิณฺณานิ ฯเปฯ.\csromanspacer{} ฌานา วุฏฺฐหิตฺวา ฌานํ ฯเปฯ.\csromanspacer{} อริยา มคฺคา วุฏฺฐหิตฺวา มคฺคํ ครุํ กตฺวา ฯเปฯ.\csromanspacer{} ผลํ ฯเปฯ.\csromanspacer{} นิพฺพานํ ฯเปฯ.\csromanspacer{} นิพฺพานํ โคตฺรภุสฺส, โวทานสฺส, มคฺคสฺส, ผลสฺส อธิปติปจฺจเยน ปจฺจโย.\csromanspacer{} \textbf{สหชาตาธิปติ}\csromandash{}นีวรณวิปฺปยุตฺตาธิปติ สมฺปยุตฺตกานํ ขนฺธานํ จิตฺตสมุฏฺฐานานญฺจ รูปานํ อธิปติปจฺจเยน ปจฺจโย. (1)}

\prose{นีวรณวิปฺปยุตฺโต ธมฺโม นีวรณสมฺปยุตฺตสฺส ธมฺมสฺส อธิปติปจฺจเยน ปจฺจโย.\csromanspacer{} \textbf{อารมฺมณาธิปติ}\csromandash{}ทานํ ฯเปฯ สีลํ ฯเปฯ.\csromanspacer{} ปุพฺเพ ฯเปฯ.\csromanspacer{} ฌานํ ฯเปฯ.\csromanspacer{} จกฺขุํ ฯเปฯ วตฺถุํ นีวรณวิปฺปยุตฺเต ขนฺเธ ครุํ กตฺวา อสฺสาเทติ อภินนฺทติ, ตํ ครุํ กตฺวา ราโค อุปฺปชฺชติ, ทิฏฺฐิ อุปฺปชฺชติ. (2)}

\csromanheader{2}{\textbf{อนนฺตร}}
\proseitem{71}{\csromanfolio{322}นีวรณสมฺปยุตฺโต ธมฺโม นีวรณสมฺปยุตฺตสฺส ธมฺมสฺส อนนฺตรปจฺจเยน ปจฺจโย.\csromanspacer{} ปุริมา ปุริมา นีวรณสมฺปยุตฺตา ขนฺธา ปจฺฉิมานํ ปจฺฉิมานํ นีวรณสมฺปยุตฺตกานํ ขนฺธานํ อนนฺตรปจฺจเยน ปจฺจโย. (มูลํ ปุจฺฉิตพฺพํ.) นีวรณสมฺปยุตฺตา ขนฺธา วุฏฺฐานสฺส อนนฺตรปจฺจเยน ปจฺจโย. (อิธ “ปุริมา ปุริมา”ติ นตฺถิ.) (2)}

\prose{(มูลํ ปุจฺฉิตพฺพํ.) ปุริมา ปุริมา นีวรณวิปฺปยุตฺตา ขนฺธา ปจฺฉิมานํ ปจฺฉิมานํ นีวรณวิปฺปยุตฺตานํ ขนฺธานํ อนนฺตรปจฺจเยน ปจฺจโย.\csromanspacer{} อนุโลมํ โคตฺรภุสฺส ฯเปฯ ผลสมาปตฺติยา อนนฺตรปจฺจเยน ปจฺจโย. (1)}

\prose{นีวรณวิปฺปยุตฺโต ธมฺโม นีวรณสมฺปยุตฺตสฺส ธมฺมสฺส อนนฺตรปจฺจเยน ปจฺจโย.\csromanspacer{} อาวชฺชนา นีวรณสมฺปยุตฺตกานํ ขนฺธานํ อนนฺตรปจฺจเยน ปจฺจโย. (2)}

\csromanheader{2}{\textbf{สมนนฺตราทิ}}
\proseitem{72}{นีวรณสมฺปยุตฺโต ธมฺโม นีวรณสมฺปยุตฺตสฺส ธมฺมสฺส สมนนฺตรปจฺจเยน ปจฺจโย.\csromanspacer{} สหชาตปจฺจเยน ปจฺจโย.\csromanspacer{} อญฺญมญฺญปจฺจเยน ปจฺจโย.\csromanspacer{} นิสฺสยปจฺจเยน ปจฺจโย.}

\csromanheader{2}{\textbf{อุปนิสฺสย}}
\proseitem{73}{นีวรณสมฺปยุตฺโต ธมฺโม นีวรณสมฺปยุตฺตสฺส ธมฺมสฺส อุปนิสฺสยปจฺจเยน ปจฺจโย.\csromanspacer{} \textbf{อารมฺมณูปนิสฺสโย, อนนฺตรูปนิสฺสโย,} \textbf{ปกตูปนิสฺสโย ฯเปฯ.}\csromanspacer{} ปกตูปนิสฺสโย\csromandash{}ราคํ อุปนิสฺสาย ปาณํ หนติ ฯเปฯ สํฆํ ภินฺทติ.\csromanspacer{} โทสํ.\csromanspacer{} โมหํ.\csromanspacer{} มานํ.\csromanspacer{} ทิฏฺฐิํ.\csromanspacer{} ปตฺถนํ อุปนิสฺสาย ปาณํ หนติ ฯเปฯ สํฆํ ภินฺทติ.\csromanspacer{} ราโค ฯเปฯ.\csromanspacer{} ปตฺถนา ราคสฺส ฯเปฯ ปตฺถนาย อุปนิสฺสยปจฺจเยน ปจฺจโย. (1)}

\prose{นีวรณสมฺปยุตฺโต ธมฺโม นีวรณวิปฺปยุตฺตสฺส ธมฺมสฺส อุปนิสฺสยปจฺจเยน ปจฺจโย.\csromanspacer{} \textbf{อนนฺตรูปนิสฺสโย, ปกตูปนิสฺสโย ฯเปฯ.}\csromanspacer{} ปกตูปนิสฺสโย\csromandash{}ราคํ อุปนิสฺสาย ทานํ เทติ ฯเปฯ สีลํ ฯเปฯ อุโปสถกมฺมํ ฯเปฯ.\csromanspacer{} ฌานํ ฯเปฯ.\csromanspacer{} วิปสฺสนํ.\csromanspacer{} มคฺคํ.\csromanspacer{} อภิญฺญํ.\csromanspacer{} สมาปตฺติํ อุปฺปาเทติ.\csromanspacer{} โทสํ ฯเปฯ.\csromanspacer{} ปตฺถนํ อุปนิสฺสาย ทานํ เทติ ฯเปฯ สมาปตฺติํ อุปฺปาเทติ.\csromanspacer{} ราโค ฯเปฯ.\csromanspacer{} ปตฺถนา}

\vspace{\tipitakaparskip}
\csromancenter{นีวรณสมฺปยุตฺตทุก สทฺธาย ฯเปฯ ปญฺญาย.\csromanspacer{} กายิกสฺส สุขสฺส, กายิกสฺส ทุกฺขสฺส, มคฺคสฺส, ผลสมาปตฺติยา อุปนิสฺสยปจฺจเยน ปจฺจโย. (2)}
\proseitem{74}{\csromanfolio{323}นีวรณวิปฺปยุตฺโต ธมฺโม นีวรณวิปฺปยุตฺตสฺส ธมฺมสฺส อุปนิสฺสยปจฺจเยน ปจฺจโย.\csromanspacer{} \textbf{อารมฺมณูปนิสฺสโย, อนนฺตรูปนิสฺสโย, ปกตูปนิสฺสโย ฯเปฯ.}\csromanspacer{} ปกตูปนิสฺสโย\csromandash{}สทฺธํ อุปนิสฺสาย ทานํ เทติ.\csromanspacer{} สีลํ ฯเปฯ มคฺคํ ฯเปฯ อภิญฺญํ ฯเปฯ สมาปตฺติํ อุปฺปาเทติ.\csromanspacer{} สีลํ ฯเปฯ.\csromanspacer{} ปญฺญํ ฯเปฯ.\csromanspacer{} เสนาสนํ อุปนิสฺสาย ทานํ เทติ ฯเปฯ สมาปตฺติํ อุปฺปาเทติ.\csromanspacer{} สทฺธา ฯเปฯ.\csromanspacer{} เสนาสนํ สทฺธาย ฯเปฯ ปญฺญาย ฯเปฯ มคฺคสฺส ผลสมาปตฺติยา อุปนิสฺสยปจฺจเยน ปจฺจโย. (1)}

\prose{นีวรณวิปฺปยุตฺโต ธมฺโม นีวรณสมฺปยุตฺตสฺส ธมฺมสฺส อุปนิสฺสยปจฺจเยน ปจฺจโย.\csromanspacer{} \textbf{อารมฺมณูปนิสฺสโย, อนนฺตรูปนิสฺสโย, ปกตูปนิสฺสโย ฯเปฯ.}\csromanspacer{} ปกตูปนิสฺสโย\csromandash{}สทฺธํ อุปนิสฺสาย มานํ ชปฺเปติ.\csromanspacer{} ทิฏฺฐิํ คณฺหาติ.\csromanspacer{} สีลํ ฯเปฯ.\csromanspacer{} ปญฺญํ.\csromanspacer{} กายิกํ สุขํ ฯเปฯ.\csromanspacer{} เสนาสนํ อุปนิสฺสาย ปาณํ หนติ ฯเปฯ.\csromanspacer{} สํฆํ ภินฺทติ.\csromanspacer{} สทฺธา ฯเปฯ.\csromanspacer{} เสนาสนํ ราคสฺส.\csromanspacer{} โทสสฺส.\csromanspacer{} โมหสฺส.\csromanspacer{} มานสฺส.\csromanspacer{} ทิฏฺฐิยา.\csromanspacer{} ปตฺถนาย อุปนิสฺสยปจฺจเยน ปจฺจโย. (2)}

\csromanheader{2}{\textbf{ปุเรชาต}}
\proseitem{75}{นีวรณวิปฺปยุตฺโต ธมฺโม นีวรณวิปฺปยุตฺตสฺส ธมฺมสฺส ปุเรชาตปจฺจเยน ปจฺจโย.\csromanspacer{} \textbf{อารมฺมณปุเรชาตํ}, วตฺถุปุเรชาตํ.\csromanspacer{}\csromanspacer{}\textbf{อารมฺมณปุเรชาตํ}\csromandash{}จกฺขุํ ฯเปฯ วตฺถุํ อนิจฺจโต ทุกฺขโต อนตฺตโต วิปฺปสฺสติ.\csromanspacer{} ทิพฺเพน จกฺขุนา รูปํ ปสฺสติ.\csromanspacer{} ทิพฺพาย โสตธาตุยา สทฺทํ สุณาติ.\csromanspacer{} รูปายตนํ จกฺขุวิญฺญาณสฺส ฯเปฯ.\csromanspacer{} โผฏฺฐพฺพายตนํ กายวิญฺญาณสฺส ฯเปฯ.\csromanspacer{} \textbf{วตฺถุปุเรชาตํ}\csromandash{}จกฺขายตนํ จกฺขุวิญฺญาณสฺส ฯเปฯ.\csromanspacer{} กายายตนํ กายวิญฺญาณสฺส ฯเปฯ.\csromanspacer{} วตฺถุ นีวรณวิปฺปยุตฺตานํ ขนฺธานํ ปุเรชาตปจฺจเยน ปจฺจโย. (1)}

\prose{นีวรณวิปฺปยุตฺโต ธมฺโม นีวรณสมฺปยุตฺตสฺส ธมฺมสฺส ปุเรชาตปจฺจเยน ปจฺจโย.\csromanspacer{} \textbf{อารมฺมณปุเรชาตํ}, วตฺถุปุเรชาตํ.\csromanspacer{}\csromanspacer{}\textbf{อารมฺมณปุเรชาตํ}\csromandash{}จกฺขุํ ฯเปฯ วตฺถุํ อสฺสาเทติ อภินนฺทติ, ตํ อารพฺภ ราโค.\csromanspacer{} โทโส.\csromanspacer{} โมโห.\csromanspacer{} \textbf{วตฺถุปุเรชาตํ}\csromandash{}วตฺถุ นีวรณสมฺปยุตฺตานํ ขนฺธานํ ปุเรชาตปจฺจเยน ปจฺจโย. (2)}

\csromanheader{2}{\textbf{ปจฺฉาชาตาเสวน}}
\proseitem{76}{\csromanfolio{324}นีวรณสมฺปยุตฺโต ธมฺโม นีวรณวิปฺปยุตฺตสฺส ธมฺมสฺส ปจฺฉาชาตปจฺจเยน ปจฺจโย.\csromanspacer{} เทฺว.\csromanspacer{} อาเสวนปจฺจเยน ปจฺจโย.\csromanspacer{} เทฺว.}

\csromanheader{2}{\textbf{กมฺมาทิ}}
\proseitem{77}{นีวรณสมฺปยุตฺโต ธมฺโม นีวรณสมฺปยุตฺตสฺส ธมฺมสฺส กมฺมปจฺจเยน ปจฺจโย.\csromanspacer{} นีวรณสมฺปยุตฺตา เจตนา สมฺปยุตฺตกานํ ขนฺธานํ กมฺมปจฺจเยน ปจฺจโย. (มูลํ กาตพฺพํ.) \textbf{สหชาตา}, นานากฺขณิกา.\csromanspacer{}\csromanspacer{}\textbf{สหชาตา} นีวรณสมฺปยุตฺตา เจตนา จิตฺตสมุฏฺฐานานํ รูปานํ กมฺมปจฺจเยน ปจฺจโย.\csromanspacer{} \textbf{นานากฺขณิกา} นีวรณสมฺปยุตฺตา เจตนา วิปากานํ ขนฺธานํ กฏตฺตา จ รูปานํ กมฺมปจฺจเยน ปจฺจโย. (มูลํ กาตพฺพํ.) นีวรณสมฺปยุตฺตา เจตนา สมฺปยุตฺตกานํ ขนฺธานํ จิตฺตสมุฏฺฐานานญฺจ รูปานํ กมฺมปจฺจเยน ปจฺจโย. (3)}

\prose{นีวรณวิปฺปยุตฺโต ธมฺโม นีวรณวิปฺปยุตฺตสฺส ธมฺมสฺส กมฺมปจฺจเยน ปจฺจโย.\csromanspacer{} \textbf{สหชาตา}, นานากฺขณิกา. (สํขิตฺตํ.) วิปากปจฺจยา เอกํ.}

\csromanheader{2}{\textbf{อาหาราทิ}}
\proseitem{78}{นีวรณสมฺปยุตฺโต ธมฺโม นีวรณสมฺปยุตฺตสฺส ธมฺมสฺส อาหารปจฺจเยน ปจฺจโย.\csromanspacer{} อินฺทฺริยปจฺจเยน ปจฺจโย.\csromanspacer{} ฌานปจฺจเยน ปจฺจโย.\csromanspacer{} มคฺคปจฺจเยน ปจฺจโย.\csromanspacer{} สมฺปยุตฺตปจฺจเยน ปจฺจโย.\csromanspacer{} เทฺว.}

\csromanheader{2}{\textbf{วิปฺปยุตฺต}}
\proseitem{79}{นีวรณสมฺปยุตฺโต ธมฺโม นีวรณวิปฺปยุตฺตสฺส ธมฺมสฺส วิปฺปยุตฺตปจฺจเยน ปจฺจโย.\csromanspacer{} สหชาตํ, ปจฺฉาชาตํ. (สํขิตฺตํ.) (1)}

\prose{นีวรณวิปฺปยุตฺโต ธมฺโม นีวรณวิปฺปยุตฺตสฺส ธมฺมสฺส วิปฺปยุตฺตปจฺจเยน ปจฺจโย.\csromanspacer{} สหชาตํ, ปุเรชาตํ, ปจฺฉาชาตํ. (สํขิตฺตํ.) (1)}

\prose{นีวรณวิปฺปยุตฺโต ธมฺโม นีวรณสมฺปยุตฺตสฺส ธมฺมสฺส วิปฺปยุตฺตปจฺจเยน ปจฺจโย.\csromanspacer{} \textbf{ปุเรชาตํ} วตฺถุ นีวรณสมฺปยุตฺตกานํ ขนฺธานํ วิปฺปยุตฺตปจฺจเยน ปจฺจโย. (2)}

\csromanheader{2}{46. นีวรณสมฺปยุตฺตทุก \textbf{อตฺถิ}}
\proseitem{80}{\csromanfolio{325}นีวรณสมฺปยุตฺโต ธมฺโม นีวรณสมฺปยุตฺตสฺส ธมฺมสฺส อตฺถิปจฺจเยน ปจฺจโย.\csromanspacer{} นีวรณสมฺปยุตฺโต เอโก ขนฺโธ ติณฺณนฺนํ ขนฺธานํ อตฺถิปจฺจเยน ปจฺจโย ฯเปฯ. (1)}

\prose{นีวรณสมฺปยุตฺโต ธมฺโม นีวรณวิปฺปยุตฺตสฺส ธมฺมสฺส อตฺถิปจฺจเยน ปจฺจโย.\csromanspacer{} สหชาตํ, ปจฺฉาชาตํ.\csromanspacer{}\csromanspacer{}\textbf{สหชาตา} นีวรณสมฺปยุตฺตา ขนฺธา จิตฺตสมุฏฺฐานานญฺจ รูปานํ อตฺถิปจฺจเยน ปจฺจโย.\csromanspacer{} \textbf{ปจฺฉาชาตา} นีวรณสมฺปยุตฺตา ขนฺธา ปุเรชาตสฺส อิมสฺส กายสฺส อตฺถิปจฺจเยน ปจฺจโย. (2)}

\prose{นีวรณสมฺปยุตฺโต ธมฺโม นีวรณสมฺปยุตฺตสฺส จ นีวรณวิปฺปยุตฺตสฺส จ ธมฺมสฺส อตฺถิปจฺจเยน ปจฺจโย.\csromanspacer{} นีวรณสมฺปยุตฺโต เอโก ขนฺโธ ติณฺณนฺนํ ขนฺธานํ จิตฺตสมุฏฺฐานานญฺจ รูปานํ อตฺถิปจฺจเยน ปจฺจโย ฯเปฯ เทฺว ขนฺธา ทฺวินฺนํ ขนฺธานํ จิตฺตสมุฏฺฐานานญฺจ รูปานํ อตฺถิปจฺจเยน ปจฺจโย. (3)}

\proseitem{81}{นีวรณวิปฺปยุตฺโต ธมฺโม นีวรณวิปฺปยุตฺตสฺส ธมฺมสฺส อตฺถิปจฺจเยน ปจฺจโย.\csromanspacer{} สหชาตํ, ปุเรชาตํ, ปจฺฉาชาตํ, อาหารํ, อินฺทฺริยํ. (สํขิตฺตํ.) (1)}

\prose{นีวรณวิปฺปยุตฺโต ธมฺโม นีวรณสมฺปยุตฺตสฺส ธมฺมสฺส อตฺถิปจฺจเยน ปจฺจโย.\csromanspacer{} \textbf{ปุเรชาตํ} จกฺขุํ ฯเปฯ วตฺถุํ อสฺสาเทติ อภินนฺทติ, ตํ อารพฺภ ราโค ฯเปฯ ทิฏฺฐิ ฯเปฯ วิจิกิจฺฉา ฯเปฯ อุทฺธจฺจํ ฯเปฯ โทมนสฺสํ อุปฺปชฺชติ.\csromanspacer{} วตฺถุ นีวรณสมฺปยุตฺตกานํ ขนฺธานํ อตฺถิปจฺจเยน ปจฺจโย. (2)}

\proseitem{82}{นีวรณสมฺปยุตฺโต จ นีวรณวิปฺปยุตฺโต จ ธมฺมา นีวรณสมฺปยุตฺตสฺส ธมฺมสฺส อตฺถิปจฺจเยน ปจฺจโย.\csromanspacer{} สหชาตํ, ปุเรชาตํ.\csromanspacer{}\csromanspacer{}สหชาโต นีวรณสมฺปยุตฺโต เอโก ขนฺโธ จ วตฺถุ จ ติณฺณนฺนํ ขนฺธานํ อตฺถิปจฺจเยน ปจฺจโย ฯเปฯ เทฺว ขนฺธา จ วตฺถุ จ ทฺวินฺนํ ขนฺธานํ อตฺถิปจฺจเยน ปจฺจโย. (1)}

\prose{นีวรณสมฺปยุตฺโต จ นีวรณวิปฺปยุตฺโต จ ธมฺมา นีวรณวิปฺปยุตฺตสฺส ธมฺมสฺส อตฺถิปจฺจเยน ปจฺจโย.\csromanspacer{} สหชาตํ, ปจฺฉาชาตํ, อาหารํ, อินฺทฺริยํ.\csromanspacer{}\csromanspacer{}\textbf{สหชาตา} นีวรณสมฺปยุตฺตา ขนฺธา จ มหาภูตา จ \csromanfolio{326}จิตฺตสมุฏฺฐานานํ รูปานํ อตฺถิปจฺจเยน ปจฺจโย.\csromanspacer{} \textbf{ปจฺฉาชาตา} นีวรณสมฺปยุตฺตา ขนฺธา จ กพฬีกาโร อาหาโร จ อิมสฺส กายสฺส อตฺถิปจฺจเยน ปจฺจโย.\csromanspacer{} \textbf{ปจฺฉาชาตา} นีวรณสมฺปยุตฺตา ขนฺธา จ รูปาชีวิตินฺทฺริยญฺจ กฏตฺตารูปานํ อตฺถิปจฺจเยน ปจฺจโย. (2)}

\csromanheader{2}{1. ปจฺจยานุโลม 2. สงฺขฺยาวาร}
\csromanheader{2}{\textbf{สุทฺธ}}
\proseitem{83}{เหตุยา จตฺตาริ, อารมฺมเณ จตฺตาริ, อธิปติยา ปญฺจ, อนนฺตเร จตฺตาริ, สมนนฺตเร จตฺตาริ, สหชาเต ปญฺจ, อญฺญมญฺเญ เทฺว, นิสฺสเย สตฺต, อุปนิสฺสเย จตฺตาริ, ปุเรชาเต เทฺว, ปจฺฉาชาเต เทฺว, อาเสวเน เทฺว, กมฺเม จตฺตาริ, วิปาเก เอกํ, อาหาเร จตฺตาริ, อินฺทฺริเย จตฺตาริ, ฌาเน จตฺตาริ, มคฺเค จตฺตาริ, สมฺปยุตฺเต เทฺว, วิปฺปยุตฺเต ตีณิ, อตฺถิยา สตฺต, นตฺถิยา จตฺตาริ, วิคเต จตฺตาริ, อวิคเต สตฺต.}

\vspace{\tipitakaparskip}
\csromancenter{อนุโลมํ.}
\csromanheader{2}{2. ปจฺจนียุทฺธาร}
\proseitem{84}{นีวรณสมฺปยุตฺโต ธมฺโม นีวรณสมฺปยุตฺตสฺส ธมฺมสฺส อารมฺมณปจฺจเยน ปจฺจโย.\csromanspacer{} สหชาตปจฺจเยน ปจฺจโย.\csromanspacer{} อุปนิสฺสยปจฺจเยน ปจฺจโย. (1)}

\prose{นีวรณสมฺปยุตฺโต ธมฺโม นีวรณวิปฺปยุตฺตสฺส ธมฺมสฺส อารมฺมณปจฺจเยน ปจฺจโย.\csromanspacer{} สหชาตปจฺจเยน ปจฺจโย.\csromanspacer{} อุปนิสฺสยปจฺจเยน ปจฺจโย.\csromanspacer{} ปจฺฉาชาตปจฺจเยน ปจฺจโย.\csromanspacer{} กมฺมปจฺจเยน ปจฺจโย. (2)}

\prose{นีวรณสมฺปยุตฺโต ธมฺโม นีวรณสมฺปยุตฺตสฺส จ นีวรณวิปฺปยุตฺตสฺส จ ธมฺมสฺส สหชาตปจฺจเยน ปจฺจโย. (3)}

\proseitem{85}{นีวรณวิปฺปยุตฺโต ธมฺโม นีวรณวิปฺปยุตฺตสฺส ธมฺมสฺส อารมฺมณปจฺจเยน ปจฺจโย.\csromanspacer{} สหชาตปจฺจเยน ปจฺจโย.\csromanspacer{} อุปนิสฺสยปจฺจเยน ปจฺจโย.\csromanspacer{} ปุเรชาตปจฺจเยน ปจฺจโย.\csromanspacer{} ปจฺฉาชาตปจฺจเยน ปจฺจโย.}

\vspace{\tipitakaparskip}
\csromancenter{นีวรณสมฺปยุตฺตทุก กมฺมปจฺจเยน ปจฺจโย.\csromanspacer{} อาหารปจฺจเยน ปจฺจโย.\csromanspacer{} อินฺทฺริยปจฺจเยน ปจฺจโย. (1)}
\prose{\csromanfolio{327}นีวรณวิปฺปยุตฺโต ธมฺโม นีวรณสมฺปยุตฺตสฺส ธมฺมสฺส อารมฺมณปจฺจเยน ปจฺจโย.\csromanspacer{} อุปนิสฺสยปจฺจเยน ปจฺจโย.\csromanspacer{} ปุเรชาตปจฺจเยน ปจฺจโย. (2)}

\prose{นีวรณสมฺปยุตฺโต จ นีวรณวิปฺปยุตฺโต จ ธมฺมา นีวรณสมฺปยุตฺตสฺส ธมฺมสฺส สหชาตํ, ปุเรชาตํ. (1)}

\prose{นีวรณสมฺปยุตฺโต จ นีวรณวิปฺปยุตฺโต จ ธมฺมา นีวรณวิปฺปยุตฺตสฺส ธมฺมสฺส สหชาตํ, ปจฺฉาชาตํ, อาหารํ, อินฺทฺริยํ. (2)}

\csromanheader{2}{2. ปจฺจยปจฺจนีย 2. สงฺขฺยาวาร}
\csromanheader{2}{\textbf{สุทฺธ}}
\proseitem{86}{นเหตุยา สตฺต, น-อารมฺมเณ สตฺต, น-อธิปติยา สตฺต, น-อนนฺตเร สตฺต, นสมนนฺตเร สตฺต, นสหชาเต ปญฺจ, น-อญฺญมญฺเญ ปญฺจ, นนิสฺสเย ปญฺจ, น-อุปนิสฺสเย สตฺต, นปุเรชาเต ฉ ฯเปฯ นมคฺเค สตฺต, นสมฺปยุตฺเต ปญฺจ, นวิปฺปยุตฺเต จตฺตาริ, โน-อตฺถิยา จตฺตาริ, โนนตฺถิยา สตฺต, โนวิคเต สตฺต, โน-อวิคเต จตฺตาริ.}

\vspace{\tipitakaparskip}
\csromancenter{ปจฺจนียํ.}
\csromanheader{2}{3. ปจฺจยานุโลมปจฺจนีย}
\csromanheader{2}{\textbf{เหตุทุก}}
\proseitem{87}{เหตุปจฺจยา น-อารมฺมเณ จตฺตาริ, น-อธิปติยา จตฺตาริ, น-อนนฺตเร จตฺตาริ, นสมนนฺตเร จตฺตาริ, น-อญฺญมญฺเญ เทฺว, น-อุปนิสฺสเย จตฺตาริ ฯเปฯ นมคฺเค จตฺตาริ, นสมฺปยุตฺเต เทฺว, นวิปฺปยุตฺเต เทฺว, โนนตฺถิยา จตฺตาริ, โนวิคเต จตฺตาริ.}

\csromanheader{2}{4. ปจฺจยปจฺจนียานุโลม}
\proseitem{88}{\csromanfolio{328}นเหตุปจฺจยา อารมฺมเณ จตฺตาริ, อธิปติยา ปญฺจ, (อนุโลมคณนา) ฯเปฯ อวิคเต สตฺต.}

\nitthitam{นีวรณสมฺปยุตฺตทุกํ นิฏฺฐิตํ.}
\csromanmatikaanchor{mtk.44}
\csromantocmarkref{subsection}{47. นีวรณนีวรณิยทุก}{mtk.44}
\bookmark[dest={mtk.44},level=2]{47. นีวรณนีวรณิยทุก}
\csromanheader{2}{47. นีวรณนีวรณิยทุก 1. ปฏิจฺจวาร}
\csromanheader{2}{1. ปจฺจยานุโลมาทิ}
\proseitem{89}{นีวรณญฺเจว นีวรณิยญฺจ ธมฺมํ ปฏิจฺจ นีวรโณ เจว นีวรณิโย จ ธมฺโม อุปฺปชฺชติ เหตุปจฺจยา.\csromanspacer{} กามจฺฉนฺทนีวรณํ ปฏิจฺจ ถินมิทฺธนีวรณํ, อุทฺธจฺจนีวรณํ, อวิชฺชานีวรณํ. (เอวํ สพฺเพ คณนา วิภชิตพฺพา.\csromanspacer{} นีวรณทุกสทิสํ.\csromanspacer{} นินฺนานากรณํ.)}

\csromanheader{2}{47. นีวรณนีวรณิยทุก 7. ปญฺหาวาร}
\csromanheader{2}{1. ปจฺจยานุโลม 1. วิภงฺควาร}
\csromanheader{2}{\textbf{เหตุ}}
\proseitem{90}{นีวรโณ เจว นีวรณิโย จ ธมฺโม นีวรณสฺส เจว นีวรณิยสฺส จ ธมฺมสฺส เหตุปจฺจเยน ปจฺจโย.\csromanspacer{} นีวรณา เจว นีวรณิยา จ เหตู สมฺปยุตฺตกานํ นีวรณานํ เหตุปจฺจเยน ปจฺจโย. (1)}

\prose{นีวรโณ เจว นีวรณิโย จ ธมฺโม นีวรณิยสฺส เจว โน จ นีวรณสฺส ธมฺมสฺส เหตุปจฺจเยน ปจฺจโย.\csromanspacer{} นีวรณา เจว นีวรณิยา จ เหตู สมฺปยุตฺตกานํ ขนฺธานํ จิตฺตสมุฏฺฐานานญฺจ รูปานํ เหตุปจฺจเยน ปจฺจโย. (2)}

\prose{นีวรโณ เจว นีวรณิโย จ ธมฺโม นีวรณสฺส เจว นีวรณิยสฺส จ นีวรณิยสฺส เจว โน จ นีวรณสฺส จ ธมฺมสฺส เหตุปจฺจเยน ปจฺจโย.}

\vspace{\tipitakaparskip}
\csromancenter{นีวรณนีวรณิยทุก นีวรณา เจว นีวรณิยา จ เหตู สมฺปยุตฺตกานํ ขนฺธานํ นีวรณานญฺจ จิตฺตสมุฏฺฐานานญฺจ รูปานํ เหตุปจฺจเยน ปจฺจโย. (3)}
\prose{\csromanfolio{329}นีวรณิโย เจว โน จ นีวรโณ ธมฺโม นีวรณิยสฺส เจว โน จ นีวรณสฺส ธมฺมสฺส เหตุปจฺจเยน ปจฺจโย.\csromanspacer{} นีวรณิยา เจว โน จ นีวรณา เหตู สมฺปยุตฺตกานํ ขนฺธานํ จิตฺตสมุฏฺฐานานญฺจ รูปานํ เหตุปจฺจเยน ปจฺจโย.\csromanspacer{} ปฏิสนฺธิกฺขเณ ฯเปฯ. (1)}

\csromanheader{2}{\textbf{อารมฺมณ}}
\proseitem{91}{นีวรโณ เจว นีวรณิโย จ ธมฺโม นีวรณสฺส เจว นีวรณิยสฺส จ ธมฺมสฺส อารมฺมณปจฺจเยน ปจฺจโย.\csromanspacer{} นีวรเณ อารพฺภ นีวรณา อุปฺปชฺชนฺติ. (มูลํ ปุจฺฉิตพฺพํ.) นีวรเณ อารพฺภ นีวรณิยา เจว โน จ นีวรณา ขนฺธา อุปฺปชฺชนฺติ. (มูลํ ปุจฺฉิตพฺพํ.) นีวรเณ อารพฺภ นีวรณา จ สมฺปยุตฺตกา จ ขนฺธา อุปฺปชฺชนฺติ. (3)}

\prose{นีวรณิโย เจว โน จ นีวรโณ ธมฺโม นีวรณิยสฺส เจว โน จ นีวรณสฺส ธมฺมสฺส อารมฺมณปจฺจเยน ปจฺจโย.\csromanspacer{} ทานํ ฯเปฯ สีลํ ฯเปฯ อุโปสถกมฺมํ กตฺวา ตํ ปจฺจเวกฺขติ, อสฺสาเทติ อภินนฺทติ, ตํ อารพฺภ ราโค อุปฺปชฺชติ.\csromanspacer{} ทิฏฺฐิ.\csromanspacer{} วิจิกิจฺฉา.\csromanspacer{} อุทฺธจฺจํ.\csromanspacer{} โทมนสฺสํ อุปฺปชฺชติ.\csromanspacer{} ปุพฺเพ สุจิณฺณานิ ฯเปฯ.\csromanspacer{} ฌานา ฯเปฯ.\csromanspacer{} อริยา โคตฺรภุํ ปจฺจเวกฺขนฺติ.\csromanspacer{} โวทานํ ฯเปฯ.\csromanspacer{} ปหีเน กิเลเส ฯเปฯ.\csromanspacer{} วิกฺขมฺภิเต กิเลเส ฯเปฯ.\csromanspacer{} ปุพฺเพ สมุทาจิณฺเณ ฯเปฯ.\csromanspacer{} จกฺขุํ ฯเปฯ วตฺถุํ นีวรณิเย เจว โน จ นีวรเณ ขนฺเธ อนิจฺจโต ฯเปฯ วิปสฺสติ, อสฺสาเทติ อภินนฺทติ ฯเปฯ โทมนสฺสํ อุปฺปชฺชติ.\csromanspacer{} ทิพฺเพน จกฺขุนา รูปํ ปสฺสติ ฯเปฯ. (ยาว อาวชฺชนา, ตาว กาตพฺพา.) (1)}

\prose{นีวรณิโย เจว โน จ นีวรโณ ธมฺโม นีวรณสฺส เจว นีวรณิยสฺส จ ธมฺมสฺส อารมฺมณปจฺจเยน ปจฺจโย.\csromanspacer{} ทานํ ฯเปฯ สีลํ ฯเปฯ อุโปสถกมฺมํ ฯเปฯ.\csromanspacer{} ปุพฺเพ สุจิณฺณานิ ฯเปฯ.\csromanspacer{} ฌานา ฯเปฯ.\csromanspacer{} จกฺขุํ ฯเปฯ วตฺถุํ นีวรณิเย เจว โน จ นีวรเณ ขนฺเธ อสฺสาเทติ อภินนฺทติ, ตํ อารพฺภ ราโค ฯเปฯ ทิฏฺฐิ.\csromanspacer{} วิจิกิจฺฉา.\csromanspacer{} อุทฺธจฺจํ.\csromanspacer{} โทมนสฺสํ อุปฺปชฺชติ. (2)}

\prose{นีวรณิโย เจว โน จ นีวรโณ ธมฺโม นีวรณสฺส เจว นีวรณิยสฺส จ นีวรณิยสฺส เจว โน จ นีวรณสฺส จ ธมฺมสฺส อารมฺมณปจฺจเยน ปจฺจโย.\csromanspacer{} ทานํ ฯเปฯ สีลํ ฯเปฯ อุโปสถกมฺมํ ฯเปฯ.\csromanspacer{} ปุพฺเพ สุจิณฺณานิ ฯเปฯ. \csromanfolio{330}ฌานา ฯเปฯ.\csromanspacer{} จกฺขุํ ฯเปฯ วตฺถุํ นีวรณิเย เจว โน จ นีวรเณ ขนฺเธ อสฺสาเทติ อภินนฺทติ, ตํ อารพฺภ นีวรณา จ สมฺปยุตฺตกา จ ขนฺธา อุปฺปชฺชนฺติ. (เอวํ อิตเรปิ ตีณิ กาตพฺพา.) (3)}

\prose{(อารมฺมณสทิสา.\csromanspacer{} อธิปติปจฺจยา.\csromanspacer{} ปุเรชาตมฺปิ อารมฺมณสทิสํ.\csromanspacer{} อุปนิสฺสเยปิ โลกุตฺตรํ น กาตพฺพํ.\csromanspacer{} สํขิตฺตํ.\csromanspacer{} เอวํ วิตฺถาเรตพฺพํ, ยถา นีวรณทุกํ, เอวํ ปจฺจเวกฺขิตฺวา กาตพฺพํ.)}

\nitthitam{นีวรณนีวรณิยทุกํ นิฏฺฐิตํ.}
\csromanmatikaanchor{mtk.45}
\csromantocmarkref{subsection}{48. นีวรณนีวรณสมฺปยุตฺตทุก}{mtk.45}
\bookmark[dest={mtk.45},level=2]{48. นีวรณนีวรณสมฺปยุตฺตทุก}
\csromanheader{2}{48. นีวรณนีวรณสมฺปยุตฺตทุก 1. ปฏิจฺจวาร}
\csromanheader{2}{1. ปจฺจยานุโลม 1. วิภงฺควาร}
\csromanheader{2}{\textbf{เหตุ}}
\proseitem{92}{นีวรณญฺเจว นีวรณสมฺปยุตฺตญฺจ ธมฺมํ ปฏิจฺจ นีวรโณ เจว นีวรณสมฺปยุตฺโต จ ธมฺโม อุปฺปชฺชติ เหตุปจฺจยา.\csromanspacer{} กามจฺฉนฺทนีวรณํ ปฏิจฺจ ถินมิทฺธนีวรณํ, อุทฺธจฺจนีวรณํ, อวิชฺชานีวรณํ. (จกฺกํ.\csromanspacer{} สพฺเพปิ นีวรณา กาตพฺพา.) (1)}

\prose{นีวรณญฺเจว นีวรณสมฺปยุตฺตญฺจ ธมฺมํ ปฏิจฺจ นีวรณสมฺปยุตฺโต เจว โน จ นีวรโณ ธมฺโม อุปฺปชฺชติ เหตุปจฺจยา.\csromanspacer{} นีวรเณ ปฏิจฺจ สมฺปยุตฺตกา ขนฺธา. (2)}

\prose{นีวรณญฺเจว นีวรณสมฺปยุตฺตญฺจ ธมฺมํ ปฏิจฺจ นีวรโณ เจว นีวรณสมฺปยุตฺโต จ นีวรณสมฺปยุตฺโต เจว โน จ นีวรโณ จ ธมฺมา อุปฺปชฺชนฺติ เหตุปจฺจยา.\csromanspacer{} กามจฺฉนฺทนีวรณํ ปฏิจฺจ ถินมิทฺธนีวรณํ, อุทฺธจฺจนีวรณํ, อวิชฺชานีวรณํ สมฺปยุตฺตกา จ ขนฺธา. (จกฺกํ.) (3)}

\proseitem{93}{นีวรณสมฺปยุตฺตญฺเจว โน จ นีวรณํ ธมฺมํ ปฏิจฺจ นีวรณสมฺปยุตฺโต เจว โน จ นีวรโณ ธมฺโม อุปฺปชฺชติ เหตุปจฺจยา.\csromanspacer{} นีวรณสมฺปยุตฺตญฺเจว โน จ นีวรณํ เอกํ ขนฺธํ ปฏิจฺจ ตโย ขนฺธา ฯเปฯ เทฺว ขนฺเธ ฯเปฯ. (1)}

\csromanheader{2}{48. นีวรณนีวรณสมฺปยุตฺตทุก}
\prose{\csromanfolio{331}นีวรณสมฺปยุตฺตญฺเจว โน จ นีวรณํ ธมฺมํ ปฏิจฺจ นีวรโณ เจว นีวรณสมฺปยุตฺโต จ ธมฺโม อุปฺปชฺชติ เหตุปจฺจยา.\csromanspacer{} นีวรณสมฺปยุตฺเต เจว โน จ นีวรเณ ขนฺเธ ปฏิจฺจ นีวรณา. (2)}

\prose{นีวรณสมฺปยุตฺตญฺเจว โน จ นีวรณํ ธมฺมํ ปฏิจฺจ นีวรโณ เจว นีวรณสมฺปยุตฺโต จ นีวรณสมฺปยุตฺโต เจว โน จ นีวรโณ จ ธมฺมา อุปฺปชฺชนฺติ เหตุปจฺจยา.\csromanspacer{} นีวรณสมฺปยุตฺตญฺเจว โน จ นีวรณํ เอกํ ขนฺธํ ปฏิจฺจ ตโย ขนฺธา นีวรณา จ ฯเปฯ เทฺว ขนฺเธ ฯเปฯ. (3)}

\proseitem{94}{นีวรณญฺเจว นีวรณสมฺปยุตฺตญฺจ นีวรณสมฺปยุตฺตญฺเจว โน จ นีวรณญฺจ ธมฺมํ ปฏิจฺจ นีวรโณ เจว นีวรณสมฺปยุตฺโต จ ธมฺโม อุปฺปชฺชติ เหตุปจฺจยา.\csromanspacer{} กามจฺฉนฺทนีวรณญฺจ สมฺปยุตฺตเก จ ขนฺเธ ปฏิจฺจ ถินมิทฺธนีวรณํ, อุทฺธจฺจนีวรณํ, อวิชฺชานีวรณํ. (จกฺกํ.) (1)}

\prose{นีวรณญฺเจว นีวรณสมฺปยุตฺตญฺจ นีวรณสมฺปยุตฺตญฺเจว โน จ นีวรณญฺจ ธมฺมํ ปฏิจฺจ นีวรณสมฺปยุตฺโต เจว โน จ นีวรโณ ธมฺโม อุปฺปชฺชติ เหตุปจฺจยา.\csromanspacer{} นีวรณสมฺปยุตฺตญฺเจว โน จ นีวรณํ เอกํ ขนฺธญฺจ นีวรณญฺจ ปฏิจฺจ ตโย ขนฺธา ฯเปฯ เทฺว ขนฺเธ จ ฯเปฯ. (2)}

\prose{นีวรณญฺเจว นีวรณสมฺปยุตฺตญฺจ นีวรณสมฺปยุตฺตญฺเจว โน จ นีวรณญฺจ ธมฺมํ ปฏิจฺจ นีวรโณ เจว นีวรณสมฺปยุตฺโต จ นีวรณสมฺปยุตฺโต เจว โน จ นีวรโณ จ ธมฺมา อุปฺปชฺชนฺติ เหตุปจฺจยา.\csromanspacer{} นีวรณสมฺปยุตฺตญฺเจว โน จ นีวรณํ เอกํ ขนฺธญฺจ กามจฺฉนฺทนีวรณญฺจ ปฏิจฺจ ตโย ขนฺธา ถินมิทฺธนีวรณํ, อุทฺธจฺจนีวรณํ, อวิชฺชานีวรณํ ฯเปฯ เทฺว ขนฺเธ จ ฯเปฯ. (จกฺกํ.\csromanspacer{} สํขิตฺตํ.) (3)}

\csromanheader{2}{1. ปจฺจยานุโลม 2. สงฺขฺยาวาร}
\proseitem{95}{เหตุยา นว, อารมฺมเณ นว, อธิปติยา นว, (สพฺพตฺถ นว.) กมฺเม นว, อาหาเร นว ฯเปฯ อวิคเต นว.}

\vspace{\tipitakaparskip}
\csromancenter{อนุโลมํ.}
\csromanheader{2}{2. ปจฺจยปจฺจนีย 1. วิภงฺควาร}
\csromanheader{2}{\textbf{นเหตุ}}
\proseitem{96}{\csromanfolio{332}นีวรณญฺเจว นีวรณสมฺปยุตฺตญฺจ ธมฺมํ ปฏิจฺจ นีวรโณ เจว นีวรณสมฺปยุตฺโต จ ธมฺโม อุปฺปชฺชติ นเหตุปจฺจยา.\csromanspacer{} วิจิกิจฺฉานีวรณํ ปฏิจฺจ อวิชฺชานีวรณํ.\csromanspacer{} อุทฺธจฺจนีวรณํ ปฏิจฺจ อวิชฺชานีวรณํ. (1)}

\prose{นีวรณสมฺปยุตฺตญฺเจว โน จ นีวรณํ ธมฺมํ ปฏิจฺจ นีวรโณ เจว นีวรณสมฺปยุตฺโต จ ธมฺโม อุปฺปชฺชติ นเหตุปจฺจยา.\csromanspacer{} วิจิกิจฺฉาสหคเต อุทฺธจฺจสหคเต ขนฺเธ ปฏิจฺจ อวิชฺชานีวรณํ. (1)}

\prose{นีวรณญฺเจว นีวรณสมฺปยุตฺตญฺจ นีวรณสมฺปยุตฺตญฺเจว โน จ นีวรณญฺจ ธมฺมํ ปฏิจฺจ นีวรโณ เจว นีวรณสมฺปยุตฺโต จ ธมฺโม อุปฺปชฺชติ นเหตุปจฺจยา.\csromanspacer{} วิจิกิจฺฉานีวรณญฺจ สมฺปยุตฺตเก จ ขนฺเธ ปฏิจฺจ อวิชฺชานีวรณํ.\csromanspacer{} อุทฺธจฺจนีวรณญฺจ สมฺปยุตฺตเก จ ขนฺเธ ปฏิจฺจ อวิชฺชานีวรณํ (สํขิตฺตํ.) (1)}

\csromanheader{2}{2. ปจฺจยปจฺจนีย 2. สงฺขฺยาวาร}
\csromanheader{2}{\textbf{สุทฺธ}}
\vspace{\tipitakaparskip}
\csromancenter{\textbf{นเหตุ}ยา ตีณิ, น-อธิปติยา นว, นปุเรชาเต นว, นปจฺฉาชาเต นว, น-อาเสวเน นว, นกมฺเม ตีณิ, นวิปาเก นว, นวิปฺปยุตฺเต นว.}
\csromanheader{2}{3. ปจฺจยานุโลมปจฺจนีย}
\vspace{\tipitakaparskip}
\csromancenter{เหตุปจฺจยา น-อธิปติยา นว ฯเปฯ นวิปฺปยุตฺเต นว.}
\csromanheader{2}{4. ปจฺจยปจฺจนียานุโลม}
\vspace{\tipitakaparskip}
\csromancenter{\textbf{นเหตุ}ปจฺจยา อารมฺมเณ ตีณิ, อนนฺตเร ตีณิ, สมนนฺตเร ตีณิ, (สพฺพตฺถ ตีณิ.) มคฺเค ตีณิ ฯเปฯ อวิคเต ตีณิ.}
\prose{(เอวํ สหชาตวาโรปิ ปจฺจยวาโรปิ นิสฺสยวาโรปิ สํสฏฺฐวาโรปิ สมฺปยุตฺตวาโรปิ ปฏิจฺจวารสทิสา, นินฺนานากรณา.)}

\vspace{\tipitakaparskip}
\csromancenter{นีวรณนีวรณสมฺปยุตฺตทุก \textbf{48.}\csromanspacer{}\textbf{ นีวรณนีวรณสมฺปยุตฺตทุก 7.}\csromanspacer{}\textbf{ ปญฺหาวาร}}
\csromanheader{2}{1. ปจฺจยานุโลม 1. วิภงฺควาร}
\csromanheader{2}{\textbf{เหตุ}}
\proseitem{100}{\csromanfolio{333}นีวรโณ เจว นีวรณสมฺปยุตฺโต จ ธมฺโม นีวรณสฺส เจว นีวรณสมฺปยุตฺตสฺส จ ธมฺมสฺส เหตุปจฺจเยน ปจฺจโย.\csromanspacer{} นีวรณา เจว นีวรณสมฺปยุตฺตา จ เหตู สมฺปยุตฺตกานํ นีวรณานํ เหตุปจฺจเยน ปจฺจโย. (1)}

\prose{นีวรโณ เจว นีวรณสมฺปยุตฺโต จ ธมฺโม นีวรณสมฺปยุตฺตสฺส เจว โน จ นีวรณสฺส ธมฺมสฺส เหตุปจฺจเยน ปจฺจโย.\csromanspacer{} นีวรณา เจว นีวรณสมฺปยุตฺตา จ เหตู สมฺปยุตฺตกานํ ขนฺธานํ เหตุปจฺจเยน ปจฺจโย. (2)}

\prose{นีวรโณ เจว นีวรณสมฺปยุตฺโต จ ธมฺโม นีวรณสฺส เจว นีวรณสมฺปยุตฺตสฺส จ นีวรณสมฺปยุตฺตสฺส เจว โน จ นีวรณสฺส จ ธมฺมสฺส เหตุปจฺจเยน ปจฺจโย.\csromanspacer{} นีวรณา เจว นีวรณสมฺปยุตฺตา จ เหตู สมฺปยุตฺตกานํ ขนฺธานํ นีวรณานญฺจ เหตุปจฺจเยน ปจฺจโย. (3)}

\csromanheader{2}{\textbf{อารมฺมณ}}
\proseitem{111}{นีวรโณ เจว นีวรณสมฺปยุตฺโต จ ธมฺโม นีวรณสฺส เจว นีวรณสมฺปยุตฺตสฺส จ ธมฺมสฺส อารมฺมณปจฺจเยน ปจฺจโย.\csromanspacer{} นีวรเณ อารพฺภ นีวรณา อุปฺปชฺชนฺติ. (มูลํ กาตพฺพํ.) นีวรเณ อารพฺภ นีวรณสมฺปยุตฺตา เจว โน จ นีวรณา ขนฺธา อุปฺปชฺชนฺติ. (มูลํ กาตพฺพํ.) นีวรเณ อารพฺภ นีวรณา จ สมฺปยุตฺตกา จ ขนฺธา อุปฺปชฺชนฺติ. (3)}

\prose{นีวรณสมฺปยุตฺโต เจว โน จ นีวรโณ ธมฺโม นีวรณสมฺปยุตฺตสฺส เจว โน จ นีวรณสฺส ธมฺมสฺส อารมฺมณปจฺจเยน ปจฺจโย.\csromanspacer{} นีวรณสมฺปยุตฺเต เจว โน จ นีวรเณ ขนฺเธ อารพฺภ นีวรณสมฺปยุตฺตา เจว โน จ นีวรณา ขนฺธา อุปฺปชฺชนฺติ. (มูลํ ปุจฺฉิตพฺพํ.) นีวรณสมฺปยุตฺเต เจว โน จ นีวรเณ ขนฺเธ อารพฺภ นีวรณา อุปฺปชฺชนฺติ. (มูลํ ปุจฺฉิตพฺพํ.) นีวรณสมฺปยุตฺเต เจว โน จ นีวรเณ ขนฺเธ อารพฺภ นีวรณา จ สมฺปยุตฺตกา จ ขนฺธา อุปฺปชฺชนฺติ. (3)}

\prose{\csromanfolio{334}นีวรโณ เจว นีวรณสมฺปยุตฺโต จ นีวรณสมฺปยุตฺโต เจว โน จ นีวรโณ จ ธมฺมา นีวรณสฺส เจว นีวรณสมฺปยุตฺตสฺส จ ธมฺมสฺส อารมฺมณปจฺจเยน ปจฺจโย.\csromanspacer{} ตีณิ. (3)}

\csromanheader{2}{\textbf{อธิปตฺยาทิ}}
\proseitem{102}{นีวรโณ เจว นีวรณสมฺปยุตฺโต จ ธมฺโม นีวรณสฺส เจว นีวรณสมฺปยุตฺตสฺส จ ธมฺมสฺส อธิปติปจฺจเยน ปจฺจโย.\csromanspacer{} \textbf{อารมฺมณาธิปติ}.\csromanspacer{} ตีณิ. (ครุการมฺมณาเยว.)}

\prose{นีวรณสมฺปยุตฺโต เจว โน จ นีวรโณ ธมฺโม นีวรณสมฺปยุตฺตสฺส เจว โน จ นีวรณสฺส ธมฺมสฺส อธิปติปจฺจเยน ปจฺจโย.\csromanspacer{} \textbf{อารมฺมณาธิปติ}, สหชาตาธิปติ.\csromanspacer{}\csromanspacer{}\textbf{อารมฺมณาธิปติ}\csromandash{}นีวรณสมฺปยุตฺเต เจว โน จ นีวรเณ ขนฺเธ ครุํ กตฺวา นีวรณสมฺปยุตฺตา เจว โน จ นีวรณา ขนฺธา อุปฺปชฺชนฺติ.\csromanspacer{} \textbf{สหชาตาธิปติ}\csromandash{}นีวรณสมฺปยุตฺตา เจว โน จ นีวรณาธิปติ สมฺปยุตฺตกานํ ขนฺธานํ อธิปติปจฺจเยน ปจฺจโย. (มูลํ กาตพฺพํ.) \textbf{อารมฺมณาธิปติ}\csromandash{}นีวรณสมฺปยุตฺเต เจว โน จ นีวรเณ ขนฺเธ ครุํ กตฺวา นีวรณา อุปฺปชฺชนฺติ.\csromanspacer{} \textbf{สหชาตาธิปติ}\csromandash{} นีวรณสมฺปยุตฺตา เจว โน จ นีวรณาธิปติ สมฺปยุตฺตกานํ นีวรณานํ อธิปติปจฺจเยน ปจฺจโย. (มูลํ กาตพฺพํ.) \textbf{อารมฺมณาธิปติ}\csromandash{}นีวารณสมฺปยุตฺเต เจว โน จ นีวรเณ ขนฺเธ ครุํ กตฺวา นีวรณา จ สมฺปยุตฺตกา จ ขนฺธา อุปฺปชฺชนฺติ.\csromanspacer{} \textbf{สหชาตาธิปติ}\csromandash{}นีวรณสมฺปยุตฺตา เจว โน จ นีวรณาธิปติ สมฺปยุตฺตกานํ ขนฺธานํ นีวรณานญฺจ อธิปติปจฺจเยน ปจฺจโย. (3)}

\proseitem{103}{นีวรโณ เจว นีวรณสมฺปยุตฺโต จ นีวรณสมฺปยุตฺโต เจว โน จ นีวรโณ จ ธมฺมา นีวรณสฺส เจว นีวรณสมฺปยุตฺตสฺส จ ธมฺมสฺส อธิปติปจฺจเยน ปจฺจโย.\csromanspacer{} \textbf{อารมฺมณาธิปติ}.\csromanspacer{} ตีณิ. (3)}

\prose{อนนฺตรปจฺจเยน ปจฺจโย. (อาวชฺชนาปิ วุฏฺฐานมฺปิ นตฺถิ, สพฺพตฺถ ปุริมา ปุริมา กาตพฺพา.) สมนนฺตรปจฺจเยน ปจฺจโย.\csromanspacer{} นว.\csromanspacer{} สหชาตปจฺจเยน ปจฺจโย.\csromanspacer{} นว.\csromanspacer{} อญฺญมญฺญปจฺจเยน ปจฺจโย.\csromanspacer{} นว.\csromanspacer{} นิสฺสยปจฺจเยน ปจฺจโย.\csromanspacer{} นว.\csromanspacer{} อุปนิสฺสยปจฺจเยน ปจฺจโย.\csromanspacer{} นว. (อารมฺมณสทิสํ, วิปาโก นตฺถิ.) อาเสวนปจฺจเยน ปจฺจโย.\csromanspacer{} ปญฺจ.}

\csromanheader{2}{48. นีวรณนีวรณสมฺปยุตฺตทุก \textbf{กมฺม}}
\proseitem{104}{\csromanfolio{335}นีวรณสมฺปยุตฺโต เจว โน จ นีวรโณ ธมฺโม นีวรณสมฺปยุตฺตสฺส เจว โน จ นีวรณสฺส ธมฺมสฺส กมฺมปจฺจเยน ปจฺจโย.\csromanspacer{} นีวรณสมฺปยุตฺตา เจว โน จ นีวรณา เจตนา สมฺปยุตฺตกานํ ขนฺธานํ กมฺมปจฺจเยน ปจฺจโย. (มูลํ กาตพฺพํ.) นีวรณสมฺปยุตฺตา เจว โน จ นีวรณา เจตนา สมฺปยุตฺตกานํ นีวรณานํ กมฺมปจฺจเยน ปจฺจโย. (มูลํ กาตพฺพํ.) นีวรณสมฺปยุตฺตา เจว โน จ นีวรณา เจตนา สมฺปยุตฺตกานํ ขนฺธานํ นีวรณานญฺจ กมฺมปจฺจเยน ปจฺจโย. (3)}

\csromanheader{2}{\textbf{อาหาราทิ}}
\proseitem{105}{นีวรณสมฺปยุตฺโต เจว โน จ นีวรโณ ธมฺโม นีวรณสมฺปยุตฺตสฺส เจว โน จ นีวรณสฺส ธมฺมสฺส อาหารปจฺจเยน ปจฺจโย.\csromanspacer{} ตีณิ.\csromanspacer{} อินฺทฺริยปจฺจเยน ปจฺจโย.\csromanspacer{} ตีณิ.\csromanspacer{} ฌานปจฺจเยน ปจฺจโย.\csromanspacer{} ตีณิ.\csromanspacer{} มคฺคปจฺจเยน ปจฺจโย.\csromanspacer{} ตีณิ.\csromanspacer{} สมฺปยุตฺตปจฺจเยน ปจฺจโย.\csromanspacer{} นว.\csromanspacer{} อตฺถิปจฺจเยน ปจฺจโย.\csromanspacer{} นว.\csromanspacer{} นตฺถิปจฺจเยน ปจฺจโย.\csromanspacer{} นว.\csromanspacer{} วิคตปจฺจเยน ปจฺจโย.\csromanspacer{} นว.\csromanspacer{} อวิคตปจฺจเยน ปจฺจโย.\csromanspacer{} นว.}

\csromanheader{2}{1. ปจฺจยานุโลม 2. สงฺขฺยาวาร}
\csromanheader{2}{\textbf{สุทฺธ}}
\proseitem{106}{เหตุยา ตีณิ, อารมฺมเณ นว, อธิปติยา นว, อนนฺตเร นว, สมนนฺตเร นว, สหชาเต นว, อญฺญมญฺเญ นว, นิสฺสเย นว, อุปนิสฺสเย นว, อาเสวเน นว, กมฺเม ตีณิ, อาหาเร ตีณิ, อินฺทฺริเย ตีณิ, ฌาเน ตีณิ, มคฺเค ตีณิ, สมฺปยุตฺเต นว, อตฺถิยา นว, นตฺถิยา นว, วิคเต นว, อวิคเต นว.}

\csromanheader{2}{2. ปจฺจนียุทฺธาร}
\proseitem{107}{นีวรโณ เจว นีวรณสมฺปยุตฺโต จ ธมฺโม นีวรณสฺส เจว นีวรณสมฺปยุตฺตสฺส จ ธมฺมสฺส อารมฺมณปจฺจเยน ปจฺจโย.\csromanspacer{} สหชาตปจฺจเยน ปจฺจโย.\csromanspacer{} อุปนิสฺสยปจฺจเยน ปจฺจโย. (เอวํ นวปิ ตีสุ ปเทสุ ปริวตฺเตตพฺพา.)}

\csromanheader{2}{2. ปจฺจยปจฺจนีย 2. สงฺขฺยาวาร}
\vspace{\tipitakaparskip}
\csromancenter{นเหตุยา นว, น-อารมฺมเณ นว ฯเปฯ โน-อวิคเต นว.}
\csromanheader{2}{3. ปจฺจยานุโลมปจฺจนีย}
\proseitem{109}{\csromanfolio{336}เหตุปจฺจยา น-อารมฺมเณ ตีณิ, น-อธิปติยา ตีณิ, น-อนนฺตเร ตีณิ, นสมนนฺตเร ตีณิ, น-อุปนิสฺสเย ตีณิ ฯเปฯ นมคฺเค ตีณิ, นวิปฺปยุตฺเต ตีณิ, โนนตฺถิยา ตีณิ, โนวิคเต ตีณิ.}

\csromanheader{2}{4. ปจฺจยปจฺจนียานุโลม}
\csromanheader{2}{110. นเหตุปจฺจยา อารมฺมเณ นว, อธิปติยา นว, (อนุโลมมาติกา) ฯเปฯ อวิคเต นว.}
\nitthitam{นีวรณนีวรณสมฺปยุตฺตทุกํ นิฏฺฐิตํ.}
\csromanmatikaanchor{mtk.46}
\csromantocmarkref{subsection}{49. นีวรณวิปฺปยุตฺตนีวรณิยทุก}{mtk.46}
\bookmark[dest={mtk.46},level=2]{49. นีวรณวิปฺปยุตฺตนีวรณิยทุก}
\csromanheader{2}{49. นีวรณวิปฺปยุตฺตนีวรณิยทุก 1. ปฏิจฺจวาร}
\csromanheader{2}{1. ปจฺจยานุโลมาทิ}
\proseitem{111}{นีวรณวิปฺปยุตฺตํ นีวรณิยํ ธมฺมํ ปฏิจฺจ นีวรณวิปฺปยุตฺโต นีวรณิโย ธมฺโม อุปฺปชฺชติ เหตุปจฺจยา.\csromanspacer{} นีวรณวิปฺปยุตฺตํ นีวรณิยํ เอกํ ขนฺธํ ปฏิจฺจ ตโย ขนฺธา จิตฺตสมุฏฺฐานญฺจ รูปํ ฯเปฯ เทฺว ขนฺเธ ฯเปฯ.\csromanspacer{} ปฏิสนฺธิกฺขเณ ฯเปฯ. (สํขิตฺตํ.)}

\prose{(ยถา จูฬนฺตรทุเก โลกิยทุกํ, เอวํ กาตพฺพํ, นินฺนานากรณํ.)}

\nitthitam{นีวรณวิปฺปยุตฺตนีวรณิยทุกํ นิฏฺฐิตํ.}
\nitthitam{นีวรณโคจฺฉกํ นิฏฺฐิตํ.}
\csromanmatikaanchor{mtk.47}
\csromanmatikaanchor{mtk.48}
\csromantocmarknopage{section}{9. ปรามาสโคจฺฉก}{mtk.47}
\bookmark[dest={mtk.47},level=1]{9. ปรามาสโคจฺฉก}
\csromantocmarkref{subsection}{50. ปรามาสทุก}{mtk.48}
\bookmark[dest={mtk.48},level=2]{50. ปรามาสทุก}
\csromanheader{2}{50. ปรามาสทุก \textbf{9. ปรามาสโคจฺฉก}}
\csromanheader{2}{50. ปรามาสทุก 1. ปฏิจฺจวาร}
\csromanheader{2}{1. ปจฺจยานุโลม 1. วิภงฺควาร}
\csromanheader{2}{\textbf{เหตุ}}
\proseitem{1}{\csromanfolio{337}ปรามาสํ ธมฺมํ ปฏิจฺจ โนปรามาโส ธมฺโม อุปฺปชฺชติ เหตุปจฺจยา.\csromanspacer{} ปรามาสํ ปฏิจฺจ สมฺปยุตฺตกา ขนฺธา จิตฺตสมุฏฺฐานญฺจ รูปํ. (1)}

\proseitem{2}{โนปรามาสํ ธมฺมํ ปฏิจฺจ โนปรามาโส ธมฺโม อุปฺปชฺชติ เหตุปจฺจยา.\csromanspacer{} โนปรามาสํ เอกํ ขนฺธํ ปฏิจฺจ ตโย ขนฺธา จิตฺตสมุฏฺฐานญฺจ รูปํ ฯเปฯ เทฺว ขนฺเธ ฯเปฯ.\csromanspacer{} ปฏิสนฺธิกฺขเณ ฯเปฯ. (ยาว อสญฺญสตฺตา มหาภูตา.) (1)}

\prose{โนปรามาสํ ธมฺมํ ปฏิจฺจ ปรามาโส ธมฺโม อุปฺปชฺชติ เหตุปจฺจยา.\csromanspacer{} โนปรามาเส ขนฺเธ ปฏิจฺจ ปรามาโส. (2)}

\prose{โนปรามาสํ ธมฺมํ ปฏิจฺจ ปรามาโส จ โนปรามาโส จ ธมฺมา อุปฺปชฺชนฺติ เหตุปจฺจยา.\csromanspacer{} โนปรามาสํ เอกํ ขนฺธํ ปฏิจฺจ ตโย ขนฺธา ปรามาโส จ จิตฺตสมุฏฺฐานญฺจ รูปํ ฯเปฯ เทฺว ขนฺเธ ฯเปฯ. (3)}

\proseitem{3}{ปรามาสญฺจ โนปรามาสญฺจ ธมฺมํ ปฏิจฺจ โนปรามาโส ธมฺโม อุปฺปชฺชติ เหตุปจฺจยา.\csromanspacer{} โนปรามาสํ เอกํ ขนฺธญฺจ ปรามาสญฺจ ปฏิจฺจ ตโย ขนฺธา จิตฺตสมุฏฺฐานญฺจ รูปํ ฯเปฯ เทฺว ขนฺเธ จ ฯเปฯ.\csromanspacer{} ปรามาเส จ มหาภูเต จ ปฏิจฺจ จิตฺตสมุฏฺฐานํ รูปํ. (1)}

\csromanheader{2}{1. ปจฺจยานุโลม 2. สงฺขฺยาวาร}
\vspace{\tipitakaparskip}
\csromancenter{\textbf{เหตุ}ยา ปญฺจ, อารมฺมเณ ปญฺจ, (สพฺพตฺถ ปญฺจ.) วิปาเก เอกํ ฯเปฯ อวิคเต ปญฺจ.}
\csromancenter{อนุโลมํ.}
\csromanheader{2}{2. ปจฺจยปจฺจนีย 1. วิภงฺควาร}
\csromanheader{2}{\textbf{นเหตุ}}
\proseitem{5}{\csromanfolio{338}โนปรามาสํ ธมฺมํ ปฏิจฺจ โนปรามาโส ธมฺโม อุปฺปชฺชติ นเหตุปจฺจยา.\csromanspacer{} อเหตุกํ โนปรามาสํ เอกํ ขนฺธํ ปฏิจฺจ ตโย ขนฺธา จิตฺตสมุฏฺฐานญฺจ รูปํ ฯเปฯ เทฺว ขนฺเธ ฯเปฯ.\csromanspacer{} อเหตุกปฏิสนฺธิกฺขเณ ฯเปฯ. (ยาว อสญฺญสตฺตา.) วิจิกิจฺฉาสหคเต อุทฺธจฺจสหคเต ขนฺเธ ปฏิจฺจ วิจิกิจฺฉาสหคโต อุทฺธจฺจสหคโต โมโห. (1)}

\csromanheader{2}{\textbf{น-อารมฺมณ}}
\proseitem{6}{ปรามาสํ ธมฺมํ ปฏิจฺจ โนปรามาโส ธมฺโม อุปฺปชฺชติ น-อารมฺมณปจฺจยา.\csromanspacer{} ปรามาสํ ปฏิจฺจ จิตฺตสมุฏฺฐานํ รูปํ. (1)}

\prose{โนปรามาสํ ธมฺมํ ปฏิจฺจ โนปรามาโส ธมฺโม อุปฺปชฺชติ น-อารมฺมณปจฺจยา.\csromanspacer{} โนปรามาเส ขนฺเธ ปฏิจฺจ จิตฺตสมุฏฺฐานํ รูปํ.\csromanspacer{} ปฏิสนฺธิกฺขเณ ฯเปฯ. (ยาว อสญฺญสตฺตา.) (1)}

\prose{ปรามาสญฺจ โนปรามาสญฺจ ธมฺมํ ปฏิจฺจ โนปรามาโส ธมฺโม อุปฺปชฺชติ.\csromanspacer{} \textbf{น-อารมฺมณ}ปจฺจยา.\csromanspacer{} ปรามาสญฺจ สมฺปยุตฺตเก จ ขนฺเธ ปฏิจฺจ จิตฺตสมุฏฺฐานํ รูปํ. (1)}

\csromanheader{2}{2. ปจฺจยปจฺจนีย 2. สงฺขฺยาวาร}
\csromanheader{2}{\textbf{สุทฺธ}}
\vspace{\tipitakaparskip}
\csromancenter{\textbf{นเหตุ}ยา เอกํ, น-อารมฺมเณ ตีณิ, น-อธิปติยา ปญฺจ, น-อนนฺตเร ตีณิ, นสมนนฺตเร ตีณิ, น-อญฺญมญฺเญ ตีณิ, น-อุปนิสฺสเย ตีณิ, นปุเรชาเต ปญฺจ, นปจฺฉาชาเต ปญฺจ, น-อาเสวเน ปญฺจ, นกมฺเม ตีณิ, นวิปาเก ปญฺจ, น-อาหาเร เอกํ, น-อินฺทฺริเย เอกํ, นฌาเน เอกํ, นมคฺเค เอกํ, นสมฺปยุตฺเต ตีณิ, นวิปฺปยุตฺเต ปญฺจ, โนนตฺถิยา ตีณิ, โนวิคเต ตีณิ.}
\csromanheader{2}{50. ปรามาสทุก \textbf{3. ปจฺจยานุโลมปจฺจนีย}}
\csromanheader{2}{\textbf{เหตุทุก}}
\proseitem{8}{\csromanfolio{339}\textbf{เหตุ}ปจฺจยา น-อารมฺมเณ ตีณิ, น-อธิปติยา ปญฺจ ฯเปฯ นวิปาเก ปญฺจ ฯเปฯ นสมฺปยุตฺเต ตีณิ, นวิปฺปยุตฺเต ปญฺจ, โนนตฺถิยา ตีณิ, โนวิคเต ตีณิ.}

\csromanheader{2}{4. ปจฺจยปจฺจนียานุโลม}
\csromanheader{2}{\textbf{นเหตุทุก}}
\vspace{\tipitakaparskip}
\csromancenter{นเหตุปจฺจยา อารมฺมเณ เอกํ, อนนฺตเร เอกํ, (สพฺพตฺถ เอกํ.) อวิคเต เอกํ.}
\csromancenter{(สหชาตวาโรปิ ปฏิจฺจวารสทิโส.)}
\csromanheader{2}{50. ปรามาสทุก 3. ปจฺจยวาร}
\csromanheader{2}{1. ปจฺจยานุโลม 1. วิภงฺควาร}
\csromanheader{2}{\textbf{เหตุ}}
\proseitem{10}{ปรามาสํ ธมฺมํ ปจฺจยา โนปรามาโส ธมฺโม อุปฺปชฺชติ เหตุปจฺจยา.\csromanspacer{} ปรามาสํ ปจฺจยา สมฺปยุตฺตกา ขนฺธา จิตฺตสมุฏฺฐานญฺจ รูปํ. (1)}

\prose{โนปรามาสํ ธมฺมํ ปจฺจยา โนปรามาโส ธมฺโม อุปฺปชฺชติ เหตุปจฺจยา.\csromanspacer{} โนปรามาสํ เอกํ ขนฺธํ ปจฺจยา ตโย ขนฺธา จิตฺตสมุฏฺฐานญฺจ รูปํ ฯเปฯ เทฺว ขนฺเธ ฯเปฯ.\csromanspacer{} ปฏิสนฺธิกฺขเณ ฯเปฯ. (ยาว อชฺฌตฺติกา มหาภูตา.) วตฺถุํ ปจฺจยา โนปรามาสา ขนฺธา. (1)}

\prose{โนปรามาสํ ธมฺมํ ปจฺจยา ปรามาโส ธมฺโม อุปฺปชฺชติ เหตุปจฺจยา.\csromanspacer{} โนปรามาเส ขนฺเธ ปจฺจยา ปรามาโส.\csromanspacer{} วตฺถุํ ปจฺจยา ปรามาโส. (2)}

\prose{โนปรามาสํ ธมฺมํ ปจฺจยา ปรามาโส จ โนปรามาโส จ ธมฺมา อุปฺปชฺชนฺติ เหตุปจฺจยา.\csromanspacer{} โนปรามาสํ เอกํ ขนฺธํ ปจฺจยา ตโย ขนฺธา \csromanfolio{340}ปรามาโส จ จิตฺตสมุฏฺฐานญฺจ รูปํ ฯเปฯ เทฺว ขนฺเธ ฯเปฯ.\csromanspacer{} วตฺถุํ ปจฺจยา ปรามาโส.\csromanspacer{} มหาภูเต ปจฺจยา จิตฺตสมุฏฺฐานํ รูปํ.\csromanspacer{} วตฺถุํ ปจฺจยา ปรามาโส จ สมฺปยุตฺตกา จ ขนฺธา. (3)}

\proseitem{11}{ปรามาสญฺจ โนปรามาสญฺจ ธมฺมํ ปจฺจยา โนปรามาโส ธมฺโม อุปฺปชฺชติ เหตุปจฺจยา.\csromanspacer{} โนปรามาสํ เอกํ ขนฺธญฺจ ปรามาสญฺจ ปจฺจยา ตโย ขนฺธา จิตฺตสมุฏฺฐานญฺจ รูปํ ฯเปฯ เทฺว ขนฺเธ จ ฯเปฯ.\csromanspacer{} ปรามาสญฺจ สมฺปยุตฺตเก จ ขนฺเธ ปจฺจยา จิตฺตสมุฏฺฐานํ รูปํ.\csromanspacer{} ปรามาสญฺจ มหาภูเต จ ปจฺจยา จิตฺตสมุฏฺฐานํ รูปํ.\csromanspacer{} ปรามาสญฺจ วตฺถุญฺจ ปจฺจยา โนปรามาสา ขนฺธา. (สํขิตฺตํ.) (1)}

\csromanheader{2}{1. ปจฺจยานุโลม 2. สงฺขฺยาวาร}
\proseitem{12}{เหตุยา ปญฺจ, อารมฺมเณ ปญฺจ, อธิปติยา ปญฺจ, (สพฺพตฺถ ปญฺจ.) วิปาเก เอกํ ฯเปฯ อวิคเต ปญฺจ.}

\csromanheader{2}{2. ปจฺจยปจฺจนีย 1. วิภงฺควาร}
\csromanheader{2}{\textbf{นเหตุ}}
\proseitem{13}{โนปรามาสํ ธมฺมํ ปจฺจยา โนปรามาโส ธมฺโม อุปฺปชฺชติ นเหตุปจฺจยา.\csromanspacer{} อเหตุกํ โนปรามาสํ เอกํ ขนฺธํ ปจฺจยา ตโย ขนฺธา จิตฺตสมุฏฺฐานญฺจ รูปํ ฯเปฯ เทฺว ขนฺเธ ฯเปฯ.\csromanspacer{} อเหตุกปฏิสนฺธิกฺขเณ ฯเปฯ. (ยาว อสญฺญสตฺตา.) จกฺขายตนํ ปจฺจยา จกฺขุวิญฺญาณํ ฯเปฯ.\csromanspacer{} กายายตนํ ปจฺจยา กายวิญฺญาณํ.\csromanspacer{} วตฺถุํ ปจฺจยา อเหตุกา โนปรามาสา ขนฺธา.\csromanspacer{} วิจิกิจฺฉาสหคเต อุทฺธจฺจสหคเต ขนฺเธ จ วตฺถุญฺจ ปจฺจยา วิจิกิจฺฉาสหคโต อุทฺธจฺจสหคโต โมโห. (สํขิตฺตํ.)}

\csromanheader{2}{2. ปจฺจยปจฺจนีย 2. สงฺขฺยาวาร}
\csromanheader{2}{\textbf{สุทฺธ}}
\vspace{\tipitakaparskip}
\csromancenter{\textbf{นเหตุ}ยา เอกํ, น-อารมฺมเณ ตีณิ, น-อธิปติยา ปญฺจ, น-อนนฺตเร ตีณิ ฯเปฯ น-อุปนิสฺสเย ตีณิ, นปุเรชาเต ปญฺจ, นปจฺฉาชาเต}
\csromancenter{ปรามาสทุก ปญฺจ, อาเสวเน ปญฺจ, นกมฺเม ตีณิ, นวิปาเก ปญฺจ, น-อาหาเร เอกํ, น-อินฺทฺริเย เอกํ, นฌาเน เอกํ, นมคฺเค เอกํ, นสมฺปยุตฺเต ตีณิ, นวิปฺปยุตฺเต ปญฺจ, โนนตฺถิยา ตีณิ, โนวิคเต ตีณิ.}
\csromanheader{2}{3. ปจฺจยานุโลมปจฺจนีย}
\proseitem{15}{\csromanfolio{341}\textbf{เหตุ}ปจฺจยา น-อารมฺมเณ ตีณิ, น-อธิปติยา ปญฺจ. (เอวํ สพฺพตฺถ กาตพฺพํ.)}

\csromanheader{2}{4. ปจฺจยปจฺจนียานุโลม}
\vspace{\tipitakaparskip}
\csromancenter{นเหตุปจฺจยา อารมฺมเณ เอกํ ฯเปฯ อวิคเต เอกํ.}
\csromancenter{(นิสฺสยวาโร ปจฺจยวารสทิโส.)}
\csromanheader{2}{50. ปรามาสทุก 5. สํสฏฺฐวาร}
\prose{ปรามาสํ ธมฺมํ สํสฏฺโฐ โนปรามาโส ธมฺโม อุปฺปชฺชติ เหตุปจฺจยา.\csromanspacer{} ปรามาสํ สํสฏฺฐา สมฺปยุตฺตกา ขนฺธา. (เอวํ ปญฺจ ปญฺหา กาตพฺพา, อรูเปเยว, สํสฏฺฐวาโรปิ สมฺปยุตฺตวาโรปิ เอวํ กาตพฺพา.)}

\csromanheader{2}{50. ปรามาสทุก 7. ปญฺหาวาร}
\csromanheader{2}{1. ปจฺจยานุโลม 1. วิภงฺควาร}
\csromanheader{2}{\textbf{เหตุ}}
\proseitem{17}{โนปรามาโส ธมฺโม โนปรามาสสฺส ธมฺมสฺส เหตุปจฺจเยน ปจฺจโย.\csromanspacer{} โนปรามาสา เหตู สมฺปยุตฺตกานํ ขนฺธานํ จิตฺตสมุฏฺฐานานญฺจ รูปานํ เหตุปจฺจเยน ปจฺจโย.\csromanspacer{} ปฏิสนฺธิกฺขเณ ฯเปฯ. (1)}

\prose{โนปรามาโส ธมฺโม ปรามาสสฺส ธมฺมสฺส เหตุปจฺจเยน ปจฺจโย.\csromanspacer{} โนปรามาโส เหตู ปรามาสสฺส เหตุปจฺจเยน ปจฺจโย. (2)}

\prose{โนปรามาโส ธมฺโม ปรามาสสฺส จ โนปรามาสสฺส จ ธมฺมสฺส เหตุปจฺจเยน ปจฺจโย.\csromanspacer{} โนปรามาสา เหตู สมฺปยุตฺตกานํ ขนฺธานํ ปรามาสสฺส จ จิตฺตสมุฏฺฐานานญฺจ รูปานํ เหตุปจฺจเยน ปจฺจโย. (3)}

\csromanheader{2}{\textbf{อารมฺมณ}}
\proseitem{18}{\csromanfolio{342}ปรามาโส ธมฺโม ปรามาสสฺส ธมฺมสฺส อารมฺมณปจฺจเยน ปจฺจโย.\csromanspacer{} ปรามาสํ อารพฺภ ปรามาโส อุปฺปชฺชติ. (มูลํ กาตพฺพํ.) ปรามาสํ อารพฺภ โนปรามาสา ขนฺธา อุปฺปชฺชนฺติ. (มูลํ กาตพฺพํ.) ปรามาสํ อารพฺภ ปรามาโส จ สมฺปยุตฺตกา จ ขนฺธา อุปฺปชฺชนฺติ. (3)}

\proseitem{19}{โนปรามาโส ธมฺโม โนปรามาสสฺส ธมฺมสฺส อารมฺมณปจฺจเยน ปจฺจโย.\csromanspacer{} ทานํ ฯเปฯ สีลํ ฯเปฯ อุโปสถกมฺมํ กตฺวา ตํ ปจฺจเวกฺขติ, อสฺสาเทติ อภินนฺทติ, ตํ อารพฺภ ราโค ฯเปฯ วิจิกิจฺฉา.\csromanspacer{} อุทฺธจฺจํ.\csromanspacer{} โทมนสฺสํ อุปฺปชฺชติ.\csromanspacer{} ปุพฺเพ สุจิณฺณานิ ฯเปฯ.\csromanspacer{} ฌานา วุฏฺฐหิตฺวา ฌานํ ฯเปฯ.\csromanspacer{} อริยา มคฺคา วุฏฺฐหิตฺวา มคฺคํ ปจฺจเวกฺขนฺติ, ผลํ ปจฺจเวกฺขนฺติ, นิพฺพานํ ปจฺจเวกฺขนฺติ.\csromanspacer{} นิพฺพานํ โคตฺรภุสฺส, โวทานสฺส, มคฺคสฺส, ผลสฺส, อาวชฺชนาย อารมฺมณปจฺจเยน ปจฺจโย.\csromanspacer{} อริยา โนปรามาเส ปหีเน กิเลเส ฯเปฯ.\csromanspacer{} วิกฺขมฺภิเต กิเลเส ฯเปฯ.\csromanspacer{} ปุพฺเพ ฯเปฯ.\csromanspacer{} จกฺขุํ ฯเปฯ วตฺถุํ โนปรามาเส ขนฺเธ อนิจฺจโต ฯเปฯ วิปสฺสติ, อสฺสาเทติ อภินนฺทติ, ตํ อารพฺภ ราโค ฯเปฯ วิจิกิจฺฉา.\csromanspacer{} อุทฺธจฺจํ.\csromanspacer{} โทมนสฺสํ อุปฺปชฺชติ.\csromanspacer{} ทิพฺเพน จกฺขุนา รูปํ ปสฺสติ.\csromanspacer{} ทิพฺพาย โสตธาตุยา สทฺทํ สุณาติ.\csromanspacer{} เจโตปริยญาเณน โนปรามาสจิตฺตสมงฺคิสฺส จิตฺตํ ชานาติ.\csromanspacer{} อากาสานญฺจายตนํ วิญฺญาณญฺจายตนสฺส ฯเปฯ.\csromanspacer{} อากิญฺจญฺญายตนํ เนวสญฺญานาสญฺญายตนสฺส ฯเปฯ.\csromanspacer{} รูปายตนํ จกฺขุวิญฺญาณสฺส ฯเปฯ.\csromanspacer{} โผฏฺฐพฺพายตนํ กายวิญฺญาณสฺส ฯเปฯ โนปรามาสา ขนฺธา อิทฺธิวิธญาณสฺส, เจโตปริยญาณสฺส, ปุพฺเพนิวาสานุสฺสติญาณสฺส, ยถากมฺมูปคญาณสฺส อนาคตํสญาณสฺส, อาวชฺชนาย อารมฺมณปจฺจเยน ปจฺจโย. (1)}

\prose{โนปรามาโส ธมฺโม ปรามาสสฺส ธมฺมสฺส อารมฺมณปจฺจเยน ปจฺจโย.\csromanspacer{} ทานํ ฯเปฯ สีลํ ฯเปฯ อุโปสถกมฺมํ ฯเปฯ ตํ อสฺสาเทติ, อภินนฺทติ, ตํ อารพฺภ ทิฏฺฐิ อุปฺปชฺชติ.\csromanspacer{} ปุพฺเพ ฯเปฯ.\csromanspacer{} ฌานา ฯเปฯ.\csromanspacer{} จกฺขุํ ฯเปฯ วตฺถุํ โนปรามาเส ขนฺเธ อสฺสาเทติ อภินนฺทติ.\csromanspacer{} ตํ อารพฺภ ทิฏฺฐิ อุปฺปชฺชติ. (2)}

\prose{โนปรามาโส ธมฺโม ปรามาสสฺส จ โนปรามาสสฺส จ ธมฺมสฺส อารมฺมณปจฺจเยน ปจฺจโย.\csromanspacer{} ทานํ ฯเปฯ สีลํ ฯเปฯ อุโปสถกมฺมํ ฯเปฯ ตํ อสฺสาเทติ อภินนฺทติ, ตํ อารพฺภ ปรามาโส จ สมฺปยุตฺตกา จ ขนฺธา อุปฺปชฺชนฺติ.\csromanspacer{} ปุพฺเพ ฯเปฯ.\csromanspacer{} ฌานา ฯเปฯ.\csromanspacer{} จกฺขุํ ฯเปฯ วตฺถุํ โนปรามาเส ขนฺเธ}

\vspace{\tipitakaparskip}
\csromancenter{ปรามาสทุก อสฺสาเทติ อภินนฺทติ, ตํ อารพฺภ ปรามาโส จ สมฺปยุตฺตกา จ ขนฺธา อุปฺปชฺชนฺติ. (3)}
\proseitem{20}{\csromanfolio{343}ปรามาโส จ โนปรามาโส จ ธมฺมา ปรามาสสฺส ธมฺมสฺส อารมฺมณปจฺจเยน ปจฺจโย.\csromanspacer{} ปรามาสญฺจ สมฺปยุตฺตเก จ ขนฺเธ อารพฺภ ปรามาโส อุปฺปชฺชติ.\csromanspacer{} ตีณิ.}

\csromanheader{2}{\textbf{อธิปติ}}
\proseitem{21}{ปรามาโส ธมฺโม ปรามาสสฺส ธมฺมสฺส อธิปติปจฺจเยน ปจฺจโย.\csromanspacer{} ปรามาสํ ครุํ กตฺวา ปรามาโส อุปฺปชฺชติ.\csromanspacer{} ตีณิ. (\textbf{อารมฺมณาธิปติ}เยว กาตพฺพา.)}

\proseitem{22}{โนปรามาโส ธมฺโม โนปรามาสสฺส ธมฺมสฺส อธิปติปจฺจเยน ปจฺจโย.\csromanspacer{} \textbf{อารมฺมณาธิปติ}, สหชาตาธิปติ.\csromanspacer{}\csromanspacer{}\textbf{อารมฺมณาธิปติ}\csromandash{}ทานํ ฯเปฯ สีลํ ฯเปฯ อุโปสถกมฺมํ ฯเปฯ ตํ ครุํ กตฺวา ปจฺจเวกฺขติ, อสฺสาเทติ อภินนฺทติ, ตํ ครุํ กตฺวา ราโค อุปฺปชฺชติ.\csromanspacer{} ปุพฺเพ ฯเปฯ.\csromanspacer{} ฌานา ฯเปฯ.\csromanspacer{} อริยา มคฺคา วุฏฺฐหิตฺวา มคฺคํ ครุํ กตฺวา ฯเปฯ.\csromanspacer{} ผลํ ครุํ กตฺวา ฯเปฯ.\csromanspacer{} นิพฺพานํ ครุํ กตฺวา ฯเปฯ.\csromanspacer{} นิพฺพานํ โคตฺรภุสฺส, โวทานสฺส, มคฺคสฺส, ผลสฺส อธิปติปจฺจเยน ปจฺจโย.\csromanspacer{} จกฺขุํ ฯเปฯ วตฺถุํ โนปรามาเส ขนฺเธ ครุํ กตฺวา อสฺสาเทติ อภินนฺทติ, ตํ ครุํ กตฺวา ราโค อุปฺปชฺชติ.\csromanspacer{} \textbf{สหชาตาธิปติ}\csromandash{} โนปรามาสาธิปติ สมฺปยุตฺตกานํ ขนฺธานํ จิตฺตสมุฏฺฐานานญฺจ รูปานํ อธิปติปจฺจเยน ปจฺจโย. (1)}

\prose{โนปรามาโส ธมฺโม ปรามาสสฺส ธมฺมสฺส อธิปติปจฺจเยน ปจฺจโย.\csromanspacer{} \textbf{อารมฺมณาธิปติ}, สหชาตาธิปติ.\csromanspacer{}\csromanspacer{}\textbf{อารมฺมณาธิปติ}\csromandash{}ทานํ ฯเปฯ สีลํ ฯเปฯ อุโปสถกมฺมํ ฯเปฯ.\csromanspacer{} ปุพฺเพ ฯเปฯ.\csromanspacer{} ฌานา ฯเปฯ.\csromanspacer{} จกฺขุํ ฯเปฯ วตฺถุํ โนปรามาเส ขนฺเธ ครุํ กตฺวา อสฺสาเทติ อภินนฺทติ, ตํ ครุํ กตฺวา ทิฏฺฐิ อุปฺปชฺชติ.\csromanspacer{} \textbf{สหชาตาธิปติ}\csromandash{}โนปรามาสาธิปติ ปรามาสสฺส อธิปติปจฺจเยน ปจฺจโย. (2)}

\prose{โนปรามาโส ธมฺโม ปรามาสสฺส จ โนปรามาสสฺส จ ธมฺมสฺส อธิปติปจฺจเยน ปจฺจโย.\csromanspacer{} \textbf{อารมฺมณาธิปติ}, สหชาตาธิปติ.\csromanspacer{}\csromanspacer{}\textbf{อารมฺมณาธิปติ}\csromandash{}ทานํ ฯเปฯ สีลํ ฯเปฯ อุโปสถกมฺมํ ฯเปฯ.\csromanspacer{} ปุพฺเพ ฯเปฯ. \csromanfolio{344}ฌานา ฯเปฯ.\csromanspacer{} จกฺขุํ ฯเปฯ วตฺถุํ โนปรามาเส ขนฺเธ ครุํ กตฺวา อสฺสาเทติ อภินนฺทติ, ตํ ครุํ กตฺวา ปรามาโส จ สมฺปยุตฺตกา จ ขนฺธา อุปฺปชฺชนฺติ.\csromanspacer{} \textbf{สหชาตาธิปติ}\csromandash{}โนปรามาสาธิปติ สมฺปยุตฺตกานํ ขนฺธานํ ปรามาสสฺส จ จิตฺตสมุฏฺฐานานญฺจ รูปานํ อธิปติปจฺจเยน ปจฺจโย. (3)}

\prose{ปรามาโส จ โนปรามาโส จ ธมฺมา ปรามาสสฺส ธมฺมสฺส อธิปติปจฺจเยน ปจฺจโย.\csromanspacer{} \textbf{อารมฺมณาธิปติ}\csromandash{}ปรามาสญฺจ สมฺปยุตฺตเก จ ขนฺเธ ครุํ กตฺวา ปรามาโส.\csromanspacer{} ตีณิ.}

\csromanheader{2}{\textbf{อนนฺตร}}
\proseitem{23}{ปรามาโส ธมฺโม ปรามาสสฺส ธมฺมสฺส อนนฺตรปจฺจเยน ปจฺจโย.\csromanspacer{} ปุริโม ปุริโม ปรามาโส ปจฺฉิมสฺส ปจฺฉิมสฺส ปรามาสสฺส อนนฺตรปจฺจเยน ปจฺจโย. (มูลํ กาตพฺพํ.) ปุริโม ปุริโม ปรามาโส ปจฺฉิมานํ ปจฺฉิมานํ โนปรามาสานํ ขนฺธานํ อนนฺตรปจฺจเยน ปจฺจโย.\csromanspacer{} ปรามาโส วุฏฺฐานสฺส อนนฺตรปจฺจเยน ปจฺจโย. (มูลํ กาตพฺพํ.) ปุริโม ปุริโม ปรามาโส ปจฺฉิมสฺส ปจฺฉิมสฺส ปรามาสสฺส สมฺปยุตฺตกานญฺจ ขนฺธานํ อนนฺตรปจฺจเยน ปจฺจโย. (3)}

\proseitem{24}{โนปรามาโส ธมฺโม โนปรามาสสฺส ธมฺมสฺส อนนฺตรปจฺจเยน ปจฺจโย.\csromanspacer{} ปุริมา ปุริมา โนปรามาสา ขนฺธา ปจฺฉิมานํ ปจฺฉิมานํ โนปรามาสานํ ขนฺธานํ อนนฺตรปจฺจเยน ปจฺจโย ฯเปฯ.\csromanspacer{} อนุโลมํ ผลสมาปตฺติยา อนนฺตรปจฺจเยน ปจฺจโย. (มูลํ กาตพฺพํ.) ปุริมา ปุริมา โนปรามาสา ขนฺธา ปจฺฉิมสฺส ปจฺฉิมสฺส ปรามาสสฺส อนนฺตรปจฺจเยน ปจฺจโย.\csromanspacer{} อาวชฺชนา ปรามาสสฺส อนนฺตรปจฺจเยน ปจฺจโย. (มูลํ กาตพฺพํ.) ปุริมา ปุริมา โนปรามาสา ขนฺธา ปจฺฉิมสฺส ปจฺฉิมสฺส ปรามาสสฺส จ สมฺปยุตฺตกานญฺจ ขนฺธานํ อนนฺตรปจฺจเยน ปจฺจโย.\csromanspacer{} อาวชฺชนา ปรามาสสฺส จ สมฺปยุตฺตกานญฺจ ขนฺธานํ อนนฺตรปจฺจเยน ปจฺจโย. (3)}

\proseitem{25}{ปรามาโส จ โนปรามาโส จ ธมฺมา ปรามาสสฺส ธมฺมสฺส อนนฺตรปจฺจเยน ปจฺจโย.\csromanspacer{} ปุริโม ปุริโม ปรามาโส จ สมฺปยุตฺตกา จ ขนฺธา ปจฺฉิมสฺส ปจฺฉิมสฺส ปรามาสสฺส อนนฺตรปจฺจเยน ปจฺจโย. (มูลํ}

\vspace{\tipitakaparskip}
\csromancenter{ปรามาสทุก กาตพฺพํ.) ปุริโม ปุริโม ปรามาโส จ สมฺปยุตฺตกา จ ขนฺธา ปจฺฉิมานํ ปจฺฉิมานํ โนปรามาสานํ ขนฺธานํ อนนฺตรปจฺจเยน ปจฺจโย.\csromanspacer{} ปรามาโส จ สมฺปยุตฺตกา จ ขนฺธา วุฏฺฐานสฺส อนนฺตรปจฺจเยน ปจฺจโย. (มูลํ กาตพฺพํ.) ปุริโม ปุริโม ปรามาโส จ สมฺปยุตฺตกา จ ขนฺธา ปจฺฉิมสฺส ปจฺฉิมสฺส ปรามาสสฺส สมฺปยุตฺตกานญฺจ ขนฺธานํ อนนฺตรปจฺจเยน ปจฺจโย. (3)}
\csromanheader{2}{\textbf{สมนนฺตราทิ}}
\proseitem{26}{\csromanfolio{345}ปรามาโส ธมฺโม ปรามาสสฺส ธมฺมสฺส สมนนฺตรปจฺจเยน ปจฺจโย.\csromanspacer{} สหชาตปจฺจเยน ปจฺจโย.\csromanspacer{} อญฺญมญฺญปจฺจเยน ปจฺจโย.\csromanspacer{} นิสฺสยปจฺจเยน ปจฺจโย.\csromanspacer{} ปญฺจ.}

\csromanheader{2}{\textbf{อุปนิสฺสย}}
\proseitem{27}{ปรามาโส ธมฺโม ปรามาสสฺส ธมฺมสฺส อุปนิสฺสยปจฺจเยน ปจฺจโย.\csromanspacer{} อารมฺมณูปนิสฺสโย, อนนฺตรูปนิสฺสโย, ปกตูปนิสฺสโย ฯเปฯ.\csromanspacer{} ปกตูปนิสฺสโย\csromandash{}ปรามาโส ปรามาสสฺส อุปนิสฺสยปจฺจเยน ปจฺจโย.\csromanspacer{} ตีณิ.}

\proseitem{28}{โนปรามาโส ธมฺโม โนปรามาสสฺส ธมฺมสฺส อุปนิสฺสยปจฺจเยน ปจฺจโย.\csromanspacer{} อารมฺมณูปนิสฺสโย, อนนฺตรูปนิสฺสโย, ปกตูปนิสฺสโย ฯเปฯ.\csromanspacer{} ปกตูปนิสฺสโย\csromandash{}สทฺธํ อุปนิสฺสาย ทานํ เทติ ฯเปฯ สมาปตฺติํ อุปฺปาเทติ.\csromanspacer{} มานํ ชปฺเปติ.\csromanspacer{} สีลํ ฯเปฯ.\csromanspacer{} ปญฺญํ.\csromanspacer{} ราคํ.\csromanspacer{} โทสํ.\csromanspacer{} โมหํ.\csromanspacer{} มานํ.\csromanspacer{} ปตฺถนํ.\csromanspacer{} กายิกํ สุขํ ฯเปฯ.\csromanspacer{} เสนาสนํ อุปนิสฺสาย ทานํ เทติ ฯเปฯ สมาปตฺติํ อุปฺปาเทติ.\csromanspacer{} ปาณํ หนติ ฯเปฯ สํฆํ ภินฺทติ.\csromanspacer{} สทฺธา ฯเปฯ.\csromanspacer{} เสนาสนํ.\csromanspacer{} สทฺธาย ฯเปฯ.\csromanspacer{} ปญฺญาย.\csromanspacer{} ราคสฺส ฯเปฯ ปตฺถนาย.\csromanspacer{} กายิกสฺส สุขสฺส ฯเปฯ มคฺคสฺส, ผลสมาปตฺติยา อุปนิสฺสยปจฺจเยน ปจฺจโย. (1)}

\prose{โนปรามาโส ธมฺโม ปรามาสสฺส ธมฺมสฺส อุปนิสฺสยปจฺจเยน ปจฺจโย.\csromanspacer{} อารมฺมณูปนิสฺสโย, อนนฺตรูปนิสฺสโย, ปกตูปนิสฺสโย ฯเปฯ.\csromanspacer{} ปกตูปนิสฺสโย\csromandash{}สทฺธํ อุปนิสฺสาย ทิฏฺฐิํ คณฺหาติ.\csromanspacer{} สีลํ ฯเปฯ.\csromanspacer{} ปญฺญํ.\csromanspacer{} ราคํ ฯเปฯ.\csromanspacer{} ปตฺถนํ.\csromanspacer{} กายิกํ สุขํ ฯเปฯ.\csromanspacer{} เสนาสนํ อุปนิสฺสาย ทิฏฺฐิํ คณฺหาติ.\csromanspacer{} สทฺธา ฯเปฯ.\csromanspacer{} เสนาสนํ ปรามาสสฺส อุปนิสฺสยปจฺจเยน ปจฺจโย. (2)}

\prose{\csromanfolio{346}โนปรามาโส ธมฺโม ปรามาสสฺส จ โนปรามาสสฺส จ ธมฺมสฺส อุปนิสฺสยปจฺจเยน ปจฺจโย.\csromanspacer{} อารมฺมณูปนิสฺสโย, อนนฺตรูปนิสฺสโย, ปกตูปนิสฺสโย ฯเปฯ.\csromanspacer{} ปกตูปนิสฺสโย\csromandash{}สทฺธํ อุปนิสฺสาย ทิฏฺฐิํ คณฺหาติ.\csromanspacer{} สีลํ ฯเปฯ.\csromanspacer{} เสนาสนํ อุปนิสฺสาย ทิฏฺฐิํ คณฺหาติ.\csromanspacer{} สทฺธา ฯเปฯ.\csromanspacer{} เสนาสนํ ปรามาสสฺส สมฺปยุตฺตกานญฺจ ขนฺธานํ อุปนิสฺสยปจฺจเยน ปจฺจโย. (3)}

\proseitem{29}{ปรามาโส จ โนปรามาโส จ ธมฺมา ปรามาสสฺส ธมฺมสฺส อุปนิสฺสยปจฺจเยน ปจฺจโย.\csromanspacer{} ตีณิ อุปนิสฺสยา ปรามาโส จ สมฺปยุตฺตกา จ ขนฺธา ปรามาสสฺส อุปนิสฺสยปจฺจเยน ปจฺจโย.\csromanspacer{} ตีณิ.}

\csromanheader{2}{\textbf{ปุเรชาต}}
\proseitem{30}{โนปรามาโส ธมฺโม โนปรามาสสฺส ธมฺมสฺส ปุเรชาตปจฺจเยน ปจฺจโย.\csromanspacer{} \textbf{อารมฺมณปุเรชาตํ}, วตฺถุปุเรชาตํ.\csromanspacer{}\csromanspacer{}\textbf{อารมฺมณปุเรชาตํ}\csromandash{}จกฺขุํ ฯเปฯ วตฺถุํ อนิจฺจโต ฯเปฯ วิปสฺสติ, อสฺสาเทติ อภินนฺทติ, ตํ อารพฺภ ราโค ฯเปฯ วิจิกิจฺฉา ฯเปฯ อุทฺธจฺจํ ฯเปฯ โทมนสฺสํ อุปฺปชฺชติ.\csromanspacer{} ทิพฺเพน จกฺขุนา รูปํ ปสฺสติ.\csromanspacer{} ทิพฺพาย โสตธาตุยา สทฺทํ สุณาติ.\csromanspacer{} รูปายตนํ จกฺขุวิญฺญาณสฺส ฯเปฯ.\csromanspacer{} โผฏฺฐพฺพายตนํ กายวิญฺญาณสฺส ฯเปฯ.\csromanspacer{} \textbf{วตฺถุปุเรชาตํ}\csromandash{}จกฺขายตนํ จกฺขุวิญฺญาณสฺส ฯเปฯ กายายตนํ กายวิญฺญาณสฺส ฯเปฯ.\csromanspacer{} วตฺถุ โนปรามาสานํ ขนฺธานํ ปุเรชาตปจฺจเยน ปจฺจโย. (1)}

\prose{โนปรามาโส ธมฺโม ปรามาสสฺส ธมฺมสฺส ปุเรชาตปจฺจเยน ปจฺจโย.\csromanspacer{} \textbf{อารมฺมณปุเรชาตํ}, วตฺถุปุเรชาตํ.\csromanspacer{}\csromanspacer{}\textbf{อารมฺมณปุเรชาตํ}\csromandash{} จกฺขุํ ฯเปฯ วตฺถุํ อสฺสาเทติ อภินนฺทติ, ตํ อารพฺภ ทิฏฺฐิ อุปฺปชฺชติ.\csromanspacer{} \textbf{วตฺถุปุเรชาตํ}\csromandash{}วตฺถุ ปรามาสสฺส ปุเรชาตปจฺจเยน ปจฺจโย. (2)}

\prose{โนปรามาโส ธมฺโม ปรามาสสฺส จ โนปรามาสสฺส จ ธมฺมสฺส ปุเรชาตปจฺจเยน ปจฺจโย.\csromanspacer{} \textbf{อารมฺมณปุเรชาตํ}, วตฺถุปุเรชาตํ.\csromanspacer{}\csromanspacer{}\textbf{อารมฺมณปุเรชาตํ}\csromandash{}จกฺขุํ ฯเปฯ วตฺถุํ อสฺสาเทติ อภินนฺทติ, ตํ อารพฺภ ปรามาโส จ สมฺปยุตฺตกา จ ขนฺธา อุปฺปชฺชนฺติ.\csromanspacer{} \textbf{วตฺถุปุเรชาตํ}\csromandash{}วตฺถุ ปรามาสสฺส จ สมฺปยุตฺตกานญฺจ ขนฺธานํ ปุเรชาตปจฺจเยน ปจฺจโย. (3)}

\csromanheader{2}{50. ปรามาสทุก \textbf{ปจฺฉาชาตาเสวน}}
\proseitem{31}{\csromanfolio{347}ปรามาโส ธมฺโม โนปรามาสสฺส ธมฺมสฺส ปจฺฉาชาตปจฺจเยน ปจฺจโย.\csromanspacer{} ตีณิ. (ปจฺฉาชาตา กาตพฺพา.) อาเสวนปจฺจเยน ปจฺจโย.\csromanspacer{} นว.}

\csromanheader{2}{\textbf{กมฺม}}
\proseitem{32}{โนปรามาโส ธมฺโม โนปรามาสสฺส ธมฺมสฺส กมฺมปจฺจเยน ปจฺจโย.\csromanspacer{} \textbf{สหชาตา}, นานากฺขณิกา.\csromanspacer{}\csromanspacer{}\textbf{สหชาตา} โนปรามาสา เจตนา สมฺปยุตฺตกานํ ขนฺธานํ จิตฺตสมุฏฺฐานานญฺจ รูปานํ กมฺมปจฺจเยน ปจฺจโย.\csromanspacer{} ปฏิสนฺธิกฺขเณ ฯเปฯ.\csromanspacer{} \textbf{นานากฺขณิกา} โนปรามาสา เจตนา วิปากานํ ขนฺธานํ กฏตฺตา จ รูปานํ กมฺมปจฺจเยน ปจฺจโย. (มูลํ กาตพฺพํ.) โนปรามาสา เจตนา ปรามาสสฺส กมฺมปจฺจเยน ปจฺจโย. (มูลํ กาตพฺพํ.) โนปรามาสา เจตนา ปรามาสสฺส จ สมฺปยุตฺตกานญฺจ ขนฺธานํ จิตฺตสมุฏฺฐานานญฺจ รูปานํ กมฺมปจฺจเยน ปจฺจโย. (3)}

\csromanheader{2}{\textbf{วิปากาทิ}}
\proseitem{33}{โนปรามาโส ธมฺโม โนปรามาสสฺส ธมฺมสฺส วิปากปจฺจเยน ปจฺจโย.\csromanspacer{} เอกํ.\csromanspacer{} อาหารปจฺจเยน ปจฺจโย.\csromanspacer{} ตีณิ.\csromanspacer{} อินฺทฺริยปจฺจเยน ปจฺจโย.\csromanspacer{} ตีณิ.\csromanspacer{} ฌานปจฺจเยน ปจฺจโย.\csromanspacer{} ตีณิ.}

\prose{ปรามาโส ธมฺโม โนปรามาสสฺส ธมฺมสฺส มคฺคปจฺจเยน ปจฺจโย.\csromanspacer{} ปรามาสานิ มคฺคงฺคานิ ฯเปฯ. (เอวํ ปญฺจ ปญฺหา กาตพฺพา.) สมฺปยุตฺตปจฺจเยน ปจฺจโย.\csromanspacer{} ปญฺจ}

\csromanheader{2}{\textbf{วิปฺปยุตฺต}}
\proseitem{34}{ปรามาโส ธมฺโม โนปรามาสสฺส ธมฺมสฺส วิปฺปยุตฺตปจฺจเยน ปจฺจโย.\csromanspacer{} สหชาตํ, ปจฺฉาชาตํ. (สํขิตฺตํ.) (1)}

\prose{โนปรามาโส ธมฺโม โนปรามาสสฺส ธมฺมสฺส วิปฺปยุตฺตปจฺจเยน ปจฺจโย.\csromanspacer{} สหชาตํ, ปุเรชาตํ, ปจฺฉาชาตํ. (สํขิตฺตํ.) (1)}

\prose{โนปรามาโส ธมฺโม ปรามาสสฺส ธมฺมสฺส วิปฺปยุตฺตปจฺจเยน ปจฺจโย.\csromanspacer{} \textbf{ปุเรชาตํ} วตฺถุ ปรามาสสฺส วิปฺปยุตฺตปจฺจเยน ปจฺจโย. (2)}

\prose{\csromanfolio{348}โนปรามาโส ธมฺโม ปรามาสสฺส จ โนปรามาสสฺส จ ธมฺมสฺส วิปฺปยุตฺตปจฺจเยน ปจฺจโย.\csromanspacer{} \textbf{ปุเรชาตํ} วตฺถุ ปรามาสสฺส จ สมฺปยุตฺตกานญฺจ ขนฺธานํ วิปฺปยุตฺตปจฺจเยน ปจฺจโย. (3)}

\prose{ปรามาโส จ โนปรามาโส จ ธมฺมา โนปรามาสสฺส ธมฺมสฺส วิปฺปยุตฺตปจฺจเยน ปจฺจโย.\csromanspacer{} สหชาตํ, ปจฺฉาชาตํ. (สํขิตฺตํ.) (1)}

\csromanheader{2}{\textbf{อตฺถิ}}
\proseitem{35}{ปรามาโส ธมฺโม โนปรามาสสฺส ธมฺมสฺส อตฺถิปจฺจเยน ปจฺจโย.\csromanspacer{} สหชาตํ, ปจฺฉาชาตํ.\csromanspacer{}\csromanspacer{}\textbf{สหชาโต} ปรามาโส สมฺปยุตฺตกานญฺจ ขนฺธานํ จิตฺตสมุฏฺฐานานญฺจ รูปานํ อตฺถิปจฺจเยน ปจฺจโย.\csromanspacer{} \textbf{ปจฺฉาชาโต} ปรามาโส ปุเรชาตสฺส อิมสฺส กายสฺส อตฺถิปจฺจเยน ปจฺจโย. (1)}

\proseitem{36}{โนปรามาโส ธมฺโม โนปรามาสสฺส ธมฺมสฺส อตฺถิปจฺจเยน ปจฺจโย.\csromanspacer{} สหชาตํ, ปุเรชาตํ, ปจฺฉาชาตํ, อาหารํ, อินฺทฺริยํ. (สํขิตฺตํ.) (1)}

\prose{โนปรามาโส ธมฺโม ปรามาสสฺส ธมฺมสฺส อตฺถิปจฺจเยน ปจฺจโย.\csromanspacer{} สหชาตํ, ปุเรชาตํ.\csromanspacer{}\csromanspacer{}\textbf{สหชาตา} โนปรามาสา ขนฺธา ปรามาสสฺส อตฺถิปจฺจเยน ปจฺจโย.\csromanspacer{} \textbf{ปุเรชาตํ} จกฺขุํ ฯเปฯ วตฺถุํ อสฺสาเทติ อภินนฺทติ, ตํ อารพฺภ ทิฏฺฐิ อุปฺปชฺชติ.\csromanspacer{} วตฺถุ ปรามาสสฺส อตฺถิปจฺจเยน ปจฺจโย. (2)}

\prose{โนปรามาโส ธมฺโม ปรามาสสฺส จ โนปรามาสสฺส จ ธมฺมสฺส อตฺถิปจฺจเยน ปจฺจโย.\csromanspacer{} สหชาตํ, ปุเรชาตํ.\csromanspacer{}\csromanspacer{}\textbf{สหชาโต} โนปรามาโส เอโก ขนฺโธ ติณฺณนฺนํ ขนฺธานํ ปรามาสสฺส จ จิตฺตสมุฏฺฐานานญฺจ รูปานํ อตฺถิปจฺจเยน ปจฺจโย ฯเปฯ.\csromanspacer{} \textbf{ปุเรชาตํ} จกฺขุํ ฯเปฯ วตฺถุํ อสฺสาเทติ อภินนฺทติ, ตํ อารพฺภ ปรามาโส จ สมฺปยุตฺตกา จ ขนฺธา อุปฺปชฺชนฺติ.\csromanspacer{} วตฺถุ ปรามาสสฺส จ สมฺปยุตฺตกานญฺจ ขนฺธานํ อตฺถิปจฺจเยน ปจฺจโย. (3)}

\proseitem{37}{ปรามาโส จ โนปรามาโส จ ธมฺมา โนปรามาสสฺส ธมฺมสฺส อตฺถิปจฺจเยน ปจฺจโย.\csromanspacer{} สหชาตํ, ปุเรชาตํ, ปจฺฉาชาตํ}

\vspace{\tipitakaparskip}
\csromancenter{ปรามาสทุก อาหารํ, อินฺทฺริยํ.\csromanspacer{}\csromanspacer{}\textbf{สหชาโต} โนปรามาโส เอโก ขนฺโธ จ ปรามาโส จ ติณฺณนฺนํ ขนฺธานํ จิตฺตสมุฏฺฐานานญฺจ รูปานํ อตฺถิปจฺจเยน ปจฺจโย ฯเปฯ เทฺว ขนฺธา จ ฯเปฯ.\csromanspacer{} ปรามาโส จ สมฺปยุตฺตกา จ ขนฺธา จิตฺตสมุฏฺฐานานํ รูปานํ อตฺถิปจฺจเยน ปจฺจโย.\csromanspacer{} ปรามาโส จ มหาภูตา จ จิตฺตสมุฏฺฐานานํ รูปานํ อตฺถิปจฺจเยน ปจฺจโย.\csromanspacer{} \textbf{ปรามาโส จ วตฺถุ จ} โนปรามาสานํ ขนฺธานํ อตฺถิปจฺจเยน ปจฺจโย.\csromanspacer{} \textbf{ปจฺฉาชาโต} ปรามาโส จ สมฺปยุตฺตกา จ ขนฺธา ปุเรชาตสฺส อิมสฺส กายสฺส อตฺถิปจฺจเยน ปจฺจโย.\csromanspacer{} \textbf{ปจฺฉาชาโต} ปรามาโส จ สมฺปยุตฺตกา จ ขนฺธา กพฬีกาโร อาหาโร จ อิมสฺส กายสฺส อตฺถิปจฺจเยน ปจฺจโย.\csromanspacer{} \textbf{ปจฺฉาชาโต} ปรามาโส จ สมฺปยุตฺตกา จ ขนฺธา รูปชีวิตินฺทฺริยญฺจ กฏตฺตารูปานํ อตฺถิปจฺจเยน ปจฺจโย. (1)}
\csromanheader{2}{1. ปจฺจยานุโลม 2. สงฺขฺยาวาร}
\csromanheader{2}{\textbf{สุทฺธ}}
\proseitem{38}{\csromanfolio{349}เหตุยา ตีณิ, อารมฺมเณ นว, อธิปติยา นว, อนนฺตเร นว, สมนนฺตเร นว, สหชาเต ปญฺจ, อญฺญมญฺเญ ปญฺจ, นิสฺสเย ปญฺจ, อุปนิสฺสเย นว, ปุเรชาเต ตีณิ, ปจฺฉาชาเต ตีณิ, อาเสวเน นว, กมฺเม ตีณิ, วิปาเก เอกํ, อาหาเร ตีณิ, อินฺทฺริเย ตีณิ, ฌาเน ตีณิ, มคฺเค ปญฺจ, สมฺปยุตฺเต ปญฺจ, วิปฺปยุตฺเต ปญฺจ, อตฺถิยา ปญฺจ, นตฺถิยา นว, วิคเต นว, อวิคเต ปญฺจ.}

\vspace{\tipitakaparskip}
\csromancenter{อนุโลมํ.}
\csromanheader{2}{2. ปจฺจนียุทฺธาร}
\proseitem{39}{ปรามาโส ธมฺโม ปรามาสสฺส ธมฺมสฺส อารมฺมณปจฺจเยน ปจฺจโย.\csromanspacer{} อุปนิสฺสยปจฺจเยน ปจฺจโย. (1)}

\prose{ปรามาโส ธมฺโม โนปรามาสสฺส ธมฺมสฺส อารมฺมณปจฺจเยน ปจฺจโย.\csromanspacer{} สหชาตปจฺจเยน ปจฺจโย.\csromanspacer{} อุปนิสฺสยปจฺจเยน ปจฺจโย.\csromanspacer{} ปจฺฉาชาตปจฺจเยน ปจฺจโย. (2)}

\prose{\csromanfolio{350}ปรามาโส ธมฺโม ปรามาสสฺส จ โนปรามาสสฺส จ ธมฺมสฺส อารมฺมณปจฺจเยน ปจฺจโย.\csromanspacer{} อุปนิสฺสยปจฺจเยน ปจฺจโย. (3)}

\proseitem{40}{โนปรามาโส ธมฺโม โนปรามาสสฺส ธมฺมสฺส อารมฺมณปจฺจเยน ปจฺจโย.\csromanspacer{} สหชาตปจฺจเยน ปจฺจโย.\csromanspacer{} อุปนิสฺสยปจฺจเยน ปจฺจโย.\csromanspacer{} ปุเรชาตปจฺจเยน ปจฺจโย.\csromanspacer{} ปจฺฉาชาตปจฺจเยน ปจฺจโย.\csromanspacer{} กมฺมปจฺจเยน ปจฺจโย.\csromanspacer{} อาหารปจฺจเยน ปจฺจโย.\csromanspacer{} อินฺทฺริยปจฺจเยน ปจฺจโย. (1)}

\prose{โนปรามาโส ธมฺโม ปรามาสสฺส ธมฺมสฺส อารมฺมณปจฺจเยน ปจฺจโย.\csromanspacer{} สหชาตปจฺจเยน ปจฺจโย.\csromanspacer{} อุปนิสฺสยปจฺจเยน ปจฺจโย.\csromanspacer{} ปุเรชาตปจฺจเยน ปจฺจโย. (2)}

\prose{โนปรามาโส ธมฺโม ปรามาสสฺส จ โนปรามาสสฺส จ ธมฺมสฺส อารมฺมณปจฺจเยน ปจฺจโย.\csromanspacer{} สหชาตปจฺจเยน ปจฺจโย.\csromanspacer{} อุปนิสฺสยปจฺจเยน ปจฺจโย.\csromanspacer{} ปุเรชาตปจฺจเยน ปจฺจโย. (3)}

\proseitem{41}{ปรามาโส จ โนปรามาโส จ ธมฺมา ปรามาสสฺส ธมฺมสฺส อารมฺมณปจฺจเยน ปจฺจโย.\csromanspacer{} อุปนิสฺสยปจฺจเยน ปจฺจโย. (1)}

\prose{ปรามาโส จ โนปรามาโส จ ธมฺมา โนปรามาสสฺส ธมฺมสฺส อารมฺมณปจฺจเยน ปจฺจโย.\csromanspacer{} สหชาตปจฺจเยน ปจฺจโย.\csromanspacer{} อุปนิสฺสยปจฺจเยน ปจฺจโย.\csromanspacer{} ปจฺฉาชาตปจฺจเยน ปจฺจโย. (2)}

\prose{ปรามาโส จ โนปรามาโส จ ธมฺมา ปรามาสสฺส จ โนปรามาสสฺส จ ธมฺมสฺส อารมฺมณปจฺจเยน ปจฺจโย.\csromanspacer{} อุปนิสฺสยปจฺจเยน ปจฺจโย. (3)}

\csromanheader{2}{2. ปจฺจยปจฺจนีย 2. สงฺขฺยาวาร}
\vspace{\tipitakaparskip}
\csromancenter{นเหตุยา นว, น-อารมฺมเณ นว, (สพฺพตฺถ นว.) โน-อวิคเต นว.}
\csromanmatikaanchor{mtk.49}
\csromanmatikaanchor{mtk.50}
\csromantocmarkref{subsection}{51. ปรามฏฺฐทุก}{mtk.49}
\bookmark[dest={mtk.49},level=2]{51. ปรามฏฺฐทุก}
\csromantocmarkref{subsection}{52. ปรามาสสมฺปยุตฺตทุก}{mtk.50}
\bookmark[dest={mtk.50},level=2]{52. ปรามาสสมฺปยุตฺตทุก}
\csromanheader{2}{52. ปรามาสสมฺปยุตฺตทุก \textbf{3. ปจฺจยานุโลมปจฺจนีย}}
\proseitem{43}{\csromanfolio{351}\textbf{เหตุ}ปจฺจยา น-อารมฺมเณ ตีณิ, น-อธิปติยา ตีณิ, น-อนนฺตเร ตีณิ, นสมนนฺตเร ตีณิ, น-อญฺญมญฺเญ เอกํ, น-อุปนิสฺสเย ตีณิ ฯเปฯ นมคฺเค ตีณิ, นสมฺปยุตฺเต เอกํ, นวิปฺปยุตฺเต ตีณิ, โนนตฺถิยา ตีณิ, โนวิคเต ตีณิ.}

\csromanheader{2}{4. ปจฺจยปจฺจนียานุโลม}
\proseitem{44}{นเหตุปจฺจยา อารมฺมเณ นว, อธิปติยา นว, (อนุโลมมาติกา) ฯเปฯ อวิคเต ปญฺจ.}

\nitthitam{ปรามาสทุกํ นิฏฺฐิตํ.}
\csromanheader{2}{51. ปรามฏฺฐทุก 1-7. ปฏิจฺจวาร}
\csromanheader{2}{1. ปจฺจยานุโลม 1. วิภงฺควาร}
\csromanheader{2}{\textbf{เหตุ}}
\proseitem{45}{ปรามฏฺฐํ ธมฺมํ ปฏิจฺจ ปรามฏฺโฐ ธมฺโม อุปฺปชฺชติ เหตุปจฺจยา.\csromanspacer{} ปรามฏฺฐํ เอกํ ขนฺธํ ปฏิจฺจ ตโย ขนฺธา จิตฺตสมุฏฺฐานญฺจ รูปํ ฯเปฯ เทฺว ขนฺเธ ฯเปฯ.\csromanspacer{} ปฏิสนฺธิกฺขเณ ฯเปฯ. (ยาว อชฺฌตฺติกา มหาภูตา.)}

\prose{ปรามฏฺฐํ ธมฺมํ ปฏิจฺจ ฯเปฯ. (ปรามฏฺฐทุกํ, ยถา จูฬนฺตรทุเก โลกิยทุกํ, เอวํ กาตพฺพํ, นินฺนานากรณํ.)}

\nitthitam{ปรามฏฺฐทุกํ นิฏฺฐิตํ.}
\csromanheader{2}{52. ปรามาสสมฺปยุตฺตทุก 1. ปฏิจฺจวาร}
\csromanheader{2}{1. ปจฺจยานุโลม 1. วิภงฺควาร}
\csromanheader{2}{\textbf{เหตุ}}
\proseitem{46}{ปรามาสสมฺปยุตฺตํ ธมฺมํ ปฏิจฺจ ปรามาสสมฺปยุตฺโต ธมฺโม อุปฺปชฺชติ เหตุปจฺจยา.\csromanspacer{} ปรามาสสมฺปยุตฺตํ เอกํ ขนฺธํ ปฏิจฺจ ตโย ขนฺธา ฯเปฯ เทฺว ขนฺเธ ฯเปฯ. (1)}

\prose{\csromanfolio{352}ปรามาสสมฺปยุตฺตํ ธมฺมํ ปฏิจฺจ ปรามาสวิปฺปยุตฺโต ธมฺโม อุปฺปชฺชติ เหตุปจฺจยา.\csromanspacer{} ปรามาสสมฺปยุตฺเต ขนฺเธ ปฏิจฺจ จิตฺตสมุฏฺฐานํ รูปํ. (2)}

\prose{ปรามาสสมฺปยุตฺตํ ธมฺมํ ปฏิจฺจ ปรามาสสมฺปยุตฺโต จ ปรามาสวิปฺปยุตฺโต จ ธมฺมา อุปฺปชฺชนฺติ เหตุปจฺจยา.\csromanspacer{} ปรามาสสมฺปยุตฺตํ เอกํ ขนฺธํ ปฏิจฺจ ตโย ขนฺธา จิตฺตสมุฏฺฐานญฺจ รูปํ.\csromanspacer{} เทฺว ขนฺเธ ฯเปฯ. (3)}

\proseitem{47}{ปรามาสวิปฺปยุตฺตํ ธมฺมํ ปฏิจฺจ ปรามาสวิปฺปยุตฺโต ธมฺโม อุปฺปชฺชติ เหตุปจฺจยา.\csromanspacer{} ปรามาสวิปฺปยุตฺตํ เอกํ ขนฺธํ ปฏิจฺจ ตโย ขนฺธา จิตฺตสมุฏฺฐานญฺจ รูปํ ฯเปฯ เทฺว ขนฺเธ ฯเปฯ.\csromanspacer{} ปฏิสนฺธิกฺขเณ ฯเปฯ.\csromanspacer{} ขนฺเธ ปฏิจฺจ วตฺถุ, วตฺถุํ ปฏิจฺจ ขนฺธา.\csromanspacer{} เอกํ มหาภูตํ ฯเปฯ. (1)}

\prose{ปรามาสสมฺปยุตฺตญฺจ ปรามาสวิปฺปยุตฺตญฺจ ธมฺมํ ปฏิจฺจ ปรามาสวิปฺปยุตฺโต ธมฺโม อุปฺปชฺชติ เหตุปจฺจยา.\csromanspacer{} ปรามาสสมฺปยุตฺเต ขนฺเธ จ มหาภูเต จ ปฏิจฺจ จิตฺตสมุฏฺฐานํ รูปํ. (สํขิตฺตํ.) (1)}

\csromanheader{2}{1. ปจฺจยานุโลม 2. สงฺขฺยาวาร}
\csromanheader{2}{\textbf{สุทฺธ}}
\proseitem{48}{เหตุยา ปญฺจ, อารมฺมเณ เทฺว, อธิปติยา ปญฺจ, อนนฺตเร เทฺว, สมนนฺตเร เทฺว, สหชาเต ปญฺจ, อญฺญมญฺเญ เทฺว, นิสฺสเย ปญฺจ, อุปนิสฺสเย เทฺว, ปุเรชาเต เทฺว, อาเสวเน เทฺว, กมฺเม ปญฺจ, วิปาเก เอกํ, อาหาเร ปญฺจ ฯเปฯ มคฺเค ปญฺจ, สมฺปยุตฺเต เทฺว, วิปฺปยุตฺเต ปญฺจ, อตฺถิยา ปญฺจ, นตฺถิยา เทฺว, วิคเต เทฺว, อวิคเต ปญฺจ.}

\vspace{\tipitakaparskip}
\csromancenter{อนุโลมํ.}
\csromanheader{2}{2. ปจฺจยปจฺจนีย 1. วิภงฺควาร}
\csromanheader{2}{\textbf{นเหตุ}}
\proseitem{49}{ปรามาสวิปฺปยุตฺตํ ธมฺมํ ปฏิจฺจ ปรามาสวิปฺปยุตฺโต ธมฺโม อุปฺปชฺชติ นเหตุปจฺจยา.\csromanspacer{} อเหตุกํ ปรามาสวิปฺปยุตฺตํ เอกํ ขนฺธํ ปฏิจฺจ ตโย}

\vspace{\tipitakaparskip}
\csromancenter{ปรามาสสมฺปยุตฺตทุก ขนฺธา จิตฺตสมุฏฺฐานญฺจ รูปํ ฯเปฯ เทฺว ขนฺเธ ฯเปฯ.\csromanspacer{} อเหตุกปฏิสนฺธิกฺขเณ ฯเปฯ. (ยาว อสญฺญสตฺตา.) วิจิกิจฺฉาสหคเต อุทฺธจฺจสหคเต ขนฺเธ ปฏิจฺจ วิจิกิจฺฉาสหคโต อุทฺธจฺจสหคโต โมโห. (1)}
\csromanheader{2}{\textbf{น-อารมฺมณ}}
\proseitem{50}{\csromanfolio{353}ปรามาสสมฺปยุตฺตํ ธมฺมํ ปฏิจฺจ ปรามาสวิปฺปยุตฺโต ธมฺโม อุปฺปชฺชติ น-อารมฺมณปจฺจยา.\csromanspacer{} ปรามาสสมฺปยุตฺเต ขนฺเธ ปฏิจฺจ จิตฺตสมุฏฺฐานํ รูปํ. (สํขิตฺตํ.) (1)}

\prose{ปรามาสวิปฺปยุตฺตํ ธมฺมํ ปฏิจฺจ ปรามาสวิปฺปยุตฺโต ธมฺโม อุปฺปชฺชติ น-อารมฺมณปจฺจยา.\csromanspacer{} ปรามาสวิปฺปยุตฺเต ขนฺเธ ปฏิจฺจ จิตฺตสมุฏฺฐานํ รูปํ. (ยาว อสญฺญสตฺตา.) (1)}

\prose{ปรามาสสมฺปยุตฺตญฺจ ปรามาสวิปฺปยุตฺตญฺจ ธมฺมํ ปฏิจฺจ ปรามาสวิปฺปยุตฺโต ธมฺโม อุปฺปชฺชติ น-อารมฺมณปจฺจยา.\csromanspacer{} ปรามาสสมฺปยุตฺเต ขนฺเธ จ มหาภูเต จ ปฏิจฺจ จิตฺตสมุฏฺฐานํ รูปํ. (สํขิตฺตํ.) (1)}

\csromanheader{2}{2. ปจฺจยปจฺจนีย 2. สงฺขฺยาวาร}
\csromanheader{2}{\textbf{สุทฺธ}}
\proseitem{51}{นเหตุยา เอกํ, น-อารมฺมเณ ตีณิ, น-อธิปติยา ปญฺจ, น-อนนฺตเร ตีณิ, นสมนนฺตเร ตีณิ, น-อญฺญมญฺเญ ตีณิ, น-อุปนิสฺสเย ตีณิ, นปุเรชาเต จตฺตาริ, นปจฺฉาชาเต ปญฺจ, น-อาเสวเน ปญฺจ, นกมฺเม เทฺว, นวิปาเก ปญฺจ, น-อาหาเร เอกํ, น-อินฺทฺริเย เอกํ, นฌาเน เอกํ, นมคฺเค เอกํ, นสมฺปยุตฺเต ตีณิ, นวิปฺปยุตฺเต เทฺว, โนนตฺถิยา ตีณิ, โนวิคเต ตีณิ.}

\csromanheader{2}{3. ปจฺจยานุโลมปจฺจนีย}
\proseitem{52}{เหตุปจฺจยา น-อารมฺมเณ ตีณิ, น-อธิปติยา ปญฺจ. (เอวํ คเณตพฺพํ.)}

\csromanheader{2}{4. ปจฺจยปจฺจนียานุโลม}
\vspace{\tipitakaparskip}
\csromancenter{นเหตุปจฺจยา อารมฺมเณ เอกํ. (สํขิตฺตํ) อวิคเต เอกํ.}
\csromancenter{(สหชาตวาโรปิ ปฏิจฺจวารสทิโส.)}
\csromanheader{2}{52. ปรามาสสมฺปยุตฺตทุก 3. ปจฺจยวาร}
\csromanheader{2}{1. ปจฺจยานุโลม 1. วิภงฺควาร}
\csromanheader{2}{\textbf{เหตุ}}
\proseitem{54}{\csromanfolio{354}ปรามาสสมฺปยุตฺตํ ธมฺมํ ปจฺจยา ปรามาสสมฺปยุตฺโต ธมฺโม อุปฺปชฺชติ เหตุปจฺจยา.\csromanspacer{} ตีณิ. (ปฏิจฺจวารสทิโส.)}

\prose{ปรามาสวิปฺปยุตฺตํ ธมฺมํ ปจฺจยา ปรามาสวิปฺปยุตฺโต ธมฺโม อุปฺปชฺชติ เหตุปจฺจยา.\csromanspacer{} ปรามาสวิปฺปยุตฺตํ เอกํ ขนฺธํ ปจฺจยา ตโย ขนฺธา จิตฺตสมุฏฺฐานญฺจ รูปํ ฯเปฯ เทฺว ขนฺเธ ฯเปฯ.\csromanspacer{} ปฏิสนฺธิกฺขเณ ฯเปฯ. (ยาว อชฺฌตฺติกา มหาภูตา.) วตฺถุํ ปจฺจยา ปรามาสวิปฺปยุตฺตา ขนฺธา. (1)}

\prose{ปรามาสวิปฺปยุตฺตํ ธมฺมํ ปจฺจยา ปรามาสสมฺปยุตฺโต ธมฺโม อุปฺปชฺชติ เหตุปจฺจยา.\csromanspacer{} วตฺถุํ ปจฺจยา ปรามาสสมฺปยุตฺตกา ขนฺธา. (2)}

\prose{ปรามาสวิปฺปยุตฺตํ ธมฺมํ ปจฺจยา ปรามาสสมฺปยุตฺโต จ ปรามาสวิปฺปยุตฺโต จ ธมฺมา อุปฺปชฺชนฺติ เหตุปจฺจยา.\csromanspacer{} วตฺถุํ ปจฺจยา ปรามาสสมฺปยุตฺตกา ขนฺธา.\csromanspacer{} มหาภูเต ปจฺจยา จิตฺตสมุฏฺฐานํ รูปํ. (3)}

\proseitem{55}{ปรามาสสมฺปยุตฺตญฺจ ปรามาสวิปฺปยุตฺตญฺจ ธมฺมํ ปจฺจยา ปรามาสสมฺปยุตฺโต ธมฺโม อุปฺปชฺชติ เหตุปจฺจยา.\csromanspacer{} ปรามาสสมฺปยุตฺตํ เอกํ ขนฺธญฺจ วตฺถุญฺจ ปจฺจยา ตโย ขนฺธา ฯเปฯ เทฺว ขนฺเธ ฯเปฯ. (1)}

\prose{ปรามาสสมฺปยุตฺตญฺจ ปรามาสวิปฺปยุตฺตญฺจ ธมฺมํ ปจฺจยา ปรามาสวิปฺปยุตฺโต ธมฺโม อุปฺปชฺชติ เหตุปจฺจยา.\csromanspacer{} ปรามาสสมฺปยุตฺเต ขนฺเธ จ มหาภูเต จ ปจฺจยา จิตฺตสมุฏฺฐานํ รูปํ. (2)}

\prose{ปรามาสสมฺปยุตฺตญฺจ ปรามาสวิปฺปยุตฺตญฺจ ธมฺมํ ปจฺจยา ปรามาสสมฺปยุตฺโต จ ปรามาสวิปฺปยุตฺโต จ ธมฺมา อุปฺปชฺชนฺติ เหตุปจฺจยา.\csromanspacer{} ปรามาสสมฺปยุตฺตํ เอกํ ขนฺธญฺจ วตฺถุญฺจ ปจฺจยา ตโย ขนฺธา ฯเปฯ เทฺว ขนฺเธ ฯเปฯ.\csromanspacer{} ปรามาสสมฺปยุตฺเต ขนฺเธ จ มหาภูเต จ ปจฺจยา จิตฺตสมุฏฺฐานํ รูปํ. (สํขิตฺตํ.) (3)}

\csromanheader{2}{52. ปรามาสสมฺปยุตฺตทุก \textbf{1. ปจฺจยานุโลม 2. สงฺขฺยาวาร}}
\csromanheader{2}{\textbf{สุทฺธ}}
\proseitem{56}{\csromanfolio{355}เหตุยา นว, อารมฺมเณ จตฺตาริ, อธิปติยา นว, อนนฺตเร จตฺตาริ, สมนนฺตเร จตฺตาริ, สหชาเต นว, อญฺญมญฺเญ จตฺตริ, นิสฺสเย นว, อุปนิสฺสเย จตฺตาริ, ปุเรชาเต จตฺตาริ, อาเสวเน จตฺตาริ, กมฺเม นว, วิปาเก เอกํ, อาหาเร นว ฯเปฯ อวิคเต นว.}

\vspace{\tipitakaparskip}
\csromancenter{อนุโลมํ.}
\csromanheader{2}{2. ปจฺจยปจฺจนีย 1. วิภงฺควาร}
\proseitem{57}{ปรามาสวิปฺปยุตฺตํ ธมฺมํ ปจฺจยา ปรามาสวิปฺปยุตฺโต ธมฺโม อุปฺปชฺชติ นเหตุปจฺจยา.\csromanspacer{} อเหตุกํ ปรามาสวิปฺปยุตฺตํ เอกํ ขนฺธํ ปจฺจยา ตโย ขนฺธา จิตฺตสมุฏฺฐานญฺจ รูปํ ฯเปฯ เทฺว ขนฺเธ ฯเปฯ.\csromanspacer{} อเหตุกปฏิสนฺธิกฺขเณ ฯเปฯ. (ยาว อสญฺญสตฺตา.) จกฺขายตนํ ปจฺจยา จกฺขุ วิญฺญาณํ ฯเปฯ.\csromanspacer{} กายายตนํ ปจฺจยา กายวิญฺญาณํ.\csromanspacer{} วตฺถุํ ปจฺจยา อเหตุกปรามาสวิปฺปยุตฺตา ขนฺธา.\csromanspacer{} วิจิกิจฺฉาสหคเต อุทฺธจฺจสหคเต ขนฺเธ จ วตฺถุญฺจ ปจฺจยา วิจิกิจฺฉาสหคโต อุทฺธจฺจสหคโต โมโห.}

\csromanheader{2}{2. ปจฺจยปจฺจนีย 2. สงฺขฺยาวาร}
\csromanheader{2}{\textbf{สุทฺธ}}
\proseitem{58}{นเหตุยา เอกํ, น-อารมฺมเณ ตีณิ, น-อธิปติยา นว, น-อนนฺตเร ตีณิ, นสมนนฺตเร ตีณิ, น-อญฺญมญฺเญ ตีณิ, น-อุปนิสฺสเย ตีณิ, นปุเรชาเต จตฺตาริ, นปจฺฉาชาเต นว, น-อาเสวเน นว, นกมฺเม จตฺตาริ, นวิปาเก นว, น-อาหาเร เอกํ, น-อินฺทฺริเย เอกํ, นฌาเน เอกํ, นมคฺเค เอกํ, นสมฺปยุตฺเต ตีณิ, นวิปฺปยุตฺเต เทฺว, โนนตฺถิยา ตีณิ โนวิคเต ตีณิ.}

\csromanheader{2}{3. ปจฺจยานุโลมปจฺจนีย}
\proseitem{59}{\csromanfolio{356}\textbf{เหตุ}ปจฺจยา น-อารมฺมเณ ตีณิ, น-อธิปติยา นว. (เอวํ คเณตพฺพํ.)}

\csromanheader{2}{4. ปจฺจยปจฺจนียานุโลม}
\vspace{\tipitakaparskip}
\csromancenter{นเหตุปจฺจยา อารมฺมเณ เอกํ ฯเปฯ อวิคเต เอกํ.}
\csromancenter{(นิสฺสยวาโร ปจฺจยวารสทิโส.)}
\csromanheader{2}{52. ปรามาสสมฺปยุตฺตทุก 5. สํสฏฺฐวาร}
\csromanheader{2}{1-4. ปจฺจยานุโลมาทิ}
\proseitem{61}{ปรามาสสมฺปยุตฺตํ ธมฺมํ สํสฏฺโฐ ปรามาสสมฺปยุตฺโต ธมฺโม อุปฺปชฺชติ เหตุปจฺจยา. (สํขิตฺตํ.)}

\csromanheader{2}{\textbf{สุทฺธ}}
\proseitem{62}{\textbf{เหตุ}ยา เทฺว, (สพฺพตฺถ เทฺว,) วิปาเก เอกํ ฯเปฯ อวิคเต เทฺว.}

\prose{นเหตุยา เอกํ, น-อธิปติยา เทฺว, นปุเรชาเต เทฺว, นปจฺฉาชาเต เทฺว, น-อาเสวเน เทฺว, นกมฺเม เทฺว, นวิปาเก เทฺว, นฌาเน เอกํ, นมคฺเค เอกํ, นวิปฺปยุตฺเต เทฺว.}

\vspace{\tipitakaparskip}
\csromancenter{(เอวํ อิตเร เทฺว คณนาปิ สมฺปยุตฺตวาราปิ กาตพฺพา.)}
\csromanheader{2}{52. ปรามาสสมฺปยุตฺตทุก 7. ปญฺหาวาร}
\csromanheader{2}{1. ปจฺจยานุโลม 1. วิภงฺควาร}
\csromanheader{2}{\textbf{เหตุ}}
\proseitem{63}{ปรามาสสมฺปยุตฺโต ธมฺโม ปรามาสสมฺปยุตฺตสฺส ธมฺมสฺส เหตุปจฺจเยน ปจฺจโย.\csromanspacer{} ปรามาสสมฺปยุตฺตา เหตู สมฺปยุตฺตกานํ}

\vspace{\tipitakaparskip}
\csromancenter{ปรามาสสมฺปยุตฺตทุก ขนฺธานํ เหตุปจฺจเยน ปจฺจโย. (มูลํ กาตพฺพํ.) ปรามาสสมฺปยุตฺตา เหตู จิตฺตสมุฏฺฐานานํ รูปานํ เหตุปจฺจเยน ปจฺจโย. (มูลํ กาตพฺพํ.) ปรามาสสมฺปยุตฺตา เหตู สมฺปยุตฺตกานํ ขนฺธานํ จิตฺตสมุฏฺฐานานญฺจ รูปานํ เหตุปจฺจเยน ปจฺจโย. (3)}
\prose{\csromanfolio{357}ปรามาสวิปฺปยุตฺโต ธมฺโม ปรามาสวิปฺปยุตฺตสฺส ธมฺมสฺส เหตุปจฺจเยน ปจฺจโย.\csromanspacer{} ปรามาสวิปฺปยุตฺโต เหตู สมฺปยุตฺตกานํ ขนฺธานํ จิตฺตสมุฏฺฐานานญฺจ รูปานํ เหตุปจฺจเยน ปจฺจโย.\csromanspacer{} ปฏิสนฺธิกฺขเณ ฯเปฯ. (1)}

\csromanheader{2}{\textbf{อารมฺมณ}}
\proseitem{64}{ปรามาสสมฺปยุตฺโต ธมฺโม ปรามาสสมฺปยุตฺตสฺส ธมฺมสฺส อารมฺมณปจฺจเยน ปจฺจโย.\csromanspacer{} ราคํ อสฺสาเทติ อภินนฺทติ, ตํ อารพฺภ ปรามาสสมฺปยุตฺโต ราโค อุปฺปชฺชติ.\csromanspacer{} ปรามาสสมฺปยุตฺเต ขนฺเธ อสฺสาเทติ อภินนฺทติ, ตํ อารพฺภ ปรามาสสมฺปยุตฺโต ราโค อุปฺปชฺชติ. (1)}

\prose{ปรามาสสมฺปยุตฺโต ธมฺโม ปรามาสวิปฺปยุตฺตสฺส ธมฺมสฺส อารมฺมณปจฺจเยน ปจฺจโย.\csromanspacer{} อริยา ปรามาสสมฺปยุตฺเต ปหีเน กิเลเส ปจฺจเวกฺขนฺติ, ปุพฺเพ สมุทาจิณฺเณ กิเลเส ชานนฺติ.\csromanspacer{} ปรามาสสมฺปยุตฺเต ขนฺเธ อนิจฺจโต ฯเปฯ วิปสฺสติ, อสฺสาเทติ อภินนฺทติ, ตํ อารพฺภ ปรามาสวิปฺปยุตฺโต ราโค ฯเปฯ.\csromanspacer{} วิจิกิจฺฉา ฯเปฯ อุทฺธจฺจํ ฯเปฯ โทมนสฺสํ อุปฺปชฺชติ.\csromanspacer{} เจโตปริยญาเณน ปรามาสสมฺปยุตฺตจิตฺตสมงฺคิสฺส จิตฺตํ ชานาติ.\csromanspacer{} ปรามาสสมฺปยุตฺตา ขนฺธา เจโตปริยญาณสฺส, ปุพฺเพนิวาสานุสฺสติญาณสฺส, ยถากมฺมูปคญาณสฺส, อนาคตํสญาณสฺส, อาวชฺชนาย อารมฺมณปจฺจเยน ปจฺจโย. (2)}

\proseitem{65}{ปรามาสวิปฺปยุตฺโต ธมฺโม ปรามาสวิปฺปยุตฺตสฺส ธมฺมสฺส อารมฺมณปจฺจเยน ปจฺจโย.\csromanspacer{} ทานํ ฯเปฯ สีลํ ฯเปฯ อุโปสถกมฺมํ กตฺวา ตํ ปจฺจเวกฺขติ, อสฺสาเทติ อภินนฺทติ, ตํ อารพฺภ ปรามาสวิปฺปยุตฺโต ราโค ฯเปฯ วิจิกิจฺฉา ฯเปฯ อุทฺธจฺจํ ฯเปฯ โทมนสฺสํ อุปฺปชฺชติ.\csromanspacer{} ปุพฺเพ สุจิณฺณานิ ฯเปฯ.\csromanspacer{} ฌานา ฯเปฯ.\csromanspacer{} อริยา มคฺคา วุฏฺฐหิตฺวา มคฺคํ ปจฺจเวกฺขนฺติ.\csromanspacer{} ผลํ ฯเปฯ.\csromanspacer{} นิพฺพานํ ฯเปฯ.\csromanspacer{} นิพฺพานํ โคตฺรภุสฺส, โวทานสฺส, มคฺคสฺส, ผลสฺส, อาวชฺชนาย อารมฺมณปจฺจเยน ปจฺจโย.\csromanspacer{} อริยา ปรามาสวิปฺปยุตฺเต ปหีเน กิเลเส ฯเปฯ.\csromanspacer{} วิกฺขมฺภิเต กิเลเส ฯเปฯ.\csromanspacer{} ปุพฺเพ ฯเปฯ.\csromanspacer{} จกฺขุํ ฯเปฯ วตฺถุํ ปรามาสวิปฺปยุตฺเต \csromanfolio{358}ขนฺเธ อนิจฺจโต ฯเปฯ วิปสฺสติ, อสฺสาเทติ อภินนฺทติ, ตํ อารพฺภ ปรามาสวิปฺปยุตฺโต ราโค ฯเปฯ วิจิกิจฺฉา ฯเปฯ อุทฺธจฺจํ ฯเปฯ โทมนสฺสํ อุปฺปชฺชติ.\csromanspacer{} ทิพฺเพน จกฺขุนา รูปํ ปสฺสติ.\csromanspacer{} ทิพฺพาย โสตธาตุยา สทฺทํ สุณาติ.\csromanspacer{} เจโตปริยญาเณน ปรามาสวิปฺปยุตฺตจิตฺตสมงฺคิสฺส จิตฺตํ ชานาติ.\csromanspacer{} อากาสานญฺจายตนํ วิญฺญาณญฺจายตนสฺส ฯเปฯ.\csromanspacer{} อากิญฺจญฺญายตนํ เนวสญฺญานาสญฺญายตนสฺส ฯเปฯ.\csromanspacer{} รูปายตนํ จกฺขุวิญฺญาณสฺส ฯเปฯ.\csromanspacer{} โผฏฺฐพฺพายตนํ กายวิญฺญาณสฺส ฯเปฯ.\csromanspacer{} ปรามาสวิปฺปยุตฺตา ขนฺธา อิทฺธิวิธญาณสฺส, เจโตปริยญาณสฺส, ปุพฺเพนิวาสานุสฺสติญาณสฺส, ยถากมฺมูปคญาณสฺส, อนาคตํสญาณสฺส, อาวชฺชนาย อารมฺมณปจฺจเยน ปจฺจโย. (1)}

\prose{ปรามาสวิปฺปยุตฺโต ธมฺโม ปรามาสสมฺปยุตฺตสฺส ธมฺมสฺส อารมฺมณปจฺจเยน ปจฺจโย.\csromanspacer{} ทานํ ฯเปฯ สีลํ ฯเปฯ อุโปสถกมฺมํ ฯเปฯ.\csromanspacer{} ปุพฺเพ ฯเปฯ.\csromanspacer{} ฌานา ฯเปฯ.\csromanspacer{} จกฺขุํ ฯเปฯ.\csromanspacer{} วตฺถุํ ปรามาสวิปฺปยุตฺเต ขนฺเธ อสฺสาเทติ อภินนฺทติ ตํ อารพฺภ ปรามาสสมฺปยุตฺโต ราโค อุปฺปชฺชติ. (2)}

\csromanheader{2}{\textbf{อธิปติ}}
\proseitem{66}{ปรามาสสมฺปยุตฺโต ธมฺโม ปรามาสสมฺปยุตฺตสฺส ธมฺมสฺส อธิปติปจฺจเยน ปจฺจโย.\csromanspacer{} \textbf{อารมฺมณาธิปติ}, สหชาตาธิปติ.\csromanspacer{}\csromanspacer{}\textbf{อารมฺมณาธิปติ\csromandash{}}ราคํ ครุํ กตฺวา อสฺสาเทติ อภินนฺทติ, ตํ ครุํ กตฺวา ปรามาสสมฺปยุตฺโต ราโค อุปฺปชฺชติ.\csromanspacer{} ปรามาสสมฺปยุตฺเต ขนฺเธ ครุํ กตฺวา อสฺสาเทติ อภินนฺทติ, ตํ ครุํ กตฺวา ปรามาสสมฺปยุตฺโต ราโค อุปฺปชฺชติ.\csromanspacer{} \textbf{สหชาตาธิปติ\csromandash{}}ปรามาสสมฺปยุตฺตาธิปติ สมฺปยุตฺตกานํ ขนฺธานํ อธิปติปจฺจเยน ปจฺจโย. (1)}

\prose{ปรามาสสมฺปยุตฺโต ธมฺโม ปรามาสวิปฺปยุตฺตสฺส ธมฺมสฺส อธิปติปจฺจเยน ปจฺจโย.\csromanspacer{} \textbf{อารมฺมณาธิปติ}, สหชาตาธิปติ.\csromanspacer{}\csromanspacer{}\textbf{อารมฺมณาธิปติ\csromandash{}}ปรามาสสมฺปยุตฺเต ขนฺเธ ครุํ กตฺวา อสฺสาเทติ อภินนฺทติ, ตํ ครุํ กตฺวา ปรามาสวิปฺปยุตฺโต ราโค อุปฺปชฺชติ.\csromanspacer{} \textbf{สหชาตาธิปติ\csromandash{}}ปรามาสสมฺปยุตฺตาธิปติ จิตฺตสมุฏฺฐานานํ รูปานํ อธิปติปจฺจเยน ปจฺจโย. (2)}

\prose{ปรามาสสมฺปยุตฺโต ธมฺโม ปรามาสสมฺปยุตฺตสฺส จ ปรามาสวิปฺปยุตฺตสฺส จ ธมฺมสฺส อธิปติปจฺจเยน ปจฺจโย.\csromanspacer{} \textbf{สหชาตาธิปติ\csromandash{}}}

\vspace{\tipitakaparskip}
\csromancenter{ปรามาสสมฺปยุตฺตทุก ปรามาสสมฺปยุตฺตาธิปติ สมฺปยุตฺตกานํ ขนฺธานํ จิตฺตสมุฏฺฐานานญฺจ รูปานํ อธิปติปจฺจเยน ปจฺจโย. (3)}
\proseitem{67}{\csromanfolio{359}ปรามาสวิปฺปยุตฺโต ธมฺโม ปรามาสวิปฺปยุตฺตสฺส ธมฺมสฺส อธิปติปจฺจเยน ปจฺจโย.\csromanspacer{} \textbf{อารมฺมณาธิปติ}, สหชาตาธิปติ.\csromanspacer{}\csromanspacer{}\textbf{อารมฺมณาธิปติ}\csromandash{}ทานํ ฯเปฯ สีลํ ฯเปฯ อุโปสถกมฺมํ กตฺวา ตํ ครุํ กตฺวา ปจฺจเวกฺขติ, อสฺสาเทติ อภินนฺทติ, ตํ ครุํ กตฺวา ปรามาสวิปฺปยุตฺโต ราโค อุปฺปชฺชติ.\csromanspacer{} ปุพฺเพ ฯเปฯ.\csromanspacer{} ฌานา ฯเปฯ.\csromanspacer{} อริยา มคฺคา วุฏฺฐหิตฺวา มคฺคํ ครุํ กตฺวา ฯเปฯ.\csromanspacer{} ผลํ ฯเปฯ นิพฺพานํ ฯเปฯ นิพฺพานํ โคตฺรภุสฺส, โวทานสฺส, มคฺคสฺส, ผลสฺส อธิปติปจฺจเยน ปจฺจโย.\csromanspacer{} จกฺขุํ ฯเปฯ วตฺถุํ ปรามาสวิปฺปยุตฺเต ขนฺเธ ครุํ กตฺวา อสฺสาเทติ อภินนฺทติ, ตํ ครุํ กตฺวา ปรามาสวิปฺปยุตฺโต ราโค อุปฺปชฺชติ.\csromanspacer{} \textbf{สหชาตาธิปติ}\csromandash{}ปรามาสวิปฺปยุตฺตาธิปติ สมฺปยุตฺตกานํ ขนฺธานํ จิตฺตสมุฏฺฐานานญฺจ รูปานํ อธิปติปจฺจเยน ปจฺจโย. (1)}

\prose{ปรามาสวิปฺปยุตฺโต ธมฺโม ปรามาสสมฺปยุตฺตสฺส ธมฺมสฺส อธิปติปจฺจเยน ปจฺจโย.\csromanspacer{} \textbf{อารมฺมณาธิปติ}\csromandash{}ทานํ ฯเปฯ สีลํ ฯเปฯ อุโปสถกมฺมํ กตฺวา ตํ ครุํ กตฺวา อสฺสาเทติ อภินนฺทติ, ตํ ครุํ กตฺวา ปรามาสสมฺปยุตฺโต ราโค อุปฺปชฺชติ.\csromanspacer{} ปุพฺเพ ฯเปฯ.\csromanspacer{} ฌานา ฯเปฯ.\csromanspacer{} จกฺขุํ ฯเปฯ วตฺถุํ ปรามาสวิปฺปยุตฺเต ขนฺเธ ครุํ กตฺวา อสฺสาเทติ อภินนฺทติ, ตํ ครุํ กตฺวา ปรามาสสมฺปยุตฺโต ราโค อุปฺปชฺชติ. (2)}

\csromanheader{2}{\textbf{อนนฺตรสมนนฺตร}}
\proseitem{68}{ปรามาสสมฺปยุตฺโต ธมฺโม ปรามาสสมฺปยุตฺตสฺส ธมฺมสฺส อนนฺตรปจฺจเยน ปจฺจโย.\csromanspacer{} ปุริมา ปุริมา ปรามาสสมฺปยุตฺตา ขนฺธา ปจฺฉิมานํ ปจฺฉิมานํ ปรามาสสมฺปยุตฺตกานํ ขนฺธานํ อนนฺตรปจฺจเยน ปจฺจโย. (1)}

\prose{ปรามาสสมฺปยุตฺโต ธมฺโม ปรามาสวิปฺปยุตฺตสฺส ธมฺมสฺส อนนฺตรปจฺจเยน ปจฺจโย.\csromanspacer{} ปรามาสสมฺปยุตฺตา ขนฺธา วุฏฺฐานสฺส อนนฺตรปจฺจเยน ปจฺจโย. (2)}

\proseitem{69}{ปรามาสวิปฺปยุตฺโต ธมฺโม ปรามาสวิปฺปยุตฺตสฺส ธมฺมสฺส อนนฺตรปจฺจเยน ปจฺจโย.\csromanspacer{} ปุริมา ปุริมา ปรามาสวิปฺปยุตฺตา ขนฺธา \csromanfolio{360}ปจฺฉิมานํ ปจฺฉิมานํ ปรามาสวิปฺปยุตฺตานํ ขนฺธานํ อนนฺตรปจฺจเยน ปจฺจโย.\csromanspacer{} อนุโลมํ ฯเปฯ ผลสมาปตฺติยา อนนฺตรปจฺจเยน ปจฺจโย. (1)}

\prose{ปรามาสวิปฺปยุตฺโต ธมฺโม ปรามาสสมฺปยุตฺตสฺส ธมฺมสฺส อนนฺตรปจฺจเยน ปจฺจโย.\csromanspacer{} อาวชฺชนา ปรามาสสมฺปยุตฺตกานํ ขนฺธานํ อนนฺตรปจฺจเยน ปจฺจโย.\csromanspacer{} สมนนฺตรปจฺจเยน ปจฺจโย.}

\csromanheader{2}{\textbf{สหชาตาทิ}}
\proseitem{70}{ปรามาสสมฺปยุตฺโต ธมฺโม ปรามาสสมฺปยุตฺตสฺส ธมฺมสฺส สหชาตปจฺจเยน ปจฺจโย.\csromanspacer{} ปญฺจ.\csromanspacer{} อญฺญมญฺญปจฺจเยน ปจฺจโย.\csromanspacer{} เทฺว.\csromanspacer{} นิสฺสยปจฺจเยน ปจฺจโย.\csromanspacer{} สตฺต.}

\csromanheader{2}{\textbf{อุปนิสฺสย}}
\proseitem{71}{ปรามาสสมฺปยุตฺโต ธมฺโม ปรามาสสมฺปยุตฺตสฺส ธมฺมสฺส อุปนิสฺสยปจฺจเยน ปจฺจโย.\csromanspacer{} อารมฺมณูปนิสฺสโย, อนนฺตรูปนิสฺสโย, ปกตูปนิสฺสโย ฯเปฯ.\csromanspacer{} ปกตูปนิสฺสโย\csromandash{}ปรามาสสมฺปยุตฺโต ราโค.\csromanspacer{} โมโห.\csromanspacer{} ปตฺถนา ปรามาสสมฺปยุตฺตสฺส ราคสฺส.\csromanspacer{} โมหสฺส.\csromanspacer{} ปตฺถนาย อุปนิสฺสยปจฺจเยน ปจฺจโย. (1)}

\prose{ปรามาสสมฺปยุตฺโต ธมฺโม ปรามาสวิปฺปยุตฺตสฺส ธมฺมสฺส อุปนิสฺสยปจฺจเยน ปจฺจโย.\csromanspacer{} อารมฺมณูปนิสฺสโย, อนนฺตรูปนิสฺสโย, ปกตูปนิสฺสโย ฯเปฯ.\csromanspacer{} ปกตูปนิสฺสโย\csromandash{}ปรามาสสมฺปยุตฺตํ ราคํ อุปนิสฺสาย ทานํ เทติ ฯเปฯ สมาปตฺติํ อุปฺปาเทติ.\csromanspacer{} มานํ ชปฺเปติ.\csromanspacer{} ปรามาสสมฺปยุตฺตํ โมหํ.\csromanspacer{} ปตฺถนํ อุปนิสฺสาย ทานํ เทติ ฯเปฯ สมาปตฺติํ อุปฺปาเทติ.\csromanspacer{} มานํ ชปฺเปติ.\csromanspacer{} ปาณํ หนติ ฯเปฯ สํฆํ ภินฺทติ.\csromanspacer{} ปรามาสสมฺปยุตฺโต ราโค โมโห.\csromanspacer{} ปตฺถนา สทฺธาย ฯเปฯ.\csromanspacer{} ปญฺญาย.\csromanspacer{} ราคสฺส.\csromanspacer{} โทสสฺส.\csromanspacer{} โมหสฺส.\csromanspacer{} มานสฺส.\csromanspacer{} ปตฺถนาย.\csromanspacer{} กายิกสฺส สุขสฺส ฯเปฯ ผลสมาปตฺติยา อุปนิสฺสยปจฺจเยน ปจฺจโย. (2)}

\proseitem{72}{ปรามาสวิปฺปยุตฺโต ธมฺโม ปรามาสวิปฺปยุตฺตสฺส ธมฺมสฺส อุปนิสฺสยปจฺจเยน ปจฺจโย.\csromanspacer{} อารมฺมณูปนิสฺสโย, อนนฺตรูปนิสฺสโย, ปกตูปนิสฺสโย ฯเปฯ.\csromanspacer{} ปกตูปนิสฺสโย\csromandash{}สทฺธํ อุปนิสฺสาย ทานํ}

\vspace{\tipitakaparskip}
\csromancenter{ปรามาสสมฺปยุตฺตทุก เทติ ฯเปฯ สมาปตฺติํ อุปฺปาเทติ.\csromanspacer{} มานํ ชปฺเปติ.\csromanspacer{} สีลํ ฯเปฯ.\csromanspacer{} ปญฺญํ.\csromanspacer{} ราคํ ฯเปฯ.\csromanspacer{} มานํ ฯเปฯ.\csromanspacer{} ปตฺถนํ อุปนิสฺสาย ทานํ เทติ ฯเปฯ สมาปตฺติํ อุปฺปาเทติ.\csromanspacer{} ปาณํ หนติ ฯเปฯ สํฆํ ภินฺทติ.\csromanspacer{} กายิกํ สุขํ ฯเปฯ.\csromanspacer{} เสนาสนํ อุปนิสฺสาย ทานํ เทติ ฯเปฯ สํฆํ ภินฺทติ.\csromanspacer{} สทฺธา ฯเปฯ.\csromanspacer{} ปญฺญา.\csromanspacer{} ราโค ฯเปฯ.\csromanspacer{} มาโน.\csromanspacer{} ปตฺถนา.\csromanspacer{} กายิกํ สุขํ ฯเปฯ.\csromanspacer{} เสนาสนํ สทฺธาย ฯเปฯ ปญฺญาย.\csromanspacer{} ราคสฺส ฯเปฯ มานสฺส.\csromanspacer{} ปตฺถนาย.\csromanspacer{} กายิกสฺส สุขสฺส ฯเปฯ ผลสมปตฺติยา อุปนิสฺสยปจฺจเยน ปจฺจโย. (1)}
\prose{\csromanfolio{361}ปรามาสวิปฺปยุตฺโต ธมฺโม ปรามาสสมฺปยุตฺตสฺส ธมฺมสฺส อุปนิสฺสยปจฺจเยน ปจฺจโย.\csromanspacer{} อารมฺมณูปนิสฺสโย, อนนฺตรูปนิสฺสโย, ปกตูปนิสฺสโย ฯเปฯ.\csromanspacer{} ปกตูปนิสฺสโย\csromandash{}สทฺธํ อุปนิสฺสาย ปรามาสสมฺปยุตฺโต ราโค อุปฺปชฺชติ.\csromanspacer{} สีลํ ฯเปฯ.\csromanspacer{} เสนาสนํ อุปนิสฺสาย ปรามาสสมฺปยุตฺโต ราโค อุปฺปชฺชติ.\csromanspacer{} สทฺธา ฯเปฯ.\csromanspacer{} เสนาสนํ ปรามาสสมฺปยุตฺตสฺส ราคสฺส ฯเปฯ ปตฺถนาย อุปนิสฺสยปจฺจเยน ปจฺจโย. (2)}

\csromanheader{2}{\textbf{ปุเรชาต}}
\proseitem{73}{ปรามาสวิปฺปยุตฺโต ธมฺโม ปรามาสวิปฺปยุตฺตสฺส ธมฺมสฺส ปุเรชาตปจฺจเยน ปจฺจโย.\csromanspacer{} \textbf{อารมฺมณปุเรชาตํ}, วตฺถุปุเรชาตํ.\csromanspacer{}\csromanspacer{}\textbf{อารมฺมณปุเรชาตํ}\csromandash{}จกฺขุํ ฯเปฯ วตฺถุํ อนิจฺจโต ฯเปฯ วิปสฺสติ, อสฺสาเทติ อภินนฺทติ, ตํ อารพฺภ ปรามาสวิปฺปยุตฺโต ราโค ฯเปฯ วิจิกิจฺฉา ฯเปฯ อุทฺธจฺจํ ฯเปฯ โทมนสฺสํ อุปฺปชฺชติ.\csromanspacer{} ทิพฺเพน จกฺขุนา รูปํ ปสฺสติ.\csromanspacer{} ทิพฺพาย โสตธาตุยา สทฺทํ สุณาติ.\csromanspacer{} รูปายตนํ จกฺขุวิญฺญาณสฺส ฯเปฯ.\csromanspacer{} โผฏฺฐพฺพายตนํ กายวิญฺญาณสฺส ปุเรชาตปจฺจเยน ปจฺจโย.\csromanspacer{} \textbf{วตฺถุปุเรชาตํ\csromandash{}}จกฺขายตนํ จกฺขุวิญฺญาณสฺส ฯเปฯ กายายตนํ กายวิญฺญาณสฺส ฯเปฯ.\csromanspacer{} วตฺถุ ปรามาสวิปฺปยุตฺตานํ ขนฺธานํ ปุเรชาตปจฺจเยน ปจฺจโย. (1)}

\prose{ปรามาสวิปฺปยุตฺโต ธมฺโม ปรามาสสมฺปยุตฺตสฺส ธมฺมสฺส ปุเรชาตปจฺจเยน ปจฺจโย.\csromanspacer{} \textbf{อารมฺมณปุเรชาตํ}, วตฺถุปุเรชาตํ.\csromanspacer{}\csromanspacer{}\textbf{อารมฺมณปุเรชาตํ}\csromandash{}จกฺขุํ ฯเปฯ วตฺถุํ อสฺสาเทติ อภินนฺทติ, ตํ อารพฺภ ปรามาสสมฺปยุตฺโต ราโค อุปฺปชฺชติ.\csromanspacer{} \textbf{วตฺถุปุเรชาตํ\csromandash{}}วตฺถุ ปรามาสสมฺปยุตฺตกานํ ขนฺธานํ ปุเรชาตปจฺจเยน ปจฺจโย. (2)}

\csromanheader{2}{\textbf{ปจฺฉาชาตาเสวน}}
\proseitem{74}{\csromanfolio{362}ปรามาสสมฺปยุตฺโต ธมฺโม ปรามาสวิปฺปยุตฺตสฺส ธมฺมสฺส ปจฺฉาชาตปจฺจเยน ปจฺจโย. (สํขิตฺตํ.) (1)}

\prose{ปรามาสวิปฺปยุตฺโต ธมฺโม ปรามาสวิปฺปยุตฺตสฺส ธมฺมสฺส ปจฺฉาชาตปจฺจเยน ปจฺจโย. (สํขิตฺตํ.) (1)}

\prose{ปรามาสสมฺปยุตฺโต ธมฺโม ปรามาสสมฺปยุตฺตสฺส ธมฺมสฺส อาเสวนปจฺจเยน ปจฺจโย.\csromanspacer{} เทฺว.}

\csromanheader{2}{\textbf{กมฺมาทิ}}
\proseitem{75}{ปรามาสสมฺปยุตฺโต ธมฺโม ปรามาสสมฺปยุตฺตสฺส ธมฺมสฺส กมฺมปจฺจเยน ปจฺจโย.\csromanspacer{} ปรามาสสมฺปยุตฺตา เจตนา สมฺปยุตฺตกานํ ขนฺธานํ กมฺมปจฺจเยน ปจฺจโย. (1)}

\prose{ปรามาสสมฺปยุตฺโต ธมฺโม ปรามาสวิปฺปยุตฺตสฺส ธมฺมสฺส กมฺมปจฺจเยน ปจฺจโย.\csromanspacer{} \textbf{สหชาตา}, นานากฺขณิกา.\csromanspacer{}\csromanspacer{}\textbf{สหชาตา} ปรามาสสมฺปยุตฺตา เจตนา จิตฺตสมุฏฺฐานานํ รูปานํ กมฺมปจฺจเยน ปจฺจโย.\csromanspacer{} \textbf{นานากฺขณิกา} ปรามาสสมฺปยุตฺตา เจตนา วิปากานํ ขนฺธานํ กฏตฺตา จ รูปานํ กมฺมปจฺจเยน ปจฺจโย. (มูลํ กาตพฺพํ.) ปรามาสสมฺปยุตฺตา เจตนา สมฺปยุตฺตกานํ ขนฺธานํ จิตฺตสมุฏฺฐานานญฺจ รูปานํ กมฺมปจฺจเยน ปจฺจโย. (3)}

\prose{ปรามาสวิปฺปยุตฺโต ธมฺโม ปรามาสวิปฺปยุตฺตสฺส ธมฺมสฺส กมฺมปจฺจเยน ปจฺจโย.\csromanspacer{} \textbf{สหชาตา}, นานากฺขณิกา.\csromanspacer{}\csromanspacer{}\textbf{สหชาตา} ปรามาสวิปฺปยุตฺตา เจตนา สมฺปยุตฺตกานํ ขนฺธานํ จิตฺตสมุฏฺฐานานญฺจ รูปานํ กมฺมปจฺจเยน ปจฺจโย.\csromanspacer{} ปฏิสนฺธิกฺขเณ ฯเปฯ.\csromanspacer{} \textbf{นานากฺขณิกา} ปรามาสวิปฺปยุตฺตา เจตนา วิปากานํ ขนฺธานํ กฏตฺตา จ รูปานํ กมฺมปจฺจเยน ปจฺจโย. (1)}

\prose{วิปากปจฺจเยน ปจฺจโย.\csromanspacer{} เอกํ.\csromanspacer{} อาหารปจฺจเยน ปจฺจโย.\csromanspacer{} จตฺตาริ.\csromanspacer{} อินฺทฺริยปจฺจเยน ปจฺจโย.\csromanspacer{} จตฺตาริ.\csromanspacer{} ฌานปจฺจเยน ปจฺจโย.\csromanspacer{} จตฺตาริ.\csromanspacer{} มคฺคปจฺจเยน ปจฺจโย.\csromanspacer{} จตฺตาริ.\csromanspacer{} สมฺปยุตฺตปจฺจเยน ปจฺจโย.\csromanspacer{} เทฺว.}

\csromanheader{2}{52. ปรามาสสมฺปยุตฺตทุก \textbf{วิปฺปยุตฺต}}
\proseitem{76}{\csromanfolio{363}ปรามาสสมฺปยุตฺโต ธมฺโม ปรามาสวิปฺปยุตฺตสฺส ธมฺมสฺส วิปฺปยุตฺตปจฺจเยน ปจฺจโย.\csromanspacer{} สหชาตํ, ปจฺฉาชาตํ. (สํขิตฺตํ.) (1)}

\prose{ปรามาสวิปฺปยุตฺโต ธมฺโม ปรามาสวิปฺปยุตฺตสฺส ธมฺมสฺส วิปฺปยุตฺตปจฺจเยน ปจฺจโย.\csromanspacer{} สหชาตํ, ปุเรชาตํ, ปจฺฉาชาตํ. (สํขิตฺตํ.) (1)}

\prose{ปรามาสวิปฺปยุตฺโต ธมฺโม ปรามาสสมฺปยุตฺตสฺส ธมฺมสฺส วิปฺปยุตฺตปจฺจเยน ปจฺจโย.\csromanspacer{} \textbf{ปุเรชาตํ} วตฺถุ ปรามาสสมฺปยุตฺตกานํ ขนฺธานํ วิปฺปยุตฺตปจฺจเยน ปจฺจโย. (1)}

\csromanheader{2}{\textbf{อตฺถิ}}
\proseitem{77}{ปรามาสสมฺปยุตฺโต ธมฺโม ปรามาสสมฺปยุตฺตสฺส ธมฺมสฺส อตฺถิปจฺจเยน ปจฺจโย.\csromanspacer{} ปรามาสสมฺปยุตฺตา เอโก ขนฺโธ ติณฺณนฺนํ ขนฺธานํ อตฺถิปจฺจเยน ปจฺจโย ฯเปฯ เทฺว ขนฺธา ฯเปฯ. (1)}

\prose{ปรามาสสมฺปยุตฺโต ธมฺโม ปรามาสวิปฺปยุตฺตสฺส ธมฺมสฺส อตฺถิปจฺจเยน ปจฺจโย.\csromanspacer{} ปรามาสสมฺปยุตฺตา ขนฺธา จิตฺตสมุฏฺฐานานํ รูปานํ อตฺถิปจฺจเยน ปจฺจโย. (มูลํ กาตพฺพํ.) (2)}

\prose{ปรามาสสมฺปยุตฺโต ธมฺโม ปรามาสสมฺปยุตฺตสฺส จ ปรามาสวิปฺปยุตฺตสฺส จ ธมฺมสฺส อตฺถิปจฺจเยน ปจฺจโย.\csromanspacer{} ปรามาสสมฺปยุตฺโต เอโก ขนฺโธ ติณฺณนฺนํ ขนฺธานํ จิตฺตสมุฏฺฐานานญฺจ รูปานํ อตฺถิปจฺจเยน ปจฺจโย ฯเปฯ เทฺว ขนฺธา ฯเปฯ. (3)}

\proseitem{78}{ปรามาสวิปฺปยุตฺโต ธมฺโม ปรามาสวิปฺปยุตฺตสฺส ธมฺมสฺส อตฺถิปจฺจเยน ปจฺจโย.\csromanspacer{} สหชาตํ, ปุเรชาตํ, ปจฺฉาชาตํ, อาหารํ, อินฺทฺริยํ. (สํขิตฺตํ.) (1)}

\prose{ปรามาสวิปฺปยุตฺโต ธมฺโม ปรามาสสมฺปยุตฺตสฺส ธมฺมสฺส อตฺถิปจฺจเยน ปจฺจโย.\csromanspacer{} \textbf{ปุเรชาตํ} จกฺขุํ ฯเปฯ วตฺถุํ อสฺสาเทติ อภินนฺทติ, ตํ อารพฺภ ปรามาสสมฺปยุตฺโต ราโค อุปฺปชฺชติ.\csromanspacer{} วตฺถุ ปรามาสสมฺปยุตฺตกานํ ขนฺธานํ อตฺถิปจฺจเยน ปจฺจโย. (2)}

\proseitem{79}{\csromanfolio{364}ปรามาสสมฺปยุตฺโต จ ปรามาสวิปฺปยุตฺโต จ ธมฺมา ปรามาสสมฺปยุตฺตสฺส ธมฺมสฺส อตฺถิปจฺจเยน ปจฺจโย.\csromanspacer{} สหชาตํ, ปุเรชาตํ.\csromanspacer{}\csromanspacer{}สหชาโต ปรามาสสมฺปยุตฺโต เอโก ขนฺโธ จ วตฺถุ จ ติณฺณนฺนํ ขนฺธานํ อตฺถิปจฺจเยน ปจฺจโย ฯเปฯ เทฺว ขนฺธา จ ฯเปฯ. (1)}

\prose{ปรามาสสมฺปยุตฺโต จ ปรามาสวิปฺปยุตฺโต จ ธมฺมา ปรามาสวิปฺปยุตฺตสฺส ธมฺมสฺส อตฺถิปจฺจเยน ปจฺจโย.\csromanspacer{} สหชาตํ, ปจฺฉาชาตํ, อาหารํ, อินฺทฺริยํ.\csromanspacer{}\csromanspacer{}\textbf{สหชาตา} ปรามาสสมฺปยุตฺตา ขนฺธา จ มหาราภูตา จ จิตฺตสมุฏฺฐานานํ รูปานํ อตฺถิปจฺจเยน ปจฺจโย.\csromanspacer{} \textbf{ปจฺฉาชาตา} ปรามาสสมฺปยุตฺตา ขนฺธา จ กพฬีกาโร อาหาโร จ อิมสฺส กายสฺส อตฺถิปจฺจเยน ปจฺจโย.\csromanspacer{} \textbf{ปจฺฉาชาตา} ปรามาสสมฺปยุตฺตา ขนฺธา จ รูปชีวิตินฺทฺริยญฺจ กฏตฺตารูปานํ อตฺถิปจฺจเยน ปจฺจโย. (2)}

\csromanheader{2}{1. ปจฺจยานุโลม 2. สงฺขฺยาวาร}
\csromanheader{2}{\textbf{สุทฺธ}}
\proseitem{80}{เหตุยา จตฺตาริ, อารมฺมเณ จตฺตาริ, อธิปติยา ปญฺจ, อนนฺตเร จตฺตาริ, สมนนฺตเร จตฺตาริ, สหชาเต ปญฺจ, อญฺญมญฺเญ เทฺว, นิสฺสเย สตฺต, อุปนิสฺสเย จตฺตาริ, ปุเรชาเต เทฺว, ปจฺฉาชาเต เทฺว, อาเสวเน เทฺว, กมฺเม จตฺตาริ, วิปาเก เอกํ, อาหาเร จตฺตาริ, อินฺทฺริเย จตฺตาริ, ฌาเน จตฺตาริ, มคฺเค จตฺตาริ, สมฺปยุตฺเต เทฺว, วิปฺปยุตฺเต ตีณิ, อตฺถิยา สตฺต, นตฺถิยา จตฺตาริ, วิคเต จตฺตาริ, อวิคเต สตฺต.}

\vspace{\tipitakaparskip}
\csromancenter{อนุโลมํ.}
\csromanheader{2}{2. ปจฺจนียุทฺธาร}
\proseitem{81}{ปรามาสสมฺปยุตฺโต ธมฺโม ปรามาสสมฺปยุตฺตสฺส ธมฺมสฺส อารมฺมณปจฺจเยน ปจฺจโย.\csromanspacer{} สหชาตปจฺจเยน ปจฺจโย.\csromanspacer{} อุปนิสฺสยปจฺจเยน ปจฺจโย. (1)}

\prose{ปรามาสสมฺปยุตฺโต ธมฺโม ปรามาสวิปฺปยุตฺตสฺส ธมฺมสฺส อารมฺมณปจฺจเยน ปจฺจโย.\csromanspacer{} สหชาตปจฺจเยน ปจฺจโย.\csromanspacer{} อุปนิสฺสยปจฺจเยน ปจฺจโย.\csromanspacer{} ปจฺฉาชาตปจฺจเยน ปจฺจโย.\csromanspacer{} กมฺมปจฺจเยน ปจฺจโย. (2)}

\csromanheader{2}{52. ปรามาสสมฺปยุตฺตทุก}
\prose{\csromanfolio{365}ปรามาสสมฺปยุตฺโต ธมฺโม ปรามาสสมฺปยุตฺตสฺส จ ปรามาสวิปฺปยุตฺตสฺส จ ธมฺมสฺส สหชาตปจฺจเยน ปจฺจโย. (3)}

\proseitem{82}{ปรามาสวิปฺปยุตฺโต ธมฺโม ปรามาสวิปฺปยุตฺตสฺส ธมฺมสฺส อารมฺมณปจฺจเยน ปจฺจโย.\csromanspacer{} สหชาตปจฺจเยน ปจฺจโย.\csromanspacer{} อุปนิสฺสยปจฺจเยน ปจฺจโย.\csromanspacer{} ปุเรชาตปจฺจเยน ปจฺจโย.\csromanspacer{} ปจฺฉาชาตปจฺจเยน ปจฺจโย.\csromanspacer{} กมฺมปจฺจเยน ปจฺจโย.\csromanspacer{} อาหารปจฺจเยน ปจฺจโย.\csromanspacer{} อินฺทฺริยปจฺจเยน ปจฺจโย. (1)}

\prose{ปรามาสวิปฺปยุตฺโต ธมฺโม ปรามาสสมฺปยุตฺตสฺส ธมฺมสฺส อารมฺมณปจฺจเยน ปจฺจโย.\csromanspacer{} อุปนิสฺสยปจฺจเยน ปจฺจโย.\csromanspacer{} ปุเรชาตปจฺจเยน ปจฺจโย. (2)}

\prose{ปรามาสสมฺปยุตฺโต จ ปรามาสวิปฺปยุตฺโต จ ธมฺมา ปรามาสสมฺปยุตฺตสฺส ธมฺมสฺส สหชาตํ, ปุเรชาตํ. (1)}

\prose{ปรามาสสมฺปยุตฺโต จ ปรามาสวิปฺปยุตฺโต จ ธมฺมา ปรามาสวิปฺปยุตฺตสฺส ธมฺมสฺส สหชาตํ, ปจฺฉาชาตํ, อาหารํ, อินฺทฺริยํ. (2)}

\csromanheader{2}{2. ปจฺจยปจฺจนีย 2. สงฺขฺยาวาร}
\csromanheader{2}{\textbf{สุทฺธ}}
\proseitem{83}{นเหตุยา สตฺต, น-อารมฺมเณ สตฺต, น-อธิปติยา สตฺต, น-อนนฺตเร สตฺต, นสมนนฺตเร สตฺต, นสหชาเต ปญฺจ, น-อญฺญมญฺเญ ปญฺจ, นนิสฺสเย ปญฺจ, น-อุปนิสฺสเย สตฺต, นปุเรชาเต ฉ, นปจฺฉาชาเต สตฺต, (สพฺพตฺถ สตฺต,) นมคฺเค สตฺต, นสมฺปยุตฺเต ปญฺจ, นวิปฺปยุตฺเต จตฺตาริ, โน-อตฺถิยา จตฺตาริ, โนนตฺถิยา สตฺต, โนวิคเต สตฺต, โน-อวิคเต จตฺตาริ.}

\csromanheader{2}{3. ปจฺจยานุโลมปจฺจนีย}
\csromanheader{2}{\textbf{เหตุทุก}}
\proseitem{84}{เหตุปจฺจยา น-อารมฺมเณ จตฺตาริ, น-อธิปติยา จตฺตาริ, น-อนนฺตเร จตฺตาริ, นสมนนฺตเร จตฺตาริ, น-อญฺญมญฺเญ เทฺว, น-อุปนิสฺสเย \csromanfolio{366}จตฺตาริ, (สพฺพตฺถ จตฺตาริ,) นมคฺเค จตฺตาริ, นสมฺปยุตฺเต เทฺว, นวิปฺปยุตฺเต เทฺว, โนนตฺถิยา จตฺตาริ, โนวิคเต จตฺตาริ.}

\csromanheader{2}{4. ปจฺจยปจฺจนียานุโลม}
\proseitem{85}{นเหตุปจฺจยา อารมฺมเณ จตฺตาริ, อธิปติยา ปญฺจ, (อนุโลมมาติกา กาตพฺพา.) ฯเปฯ อวิคเต สตฺต.}

\nitthitam{ปรามาสสมฺปยุตฺตทุกํ นิฏฺฐิตํ.}
\csromanmatikaanchor{mtk.51}
\csromantocmarkref{subsection}{53. ปรามาสปรามฏฺฐทุก}{mtk.51}
\bookmark[dest={mtk.51},level=2]{53. ปรามาสปรามฏฺฐทุก}
\csromanheader{2}{53. ปรามาสปรามฏฺฐทุก 1. ปฏิจฺจวาร}
\csromanheader{2}{1. ปจฺจยานุโลม 1. วิภงฺควาร}
\csromanheader{2}{\textbf{เหตุ}}
\proseitem{86}{ปรามาสญฺเจว ปรามฏฺฐญฺจ ธมฺมํ ปฏิจฺจ ปรามฏฺโฐ เจว โน จ ปรามาโส ธมฺโม อุปฺปชฺชติ เหตุปจฺจยา.\csromanspacer{} ปรามาสํ ปฏิจฺจ สมฺปยุตฺตกา ขนฺธา จิตฺตสมุฏฺฐานญฺจ รูปํ. (1)}

\prose{ปรามฏฺฐญฺเจว โน จ ปรามาสํ ธมฺมํ ปฏิจฺจ ปรามฏฺโฐ เจว โน จ ปรามาโส ธมฺโม อุปฺปชฺชติ เหตุปจฺจยา.\csromanspacer{} ปรามฏฺฐญฺเจว โน จ ปรามาสํ เอกํ ขนฺธํ ปฏิจฺจ ตโย ขนฺธา จิตฺตสมุฏฺฐานญฺจ รูปํ ฯเปฯ เทฺว ขนฺเธ ฯเปฯ.\csromanspacer{} ปฏิสนฺธิกฺขเณ ฯเปฯ. (ยาว อชฺฌตฺติกา มหาภูตา.) (1)}

\prose{ปรามฏฺฐญฺเจว โน จ ปรามาสํ ธมฺมํ ปฏิจฺจ ปรามาโส เจว ปรามฏฺโฐ จ ธมฺโม อุปฺปชฺชติ เหตุปจฺจยา.\csromanspacer{} ปรามฏฺเฐ เจว โน จ ปรามาเส ขนฺเธ ปฏิจฺจ ปรามาโส. (2)}

\prose{ปรามฏฺฐญฺเจว โน จ ปรามาสํ ธมฺมํ ปฏิจฺจ ปรามาโส เจว ปรามฏฺโฐ จ ปรามฏฺโฐ เจว โน จ ปรามาโส จ ธมฺมา อุปฺปชฺชนฺติ เหตุปจฺจยา.\csromanspacer{} ปรามฏฺฐญฺเจว โน จ ปรามาสํ เอกํ ขนฺธํ ปฏิจฺจ ตโย ขนฺธา ปรามาโส จ จิตฺตสมุฏฺฐานญฺจ รูปํ ฯเปฯ เทฺว ขนฺเธ ฯเปฯ. (3)}

\prose{ปรามาสญฺเจว ปรามฏฺฐญฺจ ปรามฏฺฐญฺเจว โน จ ปรามาสญฺจ ธมฺมํ ปฏิจฺจ ปรามฏฺโฐ เจว โน จ ปรามาโส ธมฺโม อุปฺปชฺชติ เหตุปจฺจยา.}

\vspace{\tipitakaparskip}
\csromancenter{ปรามาสปรามฏฺฐทุก ปรามฏฺฐญฺเจว โน จ ปรามาสํ เอกํ ขนฺธญฺจ ปรามาสญฺจ ปฏิจฺจ ตโย ขนฺธา จิตฺตสมุฏฺฐานญฺจ รูปํ ฯเปฯ เทฺว ขนฺเธ จ ฯเปฯ. (สํขิตฺตํ.)}
\csromancenter{(สพฺเพ วารา ยถา ปรามาสทุกํ, เอวํ กาตพฺพํ, นินฺนานากรณํ.)}
\csromanheader{2}{53. ปรามาสปรามฏฺฐทุก 7. ปญฺหาวาร}
\csromanheader{2}{1. ปจฺจยานุโลม 1. วิภงฺควาร}
\csromanheader{2}{\textbf{เหตุ}}
\proseitem{87}{\csromanfolio{367}ปรามฏฺโฐ เจว โน จ ปรามาโส ธมฺโม ปรามฏฺฐสฺส เจว โน จ ปรามาสสฺส ธมฺมสฺส เหตุปจฺจเยน ปจฺจโย.\csromanspacer{} ปรามฏฺฐา เจว โน จ ปรามาสา เหตู สมฺปยุตฺตกานํ ขนฺธานํ จิตฺตสมุฏฺฐานานญฺจ รูปานํ เหตุปจฺจเยน ปจฺจโย.\csromanspacer{} ปฏิสนฺธิกฺขเณ ฯเปฯ. (1)}

\prose{ปรามฏฺโฐ เจว โน จ ปรามาโส ธมฺโม ปรามาสสฺส เจว ปรามฏฺฐสฺส จ ธมฺมสฺส เหตุปจฺจเยน ปจฺจโย.\csromanspacer{} ปรามฏฺฐา เจว โน จ ปรามาสา เหตู ปรามาสสฺส เหตุปจฺจเยน ปจฺจโย. (2)}

\prose{ปรามฏฺโฐ เจว โน จ ปรามาโส ธมฺโม ปรามาสสฺส เจว ปรามฏฺฐสฺส จ ปรามฏฺฐสฺส เจว โน จ ปรามาสสฺส จ ธมฺมสฺส เหตุปจฺจเยน ปจฺจโย.\csromanspacer{} ปรามฏฺฐา เจว โน จ ปรามาสา เหตู สมฺปยุตฺตกานํ ขนฺธานํ ปรามาสสฺส จ จิตฺตสมุฏฺฐานานญฺจ รูปานํ เหตุปจฺจเยน ปจฺจโย. (3)}

\csromanheader{2}{\textbf{อารมฺมณ}}
\proseitem{88}{ปรามาโส เจว ปรามฏฺโฐ จ ธมฺโม ปรามาสสฺส เจว ปรามฏฺฐสฺส จ ธมฺมสฺส อารมฺมณปจฺจเยน ปจฺจโย.\csromanspacer{} ตีณิ. (อารพฺภ กาตพฺพานิ.\csromanspacer{} ปรามาสทุกสทิสํ.)}

\proseitem{89}{ปรามฏฺโฐ เจว โน จ ปรามาโส ธมฺโม ปรามฏฺฐสฺส เจว โน จ ปรามาสสฺส ธมฺมสฺส อารมฺมณปจฺจเยน ปจฺจโย.\csromanspacer{} ทานํ ฯเปฯ สีลํ ฯเปฯ อุโปสถกมฺมํ กตฺวา ตํ ปจฺจเวกฺขติ, อสฺสาเทติ อภินนฺทติ, ตํ อารพฺภ ราโค ฯเปฯ.\csromanspacer{} วิจิกิจฺฉา.\csromanspacer{} อุทฺธจฺจํ.\csromanspacer{} โทมนสฺสํ อุปฺปชฺชติ.\csromanspacer{} ปุพฺเพ ฯเปฯ.\csromanspacer{} ฌานา ฯเปฯ.\csromanspacer{} อริยา โคตฺรภุํ ปจฺจเวกฺขนฺติ, โวทานํ ปจฺจเวกฺขนฺติ.\csromanspacer{} ปหีเน กิเลเส ฯเปฯ. \csromanfolio{368}วิกฺขมฺภิเต กิเลเส ฯเปฯ.\csromanspacer{} ปุพฺเพ ฯเปฯ.\csromanspacer{} จกฺขุํ ฯเปฯ วตฺถุํ ปรามฏฺเฐ เจว โน จ ปรามาเส ขนฺเธ อนิจฺจโต ฯเปฯ โทมนสฺสํ อุปฺปชฺชติ.\csromanspacer{} ทิพฺเพน จกฺขุนา รูปํ ปสฺสติ. (ยาว อาวชฺชนา สพฺพํ กาตพฺพํ.) (1)}

\prose{ปรามฏฺโฐ เจว โน จ ปรามาโส ธมฺโม ปรามาสสฺส เจว ปรามฏฺฐสฺส จ ธมฺมสฺส อารมฺมณปจฺจเยน ปจฺจโย.\csromanspacer{} ทานํ ฯเปฯ สีลํ ฯเปฯ อุโปสถกมฺมํ ฯเปฯ.\csromanspacer{} ปุพฺเพ ฯเปฯ.\csromanspacer{} ฌานา ฯเปฯ.\csromanspacer{} จกฺขุํ ฯเปฯ.\csromanspacer{} วตฺถุํ ปรามฏฺเฐ เจว โน จ ปรามาเส ขนฺเธ อสฺสาเทติ อภินนฺทติ, ตํ อารพฺภ ทิฏฺฐิ อุปฺปชฺชติ. (2)}

\prose{ปรามฏฺโฐ เจว โน จ ปรามาโส ธมฺโม ปรามาสสฺส เจว ปรามฏฺฐสฺส จ ปรามฏฺฐสฺส เจว โน จ ปรามาสสฺส จ ธมฺมสฺส อารมฺมณปจฺจเยน ปจฺจโย.\csromanspacer{} ทานํ ฯเปฯ สีลํ ฯเปฯ อุโปสถกมฺมํ ฯเปฯ.\csromanspacer{} ปุพฺเพ ฯเปฯ.\csromanspacer{} ฌานา ฯเปฯ.\csromanspacer{} จกฺขุํ ฯเปฯ วตฺถุํ ปรามฏฺเฐ เจว โน จ ปรามาเส ขนฺเธ อสฺสาเทติ อภินนฺทติ, ตํ อารพฺภ ปรามาโส จ สมฺปยุตฺตกา จ ขนฺธา อุปฺปชฺชนฺติ. (3)}

\prose{(เอวํ อิตเรปิ ตีณิ.\csromanspacer{} อารพฺภ กาตพฺพานิ.\csromanspacer{} อิมํ ทุกํ ปรามาสทุกสทิสํ.\csromanspacer{} โลกุตฺตรํ ยหิํ น ลพฺภติ, ตหิํ น กาตพฺพํ.)}

\csromanmatikaanchor{mtk.52}
\csromantocmarkref{subsection}{54. ปรามาสวิปฺปยุตฺตปรามฏฺฐทุก}{mtk.52}
\bookmark[dest={mtk.52},level=2]{54. ปรามาสวิปฺปยุตฺตปรามฏฺฐทุก}
\csromanheader{2}{54. ปรามาสวิปฺปยุตฺตปรามฏฺฐทุก 1. ปฏิจฺจวาร}
\csromanheader{2}{1. ปจฺจยานุโลม 1. วิภงฺควาร}
\csromanheader{2}{\textbf{เหตุ}}
\proseitem{90}{ปรามาสวิปฺปยุตฺตํ ปรามฏฺฐํ ธมฺมํ ปฏิจฺจ ปรามาสวิปฺปยุตฺโต ปรามฏฺโฐ ธมฺโม อุปฺปชฺชติ เหตุปจฺจยา.\csromanspacer{} ปรามาสวิปฺปยุตฺตํ ปรามฏฺฐํ เอกํ ขนฺธํ ปฏิจฺจ ตโย ขนฺธา จิตฺตสมุฏฺฐานญฺจ รูปํ ฯเปฯ เทฺว ขนฺเธ ฯเปฯ.}

\prose{ปรามาสวิปฺปยุตฺตํ อปรามฏฺฐํ ธมฺมํ ปฏิจฺจ ปรามาสวิปฺปยุตฺโต อปรามฏฺโฐ ธมฺโม อุปฺปชฺชติ เหตุปจฺจยา. (สํขิตฺตํ.)}

\prose{(ยถา จูฬนฺตรทุเก โลกิยทุกํ, เอวํ กาตพฺพํ, นินฺนานากรณํ.)}

\nitthitam{ปรามาสวิปฺปยุตฺตปรามฏฺฐทุกํ นิฏฺฐิตํ.}
\nitthitam{ปรามาสโคจฺฉกํ นิฏฺฐิตํ.}
\csromanmatikaanchor{mtk.53}
\csromanmatikaanchor{mtk.54}
\csromantocmarknopage{section}{10. มหนฺตรทุก}{mtk.53}
\bookmark[dest={mtk.53},level=1]{10. มหนฺตรทุก}
\csromantocmarkref{subsection}{55. สารมฺมณทุก}{mtk.54}
\bookmark[dest={mtk.54},level=2]{55. สารมฺมณทุก}
\csromanheader{2}{55. สารมฺมณทุก \textbf{10. มหนฺตรทุก}}
\csromanheader{2}{55. สารมฺมณทุก 1. ปฏิจฺจวาร}
\csromanheader{2}{1. ปจฺจยานุโลม 1. วิภงฺควาร}
\csromanheader{2}{\textbf{เหตุ}}
\proseitem{1}{\csromanfolio{369}สารมฺมณํ ธมฺมํ ปฏิจฺจ สารมฺมโณ ธมฺโม อุปฺปชฺชติ เหตุปจฺจยา.\csromanspacer{} สารมฺมณํ เอกํ ขนฺธํ ปฏิจฺจ ตโย ขนฺธา ฯเปฯ เทฺว ขนฺเธ ฯเปฯ.\csromanspacer{} ปฏิสนฺธิกฺขเณ ฯเปฯ. (1)}

\prose{สารมฺมณํ ธมฺมํ ปฏิจฺจ อนารมฺมโณ ธมฺโม อุปฺปชฺชติ เหตุปจฺจยา.\csromanspacer{} สารมฺมเณ ขนฺเธ ปฏิจฺจ จิตฺตสมุฏฺฐานํ รูปํ.\csromanspacer{} ปฏิสนฺธิกฺขเณ ฯเปฯ. (2)}

\prose{สารมฺมณํ ธมฺมํ ปฏิจฺจ สารมฺมโณ จ อนารมฺมโณ จ ธมฺมา อุปฺปชฺชนฺติ เหตุปจฺจยา.\csromanspacer{} สารมฺมณํ เอกํ ขนฺธํ ปฏิจฺจ ตโย ขนฺธา จิตฺตสมุฏฺฐานญฺจ รูปํ ฯเปฯ เทฺว ขนฺเธ ฯเปฯ.\csromanspacer{} ปฏิสนฺธิกฺขเณ ฯเปฯ. (3)}

\proseitem{2}{อนารมฺมณํ ธมฺมํ ปฏิจฺจ อนารมฺมโณ ธมฺโม อุปฺปชฺชติ เหตุปจฺจยา.\csromanspacer{} เอกํ มหาภูตํ ฯเปฯ มหาภูเต ปฏิจฺจ จิตฺตสมุฏฺฐานํ รูปํ, กฏตฺตารูปํ, อุปาทารูปํ. (1)}

\prose{อนารมฺมณํ ธมฺมํ ปฏิจฺจ สารมฺมโณ ธมฺโม อุปฺปชฺชติ เหตุปจฺจยา.\csromanspacer{} ปฏิสนฺธิกฺขเณ วตฺถุํ ปฏิจฺจ สารมฺมณา ขนฺธา. (2)}

\prose{อนารมฺมณํ ธมฺมํ ปฏิจฺจ สารมฺมโณ จ อนารมฺมโณ จ ธมฺมา อุปฺปชฺชนฺติ เหตุปจฺจยา.\csromanspacer{} ปฏิสนฺธิกฺขเณ วตฺถุํ ปฏิจฺจ สารมฺมณา ขนฺธา, มหาภูเต ปฏิจฺจ กฏตฺตารูปํ. (3)}

\proseitem{3}{สารมฺมณญฺจ อนารมฺมณญฺจ ธมฺมํ ปฏิจฺจ สารมฺมโณ ธมฺโม อุปฺปชฺชติ เหตุปจฺจยา.\csromanspacer{} ปฏิสนฺธิกฺขเณ สารมฺมณํ เอกํ ขนฺธญฺจ วตฺถุญฺจ ปฏิจฺจ ตโย ขนฺธา ฯเปฯ เทฺว ขนฺเธ จ ฯเปฯ. (1)}

\prose{\csromanfolio{370}สารมฺมณญฺจ อนารมฺมณญฺจ ธมฺมํ ปฏิจฺจ อนารมฺมโณ ธมฺโม อุปฺปชฺชติ เหตุปจฺจยา.\csromanspacer{} สารมฺมเณ ขนฺเธ จ มหาภูเต จ ปฏิจฺจ จิตฺตสมุฏฺฐานํ รูปํ.\csromanspacer{} ปฏิสนฺธิกฺขเณ ฯเปฯ. (2)}

\prose{สารมฺมณญฺจ อนารมฺมณญฺจ ธมฺมํ ปฏิจฺจ สารมฺมโณ จ อนารมฺมโณ จ ธมฺมา อุปฺปชฺชนฺติ เหตุปจฺจยา.\csromanspacer{} ปฏิสนฺธิกฺขเณ สารมฺมณํ เอกํ ขนฺธญฺจ วตฺถุญฺจ ปฏิจฺจ ตโย ขนฺธา ฯเปฯ เทฺว ขนฺเธ จ ฯเปฯ.\csromanspacer{} สารมฺมเณ ขนฺเธ จ มหาภูเต จ ปฏิจฺจ กฏตฺตารูปํ. (3)}

\csromanheader{2}{\textbf{อารมฺมณ}}
\proseitem{4}{สารมฺมณํ ธมฺมํ ปฏิจฺจ สารมฺมโณ ธมฺโม อุปฺปชฺชติ อารมฺมณปจฺจยา.\csromanspacer{} สารมฺมณํ เอกํ ขนฺธํ ปฏิจฺจ ตโย ขนฺธา ฯเปฯ เทฺว ขนฺเธ ฯเปฯ.\csromanspacer{} ปฏิสนฺธิกฺขเณ ฯเปฯ.}

\prose{อนารมฺมณํ ธมฺมํ ปฏิจฺจ สารมฺมโณ ธมฺโม อุปฺปชฺชติ อารมฺมณปจฺจยา.\csromanspacer{} ปฏิสนฺธิกฺขเณ วตฺถุํ ปฏิจฺจ สารมฺมณา ขนฺธา.}

\prose{สารมฺมณญฺจ อนารมฺมณญฺจ ธมฺมํ ปฏิจฺจ สารมฺมโณ ธมฺโม อุปฺปชฺชติ อารมฺมณปจฺจยา.\csromanspacer{} ปฏิสนฺธิกฺขเณ สารมฺมณํ เอกํ ขนฺธญฺจ วตฺถุญฺจ ปฏิจฺจ ตโย ขนฺธา ฯเปฯ เทฺว ขนฺเธ จ ฯเปฯ. (สํขิตฺตํ.)}

\csromanheader{2}{1. ปจฺจยานุโลม 2. สญฺขฺยาวาร}
\csromanheader{2}{\textbf{สุทฺธ}}
\proseitem{5}{เหตุยา นว, อารมฺมเณ ตีณิ, อธิปติยา ปญฺจ, อนนฺตเร ตีณิ, สมนนฺตเร ตีณิ, สหชาเต นว, อญฺญมญฺเญ ฉ, นิสฺสเย นว, อุปนิสฺสเย ตีณิ, ปุเรชาเต เอกํ, อาเสวเน เอกํ, กมฺเม นว, วิปาเก นว, อาหาเร นว, อินฺทฺริเย นว, ฌาเน นว, มคฺเค นว, สมฺปยุตฺเต ตีณิ, วิปฺปยุตฺเต นว, อตฺถิยา นว, นตฺถิยา ตีณิ, วิคเต ตีณิ, อวิคเต นว.}

\csromanheader{2}{55. สารมฺมณทุก \textbf{2. ปจฺจยปจฺจนีย 1. วิภงฺควาร}}
\csromanheader{2}{\textbf{นเหตุ}}
\proseitem{6}{\csromanfolio{371}สารมฺมณํ ธมฺมํ ปฏิจฺจ สารมฺมโณ ธมฺโม อุปฺปชฺชติ นเหตุปจฺจยา.\csromanspacer{} อเหตุกํ สารมฺมณํ เอกํ ขนฺธํ ปฏิจฺจ ตโย ขนฺธา ฯเปฯ เทฺว ขนฺเธ ฯเปฯ.\csromanspacer{} อเหตุกปฏิสนฺธิกฺขเณ ฯเปฯ.\csromanspacer{} วิจิกิจฺฉาสหคเต อุทฺธจฺจสหคเต ขนฺเธ ปฏิจฺจ วิจิกิจฺฉาสหคโต อุทฺธจฺจสหคโต โมโห. (1)}

\prose{สารมฺมณํ ธมฺมํ ปฏิจฺจ อนารมฺมโณ ธมฺโม อุปฺปชฺชติ นเหตุปจฺจยา.\csromanspacer{} อเหตุเก สารมฺมเณ ขนฺเธ ปฏิจฺจ จิตฺตสมุฏฺฐานํ รูปํ.\csromanspacer{} อเหตุกปฏิสนฺธิกฺขเณ ฯเปฯ. (2)}

\prose{สารมฺมณํ ธมฺมํ ปฏิจฺจ สารมฺมโณ จ อนารมฺมโณ จ ธมฺมา อุปฺปชฺชนฺติ นเหตุปจฺจยา.\csromanspacer{} อเหตุกํ สารมฺมณํ เอกํ ขนฺธํ ปฏิจฺจ ตโย ขนฺธา จิตฺตสมุฏฺฐานญฺจ รูปํ ฯเปฯ เทฺว ขนฺเธ ฯเปฯ.\csromanspacer{} อเหตุกปฏิสนฺธิกฺขเณ ฯเปฯ. (3)}

\proseitem{7}{อนารมฺมณํ ธมฺมํ ปฏิจฺจ อนารมฺมโณ ธมฺโม อุปฺปชฺชติ นเหตุปจฺจยา.\csromanspacer{} เอกํ มหาภูตํ ฯเปฯ.\csromanspacer{} อสญฺญสตฺตานํ เอกํ มหาภูตํ ฯเปฯ. (1)}

\prose{อนารมฺมณํ ธมฺมํ ปฏิจฺจ สารมฺมโณ ธมฺโม อุปฺปชฺชติ นเหตุปจฺจยา.\csromanspacer{} อเหตุกปฏิสนฺธิกฺขเณ วตฺถุํ ปฏิจฺจ สารมฺมณา ขนฺธา. (2)}

\prose{อนารมฺมณํ ธมฺมํ ปฏิจฺจ สารมฺมโณ จ อนารมฺมโณ จ ธมฺมา อุปฺปชฺชนฺติ นเหตุปจฺจยา.\csromanspacer{} อเหตุกปฏิสนฺธิกฺขเณ วตฺถุํ ปฏิจฺจ สารมฺมณา ขนฺธา.\csromanspacer{} มหาภูเต ปฏิจฺจ กฏตฺตารูปํ. (3)}

\proseitem{8}{สารมฺมณญฺจ อนารมฺมณญฺจ ธมฺมํ ปฏิจฺจ สารมฺมโณ ธมฺโม อุปฺปชฺชติ นเหตุปจฺจยา.\csromanspacer{} อเหตุกปฏิสนฺธิกฺขเณ สารมฺมณํ เอกํ ขนฺธญฺจ วตฺถุญฺจ ปฏิจฺจ ตโย ขนฺธา ฯเปฯ เทฺว ขนฺเธ จ ฯเปฯ. (1)}

\prose{สารมฺมณญฺจ อนารมฺมณญฺจ ธมฺมํ ปฏิจฺจ อนารมฺมโณ ธมฺโม อุปฺปชฺชติ นเหตุปจฺจยา.\csromanspacer{} อเหตุเก สารมฺมเณ ขนฺเธ จ มหาภูเต จ ปฏิจฺจ จิตฺตสมุฏฺฐานํ รูปํ.\csromanspacer{} อเหตุกปฏิสนฺธิกฺขเณ ฯเปฯ. (2)}

\prose{สารมฺมณญฺจ อนารมฺมณญฺจ ธมฺมํ ปฏิจฺจ สารมฺมโณ จ อนารมฺมโณ จ ธมฺมา อุปฺปชฺชนฺติ นเหตุปจฺจยา.\csromanspacer{} อเหตุกปฏิสนฺธิกฺขเณ สารมฺมณํ เอกํ \csromanfolio{372}ขนฺธญฺจ วตฺถุญฺจ ปฏิจฺจ ตโย ขนฺธา ฯเปฯ เทฺว ขนฺเธ จ ฯเปฯ.\csromanspacer{} สารมฺมเณ ขนฺเธ จ มหาภูเต จ ปฏิจฺจ กฏตฺตารูปํ. (3)}

\csromanheader{2}{\textbf{น-อารมฺมณ}}
\proseitem{9}{สารมฺมณํ ธมฺมํ ปฏิจฺจ อนารมฺมโณ ธมฺโม อุปฺปชฺชติ น-อารมฺมณปจฺจยา.\csromanspacer{} สารมฺมเณ ขนฺเธ ปฏิจฺจ จิตฺตสมุฏฺฐานํ รูปํ.\csromanspacer{} ปฏิสนฺธิกฺขเณ ฯเปฯ. (1)}

\prose{อนารมฺมณํ ธมฺมํ ปฏิจฺจ อนารมฺมโณ ธมฺโม อุปฺปชฺชติ น-อารมฺมณปจฺจยา. (ยาว อสญฺญสตฺตา.) (1)}

\prose{สารมฺมณญฺจ อนารมฺมณญฺจ ธมฺมํ ปฏิจฺจ อนารมฺมโณ ธมฺโม อุปฺปชฺชติ น-อารมฺมณปจฺจยา.\csromanspacer{} สารมฺมเณ ขนฺเธ จ มหาภูเต จ ปฏิจฺจ จิตฺตสมุฏฺฐานํ รูปํ.\csromanspacer{} ปฏิสนฺธิกฺขเณ ฯเปฯ. (สํขิตฺตํ.) (1)}

\csromanheader{2}{2. ปจฺจยปจฺจนีย 2. สงฺขฺยาวาร}
\csromanheader{2}{\textbf{สุทฺธ}}
\proseitem{10}{นเหตุยา นว, น-อารมฺมเณ ตีณิ, น-อธิปติยา นว, น-อนนฺตเร ตีณิ, นสมนนฺตเร ตีณิ, น-อญฺญมญฺเญ ตีณิ, น-อุปนิสฺสเย ตีณิ, นปุเรชาเต นว, นปจฺฉาชาเต นว, น-อาเสวเน นว, นกมฺเม เทฺว, นวิปาเก ปญฺจ, น-อาหาเร เอกํ, น-อินฺทฺริเย เอกํ, นฌาเน เทฺว, นมคฺเค นว, นสมฺปยุตฺเต ตีณิ, นวิปฺปยุตฺเต เทฺว, โนนตฺถิยา ตีณิ, โนวิคเต ตีณิ.}
\csromansectionrule

\csromanheader{2}{3. ปจฺจยานุโลมปจฺจนีย}
\proseitem{11}{เหตุปจฺจยา น-อารมฺมเณ ตีณิ, น-อธิปติยา นว ฯเปฯ นกมฺเม เอกํ ฯเปฯ นวิปฺปยุตฺเต เอกํ. (สํขิตฺตํ.)}

\csromanheader{2}{4. ปจฺจยปจฺจนียานุโลม}
\proseitem{12}{นเหตุปจฺจยา อารมฺมเณ ตีณิ ฯเปฯ สหชาเต นว ฯเปฯ มคฺเค เอกํ ฯเปฯ อวิคเต นว.}

\vspace{\tipitakaparskip}
\csromancenter{(สหชาตวาโรปิ ปฏิจฺจวารสทิโส.)}
\csromanheader{2}{55. สารมฺมณทุก \textbf{55. สารมฺมณทุก 3. ปจฺจยวาร}}
\csromanheader{2}{1. ปจฺจยานุโลม 1. วิภงฺควาร}
\csromanheader{2}{\textbf{เหตุ}}
\proseitem{13}{\csromanfolio{373}สารมฺมณํ ธมฺมํ ปจฺจยา สารมฺมโณ ธมฺโม อุปฺปชฺชติ เหตุปจฺจยา.\csromanspacer{} ตีณิ. (ปฏิจฺจสทิสา.)}

\proseitem{14}{อนารมฺมณํ ธมฺมํ ปจฺจยา อนารมฺมโณ ธมฺโม อุปฺปชฺชติ เหตุปจฺจยา.\csromanspacer{} เอกํ มหาภูตํ ฯเปฯ มหาภูเต ปจฺจยา จิตฺตสมุฏฺฐานํ รูปํ, กฏตฺตารูปํ, อุปาทารูปํ. (1)}

\prose{อนารมฺมณํ ธมฺมํ ปจฺจยา สารมฺมโณ ธมฺโม อุปฺปชฺชติ เหตุปจฺจยา.\csromanspacer{} วตฺถุํ ปจฺจยา สารมฺมณา ขนฺธา.\csromanspacer{} ปฏิสนฺธิกฺขเณ วตฺถุํ ปจฺจยา สารมฺมณา ขนฺธา. (2)}

\prose{อนารมฺมณํ ธมฺมํ ปจฺจยา สารมฺมโณ จ อนารมฺมโณ จ ธมฺมา อุปฺปชฺชนฺติ เหตุปจฺจยา.\csromanspacer{} วตฺถุํ ปจฺจยา สารมฺมณา ขนฺธา.\csromanspacer{} มหาภูเต ปจฺจยา จิตฺตสมุฏฺฐานํ รูปํ.\csromanspacer{} ปฏิสนฺธิกฺขเณ ฯเปฯ. (3)}

\proseitem{15}{สารมฺมณญฺจ อนารมฺมณญฺจ ธมฺมํ ปจฺจยา สารมฺมโณ ธมฺโม อุปฺปชฺชติ เหตุปจฺจยา.\csromanspacer{} สารมฺมณํ เอกํ ขนฺธญฺจ วตฺถุญฺจ ปจฺจยา ตโย ขนฺธา ฯเปฯ เทฺว ขนฺเธ จ ฯเปฯ.\csromanspacer{} ปฏิสนฺธิกฺขเณ ฯเปฯ. (1)}

\prose{สารมฺมณญฺจ อนารมฺมณญฺจ ธมฺมํ ปจฺจยา อนารมฺมโณ ธมฺโม อุปฺปชฺชติ เหตุปจฺจยา.\csromanspacer{} สารมฺมเณ ขนฺเธ จ มหาภูเต จ ปจฺจยา จิตฺตสมุฏฺฐานํ รูปํ.\csromanspacer{} ปฏิสนฺธิกฺขเณ ฯเปฯ. (2)}

\prose{สารมฺมณญฺจ อนารมฺมณญฺจ ธมฺมํ ปจฺจยา สารมฺมโณ จ อนารมฺมโณ จ ธมฺมา อุปฺปชฺชนฺติ เหตุปจฺจยา.\csromanspacer{} สารมฺมณํ เอกํ ขนฺธญฺจ วตฺถุญฺจ ปจฺจยา ตโย ขนฺธา ฯเปฯ เทฺว ขนฺเธ จ ฯเปฯ.\csromanspacer{} สารมฺมเณ ขนฺเธ จ มหาภูเต จ ปจฺจยา จิตฺตสมุฏฺฐานํ รูปํ.\csromanspacer{} ปฏิสนฺธิกฺขเณ ฯเปฯ. (3)}

\csromanheader{2}{\textbf{อารมฺมณ}}
\proseitem{16}{สารมฺมณํ ธมฺมํ ปจฺจยา สารมฺมโณ ธมฺโม อุปฺปชฺชติ อารมฺมณปจฺจยา.\csromanspacer{} สารมฺมณํ เอกํ ขนฺธํ ปจฺจยา ตโย ขนฺธา ฯเปฯ เทฺว ขนฺเธ ฯเปฯ.\csromanspacer{} ปฏิสนฺธิกฺขเณ ฯเปฯ. (1)}

\prose{\csromanfolio{374}อนารมฺมณํ ธมฺมํ ปจฺจยา สารมฺมโณ ธมฺโม อุปฺปชฺชติ อารมฺมณปจฺจยา.\csromanspacer{} จกฺขายตนํ ปจฺจยา จกฺขุวิญฺญาณํ ฯเปฯ.\csromanspacer{} กายายตนํ ปจฺจยา กายวิญฺญาณํ.\csromanspacer{} วตฺถุํ ปจฺจยา สารมฺมณา ขนฺธา.\csromanspacer{} ปฏิสนฺธิกฺขเณ ฯเปฯ. (1)}

\prose{สารมฺมณญฺจ อนารมฺมณญฺจ ธมฺมํ ปจฺจยา สารมฺมโณ ธมฺโม อุปฺปชฺชติ อารมฺมณปจฺจยา.\csromanspacer{} จกฺขุวิญฺญาณสหคตํ เอกํ ขนฺธญฺจ จกฺขายตนญฺจ ปจฺจยา ตโย ขนฺธา ฯเปฯ เทฺว ขนฺเธ จ ฯเปฯ.\csromanspacer{} กายวิญฺญาณสหคตํ ฯเปฯ.\csromanspacer{} สารมฺมณํ เอกํ ขนฺธญฺจ วตฺถุญฺจ ปจฺจยา ตโย ขนฺธา ฯเปฯ เทฺว ขนฺเธ จ ฯเปฯ.\csromanspacer{} ปฏิสนฺธิกฺขเณ ฯเปฯ. (สํขิตฺตํ.) (1)}

\csromanheader{2}{1. ปจฺจยานุโลม 2. สงฺขฺยาวาร}
\csromanheader{2}{\textbf{สุทฺธ}}
\proseitem{17}{เหตุยา นว, อารมฺมเณ ตีณิ, อธิปติยา นว, อนนฺตเร ตีณิ, สมนนฺตเร ตีณิ, สหชาเต นว, อญฺญมญฺเญ ฉ, นิสฺสเย นว, อุปนิสฺสเย ตีณิ, ปุเรชาเต ตีณิ, อาเสวเน ตีณิ, กมฺเม นว, วิปาเก นว, อาหาเร นว, อินฺทฺริเย นว, ฌาเน นว, มคฺเค นว, สมฺปยุตฺเต ตีณิ, วิปฺปยุตฺเต นว, อตฺถิยา นว, นตฺถิยา ตีณิ, วิคเต ตีณิ, อวิคเต นว.}

\csromanheader{2}{2. ปจฺจยปจฺจนีย 1. วิภงฺควาร}
\csromanheader{2}{\textbf{นเหตุ}}
\proseitem{18}{สารมฺมณํ ธมฺมํ ปจฺจยา สารมฺมโณ ธมฺโม อุปฺปชฺชติ นเหตุปจฺจยา.\csromanspacer{} ตีณิ. (ปฏิจฺจวารสทิสํ.)}

\prose{อนารมฺมณํ ธมฺมํ ปจฺจยา อนารมฺมโณ ธมฺโม อุปฺปชฺชติ นเหตุปจฺจยา.\csromanspacer{} เอกํ มหาภูตํ ฯเปฯ.\csromanspacer{} อสญฺญสตฺตานํ เอกํ มหาภูตํ ฯเปฯ. (1)}

\prose{อนารมฺมณํ ธมฺมํ ปจฺจยา สารมฺมโณ ธมฺโม อุปฺปชฺชติ นเหตุปจฺจยา.\csromanspacer{} จกฺขายตนํ ปจฺจยา จกฺขุวิญฺญาณํ ฯเปฯ.\csromanspacer{} กายายตนํ ปจฺจยา กายวิญฺญาณํ.\csromanspacer{} วตฺถุํ ปจฺจยา อเหตุกา สารมฺมณา ขนฺธา.\csromanspacer{} ปฏิสนฺธิกฺขเณ ฯเปฯ.\csromanspacer{} วตฺถุํ ปจฺจยา วิจิกิจฺฉาสหคโต อุทฺธจฺจสหคโต โมโห. (2)}

\csromanheader{2}{55. สารมฺมณทุก}
\prose{\csromanfolio{375}อนารมฺมณํ ธมฺมํ ปจฺจยา สารมฺมโณ จ อนารมฺมโณ จ ธมฺมา อุปฺปชฺชนฺติ นเหตุปจฺจยา.\csromanspacer{} วตฺถุํ ปจฺจยา อเหตุกา สารมฺมณา ขนฺธา.\csromanspacer{} มหาภูเต ปจฺจยา จิตฺตสมุฏฺฐานํ รูปํ.\csromanspacer{} ปฏิสนฺธิกฺขเณ ฯเปฯ. (3)}

\proseitem{19}{สารมฺมณญฺจ อนารมฺมณญฺจ ธมฺมํ ปจฺจยา สารมฺมโณ ธมฺโม อุปฺปชฺชติ นเหตุปจฺจยา.\csromanspacer{} จกฺขุวิญฺญาณสหคตํ เอกํ ขนฺธญฺจ จกฺขายตนญฺจ ปจฺจยา ตโย ขนฺธา ฯเปฯ เทฺว ขนฺเธ จ ฯเปฯ.\csromanspacer{} กายวิญฺญาณสหคตํ ฯเปฯ.\csromanspacer{} อเหตุกํ สารมฺมณํ เอกํ ขนฺธญฺจ วตฺถุญฺจ ปจฺจยา ตโย ขนฺธา ฯเปฯ เทฺว ขนฺเธ จ ฯเปฯ.\csromanspacer{} อเหตุกปฏิสนฺธิกฺขเณ ฯเปฯ.\csromanspacer{} วิจิกิจฺฉาสหคเต อุทฺธจฺจสหคเต ขนฺเธ จ วตฺถุญฺจ ปจฺจยา วิจิกิจฺฉาสหคโต อุทฺธจฺจสหคโต โมโห. (1)}

\prose{สารมฺมณญฺจ อนารมฺมณญฺจ ธมฺมํ ปจฺจยา อนารมฺมโณ ธมฺโม อุปฺปชฺชติ นเหตุปจฺจยา.\csromanspacer{} อเหตุเก สารมฺมเณ ขนฺเธ จ มหาภูเต จ ปจฺจยา จิตฺตสมุฏฺฐานํ รูปํ.\csromanspacer{} ปฏิสนฺธิกฺขเณ ฯเปฯ. (2)}

\prose{สารมฺมณญฺจ อนารมฺมณญฺจ ธมฺมํ ปจฺจยา สารมฺมโณ จ อนารมฺมโณ จ ธมฺมา อุปฺปชฺชนฺติ นเหตุปจฺจยา.\csromanspacer{} อเหตุกํ สารมฺมณํ เอกํ ขนฺธญฺจ วตฺถุญฺจ ปจฺจยา ตโย ขนฺธา ฯเปฯ เทฺว ขนฺเธ จ ฯเปฯ.\csromanspacer{} อเหตุเก สารมฺมเณ ขนฺเธ จ มหาภูเต จ ปจฺจยา จิตฺตสมุฏฺฐานํ รูปํ.\csromanspacer{} อเหตุกปฏิสนฺธิกฺขเณ อเหตุกํ สารมฺมณํ เอกํ ขนฺธญฺจ วตฺถุญฺจ ปจฺจยา ตโย ขนฺธา ฯเปฯ เทฺว ขนฺเธ จ ฯเปฯ.\csromanspacer{} อเหตุเก สารมฺมเณ ขนฺเธ จ มหาภูเต จ ปจฺจยา กฏตฺตารูปํ ฯเปฯ. (สํขิตฺตํ.) (3)}

\csromanheader{2}{2. ปจฺจยปจฺจนีย 2. สงฺขฺยาวาร}
\csromanheader{2}{\textbf{สุทฺธ}}
\proseitem{20}{นเหตุยา นว, น-อารมฺมเณ ตีณิ, น-อธิปติยา นว, น-อนนฺตเร ตีณิ, นสมนนฺตเร ตีณิ, น-อญฺญมญฺเญ ตีณิ, น-อุปนิสฺสเย ตีณิ, นปุเรชาเต นว, นปจฺฉาชาเต นว, น-อาเสวเน นว, นกมฺเม จตฺตาริ, นวิปาเก นว, น-อาหาเร เอกํ, น-อินฺทฺริเย เอกํ, นฌาเน จตฺตาริ, นมคฺเค นว, นสมฺปยุตฺเต ตีณิ, นวิปฺปยุตฺเต เทฺว, โนนตฺถิยา ตีณิ, โนวิคเต ตีณิ.}

\csromanheader{2}{3. ปจฺจยานุโลมปจฺจนีย}
\proseitem{21}{\csromanfolio{376}เหตุปจฺจยา น-อารมฺมเณ ตีณิ, น-อธิปติยา นว ฯเปฯ นกมฺเม เอกํ ฯเปฯ นวิปฺปยุตฺเต เอกํ. (สํขิตฺตํ.)}

\csromanheader{2}{4. ปจฺจยปจฺจนียานุโลม}
\proseitem{22}{นเหตุปจฺจยา อารมฺมเณ ตีณิ, อนนฺตเร ตีณิ, สมนนฺตเร ตีณิ, สหชาเต นว ฯเปฯ มคฺเค ตีณิ ฯเปฯ อวิคเต นว.}

\vspace{\tipitakaparskip}
\csromancenter{(นิสฺสยวาโร ปจฺจยวารสทิโส.)}
\csromanheader{2}{55. สารมฺมณทุก 5. สํสฏฺฐวาร}
\csromanheader{2}{1. ปจฺจยานุโลมาทิ}
\proseitem{23}{สารมฺมณํ ธมฺมํ สํสฏฺโฐ สารมฺมโณ ธมฺโม อุปฺปชฺชติ เหตุปจฺจยา.\csromanspacer{} สารมฺมณํ เอกํ ขนฺธํ สํสฏฺฐา ตโย ขนฺธา ฯเปฯ เทฺว ขนฺเธ สํสฏฺฐา เทฺว ขนฺธา.\csromanspacer{} ปฏิสนฺธิกฺขเณ ฯเปฯ.}

\proseitem{24}{เหตุยา เอกํ, อารมฺมเณ เอกํ, อธิปติยา เอกํ, (สพฺพตฺถ เอกํ.) อวิคเต เอกํ.\csromanspacer{} อนุโลมํ.}

\proseitem{25}{สารมฺมณํ ธมฺมํ สํสฏฺโฐ สารมฺมโณ ธมฺโม อุปฺปชฺชติ นเหตุปจฺจยา.\csromanspacer{} อเหตุกํ สารมฺมณํ เอกํ ขนฺธํ สํสฏฺฐา ตโย ขนฺธา ฯเปฯ เทฺว ขนฺเธ ฯเปฯ.\csromanspacer{} วิจิกิจฺฉาสหคเต อุทฺธจฺจสหคเต ขนฺเธ สํสฏฺโฐ วิจิกิจฺฉาสหคโต อุทฺธจฺจสหคโต โมโห. (สํขิตฺตํ.)}

\proseitem{26}{นเหตุยา เอกํ, น-อธิปติยา เอกํ, นปุเรชาเต เอกํ, นปจฺฉาชาเต เอกํ, น-อาเสวเน เอกํ, นกมฺเม เอกํ, นวิปาเก เอกํ, นฌาเน เอกํ, นมคฺเค เอกํ, นวิปฺปยุตฺเต เอกํ.}

\vspace{\tipitakaparskip}
\csromancenter{ปจฺจนียํ.}
\csromancenter{(เอวํ อิตเร เทฺว คณนาปิ สมฺปยุตฺตวาโรปิ กาตพฺพา.)}
\csromanheader{2}{55. สารมฺมณทุก \textbf{55. สารมฺมณทุก 7. ปญฺหาวาร}}
\csromanheader{2}{1. ปจฺจยานุโลม 1. วิภงฺควาร}
\csromanheader{2}{\textbf{เหตุ}}
\proseitem{27}{\csromanfolio{377}สารมฺมโณ ธมฺโม สารมฺมณสฺส ธมฺมสฺส เหตุปจฺจเยน ปจฺจโย.\csromanspacer{} สารมฺมณา เหตู สมฺปยุตฺตกานํ ขนฺธานํ เหตุปจฺจเยน ปจฺจโย.\csromanspacer{} ปฏิสนฺธิกฺขเณ ฯเปฯ. (1)}

\prose{สารมฺมโณ ธมฺโม อนารมฺมณสฺส ธมฺมสฺส เหตุปจฺจเยน ปจฺจโย.\csromanspacer{} สารมฺมณา เหตู จิตฺตสมุฏฺฐานานํ รูปานํ เหตุปจฺจเยน ปจฺจโย.\csromanspacer{} ปฏิสนฺธิกฺขเณ ฯเปฯ. (2)}

\prose{สารมฺมโณ ธมฺโม สารมฺมณสฺส จ อนารมฺมณสฺส จ ธมฺมสฺส เหตุปจฺจเยน ปจฺจโย.\csromanspacer{} สารมฺมณา เหตู สมฺปยุตฺตกานํ ขนฺธานํ จิตฺตสมุฏฺฐานานญฺจ รูปานํ เหตุปจฺจเยน ปจฺจโย.\csromanspacer{} ปฏิสนฺธิกฺขเณ. (3)}

\csromanheader{2}{\textbf{อารมฺมณ}}
\proseitem{28}{สารมฺมโณ ธมฺโม สารมฺมณสฺส ธมฺมสฺส อารมฺมณปจฺจเยน ปจฺจโย.\csromanspacer{} ทานํ ฯเปฯ สีลํ ฯเปฯ อุโปสถกมฺมํ กตฺวา ตํ ปจฺจเวกฺขติ, อสฺสาเทติ อภินนฺทติ, ตํ อารพฺภ ราโค อุปฺปชฺชติ ฯเปฯ โทมนสฺสํ อุปฺปชฺชติ.\csromanspacer{} ปุพฺเพ สุจิณฺณานิ ฯเปฯ.\csromanspacer{} ฌานา ฯเปฯ.\csromanspacer{} อริยา มคฺคา วุฏฺฐหิตฺวา มคฺคํ ปจฺจเวกฺขนฺติ, ผลํ ปจฺจเวกฺขนฺติ.\csromanspacer{} ปหีเน กิเลเส ฯเปฯ.\csromanspacer{} วิกฺขมฺภิเต กิเลเส ฯเปฯ.\csromanspacer{} ปุพฺเพ สมุทาจิณฺเณ กิเสเล ชานนฺติ.\csromanspacer{} สารมฺมเณ ขนฺเธ อนิจฺจโต ฯเปฯ โทมนสฺสํ อุปฺปชฺชติ.\csromanspacer{} เจโตปริยญาเณน สารมฺมณจิตฺตสมงฺคิสฺส จิตฺตํ ชานาติ.\csromanspacer{} อากาสานญฺจายตนํ วิญฺญาณญฺจายตนสฺส ฯเปฯ.\csromanspacer{} อากิญฺจญฺญายตนํ เนวสญฺญานาสญฺญายตนสฺส ฯเปฯ.\csromanspacer{} สารมฺมณา ขนฺธา อิทฺธิวิธญาณสฺส, เจโตปริยญาณสฺส, ปุพฺเพนิวาสานุสฺสติญาณสฺส, ยถากมฺมูปคญาณสฺส, อนาคตํสญาณสฺส, อาวชฺชนาย อารมฺมณปจฺจเยน ปจฺจโย. (1)}

\proseitem{29}{อนารมฺมโณ ธมฺโม สารมฺมณสฺส ธมฺมสฺส อารมฺมณปจฺจเยน ปจฺจโย.\csromanspacer{} อริยา นิพฺพานํ ปจฺจเวกฺขนฺติ.\csromanspacer{} นิพฺพานํ โคตฺรภุสฺส, โวทานสฺส, \csromanfolio{378}มคฺคสฺส, ผลสฺส, อาวชฺชนาย อารมฺมณปจฺจเยน ปจฺจโย.\csromanspacer{} จกฺขุํ ฯเปฯ วตฺถุํ อนิจฺจโต ฯเปฯ โทมนสฺสํ อุปฺปชฺชติ.\csromanspacer{} ทิพฺเพน จกฺขุนา รูปํ ปสฺสติ.\csromanspacer{} ทิพฺพาย โสตธาตุยา สทฺทํ สุณาติ.\csromanspacer{} รูปายตนํ จกฺขุวิญฺญาณสฺส ฯเปฯ.\csromanspacer{} โผฏฺฐพฺพายตนํ กายวิญฺญาณสฺส ฯเปฯ.\csromanspacer{} อนารมฺมณา ขนฺธา อิทฺธิวิธญาณสฺส, ปุพฺเพนิวาสานุสฺสติญาณสฺส, อนาคตํสญาณสฺส, อาวชฺชนาย อารมฺมณปจฺจเยน ปจฺจโย. (1)}

\csromanheader{2}{\textbf{อธิปติ}}
\proseitem{30}{สารมฺมโณ ธมฺโม สารมฺมณสฺส ธมฺมสฺส อธิปติปจฺจเยน ปจฺจโย.\csromanspacer{} \textbf{อารมฺมณาธิปติ}, สหชาตาธิปติ.\csromanspacer{}\csromanspacer{}\textbf{อารมฺมณาธิปติ}\csromandash{}ทานํ ฯเปฯ สีลํ ฯเปฯ อุโปสถกมฺมํ กตฺวา ตํ ครุํ กตฺวา ปจฺจเวกฺขติ, อสฺสาเทติ อภินนฺทติ, ตํ ครุํ กตฺวา ราโค อุปฺปชฺชติ.\csromanspacer{} ทิฏฺฐิ อุปฺปชฺชติ.\csromanspacer{} ปุพฺเพ สุจิณฺณานิ ฯเปฯ.\csromanspacer{} ฌานา วุฏฺฐหิตฺวา ฌานํ ฯเปฯ.\csromanspacer{} อริยา มคฺคา วุฏฺฐหิตฺวา มคฺคํ ครุํ ฯเปฯ.\csromanspacer{} ผลํ ครุํ ฯเปฯ.\csromanspacer{} สารมฺมเณ ขนฺเธ ครุํ กตฺวา อสฺสาเทติ อภินนฺทติ, ตํ ครุํ กตฺวา ราโค อุปฺปชฺชติ.\csromanspacer{} ทิฏฺฐิ อุปฺปชฺชติ.\csromanspacer{} \textbf{สหชาตาธิปติ}\csromandash{}สารมฺมณาธิปติ สมฺปยุตฺตกานํ ขนฺธานํ อธิปติปจฺจเยน ปจฺจโย. (1)}

\prose{สารมฺมโณ ธมฺโม อนารมฺมณสฺส ธมฺมสฺส อธิปติปจฺจเยน ปจฺจโย.\csromanspacer{} \textbf{สหชาตาธิปติ}\csromandash{}สารมฺมณาธิปติ จิตฺตสมุฏฺฐานานํ รูปานํ อธิปติปจฺจเยน ปจฺจโย. (2)}

\prose{สารมฺมโณ ธมฺโม สารมฺมณสฺส จ อนารมฺมณสฺส จ ธมฺมสฺส อธิปติปจฺจเยน ปจฺจโย.\csromanspacer{} \textbf{สหชาตาธิปติ}\csromandash{}สารมฺมณาธิปติ สมฺปยุตฺตกานํ ขนฺธานํ จิตฺตสมุฏฺฐานานญฺจ รูปานํ อธิปติปจฺจเยน ปจฺจโย. (3)}

\proseitem{31}{อนารมฺมโณ ธมฺโม สารมฺมณสฺส ธมฺมสฺส อธิปติปจฺจเยน ปจฺจโย.\csromanspacer{} \textbf{อารมฺมณาธิปติ}\csromandash{}อริยา นิพฺพานํ ครุํ กตฺวา ปจฺจเวกฺขนฺติ.\csromanspacer{} นิพฺพานํ โคตฺรภุสฺส, โวทานสฺส, มคฺคสฺส, ผลสฺส อธิปติปจฺจเยน ปจฺจโย.\csromanspacer{} จกฺขุํ ฯเปฯ วตฺถุํ ครุํ กตฺวา อสฺสาเทติ อภินนฺทติ, ตํ ครุํ กตฺวา ราโค อุปฺปชฺชติ.\csromanspacer{} ทิฏฺฐิ อุปฺปชฺชติ. (1)}

\csromanheader{2}{55. สารมฺมณทุก \textbf{อนนฺตราทิ}}
\proseitem{32}{\csromanfolio{379}สารมฺมโณ ธมฺโม สารมฺมณสฺส ธมฺมสฺส อนนฺตรปจฺจเยน ปจฺจโย.\csromanspacer{} ปุริมา ปุริมา สารมฺมณา ฯเปฯ ผลสมาปตฺติยา อนนฺตรปจฺจเยน ปจฺจโย.\csromanspacer{} สมนนฺตรปจฺจเยน ปจฺจโย.\csromanspacer{} เอกํ.\csromanspacer{} สหชาตปจฺจเยน ปจฺจโย.\csromanspacer{} สตฺต. (ปฏิจฺจวาเร สหชาตสทิสา.) อญฺญมญฺญปจฺจเยน ปจฺจโย.\csromanspacer{} ฉ. (ปฏิจฺจวาร อญฺญมญฺญสทิสํ.) นิสฺสยปจฺจเยน ปจฺจโย.\csromanspacer{} สตฺต. (ปจฺจยวาเร นิสฺสยสทิโส.)}

\csromanheader{2}{\textbf{อุปนิสฺสย}}
\proseitem{33}{สารมฺมโณ ธมฺโม สารมฺมณสฺส ธมฺมสฺส อุปนิสฺสยปจฺจเยน ปจฺจโย.\csromanspacer{} \textbf{อารมฺมณูปนิสฺสโย, อนนฺตรูปนิสฺสโย, ปกตูปนิสฺสโย ฯเปฯ.}\csromanspacer{} ปกตูปนิสฺสโย\csromandash{}สทฺธํ อุปนิสฺสาย ทานํ เทติ ฯเปฯ.\csromanspacer{} มานํ ชปฺเปติ.\csromanspacer{} ทิฏฺฐิํ คณฺหาติ.\csromanspacer{} สีลํ ฯเปฯ.\csromanspacer{} ปญฺญํ.\csromanspacer{} ราคํ ฯเปฯ.\csromanspacer{} ปตฺถนํ.\csromanspacer{} กายิกํ สุขํ.\csromanspacer{} กายิกํ ทุกฺขํ อุปนิสฺสาย ทานํ เทติ ฯเปฯ สมาปตฺติํ อุปฺปาเทติ.\csromanspacer{} ปาณํ หนติ ฯเปฯ สํฆํ ภินฺทติ.\csromanspacer{} สทฺธา ฯเปฯ.\csromanspacer{} ปญฺญา.\csromanspacer{} ราโค ฯเปฯ.\csromanspacer{} ปตฺถนา.\csromanspacer{} กายิกํ สุขํ.\csromanspacer{} กายิกํ ทุกฺข สทฺธาย ฯเปฯ ปญฺญาย.\csromanspacer{} ราคสฺส ฯเปฯ ปตฺถนาย.\csromanspacer{} กายิกสฺส สุขสฺส.\csromanspacer{} กายิกสฺส ทุกฺขสฺส.\csromanspacer{} มคฺคสฺส, ผลสมาปตฺติยา อุปนิสฺสยปจฺจเยน ปจฺจโย. (1)}

\prose{อนารมฺมโณ ธมฺโม สารมฺมณสฺส ธมฺมสฺส อุปนิสฺสยปจฺจเยน ปจฺจโย.\csromanspacer{} \textbf{อารมฺมณูปนิสฺสโย, ปกตูปนิสฺสโย ฯเปฯ.}\csromanspacer{} ปกตูปนิสฺสโย\csromandash{} อุตุํ.\csromanspacer{} โภชนํ.\csromanspacer{} เสนาสนํ อุปนิสฺสาย ทานํ เทติ ฯเปฯ สมาปตฺติํ อุปฺปาเทติ.\csromanspacer{} ปาณํ หนติ ฯเปฯ สํฆํ ภินฺทติ.\csromanspacer{} อุตุ.\csromanspacer{} โภชนํ.\csromanspacer{} เสนาสนํ สทฺธาย ฯเปฯ ปตฺถนาย.\csromanspacer{} กายิกสฺส สุขสฺส.\csromanspacer{} กายิกสฺส ทุกฺขสฺส.\csromanspacer{} มคฺคสฺส, ผลสมาปตฺติยา อุปนิสฺสยปจฺจเยน ปจฺจโย. (1)}

\csromanheader{2}{\textbf{ปุเรชาต}}
\proseitem{34}{อนารมฺมโณ ธมฺโม สารมฺมณสฺส ธมฺมสฺส ปุเรชาตปจฺจเยน ปจฺจโย.\csromanspacer{} \textbf{อารมฺมณปุเรชาตํ}, วตฺถุปุเรชาตํ.\csromanspacer{}\csromanspacer{}\textbf{อารมฺมณปุเรชาตํ}\csromandash{}จกฺขุํ ฯเปฯ วตฺถุํ อนิจฺจโต ฯเปฯ โทมนสฺสํ อุปฺปชฺชติ.\csromanspacer{} ทิพฺเพน จกฺขุนา รูปํ \csromanfolio{380}ปสฺสติ.\csromanspacer{} ทิพฺพาย โสตธาตุยา สทฺทํ สุณาติ.\csromanspacer{} รูปายตนํ จกฺขุวิญฺญาณสฺส ฯเปฯ.\csromanspacer{} โผฏฺฐพฺพายตนํ กายวิญฺญาณสฺส ฯเปฯ.\csromanspacer{} \textbf{วตฺถุปุเรชาตํ}\csromandash{}จกฺขายตนํ จกฺขุวิญฺญาณสฺส ฯเปฯ กายายตนํ กายวิญฺญาณสฺส ฯเปฯ.\csromanspacer{} วตฺถุ สารมฺมณานํ ขนฺธานํ ปุเรชาตปจฺจเยน ปจฺจโย. (1)}

\csromanheader{2}{\textbf{ปจฺฉาชาตาเสวน}}
\proseitem{35}{สารมฺมโณ ธมฺโม อนารมฺมณสฺส ธมฺมสฺส ปจฺฉาชาตปจฺจเยน ปจฺจโย.\csromanspacer{} เอกํ.}

\prose{สารมฺมโณ ธมฺโม สารมฺมณสฺส ธมฺมสฺส อาเสวนปจฺจเยน ปจฺจโย.\csromanspacer{} เอกํ.}

\csromanheader{2}{\textbf{กมฺม}}
\proseitem{36}{สารมฺมโณ ธมฺโม สารมฺมณสฺส ธมฺมสฺส กมฺมปจฺจเยน ปจฺจโย.\csromanspacer{} \textbf{สหชาตา}, นานากฺขณิกา.\csromanspacer{}\csromanspacer{}\textbf{สหชาตา} สารมฺมณา เจตนา สมฺปยุตฺตกานํ ขนฺธานํ กมฺมปจฺจเยน ปจฺจโย.\csromanspacer{} ปฏิสนฺธิกฺขเณ ฯเปฯ.\csromanspacer{} \textbf{นานากฺขณิกา} สารมฺมณา เจตนา วิปากานํ ขนฺธานํ กมฺมปจฺจเยน ปจฺจโย. (1)}

\prose{สารมฺมโณ ธมฺโม อนารมฺมณสฺส ธมฺมสฺส กมฺมปจฺจเยน ปจฺจโย.\csromanspacer{} \textbf{สหชาตา}, นานากฺขณิกา.\csromanspacer{}\csromanspacer{}\textbf{สหชาตา} สารมฺมณา เจตนา จิตฺตสมุฏฺฐานานํ รูปานํ กมฺมปจฺจเยน ปจฺจโย.\csromanspacer{} ปฏิสนฺธิกฺขเณ ฯเปฯ.\csromanspacer{} \textbf{นานากฺขณิกา} สารมฺมณา เจตนา กฏตฺตารูปานํ กมฺมปจฺจเยน ปจฺจโย. (2)}

\prose{สารมฺมโณ ธมฺโม สารมฺมณสฺส จ อนารมฺมณสฺส จ ธมฺมสฺส กมฺมปจฺจเยน ปจฺจโย.\csromanspacer{} \textbf{สหชาตา}, นานากฺขณิกา.\csromanspacer{}\csromanspacer{}\textbf{สหชาตา} สารมฺมณา เจตนา สมฺปยุตฺตกานํ ขนฺธานํ จิตฺตสมุฏฺฐานานญฺจ รูปานํ กมฺมปจฺจเยน ปจฺจโย.\csromanspacer{} ปฏิสนฺธิกฺขเณ ฯเปฯ.\csromanspacer{} \textbf{นานากฺขณิกา} สารมฺมณา เจตนา วิปากานํ ขนฺธานํ กฏตฺตา จ รูปานํ กมฺมปจฺจเยน ปจฺจโย. (3)}

\csromanheader{2}{\textbf{วิปาก อาหาร}}
\vspace{\tipitakaparskip}
\csromancenter{สารมฺมโณ ธมฺโม สารมฺมณสฺส ธมฺมสฺส วิปากปจฺจเยน ปจฺจโย.\csromanspacer{} ตีณิ.}
\csromanheader{2}{55. สารมฺมณทุก}
\prose{\csromanfolio{381}สารมฺมโณ ธมฺโม สารมฺมณสฺส ธมฺมสฺส อาหารปจฺจเยน ปจฺจโย.\csromanspacer{} ตีณิ.}

\prose{อนารมฺมโณ ธมฺโม อนารมฺมณสฺส ธมฺมสฺส อาหารปจฺจเยน ปจฺจโย.\csromanspacer{} กพฬีกาโร อาหาโร อิมสฺส กายสฺส อาหารปจฺจเยน ปจฺจโย. (1)}

\csromanheader{2}{\textbf{อินฺทฺริย}}
\proseitem{38}{สารมฺมโณ ธมฺโม สารมฺมณสฺส ธมฺมสฺส อินฺทฺริยปจฺจเยน ปจฺจโย.\csromanspacer{} ตีณิ.}

\prose{อนารมฺมโณ ธมฺโม อนารมฺมณสฺส ธมฺมสฺส อินฺทฺริยปจฺจเยน ปจฺจโย.\csromanspacer{} รูปชีวิตินฺทฺริยํ กฏตฺตารูปานํ อินฺทฺริยปจฺจเยน ปจฺจโย. (1)}

\prose{อนารมฺมโณ ธมฺโม สารมฺมณสฺส ธมฺมสฺส อินฺทฺริยปจฺจเยน ปจฺจโย.\csromanspacer{} จกฺขุนฺทฺริยํ จกฺขุวิญฺญาณสฺส ฯเปฯ.\csromanspacer{} กายินฺทฺริยํ กายวิญฺญาณสฺส อินฺทฺริยปจฺจเยน ปจฺจโย. (1)}

\prose{สารมฺมโณ จ อนารมฺมโณ จ ธมฺมา สารมฺมณสฺส ธมฺมสฺส อินฺทฺริยปจฺจเยน ปจฺจโย.\csromanspacer{} จกฺขุนฺทฺริยญฺจ จกฺขุวิญฺญาณญฺจ จกฺขุวิญฺญาณสหคตานํ ขนฺธานํ อินฺทฺริยปจฺจเยน ปจฺจโย ฯเปฯ.\csromanspacer{} กายินฺทฺริยญฺจ กายวิญฺญาณญฺจ กายวิญฺญาณสหคตานํ ขนฺธานํ อินฺทฺริยปจฺจเยน ปจฺจโย. (1)}

\csromanheader{2}{\textbf{ฌานาทิ}}
\proseitem{39}{สารมฺมโณ ธมฺโม สารมฺมณสฺส ธมฺมสฺส ฌานปจฺจเยน ปจฺจโย.\csromanspacer{} ตีณิ.\csromanspacer{} มคฺคปจฺจเยน ปจฺจโย.\csromanspacer{} ตีณิ.\csromanspacer{} สมฺปยุตฺตปจฺจเยน ปจฺจโย.\csromanspacer{} เอกํ.}

\csromanheader{2}{\textbf{วิปฺปยุตฺต}}
\proseitem{40}{สารมฺมโณ ธมฺโม อนารมฺมณสฺส ธมฺมสฺส วิปฺปยุตฺตปจฺจเยน ปจฺจโย.\csromanspacer{} \textbf{สหชาตํ}, ปจฺฉาชาตํ. (สํขิตฺตํ.) (1)}

\prose{อนารมฺมโณ ธมฺโม สารมฺมณสฺส ธมฺมสฺส วิปฺปยุตฺตปจฺจเยน ปจฺจโย.\csromanspacer{} \textbf{สหชาตํ}, ปุเรชาตํ.\csromanspacer{}\csromanspacer{}\textbf{สหชาตํ} ปฏิสนฺธิกฺขเณ วตฺถุ สารมฺมณานํ \csromanfolio{382}ขนฺธานํ วิปฺปยุตฺตปจฺจเยน ปจฺจโย.\csromanspacer{} \textbf{ปุเรชาตํ} จกฺขายตนํ จกฺขุวิญฺญาณสฺส ฯเปฯ กายายตนํ กายวิญฺญาณสฺส วิปฺปยุตฺตปจฺจเยน ปจฺจโย.\csromanspacer{} วตฺถุ สารมฺมณานํ ขนฺธานํ วิปฺปยุตฺตปจฺจเยน ปจฺจโย. (1)}

\csromanheader{2}{\textbf{อตฺถิ}}
\proseitem{41}{สารมฺมโณ ธมฺโม สารมฺมณสฺส ธมฺมสฺส อตฺถิปจฺจเยน ปจฺจโย.\csromanspacer{} สารมฺมโณ เอโก ขนฺโธ ติณฺณนฺนํ ขนฺธานํ อตฺถิปจฺจเยน ปจฺจโย ฯเปฯ เทฺว ขนฺธา ฯเปฯ.\csromanspacer{} ปฏิสนฺธิกฺขเณ ฯเปฯ. (1)}

\prose{สารมฺมโณ ธมฺโม อนารมฺมณสฺส ธมฺมสฺส อตฺถิปจฺจเยน ปจฺจโย.\csromanspacer{} \textbf{สหชาตํ}, ปจฺฉาชาตํ. (สํขิตฺตํ.) (2)}

\prose{สารมฺมโณ ธมฺโม สารมฺมณสฺส จ อนารมฺมณสฺส จ ธมฺมสฺส อตฺถิปจฺจเยน ปจฺจโย. (ปฏิจฺจสทิสํ.) (3)}

\proseitem{42}{อนารมฺมโณ ธมฺโม อนารมฺมณสฺส ธมฺมสฺส อตฺถิปจฺจเยน ปจฺจโย.\csromanspacer{} \textbf{สหชาตํ}, อาหารํ, อินฺทฺริยํ.\csromanspacer{}\csromanspacer{}\textbf{สหชาตํ} เอกํ มหาภูตํ ฯเปฯ. (ยาว อสญฺญสตฺตา.) (1)}

\prose{อนารมฺมโณ ธมฺโม สารมฺมณสฺส ธมฺมสฺส อตฺถิปจฺจเยน ปจฺจโย.\csromanspacer{} \textbf{สหชาตํ}, ปุเรชาตํ.\csromanspacer{}\csromanspacer{}\textbf{สหชาตํ} ปฏิสนฺธิกฺขเณ วตฺถุ สารมฺมณานํ ขนฺธานํ อตฺถิปจฺจเยน ปจฺจโย.\csromanspacer{} \textbf{ปุเรชาตํ} จกฺขุํ ฯเปฯ วตฺถุํ อนิจฺจโต ฯเปฯ โทมนสฺสํ อุปฺปชฺชติ.\csromanspacer{} ทิพฺเพน จกฺขุนา รูปํ ปสฺสติ.\csromanspacer{} ทิพฺพาย โสตธาตุยา สทฺทํ สุณาติ.\csromanspacer{} รูปายตนํ จกฺขุวิญฺญาณสฺส ฯเปฯ.\csromanspacer{} โผฏฺฐพฺพายตนํ กายวิญฺญาณสฺส ฯเปฯ.\csromanspacer{} จกฺขายตนํ จกฺขุวิญฺญาณสฺส ฯเปฯ.\csromanspacer{} กายายตนํ กายวิญฺญาณสฺส ฯเปฯ.\csromanspacer{} วตฺถุ สารมฺมณานํ ขนฺธานํ อตฺถิปจฺจเยน ปจฺจโย. (2)}

\proseitem{43}{สารมฺมโณ จ อนารมฺมโณ จ ธมฺมา สารมฺมณสฺส ธมฺมสฺส อตฺถิปจฺจเยน ปจฺจโย.\csromanspacer{} \textbf{สหชาตํ}, ปุเรชาตํ.\csromanspacer{}\csromanspacer{}สหชาโต จกฺขุวิญฺญาณสหคโต เอโก ขนฺโธ จ จกฺขายตนญฺจ ติณฺณนฺนํ ขนฺธานํ อตฺถิปจฺจเยน ปจฺจโย ฯเปฯ เทฺว ขนฺธา จ ฯเปฯ.\csromanspacer{} กายวิญฺญาณสหคโต เอโก ขนฺโธ จ กายายตนญฺจ ติณฺณนฺนํ ขนฺธานํ อตฺถิปจฺจเยน ปจฺจโย ฯเปฯ เทฺว ขนฺธา จ ฯเปฯ.\csromanspacer{} สารมฺมโณ เอโก ขนฺโธ จ วตฺถุ จ ติณฺณนฺนํ ขนฺธานํ อตฺถิปจฺจเยน ปจฺจโย ฯเปฯ เทฺว ขนฺธา จ ฯเปฯ.\csromanspacer{} ปฏิสนฺธิกฺขเณ ฯเปฯ. (1)}

\csromanheader{2}{55. สารมฺมณทุก}
\prose{\csromanfolio{383}สารมฺมโณ จ อนารมฺมโณ จ ธมฺมา อนารมฺมณสฺส ธมฺมสฺส อตฺถิปจฺจเยน ปจฺจโย.\csromanspacer{} สหชาตํ, ปจฺฉาชาตํ, อาหารํ, อินฺทฺริยํ.\csromanspacer{}\csromanspacer{}\textbf{สหชาตา} สารมฺมณา ขนฺธา จ มหาภูตา จ จิตฺตสมุฏฺฐานานํ รูปานํ อตฺถิปจฺจเยน ปจฺจโย.\csromanspacer{} ปฏิสนฺธิกฺขเณ ฯเปฯ.\csromanspacer{} \textbf{ปจฺฉาชาตา} สารมฺมณา ขนฺธา จ กพฬีกาโร อาหาโร จ อิมสฺส กายสฺส อตฺถิปจฺจเยน ปจฺจโย.\csromanspacer{} \textbf{ปจฺฉาชาตา} สารมฺมณา ขนฺธา จ รูปชีวิตินฺทฺริยญฺจ กฏตฺตารูปานํ อตฺถิปจฺจเยน ปจฺจโย. (2)}

\csromanheader{2}{1. ปจฺจยานุโลม 2. สงฺขฺยาวาร}
\csromanheader{2}{\textbf{สุทฺธ}}
\proseitem{44}{เหตุยา ตีณิ, อารมฺมเณ เทฺว, อธิปติยา จตฺตาริ, อนนฺตเร เอกํ, สมนนฺตเร เอกํ, สหชาเต สตฺต, อญฺญมญฺเญ ฉ, นิสฺสเย สตฺต, อุปนิสฺสเย เทฺว, ปุเรชาเต เอกํ, ปจฺฉาชาเต เอกํ, อาเสวเน เอกํ, กมฺเม ตีณิ, วิปาเก ตีณิ, อาหาเร จตฺตาริ, อินฺทฺริเย ฉ, ฌาเน ตีณิ, มคฺเค ตีณิ, สมฺปยุตฺเต เอกํ, วิปฺปยุตฺเต เทฺว, อตฺถิยา สตฺต, นตฺถิยา เอกํ, วิคเต เอกํ, อวิคเต สตฺต.}

\csromanheader{2}{2. ปจฺจนียุทฺธาร}
\proseitem{45}{สารมฺมโณ ธมฺโม สารมฺมณสฺส ธมฺมสฺส อารมฺมณปจฺจเยน ปจฺจโย.\csromanspacer{} สหชาตปจฺจเยน ปจฺจโย.\csromanspacer{} อุปนิสฺสยปจฺจเยน ปจฺจโย.\csromanspacer{} กมฺมปจฺจเยน ปจฺจโย. (1)}

\prose{สารมฺมโณ ธมฺโม อนารมฺมณสฺส ธมฺมสฺส สหชาตปจฺจเยน ปจฺจโย.\csromanspacer{} ปจฺฉาชาตปจฺจเยน ปจฺจโย.\csromanspacer{} กมฺมปจฺจเยน ปจฺจโย. (2)}

\prose{สารมฺมโณ ธมฺโม สารมฺมณสฺส จ อนารมฺมณสฺส จ ธมฺมสฺส สหชาตปจฺจเยน ปจฺจโย.\csromanspacer{} กมฺมปจฺจเยน ปจฺจโย. (3)}

\proseitem{46}{อนารมฺมโณ ธมฺโม อนารมฺมณสฺส ธมฺมสฺส สหชาตปจฺจเยน ปจฺจโย.\csromanspacer{} อาหารปจฺจเยน ปจฺจโย.\csromanspacer{} อินฺทฺริยปจฺจเยน ปจฺจโย. (1)}

\prose{\csromanfolio{384}อนารมฺมโณ ธมฺโม สารมฺมณสฺส ธมฺมสฺส อารมฺมณปจฺจเยน ปจฺจโย.\csromanspacer{} สหชาตปจฺจเยน ปจฺจโย.\csromanspacer{} อุปนิสฺสยปจฺจเยน ปจฺจโย.\csromanspacer{} ปุเรชาตปจฺจเยน ปจฺจโย. (2)}

\prose{สารมฺมโณ จ อนารมฺมโณ จ ธมฺมา สารมฺมณสฺส ธมฺมสฺส สหชาตํ, ปุเรชาตํ. (1)}

\prose{สารมฺมโณ จ อนารมฺมโณ จ ธมฺมา อนารมฺมณสฺส ธมฺมสฺส สหชาตํ, ปจฺฉาชาตํ, อาหารํ, อินฺทฺริยํ. (2)}

\csromanheader{2}{2. ปจฺจยปจฺจนีย 2. สงฺขฺยาวาร}
\csromanheader{2}{\textbf{สุทฺธ}}
\proseitem{47}{นเหตุยา สตฺต, น-อารมฺมเณ สตฺต ฯเปฯ นสมนนฺตเร สตฺต, นสหชาเต ฉ, น-อญฺญมญฺเญ ฉ, นนิสฺสเย ฉ, น-อุปนิสฺสเย สตฺต, นปุเรชาเต สตฺต, นปจฺฉาชาเต สตฺต ฯเปฯ นมคฺเค สตฺต, นสมฺปยุตฺเต ฉ, นวิปฺปยุตฺเต ปญฺจ, โน-อตฺถิยา จตฺตาริ, โนนตฺถิยา สตฺต, โนวิคเต สตฺต, โน-อวิคเต จตฺตาริ.}

\csromanheader{2}{3. ปจฺจยานุโลมปจฺจนีย}
\proseitem{48}{เหตุปจฺจยา น-อารมฺมเณ ตีณิ, น-อธิปติยา ฯเปฯ นสมนนฺตเร ตีณิ, น-อญฺญมญฺเญ เอกํ, น-อุปนิสฺสเย ตีณิ ฯเปฯ นมคฺเค ตีณิ, นสมฺปยุตฺเต เอกํ, นวิปฺปยุตฺเต เอกํ, โนนตฺถิยา ตีณิ, โนวิคเต ตีณิ.}

\csromanheader{2}{4. ปจฺจยปจฺจนียานุโลม}
\proseitem{49}{นเหตุปจฺจยา อารมฺมเณ เทฺว, อธิปติยา จตฺตาริ, อนนฺตเร เอกํ, (อนุโลมมาติกา กาตพฺพา.) ฯเปฯ อวิคเต สตฺต.}

\nitthitam{สารมฺมณทุกํ นิฏฺฐิตํ.}
\csromanmatikaanchor{mtk.55}
\csromantocmarkref{subsection}{56. จิตฺตทุก}{mtk.55}
\bookmark[dest={mtk.55},level=2]{56. จิตฺตทุก}
\csromanheader{2}{56. จิตฺตทุก \textbf{56. จิตฺตทุก 1. ปฏิจฺจวาร}}
\csromanheader{2}{1. ปจฺจยานุโลม 1. วิภงฺควาร}
\csromanheader{2}{\textbf{เหตุ}}
\proseitem{50}{\csromanfolio{385}จิตฺตํ ธมฺมํ ปฏิจฺจ โนจิตฺโต ธมฺโม อุปฺปชฺชติ เหตุปจฺจยา.\csromanspacer{} จิตฺตํ ปฏิจฺจ สมฺปยุตฺตกา ขนฺธา จิตฺตสมุฏฺฐานญฺจ รูปํ.\csromanspacer{} ปฏิสนฺธิกฺขเณ จิตฺตํ ปฏิจฺจ สมฺปยุตฺตกา ขนฺธา กฏตฺตา จ รูปํ. (1)}

\prose{โนจิตฺตํ ธมฺมํ ปฏิจฺจ โนจิตฺโต ธมฺโม อุปฺปชฺชติ เหตุปจฺจยา.\csromanspacer{} โนจิตฺตํ เอกํ ขนฺธํ ปฏิจฺจ เทฺว ขนฺธา จิตฺตสมุฏฺฐานญฺจ รูปํ.\csromanspacer{} เทฺว ขนฺเธ ปฏิจฺจ เอโก ขนฺโธ จิตฺตสมุฏฺฐานญฺจ รูปํ.\csromanspacer{} ปฏิสนฺธิกฺขเณ ฯเปฯ.\csromanspacer{} ขนฺเธ ปฏิจฺจ วตฺถุ, วตฺถุํ ปฏิจฺจ ขนฺธา.\csromanspacer{} เอกํ มหาภูตํ ปฏิจฺจ ฯเปฯ. (1)}

\prose{โนจิตฺตํ ธมฺมํ ปฏิจฺจ จิตฺโต ธมฺโม อุปฺปชฺชติ เหตุปจฺจยา.\csromanspacer{} โนจิตฺเต ขนฺเธ ปฏิจฺจ จิตฺตํ.\csromanspacer{} ปฏิสนฺธิกฺขเณ โนจิตฺเต ขนฺเธ ปฏิจฺจ จิตฺตํ.\csromanspacer{} ปฏิสนฺธิกฺขเณ วตฺถุํ ปฏิจฺจ จิตฺตํ. (2)}

\prose{โนจิตฺตํ ธมฺมํ ปฏิจฺจ จิตฺโต จ โนจิตฺโต จ ธมฺมา อุปฺปชฺชนฺติ เหตุปจฺจยา.\csromanspacer{} โนจิตฺตํ เอกํ ขนฺธํ ปฏิจฺจ เทฺว ขนฺธา จิตฺตญฺจ จิตฺตสมุฏฺฐานญฺจ รูปํ.\csromanspacer{} เทฺว ขนฺเธ ฯเปฯ.\csromanspacer{} ปฏิสนฺธิกฺขเณ โนจิตฺตํ เอกํ ขนฺธํ ปฏิจฺจ เทฺว ขนฺธา จิตฺตญฺจ กฏตฺตา จ รูปํ.\csromanspacer{} เทฺว ขนฺเธ ปฏิจฺจ เอโก ขนฺโธ จิตฺตญฺจ กฏตฺตา จ รูปํ.\csromanspacer{} ปฏิสนฺธิกฺขเณ วตฺถุํ ปฏิจฺจ จิตฺตญฺจ สมฺปยุตฺตกา จ ขนฺธา. (3)}

\prose{จิตฺตญฺจ โนจิตฺตญฺจ ธมฺมํ ปฏิจฺจ โนจิตฺโต ธมฺโม อุปฺปชฺชติ เหตุปจฺจยา.\csromanspacer{} โนจิตฺตํ เอกํ ขนฺธญฺจ จิตฺตญฺจ ปฏิจฺจ เทฺว ขนฺธา จิตฺตสมุฏฺฐานญฺจ รูปํ.\csromanspacer{} เทฺว ขนฺเธ จ ฯเปฯ.\csromanspacer{} จิตฺตญฺจ มหาภูเต จ ปฏิจฺจ จิตฺตสมุฏฺฐานํ รูปํ.\csromanspacer{} ปฏิสนฺธิกฺขเณ โนจิตฺตํ เอกํ ขนฺธญฺจ จิตฺตญฺจ ปฏิจฺจ เทฺว ขนฺธา กฏตฺตา จ รูปํ.\csromanspacer{} เทฺว ขนฺเธ ฯเปฯ.\csromanspacer{} ปฏิสนฺธิกฺขเณ จิตฺตญฺจ วตฺถุญฺจ ปฏิจฺจ โนจิตฺตา ขนฺธา.\csromanspacer{} ปฏิสนฺธิกฺขเณ จิตฺตญฺจ มหาภูเต จ ปฏิจฺจ กฏตฺตารูปํ. (1)}

\csromanheader{2}{\textbf{อารมฺมณ}}
\proseitem{51}{จิตฺตํ ธมฺมํ ปฏิจฺจ โนจิตฺโต ธมฺโม อุปฺปชฺชติ อารมฺมณปจฺจยา.\csromanspacer{} จิตฺตํ ปฏิจฺจ สมฺปยุตฺตกา ขนฺธา.\csromanspacer{} ปฏิสนฺธิกฺขเณ ฯเปฯ. (1)}

\prose{\csromanfolio{386}โนจิตฺตํ ธมฺมํ ปฏิจฺจ โนจิตฺโต ธมฺโม อุปฺปชฺชติ อารมฺมณปจฺจยา.\csromanspacer{} โนจิตฺตํ เอกํ ขนฺธํ ปฏิจฺจ เทฺว ขนฺธา.\csromanspacer{} เทฺว ขนฺเธ ฯเปฯ.\csromanspacer{} ปฏิสนฺธิกฺขเณ ฯเปฯ.\csromanspacer{} ปฏิสนฺธิกฺขเณ วตฺถุํ ปฏิจฺจ ขนฺธา. (1)}

\prose{โนจิตฺตํ ธมฺมํ ปฏิจฺจ จิตฺโต ธมฺโม อุปฺปชฺชติ อารมฺมณปจฺจยา.\csromanspacer{} โนจิตฺเต ขนฺเธ ปฏิจฺจ จิตฺตํ.\csromanspacer{} ปฏิสนฺธิกฺขเณ โนจิตฺเต ขนฺเธ ปฏิจฺจ จิตฺตํ.\csromanspacer{} ปฏิสนฺธิกฺขเณ วตฺถุํ ปฏิจฺจ จิตฺตํ. (2)}

\prose{โนจิตฺตํ ธมฺมํ ปฏิจฺจ จิตฺโต จ โนจิตฺโต จ ธมฺมา อุปฺปชฺชนฺติ อารมฺมณปจฺจยา.\csromanspacer{} โนจิตฺตํ เอกํ ขนฺธํ ปฏิจฺจ เทฺว ขนฺธา จิตฺตญฺจ.\csromanspacer{} เทฺว ขนฺเธ ปฏิจฺจ ฯเปฯ.\csromanspacer{} ปฏิสนฺธิกฺขเณ โนจิตฺตํ เอกํ ขนฺธํ ปฏิจฺจ เทฺว ขนฺธา จิตฺตญฺจ.\csromanspacer{} เทฺว ขนฺเธ ฯเปฯ.\csromanspacer{} ปฏิสนฺธิกฺขเณ วตฺถุํ ปฏิจฺจ จิตฺตญฺจ สมฺปยุตฺตกา จ ขนฺธา. (3)}

\prose{จิตฺตญฺจ โนจิตฺตญฺจ ธมฺมํ ปฏิจฺจ โนจิตฺโต ธมฺโม อุปฺปชฺชติ อารมฺมณปจฺจยา.\csromanspacer{} โนจิตฺตํ เอกํ ขนฺธญฺจ จิตฺตญฺจ ปฏิจฺจ เทฺว ขนฺธา.\csromanspacer{} เทฺว ขนฺเธ จ ฯเปฯ.\csromanspacer{} ปฏิสนฺธิกฺขเณ ฯเปฯ.\csromanspacer{} ปฏิสนฺธิกฺขเณ จิตฺตญฺจ วตฺถุญฺจ ปฏิจฺจ โนจิตฺตา ขนฺธา. (1) (สํขิตฺตํ.)}

\csromanheader{2}{1. ปจฺจยานุโลม 2. สงฺขฺยาวาร}
\csromanheader{2}{\textbf{สุทฺธ}}
\proseitem{52}{เหตุยา ปญฺจ, อารมฺมเณ ปญฺจ, อธิปติยา ปญฺจ, อนนฺตเร ปญฺจ, สมนนฺตเร ปญฺจ, สหชาเต ปญฺจ, อญฺญมญฺเญ ปญฺจ, นิสฺสเย ปญฺจ, อุปนิสฺสเย ปญฺจ, ปุเรชาเต ปญฺจ, อาเสวเน ปญฺจ, กมฺเม ปญฺจ, วิปาเก ปญฺจ, อาหาเร ปญฺจ, อินฺทฺริเย ปญฺจ, ฌาเน ปญฺจ, มคฺเค ปญฺจ, สมฺปยุตฺเต ปญฺจ, วิปฺปยุตฺเต ปญฺจ, อตฺถิยา ปญฺจ, นตฺถิยา ปญฺจ, วิคเต ปญฺจ, อวิคเต ปญฺจ.}

\csromanheader{2}{2. ปจฺจยปจฺจนีย 1. วิภงฺควาร}
\csromanheader{2}{\textbf{นเหตุ}}
\proseitem{53}{จิตฺตํ ธมฺมํ ปฏิจฺจ โนจิตฺโต ธมฺโม อุปฺปชฺชติ นเหตุปจฺจยา.\csromanspacer{} อเหตุกํ จิตฺตํ ปฏิจฺจ สมฺปยุตฺตกา ขนฺธา จิตฺตสมุฏฺฐานญฺจ รูปํ.\csromanspacer{} อเหตุกปฏิสนฺธิกฺขเณ จิตฺตํ ฯเปฯ.\csromanspacer{} วิจิกิจฺฉาสหคตํ อุทฺธจฺจสหคตํ จิตฺตํ ปฏิจฺจ วิจิกิจฺฉาสหคโต อุทฺธจฺจสหคโต โมโห. (1)}

\csromanheader{2}{56. จิตฺตทุก}
\prose{\csromanfolio{387}โนจิตฺตํ ธมฺมํ ปฏิจฺจ โนจิตฺโต ธมฺโม อุปฺปชฺชติ นเหตุปจฺจยา.\csromanspacer{} อเหตุกํ โนจิตฺตํ เอกํ ขนฺธํ ปฏิจฺจ เทฺว ขนฺธา จิตฺตสมุฏฺฐานญฺจ รูปํ.\csromanspacer{} เทฺว ขนฺเธ ฯเปฯ.\csromanspacer{} อเหตุกปฏิสนฺธิกฺขเณ ฯเปฯ. (ยาว อสญฺญสตฺตา.) วิจิกิจฺฉาสหคเต อุทฺธจฺจสหคเต ขนฺเธ ปฏิจฺจ วิจิกิจฺฉาสหคโต อุทฺธจฺจสหคโต โมโห. (1)}

\prose{โนจิตฺตํ ธมฺมํ ปฏิจฺจ จิตฺโต ธมฺโม อุปฺปชฺชติ นเหตุปจฺจยา.\csromanspacer{} อเหตุเก โนจิตฺเต ขนฺเธ ปฏิจฺจ จิตฺตํ.\csromanspacer{} อเหตุกปฏิสนฺธิกฺขเณ โนจิตฺเต ขนฺเธ ปฏิจฺจ จิตฺตํ.\csromanspacer{} อเหตุกปฏิสนฺธิกฺขเณ วตฺถุํ ปฏิจฺจ จิตฺตํ.}

\prose{โนจิตฺตํ ธมฺมํ ปฏิจฺจ จิตฺโต จ โนจิตฺโต จ ธมฺมา อุปฺปชฺชนฺติ นเหตุปจฺจยา.\csromanspacer{} อเหตุกํ โนจิตฺตํ เอกํ ขนฺธํ ปฏิจฺจ เทฺว ขนฺธา จิตฺตญฺจ จิตฺตสมุฏฺฐานญฺจ รูปํ.\csromanspacer{} เทฺว ขนฺเธ ฯเปฯ.\csromanspacer{} อเหตุกปฏิสนฺธิกฺขเณ ฯเปฯ.\csromanspacer{} อเหตุกปฏิสนฺธิกฺขเณ วตฺถุํ ปฏิจฺจ จิตฺตญฺจ สมฺปยุตฺตกา จ ขนฺธา. (3)}

\prose{จิตฺตญฺจ โนจิตฺตญฺจ ธมฺมํ ปฏิจฺจ โนจิตฺโต ธมฺโม อุปฺปชฺชติ นเหตุปจฺจยา.\csromanspacer{} อเหตุกํ โนจิตฺตํ เอกํ ขนฺธญฺจ จิตฺตญฺจ ปฏิจฺจ เทฺว ขนฺธา จิตฺตสมุฏฺฐานญฺจ รูปํ.\csromanspacer{} เทฺว ขนฺเธ จ ฯเปฯ.\csromanspacer{} อเหตุกปฏิสนฺธิกฺขเณ โนจิตฺตํ เอกํ ขนฺธญฺจ จิตฺตญฺจ ปฏิจฺจ เทฺว ขนฺธา กฏตฺตา จ รูปํ.\csromanspacer{} เทฺว ขนฺเธ จ ฯเปฯ.\csromanspacer{} อเหตุกปฏิสนฺธิกฺขเณ จิตฺตญฺจ วตฺถุญฺจ ปฏิจฺจ โนจิตฺตา ขนฺธา.\csromanspacer{} วิจิกิจฺฉาสหคตํ อุทฺธจฺจสหคตํ จิตฺตญฺจ สมฺปยุตฺตเก จ ขนฺเธ ปฏิจฺจ วิจิกิจฺฉาสหคโต อุทฺธจฺจสหคโต โมโห. (1)}

\csromanheader{2}{\textbf{น-อารมฺมณ}}
\proseitem{54}{จิตฺตํ ธมฺมํ ปฏิจฺจ โนจิตฺโต ธมฺโม อุปฺปชฺชติ น-อารมฺมณปจฺจยา.\csromanspacer{} จิตฺตํ ปฏิจฺจ จิตฺตสมุฏฺฐานํ รูปํ.\csromanspacer{} ปฏิสนฺธิกฺขเณ ฯเปฯ. (1)}

\prose{โนจิตฺตํ ธมฺมํ ปฏิจฺจ โนจิตฺโต ธมฺโม อุปฺปชฺชติ น-อารมฺมณปจฺจยา.\csromanspacer{} โนจิตฺเต ขนฺเธ ปฏิจฺจ จิตฺตสมุฏฺฐานํ รูปํ.\csromanspacer{} ปฏิสนฺธิกฺขเณ ฯเปฯ. (ยาว อสญฺญสตฺตา.) (2)}

\prose{จิตฺตญฺจ โนจิตฺตญฺจ ธมฺมํ ปฏิจฺจ โนจิตฺโต ธมฺโม อุปฺปชฺชติ น-อารมฺมณปจฺจยา.\csromanspacer{} จิตฺตญฺจ สมฺปยุตฺตเก จ ขนฺเธ ปฏิจฺจ จิตฺตสมุฏฺฐานํ รูปํ.\csromanspacer{} จิตฺตญฺจ มหาภูเต จ ปฏิจฺจ จิตฺตสมุฏฺฐานํ รูปํ.\csromanspacer{} ปฏิสนฺธิกฺขเณ จิตฺตญฺจ สมฺปยุตฺตเก จ ขนฺเธ ปฏิจฺจ กฏตฺตารูปํ.\csromanspacer{} จิตฺตญฺจ มหาภูเต จ ปฏิจฺจ กฏตฺตารูปํ. (1)}

\csromanheader{2}{\textbf{น-อธิปตฺยาทิ}}
\proseitem{55}{\csromanfolio{388}จิตฺตํ ธมฺมํ ปฏิจฺจ โนจิตฺโต ธมฺโม อุปฺปชฺชติ น-อธิปติปจฺจยา.\csromanspacer{} ปญฺจ.\csromanspacer{} น-อนนฺตรปจฺจยา ฯเปฯ.\csromanspacer{} น-อุปนิสฺสยปจฺจยา.\csromanspacer{} ตีณิ.}

\csromanheader{2}{\textbf{นปุเรชาตาทิ}}
\proseitem{56}{จิตฺตํ ธมฺมํ ปฏิจฺจ โนจิตฺโต ธมฺโม อุปฺปชฺชติ นปุเรชาตปจฺจยา.\csromanspacer{} อรูเป จิตฺตํ ปฏิจฺจ สมฺปยุตฺตกา ขนฺธา.\csromanspacer{} จิตฺตํ ปฏิจฺจ จิตฺตสมุฏฺฐานํ รูปํ.\csromanspacer{} ปฏิสนฺธิกฺขเณ จิตฺตํ ปฏิจฺจ สมฺปยุตฺตกา ขนฺธา กฏตฺตา จ รูปํ. (1)}

\prose{โนจิตฺตํ ธมฺมํ ปฏิจฺจ โนจิตฺโต ธมฺโม อุปฺปชฺชติ นปุเรชาตปจฺจยา.\csromanspacer{} อรูเป โนจิตฺตํ เอกํ ขนฺธํ ปฏิจฺจ เทฺว ขนฺธา.\csromanspacer{} เทฺว ขนฺเธ ฯเปฯ.\csromanspacer{} โนจิตฺเต ขนฺเธ ปฏิจฺจ จิตฺตสมุฏฺฐานํ รูปํ.\csromanspacer{} ปฏิสนฺธิกฺขเณ ฯเปฯ. (ยาว อสญฺญสตฺตา.) (1)}

\prose{โนจิตฺตํ ธมฺมํ ปฏิจฺจ จิตฺโต ธมฺโม อุปฺปชฺชติ นปุเรชาตปจฺจยา.\csromanspacer{} อรูเป โนจิตฺเต ขนฺเธ ปฏิจฺจ จิตฺตํ.\csromanspacer{} ปฏิสนฺธิกฺขเณ.\csromanspacer{} วตฺถุํ ปฏิจฺจ จิตฺตํ. (2)}

\prose{โนจิตฺตํ ธมฺมํ ปฏิจฺจ จิตฺโต จ โนจิตฺโต จ ธมฺมา อุปฺปชฺชนฺติ นปุเรชาตปจฺจยา.\csromanspacer{} อรูเป โนจิตฺตํ เอกํ ขนฺธํ ปฏิจฺจ เทฺว ขนฺธา จิตฺตญฺจ.\csromanspacer{} เทฺว ขนฺเธ ฯเปฯ.\csromanspacer{} ปฏิสนฺธิกฺขเณ วตฺถุํ ปฏิจฺจ จิตฺตญฺจ สมฺปยุตฺตกา จ ขนฺธา. (3)}

\prose{จิตฺตญฺจ โนจิตฺตญฺจ ธมฺมํ ปฏิจฺจ โนจิตฺโต ธมฺโม อุปฺปชฺชติ นปุเรชาตปจฺจยา.\csromanspacer{} อรูเป โนจิตฺตํ เอกํ ขนฺธญฺจ จิตฺตญฺจ ปฏิจฺจ เทฺว ขนฺธา.\csromanspacer{} เทฺว ขนฺเธ จ ฯเปฯ.\csromanspacer{} โนจิตฺเต ขนฺเธ จ จิตฺตญฺจ ปฏิจฺจ จิตฺตสมุฏฺฐานํ รูปํ.\csromanspacer{} จิตฺตญฺจ มหาภูเต จ ปฏิจฺจ จิตฺตสมุฏฺฐานํ รูปํ.\csromanspacer{} ปฏิสนฺธิกฺขเณ จิตฺตญฺจ วตฺถุญฺจ ปฏิจฺจ โนจิตฺตา ขนฺธา.\csromanspacer{} ปฏิสนฺธิกฺขเณ จิตฺตญฺจ สมฺปยุตฺตเก จ ขนฺเธ ปฏิจฺจ กฏตฺตารูปํ.\csromanspacer{} จิตฺตญฺจ มหาภูเต จ ปฏิจฺจ กฏตฺตารูปํ. (1)}

\prose{นปจฺฉาชาตปจฺจยา.\csromanspacer{} น-อาเสวนปจฺจยา.}

\csromanheader{2}{\textbf{นกมฺม}}
\proseitem{57}{จิตฺตํ ธมฺมํ ปฏิจฺจ โนจิตฺโต ธมฺโม อุปฺปชฺชติ นกมฺมปจฺจยา.\csromanspacer{} จิตฺตํ ปฏิจฺจ สมฺปยุตฺตกา เจตนา. (1)}

\csromanheader{2}{56. จิตฺตทุก}
\prose{\csromanfolio{389}โนจิตฺตํ ธมฺมํ ปฏิจฺจ โนจิตฺโต ธมฺโม อุปฺปชฺชติ นกมฺมปจฺจยา.\csromanspacer{} โนจิตฺเต ขนฺเธ ปฏิจฺจ สมฺปยุตฺตกา เจตนา.\csromanspacer{} พาหิรํ.\csromanspacer{} อาหารสมุฏฺฐานํ.\csromanspacer{} อุตุสมุฏฺฐานํ ฯเปฯ. (1)}

\prose{จิตฺตญฺจ โนจิตฺตญฺจ ธมฺมํ ปฏิจฺจ โนจิตฺโต ธมฺโม อุปฺปชฺชติ นกมฺมปจฺจยา.\csromanspacer{} โนจิตฺเต ขนฺเธ จ จิตฺตญฺจ ปฏิจฺจ สมฺปยุตฺตกา เจตนา. (สํขิตฺตํ.) (1)}

\csromanheader{2}{2. ปจฺจยปจฺจนีย 2. สงฺขฺยาวาร}
\csromanheader{2}{\textbf{สุทฺธ}}
\proseitem{58}{นเหตุยา ปญฺจ, น-อารมฺมเณ ตีณิ, น-อธิปติยา ปญฺจ, น-อนนฺตเร ตีณิ, นสมนนฺตเร ตีณิ, น-อญฺญมญฺเญ ตีณิ, น-อุปนิสฺสเย ตีณิ, นปุเรชาเต ปญฺจ, นปจฺฉาชาเต ปญฺจ, น-อาเสวเน ปญฺจ, นกมฺเม ตีณิ, นวิปาเก ปญฺจ, น-อาหาเร เอกํ, น-อินฺทฺริเย เอกํ, นฌาเน ปญฺจ, นมคฺเค ปญฺจ, นสมฺปยุตฺเต ตีณิ, นวิปฺปยุตฺเต ปญฺจ, โนนตฺถิยา ตีณิ, โนวิคเต ตีณิ.}

\csromanheader{2}{3. ปจฺจยานุโลมปจฺจนีย}
\proseitem{59}{\textbf{เหตุ}ปจฺจยา น-อารมฺมเณ ตีณิ, น-อธิปติยา ปญฺจ. (สํขิตฺตํ.)}

\csromanheader{2}{4. ปจฺจยปจฺจนียานุโลม}
\proseitem{60}{นเหตุปจฺจยา อารมฺมเณ ปญฺจ, อนนฺตเร ปญฺจ, (สพฺพตฺถ ปญฺจ,) มคฺเค ตีณิ ฯเปฯ อวิคเต ปญฺจ.}

\vspace{\tipitakaparskip}
\csromancenter{(สหชาตวาโร ปฏิจฺจวารสทิโส.)}
\csromanheader{2}{56. จิตฺตทุก 3. ปจฺจยวาร}
\csromanheader{2}{1. ปจฺจยานุโลม 1. วิภงฺควาร}
\csromanheader{2}{\textbf{เหตุ}}
\proseitem{61}{จิตฺตํ ธมฺมํ ปจฺจยา โนจิตฺโต ธมฺโม อุปฺปชฺชติ เหตุปจฺจยา.\csromanspacer{} จิตฺตํ ปจฺจยา สมฺปยุตฺตกา ขนฺธา จิตฺตสมุฏฺฐานญฺจ รูปํ.\csromanspacer{} ปฏิสนฺธิกฺขเณ ฯเปฯ. (1)}

\prose{\csromanfolio{390}โนจิตฺตํ ธมฺมํ ปจฺจยา โนจิตฺโต ธมฺโม อุปฺปชฺชติ เหตุปจฺจยา.\csromanspacer{} โนจิตฺตํ เอกํ ขนฺธํ ปจฺจยา เทฺว ขนฺธา จิตฺตสมุฏฺฐานญฺจ รูปํ.\csromanspacer{} เทฺว ขนฺเธ ฯเปฯ.\csromanspacer{} ปฏิสนฺธิกฺขเณ ฯเปฯ. (ยาว มหาภูตา.) วตฺถุํ ปจฺจยา โนจิตฺตา ขนฺธา. (1)}

\prose{โนจิตฺตํ ธมฺมํ ปจฺจยา จิตฺโต ธมฺโม อุปฺปชฺชติ เหตุปจฺจยา.\csromanspacer{} โนจิตฺเต ขนฺเธ ปจฺจยา จิตฺตํ.\csromanspacer{} วตฺถุํ ปจฺจยา จิตฺตํ.\csromanspacer{} ปฏิสนฺธิกฺขเณ โนจิตฺเต ขนฺเธ ปจฺจยา จิตฺตํ.\csromanspacer{} ปฏิสนฺธิกฺขเณ วตฺถุํ ปจฺจยา จิตฺตํ.}

\prose{โนจิตฺตํ ธมฺมํ ปจฺจยา จิตฺโต จ โนจิตฺโต จ ธมฺมา อุปฺปชฺชนฺติ เหตุปจฺจยา.\csromanspacer{} โนจิตฺตํ เอกํ ขนฺธํ ปจฺจยา เทฺว ขนฺธา จิตฺตญฺจ จิตฺตสมุฏฺฐานญฺจ รูปํ.\csromanspacer{} เทฺว ขนฺเธ ฯเปฯ.\csromanspacer{} วตฺถุํ ปจฺจยา จิตฺตญฺจ สมฺปยุตฺตกา จ ขนฺธา.\csromanspacer{} ปฏิสนฺธิกฺขเณ โนจิตฺตํ เอกํ ขนฺธํ ปจฺจยา เทฺว ขนฺธา จิตฺตญฺจ กฏตฺตา จ รูปํ.\csromanspacer{} เทฺว ขนฺเธ ฯเปฯ.\csromanspacer{} ปฏิสนฺธิกฺขเณ วตฺถุํ ปจฺจยา จิตฺตญฺจ สมฺปยุตฺตกา จ ขนฺธา. (3)}

\prose{จิตฺตญฺจ โนจิตฺตญฺจ ธมฺมํ ปจฺจยา โนจิตฺโต ธมฺโม อุปฺปชฺชติ เหตุปจฺจยา.\csromanspacer{} โนจิตฺตํ เอกํ ขนฺธญฺจ จิตฺตญฺจ ปจฺจยา เทฺว ขนฺธา จิตฺตสมุฏฺฐานญฺจ รูปํ.\csromanspacer{} เทฺว ขนฺเธ จ ฯเปฯ.\csromanspacer{} จิตฺตญฺจ มหาภูเต จ ปจฺจยา จิตฺตสมุฏฺฐานํ รูปํ.\csromanspacer{} จิตฺตญฺจ วตฺถุญฺจ ปจฺจยา โนจิตฺตา ขนฺธา.\csromanspacer{} ปฏิสนฺธิกฺขเณ โนจิตฺตํ เอกํ ขนฺธญฺจ จิตฺตญฺจ ปจฺจยา เทฺว ขนฺธา กฏตฺตา จ รูปํ.\csromanspacer{} เทฺว ขนฺเธ จ ฯเปฯ.\csromanspacer{} ปฏิสนฺธิกฺขเณ จิตฺตญฺจ มหาภูเต จ ปจฺจยา กฏตฺตารูปํ.\csromanspacer{} ปฏิสนฺธิกฺขเณ จิตฺตญฺจ วตฺถุญฺจ ปจฺจยา โนจิตฺตา ขนฺธา. (1)}

\csromanheader{2}{\textbf{อารมฺมณ}}
\proseitem{62}{จิตฺตํ ธมฺมํ ปจฺจยา โนจิตฺโต ธมฺโม อุปฺปชฺชติ อารมฺมณปจฺจยา.\csromanspacer{} เอกํ. (ปฏิจฺจวารสทิสํ.)}

\prose{โนจิตฺตํ ธมฺมํ ปจฺจยา โนจิตฺโต ธมฺโม อุปฺปชฺชติ อารมฺมณปจฺจยา.\csromanspacer{} โนจิตฺตํ เอกํ ขนฺธํ ปจฺจยา เทฺว ขนฺธา.\csromanspacer{} เทฺว ขนฺเธ ฯเปฯ.\csromanspacer{} ปฏิสนฺธิกฺขเณ ฯเปฯ.\csromanspacer{} ปฏิสนฺธิกฺขเณ วตฺถุํ ปจฺจยา ขนฺธา.\csromanspacer{} จกฺขายตนํ ปจฺจยา จกฺขุวิญฺญาณสหคตา ขนฺธา ฯเปฯ.\csromanspacer{} กายายตนํ ปจฺจยา กายวิญฺญาณสหคตา ขนฺธา.\csromanspacer{} วตฺถุํ ปจฺจยา โนจิตฺตา ขนฺธา. (1)}

\prose{โนจิตฺตํ ธมฺมํ ปจฺจยา จิตฺโต ธมฺโม อุปฺปชฺชติ อารมฺมณปจฺจยา.\csromanspacer{} โนจิตฺเต ขนฺเธ ปจฺจยา จิตฺตํ.\csromanspacer{} วตฺถุํ ปจฺจยา จิตฺตํ.\csromanspacer{} ปฏิสนฺธิกฺขเณ}

\vspace{\tipitakaparskip}
\csromancenter{จิตฺตทุก โนจิตฺเต ขนฺเธ ปจฺจยา จิตฺตํ.\csromanspacer{} ปฏิสนฺธิกฺขเณ วตฺถุํ ปจฺจยา จิตฺตํ.\csromanspacer{} จกฺขายตนํ ปจฺจยา จกฺขุวิญฺญาณํ ฯเปฯ.\csromanspacer{} กายายตนํ ปจฺจยา กายวิญฺญาณํ. (2)}
\prose{\csromanfolio{391}โนจิตฺตํ ธมฺมํ ปจฺจยา จิตฺโต จ โนจิตฺโต จ ธมฺมา อุปฺปชฺชนฺติ อารมฺมณปจฺจยา.\csromanspacer{} โนจิตฺตํ เอกํ ขนฺธํ ปจฺจยา เทฺว ขนฺธา จิตฺตญฺจ.\csromanspacer{} เทฺว ขนฺเธ ฯเปฯ.\csromanspacer{} วตฺถุํ ปจฺจยา จิตฺตญฺจ สมฺปยุตฺตกา จ ขนฺธา.\csromanspacer{} ปฏิสนฺธิกฺขเณ ฯเปฯ.\csromanspacer{} ปฏิสนฺธิกฺขเณ วตฺถุํ ปจฺจยา จิตฺตญฺจ สมฺปยุตฺตกา จ ขนฺธา.\csromanspacer{} จกฺขายตนํ ปจฺจยา จกฺขุวิญฺญาณํ สมฺปยุตฺตกา จ ขนฺธา ฯเปฯ.\csromanspacer{} กายายตนํ ปจฺจยา กายวิญฺญาณํ สมฺปยุตฺตกา จ ขนฺธา. (3)}

\prose{จิตฺตญฺจ โนจิตฺตญฺจ ธมฺมํ ปจฺจยา โนจิตฺโต ธมฺโม อุปฺปชฺชติ อารมฺมณปจฺจยา.\csromanspacer{} โนจิตฺตํ เอกํ ขนฺธญฺจ จิตฺตญฺจ ปจฺจยา เทฺว ขนฺธา.\csromanspacer{} เทฺว ขนฺเธ จ ฯเปฯ.\csromanspacer{} จิตฺตญฺจ วตฺถุญฺจ ปจฺจยา โนจิตฺตา ขนฺธา.\csromanspacer{} ปฏิสนฺธิกฺขเณ ฯเปฯ.\csromanspacer{} ปฏิสนฺธิกฺขเณ จิตฺตญฺจ วตฺถุญฺจ ปจฺจยา โนจิตฺตา ขนฺธา.\csromanspacer{} จกฺขายตนญฺจ จกฺขุวิญฺญาณญฺจ ปจฺจยา จกฺขุวิญฺญาณสหคตา ขนฺธา ฯเปฯ.\csromanspacer{} กายายตนญฺจ. (สํขิตฺตํ.) (1)}

\csromanheader{2}{1. ปจฺจยานุโลม 2. สงฺขฺยาวาร}
\proseitem{63}{เหตุยา ปญฺจ, อารมฺมเณ ปญฺจ, อธิปติยา ปญฺจ, (สพฺพตฺถ ปญฺจ,) อวิคเต ปญฺจ.}

\csromanheader{2}{2. ปจฺจยปจฺจนีย 1. วิภงฺควาร}
\csromanheader{2}{\textbf{นเหตุ}}
\proseitem{64}{จิตฺตํ ธมฺมํ ปจฺจยา โนจิตฺโต ธมฺโม อุปฺปชฺชติ นเหตุปจฺจยา.\csromanspacer{} อเหตุกํ จิตฺตํ ปจฺจยา สมฺปยุตฺตกา ขนฺธา จิตฺตสมุฏฺฐานญฺจ รูปํ.\csromanspacer{} อเหตุกปฏิสนฺธิกฺขเณ ฯเปฯ.\csromanspacer{} วิจิกิจฺฉาสหคตํ อุทฺธจฺจสหคตํ จิตฺตํ ปจฺจยา วิจิกิจฺฉาสหคโต อุทฺธจฺจสหคโต โมโห. (1)}

\proseitem{65}{โนจิตฺตํ ธมฺมํ ปจฺจยา โนจิตฺโต ธมฺโม อุปฺปชฺชติ นเหตุปจฺจยา.\csromanspacer{} อเหตุกํ โนจิตฺตํ เอกํ ขนฺธํ ปจฺจยา เทฺว ขนฺธา จิตฺตสมุฏฺฐานญฺจ รูปํ.\csromanspacer{} เทฺว ขนฺเธ ฯเปฯ.\csromanspacer{} อเหตุกปฏิสนฺธิกฺขเณ ฯเปฯ. (ยาว อสญฺญสตฺตา.) จกฺขายตนํ ปจฺจยา จกฺขุวิญฺญาณสหคตา ขนฺธา ฯเปฯ.\csromanspacer{} กายายตนํ ปจฺจยา กายวิญฺญาณสหคตา ขนฺธา.\csromanspacer{} วตฺถุํ ปจฺจยา อเหตุกา \csromanfolio{392}โนจิตฺตา ขนฺธา.\csromanspacer{} วิจิกิจฺฉาสหคเต อุทฺธจฺจสหคเต ขนฺเธ จ วตฺถุญฺจ ปจฺจยา วิจิกิจฺฉาสหคโต อุทฺธจฺจสหคโต โมโห. (1)}

\prose{โนจิตฺตํ ธมฺมํ ปจฺจยา จิตฺโต ธมฺโม อุปฺปชฺชติ นเหตุปจฺจยา.\csromanspacer{} อเหตุเก โนจิตฺเต ขนฺเธ ปจฺจยา จิตฺตํ.\csromanspacer{} วตฺถุํ ปจฺจยา จิตฺตํ.\csromanspacer{} อเหตุกปฏิสนฺธิกฺขเณ ฯเปฯ.\csromanspacer{} อเหตุกปฏิสนฺธิกฺขเณ วตฺถุํ ปจฺจยา จิตฺตํ.\csromanspacer{} จกฺขายตนํ ปจฺจยา จกฺขุวิญฺญาณํ ฯเปฯ.\csromanspacer{} กายายตนํ ปจฺจยา กายวิญฺญาณํ. (2)}

\prose{โนจิตฺตํ ธมฺมํ ปจฺจยา จิตฺโต จ โนจิตฺโต จ ธมฺมา อุปฺปชฺชนฺติ นเหตุปจฺจยา.\csromanspacer{} อเหตุกํ โนจิตฺตํ เอกํ ขนฺธํ ปจฺจยา เทฺว ขนฺธา จิตฺตญฺจ จิตฺตสมุฏฺฐานญฺจ รูปํ.\csromanspacer{} เทฺว ขนฺเธ ฯเปฯ.\csromanspacer{} วตฺถุํ ปจฺจยา จิตฺตญฺจ สมฺปยุตฺตกา จ ขนฺธา.\csromanspacer{} อเหตุกปฏิสนฺธิกฺขเณ ฯเปฯ.\csromanspacer{} อเหตุกปฏิสนฺธิกฺขเณ วตฺถุํ ปจฺจยา จิตฺตญฺจ สมฺปยุตฺตกา จ ขนฺธา.\csromanspacer{} จกฺขายตนํ ปจฺจยา จกฺขุวิญฺญาณํ สมฺปยุตฺตกา จ ขนฺธา ฯเปฯ.\csromanspacer{} กายายตนํ ปจฺจยา กายวิญฺญาณํ ฯเปฯ. (3)}

\prose{จิตฺตญฺจ โนจิตฺตญฺจ ธมฺมํ ปจฺจยา โนจิตฺโต ธมฺโม อุปฺปชฺชติ นเหตุปจฺจยา.\csromanspacer{} อเหตุกํ โนจิตฺตํ เอกํ ขนฺธญฺจ จิตฺตญฺจ ปจฺจยา เทฺว ขนฺธา จิตฺตสมุฏฺฐานญฺจ รูปํ.\csromanspacer{} เทฺว ขนฺเธ จ ฯเปฯ.\csromanspacer{} จิตฺตญฺจ วตฺถุญฺจ ปจฺจยา โนจิตฺตา ขนฺธา.\csromanspacer{} อเหตุกปฏิสนฺธิกฺขเณ ฯเปฯ.\csromanspacer{} อเหตุกปฏิสนฺธิกฺขเณ จิตฺตญฺจ วตฺถุญฺจ ปจฺจยา โนจิตฺตา ขนฺธา.\csromanspacer{} จกฺขายตนญฺจ จกฺขุวิญฺญาณญฺจ ปจฺจยา จกฺขุวิญฺญาณสหคตา ขนฺธา ฯเปฯ.\csromanspacer{} กายายตนญฺจ ฯเปฯ.\csromanspacer{} วิจิกิจฺฉาสหคเต อุทฺธจฺจสหคเต ขนฺเธ จ จิตฺตญฺจ ปจฺจยา วิจิกิจฺฉาสหคโต อุทฺธจฺจสหคโต โมโห. (1) (สํขิตฺตํ.)}

\csromanheader{2}{2. ปจฺจยปจฺจนีย 2. สงฺขฺยาวาร}
\csromanheader{2}{\textbf{สุทฺธ}}
\proseitem{66}{นเหตุยา ปญฺจ, น-อารมฺมเณ ตีณิ, น-อธิปติยา ปญฺจ, น-อนนฺตเร ตีณิ, นสมนนฺตเร ตีณิ, น-อญฺญมญฺเญ ตีณิ, น-อุปนิสฺสเย ตีณิ, นปุเรชาเต ปญฺจ, นปจฺฉาชาเต ปญฺจ, น-อาเสวเน ปญฺจ, นกมฺเม ตีณิ, นวิปาเก ปญฺจ, น-อาหาเร เอกํ, น-อินฺทฺริเย เอกํ, นฌาเน ปญฺจ, นมคฺเค ปญฺจ, นสมฺปยุตฺเต ตีณิ, นวิปฺปยุตฺเต ปญฺจ, โนนตฺถิยา ตีณิ, โนวิคเต ตีณิ.}

\csromanheader{2}{56. จิตฺตทุก \textbf{3. ปจฺจยานุโลมปจฺจนีย}}
\proseitem{67}{\csromanfolio{393}\textbf{เหตุ}ปจฺจยา น-อารมฺมเณ ตีณิ, น-อธิปติยา ปญฺจ. (สํขิตฺตํ.)}

\csromanheader{2}{4. ปจฺจยปจฺจนียานุโลม}
\proseitem{68}{นเหตุปจฺจยา อารมฺมเณ ปญฺจ, อนนฺตเร ปญฺจ, (สพฺพตฺถ ปญฺจ,) มคฺเค ตีณิ ฯเปฯ อวิคเต ปญฺจ.}

\vspace{\tipitakaparskip}
\csromancenter{(นิสฺสยวาโร ปจฺจยวารสทิโส.)}
\csromanheader{2}{56. จิตฺตทุก 5. สํสฏฺฐวาร}
\csromanheader{2}{1. ปจฺจยานุโลมาทิ}
\csromanheader{2}{\textbf{เหตุ}}
\proseitem{69}{จิตฺตํ ธมฺมํ สํสฏฺโฐ โนจิตฺโต ธมฺโม อุปฺปชฺชติ เหตุปจฺจยา.\csromanspacer{} จิตฺตํ สํสฏฺฐา สมฺปยุตฺตกา ขนฺธา.\csromanspacer{} ปฏิสนฺธิกฺขเณ ฯเปฯ. (1)}

\prose{โนจิตฺตํ ธมฺมํ สํสฏฺโฐ โนจิตฺโต ธมฺโม อุปฺปชฺชติ เหตุปจฺจยา.\csromanspacer{} โนจิตฺตํ เอกํ ขนฺธํ สํสฏฺฐา เทฺว ขนฺธา.\csromanspacer{} เทฺว ขนฺเธ สํสฏฺโฐ เอโก ขนฺโธ.\csromanspacer{} ปฏิสนฺธิกฺขเณ ฯเปฯ. (1)}

\prose{โนจิตฺตํ ธมฺมํ สํสฏฺโฐ จิตฺโต ธมฺโม อุปฺปชฺชติ เหตุปจฺจยา.\csromanspacer{} โนจิตฺเต ขนฺเธ สํสฏฺฐํ จิตฺตํ.\csromanspacer{} ปฏิสนฺธิเณ ฯเปฯ. (2)}

\prose{โนจิตฺตํ ธมฺมํ สํสฏฺโฐ จิตฺโต จ โนจิตฺโต จ ธมฺมา อุปฺปชฺชนฺติ เหตุปจฺจยา.\csromanspacer{} โนจิตฺตํ เอกํ ขนฺธํ สํสฏฺฐา เทฺว ขนฺธา จิตฺตญฺจ.\csromanspacer{} เทฺว ขนฺเธ ฯเปฯ.\csromanspacer{} ปฏิสนฺธิกฺขเณ ฯเปฯ. (3)}

\prose{จิตฺตญฺจ โนจิตฺตญฺจ ธมฺมํ สํสฏฺโฐ โนจิตฺโต ธมฺโม อุปฺปชฺชติ เหตุปจฺจยา.\csromanspacer{} โนจิตฺตํ เอกํ ขนฺธญฺจ จิตฺตญฺจ สํสฏฺฐา เทฺว ขนฺธา.\csromanspacer{} เทฺว ขนฺเธ จ ฯเปฯ.\csromanspacer{} ปฏิสนฺธิกฺขเณ ฯเปฯ. (สํขิตฺตํ.)}

\vspace{\tipitakaparskip}
\csromancenter{\textbf{เหตุ}ยา ปญฺจ, อารมฺมเณ ปญฺจ, อธิปติยา ปญฺจ, (สพฺพตฺถ ปญฺจ,) อวิคเต ปญฺจ. (สํขิตฺตํ.)}
\prose{\csromanfolio{394}นเหตุยา ปญฺจ, น-อธิปติยา ปญฺจ, นปุเรชาเต ปญฺจ, นปจฺฉาชาเต ปญฺจ, น-อาเสวเน ปญฺจ, นกมฺเม ตีณิ, นวิปาเก ปญฺจ, นฌาเน ปญฺจ, นมคฺเค ปญฺจ, นวิปฺปยุตฺเต ปญฺจ.}

\prose{(เอวํ อิตเร เทฺว คณนาปิ สมฺปยุตฺตวาโรปิ สพฺเพ กาตพฺพา.)}

\csromanheader{2}{56. จิตฺตทุก 7. ปญฺหาวาร}
\csromanheader{2}{1. ปจฺจยานุโลม 1. วิภงฺควาร}
\csromanheader{2}{\textbf{เหตุ}}
\proseitem{70}{โนจิตฺโต ธมฺโม โนจิตฺตสฺส ธมฺมสฺส เหตุปจฺจเยน ปจฺจโย.\csromanspacer{} โนจิตฺตา เหตู สมฺปยุตฺตกานํ ขนฺธานํ จิตฺตสมุฏฺฐานานญฺจ รูปานํ เหตุปจฺจเยน ปจฺจโย.\csromanspacer{} ปฏิสนฺธิกฺขเณ ฯเปฯ. (1)}

\prose{โนจิตฺโต ธมฺโม จิตฺตสฺส ธมฺมสฺส เหตุปจฺจเยน ปจฺจโย.\csromanspacer{} โนจิตฺตา เหตู จิตฺตสฺส เหตุปจฺจเยน ปจฺจโย.\csromanspacer{} ปฏิสนฺธิกฺขเณ ฯเปฯ. (2)}

\prose{โนจิตฺโต ธมฺโม จิตฺตสฺส จ โนจิตฺตสฺส จ ธมฺมสฺส เหตุปจฺจเยน ปจฺจโย.\csromanspacer{} โนจิตฺตา เหตู สมฺปยุตฺตกานํ ขนฺธานํ จิตฺตสฺส จ จิตฺตสมุฏฺฐานานญฺจ รูปานํ เหตุปจฺจเยน ปจฺจโย.\csromanspacer{} ปฏิสนฺธิกฺขเณ ฯเปฯ. (3)}

\csromanheader{2}{\textbf{อารมฺมณ}}
\proseitem{71}{จิตฺโต ธมฺโม จิตฺตสฺส ธมฺมสฺส อารมฺมณปจฺจเยน ปจฺจโย.\csromanspacer{} จิตฺตํ อารพฺภ จิตฺตํ อุปฺปชฺชติ. (มูลํ กาตพฺพํ.) จิตฺตํ อารพฺภ โนจิตฺตา ขนฺธา อุปฺปชฺชนฺติ. (มูลํ กาตพฺพํ.) จิตฺตํ อารพฺภ จิตฺตญฺจ สมฺปยุตฺตกา จ ขนฺธา อุปฺปชฺชนฺติ. (3)}

\proseitem{72}{โนจิตฺโต ธมฺโม โนจิตฺตสฺส ธมฺมสฺส อารมฺมณปจฺจเยน ปจฺจโย.\csromanspacer{} ทานํ, ทตฺวา, สีลํ ฯเปฯ อุโปสถกมฺมํ กตฺวา ตํ ปจฺจเวกฺขติ, อสฺสาเทติ อภินนฺทติ, ตํ อารพฺภ ราโค อุปฺปชฺชติ ฯเปฯ โทมนสฺสํ อุปฺปชฺชติ.\csromanspacer{} ปุพฺเพ สุจิณฺณานิ ฯเปฯ.\csromanspacer{} ฌานา วุฏฺฐหิตฺวา ฌานํ ฯเปฯ.\csromanspacer{} อริยา มคฺคา}

\vspace{\tipitakaparskip}
\csromancenter{จิตฺตทุก วุฏฺฐหิตฺวา มคฺคํ ปจฺจเวกฺขนฺติ, ผลํ ปจฺจเวกฺขนฺติ, นิพฺพานํ ปจฺจเวกฺขนฺติ.\csromanspacer{} นิพฺพานํ โคตฺรภุสฺส, โวทานสฺส, มคฺคสฺส, ผลสฺส อาวชฺชนาย อารมฺมณปจฺจเยน ปจฺจโย.\csromanspacer{} อริยา โนจิตฺเต ปหีเน กิเลเส ปจฺจเวกฺขนฺติ, วิกฺขมฺภิเต กิเลเส ปจฺจเวกฺขนฺติ.\csromanspacer{} ปุพฺเพ สมุทาจิณฺเณ กิเลเส ชานนฺติ.\csromanspacer{} จกฺขุํ ฯเปฯ วตฺถุํ โนจิตฺเต ขนฺเธ อนิจฺจโต ฯเปฯ โทมนสฺสํ อุปฺปชฺชติ.\csromanspacer{} ทิพฺเพน จกฺขุนา รูปํ ปสฺสติ.\csromanspacer{} ทิพฺพาย โสตธาตุยา สทฺทํ สุณาติ.\csromanspacer{} เจโตปริยญาเณน โนจิตฺตสมงฺคิสฺส จิตฺตํ ชานาติ.\csromanspacer{} อากาสานญฺจายตนํ ฯเปฯ.\csromanspacer{} อากิญฺจญฺญายตนํ ฯเปฯ.\csromanspacer{} รูปายตนํ จกฺขุวิญฺญาณสหคตานํ ขนฺธานํ ฯเปฯ.\csromanspacer{} โผฏฺฐพฺพายตนํ กายวิญฺญาณสหคตานํ ขนฺธานํ ฯเปฯ.\csromanspacer{} โนจิตฺตา ขนฺธา อิทฺธิวิธญาณสฺส, เจโตปริยญาณสฺส, ปุพฺเพนิวาสานุสฺสติญาณสฺส, ยถากมฺมูปคญาณสฺส, อนาคตํสญาณสฺส, อาวชฺชนาย อารมฺมณปจฺจเยน ปจฺจโย. (1)}
\prose{\csromanfolio{395}โนจิตฺโต ธมฺโม จิตฺตสฺส ธมฺมสฺส อารมฺมณปจฺจเยน ปจฺจโย.\csromanspacer{} ทานํ ทตฺวา ฯเปฯ. (ปฐมคมนสทิสํ.\csromanspacer{} นินฺนานากรณํ.\csromanspacer{} อิมํ นานํ) รูปายตนํ จกฺขุวิญฺญาณสฺส ฯเปฯ.\csromanspacer{} โผฏฺฐพฺพายตนํ กายวิญฺญาณสฺส ฯเปฯ.\csromanspacer{} โนจิตฺตา ขนฺธา อิทฺธิวิธญาณสฺส ฯเปฯ อาวชฺชนาย อารมฺมณปจฺจเยน ปจฺจโย. (2)}

\prose{โนจิตฺโต ธมฺโม จิตฺตสฺส จ โนจิตฺตสฺส จ ธมฺมสฺส อารมฺมณปจฺจเยน ปจฺจโย.\csromanspacer{} ทานํ ทตฺวา ฯเปฯ. (ปฐมคมนสทิสํ.\csromanspacer{} นินฺนานากรณํ.\csromanspacer{} อิมํ นานํ.) รูปายตนํ จกฺขุวิญฺญาณสฺส สมฺปยุตฺตกานญฺจ ขนฺธานํ ฯเปฯ.\csromanspacer{} โผฏฺฐพฺพายตนํ กายวิญฺญาณสฺส สมฺปยุตฺตกานญฺจ ขนฺธานํ ฯเปฯ.\csromanspacer{} โนจิตฺตา ขนฺธา อิทฺธิวิธญาณสฺส ฯเปฯ อาวชฺชนาย อารมฺมณปจฺจเยน ปจฺจโย. (3)}

\prose{จิตฺโต จ โนจิตฺโต จ ธมฺมา จิตฺตสฺส ธมฺมสฺส อารมฺมณปจฺจเยน ปจฺจโย.\csromanspacer{} จิตฺตญฺจ สมฺปยุตฺตเก จ ขนฺเธ อารพฺภ จิตฺตํ อุปฺปชฺชติ.\csromanspacer{} ตีณิ.}

\csromanheader{2}{\textbf{อธิปติ}}
\proseitem{73}{จิตฺโต ธมฺโม จิตฺตสฺส ธมฺมสฺส อธิปติปจฺจเยน ปจฺจโย.\csromanspacer{} \textbf{อารมฺมณาธิปติ}\csromandash{}จิตฺตํ ครุํ กตฺวา จิตฺตํ อุปฺปชฺชติ. (1)}

\prose{จิตฺโต ธมฺโม โนจิตฺตสฺส ธมฺมสฺส อธิปติปจฺจเยน ปจฺจโย.\csromanspacer{} \textbf{อารมฺมณาธิปติ}, สหชาตาธิปติ.\csromanspacer{}\csromanspacer{}\textbf{อารมฺมณาธิปติ}\csromandash{}จิตฺตํ ครุํ \csromanfolio{396}กตฺวา โนจิตฺตา ขนฺธา อุปฺปชฺชนฺติ.\csromanspacer{} \textbf{สหชาตาธิปติ}\csromandash{}จิตฺตาธิปติ สมฺปยุตฺตกานํ ขนฺธานํ จิตฺตสมุฏฺฐานานญฺจ รูปานํ อธิปติปจฺจเยน ปจฺจโย. (2)}

\prose{จิตฺโต ธมฺโม จิตฺตสฺส จ โนจิตฺตสฺส จ ธมฺมสฺส อธิปติปจฺจเยน ปจฺจโย.\csromanspacer{} \textbf{อารมฺมณาธิปติ}\csromandash{}จิตฺตํ ครุํ กตฺวา จิตฺตญฺจ สมฺปยุตฺตกา จ ขนฺธา อุปฺปชฺชนฺติ. (3)}

\proseitem{74}{โนจิตฺโต ธมฺโม โนจิตฺตสฺส ธมฺมสฺส อธิปติปจฺจเยน ปจฺจโย.\csromanspacer{} \textbf{อารมฺมณาธิปติ}, สหชาตาธิปติ.\csromanspacer{}\csromanspacer{}\textbf{อารมฺมณาธิปติ}\csromandash{}ทานํ ฯเปฯ สีลํ ฯเปฯ อุโปสถกมฺมํ กตฺวา ตํ ครุํ กตฺวา ปจฺจเวกฺขติ, อสฺสาเทติ อภินนฺทติ, ตํ ครุํ กตฺวา ราโค อุปฺปชฺชติ.\csromanspacer{} ทิฏฺฐิ อุปฺปชฺชติ.\csromanspacer{} ปุพฺเพ ฯเปฯ.\csromanspacer{} ฌานา ฯเปฯ.\csromanspacer{} อริยา มคฺคา วุฏฺฐหิตฺวา มคฺคํ ครุํ กตฺวา ปจฺจเวกฺขนฺติ, ผลํ ครุํ กตฺวา ปจฺจเวกฺขนฺติ, นิพฺพานํ ครุํ กตฺวา ปจฺจเวกฺขนฺติ.\csromanspacer{} นิพฺพานํ โคตฺรภุสฺส, โวทานสฺส, มคฺคสฺส, ผลสฺส อธิปติปจฺจเยน ปจฺจโย.\csromanspacer{} จกฺขุํ ฯเปฯ วตฺถุํ โนจิตฺเต ขนฺเธ ครุํ กตฺวา อสฺสาเทติ อภินนฺทติ, ตํ ครุํ กตฺวา ราโค อุปฺปชฺชติ.\csromanspacer{} ทิฏฺฐิ อุปฺปชฺชติ.\csromanspacer{} \textbf{สหชาตาธิปติ}\csromandash{}โนจิตฺตาธิปติ สมฺปยุตฺตกานํ ขนฺธานํ จิตฺตสมุฏฺฐานานญฺจ รูปานํ อธิปติปจฺจเยน ปจฺจโย. (1)}

\prose{โนจิตฺโต ธมฺโม จิตฺตสฺส ธมฺมสฺส อธิปติปจฺจเยน ปจฺจโย. (เทฺวปิ คมนา ปฐมคมนสทิสํ.\csromanspacer{} นินฺนานากรณํ.\csromanspacer{} \textbf{อารมฺมณาธิปติ} สหชาตาธิปติ กาตพฺพา.) (2)}

\prose{จิตฺโต จ โนจิตฺโต จ ธมฺมา จิตฺตสฺส ธมฺมสฺส อธิปติปจฺจเยน ปจฺจโย.\csromanspacer{} \textbf{อารมฺมณาธิปติ}. (ตีณิปิ ครุการมฺมณา กาตพฺพา.\csromanspacer{} \textbf{อารมฺมณาธิปติ}เยว.)}

\csromanheader{2}{\textbf{อนนฺตร}}
\proseitem{75}{จิตฺโต ธมฺโม จิตฺตสฺส ธมฺมสฺส อนนฺตรปจฺจเยน ปจฺจโย.\csromanspacer{} ปุริมํ ปุริมํ จิตฺตํ ปจฺฉิมสฺส ปจฺฉิมสฺส จิตฺตสฺส อนนฺตรปจฺจเยน ปจฺจโย. (1)}

\prose{จิตฺโต ธมฺโม โนจิตฺตสฺส ธมฺมสฺส อนนฺตรปจฺจเยน ปจฺจโย.\csromanspacer{} ปุริมํ ปุริมํ จิตฺตํ ปจฺฉิมานํ ปจฺฉิมานํ โนจิตฺตานํ ขนฺธานํ อนนฺตรปจฺจเยน ปจฺจโย.\csromanspacer{} จิตฺตํ วุฏฺฐานสฺส อนนฺตรปจฺจเยน ปจฺจโย. (2)}

\csromanheader{2}{56. จิตฺตทุก}
\prose{\csromanfolio{397}จิตฺโต ธมฺโม จิตฺตสฺส จ โนจิตฺตสฺส จ ธมฺมสฺส อนนฺตรปจฺจเยน ปจฺจโย.\csromanspacer{} ปุริมํ ปุริมํ จิตฺตํ ปจฺฉิมสฺส ปจฺฉิมสฺส จิตฺตสฺส สมฺปยุตฺตกานญฺจ ขนฺธานํ อนนฺตรปจฺจเยน ปจฺจโย. (3)}

\proseitem{76}{โนจิตฺโต ธมฺโม โนจิตฺตสฺส ธมฺมสฺส อนนฺตรปจฺจเยน ปจฺจโย.\csromanspacer{} ปุริมา ปุริมา โนจิตฺตา ขนฺธา ฯเปฯ ผลสมาปตฺติยา อนนฺตรปจฺจเยน ปจฺจโย. (อิเม เทฺว ปูเรตุกาเมน กาตพฺพา ปุริมคมนสทิสํ.)}

\prose{จิตฺโต จ โนจิตฺโต จ ธมฺมา จิตฺตสฺส ธมฺมสฺส อนนฺตรปจฺจเยน ปจฺจโย.\csromanspacer{} ปุริมํ ปุริมํ จิตฺตญฺจ สมฺปยุตฺตกา จ ขนฺธา ปจฺฉิมสฺส ปจฺฉิมสฺส จิตฺตสฺส อนนฺตรปจฺจเยน ปจฺจโย. (มูลํ ปุจฺฉิตพฺพํ.) ปุริมํ ปุริมํ จิตฺตญฺจ สมฺปยุตฺตกา จ ขนฺธา ปจฺฉิมานํ ปจฺฉิมานํ โนจิตฺตานํ ขนฺธานํ อนนฺตรปจฺจเยน ปจฺจโย.\csromanspacer{} จิตฺตญฺจ สมฺปยุตฺตกา จ ขนฺธา วุฏฺฐานสฺส อนนฺตรปจฺจเยน ปจฺจโย. (มูลํ ปุจฺฉิตพฺพํ.) ปุริมํ ปุริมํ จิตฺตญฺจ สมฺปยุตฺตกา จ ขนฺธา ปจฺฉิมสฺส ปจฺฉิมสฺส จิตฺตสฺส สมฺปยุตฺตกานญฺจ ขนฺธานํ อนนฺตรปจฺจเยน ปจฺจโย. (3)}

\prose{สมนนฺตรปจฺจเยน ปจฺจโย.\csromanspacer{} สหชาตปจฺจเยน ปจฺจโย.\csromanspacer{} อญฺญมญฺญปจฺจเยน ปจฺจโย. (ปฏิจฺจวารสทิสา.) นิสฺสยปจฺจเยน ปจฺจโย.\csromanspacer{} ปญฺจ. (ปจฺจยวารสทิสา.)}

\csromanheader{2}{\textbf{อุปนิสฺสย}}
\proseitem{77}{จิตฺโต ธมฺโม จิตฺตสฺส ธมฺมสฺส อุปนิสฺสยปจฺจเยน ปจฺจโย.\csromanspacer{} อารมฺมณูปนิสฺสโย, อนนฺตรูปนิสฺสโย, ปกตูปนิสฺสโย ฯเปฯ.\csromanspacer{} ปกตูปนิสฺสโย\csromandash{}จิตฺตํ จิตฺตสฺส อุปนิสฺสยปจฺจเยน ปจฺจโย.\csromanspacer{} ตีณิ.}

\prose{โนจิตฺโต ธมฺโม โนจิตฺตสฺส ธมฺมสฺส อุปนิสฺสยปจฺจเยน ปจฺจโย.\csromanspacer{} อารมฺมณูปนิสฺสโย, อนนฺตรูปนิสฺสโย, ปกตูปนิสฺสโย ฯเปฯ.\csromanspacer{} ปกตูปนิสฺสโย\csromandash{}สทฺธํ อุปนิสฺสาย ทานํ เทติ ฯเปฯ.\csromanspacer{} มานํ ชปฺเปติ.\csromanspacer{} ทิฏฺฐิํ คณฺหาติ.\csromanspacer{} สีลํ ฯเปฯ.\csromanspacer{} เสนาสนํ อุปนิสฺสาย ทานํ เทติ ฯเปฯ สํฆํ ภินฺทติ.\csromanspacer{} สทฺธา ฯเปฯ.\csromanspacer{} เสนาสนํ สทฺธาย ฯเปฯ มคฺคสฺส, ผลสมาปตฺติยา อุปนิสฺสยปจฺจเยน ปจฺจโย. (1)}

\prose{โนจิตฺโต ธมฺโม จิตฺตสฺส ธมฺมสฺส อุปนิสฺสยปจฺจเยน ปจฺจโย. (อิเม เทฺวปิ ปุเรตุกาเมน สพฺพตฺถ กาตพฺพา.\csromanspacer{} ปฐมคมนสทิสํ.\csromanspacer{} นินฺนานากรณํ.)}

\prose{\csromanfolio{398}จิตฺโต จ โนจิตฺโต จ ธมฺมา จิตฺตสฺส ธมฺมสฺส อุปนิสฺสยปจฺจเยน ปจฺจโย.\csromanspacer{} อารมฺมณูปนิสฺสโย, อนนฺตรูปนิสฺสโย, ปกตูปนิสฺสโย ฯเปฯ.\csromanspacer{} ปกตูปนิสฺสโย\csromandash{}จิตฺตญฺจ สมฺปยุตฺตกา จ ขนฺธา จิตฺตสฺส อุปนิสฺสยปจฺจเยน ปจฺจโย.\csromanspacer{} ตีณิ.}

\csromanheader{2}{\textbf{ปุเรชาต}}
\proseitem{78}{โนจิตฺโต ธมฺโม โนจิตฺตสฺส ธมฺมสฺส ปุเรชาตปจฺจเยน ปจฺจโย.\csromanspacer{} \textbf{อารมฺมณปุเรชาตํ}, วตฺถุปุเรชาตํ.\csromanspacer{}\csromanspacer{}\textbf{อารมฺมณปุเรชาตํ}\csromandash{}จกฺขุํ ฯเปฯ วตฺถุํ อนิจฺจโต ฯเปฯ โทมนสฺสํ อุปฺปชฺชติ.\csromanspacer{} ทิพฺเพน จกฺขุนา รูปํ ปสฺสติ.\csromanspacer{} ทิพฺพาย โสตธาตุยา สทฺทํ สุณาติ.\csromanspacer{} รูปายตนํ จกฺขุวิญฺญาณสหคตานํ ขนฺธานํ ฯเปฯ.\csromanspacer{} โผฏฺฐพฺพายตนํ กายวิญฺญาณสหคตานํ ขนฺธานํ ปุเรชาตปจฺจเยน ปจฺจโย.\csromanspacer{} \textbf{วตฺถุปุเรชาตํ}\csromandash{}จกฺขายตนํ จกฺขุวิญฺญาณสหคตานํ ขนฺธานํ ฯเปฯ กายายตนํ ฯเปฯ.\csromanspacer{} วตฺถุ โนจิตฺตานํ ขนฺธานํ ปุเรชาตปจฺจเยน ปจฺจโย. (1)}

\prose{โนจิตฺโต ธมฺโม จิตฺตสฺส ธมฺมสฺส ปุเรชาตปจฺจเยน ปจฺจโย.\csromanspacer{} \textbf{อารมฺมณปุเรชาตํ}, วตฺถุปุเรชาตํ.\csromanspacer{}\csromanspacer{}\textbf{อารมฺมณปุเรชาตํ}\csromandash{}จกฺขุํ ฯเปฯ วตฺถุํ อนิจฺจโต ฯเปฯ โทมนสฺสํ อุปฺปชฺชติ.\csromanspacer{} ทิพฺเพน จกฺขุนา รูปํ ปสฺสติ.\csromanspacer{} ทิพฺพาย โสตธาตุยา สทฺทํ สุณาติ.\csromanspacer{} รูปายตนํ จกฺขุวิญฺญาณสฺส ฯเปฯ.\csromanspacer{} โผฏฺฐพฺพายตนํ กายวิญฺญาณสฺส ฯเปฯ.\csromanspacer{} \textbf{วตฺถุปุเรชาตํ}\csromandash{}จกฺขายตนํ จกฺขุวิญฺญาณสฺส ฯเปฯ กายายตนํ กายวิญฺญาณสฺส ฯเปฯ.\csromanspacer{} วตฺถุ จิตฺตสฺส ปุเรชาตปจฺจเยน ปจฺจโย. (2)}

\prose{โนจิตฺโต ธมฺโม จิตฺตสฺส จ โนจิตฺตสฺส จ ธมฺมสฺส ปุเรชาตปจฺจเยน ปจฺจโย.\csromanspacer{} \textbf{อารมฺมณปุเรชาตํ}, วตฺถุปุเรชาตํ.\csromanspacer{}\csromanspacer{}\textbf{อารมฺมณปุเรชาตํ}\csromandash{} จกฺขุํ ฯเปฯ วตฺถุํ อนิจฺจโต ฯเปฯ โทมนสฺสํ อุปฺปชฺชติ.\csromanspacer{} ทิพฺเพน จกฺขุนา รูปํ ปสฺสติ.\csromanspacer{} ทิพฺพาย โสตธาตุยา สทฺทํ สุณาติ.\csromanspacer{} รูปายตนํ จกฺขุวิญฺญาณสฺส จ สมฺปยุตฺตกานญฺจ ขนฺธานํ ฯเปฯ.\csromanspacer{} โผฏฺฐพฺพายตนํ ฯเปฯ.\csromanspacer{} \textbf{วตฺถุปุเรชาตํ}\csromandash{}จกฺขายตนํ จกฺขุวิญฺญาณสฺส สมฺปยุตฺตกานญฺจ ขนฺธานํ ปุเรชาตปจฺจเยน ปจฺจโย ฯเปฯ กายายตนํ กายวิญฺญาณสฺส สมฺปยุตฺตกานญฺจ ขนฺธานํ ฯเปฯ.\csromanspacer{} วตฺถุ จิตฺตสฺส จ สมฺปยุตฺตกานญฺจ ขนฺธานํ ปุเรชาตปจฺจเยน ปจฺจโย. (3)}

\csromanheader{2}{56. จิตฺตทุก \textbf{ปจฺฉาชาตาเสวน}}
\proseitem{79}{\csromanfolio{399}จิตฺโต ธมฺโม โนจิตฺตสฺส ธมฺมสฺส ปจฺฉาชาตปจฺจเยน ปจฺจโย. (สํขิตฺตํ.) (1)}

\prose{โนจิตฺโต ธมฺโม โนจิตฺตสฺส ธมฺมสฺส ปจฺฉาชาตปจฺจเยน ปจฺจโย. (สํขิตฺตํ.) (1)}

\prose{จิตฺโต จ โนจิตฺโต จ ธมฺมา โนจิตฺตสฺส ธมฺมสฺส ปจฺฉาชาตปจฺจเยน ปจฺจโย. (สํขิตฺตํ.) (1)}

\prose{จิตฺโต ธมฺโม จิตฺตสฺส ธมฺมสฺส อาเสวนปจฺจเยน ปจฺจโย.\csromanspacer{} นว.}

\csromanheader{2}{\textbf{กมฺม}}
\proseitem{80}{โนจิตฺโต ธมฺโม โนจิตฺตสฺส ธมฺมสฺส กมฺมปจฺจเยน ปจฺจโย.\csromanspacer{} \textbf{สหชาตา}, นานากฺขณิกา.\csromanspacer{}\csromanspacer{}\textbf{สหชาตา} โนจิตฺตา เจตนา สมฺปยุตฺตกานํ ขนฺธานํ จิตฺตสมุฏฺฐานานญฺจ รูปานํ กมฺมปจฺจเยน ปจฺจโย.\csromanspacer{} ปฏิสนฺธิกฺขเณ ฯเปฯ.\csromanspacer{} \textbf{นานากฺขณิกา} โนจิตฺตา เจตนา วิปากานํ ขนฺธานํ กฏตฺตา จ รูปานํ กมฺมปจฺจเยน ปจฺจโย. (1)}

\prose{โนจิตฺโต ธมฺโม จิตฺตสฺส ธมฺมสฺส กมฺมปจฺจเยน ปจฺจโย.\csromanspacer{} \textbf{สหชาตา}, นานากฺขณิกา.\csromanspacer{}\csromanspacer{}\textbf{สหชาตา} โนจิตฺตา เจตนา จิตฺตสฺส กมฺมปจฺจเยน ปจฺจโย.\csromanspacer{} ปฏิสนฺธิกฺขเณ ฯเปฯ.\csromanspacer{} \textbf{นานากฺขณิกา} โนจิตฺตา เจตนา วิปากสฺส จิตฺตสฺส กมฺมปจฺจเยน ปจฺจโย. (2)}

\prose{โนจิตฺโต ธมฺโม จิตฺตสฺส จ โนจิตฺตสฺส จ ธมฺมสฺส กมฺมปจฺจเยน ปจฺจโย.\csromanspacer{} \textbf{สหชาตา}, นานากฺขณิกา.\csromanspacer{}\csromanspacer{}\textbf{สหชาตา} โนจิตฺตา เจตนา สมฺปยุตฺตกานํ ขนฺธานํ จิตฺตสฺส จ จิตฺตสมุฏฺฐานานญฺจ รูปานํ กมฺมปจฺจเยน ปจฺจโย.\csromanspacer{} ปฏิสนฺธิกฺขเณ ฯเปฯ.\csromanspacer{} \textbf{นานากฺขณิกา} โนจิตฺตา เจตนา วิปากานํ ขนฺธานํ จิตฺตสฺส จ กฏตฺตา จ รูปานํ กมฺมปจฺจเยน ปจฺจโย. (3)}

\csromanheader{2}{\textbf{วิปากาทิ}}
\proseitem{81}{จิตฺโต ธมฺโม โนจิตฺตสฺส ธมฺมสฺส วิปากปจฺจเยน ปจฺจโย.\csromanspacer{} ปญฺจ.\csromanspacer{} อาหารปจฺจเยน ปจฺจโย.\csromanspacer{} ปญฺจ.\csromanspacer{} อินฺทฺริยปจฺจเยน ปจฺจโย.\csromanspacer{} ปญฺจ.\csromanspacer{} โนจิตฺโต ธมฺโม โนจิตฺตสฺส ธมฺมสฺส ฌานปจฺจเยน ปจฺจโย.\csromanspacer{} ตีณิ.\csromanspacer{} มคฺคปจฺจเยน ปจฺจโย.\csromanspacer{} ตีณิ.\csromanspacer{} สมฺปยุตฺตปจฺจเยน ปจฺจโย.\csromanspacer{} ปญฺจ.}

\csromanheader{2}{\textbf{วิปฺปยุตฺต}}
\proseitem{82}{\csromanfolio{400}จิตฺโต ธมฺโม โนจิตฺตสฺส ธมฺมสฺส วิปฺปยุตฺตปจฺจเยน ปจฺจโย.\csromanspacer{} \textbf{สหชาตํ}, ปจฺฉาชาตํ. (สํขิตฺตํ.) (1)}

\prose{โนจิตฺโต ธมฺโม โนจิตฺตสฺส ธมฺมสฺส วิปฺปยุตฺตปจฺจเยน ปจฺจโย.\csromanspacer{} \textbf{สหชาตํ}, ปุเรชาตํ, ปจฺฉาชาตํ.\csromanspacer{}\csromanspacer{}\textbf{สหชาตา} โนจิตฺตา ขนฺธา จิตฺตสมุฏฺฐานานํ รูปานํ วิปฺปยุตฺตปจฺจเยน ปจฺจโย.\csromanspacer{} ปฏิสนฺธิกฺขเณ โนจิตฺตา ขนฺธา กฏตฺตารูปานํ ฯเปฯ.\csromanspacer{} โนจิตฺตา ขนฺธา วตฺถุสฺส วิปฺปยุตฺตปจฺจเยน ปจฺจโย.\csromanspacer{} วตฺถุ ขนฺธานํ วิปฺปยุตฺตปจฺจเยน ปจฺจโย.\csromanspacer{} \textbf{ปุเรชาตํ} จกฺขายตนํ จกฺขุวิญฺญาณสหคตานํ ขนฺธานํ ฯเปฯ กายายตนํ กายวิญฺญาณสหคตานํ ขนฺธานํ ฯเปฯ.\csromanspacer{} วตฺถุ โนจิตฺตานํ ขนฺธานํ วิปฺปยุตฺตปจฺจเยน ปจฺจโย.\csromanspacer{} \textbf{ปจฺฉาชาตา} โนจิตฺตา ขนฺธา ปุเรชาตสฺส อิมสฺส กายสฺส วิปฺปยุตฺตปจฺจเยน ปจฺจโย. (1)}

\prose{โนจิตฺโต ธมฺโม จิตฺตสฺส ธมฺมสฺส วิปฺปยุตฺตปจฺจเยน ปจฺจโย.\csromanspacer{} \textbf{สหชาตํ}, ปุเรชาตํ.\csromanspacer{}\csromanspacer{}\textbf{สหชาตํ} ปฏิสนฺธิกฺขเณ วตฺถุ จิตฺตสฺส วิปฺปยุตฺตปจฺจเยน ปจฺจโย.\csromanspacer{} \textbf{ปุเรชาตํ} จกฺขายตนํ จกฺขุวิญฺญาณสฺส ฯเปฯ กายายตนํ กายวิญฺญาณสฺส ฯเปฯ.\csromanspacer{} วตฺถุ จิตฺตสฺส วิปฺปยุตฺตปจฺจเยน ปจฺจโย. (2)}

\prose{โนจิตฺโต ธมฺโม จิตฺตสฺส จ โนจิตฺตสฺส จ ธมฺมสฺส วิปฺปยุตฺตปจฺจเยน ปจฺจโย.\csromanspacer{} \textbf{สหชาตํ}, ปุเรชาตํ.\csromanspacer{}\csromanspacer{}\textbf{สหชาตํ} ปฏิสนฺธิกฺขเณ วตฺถุ จิตฺตสฺส สมฺปยุตฺตกานญฺจ ขนฺธานํ วิปฺปยุตฺตปจฺจเยน ปจฺจโย.\csromanspacer{} \textbf{ปุเรชาตํ} จกฺขายตนํ จกฺขุวิญฺญาณสฺส สมฺปยุตฺตกานญฺจ ขนฺธานํ วิปฺปยุตฺตปจฺจเยน ปจฺจโย ฯเปฯ กายายตนํ กายวิญฺญาณสฺส สมฺปยุตฺตกานญฺจ ขนฺธานํ ฯเปฯ.\csromanspacer{} วตฺถุ จิตฺตสฺส จ สมฺปยุตฺตกานญฺจ ขนฺธานํ วิปฺปยุตฺตปจฺจเยน ปจฺจโย. (3)}

\prose{จิตฺโต จ โนจิตฺโต จ ธมฺมา โนจิตฺตสฺส ธมฺมสฺส วิปฺปยุตฺตปจฺจเยน ปจฺจโย.\csromanspacer{} \textbf{สหชาตํ}, ปจฺฉาชาตํ. (สํขิตฺตํ.) (1)}

\csromanheader{2}{\textbf{อตฺถิ}}
\proseitem{83}{จิตฺโต ธมฺโม โนจิตฺตสฺส ธมฺมสฺส อตฺถิปจฺจเยน ปจฺจโย.\csromanspacer{} \textbf{สหชาตํ}, ปจฺฉาชาตํ. (สํขิตฺตํ.) (1)}

\prose{โนจิตฺโต ธมฺโม โนจิตฺตสฺส ธมฺมสฺส อตฺถิปจฺจเยน ปจฺจโย.\csromanspacer{} \textbf{สหชาตํ}, ปุเรชาตํ, ปจฺฉาชาตํ, อาหารํ, อินฺทฺริยํ. (สํขิตฺตํ.) (1)}

\csromanheader{2}{56. จิตฺตทุก}
\prose{\csromanfolio{401}โนจิตฺโต ธมฺโม จิตฺตสฺส ธมฺมสฺส อตฺถิปจฺจเยน ปจฺจโย.\csromanspacer{} \textbf{สหชาตํ}, ปุเรชาตํ.\csromanspacer{}\csromanspacer{}\textbf{สหชาตา} โนจิตฺตา ขนฺธา จิตฺตสฺส อตฺถิปจฺจเยน ปจฺจโย.\csromanspacer{} ปฏิสนฺธิกฺขเณ ฯเปฯ.\csromanspacer{} ปฏิสนฺธิกฺขเณ วตฺถุ จิตฺตสฺส อตฺถิปจฺจเยน ปจฺจโย.\csromanspacer{} \textbf{ปุเรชาตํ} จกฺขุํ ฯเปฯ วตฺถุํ อนิจฺจโต ฯเปฯ. (ปุเรชาตสทิสํ, สํขิตฺตํ.) (2)}

\prose{โนจิตฺโต ธมฺโม จิตฺตสฺส จ โนจิตฺตสฺส จ ธมฺมสฺส อตฺถิปจฺจเยน ปจฺจโย.\csromanspacer{} \textbf{สหชาตํ}, ปุเรชาตํ.\csromanspacer{}\csromanspacer{}\textbf{สหชาโต} โนจิตฺโต เอโก ขนฺโธ ทฺวินฺนํ ขนฺธานํ จิตฺตสฺส จ จิตฺตสมุฏฺฐานานญฺจ รูปานํ อตฺถิปจฺจเยน ปจฺจโย.\csromanspacer{} ปฏิสนฺธิกฺขเณ โนจิตฺโต เอโก ขนฺโธ ฯเปฯ.\csromanspacer{} ปฏิสนฺธิกฺขเณ วตฺถุ จิตฺตสฺส สมฺปยุตฺตกานญฺจ ขนฺธานํ อตฺถิปจฺจเยน ปจฺจโย.\csromanspacer{} \textbf{ปุเรชาตํ} จกฺขุํ ฯเปฯ วตฺถุํ ฯเปฯ. (ปุเรชาตสทิสํ, สํขิตฺตํ.) (3)}

\proseitem{84}{จิตฺโต จ โนจิตฺโต จ ธมฺมา โนจิตฺตสฺส ธมฺมสฺส อตฺถิปจฺจเยน ปจฺจโย.\csromanspacer{} \textbf{สหชาตํ}, ปุเรชาตํ, ปจฺฉาชาตํ, อาหารํ, อินฺทฺริยํ.\csromanspacer{}\csromanspacer{}\textbf{สหชาโต} โนจิตฺโต เอโก ขนฺโธ จ จิตฺตญฺจ ทฺวินฺนํ ขนฺธานํ จิตฺตสมุฏฺฐานานญฺจ รูปานํ อตฺถิปจฺจเยน ปจฺจโย.\csromanspacer{} เทฺว ขนฺธา จ ฯเปฯ.\csromanspacer{} \textbf{สหชาตํ} จิตฺตญฺจ วตฺถุญฺจ โนจิตฺตานํ ขนฺธานํ อตฺถิปจฺจเยน ปจฺจโย. (ปฏิสนฺธิกฺขเณปิ เทฺว.) \textbf{สหชาตํ} จิตฺตญฺจ สมฺปยุตฺตกา จ ขนฺธา จิตฺตสมุฏฺฐานานํ รูปานํ อตฺถิปจฺจเยน ปจฺจโย.\csromanspacer{} \textbf{สหชาตํ} จิตฺตญฺจ มหาภูตา จ จิตฺตสมุฏฺฐานานํ รูปานํ อตฺถิปจฺจเยน ปจฺจโย. (ปฏิสนฺธิกฺขเณปิ เทฺว.) \textbf{ปจฺฉาชาตํ} จิตฺตญฺจ สมฺปยุตฺตกา จ ขนฺธา ปุเรชาตสฺส อิมสฺส กายสฺส อตฺถิปจฺจเยน ปจฺจโย.\csromanspacer{} \textbf{ปจฺฉาชาตํ} จิตฺตญฺจ สมฺปยุตฺตกา จ ขนฺธา กพฬีกาโร อาหาโร จ อิมสฺส กายสฺส อตฺถิปจฺจเยน ปจฺจโย.\csromanspacer{} \textbf{ปจฺฉาชาตํ} จิตฺตญฺจ สมฺปยุตฺตกา จ ขนฺธา รูปชีวิตินฺทฺริยญฺจ กฏตฺตารูปานํ อตฺถิปจฺจเยน ปจฺจโย. (1)}

\csromanheader{2}{1. ปจฺจยานุโลม 2. สงฺขฺยาวาร}
\csromanheader{2}{\textbf{สุทฺธ}}
\proseitem{85}{เหตุยา ตีณิ, อารมฺมเณ นว, อธิปติยา นว, อนนฺตเร นว, สมนนฺตเร นว, สหชาเต ปญฺจ, อญฺญมญฺเญ ปญฺจ, นิสฺสเย ปญฺจ, อุปนิสฺสเย นว, ปุเรชาเต ตีณิ, ปจฺฉาชาเต ตีณิ, อาเสวเน \csromanfolio{402}นว, กมฺเม ตีณิ, วิปาเก ปญฺจ, อาหาเร ปญฺจ, อินฺทฺริเย ปญฺจ, ฌาเน ตีณิ, มคฺเค ตีณิ, สมฺปยุตฺเต ปญฺจ, วิปฺปยุตฺเต ปญฺจ, อตฺถิยา ปญฺจ, นตฺถิยา นว, วิคเต นว, อวิคเต ปญฺจ.}

\csromanheader{2}{2. ปจฺจนียุทฺธาร}
\proseitem{86}{จิตฺโต ธมฺโม จิตฺตสฺส ธมฺมสฺส อารมฺมณปจฺจเยน ปจฺจโย.\csromanspacer{} อุปนิสฺสยปจฺจเยน ปจฺจโย. (1)}

\prose{จิตฺโต ธมฺโม โนจิตฺตสฺส ธมฺมสฺส อารมฺมณปจฺจเยน ปจฺจโย.\csromanspacer{} สหชาตปจฺจเยน ปจฺจโย.\csromanspacer{} อุปนิสฺสยปจฺจเยน ปจฺจโย.\csromanspacer{} ปจฺฉาชาตปจฺจเยน ปจฺจโย. (2)}

\prose{จิตฺโต ธมฺโม จิตฺตสฺส จ โนจิตฺตสฺส จ ธมฺมสฺส อารมฺมณปจฺจเยน ปจฺจโย.\csromanspacer{} อุปนิสฺสยปจฺจเยน ปจฺจโย. (3)}

\proseitem{87}{โนจิตฺโต ธมฺโม โนจิตฺตสฺส ธมฺมสฺส อารมฺมณปจฺจเยน ปจฺจโย.\csromanspacer{} สหชาตปจฺจเยน ปจฺจโย.\csromanspacer{} อุปนิสฺสยปจฺจเยน ปจฺจโย.\csromanspacer{} ปุเรชาตปจฺจเยน ปจฺจโย.\csromanspacer{} ปจฺฉาชาตปจฺจเยน ปจฺจโย.\csromanspacer{} กมฺมปจฺจเยน ปจฺจโย.\csromanspacer{} อาหารปจฺจเยน ปจฺจโย.\csromanspacer{} อินฺทฺริยปจฺจเยน ปจฺจโย. (1)}

\prose{โนจิตฺโต ธมฺโม จิตฺตสฺส ธมฺมสฺส อารมฺมณปจฺจเยน ปจฺจโย.\csromanspacer{} สหชาตปจฺจเยน ปจฺจโย.\csromanspacer{} อุปนิสฺสยปจฺจเยน ปจฺจโย.\csromanspacer{} ปุเรชาตปจฺจเยน ปจฺจโย.\csromanspacer{} กมฺมปจฺจเยน ปจฺจโย. (2)}

\prose{โนจิตฺโต ธมฺโม จิตฺตสฺส จ โนจิตฺตสฺส จ ธมฺมสฺส อารมฺมณปจฺจเยน ปจฺจโย.\csromanspacer{} สหชาตปจฺจเยน ปจฺจโย.\csromanspacer{} อุปนิสฺสยปจฺจเยน ปจฺจโย.\csromanspacer{} ปุเรชาตปจฺจเยน ปจฺจโย.\csromanspacer{} กมฺมปจฺจเยน ปจฺจโย. (3)}

\proseitem{88}{จิตฺโต จ โนจิตฺโต จ ธมฺมา จิตฺตสฺส ธมฺมสฺส อารมฺมณปจฺจเยน ปจฺจโย.\csromanspacer{} อุปนิสฺสยปจฺจเยน ปจฺจโย. (1)}

\prose{จิตฺโต จ โนจิตฺโต จ ธมฺมา โนจิตฺตสฺส ธมฺมสฺส อารมฺมณปจฺจเยน ปจฺจโย.\csromanspacer{} สหชาตปจฺจเยน ปจฺจเยน ปจฺจโย.\csromanspacer{} อุปนิสฺสยปจฺจเยน ปจฺจโย.\csromanspacer{} ปจฺฉาชาตปจฺจเยน ปจฺจโย. (2)}

\csromanmatikaanchor{mtk.56}
\csromantocmarkref{subsection}{57. เจตสิกทุก}{mtk.56}
\bookmark[dest={mtk.56},level=2]{57. เจตสิกทุก}
\csromanheader{2}{57. เจตสิกทุก}
\prose{\csromanfolio{403}จิตฺโต จ โนจิตฺโต จ ธมฺมา จิตฺตสฺส จ โนจิตฺตสฺส จ ธมฺมสฺส อารมฺมณปจฺจเยน ปจฺจโย.\csromanspacer{} อุปนิสฺสยปจฺจเยน ปจฺจโย. (3)}

\csromanheader{2}{2. ปจฺจยปจฺจนีย 2. สงฺขฺยาวาร}
\vspace{\tipitakaparskip}
\csromancenter{นเหตุยา นว, น-อารมฺมเณ นว, (สพฺพตฺถ นว.) โน-อวิคเต นว.}
\csromanheader{2}{3. ปจฺจยานุโลมปจฺจนีย}
\proseitem{90}{\textbf{เหตุ}ปจฺจยา น-อารมฺมเณ ตีณิ ฯเปฯ นสมนนฺตเร ตีณิ, น-อญฺญมญฺเญ เอกํ, น-อุปนิสฺสเย ตีณิ, (สพฺพตฺถ ตีณิ.) นมคฺเค ตีณิ, นสมฺปยุตฺเต เอกํ, นวิปฺปยุตฺเต ตีณิ, โนนตฺถิยา ตีณิ, โนวิคเต ตีณิ.}

\csromanheader{2}{4. ปจฺจยปจฺจนียานุโลม}
\csromanheader{2}{91. นเหตุปจฺจยา อารมฺมเณ นว, อธิปติยา นว. (อนุโลมมาติกา กาตพฺพา.)}
\nitthitam{จิตฺตทุกํ นิฏฺฐิตํ.}
\csromanheader{2}{57. เจตสิกทุก 1. ปฏิจฺจวาร}
\csromanheader{2}{1. ปจฺจยานุโลม 1. วิภงฺควาร}
\csromanheader{2}{\textbf{เหตุ}}
\proseitem{92}{เจตสิกํ ธมฺมํ ปฏิจฺจ เจตสิโก ธมฺโม อุปฺปชฺชติ เหตุปจฺจยา.\csromanspacer{} เจตสิกํ เอกํ ขนฺธํ ปฏิจฺจ เทฺว ขนฺธา.\csromanspacer{} เทฺว ขนฺเธ ปฏิจฺจ เอโก ขนฺโธ.\csromanspacer{} ปฏิสนฺธิกฺขเณ ฯเปฯ. (1)}

\prose{เจตสิกํ ธมฺมํ ปฏิจฺจ อเจตสิโก ธมฺโม อุปฺปชฺชติ เหตุปจฺจยา.\csromanspacer{} เจตสิเก ขนฺเธ ปฏิจฺจ จิตฺตญฺจ จิตฺตสมุฏฺฐานญฺจ รูปํ.\csromanspacer{} ปฏิสนฺธิกฺขเณ เจตสิเก ขนฺเธ ปฏิจฺจ จิตฺตญฺจ กฏตฺตา จ รูปํ. (2)}

\prose{\csromanfolio{404}เจตสิกํ ธมฺมํ ปฏิจฺจ เจตสิโก จ อเจตสิโก จ ธมฺมา อุปฺปชฺชนฺติ เหตุปจฺจยา.\csromanspacer{} เจตสิกํ เอกํ ขนฺธํ ปฏิจฺจ เทฺว ขนฺธา จิตฺตญฺจ จิตฺตสมุฏฺฐานญฺจ รูปํ.\csromanspacer{} เทฺว ขนฺเธ ฯเปฯ.\csromanspacer{} ปฏิสนฺธิกฺขเณ ฯเปฯ. (3)}

\proseitem{93}{อเจตสิกํ ธมฺมํ ปฏิจฺจ อเจตสิโก ธมฺโม อุปฺปชฺชติ เหตุปจฺจยา.\csromanspacer{} จิตฺตํ ปฏิจฺจ จิตฺตสมุฏฺฐานํ รูปํ.\csromanspacer{} ปฏิสนฺธิกฺขเณ จิตฺตํ ปฏิจฺจ กฏตฺตารูปํ.\csromanspacer{} จิตฺตํ ปฏิจฺจ วตฺถุ, วตฺถุํ ปฏิจฺจ จิตฺตํ.\csromanspacer{} เอกํ มหาภูตํ ฯเปฯ มหาภูเต ปฏิจฺจ จิตฺตสมุฏฺฐานํ รูปํ, กฏตฺตารูปํ, อุปาทารูปํ. (1)}

\prose{อเจตสิกํ ธมฺมํ ปฏิจฺจ เจตสิโก ธมฺโม อุปฺปชฺชติ เหตุปจฺจยา, จิตฺตํ ปฏิจฺจ สมฺปยุตฺตกา ขนฺธา.\csromanspacer{} ปฏิสนฺธิกฺขเณ จิตฺตํ ฯเปฯ.\csromanspacer{} ปฏิสนฺธิกฺขเณ วตฺถุํ ปฏิจฺจ เจตสิกา ขนฺธา. (2)}

\prose{อเจตสิกํ ธมฺมํ ปฏิจฺจ เจตสิโก จ อเจตสิโก จ ธมฺมา อุปฺปชฺชนฺติ เหตุปจฺจยา.\csromanspacer{} จิตฺตํ ปฏิจฺจ สมฺปยุตฺตกา ขนฺธา จิตฺตสมุฏฺฐานญฺจ รูปํ.\csromanspacer{} ปฏิสนฺธิกฺขเณ จิตฺตํ ฯเปฯ.\csromanspacer{} ปฏิสนฺธิกฺขเณ วตฺถุํ ปฏิจฺจ จิตฺตญฺจ สมฺปยุตฺตกา จ ขนฺธา. (3)}

\proseitem{94}{เจตสิกญฺจ อเจตสิกญฺจ ธมฺมํ ปฏิจฺจ เจตสิโก ธมฺโม อุปฺปชฺชติ เหตุปจฺจยา.\csromanspacer{} เจตสิกํ เอกํ ขนฺธญฺจ จิตฺตญฺจ ปฏิจฺจ เทฺว ขนฺธา.\csromanspacer{} เทฺว ขนฺเธ จ ฯเปฯ.\csromanspacer{} ปฏิสนฺธิกฺขเณ เจตสิกํ เอกํ ขนฺธญฺจ จิตฺตญฺจ ปฏิจฺจ เทฺว ขนฺธา.\csromanspacer{} เทฺว ขนฺเธ จ ฯเปฯ.\csromanspacer{} ปฏิสนฺธิกฺขเณ เจตสิกํ เอกํ ขนฺธญฺจ วตฺถุญฺจ ปฏิจฺจ เทฺว ขนฺธา.\csromanspacer{} เทฺว ขนฺเธ จ ฯเปฯ. (1)}

\prose{เจตสิกญฺจ อเจตสิกญฺจ ธมฺมํ ปฏิจฺจ อเจตสิโก ธมฺโม อุปฺปชฺชติ เหตุปจฺจยา.\csromanspacer{} เจตสิเก ขนฺเธ จ จิตฺตญฺจ ปฏิจฺจ จิตฺตสมุฏฺฐานํ รูปํ.\csromanspacer{} เจตสิเก ขนฺเธ จ มหาภูเต จ ปฏิจฺจ จิตฺตสมุฏฺฐานํ รูปํ.\csromanspacer{} ปฏิสนฺธิกฺขเณ เจตสิเก ขนฺเธ จ จิตฺตญฺจ ปฏิจฺจ กฏตฺตารูปํ.\csromanspacer{} ปฏิสนฺธิกฺขเณ เจตสิเก ขนฺเธ จ มหาภูเต จ ปฏิจฺจ กฏตฺตารูปํ.\csromanspacer{} ปฏิสนฺธิกฺขเณ เจตสิเก ขนฺเธ จ วตฺถุญฺจ ปฏิจฺจ จิตฺตํ. (2)}

\prose{เจตสิกญฺจ อเจตสิกญฺจ ธมฺมํ ปฏิจฺจ เจตสิโก จ อเจตสิโก จ ธมฺมา อุปฺปชฺชนฺติ เหตุปจฺจยา.\csromanspacer{} เจตสิกํ เอกํ ขนฺธญฺจ จิตฺตญฺจ ปฏิจฺจ เทฺว ขนฺธา จิตฺตสมุฏฺฐานญฺจ รูปํ.\csromanspacer{} เทฺว ขนฺเธ จ ฯเปฯ.\csromanspacer{} ปฏิสนฺธิกฺขเณ เจตสิกํ เอกํ}

\vspace{\tipitakaparskip}
\csromancenter{เจตสิกทุก ขนฺธญฺจ จิตฺตญฺจ ปฏิจฺจ เทฺว ขนฺธา กฏตฺตา จ รูปํ.\csromanspacer{} เทฺว ขนฺเธ จ ฯเปฯ.\csromanspacer{} ปฏิสนฺธิกฺขเณ เจตสิกํ เอกํ ขนฺธญฺจ วตฺถุญฺจ ปฏิจฺจ เทฺว ขนฺธา จิตฺตญฺจ.\csromanspacer{} เทฺว ขนฺเธ จ วตฺถุญฺจ ปฏิจฺจ เอโก ขนฺโธ จิตฺตญฺจ. (3)}
\csromanheader{2}{\textbf{อารมฺมณ}}
\proseitem{95}{\csromanfolio{405}เจตสิกํ ธมฺมํ ปฏิจฺจ เจตสิโก ธมฺโม อุปฺปชฺชติ อารมฺมณปจฺจยา.\csromanspacer{} เจตสิกํ เอกํ ขนฺธํ ปฏิจฺจ เทฺว ขนฺธา.\csromanspacer{} เทฺว ขนฺเธ ฯเปฯ.\csromanspacer{} ปฏิสนฺธิกฺขเณ ฯเปฯ. (1)}

\prose{เจตสิกํ ธมฺมํ ปฏิจฺจ อเจตสิโก ธมฺโม อุปฺปชฺชติ อารมฺมณปจฺจยา.\csromanspacer{} เจตสิเก ขนฺเธ ปฏิจฺจ จิตฺตํ.\csromanspacer{} ปฏิสนฺธิกฺขเณ ฯเปฯ.}

\prose{เจตสิกํ ธมฺมํ ปฏิจฺจ เจตสิโก จ อเจตสิโก จ ธมฺมา อุปฺปชฺชนฺติ อารมฺมณปจฺจยา.\csromanspacer{} เจตสิกํ เอกํ ขนฺธํ ปฏิจฺจ เทฺว ขนฺธา จิตฺตญฺจ.\csromanspacer{} เทฺว ขนฺเธ ฯเปฯ.\csromanspacer{} ปฏิสนฺธิกฺขเณ ฯเปฯ. (3)}

\proseitem{96}{อเจตสิกํ ธมฺมํ ปฏิจฺจ อเจตสิโก ธมฺโม อุปฺปชฺชติ อารมฺมณปจฺจยา.\csromanspacer{} ปฏิสนฺธิกฺขเณ วตฺถุํ ปฏิจฺจ จิตฺตํ. (1)}

\prose{อเจตสิกํ ธมฺมํ ปฏิจฺจ เจตสิโก ธมฺโม อุปฺปชฺชติ อารมฺมณปจฺจยา.\csromanspacer{} จิตฺตํ ปฏิจฺจ สมฺปยุตฺตกา ขนฺธา.\csromanspacer{} ปฏิสนฺธิกฺขเณ จิตฺตํ ปฏิจฺจ สมฺปยุตฺตกา ขนฺธา.\csromanspacer{} ปฏิสนฺธิกฺขเณ วตฺถุํ ปฏิจฺจ เจตสิกา ขนฺธา. (2)}

\prose{อเจตสิกํ ธมฺมํ ปฏิจฺจ เจตสิโก จ อเจตสิโก จ ธมฺมา อุปฺปชฺชนฺติ อารมฺมณปจฺจยา.\csromanspacer{} ปฏิสนฺธิกฺขเณ วตฺถุํ ปฏิจฺจ จิตฺตญฺจ สมฺปยุตฺตกา จ ขนฺธา. (3)}

\prose{เจตสิกญฺจ อเจตสิกญฺจ ธมฺมํ ปฏิจฺจ เจตสิโก ธมฺโม อุปฺปชฺชติ อารมฺมณปจฺจยา.\csromanspacer{} เจตสิกํ เอกํ ขนฺธญฺจ จิตฺตญฺจ ปฏิจฺจ เทฺว ขนฺธา.\csromanspacer{} เทฺว ขนฺเธ จ ฯเปฯ.\csromanspacer{} ปฏิสนฺธิกฺขเณ เจตสิกํ เอกํ ขนฺธญฺจ จิตฺตญฺจ ปฏิจฺจ เทฺว ขนฺธา.\csromanspacer{} เทฺว ขนฺเธ จ ฯเปฯ.\csromanspacer{} ปฏิสนฺธิกฺขเณ เจตสิกํ เอกํ ขนฺธญฺจ วตฺถุญฺจ ปฏิจฺจ เทฺว ขนฺธา.\csromanspacer{} เทฺว ขนฺเธ จ ฯเปฯ. (1)}

\prose{เจตสิกญฺจ อเจตสิกญฺจ ธมฺมํ ปฏิจฺจ อเจตสิโก ธมฺโม อุปฺปชฺชติ อารมฺมณปจฺจยา.\csromanspacer{} ปฏิสนฺธิกฺขเณ เจตสิเก ขนฺเธ จ วตฺถุญฺจ ปฏิจฺจ จิตฺตํ. (2)}

\prose{\csromanfolio{406}เจตสิกญฺจ อเจตสิกญฺจ ธมฺมํ ปฏิจฺจ เจตสิโก จ อเจตสิโก จ ธมฺมา อุปฺปชฺชนฺติ อารมฺมณปจฺจยา.\csromanspacer{} ปฏิสนฺธิกฺขเณ เจตสิกํ เอกํ ขนฺธญฺจ วตฺถุญฺจ ปฏิจฺจ เทฺว ขนฺธา จิตฺตญฺจ.\csromanspacer{} เทฺว ขนฺเธ จ ฯเปฯ. (3)}

\csromanheader{2}{\textbf{อธิปติ}}
\proseitem{97}{เจตสิกํ ธมฺมํ ปฏิจฺจ เจตสิโก ธมฺโม อุปฺปชฺชติ อธิปติปจฺจยา. (สํขิตฺตํ.)}

\csromanheader{2}{1. ปจฺจยานุโลม 2. สงฺขฺยาวาร}
\csromanheader{2}{\textbf{สุทฺธ}}
\proseitem{98}{เหตุยา นว, อารมฺมเณ นว, อธิปติยา นว, อนนฺตเร นว, สมนนฺตเร นว, สหชาเต นว, อญฺญมญฺเญ นว, นิสฺสเย นว, อุปนิสฺสเย นว, ปุเรชาเต ปญฺจ, อาเสวเน ปญฺจ, กมฺเม นว, (สพฺพตฺถ นว.) อวิคเต นว.}

\csromanheader{2}{2. ปจฺจยปจฺจนีย 1. วิภงฺควาร}
\csromanheader{2}{\textbf{นเหตุ}}
\proseitem{99}{เจตสิกํ ธมฺมํ ปฏิจฺจ เจตสิโก ธมฺโม อุปฺปชฺชติ นเหตุปจฺจยา.\csromanspacer{} อเหตุกํ เจตสิกํ เอกํ ขนฺธํ ปฏิจฺจ เทฺว ขนฺธา.\csromanspacer{} เทฺว ขนฺเธ ฯเปฯ.\csromanspacer{} อเหตุกปฏิสนฺธิกฺขเณ ฯเปฯ.\csromanspacer{} วิจิกิจฺฉาสหคเต อุทฺธจฺจสหคเต ขนฺเธ ปฏิจฺจ วิจิกิจฺฉาสหคโต อุทฺธจฺจสหคโต โมโห. (1)}

\prose{เจตสิกํ ธมฺมํ ปฏิจฺจ อเจตสิโก ธมฺโม อุปฺปชฺชติ นเหตุปจฺจยา.\csromanspacer{} อเหตุเก เจตสิเก ขนฺเธ ปฏิจฺจ จิตฺตญฺจ จิตฺตสมุฏฺฐานญฺจ รูปํ.\csromanspacer{} อเหตุกปฏิสนฺธิกฺขเณ ฯเปฯ. (2)}

\prose{เจตสิกํ ธมฺมํ ปฏิจฺจ เจตสิโก จ อเจตสิโก จ ธมฺมา อุปฺปชฺชนฺติ นเหตุปจฺจยา.\csromanspacer{} อเหตุกํ เจตสิกํ เอกํ ขนฺธํ ปฏิจฺจ เทฺว ขนฺธา จิตฺตญฺจ จิตฺตสมุฏฺฐานญฺจ รูปํ.\csromanspacer{} เทฺว ขนฺเธ ฯเปฯ.\csromanspacer{} อเหตุกปฏิสนฺธิกฺขเณ ฯเปฯ. (3)}

\csromanheader{2}{57. เจตสิกทุก}
\proseitem{100}{\csromanfolio{407}อเจตสิกํ ธมฺมํ ปฏิจฺจ อเจตสิโก ธมฺโม อุปฺปชฺชติ นเหตุปจฺจยา.\csromanspacer{} อเหตุกํ จิตฺตํ ปฏิจฺจ จิตฺตสมุฏฺฐานํ รูปํ.\csromanspacer{} อเหตุกปฏิสนฺธิกฺขเณ จิตฺตํ ปฏิจฺจ กฏตฺตารูปํ.\csromanspacer{} จิตฺตํ ปฏิจฺจ วตฺถุ, วตฺถุํ ปฏิจฺจ จิตฺตํ.\csromanspacer{} เอกํ มหาภูตํ ฯเปฯ.\csromanspacer{} อสญฺญสตฺตานํ เอกํ มหาภูตํ ฯเปฯ. (1)}

\prose{อเจตสิกํ ธมฺมํ ปฏิจฺจ เจตสิโก ธมฺโม อุปฺปชฺชติ นเหตุปจฺจยา.\csromanspacer{} อเหตุกํ จิตฺตํ ปฏิจฺจ สมฺปยุตฺตกํ ขนฺธา.\csromanspacer{} อเหตุกปฏิสนฺธิกฺขเณ จิตฺตํ ปฏิจฺจ สมฺปยุตฺตกา ขนฺธา.\csromanspacer{} อเหตุกปฏิสนฺธิกฺขเณ วตฺถุํ ปฏิจฺจ เจตสิกา ขนฺธา.\csromanspacer{} วิจิกิจฺฉาสหคตํ อุทฺธจฺจสหคตํ จิตฺตํ ปฏิจฺจ วิจิกิจฺฉาสหคโต อุทฺธจฺจสหคโต โมโห. (2)}

\prose{อเจตสิกํ ธมฺมํ ปฏิจฺจ เจตสิโก จ อเจตสิโก จ ธมฺมา อุปฺปชฺชนฺติ นเหตุปจฺจยา.\csromanspacer{} อเหตุกํ จิตฺตํ ปฏิจฺจ สมฺปยุตฺตกา ขนฺธา จิตฺตสมุฏฺฐานญฺจ รูปํ.\csromanspacer{} อเหตุกปฏิสนฺธิกฺขเณ จิตฺตํ ปฏิจฺจ สมฺปยุตฺตกา ขนฺธา กฏตฺตา จ รูปํ.\csromanspacer{} อเหตุกปฏิสนฺธิกฺขเณ วตฺถุํ ปฏิจฺจ จิตฺตญฺจ สมฺปยุตฺตกา จ ขนฺธา. (3)}

\proseitem{101}{เจตสิกญฺจ อเจตสิกญฺจ ธมฺมํ ปฏิจฺจ เจตสิโก ธมฺโม อุปฺปชฺชติ นเหตุปจฺจยา.\csromanspacer{} อเหตุกํ เจตสิกํ เอกํ ขนฺธญฺจ จิตฺตญฺจ ปฏิจฺจ เทฺว ขนฺธา.\csromanspacer{} เทฺว ขนฺเธ จ ฯเปฯ.\csromanspacer{} อเหตุกปฏิสนฺธิกฺขเณ เจตสิกํ เอกํ ขนฺธญฺจ จิตฺตญฺจ ปฏิจฺจ เทฺว ขนฺธา.\csromanspacer{} เทฺว ขนฺเธ จ ฯเปฯ.\csromanspacer{} อเหตุกปฏิสนฺธิกฺขเณ เจตสิกํ เอกํ ขนฺธญฺจ วตฺถุญฺจ ปฏิจฺจ เทฺว ขนฺธา.\csromanspacer{} เทฺว ขนฺเธ จ ฯเปฯ.\csromanspacer{} วิจิกิจฺฉาสหคเต อุทฺธจฺจสหคเต ขนฺเธ จ จิตฺตญฺจ ปฏิจฺจ วิจิกิจฺฉาสหคโต อุทฺธจฺจสหคโต โมโห. (1)}

\prose{เจตสิกญฺจ อเจตสิกญฺจ ธมฺมํ ปฏิจฺจ อเจตสิโก ธมฺโม อุปฺปชฺชติ นเหตุปจฺจยา.\csromanspacer{} อเหตุเก เจตสิเก ขนฺเธ จ จิตฺตญฺจ ปฏิจฺจ จิตฺตสมุฏฺฐานํ รูปํ.\csromanspacer{} อเหตุเก เจตสิเก ขนฺเธ จ มหาภูเต จ ปฏิจฺจ จิตฺตสมุฏฺฐานํ รูปํ.\csromanspacer{} อเหตุกปฏิสนฺธิกฺขเณ เจตสิเก ขนฺเธ จ จิตฺตญฺจ ปฏิจฺจ กฏตฺตารูปํ.\csromanspacer{} อเหตุกปฏิสนฺธิกฺขเณ เจตสิเก ขนฺเธ จ มหาภูเต จ ปฏิจฺจ กฏตฺตารูปํ.\csromanspacer{} อเหตุกปฏิสนฺธิกฺขเณ เจตสิเก ขนฺเธ จ วตฺถุญฺจ ปฏิจฺจ จิตฺตํ. (2)}

\prose{เจตสิกญฺจ อเจตสิกญฺจ ธมฺมํ ปฏิจฺจ เจตสิโก จ อเจตสิโก จ ธมฺมา อุปฺปชฺชนฺติ นเหตุปจฺจยา.\csromanspacer{} อเหตุกํ เจตสิกํ เอกํ ขนฺธญฺจ จิตฺตญฺจ \csromanfolio{408}ปฏิจฺจ เทฺว ขนฺธา จิตฺตสมุฏฺฐานญฺจ รูปํ.\csromanspacer{} อเหตุกปฏิสนฺธิกฺขเณ เจตสิกํ เอกํ ขนฺธญฺจ จิตฺตญฺจ ปฏิจฺจ เทฺว ขนฺธา กฏตฺตา จ รูปํ.\csromanspacer{} อเหตุกปฏิสนฺธิกฺขเณ เจตสิกํ เอกํ ขนฺธญฺจ วตฺถุญฺจ ปฏิจฺจ เทฺว ขนฺธา จิตฺตญฺจ.\csromanspacer{} เทฺว ขนฺเธ จ ฯเปฯ. (3)}

\csromanheader{2}{\textbf{น-อารมฺมณ}}
\proseitem{102}{เจตสิกํ ธมฺมํ ปฏิจฺจ อเจตสิโก ธมฺโม อุปฺปชฺชติ น-อารมฺมณปจฺจยา.\csromanspacer{} เจตสิเก ขนฺเธ ปฏิจฺจ จิตฺตสมุฏฺฐานํ รูปํ.\csromanspacer{} ปฏิสนฺธิกฺขเณ ฯเปฯ. (1)}

\prose{อเจตสิกํ ธมฺมํ ปฏิจฺจ อเจตสิโก ธมฺโม อุปฺปชฺชติ น-อารมฺมณปจฺจยา.\csromanspacer{} จิตฺตํ ปฏิจฺจ จิตฺตสมุฏฺฐานํ รูปํ.\csromanspacer{} ปฏิสนฺธิกฺขเณ ฯเปฯ. (ยาว อสญฺญสตฺตา.) (1)}

\prose{เจตสิกญฺจ อเจตสิกญฺจ ธมฺมํ ปฏิจฺจ อเจตสิโก ธมฺโม อุปฺปชฺชติ น-อารมฺมณปจฺจยา.\csromanspacer{} เจตสิเก ขนฺเธ จ จิตฺตญฺจ ปฏิจฺจ จิตฺตสมุฏฺฐานํ รูปํ.\csromanspacer{} เจตสิเก ขนฺเธ จ มหาภูเต จ ปฏิจฺจ จิตฺตสมุฏฺฐานํ รูปํ. (ปฏิสนฺธิกฺขเณ เทฺวปิ กาตพฺพา, สํขิตฺตํ.)}

\csromanheader{2}{2. ปจฺจยปจฺจนีย 2. สงฺขฺยาวาร}
\csromanheader{2}{\textbf{สุทฺธ}}
\proseitem{103}{นเหตุยา นว, น-อารมฺมเณ ตีณิ, น-อธิปติยา นว, น-อนนฺตเร ตีณิ, นสมนนฺตเร ตีณิ, น-อญฺญมญฺเญ ตีณิ, น-อุปนิสฺสเย ตีณิ, นปุเรชาเต นว, นปจฺฉาชาเต นว, น-อาเสวเน นว, นกมฺเม จตฺตาริ, นวิปาเก นว, น-อาหาเร เอกํ, น-อินฺทฺริเย เอกํ, นฌาเน ฉ, นมคฺเค นว, นสมฺปยุตฺเต ตีณิ, นวิปฺปยุตฺเต ฉ, โนนตฺถิยา ตีณิ, โนวิคเต ตีณิ.}

\csromanheader{2}{3. ปจฺจยานุโลมปจฺจนีย}
\proseitem{104}{เหตุปจฺจยา น-อารมฺมเณ ตีณิ, น-อธิปติยา นว. (สํขิตฺตํ.)}

\csromanheader{2}{57. เจตสิกทุก \textbf{4. ปจฺจยปจฺจนียานุโลม}}
\proseitem{105}{\csromanfolio{409}นเหตุปจฺจยา อารมฺมเณ นว, อนนฺตเร นว ฯเปฯ ปุเรชาเต ปญฺจ, อาเสวเน ปญฺจ, กมฺเม นว, (สพฺพตฺถ นว.) มคฺเค ตีณิ, อวิคเต นว.}

\vspace{\tipitakaparskip}
\csromancenter{(สหชาตวาโรปิ ปฏิจฺจวารสทิโส.)}
\csromanheader{2}{57. เจตสิกทุก 3. ปจฺจยวาร}
\csromanheader{2}{1. ปจฺจยานุโลม 1. วิภงฺควาร}
\csromanheader{2}{\textbf{เหตุ}}
\proseitem{106}{เจตสิกํ ธมฺมํ ปจฺจยา เจตสิโก ธมฺโม อุปฺปชฺชติ เหตุปจฺจยา.\csromanspacer{} ตีณิ. (ปฏิจฺจสทิสา.)}

\prose{อเจตสิกํ ธมฺมํ ปจฺจยา อเจตสิโก ธมฺโม อุปฺปชฺชติ เหตุปจฺจยา.\csromanspacer{} จิตฺตํ ปจฺจยา จิตฺตสมุฏฺฐานํ รูปํ.\csromanspacer{} วตฺถุํ ปจฺจยา จิตฺตํ.\csromanspacer{} ปฏิสนฺธิกฺขเณ จิตฺตํ ปจฺจยา กฏตฺตารูปํ.\csromanspacer{} จิตฺตํ ปจฺจยา วตฺถุ, วตฺถุํ ปจฺจยา จิตฺตํ.\csromanspacer{} เอกํ มหาภูตํ ฯเปฯ. (1)}

\prose{อเจตสิกํ ธมฺมํ ปจฺจยา เจตสิโก ธมฺโม อุปฺปชฺชติ เหตุปจฺจยา.\csromanspacer{} จิตฺตํ ปจฺจยา สมฺปยุตฺตกา ขนฺธา, วตฺถุํ ปจฺจยา เจตสิกา ขนฺธา. (ปฏิสนฺธิกฺขเณ เทฺวปิ กาตพฺพา.) (2)}

\prose{อเจตสิกํ ธมฺมํ ปจฺจยา เจตสิโก จ อเจตสิโก จ ธมฺมา อุปฺปชฺชนฺติ เหตุปจฺจยา.\csromanspacer{} จิตฺตํ ปจฺจยา สมฺปยุตฺตกา ขนฺธา จิตฺตสมุฏฺฐานญฺจ รูปํ.\csromanspacer{} วตฺถุํ ปจฺจยา จิตฺตญฺจ สมฺปยุตฺตกา จ ขนฺธา. (ปฏิสนฺธิกฺขเณ เทฺวปิ กาตพฺพา.) (3)}

\proseitem{107}{เจตสิกญฺจ อเจตสิกญฺจ ธมฺมํ ปจฺจยา เจตสิโก ธมฺโม อุปฺปชฺชติ เหตุปจฺจยา.\csromanspacer{} เจตสิกํ เอกํ ขนฺธญฺจ จิตฺตญฺจ ปจฺจยา เทฺว ขนฺธา.\csromanspacer{} เทฺว ขนฺเธ จ ฯเปฯ.\csromanspacer{} เจตสิกํ เอกํ ขนฺธญฺจ วตฺถุญฺจ ปจฺจยา เทฺว ขนฺธา.\csromanspacer{} เทฺว ขนฺเธ จ ฯเปฯ. (ปฏิสนฺธิกฺขเณ เทฺวปิ กาตพฺพา.) (1)}

\prose{\csromanfolio{410}เจตสิกญฺจ อเจตสิกญฺจ ธมฺมํ ปจฺจยา อเจตสิโก ธมฺโม อุปฺปชฺชติ เหตุปจฺจยา.\csromanspacer{} เจตสิเก ขนฺเธ จ จิตฺตญฺจ ปจฺจยา จิตฺตสมุฏฺฐานํ รูปํ.\csromanspacer{} เจตสิเก ขนฺเธ จ มหาภูเต จ ปจฺจยา จิตฺตสมุฏฺฐานํ รูปํ.\csromanspacer{} เจตสิเก ขนฺเธ จ วตฺถุญฺจ ปจฺจยา จิตฺตํ. (ปฏิสนฺธิกฺขเณ ตีณิปิ กาตพฺพา.) (2)}

\prose{เจตสิกญฺจ อเจตสิกญฺจ ธมฺมํ ปจฺจยา เจตสิโก จ อเจตสิโก จ ธมฺมา อุปฺปชฺชนฺติ เหตุปจฺจยา.\csromanspacer{} เจตสิกํ เอกํ ขนฺธญฺจ จิตฺตญฺจ ปจฺจยา เทฺว ขนฺธา จิตฺตสมุฏฺฐานญฺจ รูปํ.\csromanspacer{} เทฺว ขนฺเธ จ ฯเปฯ.\csromanspacer{} เจตสิกํ เอกํ ขนฺธญฺจ วตฺถุญฺจ ปจฺจยา เทฺว ขนฺธา จิตฺตญฺจ.\csromanspacer{} เทฺว ขนฺเธ จ ฯเปฯ. (ปฏิสนฺธิกฺขเณ เทฺวปิ กาตพฺพา.) (3)}

\csromanheader{2}{\textbf{อรมฺมณ}}
\proseitem{108}{เจตสิกํ ธมฺมํ ปจฺจยา เจตสิโก ธมฺโม อุปฺปชฺชติ อารมฺมณปจฺจยา.\csromanspacer{} ตีณิ. (ปฏิจฺจสทิสา.)}

\prose{อเจตสิกํ ธมฺมํ ปจฺจยา อเจตสิโก ธมฺโม อุปฺปชฺชติ อารมฺมณปจฺจยา.\csromanspacer{} จกฺขายตนํ ปจฺจยา จกฺขุวิญฺญาณํ ฯเปฯ.\csromanspacer{} กายายตนํ ปจฺจยา กายวิญฺญาณํ.\csromanspacer{} วตฺถุํ ปจฺจยา จิตฺตํ.\csromanspacer{} ปฏิสนฺธิกฺขเณ ฯเปฯ. (1)}

\prose{อเจตสิกํ ธมฺมํ ปจฺจยา เจตสิโก ธมฺโม อุปฺปชฺชติ อารมฺมณปจฺจยา.\csromanspacer{} จกฺขายตนํ ปจฺจยา จกฺขุวิญฺญาณสหคตา ขนฺธา ฯเปฯ.\csromanspacer{} กายายตนํ ปจฺจยา ฯเปฯ.\csromanspacer{} จิตฺตํ ปจฺจยา สมฺปยุตฺตกา ขนฺธา.\csromanspacer{} วตฺถุํ ปจฺจยา เจตสิกา ขนฺธา. (ปฏิสนฺธิกฺขเณ เทฺวปิ กาตพฺพา.) (2)}

\prose{อเจตสิกํ ธมฺมํ ปจฺจยา เจตสิโก จ อเจตสิโก จ ธมฺมา อุปฺปชฺชนฺติ อารมฺมณปจฺจยา.\csromanspacer{} จกฺขายตนํ ปจฺจยา จกฺขุวิญฺญาณํ สมฺปยุตฺตกา จ ขนฺธา ฯเปฯ.\csromanspacer{} กายายตนํ ปจฺจยา.\csromanspacer{} วตฺถุํ ปจฺจยา จิตฺตํ สมฺปยุตฺตกา จ ขนฺธา. (ปฏิสนฺธิกฺขเณ เอกํ.) (3)}

\proseitem{109}{เจตสิกญฺจ อเจตสิกญฺจ ธมฺมํ ปจฺจยา เจตสิโก ธมฺโม อุปฺปชฺชติ อารมฺมณปจฺจยา.\csromanspacer{} จกฺขุวิญฺญาณสหคตํ เอกํ ขนฺธญฺจ จกฺขุวิญฺญาณญฺจ ปจฺจยา เทฺว ขนฺธา.\csromanspacer{} เทฺว ขนฺเธ จ ฯเปฯ.\csromanspacer{} จกฺขุวิญฺญาณสหคตํ เอกํ ขนฺธญฺจ จกฺขายตนญฺจ ปจฺจยา เทฺว ขนฺธา.\csromanspacer{} เทฺว ขนฺเธ จ ฯเปฯ.\csromanspacer{} กายวิญฺญาณสหคตํ ฯเปฯ.\csromanspacer{} เจตสิกํ เอกํ ขนฺธญฺจ จิตฺตญฺจ ปจฺจยา เทฺว ขนฺธา.\csromanspacer{} เทฺว ขนฺเธ จ ฯเปฯ.\csromanspacer{} เจตสิกํ เอกํ ขนฺธญฺจ วตฺถุญฺจ ปจฺจยา เทฺว ขนฺธา.\csromanspacer{} เทฺว ขนฺเธ จ ฯเปฯ. (ปฏิสนฺธิกฺขเณ เทฺว.) (1)}

\csromanheader{2}{57. เจตสิกทุก}
\prose{\csromanfolio{411}เจตสิกญฺจ อเจตสิกญฺจ ธมฺมํ ปจฺจยา อเจตสิโก ธมฺโม อุปฺปชฺชติ อารมฺมณปจฺจยา.\csromanspacer{} จกฺขุวิญฺญาณสหคเต ขนฺเธ จ จกฺขายตนญฺจ ปจฺจยา จกฺขุวิญฺญาณํ ฯเปฯ.\csromanspacer{} กายวิญฺญาณสหคเต ฯเปฯ.\csromanspacer{} เจตสิเก ขนฺเธ จ วตฺถุญฺจ ปจฺจยา จิตฺตํ. (ปฏิสนฺธิกฺขเณ เอกํ.)}

\prose{เจตสิกญฺจ อเจตสิกญฺจ ธมฺมํ ปจฺจยา เจตสิโก จ อเจตสิโก จ ธมฺมา อุปฺปชฺชนฺติ อารมฺมณปจฺจยา.\csromanspacer{} จกฺขุวิญฺญาณสหคตํ เอกํ ขนฺธญฺจ จกฺขายตนญฺจ ปจฺจยา เทฺว ขนฺธา จกฺขุวิญฺญาณญฺจ.\csromanspacer{} เทฺว ขนฺเธ จ ฯเปฯ.\csromanspacer{} เจตสิกํ เอกํ ขนฺธญฺจ วตฺถุญฺจ ปจฺจยา เทฺว ขนฺธา จิตฺตญฺจ.\csromanspacer{} เทฺว ขนฺเธ จ ฯเปฯ. (ปฏิสนฺธิกฺขเณ เอกํ.\csromanspacer{} สํ ขิตฺตํ.) (3)}

\csromanheader{2}{1. ปจฺจยานุโลม 2. สงฺขฺยาวาร}
\proseitem{110}{เหตุยา นว, อารมฺมเณ นว, อธิปติยา นว, (สพฺพตฺถ นว.) ปุเรชาเต นว, อาเสวเน นว ฯเปฯ อวิคเต นว.}

\csromanheader{2}{2. ปจฺจยปจฺจนีย 1. วิภงฺควาร}
\csromanheader{2}{\textbf{นเหตุ}}
\proseitem{111}{เจตสิกํ ธมฺมํ ปจฺจยา เจตสิโก ธมฺโม อุปฺปชฺชติ นเหตุปจฺจยา.\csromanspacer{} อเหตุกํ เจตสิกํ. (สํขิตฺตํ.)}

\prose{(นว ปญฺหา, ปญฺจวิญฺญาณมฺปิ ยถา อารมฺมณปจฺจยา เอวํ กาตพฺพํ.\csromanspacer{} ตีสุเยว โมโห.\csromanspacer{} สพฺเพ ปญฺหา ปวตฺติปฏิสนฺธิยา กาตพฺพา อสมฺโมหนฺเตน.)}

\csromanheader{2}{2. ปจฺจยปจฺจนีย 2. สงฺขฺยาวาร}
\csromanheader{2}{\textbf{สุทฺธ}}
\vspace{\tipitakaparskip}
\csromancenter{\textbf{นเหตุ}ยา นว, น-อารมฺมเณ ตีณิ, น-อธิปติยา นว, น-อนนฺตเร ตีณิ, นสมนนฺตเร ตีณิ, น-อญฺญมญฺเญ ตีณิ, น-อุปนิสฺสเย ตีณิ, นปุเรชาเต นว, นปจฺฉาชาเต นว, น-อาเสวเน นว, นกมฺเม}
\prosecont{\csromanfolio{412}จตฺตาริ, นวิปาเก นว, น-อาหาเร เอกํ, น-อินฺทฺริเย เอกํ, นฌาเน นว, นมคฺเค นว, นสมฺปยุตฺเต ตีณิ, นวิปฺปยุตฺเต ฉ, โนนตฺถิยา ตีณิ, โนวิคเต ตีณิ.}

\csromanheader{2}{3. ปจฺจยานุโลมปจฺจนีย}
\proseitem{113}{\textbf{เหตุ}ปจฺจยา น-อารมฺมเณ ตีณิ, น-อธิปติยา นว. (สํขิตฺตํ.)}

\csromanheader{2}{4. ปจฺจยปจฺจนียานุโลม}
\proseitem{114}{นเหตุปจฺจยา อารมฺมเณ นว, อนนฺตเร นว, (สพฺพตฺถ นว.) มคฺเค ตีณิ ฯเปฯ อวิคเต นว.}

\vspace{\tipitakaparskip}
\csromancenter{(นิสฺสยวาโร ปจฺจยวารสทิโส.)}
\csromanheader{2}{57. เจตสิกทุก 5. สํสฏฺฐวาร}
\csromanheader{2}{1. ปจฺจยานุโลม 1. วิภงฺควาร}
\csromanheader{2}{\textbf{เหตุ}}
\proseitem{115}{เจตสิกํ ธมฺมํ สํสฏฺโฐ เจตสิโก ธมฺโม อุปฺปชฺชติ เหตุปจฺจยา.\csromanspacer{} เจตสิกํ เอกํ ขนฺธํ สํสฏฺฐา เทฺว ขนฺธา.\csromanspacer{} เทฺว ขนฺเธ ฯเปฯ.\csromanspacer{} ปฏิสนฺธิกฺขเณ ฯเปฯ. (1)}

\prose{เจตสิกํ ธมฺมํ สํสฏฺโฐ อเจตสิโก ธมฺโม อุปฺปชฺชติ เหตุปจฺจยา.\csromanspacer{} เจตสิเก ขนฺเธ สํสฏฺฐํ จิตฺตํ.\csromanspacer{} ปฏิสนฺธิกฺขเณ ฯเปฯ.}

\prose{เจตสิกํ ธมฺมํ สํสฏฺโฐ เจตสิโก จ อเจตสิโก จ ธมฺมา อุปฺปชฺชนฺติ เหตุปจฺจยา.\csromanspacer{} เจตสิกํ เอกํ ขนฺธํ สํสฏฺฐา เทฺว ขนฺธา จิตฺตญฺจ.\csromanspacer{} เทฺว ขนฺเธ ฯเปฯ.\csromanspacer{} ปฏิสนฺธิกฺขเณ ฯเปฯ. (3)}

\prose{อเจตสิกํ ธมฺมํ สํสฏฺโฐ เจตสิโก ธมฺโม อุปฺปชฺชติ เหตุปจฺจยา.\csromanspacer{} จิตฺตํ สํสฏฺฐา สมฺปยุตฺตกา ขนฺธา.\csromanspacer{} ปฏิสนฺธิกฺขเณ ฯเปฯ. (1)}

\csromanheader{2}{57. เจตสิกทุก}
\prose{\csromanfolio{413}เจตสิกญฺจ อเจตสิกญฺจ ธมฺมํ สํสฏฺโฐ เจตสิโก ธมฺโม อุปฺปชฺชติ เหตุปจฺจยา.\csromanspacer{} เจตสิกํ เอกํ ขนฺธญฺจ จิตฺตญฺจ สํสฏฺฐา เทฺว ขนฺธา.\csromanspacer{} เทฺว ขนฺเธ จ ฯเปฯ.\csromanspacer{} ปฏิสนฺธิกฺขเณ ฯเปฯ. (1) (สํขิตฺตํ.)}

\csromanheader{2}{1. ปจฺจยานุโลม 2. สงฺขฺยาวาร}
\proseitem{116}{\textbf{เหตุ}ยา ปญฺจ, อารมฺมเณ ปญฺจ, อธิปติยา ปญฺจ, (สพฺพตฺถ ปญฺจ.) อวิคเต ปญฺจ.}

\csromanheader{2}{2. ปจฺจยปจฺจนีย 1. วิภงฺควาร}
\proseitem{117}{เจตสิกํ ธมฺมํ สํสฏฺโฐ เจตสิโก ธมฺโม อุปฺปชฺชติ นเหตุปจฺจยา. (เอวํ ปญฺจปิ ปญฺหา กาตพฺพา, ตีณิเยว โมโห.\csromanspacer{} สํขิตฺตํ.)}

\csromanheader{2}{2. ปจฺจยปจฺจนีย 2. สงฺขฺยาวาร}
\proseitem{118}{นเหตุยา ปญฺจ, น-อธิปติยา ปญฺจ, นปุเรชาเต ปญฺจ, นปจฺฉาชาเต ปญฺจ, น-อาเสวเน ปญฺจ, นกมฺเม ตีณิ, นวิปาเก ปญฺจ, นฌาเน ปญฺจ, นมคฺเค ปญฺจ, นวิปฺปยุตฺเต ปญฺจ. (เอวํ อิตเร เทฺว คณนาปิ สมฺปยุตฺตวาโรปิ กาตพฺโพ.)}

\csromanheader{2}{57. เจตสิกทุก 7. ปญฺหาวาร}
\csromanheader{2}{1. ปจฺจยานุโลม 1. วิภงฺควาร}
\csromanheader{2}{\textbf{เหตุ}}
\proseitem{119}{เจตสิโก ธมฺโม เจตสิกสฺส ธมฺมสฺส เหตุปจฺจเยน ปจฺจโย.\csromanspacer{} เจตสิกา เหตู สมฺปยุตฺตกานํ ขนฺธานํ เหตุปจฺจเยน ปจฺจโย.\csromanspacer{} ปฏิสนฺธิกฺขเณ ฯเปฯ. (1)}

\prose{เจตสิโก ธมฺโม อเจตสิกสฺส ธมฺมสฺส เหตุปจฺจเยน ปจฺจโย.\csromanspacer{} เจตสิกา เหตู จิตฺตสฺส จิตฺตสมุฏฺฐานานญฺจ รูปานํ เหตุปจฺจเยน ปจฺจโย.\csromanspacer{} ปฏิสนฺธิกฺขเณ ฯเปฯ. (2)}

\prose{\csromanfolio{414}เจตสิโก ธมฺโม เจตสิกสฺส จ อเจตสิกสฺส จ ธมฺมสฺส เหตุปจฺจเยน ปจฺจโย.\csromanspacer{} เจตสิกา เหตู สมฺปยุตฺตกานํ ขนฺธานํ จิตฺตสฺส จ จิตฺตสมุฏฺฐานาญฺจ รูปานํ เหตุปจฺจเยน ปจฺจโย.\csromanspacer{} ปฏิสนฺธิกฺขเณ ฯเปฯ. (3)}

\csromanheader{2}{\textbf{อารมฺมณ}}
\proseitem{120}{เจตสิโก ธมฺโม เจตสิกสฺส ธมฺมสฺส อารมฺมณปจฺจเยน ปจฺจโย.\csromanspacer{} เจตสิเก ขนฺเธ อารพฺภ เจตสิกา ขนฺธา อุปฺปชฺชนฺติ. (มูลํ ปุจฺฉิตพฺพํ.) เจตสิเก ขนฺเธ อารพฺภ จิตฺตํ อุปฺปชฺชติ. (มูลํ ปุจฺฉิตพฺพํ.) เจตสิเก ขนฺเธ อารพฺภ เจตสิกา ขนฺธา จ จิตฺตญฺจ อุปฺปชฺชนฺติ. (3)}

\proseitem{121}{อเจตสิโก ธมฺโม อเจตสิกสฺส ธมฺมสฺส อารมฺมณปจฺจเยน ปจฺจโย.\csromanspacer{} อริยา มคฺคา วุฏฺฐหิตฺวา ฯเปฯ.\csromanspacer{} ผลํ ปจฺจเวกฺขนฺติ, นิพฺพานํ ปจฺจเวกฺขนฺติ.\csromanspacer{} นิพฺพานํ โคตฺรภุสฺส, โวทานสฺส, มคฺคสฺส, ผลสฺส, อาวชฺชนาย อารมฺมณปจฺจเยน ปจฺจโย.\csromanspacer{} จกฺขุํ ฯเปฯ วตฺถุํ อเจตสิเก ขนฺเธ อนิจฺจโต ฯเปฯ วิปสฺสติ, อสฺสาเทติ อภินนฺทติ, ตํ อารพฺภ จิตฺตํ อุปฺปชฺชติ.\csromanspacer{} ทิพฺเพน จกฺขุนา รูปํ ปสฺสติ.\csromanspacer{} ทิพฺพาย โสตธาตุยา สทฺทํ สุณาติ.\csromanspacer{} เจโตปริยญาเณน อเจตสิกจิตฺตสมงฺคิสฺส จิตฺตํ ชานาติ.\csromanspacer{} อากาสานญฺจายตนํ ฯเปฯ.\csromanspacer{} อากิญฺจญฺญายตนํ ฯเปฯ.\csromanspacer{} รูปายตนํ จกฺขุวิญฺญาณสฺส ฯเปฯ.\csromanspacer{} โผฏฺฐพฺพายตนํ กายวิญฺญาณสฺส ฯเปฯ.\csromanspacer{} อเจตสิกา ขนฺธา อิทฺธิวิธญาณสฺส, เจโตปริยญาณสฺส, ปุพฺเพนิวาสานุสฺสติญาณสฺส, อนาคตํสญาณสฺส, อาวชฺชนาย อารมฺมณปจฺจเยน ปจฺจโย. (1)}

\prose{อเจตสิโก ธมฺโม เจตสิกสฺส ธมฺมสฺส อารมฺมณปจฺจเยน ปจฺจโย.\csromanspacer{} อริยา มคฺคา วุฏฺฐหิตฺวา ฯเปฯ.\csromanspacer{} นิพฺพานํ ปจฺจเวกฺขนฺติ. (ปฐมคมนสทิสํ.) จกฺขุํ ฯเปฯ วตฺถุํ อเจตสิเก ขนฺเธ อนิจฺจโต ฯเปฯ โทมนสฺสํ อุปฺปชฺชติ.\csromanspacer{} ทิพฺเพน จกฺขุนา รูปํ ปสฺสติ.\csromanspacer{} ทิพฺพาย โสตธาตุยา สทฺทํ สุณาติ.\csromanspacer{} เจโตปริยญาเณน อเจตสิกจิตฺตสมงฺคิสฺส จิตฺตํ ชานาติ.\csromanspacer{} อากาสานญฺจายตนํ ฯเปฯ.\csromanspacer{} อากิญฺจญฺญายตนํ เนวสญฺญานาสญฺญายตนสฺส ฯเปฯ.\csromanspacer{} รูปายตนํ จกฺขุวิญฺญาณสหคตานํ ขนฺธานํ ฯเปฯ.\csromanspacer{} โผฏฺฐพฺพายตนํ กายวิญฺญาณสหคตานํ ขนฺธานํ ฯเปฯ.\csromanspacer{} อเจตสิกา ขนฺธา อิทฺธิวิธญาณสฺส, เจโตปริยญาณสฺส, ปุพฺเพนิวาสานุสฺสติญาณสฺส, อนาคตํสญาณสฺส, อาวชฺชนาย อารมฺมณปจฺจเยน ปจฺจโย. (2)}

\csromanheader{2}{57. เจตสิกทุก}
\prose{\csromanfolio{415}อเจตสิโก ธมฺโม เจตสิกสฺส จ อเจตสิกสฺส จ ธมฺมสฺส อารมฺมณปจฺจเยน ปจฺจโย.\csromanspacer{} อริยา มคฺคา วุฏฺฐหิตฺวา ฯเปฯ.\csromanspacer{} นิพฺพานํ ปจฺจเวกฺขนฺติ. (ปฐมคมนสทิสํ.) จกฺขุํ ฯเปฯ วตฺถุํ อเจตสิเก ขนฺเธ อนิจฺจโต ฯเปฯ วิปสฺสติ, อสฺสาเทติ อภินนฺทติ, ตํ อารพฺภ จิตฺตญฺจ สมฺปยุตฺตกา จ ขนฺธา อุปฺปชฺชนฺติ.\csromanspacer{} ทิพฺเพน จกฺขุนา รูปํ ปสฺสติ ฯเปฯ.\csromanspacer{} รูปายตนํ จกฺขุวิญฺญาณสฺส สมฺปยุตฺตกานญฺจ ขนฺธานํ อารมฺมณปจฺจเยน ปจฺจโย ฯเปฯ.\csromanspacer{} โผฏฺฐพฺพายตนํ กายวิญฺญาณสฺส สมฺปยุตฺตกานญฺจ ขนฺธานํ ฯเปฯ.\csromanspacer{} อเจตสิกา ขนฺธา อิทฺธิวิธญาณสฺส, เจโตปริยญาณสฺส, ปุพฺเพนิวาสานุสฺสติญาณสฺส, อนาคตํสญาณสฺส, อาวชฺชนาย อารมฺมณปจฺจเยน ปจฺจโย. (3)}

\proseitem{112}{เจตสิโก จ อเจตสิโก จ ธมฺมา เจตสิกสฺส ธมฺมสฺส อารมฺมณปจฺจเยน ปจฺจโย.\csromanspacer{} เจตสิเก ขนฺเธ จ จิตฺตญฺจ อารพฺภ เจตสิกา ขนฺธา อุปฺปชฺชนฺติ. (มูลํ ปุจฺฉิตพฺพํ.) เจตสิเก ขนฺเธ จ จิตฺตญฺจ อารพฺภ จิตฺตํ อุปฺปชฺชติ. (มูลํ ปุจฺฉิตพฺพํ.) เจตสิเก ขนฺเธ จ จิตฺตญฺจ อารพฺภ เจตสิกา ขนฺธา จ จิตฺตญฺจ อุปฺปชฺชนฺติ. (3)}

\csromanheader{2}{\textbf{อธิปติ}}
\proseitem{123}{เจตสิโก ธมฺโม เจตสิกสฺส ธมฺมสฺส อธิปติปจฺจเยน ปจฺจโย.\csromanspacer{} \textbf{อารมฺมณาธิปติ}, สหชาตาธิปติ.\csromanspacer{}\csromanspacer{}\textbf{อารมฺมณธิปติ}\csromandash{}เจตสิเก ขนฺเธ ครุํ กตฺวา เจตสิกา ขนฺธา อุปฺปชฺชนฺติ.\csromanspacer{} \textbf{สหชาตาธิปติ}\csromandash{} เจตสิกาธิปติ สมฺปยุตฺตกานํ ขนฺธานํ อธิปติปจฺจเยน ปจฺจโย. (มูลํ ปุจฺฉิตพฺพํ.) \textbf{อารมฺมณาธิปติ}\csromandash{}เจตสิเก ขนฺเธ ครุํ กตฺวา จิตฺตํ อุปฺปชฺชติ.\csromanspacer{} \textbf{สหชาตาธิปติ}\csromandash{}เจตสิกาธิปติ จิตฺตสฺส จิตฺตสมุฏฺฐานานญฺจ รูปานํ อธิปติปจฺจเยน ปจฺจโย. (มูลํ ปุจฺฉิตพฺพํ.) \textbf{อารมฺมณาธิปติ}\csromandash{}เจตสิเก ขนฺเธ ครุํ กตฺวา เจตสิกา ขนฺธา จ จิตฺตญฺจ อุปฺปชฺชนฺติ.\csromanspacer{} \textbf{สหชาตาธิปติ}\csromandash{}เจตสิกาธิปติ สมฺปยุตฺตกานํ ขนฺธานํ จิตฺตสฺส จ จิตฺตสมุฏฺฐานานญฺจ รูปานํ อธิปติปจฺจเยน ปจฺจโย. (3)}

\proseitem{124}{อเจตสิโก ธมฺโม อเจตสิกสฺส ธมฺมสฺส อธิปติปจฺจเยน ปจฺจโย.\csromanspacer{} \textbf{อารมฺมณาธิปติ}, สหชาตาธิปติ.\csromanspacer{}\csromanspacer{}\textbf{อารมฺมณาธิปติ}\csromandash{}อริยา มคฺคา วุฏฺฐหิตฺวา ฯเปฯ.\csromanspacer{} นิพฺพานํ ครุํ กตฺวา ปจฺจเวกฺขนฺติ.\csromanspacer{} นิพฺพานํ โคตฺรภุสฺส, โวทานสฺส, มคฺคสฺส, ผลสฺส อธิปติปจฺจเยน ปจฺจโย. \csromanfolio{416}จกฺขุํ ฯเปฯ วตฺถุํ อเจตสิเก ขนฺเธ ครุํ กตฺวา อสฺสาเทติ อภินนฺทติ, ตํ ครุํ กตฺวา จิตฺตํ อุปฺปชฺชติ.\csromanspacer{} \textbf{สหชาตาธิปติ}\csromandash{}อเจตสิกาธิปติ จิตฺตสมุฏฺฐานานํ รูปานํ อธิปติปจฺจเยน ปจฺจโย. (1)}

\prose{อเจตสิโก ธมฺโม เจตสิกสฺส ธมฺมสฺส อธิปติปจฺจเยน ปจฺจโย.\csromanspacer{} อารมฺมณาธิปติ, สหชาตาธิปติ.\csromanspacer{}\csromanspacer{}\textbf{อารมฺมณาธิปติ\csromandash{}}อริยา มคฺคา วุฏฺฐหิตฺวา ฯเปฯ.\csromanspacer{} นิพฺพานํ ครุํ กตฺวา ฯเปฯ. (ปฐมคมนสทิสํ.) จกฺขุํ ฯเปฯ วตฺถุํ อเจตสิเก ขนฺเธ ครุํ กตฺวา อสฺสาเทติ อภินนฺทติ, ตํ ครุํ กตฺวา ราโค อุปฺปชฺชติ, ทิฏฺฐิ อุปฺปชฺชติ.\csromanspacer{} \textbf{สหชาตาธิปติ}\csromandash{} อเจตสิกาธิปติ สมฺปยุตฺตกานํ ขนฺธานํ อธิปติปจฺจเยน ปจฺจโย.}

\prose{อเจตสิโก ธมฺโม เจตสิกสฺส จ อเจตสิกสฺส จ ธมฺมสฺส อธิปติปจฺจเยน ปจฺจโย.\csromanspacer{} อารมฺมณาธิปติ, สหชาตาธิปติ.\csromanspacer{}\csromanspacer{}\textbf{อารมฺมณาธิปติ\csromandash{}}อริยา มคฺคา วุฏฺฐหิตฺวา ฯเปฯ.\csromanspacer{} นิพฺพานํ ฯเปฯ. (ปฐมคมนํ.) อเจตสิเก ขนฺเธ ครุํ กตฺวา อสฺสาเทติ อภินนฺทติ, ตํ ครุํ กตฺวา เจตสิกา ขนฺธา จ จิตฺตญฺจ อุปฺปชฺชนฺติ.\csromanspacer{} \textbf{สหชาตาธิปติ}\csromandash{} อเจตสิกาธิปติ สมฺปยุตฺตกานํ ขนฺธานํ จิตฺตสมุฏฺฐานานญฺจ รูปานํ อธิปติปจฺจเยน ปจฺจโย. (3)}

\prose{เจตสิโก จ อเจตสิโก จ ธมฺมา เจตสิกสฺส ธมฺมสฺส อธิปติปจฺจเยน ปจฺจโย.\csromanspacer{} อารมฺมณาธิปติ.\csromanspacer{} ตีณิ. (อารมฺมณาธิปติเยว.)}

\csromanheader{2}{\textbf{อนนฺตราทิ}}
\proseitem{125}{เจตสิโก ธมฺโม เจตสิกสฺส ธมฺมสฺส อนนฺตรปจฺจเยน ปจฺจโย.\csromanspacer{} ปุริมา ปุริมา เจตสิกา ขนฺธา ปจฺฉิมานํ ปจฺฉิมานํ เจตสิกานํ ขนฺธานํ อนนฺตรปจฺจเยน ปจฺจโย. (มูลํ ปุจฺฉิตพฺพํ.) ปุริมา ปุริมา เจตสิกา ขนฺธา ปจฺฉิมสฺส ปจฺฉิมสฺส จิตฺตสฺส อนนฺตรปจฺจเยน ปจฺจโย. (มูลํ ปุจฺฉิตพฺพํ.) ปุริมา ปุริมา เจตสิกา ขนฺธา ปจฺฉิมานํ ปจฺฉิมานํ เจตสิกานํ ขนฺธานํ จิตฺตสฺส จ นนฺตรปจฺจเยน ปจฺจโย. (3)}

\proseitem{126}{อเจตสิโก ธมฺโม อเจตสิกสฺส ธมฺมสฺส อนนฺตรปจฺจเยน ปจฺจโย.\csromanspacer{} ปุริมํ ปุริมํ จิตฺตํ ปจฺฉิมสฺส ปจฺฉิมสฺส จิตฺตสฺส อนนฺตรปจฺจเยน ปจฺจโย.\csromanspacer{} อนุโลมํ โคตฺรภุสฺส ฯเปฯ ผลสมาปตฺติยา อนนฺตรปจฺจเยน ปจฺจโย. (1)}

\csromanheader{2}{57. เจตสิกทุก}
\prose{\csromanfolio{417}อเจตสิโก ธมฺโม เจตสิกสฺส ธมฺมสฺส อนนฺตรปจฺจเยน ปจฺจโย. (เอวํ ตีณิ กาตพฺพานิ, ปฐมคมนสทิสํ.\csromanspacer{} ปูริตฺวา กาตพฺพํ, นินฺนานากรณํ.)}

\prose{เจตสิโก จ อเจตสิโก จ ธมฺมา เจตสิกสฺส ธมฺมสฺส อนนฺตรปจฺจเยน ปจฺจโย.\csromanspacer{} ตีณิ. (อาวชฺชนาปิ วุฏฺฐานมฺปิ นตฺถิ.)}

\prose{สมนนฺตรปจฺจเยน ปจฺจโย.\csromanspacer{} นว.\csromanspacer{} สหชาตปจฺจเยน ปจฺจโย.\csromanspacer{} นว. (ปฏิจฺจวารสทิสํ.) อญฺญมญฺญปจฺจเยน ปจฺจโย.\csromanspacer{} นว. (ปฏิจฺจวารสทิสํ.) นิสฺสยปจฺจเยน ปจฺจโย.\csromanspacer{} นว. (ปจฺจยวารสทิสํ.)}

\csromanheader{2}{\textbf{อุปนิสฺสย}}
\proseitem{127}{เจตสิโก ธมฺโม เจตสิกสฺส ธมฺมสฺส อุปนิสฺสยปจฺจเยน ปจฺจโย.\csromanspacer{} อารมฺมณูปนิสฺสโย, อนนฺตรูปนิสฺสโย, ปกตูปนิสฺสโย ฯเปฯ.\csromanspacer{} ปกตูปนิสฺสโย\csromandash{}เจตสิกา ขนฺธา เจตสิกานํ ขนฺธานํ อุปนิสฺสยปจฺจเยน ปจฺจโย. (มูลํ ปุจฺฉิตพฺพํ.\csromanspacer{} ตีณิ อุปนิสฺสยา.) เจตสิกา ขนฺธา จิตฺตสฺส อุปนิสฺสยปจฺจเยน ปจฺจโย. (มูลํ ปุจฺฉิตพฺพํ.\csromanspacer{} ตีณิ อุปนิสฺสยา.) เจตสิกา ขนฺธา เจตสิกานํ ขนฺธานํ จิตฺตสฺส จ อุปนิสฺสยปจฺจเยน ปจฺจโย. (3)}

\proseitem{128}{อเจตสิโก ธมฺโม อเจตสิกสฺส ธมฺมสฺส อุปนิสฺสยปจฺจเยน ปจฺจโย.\csromanspacer{} อารมฺมณูปนิสฺสโย, อนนฺตรูปนิสฺสโย, ปกตูปนิสฺสโย ฯเปฯ.\csromanspacer{} ปกตูปนิสฺสโย\csromandash{}อุตุํ.\csromanspacer{} โภชนํ.\csromanspacer{} เสนาสนํ จิตฺตํ อุปนิสฺสาย ทานํ เทติ ฯเปฯ สํฆํ ภินฺทติ.\csromanspacer{} อุตุํ.\csromanspacer{} โภชนํ.\csromanspacer{} เสนาสนํ จิตฺตํ สทฺธาย ฯเปฯ ปญฺญาย.\csromanspacer{} ราคสฺส ฯเปฯ ปตฺถนาย กายิกสฺส สุขสฺส, กายิกสฺส ทุกฺขสฺส, มคฺคสฺส, ผลสมาปตฺติยา อุปนิสฺสยปจฺจเยน ปจฺจโย. (1)}

\prose{อเจตสิโก ธมฺโม เจตสิกสฺส ธมฺมสฺส อุปนิสฺสยปจฺจเยน ปจฺจโย.\csromanspacer{} อารมฺมณูปนิสฺสโย, อนนฺตรูปนิสฺสโย, ปกตูปนิสฺสโย ฯเปฯ.\csromanspacer{} ปกตูปนิสฺสโย\csromandash{}อุตุํ.\csromanspacer{} โภชนํ.\csromanspacer{} เสนาสนํ จิตฺตํ อุปนิสฺสาย ทานํ เทติ ฯเปฯ. (ตีณิ.\csromanspacer{} ปฐมคมนสทิสํ.\csromanspacer{} นินฺนานากรณํ.)}

\prose{เจตสิโก จ อเจตสิโก จ ธมฺมา เจตสิกสฺส ธมฺมสฺส อุปนิสฺสยปจฺจเยน ปจฺจโย.\csromanspacer{} ตีณิ.}

\csromanheader{2}{\textbf{ปุเรชาต}}
\proseitem{129}{\csromanfolio{418}อเจตสิโก ธมฺโม อเจตสิกสฺส ธมฺมสฺส ปุเรชาตปจฺจเยน ปจฺจโย.\csromanspacer{} \textbf{อารมฺมณปุเรชาตํ}, วตฺถุปุเรชาตํ.\csromanspacer{}\csromanspacer{}\textbf{อารมฺมณปุเรชาตํ}\csromandash{} จกฺขุํ ฯเปฯ วตฺถุํ อนิจฺจโต ฯเปฯ วิปสฺสติ, อสฺสาเทติ อภินนฺทติ, ตํ อารพฺภ จิตฺตํ อุปฺปชฺชติ.\csromanspacer{} ทิพฺเพน จกฺขุนา รูปํ ปสฺสติ.\csromanspacer{} ทิพฺพาย โสตธาตุยา สทฺทํ สุณาติ.\csromanspacer{} รูปายตนํ จกฺขุวิญฺญาณสฺส ฯเปฯ.\csromanspacer{} โผฏฺฐพฺพายตนํ กายวิญฺญาณสฺส ฯเปฯ.\csromanspacer{} \textbf{วตฺถุปุเรชาตํ}\csromandash{}จกฺขายตนํ จกฺขุวิญฺญาณสฺส ฯเปฯ กายายตนํ กายวิญฺญาณสฺส ฯเปฯ.\csromanspacer{} วตฺถุ จิตฺตสฺส ปุเรชาตปจฺจเยน ปจฺจโย. (1)}

\prose{อเจตสิโก ธมฺโม เจตสิกสฺส ธมฺมสฺส ปุเรชาตปจฺจเยน ปจฺจโย.\csromanspacer{} \textbf{อารมฺมณปุเรชาตํ}, วตฺถุปุเรชาตํ.\csromanspacer{}\csromanspacer{}\textbf{อารมฺมณปุเรชาตํ}\csromandash{}จกฺขุํ ฯเปฯ.\csromanspacer{} วตฺถุํ อนิจฺจโต ฯเปฯ โทมนสฺสํ อุปฺปชฺชติ.\csromanspacer{} ทิพฺเพน จกฺขุนา รูปํ ปสฺสติ.\csromanspacer{} ทิพฺพาย โสตธาตุยา สทฺทํ สุณาติ.\csromanspacer{} รูปายตนํ จกฺขุวิญฺญาณสหคตานํ ขนฺธานํ ฯเปฯ.\csromanspacer{} โผฏฺฐพฺพายตนํ กายวิญฺญาณสหคตานํ ขนฺธานํ ฯเปฯ.\csromanspacer{} \textbf{วตฺถุปุเรชาตํ}\csromandash{} จกฺขายตนํ จกฺขุวิญฺญาณสหคตานํ ขนฺธานํ ฯเปฯ กายายตนํ ฯเปฯ.\csromanspacer{} วตฺถุ เจตสิกานํ ขนฺธานํ ปุเรชาตปจฺจเยน ปจฺจโย. (2)}

\prose{อเจตสิโก ธมฺโม เจตสิกสฺส จ อเจตสิกสฺส จ ธมฺมสฺส ปุเรชาตปจฺจเยน ปจฺจโย.\csromanspacer{} \textbf{อารมฺมณปุเรชาตํ}, วตฺถุปุเรชาตํ.\csromanspacer{}\csromanspacer{}\textbf{อารมฺมณปุเรชาตํ}\csromandash{}จกฺขุํ ฯเปฯ วตฺถุํ อนิจฺจโต ฯเปฯ วิปสฺสติ, อสฺสาเทติ อภินนฺทติ, ตํ อารพฺภ จิตฺตญฺจ สมฺปยุตฺตกา จ ขนฺธา อุปฺปชฺชนฺติ.\csromanspacer{} ทิพฺเพน จกฺขุนา รูปํ ปสฺสติ.\csromanspacer{} ทิพฺพาย โสตธาตุยา สทฺทํ สุณาติ.\csromanspacer{} รูปายตนํ จกฺขุวิญฺญาณสฺส สมฺปยุตฺตกานญฺจ ขนฺธานํ ฯเปฯ.\csromanspacer{} โผฏฺฐพฺพายตนํ ฯเปฯ.\csromanspacer{} \textbf{วตฺถุปุเรชาตํ}\csromandash{}จกฺขายตนํ จกฺขุวิญฺญาณสฺส สมฺปยุตฺตกานญฺจ ขนฺธานํ ฯเปฯ กายายตนํ ฯเปฯ.\csromanspacer{} วตฺถุ จิตฺตสฺส สมฺปยุตฺตกานญฺจ ขนฺธานํ ปุเรชาตปจฺจเยน ปจฺจโย. (3)}

\csromanheader{2}{\textbf{ปจฺฉาชาตาเสวน}}
\proseitem{130}{เจตสิโก ธมฺโม อเจตสิกสฺส ธมฺมสฺส ปจฺฉาชาตปจฺจเยน ปจฺจโย. (สํขิตฺตํ.)}

\csromanheader{2}{57. เจตสิกทุก}
\prose{\csromanfolio{419}อเจตสิโก ธมฺโม อเจตสิกสฺส ธมฺมสฺส ปจฺฉาชาตปจฺจเยน ปจฺจโย. (สํขิตฺตํ.)}

\prose{เจตสิโก จ อเจตสิโก จ ธมฺมา อเจตสิกสฺส ธมฺมสฺส ปจฺฉาชาตปจฺจเยน ปจฺจโย. (สํขิตฺตํ.)}

\prose{เจตสิโก ธมฺโม เจตสิกสฺส ธมฺมสฺส อาเสวนปจฺจเยน ปจฺจโย. (สํขิตฺตํ.)}

\csromanheader{2}{\textbf{กมฺม}}
\proseitem{131}{เจตสิโก ธมฺโม เจตสิกสฺส ธมฺมสฺส กมฺมปจฺจเยน ปจฺจโย.\csromanspacer{} \textbf{สหชาตา}, นานากฺขณิกา.\csromanspacer{}\csromanspacer{}\textbf{สหชาตา} เจตสิกา เจตนา สมฺปยุตฺตกานํ ขนฺธานํ กมฺมปจฺจเยน ปจฺจโย.\csromanspacer{} ปฏิสนฺธิกฺขเณ ฯเปฯ.\csromanspacer{} \textbf{นานากฺขณิกา} เจตสิกา เจตนา วิปากานํ ขนฺธานํ กมฺมปจฺจเยน ปจฺจโย. (1)}

\prose{เจตสิโก ธมฺโม อเจตสิกสฺส ธมฺมสฺส กมฺมปจฺจเยน ปจฺจโย.\csromanspacer{} \textbf{สหชาตา}, นานากฺขณิกา.\csromanspacer{}\csromanspacer{}\textbf{สหชาตา} เจตสิกา เจตนา จิตฺตสฺส จิตฺตสมุฏฺฐานานญฺจ รูปานํ กมฺมปจฺจเยน ปจฺจโย.\csromanspacer{} ปฏิสนฺธิกฺขเณ ฯเปฯ.\csromanspacer{} \textbf{นานากฺขณิกา} เจตสิกา เจตนา วิปากสฺส จิตฺตสฺส กฏตฺตา จ รูปานํ กมฺมปจฺจเยน ปจฺจโย. (2)}

\prose{เจตสิโก ธมฺโม เจตสิกสฺส จ อเจตสิกสฺส จ ธมฺมสฺส กมฺมปจฺจเยน ปจฺจโย.\csromanspacer{} \textbf{สหชาตา}, นานากฺขณิกา.\csromanspacer{}\csromanspacer{}\textbf{สหชาตา} เจตสิกา เจตนา สมฺปยุตฺตกานํ ขนฺธานํ จิตฺตสฺส จ จิตฺตสมุฏฺฐานญฺจ รูปานํ กมฺมปจฺจเยน ปจฺจโย.\csromanspacer{} ปฏิสนฺธิกฺขเณ ฯเปฯ.\csromanspacer{} \textbf{นานากฺขณิกา} เจตสิกา เจตนา วิปากานํ ขนฺธานํ จิตฺตสฺส กฏตฺตา จ รูปานํ กมฺมปจฺจเยน ปจฺจโย. (3)}

\csromanheader{2}{\textbf{วิปากาทิ}}
\proseitem{132}{เจตสิโก ธมฺโม เจตสิกสฺส ธมฺมสฺส วิปากปจฺจเยน ปจฺจโย.\csromanspacer{} นว.\csromanspacer{} อาหารปจฺจเยน ปจฺจโย.\csromanspacer{} นว.\csromanspacer{} อินฺทฺริยปจฺจเยน ปจฺจโย.\csromanspacer{} นว.\csromanspacer{} ฌานปจฺจเยน ปจฺจโย.\csromanspacer{} ตีณิ.\csromanspacer{} มคฺคปจฺจเยน ปจฺจโย.\csromanspacer{} ตีณิ.\csromanspacer{} สมฺปยุตฺตปจฺจเยน ปจฺจโย.\csromanspacer{} ปญฺจ.}

\csromanheader{2}{\textbf{วิปฺปยุตฺต}}
\proseitem{133}{\csromanfolio{420}เจตสิโก ธมฺโม อเจตสิกสฺส ธมฺมสฺส วิปฺปยุตฺตปจฺจเยน ปจฺจโย.\csromanspacer{} \textbf{สหชาตํ}, ปจฺฉาชาตํ. (สํขิตฺตํ.) (1)}

\prose{อเจตสิโก ธมฺโม อเจตสิกสฺส ธมฺมสฺส วิปฺปยุตฺตปจฺจเยน ปจฺจโย.\csromanspacer{} \textbf{สหชาตํ}, ปุเรชาตํ, ปจฺฉาชาตํ.\csromanspacer{}\csromanspacer{}\textbf{สหชาตํ} จิตฺตํ จิตฺตสมุฏฺฐานานํ รูปานํ วิปฺปยุตฺตปจฺจเยน ปจฺจโย.\csromanspacer{} ปฏิสนฺธิกฺขเณ จิตฺตํ กฏตฺตารูปานํ วิปฺปยุตฺตปจฺจเยน ปจฺจโย.\csromanspacer{} จิตฺตํ วตฺถุสฺส วิปฺปยุตฺตปจฺจเยน ปจฺจโย.\csromanspacer{} วตฺถุ จิตฺตสฺส วิปฺปยุตฺตปจฺจเยน ปจฺจโย.\csromanspacer{} \textbf{ปุเรชาตํ} จกฺขายตนํ จกฺขุวิญฺญาณสฺส วิปฺปยุตฺตปจฺจเยน ปจฺจโย ฯเปฯ กายายตนํ กายวิญฺญาณสฺส ฯเปฯ.\csromanspacer{} วตฺถุ จิตฺตสฺส วิปฺปยุตฺตปจฺจเยน ปจฺจโย.\csromanspacer{} \textbf{ปจฺฉาชาตํ} จิตฺตํ ปุเรชาตสฺส อิมสฺส กายสฺส วิปฺปยุตฺตปจฺจเยน ปจฺจโย. (1)}

\prose{อเจตสิโก ธมฺโม เจตสิกสฺส ธมฺมสฺส วิปฺปยุตฺตปจฺจเยน ปจฺจโย.\csromanspacer{} \textbf{สหชาตํ}, ปุเรชาตํ.\csromanspacer{}\csromanspacer{}\textbf{สหชาตํ} ปฏิสนฺธิกฺขเณ วตฺถุ เจตสิกานํ ขนฺธานํ วิปฺปยุตฺตปจฺจเยน ปจฺจโย.\csromanspacer{} \textbf{ปุเรชาตํ} จกฺขายตนํ จกฺขุวิญฺญาณสหคตานํ ขนฺธานํ ฯเปฯ กายายตนํ กายวิญฺญาณสหคตานํ ขนฺธานํ ฯเปฯ.\csromanspacer{} วตฺถุ เจตสิกานํ ขนฺธานํ วิปฺปยุตฺตปจฺจเยน ปจฺจโย. (2)}

\prose{อเจตสิโก ธมฺโม เจตสิกสฺส จ อเจตสิกสฺส จ ธมฺมสฺส วิปฺปยุตฺตปจฺจเยน ปจฺจโย.\csromanspacer{} \textbf{สหชาตํ}, ปุเรชาตํ.\csromanspacer{}\csromanspacer{}\textbf{สหชาตํ} ปฏิสนฺธิกฺขเณ วตฺถุ เจตสิกานํ ขนฺธานํ จิตฺตสฺส จ วิปฺปยุตฺตปจฺจเยน ปจฺจโย.\csromanspacer{} \textbf{ปุเรชาตํ} จกฺขายตนํ จกฺขุวิญฺญาณสฺส สมฺปยุตฺตกานญฺจ ขนฺธานํ ฯเปฯ กายายตนํ กายวิญฺญาณสฺส สมฺปยุตฺตกานญฺจ ขนฺธานํ ฯเปฯ.\csromanspacer{} วตฺถุ จิตฺตสฺส สมฺปยุตฺตกานญฺจ ขนฺธานํ วิปฺปยุตฺตปจฺจเยน ปจฺจโย. (3)}

\prose{เจตสิโก จ อเจตสิโก จ ธมฺมา อเจตสิกสฺส ธมฺมสฺส วิปฺปยุตฺตปจฺจเยน ปจฺจโย.\csromanspacer{} \textbf{สหชาตํ}, ปจฺฉาชาตํ. (สํขิตฺตํ.)}

\csromanheader{2}{\textbf{อตฺถิ}}
\proseitem{134}{เจตสิโก ธมฺโม เจตสิกสฺส ธมฺมสฺส อตฺถิปจฺจเยน ปจฺจโย.\csromanspacer{} เอกํ. (ปฏิจฺจวารสทิสํ.)}

\prose{เจตสิโก ธมฺโม อเจตสิกสฺส ธมฺมสฺส อตฺถิปจฺจเยน ปจฺจโย.\csromanspacer{} \textbf{สหชาตํ}, ปจฺฉาชาตํ. (สํขิตฺตํ.) (2)}

\csromanheader{2}{57. เจตสิกทุก}
\prose{\csromanfolio{421}เจตสิโก ธมฺโม เจตสิกสฺส จ อเจตสิกสฺส จ ธมฺมสฺส อตฺถิปจฺจเยน ปจฺจโย.\csromanspacer{} เอกํ. (ปฏิจฺจวารสทิสํ.) (3)}

\proseitem{135}{อเจตสิโก ธมฺโม อเจตสิกสฺส ธมฺมสฺส อตฺถิปจฺจเยน ปจฺจโย.\csromanspacer{} \textbf{สหชาตํ}, ปุเรชาตํ, ปจฺฉาชาตํ, อาหารํ, อินฺทฺริยํ. (สํขิตฺตํ.) (1)}

\prose{อเจตสิโก ธมฺโม เจตสิกสฺส ธมฺมสฺส อตฺถิปจฺจเยน ปจฺจโย.\csromanspacer{} \textbf{สหชาตํ}, ปุเรชาตํ.\csromanspacer{}\csromanspacer{}\textbf{สหชาตํ} จิตฺตํ เจตสิกานํ ขนฺธานํ อตฺถิปจฺจเยน ปจฺจโย.\csromanspacer{} ปฏิสนฺธิกฺขเณ จิตฺตํ ฯเปฯ.\csromanspacer{} ปฏิสนฺธิกฺขเณ วตฺถุ เจตสิกา ขนฺธานํ อตฺถิปจฺจเยน ปจฺจโย.\csromanspacer{} \textbf{ปุเรชาตํ} จกฺขุํ ฯเปฯ วตฺถุํ อนิจฺจโต ฯเปฯ. (ปุเรชาตสทิสํ, นินฺนานากรณํ.) (2)}

\prose{อเจตสิโก ธมฺโม เจตสิกสฺส จ อเจตสิกสฺส จ ธมฺมสฺส อตฺถิปจฺจเยน ปจฺจโย.\csromanspacer{} \textbf{สหชาตํ}, ปุเรชาตํ.\csromanspacer{}\csromanspacer{}\textbf{สหชาตํ} จิตฺตํ สมฺปยุตฺตกานํ ขนฺธานํ จิตฺตสมุฏฺฐานานญฺจ รูปานํ อตฺถิปจฺจเยน ปจฺจโย.\csromanspacer{} ปฏิสนฺธิกฺขเณ จิตฺตํ ฯเปฯ.\csromanspacer{} ปฏิสนฺธิกฺขเณ วตฺถุ เจตสิกานํ ขนฺธานํ จิตฺตสฺส จ อตฺถิปจฺจเยน ปจฺจโย.\csromanspacer{} \textbf{ปุเรชาตํ} จกฺขุํ ฯเปฯ วตฺถุํ อนิจฺจโต ฯเปฯ. (ปุเรชาตสทิสํ นินฺนานํ.) (3)}

\proseitem{136}{เจตสิโก จ อเจตสิโก จ ธมฺมา เจตสิกสฺส ธมฺมสฺส อตฺถิปจฺจเยน ปจฺจโย.\csromanspacer{} \textbf{สหชาตํ}, ปุเรชาตํ.\csromanspacer{}\csromanspacer{}\textbf{สหชาโต} จกฺขุวิญฺญาณสหคโต เอโก ขนฺโธ จ จกฺขายตนญฺจ ทฺวินฺนํ ขนฺธานํ ฯเปฯ.\csromanspacer{} กายวิญฺญาณสหคโต ฯเปฯ.\csromanspacer{} เจตสิโก เอโก ขนฺโธ จ วตฺถุ จ ทฺวินฺนํ ขนฺธานํ อตฺถิปจฺจเยน ปจฺจโย.\csromanspacer{} เทฺว ขนฺธา จ ฯเปฯ.\csromanspacer{} \textbf{สหชาโต} เจตสิโก เอโก ขนฺโธ จ จิตฺตญฺจ ทฺวินฺนํ ขนฺธานํ อตฺถิปจฺจเยน ปจฺจโย.\csromanspacer{} เทฺว ขนฺธา จ ฯเปฯ. (ปฏิสนฺธิกฺขเณ เทฺวปิ กาตพฺพา.) (1)}

\prose{เจตสิโก จ อเจตสิโก จ ธมฺมา อเจตสิกสฺส ธมฺมสฺส อตฺถิปจฺจเยน ปจฺจโย.\csromanspacer{} \textbf{สหชาตํ}, ปุเรชาตํ, ปจฺฉาชาตํ, อาหารํ, อินฺทฺริยํ.\csromanspacer{}\csromanspacer{}\textbf{สหชาตา} จกฺขุวิญฺญาณสหคตา ขนฺธา จ จกฺขายตนญฺจ จกฺขุวิญฺญาณสฺส ฯเปฯ.\csromanspacer{} กายวิญฺญาณสฺส ฯเปฯ.\csromanspacer{} เจตสิกา ขนฺธา จ จิตฺตญฺจ จิตฺตสมุฏฺฐานานํ รูปานํ อตฺถิปจฺจเยน ปจฺจโย.\csromanspacer{} เจตสิกา ขนฺธา จ จิตฺตญฺจ มหาภูตา จ จิตฺตสมุฏฺฐานานํ รูปานํ อตฺถิปจฺจเยน ปจฺจโย. \csromanfolio{422}\textbf{สหชาตา} เจตสิกา ขนฺธา จ วตฺถุ จ จิตฺตสฺส อตฺถิปจฺจเยน ปจฺจโย. (ปฏิสนฺธิกฺขเณ ตีณิปิ กาตพฺพา.) \textbf{ปจฺฉาชาตา} เจตสิกา ขนฺธา จ จิตฺตญฺจ ปุเรชาตสฺส อิมสฺส กายสฺส อตฺถิปจฺจเยน ปจฺจโย.\csromanspacer{} \textbf{ปจฺฉาชาตา} เจตสิกา ขนฺธา จ จิตฺตญฺจ กพฬีกาโร อาหาโร จ อิมสฺส กายสฺส อตฺถิปจฺจเยน ปจฺจโย.\csromanspacer{} \textbf{ปจฺฉาชาตา} เจตสิกา ขนฺธา จ จิตฺตญฺจ รูปชีวิตินฺทฺริยญฺจ กฏตฺตารูปานํ อตฺถิปจฺจเยน ปจฺจโย. (2)}

\prose{เจตสิโก จ อเจตสิโก จ ธมฺมา เจตสิกสฺส จ อเจตสิกสฺส จ ธมฺมสฺส อตฺถิปจฺจเยน ปจฺจโย.\csromanspacer{} สหชาตํ, ปุเรชาตํ.\csromanspacer{}\csromanspacer{}\textbf{สหชาโต} จกฺขุวิญฺญาณสหคโต เอโก ขนฺโธ จ จกฺขายตนญฺจ ทฺวินฺนํ ขนฺธานํ จกฺขุวิญฺญาณสฺส จ อตฺถิปจฺจเยน ปจฺจโย ฯเปฯ.\csromanspacer{} กายวิญฺญาณสหคโต เอโก ขนฺโธ จ กายายตนญฺจ ทฺวินฺนํ ขนฺธานํ กายวิญฺญาณสฺส จ อตฺถิปจฺจเยน ปจฺจโย.\csromanspacer{} เทฺว ขนฺธา จ ฯเปฯ.\csromanspacer{} \textbf{สหชาโต} เจตสิโก เอโก ขนฺโธ จ วตฺถุ จ ทฺวินฺนํ ขนฺธานํ จิตฺตสฺส จ อตฺถิปจฺจเยน ปจฺจโย.\csromanspacer{} เทฺว ขนฺธา จ ฯเปฯ. (ปฏิสนฺธิยา เทฺว กาตพฺพา.) (3)}

\csromanheader{2}{1. ปจฺจยานุโลม 2. สงฺขฺยาวาร}
\csromanheader{2}{\textbf{สุทฺธ}}
\proseitem{137}{เหตุยา ตีณิ, อารมฺมเณ นว, อธิปติยา นว, อนนฺตเร นว, สมนนฺตเร นว, สหชาเต นว, อญฺญมญฺเญ นว, นิสฺสเย นว, อุปนิสฺสเย อนว, ปุเรชาเต ตีณิ, ปจฺฉาชาเต ตีณิ, อาเสวเน นว, กมฺเม ตีณิ, วิปาเก นว, อาหาเร นว, อินฺทฺริเย นว, ฌาเน ตีณิ, มคฺเค ตีณิ, สมฺปยุตฺเต ปญฺจ, วิปฺปยุตฺเต ปญฺจ, อตฺถิยา นว, นตฺถิยา นว, วิคเต นว, อวิคเต นว.}

\vspace{\tipitakaparskip}
\csromancenter{อนุโลมํ.}
\csromanheader{2}{2. ปจฺจนียุทฺธาร}
\proseitem{138}{เจตสิโก ธมฺโม เจตสิกสฺส ธมฺมสฺส อารมฺมณปจฺจเยน ปจฺจโย.\csromanspacer{} สหชาตปจฺจเยน ปจฺจโย.\csromanspacer{} อุปนิสฺสยปจฺจเยน ปจฺจโย.\csromanspacer{} กมฺมปจฺจเยน ปจฺจโย. (1)}

\csromanheader{2}{57. เจตสิกทุก}
\prose{\csromanfolio{423}เจตสิโก ธมฺโม อเจตสิกสฺส ธมฺมสฺส อารมฺมณปจฺจเยน ปจฺจโย.\csromanspacer{} สหชาตปจฺจเยน ปจฺจโย.\csromanspacer{} อุปนิสฺสยปจฺจเยน ปจฺจโย.\csromanspacer{} ปจฺฉาชาตปจฺจเยน ปจฺจโย.\csromanspacer{} กมฺมปจฺจเยน ปจฺจโย. (2)}

\prose{เจตสิโก ธมฺโม เจตสิกสฺส จ อเจตสิกสฺส จ ธมฺมสฺส อารมฺมณปจฺจเยน ปจฺจโย.\csromanspacer{} สหชาตปจฺจเยน ปจฺจโย.\csromanspacer{} อุปนิสฺสยปจฺจเยน ปจฺจโย.\csromanspacer{} กมฺมปจฺจเยน ปจฺจโย. (3)}

\proseitem{139}{อเจตสิโก ธมฺโม อเจตสิกสฺส ธมฺมสฺส อารมฺมณปจฺจเยน ปจฺจโย.\csromanspacer{} สหชาตปจฺจเยน ปจฺจโย.\csromanspacer{} อุปนิสฺสยปจฺจเยน ปจฺจโย.\csromanspacer{} ปุเรชาตปจฺจเยน ปจฺจโย.\csromanspacer{} ปจฺฉาชาตปจฺจเยน ปจฺจโย.\csromanspacer{} อาหารปจฺจเยน ปจฺจโย.\csromanspacer{} อินฺทฺริยปจฺจเยน ปจฺจโย. (1)}

\prose{อเจตสิโก ธมฺโม เจตสิกสฺส ธมฺมสฺส อารมฺมณปจฺจเยน ปจฺจโย.\csromanspacer{} สหชาตปจฺจเยน ปจฺจโย.\csromanspacer{} อุปนิสฺสยปจฺจเยน ปจฺจโย.\csromanspacer{} ปุเรชาตปจฺจเยน ปจฺจโย. (2)}

\prose{อเจตสิโก ธมฺโม เจตสิกสฺส จ อเจตสิกสฺส จ ธมฺมสฺส อารมฺมณปจฺจเยน ปจฺจโย.\csromanspacer{} สหชาตปจฺจเยน ปจฺจโย.\csromanspacer{} อุปนิสฺสยปจฺจเยน ปจฺจโย.\csromanspacer{} ปุเรชาตปจฺจเยน ปจฺจโย. (3)}

\proseitem{140}{เจตสิโก จ อเจตสิโก จ ธมฺมา เจตสิกสฺส ธมฺมสฺส อารมฺมณปจฺจเยน ปจฺจโย.\csromanspacer{} สหชาตปจฺจเยน ปจฺจโย.\csromanspacer{} อุปนิสฺสยปจฺจเยน ปจฺจโย. (1)}

\prose{เจตสิโก จ อเจตสิโก จ ธมฺมา อเจตสิกสฺส ธมฺมสฺส อารมฺมณปจฺจเยน ปจฺจโย.\csromanspacer{} สหชาตปจฺจเยน ปจฺจโย.\csromanspacer{} อุปนิสฺสยปจฺจเยน ปจฺจโย.\csromanspacer{} ปจฺฉาชาตปจฺจเยน ปจฺจโย. (2)}

\prose{เจตสิโก จ อเจตสิโก จ ธมฺมา เจตสิกสฺส จ อเจตสิกสฺส จ ธมฺมสฺส อารมฺมณปจฺจเยน ปจฺจโย.\csromanspacer{} สหชาตปจฺจเยน ปจฺจโย.\csromanspacer{} อุปนิสฺสยปจฺจเยน ปจฺจโย. (3)}

\csromanheader{2}{2. ปจฺจยปจฺจนีย 2. สงฺขฺยาวาร}
\vspace{\tipitakaparskip}
\csromancenter{นเหตุยา นว, น-อารมฺมเณ นว, (สพฺพตฺถ นว,) โน-อวิคเต นว.}
\csromanheader{2}{3. ปจฺจยานุโลมปจฺจนีย}
\proseitem{142}{\csromanfolio{424}\textbf{เหตุ}ปจฺจยา น-อารมฺมเณ ตีณิ, น-อธิปติยา ตีณิ, น-อนนฺตเร ตีณิ, นสมนนฺตเร ตีณิ, น-อญฺญมญฺเญ เอกํ, น-อุปนิสฺสเย ตีณิ, (สพฺพตฺถ ตีณิ,) นมคฺเค ตีณิ, นสมฺปยุตฺเต เอกํ, นวิปฺปยุตฺเต ตีณิ, โนนตฺถิยา ตีณิ, โนวิคเต ตีณิ.}

\csromanheader{2}{4. ปจฺจยปจฺจนียานุโลม}
\proseitem{143}{นเหตุปจฺจยา อารมฺมเณ นว, อธิปติยา นว, (อนุโลมมาติกา กาตพฺพา.) ฯเปฯ อวิคเต นว.}

\nitthitam{เจตสิกทุกํ นิฏฺฐิตํ.}
\csromanmatikaanchor{mtk.57}
\csromantocmarkref{subsection}{58. จิตฺตสมฺปยุตฺตทุก}{mtk.57}
\bookmark[dest={mtk.57},level=2]{58. จิตฺตสมฺปยุตฺตทุก}
\csromanheader{2}{58. จิตฺตสมฺปยุตฺตทุก 1. ปฏิจฺจวาร}
\csromanheader{2}{1. ปจฺจยานุโลม 1. วิภงฺควาร}
\csromanheader{2}{\textbf{เหตุ}}
\proseitem{144}{จิตฺตสมฺปยุตฺตํ ธมฺมํ ปฏิจฺจ จิตฺตสมฺปยุตฺโต ธมฺโม อุปฺปชฺชติ เหตุปจฺจยา.\csromanspacer{} จิตฺตสมฺปยุตฺตํ เอกํ ขนฺธํ ปฏิจฺจ เทฺว ขนฺธา, เทฺว ขนฺเธ ปฏิจฺจ เอโก ขนฺโธ ปฏิสนฺธิกฺขเณ ฯเปฯ. (1)}

\prose{จิตฺตสมฺปยุตฺตํ ธมฺมํ ปฏิจฺจ จิตฺตวิปฺปยุตฺโต ธมฺโม อุปฺปชฺชติ เหตุปจฺจยา.\csromanspacer{} จิตฺตสมฺปยุตฺเต ขนฺเธ จิตฺตสมุฏฺฐานํ รูปํ.\csromanspacer{} ปฏิสนฺธิกฺขเณ ฯเปฯ. (2)}

\prose{จิตฺตสมฺปยุตฺตํ ธมฺมํ ปฏิจฺจ จิตฺตสมฺปยุตฺโต จ จิตฺตวิปฺปยุตฺโต จ ธมฺมา อุปฺปชฺชนฺติ เหตุปจฺจยา.\csromanspacer{} จิตฺตสมฺปยุตฺตํ เอกํ ขนฺธํ ปฏิจฺจ เทฺว ขนฺธา จิตฺตสมุฏฺฐานญฺจ รูปํ.\csromanspacer{} เทฺว ขนฺเธ ฯเปฯ.\csromanspacer{} ปฏิสนฺธิกฺขเณ ฯเปฯ. (3)}

\proseitem{145}{จิตฺตวิปฺปยุตฺตํ ธมฺมํ ปฏิจฺจ จิตฺตวิปฺปยุตฺโต ธมฺโม อุปฺปชฺชติ เหตุปจฺจยา.\csromanspacer{} เอกํ มหาภูตํ ฯเปฯ มหาภูเต ปฏิจฺจ จิตฺตสมุฏฺฐานํ รูปํ ฯเปฯ.\csromanspacer{} ปฏิสนฺธิกฺขเณ ฯเปฯ.\csromanspacer{} เอกํ มหาภูตํ ฯเปฯ มหาภูเต ปฏิจฺจ กฏตฺตารูปํ, อุปาทารูปํ. (1)}

\csromanheader{2}{58. จิตฺตสมฺปยุตฺตทุก}
\prose{\csromanfolio{425}จิตฺตวิปฺปยุตฺตํ ธมฺมํ ปฏิจฺจ จิตฺตสมฺปยุตฺโต ธมฺโม อุปฺปชฺชติ เหตุปจฺจยา.\csromanspacer{} ปฏิสนฺธิกฺขเณ วตฺถุํ ปฏิจฺจ จิตฺตสมฺปยุตฺตกา ขนฺธา.}

\prose{จิตฺตวิปฺปยุตฺตํ ธมฺมํ ปฏิจฺจ จิตฺตสมฺปยุตฺโต จ จิตฺตวิปฺปยุตฺโต จ ธมฺมา อุปฺปชฺชนฺติ เหตุปจฺจยา.\csromanspacer{} ปฏิสนฺธิกฺขเณ วตฺถุํ ปฏิจฺจ จิตฺตสมฺปยุตฺตกา ขนฺธา.\csromanspacer{} มหาภูเต ปฏิจฺจ กฏตฺตารูปํ. (3)}

\proseitem{146}{จิตฺตสมฺปยุตฺตญฺจ จิตฺตวิปฺปยุตฺตญฺจ ธมฺมํ ปฏิจฺจ จิตฺตสมฺปยุตฺโต ธมฺโม อุปฺปชฺชติ เหตุปจฺจยา.\csromanspacer{} ปฏิสนฺธิกฺขเณ จิตฺตสมฺปยุตฺตํ เอกํ ขนฺธญฺจ วตฺถุญฺจ ปฏิจฺจ เทฺว ขนฺธา.\csromanspacer{} เทฺว ขนฺเธ จ ฯเปฯ. (1)}

\prose{จิตฺตสมฺปยุตฺตญฺจ จิตฺตวิปฺปยุตฺตญฺจ ธมฺมํ ปฏิจฺจ จิตฺตวิปฺปยุตฺโต ธมฺโม อุปฺปชฺชติ เหตุปจฺจยา.\csromanspacer{} จิตฺตสมฺปยุตฺเต ขนฺเธ จ มหาภูเต จ ปฏิจฺจ จิตฺตสมุฏฺฐานํ รูปํ.\csromanspacer{} ปฏิสนฺธิกฺขเณ จิตฺตสมฺปยุตฺเต ขนฺเธ จ มหาภูเต จ ปฏิจฺจ กฏตฺตารูปํ. (2)}

\prose{จิตฺตสมฺปยุตฺตญฺจ จิตฺตวิปฺปยุตฺตญฺจ ธมฺมํ ปฏิจฺจ จิตฺตสมฺปยุตฺโต จ จิตฺตวิปฺปยุตฺโต จ ธมฺมา อุปฺปชฺชนฺติ เหตุปจฺจยา.\csromanspacer{} ปฏิสนฺธิกฺขเณ จิตฺตสมฺปยุตฺตํ เอกํ ขนฺธญฺจ วตฺถุญฺจ ปฏิจฺจ เทฺว ขนฺธา.\csromanspacer{} เทฺว ขนฺเธ จ ฯเปฯ.\csromanspacer{} จิตฺตสมฺปยุตฺเต ขนฺเธ จ มหาภูเต จ ปฏิจฺจ กฏตฺตารูปํ. (3)}

\csromanheader{2}{\textbf{อารมฺมณ}}
\proseitem{147}{จิตฺตสมฺปยุตฺตํ ธมฺมํ ปฏิจฺจ จิตฺตสมฺปยุตฺโต ธมฺโม อุปฺปชฺชติ อารมฺมณปจฺจยา.\csromanspacer{} จิตฺตสมฺปยุตฺตํ เอกํ ขนฺธํ ปฏิจฺจ เทฺว ขนฺธา.\csromanspacer{} เทฺว ขนฺเธ ฯเปฯ.\csromanspacer{} ปฏิสนฺธิกฺขเณ ฯเปฯ. (1)}

\prose{จิตฺตวิปฺปยุตฺตํ ธมฺมํ ปฏิจฺจ จิตฺตสมฺปยุตฺโต ธมฺโม อุปฺปชฺชติ อารมฺมณปจฺจยา.\csromanspacer{} ปฏิสนฺธิกฺขเณ วตฺถุํ ปฏิจฺจ จิตฺตสมฺปยุตฺตกา ขนฺธา. (1)}

\prose{จิตฺตสมฺปยุตฺตญฺจ จิตฺตวิปฺปยุตฺตญฺจ ธมฺมํ ปฏิจฺจ จิตฺตสมฺปยุตฺโต ธมฺโม อุปฺปชฺชติ อารมฺมณปจฺจยา.\csromanspacer{} ปฏิสนฺธิกฺขเณ จิตฺตสมฺปยุตฺตํ เอกํ ขนฺธญฺจ วตฺถุญฺจ ปฏิจฺจ เทฺว ขนฺธา.\csromanspacer{} เทฺว ขนฺเธ จ ฯเปฯ. (สํขิตฺตํ.) (1)}

\csromanheader{2}{1. ปจฺจยานุโลม 2. สงฺขฺยาวาร}
\csromanheader{2}{\textbf{สุทฺธ}}
\proseitem{148}{\csromanfolio{426}เหตุยา นว, อารมฺมเณ ตีณิ, อธิปติยา ปญฺจ, อนนฺตเร ตีณิ, สมนนฺตเร ตีณิ, สหชาเต นว, อญฺญมญฺเญ ฉ, นิสฺสเย นว, อุปนิสฺสเย ตีณิ, ปุเรชาเต เอกํ, อาเสวเน เอกํ, กมฺเม นว, วิปาเก นว, อาหาเร นว, อินฺทฺริเย นว, ฌาเน นว, มคฺเค นว, สมฺปยุตฺเต ตีณิ, วิปฺปยุตฺเต นว, อตฺถิยา นว, นตฺถิยา ตีณิ, วิคเต ตีณิ, อวิคเต นว.}

\csromanheader{2}{2. ปจฺจยปจฺจนีย 1. วิภงฺควาร}
\csromanheader{2}{\textbf{นเหตุ}}
\proseitem{149}{จิตฺตสมฺปยุตฺตํ ธมฺมํ ปฏิจฺจ จิตฺตสมฺปยุตฺโต ธมฺโม อุปฺปชฺชติ นเหตุปจฺจยา.\csromanspacer{} อเหตุกํ จิตฺตสมฺปยุตฺตํ เอกํ ขนฺธํ ปฏิจฺจ เทฺว ขนฺธา.\csromanspacer{} เทฺว ขนฺเธ ฯเปฯ.\csromanspacer{} อเหตุกปฏิสนฺธิกฺขเณ ฯเปฯ.\csromanspacer{} วิจิกิจฺฉาสหคเต อุทฺธจฺจสหคเต ขนฺเธ ปฏิจฺจ วิจิกิจฺฉาสหคโต อุทฺธจฺจสหคโต โมโห. (1)}

\prose{(เอวํ นวปิ ปญฺหา กาตพฺพา. “อเหตุกนฺ”ติ สพฺพตฺถ นิยาเมตพฺพํ, เอกํเยว โมหํ มูลปเท.)}

\csromanheader{2}{\textbf{น-อารมฺมณ}}
\proseitem{150}{จิตฺตสมฺปยุตฺตํ ธมฺมํ ปฏิจฺจ จิตฺตวิปฺปยุตฺโต ธมฺโม อุปฺปชฺชติ น-อารมฺมณปจฺจยา.\csromanspacer{} จิตฺตสมฺปยุตฺเต ขนฺเธ ปฏิจฺจ จิตฺตสมุฏฺฐานํ รูปํ.\csromanspacer{} ปฏิสนฺธิกฺขเณ ฯเปฯ. (1)}

\prose{จิตฺตวิปฺปยุตฺตํ ธมฺมํ ปฏิจฺจ จิตฺตวิปฺปยุตฺโต ธมฺโม อุปฺปชฺชติ น-อารมฺมณปจฺจยา.\csromanspacer{} เอกํ มหาภูตํ ฯเปฯ. (ยาว อสญฺญสตฺตา.) (1)}

\prose{จิตฺตสมฺปยุตฺตญฺจ จิตฺตวิปฺปยุตฺตญฺจ ธมฺมํ ปฏิจฺจ จิตฺตวิปฺปยุตฺโต ธมฺโม อุปฺปชฺชติ น-อารมฺมณปจฺจยา.\csromanspacer{} จิตฺตสมฺปยุตฺเต ขนฺเธ จ มหาภูเต จ ปฏิจฺจ จิตฺตสมุฏฺฐานํ รูปํ.\csromanspacer{} ปฏิสนฺธิกฺขเณ เอกํ. (สํขิตฺตํ.) (1)}

\csromanheader{2}{58. จิตฺตสมฺปยุตฺตทุก \textbf{2. ปจฺจยปจฺจนีย 2. สงฺขฺยาวาร}}
\csromanheader{2}{\textbf{สุทฺธ}}
\proseitem{151}{\csromanfolio{427}นเหตุยา นว, น-อารมฺมเณ ตีณิ, น-อธิปติยา นว, น-อนนฺตเร ตีณิ, นสมนนฺตเร ตีณิ, น-อญฺญมญฺเญ ตีณิ, น-อุปนิสฺสเย ตีณิ, นปุเรชาเต นว, นปจฺฉาชาเต นว, น-อาเสวเน นว, นกมฺเม เทฺว, นวิปาเก ปญฺจ, น-อาหาเร เอกํ, น-อินฺทฺริเย เอกํ, นฌาเน เทฺว, นมคฺเค นว, นสมฺปยุตฺเต ตีณิ, นวิปฺปยุตฺเต เทฺว, โนนตฺถิยา ตีณิ, โนวิคเต ตีณิ.}

\csromanheader{2}{3. ปจฺจยานุโลมปจฺจนีย}
\proseitem{152}{\textbf{เหตุ}ปจฺจยา น-อารมฺมเณ ตีณิ, น-อธิปติยา นว ฯเปฯ นกมฺเม เอกํ, นวิปาเก ปญฺจ, นสมฺปยุตฺเต ตีณิ, นวิปฺปยุตฺเต เอกํ, โนนตฺถิยา ตีณิ, โนวิคเต ตีณิ.}

\csromanheader{2}{4. ปจฺจยปจฺจนียานุโลม}
\proseitem{153}{นเหตุปจฺจยา อารมฺมเณ ตีณิ, อนนฺตเร ตีณิ, สมนนฺตเร ตีณิ ฯเปฯ อญฺญมญฺเญ ฉ ฯเปฯ ปุเรชาเต เอกํ, อาเสวเน เอกํ, กมฺเม นว ฯเปฯ มคฺเค เอกํ ฯเปฯ อวิคเต นว.}

\vspace{\tipitakaparskip}
\csromancenter{(สหชาตวาโร ปฏิจฺจวารสทิโส.)}
\csromanheader{2}{58. จิตฺตสมฺปยุตฺตทุก 3. ปจฺจยวาร}
\csromanheader{2}{1. ปจฺจยานุโลม 1. วิภงฺควาร}
\csromanheader{2}{\textbf{เหตุ}}
\proseitem{154}{จิตฺตสมฺปยุตฺตํ ธมฺมํ ปจฺจยา จิตฺตสมฺปยุตฺโต ธมฺโม อุปฺปชฺชติ เหตุปจฺจยา.\csromanspacer{} ตีณิ. (ปฏิจฺจสทิสา.)}

\prose{จิตฺตวิปฺปยุตฺตํ ธมฺมํ ปจฺจยา จิตฺตวิปฺปยุตฺโต ธมฺโม อุปฺปชฺชติ เหตุปจฺจยา.\csromanspacer{} เอกํ. (ปฏิจฺจสทิสํ) (1)}

\prose{\csromanfolio{428}จิตฺตวิปฺปยุตฺตํ ธมฺมํ ปจฺจยา จิตฺตสมฺปยุตฺโต ธมฺโม อุปฺปชฺชติ เหตุปจฺจยา.\csromanspacer{} วตฺถุํ ปจฺจยา จิตฺตสมฺปยุตฺตกา ขนฺธา.\csromanspacer{} ปฏิสนฺธิกฺขเณ ฯเปฯ. (2)}

\prose{จิตฺตวิปฺปยุตฺตํ ธมฺมํ ปจฺจยา จิตฺตสมฺปยุตฺโต จ จิตฺตวิปฺปยุตฺโต จ ธมฺมา อุปฺปชฺชนฺติ เหตุปจฺจยา.\csromanspacer{} วตฺถุํ ปจฺจยา จิตฺตสมฺปยุตฺตกา ขนฺธา.\csromanspacer{} มหาภูเต ปจฺจยา จิตฺตสมุฏฺฐานํ รูปํ.\csromanspacer{} ปฏิสนฺธิกฺขเณ ฯเปฯ. (3)}

\proseitem{155}{จิตฺตสมฺปยุตฺตญฺจ จิตฺตวิปฺปยุตฺตญฺจ ธมฺมํ ปจฺจยา จิตฺตสมฺปยุตฺโต ธมฺโม อุปฺปชฺชติ เหตุปจฺจยา.\csromanspacer{} จิตฺตสมฺปยุตฺตํ เอกํ ขนฺธญฺจ วตฺถุญฺจ ปจฺจยา เทฺว ขนฺธา.\csromanspacer{} เทฺว ขนฺเธ จ ฯเปฯ.\csromanspacer{} ปฏิสนฺธิกฺขเณ ฯเปฯ. (1)}

\prose{จิตฺตสมฺปยุตฺตญฺจ จิตฺตวิปฺปยุตฺตญฺจ ธมฺมํ ปจฺจยา จิตฺตวิปฺปยุตฺโต ธมฺโม อุปฺปชฺชติ เหตุปจฺจยา.\csromanspacer{} จิตฺตสมฺปยุตฺเต ขนฺเธ จ มหาภูเต จ ปจฺจยา จิตฺตสมุฏฺฐานํ รูปํ.\csromanspacer{} ปฏิสนฺธิกฺขเณ ฯเปฯ. (2)}

\prose{จิตฺตสมฺปยุตฺตญฺจ จิตฺตวิปฺปยุตฺตญฺจ ธมฺมํ ปจฺจยา จิตฺตสมฺปยุตฺโต จ จิตฺตวิปฺปยุตฺโต จ ธมฺมา อุปฺปชฺชนฺติ เหตุปจฺจยา.\csromanspacer{} จิตฺตสมฺปยุตฺตํ เอกํ ขนฺธญฺจ วตฺถุญฺจ ปจฺจยา เทฺว ขนฺธา.\csromanspacer{} เทฺว ขนฺเธ จ ฯเปฯ.\csromanspacer{} จิตฺตสมฺปยุตฺเต ขนฺเธ จ มหาภูเต จ ปจฺจยา จิตฺตสมุฏฺฐานํ รูปํ.\csromanspacer{} ปฏิสนฺธิกฺขเณ ฯเปฯ. (3)}

\csromanheader{2}{\textbf{อารมฺมณ}}
\proseitem{156}{จิตฺตสมฺปยุตฺตํ ธมฺมํ ปจฺจยา จิตฺตสมฺปยุตฺโต ธมฺโม อุปฺปชฺชติ อารมฺมณปจฺจยา.\csromanspacer{} เอกํ. (ปฏิจฺจสทิสํ.) (1)}

\prose{จิตฺตวิปฺปยุตฺตํ ธมฺมํ ปจฺจยา จิตฺตสมฺปยุตฺโต ธมฺโม อุปฺปชฺชติ อารมฺมณปจฺจยา.\csromanspacer{} จกฺขายตนํ ปจฺจยา จกฺขุวิญฺญาณสหคตา ขนฺธา ฯเปฯ.\csromanspacer{} กายายตนํ ปจฺจยา กายวิญฺญาณสหคตา ขนฺธา ฯเปฯ.\csromanspacer{} วตฺถุํ ปจฺจยา จิตฺตสมฺปยุตฺตกา ขนฺธา.\csromanspacer{} ปฏิสนฺธิกฺขเณ ฯเปฯ. (1)}

\prose{จิตฺตสมฺปยุตฺตญฺจ จิตฺตวิปฺปยุตฺตญฺจ ธมฺมํ ปจฺจยา จิตฺตสมฺปยุตฺโต ธมฺโม อุปฺปชฺชติ อารมฺมณปจฺจยา.\csromanspacer{} จกฺขุวิญฺญาณสหคตํ เอกํ ขนฺธญฺจ จกฺขายตนญฺจ ปจฺจยา เทฺว ขนฺธา.\csromanspacer{} เทฺว ขนฺเธ จ ฯเปฯ.\csromanspacer{} กายวิญฺญาณสหคตํ เอกํ ขนฺธญฺจ กายายตนญฺจ ปจฺจยา เทฺว ขนฺธา.\csromanspacer{} เทฺว ขนฺเธ จ ฯเปฯ.\csromanspacer{} จิตฺตสมฺปยุตฺตํ เอกํ ขนฺธญฺจ วตฺถุญฺจ ปจฺจยา เทฺว ขนฺธา.\csromanspacer{} เทฺว ขนฺเธ จ ฯเปฯ.\csromanspacer{} ปฏิสนฺธิกฺขเณ. (สํขิตฺตํ.) (1)}

\csromanheader{2}{58. จิตฺตสมฺปยุตฺตทุก \textbf{1. ปจฺจยานุโลม 2. สงฺขฺยาวาร}}
\csromanheader{2}{\textbf{สุทฺธ}}
\proseitem{157}{\csromanfolio{429}เหตุยา นว, อารมฺมเณ ตีณิ, อธิปติยา นว, อนนฺตเร ตีณิ, สมนนฺตเร ตีณิ, สหชาเต นว, อญฺญมญฺเญ ฉ, นิสฺสเย นว, อุปนิสฺสเย ตีณิ, ปุเรชาเต ตีณิ, อาเสวเน ตีณิ, กมฺเม นว ฯเปฯ อวิคเต นว.}

\csromanheader{2}{2. ปจฺจยปจฺจนีย 1. วิภงฺควาร}
\csromanheader{2}{\textbf{นเหตุ}}
\proseitem{158}{จิตฺตสมฺปยุตฺตํ ธมฺมํ ปจฺจยา จิตฺตสมฺปยุตฺโต ธมฺโม อุปฺปชฺชติ นเหตุปจฺจยา.\csromanspacer{} ตีณิ. (ปฏิจฺจสทิสา.)}

\prose{จิตฺตวิปฺปยุตฺตํ ธมฺมํ ปจฺจยา จิตฺตวิปฺปยุตฺโต ธมฺโม อุปฺปชฺชติ นเหตุปจฺจยา.\csromanspacer{} เอกํ มหาภูตํ ฯเปฯ. (ยาว อสญฺญสตฺตา.) (1)}

\prose{จิตฺตวิปฺปยุตฺตํ ธมฺมํ ปจฺจยา จิตฺตสมฺปยุตฺโต ธมฺโม อุปฺปชฺชติ นเหตุปจฺจยา.\csromanspacer{} จกฺขายตนํ ปจฺจยา จกฺขุวิญฺญาณสหคตา ขนฺธา ฯเปฯ.\csromanspacer{} กายายตนํ ฯเปฯ.\csromanspacer{} วตฺถุํ ปจฺจยา อเหตุกา จิตฺตสมฺปยุตฺตกา ขนฺธา.\csromanspacer{} ปฏิสนฺธิกฺขเณ ฯเปฯ.\csromanspacer{} วตฺถุํ ปจฺจยา วิจิกิจฺฉาสหคโต อุทฺธจฺจสหคโต โมโห. (2)}

\prose{จิตฺตวิปฺปยุตฺตํ ธมฺมํ ปจฺจยา จิตฺตสมฺปยุตฺโต จ จิตฺตวิปฺปยุตฺโต จ ธมฺมา อุปฺปชฺชนฺติ นเหตุปจฺจยา.\csromanspacer{} วตฺถุํ ปจฺจยา อเหตุกา จิตฺตสมฺปยุตฺตกา ขนฺธา.\csromanspacer{} มหาภูเต ปจฺจยา จิตฺตสมุฏฺฐานํ รูปํ.\csromanspacer{} ปฏิสนฺธิกฺขเณ ฯเปฯ. (3)}

\prose{จิตฺตสมฺปยุตฺตญฺจ จิตฺตวิปฺปยุตฺตญฺจ ธมฺมํ ปจฺจยา จิตฺตสมฺปยุตฺโต ธมฺโม อุปฺปชฺชติ นเหตุปจฺจยา.\csromanspacer{} จกฺขุวิญฺญาณสหคตํ เอกํ ขนฺธญฺจ จกฺขายตนญฺจ ปจฺจยา เทฺว ขนฺธา.\csromanspacer{} เทฺว ขนฺเธ จ ฯเปฯ.\csromanspacer{} กายวิญฺญาณสหคตํ เอกํ ขนฺธญฺจ กายายตนญฺจ ปจฺจยา เทฺว ขนฺธา.\csromanspacer{} เทฺว ขนฺเธ จ ฯเปฯ.\csromanspacer{} อเหตุกํ จิตฺตสมฺปยุตฺตํ เอกํ ขนฺธญฺจ วตฺถุญฺจ ปจฺจยา เทฺว ขนฺธา.\csromanspacer{} เทฺว ขนฺเธ จ ฯเปฯ.\csromanspacer{} อเหตุกปฏิสนฺธิกฺขเณ ฯเปฯ.\csromanspacer{} วิจิกิจฺฉาสหคเต อุทฺธจฺจสหคเต ขนฺเธ จ วตฺถุญฺจ ปจฺจยา วิจิกิจฺฉาสหคโต อุทฺธจฺจสหคโต โมโห. (เอวํ เทฺว ปญฺหา ปวตฺติปฏิสนฺธิ กาตพฺพา.\csromanspacer{} สํขิตฺตํ.)}

\csromanheader{2}{2. ปจฺจยปจฺจนีย 2. สงฺขฺยาวาร}
\csromanheader{2}{\textbf{สุทฺธ}}
\proseitem{159}{\csromanfolio{430}นเหตุยา นว, น-อารมฺมเณ ตีณิ, น-อธิปติยา นว, น-อนนฺตเร ตีณิ, นสมนนฺตเร ตีณิ, น-อญฺญมญฺเญ ตีณิ, น-อุปนิสฺสเย ตีณิ, นปุเรชาเต นว, นปจฺฉาชาเต นว, น-อาเสวเน นว, นกมฺเม จตฺตาริ, นวิปาเก นว, น-อาหาเร เอกํ, น-อินฺทฺริเย เอกํ, นฌาเน จตฺตาริ, นมคฺเค นว, นสมฺปยุตฺเต ตีณิ, นวิปฺปยุตฺเต เทฺว, โนนตฺถิยา ตีณิ, โนวิคเต ตีณิ.}

\csromanheader{2}{3. ปจฺจยานุโลมปจฺจนีย}
\proseitem{160}{เหตุปจฺจยา น-อารมฺมเณ ตีณิ ฯเปฯ นกมฺเม ตีณิ ฯเปฯ นวิปฺปยุตฺเต เอกํ, โนนตฺถิยา ตีณิ, โนวิคเต ตีณิ.}

\csromanheader{2}{4. ปจฺจยปจฺจนียานุโลม}
\vspace{\tipitakaparskip}
\csromancenter{นเหตุปจฺจยา อารมฺมเณ ตีณิ ฯเปฯ มคฺเค ตีณิ ฯเปฯ อวิคเต นว.}
\csromancenter{(นิสฺสยวาโร ปจฺจยวารสทิโส.)}
\csromanheader{2}{58. จิตฺตสมฺปยุตฺตทุก 5. สํสฏฺฐวาร}
\csromanheader{2}{1. ปจฺจยานุโลมาทิ}
\proseitem{162}{จิตฺตสมฺปยุตฺตํ ธมฺมํ สํสฏฺโฐ จิตฺตสมฺปยุตฺโต ธมฺโม อุปฺปชฺชติ เหตุปจฺจยา.\csromanspacer{} จิตฺตสมฺปยุตฺตํ เอกํ ขนฺธํ สํสฏฺฐา เทฺว ขนฺธา.\csromanspacer{} เทฺว ขนฺเธ ฯเปฯ.\csromanspacer{} ปฏิสนฺธิกฺขเณ ฯเปฯ.}

\prose{เหตุยา เอกํ, อารมฺมเณ เอกํ, อธิปติยา เอกํ, (สพฺพตฺถ เอกํ,) อวิคเต เอกํ.}

\csromanheader{2}{58. จิตฺตสมฺปยุตฺตทุก}
\prose{\csromanfolio{431}นเหตุยา เอกํ, น-อธิปติยา เอกํ, นปุเรชาเต เอกํ, นปจฺฉาชาเต เอกํ, น-อาเสวเน เอกํ, นกมฺเม เอกํ, นวิปาเก เอกํ, นฌาเน เอกํ, นมคฺเค เอกํ, นวิปฺปยุตฺเต เอกํ.}

\prose{(เอวํ อิตเร เทฺว คณนาปิ สมฺปยุตฺตวาโรปิ กาตพฺโพ.)}

\csromanheader{2}{58. จิตฺตสมฺปยุตฺตทุก 7. ปญฺหาวาร}
\csromanheader{2}{1. ปจฺจยานุโลม 1. วิภงฺควาร}
\csromanheader{2}{\textbf{เหตุ}}
\proseitem{163}{จิตฺตสมฺปยุตฺโต ธมฺโม จิตฺตสมฺปยุตฺตสฺส ธมฺมสฺส เหตุปจฺจเยน ปจฺจโย.\csromanspacer{} จิตฺตสมฺปยุตฺตา เหตู สมฺปยุตฺตกานํ ขนฺธานํ เหตุปจฺจเยน ปจฺจโย.\csromanspacer{} ปฏิสนฺธิกฺขเณ. (1)}

\prose{จิตฺตสมฺปยุตฺโต ธมฺโม จิตฺตวิปฺปยุตฺตสฺส ธมฺมสฺส เหตุปจฺจเยน ปจฺจโย.\csromanspacer{} จิตฺตสมฺปยุตฺตา เหตู จิตฺตสมุฏฺฐานานํ รูปานํ เหตุปจฺจเยน ปจฺจโย.\csromanspacer{} ปฏิสนฺธิกฺขเณ ฯเปฯ. (2)}

\prose{จิตฺตสมฺปยุตฺโต ธมฺโม จิตฺตสมฺปยุตฺตสฺส จ จิตฺตวิปฺปยุตฺตสฺส จ ธมฺมสฺส เหตุปจฺจเยน ปจฺจโย.\csromanspacer{} จิตฺตสมฺปยุตฺตา เหตู สมฺปยุตฺตกานํ ขนฺธานํ จิตฺตสมุฏฺฐานานญฺจ รูปานํ เหตุปจฺจเยน ปจฺจโย.\csromanspacer{} ปฏิสนฺธิกฺขเณ ฯเปฯ. (3)}

\csromanheader{2}{\textbf{อารมฺมณ}}
\proseitem{164}{จิตฺตสมฺปยุตฺโต ธมฺโม จิตฺตสมฺปยุตฺตสฺส ธมฺมสฺส อารมฺมณปจฺจเยน ปจฺจโย.\csromanspacer{} ทานํ ฯเปฯ สีลํ ฯเปฯ อุโปสถกมฺมํ กตฺวา ตํ ปจฺจเวกฺขติ, อสฺสาเทติ อภินนฺทติ, ตํ อารพฺภ ราโค อุปฺปชฺชติ ฯเปฯ โทมนสฺสํ อุปฺปชฺชติ.\csromanspacer{} ปุพฺเพ สุจิณฺณานิ ฯเปฯ.\csromanspacer{} ฌานา วุฏฺฐหิตฺวา ฌานํ ฯเปฯ.\csromanspacer{} อริยา มคฺคา วุฏฺฐหิตฺวา มคฺคํ ปจฺจเวกฺขนฺติ, ผลํ ปจฺจเวกฺขนฺติ.\csromanspacer{} ปหีเน กิเลเส ฯเปฯ.\csromanspacer{} วิกฺขมฺภิเต กิเลเส ปจฺจเวกฺขนฺติ.\csromanspacer{} ปุพฺเพ สมุทาจิณฺเณ กิเลเส ชานนฺติ.\csromanspacer{} จิตฺตสมฺปยุตฺเต ขนฺเธ อนิจฺจโต ฯเปฯ โทมนสฺสํ อุปฺปชฺชติ.\csromanspacer{} เจโตปริยญาเณน จิตฺตสมฺปยุตฺตสมงฺคิสฺส จิตฺตํ ชานนฺติ.\csromanspacer{} อากาสานญฺจายตนํ วิญฺญาณญฺจายตนสฺส ฯเปฯ. \csromanfolio{432}อากิญฺจญฺญายตนํ เนวสญฺญานาสญฺญายตนสฺส ฯเปฯ.\csromanspacer{} จิตฺตสมฺปยุตฺตา ขนฺธา อิทฺธิวิธญาณสฺส, เจโตปริยญาณสฺส, ปุพฺเพนิวาสานุสฺสติญาณสฺส, ยถากมฺมูปคญาณสฺส, อนาคตํสญาณสฺส, อาวชฺชนาย อารมฺมณปจฺจเยน ปจฺจโย. (1)}

\proseitem{165}{จิตฺตวิปฺปยุตฺโต ธมฺโม จิตฺตสมฺปยุตฺตสฺส ธมฺมสฺส อารมฺมณปจฺจเยน ปจฺจโย.\csromanspacer{} อริยา นิพฺพานํ ปจฺจเวกฺขนฺติ.\csromanspacer{} นิพฺพานํ โคตฺรภุสฺส, โวทานสฺส, มคฺคสฺส, ผลสฺส, อาวชฺชนาย อารมฺมณปจฺจเยน ปจฺจโย.\csromanspacer{} จกฺขุํ ฯเปฯ วตฺถุํ จิตฺตวิปฺปยุตฺเต ขนฺเธ อนิจฺจโต ฯเปฯ โทมนสฺสํ อุปฺปชฺชติ.\csromanspacer{} ทิพฺเพน จกฺขุนา รูปํ ปสฺสติ.\csromanspacer{} ทิพฺพาย โสตธาตุยา สทฺทํ สุณาติ.\csromanspacer{} รูปายตนํ จกฺขุวิญฺญาณสหคตานํ ขนฺธานํ ฯเปฯ.\csromanspacer{} โผฏฺฐพฺพายตนํ ฯเปฯ.\csromanspacer{} จิตฺตวิปฺปยุตฺตา ขนฺธา อิทฺธิวิธญาณสฺส ปุพฺเพนิวาสานุสฺสติญาณสฺส, อนาคตํสญาณสฺส, อาวชฺชนาย อารมฺมณปจฺจเยน ปจฺจโย. (1)}

\csromanheader{2}{\textbf{อธิปติ}}
\proseitem{166}{จิตฺตสมฺปยุตฺโต ธมฺโม จิตฺตสมฺปยุตฺตสฺส ธมฺมสฺส อธิปติปจฺจเยน ปจฺจโย.\csromanspacer{} \textbf{อารมฺมณาธิปติ}, สหชาตาธิปติ.\csromanspacer{}\csromanspacer{}\textbf{อารมฺมณาธิปติ}\csromandash{}ทานํ ฯเปฯ สีลํ ฯเปฯ อุโปสถกมฺมํ กตฺวา ตํ ครุํ กตฺวา ปจฺจเวกฺขติ, อสฺสาเทติ อภินนฺทติ, ตํ ครุํ กตฺวา ราโค อุปฺปชฺชติ.\csromanspacer{} ทิฏฺฐิ อุปฺปชฺชติ.\csromanspacer{} ปุพฺเพ ฯเปฯ.\csromanspacer{} ฌานา ฯเปฯ.\csromanspacer{} อริยา มคฺคา วุฏฺฐหิตฺวา มคฺคํ ครุํ กตฺวา ปจฺจเวกฺขนฺติ.\csromanspacer{} ผลํ ครุํ กตฺวา ปจฺจเวกฺขนฺติ.\csromanspacer{} จิตฺตสมฺปยุตฺเต ขนฺเธ ครุํ กตฺวา อสฺสาเทติ อภินนฺทติ, ตํ ครุํ กตฺวา ราโค อุปฺปชฺชติ.\csromanspacer{} ทิฏฺฐิ อุปฺปชฺชติ.\csromanspacer{} \textbf{สหชาตาธิปติ}\csromandash{} จิตฺตสมฺปยุตฺตาธิปติ สมฺปยุตฺตกานํ ขนฺธานํ อธิปติปจฺจเยน ปจฺจโย. (1)}

\prose{จิตฺตสมฺปยุตฺโต ธมฺโม จิตฺตวิปฺปยุตฺตสฺส ธมฺมสฺส อธิปติปจฺจเยน ปจฺจโย.\csromanspacer{} \textbf{สหชาตาธิปติ}\csromandash{}จิตฺตสมฺปยุตฺตาธิปติ จิตฺตสมุฏฺฐานานํ รูปานํ อธิปติปจฺจเยน ปจฺจโย. (2)}

\prose{จิตฺตสมฺปยุตฺโต ธมฺโม จิตฺตสมฺปยุตฺตสฺส จ จิตฺตวิปฺปยุตฺตสฺส จ ธมฺมสฺส อธิปติปจฺจเยน ปจฺจโย.\csromanspacer{} \textbf{สหชาตาธิปติ}\csromandash{} จิตฺตสมฺปยุตฺตาธิปติ สมฺปยุตฺตกานํ ขนฺธานํ จิตฺตสมุฏฺฐานานญฺจ รูปานํ อธิปติปจฺจเยน ปจฺจโย. (3)}

\csromanheader{2}{58. จิตฺตสมฺปยุตฺตทุก}
\proseitem{167}{\csromanfolio{433}จิตฺตวิปฺปยุตฺโต ธมฺโม จิตฺตสมฺปยุตฺตสฺส ธมฺมสฺส อธิปติปจฺจเยน ปจฺจโย.\csromanspacer{} \textbf{อารมฺมณาธิปติ}\csromandash{}อริยา นิพฺพานํ ครุํ กตฺวา ปจฺจเวกฺขนฺติ.\csromanspacer{} นิพฺพานํ โคตฺรภุสฺส, โวทานสฺส, มคฺคสฺส, ผลสฺส อธิปติปจฺจเยน ปจฺจโย.\csromanspacer{} จกฺขุํ ฯเปฯ วตฺถุํ จิตฺตวิปฺปยุตฺเต ขนฺเธ ครุํ กตฺวา อสฺสาเทติ อภินนฺทติ, ตํ ครุํ กตฺวา ราโค อุปฺปชฺชติ.\csromanspacer{} ทิฏฺฐิ อุปฺปชฺชติ. (1)}

\csromanheader{2}{\textbf{อนนฺตราทิ}}
\proseitem{168}{จิตฺตสมฺปยุตฺโต ธมฺโม จิตฺตสมฺปยุตฺตสฺส ธมฺมสฺส อนนฺตรปจฺจเยน ปจฺจโย.\csromanspacer{} ปุริมา ปุริมา จิตฺตสมฺปยุตฺตา ขนฺธา ฯเปฯ ผลสมาปตฺติยา อนนฺตรปจฺจเยน ปจฺจโย.\csromanspacer{} สมนนฺตรปจฺจเยน ปจฺจโย.\csromanspacer{} สหชาตปจฺจเยน ปจฺจโย.\csromanspacer{} สตฺต. (ปฏิจฺจสทิสา, ปญฺหาฆฏนา นตฺถิ.) อญฺญมญฺญปจฺจเยน ปจฺจโย.\csromanspacer{} ฉ. (ปฏิจฺจสทิสา.) นิสฺสยปจฺจเยน ปจฺจโย.\csromanspacer{} สตฺต. (ปจฺจยวารสทิสา, ปญฺหาฆฏนา นตฺถิ.)}

\csromanheader{2}{\textbf{อุปนิสฺสย}}
\proseitem{169}{จิตฺตสมฺปยุตฺโต ธมฺโม จิตฺตสมฺปยุตฺตสฺส ธมฺมสฺส อุปนิสฺสยปจฺจเยน ปจฺจโย.\csromanspacer{} \textbf{อารมฺมณูปนิสฺสโย, อนนฺตรูปนิสฺสโย,} \textbf{ปกตูปนิสฺสโย ฯเปฯ.}\csromanspacer{} ปกตูปนิสฺสโย\csromandash{}สทฺธํ อุปนิสฺสาย ทานํ เทติ ฯเปฯ มานํ ชปฺเปติ, ทิฏฺฐิํ คณฺหาติ, สีลํ ฯเปฯ.\csromanspacer{} ปตฺถนํ.\csromanspacer{} กายิกํ สุขํ.\csromanspacer{} กายิกํ ทุกฺขํ อุปนิสฺสาย ทานํ เทติ ฯเปฯ.\csromanspacer{} ปาณํ หนติ ฯเปฯ สํฆํ ภินฺทติ.\csromanspacer{} สทฺธา ฯเปฯ.\csromanspacer{} กายิกํ ทุกฺขํ สทฺธาย ฯเปฯ มคฺคสฺส, ผลสมาปตฺติยา อุปนิสฺสยปจฺจเยน ปจฺจโย. (1)}

\prose{จิตฺตวิปฺปยุตฺโต ธมฺโม จิตฺตสมฺปยุตฺตสฺส ธมฺมสฺส อุปนิสฺสยปจฺจเยน ปจฺจโย.\csromanspacer{} \textbf{อารมฺมณูปนิสฺสโย, ปกตูปนิสฺสโย ฯเปฯ.}\csromanspacer{} ปกตูปนิสฺสโย\csromandash{}อุตุํ.\csromanspacer{} โภชนํ.\csromanspacer{} เสนาสนํ อุปนิสฺสาย ทานํ เทติ ฯเปฯ สํฆํ ภินฺทติ.\csromanspacer{} อุตุ.\csromanspacer{} โภชนํ.\csromanspacer{} เสนาสนํ สทฺธาย ฯเปฯ ผลสมาปตฺติยา อุปนิสฺสยปจฺจเยน ปจฺจโย. (1)}

\csromanheader{2}{\textbf{ปุเรชาต}}
\proseitem{170}{จิตฺตวิปฺปยุตฺโต ธมฺโม จิตฺตสมฺปยุตฺตสฺส ธมฺมสฺส ปุเรชาตปจฺจเยน ปจฺจโย.\csromanspacer{} \textbf{อารมฺมณปุเรชาตํ}, วตฺถุปุเรชาตํ.\csromanspacer{}\csromanspacer{}\textbf{อารมฺมณปุเรชาตํ}\csromandash{}จกฺขุํ ฯเปฯ วตฺถุํ อนิจฺจโต ฯเปฯ โทมนสฺสํ อุปฺปชฺชติ.\csromanspacer{} ทิพฺเพน \csromanfolio{434}จกฺขุนา รูปํ ปสฺสติ.\csromanspacer{} ทิพฺพาย โสตธาตุยา สทฺทํ สุณาติ.\csromanspacer{} รูปายตนํ จกฺขุวิญฺญาณสหคตานํ ขนฺธานํ ฯเปฯ โผฏฺฐพฺพายตนํ ฯเปฯ.\csromanspacer{} \textbf{วตฺถุปุเรชาตํ}\csromandash{}จกฺขายตนํ จกฺขุวิญฺญาณสหคตานํ ขนฺธานํ ฯเปฯ กายายตนํ ฯเปฯ.\csromanspacer{} วตฺถุ จิตฺตสมฺปยุตฺตกานํ ขนฺธานํ ปุเรชาตปจฺจเยน ปจฺจโย. (1)}

\csromanheader{2}{\textbf{ปจฺฉาชาตาเสวน}}
\proseitem{171}{จิตฺตสมฺปยุตฺโต ธมฺโม จิตฺตวิปฺปยุตฺตสฺส ธมฺมสฺส ปจฺฉาชาตปจฺจเยน ปจฺจโย. (สํขิตฺตํ.) เอกํ.\csromanspacer{} อาเสวนปจฺจเยน ปจฺจโย.\csromanspacer{} เอกํ.}

\csromanheader{2}{\textbf{กมฺม}}
\proseitem{172}{จิตฺตสมฺปยุตฺโต ธมฺโม จิตฺตสมฺปยุตฺตสฺส ธมฺมสฺส กมฺมปจฺจเยน ปจฺจโย.\csromanspacer{} \textbf{สหชาตา}, นานากฺขณิกา.\csromanspacer{}\csromanspacer{}\textbf{สหชาตา} จิตฺตสมฺปยุตฺตา เจตนา สมฺปยุตฺตกานํ ขนฺธานํ กมฺมปจฺจเยน ปจฺจโย.\csromanspacer{} ปฏิสนฺธิกฺขเณ ฯเปฯ.\csromanspacer{} \textbf{นานากฺขณิกา} จิตฺตสมฺปยุตฺตา เจตนา วิปากานํ ขนฺธานํ กมฺมปจฺจเยน ปจฺจโย. (1)}

\prose{จิตฺตสมฺปยุตฺโต ธมฺโม จิตฺตวิปฺปยุตฺตสฺส ธมฺมสฺส กมฺมปจฺจเยน ปจฺจโย.\csromanspacer{} \textbf{สหชาตา}, นานากฺขณิกา.\csromanspacer{}\csromanspacer{}\textbf{สหชาตา} จิตฺตสมฺปยุตฺตา เจตนา จิตฺตสมุฏฺฐานานํ รูปานํ กมฺมปจฺจเยน ปจฺจโย.\csromanspacer{} ปฏิสนฺธิกฺขเณ ฯเปฯ.\csromanspacer{} \textbf{นานากฺขณิกา} จิตฺตสมฺปยุตฺตา เจตนา กฏตฺตารูปานํ กมฺมปจฺจเยน ปจฺจโย. (2)}

\prose{จิตฺตสมฺปยุตฺโต ธมฺโม จิตฺตสมฺปยุตฺตสฺส จ จิตฺตวิปฺปยุตฺตสฺส จ ธมฺมสฺส กมฺมปจฺจเยน ปจฺจโย.\csromanspacer{} \textbf{สหชาตา}, นานากฺขณิกา.\csromanspacer{}\csromanspacer{}\textbf{สหชาตา} จิตฺตสมฺปยุตฺตา เจตนา สมฺปยุตฺตกานํ ขนฺธานํ จิตฺตสมุฏฺฐานานญฺจ รูปานํ กมฺมปจฺจเยน ปจฺจโย.\csromanspacer{} ปฏิสนฺธิกฺขเณ ฯเปฯ.\csromanspacer{} \textbf{นานากฺขณิกา} จิตฺตสมฺปยุตฺตา เจตนา วิปากานํ ขนฺธานํ กฏตฺตา จ รูปานํ กมฺมปจฺจเยน ปจฺจโย. (3)}

\csromanheader{2}{\textbf{วิปากาหาร}}
\proseitem{173}{จิตฺตสมฺปยุตฺโต ธมฺโม จิตฺตสมฺปยุตฺตสฺส ธมฺมสฺส วิปากปจฺจเยน ปจฺจโย.\csromanspacer{} ตีณิ.}

\prose{จิตฺตสมฺปยุตฺโต ธมฺโม จิตฺตสมฺปยุตฺตสฺส ธมฺมสฺส อาหารปจฺจเยน ปจฺจโย.\csromanspacer{} ตีณิ.}

\csromanheader{2}{58. จิตฺตสมฺปยุตฺตทุก}
\prose{\csromanfolio{435}จิตฺตวิปฺปยุตฺโต ธมฺโม จิตฺตวิปฺปยุตฺตสฺส ธมฺมสฺส อาหารปจฺจเยน ปจฺจโย.\csromanspacer{} กพฬีกาโร อาหาโร อิมสฺส กายสฺส อาหารปจฺจเยน ปจฺจโย. (1)}

\csromanheader{2}{\textbf{อินฺทฺริยาทิ}}
\proseitem{174}{จิตฺตสมฺปยุตฺโต ธมฺโม จิตฺตสมฺปยุตฺตสฺส ธมฺมสฺส อินฺทฺริยปจฺจเยน ปจฺจโย.\csromanspacer{} ตีณิ.}

\prose{จิตฺตวิปฺปยุตฺโต ธมฺโม จิตฺตวิปฺปยุตฺตสฺส ธมฺมสฺส อินฺทฺริยปจฺจเยน ปจฺจโย.\csromanspacer{} รูปชีวิตินฺทฺริยํ กฏตฺตารูปานํ อินฺทฺริยปจฺจเยน ปจฺจโย. (1)}

\prose{จิตฺตวิปฺปยุตฺโต ธมฺโม จิตฺตสมฺปยุตฺตสฺส ธมฺมสฺส อินฺทฺริยปจฺจเยน ปจฺจโย.\csromanspacer{} จกฺขุนฺทฺริยํ จกฺขุวิญฺญาณสหคตานํ ขนฺธานํ อินฺทฺริยปจฺจเยน ปจฺจโย ฯเปฯ.\csromanspacer{} กายินฺทฺริยํ. (2)}

\prose{จิตฺตสมฺปยุตฺโต จ จิตฺตวิปฺปยุตฺโต จ ธมฺมา จิตฺตสมฺปยุตฺตสฺส ธมฺมสฺส อินฺทฺริยปจฺจเยน ปจฺจโย.\csromanspacer{} จกฺขุนฺทฺริยญฺจ อุเปกฺขินฺทฺริยญฺจ จกฺขุวิญฺญาณสหคตานํ ขนฺธานํ อินฺทฺริยปจฺจเยน ปจฺจโย ฯเปฯ.\csromanspacer{} กายินฺทฺริยญฺจ สุขินฺทฺริยญฺจ ฯเปฯ.\csromanspacer{} กายินฺทฺริยญฺจ ทุกฺขินฺทฺริยญฺจ กายวิญฺญาณสหคตานํ ขนฺธานํ อินฺทฺริยปจฺจเยน ปจฺจโย. (1)}

\prose{ฌานปจฺจเยน ปจฺจโย.\csromanspacer{} ตีณิ.\csromanspacer{} มคฺคปจฺจเยน ปจฺจโย.\csromanspacer{} ตีณิ.\csromanspacer{} สมฺปยุตฺตปจฺจเยน ปจฺจโย.\csromanspacer{} เอกํ.}

\csromanheader{2}{\textbf{วิปฺปยุตฺต}}
\proseitem{175}{จิตฺตสมฺปยุตฺโต ธมฺโม จิตฺตวิปฺปยุตฺตสฺส ธมฺมสฺส วิปฺปยุตฺตปจฺจเยน ปจฺจโย.\csromanspacer{} \textbf{สหชาตํ}, ปจฺฉาชาตํ. (สํขิตฺตํ.) (1)}

\prose{จิตฺตวิปฺปยุตฺโต ธมฺโม จิตฺตสมฺปยุตฺตสฺส ธมฺมสฺส วิปฺปยุตฺตปจฺจเยน ปจฺจโย.\csromanspacer{} \textbf{สหชาตํ}, ปุเรชาตํ.\csromanspacer{}\csromanspacer{}\textbf{สหชาตํ} ปฏิสนฺธิกฺขเณ วตฺถุ จิตฺตสมฺปยุตฺตกานํ ขนฺธานํ วิปฺปยุตฺตปจฺจเยน ปจฺจโย.\csromanspacer{} \textbf{ปุเรชาตํ} จกฺขายตนํ จกฺขุวิญฺญาณสหคตานํ ขนฺธานํ วิปฺปยุตฺตปจฺจเยน ปจฺจโย ฯเปฯ.\csromanspacer{} กายายตนํ ฯเปฯ.\csromanspacer{} วตฺถุ จิตฺตสมฺปยุตฺตกานํ ขนฺธานํ วิปฺปยุตฺตปจฺจเยน ปจฺจโย. (1)}

\csromanheader{2}{\textbf{อตฺถิ}}
\proseitem{176}{\csromanfolio{436}จิตฺตสมฺปยุตฺโต ธมฺโม จิตฺตสมฺปยุตฺตสฺส ธมฺมสฺส อตฺถิปจฺจเยน ปจฺจโย.\csromanspacer{} เอกํ. (ปฏิจฺจสทิสํ.) (1)}

\prose{จิตฺตสมฺปยุตฺโต ธมฺโม จิตฺตวิปฺปยุตฺตสฺส ธมฺมสฺส อตฺถิปจฺจเยน ปจฺจโย.\csromanspacer{} \textbf{สหชาตํ}, ปจฺฉาชาตํ. (สํขิตฺตํ.) (2)}

\prose{จิตฺตสมฺปยุตฺโต ธมฺโม จิตฺตสมฺปยุตฺตสฺส จ จิตฺตวิปฺปยุตฺตสฺส จ ธมฺมสฺส อตฺถิปจฺจเยน ปจฺจโย. (ปฏิจฺจสทิสํ.) (3)}

\proseitem{177}{จิตฺตวิปฺปยุตฺโต ธมฺโม จิตฺตวิปฺปยุตฺตสฺส ธมฺมสฺส อตฺถิปจฺจเยน ปจฺจโย.\csromanspacer{} \textbf{สหชาตํ}, อาหารํ, อินฺทฺริยํ. (สํขิตฺตํ.) (1)}

\prose{จิตฺตวิปฺปยุตฺโต ธมฺโม จิตฺตสมฺปยุตฺตสฺส ธมฺมสฺส อตฺถิปจฺจเยน ปจฺจโย.\csromanspacer{} \textbf{สหชาตํ, ปุเรชาตํ}.\csromanspacer{}\csromanspacer{}\textbf{สหชาตํ} ปฏิสนฺธิกฺขเณ วตฺถุ จิตฺตสมฺปยุตฺตกานํ ขนฺธานํ อตฺถิปจฺจเยน ปจฺจโย.\csromanspacer{} \textbf{ปุเรชาตํ} จกฺขุํ ฯเปฯ วตฺถุํ อนิจฺจโต ฯเปฯ. (ปุเรชาตสทิสํ.) (2)}

\proseitem{178}{จิตฺตสมฺปยุตฺโต จ จิตฺตวิปฺปยุตฺโต จ ธมฺมา จิตฺตสมฺปยุตฺตสฺส ธมฺมสฺส อตฺถิปจฺจเยน ปจฺจโย.\csromanspacer{} \textbf{สหชาตํ, ปุเรชาตํ}.\csromanspacer{}\csromanspacer{}สหชาโต จกฺขุวิญฺญาณสหคโต เอโก ขนฺโธ จ จกฺขายตนญฺจ ทฺวินฺนํ ขนฺธานํ ฯเปฯ.\csromanspacer{} กายวิญฺญาณสหคโต เอโก ขนฺโธ จ กายายตนญฺจ ทฺวินฺนํ ขนฺธานํ ฯเปฯ.\csromanspacer{} จิตฺตสมฺปยุตฺโต เอโก ขนฺโธ จ วตฺถุ จ ทฺวินฺนํ ขนฺธานํ ฯเปฯ เทฺว ขนฺธา จ ฯเปฯ.\csromanspacer{} ปฏิสนฺธิกฺขเณ ฯเปฯ. (1)}

\prose{จิตฺตสมฺปยุตฺโต จ จิตฺตวิปฺปยุตฺโต จ ธมฺมา จิตฺตวิปฺปยุตฺตสฺส ธมฺมสฺส อตฺถิปจฺจเยน ปจฺจโย.\csromanspacer{} \textbf{สหชาตํ}, ปจฺฉาชาตํ, อาหารํ, อินฺทฺริยํ.\csromanspacer{}\csromanspacer{}\textbf{สหชาตา} จิตฺตสมฺปยุตฺตา ขนฺธา จ มหาภูตา จ จิตฺตสมุฏฺฐานานํ รูปานํ อตฺถิปจฺจเยน ปจฺจโย.\csromanspacer{} ปฏิสนฺธิกฺขเณ ฯเปฯ.\csromanspacer{} \textbf{ปจฺฉาชาตา} จิตฺตสมฺปยุตฺตกา ขนฺธา จ กพฬีกาโร อาหาโร จ อิมสฺส กายสฺส อตฺถิปจฺจเยน ปจฺจโย.\csromanspacer{} \textbf{ปจฺฉาชาตา} จิตฺตสมฺปยุตฺตกา ขนฺธา จ รูปชีวิตินฺทฺริยญฺจ กฏตฺตารูปานํ อตฺถิปจฺจเยน ปจฺจโย. (2)}

\csromanheader{2}{58. จิตฺตสมฺปยุตฺตทุก \textbf{1. ปจฺจยานุโลม 2. สงฺขฺยาวาร}}
\csromanheader{2}{\textbf{สุทฺธ}}
\proseitem{179}{\csromanfolio{437}เหตุยา ตีณิ, อารมฺมเณ เทฺว, อธิปติยา จตฺตาริ, อนนฺตเร เอกํ, สมนนฺตเร เอกํ, สหชาเต สตฺต, อญฺญมญฺเญ ฉ, นิสฺสเย สตฺต, อุปนิสฺสเย เทฺว, ปุเรชาเต เอกํ, ปจฺฉาชาเต เอกํ, อาเสวเน เอกํ, กมฺเม ตีณิ, วิปาเก ตีณิ, อาหาเร จตฺตาริ, อินฺทฺริเย ฉ, ฌาเน ตีณิ, มคฺเค ตีณิ, สมฺปยุตฺเต เอกํ, วิปฺปยุตฺเต เทฺว, อตฺถิยา สตฺต, นตฺถิยา เอกํ, วิคเต เอกํ, อวิคเต สตฺต.}

\vspace{\tipitakaparskip}
\csromancenter{อนุโลมํ.}
\csromanheader{2}{2. ปจฺจนียุทฺธาร}
\proseitem{180}{จิตฺตสมฺปยุตฺโต ธมฺโม จิตฺตสมฺปยุตฺตสฺส ธมฺมสฺส อารมฺมณปจฺจเยน ปจฺจโย.\csromanspacer{} สหชาตปจฺจเยน ปจฺจโย.\csromanspacer{} อุปนิสฺสยปจฺจเยน ปจฺจโย.\csromanspacer{} กมฺมปจฺจเยน ปจฺจโย. (1)}

\prose{จิตฺตสมฺปยุตฺโต ธมฺโม จิตฺตวิปฺปยุตฺตสฺส ธมฺมสฺส สหชาตปจฺจเยน ปจฺจโย.\csromanspacer{} ปจฺฉาชาตปจฺจเยน ปจฺจโย.\csromanspacer{} กมฺมปจฺจเยน ปจฺจโย. (2)}

\prose{จิตฺตสมฺปยุตฺโต ธมฺโม จิตฺตสมฺปยุตฺตสฺส จ จิตฺตวิปฺปยุตฺตสฺส จ ธมฺมสฺส สหชาตปจฺจเยน ปจฺจโย.\csromanspacer{} กมฺมปจฺจเยน ปจฺจโย. (3)}

\proseitem{181}{จิตฺตวิปฺปยุตฺโต ธมฺโม จิตฺตวิปฺปยุตฺตสฺส ธมฺมสฺส สหชาตปจฺจเยน ปจฺจโย.\csromanspacer{} อาหารปจฺจเยน ปจฺจโย.\csromanspacer{} อินฺทฺริยปจฺจเยน ปจฺจโย. (1)}

\prose{จิตฺตวิปฺปยุตฺโต ธมฺโม จิตฺตสมฺปยุตฺตสฺส ธมฺมสฺส อารมฺมณปจฺจเยน ปจฺจโย.\csromanspacer{} สหชาตปจฺจเยน ปจฺจโย.\csromanspacer{} อุปนิสฺสยปจฺจเยน ปจฺจโย.\csromanspacer{} ปุเรชาตปจฺจเยน ปจฺจโย. (2)}

\prose{จิตฺตสมฺปยุตฺโต จ จิตฺตวิปฺปยุตฺโต จ ธมฺมา จิตฺตสมฺปยุตฺตสฺส ธมฺมสฺส สหชาตํ, ปุเรชาตํ. (1)}

\prose{จิตฺตสมฺปยุตฺโต จ จิตฺตวิปฺปยุตฺโต จ ธมฺมา จิตฺตวิปฺปยุตฺตสฺส ธมฺมสฺส สหชาตํ, ปจฺฉาชาตํ, อาหารํ, อินฺทฺริยํ. (2)}

\csromanheader{2}{2. ปจฺจยปจฺจนีย 2. สงฺขฺยาวาร}
\csromanheader{2}{\textbf{สุทฺธ}}
\proseitem{182}{\csromanfolio{438}นเหตุยา สตฺต, น-อารมฺมเณ สตฺต, น-อธิปติยา สตฺต, น-อนนฺตเร สตฺต, นสมนนฺตเร สตฺต, นสหชาเต ฉ, น-อญฺญมญฺเญ ฉ, นนิสฺสเย ฉ, น-อุปนิสฺสเย สตฺต, นปุเรชาเต สตฺต, (สพฺพตฺถ สตฺต.) นสมฺปยุตฺเต ฉ, นวิปฺปยุตฺเต ปญฺจ, โน-อตฺถิยา จตฺตาริ, โนนตฺถิยา สตฺต, โนวิคเต สตฺต, โน-อวิคเต จตฺตาริ.}

\csromanheader{2}{3. ปจฺจยานุโลมปจฺจนีย}
\proseitem{183}{เหตุปจฺจยา น-อารมฺมเณ ตีณิ, น-อธิปติยา ตีณิ, น-อนนฺตเร ตีณิ, นสมนนฺตเร ตีณิ, น-อญฺญมญฺเญ เอกํ, น-อุปนิสฺสเย ตีณิ, (สพฺพตฺถ ตีณิ.) นสมฺปยุตฺเต เอกํ, นวิปฺปยุตฺเต เอกํ, โนนตฺถิยา ตีณิ, โนวิคเต ตีณิ.}

\csromanheader{2}{4. ปจฺจยปจฺจนียานุโลม}
\proseitem{184}{นเหตุปจฺจยา อารมฺมเณ เทฺว, อธิปติยา จตฺตาริ, (อนุโลมมาติกา กาตพฺพา.) ฯเปฯ อวิคเต สตฺต.}

\nitthitam{จิตฺตสมฺปยุตฺตทุกํ นิฏฺฐิตํ.}
\csromanmatikaanchor{mtk.58}
\csromantocmarkref{subsection}{59. จิตฺตสํสฏฺฐทุก}{mtk.58}
\bookmark[dest={mtk.58},level=2]{59. จิตฺตสํสฏฺฐทุก}
\csromanheader{2}{59. จิตฺตสํสฏฺฐทุก 1. ปฏิจฺจวาร}
\csromanheader{2}{1. ปจฺจยานุโลม 1. วิภงฺควาร}
\proseitem{185}{จิตฺตสํสฏฺฐํ ธมฺมํ ปฏิจฺจ จิตฺตสํสฏฺโฐ ธมฺโม อุปฺปชฺชติ เหตุปจฺจยา.\csromanspacer{} จิตฺตสํสฏฺฐํ เอกํ ขนฺธํ ปฏิจฺจ เทฺว ขนฺธา.\csromanspacer{} เทฺว ขนฺเธ ฯเปฯ.\csromanspacer{} ปฏิสนฺธิกฺขเณ ฯเปฯ. (1)}

\prose{จิตฺตสํสฏฺฐํ ธมฺมํ ปฏิจฺจ จิตฺตวิสํสฏฺโฐ ธมฺโม อุปฺปชฺชติ เหตุปจฺจยา.\csromanspacer{} จิตฺตสํสฏฺเฐ ขนฺเธ ปฏิจฺจ จิตฺตสมุฏฺฐานํ รูปํ.\csromanspacer{} ปฏิสนฺธิกฺขเณ ฯเปฯ. (2)}

\csromanmatikaanchor{mtk.59}
\csromantocmarkref{subsection}{60. จิตฺตสมุฏฺฐานทุก}{mtk.59}
\bookmark[dest={mtk.59},level=2]{60. จิตฺตสมุฏฺฐานทุก}
\csromanheader{2}{60. จิตฺตสมุฏฺฐานทุก}
\vspace{\tipitakaparskip}
\csromancenter{(จิตฺตสํสฏฺฐทุกํ ยถา จิตฺตสมฺปยุตฺตทุกํ, เอวํ กาตพฺพํ, นินฺนานากรณํ.)}
\nitthitam{จิตฺตสํสฏฺฐทุกํ นิฏฺฐิตํ.}
\csromanheader{2}{60. จิตฺตสมุฏฺฐานทุก 1. ปฏิจฺจวาร}
\csromanheader{2}{1. ปจฺจยานุโลม 1. วิภงฺควาร}
\csromanheader{2}{\textbf{เหตุ}}
\proseitem{186}{\csromanfolio{439}จิตฺตสมุฏฺฐานํ ธมฺมํ ปฏิจฺจ จิตฺตสมุฏฺฐาโน ธมฺโม อุปฺปชฺชติ เหตุปจฺจยา.\csromanspacer{} จิตฺตสมุฏฺฐานํ เอกํ ขนฺธํ ปฏิจฺจ เทฺว ขนฺธา จิตฺตสมุฏฺฐานญฺจ รูปํ.\csromanspacer{} เทฺว ขนฺเธ ฯเปฯ.\csromanspacer{} ปฏิสนฺธิกฺขเณ จิตฺตสมุฏฺฐานํ เอกํ ขนฺธํ ปฏิจฺจ เทฺว ขนฺธา.\csromanspacer{} เทฺว ขนฺเธ ฯเปฯ.\csromanspacer{} เอกํ มหาภูตํ ฯเปฯ มหาภูเต ปฏิจฺจ จิตฺตสมุฏฺฐานํ รูปํ, อุปาทารูปํ. (1)}

\prose{จิตฺตสมุฏฺฐานํ ธมฺมํ ปฏิจฺจ โนจิตฺตสมุฏฺฐาโน ธมฺโม อุปฺปชฺชติ เหตุปจฺจยา.\csromanspacer{} จิตฺตสมุฏฺฐาเน ขนฺเธ ปฏิจฺจ จิตฺตํ.\csromanspacer{} ปฏิสนฺธิกฺขเณ จิตฺตสมุฏฺฐาเน ขนฺเธ ปฏิจฺจ จิตฺตํ กฏตฺตา จ รูปํ.}

\prose{จิตฺตสมุฏฺฐานํ ธมฺมํ ปฏิจฺจ จิตฺตสมุฏฺฐาโน จ โนจิตฺตสมุฏฺฐาโน จ ธมฺมา อุปฺปชฺชนฺติ เหตุปจฺจยา.\csromanspacer{} จิตฺตสมุฏฺฐานํ เอกํ ขนฺธํ ปฏิจฺจ เทฺว ขนฺธา จิตฺตญฺจ จิตฺตสมุฏฺฐานญฺจ รูปํ.\csromanspacer{} เทฺว ขนฺเธ ฯเปฯ.\csromanspacer{} ปฏิสนฺธิกฺขเณ จิตฺตสมุฏฺฐานํ เอกํ ขนฺธํ ปฏิจฺจ เทฺว ขนฺธา จิตฺตญฺจ กฏตฺตา จ รูปํ.\csromanspacer{} เทฺว ขนฺเธ ฯเปฯ. (3)}

\proseitem{187}{โนจิตฺตสมุฏฺฐานํ ธมฺมํ ปฏิจฺจ โนจิตฺตสมุฏฺฐาโน ธมฺโม อุปฺปชฺชติ เหตุปจฺจยา.\csromanspacer{} ปฏิสนฺธิกฺขเณ จิตฺตํ ปฏิจฺจ กฏตฺตารูปํ.\csromanspacer{} จิตฺตํ ปฏิจฺจ วตฺถุ, วตฺถุ ปฏิจฺจ จิตฺตํ.\csromanspacer{} เอกํ มหาภูตํ ฯเปฯ มหาภูเต ปฏิจฺจ กฏตฺตารูปํ, อุปาทารูปํ. (1)}

\prose{โนจิตฺตสมุฏฺฐานํ ธมฺมํ ปฏิจฺจ จิตฺตสมุฏฺฐาโน ธมฺโม อุปฺปชฺชติ เหตุปจฺจยา.\csromanspacer{} จิตฺตํ ปฏิจฺจ สมฺปยุตฺตกา ขนฺธา จิตฺตสมุฏฺฐานญฺจ รูปํ.\csromanspacer{} ปฏิสนฺธิกฺขเณ จิตฺตํ ปฏิจฺจ สมฺปยุตฺตกา ขนฺธา.\csromanspacer{} ปฏิสนฺธิกฺขเณ วตฺถุํ ปฏิจฺจ จิตฺตสมุฏฺฐานา ขนฺธา. (2)}

\prose{\csromanfolio{440}โนจิตฺตสมุฏฺฐานํ ธมฺมํ ปฏิจฺจ จิตฺตสมุฏฺฐาโน จ โนจิตฺตสมุฏฺฐาโน จ ธมฺมา อุปฺปชฺชนฺติ เหตุปจฺจยา.\csromanspacer{} ปฏิสนฺธิกฺขเณ จิตฺตํ ปฏิจฺจ สมฺปยุตฺตกา ขนฺธา กฏตฺตา จ รูปํ.\csromanspacer{} ปฏิสนฺธิกฺขเณ วตฺถุํ ปฏิจฺจ จิตฺตํ สมฺปยุตฺตกา จ ขนฺธา. (3)}

\proseitem{188}{จิตฺตสมุฏฺฐานญฺจ โนจิตฺตสมุฏฺฐานญฺจ ธมฺมํ ปฏิจฺจ จิตฺตสมุฏฺฐาโน ธมฺโม อุปฺปชฺชติ เหตุปจฺจยา.\csromanspacer{} จิตฺตสมุฏฺฐานํ เอกํ ขนฺธญฺจ จิตฺตญฺจ ปฏิจฺจ เทฺว ขนฺธา จิตฺตสมุฏฺฐานญฺจ รูปํ.\csromanspacer{} เทฺว ขนฺเธ จ ฯเปฯ.\csromanspacer{} ปฏิสนฺธิกฺขเณ จิตฺตสมุฏฺฐานํ เอกํ ขนฺธญฺจ จิตฺตญฺจ ปฏิจฺจ เทฺว ขนฺธา.\csromanspacer{} เทฺว ขนฺเธ จ ฯเปฯ.\csromanspacer{} ปฏิสนฺธิกฺขเณ จิตฺตสมุฏฺฐานํ เอกํ ขนฺธญฺจ วตฺถุญฺจ ปฏิจฺจ เทฺว ขนฺธา.\csromanspacer{} เทฺว ขนฺเธ จ ฯเปฯ. (1)}

\prose{จิตฺตสมุฏฺฐานญฺจ โนจิตฺตสมุฏฺฐานญฺจ ธมฺมํ ปฏิจฺจ โนจิตฺตสมุฏฺฐาโน ธมฺโม อุปฺปชฺชติ เหตุปจฺจยา.\csromanspacer{} ปฏิสนฺธิกฺขเณ จิตฺตสมุฏฺฐาเน ขนฺเธ จ จิตฺตญฺจ ปฏิจฺจ กฏตฺตารูปํ.\csromanspacer{} ปฏิสนฺธิกฺขเณ จิตฺตสมุฏฺฐาเน ขนฺเธ จ มหาภูเต จ ปฏิจฺจ กฏตฺตารูปํ.\csromanspacer{} ปฏิสนฺธิกฺขเณ จิตฺตสมุฏฺฐาเน ขนฺเธ จ วตฺถุญฺจ ปฏิจฺจ จิตฺตํ. (2)}

\prose{จิตฺตสมุฏฺฐานญฺจ โนจิตฺตสมุฏฺฐานญฺจ ธมฺมํ ปฏิจฺจ จิตฺตสมุฏฺฐาโน จ โนจิตฺตสมุฏฺฐาโน จ ธมฺมา อุปฺปชฺชนฺติ เหตุปจฺจยา.\csromanspacer{} ปฏิสนฺธิกฺขเณ จิตฺตสมุฏฺฐานํ เอกํ ขนฺธญฺจ จิตฺตญฺจ ปฏิจฺจ เทฺว ขนฺธา กฏตฺตา จ รูปํ.\csromanspacer{} เทฺว ขนฺเธ จ ฯเปฯ.\csromanspacer{} ปฏิสนฺธิกฺขเณ จิตฺตสมุฏฺฐานํ เอกํ ขนฺธญฺจ วตฺถุญฺจ ปฏิจฺจ เทฺว ขนฺธา จิตฺตญฺจ.\csromanspacer{} เทฺว ขนฺเธ จ ฯเปฯ. (3)}

\csromanheader{2}{\textbf{อารมฺมณ}}
\proseitem{189}{จิตฺตสมุฏฺฐานํ ธมฺมํ ปฏิจฺจ จิตฺตสมุฏฺฐาโน ธมฺโม อุปฺปชฺชติ อารมฺมณปจฺจยา.\csromanspacer{} จิตฺตสมุฏฺฐานํ เอกํ ขนฺธํ ปฏิจฺจ เทฺว ขนฺธา.\csromanspacer{} เทฺว ขนฺเธ ฯเปฯ.\csromanspacer{} ปฏิสนฺธิกฺขเณ ฯเปฯ. (1)}

\prose{จิตฺตสมุฏฺฐานํ ธมฺมํ ปฏิจฺจ โนจิตฺตสมุฏฺฐาโน ธมฺโม อุปฺปชฺชติ อารมฺมณปจฺจยา.\csromanspacer{} จิตฺตสมุฏฺฐาเน ขนฺเธ ปฏิจฺจ จิตฺตํ.\csromanspacer{} ปฏิสนฺธิกฺขเณ ฯเปฯ. (2)}

\prose{จิตฺตสมุฏฺฐานํ ธมฺมํ ปฏิจฺจ จิตฺตสมุฏฺฐาโน จ โนจิตฺตสมุฏฺฐาโน จ ธมฺมา อุปฺปชฺชนฺติ อารมฺมณปจฺจยา.\csromanspacer{} จิตฺตสมุฏฺฐานํ เอกํ ขนฺธํ ปฏิจฺจ เทฺว ขนฺธา จิตฺตญฺจ.\csromanspacer{} เทฺว ขนฺเธ ฯเปฯ.\csromanspacer{} ปฏิสนฺธิกฺขเณ ฯเปฯ. (3)}

\csromanheader{2}{60. จิตฺตสมุฏฺฐานทุก}
\proseitem{190}{\csromanfolio{441}โนจิตฺตสมุฏฺฐานํ ธมฺมํ ปฏิจฺจ โนจิตฺตสมุฏฺฐาโน ธมฺโม อุปฺปชฺชติ อารมฺมณปจฺจยา.\csromanspacer{} ปฏิสนฺธิกฺขเณ วตฺถุํ ปฏิจฺจ จิตฺตํ. (1)}

\prose{โนจิตฺตสมุฏฺฐานํ ธมฺมํ ปฏิจฺจ จิตฺตสมุฏฺฐาโน ธมฺโม อุปฺปชฺชติ อารมฺมณปจฺจยา.\csromanspacer{} จิตฺตํ ปฏิจฺจ สมฺปยุตฺตกา ขนฺธา.\csromanspacer{} ปฏิสนฺธิกฺขเณ จิตฺตํ ปฏิจฺจ สมฺปยุตฺตกา ขนฺธา.\csromanspacer{} ปฏิสนฺธิกฺขเณ วตฺถุํ ปฏิจฺจ จิตฺตสมุฏฺฐานา ขนฺธา. (2)}

\prose{โนจิตฺตสมุฏฺฐานํ ธมฺมํ ปฏิจฺจ จิตฺตสมุฏฺฐาโน จ โนจิตฺตสมุฏฺฐาโน จ ธมฺมา อุปฺปชฺชนฺติ อารมฺมณปจฺจยา.\csromanspacer{} ปฏิสนฺธิกฺขเณ วตฺถุํ ปฏิจฺจ จิตฺตํ สมฺปยุตฺตกา จ ขนฺธา. (3)}

\proseitem{191}{จิตฺตสมุฏฺฐานญฺจ โนจิตฺตสมุฏฺฐานญฺจ ธมฺมํ ปฏิจฺจ จิตฺตสมุฏฺฐาโน ธมฺโม อุปฺปชฺชติ อารมฺมณปจฺจยา.\csromanspacer{} จิตฺตสมุฏฺฐานํ เอกํ ขนฺธญฺจ จิตฺตญฺจ ปฏิจฺจ เทฺว ขนฺธา.\csromanspacer{} เทฺว ขนฺเธ ฯเปฯ.\csromanspacer{} ปฏิสนฺธิกฺขเณ จิตฺตสมุฏฺฐานํ เอกํ ขนฺธญฺจ จิตฺตญฺจ ปฏิจฺจ เทฺว ขนฺธา.\csromanspacer{} เทฺว ขนฺเธ ฯเปฯ.\csromanspacer{} ปฏิสนฺธิกฺขเณ จิตฺตสมุฏฺฐานํ เอกํ ขนฺธญฺจ วตฺถุญฺจ ปฏิจฺจ เทฺว ขนฺธา.\csromanspacer{} เทฺว ขนฺเธ ฯเปฯ. (1)}

\prose{จิตฺตสมุฏฺฐานญฺจ โนจิตฺตสมุฏฺฐานญฺจ ธมฺมํ ปฏิจฺจ โนจิตฺตสมุฏฺฐาโน ธมฺโม อุปฺปชฺชติ อารมฺมณปจฺจยา.\csromanspacer{} ปฏิสนฺธิกฺขเณ จิตฺตสมุฏฺฐาเน ขนฺเธ จ วตฺถุญฺจ ปฏิจฺจ จิตฺตํ. (2)}

\prose{จิตฺตสมุฏฺฐานญฺจ โนจิตฺตสมุฏฺฐานญฺจ ธมฺมํ ปฏิจฺจ จิตฺตสมุฏฺฐาโน จ โนจิตฺตสมุฏฺฐาโน จ ธมฺมา อุปฺปชฺชนฺติ อารมฺมณปจฺจยา.\csromanspacer{} ปฏิสนฺธิกฺขเณ จิตฺตสมุฏฺฐานํ เอกํ ขนฺธญฺจ วตฺถุญฺจ ปฏิจฺจ เทฺว ขนฺธา จิตฺตญฺจ.\csromanspacer{} เทฺว ขนฺเธ ฯเปฯ. (3)}

\csromanheader{2}{1. ปจฺจยานุโลม 2. สงฺขฺยาวาร}
\csromanheader{2}{\textbf{สุทฺธ}}
\proseitem{192}{เหตุยา นว, อารมฺมเณ นว, อธิปติยา ปญฺจ, อนนฺตเร นว, สมนนฺตเร นว, สหชาเต นว, อญฺญมญฺเญ นว, นิสฺสเย นว, อุปนิสฺสเย นว, ปุเรชาเต ปญฺจ, อาเสวเน ปญฺจ, กมฺเม นว, วิปาเก นว, อาหาเร นว, อินฺทฺริเย นว, ฌาเน นว, มคฺเค นว, สมฺปยุตฺเต นว, (สพฺพตฺถ นว.) อวิคเต นว.}

\vspace{\tipitakaparskip}
\csromancenter{อนุโลมํ.}
\csromanheader{2}{2. ปจฺจยปจฺจนีย 1. วิภงฺควาร}
\csromanheader{2}{\textbf{นเหตุ}}
\proseitem{193}{\csromanfolio{442}จิตฺตสมุฏฺฐานํ ธมฺมํ ปฏิจฺจ จิตฺตสมุฏฺฐาโน ธมฺโม อุปฺปชฺชติ นเหตุปจฺจยา.\csromanspacer{} อเหตุกํ จิตฺตสมุฏฺฐานํ เอกํ ขนฺธํ ปฏิจฺจ เทฺว ขนฺธา จิตฺตสมุฏฺฐานญฺจ รูปํ.\csromanspacer{} เทฺว ขนฺเธ ฯเปฯ.\csromanspacer{} อเหตุกปฏิสนฺธิกฺขเณ จิตฺตสมุฏฺฐานํ เอกํ ขนฺธํ ปฏิจฺจ เทฺว ขนฺธา.\csromanspacer{} เทฺว ขนฺเธ ฯเปฯ.\csromanspacer{} จิตฺตสมุฏฺฐานํ เอกํ มหาภูตํ ฯเปฯ.\csromanspacer{} มหาภูเต ปฏิจฺจ จิตฺตสมุฏฺฐานํ รูปํ, อุปาทารูปํ.\csromanspacer{} วิจิกิจฺฉาสหคเต อุทฺธจฺจสหคเต ขนฺเธ ปฏิจฺจ วิจิกิจฺฉาสหคโต อุทฺธจฺจสหคโต โมโห. (1)}

\prose{จิตฺตสมุฏฺฐานํ ธมฺมํ ปฏิจฺจ โนจิตฺตสมุฏฺฐาโน ธมฺโม อุปฺปชฺชติ นเหตุปจฺจยา.\csromanspacer{} อเหตุเก จิตฺตสมุฏฺฐาเน ขนฺเธ ปฏิจฺจ จิตฺตํ.\csromanspacer{} อเหตุกปฏิสนฺธิกฺขเณ จิตฺตสมุฏฺฐาเน ขนฺเธ ปฏิจฺจ จิตฺตํ กฏตฺตา จ รูปํ. (2)}

\prose{จิตฺตสมุฏฺฐานํ ธมฺมํ ปฏิจฺจ จิตฺตสมุฏฺฐาโน จ โนจิตฺตสมุฏฺฐาโน จ ธมฺมา อุปฺปชฺชนฺติ นเหตุปจฺจยา.\csromanspacer{} อเหตุกํ จิตฺตสมุฏฺฐานํ เอกํ ขนฺธํ ปฏิจฺจ เทฺว ขนฺธา จิตฺตญฺจ จิตฺตสมุฏฺฐานญฺจ รูปํ.\csromanspacer{} เทฺว ขนฺเธ ฯเปฯ.\csromanspacer{} อเหตุกปฏิสนฺธิกฺขเณ ฯเปฯ. (3)}

\proseitem{194}{โนจิตฺตสมุฏฺฐานํ ธมฺมํ ปฏิจฺจ โนจิตฺตสมุฏฺฐาโน ธมฺโม อุปฺปชฺชติ นเหตุปจฺจยา.\csromanspacer{} อเหตุกปฏิสนฺธิกฺขเณ จิตฺตํ ปฏิจฺจ กฏตฺตารูปํ.\csromanspacer{} จิตฺตํ ปฏิจฺจ วตฺถุ, วตฺถุํ ปฏิจฺจ จิตฺตํ.\csromanspacer{} เอกํ มหาภูตํ ฯเปฯ. (ยาว อสญฺญสตฺตา กาตพฺพา.) (1)}

\prose{โนจิตฺตสมุฏฺฐานํ ธมฺมํ ปฏิจฺจ จิตฺตสมุฏฺฐาโน ธมฺโม อุปฺปชฺชติ นเหตุปจฺจยา.\csromanspacer{} อเหตุกํ จิตฺตํ ปฏิจฺจ สมฺปยุตฺตกา ขนฺธา จิตฺตสมุฏฺฐานญฺจ รูปํ.\csromanspacer{} อเหตุกปฏิสนฺธิกฺขเณ จิตฺตํ ฯเปฯ.\csromanspacer{} อเหตุกปฏิสนฺธิกฺขเณ วตฺถุํ ปฏิจฺจ จิตฺตสมุฏฺฐานา ขนฺธา.\csromanspacer{} วิจิกิจฺฉาสหคเต อุทฺธจฺจสหคเต ขนฺเธ ปฏิจฺจ วิจิกิจฺฉาสหคโต อุทฺธจฺจสหคโต โมโห. (2)}

\prose{โนจิตฺตสมุฏฺฐานํ ธมฺมํ ปฏิจฺจ จิตฺตสมุฏฺฐาโน จ โนจิตฺตสมุฏฺฐาโน จ ธมฺมา อุปฺปชฺชนฺติ นเหตุปจฺจยา.\csromanspacer{} อเหตุกปฏิสนฺธิกฺขเณ จิตฺตํ ปฏิจฺจ สมฺปยุตฺตกา ขนฺธา กฏตฺตา จ รูปํ.\csromanspacer{} อเหตุกปฏิสนฺธิกฺขเณ วตฺถุํ ปฏิจฺจ จิตฺตํ สมฺปยุตฺตกา ขนฺธา จ. (3)}

\csromanheader{2}{60. จิตฺตสมุฏฺฐานทุก}
\proseitem{195}{\csromanfolio{443}จิตฺตสมุฏฺฐานญฺจ โนจิตฺตสมุฏฺฐานญฺจ ธมฺมํ ปฏิจฺจ จิตฺตสมุฏฺฐาโน ธมฺโม อุปฺปชฺชติ นเหตุปจฺจยา.\csromanspacer{} อเหตุกํ จิตฺตสมุฏฺฐานํ เอกํ ขนฺธญฺจ จิตฺตญฺจ ปฏิจฺจ เทฺว ขนฺธา จิตฺตสมุฏฺฐานญฺจ รูปํ.\csromanspacer{} เทฺว ขนฺเธ ฯเปฯ.\csromanspacer{} อเหตุกปฏิสนฺธิกฺขเณ จิตฺตสมุฏฺฐานํ เอกํ ขนฺธญฺจ จิตฺตญฺจ ปฏิจฺจ เทฺว ขนฺธา.\csromanspacer{} เทฺว ขนฺเธ ฯเปฯ.\csromanspacer{} อเหตุกปฏิสนฺธิกฺขเณ จิตฺตสมุฏฺฐานํ เอกํ ขนฺธญฺจ วตฺถุญฺจ ปฏิจฺจ เทฺว ขนฺธา.\csromanspacer{} เทฺว ขนฺเธ ฯเปฯ.\csromanspacer{} วิจิกิจฺฉาสหคเต อุทฺธจฺจสหคเต ขนฺเธ จ จิตฺตญฺจ ปฏิจฺจ วิจิกิจฺฉาสหคโต อุทฺธจฺจสหคโต โมโห. (1)}

\prose{จิตฺตสมุฏฺฐานญฺจ โนจิตฺตสมุฏฺฐานญฺจ ธมฺมํ ปฏิจฺจ โนจิตฺตสมุฏฺฐาโน ธมฺโม อุปฺปชฺชติ นเหตุปจฺจยา.\csromanspacer{} อเหตุกปฏิสนฺธิกฺขเณ จิตฺตสมุฏฺฐาเน ขนฺเธ จ จิตฺตญฺจ ปฏิจฺจ กฏตฺตารูปํ.\csromanspacer{} อเหตุกปฏิสนฺธิกฺขเณ จิตฺตสมุฏฺฐาเน ขนฺเธ จ มหาภูเต จ ปฏิจฺจ กฏตฺตารูปํ.\csromanspacer{} อเหตุกปฏิสนฺธิกฺขเณ จิตฺตสมุฏฺฐาเน ขนฺเธ จ วตฺถุญฺจ ปฏิจฺจ จิตฺตํ. (2)}

\prose{จิตฺตสมุฏฺฐานญฺจ โนจิตฺตสมุฏฺฐานญฺจ ธมฺมํ ปฏิจฺจ จิตฺตสมุฏฺฐาโน จ โนจิตฺตสมุฏฺฐาโน จ ธมฺมา อุปฺปชฺชนฺติ นเหตุปจฺจยา.\csromanspacer{} อเหตุกปฏิสนฺธิกฺขเณ จิตฺตสมุฏฺฐานํ เอกํ ขนฺธญฺจ จิตฺตญฺจ ปฏิจฺจ เทฺว ขนฺธา กฏตฺตา จ รูปํ.\csromanspacer{} เทฺว ขนฺเธ ฯเปฯ.\csromanspacer{} อเหตุกปฏิสนฺธิกฺขเณ จิตฺตสมุฏฺฐานํ เอกํ ขนฺธญฺจ วตฺถุญฺจ ปฏิจฺจ เทฺว ขนฺธา จิตฺตญฺจ.\csromanspacer{} เทฺว ขนฺเธ ฯเปฯ. (3)}

\csromanheader{2}{\textbf{น-อารมฺมณ}}
\proseitem{196}{จิตฺตสมุฏฺฐานํ ธมฺมํ ปฏิจฺจ จิตฺตสมุฏฺฐาโน ธมฺโม อุปฺปชฺชติ น-อารมฺมณปจฺจยา.\csromanspacer{} จิตฺตสมุฏฺฐาเน ขนฺเธ ปฏิจฺจ จิตฺตสมุฏฺฐานํ รูปํ.\csromanspacer{} จิตฺตสมุฏฺฐานํ เอกํ มหาภูตํ ฯเปฯ.\csromanspacer{} มหาภูเต ปฏิจฺจ จิตฺตสมุฏฺฐานํ รูปํ, อุปาทารูปํ. (1)}

\prose{จิตฺตสมุฏฺฐานํ ธมฺมํ ปฏิจฺจ โนจิตฺตสมุฏฺฐาโน ธมฺโม อุปฺปชฺชติ น-อารมฺมณปจฺจยา.\csromanspacer{} ปฏิสนฺธิกฺขเณ จิตฺตสมุฏฺฐาเน ขนฺเธ ปฏิจฺจ กฏตฺตารูปํ. (2)}

\proseitem{197}{โนจิตฺตสมุฏฺฐานํ ธมฺมํ ปฏิจฺจ โนจิตฺตสมุฏฺฐาโน ธมฺโม อุปฺปชฺชติ น-อารมฺมณปจฺจยา.\csromanspacer{} ปฏิสนฺธิกฺขเณ จิตฺตํ ปฏิจฺจ กฏตฺตารูปํ.\csromanspacer{} จิตฺตํ ปฏิจฺจ วตฺถุ.\csromanspacer{} เอกํ มหาภูตํ ฯเปฯ. (ยาว อสญฺญสตฺตา.) (1)}

\prose{โนจิตฺตสมุฏฺฐานํ ธมฺมํ ปฏิจฺจ จิตฺตสมุฏฺฐาโน ธมฺโม อุปฺปชฺชติ น-อารมฺมณปจฺจยา.\csromanspacer{} จิตฺตํ ปฏิจฺจ จิตฺตสมุฏฺฐานํ รูปํ. (2)}

\prose{\csromanfolio{444}จิตฺตสมุฏฺฐานญฺจ โนจิตฺตสมุฏฺฐานญฺจ ธมฺมํ ปฏิจฺจ จิตฺตสมุฏฺฐาโน ธมฺโม อุปฺปชฺชติ น-อารมฺมณปจฺจยา.\csromanspacer{} จิตฺตสมุฏฺฐาเน ขนฺเธ จ จิตฺตญฺจ ปฏิจฺจ จิตฺตสมุฏฺฐานํ รูปํ.\csromanspacer{} จิตฺตญฺจ มหาภูเต จ ปฏิจฺจ จิตฺตสมุฏฺฐานํ รูปํ. (1)}

\prose{จิตฺตสมุฏฺฐานญฺจ โนจิตฺตสมุฏฺฐานญฺจ ธมฺมํ ปฏิจฺจ โนจิตฺตสมุฏฺฐาโน ธมฺโม อุปฺปชฺชติ น-อารมฺมณปจฺจยา.\csromanspacer{} ปฏิสนฺธิกฺขเณ จิตฺตสมุฏฺฐาเน ขนฺเธ จ จิตฺตญฺจ ปฏิจฺจ กฏตฺตารูปํ.\csromanspacer{} ปฏิสนฺธิกฺขเณ จิตฺตสมุฏฺฐาเน ขนฺเธ จ มหาภูเต จ ปฏิจฺจ กฏตฺตารูปํ. (สํขิตฺตํ.) (2)}

\csromanheader{2}{2. ปจฺจยปจฺจนีย 2. สงฺขฺยาวาร}
\csromanheader{2}{\textbf{สุทฺธ}}
\proseitem{198}{นเหตุยา นว, น-อารมฺมเณ ฉ, น-อธิปติยา นว, น-อนนฺตเร ฉ, นสมนนฺตเร ฉ, น-อญฺญมญฺเญ ฉ, น-อุปนิสฺสเย ฉ, นปุเรชาเต นว, นปจฺฉาชาเต นว, น-อาเสวเน นว, นกมฺเม จตฺตาริ, นวิปาเก ปญฺจ, น-อาหาเร เอกํ, น-อินฺทฺริเย เอกํ, นฌาเน ฉ, นมคฺเค นว, นสมฺปยุตฺเต ฉ, นวิปฺปยุตฺเต ฉ, โนนตฺถิยา ฉ, โนวิคเต ฉ.}

\csromanheader{2}{3. ปจฺจยานุโลมปจฺจนีย}
\proseitem{199}{เหตุปจฺจยา น-อารมฺมเณ ฉ, น-อธิปติยา นว ฯเปฯ นกมฺเม ตีณิ, นวิปาเก ปญฺจ, นสมฺปยุตฺเต ฉ, นวิปฺปยุตฺเต ปญฺจ, โนนตฺถิยา ฉ, โนวิคเต ฉ.}

\csromanheader{2}{4. ปจฺจยปจฺจนียานุโลม}
\proseitem{200}{นเหตุปจฺจยา อารมฺมเณ นว, อนนฺตเร นว, สมนนฺตเร นว ฯเปฯ ปุเรชาเต ปญฺจ, อาเสวเน ปญฺจ ฯเปฯ ฌาเน นว, มคฺเค ตีณิ ฯเปฯ อวิคเต นว.}

\vspace{\tipitakaparskip}
\csromancenter{(สหชาตวาโร ปฏิจฺจวารสทิโส.)}
\csromanheader{2}{60. จิตฺตสมุฏฺฐานทุก \textbf{60. จิตฺตสมุฏฺฐานทุก 3. ปจฺจยวาร}}
\csromanheader{2}{1. ปจฺจยานุโลม 1. วิภงฺควาร}
\csromanheader{2}{\textbf{เหตุ}}
\proseitem{201}{\csromanfolio{445}จิตฺตสมุฏฺฐานํ ธมฺมํ ปจฺจยา จิตฺตสมุฏฺฐาโน ธมฺโม อุปฺปชฺชติ เหตุปจฺจยา.\csromanspacer{} ตีณิ. (ปฏิจฺจวารสทิสํ.)}

\prose{โนจิตฺตสมุฏฺฐานํ ธมฺมํ ปจฺจยา โนจิตฺตสมุฏฺฐาโน ธมฺโม อุปฺปชฺชติ เหตุปจฺจยา.\csromanspacer{} วตฺถุํ ปจฺจยา จิตฺตํ.\csromanspacer{} ปฏิสนฺธิกฺขเณ ฯเปฯ. (ปฏิจฺจวารสทิสา.) (1)}

\prose{โนจิตฺตสมุฏฺฐานํ ธมฺมํ ปจฺจยา จิตฺตสมุฏฺฐาโน ธมฺโม อุปฺปชฺชติ เหตุปจฺจยา.\csromanspacer{} จิตฺตํ ปจฺจยา สมฺปยุตฺตกา ขนฺธา จิตฺตสมุฏฺฐานญฺจ รูปํ.\csromanspacer{} วตฺถุํ ปจฺจยา จิตฺตสมุฏฺฐานา ขนฺธา. (ปฏิสนฺธิกฺขเณ เทฺวปิ กาตพฺพา.) (2)}

\prose{โนจิตฺตสมุฏฺฐานํ ธมฺมํ ปจฺจยา จิตฺตสมุฏฺฐาโน จ โนจิตฺตสมุฏฺฐาโน จ ธมฺมา อุปฺปชฺชนฺติ เหตุปจฺจยา.\csromanspacer{} วตฺถุํ ปจฺจยา จิตฺตํ สมฺปยุตฺตกา จ ขนฺธา.\csromanspacer{} ปฏิสนฺธิกฺขเณ เทฺว. (ปฏิจฺจวารสทิสํ.) (3)}

\proseitem{202}{จิตฺตสมุฏฺฐานญฺจ โนจิตฺตสมุฏฺฐานญฺจ ธมฺมํ ปจฺจยา จิตฺตสมุฏฺฐาโน ธมฺโม อุปฺปชฺชติ เหตุปจฺจยา.\csromanspacer{} จิตฺตสมุฏฺฐานํ เอกํ ขนฺธญฺจ จิตฺตญฺจ ปจฺจยา เทฺว ขนฺธา.\csromanspacer{} เทฺว ขนฺเธ ฯเปฯ.\csromanspacer{} จิตฺตสมุฏฺฐานํ เอกํ ขนฺธญฺจ วตฺถุญฺจ ปจฺจยา เทฺว ขนฺธา.\csromanspacer{} เทฺว ขนฺเธ ฯเปฯ. (ปฏิสนฺธิกฺขเณ เทฺวปิ ปฏิจฺจวารสทิสา.) (1)}

\prose{จิตฺตสมุฏฺฐานญฺจ โนจิตฺตสมุฏฺฐานญฺจ ธมฺมํ ปจฺจยา โนจิตฺตสมุฏฺฐาโน ธมฺโม อุปฺปชฺชติ เหตุปจฺจยา.\csromanspacer{} จิตฺตสมุฏฺฐาเน ขนฺเธ จ วตฺถุญฺจ ปจฺจยา จิตฺตํ. (ปฏิสนฺธิกฺขเณ ตีณิปิ กาตพฺพา.\csromanspacer{} ปฏิจฺจวารสทิสา.) (2)}

\prose{จิตฺตสมุฏฺฐานญฺจ โนจิตฺตสมุฏฺฐานญฺจ ธมฺมํ ปจฺจยา จิตฺตสมุฏฺฐาโน จ โนจิตฺตสมุฏฺฐาโน จ ธมฺมา อุปฺปชฺชนฺติ เหตุปจฺจยา.\csromanspacer{} จิตฺตสมุฏฺฐานํ เอกํ ขนฺธญฺจ วตฺถุญฺจ ปจฺจยา เทฺว ขนฺธา จิตฺตญฺจ.\csromanspacer{} เทฺว ขนฺเธ ฯเปฯ. (ปฏิสนฺธิกฺขเณ เทฺวปิ กาตพฺพา.\csromanspacer{} ปฏิจฺจวารสทิสา.) (3)}

\csromanheader{2}{\textbf{อารมฺมณ}}
\proseitem{203}{จิตฺตสมุฏฺฐานํ ธมฺมํ ปจฺจยา จิตฺตสมุฏฺฐาโน ธมฺโม อุปฺปชฺชติ อารมฺมณปจฺจยา.\csromanspacer{} ตีณิ. (ปฏิจฺจวารสทิสา.)}

\prose{\csromanfolio{446}โนจิตฺตสมุฏฺฐานํ ธมฺมํ ปจฺจยา โนจิตฺตสมุฏฺฐาโน ธมฺโม อุปฺปชฺชติ อารมฺมณปจฺจยา.\csromanspacer{} จกฺขายตนํ ปจฺจยา จกฺขุวิญฺญาณํ ฯเปฯ.\csromanspacer{} กายายตนํ ปจฺจยา กายวิญฺญาณํ ฯเปฯ.\csromanspacer{} วตฺถุํ ปจฺจยา จิตฺตํ. (1)}

\prose{โนจิตฺตสมุฏฺฐานํ ธมฺมํ ปจฺจยา จิตฺตสมุฏฺฐาโน ธมฺโม อุปฺปชฺชติ อารมฺมณปจฺจยา.\csromanspacer{} จกฺขายตนํ ปจฺจยา จกฺขุวิญฺญาณสหคตา ขนฺธา ฯเปฯ.\csromanspacer{} กายายตนํ ปจฺจยา กายวิญฺญาณสหคตา ขนฺธา.\csromanspacer{} จิตฺตํ ปจฺจยา สมฺปยุตฺตกา ขนฺธา.\csromanspacer{} วตฺถุํ ปจฺจยา จิตฺตสมุฏฺฐานา ขนฺธา. (ปฏิสนฺธิกฺขเณ เทฺวปิ.) (2)}

\prose{โนจิตฺตสมุฏฺฐานํ ธมฺมํ ปจฺจยา จิตฺตสมุฏฺฐาโน จ โนจิตฺตสมุฏฺฐาโน จ ธมฺมา อุปฺปชฺชนฺติ อารมฺมณปจฺจยา.\csromanspacer{} จกฺขายตนํ ปจฺจยา จกฺขุวิญฺญาณํ สมฺปยุตฺตกา จ ขนฺธา ฯเปฯ.\csromanspacer{} กายายตนํ ฯเปฯ.\csromanspacer{} วตฺถุํ ปจฺจยา จิตฺตํ สมฺปยุตฺตกา จ ขนฺธา.\csromanspacer{} ปฏิสนฺธิกฺขเณ ฯเปฯ. (3)}

\proseitem{204}{จิตฺตสมุฏฺฐานญฺจ โนจิตฺตสมุฏฺฐานญฺจ ธมฺมํ ปจฺจยา จิตฺตสมุฏฺฐาโน ธมฺโม อุปฺปชฺชติ อารมฺมณปจฺจยา.\csromanspacer{} จิตฺตสมุฏฺฐานํ เอกํ ขนฺธญฺจ จิตฺตญฺจ ปจฺจยา เทฺว ขนฺธา.\csromanspacer{} เทฺว ขนฺเธ ฯเปฯ.\csromanspacer{} จิตฺตสมุฏฺฐานํ เอกํ ขนฺธญฺจ วตฺถุญฺจ ปจฺจยา เทฺว ขนฺธา.\csromanspacer{} เทฺว ขนฺเธ ฯเปฯ. (ปฏิสนฺธิกฺขเณ เทฺวปิ กาตพฺพา.) (1)}

\prose{จิตฺตสมุฏฺฐานญฺจ โนจิตฺตสมุฏฺฐานญฺจ ธมฺมํ ปจฺจยา โนจิตฺตสมุฏฺฐาโน ธมฺโม อุปฺปชฺชติ อารมฺมณปจฺจยา.\csromanspacer{} จกฺขุวิญฺญาณสหคเต ขนฺเธ จ จกฺขายตนญฺจ ปจฺจยา จกฺขุวิญฺญาณํ ฯเปฯ.\csromanspacer{} กายวิญฺญาณสหคเต ฯเปฯ.\csromanspacer{} จิตฺตสมุฏฺฐาเน ขนฺเธ จ วตฺถุญฺจ ปจฺจยา จิตฺตํ.\csromanspacer{} ปฏิสนฺธิกฺขเณ ฯเปฯ. (2)}

\prose{จิตฺตสมุฏฺฐานญฺจ โนจิตฺตสมุฏฺฐานญฺจ ธมฺมํ ปจฺจยา จิตฺตสมุฏฺฐาโน จ โนจิตฺตสมุฏฺฐาโน จ ธมฺมา อุปฺปชฺชนฺติ อารมฺมณปจฺจยา.\csromanspacer{} จกฺขุวิญฺญาณสหคตํ เอกํ ขนฺธญฺจ จกฺขายตนญฺจ ปจฺจยา เทฺว ขนฺธา จกฺขุวิญฺญาณญฺจ.\csromanspacer{} เทฺว ขนฺเธ ฯเปฯ.\csromanspacer{} กายวิญฺญาณสหคตํ ฯเปฯ.\csromanspacer{} จิตฺตสมุฏฺฐานํ เอกํ ขนฺธญฺจ วตฺถุญฺจ ปจฺจยา เทฺว ขนฺธา จิตฺตญฺจ.\csromanspacer{} เทฺว ขนฺเธ ฯเปฯ.\csromanspacer{} ปฏิสนฺธิกฺขเณ ฯเปฯ. (3) (สํขิตฺตํ.)}

\csromanheader{2}{1. ปจฺจยานุโลม 2. สงฺขฺยาวาร}
\vspace{\tipitakaparskip}
\csromancenter{เหตุยา นว, อารมฺมเณ นว, อธิปติยา นว, (สพฺพตฺถ นว.) อวิคเต นว.}
\csromanheader{2}{60. จิตฺตสมุฏฺฐานทุก \textbf{2. ปจฺจยปจฺจนีย 1. วิภงฺควาร}}
\csromanheader{2}{\textbf{นเหตุ}}
\proseitem{206}{\csromanfolio{447}จิตฺตสมุฏฺฐานํ ธมฺมํ ปจฺจยา จิตฺตสมุฏฺฐาโน ธมฺโม อุปฺปชฺชติ นเหตุปจฺจยา. (สพฺเพ นวปิ ปญฺหา กาตพฺพา, ปฏิจฺจวารสทิสา.\csromanspacer{} ปญฺจวิญฺญาณมฺปิ กาตพฺพํ.\csromanspacer{} ตีณิเยว, โมโห.)}

\csromanheader{2}{2. ปจฺจยปจฺจนีย 2. สงฺขฺยาวาร}
\csromanheader{2}{\textbf{สุทฺธ}}
\vspace{\tipitakaparskip}
\csromancenter{\textbf{นเหตุ}ยา นว, น-อารมฺมเณ ฉ, น-อธิปติยา นว, น-อนนฺตเร ฉ, นสมนนฺตเร ฉ, น-อญฺญมญฺเญ ฉ, น-อุปนิสฺสเย ฉ, นปุเรชาเต นว, นปจฺฉาชาเต นว, น-อาเสวเน นว, นกมฺเม จตฺตาริ, นวิปาเก นว, น-อาหาเร เอกํ, น-อินฺทฺริเย เอกํ, นฌาเน นว, นมคฺเค นว, นสมฺปยุตฺเต ฉ, นวิปฺปยุตฺเต ฉ, โนนตฺถิยา ฉ, โนวิคเต ฉ.}
\csromancenter{(เอวํ อิตเร เทฺว คณนาปิ นิสฺสยวาโรปิ กาตพฺโพ.)}
\csromanheader{2}{60. จิตฺตสมุฏฺฐานทุก 5. สํสฏฺฐวาร}
\csromanheader{2}{1. ปจฺจยานุโลมาทิ}
\proseitem{208}{จิตฺตสมุฏฺฐานํ ธมฺมํ สํสฏฺโฐ จิตฺตสมุฏฺฐาโน ธมฺโม อุปฺปชฺชติ เหตุปจฺจยา.\csromanspacer{} จิตฺตสมุฏฺฐานํ เอกํ ขนฺธํ สํสฏฺฐา เทฺว ขนฺธา.\csromanspacer{} เทฺว ขนฺเธ ฯเปฯ.\csromanspacer{} ปฏิสนฺธิกฺขเณ ฯเปฯ. (1)}

\prose{จิตฺตสมุฏฺฐานํ ธมฺมํ สํสฏฺโฐ โนจิตฺตสมุฏฺฐาโน ธมฺโม อุปฺปชฺชติ เหตุปจฺจยา.\csromanspacer{} จิตฺตสมุฏฺฐาเน ขนฺเธ สํสฏฺฐํ จิตฺตํ.\csromanspacer{} ปฏิสนฺธิกฺขเณ ฯเปฯ. (2)}

\prose{จิตฺตสมุฏฺฐานํ ธมฺมํ สํสฏฺโฐ จิตฺตสมุฏฺฐาโน จ โนจิตฺตสมุฏฺฐาโน จ ธมฺมา อุปฺปชฺชนฺติ เหตุปจฺจยา.\csromanspacer{} จิตฺตสมุฏฺฐานํ เอกํ ขนฺธํ สํสฏฺฐา เทฺว ขนฺธา จิตฺตญฺจ.\csromanspacer{} เทฺว ขนฺเธ ฯเปฯ.\csromanspacer{} ปฏิสนฺธิกฺขเณ ฯเปฯ. (3)}

\proseitem{209}{\csromanfolio{448}โนจิตฺตสมุฏฺฐานํ ธมฺมํ สํสฏฺโฐ จิตฺตสมุฏฺฐาโน ธมฺโม อุปฺปชฺชติ เหตุปจฺจยา.\csromanspacer{} จิตฺตสํสฏฺฐา สมฺปยุตฺตกา ขนฺธา.\csromanspacer{} ปฏิสนฺธิกฺขเณ ฯเปฯ. (1)}

\prose{จิตฺตสมุฏฺฐานญฺจ โนจิตฺตสมุฏฺฐานญฺจ ธมฺมํ สํสฏฺโฐ จิตฺตสมุฏฺฐาโน ธมฺโม อุปฺปชฺชติ เหตุปจฺจยา.\csromanspacer{} จิตฺตสมุฏฺฐานํ เอกํ ขนฺธญฺจ จิตฺตญฺจ สํสฏฺฐา เทฺว ขนฺธา.\csromanspacer{} เทฺว ขนฺเธ ฯเปฯ.\csromanspacer{} ปฏิสนฺธิกฺขเณ ฯเปฯ. (สํขิตฺตํ.) (1)}

\prose{\textbf{เหตุ}ยา ปญฺจ, อารมฺมเณ ปญฺจ, (สพฺพตฺถ ปญฺจ.) อวิคเต ปญฺจ.}

\prose{จิตฺตสมุฏฺฐานํ ธมฺมํ สํสฏฺโฐ จิตฺตสมุฏฺฐาโน ธมฺโม อุปฺปชฺชติ นเหตุปจฺจยา. (ปญฺจ ปญฺหา กาตพฺพา.\csromanspacer{} ตีณิ.\csromanspacer{} โมโห.)}

\prose{นเหตุยา ปญฺจ, น-อธิปติยา ปญฺจ, นปุเรชาเต ปญฺจ, นปจฺฉาชาเต ปญฺจ, น-อาเสวเน ปญฺจ, นกมฺเม ตีณิ, นวิปาเก ปญฺจ, นฌาเน ปญฺจ, นมคฺเค ปญฺจ, นวิปฺปยุตฺเต ปญฺจ.}

\vspace{\tipitakaparskip}
\csromancenter{(เอวํ อิตเร เทฺว คณนาปิ สมฺปยุตฺตวาโรปิ กาตพฺโพ.)}
\csromanheader{2}{60. จิตฺตสมุฏฺฐานทุก 7. ปญฺหาวาร}
\csromanheader{2}{1. ปจฺจยานุโลม 1. วิภงฺควาร}
\csromanheader{2}{\textbf{เหตุ}}
\proseitem{210}{จิตฺตสมุฏฺฐาโน ธมฺโม จิตฺตสมุฏฺฐานสฺส ธมฺมสฺส เหตุปจฺจเยน ปจฺจโย.\csromanspacer{} จิตฺตสมุฏฺฐานา เหตู สมฺปยุตฺตกานํ ขนฺธานํ จิตฺตสมุฏฺฐานานญฺจ รูปานํ เหตุปจฺจเยน ปจฺจโย.\csromanspacer{} ปฏิสนฺธิกฺขเณ ฯเปฯ. (1)}

\prose{จิตฺตสมุฏฺฐาโน ธมฺโม โนจิตฺตสมุฏฺฐานสฺส ธมฺมสฺส เหตุปจฺจเยน ปจฺจโย.\csromanspacer{} จิตฺตสมุฏฺฐานา เหตู จิตฺตสฺส เหตุปจฺจเยน ปจฺจโย.\csromanspacer{} ปฏิสนฺธิกฺขเณ จิตฺตสมุฏฺฐานา เหตู จิตฺตสฺส กฏตฺตา จ รูปานํ เหตุปจฺจเยน ปจฺจโย. (2)}

\prose{จิตฺตสมุฏฺฐาโน ธมฺโม จิตฺตสมุฏฺฐานสฺส จ โนจิตฺตสมุฏฺฐานสฺส จ ธมฺมสฺส เหตุปจฺจเยน ปจฺจโย.\csromanspacer{} จิตฺตสมุฏฺฐานา เหตู สมฺปยุตฺตกานํ ขนฺธานํ จิตฺตสฺส จ จิตฺตสมุฏฺฐานานญฺจ รูปานํ เหตุปจฺจเยน ปจฺจโย.\csromanspacer{} ปฏิสนฺธิกฺขเณ ฯเปฯ. (3)}

\csromanheader{2}{60. จิตฺตสมุฏฺฐานทุก \textbf{อารมฺมณ}}
\proseitem{211}{\csromanfolio{449}จิตฺตสมุฏฺฐาโน ธมฺโม จิตฺตสมุฏฺฐานสฺส ธมฺมสฺส อารมฺมณปจฺจเยน ปจฺจโย.\csromanspacer{} จิตฺตสมุฏฺฐาเน ขนฺเธ อารพฺภ จิตฺตสมุฏฺฐานา ขนฺธา อุปฺปชฺชนฺติ. (มูลํ กาตพฺพํ.) จิตฺตสมุฏฺฐาเน ขนฺเธ อารพฺภ จิตฺตํ อุปฺปชฺชติ. (มูลํ กาตพฺพํ.) จิตฺตสมุฏฺฐาเน ขนฺเธ อารพฺภ จิตฺตสมุฏฺฐานา ขนฺธา จ จิตฺตญฺจ อุปฺปชฺชนฺติ. (3)}

\proseitem{212}{โนจิตฺตสมุฏฺฐาโน ธมฺโม โนจิตฺตสมุฏฺฐานสฺส ธมฺมสฺส อารมฺมณปจฺจเยน ปจฺจโย.\csromanspacer{} อริยา มคฺคา วุฏฺฐหิตฺวา มคฺคํ ปจฺจเวกฺขนฺติ.\csromanspacer{} ผลํ ฯเปฯ.\csromanspacer{} นิพฺพานํ ปจฺจเวกฺขนฺติ.\csromanspacer{} นิพฺพานํ โคตฺรภุสฺส, โวทานสฺส, มคฺคสฺส, ผลสฺส, อาวชฺชนาย อารมฺมณปจฺจเยน ปจฺจโย.\csromanspacer{} จกฺขุํ ฯเปฯ.\csromanspacer{} วตฺถุํ โนจิตฺตสมุฏฺฐาเน ขนฺเธ อนิจฺจโต ฯเปฯ วิปสฺสติ, อสฺสาเทติ อภินนฺทติ, ตํ อารพฺภ จิตฺตํ อุปฺปชฺชติ.\csromanspacer{} ทิพฺเพน จกฺขุนา รูปํ ปสฺสติ.\csromanspacer{} ทิพฺพาย โสตธาตุยา สทฺทํ สุณาติ.\csromanspacer{} เจโตปริยญาเณน โนจิตฺตสมุฏฺฐานจิตฺตสมงฺคิสฺส จิตฺตํ ชานาติ.\csromanspacer{} อากาสานญฺจายตนํ ฯเปฯ.\csromanspacer{} อากิญฺจญฺญายตนํ ฯเปฯ.\csromanspacer{} รูปายตนํ จกฺขุวิญฺญาณสฺส ฯเปฯ.\csromanspacer{} โผฏฺฐพฺพายตนํ ฯเปฯ.\csromanspacer{} โนจิตฺตสมุฏฺฐานา ขนฺธา อิทฺธิวิธญาณสฺส, เจโตปริยญาณสฺส, ปุพฺเพนิวาสานุสฺสติญาณสฺส, อนาคตํสญาณสฺส, อาวชฺชนาย อารมฺมณปจฺจเยน ปจฺจโย. (1)}

\prose{โนจิตฺตสมุฏฺฐาโน ธมฺโม จิตฺตสมุฏฺฐานสฺส ธมฺมสฺส อารมฺมณปจฺจเยน ปจฺจโย.\csromanspacer{} อริยา มคฺคา ฯเปฯ.\csromanspacer{} นิพฺพานํ ปจฺจเวกฺขนฺติ. (ปฐมคมนสทิสํ.) จกฺขุํ ฯเปฯ.\csromanspacer{} วตฺถุํ โนจิตฺตสมุฏฺฐาเน ขนฺเธ อนิจฺจโต ฯเปฯ โทมนสฺสํ อุปฺปชฺชติ.\csromanspacer{} ทิพฺเพน จกฺขุนา รูปํ ปสฺสติ ฯเปฯ.\csromanspacer{} รูปายตนํ จกฺขุวิญฺญาณสหคตานํ ขนฺธานํ ฯเปฯ.\csromanspacer{} โผฏฺฐพฺพายตนํ ฯเปฯ.\csromanspacer{} โนจิตฺตสมุฏฺฐานา ขนฺธา อิทฺธิวิธญาณสฺส, เจโตปริยญาณสฺส, ปุพฺเพนิวาสานุสฺสติญาณสฺส, อนาคตํสญาณสฺส, อาวชฺชนาย อารมฺมณปจฺจเยน ปจฺจโย. (2)}

\prose{โนจิตฺตสมุฏฺฐาโน ธมฺโม จิตฺตสมุฏฺฐานสฺส จ โนจิตฺตสมุฏฺฐานสฺส จ ธมฺมสฺส อารมฺมณปจฺจเยน ปจฺจโย.\csromanspacer{} อริยา มคฺคา ฯเปฯ.\csromanspacer{} นิพฺพานํ ปจฺจเวกฺขนฺติ. (ปฐมคมนสทิสํ.) โนจิตฺตสมุฏฺฐาเน ขนฺเธ อนิจฺจโต ฯเปฯ วิปสฺสติ, อสฺสาเทติ, อภินนฺทติ, ตํ อารพฺภ จิตฺตญฺจ สมฺปยุตฺตกา จ ขนฺธา อุปฺปชฺชนฺติ.\csromanspacer{} ทิพฺเพน จกฺขุนา รูปํ ปสฺสติ.\csromanspacer{} ทิพฺพาย โสตธาตุยา สทฺทํ สุณาติ.\csromanspacer{} รูปายตนํ จกฺขุวิญฺญาณสฺส สมฺปยุตฺตกานญฺจ ขนฺธานํ ฯเปฯ.\csromanspacer{} โผฏฺฐพฺพายตนํ ฯเปฯ. \csromanfolio{450}โนจิตฺตสมุฏฺฐานา ขนฺธา อิทฺธิวิธญาณสฺส, เจโตปริยญาณสฺส, ปุพฺเพนิวาสานุสฺสติญาณสฺส, อนาคตํสญาณสฺส, อาวชฺชนาย อารมฺมณปจฺจเยน ปจฺจโย. (3)}

\prose{จิตฺตสมุฏฺฐาโน จ โนจิตฺตสมุฏฺฐาโน จ ธมฺมา จิตฺตสมุฏฺฐานสฺส ธมฺมสฺส อารมฺมณปจฺจเยน ปจฺจโย.\csromanspacer{} ตีณิ. (อารพฺภ กาตพฺพา.)}

\csromanheader{2}{\textbf{อธิปติ}}
\proseitem{213}{จิตฺตสมุฏฺฐาโน ธมฺโม จิตฺตสมุฏฺฐานสฺส ธมฺมสฺส อธิปติปจฺจเยน ปจฺจโย.\csromanspacer{} \textbf{อารมฺมณาธิปติ}, สหชาตาธิปติ.\csromanspacer{}\csromanspacer{}\textbf{อารมฺมณาธิปติ}\csromandash{}จิตฺตสมุฏฺฐาเน ขนฺเธ ครุํ กตฺวา จิตฺตสมุฏฺฐานา ขนฺธา อุปฺปชฺชนฺติ.\csromanspacer{} \textbf{สหชาตาธิปติ}\csromandash{}จิตฺตสมุฏฺฐานาธิปติ สมฺปยุตฺตกานํ ขนฺธานํ จิตฺตสมุฏฺฐานานญฺจ รูปานํ อธิปติปจฺจเยน ปจฺจโย. (ตีณิปิ.\csromanspacer{} \textbf{อารมฺมณาธิปติ} สหชาตาธิปติปิ กาตพฺพา.) (3)}

\prose{โนจิตฺตสมุฏฺฐาโน ธมฺโม โนจิตฺตสมุฏฺฐานสฺส ธมฺมสฺส อธิปติปจฺจเยน ปจฺจโย.\csromanspacer{} \textbf{อารมฺมณาธิปติ}\csromandash{}อริยา มคฺคา ฯเปฯ.\csromanspacer{} นิพฺพานํ ครุํ กตฺวา ฯเปฯ.\csromanspacer{} โนจิตฺตสมุฏฺฐาเน ขนฺเธ ครุํ กตฺวา อสฺสาเทติ อภินนฺทติ, ตํ ครุํ กตฺวา จิตฺตํ อุปฺปชฺชติ. (1)}

\prose{โนจิตฺตสมุฏฺฐาโน ธมฺโม จิตฺตสมุฏฺฐานสฺส ธมฺมสฺส อธิปติปจฺจเยน ปจฺจโย.\csromanspacer{} \textbf{อารมฺมณาธิปติ}, สหชาตาธิปติ.\csromanspacer{}\csromanspacer{}\textbf{อารมฺมณาธิปติ}\csromandash{}อริยา มคฺคา ฯเปฯ.\csromanspacer{} นิพฺพานํ ครุํ กตฺวา ปจฺจเวกฺขนฺติ ฯเปฯ.\csromanspacer{} โนจิตฺตสมุฏฺฐาเน ขนฺเธ ครุํ กตฺวา อสฺสาเทติ อภินนฺทติ, ตํ ครุํ กตฺวา ราโค อุปฺปชฺชติ.\csromanspacer{} ทิฏฺฐิ อุปฺปชฺชติ.\csromanspacer{} \textbf{สหชาตาธิปติ}\csromandash{}โนจิตฺตสมุฏฺฐานาธิปติ สมฺปยุตฺตกานํ ขนฺธานํ จิตฺตสมุฏฺฐานานญฺจ รูปานํ อธิปติปจฺจเยน ปจฺจโย. (2)}

\prose{โนจิตฺตสมุฏฺฐาโน ธมฺโม จิตฺตสมุฏฺฐานสฺส จ โนจิตฺตสมุฏฺฐานสฺส จ ธมฺมสฺส อธิปติปจฺจเยน ปจฺจโย.\csromanspacer{} \textbf{อารมฺมณาธิปติ}\csromandash{}อริยา มคฺคา ฯเปฯ.\csromanspacer{} นิพฺพานํ ครุํ กตฺวา ปจฺจเวกฺขนฺติ ฯเปฯ.\csromanspacer{} โนจิตฺตสมุฏฺฐาเน ขนฺเธ ครุํ กตฺวา ฯเปฯ.\csromanspacer{} จิตฺตญฺจ สมฺปยุตฺตกา จ ขนฺธา อุปฺปชฺชนฺติ. (3)}

\prose{จิตฺตสมุฏฺฐาโน จ โนจิตฺตสมุฏฺฐาโน จ ธมฺมา จิตฺตสมุฏฺฐานสฺส ธมฺมสฺส อธิปติปจฺจเยน ปจฺจโย.\csromanspacer{} \textbf{อารมฺมณาธิปติ}.\csromanspacer{} ตีณิ. (\textbf{อารมฺมณาธิปติ}เยว.)}

\csromanheader{2}{60. จิตฺตสมุฏฺฐานทุก \textbf{อนนฺตรสมนนฺตร}}
\proseitem{214}{\csromanfolio{451}จิตฺตสมุฏฺฐาโน ธมฺโม จิตฺตสมุฏฺฐานสฺส ธมฺมสฺส อนนฺตรปจฺจเยน ปจฺจโย.\csromanspacer{} ตีณิ. (วุฏฺฐานํ นตฺถิ.)}

\prose{โนจิตฺตสมุฏฺฐาโน ธมฺโม โนจิตฺตสมุฏฺฐานสฺส ธมฺมสฺส อนนฺตรปจฺจเยน ปจฺจโย.\csromanspacer{} ปุริมํ ปุริมํ จิตฺตํ ปจฺฉิมสฺส ปจฺฉิมสฺส ฯเปฯ.\csromanspacer{} นิโรธา วุฏฺฐหนฺตสฺส เนวสญฺญานาสญฺญายตนํ ผลสมาปตฺติยา อนนฺตรปจฺจเยน ปจฺจโย. (อิตเร เทฺว คณนา, อิมสฺส สทิสาเยว กาตพฺพา.)}

\prose{จิตฺตสมุฏฺฐาโน จ โนจิตฺตสมุฏฺฐาโน จ ธมฺมา จิตฺตสมุฏฺฐานสฺส ธมฺมสฺส อนนฺตรปจฺจเยน ปจฺจโย. (ตีณิ กาตพฺพา.\csromanspacer{} วุฏฺฐานํ นตฺถิ.) สมนนฺตรปจฺจเยน ปจฺจโย.}

\csromanheader{2}{\textbf{สหชาตาทิ}}
\proseitem{215}{จิตฺตสมุฏฺฐาโน ธมฺโม จิตฺตสมุฏฺฐานสฺส ธมฺมสฺส สหชาตปจฺจเยน ปจฺจโย. (ปฏิจฺจวารสทิสํ.) อญฺญมญฺญปจฺจเยน ปจฺจโย. (ปฏิจฺจวารสทิสํ.) นิสฺสยปจฺจเยน ปจฺจโย. (ปจฺจยวารสทิสํ)}

\csromanheader{2}{\textbf{อุปนิสฺสย}}
\proseitem{216}{จิตฺตสมุฏฺฐาโน ธมฺโม จิตฺตสมุฏฺฐานสฺส ธมฺมสฺส อุปนิสฺสยปจฺจเยน ปจฺจโย.\csromanspacer{} อารมฺมณูปนิสฺสโย, อนนฺตรูปนิสฺสโย, ปกตูปนิสฺสโย ฯเปฯ.\csromanspacer{} ปกตูปนิสฺสโย\csromandash{}(ตีณิ ปญฺหา กาตพฺพา.) (3)}

\prose{โนจิตฺตสมุฏฺฐาโน ธมฺโม โนจิตฺตสมุฏฺฐานสฺส ธมฺมสฺส อุปนิสฺสยปจฺจเยน ปจฺจโย.\csromanspacer{} อารมฺมณูปนิสฺสโย, อนนฺตรูปนิสฺสโย, ปกตูปนิสฺสโย ฯเปฯ.\csromanspacer{} ปกตูปนิสฺสโย\csromandash{}อุตุํ.\csromanspacer{} โภชนํ.\csromanspacer{} เสนาสนํ จิตฺตํ อุปนิสฺสาย ทานํ เทติ ฯเปฯ สํฆํ ภินฺทติ.\csromanspacer{} อุตุ.\csromanspacer{} โภชนํ.\csromanspacer{} เสนาสนํ จิตฺตํ จิตฺตสฺส อุปนิสฺสยปจฺจเยน ปจฺจโย. (1)}

\prose{โนจิตฺตสมุฏฺฐาโน ธมฺโม จิตฺตสมุฏฺฐานสฺส ธมฺมสฺส อุปนิสฺสยปจฺจเยน ปจฺจโย.\csromanspacer{} อารมฺมณูปนิสฺสโย, อนนฺตรูปนิสฺสโย, ปกตูปนิสฺสโย ฯเปฯ.\csromanspacer{} ปกตูปนิสฺสโย\csromandash{}อุตุํ.\csromanspacer{} โภชนํ.\csromanspacer{} เสนาสนํ จิตฺตํ \csromanfolio{452}อุปนิสฺสาย ทานํ เทติ ฯเปฯ สํฆํ ภินฺทติ.\csromanspacer{} อุตุ.\csromanspacer{} โภชนํ.\csromanspacer{} เสนาสนํ จิตฺตํ สทฺธาย ฯเปฯ มคฺคสฺส ผลสมาปตฺติยา อุปนิสฺสยปจฺจเยน ปจฺจโย. (2)}

\prose{โนจิตฺตสมุฏฺฐาโน ธมฺโม จิตฺตสมุฏฺฐานสฺส จ โนจิตฺตสมุฏฺฐานสฺส จ ธมฺมสฺส อุปนิสฺสยปจฺจเยน ปจฺจโย.\csromanspacer{} อารมฺมณูปนิสฺสโย, อนนฺตรูปนิสฺสโย, ปกตูปนิสฺสโย ฯเปฯ.\csromanspacer{} ปกตูปนิสฺสโย\csromandash{}อุตุํ.\csromanspacer{} โภชนํ.\csromanspacer{} เสนาสนํ จิตฺตํ อุปนิสฺสาย ทานํ เทติ ฯเปฯ สํฆํ ภินฺทติ.\csromanspacer{} อุตุ.\csromanspacer{} โภชนํ.\csromanspacer{} เสนาสนํ จิตฺตํ จิตฺตสมุฏฺฐานานํ ขนฺธานํ จิตฺตสฺส จ อุปนิสฺสยปจฺจเยน ปจฺจโย. (3)}

\prose{จิตฺตสมุฏฺฐาโน จ โนจิตฺตสมุฏฺฐาโน จ ธมฺมา จิตฺตสมุฏฺฐานสฺส ธมฺมสฺส อุปนิสฺสยปจฺจเยน ปจฺจโย.\csromanspacer{} อารมฺมณูปนิสฺสโย, อนนฺตรูปนิสฺสโย, ปกตูปนิสฺสโย ฯเปฯ.\csromanspacer{} ปกตูปนิสฺสโย\csromandash{}ตีณิ.}

\csromanheader{2}{\textbf{ปุเรชาต}}
\proseitem{217}{จิตฺตสมุฏฺฐาโน ธมฺโม จิตฺตสมุฏฺฐานสฺส ธมฺมสฺส ปุเรชาตปจฺจเยน ปจฺจโย.\csromanspacer{} \textbf{อารมฺมณปุเรชาตํ}\csromandash{}จิตฺตสมุฏฺฐาเน รูเป ฯเปฯ.\csromanspacer{} โผฏฺฐพฺเพ อนิจฺจโต ฯเปฯ วิปสฺสติ ฯเปฯ โทมนสฺสํ อุปฺปชฺชติ.\csromanspacer{} ทิพฺเพน จกฺขุนา รูปํ ปสฺสติ.\csromanspacer{} ทิพฺพาย โสตธาตุยา สทฺทํ สุณาติ.\csromanspacer{} รูปายตนํ จกฺขุวิญฺญาณสหคตานํ ขนฺธานํ ปุเรชาตปจฺจเยน ปจฺจโย ฯเปฯ.\csromanspacer{} โผฏฺฐพฺพายตนํ กายวิญฺญาณสหคตานํ ขนฺธานํ ปุเรชาตปจฺจเยน ปจฺจโย. (1)}

\prose{จิตฺตสมุฏฺฐาโน ธมฺโม โนจิตฺตสมุฏฺฐานสฺส ธมฺมสฺส ปุเรชาตปจฺจเยน ปจฺจโย.\csromanspacer{} \textbf{อารมฺมณปุเรชาตํ}\csromandash{}จิตฺตสมุฏฺฐาเน รูเป ฯเปฯ โผฏฺฐพฺเพ อนิจฺจโต ฯเปฯ วิปสฺสติ, อสฺสาเทติ อภินนฺทติ, ตํ อารพฺภ จิตฺตํ อุปฺปชฺชติ.\csromanspacer{} ทิพฺเพน จกฺขุนา รูปํ ปสฺสติ.\csromanspacer{} ทิพฺพาย โสตธาตุยา สทฺทํ สุณาติ.\csromanspacer{} รูปายตนํ จกฺขุวิญฺญาณสฺส ฯเปฯ.\csromanspacer{} โผฏฺฐพฺพายตนํ กายวิญฺญาณสฺส ปุเรชาตปจฺจเยน ปจฺจโย. (2)}

\prose{จิตฺตสมุฏฺฐาโน ธมฺโม จิตฺตสมุฏฺฐานสฺส จ โนจิตฺตสมุฏฺฐานสฺส จ ธมฺมสฺส ปุเรชาตปจฺจเยน ปจฺจโย.\csromanspacer{} \textbf{อารมฺมณปุเรชาตํ}\csromandash{} จิตฺตสมุฏฺฐาเน รูเป ฯเปฯ โผฏฺฐพฺเพ อนิจฺจโต ฯเปฯ วิปสฺสติ, อสฺสาเทติ อภินนฺทติ, ตํ อารพฺภ จิตฺตญฺจ สมฺปยุตฺตกา จ ขนฺธา อุปฺปชฺชนฺติ.\csromanspacer{} ทิพฺเพน จกฺขุนา รูปํ ปสฺสติ.}

\vspace{\tipitakaparskip}
\csromancenter{จิตฺตสมุฏฺฐานทุก ทิพฺพาย โสตธาตุยา สทฺทํ สุณาติ.\csromanspacer{} รูปายตนํ จกฺขุวิญฺญาณสฺส สมฺปยุตฺตกานญฺจ ขนฺธานํ ปุเรชาตปจฺจเยน ปจฺจโย ฯเปฯ.\csromanspacer{} โผฏฺฐพฺพายตนํ. (3)}
\proseitem{218}{\csromanfolio{453}โนจิตฺตสมุฏฺฐาโน ธมฺโม โนจิตฺตสมุฏฺฐานสฺส ธมฺมสฺส ปุเรชาตปจฺจเยน ปจฺจโย.\csromanspacer{} \textbf{อารมฺมณปุเรชาตํ, วตฺถุปุเรชาตํ.}\csromanspacer{}\csromanspacer{}\textbf{อารมฺมณปุเรชาตํ}\csromandash{}จกฺขุํ ฯเปฯ กายํ.\csromanspacer{} รูเป ฯเปฯ เผฏฺฐพฺเพ วตฺถุํ อนิจฺจโต ฯเปฯ ตํ อารพฺภ จิตฺตํ อุปฺปชฺชติ.\csromanspacer{} ทิพฺเพน จกฺขุนา รูปํ ปสฺสติ.\csromanspacer{} ทิพฺพาย โสตธาตุยา สทฺทํ สุณาติ.\csromanspacer{} รูปายตนํ จกฺขุวิญฺญาณสฺส ฯเปฯ.\csromanspacer{} โผฏฺฐพฺพายตนํ ฯเปฯ.\csromanspacer{} \textbf{วตฺถุปุเรชาตํ}\csromandash{} จกฺขายตนํ จกฺขุวิญฺญาณสฺส ฯเปฯ กายายตนํ ฯเปฯ.\csromanspacer{} วตฺถุ จิตฺตสฺส ปุเรชาตปจฺจเยน ปจฺจโย. (1)}

\prose{โนจิตฺตสมุฏฺฐาโน ธมฺโม จิตฺตสมุฏฺฐานสฺส ธมฺมสฺส ปุเรชาตปจฺจเยน ปจฺจโย.\csromanspacer{} \textbf{อารมฺมณปุเรชาตํ, วตฺถุปุเรชาตํ.}\csromanspacer{}\csromanspacer{}\textbf{อารมฺมณปุเรชาตํ}\csromandash{}จกฺขุํ ฯเปฯ วตฺถุํ อนิจฺจโต ฯเปฯ โทมนสฺสํ อุปชฺชติ.\csromanspacer{} ทิพฺเพน จกฺขุนา รูปํ ปสฺสติ.\csromanspacer{} ทิพฺพาย โสตธาตุยา สทฺทํ สุณาติ.\csromanspacer{} รูปายตนํ จกฺขุวิญฺญาณสหคตานํ ขนฺธานํ ฯเปฯ โผฏฺฐพฺพายตนํ ฯเปฯ.\csromanspacer{} \textbf{วตฺถุปุเรชาตํ}\csromandash{}จกฺขายตนํ จกฺขุวิญฺญาณสหคตานํ ขนฺธานํ ฯเปฯ กายายตนํ ฯเปฯ.\csromanspacer{} วตฺถุ จิตฺตสมุฏฺฐานานํ ขนฺธานํ ปุเรชาตปจฺจเยน ปจฺจโย. (2)}

\prose{โนจิตฺตสมุฏฺฐาโน ธมฺโม จิตฺตสมุฏฺฐานสฺส จ โนจิตฺตสมุฏฺฐานสฺส จ ธมฺมสฺส ปุเรชาตปจฺจเยน ปจฺจโย.\csromanspacer{} \textbf{อารมฺมณปุเรชาตํ, วตฺถุปุเรชาตํ.}\csromanspacer{}\csromanspacer{}\textbf{อารมฺมณปุเรชาตํ}\csromandash{}จกฺขุํ ฯเปฯ วตฺถุํ อนิจฺจโต ฯเปฯ ตํ อารพฺภ จิตฺตญฺจ สมฺปยุตฺตกา จ ขนฺธา อุปฺปชฺชนฺติ.\csromanspacer{} ทิพฺเพน จกฺขุนา รูปํ ปสฺสติ ฯเปฯ.\csromanspacer{} ทิพฺพาย โสตธาตุยา สทฺทํ สุณาติ.\csromanspacer{} รูปายตนํ จกฺขุวิญฺญาณสฺส สมฺปยุตฺตกานญฺจ ขนฺธานํ ฯเปฯ.\csromanspacer{} โผฏฺฐพฺพายตนํ ฯเปฯ.\csromanspacer{} \textbf{วตฺถุปุเรชาตํ}\csromandash{}จกฺขายตนํ จกฺขุวิญฺญาณสฺส สมฺปยุตฺตกานญฺจ ขนฺธานํ ฯเปฯ กายายตนํ ฯเปฯ.\csromanspacer{} วตฺถุ จิตฺตสฺส สมฺปยุตฺตกานญฺจ ขนฺธานํ ปุเรชาตปจฺจเยน ปจฺจโย. (3)}

\proseitem{219}{จิตฺตสมุฏฺฐาโน จ โนจิตฺตสมุฏฺฐาโน จ ธมฺมา จิตฺตสมุฏฺฐานสฺส ธมฺมสฺส ปุเรชาตปจฺจเยน ปจฺจโย.\csromanspacer{} \textbf{อารมฺมณปุเรชาตํ, วตฺถุปุเรชาตํ.}\csromanspacer{} จิตฺตสมุฏฺฐานํ รูปายตนญฺจ จกฺขายตนญฺจ จกฺขุวิญฺญาณสหคตานํ \csromanfolio{454}ขนฺธานํ ฯเปฯ.\csromanspacer{} จิตฺตสมุฏฺฐานํ โผฏฺฐพฺพายตนญฺจ กายายตนญฺจ กายวิญฺญาณสหคตานํ ขนฺธานํ ปุเรชาตปจฺจเยน ปจฺจโย.\csromanspacer{} จิตฺตสมุฏฺฐานํ รูปายตนญฺจ วตฺถุ จ จิตฺตสมุฏฺฐานานํ ขนฺธานํ ปุเรชาตปจฺจเยน ปจฺจโย ฯเปฯ.\csromanspacer{} จิตฺตสมุฏฺฐานํ โผฏฺฐพฺพายตนญฺจ วตฺถุ จ ฯเปฯ. (1)}

\prose{จิตฺตสมุฏฺฐาโน จ โนจิตฺตสมุฏฺฐาโน จ ธมฺมา โนจิตฺตสมุฏฺฐานสฺส ธมฺมสฺส ปุเรชาตปจฺจเยน ปจฺจโย.\csromanspacer{} \textbf{อารมฺมณปุเรชาตํ, วตฺถุปุเรชาตํ.}\csromanspacer{} จิตฺตสมุฏฺฐานํ รูปายตนญฺจ จกฺขายตนญฺจ จกฺขุวิญฺญาณสฺส ฯเปฯ.\csromanspacer{} จิตฺตสมุฏฺฐานํ โผฏฺฐพฺพายตนญฺจ กายายตนญฺจ กายวิญฺญาณสฺส ปุเรชาตปจฺจเยน ปจฺจโย.\csromanspacer{} จิตฺตสมุฏฺฐานํ รูปายตนญฺจ วตฺถุ จ จิตฺตสฺส ปุเรชาตปจฺจเยน ปจฺจโย ฯเปฯ.\csromanspacer{} จิตฺตสมุฏฺฐานํ โผฏฺฐพฺพายตนญฺจ วตฺถุ จ ฯเปฯ. (2)}

\prose{จิตฺตสมุฏฺฐาโน จ โนจิตฺตสมุฏฺฐาโน จ ธมฺมา จิตฺตสมุฏฺฐานสฺส จ โนจิตฺตสมุฏฺฐานสฺส จ ธมฺมสฺส ปุเรชาตปจฺจเยน ปจฺจโย.\csromanspacer{} \textbf{อารมฺมณปุเรชาตํ, วตฺถุปุเรชาตํ.}\csromanspacer{} จิตฺตสมุฏฺฐานํ รูปายตนญฺจ จกฺขายตนญฺจ จกฺขุวิญฺญาณสฺส สมฺปยุตฺตกานญฺจ ขนฺธานํ ปุเรชาตปจฺจเยน ปจฺจโย ฯเปฯ.\csromanspacer{} จิตฺตสมุฏฺฐานํ โผฏฺฐพฺพายตนญฺจ ฯเปฯ.\csromanspacer{} จิตฺตสมุฏฺฐานํ รูปายตนญฺจ วตฺถุ จ จิตฺตสฺส สมฺปยุตฺตกานญฺจ ขนฺธานํ ปุเรชาตปจฺจเยน ปจฺจโย ฯเปฯ.\csromanspacer{} จิตฺตสมุฏฺฐานํ โผฏฺฐพฺพายตนญฺจ วตฺถุ จ ฯเปฯ. (3)}

\csromanheader{2}{\textbf{ปจฺฉาชาตาเสวน}}
\proseitem{220}{จิตฺตสมุฏฺฐาโน ธมฺโม จิตฺตสมุฏฺฐานสฺส ธมฺมสฺส ปจฺฉาชาตปจฺจเยน ปจฺจโย.\csromanspacer{} ปจฺฉาชาตา จิตฺตสมุฏฺฐานา ขนฺธา ปุเรชาตสฺส อิมสฺส จิตฺตสมุฏฺฐานสฺส กายสฺส ปจฺฉาชาปจฺจเยน ปจฺจโย. (อิมินากาเรเนว ปจฺฉาชาโต วิตฺถาเรตพฺโพ.) อาเสวนปจฺจเยน ปจฺจโย.\csromanspacer{} นว.}

\csromanheader{2}{\textbf{กมฺมวิปาก}}
\proseitem{221}{จิตฺตสมุฏฺฐาโน ธมฺโม จิตฺตสมุฏฺฐานสฺส ธมฺมสฺส กมฺมปจฺจเยน ปจฺจโย.\csromanspacer{} \textbf{สหชาตา}, นานากฺขณิกา.\csromanspacer{}\csromanspacer{}\textbf{สหชาตา} จิตฺตสมุฏฺฐานา เจตนา สมฺปยุตฺตกานํ ขนฺธานํ จิตฺตสมุฏฺฐานานญฺจ รูปานํ กมฺมปจฺจเยน ปจฺจโย.\csromanspacer{} ปฏิสนฺธิกฺขเณ ฯเปฯ.\csromanspacer{} \textbf{นานากฺขณิกา} จิตฺตสมุฏฺฐานา เจตนา วิปากานํ ขนฺธานํ กมฺมปจฺจเยน ปจฺจโย. (1)}

\csromanheader{2}{60. จิตฺตสมุฏฺฐานทุก}
\prose{\csromanfolio{455}จิตฺตสมุฏฺฐาโน ธมฺโม โนจิตฺตสมุฏฺฐานสฺส ธมฺมสฺส กมฺมปจฺจเยน ปจฺจโย.\csromanspacer{} \textbf{สหชาตา}, นานากฺขณิกา.\csromanspacer{}\csromanspacer{}\textbf{สหชาตา} จิตฺตสมุฏฺฐานา เจตนา จิตฺตสฺส กมฺมปจฺจเยน ปจฺจโย.\csromanspacer{} ปฏิสนฺธิกฺขเณ ฯเปฯ.\csromanspacer{} \textbf{นานากฺขณิกา} จิตฺตสมุฏฺฐานา เจตนา วิปากสฺส จิตฺตสฺส กฏตฺตา จ รูปานํ กมฺมปจฺจเยน ปจฺจโย. (2)}

\prose{จิตฺตสมุฏฺฐาโน ธมฺโม จิตฺตสมุฏฺฐานสฺส จ โนจิตฺตสมุฏฺฐานสฺส จ ธมฺมสฺส กมฺมปจฺจเยน ปจฺจโย.\csromanspacer{} \textbf{สหชาตา}, นานากฺขณิกา.\csromanspacer{}\csromanspacer{}\textbf{สหชาตา} จิตฺตสมุฏฺฐานา เจตนา สมฺปยุตฺตกานํ ขนฺธานํ จิตฺตสฺส จ จิตฺตสมุฏฺฐานานญฺจ รูปานํ กมฺมปจฺจเยน ปจฺจโย.\csromanspacer{} ปฏิสนฺธิกฺขเณ ฯเปฯ.\csromanspacer{} \textbf{นานากฺขณิกา} จิตฺตสมุฏฺฐานา เจตนา วิปากานํ ขนฺธานํ จิตฺตสฺส จ กฏตฺตา จ รูปานํ กมฺมปจฺจเยน ปจฺจโย. (3)}

\prose{จิตฺตสมุฏฺฐาโน ธมฺโม จิตฺตสมุฏฺฐานสฺส ธมฺมสฺส วิปากปจฺจเยน ปจฺจโย.\csromanspacer{} นว.}

\csromanheader{2}{\textbf{อาหาร}}
\proseitem{222}{จิตฺตสมุฏฺฐาโน ธมฺโม จิตฺตสมุฏฺฐานสฺส ธมฺมสฺส อาหารปจฺจเยน ปจฺจโย.\csromanspacer{} จิตฺตสมุฏฺฐานา อาหารา สมฺปยุตฺตกานํ ขนฺธานํ จิตฺตสมุฏฺฐานานญฺจ รูปานํ อาหารปจฺจเยน ปจฺจโย.\csromanspacer{} ปฏิสนฺธิกฺขเณ ฯเปฯ. (มูลํ กาตพฺพํ.) จิตฺตสมุฏฺฐานา อาหารา จิตฺตสฺส อาหารปจฺจเยน ปจฺจโย.\csromanspacer{} ปฏิสนฺธิกฺขเณ ฯเปฯ.\csromanspacer{} จิตฺตสมุฏฺฐาโน กพฬีกาโร อาหาโร อิมสฺส โนจิตฺตสมุฏฺฐานสฺส กายสฺส อาหารปจฺจเยน ปจฺจโย. (มูลํ กาตพฺพํ.) จิตฺตสมุฏฺฐานา อาหารา สมฺปยุตฺตกานํ ขนฺธานํ จิตฺตสฺส จ จิตฺตสมุฏฺฐานานญฺจ รูปานํ อาหารปจฺจเยน ปจฺจโย.\csromanspacer{} ปฏิสนฺธิกฺขเณ ฯเปฯ. (3)}

\proseitem{223}{โนจิตฺตสมุฏฺฐาโน ธมฺโม โนจิตฺตสมุฏฺฐานสฺส ธมฺมสฺส อาหารปจฺจเยน ปจฺจโย.\csromanspacer{} ปฏิสนฺธิกฺขเณ โนจิตฺตสมุฏฺฐานา อาหารา กฏตฺตารูปานํ อาหารปจฺจเยน ปจฺจโย.\csromanspacer{} โนจิตฺตสมุฏฺฐาโน กพฬีกาโร อาหาโร อิมสฺส โนจิตฺตสมุฏฺฐานสฺส กายสฺส อาหารปจฺจเยน ปจฺจโย. (มูลํ กาตพฺพํ.) โนจิตฺตสมุฏฺฐานา อาหารา สมฺปยุตฺตกานํ ขนฺธานํ จิตฺตสมุฏฺฐานานญฺจ รูปานํ อาหารปจฺจเยน ปจฺจโย.\csromanspacer{} ปฏิสนฺธิกฺขเณ ฯเปฯ. (มูลํ กาตพฺพํ.) ปฏิสนฺธิกฺขเณ โนจิตฺตสมุฏฺฐานา \csromanfolio{456}อาหารา สมฺปยุตฺตกานํ ขนฺธานํ กฏตฺตา จ รูปานํ อาหารปจฺจเยน ปจฺจโย. (3)}

\proseitem{224}{จิตฺตสมุฏฺฐาโน จ โนจิตฺตสมุฏฺฐาโน จ ธมฺมา จิตฺตสมุฏฺฐานสฺส ธมฺมสฺส อาหารปจฺจเยน ปจฺจโย.\csromanspacer{} จิตฺตสมุฏฺฐานา จ โนจิตฺตสมุฏฺฐานา จ อาหารา สมฺปยุตฺตกานํ ขนฺธานํ จิตฺตสมุฏฺฐานานญฺจ รูปานํ อาหารปจฺจเยน ปจฺจโย.\csromanspacer{} ปฏิสนฺธิกฺขเณ ฯเปฯ. (มูลํ กาตพฺพํ.) ปฏิสนฺธิกฺขเณ จิตฺตสมุฏฺฐานา จ โนจิตฺตสมุฏฺฐานา จ อาหารา กฏตฺตารูปานํ อาหารปจฺจเยน ปจฺจโย.\csromanspacer{} จิตฺตสมุฏฺฐาโน จ โนจิตฺตสมุฏฺฐาโน จ กพฬีกาโร อาหาโร อิมสฺส โนจิตฺตสมุฏฺฐานสฺส กายสฺส อาหารปจฺจเยน ปจฺจโย. (มูลํ กาตพฺพํ.) ปฏิสนฺธิกฺขเณ จิตฺตสมุฏฺฐานา จ โนจิตฺตสมุฏฺฐานา จ อาหารา สมฺปยุตฺตกานํ ขนฺธานํ กฏตฺตา จ รูปานํ อาหารปจฺจเยน ปจฺจโย. (3)}

\csromanheader{2}{\textbf{อินฺทฺริยาทิ}}
\proseitem{225}{จิตฺตสมุฏฺฐาโน ธมฺโม จิตฺตสมุฏฺฐานสฺส ธมฺมสฺส อินฺทฺริยปจฺจเยน ปจฺจโย.\csromanspacer{} ตีณิ.}

\prose{โนจิตฺตสมุฏฺฐาโน ธมฺโม โนจิตฺตสมุฏฺฐานสฺส ธมฺมสฺส อินฺทฺริยปจฺจเยน ปจฺจโย.\csromanspacer{} ปฏิสนฺธิกฺขเณ โนจิตฺตสมุฏฺฐานา อินฺทฺริยา กฏตฺตารูปานํ อินฺทฺริยปจฺจเยน ปจฺจโย.\csromanspacer{} ปฏิสนฺธิกฺขเณ ฯเปฯ.\csromanspacer{} จกฺขุนฺทฺริยํ จกฺขุวิญฺญาณสฺส ฯเปฯ.\csromanspacer{} กายินฺทฺริยํ กายวิญฺญาณสฺส ฯเปฯ.\csromanspacer{} รูปชีวิตินฺทฺริยํ กฏตฺตารูปานํ อินฺทฺริยปจฺจเยน ปจฺจโย. (มูลํ กาตพฺพํ.) โนจิตฺตสมุฏฺฐานา อินฺทฺริยา สมฺปยุตฺตกานํ ขนฺธานํ จิตฺตสมุฏฺฐานานญฺจ รูปานํ อินฺทฺริยปจฺจเยน ปจฺจโย.\csromanspacer{} ปฏิสนฺธิกฺขเณ ฯเปฯ.\csromanspacer{} จกฺขุนฺทฺริยํ จกฺขุวิญฺญาณสหคตานํ ขนฺธานํ ฯเปฯ.\csromanspacer{} กายินฺทฺริยํ ฯเปฯ. (มูลํ กาตพฺพํ.) ปฏิสนฺธิกฺขเณ โนจิตฺตสมุฏฺฐานา อินฺทฺริยา สมฺปยุตฺตกานํ ขนฺธานํ กฏตฺตา จ รูปานํ อินฺทฺริยปจฺจเยน ปจฺจโย.\csromanspacer{} จกฺขุนฺทฺริยํ จกฺขุวิญฺญาณสฺส จกฺขุวิญฺญาณสหคตานญฺจ ขนฺธานํ ฯเปฯ.\csromanspacer{} กายินฺทฺริยํ ฯเปฯ. (3)}

\proseitem{226}{จิตฺตสมุฏฺฐาโน จ โนจิตฺตสมุฏฺฐาโน จ ธมฺมา จิตฺตสมุฏฺฐานุสฺส ธมฺมสฺส อินฺทฺริยปจฺจเยน ปจฺจโย.\csromanspacer{} จิตฺตสมุฏฺฐานา จ โนจิตฺตสมุฏฺฐานา จ อินฺทฺริยา สมฺปยุตฺตกานํ ขนฺธานํ จิตฺตสมุฏฺฐานานญฺจ รูปานํ อินฺทฺริยปจฺจเยน}

\vspace{\tipitakaparskip}
\csromancenter{จิตฺตสมุฏฺฐานทุก ปจฺจโย.\csromanspacer{} ปฏิสนฺธิกฺขเณ ฯเปฯ.\csromanspacer{} จกฺขุนฺทฺริยญฺจ อุเปกฺขินฺทฺริยญฺจ จกฺขุวิญฺญาณสหคตานํ ขนฺธานํ ฯเปฯ.\csromanspacer{} กายินฺทฺริยญฺจ สุขินฺทฺริยญฺจ ฯเปฯ.\csromanspacer{} กายินฺทฺริยญฺจ ทุกฺขินฺทฺริยญฺจ กายวิญฺญาณสหคตานํ ขนฺธานํ อินฺทฺริยปจฺจเยน ปจฺจโย. (มูลํ กาตพฺพํ.) ปฏิสนฺธิกฺขเณ จิตฺตสมุฏฺฐานา จ โนจิตฺตสมุฏฺฐานา จ อินฺทฺริยา กฏตฺตารูปานํ อินฺทฺริยปจฺจเยน ปจฺจโย.\csromanspacer{} จกฺขุนฺทฺริยญฺจ อุเปกฺขินฺทฺริยญฺจ จกฺขุวิญฺญาณสฺส ฯเปฯ.\csromanspacer{} กายินฺทฺริยญฺจ ฯเปฯ. (มูลํ กาตพฺพํ.) ปฏิสนฺธิกฺขเณ จิตฺตสมุฏฺฐานา จ โนจิตฺตสมุฏฺฐานา จ อินฺทฺริยา สมฺปยุตฺตกานํ ขนฺธานํ กฏตฺตา จ รูปานํ อินฺทฺริยปจฺจเยน ปจฺจโย.\csromanspacer{} จกฺขุนฺทฺริยญฺจ อุเปกฺขินฺทฺริยญฺจ จกฺขุวิญฺญาณสฺส สมฺปยุตฺตกานญฺจ ขนฺธานํ อินฺทฺริยปจฺจเยน ปจฺจโย.\csromanspacer{} กายินฺทฺริยญฺจ ฯเปฯ. (3)}
\prose{\csromanfolio{457}ฌานปจฺจเยน ปจฺจโย.\csromanspacer{} ตีณิ, มคฺคปจฺจเยน ปจฺจโย.\csromanspacer{} ตีณิ, สมฺปยุตฺตปจฺจเยน ปจฺจโย.\csromanspacer{} ปญฺจ ฯเปฯ.}

\csromanheader{2}{\textbf{วิปฺปยุตฺต}}
\proseitem{227}{จิตฺตสมุฏฺฐาโน ธมฺโม จิตฺตสมุฏฺฐานสฺส ธมฺมสฺส วิปฺปยุตฺตปจฺจเยน ปจฺจโย.\csromanspacer{} \textbf{สหชาตํ}, ปจฺฉาชาตํ. (สํขิตฺตํ.) (1)}

\prose{จิตฺตสมุฏฺฐาโน ธมฺโม โนจิตฺตสมุฏฺฐานสฺส ธมฺมสฺส วิปฺปยุตฺตปจฺจเยน ปจฺจโย.\csromanspacer{} \textbf{สหชาตํ}, ปจฺฉาชาตํ.\csromanspacer{}\csromanspacer{}\textbf{สหชาตํ} ปฏิสนฺธิกฺขเณ. (สํขิตฺตํ.) (2)}

\prose{จิตฺตสมุฏฺฐาโน ธมฺโม จิตฺตสมุฏฺฐานสฺส จ โนจิตฺตสมุฏฺฐานสฺส จ ธมฺมสฺส วิปฺปยุตฺตปจฺจเยน ปจฺจโย.\csromanspacer{} ปจฺฉาชาตํ. (สํขิตฺตํ.) (3)}

\proseitem{228}{โนจิตฺตสมุฏฺฐาโน ธมฺโม โนจิตฺตสมุฏฺฐานสฺส ธมฺมสฺส วิปฺปยุตฺตปจฺจเยน ปจฺจโย.\csromanspacer{} \textbf{สหชาตํ}, ปุเรชาตํ, ปจฺฉาชาตํ.\csromanspacer{}\csromanspacer{}\textbf{สหชาตํ} ปฏิสนฺธิกฺขเณ จิตฺตํ กฏตฺตารูปานํ วิปฺปยุตฺตปจฺจเยน ปจฺจโย.\csromanspacer{} จิตฺตํ วตฺถุสฺส วิปฺปยุตฺตปจฺจเยน ปจฺจโย.\csromanspacer{} วตฺถุ จิตฺตสฺส วิปฺปยุตฺตปจฺจเยน ปจฺจโย.\csromanspacer{} \textbf{ปุเรชาตํ} จกฺขายตนํ จกฺขุวิญฺญาณสฺส ฯเปฯ.\csromanspacer{} กายายตนํ กายวิญฺญาณสฺส ฯเปฯ.\csromanspacer{} วตฺถุ จิตฺตสฺส วิปฺปยุตฺตปจฺจเยน ปจฺจโย.\csromanspacer{} \textbf{ปจฺฉาชาตา} โนจิตฺตสมุฏฺฐานา ขนฺธา ปุเรชาตสฺส อิมสฺส โนจิตฺตสมุฏฺฐานสฺส กายสฺส วิปฺปยุตฺตปจฺจเยน ปจฺจโย. (1)}

\prose{โนจิตฺตสมุฏฺฐาโน ธมฺโม จิตฺตสมุฏฺฐานสฺส ธมฺมสฺส วิปฺปยุตฺตปจฺจเยน ปจฺจโย.\csromanspacer{} \textbf{สหชาตํ}, ปุเรชาตํ, ปจฺฉาชาตํ. (สํขิตฺตํ.) (2)}

\prose{\csromanfolio{458}โนจิตฺตสมุฏฺฐาโน ธมฺโม จิตฺตสมุฏฺฐานสฺส จ โนจิตฺตสมุฏฺฐานสฺส จ ธมฺมสฺส วิปฺปยุตฺตปจฺจเยน ปจฺจโย.\csromanspacer{} สหชาตํ, ปุเรชาตํ, ปจฺฉาชาตํ. (สํขิตฺตํ.) (3)}

\proseitem{229}{จิตฺตสมุฏฺฐาโน จ โนจิตฺตสมุฏฺฐาโน จ ธมฺมา จิตฺตสมุฏฺฐานสฺส ธมฺมสฺส วิปฺปยุตฺตปจฺจเยน ปจฺจโย.\csromanspacer{} สหชาตํ, ปจฺฉาชาตํ. (สํขิตฺตํ.) (1)}

\prose{จิตฺตสมุฏฺฐาโน จ โนจิตฺตสมุฏฺฐาโน จ ธมฺมา โนจิตฺตสมุฏฺฐานสฺส ธมฺมสฺส วิปฺปยุตฺตปจฺจเยน ปจฺจโย.\csromanspacer{} สหชาตํ, ปจฺฉาชาตํ. (สํขิตฺตํ.) (2)}

\prose{จิตฺตสมุฏฺฐาโน จ โนจิตฺตสมุฏฺฐาโน จ ธมฺมา จิตฺตสมุฏฺฐานสฺส จ โนจิตฺตสมุฏฺฐานสฺส จ ธมฺมสฺส วิปฺปยุตฺตปจฺจเยน ปจฺจโย.\csromanspacer{} ปจฺฉาชาตํ. (สํขิตฺตํ.) (3)}

\csromanheader{2}{\textbf{อตฺถิ}}
\proseitem{230}{จิตฺตสมุฏฺฐาโน ธมฺโม จิตฺตสมุฏฺฐานสฺส ธมฺมสฺส อตฺถิปจฺจเยน ปจฺจโย.\csromanspacer{} สหชาตํ, ปุเรชาตํ, ปจฺฉาชาตํ. (สํขิตฺตํ.) (1)}

\prose{จิตฺตสมุฏฺฐาโน ธมฺโม โนจิตฺตสมุฏฺฐานสฺส ธมฺมสฺส อตฺถิปจฺจเยน ปจฺจโย.\csromanspacer{} สหชาตํ, ปุเรชาตํ, ปจฺฉาชาตํ, อาหารํ. (สํขิตฺตํ.) (2)}

\prose{จิตฺตสมุฏฺฐาโน ธมฺโม จิตฺตสมุฏฺฐานสฺส จ โนจิตฺตสมุฏฺฐานสฺส จ ธมฺมสฺส อตฺถิปจฺจเยน ปจฺจโย.\csromanspacer{} สหชาตํ, ปุเรชาตํ, ปจฺฉาชาตํ. (สํขิตฺตํ.) (3)}

\prose{โนจิตฺตสมุฏฺฐาโน ธมฺโม โนจิตฺตสมุฏฺฐานสฺส ธมฺมสฺส อตฺถิปจฺจเยน ปจฺจโย.\csromanspacer{} สหชาตํ, ปุเรชาตํ, ปจฺฉาชาตํ, อาหารํ, อินฺทฺริยํ. (สํขิตฺตํ.) (1)}

\prose{โนจิตฺตสมุฏฺฐาโน ธมฺโม จิตฺตสมุฏฺฐานสฺส ธมฺมสฺส อตฺถิปจฺจเยน ปจฺจโย.\csromanspacer{} สหชาตํ, ปุเรชาตํ, ปจฺฉาชาตํ. (สํขิตฺตํ.)}

\prose{โนจิตฺตสมุฏฺฐาโน ธมฺโม จิตฺตสมุฏฺฐานสฺส จ โนจิตฺตสมุฏฺฐานสฺส จ ธมฺมสฺส อตฺถิปจฺจเยน ปจฺจโย.\csromanspacer{} สหชาตํ, ปุเรชาตํ, ปจฺฉาชาตํ, อาหารํ. (สํขิตฺตํ.) (3)}

\csromanheader{2}{60. จิตฺตสมุฏฺฐานทุก}
\proseitem{231}{\csromanfolio{459}จิตฺตสมุฏฺฐาโน จ โนจิตฺตสมุฏฺฐาโน จ ธมฺมา จิตฺตสมุฏฺฐานสฺส ธมฺมสฺส อตฺถิปจฺจเยน ปจฺจโย.\csromanspacer{} สหชาตํ, ปุเรชาตํ, ปจฺฉาชาตํ.\csromanspacer{}\csromanspacer{}\textbf{สหชาโต} จกฺขุวิญฺญาณสหคโต เอโก ขนฺโธ ฯเปฯ. (สํขิตฺตํ.) (1)}

\prose{จิตฺตสมุฏฺฐาโน จ โนจิตฺตสมุฏฺฐาโน จ ธมฺมา โนจิตฺตสมุฏฺฐานสฺส ธมฺมสฺส อตฺถิปจฺจเยน ปจฺจโย.\csromanspacer{} สหชาตํ, ปุเรชาตํ, ปจฺฉาชาตํ, อาหารํ, อินฺทฺริยํ.\csromanspacer{}\csromanspacer{}\textbf{สหชาตา} จกฺขุวิญฺญาณสหคตา ขนฺธา จ จกฺขายตนญฺจ จกฺขุวิญฺญาณสฺส อตฺถิปจฺจเยน ปจฺจโย ฯเปฯ.\csromanspacer{} กายวิญฺญาณสหคตา ฯเปฯ.\csromanspacer{} \textbf{สหชาตา} จิตฺตสมุฏฺฐานา ฯเปฯ. (ปจฺจยวารสทิสํ ปฏิสนฺธิปิ ปวตฺติปิ กาตพฺพา สพฺเพสมฺปิ ปญฺหานํ.) \textbf{ปจฺฉาชาตา} จิตฺตสมุฏฺฐานา ขนฺธา จ จิตฺตญฺจ ปุเรชาตสฺส อิมสฺส โนจิตฺตสมุฏฺฐานสฺส กายสฺส อตฺถิปจฺจเยน ปจฺจโย.\csromanspacer{} \textbf{ปจฺฉาชาตา} จิตฺตสมุฏฺฐานา ขนฺธา จ จิตฺตญฺจ กพฬีกาโร อาหาโร จ อิมสฺส โนจิตฺตสมุฏฺฐานสฺส กายสฺส อตฺถิปจฺจเยน ปจฺจโย.\csromanspacer{} \textbf{ปจฺฉาชาตา} จิตฺตสมุฏฺฐานา ขนฺธา จ จิตฺตญฺจ รูปชีวิตินฺทฺริยญฺจ กฏตฺตารูปานํ อตฺถิปจฺจเยน ปจฺจโย. (2)}

\prose{จิตฺตสมุฏฺฐาโน จ โนจิตฺตสมุฏฺฐาโน จ ธมฺมา จิตฺตสมุฏฺฐานสฺส จ โนจิตฺตสมุฏฺฐานสฺส จ ธมฺมสฺส อตฺถิปจฺจเยน ปจฺจโย.\csromanspacer{} สหชาตํ, ปุเรชาตํ, ปจฺฉาชาตํ.\csromanspacer{}\csromanspacer{}\textbf{สหชาโต} จกฺขุวิญฺญาณสหคโต ฯเปฯ. (สํขิตฺตํ.) (3) นตฺถิปจฺจเยน ปจฺจโย.\csromanspacer{} วิคตปจฺจเยน ปจฺจโย.\csromanspacer{} อวิคตปจฺจเยน ปจฺจโย.}

\csromanheader{2}{1. ปจฺจยานุโลม 2. สงฺขฺยาวาร}
\csromanheader{2}{\textbf{สุทฺธ}}
\proseitem{232}{เหตุยา ตีณิ, อารมฺมเณ นว, อธิปติยา นว, อนนฺตเร นว, สมนนฺตเร นว, สหชาเต นว, อญฺญมญฺเญ นว, นิสฺสเย นว, อุปนิสฺสเย นว, ปุเรชาเต นว, ปจฺฉาชาเต นว, อาเสวเน นว, กมฺเม ตีณิ, วิปาเก นว, (สพฺพตฺถ นว,) อินฺทฺริเย นว, ฌาเน ตีณิ, มคฺเค ตีณิ, สมฺปยุตฺเต ปญฺจ, วิปฺปยุตฺเต นว ฯเปฯ อวิคเต นว.}

\vspace{\tipitakaparskip}
\csromancenter{อนุโลมํ.}
\csromanheader{2}{2. ปจฺจนียุทฺธาร}
\proseitem{233}{\csromanfolio{460}จิตฺตสมุฏฺฐาโน ธมฺโม จิตฺตสมุฏฺฐานสฺส ธมฺมสฺส อารมฺมณปจฺจเยน ปจฺจโย.\csromanspacer{} สหชาตปจฺจเยน ปจฺจโย.\csromanspacer{} อุปนิสฺสยปจฺจเยน ปจฺจโย.\csromanspacer{} ปุเรชาตปจฺจเยน ปจฺจโย.\csromanspacer{} ปจฺฉาชาตปจฺจเยน ปจฺจโย.\csromanspacer{} กมฺมปจฺจเยน ปจฺจโย. (1)}

\prose{จิตฺตสมุฏฺฐาโน ธมฺโม โนจิตฺตสมุฏฺฐานสฺส ธมฺมสฺส อารมฺมณปจฺจเยน ปจฺจโย.\csromanspacer{} สหชาตปจฺจเยน ปจฺจโย.\csromanspacer{} อุปนิสฺสยปจฺจเยน ปจฺจโย.\csromanspacer{} ปุเรชาตปจฺจเยน ปจฺจโย.\csromanspacer{} ปจฺฉาชาตปจฺจเยน ปจฺจโย.\csromanspacer{} กมฺมปจฺจเยน ปจฺจโย.\csromanspacer{} อาหารปจฺจเยน ปจฺจโย. (2)}

\prose{จิตฺตสมุฏฺฐาโน ธมฺโม จิตฺตสมุฏฺฐานสฺส จ โนจิตฺตสมุฏฺฐานสฺส จ ธมฺมสฺส อารมฺมณปจฺจเยน ปจฺจโย.\csromanspacer{} สหชาตปจฺจเยน ปจฺจโย.\csromanspacer{} อุปนิสฺสยปจฺจเยน ปจฺจโย.\csromanspacer{} ปุเรชาตปจฺจเยน ปจฺจโย.\csromanspacer{} ปจฺฉาชาตปจฺจเยน ปจฺจโย.\csromanspacer{} กมฺมปจฺจเยน ปจฺจโย. (3)}

\proseitem{234}{โนจิตฺตสมุฏฺฐาโน ธมฺโม โนจิตฺตสมุฏฺฐานสฺส ธมฺมสฺส อารมฺมณปจฺจเยน ปจฺจโย.\csromanspacer{} สหชาตปจฺจเยน ปจฺจโย.\csromanspacer{} อุปนิสฺสยปจฺจเยน ปจฺจโย.\csromanspacer{} ปุเรชาตปจฺจเยน ปจฺจโย.\csromanspacer{} ปจฺฉาชาตปจฺจเยน ปจฺจโย.\csromanspacer{} อาหารปจฺจเยน ปจฺจโย.\csromanspacer{} อินฺทฺริยปจฺจเยน ปจฺจโย. (1)}

\prose{โนจิตฺตสมุฏฺฐาโน ธมฺโม จิตฺตสมุฏฺฐานสฺส ธมฺมสฺส อารมฺมณปจฺจเยน ปจฺจโย.\csromanspacer{} สหชาตปจฺจเยน ปจฺจโย.\csromanspacer{} อุปนิสฺสยปจฺจเยน ปจฺจโย.\csromanspacer{} ปุเรชาตปจฺจเยน ปจฺจโย.\csromanspacer{} ปจฺฉาชาตปจฺจเยน ปจฺจโย. (2)}

\prose{โนจิตฺตสมุฏฺฐาโน ธมฺโม จิตฺตสมุฏฺฐานสฺส จ โนจิตฺตสมุฏฺฐานสฺส จ ธมฺมสฺส อารมฺมณปจฺจเยน ปจฺจโย.\csromanspacer{} สหชาตปจฺจเยน ปจฺจโย.\csromanspacer{} อุปนิสฺสยปจฺจเยน ปจฺจโย.\csromanspacer{} ปุเรชาตปจฺจเยน ปจฺจโย.\csromanspacer{} ปจฺฉาชาตปจฺจเยน ปจฺจโย. (3)}

\proseitem{235}{จิตฺตสมุฏฺฐาโน จ โนจิตฺตสมุฏฺฐาโน จ ธมฺมา จิตฺตสมุฏฺฐานสฺส ธมฺมสฺส อารมฺมณปจฺจเยน ปจฺจโย.\csromanspacer{} สหชาตปจฺจเยน ปจฺจโย.\csromanspacer{} อุปนิสฺสยปจฺจเยน ปจฺจโย.\csromanspacer{} ปุเรชาตปจฺจเยน ปจฺจโย.\csromanspacer{} ปจฺฉาชาตปจฺจเยน ปจฺจโย. (1)}

\csromanmatikaanchor{mtk.60}
\csromantocmarkref{subsection}{61. จิตฺตสหภูทุก}{mtk.60}
\bookmark[dest={mtk.60},level=2]{61. จิตฺตสหภูทุก}
\csromanheader{2}{61. จิตฺตสหภูทุก}
\prose{\csromanfolio{461}จิตฺตสมุฏฺฐาโน จ โนจิตฺตสมุฏฺฐาโน จ ธมฺมา โนจิตฺตสมุฏฺฐานสฺส ธมฺมสฺส อารมฺมณปจฺจเยน ปจฺจโย.\csromanspacer{} สหชาตปจฺจเยน ปจฺจโย.\csromanspacer{} อุปนิสฺสยปจฺจเยน ปจฺจโย.\csromanspacer{} ปุเรชาตปจฺจเยน ปจฺจโย.\csromanspacer{} ปจฺฉาชาตปจฺจเยน ปจฺจโย.\csromanspacer{} อาหารปจฺจเยน ปจฺจโย. (2)}

\prose{จิตฺตสมุฏฺฐาโน จ โนจิตฺตสมุฏฺฐาโน จ ธมฺมา จิตฺตสมุฏฺฐานสฺส จ โนจิตฺตสมุฏฺฐานสฺส จ ธมฺมสฺส อารมฺมณปจฺจเยน ปจฺจโย.\csromanspacer{} สหชาตปจฺจเยน ปจฺจโย.\csromanspacer{} อุปนิสฺสยปจฺจเยน ปจฺจโย.\csromanspacer{} ปุเรชาตปจฺจเยน ปจฺจโย.\csromanspacer{} ปจฺฉาชาตปจฺจเยน ปจฺจโย. (3)}

\csromanheader{2}{2. ปจฺจยปจฺจนีย 2. สงฺขฺยาวาร}
\vspace{\tipitakaparskip}
\csromancenter{นเหตุยา นว, น-อารมฺมเณ นว, (สพฺพตฺถ นว,) โน-อวิคเต นว.}
\csromanheader{2}{3. ปจฺจยานุโลมปจฺจนีย}
\proseitem{237}{\textbf{เหตุ}ปจฺจยา น-อารมฺมเณ ตีณิ, น-อธิปติยา ตีณิ, น-อนนฺตเร ตีณิ, นสมนนฺตเร ตีณิ, น-อญฺญมญฺเญ เทฺว, น-อุปนิสฺสเย ตีณิ ฯเปฯ นมคฺเค ตีณิ, นสมฺปยุตฺเต เทฺว, นวิปฺปยุตฺเต ตีณิ, โนนตฺถิยา ตีณิ, โนวิคเต ตีณิ.}

\csromanheader{2}{4. ปจฺจยปจฺจนียานุโลม}
\proseitem{238}{นเหตุปจฺจยา อารมฺมเณ นว, อธิปติยา นว, (อนุโลมคณนา กาตพฺพา.) ฯเปฯ อวิคเต นว.}

\nitthitam{จิตฺตสมุฏฺฐานทุกํ นิฏฺฐิตํ.}
\csromanheader{2}{61. จิตฺตสหภูทุก 1. ปฏิจฺจวาร}
\csromanheader{2}{1. ปจฺจยานุโลม 1. วิภงฺควาร}
\csromanheader{2}{\textbf{เหตุ}}
\proseitem{239}{จิตฺตสหภุํ ธมฺมํ ปฏิจฺจ จิตฺตสหภู ธมฺโม อุปฺปชฺชติ เหตุปจฺจยา.\csromanspacer{} จิตฺตสหภุํ เอกํ ขนฺธํ ปฏิจฺจ เทฺว ขนฺธา จิตฺตสหภุ จิตฺตสมุฏฺฐานญฺจ \csromanfolio{462}รูปํ.\csromanspacer{} เทฺว ขนฺเธ ฯเปฯ.\csromanspacer{} ปฏิสนฺธิกฺขเณ จิตฺตสหภุํ เอกํ ขนฺธํ ปฏิจฺจ เทฺว ขนฺธา.\csromanspacer{} เทฺว ขนฺเธ ฯเปฯ. (1)}

\prose{จิตฺตสหภุํ ธมฺมํ ปฏิจฺจ โนจิตฺตสหภู ธมฺโม อุปฺปชฺชติ เหตุปจฺจยา.\csromanspacer{} จิตฺตสหภู ขนฺเธ ปฏิจฺจ จิตฺตํ โนจิตฺตสหภุ จิตฺตสมุฏฺฐานญฺจ รูปํ.\csromanspacer{} ปฏิสนฺธิกฺขเณ จิตฺตสหภู ขนฺเธ ปฏิจฺจ จิตฺตํ กฏตฺตา จ รูปํ. (2)}

\prose{จิตฺตสหภุํ ธมฺมํ ปฏิจฺจ จิตฺตสหภู จ โนจิตฺตสหภู จ ธมฺมา อุปฺปชฺชนฺติ เหตุปจฺจยา.\csromanspacer{} จิตฺตสหภุํ เอกํ ขนฺธํ ปฏิจฺจ เทฺว ขนฺธา จิตฺตญฺจ จิตฺตสหภุ จ โนจิตฺตสหภุ จ จิตฺตสมุฏฺฐานํ รูปํ.\csromanspacer{} เทฺว ขนฺเธ ฯเปฯ.\csromanspacer{} ปฏิสนฺธิกฺขเณ จิตฺตสหภุํ เอกํ ขนฺธํ ปฏิจฺจ เทฺว ขนฺธา จิตฺตญฺจ กฏตฺตา จ รูปํ.\csromanspacer{} เทฺว ขนฺเธ ฯเปฯ. (3)}

\proseitem{240}{โนจิตฺตสหภุํ ธมฺมํ ปฏิจฺจ โนจิตฺตสหภู ธมฺโม อุปฺปชฺชติ เหตุปจฺจยา.\csromanspacer{} จิตฺตํ ปฏิจฺจ โนจิตฺตสหภุ จิตฺตสมุฏฺฐานํ รูปํ.\csromanspacer{} ปฏิสนฺธิกฺขเณ จิตฺตํ ปฏิจฺจ กฏตฺตารูปํ.\csromanspacer{} จิตฺตํ ปฏิจฺจ วตฺถุ, วตฺถุํ ปฏิจฺจ จิตฺตํ.\csromanspacer{} เอกํ มหาภูตํ ฯเปฯ.\csromanspacer{} มหาภูเต ปฏิจฺจ โนจิตฺตสหภุ จิตฺตสมุฏฺฐานํ รูปํ, กฏตฺตารูปํ, อุปาทารูปํ. (1)}

\prose{โนจิตฺตสหภุํ ธมฺมํ ปฏิจฺจ จิตฺตสหภู ธมฺโม อุปฺปชฺชติ เหตุปจฺจยา.\csromanspacer{} จิตฺตํ ปฏิจฺจ สมฺปยุตฺตกา ขนฺธา จิตฺตสหภุ จิตฺตสมุฏฺฐานญฺจ รูปํ.\csromanspacer{} ปฏิสนฺธิกฺขเณ จิตฺตํ ปฏิจฺจ สมฺปยุตฺตกา ขนฺธา.\csromanspacer{} ปฏิสนฺธิกฺขเณ วตฺถุํ ปฏิจฺจ จิตฺตสหภู ขนฺธา.\csromanspacer{} มหาภูเต ปฏิจฺจ จิตฺตสหภุ จิตฺตสมุฏฺฐานํ รูปํ, อุปาทารูปํ. (2)}

\prose{โนจิตฺตสหภุํ ธมฺมํ ปฏิจฺจ จิตฺตสหภู จ โนจิตฺตสหภู จ ธมฺมา อุปฺปชฺชนฺติ เหตุปจฺจยา.\csromanspacer{} จิตฺตํ ปฏิจฺจ สมฺปยุตฺตกา ขนฺธา จิตฺตสหภุ จ โนจิตฺตสหภุ จ จิตฺตสมุฏฺฐานํ รูปํ.\csromanspacer{} ปฏิสนฺธิกฺขเณ จิตฺตํ ปฏิจฺจ สมฺปยุตฺตกา ขนฺธา กฏตฺตา จ รูปํ.\csromanspacer{} ปฏิสนฺธิกฺขเณ วตฺถุํ ปฏิจฺจ จิตฺตญฺจ สมฺปยุตฺตกา จ ขนฺธา.\csromanspacer{} มหาภูเต ปฏิจฺจ จิตฺตสหภุ จ โนจิตฺตสหภุ จ จิตฺตสมุฏฺฐานํ รูปํ, อุปาทารูปํ. (3)}

\proseitem{241}{จิตฺตสหภุญฺจ โนจิตฺตสหภุญฺจ ธมฺมํ ปฏิจฺจ จิตฺตสหภู ธมฺโม อุปฺปชฺชติ เหตุปจฺจยา.\csromanspacer{} จิตฺตสหภุํ เอกํ ขนฺธญฺจ จิตฺตญฺจ ปฏิจฺจ เทฺว ขนฺธา จิตฺตสหภุ จิตฺตสมุฏฺฐานญฺจ รูปํ.\csromanspacer{} เทฺว ขนฺเธ จ ฯเปฯ.\csromanspacer{} ปฏิสนฺธิกฺขเณ จิตฺตสหภุํ เอกํ ขนฺธญฺจ จิตฺตญฺจ ปฏิจฺจ เทฺว ขนฺธา.\csromanspacer{} เทฺว ขนฺเธ จ ฯเปฯ.\csromanspacer{} ปฏิสนฺธิกฺขเณ จิตฺตสหภุํ}

\vspace{\tipitakaparskip}
\csromancenter{จิตฺตสหภูทุก เอกํ ขนฺธญฺจ วตฺถุญฺจ ปฏิจฺจ เทฺว ขนฺธา.\csromanspacer{} เทฺว ขนฺเธ จ ฯเปฯ.\csromanspacer{} จิตฺตสหภู ขนฺเธ จ มหาภูเต จ ปฏิจฺจ จิตฺตสหภุ จิตฺตสมุฏฺฐานํ รูปํ, อุปาทารูปํ. (1)}
\prose{\csromanfolio{463}จิตฺตสหภุญฺจ โนจิตฺตสหภุญฺจ ธมฺมํ ปฏิจฺจ โนจิตฺตสหภู ธมฺโม อุปฺปชฺชติ เหตุปจฺจยา.\csromanspacer{} จิตฺตสหภู ขนฺเธ จ จิตฺตญฺจ ปฏิจฺจ โนจิตฺตสหภุ จิตฺตสมุฏฺฐานํ รูปํ.\csromanspacer{} ปฏิสนฺธิกฺขเณ จิตฺตสหภู ขนฺเธ จ จิตฺตญฺจ ปฏิจฺจ กฏตฺตารูปํ.\csromanspacer{} ปฏิสนฺธิกฺขเณ จิตฺตสหภู ขนฺเธ จ วตฺถุญฺจ ปฏิจฺจ จิตฺตํ.\csromanspacer{} จิตฺตสหภู ขนฺเธ จ มหาภูเต จ ปฏิจฺจ โนจิตฺตสหภุ จิตฺตสมุฏฺฐานํ รูปํ, กฏตฺตารูปํ, อุปาทารูปํ.}

\prose{จิตฺตสหภุญฺจ โนจิตฺตสหภุญฺจ ธมฺมํ ปฏิจฺจ จิตฺตสหภู จ โนจิตฺตสหภู จ ธมฺมา อุปฺปชฺชนฺติ เหตุปจฺจยา.\csromanspacer{} จิตฺตสหภุํ เอกํ ขนฺธญฺจ จิตฺตญฺจ ปฏิจฺจ เทฺว ขนฺธา จิตฺตสหภุ จ โนจิตฺตสหภุ จ จิตฺตสมุฏฺฐานํ รูปํ.\csromanspacer{} เทฺว ขนฺเธ จ ฯเปฯ.\csromanspacer{} ปฏิสนฺธิกฺขเณ จิตฺตสหภุํ เอกํ ขนฺธญฺจ จิตฺตญฺจ ปฏิจฺจ เทฺว ขนฺธา กฏตฺตา จ รูปํ.\csromanspacer{} เทฺว ขนฺเธ จ ฯเปฯ.\csromanspacer{} ปฏิสนฺธิกฺขเณ จิตฺตสหภุํ เอกํ ขนฺธญฺจ วตฺถุญฺจ ปฏิจฺจ เทฺว ขนฺธา จิตฺตญฺจ.\csromanspacer{} เทฺว ขนฺเธ จ ฯเปฯ.\csromanspacer{} จิตฺตสหภู ขนฺเธ จ มหาภูเต จ ปฏิจฺจ กฏตฺตารูปํ, อุปาทารูปํ. (สํขิตฺตํ.) (3)}

\csromanheader{2}{1. ปจฺจยานุโลม 2. สงฺขฺยาวาร}
\proseitem{242}{เหตุยา นว, อารมฺมเณ นว, (อรูปํ สพฺพํ อุทฺธริตพฺพํ.\csromanspacer{} จิตฺตสมุฏฺฐานทุกสทิสํ.) อธิปติยา นว, (มหาภูตา ฉสุปิ ปเญฺหสุ กาตพฺพา.\csromanspacer{} อธิปติยา ตีสุ นตฺถิ.) อนนฺตเร นว, สมนนฺตเร นว, สหชาเต นว, อญฺญมญฺเญ นว, นิสฺสเย นว, อุปนิสฺสเย นว, ปุเรชาเต ปญฺจ, อาเสวเน ปญฺจ, กมฺเม นว, (สพฺพตฺถ นว,) อวิคเต นว.}

\csromanheader{2}{2. ปจฺจยปจฺจนีย 2. วิภงฺควาร}
\csromanheader{2}{\textbf{นเหตุ}}
\proseitem{243}{จิตฺตสหภุํ ธมฺมํ ปฏิจฺจ จิตฺตสหภู ธมฺโม อุปฺปชฺชติ นเหตุปจฺจยา.\csromanspacer{} อเหตุกํ จิตฺตสหภุํ เอกํ ขนฺธํ ปฏิจฺจ เทฺว ขนฺธา จิตฺตสหภุ จิตฺตสมุฏฺฐานญฺจ รูปํ.\csromanspacer{} เทฺว ขนฺเธ ฯเปฯ.\csromanspacer{} อเหตุกปฏิสนฺธิกฺขเณ ฯเปฯ.\csromanspacer{} วิจิกิจฺฉาสหคเต \csromanfolio{464}อุทฺธจฺจสหคเต ขนฺเธ ปฏิจฺจ วิจิกิจฺฉาสหคโต อุทฺธจฺจสหคโต โมโห. (เอวํ นวปิ ปญฺหา กาตพฺพา. “อเหตุกนฺ”ติ นิยาเมตพฺพํ.\csromanspacer{} ยถา อนุโลเม ลพฺภติ, เอวํ กาตพฺพํ.\csromanspacer{} ตีณิ, โมโห.\csromanspacer{} ยถา จิตฺตสมุฏฺฐานทุเก, เอวเมว กาตพฺพา.)}

\csromanheader{2}{\textbf{นกมฺม}}
\proseitem{244}{จิตฺตสหภุํ ธมฺมํ ปฏิจฺจ จิตฺตสหภู ธมฺโม อุปฺปชฺชติ นกมฺมปจฺจยา.\csromanspacer{} จิตฺตสหภู ขนฺเธ ปฏิจฺจ จิตฺตสหภู เจตนา.}

\prose{โนจิตฺตสหภุํ ธมฺมํ ปฏิจฺจ โนจิตฺตสหภู ธมฺโม อุปฺปชฺชติ นกมฺมปจฺจยา.\csromanspacer{} พาหิรํ.\csromanspacer{} อาหารสมุฏฺฐานํ.\csromanspacer{} อุตุสมุฏฺฐานํ ฯเปฯ.}

\prose{โนจิตฺตสหภุํ ธมฺมํ ปฏิจฺจ จิตฺตสหภู ธมฺโม อุปฺปชฺชติ นกมฺมปจฺจยา.\csromanspacer{} จิตฺตํ ปฏิจฺจ สมฺปยุตฺตกา เจตนา.}

\prose{จิตฺตสหภุญฺจ โนจิตฺตสหภุญฺจ ธมฺมํ ปฏิจฺจ จิตฺตสหภู ธมฺโม อุปฺปชฺชติ นกมฺมปจฺจยา.\csromanspacer{} จิตฺตสหภู ขนฺเธ จ จิตฺตญฺจ ปฏิจฺจ สมฺปยุตฺตกา เจตนา.}

\csromanheader{2}{\textbf{นฌาน}}
\proseitem{245}{จิตฺตสหภุํ ธมฺมํ ปฏิจฺจ จิตฺตสหภู ธมฺโม อุปฺปชฺชติ นฌานปจฺจยา.\csromanspacer{} ปญฺจวิญฺญาณสหคตํ ฯเปฯ.}

\csromanheader{2}{2. ปจฺจยปจฺจนีย 2. สงฺขฺยาวาร}
\csromanheader{2}{\textbf{สุทฺธ}}
\proseitem{246}{นเหตุยา นว, น-อารมฺมเณ นว, น-อธิปติยา นว, น-อนนฺตเร นว, นสมนนฺตเร นว, น-อญฺญมญฺเญ นว, น-อุปนิสฺสเย นว, นปุเรชาเต นว, นปจฺฉาชาเต นว, น-อาเสวเน นว.\csromanspacer{} นกมฺเม จตฺตาริ, นวิปาเก นว, น-อาหาเร เอกํ, น-อินฺทฺริเย เอกํ, นฌาเน ฉ, นมคฺเค นว, นสมฺปยุตฺเต นว, นวิปฺปยุตฺเต ฉ, โนนตฺถิยา นว, โนวิคเต นว. (เอวํ อิตเร เทฺวปิ คณนา กาตพฺพา.)}

\vspace{\tipitakaparskip}
\csromancenter{(สหชาตวาโรปิ ปฏิจฺจวารสทิโส.)}
\csromanheader{2}{61. จิตฺตสหภูทุก \textbf{61. จิตฺตสหภูทุก 3. ปจฺจยวาร}}
\csromanheader{2}{1. ปจฺจยานุโลม 1. วิภงฺควาร}
\csromanheader{2}{\textbf{เหตุ}}
\proseitem{247}{\csromanfolio{465}จิตฺตสหภุํ ธมฺมํ ปจฺจยา จิตฺตสหภู ธมฺโม อุปฺปชฺชติ เหตุปจฺจยา.\csromanspacer{} ตีณิ. (ปฏิจฺจสทิสา.)}

\prose{โนจิตฺตสหภุํ ธมฺมํ ปจฺจยา โนจิตฺตสหภู ธมฺโม อุปฺปชฺชติ เหตุปจฺจยา.\csromanspacer{} วตฺถุํ ปจฺจยา จิตฺตํ.\csromanspacer{} จิตฺตํ ปจฺจยา โนจิตฺตสหภุ จิตฺตสมุฏฺฐานํ รูปํ.\csromanspacer{} ปฏิสนฺธิกฺขเณ ฯเปฯ. (ปฏิจฺจวารสทิสํ.\csromanspacer{} สพฺเพ มหาภูตา.) (1)}

\prose{โนจิตฺตสหภุํ ธมฺมํ ปจฺจยา จิตฺตสหภู ธมฺโม อุปฺปชฺชติ เหตุปจฺจยา.\csromanspacer{} จิตฺตํ ปจฺจยา สมฺปยุตฺตกา ขนฺธา จิตฺตสหภุ จิตฺตสมุฏฺฐานญฺจ รูปํ.\csromanspacer{} วตฺถุํ ปจฺจยา จิตฺตสหภู ขนฺธา.\csromanspacer{} ปฏิสนฺธิกฺขเณ ฯเปฯ. (สพฺเพ มหาภูตา.\csromanspacer{} ปฏิจฺจสทิสํ.) (2)}

\prose{โนจิตฺตสหภุํ ธมฺมํ ปจฺจยา จิตฺตสหภู จ โนจิตฺตสหภู จ ธมฺมา อุปฺปชฺชนฺติ เหตุปจฺจยา.\csromanspacer{} จิตฺตํ ปจฺจยา สมฺปยุตฺตกา ขนฺธา จิตฺตสหภุ จ โนจิตฺตสหภุ จ จิตฺตสมุฏฺฐานํ รูปํ.\csromanspacer{} วตฺถุํ ปจฺจยา จิตฺตญฺจ สมฺปยุตฺตกา จ ขนฺธา.\csromanspacer{} ปฏิสนฺธิกฺขเณ ฯเปฯ. (ปฏิจฺจสทิสํ.\csromanspacer{} สพฺเพ มหาภูตา.) (3)}

\proseitem{248}{จิตฺตสหภุญฺจ โนจิตฺตสหภุญฺจ ธมฺมํ ปจฺจยา จิตฺตสหภู ธมฺโม อุปฺปชฺชติ เหตุปจฺจยา.\csromanspacer{} จิตฺตสหตุํ เอกํ ขนฺธญฺจ จิตฺตญฺจ ปจฺจยา เทฺว ขนฺธา จิตฺตสหภุ จิตฺตสมุฏฺฐานญฺจ รูปํ.\csromanspacer{} เทฺว ขนฺเธ จ ฯเปฯ.\csromanspacer{} จิตฺตสหภุํ เอกํ ขนฺธญฺจ วตฺถุญฺจ ปจฺจยา เทฺว ขนฺธา.\csromanspacer{} เทฺว ขนฺเธ จ ฯเปฯ.\csromanspacer{} ปฏิสนฺธิกฺขเณ ฯเปฯ. (ปฏิจฺจสทิสํ.\csromanspacer{} สพฺเพ มหาภูตา.) (1)}

\prose{จิตฺตสหภุญฺจ โนจิตฺตสหภุญฺจ ธมฺมํ ปจฺจยา โนจิตฺตสหภู ธมฺโม อุปฺปชฺชติ เหตุปจฺจยา.\csromanspacer{} จิตฺตสหภู ขนฺเธ จ จิตฺตญฺจ ปจฺจยา โนจิตฺตสหภุ จิตฺตสมุฏฺฐานํ รูปํ.\csromanspacer{} จิตฺตสหภู ขนฺเธ จ วตฺถุญฺจ ปจฺจยา จิตฺตํ.\csromanspacer{} จิตฺตสหภู ขนฺเธ จ มหาภูเต จ ปจฺจยา โนจิตฺตสหภุ จิตฺตสมุฏฺฐานํ รูปํ.\csromanspacer{} ปฏิสนฺธิกฺขเณ ฯเปฯ. (ปฏิจฺจสทิสํ.\csromanspacer{} สพฺเพ มหาภูตา.) (2)}

\prose{จิตฺตสหภุญฺจ โนจิตฺตสหภุญฺจ ธมฺมํ ปจฺจยา จิตฺตสหภู จ โนจิตฺตสหภู จ ธมฺมา อุปฺปชฺชนฺติ เหตุปจฺจยา.\csromanspacer{} จิตฺตสหภุํ เอกํ ขนฺธญฺจ จิตฺตญฺจ \csromanfolio{466}ปจฺจยา เทฺว ขนฺธา จิตฺตสหภุ จ โนจิตฺตสหภุ จ จิตฺตสมุฏฺฐานํ รูปํ.\csromanspacer{} จิตฺตสหภุํ เอกํ ขนฺธญฺจ วตฺถุญฺจ ปจฺจยา เทฺว ขนฺธา จิตฺตญฺจ.\csromanspacer{} เทฺว ขนฺเธ จ ฯเปฯ.\csromanspacer{} ปฏิสนฺธิกฺขเณ ฯเปฯ. (ปฏิจฺจสทิสํ.\csromanspacer{} สพฺเพ มหาภูตา.) (3)}

\csromanheader{2}{\textbf{อารมฺมณ}}
\proseitem{249}{จิตฺตสหภุํ ธมฺมํ ปจฺจยา จิตฺตสหภู ธมฺโม อุปฺปชฺชติ อารมฺมณปจฺจยา.\csromanspacer{} ตีณิ. (ปฏิจฺจสทิสํ.)}

\prose{โนจิตฺตสหภุํ ธมฺมํ ปจฺจยา โนจิตฺตสหภู ธมฺโม อุปฺปชฺชติ อารมฺมณปจฺจยา ฯเปฯ.\csromanspacer{} จกฺขายตนํ ปจฺจยา จกฺขุวิญฺญาณํ ฯเปฯ.\csromanspacer{} กายายตนํ ฯเปฯ. (อิมํ จิตฺตสมุฏฺฐานทุกํ ปจฺจยวาเร อารมฺมณสทิสํ ฉนฺนมฺปิ.\csromanspacer{} อิเมสํ ปญฺจวิญฺญาณมูลา กาตพฺพา.\csromanspacer{} สํขิตฺตํ.)}

\csromanheader{2}{1. ปจฺจยานุโลม 2. สงฺขฺยาวาร}
\vspace{\tipitakaparskip}
\csromancenter{เหตุยา นว, อารมฺมเณ นว, อธิปติยา นว, (สพฺพตฺถ นว.) อวิคเต นว.}
\csromanheader{2}{2. ปจฺจยปจฺจนีย 1. วิภงฺควาร}
\csromanheader{2}{\textbf{นเหตุ}}
\proseitem{251}{จิตฺตสหภุํ ธมฺมํ ปจฺจยา จิตฺตสหภู ธมฺโม อุปฺปชฺชติ นเหตุปจฺจยา.\csromanspacer{} อเหตุกํ จิตฺตสหภุํ เอกํ ขนฺธํ ฯเปฯ. (สํขิตฺตํ.\csromanspacer{} สพฺพํ กาตพฺพํ.) (ปจฺจยวารสฺส ปญฺจวิญฺญาณํ ฉนฺนมฺปิ มูลา กาตพฺพา.\csromanspacer{} สพฺเพ มหาภูเต ตีณิเยว.\csromanspacer{} โมโห.\csromanspacer{} สํขิตฺตํ.)}

\csromanheader{2}{2. ปจฺจยปจฺจนีย 2. สงฺขฺยาวาร}
\csromanheader{2}{\textbf{สุทฺธ}}
\vspace{\tipitakaparskip}
\csromancenter{\textbf{นเหตุ}ยา นว, น-อารมฺมเณ นว, น-อธิปติยา นว, น-อนนฺตเร นว, นสมนนฺตเร นว, น-อญฺญมญฺเญ นว, น-อุปนิสฺสเย นว, นปุเรชาเต}
\csromancenter{จิตฺตสหภูทุก นว, นปจฺฉาชาเต นว, น-อาเสวเน นว, นกมฺเม จตฺตาริ, นวิปาเก นว, น-อาหาเร เอกํ, น-อินฺทฺริเย เอกํ, นฌาเน นว, นมคฺเค นว, นสมฺปยุตฺเต นว, นวิปฺปยุตฺเต ฉ, โนนตฺถิยา นว, โนวิคเต นว.}
\csromanheader{2}{3. ปจฺจยานุโลมปจฺจนีย}
\proseitem{253}{\csromanfolio{467}เหตุปจฺจยา น-อารมฺมเณ นว, (สพฺพตฺถ นว.) นกมฺเม ตีณิ, นวิปาเก นว, นสมฺปยุตฺเต นว, นวิปฺปยุตฺเต ปญฺจ, โนนตฺถิยา นว, โนวิคเต นว.}

\csromanheader{2}{4. ปจฺจยปจฺจนียานุโลม}
\proseitem{254}{นเหตุปจฺจยา อารมฺมเณ นว, อนนฺตเร นว, (สพฺพตฺถ นว.) มคฺเค ตีณิ ฯเปฯ อวิคเต นว.}

\vspace{\tipitakaparskip}
\csromancenter{(นิสฺสยวาโร ปจฺจยวารสทิโส.)}
\csromanheader{2}{61. จิตฺตสหภูทุก 5. สํสฏฺฐวาร}
\csromanheader{2}{1. ปจฺจยานุโลมาทิ}
\proseitem{255}{จิตฺตสหภุํ ธมฺมํ สํสฏฺโฐ จิตฺตสหภู ธมฺโม อุปฺปชฺชติ เหตุปจฺจยา.\csromanspacer{} จิตฺตสหภุํ เอกํ ขนฺธํ สํสฏฺฐา เทฺว ขนฺธา.\csromanspacer{} เทฺว ขนฺเธ ฯเปฯ.\csromanspacer{} ปฏิสนฺธิกฺขเณ ฯเปฯ. (1)}

\prose{จิตฺตสหภุํ ธมฺมํ สํสฏฺโฐ โนจิตฺตสหภู ธมฺโม อุปฺปชฺชติ เหตุปจฺจยา.\csromanspacer{} จิตฺตสหภู ขนฺเธ สํสฏฺฐํ จิตฺตํ.\csromanspacer{} ปฏิสนฺธิกฺขเณ ฯเปฯ. (2)}

\prose{จิตฺตสหภุํ ธมฺมํ สํสฏฺโฐ จิตฺตสหภู จ โนจิตฺตสหภู จ ธมฺมา อุปฺปชฺชนฺติ เหตุปจฺจยา.\csromanspacer{} จิตฺตสหภุํ เอกํ ขนฺธํ สํสฏฺฐา เทฺว ขนฺธา จิตฺตญฺจ.\csromanspacer{} เทฺว ขนฺเธ ฯเปฯ.\csromanspacer{} ปฏิสนฺธิกฺขเณ ฯเปฯ. (3)}

\proseitem{256}{โนจิตฺตสหภุํ ธมฺมํ สํสฏฺโฐ จิตฺตสหภู ธมฺโม อุปฺปชฺชติ เหตุปจฺจยา.\csromanspacer{} จิตฺตํ สํสฏฺฐา สมฺปยุตฺตกา ขนฺธา.\csromanspacer{} ปฏิสนฺธิกฺขเณ ฯเปฯ. (1)}

\prose{\csromanfolio{468}จิตฺตสหภุญฺจ โนจิตฺตสหภุญฺจ ธมฺมํ สํสฏฺโฐ จิตฺตสหภู ธมฺโม อุปฺปชฺชติ เหตุปจฺจยา.\csromanspacer{} จิตฺตสหภุํ เอกํ ขนฺธญฺจ จิตฺตญฺจ สํสฏฺฐา เทฺว ขนฺธา.\csromanspacer{} เทฺว ขนฺเธ จ ฯเปฯ.\csromanspacer{} ปฏิสนฺธิกฺขเณ ฯเปฯ. (สํขิตฺตํ.) (2)}

\prose{\textbf{เหตุ}ยา ปญฺจ, อารมฺมเณ ปญฺจ, (สพฺพตฺถ ปญฺจ.) อวิคเต ปญฺจ.\csromanspacer{} อนุโลมํ.}

\prose{นเหตุยา ปญฺจ, (ตีณิ, โมโห.) น-อธิปติยา ปญฺจ, นปุเรชาเต ปญฺจ, นปจฺฉาชาเต ปญฺจ, น-อาเสวเน ปญฺจ, นกมฺเม ตีณิ, นวิปาเก ปญฺจ, นฌาเน ปญฺจ, นมคฺเค ปญฺจ, นวิปฺปยุตฺเต ปญฺจ. (อิตเร เทฺว คณนาปิ สมฺปยุตฺตวาโรปิ สพฺพํ กาตพฺพํ.)}

\csromanheader{2}{61. จิตฺตสหภูทุก 7. ปญฺหาวาร}
\csromanheader{2}{1. ปจฺจยานุโลม 1. วิภงฺควาร}
\csromanheader{2}{\textbf{เหตุ}}
\proseitem{257}{จิตฺตสหภู ธมฺโม จิตฺตสหภุสฺส ธมฺมสฺส เหตุปจฺจเยน ปจฺจโย.\csromanspacer{} จิตฺตสหภู เหตู สมฺปยุตฺตกานํ ขนฺธานํ จิตฺตสหภูนํ จิตฺตสมุฏฺฐานานญฺจ รูปานํ เหตุปจฺจเยน ปจฺจโย.\csromanspacer{} ปฏิสนฺธิกฺขเณ ฯเปฯ. (1)}

\prose{จิตฺตสหภู ธมฺโม โนจิตฺตสหภุสฺส ธมฺมสฺส เหตุปจฺจเยน ปจฺจโย.\csromanspacer{} จิตฺตสหภู เหตู จิตฺตสฺส โนจิตฺตสหภูนํ จิตฺตสมุฏฺฐานานญฺจ รูปานํ เหตุปจฺจเยน ปจฺจโย.\csromanspacer{} ปฏิสนฺธิกฺขเณ จิตฺตสหภู เหตู จิตฺตสฺส กฏตฺตา จ รูปานํ เหตุปจฺจเยน ปจฺจโย. (2)}

\prose{จิตฺตสหภู ธมฺโม จิตฺตสหภุสฺส จ โนจิตฺตสหภุสฺส จ ธมฺมสฺส เหตุปจฺจเยน ปจฺจโย.\csromanspacer{} จิตฺตสหภู เหตู สมฺปยุตฺตกานํ ขนฺธานํ จิตฺตสฺส จ จิตฺตสหภูนญฺจ โนจิตฺตสหภูนญฺจ จิตฺตสมุฏฺฐานานํ รูปานํ เหตุปจฺจเยน ปจฺจโย. (3)}

\csromanheader{2}{\textbf{อรมฺมณ}}
\proseitem{258}{จิตฺตสหภู ธมฺโม จิตฺตสหภุสฺส ธมฺมสฺส อารมฺมณปจฺจเยน ปจฺจโย.\csromanspacer{} นว. (จิตฺตสมุฏฺฐานทุกสทิสํ นินฺนานากรณํ.)}

\csromanheader{2}{61. จิตฺตสหภูทุก \textbf{อธิปติ}}
\proseitem{259}{\csromanfolio{469}จิตฺตสหภู ธมฺโม จิตฺตสหภุสฺส ธมฺมสฺส อธิปติปจฺจเยน ปจฺจโย.\csromanspacer{} ตีณิ. (อารมฺมณาธิปติปิ สหชาตาธิปติปิ กาตพฺพา.)}

\prose{โนจิตฺตสหภู ธมฺโม โนจิตฺตสหภุสฺส ธมฺมสฺส อธิปติปจฺจเยน ปจฺจโย.\csromanspacer{} ตีณิ. (อารมฺมณาธิปติปิ สหชาตาธิปติปิ อิเมสมฺปิ ติณฺณํ กาตพฺพา.\csromanspacer{} นวปิ ปญฺหา จิตฺตสมุฏฺฐานทุกสทิสา.\csromanspacer{} อนฺเต ตีณิ อารมฺมณาธิปติเยว.)}

\csromanheader{2}{\textbf{อนนฺตราทิ}}
\proseitem{260}{จิตฺตสหภู ธมฺโม จิตฺตสหภุสฺส ธมฺมสฺส อนนฺตรปจฺจเยน ปจฺจโย.\csromanspacer{} นว. (จิตฺตสมุฏฺฐานทุกสทิสํ นินฺนานากรณํ.) สมนนฺตรปจฺจเยน ปจฺจโย.\csromanspacer{} นว. (ปฏิจฺจสทิสา.) สหชาตปจฺจเยน ปจฺจโย.\csromanspacer{} นว. (ปฏิจฺจสทิสา.) อญฺญมญฺญปจฺจเยน ปจฺจโย.\csromanspacer{} นว. (ปฏิจฺจสทิสา.) นิสฺสยปจฺจเยน ปจฺจโย.\csromanspacer{} นว. (ปจฺจยวารสทิสา.) อุปนิสฺสยปจฺจเยน ปจฺจโย.\csromanspacer{} นว. (จิตฺตสมุฏฺฐานทุกสทิสา.)}

\csromanheader{2}{\textbf{ปุเรชาต}}
\proseitem{261}{โนจิตฺตสหภู ธมฺโม โนจิตฺตสหภุสฺส ธมฺมสฺส ปุเรชาตปจฺจเยน ปจฺจโย.\csromanspacer{} อารมฺมณปุเรชาตํ, วตฺถุปุเรชาตํ.\csromanspacer{} ตีณิ. (โนจิตฺตสหภุมูลํเยว ลพฺภติ.\csromanspacer{} จิตฺตสมุฏฺฐานทุกสทิสา.\csromanspacer{} ตีณิปิ นินฺนานากรณํ.)}

\csromanheader{2}{\textbf{ปจฺฉาชาตาเสวน}}
\proseitem{262}{จิตฺตสหภู ธมฺโม โนจิตฺตสหภุสฺส ธมฺมสฺส ปจฺฉาชาตปจฺจเยน ปจฺจโย.\csromanspacer{} ปจฺฉาชาตา จิตฺตสหภู ขนฺธา ปุเรชาตสฺส อิมสฺส โนจิตฺตสหภุสฺส กายสฺส ปจฺฉาชาตปจฺจเยน ปจฺจโย. (1)}

\prose{โนจิตฺตสหภู ธมฺโม โนจิตฺตสหภุสฺส ธมฺมสฺส ปจฺฉาชาตปจฺจเยน ปจฺจโย ฯเปฯ. (1)}

\prose{จิตฺตสหภู จ โนจิตฺตสหภู จ ธมฺมา โนจิตฺตสหภุสฺส ธมฺมสฺส ปจฺฉาชาตปจฺจเยน ปจฺจโย. (สํขิตฺตํ.) อาเสวนปจฺจเยน ปจฺจโย.\csromanspacer{} นว.}

\csromanheader{2}{\textbf{กมฺม}}
\proseitem{263}{\csromanfolio{470}จิตฺตสหภู ธมฺโม จิตฺตสหภุสฺส ธมฺมสฺส กมฺมปจฺจเยน ปจฺจโย.\csromanspacer{} \textbf{สหชาตา}, นานากฺขณิกา.\csromanspacer{}\csromanspacer{}\textbf{สหชาตา} จิตฺตสหภู เจตนา สมฺปยุตฺตกานํ ขนฺธานํ จิตฺตสหภูนํ จิตฺตสมุฏฺฐานานญฺจ รูปานํ กมฺมปจฺจเยน ปจฺจโย.\csromanspacer{} ปฏิสนฺธิกฺขเณ ฯเปฯ.\csromanspacer{} \textbf{นานากฺขณิกา} จิตฺตสหภู เจตนา วิปากานํ จิตฺตสหภูนํ ขนฺธานํ กมฺมปจฺจเยน ปจฺจโย. (1)}

\prose{จิตฺตสหภู ธมฺโม โนจิตฺตสหภุสฺส ธมฺมสฺส กมฺมปจฺจเยน ปจฺจโย.\csromanspacer{} \textbf{สหชาตา}, นานากฺขณิกา.\csromanspacer{}\csromanspacer{}\textbf{สหชาตา} จิตฺตสหภู เจตนา จิตฺตสฺส โนจิตฺตสหภูนํ จิตฺตสมุฏฺฐานานญฺจ รูปานํ กมฺมปจฺจเยน ปจฺจโย.\csromanspacer{} ปฏิสนฺธิกฺขเณ ฯเปฯ.\csromanspacer{} \textbf{นานากฺขณิกา} จิตฺตสหภู เจตนา วิปากสฺส จิตฺตสฺส กฏตฺตา จ รูปานํ กมฺมปจฺจเยน ปจฺจโย. (2)}

\prose{จิตฺตสหภู ธมฺโม จิตฺตสหภุสฺส จ โนจิตฺตสหภุสฺส จ ธมฺมสฺส กมฺมปจฺจเยน ปจฺจโย.\csromanspacer{} \textbf{สหชาตา}, นานากฺขณิกา.\csromanspacer{}\csromanspacer{}\textbf{สหชาตา} จิตฺตสหภู เจตนา สมฺปยุตฺตกานํ ขนฺธานํ จิตฺตสฺส จ จิตฺตสหภูนญฺจ โนจิตฺตสหภูนญฺจ จิตฺตสมุฏฺฐานานญฺจ รูปานํ กมฺมปจฺจเยน ปจฺจโย.\csromanspacer{} ปฏิสนฺธิกฺขเณ ฯเปฯ.\csromanspacer{} \textbf{นานากฺขณิกา} จิตฺตสหภู เจตนา วิปากานํ ขนฺธานํ จิตฺตสฺส กฏตฺตา จ รูปานํ กมฺมปจฺจเยน ปจฺจโย. (3)}

\csromanheader{2}{\textbf{วิปากาทิ}}
\proseitem{264}{จิตฺตสหภู ธมฺโม จิตฺตสหภุสฺส ธมฺมสฺส วิปากปจฺจเยน ปจฺจโย. (จิตฺตสมุฏฺฐานทุกสทิสํ.) อาหารปจฺจเยน ปจฺจโย.\csromanspacer{} นว. (จิตฺตสมุฏฺฐานทุกสทิสา.\csromanspacer{} อิมมฺปิ เอกํ กพฬีการ-อาหารสทิสํ.) จิตฺตสหภู ธมฺโม จิตฺตสหภุสฺส ธมฺมสฺส อินฺทฺริยปจฺจเยน ปจฺจโย.\csromanspacer{} นว. (จิตฺตสมุฏฺฐานทุกสทิสํ นินฺนานากรณํ.) ฌานปจฺจเยน ปจฺจโย.\csromanspacer{} ตีณิ.\csromanspacer{} มคฺคปจฺจเยน ปจฺจโย.\csromanspacer{} ตีณิ.\csromanspacer{} สมฺปยุตฺตปจฺจเยน ปจฺจโย.\csromanspacer{} ปญฺจ.}

\csromanheader{2}{\textbf{วิปฺปยุตฺต}}
\proseitem{265}{จิตฺตสหภู ธมฺโม จิตฺตสหภุสฺส ธมฺมสฺส วิปฺปยุตฺตปจฺจเยน ปจฺจโย.\csromanspacer{} \textbf{สหชาตา} จิตฺตสหภู ขนฺธา จิตฺตสหภูนํ จิตฺตสมุฏฺฐานานํ รูปานํ วิปฺปยุตฺตปจฺจเยน ปจฺจโย. (1)}

\csromanheader{2}{61. จิตฺตสหภูทุก}
\prose{\csromanfolio{471}จิตฺตสหภู ธมฺโม โนจิตฺตสหภุสฺส ธมฺมสฺส วิปฺปยุตฺตปจฺจเยน ปจฺจโย.\csromanspacer{} \textbf{สหชาตํ}, ปจฺฉาชาตํ.\csromanspacer{}\csromanspacer{}\textbf{สหชาตา} จิตฺตสหภู ขนฺธา โนจิตฺตสหภูนํ จิตฺตสมุฏฺฐานานํ รูปานํ วิปฺปยุตฺตปจฺจเยน ปจฺจโย.\csromanspacer{} ปฏิสนฺธิกฺขเณ ฯเปฯ.\csromanspacer{} \textbf{ปจฺฉาชาตา} จิตฺตสหภู ขนฺธา ปุเรชาตสฺส อิมสฺส โนจิตฺตสหภุสฺส กายสฺส วิปฺปยุตฺตปจฺจเยน ปจฺจโย. (2)}

\prose{จิตฺตสหภู ธมฺโม จิตฺตสหภุสฺส จ โนจิตฺตสหภุสฺส จ ธมฺมสฺส วิปฺปยุตฺตปจฺจเยน ปจฺจโย.\csromanspacer{} \textbf{สหชาตา} จิตฺตสหภู ขนฺธา จิตฺตสหภูนญฺจ โนจิตฺตสหภูนญฺจ จิตฺตสมุฏฺฐานานํ รูปานํ วิปฺปยุตฺตปจฺจเยน ปจฺจโย. (3)}

\proseitem{266}{โนจิตฺตสหภู ธมฺโม โนจิตฺตสหภุสฺส ธมฺมสฺส วิปฺปยุตฺตปจฺจเยน ปจฺจโย.\csromanspacer{} \textbf{สหชาตํ}, ปุเรชาตํ, ปจฺฉาชาตํ.\csromanspacer{}\csromanspacer{}\textbf{สหชาตํ} จิตฺตํ โนจิตฺตสหภูนํ จิตฺตสมุฏฺฐานานํ รูปานํ วิปฺปยุตฺตปจฺจเยน ปจฺจโย.\csromanspacer{} ปฏิสนฺธิกฺขเณ จิตฺตํ กฏตฺตารูปานํ วิปฺปยุตฺตปจฺจเยน ปจฺจโย.\csromanspacer{} จิตฺตํ วตฺถุสฺส วิปฺปยุตฺตปจฺจเยน ปจฺจโย.\csromanspacer{} วตฺถุ จิตฺตสฺส วิปฺปยุตฺตปจฺจเยน ปจฺจโย.\csromanspacer{} \textbf{ปุเรชาตํ} จกฺขายตนํ จกฺขุวิญฺญาณสฺส วิปฺปยุตฺตปจฺจเยน ปจฺจโย ฯเปฯ.\csromanspacer{} กายายตนํ ฯเปฯ.\csromanspacer{} วตฺถุ จิตฺตสฺส วิปฺปยุตฺตปจฺจเยน ปจฺจโย.\csromanspacer{} \textbf{ปจฺฉาชาตํ} จิตฺตํ ปุเรชาตสฺส อิมสฺส โนจิตฺตสหภุสฺส กายสฺส วิปฺปยุตฺตปจฺจเยน ปจฺจโย. (1)}

\prose{โนจิตฺตสหภู ธมฺโม จิตฺตสหภุสฺส ธมฺมสฺส วิปฺปยุตฺตปจฺจเยน ปจฺจโย.\csromanspacer{} \textbf{สหชาตํ}, ปุเรชาตํ.\csromanspacer{}\csromanspacer{}\textbf{สหชาตํ} จิตฺตํ จิตฺตสหภูนํ จิตฺตสมุฏฺฐานานํ รูปานํ วิปฺปยุตฺตปจฺจเยน ปจฺจโย.\csromanspacer{} \textbf{ปุเรชาตํ} จกฺขายตนํ จกฺขุวิญฺญาณสหคตานํ ขนฺธานํ วิปฺปยุตฺตปจฺจเยน ปจฺจโย ฯเปฯ.\csromanspacer{} กายายตนํ ฯเปฯ.\csromanspacer{} วตฺถุ จิตฺตสหภูนํ ขนฺธานํ วิปฺปยุตฺตปจฺจเยน ปจฺจโย. (2)}

\prose{โนจิตฺตสหภู ธมฺโม จิตฺตสหภุสฺส จ โนจิตฺตสหภุสฺส จ ธมฺมสฺส วิปฺปยุตฺตปจฺจเยน ปจฺจโย.\csromanspacer{} \textbf{สหชาตํ}, ปุเรชาตํ.\csromanspacer{}\csromanspacer{}\textbf{สหชาตํ} จิตฺตํ จิตฺตสหภูนญฺจ โนจิตฺตสหภูนญฺจ จิตฺตสมุฏฺฐานานํ รูปานํ วิปฺปยุตฺตปจฺจเยน ปจฺจโย.\csromanspacer{} ปฏิสนฺธิกฺขเณ ฯเปฯ.\csromanspacer{} \textbf{วตฺถุปุเรชาตํ} จกฺขายตนํ จกฺขุวิญฺญาณสฺส สมฺปยุตฺตกานญฺจ ขนฺธานํ วิปฺปยุตฺตปจฺจเยน ปจฺจโย ฯเปฯ.\csromanspacer{} กายายตนํ ฯเปฯ.\csromanspacer{} วตฺถุ จิตฺตสฺส สมฺปยุตฺตกานญฺจ ขนฺธานํ วิปฺปยุตฺตปจฺจเยน ปจฺจโย. (3)}

\proseitem{267}{\csromanfolio{472}จิตฺตสหภู จ โนจิตฺตสหภู จ ธมฺมา จิตฺตสหภุสฺส ธมฺมสฺส วิปฺปยุตฺตปจฺจเยน ปจฺจโย.\csromanspacer{} \textbf{สหชาตา} จิตฺตสหภู ขนฺธา จ จิตฺตญฺจ จิตฺตสหภูนํ จิตฺตสมุฏฺฐานานํ รูปานํ วิปฺปยุตฺตปจฺจเยน ปจฺจโย. (1)}

\prose{จิตฺตสหภู จ โนจิตฺตสหภู จ ธมฺมา โนจิตฺตสหภุสฺส ธมฺมสฺส วิปฺปยุตฺตปจฺจเยน ปจฺจโย.\csromanspacer{} สหชาตํ, ปจฺฉาชาตํ.\csromanspacer{}\csromanspacer{}\textbf{สหชาตา} จิตฺตสหภู ขนฺธา จ จิตฺตญฺจ โนจิตฺตสหภูนํ จิตฺตสมุฏฺฐานานํ รูปานํ วิปฺปยุตฺตปจฺจเยน ปจฺจโย.\csromanspacer{} ปฏิสนฺธิกฺขเณ ฯเปฯ.\csromanspacer{} \textbf{ปจฺฉาชาตา} จิตฺตสหภู ขนฺธา จ จิตฺตญฺจ ปุเรชาตสฺส อิมสฺส โนจิตฺตสหภุสฺส กายสฺส วิปฺปยุตฺตปจฺจเยน ปจฺจโย.}

\prose{จิตฺตสหภู จ โนจิตฺตสหภู จ ธมฺมา จิตฺตสหภุสฺส จ โนจิตฺตสหภุสฺส จ ธมฺมสฺส วิปฺปยุตฺตปจฺจเยน ปจฺจโย.\csromanspacer{} จิตฺตสหภู ขนฺธา จ จิตฺตญฺจ จิตฺตสหภูนญฺจ โนจิตฺตสหภูนญฺจ จิตฺตสมุฏฺฐานานํ รูปานํ วิปฺปยุตฺตปจฺจเยน ปจฺจโย. (3)}

\csromanheader{2}{\textbf{อตฺถิ}}
\proseitem{268}{จิตฺตสหภู ธมฺโม จิตฺตสหภุสฺส ธมฺมสฺส อตฺถิปจฺจเยน ปจฺจโย.\csromanspacer{} จิตฺตสหภู เอโก ขนฺโธ ฯเปฯ. (ปฏิจฺจสทิสํ.) (1)}

\prose{จิตฺตสหภู ธมฺโม โนจิตฺตสหภุสฺส ธมฺมสฺส อตฺถิปจฺจเยน ปจฺจโย.\csromanspacer{} สหชาตํ, ปจฺฉาชาตํ. (สํขิตฺตํ.) (2)}

\prose{จิตฺตสหภู ธมฺโม จิตฺตสหภุสฺส จ โนจิตฺตสหภุสฺส จ ธมฺมสฺส อตฺถิปจฺจเยน ปจฺจโย.\csromanspacer{} จิตฺตสหภู เอโก ขนฺโธ ฯเปฯ. (ปฏิจฺจสทิสํ.) (3)}

\prose{โนจิตฺตสหภู ธมฺโม โนจิตฺตสหภุสฺส ธมฺมสฺส อตฺถิปจฺจเยน ปจฺจโย.\csromanspacer{} สหชาตํ, ปุเรชาตํ, ปจฺฉาชาตํ, อาหารํ, อินฺทฺริยํ. (สํขิตฺตํ.) (1)}

\prose{โนจิตฺตสหภู ธมฺโม โนจิตฺตสหภุสฺส ธมฺมสฺส อตฺถิปจฺจเยน ปจฺจโย.\csromanspacer{} สหชาตํ, ปุเรชาตํ. (สํขิตฺตํ.) (2)}

\prose{โนจิตฺตสหภู ธมฺโม จิตฺตสหภุสฺส จ โนจิตฺตสหภุสฺส จ ธมฺมสฺส อตฺถิปจฺจเยน ปจฺจโย.\csromanspacer{} สหชาตํ, ปุเรชาตํ. (สํขิตฺตํ.) (3)}

\csromanheader{2}{61. จิตฺตสหภูทุก}
\proseitem{269}{\csromanfolio{473}จิตฺตสหภู จ โนจิตฺตสหภู จ ธมฺมา จิตฺตสหภุสฺส ธมฺมสฺส อตฺถิปจฺจเยน ปจฺจโย.\csromanspacer{} \textbf{สหชาตํ, ปุเรชาตํ}.\csromanspacer{}\csromanspacer{}สหชาโต จกฺขุวิญฺญาณสหคโต เอโก ขนฺโธ จ จกฺขายตนญฺจ จกฺขุวิญฺญาณญฺจ ทฺวินฺนํ ขนฺธานํ ฯเปฯ.\csromanspacer{} เทฺว ขนฺธา จ ฯเปฯ. (สพฺพํ ปฏิสนฺธิยํ กาตพฺพํ.\csromanspacer{} สหชาตํ ปุเรชาตมฺปิ.) (1)}

\prose{จิตฺตสหภู จ โนจิตฺตสหภู จ ธมฺมา โนจิตฺตสหภุสฺส ธมฺมสฺส อตฺถิปจฺจเยน ปจฺจโย.\csromanspacer{} \textbf{สหชาตํ, ปุเรชาตํ}, ปจฺฉาชาตํ, อาหารํ, อินฺทฺริยํ.\csromanspacer{}\csromanspacer{}\textbf{สหชาตา} จกฺขุวิญฺญาณสหคตา ขนฺธา จ จกฺขายตนญฺจ จกฺขุวิญฺญาณสฺส อตฺถิปจฺจเยน ปจฺจโย ฯเปฯ.\csromanspacer{} กายวิญฺญาณสหคตา ฯเปฯ.\csromanspacer{} จิตฺตสหภู ขนฺธา จ จิตฺตญฺจ โนจิตฺตสหภูนํ จิตฺตสมุฏฺฐานานํ รูปานํ อตฺถิปจฺจเยน ปจฺจโย.\csromanspacer{} \textbf{สหชาตา} จิตฺตสหภู ขนฺธา จ วตฺถุ จ จิตฺตสฺส อตฺถิปจฺจเยน ปจฺจโย.\csromanspacer{} \textbf{สหชาตา} จิตฺตสหภู ขนฺธา จ มหาภูตา จ โนจิตฺตสหภูนํ จิตฺตสมุฏฺฐานานํ รูปานํ อตฺถิปจฺจเยน ปจฺจโย. (ปฏิสนฺธิกฺขเณ ตีณิปิ กาตพฺพา.) \textbf{ปจฺฉาชาตา} จิตฺตสหภู ขนฺธา จ จิตฺตญฺจ ปุเรชาตสฺส อิมสฺส โนจิตฺตสหภุสฺส กายสฺส อตฺถิปจฺจเยน ปจฺจโย.\csromanspacer{} \textbf{ปจฺฉาชาตา} จิตฺตสหภู ขนฺธา จ กพฬีกาโร อาหาโร จ อิมสฺส โนจิตฺตสหภุสฺส กายสฺส อตฺถิปจฺจเยน ปจฺจโย.\csromanspacer{} \textbf{ปจฺฉาชาตา} จิตฺตสหภู ขนฺธา จ รูปชีวิตินฺทฺริยญฺจ กฏตฺตารูปานํ อตฺถิปจฺจเยน ปจฺจโย. (2)}

\prose{จิตฺตสหภู จ โนจิตฺตสหภู จ ธมฺมา จิตฺตสหภุสฺส จ โนจิตฺตสหภุสฺส จ ธมฺมสฺส อตฺถิปจฺจเยน ปจฺจโย.\csromanspacer{} \textbf{สหชาตํ, ปุเรชาตํ}.\csromanspacer{}\csromanspacer{}สหชาโต จกฺขุวิญฺญาณสหคโต เอโก ขนฺโธ จ จกฺขายตนญฺจ ฯเปฯ. (ปจฺจยวารสทิสํ.) (3)}

\csromanheader{2}{1. ปจฺจยานุโลม 2. สงฺขฺยาวาร}
\csromanheader{2}{\textbf{สุทฺธ}}
\proseitem{270}{เหตุยา ตีณิ, อารมฺมเณ นว, อธิปติยา นว, อนนฺตเร นว, สมนนฺตเร นว, สหชาเต นว, อญฺญมญฺเญ นว, นิสฺสเย นว, อุปนิสฺสเย นว, ปุเรชาเต ตีณิ, ปจฺฉาชาเต ตีณิ, อาเสวเน นว, \csromanfolio{474}กมฺเม ตีณิ, วิปาเก นว, อาหาเร นว, อินฺทฺริเย นว, ฌาเน ตีณิ, มคฺเค ตีณิ, สมฺปยุตฺเต ปญฺจ, วิปฺปยุตฺเต นว, อตฺถิยา นว, นตฺถิยา นว, วิคเต นว, อวิคเต นว.}

\vspace{\tipitakaparskip}
\csromancenter{อนุโลมํ.}
\csromanheader{2}{2. ปจฺจนียุทฺธาร}
\proseitem{271}{จิตฺตสหภู ธมฺโม จิตฺตสหภุสฺส ธมฺมสฺส อารมฺมณปจฺจเยน ปจฺจโย.\csromanspacer{} สหชาตปจฺจเยน ปจฺจโย.\csromanspacer{} อุปนิสฺสยปจฺจเยน ปจฺจโย.\csromanspacer{} กมฺมปจฺจเยน ปจฺจโย. (1)}

\prose{จิตฺตสหภู ธมฺโม โนจิตฺตสหภุสฺส ธมฺมสฺส อารมฺมณปจฺจเยน ปจฺจโย.\csromanspacer{} สหชาตปจฺจเยน ปจฺจโย.\csromanspacer{} อุปนิสฺสยปจฺจเยน ปจฺจโย.\csromanspacer{} ปจฺฉาชาตปจฺจเยน ปจฺจโย.\csromanspacer{} กมฺมปจฺจเยน ปจฺจโย. (2)}

\prose{จิตฺตสหภู ธมฺโม จิตฺตสหภุสฺส จ โนจิตฺตสหภุสฺส จ ธมฺมสฺส อารมฺมณปจฺจเยน ปจฺจโย.\csromanspacer{} สหชาตปจฺจเยน ปจฺจโย.\csromanspacer{} อุปนิสฺสยปจฺจเยน ปจฺจโย.\csromanspacer{} กมฺมปจฺจเยน ปจฺจโย. (3)}

\proseitem{272}{โนจิตฺตสหภู ธมฺโม โนจิตฺตสหภุสฺส ธมฺมสฺส อารมฺมณปจฺจเยน ปจฺจโย.\csromanspacer{} สหชาตปจฺจเยน ปจฺจโย.\csromanspacer{} อุปนิสฺสยปจฺจเยน ปจฺจโย.\csromanspacer{} ปุเรชาตปจฺจเยน ปจฺจโย.\csromanspacer{} ปจฺฉาชาตปจฺจเยน ปจฺจโย.\csromanspacer{} อาหารปจฺจเยน ปจฺจโย.\csromanspacer{} อินฺทฺริยปจฺจเยน ปจฺจโย. (1)}

\prose{โนจิตฺตสหภู ธมฺโม จิตฺตสหภุสฺส ธมฺมสฺส อารมฺมณปจฺจเยน ปจฺจโย.\csromanspacer{} สหชาตปจฺจเยน ปจฺจโย.\csromanspacer{} อุปนิสฺสยปจฺจเยน ปจฺจโย.\csromanspacer{} ปุเรชาตปจฺจเยน ปจฺจโย. (2)}

\prose{โนจิตฺตสหภู ธมฺโม จิตฺตสหภุสฺส จ โนจิตฺตสหภุสฺส จ ธมฺมสฺส อารมฺมณปจฺจเยน ปจฺจโย.\csromanspacer{} สหชาตปจฺจเยน ปจฺจโย.\csromanspacer{} อุปนิสฺสยปจฺจเยน ปจฺจโย.\csromanspacer{} ปุเรชาตปจฺจเยน ปจฺจโย. (3)}

\csromanmatikaanchor{mtk.61}
\csromantocmarkref{subsection}{62. จิตฺตานุปริวตฺติทุก}{mtk.61}
\bookmark[dest={mtk.61},level=2]{62. จิตฺตานุปริวตฺติทุก}
\csromanheader{2}{62. จิตฺตานุปริวตฺติทุก}
\proseitem{273}{\csromanfolio{475}จิตฺตสหภู จ โนจิตฺตสหภู จ ธมฺมา จิตฺตสหภุสฺส ธมฺมสฺส อารมฺมณปจฺจเยน ปจฺจโย.\csromanspacer{} สหชาตปจฺจเยน ปจฺจโย.\csromanspacer{} อุปนิสฺสยปจฺจเยน ปจฺจโย. (1)}

\prose{จิตฺตสหภู จ โนจิตฺตสหภู จ ธมฺมา โนจิตฺตสหภุสฺส ธมฺมสฺส อารมฺมณปจฺจเยน ปจฺจโย.\csromanspacer{} สหชาตปจฺจเยน ปจฺจโย.\csromanspacer{} อุปนิสฺสยปจฺจเยน ปจฺจโย.\csromanspacer{} ปจฺฉาชาตปจฺจเยน ปจฺจโย. (2)}

\prose{จิตฺตสหภู จ โนจิตฺตสหภู จ ธมฺมา จิตฺตสหภุสฺส จ โนจิตฺตสหภุสฺส จ ธมฺมสฺส อารมฺมณปจฺจเยน ปจฺจโย.\csromanspacer{} สหชาตปจฺจเยน ปจฺจโย.\csromanspacer{} อุปนิสฺสยปจฺจเยน ปจฺจโย. (3)}

\csromanheader{2}{2. ปจฺจยปจฺจนีย 2. สงฺขฺยาวาร}
\vspace{\tipitakaparskip}
\csromancenter{นเหตุยา นว, น-อารมฺมเณ นว, (สพฺพตฺถ นว.) โน-อวิคเต นว.}
\csromanheader{2}{3. ปจฺจยานุโลมปจฺจนีย}
\proseitem{275}{เหตุปจฺจยา น-อารมฺมเณ ตีณิ, น-อธิปติยา ตีณิ, น-อนนฺตเร ตีณิ, นสมนนฺตเร ตีณิ, น-อญฺญมญฺเญ ตีณิ, น-อุปนิสฺสเย ตีณิ, (สพฺพตฺถ ตีณิ.) นสมฺปยุตฺเต ตีณิ, นวิปฺปยุตฺเต ตีณิ, โนนตฺถิยา ตีณิ, โนวิคเต ตีณิ.}

\csromanheader{2}{4. ปจฺจยปจฺจนียานุโลม}
\csromanheader{2}{276. นเหตุปจฺจยา อารมฺมเณ นว, อธิปติยา นว. (อนุโลมมาติกา กาตพฺพา.)}
\nitthitam{จิตฺตสหภูทุกํ นิฏฺฐิตํ.}
\csromanheader{2}{62. จิตฺตานุปริวตฺติทุก 1. ปฏิจฺจวาร}
\proseitem{277}{จิตฺตานุปริวตฺติํ ธมฺมํ ปฏิจฺจ จิตฺตานุปริวตฺตี ธมฺโม อุปฺปชฺชติ เหตุปจฺจยา.\csromanspacer{} จิตฺตานุปริวตฺติํ เอกํ ขนฺธํ ปฏิจฺจ เทฺว ขนฺธา จิตฺตานุปริวตฺติํ \csromanfolio{476}จิตฺตสมุฏฺฐานญฺจ รูปํ.\csromanspacer{} เทฺว ขนฺเธ ฯเปฯ.\csromanspacer{} ปฏิสนฺธิกฺขเณ ฯเปฯ. (ยถา จิตฺตสหภูทุกํ, เอวํ อิมํ ทุกํ กาตพฺพํ, นินฺนานากรณํ.)}

\nitthitam{จิตฺตานุปริวตฺติทุกํ นิฏฺฐิตํ.}
\csromanmatikaanchor{mtk.62}
\csromantocmarkref{subsection}{63. จิตฺตสํสฏฺฐสมุฏฺฐานทุก}{mtk.62}
\bookmark[dest={mtk.62},level=2]{63. จิตฺตสํสฏฺฐสมุฏฺฐานทุก}
\csromanheader{2}{63. จิตฺตสํสฏฺฐสมุฏฺฐานทุก 1. ปฏิจฺจวาร}
\csromanheader{2}{1. ปจฺจยานุโลม 1. วิภงฺควาร}
\csromanheader{2}{\textbf{เหตุ}}
\proseitem{278}{จิตฺตสํสฏฺฐสมุฏฺฐานํ ธมฺมํ ปฏิจฺจ จิตฺตสํสฏฺฐสมุฏฺฐาโน ธมฺโม อุปฺปชฺชติ เหตุปจฺจยา.\csromanspacer{} จิตฺตสํสฏฺฐสมุฏฺฐานํ เอกํ ขนฺธํ ปฏิจฺจ เทฺว ขนฺธา.\csromanspacer{} เทฺว ขนฺเธ ปฏิจฺจ เอโก ขนฺโธ.\csromanspacer{} ปฏิสนฺธิกฺขเณ ฯเปฯ. (1)}

\prose{จิตฺตสํสฏฺฐสมุฏฺฐานํ ธมฺมํ ปฏิจฺจ โนจิตฺตสํสฏฺฐสมุฏฺฐาโน ธมฺโม อุปฺปชฺชติ เหตุปจฺจยา.\csromanspacer{} จิตฺตสํสฏฺฐสมุฏฺฐาเน ขนฺเธ ปฏิจฺจ จิตฺตํ จิตฺตสมุฏฺฐานญฺจ รูปํ.\csromanspacer{} ปฏิสนฺธิกฺขเณ จิตฺตสํสฏฺฐสมุฏฺฐาเน ขนฺเธ ปฏิจฺจ จิตฺตํ กฏตฺตา จ รูปํ. (2)}

\prose{จิตฺตสํสฏฺฐสมุฏฺฐานํ ธมฺมํ ปฏิจฺจ จิตฺตสํสฏฺฐสมุฏฺฐาโน จ โนจิตฺตสํสฏฺฐสมุฏฺฐาโน จ ธมฺมา อุปฺปชฺชนฺติ เหตุปจฺจยา.\csromanspacer{} จิตฺตสํสฏฺฐสมุฏฺฐานํ เอกํ ขนฺธํ ปฏิจฺจ เทฺว ขนฺธา จิตฺตญฺจ จิตฺตสมุฏฺฐานญฺจ รูปํ.\csromanspacer{} เทฺว ขนฺเธ ฯเปฯ.\csromanspacer{} ปฏิสนฺธิกฺขเณ ฯเปฯ. (3)}

\proseitem{279}{โนจิตฺตสํสฏฺฐสมุฏฺฐานํ ธมฺมํ ปฏิจฺจ โนจิตฺตสํสฏฺฐสมุฏฺฐาโน ธมฺโม อุปฺปชฺชติ เหตุปจฺจยา.\csromanspacer{} จิตฺตํ ปฏิจฺจ จิตฺตสมุฏฺฐานํ รูปํ.\csromanspacer{} ปฏิสนฺธิกฺขเณ จิตฺตํ ปฏิจฺจ กฏตฺตารูปํ.\csromanspacer{} จิตฺตํ ปฏิจฺจ วตฺถุ, วตฺถุํ ปฏิจฺจ จิตฺตํ.\csromanspacer{} เอกํ มหาภูตํ ฯเปฯ มหาภูเต ปฏิจฺจ จิตฺตสมุฏฺฐานํ รูปํ, กฏตฺตารูปํ, อุปาทารูปํ. (1)}

\prose{โนจิตฺตสํสฏฺฐสมุฏฺฐานํ ธมฺมํ ปฏิจฺจ จิตฺตสํสฏฺฐสมุฏฺฐาโน ธมฺโม อุปฺปชฺชติ เหตุปจฺจยา.\csromanspacer{} จิตฺตํ ปฏิจฺจ สมฺปยุตฺตกา ขนฺธา.\csromanspacer{} ปฏิสนฺธิกฺขเณ จิตฺตํ ปฏิจฺจ สมฺปยุตฺตกา ขนฺธา.\csromanspacer{} ปฏิสนฺธิกฺขเณ วตฺถุํ ปฏิจฺจ จิตฺตสํสฏฺฐสมุฏฺฐานา ขนฺธา. (2)}

\csromanheader{2}{63. จิตฺตสํสฏฺฐสมุฏฺฐานทุก}
\prose{\csromanfolio{477}โนจิตฺตสํสฏฺฐสมุฏฺฐานํ ธมฺมํ ปฏิจฺจ จิตฺตสํสฏฺฐสมุฏฺฐาโน จ โนจิตฺตสํสฏฺฐสมุฏฺฐาโน จ ธมฺมา อุปฺปชฺชนฺติ เหตุปจฺจยา.\csromanspacer{} จิตฺตํ ปฏิจฺจ สมฺปยุตฺตกา ขนฺธา จิตฺตสมุฏฺฐานญฺจ รูปํ.\csromanspacer{} ปฏิสนฺธิกฺขเณ จิตฺตํ ปฏิจฺจ สมฺปยุตฺตกา ขนฺธา กฏตฺตา จ รูปํ.\csromanspacer{} ปฏิสนฺธิกฺขเณ วตฺถุํ ปฏิจฺจ จิตฺตํ สมฺปยุตฺตกา จ ขนฺธา. (3)}

\proseitem{280}{จิตฺตสํสฏฺฐสมุฏฺฐานญฺจ โนจิตฺตสํสฏฺฐสมุฏฺฐานญฺจ ธมฺมํ ปฏิจฺจ จิตฺตสํสฏฺฐสมุฏฺฐาโน ธมฺโม อุปฺปชฺชติ เหตุปจฺจยา.\csromanspacer{} จิตฺตสํสฏฺฐสมุฏฺฐานํ เอกํ ขนฺธญฺจ จิตฺตญฺจ ปฏิจฺจ เทฺว ขนฺธา.\csromanspacer{} เทฺว ขนฺเธ ฯเปฯ.\csromanspacer{} ปฏิสนฺธิกฺขเณ จิตฺตสํสฏฺฐสมุฏฺฐานํ เอกํ ขนฺธญฺจ จิตฺตญฺจ ปฏิจฺจ เทฺว ขนฺธา.\csromanspacer{} เทฺว ขนฺเธ ฯเปฯ.\csromanspacer{} ปฏิสนฺธิกฺขเณ จิตฺตสํสฏฺฐสมุฏฺฐานํ เอกํ ขนฺธญฺจ วตฺถุญฺจ ปฏิจฺจ เทฺว ขนฺธา.\csromanspacer{} เทฺว ขนฺเธ ฯเปฯ. (1)}

\prose{จิตฺตสํสฏฺฐสมุฏฺฐานญฺจ โนจิตฺตสํสฏฺฐสมุฏฺฐานญฺจ ธมฺมํ ปฏิจฺจ โนจิตฺตสํสฏฺฐสมุฏฺฐาโน ธมฺโม อุปฺปชฺชติ เหตุปจฺจยา.\csromanspacer{} จิตฺตสํสฏฺฐสมุฏฺฐาเน ขนฺเธ จ จิตฺตญฺจ ปฏิจฺจ จิตฺตสมุฏฺฐานํ รูปํ.\csromanspacer{} จิตฺตสํสฏฺฐสมุฏฺฐาเน ขนฺเธ จ มหาภูเต จ ปฏิจฺจ จิตฺตสมุฏฺฐานํ รูปํ.\csromanspacer{} ปฏิสนฺธิกฺขเณ จิตฺตสํสฏฺฐสมุฏฺฐาเน ขนฺเธ จ จิตฺตญฺจ ปฏิจฺจ กฏตฺตารูปํ.\csromanspacer{} ปฏิสนฺธิกฺขเณ จิตฺตสํสฏฺฐสมุฏฺฐาเน ขนฺเธ จ มหาภูเต จ ปฏิจฺจ กฏตฺตารูปํ.\csromanspacer{} ปฏิสนฺธิกฺขเณ จิตฺตสํสฏฺฐสมุฏฺฐาเน ขนฺเธ จ วตฺถุญฺจ ปฏิจฺจ จิตฺตํ. (2)}

\prose{จิตฺตสํสฏฺฐสมุฏฺฐานญฺจ โนจิตฺตสํสฏฺฐสมุฏฺฐานญฺจ ธมฺมํ ปฏิจฺจ จิตฺตสํสฏฺฐสมุฏฺฐาโน จ โนจิตฺตสํสฏฺฐสมุฏฺฐาโน จ ธมฺมา อุปฺปชฺชนฺติ เหตุปจฺจยา.\csromanspacer{} จิตฺตสํสฏฺฐสมุฏฺฐานํ เอกํ ขนฺธญฺจ จิตฺตญฺจ ปฏิจฺจ เทฺว ขนฺธา จิตฺตสมุฏฺฐานญฺจ รูปํ.\csromanspacer{} เทฺว ขนฺเธ ฯเปฯ.\csromanspacer{} ปฏิสนฺธิกฺขเณ จิตฺตสํสฏฺฐสมุฏฺฐานํ เอกํ ขนฺธญฺจ จิตฺตญฺจ ปฏิจฺจ เทฺว ขนฺธา กฏตฺตา จ รูปํ.\csromanspacer{} เทฺว ขนฺเธ ฯเปฯ.\csromanspacer{} ปฏิสนฺธิกฺขเณ จิตฺตสํสฏฺฐสมุฏฺฐานํ เอกํ ขนฺธญฺจ วตฺถุญฺจ ปฏิจฺจ เทฺว ขนฺธา จิตฺตญฺจ.\csromanspacer{} เทฺว ขนฺเธ ฯเปฯ. (สํขิตฺตํ.) (3)}

\csromanheader{2}{1. ปจฺจยานุโลม 2. สงฺขฺยาวาร}
\csromanheader{2}{\textbf{สุทฺธ}}
\proseitem{281}{เหตุยา นว, อารมฺมเณ นว, อธิปติยา นว, อนนฺตเร นว, สมนนฺตเร นว, สหชาเต นว, อญฺญมญฺเญ นว, นิสฺสเย นว, อุปนิสฺสเย นว, ปุเรชาเต ปญฺจ, อาเสวเน ปญฺจ, กมฺเม นว, วิปาเก นว, (สพฺพตฺถ นว.) อวิคเต นว.}

\vspace{\tipitakaparskip}
\csromancenter{อนุโลมํ.}
\csromanheader{2}{2. ปจฺจยปจฺจนีย 1. วิภงฺควาร}
\csromanheader{2}{\textbf{นเหตุ}}
\proseitem{282}{\csromanfolio{478}จิตฺตสํสฏฺฐสมุฏฺฐานํ ธมฺมํ ปฏิจฺจ จิตฺตสํสฏฺฐสมุฏฺฐาโน ธมฺโม อุปฺปชฺชติ นเหตุปจฺจยา.\csromanspacer{} อเหตุกํ จิตฺตสํสฏฺฐสมุฏฺฐานํ เอกํ ขนฺธํ ปฏิจฺจ เทฺว ขนฺธา.\csromanspacer{} เทฺว ขนฺเธ ฯเปฯ.\csromanspacer{} อเหตุกปฏิสนฺธิกฺขเณ ฯเปฯ.\csromanspacer{} วิจิกิจฺฉาสหคเต อุทฺธจฺจสหคเต ขนฺเธ ปฏิจฺจ วิจิกิจฺฉาสหคโต อุทฺธจฺจสหคโต โมโห. (1)}

\prose{จิตฺตสํสฏฺฐสมุฏฺฐานํ ธมฺมํ ปฏิจฺจ โนจิตฺตสํสฏฺฐสมุฏฺฐาโน ธมฺโม อุปฺปชฺชติ นเหตุปจฺจยา.\csromanspacer{} อเหตุเก จิตฺตสํสฏฺฐสมุฏฺฐาเน ขนฺเธ ปฏิจฺจ จิตฺตํ จิตฺตสมุฏฺฐานญฺจ รูปํ.\csromanspacer{} อเหตุกปฏิสนฺธิกฺขเณ ฯเปฯ. (2)}

\prose{จิตฺตสํสฏฺฐสมุฏฺฐานํ ธมฺมํ ปฏิจฺจ จิตฺตสํสฏฺฐสมุฏฺฐาโน จ โนจิตฺตสํสฏฺฐสมุฏฺฐาโน จ ธมฺมา อุปฺปชฺชนฺติ นเหตุปจฺจยา.\csromanspacer{} อเหตุกํ จิตฺตสํสฏฺฐสมุฏฺฐานํ เอกํ ขนฺธํ ปฏิจฺจ เทฺว ขนฺธา จิตฺตญฺจ จิตฺตสมุฏฺฐานญฺจ รูปํ.\csromanspacer{} เทฺว ขนฺเธ ฯเปฯ.\csromanspacer{} อเหตุกปฏิสนฺธิกฺขเณ ฯเปฯ. (3)}

\proseitem{283}{โนจิตฺตสํสฏฺฐสมุฏฺฐานํ ธมฺมํ ปฏิจฺจ โนจิตฺตสํสฏฺฐสมุฏฺฐาโน ธมฺโม อุปฺปชฺชติ นเหตุปจฺจยา.\csromanspacer{} อเหตุกํ จิตฺตํ ปฏิจฺจ จิตฺตสมุฏฺฐานํ รูปํ.\csromanspacer{} อเหตุกปฏิสนฺธิกฺขเณ จิตฺตํ ปฏิจฺจ กฏตฺตารูปํ.\csromanspacer{} จิตฺตํ ปฏิจฺจ วตฺถุ, วตฺถุํ ปฏิจฺจ จิตฺตํ.\csromanspacer{} เอกํ มหาภูตํ ฯเปฯ. (ยาว อสญฺญสตฺตา.) (1)}

\prose{โนจิตฺตสํสฏฺฐสมุฏฺฐานํ ธมฺมํ ปฏิจฺจ จิตฺตสํสฏฺฐสมุฏฺฐาโน ธมฺโม อุปฺปชฺชติ นเหตุปจฺจยา.\csromanspacer{} อเหตุกํ จิตฺตํ ปฏิจฺจ สมฺปยุตฺตกา ขนฺธา.\csromanspacer{} อเหตุกปฏิสนฺธิกฺขเณ จิตฺตํ ปฏิจฺจ สมฺปยุตฺตกา ขนฺธา.\csromanspacer{} อเหตุกปฏิสนฺธิกฺขเณ วตฺถุํ ปฏิจฺจ จิตฺตสํสฏฺฐสมุฏฺฐานา ขนฺธา.\csromanspacer{} วิจิกิจฺฉาสหคตํ อุทฺธจฺจสหคตํ จิตฺตํ ปฏิจฺจ วิจิกิจฺฉาสหคโต อุทฺธจฺจสหคโต โมโห. (2)}

\prose{โนจิตฺตสํสฏฺฐสมุฏฺฐานํ ธมฺมํ ปฏิจฺจ จิตฺตสํสฏฺฐสมุฏฺฐาโน จ โนจิตฺตสํสฏฺฐสมุฏฺฐาโน จ ธมฺมา อุปฺปชฺชนฺติ นเหตุปจฺจยา.\csromanspacer{} อเหตุกํ จิตฺตํ ปฏิจฺจ สมฺปยุตฺตกา ขนฺธา จิตฺตสมุฏฺฐานญฺจ รูปํ.\csromanspacer{} อเหตุกปฏิสนฺธิกฺขเณ จิตฺตํ ปฏิจฺจ สมฺปยุตฺตกา ขนฺธา กฏตฺตา จ รูปํ.\csromanspacer{} อเหตุกปฏิสนฺธิกฺขเณ วตฺถุํ ปฏิจฺจ จิตฺตํ สมฺปยุตฺตกา จ ขนฺธา. (3)}

\csromanheader{2}{63. จิตฺตสํสฏฺฐสมุฏฺฐานทุก}
\proseitem{284}{\csromanfolio{479}จิตฺตสํสฏฺฐสมุฏฺฐานญฺจ โนจิตฺตสํสฏฺฐสมุฏฺฐานญฺจ ธมฺมํ ปฏิจฺจ จิตฺตสํสฏฺฐสมุฏฺฐาโน ธมฺโม อุปฺปชฺชติ นเหตุปจฺจยา.\csromanspacer{} อเหตุกํ จิตฺตสํสฏฺฐสมุฏฺฐานํ เอกํ ขนฺธญฺจ จิตฺตญฺจ ปฏิจฺจ เทฺว ขนฺธา.\csromanspacer{} เทฺว ขนฺเธ ฯเปฯ.\csromanspacer{} อเหตุกปฏิสนฺธิกฺขเณ จิตฺตสํสฏฺฐสมุฏฺฐานํ เอกํ ขนฺธญฺจ จิตฺตญฺจ ฯเปฯ.\csromanspacer{} อเหตุกปฏิสนฺธิกฺขเณ จิตฺตสํสฏฺฐสมุฏฺฐานํ เอกํ ขนฺธญฺจ วตฺถุญฺจ ปฏิจฺจ เทฺว ขนฺธา.\csromanspacer{} เทฺว ขนฺเธ ฯเปฯ.\csromanspacer{} วิจิกิจฺฉาสหคเต อุทฺธจฺจสหคเต ขนฺเธ จ จิตฺตญฺจ ปฏิจฺจ วิจิกิจฺฉาสคหโต อุทฺธจฺจสหคโต โมโห. (1)}

\prose{จิตฺตสํสฏฺฐสมุฏฺฐานญฺจ โนจิตฺตสํสฏฺฐสมุฏฺฐานญฺจ ธมฺมํ ปฏิจฺจ โนจิตฺตสํสฏฺฐสมุฏฺฐาโน ธมฺโม อุปฺปชฺชติ นเหตุปจฺจยา.\csromanspacer{} อเหตุเก จิตฺตสํสฏฺฐสมุฏฺฐาเน ขนฺเธ จ จิตฺตญฺจ ปฏิจฺจ จิตฺตสมุฏฺฐานํ รูปํ.\csromanspacer{} อเหตุเก จิตฺตสํสฏฺฐสมุฏฺฐาเน ขนฺเธ จ มหาภูเต จ ปฏิจฺจ จิตฺตสมุฏฺฐานํ รูปํ.\csromanspacer{} อเหตุกปฏิสนฺธิกฺขเณ จิตฺตสํสฏฺฐสมุฏฺฐาเน ขนฺเธ จ จิตฺตญฺจ ปฏิจฺจ กฏตฺตารูปํ.\csromanspacer{} อเหตุกปฏิสนฺธิกฺขเณ จิตฺตสํสฏฺฐสมุฏฺฐาเน ขนฺเธ จ มหาภูเต จ ปฏิจฺจ กฏตฺตารูปํ.\csromanspacer{} อเหตุกปฏิสนฺธิกฺขเณ จิตฺตสํสฏฺฐสมุฏฺฐาเน ขนฺเธ จ วตฺถุญฺจ ปฏิจฺจ จิตฺตํ. (2)}

\prose{จิตฺตสํสฏฺฐสมุฏฺฐานญฺจ โนจิตฺตสํสฏฺฐสมุฏฺฐานญฺจ ธมฺมํ ปฏิจฺจ จิตฺตสํสฏฺฐสมุฏฺฐาโน จ โนจิตฺตสํสฏฺฐสมุฏฺฐาโน จ ธมฺมา อุปฺปชฺชนฺติ นเหตุปจฺจยา.\csromanspacer{} อเหตุกํ จิตฺตสํสฏฺฐสมุฏฺฐานํ เอกํ ขนฺธญฺจ จิตฺตญฺจ ปฏิจฺจ เทฺว ขนฺธา จิตฺตสมุฏฺฐานญฺจ รูปํ.\csromanspacer{} เทฺว ขนฺเธ ฯเปฯ. (อเหตุกปฏิสนฺธิกฺขเณ เทฺวปิ กาตพฺพา.) (3)}

\csromanheader{2}{\textbf{น-อารมฺมณ}}
\proseitem{285}{จิตฺตสํสฏฺฐสมุฏฺฐานํ ธมฺมํ ปฏิจฺจ โนจิตฺตสํสฏฺฐสมุฏฺฐาโน ธมฺโม อุปฺปชฺชติ น-อารมฺมณปจฺจยา.\csromanspacer{} จิตฺตสํสฏฺฐสมุฏฺฐาเน ขนฺเธ ปฏิจฺจ จิตฺตสมุฏฺฐานํ รูปํ.\csromanspacer{} ปฏิสนฺธิกฺขเณ ฯเปฯ. (1)}

\prose{โนจิตฺตสํสฏฺฐสมุฏฺฐานํ ธมฺมํ ปฏิจฺจ โนจิตฺตสํสฏฺฐสมุฏฺฐาโน ธมฺโม อุปฺปชฺชติ น-อารมฺมณปจฺจยา.\csromanspacer{} จิตฺตํ ปฏิจฺจ จิตฺตสมุฏฺฐานํ รูปํ.\csromanspacer{} ปฏิสนฺธิกฺขเณ จิตฺตํ ปฏิจฺจ กฏตฺตารูปํ.\csromanspacer{} จิตฺตํ ปฏิจฺจ วตฺถุ.\csromanspacer{} เอกํ มหาภูตํ ฯเปฯ. (ยาว อสญฺญสตฺตา.) (1)}

\prose{จิตฺตสํสฏฺฐสมุฏฺฐานญฺจ โนจิตฺตสํสฏฺฐสมุฏฺฐานญฺจ ธมฺมํ ปฏิจฺจ โนจิตฺตสํสฏฺฐสมุฏฺฐาโน ธมฺโม อุปฺปชฺชติ น-อารมฺมณปจฺจยา.\csromanspacer{} จิตฺตสํสฏฺฐสมุฏฺฐาเน \csromanfolio{480}ขนฺเธ จ จิตฺตญฺจ ปฏิจฺจ จิตฺตสมุฏฺฐานํ รูปํ.\csromanspacer{} จิตฺตสํสฏฺฐสมุฏฺฐาเน ขนฺเธ จ มหาภูเต จ ปฏิจฺจ จิตฺตสมุฏฺฐานํ รูปํ. (ปฏิสนฺธิกฺขเณ เทฺว.\csromanspacer{} สํขิตฺตํ.) (1)}

\csromanheader{2}{2. ปจฺจยปจฺจนีย 2. สงฺขฺยาวาร}
\csromanheader{2}{\textbf{สุทฺธ}}
\proseitem{286}{นเหตุยา นว, น-อารมฺมเณ ตีณิ, น-อธิปติยา นว, น-อนนฺตเร ตีณิ, นสมนนฺตเร ตีณิ, น-อญฺญมญฺเญ ตีณิ, น-อุปนิสฺสเย ตีณิ, นปุเรชาเต นว, นปจฺฉาชาเต นว, น-อาเสวเน นว, นกมฺเม จตฺตาริ, นวิปาเก นว, น-อาหาเร เอกํ, น-อินฺทฺริเย เอกํ, นฌาเน ฉ, นมคฺเค นว, นสมฺปยุตฺเต ตีณิ, นวิปฺปยุตฺเต ฉ, โนนตฺถิยา ตีณิ, โนวิคเต ตีณิ.}

\csromanheader{2}{3. ปจฺจยานุโลมปจฺจนีย}
\csromanheader{2}{287. \textbf{เหตุ}ปจฺจยา น-อารมฺมเณ ตีณิ. (สํขิตฺตํ.)}
\csromanheader{2}{4. ปจฺจยปจฺจนียานุโลม}
\proseitem{288}{นเหตุปจฺจยา อารมฺมเณ นว, อนนฺตเร นว. (สํขิตฺตํ.)}

\vspace{\tipitakaparskip}
\csromancenter{(สหชาตวาโร ปฏิจฺจวารสทิโส.)}
\csromanheader{2}{63. จิตฺตสํสฏฺฐสมุฏฺฐานทุก 3. ปจฺจยวาร}
\csromanheader{2}{1. ปจฺจยานุโลม 1. วิภงฺควาร}
\csromanheader{2}{\textbf{เหตุ}}
\proseitem{289}{จิตฺตสํสฏฺฐสมุฏฺฐานํ ธมฺมํ ปจฺจยา จิตฺตสํสฏฺฐสมุฏฺฐาโน ธมฺโม อุปฺปชฺชติ เหตุปจฺจยา. (สํขิตฺตํ.) ตีณิ. (ปฏิจฺจวารสทิสา.)}

\csromanheader{2}{63. จิตฺตสํสฏฺฐสมุฏฺฐานทุก}
\prose{\csromanfolio{481}โนจิตฺตสํสฏฺฐสมุฏฺฐานํ ธมฺมํ ปจฺจยา โนจิตฺตสํสฏฺฐสมุฏฺฐาโน ธมฺโม อุปฺปชฺชติ เหตุปจฺจยา.\csromanspacer{} จิตฺตํ ปจฺจยา จิตฺตสมุฏฺฐานํ รูปํ.\csromanspacer{} วตฺถุํ ปจฺจยา จิตฺตํ.\csromanspacer{} ปฏิสนฺธิกฺขเณ ฯเปฯ. (ยาว มหาภูตา.) (1)}

\prose{โนจิตฺตสํสฏฺฐสมุฏฺฐานํ ธมฺมํ ปจฺจยา จิตฺตสํสฏฺฐสมุฏฺฐาโน ธมฺโม อุปฺปชฺชติ เหตุปจฺจยา.\csromanspacer{} จิตฺตํ ปจฺจยา สมฺปยุตฺตกา ขนฺธา.\csromanspacer{} วตฺถุํ ปจฺจยา จิตฺตสํสฏฺฐสมุฏฺฐานา ขนฺธา. (ปฏิสนฺธิกฺขเณ เทฺวปิ กาตพฺพา.) ตีณิ. (3)}

\proseitem{290}{จิตฺตสํสฏฺฐสมุฏฺฐานญฺจ โนจิตฺตสํสฏฺฐสมุฏฺฐานญฺจ ธมฺมํ ปจฺจยา จิตฺตสํสฏฺฐสมุฏฺฐาโน ธมฺโม อุปฺปชฺชติ เหตุปจฺจยา.\csromanspacer{} จิตฺตสํสฏฺฐสมุฏฺฐานํ เอกํ ขนฺธญฺจ จิตฺตญฺจ ปจฺจยา เทฺว ขนฺธา.\csromanspacer{} เทฺว ขนฺเธ จ ฯเปฯ.\csromanspacer{} จิตฺตสํสฏฺฐสมุฏฺฐานํ เอกํ ขนฺธญฺจ วตฺถุญฺจ ปจฺจยา เทฺว ขนฺธา.\csromanspacer{} เทฺว ขนฺเธ จ ฯเปฯ. (ปฏิสนฺธิกฺขเณ เทฺวปิ กาตพฺพา.) (1)}

\prose{จิตฺตสํสฏฺฐสมุฏฺฐานญฺจ โนจิตฺตสํสฏฺฐสมุฏฺฐานญฺจ ธมฺมํ ปจฺจยา โนจิตฺตสํสฏฺฐสมุฏฺฐาโน ธมฺโม อุปฺปชฺชติ เหตุปจฺจยา.\csromanspacer{} จิตฺตสํสฏฺฐสมุฏฺฐาเน ขนฺเธ จ จิตฺตญฺจ ปจฺจยา โนจิตฺตสํสฏฺฐสมุฏฺฐานํ รูปํ.\csromanspacer{} จิตฺตสํสฏฺฐสมุฏฺฐาเน ขนฺเธ จ มหาภูเต จ ปจฺจยา โนจิตฺตสํสฏฺฐสมุฏฺฐานํ รูปํ.\csromanspacer{} จิตฺตสํสฏฺฐสมุฏฺฐาเน ขนฺเธ จ วตฺถุญฺจ ปจฺจยา จิตฺตํ. (ปฏิสนฺธิกฺขเณ ตีณิปิ กาตพฺพา.) (2)}

\prose{จิตฺตสํสฏฺฐสมุฏฺฐานญฺจ โนจิตฺตสํสฏฺฐสมุฏฺฐานญฺจ ธมฺมํ ปจฺจยา จิตฺตสํสฏฺฐสมุฏฺฐาโน จ โนจิตฺตสํสฏฺฐสมุฏฺฐาโน จ ธมฺมา อุปฺปชฺชนฺติ เหตุปจฺจยา.\csromanspacer{} จิตฺตสํสฏฺฐสมุฏฺฐานํ เอกํ ขนฺธญฺจ จิตฺตญฺจ ปจฺจยา เทฺว ขนฺธา จิตฺตสํสฏฺฐสมุฏฺฐานญฺจ รูปํ.\csromanspacer{} เทฺว ขนฺเธ ฯเปฯ.\csromanspacer{} จิตฺตสํสฏฺฐสมุฏฺฐานํ เอกํ ขนฺธญฺจ วตฺถุญฺจ ปจฺจยา เทฺว ขนฺธา จิตฺตญฺจ.\csromanspacer{} เทฺว ขนฺเธ ฯเปฯ. (ปฏิสนฺธิกฺขเณ เทฺวปิ กาตพฺพา. (3)}

\csromanheader{2}{\textbf{อรมฺมณ}}
\proseitem{291}{จิตฺตสํสฏฺฐสมุฏฺฐานํ ธมฺมํ ปจฺจยา จิตฺตสํสฏฺฐสมุฏฺฐาโน ธมฺโม อุปฺปชฺชติ อารมฺมณปจฺจยา.\csromanspacer{} ตีณิ. (ปฏิจฺจสทิสา.)}

\prose{โนจิตฺตสํสฏฺฐสมุฏฺฐานํ ธมฺมํ ปจฺจยา โนจิตฺตสํสฏฺฐสมุฏฺฐาโน ธมฺโม อุปฺปชฺชติ อารมฺมณปจฺจยา.\csromanspacer{} จกฺขายตนํ ปจฺจยา จกฺขุวิญฺญาณํ ฯเปฯ.\csromanspacer{} กายายตนํ ฯเปฯ.\csromanspacer{} วตฺถุํ ปจฺจยา จิตฺตํ.\csromanspacer{} ปฏิสนฺธิกฺขเณ ฯเปฯ. (1)}

\prose{\csromanfolio{482}โนจิตฺตสํสฏฺฐสมุฏฺฐานํ ธมฺมํ ปจฺจยา จิตฺตสํสฏฺฐสมุฏฺฐาโน ธมฺโม อุปฺปชฺชติ อารมฺมณปจฺจยา.\csromanspacer{} จกฺขายตนํ ปจฺจยา จกฺขุวิญฺญาณสหคตา ขนฺธา ฯเปฯ.\csromanspacer{} กายายตนํ ฯเปฯ.\csromanspacer{} จิตฺตํ ปจฺจยา สมฺปยุตฺตกา ขนฺธา.\csromanspacer{} วตฺถุํ ปจฺจยา จิตฺตสํสฏฺฐสมุฏฺฐานา ขนฺธา. (ปฏิสนฺธิกฺขเณ เทฺวปิ กาตพฺพา.) (2)}

\prose{โนจิตฺตสํสฏฺฐสมุฏฺฐานํ ธมฺมํ ปจฺจยา จิตฺตสํสฏฺฐสมุฏฺฐาโน จ โนจิตฺตสํสฏฺฐสมุฏฺฐาโน จ ธมฺมา อุปฺปชฺชนฺติ อารมฺมณปจฺจยา.\csromanspacer{} จกฺขายตนํ ปจฺจยา จกฺขุวิญฺญาณํ สมฺปยุตฺตกา จ ขนฺธา ฯเปฯ.\csromanspacer{} กายายตนํ ฯเปฯ.\csromanspacer{} วตฺถุํ ปจฺจยา จิตฺตํ สมฺปยุตฺตกา จ ขนฺธา. (ปฏิสนฺธิกฺขเณ เอกํ กาตพฺพํ.) (3)}

\proseitem{292}{จิตฺตสํสฏฺฐสมุฏฺฐานญฺจ โนจิตฺตสํสฏฺฐสมุฏฺฐานญฺจ ธมฺมํ ปจฺจยา จิตฺตสํสฏฺฐสมุฏฺฐาโน ธมฺโม อุปฺปชฺชติ อารมฺมณปจฺจยา.\csromanspacer{} จกฺขุวิญฺญาณสหคตํ เอกํ ขนฺธญฺจ จกฺขายตนญฺจ จกฺขุวิญฺญาณญฺจ ปจฺจยา เทฺว ขนฺธา.\csromanspacer{} เทฺว ขนฺเธ ฯเปฯ.\csromanspacer{} กายวิญฺญาณสหคตํ ฯเปฯ.\csromanspacer{} จิตฺตสํสฏฺฐสมุฏฺฐานํ เอกํ ขนฺธญฺจ จิตฺตญฺจ ปจฺจยา เทฺว ขนฺธา.\csromanspacer{} เทฺว ขนฺเธ ฯเปฯ.\csromanspacer{} จิตฺตสํสฏฺฐสมุฏฺฐานํ เอกํ ขนฺธญฺจ วตฺถุญฺจ ปจฺจยา เทฺว ขนฺธา.\csromanspacer{} เทฺว ขนฺเธ ฯเปฯ. (ปฏิสนฺธิกฺขเณ เทฺว กาตพฺพา.) (1)}

\prose{จิตฺตสํสฏฺฐสมุฏฺฐานญฺจ โนจิตฺตสํสฏฺฐสมุฏฺฐานญฺจ ธมฺมํ ปจฺจยา โนจิตฺตสํสฏฺฐสมุฏฺฐาโน ธมฺโม อุปฺปชฺชติ อารมฺมณปจฺจยา.\csromanspacer{} จกฺขุวิญฺญาณสหคเต ขนฺเธ จ จกฺขายตนญฺจ ปจฺจยา จกฺขุวิญฺญาณํ ฯเปฯ.\csromanspacer{} กายวิญฺญาณสหคเต ฯเปฯ.\csromanspacer{} จิตฺตสํสฏฺฐสมุฏฺฐาเน ขนฺเธ จ วตฺถุญฺจ ปจฺจยา จิตฺตํ. (ปฏิสนฺธิกฺขเณ เอกํ กาตพฺพํ.)}

\prose{จิตฺตสํสฏฺฐสมุฏฺฐานญฺจ โนจิตฺตสํสฏฺฐสมุฏฺฐานญฺจ ธมฺมํ ปจฺจยา จิตฺตสํสฏฺฐสมุฏฺฐาโน จ โนจิตฺตสํสฏฺฐสมุฏฺฐาโน จ ธมฺมา อุปฺปชฺชนฺติ อารมฺมณปจฺจยา.\csromanspacer{} จิตฺตสํสฏฺฐสมุฏฺฐานํ เอกํ ขนฺธญฺจ วตฺถุญฺจ ปจฺจยา เทฺว ขนฺธา จิตฺตญฺจ.\csromanspacer{} เทฺว ขนฺเธ ฯเปฯ. (ปฏิสนฺธิกฺขเณ เอกํ กาตพฺพํ.\csromanspacer{} สํขิตฺตํ.) (3)}

\csromanheader{2}{1. ปจฺจยานุโลม 2. สงฺขฺยาวาร}
\vspace{\tipitakaparskip}
\csromancenter{เหตุยา นว, อารมฺมเณ นว, อธิปติยา นว, (สพฺพตฺถ นว.) อวิคเต นว.}
\csromancenter{จิตฺตสํสฏฺฐสมุฏฺฐานทุก \textbf{2.}\csromanspacer{}\textbf{ ปจฺจยปจฺจนีย 1.}\csromanspacer{}\textbf{ วิภงฺควาร}}
\proseitem{294}{\csromanfolio{483}จิตฺตสํสฏฺฐสมุฏฺฐานํ ธมฺมํ ปจฺจยา จิตฺตสํสฏฺฐสมุฏฺฐาโน ธมฺโม อุปฺปชฺชติ นเหตุปจฺจยา. (เอวํ นว ปญฺหา กาตพฺพา.\csromanspacer{} ปจฺจยวาเร ปญฺจวิญฺญาณมฺปิ กาตพฺพํ.\csromanspacer{} ตีณิเยว, โมโห.)}

\csromanheader{2}{2. ปจฺจยปจฺจนีย 2. สงฺขฺยาวาร}
\csromanheader{2}{\textbf{สุทฺธ}}
\proseitem{295}{นเหตุยา นว, น-อารมฺมเณ ตีณิ, น-อธิปติยา นว, น-อนนฺตเร ตีณิ, นสมนนฺตเร ตีณิ, น-อญฺญมญฺเญ ตีณิ, น-อุปนิสฺสเย ตีณิ, นปุเรชาเต นว, นปจฺฉาชาเต นว, น-อาเสวเน นว, นกมฺเม จตฺตาริ, นวิปาเก นว, น-อาหาเร เอกํ, น-อินฺทฺริเย เอกํ, นฌาเน นว, นมคฺเค นว, นสมฺปยุตฺเต ตีณิ, นวิปฺปยุตฺเต ฉ, โนนตฺถิยา ตีณิ, โนวิคเต ตีณิ.}

\vspace{\tipitakaparskip}
\csromancenter{(เอวํ อิตเร เทฺว คณนาปิ นิสฺสยวาโรปิ กาตพฺโพ.)}
\csromanheader{2}{63. จิตฺตสํสฏฺฐสมุฏฺฐานทุก 5. สํสฏฺฐวาร}
\csromanheader{2}{1. ปจฺจยานุโลม 1. วิภงฺควาร}
\csromanheader{2}{\textbf{เหตุ}}
\proseitem{296}{จิตฺตสํสฏฺฐสมุฏฺฐานํ ธมฺมํ สํสฏฺโฐ จิตฺตสํสฏฺฐสมุฏฺฐาโน ธมฺโม อุปฺปชฺชติ เหตุปจฺจยา.\csromanspacer{} จิตฺตสํสฏฺฐสมุฏฺฐานํ เอกํ ขนฺธํ สํสฏฺฐา เทฺว ขนฺธา.\csromanspacer{} เทฺว ขนฺเธ ฯเปฯ.\csromanspacer{} ปฏิสนฺธิกฺขเณ ฯเปฯ. (1)}

\prose{จิตฺตสํสฏฺฐสมุฏฺฐานํ ธมฺมํ สํสฏฺโฐ โนจิตฺตสํสฏฺฐสมุฏฺฐาโน ธมฺโม อุปฺปชฺชติ เหตุปจฺจยา.\csromanspacer{} จิตฺตสํสฏฺฐสมุฏฺฐาเน ขนฺเธ สํสฏฺฐํ จิตฺตํ.\csromanspacer{} ปฏิสนฺธิกฺขเณ ฯเปฯ. (2)}

\prose{จิตฺตสํสฏฺฐสมุฏฺฐานํ ธมฺมํ สํสฏฺโฐ จิตฺตสํสฏฺฐสมุฏฺฐาโน จ โนจิตฺตสํสฏฺฐสมุฏฺฐาโน จ ธมฺมา อุปฺปชฺชนฺติ เหตุปจฺจยา.\csromanspacer{} จิตฺตสํสฏฺฐสมุฏฺฐานํ เอกํ ขนฺธํ สํสฏฺฐา เทฺว ขนฺธา จิตฺตญฺจ.\csromanspacer{} เทฺว ขนฺเธ ฯเปฯ.\csromanspacer{} ปฏิสนฺธิกฺขเณ ฯเปฯ. (3)}

\prose{\csromanfolio{484}โนจิตฺตสํสฏฺฐสมุฏฺฐานํ ธมฺมํ สํสฏฺโฐ จิตฺตสํสฏฺฐสมุฏฺฐาโน ธมฺโม อุปฺปชฺชติ เหตุปจฺจยา.\csromanspacer{} จิตฺตํ สํสฏฺฐา สมฺปยุตฺตกา ขนฺธา.\csromanspacer{} ปฏิสนฺธิกฺขเณ ฯเปฯ. (1)}

\prose{จิตฺตสํสฏฺฐสมุฏฺฐานญฺจ โนจิตฺตสํสฏฺฐสมุฏฺฐานญฺจ ธมฺมํ สํสฏฺโฐ จิตฺตสํสฏฺฐสมุฏฺฐาโน ธมฺโม อุปฺปชฺชติ เหตุปจฺจยา.\csromanspacer{} จิตฺตสํสฏฺฐสมุฏฺฐานํ เอกํ ขนฺธญฺจ จิตฺตญฺจ สํสฏฺฐา เทฺว ขนฺธา.\csromanspacer{} เทฺว ขนฺเธ ฯเปฯ.\csromanspacer{} ปฏิสนฺธิกฺขเณ ฯเปฯ. (1) (สํขิตฺตํ.)}

\csromanheader{2}{1. ปจฺจยานุโลม 2. สงฺขฺยาวาร}
\proseitem{297}{\textbf{เหตุ}ยา ปญฺจ, อารมฺมเณ ปญฺจ, อธิปติยา ปญฺจ, (สพฺพตฺถ ปญฺจ,) อวิคเต ปญฺจ.}

\csromanheader{2}{2. ปจฺจยปจฺจนีย 1. วิภงฺควาร}
\proseitem{298}{จิตฺตสํสฏฺฐสมุฏฺฐานํ ธมฺมํ สํสฏฺโฐ จิตฺตสํสฏฺฐสมุฏฺฐาโน ธมฺโม อุปฺปชฺชติ นเหตุปจฺจยา. (สํขิตฺตํ.\csromanspacer{} ตีณิเยว, โมโห.)}

\csromanheader{2}{2. ปจฺจยปจฺจนีย 1. สงฺขฺยาวาร}
\csromanheader{2}{\textbf{สุทฺธ}}
\proseitem{299}{นเหตุยา ปญฺจ, น-อธิปติยา ปญฺจ, นปุเรชาเต ปญฺจ, นปจฺฉาชาเต ปญฺจ, น-อาเสวเน ปญฺจ, นกมฺเม ตีณิ, นวิปาเก ปญฺจ, นฌาเน ปญฺจ, นมคฺเค ปญฺจ, นวิปฺปยุตฺเต ปญฺจ.}

\vspace{\tipitakaparskip}
\csromancenter{(เอวํ อิตเร เทฺว คณนาปิ สมฺปยุตฺตวาโรปิ กาตพฺโพ.)}
\csromanheader{2}{63. จิตฺตสํสฏฺฐสมุฏฺฐานทุก 7. ปญฺหาวาร}
\csromanheader{2}{1. ปจฺจยานุโลม 1. วิภงฺควาร}
\csromanheader{2}{\textbf{เหตุ}}
\proseitem{300}{จิตฺตสํสฏฺฐสมุฏฺฐาโน ธมฺโม จิตฺตสํสฏฺฐสมุฏฺฐานสฺส ธมฺมสฺส เหตุปจฺจเยน ปจฺจโย.\csromanspacer{} จิตฺตสํสฏฺฐสมุฏฺฐานา เหตู สมฺปยุตฺตกานํ}

\vspace{\tipitakaparskip}
\csromancenter{จิตฺตสํสฏฺฐสมุฏฺฐานทุก ขนฺธานํ เหตุปจฺจเยน ปจฺจโย.\csromanspacer{} ปฏิสนฺธิกฺขเณ ฯเปฯ. (มูลํ กาตพฺพํ.) จิตฺตสํสฏฺฐสมุฏฺฐานา เหตู จิตฺตสฺส จิตฺตสมุฏฺฐานานญฺจ รูปานํ เหตุปจฺจเยน ปจฺจโย.\csromanspacer{} ปฏิสนฺธิกฺขเณ ฯเปฯ. (มูลํ กาตพฺพํ.) จิตฺตสํสฏฺฐสมุฏฺฐานา เหตู สมฺปยุตฺตกานํ ขนฺธานํ จิตฺตสฺส จ จิตฺตสมุฏฺฐานานญฺจ รูปานํ เหตุปจฺจเยน ปจฺจโย.\csromanspacer{} ปฏิสนฺธิกฺขเณ ฯเปฯ. (3)}
\csromanheader{2}{\textbf{อารมฺมณ}}
\proseitem{301}{\csromanfolio{485}จิตฺตสํสฏฺฐสมุฏฺฐาโน ธมฺโม จิตฺตสํสฏฺฐสมุฏฺฐานสฺส ธมฺมสฺส อารมฺมณปจฺจเยน ปจฺจโย.\csromanspacer{} จิตฺตสํสฏฺฐสมุฏฺฐาเน ขนฺเธ อารพฺภ จิตฺตสํสฏฺฐสมุฏฺฐานา ขนฺธา อุปฺปชฺชนฺติ. (มูลํ กาตพฺพํ.) จิตฺตสํสฏฺฐสมุฏฺฐาเน ขนฺเธ อารพฺภ จิตฺตํ อุปฺปชฺชติ. (มูลํ กาตพฺพํ.) จิตฺตสํสฏฺฐสมุฏฺฐาเน ขนฺเธ อารพฺภ จิตฺตญฺจ สมฺปยุตฺตกา จ ขนฺธา อุปฺปชฺชนฺติ. (3)}

\prose{โนจิตฺตสํสฏฺฐาสมุฏฺฐาโน ธมฺโม โนจิตฺตสํสฏฺฐสมุฏฺฐานสฺส ธมฺมสฺส อารมฺมณปจฺจเยน ปจฺจโย.\csromanspacer{} อริยา มคฺคา วุฏฺฐหิตฺวา มคฺคํ ปจฺจเวกฺขนฺติ ฯเปฯ.\csromanspacer{} นิพฺพานํ ปจฺจเวกฺขนฺติ.\csromanspacer{} นิพฺพานํ โคตฺรภุสฺส ฯเปฯ. (สํขิตฺตํ.\csromanspacer{} ยถา จิตฺตสหภูทุเก อารมฺมณํ, เอวํ กาตพฺพํ.\csromanspacer{} นินฺนานากรณํ.\csromanspacer{} นวปิ ปญฺหา.)}

\csromanheader{2}{\textbf{อธิปติ}}
\proseitem{302}{จิตฺตสํสฏฺฐสมุฏฺฐาโน ธมฺโม จิตฺตสํสฏฺฐสมุฏฺฐานสฺส ธมฺมสฺส อธิปติปจฺจเยน ปจฺจโย.\csromanspacer{} อารมฺมณาธิปติ, สหชาตาธิปติ.\csromanspacer{} ตีณิ. (เทฺวปิ อธิปตี กาตพฺพา.) (3)}

\prose{โนจิตฺตสํสฏฺฐสมุฏฺฐาโน ธมฺโม โนจิตฺตสํสฏฺฐสมุฏฺฐานสฺส ธมฺมสฺส อธิปติปจฺจเยน ปจฺจโย.\csromanspacer{} อารมฺมณาธิปติ, สหชาตาธิปติ.\csromanspacer{} ตีณิ. (เทฺวปิ อธิปตี กาตพฺพา.) (3)}

\prose{จิตฺตสํสฏฺฐสมุฏฺฐาโน จ โนจิตฺตสํสฏฺฐสมุฏฺฐาโน จ ธมฺมา จิตฺตสํสฏฺฐสมุฏฺฐานสฺส ธมฺมสฺส อธิปติปจฺจเยน ปจฺจโย.\csromanspacer{} อารมฺมณาธิปติ. (เอกาเยว อธิปติ กาตพฺพา, นวปิ ปญฺหา.\csromanspacer{} ยถา จิตฺตสหภูทุกํ, เอวํ กาตพฺพํ.\csromanspacer{} นินฺนานากรณํ.)}

\csromanheader{2}{\textbf{อนนฺตราทิ}}
\proseitem{303}{\csromanfolio{486}จิตฺตสํสฏฺฐสมุฏฺฐาโน ธมฺโม จิตฺตสํสฏฺฐสมุฏฺฐานสฺส ธมฺมสฺส อนนฺตรปจฺจเยน ปจฺจโย. (นวปิ ปญฺหา.\csromanspacer{} จิตฺตสหภูทุกสทิสา.) สมนนฺตรปจฺจเยน ปจฺจโย.\csromanspacer{} นว.\csromanspacer{} สหชาตปจฺจเยน ปจฺจโย.\csromanspacer{} นว. (ปฏิจฺจสทิสา.) อญฺญมญฺญปจฺจเยน ปจฺจโย.\csromanspacer{} นว. (ปฏิจฺจสทิสา.) นิสฺสยปจฺจเยน ปจฺจโย.\csromanspacer{} นว. (ปจฺจยสทิสา.) อุปนิสฺสยปจฺจเยน ปจฺจโย. (นวปิ ปญฺหา.\csromanspacer{} จิตฺตสหภูทุกสทิสา.\csromanspacer{} นินฺนานากรณํ.)}

\csromanheader{2}{\textbf{ปุเรชาตาทิ}}
\proseitem{304}{โนจิตฺตสํสฏฺฐสมุฏฺฐาโน ธมฺโม โนจิตฺตสํสฏฺฐสมุฏฺฐานสฺส ธมฺมสฺส ปุเรชาตปจฺจเยน ปจฺจโย.\csromanspacer{} อารมฺมณปุเรชาตํ, วตฺถุปุเรชาตํ.\csromanspacer{} ตีณิ. (จิตฺตสหภูทุกสทิสา.\csromanspacer{} นินฺนานากรณํ.)}

\prose{จิตฺตสํสฏฺฐสมุฏฺฐาโน ธมฺโม โนจิตฺตสํสฏฺฐสมุฏฺฐานสฺส ธมฺมสฺส ปจฺฉาชาตปจฺจเยน ปจฺจโย. (จิตฺตสหภูทุกสทิสา.\csromanspacer{} นินฺนานากรณํ.\csromanspacer{} ตีณิปิ ปจฺฉาชาตา.\csromanspacer{} เทฺว.\csromanspacer{} เอกมูลานํ เอกา ฆฏนา.) อาเสวนปจฺจเยน ปจฺจโย.\csromanspacer{} นว.}

\csromanheader{2}{\textbf{กมฺมาทิ}}
\proseitem{305}{จิตฺตสํสฏฺฐสมุฏฺฐาโน ธมฺโม จิตฺตสํสฏฺฐสมุฏฺฐานสฺส ธมฺมสฺส กมฺมปจฺจเยน ปจฺจโย.\csromanspacer{} ตีณิ. (จิตฺตสหภูทุกสทิสา.\csromanspacer{} นินฺนานากรณา.\csromanspacer{} ตีณิปิ สหชาตา, นานากฺขณิกา.) วิปากปจฺจเยน ปจฺจโย.\csromanspacer{} นว.\csromanspacer{} อาหารปจฺจเยน ปจฺจโย.\csromanspacer{} นว.\csromanspacer{} จิตฺตสหภูคมนสทิสา.\csromanspacer{} เอกํเยว กพฬีการํ อาหารํ.) อินฺทฺริยปจฺจเยน ปจฺจโย.\csromanspacer{} นว.\csromanspacer{} ฌานปจฺจเยน ปจฺจโย.\csromanspacer{} ตีณิ.\csromanspacer{} มคฺคปจฺจเยน ปจฺจโย.\csromanspacer{} ตีณิ.\csromanspacer{} สมฺปยุตฺตปจฺจเยน ปจฺจโย.\csromanspacer{} ปญฺจ.}

\csromanheader{2}{\textbf{วิปฺปยุตฺต}}
\proseitem{306}{จิตฺตสํสฏฺฐสมุฏฺฐาโน ธมฺโม โนจิตฺตสํสฏฺฐสมุฏฺฐานสฺส ธมฺมสฺส วิปฺปยุตฺตปจฺจเยน ปจฺจโย.\csromanspacer{} สหชาตํ, ปจฺฉาชาตํ. (สํขิตฺตํ.) (1)}

\csromanheader{2}{63. จิตฺตสํสฏฺฐสมุฏฺฐานทุก}
\prose{\csromanfolio{487}โนจิตฺตสํสฏฺฐสมุฏฺฐาโน ธมฺโม โนจิตฺตสํสฏฺฐสมุฏฺฐานสฺส ธมฺมสฺส วิปฺปยุตฺตปจฺจเยน ปจฺจโย.\csromanspacer{} \textbf{สหชาตํ}, ปุเรชาตํ, ปจฺฉาชาตํ. (สํขิตฺตํ.) (1)}

\prose{โนจิตฺตสํสฏฺฐสมุฏฺฐาโน ธมฺโม จิตฺตสํสฏฺฐสมุฏฺฐานสฺส ธมฺมสฺส วิปฺปยุตฺตปจฺจเยน ปจฺจโย.\csromanspacer{} \textbf{สหชาตํ}, ปุเรชาตํ.\csromanspacer{}\csromanspacer{}\textbf{สหชาตํ} ปฏิสนฺธิกฺขเณ วตฺถุ จิตฺตสํสฏฺฐสมุฏฺฐานานํ ขนฺธานํ วิปฺปยุตฺตปจฺจเยน ปจฺจโย.\csromanspacer{} \textbf{ปุเรชาตํ} จกฺขายตนํ จกฺขุวิญฺญาณสหคตานํ ขนฺธานํ วิปฺปยุตฺตปจฺจเยน ปจฺจโย ฯเปฯ.\csromanspacer{} กายายตนํ ฯเปฯ.\csromanspacer{} วตฺถุ จิตฺตสํสฏฺฐสมุฏฺฐานานํ ขนฺธานํ วิปฺปยุตฺตปจฺจเยน ปจฺจโย. (2)}

\prose{โนจิตฺตสํสฏฺฐสมุฏฺฐาโน ธมฺโม จิตฺตสํสฏฺฐสมุฏฺฐานสฺส จ โนจิตฺตสํสฏฺฐสมุฏฺฐานสฺส จ ธมฺมสฺส วิปฺปยุตฺตปจฺจเยน ปจฺจโย.\csromanspacer{} \textbf{สหชาตํ}, ปุเรชาตํ.\csromanspacer{}\csromanspacer{}\textbf{สหชาตํ} ปฏิสนฺธิกฺขเณ วตฺถุ จิตฺตสฺส สมฺปยุตฺตกานญฺจ ขนฺธานํ วิปฺปยุตฺตปจฺจเยน ปจฺจโย.\csromanspacer{} \textbf{ปุเรชาตํ} จกฺขายตนํ จกฺขุวิญฺญาณสฺส สมฺปยุตฺตกานญฺจ ขนฺธานํ วิปฺปยุตฺตปจฺจเยน ปจฺจโย ฯเปฯ.\csromanspacer{} กายายตนํ ฯเปฯ.\csromanspacer{} วตฺถุ จิตฺตสฺส สมฺปยุตฺตกานญฺจ ขนฺธานํ วิปฺปยุตฺตปจฺจเยน ปจฺจโย. (3)}

\prose{จิตฺตสํสฏฺฐสมุฏฺฐาโน จ โนจิตฺตสํสฏฺฐสมุฏฺฐาโน จ ธมฺมา โนจิตฺตสํสฏฺฐสมุฏฺฐานสฺส ธมฺมสฺส วิปฺปยุตฺตปจฺจเยน ปจฺจโย.\csromanspacer{} \textbf{สหชาตํ}, ปจฺฉาชาตํ. (สํขิตฺตํ.)}

\csromanheader{2}{\textbf{อตฺถิ}}
\proseitem{307}{จิตฺตสํสฏฺฐสมุฏฺฐาโน ธมฺโม จิตฺตสํสฏฺฐสมุฏฺฐานสฺส ธมฺมสฺส อตฺถิปจฺจเยน ปจฺจโย. (ปฏิจฺจสทิสา.) จิตฺตสํสฏฺฐสมุฏฺฐาโน ธมฺโม โนจิตฺตสํสฏฺฐสมุฏฺฐานสฺส ธมฺมสฺส อตฺถิปจฺจเยน ปจฺจโย.\csromanspacer{} \textbf{สหชาตํ}, ปจฺฉาชาตํ. (สํขิตฺตํ.) จิตฺตสํสฏฺฐสมุฏฺฐาโน ธมฺโม จิตฺตสํสฏฺฐสมุฏฺฐานสฺส จ โนจิตฺตสํสฏฺฐสมุฏฺฐานสฺส จ ธมฺมสฺส อตฺถิปจฺจเยน ปจฺจโย. (ปฏิจฺจสทิสา.) (3)}

\prose{โนจิตฺตสํสฏฺฐสมุฏฺฐาโน ธมฺโม โนจิตฺตสํสฏฺฐสมุฏฺฐานสฺส ธมฺมสฺส อตฺถิปจฺจเยน ปจฺจโย, สหชาตํ, ปุเรชาตํ, ปจฺฉาชาตํ, อาหารํ, อินฺทฺริยํ. (ปุเรชาตสทิสํ, ปุเรชาตํ กาตพฺพํ.\csromanspacer{} สพฺพํ สํขิตฺตํ.\csromanspacer{} วิตฺถาเรตพฺพํ. (1)}

\prose{\csromanfolio{488}โนจิตฺตสํสฏฺฐสมุฏฺฐาโน ธมฺโม จิตฺตสํสฏฺฐสมุฏฺฐานสฺส ธมฺมสฺส อตฺถิปจฺจเยน ปจฺจโย.\csromanspacer{} \textbf{สหชาตํ}, ปุเรชาตํ.\csromanspacer{}\csromanspacer{}\textbf{สหชาตํ} จิตฺตํ สมฺปยุตฺตกานํ ขนฺธานํ อตฺถิปจฺจเยน ปจฺจโย.\csromanspacer{} ปฏิสนฺธิกฺขเณ จิตฺตํ สมฺปยุตฺตกานํ ขนฺธานํ อตฺถิปจฺจเยน ปจฺจโย.\csromanspacer{} ปฏิสนฺธิกฺขเณ วตฺถุ จิตฺตสํสฏฺฐสมุฏฺฐานานํ ขนฺธานํ อตฺถิปจฺจเยน ปจฺจโย. \textbf{(ปุเรชาตํ} ปุเรชาตสทิสํ.\csromanspacer{} นินฺนานากรณํ.\textbf{)} (2\textbf{)}}

\prose{โนจิตฺตสํสฏฺฐสมุฏฺฐาโน ธมฺโม จิตฺตสํสฏฺฐสมุฏฺฐานสฺส จ โนจิตฺตสํสฏฺฐสมุฏฺฐานสฺส จ ธมฺมสฺส อตฺถิปจฺจเยน ปจฺจโย.\csromanspacer{} \textbf{สหชาตํ}, ปุเรชาตํ.\csromanspacer{}\csromanspacer{}\textbf{สหชาตํ} จิตฺตํ สมฺปยุตฺตกานํ ขนฺธานํ จิตฺตสมุฏฺฐานานญฺจ รูปานํ อตฺถิปจฺจเยน ปจฺจโย.\csromanspacer{} ปฏิสนฺธิกฺขเณ จิตฺตํ สมฺปยุตฺตกานํ ขนฺธานํ กฏตฺตา จ รูปานํ อตฺถิปจฺจเยน ปจฺจโย.\csromanspacer{} ปฏิสนฺธิกฺขเณ วตฺถุ จิตฺตสฺส สมฺปยุตฺตกานญฺจ ขนฺธานํ อตฺถิปจฺจเยน ปจฺจโย. \textbf{(ปุเรชาตํ} ปุเรชาตสทิสํ.\textbf{)} (3\textbf{)}}

\proseitem{308}{จิตฺตสํสฏฺฐสมุฏฺฐาโน จ โนจิตฺตสํสฏฺฐสมุฏฺฐาโน จ ธมฺมา จิตฺตสํสฏฺฐสมุฏฺฐานสฺส ธมฺมสฺส อตฺถิปจฺจเยน ปจฺจโย.\csromanspacer{} \textbf{สหชาตํ}, ปุเรชาตํ.\csromanspacer{}\csromanspacer{}สหชาโต จกฺขุวิญฺญาณสหคโต เอโก ขนฺโธ จ จกฺขุวิญฺญาณญฺจ ทฺวินฺนํ ขนฺธานํ ฯเปฯ เทฺว ขนฺธา จ ฯเปฯ.\csromanspacer{} จกฺขุวิญฺญาณสหคโต เอโก ขนฺโธ จ จกฺขายตนญฺจ ทฺวินฺนํ ขนฺธานํ ฯเปฯ เทฺว ขนฺธา จ ฯเปฯ.\csromanspacer{} กายวิญฺญาณสหคโต ฯเปฯ.\csromanspacer{} จิตฺตสํสฏฺฐสมุฏฺฐาโน เอโก ขนฺโธ จ จิตฺตญฺจ ทฺวินฺนํ ขนฺธานํ อตฺถิปจฺจเยน ปจฺจโย.\csromanspacer{} เทฺว ขนฺธา จ ฯเปฯ.\csromanspacer{} จิตฺตสํสฏฺฐสมุฏฺฐาโน เอโก ขนฺโธ จ วตฺถุ จ ทฺวินฺนํ ขนฺธานํ อตฺถิปจฺจเยน ปจฺจโย ฯเปฯ เทฺว ขนฺธา จ ฯเปฯ. (ปฏิสนฺธิกฺขเณ เทฺวปิ กาตพฺพา.\textbf{)} (1\textbf{)}}

\prose{จิตฺตสํสฏฺฐสมุฏฺฐาโน จ โนจิตฺตสํสฏฺฐสมุฏฺฐาโน จ ธมฺมา โนจิตฺตสํสฏฺฐสมุฏฺฐานสฺส ธมฺมสฺส อตฺถิปจฺจเยน ปจฺจโย.\csromanspacer{} \textbf{สหชาตํ}, ปุเรชาตํ, ปจฺฉาชาตํ, อาหารํ, อินฺทฺริยํ.\csromanspacer{}\csromanspacer{}\textbf{สหชาตา} จกฺขุวิญฺญาณสหคตา ขนฺธา จ จกฺขายตนญฺจ จกฺขุวิญฺญาณสฺส อตฺถิปจฺจเยน ปจฺจโย ฯเปฯ.\csromanspacer{} กายวิญฺญาณสหคตา ฯเปฯ.\csromanspacer{} จิตฺตสํสฏฺฐสมุฏฺฐานา ขนฺธา จ จิตฺตญฺจ จิตฺตสํสฏฺฐสมุฏฺฐานานํ รูปานํ อตฺถิปจฺจเยน ปจฺจโย.\csromanspacer{} จิตฺตสํสฏฺฐสมุฏฺฐานา ขนฺธา จ มหาภูตา จ จิตฺตสํสฏฺฐสมุฏฺฐานานํ รูปานํ อตฺถิปจฺจเยน ปจฺจโย.\csromanspacer{} จิตฺตสํสฏฺฐสมุฏฺฐานา ขนฺธา จ วตฺถุ จ จิตฺตสฺส อตฺถิปจฺจเยน ปจฺจโย.}

\vspace{\tipitakaparskip}
\csromancenter{จิตฺตสํสฏฺฐสมุฏฺฐานทุก (ปฏิสนฺธิกฺขเณ ตีณิ กาตพฺพา.) \textbf{ปจฺฉาชาตา} จิตฺตสํสฏฺฐสมุฏฺฐานา ขนฺธา จ จิตฺตญฺจ ปุเรชาตสฺส อิมสฺส โนจิตฺตสํสฏฺฐสมุฏฺฐานสฺส กายสฺส อตฺถิปจฺจเยน ปจฺจโย.\csromanspacer{} \textbf{ปจฺฉาชาตา} จิตฺตสํสฏฺฐสมุฏฺฐานา ขนฺธา จ จิตฺตญฺจ กพฬีกาโร อาหาโร จ อิมสฺส โนจิตฺตสํสฏฺฐสมุฏฺฐานสฺส กายสฺส อตฺถิปจฺจเยน ปจฺจโย.\csromanspacer{} \textbf{ปจฺฉาชาตา} จิตฺตสํสฏฺฐสมุฏฺฐานา ขนฺธา จ จิตฺตญฺจ รูปชีวิตินฺทฺริยญฺจ กฏตฺตารูปานํ อตฺถิปจฺจเยน ปจฺจโย. (2)}
\prose{\csromanfolio{489}จิตฺตสํสฏฺฐสมุฏฺฐาโน จ โนจิตฺตสํสฏฺฐสมุฏฺฐาโน จ ธมฺมา จิตฺตสํสฏฺฐสมุฏฺฐานสฺส จ โนจิตฺตสํสฏฺฐสมุฏฺฐานสฺส จ ธมฺมสฺส อตฺถิปจฺจเยน ปจฺจโย.\csromanspacer{} สหชาตํ, ปุเรชาตํ.\csromanspacer{}\csromanspacer{}\textbf{สหชาโต} จกฺขุวิญฺญาณสหคโต เอโก ขนฺโธ จ จกฺขายตนญฺจ ทฺวินฺนํ ขนฺธานํ จกฺขุวิญฺญาณสฺส จ อตฺถิปจฺจเยน ปจฺจโย.\csromanspacer{} เทฺว ขนฺธา จ ฯเปฯ.\csromanspacer{} กายวิญฺญาณสหคโต ฯเปฯ.\csromanspacer{} \textbf{สหชาโต} จิตฺตสํสฏฺฐสมุฏฺฐาโน เอโก ขนฺโธ จ จิตฺตญฺจ ทฺวินฺนํ ขนฺธานํ จิตฺตสมุฏฺฐานานญฺจ รูปานํ อตฺถิปจฺจเยน ปจฺจโย.\csromanspacer{} เทฺว ขนฺธา จ ฯเปฯ.\csromanspacer{} จิตฺตํสํสฏฺฐสมุฏฺฐาโน เอโก ขนฺโธ จ วตฺถุ จ ทฺวินฺนํ ขนฺธานํ จิตฺตสฺส จ อตฺถิปจฺจเยน ปจฺจโย.\csromanspacer{} เทฺว ขนฺธา จ ฯเปฯ. (ปฏิสนฺธิกฺขเณ เทฺวปิ กาตพฺพา.) (3)}

\csromanheader{2}{1. ปจฺจยานุโลม 2. สงฺขฺยาวาร}
\csromanheader{2}{\textbf{สุทฺธ}}
\proseitem{309}{เหตุยา ตีณิ, อารมฺมเณ นว, อธิปติยา นว, อนนฺตเร นว, สมนนฺตเร นว, สหชาเต นว, อญฺญมญฺเญ นว, นิสฺสเย นว, อุปนิสฺสเย นว, ปุเรชาเต ตีณิ, ปจฺฉาชาเต ตีณิ, อาเสวเน นว, กมฺเม ตีณิ, วิปาเก นว, อาหาเร นว, อินฺทฺริเย นว, ฌาเน ตีณิ, มคฺเค ตีณิ, สมฺปยุตฺเต ปญฺจ, วิปฺปยุตฺเต ปญฺจ, อตฺถิยา นว, นตฺถิยา นว, วิคเต นว, อวิคเต นว.}

\vspace{\tipitakaparskip}
\csromancenter{อนุโลมํ.}
\csromanheader{2}{2. ปจฺจนียุทฺธาร}
\proseitem{310}{จิตฺตสํสฏฺฐสมุฏฺฐาโน ธมฺโม จิตฺตสํสฏฺฐสมุฏฺฐานสฺส ธมฺมสฺส อารมฺมณปจฺจเยน ปจฺจโย.\csromanspacer{} สหชาตปจฺจเยน ปจฺจโย.\csromanspacer{} อุปนิสฺสยปจฺจเยน ปจฺจโย.\csromanspacer{} กมฺมปจฺจเยน ปจฺจโย. (1)}

\prose{\csromanfolio{490}จิตฺตสํสฏฺฐสมุฏฺฐาโน ธมฺโม โนจิตฺตสํสฏฺฐสมุฏฺฐานสฺส ธมฺมสฺส อารมฺมณปจฺจเยน ปจฺจโย.\csromanspacer{} สหชาตปจฺจเยน ปจฺจโย.\csromanspacer{} อุปนิสฺสยปจฺจเยน ปจฺจโย.\csromanspacer{} ปจฺฉาชาตปจฺจเยน ปจฺจโย.\csromanspacer{} กมฺมปจฺจเยน ปจฺจโย. (2)}

\prose{จิตฺตสํสฏฺฐสมุฏฺฐาโน ธมฺโม จิตฺตสํสฏฺฐสมุฏฺฐานสฺส จ โนจิตฺตสํสฏฺฐสมุฏฺฐานสฺส จ ธมฺมสฺส อารมฺมณปจฺจเยน ปจฺจโย.\csromanspacer{} สหชาตปจฺจเยน ปจฺจโย.\csromanspacer{} อุปนิสฺสยปจฺจเยน ปจฺจโย.\csromanspacer{} กมฺมปจฺจเยน ปจฺจโย. (3)}

\proseitem{311}{โนจิตฺตสํสฏฺฐสมุฏฺฐาโน ธมฺโม โนจิตฺตสํสฏฺฐสมุฏฺฐานสฺส ธมฺมสฺส อารมฺมณปจฺจเยน ปจฺจโย.\csromanspacer{} สหชาตปจฺจเยน ปจฺจโย.\csromanspacer{} อุปนิสฺสยปจฺจเยน ปจฺจโย.\csromanspacer{} ปุเรชาตปจฺจเยน ปจฺจโย.\csromanspacer{} ปจฺฉาชาตปจฺจเยน ปจฺจโย.\csromanspacer{} อาหารปจฺจเยน ปจฺจโย.\csromanspacer{} อินฺทฺริยปจฺจเยน ปจฺจโย. (1)}

\prose{โนจิตฺตสํสฏฺฐสมุฏฺฐาโน ธมฺโม จิตฺตสํสฏฺฐสมุฏฺฐานสฺส ธมฺมสฺส อารมฺมณปจฺจเยน ปจฺจโย.\csromanspacer{} สหชาตปจฺจเยน ปจฺจโย.\csromanspacer{} อุปนิสฺสยปจฺจเยน ปจฺจโย.\csromanspacer{} ปุเรชาตปจฺจเยน ปจฺจโย. (2)}

\prose{โนจิตฺตสํสฏฺฐสมุฏฺฐาโน ธมฺโม จิตฺตสํสฏฺฐสมุฏฺฐานสฺส จ โนจิตฺตสํสฏฺฐสมุฏฺฐานสฺส จ ธมฺมสฺส อารมฺมณปจฺจเยน ปจฺจโย.\csromanspacer{} สหชาตปจฺจเยน ปจฺจโย.\csromanspacer{} อุปนิสฺสยปจฺจเยน ปจฺจโย.\csromanspacer{} ปุเรชาตปจฺจเยน ปจฺจโย. (3)}

\proseitem{312}{จิตฺตสํสฏฺฐสมุฏฺฐาโน จ โนจิตฺตสํสฏฺฐสมุฏฺฐาโน จ ธมฺมา จิตฺตสํสฏฺฐสมุฏฺฐานสฺส ธมฺมสฺส อารมฺมณปจฺจเยน ปจฺจโย.\csromanspacer{} สหชาตปจฺจเยน ปจฺจโย.\csromanspacer{} อุปนิสฺสยปจฺจเยน ปจฺจโย. (1)}

\prose{จิตฺตสํสฏฺฐสมุฏฺฐาโน จ โนจิตฺตสํสฏฺฐสมุฏฺฐาโน จ ธมฺมา โนจิตฺตสํสฏฺฐสมุฏฺฐานสฺส ธมฺมสฺส อารมฺมณปจฺจเยน ปจฺจโย.\csromanspacer{} สหชาตปจฺจเยน ปจฺจโย.\csromanspacer{} อุปนิสฺสยปจฺจเยน ปจฺจโย.\csromanspacer{} ปจฺฉาชาตปจฺจเยน ปจฺจโย. (2)}

\prose{จิตฺตสํสฏฺฐสมุฏฺฐาโน จ โนจิตฺตสํสฏฺฐสมุฏฺฐาโน จ ธมฺมา จิตฺตสํสฏฺฐสมุฏฺฐานสฺส จ โนจิตฺตสํสฏฺฐสมุฏฺฐานสฺส จ ธมฺมสฺส อารมฺมณปจฺจเยน ปจฺจโย.\csromanspacer{} สหชาตปจฺจเยน ปจฺจโย.\csromanspacer{} อุปนิสฺสยปจฺจเยน ปจฺจโย. (3)}

\vspace{\tipitakaparskip}
\csromancenter{จิตฺตสํสฏฺฐสมุฏฺฐานานุปริวตฺติทุก \textbf{2.}\csromanspacer{}\textbf{ ปจฺจยปจฺจนีย 2.}\csromanspacer{}\textbf{ สงฺขฺยาวาร}}
\csromancenter{นเหตุยา นว, น-อารมฺมเณ นว, (สพฺพตฺถ นว,) โน-อวิคเต นว.}
\csromanheader{2}{3. ปจฺจยานุโลมปจฺจนีย}
\proseitem{314}{\csromanfolio{491}\textbf{เหตุ}ปจฺจยา น-อารมฺมเณ ตีณิ, น-อธิปติยา นว, น-อนนฺตเร ตีณิ, นสมนนฺตเร ตีณิ, น-อญฺญมญฺเญ เอกํ, น-อุปนิสฺสเย ตีณิ, (สพฺพตฺถ ตีณิ,) นสมฺปยุตฺเต เอกํ, นวิปฺปยุตฺเต ตีณิ, โนนตฺถิยา ตีณิ, โนวิคเต ตีณิ.}

\csromanheader{2}{4. ปจฺจยปจฺจนียานุโลม}
\csromanheader{2}{315. นเหตุปจฺจยา อารมฺมเณ นว, อธิปติยา นว. (สพฺพตฺถ นว, อนุโลมมาติกา.)}
\nitthitam{จิตฺตสํสฏฺฐสมุฏฺฐานทุกํ นิฏฺฐิตํ.}
\csromanmatikaanchor{mtk.63}
\csromantocmarkref{subsection}{64. จิตฺตสํสฏฺฐสมุฏฺฐานสหภูทุก}{mtk.63}
\bookmark[dest={mtk.63},level=2]{64. จิตฺตสํสฏฺฐสมุฏฺฐานสหภูทุก}
\csromanheader{2}{64. จิตฺตสํสฏฺฐสมุฏฺฐานสหภูทุก 1. ปฏิจฺจวาร}
\csromanheader{2}{\textbf{เหตุ}}
\proseitem{316}{จิตฺตสํสฏฺฐสมุฏฺฐานสหภุํ ธมฺมํ ปฏิจฺจ จิตฺตสํสฏฺฐสมุฏฺฐานสหภู ธมฺโม อุปฺปชฺชติ เหตุปจฺจยา.\csromanspacer{} จิตฺตสํสฏฺฐสมุฏฺฐานสหภุํ เอกํ ขนฺธํ ปฏิจฺจ เทฺว ขนฺธา.\csromanspacer{} เทฺว ขนฺเธ ฯเปฯ.\csromanspacer{} ปฏิสนฺธิกฺขเณ ฯเปฯ. (ยถา จิตฺตสํสฏฺฐสมุฏฺฐานทุกํ, เอวํ อิมมฺปิ ทุกํ, นินฺนานากรณํ.)}

\nitthitam{จิตฺตสํสฏฺฐสมุฏฺฐานสหภูทุกํ นิฏฺฐิตํ.}
\csromanmatikaanchor{mtk.64}
\csromantocmarkref{subsection}{65. จิตฺตสํสฏฺฐสมุฏฺฐานานุปริวตฺติทุก}{mtk.64}
\bookmark[dest={mtk.64},level=2]{65. จิตฺตสํสฏฺฐสมุฏฺฐานานุปริวตฺติทุก}
\csromanheader{2}{65. จิตฺตสํสฏฺฐสมุฏฺฐานานุปริวตฺติทุก 1. ปฏิจฺจวาร}
\csromanheader{2}{\textbf{เหตุ}}
\vspace{\tipitakaparskip}
\csromancenter{จิตฺตสํสฏฺฐสมุฏฺฐานานุปริวตฺติํ ธมฺมํ ปฏิจฺจ จิตฺตสํสฏฺฐสมุฏฺฐานานุปริวตฺตี ธมฺโม อุปฺปชฺชติ เหตุปจฺจยา.\csromanspacer{} จิตฺตสํสฏฺฐสมุฏฺฐานานุปริวตฺติํ}
\prosecont{\csromanfolio{492}เอกํ ขนฺธํ ปฏิจฺจ เทฺว ขนฺธา, เทฺว ขนฺเธ ฯเปฯ.\csromanspacer{} ปฏิสนฺธิกฺขเณ ฯเปฯ. (ยถา จิตฺตสํสฏฺฐสมุฏฺฐานทุกสทิสํ, นินฺนานากรณํ.)}

\nitthitam{จิตฺตสํสฏฺฐสมุฏฺฐานานุปริวตฺติทุกํ นิฏฺฐิตํ.}
\csromanmatikaanchor{mtk.65}
\csromantocmarkref{subsection}{66. อชฺฌตฺติกทุก}{mtk.65}
\bookmark[dest={mtk.65},level=2]{66. อชฺฌตฺติกทุก}
\csromanheader{2}{66. อชฺฌตฺติกทุก 1. ปฏิจฺจวาร}
\csromanheader{2}{1. ปจฺจยานุโลม 1. วิภงฺควาร}
\csromanheader{2}{\textbf{เหตุ}}
\proseitem{318}{อชฺฌตฺติกํ ธมฺมํ ปฏิจฺจ อชฺฌตฺติโก ธมฺโม อุปฺปชฺชติ เหตุปจฺจยา.\csromanspacer{} ปฏิสนฺธิกฺขเณ จิตฺตํ ปฏิจฺจ อชฺฌตฺติกํ กฏตฺตารูปํ. (1)}

\prose{อชฺฌตฺติกํ ธมฺมํ ปฏิจฺจ พาหิโร ธมฺโม อุปฺปชฺชติ เหตุปจฺจยา.\csromanspacer{} จิตฺตํ ปฏิจฺจ สมฺปยุตฺตกา ขนฺธา จิตฺตสมุฏฺฐานญฺจ รูปํ.\csromanspacer{} ปฏิสนฺธิกฺขเณ จิตฺตํ ปฏิจฺจ สมฺปยุตฺตกา ขนฺธา พาหิรํ กฏตฺตา จ รูปํ. (2)}

\prose{อชฺฌตฺติกํ ธมฺมํ ปฏิจฺจ อชฺฌตฺติโก จ พาหิโร จ ธมฺมา อุปฺปชฺชนฺติ เหตุปจฺจยา.\csromanspacer{} ปฏิสนฺธิกฺขเณ จิตฺตํ ปฏิจฺจ สมฺปยุตฺตกา ขนฺธา อชฺฌตฺติกญฺจ พาหิรญฺจ กฏตฺตารูปํ. (3)}

\proseitem{319}{พาหิรํ ธมฺมํ ปฏิจฺจ พาหิโร ธมฺโม อุปฺปชฺชติ เหตุปจฺจยา.\csromanspacer{} พาหิรํ เอกํ ขนฺธํ ปฏิจฺจ เทฺว ขนฺธา จิตฺตสมุฏฺฐานญฺจ รูปํ.\csromanspacer{} เทฺว ขนฺเธ ฯเปฯ.\csromanspacer{} ปฏิสนฺธิกฺขเณ พาหิรํ เอกํ ขนฺธํ ปฏิจฺจ เทฺว ขนฺธา พาหิรํ กฏตฺตา จ รูปํ.\csromanspacer{} เทฺว ขนฺเธ ฯเปฯ.\csromanspacer{} ขนฺเธ ปฏิจฺจ วตฺถุ, วตฺถุํ ปฏิจฺจ ขนฺธา.\csromanspacer{} เอกํ มหาภูตํ ฯเปฯ มหาภูเต ปฏิจฺจ จิตฺตสมุฏฺฐานํ รูปํ, กฏตฺตารูปํ, อุปาทารูปํ. (1)}

\prose{พาหิรํ ธมฺมํ ปฏิจฺจ อชฺฌตฺติโก ธมฺโม อุปฺปชฺชติ เหตุปจฺจยา.\csromanspacer{} พาหิเร ขนฺเธ ปฏิจฺจ จิตฺตํ.\csromanspacer{} ปฏิสนฺธิกฺขเณ พาหิเร ขนฺเธ ปฏิจฺจ จิตฺตํ อชฺฌตฺติกํ กฏตฺตา จ รูปํ.\csromanspacer{} ปฏิสนฺธิกฺขเณ พาหิรํ วตฺถุํ ปฏิจฺจ จิตฺตํ. (2)}

\prose{พาหิรํ ธมฺมํ ปฏิจฺจ อชฺฌตฺติโก จ พาหิโร จ ธมฺมา อุปฺปชฺชนฺติ เหตุปจฺจยา.\csromanspacer{} พาหิรํ เอกํ ขนฺธํ ปฏิจฺจ เทฺว ขนฺธา จิตฺตญฺจ จิตฺตสมุฏฺฐานญฺจ รูปํ.\csromanspacer{} เทฺว ขนฺเธ ฯเปฯ.\csromanspacer{} ปฏิสนฺธิกฺขเณ พาหิรํ เอกํ ขนฺธํ ปฏิจฺจ เทฺว ขนฺธา จิตฺตญฺจ อชฺฌตฺติกญฺจ พาหิรญฺจ}

\vspace{\tipitakaparskip}
\csromancenter{อชฺฌตฺติกทุก กฏตฺตารูปํ.\csromanspacer{} เทฺว ขนฺเธ ฯเปฯ.\csromanspacer{} ปฏิสนฺธิกฺขเณ วตฺถุํ ปฏิจฺจ จิตฺตํ สมฺปยุตฺตกา จ ขนฺธา. (3)}
\proseitem{320}{\csromanfolio{493}อชฺฌตฺติกญฺจ พาหิรญฺจ ธมฺมํ ปฏิจฺจ อชฺฌตฺติโก ธมฺโม อุปฺปชฺชติ เหตุปจฺจยา.\csromanspacer{} ปฏิสนฺธิกฺขเณ จิตฺตญฺจ สมฺปยุตฺตเก จ ขนฺเธ ปฏิจฺจ อชฺฌตฺติกํ กฏตฺตารูปํ. (1)}

\prose{อชฺฌตฺติกญฺจ พาหิรญฺจ ธมฺมํ ปฏิจฺจ พาหิโร ธมฺโม อุปฺปชฺชติ เหตุปจฺจยา.\csromanspacer{} พาหิรํ เอกํ ขนฺธญฺจ จิตฺตญฺจ ปฏิจฺจ เทฺว ขนฺธา จิตฺตสมุฏฺฐานญฺจ รูปํ.\csromanspacer{} เทฺว ขนฺเธ จ ฯเปฯ.\csromanspacer{} จิตฺตญฺจ มหาภูเต จ ปฏิจฺจ จิตฺตสมุฏฺฐานํ รูปํ.\csromanspacer{} ปฏิสนฺธิกฺขเณ พาหิรํ เอกํ ขนฺธญฺจ จิตฺตญฺจ ปฏิจฺจ เทฺว ขนฺธา พาหิรํ กฏตฺตา จ รูปํ.\csromanspacer{} เทฺว ขนฺเธ จ ฯเปฯ.\csromanspacer{} จิตฺตญฺจ มหาภูเต จ ปฏิจฺจ พาหิรํ กฏตฺตารูปํ.\csromanspacer{} ปฏิสนฺธิกฺขเณ จิตฺตญฺจ วตฺถุญฺจ ปฏิจฺจ พาหิรา ขนฺธา. (2)}

\prose{อชฺฌตฺติกญฺจ พาหิรญฺจ ธมฺมํ ปฏิจฺจ อชฺฌตฺติโก จ พาหิโร จ ธมฺมา อุปฺปชฺชนฺติ เหตุปจฺจยา.\csromanspacer{} ปฏิสนฺธิกฺขเณ พาหิรํ เอกํ ขนฺธญฺจ จิตฺตญฺจ ปฏิจฺจ เทฺว ขนฺธา อชฺฌตฺติกญฺจ พาหิรญฺจ กฏตฺตารูปํ.\csromanspacer{} เทฺว ขนฺเธ จ ฯเปฯ. (3)}

\csromanheader{2}{\textbf{อารมฺมณ}}
\proseitem{321}{อชฺฌตฺติกํ ธมฺมํ ปฏิจฺจ พาหิโร ธมฺโม อุปฺปชฺชติ อารมฺมณปจฺจยา.\csromanspacer{} จิตฺตํ ปฏิจฺจ สมฺปยุตฺตกา ขนฺธา.\csromanspacer{} ปฏิสนฺธิกฺขเณ จิตฺตํ ปฏิจฺจ สมฺปยุตฺตกา ขนฺธา. (1)}

\prose{พาหิรํ ธมฺมํ ปฏิจฺจ พาหิโร ธมฺโม อุปฺปชฺชติ อารมฺมณปจฺจยา.\csromanspacer{} พาหิรํ เอกํ ขนฺธํ ปฏิจฺจ เทฺว ขนฺธา.\csromanspacer{} เทฺว ขนฺเธ ฯเปฯ.\csromanspacer{} ปฏิสนฺธิกฺขเณ ฯเปฯ.\csromanspacer{} วตฺถุํ ปฏิจฺจ ขนฺธา. (1)}

\prose{พาหิรํ ธมฺมํ ปฏิจฺจ อชฺฌตฺติโก ธมฺโม อุปฺปชฺชติ อารมฺมณปจฺจยา.\csromanspacer{} พาหิเร ขนฺเธ ปฏิจฺจ จิตฺตํ.\csromanspacer{} ปฏิสนฺธิกฺขเณ พาหิเร ขนฺเธ ปฏิจฺจ จิตฺตํ.\csromanspacer{} ปฏิสนฺธิกฺขเณ วตฺถุํ ปฏิจฺจ จิตฺตํ. (2)}

\prose{พาหิรํ ธมฺมํ ปฏิจฺจ อชฺฌตฺติโก จ พาหิโร จ ธมฺมา อุปฺปชฺชนฺติ อารมฺมณปจฺจยา.\csromanspacer{} พาหิรํ เอกํ ขนฺธํ ปฏิจฺจ เทฺว ขนฺธา จิตฺตญฺจ.\csromanspacer{} เทฺว ขนฺเธ ฯเปฯ.\csromanspacer{} ปฏิสนฺธิกฺขเณ พาหิรํ เอกํ ขนฺธํ ปฏิจฺจ เทฺว ขนฺธา จิตฺตญฺจ.\csromanspacer{} เทฺว ขนฺเธ ฯเปฯ.\csromanspacer{} ปฏิสนฺธิกฺขเณ วตฺถุํ ปฏิจฺจ จิตฺตํ สมฺปยุตฺตกา จ ขนฺธา. (3)}

\prose{\csromanfolio{494}อชฺฌตฺติกญฺจ พาหิรญฺจ ธมฺมํ ปฏิจฺจ พาหิโร ธมฺโม อุปฺปชฺชติ อารมฺมณปจฺจยา.\csromanspacer{} พาหิรํ เอกํ ขนฺธญฺจ จิตฺตญฺจ ปฏิจฺจ เทฺว ขนฺธา.\csromanspacer{} เทฺว ขนฺเธ จ ฯเปฯ.\csromanspacer{} ปฏิสนฺธิกฺขเณ พาหิรํ เอกํ ขนฺธญฺจ จิตฺตญฺจ ปฏิจฺจ เทฺว ขนฺธา.\csromanspacer{} เทฺว ขนฺเธ จ ฯเปฯ.\csromanspacer{} ปฏิสนฺธิกฺขเณ พาหิรํ เอกํ ขนฺธญฺจ วตฺถุญฺจ ปฏิจฺจ เทฺว ขนฺธา.\csromanspacer{} เทฺว ขนฺเธ จ ฯเปฯ. (สํขิตฺตํ.) (1)}

\csromanheader{2}{1. ปจฺจยานุโลม 2. สงฺขฺยาวาร}
\csromanheader{2}{\textbf{สุทฺธ}}
\proseitem{322}{เหตุยา นว, อารมฺมเณ ปญฺจ, อธิปติยา ปญฺจ, อนนฺตเร ปญฺจ, สมนนฺตเร ปญฺจ, สหชาเต นว, อญฺญมญฺเญ ปญฺจ, นิสฺสเย นว, อุปนิสฺสเย ปญฺจ, ปุเรชาเต ปญฺจ, อาเสวเน ปญฺจ, กมฺเม นว, วิปาเก นว, (สพฺพตฺถ นว,) สมฺปยุตฺเต ปญฺจ, วิปฺปยุตฺเต นว, อตฺถิยา นว, นตฺถิยา ปญฺจ, วิคเต ปญฺจ, อวิคเต นว.}

\csromanheader{2}{2. ปจฺจยปจฺจนีย 1. วิภงฺควาร}
\csromanheader{2}{\textbf{นเหตุ}}
\proseitem{323}{อชฺฌตฺติกํ ธมฺมํ ปฏิจฺจ อชฺฌตฺติโก ธมฺโม อุปฺปชฺชติ นเหตุปจฺจยา.\csromanspacer{} อเหตุกปฏิสนฺธิกฺขเณ จิตฺตํ ปฏิจฺจ อชฺฌตฺติกํ กฏตฺตารูปํ. (1)}

\prose{อชฺฌตฺติกํ ธมฺมํ ปฏิจฺจ พาหิโร ธมฺโม อุปฺปชฺชติ นเหตุปจฺจยา.\csromanspacer{} อเหตุกํ จิตฺตํ ปฏิจฺจ สมฺปยุตฺตกา ขนฺธา จิตฺตสมุฏฺฐานญฺจ รูปํ.\csromanspacer{} อเหตุกปฏิสนฺธิกฺขเณ จิตฺตํ ปฏิจฺจ สมฺปยุตฺตกา ขนฺธา พาหิรํ กฏตฺตา จ รูปํ.\csromanspacer{} วิจิกิจฺฉาสหคตํ อุทฺธจฺจสหคตํ จิตฺตํ ปฏิจฺจ วิจิกิจฺฉาสหคโต อุทฺธจฺจสหคโต โมโห. (2)}

\prose{อชฺฌตฺติกํ ธมฺมํ ปฏิจฺจ อชฺฌตฺติโก จ พาหิโร จ ธมฺมา อุปฺปชฺชนฺติ นเหตุปจฺจยา.\csromanspacer{} อเหตุกปฏิสนฺธิกฺขเณ จิตฺตํ ปฏิจฺจ สมฺปยุตฺตกา ขนฺธา อชฺฌตฺติกญฺจ พาหิรญฺจ กฏตฺตารูปํ. (3)}

\proseitem{324}{พาหิรํ ธมฺมํ ปฏิจฺจ พาหิโร ธมฺโม อุปฺปชฺชติ นเหตุปจฺจยา.\csromanspacer{} อเหตุกํ พาหิรํ เอกํ ขนฺธํ ปฏิจฺจ เทฺว ขนฺธา จิตฺตสมุฏฺฐานญฺจ รูปํ.\csromanspacer{} เทฺว ขนฺเธ ฯเปฯ.}

\vspace{\tipitakaparskip}
\csromancenter{อชฺฌตฺติกทุก อเหตุกปฏิสนฺธิกฺขเณ ฯเปฯ. (ยาว อสญฺญสตฺตา.) วิจิกิจฺฉาสหคเต อุทฺธจฺจสหคเต ขนฺเธ ปฏิจฺจ วิจิกิจฺฉาสหคโต อุทฺธจฺจสหคโต โมโห. (1)}
\prose{\csromanfolio{495}พาหิรํ ธมฺมํ ปฏิจฺจ อชฺฌตฺติโก ธมฺโม อุปฺปชฺชติ นเหตุปจฺจยา.\csromanspacer{} อเหตุเก พาหิเร ขนฺเธ ปฏิจฺจ จิตฺตํ.\csromanspacer{} อเหตุกปฏิสนฺธิกฺขเณ พาหิเร ขนฺเธ ปฏิจฺจ จิตฺตํ อชฺฌตฺติกํ กฏตฺตา จ รูปํ.\csromanspacer{} อเหตุกปฏิสนฺธิกฺขเณ วตฺถุํ ปฏิจฺจ จิตฺตํ. (2)}

\prose{พาหิรํ ธมฺมํ ปฏิจฺจ อชฺฌตฺติโก จ พาหิโร จ ธมฺมา อุปฺปชฺชนฺติ นเหตุปจฺจยา.\csromanspacer{} อเหตุกํ พาหิรํ เอกํ ขนฺธํ ปฏิจฺจ เทฺว ขนฺธา จิตฺตญฺจ จิตฺตสมุฏฺฐานญฺจ รูปํ.\csromanspacer{} เทฺว ขนฺเธ ฯเปฯ.\csromanspacer{} อเหตุกปฏิสนฺธิกฺขเณ พาหิรํ เอกํ ขนฺธํ ปฏิจฺจ เทฺว ขนฺธา จิตฺตญฺจ อชฺฌตฺติกญฺจ พาหิรญฺจ กฏตฺตารูปํ.\csromanspacer{} อเหตุกปฏิสนฺธิกฺขเณ วตฺถุํ ปฏิจฺจ จิตฺตํ สมฺปยุตฺตกา จ ขนฺธา. (3)}

\proseitem{325}{อชฺฌตฺติกญฺจ พาหิรญฺจ ธมฺมํ ปฏิจฺจ อชฺฌตฺติโก ธมฺโม อุปฺปชฺชติ นเหตุปจฺจยา.\csromanspacer{} อเหตุกปฏิสนฺธิกฺขเณ จิตฺตญฺจ สมฺปยุตฺตเก จ ขนฺเธ ปฏิจฺจ อชฺฌตฺติกํ กฏตฺตารูปํ. (1)}

\prose{อชฺฌตฺติกญฺจ พาหิรญฺจ ธมฺมํ ปฏิจฺจ พาหิโร ธมฺโม อุปฺปชฺชติ นเหตุปจฺจยา.\csromanspacer{} อเหตุกํ พาหิรํ เอกํ ขนฺธญฺจ จิตฺตญฺจ ปฏิจฺจ เทฺว ขนฺธา จิตฺตสมุฏฺฐานญฺจ รูปํ.\csromanspacer{} เทฺว ขนฺเธ ฯเปฯ.\csromanspacer{} อเหตุกํ จิตฺตญฺจ มหาภูเต จ ปฏิจฺจ จิตฺตสมุฏฺฐานํ รูปํ.\csromanspacer{} อเหตุกปฏิสนฺธิกฺขเณ พาหิรํ เอกํ ขนฺธญฺจ จิตฺตญฺจ ปฏิจฺจ เทฺว ขนฺธา พาหิรํ กฏตฺตา จ รูปํ.\csromanspacer{} อเหตุกปฏิสนฺธิกฺขเณ จิตฺตญฺจ มหาภูเต จ ปฏิจฺจ พาหิรํ กฏตฺตารูปํ.\csromanspacer{} อเหตุกปฏิสนฺธิกฺขเณ จิตฺตญฺจ วตฺถุญฺจ ปฏิจฺจ พาหิรา ขนฺธา.\csromanspacer{} วิจิกิจฺฉาสหคเต อุทฺธจฺจสหคเต ขนฺเธ จ จิตฺตญฺจ ปฏิจฺจ วิจิกิจฺฉาสหคโต อุทฺธจฺจสหคโต โมโห. (2)}

\prose{อชฺฌตฺติกญฺจ พาหิรญฺจ ธมฺมํ ปฏิจฺจ อชฺฌตฺติโก จ พาหิโร จ ธมฺมา อุปฺปชฺชนฺติ นเหตุปจฺจยา.\csromanspacer{} อเหตุกปฏิสนฺธิกฺขเณ พาหิรํ เอกํ ขนฺธญฺจ จิตฺตญฺจ ปฏิจฺจ เทฺว ขนฺธา อชฺฌตฺติกญฺจ พาหิรญฺจ กฏตฺตารูปํ.\csromanspacer{} เทฺว ขนฺเธ จ ฯเปฯ. (3)}

\csromanheader{2}{\textbf{น-อารมฺมณ}}
\proseitem{326}{\csromanfolio{496}อชฺฌตฺติกํ ธมฺมํ ปฏิจฺจ อชฺฌตฺติโก ธมฺโม อุปฺปชฺชติ น-อารมฺมณปจฺจยา.\csromanspacer{} ปฏิสนฺธิกฺขเณ จิตฺตํ ปฏิจฺจ อชฺฌตฺติกํ กฏตฺตารูปํ. (1)}

\prose{อชฺฌตฺติกํ ธมฺมํ ปฏิจฺจ พาหิโร ธมฺโม อุปฺปชฺชติ น-อารมฺมณปจฺจยา.\csromanspacer{} จิตฺตํ ปฏิจฺจ จิตฺตสมุฏฺฐานํ รูปํ.\csromanspacer{} ปฏิสนฺธิกฺขเณ จิตฺตํ ปฏิจฺจ พาหิรํ กฏตฺตารูปํ. (2)}

\prose{อชฺฌตฺติกํ ธมฺมํ ปฏิจฺจ อชฺฌตฺติโก จ พาหิโร จ ธมฺมา อุปฺปชฺชนฺติ น-อารมฺมณปจฺจยา.\csromanspacer{} ปฏิสนฺธิกฺขเณ จิตฺตํ ปฏิจฺจ อชฺฌตฺติกญฺจ พาหิรญฺจ กฏตฺตารูปํ. (3)}

\prose{พาหิรํ ธมฺมํ ปฏิจฺจ พาหิโร ธมฺโม อุปฺปชฺชติ น-อารมฺมณปจฺจยา.\csromanspacer{} พาหิเร ขนฺเธ ปฏิจฺจ จิตฺตสมุฏฺฐานํ รูปํ.\csromanspacer{} ปฏิสนฺธิกฺขเณ พาหิเร ขนฺเธ ปฏิจฺจ พาหิรํ กฏตฺตารูปํ.\csromanspacer{} พาหิเร ขนฺเธ ปฏิจฺจ วตฺถุ.\csromanspacer{} เอกํ มหาภูตํ ฯเปฯ. (ยาว อสญฺญสตฺตา.) (1)}

\prose{พาหิรํ ธมฺมํ ปฏิจฺจ อชฺฌตฺติโก ธมฺโม อุปฺปชฺชติ น-อารมฺมณปจฺจยา.\csromanspacer{} ปฏิสนฺธิกฺขเณ พาหิเร ขนฺเธ ปฏิจฺจ อชฺฌตฺติกํ กฏตฺตารูปํ. (2)}

\prose{พาหิรํ ธมฺมํ ปฏิจฺจ อชฺฌตฺติโก จ พาหิโร จ ธมฺมา อุปฺปชฺชนฺติ น-อารมฺมณปจฺจยา.\csromanspacer{} ปฏิสนฺธิกฺขเณ พาหิเร ขนฺเธ ปฏิจฺจ อชฺฌตฺติกญฺจ พาหิรญฺจ กฏตฺตารูปํ. (3)}

\proseitem{327}{อชฺฌตฺติกญฺจ พาหิรญฺจ ธมฺมํ ปฏิจฺจ อชฺฌตฺติโก ธมฺโม อุปฺปชฺชติ น-อารมฺมณปจฺจยา.\csromanspacer{} ปฏิสนฺธิกฺขเณ จิตฺตญฺจ สมฺปยุตฺตเก จ ขนฺเธ ปฏิจฺจ อชฺฌตฺติกํ กฏตฺตารูปํ. (1)}

\prose{อชฺฌตฺติกญฺจ พาหิรญฺจ ธมฺมํ ปฏิจฺจ พาหิโร ธมฺโม อุปฺปชฺชติ น-อารมฺมณปจฺจยา.\csromanspacer{} พาหิเร ขนฺเธ จ จิตฺตญฺจ ปฏิจฺจ จิตฺตสมุฏฺฐานํ รูปํ.\csromanspacer{} จิตฺตญฺจ มหาภูเต จ ปฏิจฺจ จิตฺตสมุฏฺฐานํ รูปํ. (ปฏิสนฺธิกฺขเณ เทฺว กาตพฺพา.)}

\prose{อชฺฌตฺติกญฺจ พาหิรญฺจ ธมฺมํ ปฏิจฺจ อชฺฌตฺติโก จ พาหิโร จ ธมฺมา อุปฺปชฺชนฺติ น-อารมฺมณปจฺจยา.\csromanspacer{} ปฏิสนฺธิกฺขเณ จิตฺตญฺจ สมฺปยุตฺตเก จ ขนฺเธ ปฏิจฺจ อชฺฌตฺติกญฺจ พาหิรญฺจ กฏตฺตารูปํ. (สํขิตฺตํ.) (3)}

\csromanheader{2}{66. อชฺฌตฺติกทุก \textbf{นฌาน}}
\proseitem{328}{\csromanfolio{497}อชฺฌตฺติกํ ธมฺมํ ปฏิจฺจ พาหิโร ธมฺโม อุปฺปชฺชติ นฌานปจฺจยา.\csromanspacer{} จกฺขุวิญฺญาณํ ปฏิจฺจ สมฺปยุตฺตกา ขนฺธา ฯเปฯ.\csromanspacer{} กายวิญฺญาณํ ฯเปฯ.}

\prose{พาหิรํ ธมฺมํ ปฏิจฺจ พาหิโร ธมฺโม อุปฺปชฺชติ นฌานปจฺจยา.\csromanspacer{} จกฺขุวิญฺญาณสหคตํ เอกํ ขนฺธํ ปฏิจฺจ เทฺว ขนฺธา.\csromanspacer{} เทฺว ขนฺเธ ฯเปฯ.\csromanspacer{} กายวิญฺญาณสหคตํ ฯเปฯ.\csromanspacer{} พาหิรํ.\csromanspacer{} อาหารสมุฏฺฐานํ.\csromanspacer{} อุตุสมุฏฺฐานํ.\csromanspacer{} อสญฺญสตฺตานํ ฯเปฯ.\csromanspacer{} พาหิรํ ธมฺมํ ปฏิจฺจ อชฺฌตฺติโก ธมฺโม อุปฺปชฺชติ นฌานปจฺจยา.\csromanspacer{} จกฺขุวิญฺญาณสหคเต ขนฺเธ ปฏิจฺจ จกฺขุวิญฺญาณํ ฯเปฯ.\csromanspacer{} กายวิญฺญาณสหคเต ขนฺเธ ปฏิจฺจ กายวิญฺญาณํ.\csromanspacer{}\csromanspacer{}พาหิรํ ธมฺมํ ปฏิจฺจ อชฺฌตฺติโก จ พาหิโร จ ธมฺมา อุปฺปชฺชนฺติ นฌานปจฺจยา.\csromanspacer{} จกฺขุวิญฺญาณสหคตํ เอกํ ขนฺธํ ปฏิจฺจ เทฺว ขนฺธา จกฺขุวิญฺญาณญฺจ.\csromanspacer{} เทฺว ขนฺเธ ฯเปฯ.\csromanspacer{} กายวิญฺญาณสหคตํ เอกํ ขนฺธํ ฯเปฯ. (3)}

\prose{อชฺฌตฺติกญฺจ พาหิรญฺจ ธมฺมํ ปฏิจฺจ พาหิโร ธมฺโม อุปฺปชฺชติ นฌานปจฺจยา.\csromanspacer{} จกฺขุวิญฺญาณสหคตํ เอกํ ขนฺธญฺจ จกฺขุวิญฺญาณญฺจ ปฏิจฺจ เทฺว ขนฺธา.\csromanspacer{} เทฺว ขนฺเธ ฯเปฯ.\csromanspacer{} กายวิญฺญาณํ. (จกฺกํ.)}

\csromanheader{2}{2. ปจฺจยปจฺจนีย 2. สงฺขฺยาวาร}
\proseitem{329}{นเหตุยา นว, น-อารมฺมเณ นว, น-อธิปติยา นว, น-อนนฺตเร นว, (สพฺพตฺถ นว.) นกมฺเม ตีณิ, นวิปาเก ปญฺจ, น-อาหาเร เอกํ, น-อินฺทฺริเย เอกํ, นฌาเน ปญฺจ, นมคฺเค นว, นสมฺปยุตฺเต นว, นวิปฺปยุตฺเต ปญฺจ, โนนตฺถิยา นว, โนวิคเต นว.}

\csromanheader{2}{3. ปจฺจยานุโลมปจฺจนีย}
\proseitem{330}{เหตุปจฺจยา น-อารมฺมเณ นว, น-อธิปติยา นว. (สํขิตฺตํ.)}

\csromanheader{2}{4. ปจฺจยปจฺจนียานุโลม}
\proseitem{331}{นเหตุปจฺจยา อารมฺมเณ ปญฺจ, อนนฺตเร ปญฺจ, สมนนฺตเร ปญฺจ, สหชาเต นว ฯเปฯ มคฺเค ตีณิ. (สํขิตฺตํ.)}

\vspace{\tipitakaparskip}
\csromancenter{(สหชาตวาโรปิ ปฏิจฺจวารสทิโส.)}
\csromanheader{2}{66. อชฺฌตฺติกทุก 3. ปจฺจยวาร}
\csromanheader{2}{1. ปจฺจยานุโลม 1. วิภงฺควาร}
\csromanheader{2}{\textbf{เหตุ}}
\proseitem{332}{\csromanfolio{498}อชฺฌตฺติกํ ธมฺมํ ปจฺจยา อชฺฌตฺติโก ธมฺโม อุปฺปชฺชติ เหตุปจฺจยา.\csromanspacer{} ตีณิ. (ปฏิจฺจสทิสา.)}

\prose{พาหิรํ ธมฺมํ ปจฺจยา พาหิโร ธมฺโม อุปฺปชฺชติ เหตุปจฺจยา.\csromanspacer{} พาหิรํ เอกํ ขนฺธํ ปจฺจยา เทฺว ขนฺธา จิตฺตสมุฏฺฐานญฺจ รูปํ.\csromanspacer{} เทฺว ขนฺเธ ฯเปฯ. (ปฏิสนฺธิกฺขเณ เทฺวปิ กาตพฺพา, ยาว อชฺฌตฺติกา มหาภูตา.) วตฺถุํ ปจฺจยา พาหิรา ขนฺธา.\csromanspacer{}\csromanspacer{}พาหิรํ ธมฺมํ ปจฺจยา อชฺฌตฺติโก ธมฺโม อุปฺปชฺชติ เหตุปจฺจยา.\csromanspacer{} พาหิเร ขนฺเธ ปจฺจยา จิตฺตํ.\csromanspacer{} วตฺถุํ ปจฺจยา จิตฺตํ. (ปฏิสนฺธิกฺขเณ เทฺวปิ กาตพฺพา.).\csromanspacer{} พาหิรํ ธมฺมํ ปจฺจยา อชฺฌตฺติโก จ พาหิโร จ ธมฺมา อุปฺปชฺชนฺติ เหตุปจฺจยา.\csromanspacer{} พาหิรํ เอกํ ขนฺธํ ปจฺจยา เทฺว ขนฺธา จิตฺตญฺจ จิตฺตสมุฏฺฐานญฺจ รูปํ.\csromanspacer{} เทฺว ขนฺเธ ฯเปฯ.\csromanspacer{} วตฺถุํ ปจฺจยา จิตฺตํ สมฺปยุตฺตกา จ ขนฺธา. (ปฏิสนฺธิกฺขเณ เทฺวปิ กาตพฺพา.) (3)}

\proseitem{333}{อชฺฌตฺติกญฺจ พาหิรญฺจ ธมฺมํ ปจฺจยา อชฺฌตฺติโก ธมฺโม อุปฺปชฺชติ เหตุปจฺจยา.\csromanspacer{} ปฏิสนฺธิกฺขเณ จิตฺตญฺจ สมฺปยุตฺตเก จ ขนฺเธ ปจฺจยา อชฺฌตฺติกํ กฏตฺตารูปํ.\csromanspacer{}\csromanspacer{}อชฺฌตฺติกญฺจ พาหิรญฺจ ธมฺมํ ปจฺจยา พาหิโร ธมฺโม อุปฺปชฺชติ เหตุปจฺจยา.\csromanspacer{} พาหิรํ เอกํ ขนฺธญฺจ จิตฺตญฺจ ปจฺจยา เทฺว ขนฺธา จิตฺตสมุฏฺฐานญฺจ รูปํ.\csromanspacer{} เทฺว ขนฺเธ ฯเปฯ.\csromanspacer{} จิตฺตญฺจ มหาภูเต จ ปจฺจยา จิตฺตสมุฏฺฐานํ รูปํ.\csromanspacer{} จิตฺตญฺจ วตฺถุญฺจ ปจฺจยา พาหิรา ขนฺธา. (ปฏิสนฺธิกฺขเณ ตีณิปิ กาตพฺพา.).\csromanspacer{} อชฺฌตฺติกญฺจ พาหิรญฺจ ธมฺมํ ปจฺจยา อชฺฌตฺติโก จ พาหิโร จ ธมฺมา อุปฺปชฺชนฺติ เหตุปจฺจยา.\csromanspacer{} ปฏิสนฺธิกฺขเณ พาหิรํ เอกํ ขนฺธญฺจ จิตฺตญฺจ ปจฺจยา เทฺว ขนฺธา อชฺฌตฺติกญฺจ พาหิรญฺจ กฏตฺตารูปํ.\csromanspacer{} เทฺว ขนฺเธ ฯเปฯ. (3)}

\csromanheader{2}{\textbf{อรมฺมณ}}
\proseitem{334}{อชฺฌตฺติกํ ธมฺมํ ปจฺจยา อชฺฌตฺติโก ธมฺโม อุปฺปชฺชติ อารมฺมณปจฺจยา.\csromanspacer{} จกฺขายตนํ ปจฺจยา จกฺขุวิญฺญาณํ ฯเปฯ.\csromanspacer{} กายายตนํ ฯเปฯ.\csromanspacer{}\csromanspacer{}อชฺฌตฺติกํ ธมฺมํ ปจฺจยา พาหิโร ธมฺโม อุปฺปชฺชติ อารมฺมณปจฺจโย.}

\vspace{\tipitakaparskip}
\csromancenter{อชฺฌตฺติกทุก จกฺขายตนญฺจ จกฺขุวิญฺญาณญฺจ ปจฺจยา จกฺขุวิญฺญาณสหคตา ขนฺธา ฯเปฯ.\csromanspacer{} กายายตนญฺจ ฯเปฯ.\csromanspacer{} จิตฺตํ ปจฺจยา สมฺปยุตฺตกา ขนฺธา.\csromanspacer{} ปฏิสนฺธิกฺขเณ ฯเปฯ.\csromanspacer{} อชฺฌตฺติกํ ธมฺมํ ปจฺจยา อชฺฌตฺติโก จ พาหิโร จ ธมฺมา อุปฺปชฺชนฺติ อารมฺมณปจฺจยา.\csromanspacer{} จกฺขายตนํ ปจฺจยา จกฺขุวิญฺญาณํ สมฺปยุตฺตกา จ ขนฺธา ฯเปฯ.\csromanspacer{} กายายตนํ. (3)}
\prose{\csromanfolio{499}พาหิรํ ธมฺมํ ปจฺจยา พาหิโร ธมฺโม อุปฺปชฺชติ อารมฺมณปจฺจยา.\csromanspacer{} พาหิรํ เอกํ ขนฺธํ ปจฺจยา เทฺว ขนฺธา.\csromanspacer{} เทฺว ขนฺเธ ฯเปฯ.\csromanspacer{} ปฏิสนฺธิกฺขเณ ฯเปฯ.\csromanspacer{} วตฺถุํ ปจฺจยา พาหิรา ขนฺธา.\csromanspacer{}\csromanspacer{}พาหิรํ ธมฺมํ ปจฺจยา อชฺฌตฺติโก ธมฺโม อุปฺปชฺชติ อารมฺมณปจฺจยา.\csromanspacer{} พาหิเร ขนฺเธ ปจฺจยา จิตฺตํ.\csromanspacer{} วตฺถุํ ปจฺจยา จิตฺตํ. (ปฏิสนฺธิกฺขเณ เทฺวปิ กาตพฺพา.).\csromanspacer{} พาหิรํ ธมฺมํ ปจฺจยา อชฺฌตฺติโก จ พาหิโร จ ธมฺมา อุปฺปชฺชนฺติ อารมฺมณปจฺจยา.\csromanspacer{} วตฺถุํ ปจฺจยา จิตฺตญฺจ สมฺปยุตฺตกา จ ขนฺธา. (ปฏิสนฺธิกฺขเณ เอกํ กาตพฺพํ.) (3)}

\proseitem{335}{อชฺฌตฺติกญฺจ พาหิรญฺจ ธมฺมํ ปจฺจยา อชฺฌตฺติโก ธมฺโม อุปฺปชฺชติ อารมฺมณปจฺจยา.\csromanspacer{} จกฺขุวิญฺญาณสหคเต ขนฺเธ จ จกฺขายตนญฺจ ปจฺจยา จกฺขุวิญฺญาณํ ฯเปฯ.\csromanspacer{} กายวิญฺญาณสหคเต ฯเปฯ.\csromanspacer{} อชฺฌตฺติกญฺจ พาหิรญฺจ ธมฺมํ ปจฺจยา พาหิโร ธมฺโม อุปฺปชฺชติ อารมฺมณปจฺจยา.\csromanspacer{} จกฺขุวิญฺญาณสหคตํ เอกํ ขนฺธญฺจ จกฺขายตนญฺจ จกฺขุวิญฺญาณญฺจ ปจฺจยา เทฺว ขนฺธา.\csromanspacer{} เทฺว ขนฺเธ ฯเปฯ.\csromanspacer{} กายวิญฺญาณสหคตํ ฯเปฯ.\csromanspacer{} พาหิรํ เอกํ ขนฺธญฺจ จิตฺตญฺจ ปจฺจยา เทฺว ขนฺธา.\csromanspacer{} เทฺว ขนฺเธ ฯเปฯ.\csromanspacer{} จิตฺตญฺจ วตฺถุญฺจ ปจฺจยา พาหิรา ขนฺธา. (ปฏิสนฺธิกฺขเณ เทฺวปิ กาตพฺพา.) อชฺฌตฺติกญฺจ พาหิรญฺจ ธมฺมํ ปจฺจยา อชฺฌตฺติโก จ พาหิโร จ ธมฺมา อุปฺปชฺชนฺติ อารมฺมณปจฺจยา.\csromanspacer{} จกฺขุวิญฺญาณสหคตํ เอกํ ขนฺธญฺจ จกฺขายตนญฺจ ปจฺจยา เทฺว ขนฺธา จกฺขุวิญฺญาณญฺจ.\csromanspacer{} เทฺว ขนฺเธ ฯเปฯ. (สํขิตฺตํ.) (3)}

\csromanheader{2}{1. ปจฺจยานุโลม 2. สงฺขฺยาวาร}
\proseitem{336}{เหตุยา นว, อารมฺมเณ นว, อธิปติยา ปญฺจ, อนนฺตเร นว, สมนนฺตเร นว, สหชาเต นว, อญฺญมญฺเญ นว, นิสฺสเย นว, อุปนิสฺสเย นว, ปุเรชาเต นว, อาเสวเน นว, กมฺเม นว, (สพฺพตฺถ นว.) อวิคเต นว.}

\csromanheader{2}{2. ปจฺจยปจฺจนีย 1. วิภงฺควาร}
\csromanheader{2}{\textbf{นเหตุ}}
\proseitem{337}{\csromanfolio{500}อชฺฌตฺติกํ ธมฺมํ ปจฺจยา อชฺฌตฺติโก ธมฺโม อุปฺปชฺชติ นเหตุปจฺจยา.\csromanspacer{} อเหตุกปฏิสนฺธิกฺขเณ จิตฺตํ ปจฺจยา อชฺฌตฺติกํ กฏตฺตารูปํ.\csromanspacer{} จกฺขายตนํ ปจฺจยา จกฺขุวิญฺญาณํ. (สํขิตฺตํ.) (เอวํ นวปิ ปญฺหา กาตพฺพา, ปญฺจวิญฺญาณมฺปิ ปเวเสตฺวา ตีณิเยว โมโห.)}

\csromanheader{2}{2. ปจฺจยปจฺจนีย 2. สงฺขฺยาวาร}
\csromanheader{2}{\textbf{สุทฺธ}}
\vspace{\tipitakaparskip}
\csromancenter{\textbf{นเหตุ}ยา นว, น-อารมฺมเณ นว, น-อธิปติยา นว, น-อนนฺตเร นว, นสมนนฺตเร นว, น-อญฺญมญฺเญ นว, น-อุปนิสฺสเย นว, นปุเรชาเต นว, นปจฺฉาชาเต นว, น-อาเสวเน นว, นกมฺเม ตีณิ, นวิปาเก ปญฺจ, น-อาหาเร เอกํ, น-อินฺทฺริเย เอกํ, นฌาเน นว, นมคฺเค นว, นสมฺปยุตฺเต นว, นวิปฺปยุตฺเต ปญฺจ, โนนตฺถิยา นว, โนวิคเต นว.}
\csromancenter{(เอวํ อิตเร เทฺว คณนาปิ นิสฺสยวาโรปิ กาตพฺโพ.)}
\csromanheader{2}{66. อชฺฌตฺติกทุก 5. สํสฏฺฐวาร}
\csromanheader{2}{1. ปจฺจยานุโลมาทิ}
\proseitem{339}{อชฺฌตฺติกํ ธมฺมํ สํสฏฺโฐ พาหิโร ธมฺโม อุปฺปชฺชติ เหตุปจฺจยา.\csromanspacer{} จิตฺตํ สํสฏฺฐา สมฺปยุตฺตกา ขนฺธา.\csromanspacer{} ปฏิสนฺธิกฺขเณ จิตฺตํ สํสฏฺฐา สมฺปยุตฺตกา ขนฺธา. (1)}

\prose{พาหิรํ ธมฺมํ สํสฏฺโฐ พาหิโร ธมฺโม อุปฺปชฺชติ เหตุปจฺจยา.\csromanspacer{} พาหิรํ เอกํ ขนฺธํ สํสฏฺฐา เทฺว ขนฺธา.\csromanspacer{} เทฺว ขนฺเธ ฯเปฯ.\csromanspacer{} ปฏิสนฺธิกฺขเณ ฯเปฯ.\csromanspacer{} พาหิรํ ธมฺมํ สํสฏฺโฐ อชฺฌตฺติโก ธมฺโม อุปฺปชฺชติ เหตุปจฺจยา.\csromanspacer{} พาหิเร ขนฺเธ สํสฏฺฐํ จิตฺตํ.\csromanspacer{} ปฏิสนฺธิกฺขเณ ฯเปฯ.\csromanspacer{} พาหิรํ ธมฺมํ สํสฏฺโฐ อชฺฌตฺติโก จ พาหิโร จ ธมฺมา อุปฺปชฺชนฺติ เหตุปจฺจยา.\csromanspacer{} พาหิรํ เอกํ ขนฺธํ สํสฏฺฐา เทฺว ขนฺธา จิตฺตญฺจ.\csromanspacer{} เทฺว ขนฺเธ ฯเปฯ.\csromanspacer{} ปฏิสนฺธิกฺขเณ ฯเปฯ. (3)}

\csromanheader{2}{66. อชฺฌตฺติกทุก}
\prose{\csromanfolio{501}อชฺฌตฺติกญฺจ พาหิรญฺจ ธมฺมํ สํสฏฺโฐ พาหิโร ธมฺโม อุปฺปชฺชติ เหตุปจฺจยา.\csromanspacer{} พาหิรํ เอกํ ขนฺธญฺจ จิตฺตญฺจ สํสฏฺฐา เทฺว ขนฺธา.\csromanspacer{} เทฺว ขนฺเธ ฯเปฯ.\csromanspacer{} ปฏิสนฺธิกฺขเณ ฯเปฯ. (สํขิตฺตํ.)}

\prose{\textbf{เหตุ}ยา ปญฺจ, อารมฺมเณ ปญฺจ, อธิปติยา ปญฺจ, (สพฺพตฺถ ปญฺจ.) อวิคเต ปญฺจ. (อนุโลมํ.)}

\prose{อชฺฌตฺติกํ ธมฺมํ สํสฏฺโฐ พาหิโร ธมฺโม อุปฺปชฺชติ นเหตุปจฺจยา. (เอวํ ปญฺจ กาตพฺพา, ตีณิเยว โมโห.)}

\prose{นเหตุยา ปญฺจ, น-อธิปติยา ปญฺจ, นปุเรชาเต ปญฺจ, นปจฺฉาชาเต ปญฺจ, น-อาเสวเน ปญฺจ, นกมฺเม ตีณิ, นวิปาเก ปญฺจ, นฌาเน ปญฺจ, นมคฺเค ปญฺจ, นวิปฺปยุตฺเต ปญฺจ. (ปจฺจนียํ.)}

\prose{(เอวํ อิตเร เทฺว ค คณนาปิ สมฺปยุตฺตวาโรปิ กาตพฺโพ.)}

\csromanheader{2}{66. อชฺฌตฺติกทุก 7. ปญฺหาวาร}
\csromanheader{2}{1. ปจฺจยานุโลม 1. วิภงฺควาร}
\csromanheader{2}{\textbf{เหตุ}}
\proseitem{340}{พาหิโร ธมฺโม พาหิรสฺส ธมฺมสฺส เหตุปจฺจเยน ปจฺจโย.\csromanspacer{} พาหิรา เหตู สมฺปยุตฺตกานํ ขนฺธานํ จิตฺตสมุฏฺฐานานญฺจ รูปานํ เหตุปจฺจเยน ปจฺจโย.\csromanspacer{} ปฏิสนฺธิกฺขเณ พาหิรา เหตู สมฺปยุตฺตกานํ ขนฺธานํ พาหิรานญฺจ กฏตฺตารูปานํ เหตุปจฺจเยน ปจฺจโย. (1)}

\prose{พาหิโร ธมฺโม อชฺฌตฺติกสฺส ธมฺมสฺส เหตุปจฺจเยน ปจฺจโย.\csromanspacer{} พาหิรา เหตู จิตฺตสฺส เหตุปจฺจเยน ปจฺจโย.\csromanspacer{} ปฏิสนฺธิกฺขเณ พาหิรา เหตู จิตฺตสฺส อชฺฌตฺติกานญฺจ กฏตฺตารูปานํ เหตุปจฺจเยน ปจฺจโย. (2)}

\prose{พาหิโร ธมฺโม อชฺฌตฺติกสฺส จ พาหิรสฺส จ ธมฺมสฺส เหตุปจฺจเยน ปจฺจโย.\csromanspacer{} พาหิรา เหตู สมฺปยุตฺตกานํ ขนฺธานํ จิตฺตสฺส จ จิตฺตสมุฏฺฐานานญฺจ รูปานํ เหตุปจฺจเยน ปจฺจโย.\csromanspacer{} ปฏิสนฺธิกฺขเณ พาหิรา เหตู สมฺปยุตฺตกานํ ขนฺธานํ จิตฺตสฺส จ อชฺฌตฺติกานญฺจ พาหิรานญฺจ กฏตฺตารูปานํ เหตุปจฺจเยน ปจฺจโย. (3)}

\csromanheader{2}{\textbf{อารมฺมณ}}
\proseitem{341}{\csromanfolio{502}อชฺฌตฺติโก ธมฺโม อชฺฌตฺติกสฺส ธมฺมสฺส อารมฺมณปจฺจเยน ปจฺจโย.\csromanspacer{} จิตฺตํ อารพฺภ จิตฺตํ อุปฺปชฺชติ. (มูลํ ปุจฺฉิตพฺพํ.) จิตฺตํ อารพฺภ พาหิรา ขนฺธา อุปฺปชฺชนฺติ. (มูลํ ปุจฺฉิตพฺพํ.) จิตฺตํ อารพฺภ จิตฺตญฺจ สมฺปยุตฺตกา จ ขนฺธา อุปฺปชฺชนฺติ. (3)}

\prose{พาหิโร ธมฺโม พาหิรสฺส ธมฺมสฺส อารมฺมณปจฺจเยน ปจฺจโย.\csromanspacer{} ทานํ ฯเปฯ สีลํ ฯเปฯ อุโปสถกมฺมํ กตฺวา ตํ ปจฺจเวกฺขติ, อสฺสาเทติ อภินนฺทติ, ตํ อารพฺภ ราโค ฯเปฯ โทมนสฺสํ อุปฺปชฺชติ.\csromanspacer{} ปุพฺเพ สุจิณฺณานิ ปจฺจเวกฺขติ.\csromanspacer{} ฌานา วุฏฺฐหิตฺวา ฌานํ ฯเปฯ.\csromanspacer{} อริยา มคฺคา วุฏฺฐหิตฺวา มคฺคํ ฯเปฯ.\csromanspacer{} ผลํ ฯเปฯ.\csromanspacer{} นิพฺพานํ ปจฺจเวกฺขนฺติ.\csromanspacer{} นิพฺพานํ โคตฺรภุสฺส, โวทานสฺส, มคฺคสฺส, ผลสฺส, อาวชฺชนาย อารมฺมณปจฺจเยน ปจฺจโย.\csromanspacer{} อริยา พาหิเร ปหีเน กิเลเส ปจฺจเวกฺขนฺติ.\csromanspacer{} วิกฺขมฺภิเต กิเลเส ฯเปฯ.\csromanspacer{} ปุพฺเพ สมุทาจิณฺเณ กิเลเส ชานนฺติ.\csromanspacer{} รูเป ฯเปฯ.\csromanspacer{} วตฺถุํ พาหิเร ขนฺเธ อนิจฺจโต ฯเปฯ โทมนสฺสํ อุปฺปชฺชติ.\csromanspacer{} ทิพฺเพน จกฺขุนา รูปํ ปสฺสติ.\csromanspacer{} ทิพฺพาย โสตธาตุยา สทฺทํ สุณาติ.\csromanspacer{} เจโตปริยญาเณน พาหิรจิตฺตสมงฺคสฺส จิตฺตํ ชานาติ.\csromanspacer{} อากาสานญฺจายตนํ วิญฺญาณญฺจายตนสฺส ฯเปฯ.\csromanspacer{} อากิญฺจญฺญายตนํ เนวสญฺญานาสญฺญายตนสฺส ฯเปฯ.\csromanspacer{} รูปายตนํ จกฺขุวิญฺญาณสหคตานํ ขนฺธานํ อารมฺมณปจฺจเยน ปจฺจโย ฯเปฯ.\csromanspacer{} โผฏฺฐพฺพายตนํ.\csromanspacer{} พาหิรา ขนฺธา อิทฺธิวิธญาณสฺส, เจโตปริยญาณสฺส, ปุพฺเพนิวาสานุสฺสติญาณสฺส, ยถากมฺมูปคญาณสฺส, อนาคตํสญาณสฺส, อาวชฺชนาย อารมฺมณปจฺจเยน ปจฺจโย. (1)}

\prose{พาหิโร ธมฺโม อชฺฌตฺติกสฺส ธมฺมสฺส อารมฺมณปจฺจเยน ปจฺจโย.\csromanspacer{} ทานํ ฯเปฯ สีลํ ฯเปฯ อุโปสถกมฺมํ กตฺวา ตํ ปจฺจเวกฺขติ, อสฺสาเทติ อภินนฺทติ, ตํ อารพฺภ จิตฺตํ อุปฺปชฺชติ.\csromanspacer{} ปุพฺเพ สุจิณฺณานิ ฯเปฯ.\csromanspacer{} ฌานํ ฯเปฯ. (สํขิตฺตํ.\csromanspacer{} สพฺพํ กาตพฺพํ.) ปุพฺเพ สมุทาจิณฺเณ ฯเปฯ.\csromanspacer{} รูเป ฯเปฯ.\csromanspacer{} วตฺถุํ พาหิเร ขนฺเธ อนิจฺจโต ฯเปฯ วิปสฺสติ, อสฺสาเทติ อภินนฺทติ, ตํ อารพฺภ จิตฺตํ อุปฺปชฺชติ.\csromanspacer{} ทิพฺเพน จกฺขุนา รูปํ ปสฺสติ ฯเปฯ.\csromanspacer{} รูปายตนํ จกฺขุวิญฺญาณสฺส ฯเปฯ.\csromanspacer{} โผฏฺฐพฺพายตนํ ฯเปฯ.\csromanspacer{} พาหิรา ขนฺธา อิทฺธิวิธญาณสฺส, เจโตปริยญาณสฺส, ปุพฺเพนิวาสานุสฺสติญาณสฺส, ยถากมฺมูปคญาณสฺส, อนาคตํสญาณสฺส, อาวชฺชนาย อารมฺมณปจฺจเยน ปจฺจโย. (2)}

\csromanheader{2}{66. อชฺฌตฺติกทุก}
\prose{\csromanfolio{503}พาหิโร ธมฺโม อชฺฌตฺติกสฺส จ พาหิรสฺส จ ธมฺมสฺส อารมฺมณปจฺจเยน ปจฺจโย.\csromanspacer{} ทานํ ฯเปฯ สีลํ ฯเปฯ อุโปสถกมฺมํ กตฺวา ตํ ปจฺจเวกฺขติ, อสฺสาเทติ อภินนฺทติ, ตํ อารพฺภ จิตฺตญฺจ สมฺปยุตฺตกา จ ขนฺธา อุปฺปชฺชนฺติ. (สํขิตฺตํ.\csromanspacer{} สพฺพํ กาตพฺพํ.) พาหิเร ขนฺเธ อนิจฺจโต ฯเปฯ วิปสฺสติ, อสฺสาเทติ อภินนฺทติ, ตํ อารพฺภ จิตฺตญฺจ สมฺปยุตฺตกา จ ขนฺธา อุปฺปชฺชนฺติ.\csromanspacer{} ทิพฺเพน จกฺขุนา รูปํ ปสฺสติ ฯเปฯ.\csromanspacer{} รูปายตนํ จกฺขุวิญฺญาณสฺส จ สมฺปยุตฺตกานญฺจ ขนฺธานํ ฯเปฯ.\csromanspacer{} โผฏฺฐพฺพายตนํ ฯเปฯ.\csromanspacer{} พาหิรา ขนฺธา อิทฺธิวิธญาณสฺส, เจโตปริยญาณสฺส, ปุพฺเพนิวาสานุสฺสติญาณสฺส, ยถากมฺมูปคญาณสฺส, อนาคตํสญาณสฺส, อาวชฺชนาย อารมฺมณปจฺจเยน ปจฺจโย. (3)}

\prose{อชฺฌตฺติโก จ พาหิโร จ ธมฺมา อชฺฌตฺติกสฺส ธมฺมสฺส อารมฺมณปจฺจเยน ปจฺจโย.\csromanspacer{} ตีณิ.}

\csromanheader{2}{\textbf{อธิปติ}}
\proseitem{342}{อชฺฌตฺติโก ธมฺโม อชฺฌตฺติกสฺส ธมฺมสฺส อธิปติปจฺจเยน ปจฺจโย.\csromanspacer{} \textbf{อารมฺมณาธิปติ}\csromandash{}จิตฺตํ ครุํ กตฺวา จิตฺตํ อุปฺปชฺชติ. (1)}

\prose{อชฺฌตฺติโก ธมฺโม พาหิรสฺส ธมฺมสฺส อธิปติปจฺจเยน ปจฺจโย.\csromanspacer{} \textbf{อารมฺมณาธิปติ}, สหชาตาธิปติ.\csromanspacer{}\csromanspacer{}\textbf{อารมฺมณาธิปติ}\csromandash{}จิตฺตํ ครุํ กตฺวา พาหิรา ขนฺธา อุปฺปชฺชนฺติ.\csromanspacer{} \textbf{สหชาตาธิปติ}\csromandash{}จิตฺตาธิปติ สมฺปยุตฺตกานํ ขนฺธานํ จิตฺตสมุฏฺฐานานญฺจ รูปานํ อธิปติปจฺจเยน ปจฺจโย. (มูลํ.) \textbf{อารมฺมณาธิปติ}\csromandash{}อชฺฌตฺติกํ จิตฺตํ ครุํ กตฺวา จิตฺตญฺจ สมฺปยุตฺตกา จ ขนฺธา อุปฺปชฺชนฺติ. (3)}

\prose{พาหิโร ธมฺโม พาหิรสฺส ธมฺมสฺส อธิปติปจฺจเยน ปจฺจโย.\csromanspacer{} \textbf{อารมฺมณาธิปติ}, สหชาตาธิปติ.\csromanspacer{}\csromanspacer{}\textbf{อารมฺมณาธิปติ}\csromandash{}ทานํ ทตฺวา ฯเปฯ.\csromanspacer{} ตีณิ. (เทฺว อธิปตี, ติณฺณมฺปิ กาตพฺพา.) (3)}

\prose{อชฺฌตฺติโก จ พาหิโร จ ธมฺมา อชฺฌตฺติกสฺส ธมฺมสฺส อธิปติปจฺจเยน ปจฺจโย.\csromanspacer{} ตีณิ. (ติณฺณมฺปิ เอกาเยว อธิปติ.)}

\csromanheader{2}{\textbf{อนนฺตราทิ}}
\proseitem{343}{\csromanfolio{504}อชฺฌตฺติโก ธมฺโม อชฺฌตฺติกสฺส ธมฺมสฺส อนนฺตรปจฺจเยน ปจฺจโย.\csromanspacer{} ปุริมํ ปุริมํ จิตฺตํ ปจฺฉิมสฺส ปจฺฉิมสฺส จิตฺตสฺส อนนฺตรปจฺจเยน ปจฺจโย.\csromanspacer{} ตีณิ.}

\prose{พาหิโร ธมฺโม พาหิรสฺส ธมฺมสฺส อนนฺตรปจฺจเยน ปจฺจโย.\csromanspacer{} ปุริมา ปุริมา พาหิรา ขนฺธา ปจฺฉิมานํ ปจฺฉิมานํ ขนฺธานํ อนนฺตรปจฺจเยน ปจฺจโย.\csromanspacer{} อนุโลมํ โคตฺรภุสฺส ฯเปฯ.\csromanspacer{} ตีณิ. (ติณฺณมฺปิ เอกสทิสา.)}

\prose{อชฺฌตฺติโก จ พาหิโร จ ธมฺมา อชฺฌตฺติกสฺส ธมฺมสฺส อนนฺตรปจฺจเยน ปจฺจโย.\csromanspacer{} ตีณิ.\csromanspacer{} สมนนฺตรปจฺจเยน ปจฺจโย.\csromanspacer{} นว.\csromanspacer{} สหชาตปจฺจเยน ปจฺจโย.\csromanspacer{} นว. (ปฏิจฺจสทิสา.) อญฺญมญฺญปจฺจเยน ปจฺจโย.\csromanspacer{} ปญฺจ. (ปฏิจฺจสทิสา.) นิสฺสยปจฺจเยน ปจฺจโย.\csromanspacer{} นว. (ปจฺจยวารสทิสา.)}

\csromanheader{2}{\textbf{อุปนิสฺสย}}
\proseitem{344}{อชฺฌตฺติโก ธมฺโม อชฺฌตฺติกสฺส ธมฺมสฺส อุปนิสฺสยปจฺจเยน ปจฺจโย.\csromanspacer{} อารมฺมณูปนิสฺสโย, อนนฺตรูปนิสฺสโย.\csromanspacer{} ปกตูปนิสฺสโย ฯเปฯ.\csromanspacer{} ปกตูปนิสฺสโย\csromandash{}จิตฺตํ จิตฺตสฺส อุปนิสฺสยปจฺจเยน ปจฺจโย.\csromanspacer{} ตีณิ.}

\prose{พาหิโร ธมฺโม พาหิรสฺส ธมฺมสฺส อุปนิสฺสยปจฺจเยน ปจฺจโย.\csromanspacer{} อารมฺมณูปนิสฺสโย, อนนฺตรูปนิสฺสโย, ปกตูปนิสฺสโย ฯเปฯ.\csromanspacer{} ปกตูปนิสฺสโย\csromandash{}สทฺธํ อุปนิสฺสาย ทานํ เทติ ฯเปฯ มานํ ชปฺเปติ, ทิฏฺฐิํ คณฺหาติ.\csromanspacer{} สีลํ ฯเปฯ.\csromanspacer{} เสนาสนํ อุปนิสฺสาย ทานํ เทติ ฯเปฯ สํฆํ ภินฺทติ.\csromanspacer{} สทฺธา ฯเปฯ.\csromanspacer{} เสนาสนํ สทฺธาย ฯเปฯ ผลสมาปตฺติยา อุปนิสฺสยปจฺจเยน ปจฺจโย. (ตีณิปิ ปูเรตฺวา กาตพฺพา. “จิตฺตสฺสา”ติ กาตพฺพา. “สมฺปยุตฺตกานญฺจา”ติ กาตพฺพา.)}

\prose{อชฺฌตฺติโก จ พาหิโร จ ธมฺมา อชฺฌตฺติกสฺส ธมฺมสฺส อุปนิสฺสยปจฺจเยน ปจฺจโย.\csromanspacer{} ตีณิ.}

\csromanheader{2}{\textbf{ปุเรชาต}}
\proseitem{345}{อชฺฌตฺติโก ธมฺโม อชฺฌตฺติกสฺส ธมฺมสฺส ปุเรชาตปจฺจเยน ปจฺจโย.\csromanspacer{} อารมฺมณปุเรชาตํ, วตฺถุปุเรชาตํ.\csromanspacer{}\csromanspacer{}\textbf{อารมฺมณปุเรชาตํ\csromandash{}}}

\vspace{\tipitakaparskip}
\csromancenter{อชฺฌตฺติกทุก จกฺขุํ ฯเปฯ กายํ อนิจฺจโต ฯเปฯ วิปสฺสติ, อสฺสาเทติ อภินนฺทติ, ตํ อารพฺภ จิตฺตํ อุปฺปชฺชติ.\csromanspacer{} \textbf{วตฺถุปุเรชาตํ\csromandash{}}จกฺขายตนํ จกฺขุวิญฺญาณสฺส ฯเปฯ กายายตนํ กายวิญฺญาณสฺส ปุเรชาตปจฺจเยน ปจฺจโย. (1)}
\prose{\csromanfolio{505}อชฺฌตฺติโก ธมฺโม พาหิรสฺส ธมฺมสฺส ปุเรชาตปจฺจเยน ปจฺจโย.\csromanspacer{} \textbf{อารมฺมณปุเรชาตํ}, วตฺถุปุเรชาตํ.\csromanspacer{}\csromanspacer{}\textbf{อารมฺมณปุเรชาตํ}\csromandash{}จกฺขุํ ฯเปฯ กายํ อนิจฺจโต ฯเปฯ วิปสฺสติ, อสฺสาเทติ ฯเปฯ โทมนสฺสํ อุปฺปชฺชติ.\csromanspacer{} \textbf{วตฺถุปุเรชาตํ\csromandash{}}จกฺขายตนํ จกฺขุวิญฺญาณสหคตานํ ขนฺธานํ ฯเปฯ กายายตนํ กายวิญฺญาณสหคตานํ ขนฺธานํ ปุเรชาตปจฺจเยน ปจฺจโย. (2)}

\prose{อชฺฌตฺติโก ธมฺโม อชฺฌตฺติกสฺส จ พาหิรสฺส จ ธมฺมสฺส ปุเรชาตปจฺจเยน ปจฺจโย.\csromanspacer{} \textbf{อารมฺมณปุเรชาตํ}, วตฺถุปุเรชาตํ.\csromanspacer{}\csromanspacer{}\textbf{อารมฺมณปุเรชาตํ}\csromandash{}จกฺขุํ ฯเปฯ กายํ อนิจฺจโต ฯเปฯ วิปสฺสติ, ตํ อารพฺภ จิตฺตญฺจ สมฺปยุตฺตกา ขนฺธา จ อุปฺปชฺชนฺติ.\csromanspacer{} \textbf{วตฺถุปุเรชาตํ\csromandash{}} จกฺขายตนํ จกฺขุวิญฺญาณสฺส สมฺปยุตฺตกานญฺจ ขนฺธานํ ฯเปฯ กายายตนํ กายวิญฺญาณสฺส สมฺปยุตฺตกานญฺจ ขนฺธานํ ปุเรชาตปจฺจเยน ปจฺจโย. (3)}

\proseitem{346}{พาหิโร ธมฺโม พาหิรสฺส ธมฺมสฺส ปุเรชาตปจฺจเยน ปจฺจโย.\csromanspacer{} \textbf{อารมฺมณปุเรชาตํ}, วตฺถุปุเรชาตํ.\csromanspacer{}\csromanspacer{}\textbf{อารมฺมณปุเรชาตํ}\csromandash{}รูเป ฯเปฯ โผฏฺฐพฺเพ.\csromanspacer{} วตฺถุํ อนิจฺจโต ฯเปฯ โทมนสฺสํ อุปฺปชฺชติ.\csromanspacer{} ทิพฺเพน จกฺขุนา รูปํ ปสฺสติ.\csromanspacer{} ทิพฺพาย โสตธาตุยา สทฺทํ สุณาติ.\csromanspacer{} \textbf{วตฺถุปุเรชาตํ\csromandash{}}วตฺถุ พาหิรานํ ขนฺธานํ ปุเรชาตปจฺจเยน ปจฺจโย. (1)}

\prose{พาหิโร ธมฺโม อชฺฌตฺติกสฺส ธมฺมสฺส ปุเรชาตปจฺจเยน ปจฺจโย.\csromanspacer{} \textbf{อารมฺมณปุเรชาตํ}, วตฺถุปุเรชาตํ.\csromanspacer{}\csromanspacer{}\textbf{อารมฺมณปุเรชาตํ}\csromandash{}รูเป ฯเปฯ โผฏฺฐพฺเพ.\csromanspacer{} วตฺถุํ อนิจฺจโต ฯเปฯ ตํ อารพฺภ จิตฺตํ อุปฺปชฺชติ.\csromanspacer{} ทิพฺเพน จกฺขุนา รูปํ ปสฺสติ.\csromanspacer{} ทิพฺพาย โสตธาตุยา สทฺทํ สุณาติ.\csromanspacer{} \textbf{วตฺถุปุเรชาตํ\csromandash{}}วตฺถุ จิตฺตสฺส ปุเรชาตปจฺจเยน ปจฺจโย. (2)}

\prose{พาหิโร ธมฺโม อชฺฌตฺติกสฺส จ พาหิรสฺส จ ธมฺมสฺส ปุเรชาตปจฺจเยน ปจฺจโย.\csromanspacer{} \textbf{อารมฺมณปุเรชาตํ}, วตฺถุปุเรชาตํ.\csromanspacer{}\csromanspacer{}\textbf{อารมฺมณปุเรชาตํ}\csromandash{}รูเป ฯเปฯ โผฏฺฐพฺเพ.\csromanspacer{} วตฺถุํ อนิจฺจโต ฯเปฯ ตํ อารพฺภ จิตฺตญฺจ สมฺปยุตฺตกา ขนฺธา จ อุปฺปชฺชนฺติ.\csromanspacer{} ทิพฺเพน จกฺขุนา รูปํ ปสฺสติ.\csromanspacer{} ทิพฺพาย โสตธาตุยา \csromanfolio{506}สทฺทํ สุณาติ.\csromanspacer{} \textbf{วตฺถุปุเรชาตํ}\csromandash{}วตฺถุ จิตฺตสฺส สมฺปยุตฺตกานญฺจ ขนฺธานํ ปุเรชาตปจฺจเยน ปจฺจโย. (3)}

\proseitem{347}{อชฺฌตฺติโก จ พาหิโร จ ธมฺมา อชฺฌตฺติกสฺส ธมฺมสฺส ปุเรชาตปจฺจเยน ปจฺจโย.\csromanspacer{} \textbf{อารมฺมณปุเรชาตํ, วตฺถุปุเรชาตํ.}\csromanspacer{} จกฺขายตนญฺจ วตฺถุ จ จิตฺตสฺส ฯเปฯ.\csromanspacer{} กายายตนญฺจ วตฺถุ จ จิตฺตสฺส ปุเรชาตปจฺจเยน ปจฺจโย.\csromanspacer{} รูปายตนญฺจ จกฺขายตนญฺจ จกฺขุวิญฺญาณสฺส ฯเปฯ.\csromanspacer{} โผฏฺฐพฺพายตนญฺจ กายายตนญฺจ กายวิญฺญาณสฺส ปุเรชาตปจฺจเยน ปจฺจโย. (1)}

\prose{อชฺฌตฺติโก จ พาหิโร จ ธมฺมา พาหิรสฺส ธมฺมสฺส ปุเรชาตปจฺจเยน ปจฺจโย.\csromanspacer{} \textbf{อารมฺมณปุเรชาตํ, วตฺถุปุเรชาตํ.}\csromanspacer{} จกฺขายตนญฺจ วตฺถุ จ พาหิรานํ ขนฺธานํ ปุเรชาตปจฺจเยน ปจฺจโย ฯเปฯ.\csromanspacer{} กายายตนญฺจ วตฺถุ จ พาหิรานํ ขนฺธานํ ปุเรชาตปจฺจเยน ปจฺจโย.\csromanspacer{} รูปายตนญฺจ จกฺขายตนญฺจ จกฺขุวิญฺญาณสหคตานํ ขนฺธานํ ปุเรชาตปจฺจเยน ปจฺจโย ฯเปฯ.\csromanspacer{} โผฏฺฐพฺพายตนญฺจ กายายตนญฺจ กายวิญฺญาณสหคตานํ ขนฺธานํ ปุเรชาตปจฺจเยน ปจฺจโย. (2)}

\prose{อชฺฌตฺติโก จ พาหิโร จ ธมฺมา อชฺฌตฺติกสฺส จ พาหิรสฺส จ ธมฺมสฺส ปุเรชาตปจฺจเยน ปจฺจโย.\csromanspacer{} \textbf{อารมฺมณปุเรชาตํ, วตฺถุปุเรชาตํ.}\csromanspacer{} จกฺขายตนญฺจ วตฺถุ จ จิตฺตสฺส สมฺปยุตฺตกานญฺจ ขนฺธานํ ปุเรชาตปจฺจเยน ปจฺจโย ฯเปฯ.\csromanspacer{} กายายตนญฺจ วตฺถุ จ ฯเปฯ.\csromanspacer{} รูปายตนญฺจ จกฺขายตนญฺจ จกฺขุวิญฺญาณสฺส สมฺปยุตฺตกานญฺจ ขนฺธานํ ปุเรชาตปจฺจเยน ปจฺจโย ฯเปฯ.\csromanspacer{} โผฏฺฐพฺพายตนญฺจ ฯเปฯ. (3)}

\csromanheader{2}{\textbf{ปจฺฉาชาตาเสวน}}
\proseitem{348}{อชฺฌตฺติโก ธมฺโม อชฺฌตฺติกสฺส ธมฺมสฺส ปจฺฉาชาตปจฺจเยน ปจฺจโย.\csromanspacer{} ปจฺฉาชาตา อชฺฌตฺติกา ขนฺธา ปุเรชาตสฺส อิมสฺส อชฺฌตฺติกสฺส กายสฺส ปจฺฉาชาตปจฺจเยน ปจฺจโย. (มูลํ กาตพฺพํ.) ปจฺฉาชาตา อชฺฌตฺติกา ขนฺธา ปุเรชาตสฺส อิมสฺส พาหิรสฺส กายสฺส ปจฺฉาชาตปจฺจเยน ปจฺจโย. (มูลํ กาตพฺพํ.) ปจฺฉาชาตา อชฺฌตฺติกา ขนฺธา ปุเรชาตสฺส อิมสฺส อชฺฌตฺติกสฺส จ พาหิรสฺส จ}

\vspace{\tipitakaparskip}
\csromancenter{อชฺฌตฺติกทุก กายสฺส ปจฺฉาชาตปจฺจเยน ปจฺจโย. (เอวํ นวปิ ปญฺหา กาตพฺพา.) อาเสวนปจฺจเยน ปจฺจโย. (นว ปญฺหา กาตพฺพา.)}
\csromanheader{2}{\textbf{กมฺม วิปาก}}
\proseitem{349}{\csromanfolio{507}พาหิโร ธมฺโม พาหิรสฺส ธมฺมสฺส กมฺมปจฺจเยน ปจฺจโย.\csromanspacer{} \textbf{สหชาตา}, นานากฺขณิกา.\csromanspacer{}\csromanspacer{}\textbf{สหชาตา} พาหิรา เจตนา สมฺปยุตฺตกานํ ขนฺธานํ จิตฺตสมุฏฺฐานานญฺจ รูปานํ กมฺมปจฺจเยน ปจฺจโย.\csromanspacer{} ปฏิสนฺธิกฺขเณ ฯเปฯ.\csromanspacer{} \textbf{นานากฺขณิกา} พาหิรา เจตนา วิปากานํ พาหิรานํ ขนฺธานํ กฏตฺตา จ รูปานํ กมฺมปจฺจเยน ปจฺจโย. (1)}

\prose{พาหิโร ธมฺโม อชฺฌตฺติกสฺส ธมฺมสฺส กมฺมปจฺจเยน ปจฺจโย.\csromanspacer{} \textbf{สหชาตา}, นานากฺขณิกา.\csromanspacer{}\csromanspacer{}\textbf{สหชาตา} พาหิรา เจตนา จิตฺตสฺส กมฺมปจฺจเยน ปจฺจโย.\csromanspacer{} ปฏิสนฺธิกฺขเณ ฯเปฯ.\csromanspacer{} \textbf{นานากฺขณิกา} พาหิรา เจตนา วิปากสฺส จิตฺตสฺส อชฺฌตฺติกานญฺจ กฏตฺตารูปานํ กมฺมปจฺจเยน ปจฺจโย. (2)}

\prose{พาหิโร ธมฺโม อชฺฌตฺติกสฺส จ พาหิรสฺส จ ธมฺมสฺส กมฺมปจฺจเยน ปจฺจโย.\csromanspacer{} \textbf{สหชาตา}, นานากฺขณิกา.\csromanspacer{}\csromanspacer{}\textbf{สหชาตา} พาหิรา เจตนา สมฺปยุตฺตกานํ ขนฺธานํ จิตฺตสฺส จ จิตฺตสมุฏฺฐานานญฺจ รูปานํ กมฺมปจฺจเยน ปจฺจโย.\csromanspacer{} ปฏิสนฺธิกฺขเณ ฯเปฯ.\csromanspacer{} \textbf{นานากฺขณิกา} พาหิรา เจตนา วิปากานํ ขนฺธานํ จิตฺตสฺส จ อชฺฌตฺติกานญฺจ พาหิรานญฺจ กฏตฺตารูปานํ กมฺมปจฺจเยน ปจฺจโย. (3) วิปากปจฺจเยน ปจฺจโย.\csromanspacer{} นว.}

\csromanheader{2}{\textbf{อาหาร}}
\proseitem{350}{อชฺฌตฺติโก ธมฺโม อชฺฌตฺติกสฺส ธมฺมสฺส อาหารปจฺจเยน ปจฺจโย.\csromanspacer{} ปฏิสนฺธิกฺขเณ อชฺฌตฺติกา อาหารา อชฺฌตฺติกานํ กฏตฺตารูปานํ อาหารปจฺจเยน ปจฺจโย. (1)}

\prose{อชฺฌตฺติโก ธมฺโม พาหิรสฺส ธมฺมสฺส อาหารปจฺจเยน ปจฺจโย.\csromanspacer{} อชฺฌตฺติกา อาหารา สมฺปยุตฺตกานํ ขนฺธานํ จิตฺตสมุฏฺฐานานญฺจ รูปานํ อาหารปจฺจเยน ปจฺจโย.\csromanspacer{} ปฏิสนฺธิกฺขเณ อชฺฌตฺติกา อาหารา สมฺปยุตฺตกานํ ขนฺธานํ พาหิรานญฺจ กฏตฺตารูปานํ อาหารปจฺจเยน ปจฺจโย. (มูลํ กาตพฺพํ.) ปฏิสนฺธิกฺขเณ อชฺฌตฺติกา อาหารา \csromanfolio{508}สมฺปยุตฺตกานํ ขนฺธานํ อชฺฌตฺติกานญฺจ พาหิรานญฺจ กฏตฺตารูปานํ อาหารปจฺจเยน ปจฺจโย. (3)}

\proseitem{351}{พาหิโร ธมฺโม พาหิรสฺส ธมฺมสฺส อาหารปจฺจเยน ปจฺจโย.\csromanspacer{} พาหิรา อาหารา สมฺปยุตฺตกานํ ขนฺธานํ จิตฺตสมุฏฺฐานานญฺจ รูปานํ อาหารปจฺจเยน ปจฺจโย.\csromanspacer{} ปฏิสนฺธิกฺขเณ ฯเปฯ.\csromanspacer{} พาหิโร กพฬีกาโร อาหาโร พาหิรสฺส กายสฺส อาหารปจฺจเยน ปจฺจโย. (1)}

\prose{พาหิโร ธมฺโม อชฺฌตฺติกสฺส ธมฺมสฺส อาหารปจฺจเยน ปจฺจโย.\csromanspacer{} พาหิรา อาหารา จิตฺตสฺส อาหารปจฺจเยน ปจฺจโย.\csromanspacer{} ปฏิสนฺธิกฺขเณ พาหิรา อาหารา จิตฺตสฺส อชฺฌตฺติกานญฺจ กฏตฺตารูปานํ อาหารปจฺจเยน ปจฺจโย.\csromanspacer{} พาหิโร กพฬีกาโร อาหาโร อชฺฌตฺติกสฺส กายสฺส อาหารปจฺจเยน ปจฺจโย. (2)}

\prose{พาหิโร ธมฺโม อชฺฌตฺติกสฺส จ พาหิรสฺส จ ธมฺมสฺส อาหารปจฺจเยน ปจฺจโย.\csromanspacer{} พาหิรา อาหารา สมฺปยุตฺตกานํ ขนฺธานํ จิตฺตสฺส จ จิตฺตสมุฏฺฐานานญฺจ รูปานํ อาหารปจฺจเยน ปจฺจโย.\csromanspacer{} ปฏิสนฺธิกฺขเณ พาหิรา อาหารา สมฺปยุตฺตกานํ ขนฺธานํ จิตฺตสฺส จ อชฺฌตฺติกานญฺจ พาหิรานญฺจ กฏตฺตารูปานํ อาหารปจฺจเยน ปจฺจโย.\csromanspacer{} พาหิโร กพฬีกาโร อาหาโร อชฺฌตฺติกสฺส จ พาหิรสฺส จ กายสฺส อาหารปจฺจเยน ปจฺจโย. (3)}

\proseitem{352}{อชฺฌตฺติโก จ พาหิโร จ ธมฺมา อชฺฌตฺติกสฺส ธมฺมสฺส อาหารปจฺจเยน ปจฺจโย.\csromanspacer{} ปฏิสนฺธิกฺขเณ อชฺฌตฺติกา จ พาหิรา จ อาหารา อชฺฌตฺติกานํ กฏตฺตารูปานํ อาหารปจฺจเยน ปจฺจโย. (1)}

\prose{อชฺฌตฺติโก จ พาหิโร จ ธมฺมา พาหิรสฺส ธมฺมสฺส อาหารปจฺจเยน ปจฺจโย.\csromanspacer{} อชฺฌตฺติกา จ พาหิรา จ อาหารา สมฺปยุตฺตกานํ ขนฺธานํ จิตฺตสมุฏฺฐานานญฺจ รูปานํ อาหารปจฺจเยน ปจฺจโย.\csromanspacer{} ปฏิสนฺธิกฺขเณ อชฺฌตฺติกา จ พาหิรา จ อาหารา สมฺปยุตฺตกานํ ขนฺธานํ พาหิรานญฺจ กฏตฺตารูปานํ อาหารปจฺจเยน ปจฺจโย. (2)}

\prose{อชฺฌตฺติโก จ พาหิโร จ ธมฺมา อชฺฌตฺติกสฺส จ พาหิรสฺส จ ธมฺมสฺส อาหารปจฺจเยน ปจฺจโย.\csromanspacer{} ปฏิสนฺธิกฺขเณ อชฺฌตฺติกา จ พาหิรา จ}

\vspace{\tipitakaparskip}
\csromancenter{อชฺฌตฺติกทุก อาหารา สมฺปยุตฺตกานํ ขนฺธานํ อชฺฌตฺติกานญฺจ พาหิรานญฺจ กฏตฺตารูปานํ อาหารปจฺจเยน ปจฺจโย. (3)}
\csromanheader{2}{\textbf{อินฺทฺริย}}
\proseitem{353}{\csromanfolio{509}อชฺฌตฺติโก ธมฺโม อชฺฌตฺติกสฺส ธมฺมสฺส อินฺทฺริยปจฺจเยน ปจฺจโย.\csromanspacer{} ปฏิสนฺธิกฺขเณ อชฺฌตฺติกา อินฺทฺริยา อชฺฌตฺติกานํ กฏตฺตารูปานํ อินฺทฺริยปจฺจเยน ปจฺจโย.\csromanspacer{} จกฺขุนฺทฺริยํ จกฺขุวิญฺญาณสฺส ฯเปฯ.\csromanspacer{} กายินฺทฺริยํ กายวิญฺญาณสฺส อินฺทฺริยปจฺจเยน ปจฺจโย. (1)}

\prose{อชฺฌตฺติโก ธมฺโม พาหิรสฺส ธมฺมสฺส อินฺทฺริยปจฺจเยน ปจฺจโย.\csromanspacer{} อชฺฌตฺติกา อินฺทฺริยา สมฺปยุตฺตกานํ ขนฺธานํ จิตฺตสมุฏฺฐานานญฺจ รูปานํ อินฺทฺริยปจฺจเยน ปจฺจโย.\csromanspacer{} ปฏิสนฺธิกฺขเณ อชฺฌตฺติกา อินฺทฺริยา สมฺปยุตฺตกานํ ขนฺธานํ พาหิรานญฺจ กฏตฺตารูปานํ อินฺทฺริยปจฺจเยน ปจฺจโย.\csromanspacer{} จกฺขุนฺทฺริยํ จกฺขุวิญฺญาณสหคตานํ ขนฺธานํ ฯเปฯ.\csromanspacer{} กายินฺทฺริยํ กายวิญฺญาณสหคตานํ ขนฺธานํ อินฺทฺริยปจฺจเยน ปจฺจโย. (2)}

\prose{อชฺฌตฺติโก ธมฺโม อชฺฌตฺติกสฺส จ พาหิรสฺส จ ธมฺมสฺส อินฺทฺริยปจฺจเยน ปจฺจโย.\csromanspacer{} ปฏิสนฺธิกฺขเณ อชฺฌตฺติกา อินฺทฺริยา สมฺปยุตฺตกานํ ขนฺธานํ อชฺฌตฺติกานญฺจ พาหิรานญฺจ กฏตฺตารูปานํ อินฺทฺริยปจฺจเยน ปจฺจโย.\csromanspacer{} จกฺขุนฺทฺริยํ จกฺขุวิญฺญาณสฺส สมฺปยุตฺตกานญฺจ ขนฺธานํ ฯเปฯ.\csromanspacer{} กายินฺทฺริยํ กายวิญฺญาณสฺส สมฺปยุตฺตกานญฺจ ขนฺธานํ อินฺทฺริยปจฺจเยน ปจฺจโย. (3)}

\proseitem{354}{พาหิโร ธมฺโม พาหิรสฺส ธมฺมสฺส อินฺทฺริยปจฺจเยน ปจฺจโย.\csromanspacer{} พาหิรา อินฺทฺริยา สมฺปยุตฺตกานํ ขนฺธานํ จิตฺตสมุฏฺฐานานญฺจ รูปานํ อินฺทฺริยปจฺจเยน ปจฺจโย.\csromanspacer{} ปฏิสนฺธิกฺขเณ พาหิรา อินฺทฺริยา สมฺปยุตฺตกานํ ขนฺธานํ พาหิรานญฺจ กฏตฺตารูปานํ อินฺทฺริยปจฺจเยน ปจฺจโย.\csromanspacer{} รูปชีวิตินฺทฺริยํ พาหิรานํ กฏตฺตารูปานํ อินฺทฺริยปจฺจเยน ปจฺจโย. (1)}

\prose{พาหิโร ธมฺโม อชฺฌตฺติกสฺส ธมฺมสฺส อินฺทฺริยปจฺจเยน ปจฺจโย.\csromanspacer{} พาหิรา อินฺทฺริยา จิตฺตสฺส อินฺทฺริยปจฺจเยน ปจฺจโย.\csromanspacer{} ปฏิสนฺธิกฺขเณ พาหิรา อินฺทฺริยา จิตฺตสฺส อชฺฌตฺติกานญฺจ กฏตฺตารูปานํ อินฺทฺริยปจฺจเยน ปจฺจโย.\csromanspacer{} รูปชีวิตินฺทฺริยํ อชฺฌตฺติกานํ กฏตฺตารูปานํ อินฺทฺริยปจฺจเยน ปจฺจโย.}

\prose{\csromanfolio{510}พาหิโร ธมฺโม อชฺฌตฺติกสฺส จ พาหิรสฺส จ ธมฺมสฺส อินฺทฺริยปจฺจเยน ปจฺจโย.\csromanspacer{} พาหิรา อินฺทฺริยา สมฺปยุตฺตกานํ ขนฺธานํ จิตฺตสฺส จ จิตฺตสมุฏฺฐานานญฺจ รูปานํ อินฺทฺริยปจฺจเยน ปจฺจโย.\csromanspacer{} ปฏิสนฺธิกฺขเณ พาหิรา อินฺทฺริยา สมฺปยุตฺตกานํ ขนฺธานํ จิตฺตสฺส จ อชฺฌตฺติกานญฺจ พาหิรานญฺจ กฏตฺตารูปานํ อินฺทฺริยปจฺจเยน ปจฺจโย.\csromanspacer{} รูปชีวิตินฺทฺริยํ อชฺฌตฺติกานญฺจ พาหิรานญฺจ กฏตฺตารูปานํ อินฺทฺริยปจฺจเยน ปจฺจโย. (3)}

\proseitem{355}{อชฺฌตฺติโก จ พาหิโร จ ธมฺมา อชฺฌตฺติกสฺส ธมฺมสฺส อินฺทฺริยปจฺจเยน ปจฺจโย.\csromanspacer{} ปฏิสนฺธิกฺขเณ อชฺฌตฺติกา จ พาหิรา จ อินฺทฺริยา อชฺฌตฺติกานํ กฏตฺตารูปานํ อินฺทฺริยปจฺจเยน ปจฺจโย.\csromanspacer{} จกฺขุนฺทฺริยญฺจ อุเปกฺขินฺทฺริยญฺจ จกฺขุวิญฺญาณสฺส อินฺทฺริยปจฺจเยน ปจฺจโย ฯเปฯ.\csromanspacer{} กายินฺทฺริยญฺจ สุขินฺทฺริยญฺจ ฯเปฯ.\csromanspacer{} กายินฺทฺริยญฺจ ทุกฺขินฺทฺริยญฺจ กายวิญฺญาณสฺส อินฺทฺริยปจฺจเยน ปจฺจโย. (1)}

\prose{อชฺฌตฺติโก จ พาหิโร จ ธมฺมา พาหิรสฺส ธมฺมสฺส อินฺทฺริยปจฺจเยน ปจฺจโย.\csromanspacer{} อชฺฌตฺติกา จ พาหิรา จ อินฺทฺริยา สมฺปยุตฺตกานํ ขนฺธานํ จิตฺตสมุฏฺฐานานญฺจ รูปานํ อินฺทฺริยปจฺจเยน ปจฺจโย.\csromanspacer{} ปฏิสนฺธิกฺขเณ อชฺฌตฺติกา จ พาหิรา จ อินฺทฺริยา สมฺปยุตฺตกานํ ขนฺธานํ พาหิรานญฺจ กฏตฺตารูปานํ อินฺทฺริยปจฺจเยน ปจฺจโย.\csromanspacer{} จกฺขุนฺทฺริยญฺจ อุเปกฺขินฺทฺริยญฺจ จกฺขุวิญฺญาณสหคตานํ ขนฺธานํ อินฺทฺริยปจฺจเยน ปจฺจโย ฯเปฯ.\csromanspacer{} กายินฺทฺริยญฺจ สุขินฺทฺริยญฺจ ฯเปฯ.\csromanspacer{} กายินฺทฺริยญฺจ ทุกฺขินฺทฺริยญฺจ กายวิญฺญาณสหคตานํ ขนฺธานํ อินฺทฺริยปจฺจเยน ปจฺจโย. (2)}

\prose{อชฺฌตฺติโก จ พาหิโร จ ธมฺมา อชฺฌตฺติกสฺส จ พาหิรสฺส จ ธมฺมสฺส อินฺทฺริยปจฺจเยน ปจฺจโย.\csromanspacer{} ปฏิสนฺธิกฺขเณ อชฺฌตฺติกา จ พาหิรา จ อินฺทฺริยา สมฺปยุตฺตกานํ ขนฺธานํ อชฺฌตฺติกานญฺจ พาหิรานญฺจ กฏตฺตารูปานํ อินฺทฺริยปจฺจเยน ปจฺจโย.\csromanspacer{} จกฺขุนฺทฺริยญฺจ อุเปกฺขินฺทฺริยญฺจ จกฺขุวิญฺญาณสฺส สมฺปยุตฺตกานญฺจ ขนฺธานํ อินฺทฺริยปจฺจเยน ปจฺจโย.\csromanspacer{} กายินฺทฺริยญฺจ ฯเปฯ. (3)}

\csromanheader{2}{\textbf{ฌานาทิ}}
\proseitem{356}{พาหิโร ธมฺโม พาหิรสฺส ธมฺมสฺส ฌานปจฺจเยน ปจฺจโย.\csromanspacer{} ตีณิ.\csromanspacer{} มคฺคปจฺจเยน ปจฺจโย.\csromanspacer{} ตีณิ.\csromanspacer{} สมฺปยุตฺตปจฺจเยน ปจฺจโย.\csromanspacer{} ปญฺจ.}

\csromanheader{2}{66. อชฺฌตฺติกทุก \textbf{วิปฺปยุตฺต}}
\proseitem{357}{\csromanfolio{511}อชฺฌตฺติโก ธมฺโม อชฺฌตฺติกสฺส ธมฺมสฺส วิปฺปยุตฺตปจฺจเยน ปจฺจโย.\csromanspacer{} \textbf{สหชาตํ}, ปุเรชาตํ, ปจฺฉาชาตํ.\csromanspacer{}\csromanspacer{}\textbf{สหชาตํ} ปฏิสนฺธิกฺขเณ จิตฺตํ อชฺฌตฺติกานํ กฏตฺตารูปานํ วิปฺปยุตฺตปจฺจเยน ปจฺจโย.\csromanspacer{} \textbf{ปุเรชาตํ} จกฺขายตนํ จกฺขุวิญฺญาณสฺส ฯเปฯ.\csromanspacer{} กายายตนํ กายวิญฺญาณสฺส วิปฺปยุตฺตปจฺจเยน ปจฺจโย.\csromanspacer{} \textbf{ปจฺฉาชาตา} อชฺฌตฺติกา ขนฺธา ปุเรชาตสฺส อิมสฺส อชฺฌตฺติกสฺส กายสฺส วิปฺปยุตฺตปจฺจเยน ปจฺจโย. (1)}

\prose{อชฺฌตฺติโก ธมฺโม พาหิรสฺส ธมฺมสฺส วิปฺปยุตฺตปจฺจเยน ปจฺจโย.\csromanspacer{} \textbf{สหชาตํ}, ปุเรชาตํ, ปจฺฉาชาตํ.\csromanspacer{}\csromanspacer{}\textbf{สหชาตา} อชฺฌตฺติกา ขนฺธา จิตฺตสมุฏฺฐานานํ รูปานํ วิปฺปยุตฺตปจฺจเยน ปจฺจโย.\csromanspacer{} ปฏิสนฺธิกฺขเณ ฯเปฯ.\csromanspacer{} \textbf{ปุเรชาตํ} จกฺขายตนํ จกฺขุวิญฺญาณสหคตานํ ขนฺธานํ ฯเปฯ กายายตนํ กายวิญฺญาณสหคตานํ ขนฺธานํ วิปฺปยุตฺตปจฺจเยน ปจฺจโย.\csromanspacer{} \textbf{ปจฺฉาชาตา} อชฺฌตฺติกา ขนฺธา ปุเรชาตสฺส อิมสฺส พาหิรสฺส กายสฺส วิปฺปยุตฺตปจฺจเยน ปจฺจโย. (2)}

\prose{อชฺฌตฺติโก ธมฺโม อชฺฌตฺติกสฺส จ พาหิรสฺส จ ธมฺมสฺส วิปฺปยุตฺตปจฺจเยน ปจฺจโย.\csromanspacer{} \textbf{สหชาตํ}, ปุเรชาตํ, ปจฺฉาชาตํ.\csromanspacer{}\csromanspacer{}\textbf{สหชาตา} ปฏิสนฺธิกฺขเณ อชฺฌตฺติกา ขนฺธา อชฺฌตฺติกานญฺจ พาหิรานญฺจ กฏตฺตารูปานํ วิปฺปยุตฺตปจฺจเยน ปจฺจโย.\csromanspacer{} \textbf{ปุเรชาตํ} จกฺขายตนํ จกฺขุวิญฺญาณสฺส สมฺปยุตฺตกานญฺจ ขนฺธานํ ฯเปฯ.\csromanspacer{} กายายตนํ กายวิญฺญาณสฺส สมฺปยุตฺตกานํ ขนฺธานํ วิปฺปยุตฺตปจฺจเยน ปจฺจโย.\csromanspacer{} \textbf{ปจฺฉาชาตา} อชฺฌตฺติกา ขนฺธา ปุเรชาตสฺส อิมสฺส อชฺฌตฺติกสฺส จ พาหิรสฺส จ กายสฺส วิปฺปยุตฺตปจฺจเยน ปจฺจโย. (3)}

\proseitem{358}{พาหิโร ธมฺโม พาหิรสฺส ธมฺมสฺส วิปฺปยุตฺตปจฺจเยน ปจฺจโย.\csromanspacer{} \textbf{สหชาตํ}, ปุเรชาตํ, ปจฺฉาชาตํ.\csromanspacer{}\csromanspacer{}\textbf{สหชาตา} พาหิรา ขนฺธา จิตฺตสมุฏฺฐานานํ รูปานํ วิปฺปยุตฺตปจฺจเยน ปจฺจโย.\csromanspacer{} ปฏิสนฺธิกฺขเณ ฯเปฯ.\csromanspacer{} ขนฺธา วตฺถุสฺส วิปฺปยุตฺตปจฺจเยน ปจฺจโย.\csromanspacer{} วตฺถุ ขนฺธานํ วิปฺปยุตฺตปจฺจเยน ปจฺจโย.\csromanspacer{} \textbf{ปุเรชาตํ} วตฺถุ พาหิรานํ ขนฺธานํ วิปฺปยุตฺตปจฺจเยน ปจฺจโย.\csromanspacer{} \textbf{ปจฺฉาชาตา} พาหิรา ขนฺธา ปุเรชาตสฺส อิมสฺส พาหิรสฺส กายสฺส วิปฺปยุตฺตปจฺจเยน ปจฺจโย. (1)}

\prose{\csromanfolio{512}พาหิโร ธมฺโม อชฺฌตฺติกสฺส ธมฺมสฺส วิปฺปยุตฺตปจฺจเยน ปจฺจโย.\csromanspacer{} \textbf{สหชาตํ}, ปุเรชาตํ, ปจฺฉาชาตํ.\csromanspacer{}\csromanspacer{}\textbf{สหชาตา} ปฏิสนฺธิกฺขเณ พาหิรา ขนฺธา อชฺฌตฺติกานํ กฏตฺตารูปานํ วิปฺปยุตฺตปจฺจเยน ปจฺจโย.\csromanspacer{} ปฏิสนฺธิกฺขเณ วตฺถุ จิตฺตสฺส วิปฺปยุตฺตปจฺจเยน ปจฺจโย.\csromanspacer{} \textbf{ปุเรชาตํ} วตฺถุ จิตฺตสฺส วิปฺปยุตฺตปจฺจเยน ปจฺจโย.\csromanspacer{} \textbf{ปจฺฉาชาตา} พาหิรา ขนฺธา ปุเรชาตสฺส อิมสฺส อชฺฌตฺติกสฺส กายสฺส วิปฺปยุตฺตปจฺจเยน ปจฺจโย. (2)}

\prose{พาหิโร ธมฺโม อชฺฌตฺติกสฺส จ พาหิรสฺส จ ธมฺมสฺส วิปฺปยุตฺตปจฺจเยน ปจฺจโย.\csromanspacer{} \textbf{สหชาตํ}, ปุเรชาตํ, ปจฺฉาชาตํ. (สํขิตฺตํ.) (3)}

\proseitem{359}{อชฺฌตฺติโก จ พาหิโร จ ธมฺมา อชฺฌตฺติกสฺส ธมฺมสฺส วิปฺปยุตฺตปจฺจเยน ปจฺจโย.\csromanspacer{} \textbf{สหชาตํ}, ปจฺฉาชาตํ.\csromanspacer{}\csromanspacer{}\textbf{สหชาตา} ปฏิสนฺธิกฺขเณ อชฺฌตฺติกา จ พาหิรา จ ขนฺธา อชฺฌตฺติกานํ กฏตฺตารูปานํ วิปฺปยุตฺตปจฺจเยน ปจฺจโย.\csromanspacer{} \textbf{ปจฺฉาชาตํ}\csromandash{}ปจฺฉาชาตา ฯเปฯ. (สํขิตฺตํ.) (1)}

\prose{อชฺฌตฺติโก จ พาหิโร จ ธมฺมา พาหิรสฺส ธมฺมสฺส วิปฺปยุตฺตปจฺจเยน ปจฺจโย.\csromanspacer{} \textbf{สหชาตํ}, ปจฺฉาชาตํ. (สํขิตฺตํ.) (2)}

\prose{อชฺฌตฺติโก จ พาหิโร จ ธมฺมา อชฺฌตฺติกสฺส จ พาหิรสฺส จ ธมฺมสฺส วิปฺปยุตฺตปจฺจเยน ปจฺจโย.\csromanspacer{} \textbf{สหชาตํ}, ปจฺฉาชาตํ.\csromanspacer{}\csromanspacer{}\textbf{สหชาตา} ปฏิสนฺธิกฺขเณ อชฺฌตฺติกา จ พาหิรา จ ขนฺธา ฯเปฯ. (สํขิตฺตํ.) (3)}

\csromanheader{2}{\textbf{อตฺถิ-อาทิ}}
\proseitem{360}{อชฺฌตฺติโก ธมฺโม อชฺฌตฺติกสฺส ธมฺมสฺส อตฺถิปจฺจเยน ปจฺจโย.\csromanspacer{} \textbf{สหชาตํ}, ปุเรชาตํ, ปจฺฉาชาตํ.\csromanspacer{}\csromanspacer{}\textbf{สหชาตํ} ปฏิสนฺธิกฺขเณ จิตฺตํ อชฺฌตฺติกานํ กฏตฺตารูปานํ อตฺถิปจฺจเยน ปจฺจโย.\csromanspacer{} \textbf{ปุเรชาตํ} จกฺขุํ ฯเปฯ กายํ อนิจฺจโต ฯเปฯ. (ปุเรชาตสทิสํ นินฺนานากรณํ.) \textbf{ปจฺฉาชาตํ}\csromandash{}(ปจฺฉาชาตสทิสํ กาตพฺพํ.) (1)}

\prose{อชฺฌตฺติโก ธมฺโม พาหิรสฺส ธมฺมสฺส อตฺถิปจฺจเยน ปจฺจโย.\csromanspacer{} \textbf{สหชาตํ}, ปุเรชาตํ, ปจฺฉาชาตํ.\csromanspacer{}\csromanspacer{}\textbf{สหชาตํ}\csromandash{}สหชาตา อชฺฌตฺติกา ขนฺธา สมฺปยุตฺตกานํ ขนฺธานํ จิตฺตสมุฏฺฐานานญฺจ รูปานํ อตฺถิปจฺจเยน ปจฺจโย. (สํขิตฺตํ.) (2)}

\csromanheader{2}{66. อชฺฌตฺติกทุก}
\prose{\csromanfolio{513}(อิธ อตฺถิ สพฺพฏฺฐาเน สหชาตํ ปจฺจยวารสทิสํ, ปุเรชาตํ ปุเรชาตสทิสํ, ปจฺฉาชาตํ ปจฺฉาชาตสทิสํ กาตพฺพํ.\csromanspacer{} นินฺนานากรณํ.)}

\prose{อชฺฌตฺติโก ธมฺโม อชฺฌตฺติกสฺส จ พาหิรสฺส จ ธมฺมสฺส อตฺถิปจฺจเยน ปจฺจโย.\csromanspacer{} สหชาตํ, ปุเรชาตํ, ปจฺฉาชาตํ. (สํขิตฺตํ.) (3)}

\proseitem{361}{พาหิโร ธมฺโม พาหิรสฺส ธมฺมสฺส อตฺถิปจฺจเยน ปจฺจโย.\csromanspacer{} สหชาตํ, ปุเรชาตํ, ปจฺฉาชาตํ, อาหารํ, อินฺทฺริยํ. (สพฺพํ วิตฺถาเรตพฺพํ.) (1)}

\prose{พาหิโร ธมฺโม อชฺฌตฺติกสฺส ธมฺมสฺส อตฺถิปจฺจเยน ปจฺจโย.\csromanspacer{} สหชาตํ, ปุเรชาตํ, ปจฺฉาชาตํ, อาหารํ, อินฺทฺริยํ. (สํขิตฺตํ.) (2)}

\prose{พาหิโร ธมฺโม อชฺฌตฺติกสฺส จ พาหิรสฺส จ ธมฺมสฺส อตฺถิปจฺจเยน ปจฺจโย.\csromanspacer{} สหชาตํ, ปุเรชาตํ, ปจฺฉาชาตํ, อาหารํ, อินฺทฺริยํ. (สํขิตฺตํ.) (3)}

\proseitem{362}{อชฺฌตฺติโก จ พาหิโร จ ธมฺมา อชฺฌตฺติกสฺส ธมฺมสฺส อตฺถิปจฺจเยน ปจฺจโย.\csromanspacer{} สหชาตํ, ปุเรชาตํ, ปจฺฉาชาตํ, อาหารํ, อินฺทฺริยํ.\csromanspacer{}\csromanspacer{}\textbf{สหชาตา} จกฺขุวิญฺญาณสหคตา ขนฺธา จ จกฺขายตนญฺจ จกฺขุวิญฺญาณสฺส อตฺถิปจฺจเยน ปจฺจโย ฯเปฯ.\csromanspacer{} กายวิญฺญาณสหคตา ขนฺธา จ ฯเปฯ.\csromanspacer{} ปฏิสนฺธิกฺขเณ อชฺฌตฺติกา จ พาหิรา จ ขนฺธา อชฺฌตฺติกานํ กฏตฺตารูปานํ อตฺถิปจฺจเยน ปจฺจโย.\csromanspacer{} \textbf{ปุเรชาตํ} จกฺขายตนญฺจ วตฺถุ จ จิตฺตสฺส อตฺถิปจฺจเยน ปจฺจโย ฯเปฯ.\csromanspacer{} กายายตนญฺจ วตฺถุ จ จิตฺตสฺส อตฺถิปจฺจเยน ปจฺจโย.\csromanspacer{} รูปายตนญฺจ จกฺขายตนญฺจ จกฺขุวิญฺญาณสฺส ฯเปฯ.\csromanspacer{} โผฏฺฐพฺพายตนญฺจ กายายตนญฺจ กายวิญฺญาณสฺส อตฺถิปจฺจเยน ปจฺจโย.\csromanspacer{} \textbf{ปจฺฉาชาตา} อชฺฌตฺติกา จ พาหิรา จ ขนฺธา ปุเรชาตสฺส อิมสฺส อชฺฌตฺติกสฺส กายสฺส อตฺถิปจฺจเยน ปจฺจโย.\csromanspacer{} \textbf{ปจฺฉาชาตา} อชฺฌตฺติกา จ พาหิรา จ ขนฺธา กพฬีกาโร อาหาโร จ อิมสฺส อชฺฌตฺติกสฺส กายสฺส อตฺถิปจฺจเยน ปจฺจโย.\csromanspacer{} \textbf{ปจฺฉาชาตา} อชฺฌตฺติกา จ พาหิรา จ ขนฺธา รูปชีวิตินฺทฺริยญฺจ อชฺฌตฺติกานํ กฏตฺตารูปานํ อตฺถิปจฺจเยน ปจฺจโย. (1)}

\prose{อชฺฌตฺติโก จ พาหิโร จ ธมฺมา พาหิรสฺส ธมฺมสฺส อตฺถิปจฺจเยน ปจฺจโย.\csromanspacer{} สหชาตํ, ปุเรชาตํ, อาหารํ, อินฺทฺริยํ.\csromanspacer{}\csromanspacer{} \csromanfolio{514}\textbf{สหชาโต} จกฺขุวิญฺญาณสหคโต เอโก ขนฺโธ จ จกฺขายตนญฺจ จกฺขุวิญฺญาณญฺจ ทฺวินฺนํ ขนฺธานํ อตฺถิปจฺจเยน ปจฺจโย ฯเปฯ. (สหชาตปจฺจยวารสทิสํ นินฺนานากรณํ, ปฐมคมนสทิสํเยว.\csromanspacer{} สพฺเพ ปทา ปฐมฆฏนานเยน วิภชิตพฺพา.) (2)}

\prose{อชฺฌตฺติโก จ พาหิโร จ ธมฺมา อชฺฌตฺติกสฺส จ พาหิรสฺส จ ธมฺมสฺส อตฺถิปจฺจเยน ปจฺจโย.\csromanspacer{} สหชาตํ, ปุเรชาตํ, ปจฺฉาชาตํ, อาหารํ, อินฺทฺริยํ.\csromanspacer{}\csromanspacer{}\textbf{สหชาโต} จกฺขุวิญฺญาณสหคโต เอโก ขนฺโธ จ จกฺขายตนญฺจ ทฺวินฺนํ ขนฺธานํ จกฺขุวิญฺญาณสฺส จ อตฺถิปจฺจเยน ปจฺจโย. (สํขิตฺตํ.\csromanspacer{} สพฺเพ ปทา วิภชิตพฺพา, ปฐมฆฏนานเยน.) (3)}

\prose{นตฺถิปจฺจเยน ปจฺจโย.\csromanspacer{} วิคตปจฺจเยน ปจฺจโย.\csromanspacer{} อวิคตปจฺจเยน ปจฺจโย.}

\csromanheader{2}{1. ปจฺจยานุโลม 2. สงฺขฺยาวาร}
\csromanheader{2}{\textbf{สุทฺธ}}
\proseitem{363}{เหตุยา ตีณิ, อารมฺมเณ นว, อธิปติยา นว, อนนฺตเร นว, สมนนฺตเร นว, สหชาเต นว, อญฺญมญฺเญ ปญฺจ, นิสฺสเย นว, อุปนิสฺสเย นว, ปุเรชาเต นว, ปจฺฉาชาเต นว, อาเสวเน นว, กมฺเม ตีณิ, วิปาเก นว, อาหาเร นว, อินฺทฺริเย นว, ฌาเน ตีณิ, มคฺเค ตีณิ, สมฺปยุตฺเต ปญฺจ, วิปฺปยุตฺเต นว, อตฺถิยา นว, นตฺถิยา นว, วิคเต นว, อวิคเต นว.}

\vspace{\tipitakaparskip}
\csromancenter{อนุโลมํ.}
\csromanheader{2}{2. ปจฺจนียุทฺธาร}
\proseitem{364}{อชฺฌตฺติโก ธมฺโม อชฺฌตฺติกสฺส ธมฺมสฺส อารมฺมณปจฺจเยน ปจฺจโย.\csromanspacer{} สหชาตปจฺจเยน ปจฺจโย.\csromanspacer{} อุปนิสฺสยปจฺจเยน ปจฺจโย.\csromanspacer{} ปุเรชาตปจฺจเยน ปจฺจโย.\csromanspacer{} ปจฺฉาชาตปจฺจเยน ปจฺจโย. (1)}

\prose{อชฺฌตฺติโก ธมฺโม พาหิรสฺส ธมฺมสฺส อารมฺมณปจฺจเยน ปจฺจโย.\csromanspacer{} สหชาตปจฺจเยน ปจฺจโย.\csromanspacer{} อุปนิสฺสยปจฺจเยน ปจฺจโย.\csromanspacer{} ปุเรชาตปจฺจเยน ปจฺจโย.\csromanspacer{} ปจฺฉาชาตปจฺจเยน ปจฺจโย. (2)}

\csromanheader{2}{66. อชฺฌตฺติกทุก}
\prose{\csromanfolio{515}อชฺฌตฺติโก ธมฺโม อชฺฌตฺติกสฺส จ พาหิรสฺส จ ธมฺมสฺส อารมฺมณปจฺจเยน ปจฺจโย.\csromanspacer{} สหชาตปจฺจเยน ปจฺจโย.\csromanspacer{} อุปนิสฺสยปจฺจเยน ปจฺจโย.\csromanspacer{} ปุเรชาตปจฺจเยน ปจฺจโย.\csromanspacer{} ปจฺฉาชาตปจฺจเยน ปจฺจโย. (3)}

\proseitem{365}{พาหิโร ธมฺโม พาหิรสฺส ธมฺมสฺส อารมฺมณปจฺจเยน ปจฺจโย.\csromanspacer{} สหชาตปจฺจเยน ปจฺจโย.\csromanspacer{} อุปนิสฺสยปจฺจเยน ปจฺจโย.\csromanspacer{} ปุเรชาตปจฺจเยน ปจฺจโย.\csromanspacer{} ปจฺฉาชาตปจฺจเยน ปจฺจโย.\csromanspacer{} กมฺมปจฺจเยน ปจฺจโย.\csromanspacer{} อาหารปจฺจเยน ปจฺจโย.\csromanspacer{} อินฺทฺริยปจฺจเยน ปจฺจโย. (1)}

\prose{พาหิโร ธมฺโม อชฺฌตฺติกสฺส ธมฺมสฺส อารมฺมณปจฺจเยน ปจฺจโย.\csromanspacer{} สหชาตปจฺจเยน ปจฺจโย.\csromanspacer{} อุปนิสฺสยปจฺจเยน ปจฺจโย.\csromanspacer{} ปุเรชาตปจฺจเยน ปจฺจโย.\csromanspacer{} ปจฺฉาชาตปจฺจเยน ปจฺจโย.\csromanspacer{} กมฺมปจฺจเยน ปจฺจโย.\csromanspacer{} อาหารปจฺจเยน ปจฺจโย.\csromanspacer{} อินฺทฺริยปจฺจเยน ปจฺจโย. (2)}

\prose{พาหิโร ธมฺโม อชฺฌตฺติกสฺส จ พาหิรสฺส จ ธมฺมสฺส อารมฺมณปจฺจเยน ปจฺจโย.\csromanspacer{} สหชาตปจฺจเยน ปจฺจโย.\csromanspacer{} อุปนิสฺสยปจฺจเยน ปจฺจโย.\csromanspacer{} ปุเรชาตปจฺจเยน ปจฺจโย.\csromanspacer{} ปจฺฉาชาตปจฺจเยน ปจฺจโย.\csromanspacer{} กมฺมปจฺจเยน ปจฺจโย. (3)}

\proseitem{366}{อชฺฌตฺติโก จ พาหิโร จ ธมฺมา อชฺฌตฺติกสฺส ธมฺมสฺส อารมฺมณปจฺจเยน ปจฺจโย.\csromanspacer{} สหชาตปจฺจเยน ปจฺจโย.\csromanspacer{} อุปนิสฺสยปจฺจเยน ปจฺจโย.\csromanspacer{} ปุเรชาตปจฺจเยน ปจฺจโย.\csromanspacer{} ปจฺฉาชาตปจฺจเยน ปจฺจโย. (1)}

\prose{อชฺฌตฺติโก จ พาหิโร จ ธมฺมา พาหิรสฺส ธมฺมสฺส อารมฺมณปจฺจเยน ปจฺจโย.\csromanspacer{} สหชาตปจฺจเยน ปจฺจโย.\csromanspacer{} อุปนิสฺสยปจฺจเยน ปจฺจโย.\csromanspacer{} ปุเรชาตปจฺจเยน ปจฺจโย.\csromanspacer{} ปจฺฉาชาตปจฺจเยน ปจฺจโย. (2)}

\prose{อชฺฌตฺติโก จ พาหิโร จ ธมฺมา อชฺฌตฺติตสฺส จ พาหิรสฺส จ ธมฺมสฺส อารมฺมณปจฺจเยน ปจฺจโย.\csromanspacer{} สหชาตปจฺจเยน ปจฺจโย.\csromanspacer{} อุปนิสฺสยปจฺจเยน ปจฺจโย.\csromanspacer{} ปุเรชาตปจฺจเยน ปจฺจโย.\csromanspacer{} ปจฺฉาชาตปจฺจเยน ปจฺจโย. (3)}

\csromanheader{2}{2. ปจฺจยปจฺจนีย 2. สงฺขฺยาวาร}
\vspace{\tipitakaparskip}
\csromancenter{นเหตุยา นว, น-อารมฺมเณ นว, (สพฺพตฺถ นว,) โน-อวิคเต นว.}
\csromanheader{2}{3. ปจฺจยานุโลมปจฺจนีย}
\proseitem{368}{\csromanfolio{516}\textbf{เหตุ}ปจฺจยา น-อารมฺมเณ ตีณิ, น-อธิปติยา ตีณิ, น-อนนฺตเร ตีณิ, นสมนนฺตเร ตีณิ, น-อญฺญมญฺเญ ตีณิ, น-อุปนิสฺสเย ตีณิ, (สพฺพตฺถ ตีณิ,) นสมฺปยุตฺเต ตีณิ, นวิปฺปยุตฺเต ตีณิ, โนนตฺถิยา ตีณิ, โนวิคเต ตีณิ.}

\csromanheader{2}{4. ปจฺจยปจฺจนียานุโลม}
\proseitem{369}{นเหตุปจฺจยา อารมฺมเณ นว, อธิปติยา นว, (อนุโลมมาติกา กาตพฺพา.) ฯเปฯ อวิคเต นว.}

\nitthitam{อชฺฌตฺติกทุกํ นิฏฺฐิตํ.}
\csromanmatikaanchor{mtk.66}
\csromantocmarkref{subsection}{67. อุปาทาทุก}{mtk.66}
\bookmark[dest={mtk.66},level=2]{67. อุปาทาทุก}
\csromanheader{2}{67. อุปาทาทุก 1. ปฏิจฺจวาร}
\csromanheader{2}{1. ปจฺจยานุโลม 1. วิภงฺควาร}
\csromanheader{2}{\textbf{เหตุ}}
\proseitem{370}{อุปาทา ธมฺมํ ปฏิจฺจ โน-อุปาทา ธมฺโม อุปฺปชฺชติ เหตุปจฺจยา.\csromanspacer{} ปฏิสนฺธิกฺขเณ วตฺถุํ ปฏิจฺจ โน-อุปาทา ขนฺธา. (1)}

\prose{โน-อุปาทา ธมฺมํ ปฏิจฺจ โน-อุปาทา ธมฺโม อุปฺปชฺชติ เหตุปจฺจยา.\csromanspacer{} โน-อุปาทา เอกํ ขนฺธํ ปฏิจฺจ ตโย ขนฺธา โน-อุปาทา จ จิตฺตสมุฏฺฐานํ รูปํ ฯเปฯ.\csromanspacer{} เทฺว ขนฺเธ ฯเปฯ.\csromanspacer{} ปฏิสนฺธิกฺขเณ โน-อุปาทา เอกํ ขนฺธํ ปฏิจฺจ ตโย ขนฺธา โน-อุปาทา จ กฏตฺตารูปํ ฯเปฯ เทฺว ขนฺเธ ฯเปฯ.\csromanspacer{} เอกํ มหาภูตํ ฯเปฯ เทฺว มหาภูเต ปฏิจฺจ เทฺว มหาภูตา. (1)}

\prose{โน-อุปาทา ธมฺมํ ปฏิจฺจ อุปาทา ธมฺโม อุปฺปชฺชติ เหตุปจฺจยา.\csromanspacer{} โน-อุปาทา ขนฺเธ ปฏิจฺจ อุปาทา จิตฺตสมุฏฺฐานํ รูปํ.\csromanspacer{} ปฏิสนฺธิกฺขเณ ฯเปฯ มหภูเต ปฏิจฺจ อุปาทา จิตฺตสมุฏฺฐานํ รูปํ, กฏตฺตารูปํ, อุปาทารูปํ. (2)}

\prose{โน-อุปาทา ธมฺมํ ปฏิจฺจ อุปาทา จ โน-อุปาทา จ ธมฺมา อุปฺปชฺชนฺติ เหตุปจฺจยา.\csromanspacer{} โน-อุปาทา เอกํ ขนฺธํ ปฏิจฺจ ตโย ขนฺธา อุปาทา จ โน-อุปาทา จ จิตฺตสมุฏฺฐานํ รูปํ ฯเปฯ เทฺว ขนฺเธ ฯเปฯ.\csromanspacer{} ปฏิสนฺธิกฺขเณ ฯเปฯ. (3)}

\csromanheader{2}{67. อุปาทาทุก}
\prose{\csromanfolio{517}อุปาทา จ โน-อุปาทา จ ธมฺมํ ปฏิจฺจ โน-อุปาทา ธมฺโม อุปฺปชฺชติ เหตุปจฺจยา.\csromanspacer{} ปฏิสนฺธิกฺขเณ โน-อุปาทา เอกํ ขนฺธญฺจ วตฺถุญฺจ ปฏิจฺจ ตโย ขนฺธา.\csromanspacer{} เทฺว ขนฺเธ จ ฯเปฯ. (1)}

\csromanheader{2}{\textbf{อารมฺมณ}}
\proseitem{371}{อุปาทา ธมฺมํ ปฏิจฺจ โน-อุปาทา ธมฺโม อุปฺปชฺชติ อารมฺมณปจฺจยา.\csromanspacer{} ปฏิสนฺธิกฺขเณ วตฺถุํ ปฏิจฺจ โน-อุปาทา ขนฺธา. (1)}

\prose{โน-อุปาทา ธมฺมํ ปฏิจฺจ โน-อุปาทา ธมฺโม อุปฺปชฺชติ อารมฺมณปจฺจยา.\csromanspacer{} โน-อุปาทา เอกํ ขนฺธํ ปฏิจฺจ ตโย ขนฺธา ฯเปฯ เทฺว ขนฺเธ ฯเปฯ.\csromanspacer{} ปฏิสนฺธิกฺขเณ ฯเปฯ. (1)}

\prose{อุปาทา จ โน-อุปาทา จ ธมฺมํ ปฏิจฺจ โน-อุปาทา ธมฺโม อุปฺปชฺชติ อารมฺมณปจฺจยา.\csromanspacer{} ปฏิสนฺธิกฺขเณ โน-อุปาทา เอกํ ขนฺธญฺจ วตฺถุญฺจ ปฏิจฺจ ตโย ขนฺธา ฯเปฯ เทฺว ขนฺเธ จ ฯเปฯ. (1)}

\csromanheader{2}{\textbf{อธิปตฺยาทิ}}
\proseitem{372}{โน-อุปาทา ธมฺมํ ปฏิจฺจ โน-อุปาทา ธมฺโม อุปฺปชฺชติ อธิปติปจฺจยา.\csromanspacer{} โน-อุปาทา เอกํ ขนฺธํ ปฏิจฺจ ตโย ขนฺธา โน-อุปาทา จ จิตฺตสมุฏฺฐานํ รูปํ ฯเปฯ เทฺว ขนฺเธ ฯเปฯ. (1)}

\prose{โน-อุปาทา ธมฺมํ ปฏิจฺจ อุปาทา ธมฺโม อุปฺปชฺชติ อธิปติปจฺจยา.\csromanspacer{} โน-อุปาทา ขนฺเธ ปฏิจฺจ อุปาทา จิตฺตสมุฏฺฐานํ รูปํ.\csromanspacer{} มหาภูเต ปฏิจฺจ อุปาทา จิตฺตสมุฏฺฐานํ รูปํ. (2)}

\prose{โน-อุปาทา ธมฺมํ ปฏิจฺจ อุปาทา จ โน-อุปาทา จ ธมฺมา อุปฺปชฺชนฺติ อธิปติปจฺจยา.\csromanspacer{} โน-อุปาทา เอกํ ขนฺธํ ปฏิจฺจ ตโย ขนฺธา อุปาทา จ โน-อุปาทา จ จิตฺตสมุฏฺฐานํ รูปํ ฯเปฯ เทฺว ขนฺเธ ฯเปฯ. (3)}

\prose{อนนฺตรปจฺจยา ตีณิ.\csromanspacer{} สมนนฺตรปจฺจยา ตีณิ.\csromanspacer{} สหชาตปจฺจยา ปญฺจ.}

\csromanheader{2}{\textbf{อญฺญมญฺญ}}
\proseitem{373}{อุปาทา ธมฺมํ ปฏิจฺจ โน-อุปาทา ธมฺโม อุปฺปชฺชติ อญฺญมญฺญปจฺจยา.\csromanspacer{} ปฏิสนฺธิกฺขเณ วตฺถุํ ปฏิจฺจ โน-อุปาทา ขนฺธา. (1)}

\prose{โน-อุปาทา ธมฺมํ ปฏิจฺจ โน-อุปาทา ธมฺโม อุปฺปชฺชติ อญฺญมญฺญปจฺจยา.\csromanspacer{} โน-อุปาทา เอกํ ขนฺธํ ปฏิจฺจ ตโย ขนฺธา ฯเปฯ เทฺว ขนฺเธ ฯเปฯ.\csromanspacer{} ปฏิสนฺธิกฺขเณ ฯเปฯ. \csromanfolio{518}เอกํ มหาภูตํ ฯเปฯ.\csromanspacer{} อสญฺญสตฺตานํ เอกํ มหาภูตํ ปฏิจฺจ ฯเปฯ เทฺว มหาภูเต ปฏิจฺจ เทฺว มหาภูตา. (1)}

\prose{โน-อุปาทา ธมฺมํ ปฏิจฺจ อุปาทา ธมฺโม อุปฺปชฺชติ อญฺญมญฺญปจฺจยา.\csromanspacer{} ปฏิสนฺธิกฺขเณ โน-อุปาทา ขนฺเธ ปฏิจฺจ วตฺถุ. (2)}

\prose{โน-อุปาทา ธมฺมํ ปฏิจฺจ อุปาทา จ โน-อุปาทา จ ธมฺมา อุปฺปชฺชนฺติ อญฺญมญฺญปจฺจยา.\csromanspacer{} ปฏิสนฺธิกฺขเณ โน-อุปาทา เอกํ ขนฺธํ ปฏิจฺจ ตโย ขนฺธา วตฺถุ จ ฯเปฯ เทฺว ขนฺเธ ฯเปฯ. (3)}

\prose{อุปาทา จ โน-อุปาทา จ ธมฺมํ ปฏิจฺจ โน-อุปาทา ธมฺโม อุปฺปชฺชติ อญฺญมญฺญปจฺจยา.\csromanspacer{} ปฏิสนฺธิกฺขเณ โน-อุปาทา เอกํ ขนฺธญฺจ วตฺถุญฺจ ปฏิจฺจ ตโย ขนฺธา ฯเปฯ เทฺว ขนฺเธ จ ฯเปฯ. (สํขิตฺตํ.) (1)}

\csromanheader{2}{1. ปจฺจยานุโลม 2. สงฺขฺยาวาร}
\csromanheader{2}{\textbf{สุทฺธ}}
\proseitem{374}{เหตุยา ปญฺจ, อารมฺมเณ ตีณิ, อธิปติยา ตีณิ, อนนฺตเร ตีณิ, สมนนฺตเร ตีณิ, สหชาเต ปญฺจ, อญฺญมญฺเญ ปญฺจ, นิสฺสเย ปญฺจ, อุปนิสฺสเย ตีณิ, ปุเรชาเต เอกํ, อาเสวเน เอกํ, กมฺเม ปญฺจ, วิปาเก ปญฺจ, (สพฺพตฺถ ปญฺจ,) สมฺปยุตฺเต ตีณิ, วิปฺปยุตฺเต ปญฺจ, อตฺถิยา ปญฺจ, นตฺถิยา ตีณิ, วิคเต ตีณิ, อวิคเต ปญฺจ.}

\csromanheader{2}{2. ปจฺจยปจฺจนีย 1. วิภงฺควาร}
\csromanheader{2}{\textbf{นเหตุ}}
\proseitem{375}{อุปาทา ธมฺมํ ปฏิจฺจ โน-อุปาทา ธมฺโม อุปฺปชฺชติ นเหตุปจฺจยา.\csromanspacer{} อเหตุกปฏิสนฺธิกฺขเณ วตฺถุํ ปฏิจฺจ โน-อุปาทา ขนฺธา. (1)}

\prose{โน-อุปาทา ธมฺมํ ปฏิจฺจ โน-อุปาทา ธมฺโม อุปฺปชฺชติ นเหตุปจฺจยา.\csromanspacer{} อเหตุกํ โน-อุปาทา เอกํ ขนฺธํ ปฏิจฺจ ตโย ขนฺธา โน-อุปาทา จ จิตฺตสมุฏฺฐานํ รูปํ ฯเปฯ เทฺว ขนฺเธ ฯเปฯ.\csromanspacer{} อเหตุกปฏิสนฺธิกฺขเณ ฯเปฯ.\csromanspacer{} เอกํ มหาภูตํ ฯเปฯ.\csromanspacer{} อสญฺญสตฺตานํ เอกํ มหาภูตํ ปฏิจฺจ ตโย มหาภูตา ฯเปฯ เทฺว มหาภูเต ปฏิจฺจ เทฺว มหาภูตา.\csromanspacer{} วิจิกิจฺฉาสหคเต}

\vspace{\tipitakaparskip}
\csromancenter{อุปาทาทุก อุทฺธจฺจสหคเต ขนฺเธ ปฏิจฺจ วิจิกิจฺฉาสหคโต อุทฺธจฺจสหคโต โมโห. (1)}
\prose{\csromanfolio{519}โน-อุปาทา ธมฺมํ ปฏิจฺจ อุปาทา ธมฺโม อุปฺปชฺชติ นเหตุปจฺจยา.\csromanspacer{} อเหตุเก โน-อุปาทา ขนฺเธ ปฏิจฺจ อุปาทา จิตฺตสมุฏฺฐานํ รูปํ.\csromanspacer{} อเหตุกปฏิสนฺธิกฺขเณ ฯเปฯ มหาภูเต ปฏิจฺจ อุปาทา จิตฺตสมุฏฺฐานํ รูปํ, กฏตฺตารูปํ อุปาทารูปํ ฯเปฯ. (ยาว อสญฺญสตฺตา.) (2)}

\prose{โน-อุปาทา ธมฺมํ ปฏิจฺจ อุปาทา จ โน-อุปาทา จ ธมฺมา อุปฺปชฺชนฺติ นเหตุปจฺจยา.\csromanspacer{} อเหตุกํ โน-อุปาทา เอกํ ขนฺธํ ปฏิจฺจ ตโย ขนฺธา อุปาทา จ โน-อุปาทา จ จิตฺตสมุฏฺฐานํ รูปํ ฯเปฯ เทฺว ขนฺเธ ฯเปฯ.\csromanspacer{} อเหตุกปฏิสนฺธิกฺขเณ ฯเปฯ. (3)}

\prose{อุปาทา จ โน-อุปาทา จ ธมฺมํ ปฏิจฺจ โน-อุปาทา ธมฺโม อุปฺปชฺชติ นเหตุปจฺจยา.\csromanspacer{} อเหตุกปฏิสนฺธิกฺขเณ โน-อุปาทา เอกํ ขนฺธญฺจ วตฺถุญฺจ ปฏิจฺจ ตโย ขนฺธา ฯเปฯ เทฺว ขนฺเธ จ ฯเปฯ. (1)}

\csromanheader{2}{\textbf{น-อารมฺมณาทิ}}
\proseitem{376}{โน-อุปาทา ธมฺมํ ปฏิจฺจ โน-อุปาทา ธมฺโม อุปฺปชฺชติ น-อารมฺมณปจฺจยา.\csromanspacer{} โน-อุปาทา ขนฺเธ ปฏิจฺจ โน-อุปาทา จิตฺตสมุฏฺฐานํ รูปํ.\csromanspacer{} ปฏิสนฺธิกฺขเณ ฯเปฯ.\csromanspacer{} เอกํ มหาภูตํ ฯเปฯ. (ยาว อสญฺญสตฺตา.) เทฺว มหาภูเต ปฏิจฺจ เทฺว มหาภูตา. (1)}

\prose{โน-อุปาทา ธมฺมํ ปฏิจฺจ อุปาทา ธมฺโม อุปฺปชฺชติ น-อารมฺมณปจฺจยา.\csromanspacer{} โน-อุปาทา ขนฺเธ ปฏิจฺจ อุปาทา จิตฺตสมุฏฺฐานํ รูปํ.\csromanspacer{} ปฏิสนฺธิกฺขเณ ฯเปฯ มหาภูเต ปฏิจฺจ อุปาทา จิตฺตสมุฏฺฐานํ รูปํ, กฏตฺตารูปํ, อุปาทารูปํ. (ยาว อสญฺญสตฺตา.) (2)}

\prose{โน-อุปาทา ธมฺมํ ปฏิจฺจ อุปาทา จ โน-อุปาทา จ ธมฺมา อุปฺปชฺชนฺติ น-อารมฺมณปจฺจยา.\csromanspacer{} โน-อุปาทา ขนฺเธ ปฏิจฺจ อุปาทา จ โน-อุปาทา จ จิตฺตสมุฏฺฐานํ รูปํ.\csromanspacer{} ปฏิสนฺธิกฺขเณ ฯเปฯ. (3)}

\prose{น-อธิปติปจฺจยา ปญฺจ.\csromanspacer{} น-อนนฺตรปจฺจยา ตีณิ ฯเปฯ.\csromanspacer{} น-อุปนิสฺสยปจฺจยา ปญฺจ.\csromanspacer{} นปุเรชาตปจฺจยา ปญฺจ.\csromanspacer{} นปจฺฉาชาตปจฺจยา ปญฺจ.\csromanspacer{} น-อาเสวนปจฺจยา ปญฺจ.}

\csromanheader{2}{\textbf{นกมฺม}}
\proseitem{377}{\csromanfolio{520}โน-อุปาทา ธมฺมํ ปฏิจฺจ โน-อุปาทา ธมฺโม อุปฺปชฺชติ นกมฺมปจฺจยา.\csromanspacer{} โน-อุปาทา ขนฺเธ ปฏิจฺจ สมฺปยุตฺตกา เจตนา.\csromanspacer{} พาหิรํ อาหารสมุฏฺฐานํ อุตุสมุฏฺฐานํ ฯเปฯ เทฺว มหาภูเต ปฏิจฺจ เทฺว มหาภูตา. (1)}

\prose{โน-อุปาทา ธมฺมํ ปฏิจฺจ อุปาทา ธมฺโม อุปฺปชฺชติ นกมฺมปจฺจยา.\csromanspacer{} พาหิเร อาหารสมุฏฺฐาเน อุตุสมุฏฺฐาเน มหาภูเต ปฏิจฺจ อุปาทารูปํ. (2)}

\prose{โน-อุปาทา ธมฺมํ ปฏิจฺจ อุปาทา จ โน-อุปาทา จ ธมฺมา อุปฺปชฺชนฺติ นกมฺมปจฺจยา.\csromanspacer{} พาหิรํ อาหารสมุฏฺฐานํ อุตุสมุฏฺฐานํ เอกํ มหาภูตํ ปฏิจฺจ ตโย มหาภูตา อุปาทา จ รูปํ ฯเปฯ เทฺว มหาภูเต ฯเปฯ. (3)}

\csromanheader{2}{\textbf{นวิปาก}}
\proseitem{378}{โน-อุปาทา ธมฺมํ ปฏิจฺจ โน-อุปาทา ธมฺโม อุปฺปชฺชติ นวิปากปจฺจยา.\csromanspacer{} โน-อุปาทา เอกํ ขนฺธํ ปฏิจฺจ ตโย ขนฺธา โน-อุปาทา จ จิตฺตสมุฏฺฐานํ รูปํ ฯเปฯ เทฺว ขนฺเธ ฯเปฯ.\csromanspacer{} เอกํ มหาภูตํ ฯเปฯ. (ยาว อสญฺญสตฺตา, เอวํ ตีณิ, โน-อุปาทามูลเก.)}

\csromanheader{2}{\textbf{น-อาหาร}}
\proseitem{379}{โน-อุปาทา ธมฺมํ ปฏิจฺจ โน-อุปาทา ธมฺโม อุปฺปชฺชติ น-อาหารปจฺจยา.\csromanspacer{} พาหิรํ อุตุสมุฏฺฐานํ อสญฺญสตฺตานํ เอกํ มหาภูตํ ฯเปฯ เทฺว มหาภูเต ปฏิจฺจ มหาภูตา. (1)}

\prose{โน-อุปาทา ธมฺมํ ปฏิจฺจ อุปาทา ธมฺโม อุปฺปชฺชติ น-อาหารปจฺจยา.\csromanspacer{} พาหิรํ อุตุสมุฏฺฐานํ อสญฺญสตฺตานํ มหาภูเต ปฏิจฺจ อุปาทารูปํ.}

\prose{โน-อุปาทา ธมฺมํ ปฏิจฺจ อุปาทา จ โน-อุปาทา จ ธมฺมา อุปฺปชฺชนฺติ น-อาหารปจฺจยา.\csromanspacer{} พาหิรํ อุตุสมุฏฺฐานํ อสญฺญสตฺตานํ เอกํ มหาภูตํ ปฏิจฺจ ตโย มหาภูตา อุปาทา จ รูปํ ฯเปฯ เทฺว มหาภูเต ฯเปฯ. (3)}

\csromanheader{2}{\textbf{น-อินฺทฺริย}}
\proseitem{380}{โน-อุปาทา ธมฺมํ ปฏิจฺจ โน-อุปาทา ธมฺโม อุปฺปชฺชติ น-อินฺทฺริยปจฺจยา.\csromanspacer{} พาหิรํ อาหารสมุฏฺฐานํ อุตุสมุฏฺฐานํ เอกํ มหาภูตํ ฯเปฯ. (1)}

\csromanheader{2}{67. อุปาทาทุก}
\prose{\csromanfolio{521}โน-อุปาทา ธมฺมํ ปฏิจฺจ อุปาทา ธมฺโม อุปฺปชฺชติ น-อินฺทฺริยปจฺจยา.\csromanspacer{} พาหิเร อาหารสมุฏฺฐาเน อุตุสมุฏฺฐาเน มหาภูเต ปฏิจฺจ อุปาทารูปํ.\csromanspacer{} อสญฺญสตฺตานํ มหาภูเต ปฏิจฺจ รูปชีวิตินฺทฺริยํ.}

\prose{โน-อุปาทา ธมฺมํ ปฏิจฺจ อุปาทา จ โน-อุปาทา จ ธมฺมา อุปฺปชฺชนฺติ น-อินฺทฺริยปจฺจยา.\csromanspacer{} พาหิรํ อาหารสมุฏฺฐานํ อุตุสมุฏฺฐานํ เอกํ มหาภูตํ ฯเปฯ. (3)}

\csromanheader{2}{\textbf{นฌานาทิ}}
\proseitem{381}{โน-อุปาทา ธมฺมํ ปฏิจฺจ โน-อุปาทา ธมฺโม อุปฺปชฺชติ นฌานปจฺจยา.\csromanspacer{} ปญฺจวิญฺญาณสหคตํ เอกํ ขนฺธํ ปฏิจฺจ ตโย ขนฺธา ฯเปฯ เทฺว ขนฺเธ ฯเปฯ.\csromanspacer{} พาหิรํ อาหารสมุฏฺฐานํ อุตุสมุฏฺฐานํ อสญฺญสตฺตานํ ฯเปฯ เทฺว มหาภูเต ปฏิจฺจ เทฺว มหาภูตา. (1)}

\prose{โน-อุปาทา ธมฺมํ ปฏิจฺจ อุปาทา ธมฺโม อุปฺปชฺชติ นฌานปจฺจยา.\csromanspacer{} พาหิเร อาหารสมุฏฺฐาเน อุตุสมุฏฺฐาเน อสญฺญสตฺตานํ มหาภูเต ปฏิจฺจ อุปาทา กฏตฺตารูปํ. (2)}

\prose{โน-อุปาทา ธมฺมํ ปฏิจฺจ อุปาทา จ โน-อุปาทา จ ธมฺมา อุปฺปชฺชนฺติ นฌานปจฺจยา.\csromanspacer{} พาหิรํ อาหารสมุฏฺฐานํ อุตุสมุฏฺฐานํ อสญฺญสตฺตานํ เอกํ มหาภูตํ ปฏิจฺจ ตโย มหาภูตา ฯเปฯ เทฺว มหาภูเต ฯเปฯ มหาภูเต ปฏิจฺจ อุปาทา กฏตฺตารูปํ, อุปาทารูปํ. (3)}

\prose{อุปาทา ธมฺมํ ปฏิจฺจ โน-อุปาทา ธมฺโม อุปฺปชฺชติ นมคฺคปจฺจยา ปญฺจ.\csromanspacer{} นสมฺปยุตฺตปจฺจยา.\csromanspacer{} ตีณิ.}

\csromanheader{2}{\textbf{นวิปฺปยุตฺตาทิ}}
\proseitem{382}{โน-อุปาทา ธมฺมํ ปฏิจฺจ โน-อุปาทา ธมฺโม อุปฺปชฺชติ นวิปฺปยุตฺตปจฺจยา.\csromanspacer{} อรูเป โน-อุปาทา เอกํ ขนฺธํ ปฏิจฺจ ตโย ขนฺธา ฯเปฯ เทฺว ขนฺเธ ฯเปฯ.\csromanspacer{} พาหิรํ อาหารสมุฏฺฐานํ อุตุสมุฏฺฐานํ อสญฺญสตฺตานํ เอกํ มหาภูตํ ฯเปฯ. (1)}

\prose{โน-อุปาทา ธมฺมํ ปฏิจฺจ อุปาทา ธมฺโม อุปฺปชฺชติ นวิปฺปยุตฺตปจฺจยา.\csromanspacer{} พาหิเร อาหารสมุฏฺฐาเน อุตุสมุฏฺฐาเน อสญฺญสตฺตานํ มหาภูเต ปฏิจฺจ อุปาทา กฏตฺตารูปํ, อุปาทารูปํ. (2)}

\prose{\csromanfolio{522}โน-อุปาทา ธมฺมํ ปฏิจฺจ อุปาทา จ โน-อุปาทา จ ธมฺมา อุปฺปชฺชนฺติ นวิปฺปยุตฺตปจฺจยา.\csromanspacer{} พาหิรํ อาหารสมุฏฺฐานํ อุตุสมุฏฺฐานํ เอกํ มหาภูตํ ปฏิจฺจ ตโย มหาภูตา อุปาทา จ รูปํ ฯเปฯ เทฺว มหาภูเต ปฏิจฺจ เทฺว มหาภูตา อุปาทา จ รูปํ.\csromanspacer{} อสญฺญสตฺตานํ เอกํ มหาภูตํ ปฏิจฺจ ตโย มหาภูตา กฏตฺตา จ รูปํ, อุปาทารูปํ ฯเปฯ เทฺว มหาภูเต ปฏิจฺจ เทฺว มหาภูตา กฏตฺตา จ รูปํ อุปาทารูปํ.\csromanspacer{} โนนตฺถิปจฺจยา.\csromanspacer{} โนวิคตปจฺจยา. (3)}

\csromanheader{2}{2. ปจฺจยปจฺจนีย 2. สงฺขฺยาวาร}
\csromanheader{2}{\textbf{สุทฺธ}}
\proseitem{383}{นเหตุยา ปญฺจ, น-อารมฺมเณ ตีณิ, น-อธิปติยา ปญฺจ, น-อนนฺตเร ตีณิ, นสมนนฺตเร ตีณิ, น-อญฺญมญฺเญ ตีณิ, น-อุปนิสฺสเย ตีณิ, นปุเรชาเต ปญฺจ, นปจฺฉาชาเต ปญฺจ, น-อาเสวเน ปญฺจ, นกมฺเม ตีณิ, นวิปาเก ตีณิ, น-อาหาเร ตีณิ, น-อินฺทฺริเย ตีณิ, นฌาเน ตีณิ, นมคฺเค ปญฺจ, นสมฺปยุตฺเต ตีณิ, นวิปฺปยุตฺเต ตีณิ, โนนตฺถิยา ตีณิ, โนวิคเต ตีณิ.}

\csromanheader{2}{3. ปจฺจยานุโลมปจฺจนีย}
\proseitem{384}{เหตุปจฺจยา น-อารมฺมเณ ตีณิ, น-อธิปติยา ปญฺจ ฯเปฯ นกมฺเม เอกํ, นวิปาเก ตีณิ ฯเปฯ นสมฺปยุตฺเต ตีณิ, นวิปฺปยุตฺเต เอกํ, โนนตฺถิยา ตีณิ, โนวิคเต ตีณิ.}

\csromanheader{2}{4. ปจฺจยปจฺจนียานุโลม}
\proseitem{385}{นเหตุปจฺจยา อารมฺมเณ ตีณิ, อนนฺตเร ตีณิ, สมนนฺตเร ตีณิ, สหชาเต ปญฺจ, อญฺญมญฺเญ ปญฺจ, นิสฺสเย ปญฺจ, อุปนิสฺสเย ตีณิ, ปุเรชาเต เอกํ, อาเสวเน เอกํ ฯเปฯ มคฺเค เอกํ ฯเปฯ อวิคเต ปญฺจ.}

\vspace{\tipitakaparskip}
\csromancenter{(สหชาตวาโร ปฏิจฺจวารสทิโส.)}
\csromanheader{2}{67. อุปาทาทุก \textbf{67. อุปาทาทุก 3. ปจฺจยวาร}}
\csromanheader{2}{1. ปจฺจยานุโลม 1. วิภงฺควาร}
\csromanheader{2}{\textbf{เหตุ}}
\proseitem{386}{\csromanfolio{523}อุปาทา ธมฺมํ ปจฺจยา โน-อุปาทา ธมฺโม อุปฺปชฺชติ เหตุปจฺจยา.\csromanspacer{} วตฺถุํ ปจฺจยา โน-อุปาทา ขนฺธา.\csromanspacer{} ปฏิสนฺธิกฺขเณ ฯเปฯ. (1)}

\prose{โน-อุปาทา ธมฺมํ ปจฺจยา โน-อุปาทา ธมฺโม อุปฺปชฺชติ เหตุปจฺจยา.\csromanspacer{} ตีณิ. (โน-อุปาทามูลเก ตีณิปิ ปฏิจฺจสทิสา, นินฺนานากรณา.) (3)}

\prose{อุปาทา จ โน-อุปาทา จ ธมฺมํ ปจฺจยา โน-อุปาทา ธมฺโม อุปฺปชฺชติ เหตุปจฺจยา.\csromanspacer{} โน-อุปาทา เอกํ ขนฺธญฺจ วตฺถุญฺจ ปจฺจยา ตโย ขนฺธา ฯเปฯ เทฺว ขนฺเธ จ ฯเปฯ.\csromanspacer{} ปฏิสนฺธิกฺขเณ ฯเปฯ. (1)}

\csromanheader{2}{\textbf{อารมฺมณ}}
\proseitem{387}{อุปาทา ธมฺมํ ปจฺจยา โน-อุปาทา ธมฺโม อุปฺปชฺชติ อารมฺมณปจฺจยา.\csromanspacer{} จกฺขายตนํ ปจฺจยา จกฺขุวิญฺญาณํ ฯเปฯ.\csromanspacer{} กายายตนํ ปจฺจยา กายวิญฺญาณํ.\csromanspacer{} วตฺถุํ ปจฺจยา โน-อุปาทา ขนฺธา.\csromanspacer{} ปฏิสนฺธิกฺขเณ ฯเปฯ. (1)}

\prose{โน-อุปาทา ธมฺมํ ปจฺจยา โน-อุปาทา ธมฺโม อุปฺปชฺชติ อารมฺมณปจฺจยา.\csromanspacer{} เอกํ. (ปฏิจฺจสทิสํ.) (1)}

\prose{อุปาทา ธมฺมญฺจ โน-อุปาทา ธมฺมญฺจ ธมฺมํ ปจฺจยา โน-อุปาทา ธมฺโม อุปฺปชฺชติ อารมฺมณปจฺจยา.\csromanspacer{} จกฺขุวิญฺญาณสหคตํ เอกํ ขนฺธญฺจ จกฺขายตนญฺจ ปจฺจยา ตโย ขนฺธา ฯเปฯ เทฺว ขนฺเธ จ ฯเปฯ.\csromanspacer{} กายวิญฺญาณสหคตํ ฯเปฯ.\csromanspacer{} โน-อุปาทา เอกํ ขนฺธญฺจ วตฺถุญฺจ ปจฺจยา ตโย ขนฺธา ฯเปฯ เทฺว ขนฺเธ จ ฯเปฯ.\csromanspacer{} ปฏิสนฺธิกฺขเณ ฯเปฯ. (สํขิตฺตํ.) (1)}

\csromanheader{2}{1. ปจฺจยานุโลม 2. สงฺขฺยาวาร}
\csromanheader{2}{\textbf{สุทฺธ}}
\vspace{\tipitakaparskip}
\csromancenter{\textbf{เหตุ}ยา ปญฺจ, อารมฺมเณ ตีณิ, อธิปติยา ปญฺจ, อนนฺตเร ตีณิ, สมนนฺตเร ตีณิ, สหชาเต ปญฺจ, อญฺญมญฺเญ ปญฺจ, นิสฺสเย ปญฺจ, อุปนิสฺสเย ตีณิ, ปุเรชาเต ตีณิ, อาเสวเน ตีณิ, กมฺเม ปญฺจ ฯเปฯ อวิคเต ปญฺจ. (เอวํ คเณตพฺพํ.)}
\csromanheader{2}{2. ปจฺจยปจฺจนีย 1. วิภงฺควาร}
\csromanheader{2}{\textbf{นเหตุ}}
\proseitem{389}{\csromanfolio{524}อุปาทา ธมฺมํ ปจฺจยา โน-อุปาทา ธมฺโม อุปฺปชฺชติ นเหตุปจฺจยา.\csromanspacer{} จกฺขายตนํ ปจฺจยา จกฺขุวิญฺญาณํ ฯเปฯ.\csromanspacer{} กายายตนํ ปจฺจยา กายวิญฺญาณํ.\csromanspacer{} วตฺถุํ ปจฺจยา อเหตุกา โน-อุปาทา ขนฺธา.\csromanspacer{} อเหตุกปฏิสนฺธิกฺขเณ ฯเปฯ.\csromanspacer{} วตฺถุํ ปจฺจยา วิจิกิจฺฉาสหคโต อุทฺธจฺจสหคโต โมโห. (1)}

\prose{โน-อุปาทา ธมฺมํ ปจฺจยา โน-อุปาทา ธมฺโม อุปฺปชฺชติ นเหตุปจฺจยา.\csromanspacer{} อเหตุกํ โน-อุปาทา เอกํ ขนฺธํ ปจฺจยา ตโย ขนฺธา โน-อุปาทา จ จิตฺตสมุฏฺฐานํ รูปํ ฯเปฯ เทฺว ขนฺเธ ฯเปฯ.\csromanspacer{} อเหตุกปฏิสนฺธิกฺขเณ ฯเปฯ.\csromanspacer{} เอกํ มหาภูตํ ฯเปฯ. (ยาว อสญฺญสตฺตา.) วิจิกิจฺฉาสหคเต อุทฺธจฺจสหคเต ขนฺเธ ปจฺจยา วิจิกิจฺฉาสหคโต อุทฺธจฺจสหคโต โมโห. (ตีณิปิ.\csromanspacer{} ปฏิจฺจสทิสา นินฺนานากรณา.) (3)}

\prose{อุปาทา จ โน-อุปาทา จ ธมฺมํ ปจฺจยา โน-อุปาทา ธมฺโม อุปฺปชฺชติ นเหตุปจฺจยา.\csromanspacer{} จกฺขุวิญฺญาณสหคตํ เอกํ ขนฺธญฺจ จกฺขายตนญฺจ ปจฺจยา ตโย ขนฺธา ฯเปฯ เทฺว ขนฺเธ จ ฯเปฯ.\csromanspacer{} กายวิญฺญาณสหคตํ ฯเปฯ.\csromanspacer{} อเหตุกํ โน-อุปาทา เอกํ ขนฺธญฺจ วตฺถุญฺจ ปจฺจยา ตโย ขนฺธา ฯเปฯ เทฺว ขนฺเธ จ ฯเปฯ.\csromanspacer{} ปฏิสนฺธิกฺขเณ ฯเปฯ.\csromanspacer{} วิจิกิจฺฉาสหคเต อุทฺธจฺจสหคเต ขนฺเธ จ วตฺถุญฺจ ปจฺจยา วิจิกิจฺฉาสหคโต อุทฺธจฺจสหคโต โมโห. (1)}

\prose{น-อารมฺมณปจฺจยา ตีณิ.\csromanspacer{} น-อาเสวนปจฺจยา ปญฺจ.}

\csromanheader{2}{\textbf{นกมฺมาทิ}}
\proseitem{390}{อุปาทา ธมฺมํ ปจฺจยา โน-อุปาทา ธมฺโม อุปฺปชฺชติ นกมฺมปจฺจยา.\csromanspacer{} วตฺถุํ ปจฺจยา โน-อุปาทา เจตนา. (1)}

\prose{โน-อุปาทา ธมฺมํ ปจฺจยา โน-อุปาทา ธมฺโม อุปฺปชฺชติ นกมฺมปจฺจยา.\csromanspacer{} โน-อุปาทา ขนฺเธ ปจฺจยา สมฺปยุตฺตกา เจตนา.\csromanspacer{} พาหิรํ อาหารสมุฏฺฐานํ อุตุสมุฏฺฐานํ ฯเปฯ เทฺว มหาภูเต ปจฺจยา เทฺว มหาภูตา. (1)}

\prose{โน-อุปาทา ธมฺมํ ปจฺจยา อุปาทา ธมฺโม อุปฺปชฺชติ นกมฺมปจฺจยา.\csromanspacer{} พาหิเร อาหารสมุฏฺฐาเน อุตุสมุฏฺฐาเน มหาภูเต ปจฺจยา อุปาทารูปํ. (2)}

\csromanheader{2}{67. อุปาทาทุก}
\prose{\csromanfolio{525}โน-อุปาทา ธมฺมํ ปจฺจยา อุปาทา จ โน-อุปาทา จ ธมฺมา อุปฺปชฺชนฺติ นกมฺมปจฺจยา.\csromanspacer{} พาหิรํ อาหารสมุฏฺฐานํ อุตุสมุฏฺฐานํ เอกํ มหาภูตํ ปจฺจยา ตโย มหาภูตา อุปาทา จ รูปํ ฯเปฯ เทฺว มหาภูเต ฯเปฯ. (3)}

\prose{อุปาทาจ โน-อุปาทา จ ธมฺมํ ปจฺจยา โน-อุปาทา ธมฺโม อุปฺปชฺชติ นกมฺมปจฺจยา.\csromanspacer{} โน-อุปาทา ขนฺเธ จ วตฺถุญฺจ ปจฺจยา โน-อุปาทา เจตนา. (1)}

\prose{นวิปากปจฺจยา ปญฺจ.\csromanspacer{} น-อาหารปจฺจยา ตีณิ.\csromanspacer{} น-อินฺทฺริยปจฺจยา ตีณิ.}

\csromanheader{2}{\textbf{นฌานาทิ}}
\proseitem{391}{อุปาทา ธมฺมํ ปจฺจยา โน-อุปาทา ธมฺโม อุปฺปชฺชติ นฌานปจฺจยา.\csromanspacer{} จกฺขายตนํ ปจฺจยา จกฺขุวิญฺญาณํ ฯเปฯ.\csromanspacer{} กายายตนํ ปจฺจยา กายวิญฺญาณํ. (1)}

\prose{โน-อุปาทา ธมฺมํ ปจฺจยา โน-อุปาทา ธมฺโม อุปฺปชฺชติ นฌานปจฺจยา.\csromanspacer{} ปญฺจวิญฺญาณสหคตํ เอกํ ขนฺธํ ปจฺจยา ตโย ขนฺธา ฯเปฯ เทฺว ขนฺเธ ฯเปฯ.\csromanspacer{} พาหิรํ อาหารสมุฏฺฐานํ อุตุสมุฏฺฐานํ อสญฺญสตฺตานํ ฯเปฯ เทฺว มหาภูเต ปจฺจยา เทฺว มหาภูตา. (1)}

\prose{โน-อุปาทา ธมฺมํ ปจฺจยา อุปาทา ธมฺโม อุปฺปชฺชติ นฌานปจฺจยา.\csromanspacer{} พาหิเร อาหารสมุฏฺฐาเน อุตุสมุฏฺฐาเน อสญฺญสตฺตานํ มหาภูเต ปจฺจยา อุปาทา กฏตฺตารูปํ. (2)}

\prose{โน-อุปาทา ธมฺมํ ปจฺจยา อุปาทา จ โน-อุปาทา จ ธมฺมา อุปฺปชฺชนฺติ นฌานปจฺจยา.\csromanspacer{} พาหิรํ อาหารสมุฏฺฐานํ อุตุสมุฏฺฐานํ อสญฺญสตฺตานํ เอกํ มหาภูตํ ปจฺจยา ตโย มหาภูตา อุปาทา จ กฏตฺตารูปํ ฯเปฯ เทฺว ฯเปฯ. (3)}

\prose{อุปาทา จ โน-อุปาทา จ ธมฺมํ ปจฺจยา โน-อุปาทา ธมฺโม อุปฺปชฺชติ นฌานปจฺจยา.\csromanspacer{} จกฺขุวิญฺญาณสหคตํ เอกํ ขนฺธญฺจ จกฺขายตนญฺจ ปจฺจยา ตโย ขนฺธา ฯเปฯ เทฺว ขนฺเธ จ ฯเปฯ. (1)}

\prose{นมคฺคปจฺจยา ปญฺจ.\csromanspacer{} นสมฺปยุตฺตปจฺจยา ตีณิ.\csromanspacer{} นวิปฺปยุตฺตปจฺจยา ตีณิ.\csromanspacer{} โนนตฺถิปจฺจยา ตีณิ.\csromanspacer{} โนวิคตปจฺจยา ตีณิ.}

\csromanheader{2}{2. ปจฺจยปจฺจนีย 2. สงฺขฺยาวาร}
\csromanheader{2}{\textbf{สุทฺธ}}
\proseitem{392}{\csromanfolio{526}นเหตุยา ปญฺจ, น-อารมฺมเณ ตีณิ, น-อธิปติยา ปญฺจ ฯเปฯ นกมฺเม ปญฺจ, นวิปาเก ปญฺจ, น-อาหาเร ตีณิ, น-อินฺทฺริเย ตีณิ, นฌาเน ปญฺจ, นมคฺเค ปญฺจ, นสมฺปยุตฺเต ตีณิ, นวิปฺปยุตฺเต ตีณิ, โนนตฺถิยา ตีณิ, โนวิคเต ตีณิ.}

\vspace{\tipitakaparskip}
\csromancenter{(เอวํ อิตเร เทฺว คณนาปิ นิสฺสวาโรปิ กาตพฺโพ.)}
\csromanheader{2}{67. อุปาทาทุก 5. สํสฏฺฐวาร}
\csromanheader{2}{1. ปจฺจยานุโลมาทิ}
\proseitem{393}{โน-อุปาทา ธมฺมํ สํสฏฺโฐ โน-อุปาทา ธมฺโม อุปฺปชฺชติ เหตุปจฺจยา.\csromanspacer{} โน-อุปาทา เอกํ ขนฺธํ สํสฏฺฐา ตโย ขนฺธา ฯเปฯ เทฺว ขนฺเธ สํสฏฺฐา เทฺว ขนฺธา.\csromanspacer{} ปฏิสนฺธิกฺขเณ ฯเปฯ. (สํขิตฺตํ.)}

\prose{เหตุยา เอกํ, อารมฺมเณ เอกํ, อธิปติยา เอกํ, (สพฺพตฺถ เอกํ.) อวิคเต เอกํ.\csromanspacer{} อนุโลมํ.}

\prose{โน-อุปาทา ธมฺมํ สํสฏฺโฐ โน-อุปาทา ธมฺโม อุปฺปชฺชติ นเหตุปจฺจยา.\csromanspacer{} อเหตุกํ โน-อุปาทา เอกํ ขนฺธํ สํสฏฺฐา ตโย ขนฺธา ฯเปฯ เทฺว ขนฺเธ ฯเปฯ.\csromanspacer{} อเหตุกปฏิสนฺธิกฺขเณ ฯเปฯ.\csromanspacer{} วิจิกิจฺฉาสหคเต อุทฺธจฺจสหคเต ขนฺเธ สํสฏฺโฐ วิจิกิจฺฉาสหคโต อุทฺธจฺจสหคโต โมโห. (สํขิตฺตํ.)}

\prose{นเหตุยา เอกํ, น-อธิปติยา เอกํ, นปุเรชาเต เอกํ, นปจฺฉาชาเต เอกํ, น-อาเสวเน เอกํ, นกมฺเม เอกํ, นวิปาเก เอกํ, นฌาเน เอกํ, นมคฺเค เอกํ, นวิปฺปยุตฺเต เอกํ.}

\vspace{\tipitakaparskip}
\csromancenter{ปจฺจนียํ.}
\csromancenter{(อิตเร เทฺว คณนาปิ สมฺปยุตฺตวาโรปิ กาตพฺโพ.)}
\csromanheader{2}{67. อุปาทาทุก \textbf{67. อุปาทาทุก 7. ปญฺหาวาร}}
\csromanheader{2}{1. ปจฺจยานุโลม 1. วิภงฺควาร}
\csromanheader{2}{\textbf{เหตุ}}
\proseitem{394}{\csromanfolio{527}โน-อุปาทา ธมฺโม โน-อุปาทา ธมฺมสฺส เหตุปจฺจเยน ปจฺจโย.\csromanspacer{} โน-อุปาทา เหตู สมฺปยุตฺตกานํ ขนฺธานํ โน-อุปาทา จ จิตฺตสมุฏฺฐานานํ รูปานํ เหตุปจฺจเยน ปจฺจโย.\csromanspacer{} ปฏิสนฺธิกฺขเณ ฯเปฯ. (1)}

\prose{โน-อุปาทา ธมฺโม อุปาทา ธมฺมสฺส เหตุปจฺจเยน ปจฺจโย.\csromanspacer{} โน-อุปาทา เหตู อุปาทา จิตฺตสมุฏฺฐานานํ รูปานํ เหตุปจฺจเยน ปจฺจโย.\csromanspacer{} ปฏิสนฺธิกฺขเณ ฯเปฯ. (2)}

\prose{โน-อุปาทา ธมฺโม อุปาทา จ โน-อุปาทา จ ธมฺมสฺส เหตุปจฺจเยน ปจฺจโย.\csromanspacer{} โน-อุปาทา เหตู สมฺปยุตฺตกานํ ขนฺธานํ อุปาทา จ โน-อุปาทา จ จิตฺตสมุฏฺฐานานํ รูปานํ เหตุปจฺจเยน ปจฺจโย.\csromanspacer{} ปฏิสนฺธิกฺขเณ ฯเปฯ. (3)}

\csromanheader{2}{\textbf{อารมฺมณ}}
\proseitem{395}{อุปาทา ธมฺโม โน-อุปาทา ธมฺมสฺส อารมฺมณปจฺจเยน ปจฺจโย.\csromanspacer{} จกฺขุํ ฯเปฯ กายํ.\csromanspacer{} รูเป ฯเปฯ รเส วตฺถุํ อนิจฺจโต ฯเปฯ โทมนสฺสํ อุปฺปชฺชติ.\csromanspacer{} ทิพฺเพน จกฺขุนา รูปํ ปสฺสติ.\csromanspacer{} ทิพฺพาย โสตธาตุยา สทฺทํ สุณาติ.\csromanspacer{} รูปายตนํ จกฺขุวิญฺญาณสฺส ฯเปฯ.\csromanspacer{} รสายตนํ ชิวฺหาวิญฺญาณสฺส อารมฺมณปจฺจเยน ปจฺจโย.\csromanspacer{} อุปาทา ขนฺธา อิทฺธิวิธญาณสฺส, ปุพฺเพนิวาสานุสฺสติญาณสฺส, อนาคตํสญาณสฺส, อาวชฺชนาย อารมฺมณปจฺจเยน ปจฺจโย. (1)}

\prose{โน-อุปาทา ธมฺโม โน-อุปาทา ธมฺมสฺส อารมฺมณปจฺจเยน ปจฺจโย.\csromanspacer{} ทานํ ฯเปฯ สีลํ ฯเปฯ อุโปสถกมฺมํ กตฺวา ตํ ปจฺจเวกฺขติ, อสฺสาเทติ อภินนฺทติ, ตํ อารพฺภ ราโค อุปฺปชฺชติ ฯเปฯ โทมนสฺสํ อุปฺปชฺชติ.\csromanspacer{} ปุพฺเพ ฯเปฯ.\csromanspacer{} ฌานา ฯเปฯ.\csromanspacer{} อริยา มคฺคา วุฏฺฐหิตฺวา มคฺคํ ปจฺจเวกฺขนฺติ.\csromanspacer{} ผลํ ฯเปฯ.\csromanspacer{} นิพฺพานํ ฯเปฯ.\csromanspacer{} นิพฺพานํ โคตฺรภุสฺส, โวทานสฺส, มคฺคสฺส, ผลสฺส, อาวชฺชนาย อารมฺมณปจฺจเยน ปจฺจโย.\csromanspacer{} อริยา ปหีเน กิเลเส ปจฺจเวกฺขนฺติ.\csromanspacer{} วิกฺขมฺภิเต กิเลเส ฯเปฯ.\csromanspacer{} ปุพฺเพ ฯเปฯ.\csromanspacer{} โผฏฺฐพฺเพ โน-อุปาทา ขนฺเธ อนิจฺจโต ฯเปฯ โทมนสฺสํ อุปฺปชฺชติ.\csromanspacer{} เจโตปริยญาเณน โน-อุปาทาจิตฺตสมงฺคิสฺส จิตฺตํ \csromanfolio{528}ชานาติ.\csromanspacer{} อากาสานญฺจายตนํ วิญฺญาณญฺจายตนสฺส ฯเปฯ.\csromanspacer{} อากิญฺจญฺญายตนํ เนวสญฺญานาสญฺญายตนสฺส ฯเปฯ.\csromanspacer{} โผฏฺฐพฺพายตนํ กายวิญฺญาณสฺส, อารมฺมณปจฺจเยน ปจฺจโย.\csromanspacer{} โน-อุปาทา ขนฺธา อิทฺธิวิธญาณสฺส เจโตปริยญาณสฺส, ปุพฺเพนิวาสานุสฺสติญาณสฺส, ยถากมฺมูปคญาณสฺส, อนาคตํสญาณสฺส, อาวชฺชนาย อารมฺมณปจฺจเยน ปจฺจโย. (1)}

\csromanheader{2}{\textbf{อธิปติ}}
\proseitem{396}{อุปาทา ธมฺโม โน-อุปาทา ธมฺมสฺส อธิปติปจฺจเยน ปจฺจโย.\csromanspacer{} \textbf{อารมฺมณาธิปติ}\csromandash{}จกฺขุํ ฯเปฯ กายํ.\csromanspacer{} รูเป ฯเปฯ รเส วตฺถุํ ครุํ กตฺวา อสฺสาเทติ, อภินนฺทติ, ตํ ครุํ กตฺวา ราโค อุปฺปชฺชติ.\csromanspacer{} ทิฏฺฐิ อุปฺปชฺชติ. (1)}

\prose{โน-อุปาทา ธมฺโม โน-อุปาทา ธมฺมสฺส อธิปติปจฺจเยน ปจฺจโย.\csromanspacer{} \textbf{อารมฺมณาธิปติ}, สหชาตาธิปติ.\csromanspacer{}\csromanspacer{}\textbf{อารมฺมณาธิปติ}\csromandash{}ทานํ ฯเปฯ สีลํ ฯเปฯ อุโปสถกมฺมํ กตฺวา ตํ ครุํ กตฺวา ปจฺจเวกฺขติ, อสฺสาเทติ อภินนฺทติ, ตํ ครุํ กตฺวา ราโค อุปฺปชฺชติ.\csromanspacer{} ทิฏฺฐิ อุปฺปชฺชติ.\csromanspacer{} ปุพฺเพ สุจิณฺณานิ ฯเปฯ.\csromanspacer{} ฌานา วุฏฺฐหิตฺวา ฌานํ ฯเปฯ.\csromanspacer{} อริยา มคฺคา วุฏฺฐหิตฺวา มคฺคํ ครุํ กตฺวา ฯเปฯ.\csromanspacer{} ผลํ ฯเปฯ.\csromanspacer{} นิพฺพานํ ครุํ กตฺวา ปจฺจเวกฺขนฺติ.\csromanspacer{} นิพฺพานํ โคตฺรภุสฺส, โวทานสฺส, มคฺคสฺส, ผลสฺส อธิปติปจฺจเยน ปจฺจโย.\csromanspacer{} โผฏฺฐพฺเพ โน-อุปาทา ขนฺเธ ครุํ กตฺวา อสฺสาเทติ อภินนฺทติ, ตํ ครุํ กตฺวา ราโค อุปฺปชฺชติ.\csromanspacer{} ทิฏฺฐิ อุปฺปชฺชติ.\csromanspacer{} \textbf{สหชาตาธิปติ}\csromandash{}โน-อุปาทาธิปติ สมฺปยุตฺตกานํ ขนฺธานํ โน-อุปาทา จ จิตฺตสมุฏฺฐานานํ รูปานํ อธิปติปจฺจเยน ปจฺจโย. (1)}

\prose{โน-อุปาทา ธมฺโม อุปาทา ธมฺมสฺส อธิปติปจฺจเยน ปจฺจโย.\csromanspacer{} \textbf{สหชาตาธิปติ}\csromandash{}โน-อุปาทาธิปติ อุปาทา จิตฺตสมุฏฺฐานานํ รูปานํ อธิปติปจฺจเยน ปจฺจโย. (2)}

\prose{โน-อุปาทา ธมฺโม อุปาทา จ โน-อุปาทา จ ธมฺมสฺส อธิปติปจฺจเยน ปจฺจโย.\csromanspacer{} \textbf{สหชาตาธิปติ}\csromandash{}โน-อุปาทาธิปติ สมฺปยุตฺตกานํ ขนฺธานํ อุปาทา จ โน-อุปาทา จ จิตฺตสมุฏฺฐานานํ รูปานํ อธิปติปจฺจเยน ปจฺจโย. (3)}

\csromanheader{2}{67. อุปาทาทุก \textbf{อนนฺตราทิ}}
\proseitem{397}{\csromanfolio{529}โน-อุปาทา ธมฺโม โน-อุปาทา ธมฺมสฺส อนนฺตรปจฺจเยน ปจฺจโย.\csromanspacer{} ปุริมา ปุริมา โน-อุปาทา ขนฺธา ปจฺฉิมานํ ปจฺฉิมานํ โน-อุปาทา ขนฺธานํ อนนฺตรปจฺจเยน ปจฺจโย.\csromanspacer{} อนุโลมํ โคตฺรภุสฺส ฯเปฯ ผลสมาปตฺติยา อนนฺตรปจฺจเยน ปจฺจโย. (1)}

\prose{สมนนฺตรปจฺจเยน ปจฺจโย.\csromanspacer{} สหชาตปจฺจเยน ปจฺจโย. (ปฏิจฺจสทิสํ.) อญฺญมญฺญปจฺจเยน ปจฺจโย. (ปฏิจฺจสทิสํ.) นิสฺสยปจฺจเยน ปจฺจโย. (ปจฺจวาเร นิสฺสยสทิสํ.)}

\csromanheader{2}{\textbf{อุปนิสฺสย}}
\proseitem{398}{อุปาทา ธมฺโม โน-อุปาทา ธมฺมสฺส อุปนิสฺสยปจฺจเยน ปจฺจโย.\csromanspacer{} \textbf{อารมฺมณูปนิสฺสโย, ปกตูปนิสฺสโย ฯเปฯ.}\csromanspacer{} ปกตูปนิสฺสโย\csromandash{} จกฺขุสมฺปทํ ฯเปฯ กายสมฺปทํ.\csromanspacer{} วณฺณสมฺปทํ.\csromanspacer{} สทฺทสมฺปทํ.\csromanspacer{} คนฺธสมฺปทํ.\csromanspacer{} รสสมฺปทํ.\csromanspacer{} โภชนํ อุปนิสฺสาย ทานํ เทติ ฯเปฯ สํฆํ ภินฺทติ.\csromanspacer{} จกฺขุสมฺปทา ฯเปฯ.\csromanspacer{} กายสมฺปทา.\csromanspacer{} วณฺณสมฺปทา.\csromanspacer{} สทฺทสมฺปทา.\csromanspacer{} คนฺธสมฺปทา.\csromanspacer{} รสสมฺปทา.\csromanspacer{} โภชนํ สทฺธาย ฯเปฯ ผลสมาปตฺติยา อุปนิสฺสยปจฺจเยน ปจฺจโย. (1)}

\prose{โน-อุปาทา ธมฺโม โน-อุปาทา ธมฺมสฺส อุปนิสฺสยปจฺจเยน ปจฺจโย.\csromanspacer{} \textbf{อารมฺมณูปนิสฺสโย, อนนฺตรูปนิสฺสโย, ปกตูปนิสฺสโย ฯเปฯ.}\csromanspacer{} ปกตูปนิสฺสโย\csromandash{}สทฺธํ อุปนิสฺสาย ทานํ เทติ ฯเปฯ สมาปตฺติํ อุปฺปาเทติ.\csromanspacer{} มานํ ชปฺเปติ.\csromanspacer{} ทิฏฺฐิํ คณฺหาติ.\csromanspacer{} สีลํ ฯเปฯ.\csromanspacer{} ปญฺญํ.\csromanspacer{} ราคํ ฯเปฯ.\csromanspacer{} ปตฺถนํ.\csromanspacer{} กายิกํ สุขํ.\csromanspacer{} กายิกํ ทุกฺขํ.\csromanspacer{} อุตุํ.\csromanspacer{} เสนาสนํ อุปนิสฺสาย ทานํ เทติ ฯเปฯ สํฆํ ภินฺทติ.\csromanspacer{} สทฺธา ฯเปฯ.\csromanspacer{} เสนาสนํ สทฺธาย ฯเปฯ ผลสมาปตฺติยา อุปนิสฺสยปจฺจเยน ปจฺจโย. (1)}

\csromanheader{2}{\textbf{ปุเรชาต}}
\proseitem{399}{อุปาทา ธมฺโม โน-อุปาทา ธมฺมสฺส ปุเรชาตปจฺจเยน ปจฺจโย.\csromanspacer{} \textbf{อารมฺมณปุเรชาตํ}, วตฺถุปุเรชาตํ.\csromanspacer{}\csromanspacer{}\textbf{อารมฺมณปุเรชาตํ}\csromandash{}จกฺขุํ ฯเปฯ.\csromanspacer{} วตฺถุํ อนิจฺจโต ฯเปฯ โทมนสฺสํ อุปฺปชฺชติ.\csromanspacer{} ทิพฺเพน จกฺขุนา รูปํ ปสฺสติ.\csromanspacer{} ทิพฺพาย โสตธาตุยา สทฺทํ สุณาติ.\csromanspacer{} รูปายตนํ จกฺขุวิญฺญาณสฺส ฯเปฯ. \csromanfolio{530}รสายตนํ ชิวฺหาวิญฺญาณสฺส ฯเปฯ.\csromanspacer{} \textbf{วตฺถุปุเรชาตํ}\csromandash{}จกฺขายตนํ จกฺขุวิญฺญาณสฺส ฯเปฯ กายายตนํ ฯเปฯ.\csromanspacer{} วตฺถุ โน-อุปาทา ขนฺธานํ ปุเรชาตปจฺจเยน ปจฺจโย. (1)}

\prose{โน-อุปาทา ธมฺโม โน-อุปาทา ธมฺมสฺส ปุเรชาตปจฺจเยน ปจฺจโย.\csromanspacer{} \textbf{อารมฺมณปุเรชาตํ}\csromandash{}โผฏฺฐพฺเพ อนิจฺจโต ฯเปฯ โทมนสฺสํ อุปฺปชฺชติ.\csromanspacer{} โผฏฺฐพฺพายตนํ กายวิญฺญาณสฺส ปุเรชาตปจฺจเยน ปจฺจโย. (1)}

\prose{อุปาทา จ โน-อุปาทา จ ธมฺมา โน-อุปาทา ธมฺมสฺส ปุเรชาตปจฺจเยน ปจฺจโย.\csromanspacer{} \textbf{อารมฺมณปุเรชาตํ, วตฺถุปุเรชาตํ.}\csromanspacer{} โผฏฺฐพฺพายตนญฺจ กายายตนญฺจ กายวิญฺญาณสฺส ปุเรชาตปจฺจเยน ปจฺจโย.\csromanspacer{} โผฏฺฐพฺพายตนญฺจ วตฺถุ จ โน-อุปาทา ขนฺธานํ ปุเรชาตปจฺจเยน ปจฺจโย. (1)}

\csromanheader{2}{\textbf{ปจฺฉาชาตเสวน}}
\proseitem{400}{โน-อุปาทา ธมฺโม โน-อุปาทา ธมฺมสฺส ปจฺฉาชาตปจฺจเยน ปจฺจโย.\csromanspacer{} ปจฺฉาชาตา โน-อุปาทา ขนฺธา ปุเรชาตสฺส อิมสฺส โน-อุปาทา กายสฺส ปจฺฉาชาตปจฺจเยน ปจฺจโย. (มูลํ ปุจฺฉิตพฺพํ.) ปจฺฉาชาตา โน-อุปาทา ขนฺธา ปุเรชาตสฺส อิมสฺส อุปาทา กายสฺส ปจฺฉาชาตปจฺจเยน ปจฺจโย. (มูลํ ปุจฺฉิตพฺพํ.) ปจฺฉาชาตา โน-อุปาทา ขนฺธา ปุเรชาตสฺส อิมสฺส อุปาทา กายสฺส จ โน-อุปาทา กายสฺส จ ปจฺฉาชาตปจฺจเยน ปจฺจโย. (3) อาเสวนปจฺจเยน ปจฺจโย.\csromanspacer{} เอกํ.}

\csromanheader{2}{\textbf{กมฺม}}
\proseitem{401}{โน-อุปาทา ธมฺโม โน-อุปาทา ธมฺมสฺส กมฺมปจฺจเยน ปจฺจโย.\csromanspacer{} \textbf{สหชาตา}, นานากฺขณิกา.\csromanspacer{}\csromanspacer{}\textbf{สหชาตา} โน-อุปาทา เจตนา สมฺปยุตฺตกานํ ขนฺธานํ โน-อุปาทา จ จิตฺตสมุฏฺฐานานํ รูปานํ กมฺมปจฺจเยน ปจฺจโย.\csromanspacer{} ปฏิสนฺธิกฺขเณ ฯเปฯ.\csromanspacer{} \textbf{นานากฺขณิกา} โน-อุปาทา เจตนา วิปากานํ ขนฺธานํ โน-อุปาทา จ กฏตฺตารูปานํ กมฺมปจฺจเยน ปจฺจโย. (1)}

\prose{โน-อุปาทา ธมฺโม อุปาทา ธมฺมสฺส กมฺมปจฺจเยน ปจฺจโย.\csromanspacer{} \textbf{สหชาตา}, นานากฺขณิกา.\csromanspacer{}\csromanspacer{}\textbf{สหชาตา} โน-อุปาทา เจตนา อุปาทา จิตฺตสมุฏฺฐานานํ รูปานํ กมฺมปจฺจเยน ปจฺจโย.\csromanspacer{} ปฏิสนฺธิกฺขเณ ฯเปฯ.\csromanspacer{} \textbf{นานากฺขณิกา} โน-อุปาทา เจตนา อุปาทา กฏตฺตารูปานํ กมฺมปจฺจเยน ปจฺจโย. (2)}

\csromanheader{2}{67. อุปาทาทุก}
\prose{\csromanfolio{531}โน-อุปาทา ธมฺโม อุปาทา จ โน-อุปาทา จ ธมฺมสฺส กมฺมปจฺจเยน ปจฺจโย.\csromanspacer{} \textbf{สหชาตา}, นานากฺขณิกา.\csromanspacer{}\csromanspacer{}\textbf{สหชาตา} โน-อุปาทา เจตนา สมฺปยุตฺตกานํ ขนฺธานํ อุปาทา จ โน-อุปาทา จ จิตฺตสมุฏฺฐานานํ รูปานํ กมฺมปจฺจเยน ปจฺจโย.\csromanspacer{} ปฏิสนฺธิกฺขเณ ฯเปฯ.\csromanspacer{} \textbf{นานากฺขณิกา} โน-อุปาทา เจตนา วิปากานํ ขนฺธานํ อุปาทา จ โน-อุปาทา จ กฏตฺตารูปานํ กมฺมปจฺจเยน ปจฺจโย. (3)}

\csromanheader{2}{\textbf{วิปาก}}
\proseitem{402}{โน-อุปาทา ธมฺโม โน-อุปาทา ธมฺมสฺส วิปากปจฺจเยน ปจฺจโย.\csromanspacer{} วิปาโก โน-อุปาทา เอโก ขนฺโธ ติณฺณนฺนํ ขนฺธานํ ฯเปฯ เทฺว ขนฺธา ฯเปฯ.\csromanspacer{} ปฏิสนฺธิกฺขเณ ฯเปฯ.\csromanspacer{} ตีณิ.}

\csromanheader{2}{\textbf{อาหาร}}
\proseitem{403}{อุปาทา ธมฺโม อุปาทา ธมฺมสฺส อาหารปจฺจเยน ปจฺจโย.\csromanspacer{} กพฬีกาโร อาหาโร อิมสฺส อุปาทา กายสฺส อาหารปจฺจเยน ปจฺจโย. (มูลํ ปุจฺฉิตพฺพํ.) กพฬีกาโร อาหาโร อิมสฺส โน-อุปาทา กายสฺส อาหารปจฺจเยน ปจฺจโย. (มูลํ ปุจฺฉิตพฺพํ.) กพฬีกาโร อาหาโร อิมสฺส อุปาทา กายสฺส จ โน-อุปาทา กายสฺส จ อาหารปจฺจเยน ปจฺจโย. (3)}

\prose{โน-อุปาทา ธมฺโม โน-อุปาทา ธมฺมสฺส อาหารปจฺจเยน ปจฺจโย.\csromanspacer{} โน-อุปาทา อาหารา สมฺปยุตฺตกานํ ขนฺธานํ โน-อุปาทา จ จิตฺตสมุฏฺฐานานํ รูปานํ อาหารปจฺจเยน ปจฺจโย.\csromanspacer{} ปฏิสนฺธิกฺขเณ ฯเปฯ. (โน-อุปาทามูลเก ตีณิ.\csromanspacer{} ปฏิสนฺธิกฺขเณ ตีณิปิ กาตพฺพา.) (3)}

\csromanheader{2}{\textbf{อินฺทฺริย}}
\proseitem{404}{อุปาทา ธมฺโม อุปาทา ธมฺมสฺส อินฺทฺริยปจฺจเยน ปจฺจโย.\csromanspacer{} รูปชีวิตินฺทฺริยํ อุปาทา กฏตฺตารูปานํ อินฺทฺริยปจฺจเยน ปจฺจโย. (มูลํ กาตพฺพํ.) จกฺขุนฺทฺริยํ จกฺขุวิญฺญาณสฺส ฯเปฯ.\csromanspacer{} กายินฺทฺริยํ กายวิญฺญาณสฺส รูปชีวิตินฺทฺริยํ โน-อุปาทา กฏตฺตารูปานํ อินฺทฺริยปจฺจเยน ปจฺจโย. (มูลํ กาตพฺพํ.) รูปชีวิตินฺทฺริยํ อุปาทา จ โน-อุปาทา จ กฏตฺตารูปานํ อินฺทฺริยปจฺจเยน ปจฺจโย. (3)}

\prose{\csromanfolio{532}โน-อุปาทา ธมฺโม โน-อุปาทา ธมฺมสฺส อินฺทฺริยปจฺจเยน ปจฺจโย.\csromanspacer{} ตีณิ.\csromanspacer{} ปฏิสนฺธิกฺขเณ ฯเปฯ.}

\prose{อุปาทา จ โน-อุปาทา จ ธมฺมา โน-อุปาทา ธมฺมสฺส อินฺทฺริยปจฺจเยน ปจฺจโย.\csromanspacer{} จกฺขุนฺทฺริยญฺจ จกฺขุวิญฺญาณญฺจ จกฺขุวิญฺญาณสหคตานํ ขนฺธานํ อินฺทฺริยปจฺจเยน ปจฺจโย ฯเปฯ.\csromanspacer{} กยินฺทฺริยญฺจ ฯเปฯ.}

\csromanheader{2}{\textbf{ฌานาทิ}}
\proseitem{405}{โน อุปาทา ธมฺโม โน-อุปาทา ธมฺมสฺส ฌานปจฺจเยน ปจฺจโย.\csromanspacer{} ตีณิ.\csromanspacer{} มคฺคปจฺจเยน ปจฺจโย.\csromanspacer{} ตีณิ.\csromanspacer{} ปฏิสนฺธิกฺขเณ ฯเปฯ.\csromanspacer{} สมฺปยุตฺตปจฺจเยน ปจฺจโย.\csromanspacer{} เอกํ.}

\csromanheader{2}{\textbf{วิปฺปยุตฺต}}
\proseitem{406}{อุปาทา ธมฺโม โน-อุปาทา ธมฺมสฺส วิปฺปยุตฺตปจฺจเยน ปจฺจโย.\csromanspacer{} \textbf{สหชาตํ}, ปุเรชาตํ.\csromanspacer{}\csromanspacer{}\textbf{สหชาตํ} ปฏิสนฺธิกฺขเณ วตฺถุ โน-อุปาทา ขนฺธานํ วิปฺปยุตฺตปจฺจเยน ปจฺจโย.\csromanspacer{} \textbf{ปุเรชาตํ} จกฺขายตนํ จกฺขุวิญฺญาณสฺส ฯเปฯ กายายตนํ ฯเปฯ.\csromanspacer{} วตฺถุ โน-อุปาทา ขนฺธานํ วิปฺปยุตฺตปจฺจเยน ปจฺจโย. (1)}

\prose{โน-อุปาทา ธมฺโม โน-อุปาทา ธมฺมสฺส วิปฺปยุตฺตปจฺจเยน ปจฺจโย.\csromanspacer{} \textbf{สหชาตํ}, ปจฺฉาชาตํ.\csromanspacer{}\csromanspacer{}\textbf{สหชาตา} โน-อุปาทา ขนฺธา โน-อุปาทา จิตฺตสมุฏฺฐานานํ รูปานํ วิปฺปยุตฺตปจฺจเยน ปจฺจโย.\csromanspacer{} ปฏิสนฺธิกฺขเณ โน-อุปาทา ขนฺธา โน-อุปาทา กฏตฺตารูปานํ วิปฺปยุตฺตปจฺจเยน ปจฺจโย.\csromanspacer{} \textbf{ปจฺฉาชาตา} โน-อุปาทา ขนฺธา ปุเรชาตสฺส อิมสฺส โน-อุปาทา กายสฺส วิปฺปยุตฺตปจฺจเยน ปจฺจโย. (1)}

\prose{โน-อุปาทา ธมฺโม อุปาทา ธมฺมสฺส วิปฺปยุตฺตปจฺจเยน ปจฺจโย.\csromanspacer{} \textbf{สหชาตํ}, ปจฺฉาชาตํ.\csromanspacer{}\csromanspacer{}\textbf{สหชาตา} โน-อุปาทา ขนฺธา อุปาทา จิตฺตสมุฏฺฐานานํ รูปานํ วิปฺปยุตฺตปจฺจเยน ปจฺจโย.\csromanspacer{} ปฏิสนฺธิกฺขเณ ฯเปฯ.\csromanspacer{} \textbf{ปจฺฉาชาตา} โน-อุปาทา ขนฺธา ปุเรชาตสฺส อิมสฺส อุปาทา กายสฺส วิปฺปยุตฺตปจฺจเยน ปจฺจโย. (2)}

\prose{โน-อุปาทา ธมฺโม อุปาทา ธมฺมสฺส จ โน-อุปาทา ธมฺมสฺส จ วิปฺปยุตฺตปจฺจเยน ปจฺจโย.\csromanspacer{} \textbf{สหชาตํ}, ปจฺฉาชาตํ.\csromanspacer{}\csromanspacer{}\textbf{สหชาตา} โน-อุปาทา ขนฺธา}

\vspace{\tipitakaparskip}
\csromancenter{อุปาทาทุก อุปาทา จ โน-อุปาทา จ จิตฺตสมุฏฺฐานานํ รูปานํ วิปฺปยุตฺตปจฺจเยน ปจฺจโย.\csromanspacer{} ปฏิสนฺธิกฺขเณ ฯเปฯ.\csromanspacer{} \textbf{ปจฺฉาชาตา} โน-อุปาทา ขนฺธา ปุเรชาตสฺส อิมสฺส อุปาทา กายสฺส จ โน-อุปาทา กายสฺส จ วิปฺปยุตฺตปจฺจเยน ปจฺจโย. (3)}
\csromanheader{2}{\textbf{อตฺถิ}}
\proseitem{407}{\csromanfolio{533}อุปาทา ธมฺโม อุปาทา ธมฺมสฺส อตฺถิปจฺจเยน ปจฺจโย.\csromanspacer{} อาหารํ, อินฺทฺริยํ.\csromanspacer{} \textbf{กพฬีกาโร อาหาโร} อิมสฺส อุปาทา กายสฺส อตฺถิปจฺจเยน ปจฺจโย.\csromanspacer{} \textbf{รูปชีวิตินฺทฺริยํ} อุปาทา กฏตฺตารูปานํ อตฺถิปจฺจเยน ปจฺจโย. (1)}

\prose{อุปาทา ธมฺโม โน-อุปาทา ธมฺมสฺส อตฺถิปจฺจเยน ปจฺจโย.\csromanspacer{} \textbf{สหชาตํ}, ปุเรชาตํ, อาหารํ, อินฺทฺริยํ.\csromanspacer{}\csromanspacer{}\textbf{สหชาตํ} ปฏิสนฺธิกฺขเณ วตฺถุ โน-อุปาทา ขนฺธานํ อตฺถิปจฺจเยน ปจฺจโย.\csromanspacer{} \textbf{ปุเรชาตํ} จกฺขุํ อนิจฺจโต ฯเปฯ. (สํขิตฺตํ.\csromanspacer{} ปุเรชาตสทิสํ.\csromanspacer{} นินฺนานากรณํ.) \textbf{กพฬีกาโร อาหาโร} อิมสฺส โน-อุปาทา กายสฺส อตฺถิปจฺจเยน ปจฺจโย.\csromanspacer{} \textbf{รูปชีวิตินฺทฺริยํ} โน-อุปาทา กฏตฺตารูปานํ อตฺถิปจฺจเยน ปจฺจโย. (2)}

\prose{อุปาทา ธมฺโม อุปาทา ธมฺมสฺส จ โน-อุปาทา ธมฺมสฺส จ อตฺถิปจฺจเยน ปจฺจโย.\csromanspacer{} อาหารํ, อินฺทฺริยํ.\csromanspacer{} \textbf{กพฬีกาโร อาหาโร} อิมสฺส อุปาทา กายสฺส จ โน-อุปาทา กายสฺส จ อตฺถิปจฺจเยน ปจฺจโย.\csromanspacer{} \textbf{รูปชีวิตินฺทฺริยํ} อุปาทา จ โน-อุปาทา จ กฏตฺตารูปานํ อตฺถิปจฺจเยน ปจฺจโย. (3)}

\proseitem{408}{โน-อุปาทา ธมฺโม โน-อุปาทา ธมฺมสฺส อตฺถิปจฺจเยน ปจฺจโย.\csromanspacer{} \textbf{สหชาตํ}, ปุเรชาตํ, ปจฺฉาชาตํ.\csromanspacer{}\csromanspacer{}\textbf{สหชาโต} โน-อุปาทา เอโก ขนฺโธ ติณฺณนฺนํ ขนฺธานํ โน-อุปาทา จ จิตฺตสมุฏฺฐานานํ รูปานํ อตฺถิปจฺจเยน ปจฺจโย ฯเปฯ เทฺว ขนฺธา ฯเปฯ.\csromanspacer{} ปฏิสนฺธิกฺขเณ ฯเปฯ.\csromanspacer{} เอกํ มหาภูตํ ติณฺณนฺนํ มหาภูตานํ อตฺถิปจฺจเยน ปจฺจโย ฯเปฯ เทฺว มหาภูตา ทฺวินฺนํ มหาภูตานํ อตฺถิปจฺจเยน ปจฺจโย. (ยาว อสญฺญสตฺตา.) \textbf{ปุเรชาตํ}\csromandash{}โผฏฺฐพฺเพ อนิจฺจโต ฯเปฯ โทมนสฺสํ อุปฺปชฺชติ.\csromanspacer{} โผฏฺฐพฺพายตนํ กายวิญฺญาณสฺส อตฺถิปจฺจเยน ปจฺจโย.\csromanspacer{} \textbf{ปจฺฉาชาตา} โน-อุปาทา ขนฺธา ปุเรชาตสฺส อิมสฺส โน-อุปาทา กายสฺส อตฺถิปจฺจเยน ปจฺจโย. (1)}

\prose{โน-อุปาทา ธมฺโม อุปาทา ธมฺมสฺส อตฺถิปจฺจเยน ปจฺจโย.\csromanspacer{} \textbf{สหชาตํ}, ปจฺฉาชาตํ.\csromanspacer{}\csromanspacer{}\textbf{สหชาตา} โน-อุปาทา ขนฺธา อุปาทา จิตฺตสมุฏฺฐานานํ \csromanfolio{534}รูปานํ อตฺถิปจฺจเยน ปจฺจโย.\csromanspacer{} ปฏิสนฺธิกฺขเณ ฯเปฯ.\csromanspacer{} \textbf{ปจฺฉาชาตา} โน-อุปาทา ขนฺธา ปุเรชาตสฺส อิมสฺส อุปาทา กายสฺส อตฺถิปจฺจเยน ปจฺจโย. (2)}

\prose{โน-อุปาทา ธมฺโม อุปาทา จ โน-อุปาทา จ ธมฺมสฺส อตฺถิปจฺจเยน ปจฺจโย.\csromanspacer{} สหชาตํ, ปจฺฉาชาตํ.\csromanspacer{}\csromanspacer{}\textbf{สหชาโต} โน-อุปาทา เอโก ขนฺโธ ติณฺณนฺนํ ขนฺธานํ อุปาทา จ โน-อุปาทา จ จิตฺตสมุฏฺฐานานํ รูปานํ อตฺถิปจฺจเยน ปจฺจโย.\csromanspacer{} เทฺว ขนฺธา ฯเปฯ.\csromanspacer{} ปฏิสนฺธิกฺขเณ ฯเปฯ.\csromanspacer{} \textbf{ปจฺฉาชาตา} โน-อุปาทา ขนฺธา ปุเรชาตสฺส อิมสฺส อุปาทา กายสฺส จ โน-อุปาทา กายสฺส จ อตฺถิปจฺจเยน ปจฺจโย. (3)}

\proseitem{409}{อุปาทา จ โน-อุปาทา จ ธมฺมา อุปาทา ธมฺมสฺส อตฺถิปจฺจเยน ปจฺจโย.\csromanspacer{} ปจฺฉาชาตํ, อาหารํ, อินฺทฺริยํ.\csromanspacer{}\csromanspacer{}\textbf{ปจฺฉาชาตา} โน-อุปาทา ขนฺธา จ กพฬีกาโร อาหาโร จ อิมสฺส อุปาทา กายสฺส อตฺถิปจฺจเยน ปจฺจโย.\csromanspacer{} \textbf{ปจฺฉาชาตา} โน-อุปาทา ขนฺธา จ รูปชีวิตินฺทฺริยญฺจ อุปาทา กฏตฺตารูปานํ อตฺถิปจฺจเยน ปจฺจโย. (1)}

\prose{อุปาทา จ โน-อุปาทา จ ธมฺมา โน-อุปาทา ธมฺมสฺส อตฺถิปจฺจเยน ปจฺจโย.\csromanspacer{} สหชาตํ, ปุเรชาตํ, ปจฺฉาชาตํ, อาหารํ, อินฺทฺริยํ.\csromanspacer{}\csromanspacer{}\textbf{สหชาโต} จกฺขุวิญฺญาณสหคโต เอโก ขนฺโธ จ จกฺขายตนญฺจ ติณฺณนฺนํ ขนฺธานํ อตฺถิปจฺจเยน ปจฺจโย ฯเปฯ เทฺว ขนฺธา จ ฯเปฯ.\csromanspacer{} โน-อุปาทา เอโก ขนฺโธ จ วตฺถุ จ ติณฺณนฺนํ ขนฺธานํ อตฺถิปจฺจเยน ปจฺจโย ฯเปฯ เทฺว ขนฺธา จ ฯเปฯ.\csromanspacer{} ปฏิสนฺธิกฺขเณ โน-อุปาทา เอโก ขนฺโธ จ วตฺถุ จ ติณฺณนฺนํ ขนฺธานํ อตฺถิปจฺจเยน ปจฺจโย ฯเปฯ เทฺว ขนฺธา จ ฯเปฯ.\csromanspacer{} \textbf{ปุเรชาตํ} โผฏฺฐพฺพายตนญฺจ กายายตนญฺจ กายวิญฺญาณสฺส อตฺถิปจฺจเยน ปจฺจโย.\csromanspacer{} โผฏฺฐพฺพายตนญฺจ วตฺถุ จ โน-อุปาทา ขนฺธานํ อตฺถิปจฺจเยน ปจฺจโย.\csromanspacer{} \textbf{ปจฺฉาชาตา} โน-อุปาทา ขนฺธา จ กพฬีกาโร อาหาโร จ อิมสฺส โน-อุปาทา กายสฺส อตฺถิปจฺจเยน ปจฺจโย.\csromanspacer{} \textbf{ปจฺฉาชาตา} โน-อุปาทา ขนฺธา จ รูปชีวิตินฺทฺริยญฺจ โน-อุปาทา กฏตฺตารูปานํ อตฺถิปจฺจเยน ปจฺจโย. (2)}

\prose{อุปาทา จ โน-อุปาทา จ ธมฺมา อุปาทา ธมฺมสฺส จ โน-อุปาทา ธมฺมสฺส จ อตฺถิปจฺจเยน ปจฺจโย.\csromanspacer{} ปจฺฉาชาตํ, อาหารํ, อินฺทฺริยํ.\csromanspacer{}\csromanspacer{}\textbf{ปจฺฉาชาตา} โน-อุปาทา ขนฺธา จ กพฬีกาโร อาหาโร จ อิมสฺส อุปาทา กายสฺส จ โน-อุปาทา กายสฺส จ อตฺถิปจฺจเยน ปจฺจโย.}

\vspace{\tipitakaparskip}
\csromancenter{อุปาทาทุก \textbf{ปจฺฉาชาตา} โน-อุปาทา ขนฺธา จ รูปชีวิตินฺทฺริยญฺจ อุปาทา จ โน-อุปาทา จ กฏตฺตารูปานํ อตฺถิปจฺจเยน ปจฺจโย. (3)}
\csromanheader{2}{2. ปจฺจยานุโลม 2. สงฺขฺยาวาร}
\csromanheader{2}{\textbf{สุทฺธ}}
\proseitem{410}{\csromanfolio{535}เหตุยา ตีณิ, อารมฺมเณ เทฺว, อธิปติยา จตฺตาริ, อนนฺตเร เอกํ, สมนนฺตเร เอกํ, สหชาเต ปญฺจ, อญฺญมญฺเญ ปญฺจ, นิสฺสเย ปญฺจ, อุปนิสฺสเย เทฺว, ปุเรชาเต ตีณิ, ปจฺฉาชาเต ตีณิ, อาเสวเน เอกํ, กมฺเม ตีณิ, วิปาเก ตีณิ, อาหาเร ฉ, อินฺทฺริเย สตฺต, ฌาเน ตีณิ, มคฺเค ตีณิ, สมฺปยุตฺเต เอกํ, วิปฺปยุตฺเต จตฺตาริ, อตฺถิยา นว, นตฺถิยา เอกํ, วิคเต เอกํ, อวิคเต นว.}

\vspace{\tipitakaparskip}
\csromancenter{อนุโลมํ.}
\csromanheader{2}{2. ปจฺจนียุทฺธาร}
\proseitem{411}{อุปาทา ธมฺโม อุปาทา ธมฺมสฺส อาหารปจฺจเยน ปจฺจโย.\csromanspacer{} อินฺทฺริยปจฺจเยน ปจฺจโย. (1)}

\prose{อุปาทา ธมฺโม โน-อุปาทา ธมฺมสฺส อารมฺมณปจฺจเยน ปจฺจโย.\csromanspacer{} สหชาตปจฺจเยน ปจฺจโย.\csromanspacer{} อุปนิสฺสยปจฺจเยน ปจฺจโย.\csromanspacer{} ปุเรชาตปจฺจเยน ปจฺจโย.\csromanspacer{} อาหารปจฺจเยน ปจฺจโย.\csromanspacer{} อินฺทฺริยปจฺจเยน ปจฺจโย. (2)}

\prose{อุปาทา ธมฺโม อุปาทา ธมฺมสฺส จ โน-อุปาทา ธมฺมสฺส จ อาหารปจฺจเยน ปจฺจโย.\csromanspacer{} อินฺทฺริยปจฺจเยน ปจฺจโย. (3)}

\proseitem{412}{โน-อุปาทา ธมฺโม โน-อุปาทา ธมฺมสฺส อารมฺมณปจฺจเยน ปจฺจโย.\csromanspacer{} สหชาตปจฺจเยน ปจฺจโย.\csromanspacer{} อุปนิสฺสยปจฺจเยน ปจฺจโย.\csromanspacer{} ปุเรชาตปจฺจเยน ปจฺจโย.\csromanspacer{} ปจฺฉาชาตปจฺจเยน ปจฺจโย.\csromanspacer{} กมฺมปจฺจเยน ปจฺจโย. (1)}

\prose{โน-อุปาทา ธมฺโม อุปาทา ธมฺมสฺส สหชาตปจฺจเยน ปจฺจโย.\csromanspacer{} ปจฺฉาชาตปจฺจเยน ปจฺจโย.\csromanspacer{} กมฺมปจฺจเยน ปจฺจโย. (2)}

\prose{\csromanfolio{536}โน-อุปาทา ธมฺโม อุปาทา ธมฺมสฺส จ โน-อุปาทา ธมฺมสฺส จ สหชาตปจฺจเยน ปจฺจโย.\csromanspacer{} ปจฺฉาชาตปจฺจเยน ปจฺจโย.\csromanspacer{} กมฺมปจฺจเยน ปจฺจโย. (3)}

\prose{อุปาทา จ โน-อุปาทา จ ธมฺมา อุปาทา ธมฺมสฺส ปจฺฉาชาตํ, อาหารํ, อินฺทฺริยํ. (1)}

\prose{อุปาทา จ โน-อุปาทา จ ธมฺมา โน-อุปาทา ธมฺมสฺส สหชาตํ, ปุเรชาตํ, ปจฺฉาชาตํ, อาหารํ, อินฺทฺริยํ. (2)}

\prose{อุปาทา จ โน-อุปาทา จ ธมฺมา อุปาทา ธมฺมสฺส จ โน-อุปาทา ธมฺมสฺส จ ปจฺฉาชาตํ, อาหารํ, อินฺทฺริยํ. (3)}

\csromanheader{2}{2. ปจฺจยปจฺจนีย 2. สงฺขฺยาวาร}
\csromanheader{2}{\textbf{สุทฺธ}}
\proseitem{413}{นเหตุยา นว, น-อารมฺมเณ นว, (สพฺพตฺถ นว.) นสมฺปยุตฺเต นว, นวิปฺปยุตฺเต ฉ, โน-อตฺถิยา จตฺตาริ, โนนตฺถิยา นว, โนวิคเต นว, โน-อวิคเต จตฺตาริ.}

\csromanheader{2}{3. ปจฺจยานุโลมปจฺจนีย}
\proseitem{414}{เหตุปจฺจยา น-อารมฺมเณ ตีณิ, น-อธิปติยา ตีณิ, น-อนนฺตเร ตีณิ, นสมนนฺตเร ตีณิ, น-อญฺญมญฺเญ ตีณิ, น-อุปนิสฺสเย ตีณิ, (สพฺพตฺถ ตีณิ.) นสมฺปยุตฺเต ตีณิ, นวิปฺปยุตฺเต เอกํ, โนนตฺถิยา ตีณิ, โนวิคเต ตีณิ.}

\csromanheader{2}{4. ปจฺจยปจฺจนียานุโลม}
\proseitem{415}{นเหตุปจฺจยา อารมฺมเณ เทฺว, อธิปติยา จตฺตาริ ฯเปฯ. (อนุโลมมาติกา วิตฺถาเรตพฺพา.) ฯเปฯ อวิคเต นว.}

\nitthitam{อุปาทาทุกํ นิฏฺฐิตํ.}
\csromanmatikaanchor{mtk.67}
\csromantocmarkref{subsection}{68. อุปาทินฺนทุก}{mtk.67}
\bookmark[dest={mtk.67},level=2]{68. อุปาทินฺนทุก}
\csromanheader{2}{68. อุปาทินฺนทุก \textbf{68. อุปาทินฺนทุก 1. ปฏิจฺจวาร}}
\csromanheader{2}{1. ปจฺจยานุโลม 1. วิภงฺควาร}
\csromanheader{2}{\textbf{เหตุ}}
\proseitem{416}{\csromanfolio{537}อุปาทินฺนํ ธมฺมํ ปฏิจฺจ อุปาทินฺโน ธมฺโม อุปฺปชฺชติ เหตุปจฺจยา.\csromanspacer{} อุปาทินฺนํ เอกํ ขนฺธํ ปฏิจฺจ ตโย ขนฺธา ฯเปฯ เทฺว ขนฺเธ ฯเปฯ.\csromanspacer{} ปฏิสนฺธิกฺขเณ อุปาทินฺนํ เอกํ ขนฺธํ ปฏิจฺจ ตโย ขนฺธา กฏตฺตา จ รูปํ ฯเปฯ เทฺว ขนฺเธ ฯเปฯ.\csromanspacer{} ขนฺเธ ปฏิจฺจ วตฺถุ, วตฺถุํ ปฏิจฺจ ขนฺธา.\csromanspacer{} เอกํ มหาภูตํ ฯเปฯ มหาภูเต ปฏิจฺจ กฏตฺตารูปํ, อุปาทารูปํ. (1)}

\prose{อุปาทินฺนํ ธมฺมํ ปฏิจฺจ อนุปาทินฺโน ธมฺโม อุปฺปชฺชติ เหตุปจฺจยา.\csromanspacer{} อุปาทินฺเน ขนฺเธ ปฏิจฺจ จิตฺตสมุฏฺฐานํ รูปํ. (2)}

\prose{อุปาทินฺนํ ธมฺมํ ปฏิจฺจ อุปาทินฺโน จ อนุปาทินฺโน จ ธมฺมา อุปฺปชฺชนฺติ เหตุปจฺจยา.\csromanspacer{} อุปาทินฺนํ เอกํ ขนฺธํ ปฏิจฺจ ตโย ขนฺธา จิตฺตสมุฏฺฐานญฺจ รูปํ ฯเปฯ เทฺว ขนฺเธ ฯเปฯ. (3)}

\prose{อนุปาทินฺนํ ธมฺมํ ปฏิจฺจ อนุปาทินฺโน ธมฺโม อุปฺปชฺชติ เหตุปจฺจยา.\csromanspacer{} อนุปาทินฺนํ เอกํ ขนฺธํ ปฏิจฺจ ตโย ขนฺธา จิตฺตสมุฏฺฐานญฺจ รูปํ ฯเปฯ เทฺว ขนฺเธ ฯเปฯ.\csromanspacer{} เอกํ มหาภูตํ ปฏิจฺจ ฯเปฯ มหาภูเต ปฏิจฺจ จิตฺตสมุฏฺฐานํ รูปํ, อุปาทารูปํ. (1)}

\prose{อุปาทินฺนญฺจ อนุปาทินฺนญฺจ ธมฺมํ ปฏิจฺจ อนุปาทินฺโน ธมฺโม อุปฺปชฺชติ เหตุปจฺจยา.\csromanspacer{} อุปาทินฺเน ขนฺเธ จ มหาภูเต จ ปฏิจฺจ จิตฺตสมุฏฺฐานํ รูปํ. (1)}

\csromanheader{2}{\textbf{อารมฺมณ}}
\proseitem{417}{อุปาทินฺนํ ธมฺมํ ปฏิจฺจ อุปาทินฺโน ธมฺโม อุปฺปชฺชติ อารมฺมณปจฺจยา.\csromanspacer{} อุปาทินฺนํ เอกํ ขนฺธํ ปฏิจฺจ ตโย ขนฺธา ฯเปฯ เทฺว ขนฺเธ ฯเปฯ.\csromanspacer{} ปฏิสนฺธิกฺขเณ ฯเปฯ.\csromanspacer{} วตฺถุํ ปฏิจฺจ ขนฺธา. (1)}

\prose{อนุปาทินฺนํ ธมฺมํ ปฏิจฺจ อนุปาทินฺโน ธมฺโม อุปฺปชฺชติ อารมฺมณปจฺจยา.\csromanspacer{} อนุปาทินฺนํ เอกํ ขนฺธํ ปฏิจฺจ ตโย ขนฺธา ฯเปฯ เทฺว ขนฺเธ ฯเปฯ. (1)}

\csromanheader{2}{\textbf{อธิปติ}}
\proseitem{418}{\csromanfolio{538}อนุปาทินฺนํ ธมฺมํ ปฏิจฺจ อนุปาทินฺโน ธมฺโม อุปฺปชฺชติ อธิปติปจฺจยา.\csromanspacer{} อนุปาทินฺนํ เอกํ ขนฺธํ ปฏิจฺจ ตโย ขนฺธา จิตฺตสมุฏฺฐานญฺจ รูปํ ฯเปฯ เทฺว ขนฺเธ ฯเปฯ.\csromanspacer{} เอกํ มหาภูตํ ปฏิจฺจ ฯเปฯ มหาภูเต ปฏิจฺจ จิตฺตสมุฏฺฐานํ รูปํ, อุปาทารูปํ. (สํขิตฺตํ.)}

\csromanheader{2}{1. ปจฺจยานุโลม 2. สงฺขฺยาวาร}
\csromanheader{2}{\textbf{สุทฺธ}}
\proseitem{419}{เหตุยา ปญฺจ, อารมฺมเณ เทฺว, อธิปติยา เอกํ, อนนฺตเร เทฺว, สมนนฺตเร เทฺว, สหชาเต ปญฺจ, อญฺญมญฺเญ เทฺว, นิสฺสเย ปญฺจ, อุปนิสฺสเย เทฺว, ปุเรชาเต เทฺว, อาเสวเน เอกํ, กมฺเม ปญฺจ, วิปาเก ปญฺจ, อาหาเร ปญฺจ, อินฺทฺริเย ปญฺจ, ฌาเน ปญฺจ, มคฺเค ปญฺจ, สมฺปยุตฺเต เทฺว, วิปฺปยุตฺเต ปญฺจ, อตฺถิยา ปญฺจ, นตฺถิยา เทฺว, วิคเต เทฺว, อวิคเต ปญฺจ.}

\csromanheader{2}{2. ปจฺจยปจฺจนีย 1. วิภงฺควาร}
\csromanheader{2}{\textbf{นเหตุ}}
\proseitem{420}{อุปาทินฺนํ ธมฺมํ ปฏิจฺจ อุปาทินฺโน ธมฺโม อุปฺปชฺชติ นเหตุปจฺจยา.\csromanspacer{} อเหตุกํ อุปาทินฺนํ เอกํ ขนฺธํ ปฏิจฺจ ตโย ขนฺธา ฯเปฯ เทฺว ขนฺเธ ฯเปฯ.\csromanspacer{} อเหตุกปฏิสนฺธิกฺขเณ อุปาทินฺนํ เอกํ ขนฺธํ ปฏิจฺจ ตโย ขนฺธา กฏตฺตา จ รูปํ ฯเปฯ เทฺว ขนฺเธ ฯเปฯ.\csromanspacer{} ขนฺเธ ปฏิจฺจ วตฺถุ, วตฺถุํ ปฏิจฺจ ขนฺธา.\csromanspacer{} เอกํ มหาภูตํ ปฏิจฺจ ฯเปฯ มหาภูเต ปฏิจฺจ กฏตฺตารูปํ, อุปาทารูปํ.\csromanspacer{} อสญฺญสตฺตานํ เอกํ มหาภูตํ ปฏิจฺจ ฯเปฯ มหาภูเต ปฏิจฺจ กฏตฺตารูปํ, อุปาทารูปํ. (1)}

\prose{อุปาทินฺนํ ธมฺมํ ปฏิจฺจ อนุปาทินฺโน ธมฺโม อุปฺปชฺชติ นเหตุปจฺจยา.\csromanspacer{} อเหตุเก อุปาทินฺเน ขนฺเธ ปฏิจฺจ จิตฺตสมุฏฺฐานํ รูปํ.}

\prose{อุปาทินฺนํ ธมฺมํ ปฏิจฺจ อุปาทินฺโน จ อนุปาทินฺโน จ ธมฺมา อุปฺปชฺชนฺติ นเหตุปจฺจยา.\csromanspacer{} อเหตุกํ อุปาทินฺนํ เอกํ ขนฺธํ ปฏิจฺจ ตโย ขนฺธา จิตฺตสมุฏฺฐานญฺจ รูปํ ฯเปฯ เทฺว ขนฺเธ ฯเปฯ. (3)}

\csromanheader{2}{68. อุปาทินฺนทุก}
\proseitem{421}{\csromanfolio{539}อนุปาทินฺนํ ธมฺมํ ปฏิจฺจ อนุปาทินฺโน ธมฺโม อุปฺปชฺชติ นเหตุปจฺจยา.\csromanspacer{} อเหตุกํ อนุปาทินฺนํ เอกํ ขนฺธํ ปฏิจฺจ ตโย ขนฺธา จิตฺตสมุฏฺฐานญฺจ รูปํ ฯเปฯ เทฺว ขนฺเธ ฯเปฯ.\csromanspacer{} เอกํ มหาภูตํ ฯเปฯ มหาภูเต ปฏิจฺจ จิตฺตสมุฏฺฐานํ รูปํ, อุปาทารูปํ.\csromanspacer{} พาหิรํ.\csromanspacer{} อาหารสมุฏฺฐานํ.\csromanspacer{} อุตุสมุฏฺฐานํ เอกํ มหาภูตํ ฯเปฯ มหาภูเต ปฏิจฺจ อุปาทารูปํ.\csromanspacer{} วิจิกิจฺฉาสหคเต อุทฺธจฺจสหคเต ขนฺเธ ปฏิจฺจ วิจิกิจฺฉาสหคโต อุทฺธจฺจสหคโต โมโห. (1)}

\prose{อุปาทินฺนญฺจ อนุปาทินฺนญฺจ ธมฺมํ ปฏิจฺจ อนุปาทินฺโน ธมฺโม อุปฺปชฺชติ นเหตุปจฺจยา.\csromanspacer{} อเหตุเก อุปาทินฺเน ขนฺเธ จ มหาภูเต จ ปฏิจฺจ จิตฺตสมุฏฺฐานํ รูปํ. (1)}

\csromanheader{2}{\textbf{น-อารมฺมณ}}
\proseitem{422}{อุปาทินฺนํ ธมฺมํ ปฏิจฺจ อุปาทินฺโน ธมฺโม อุปฺปชฺชติ น-อารมฺมณปจฺจยา.\csromanspacer{} ปฏิสนฺธิกฺขเณ อุปาทินฺเน ขนฺเธ ปฏิจฺจ กฏตฺตารูปํ.\csromanspacer{} ขนฺเธ ปฏิจฺจ วตฺถุ.\csromanspacer{} เอกํ มหาภูตํ ฯเปฯ มหาภูเต ปฏิจฺจ กฏตฺตารูปํ, อุปาทารูปํ.\csromanspacer{} อสญฺญสตฺตานํ เอกํ มหาภูตํ ฯเปฯ มหาภูเต ปฏิจฺจ กฏตฺตารูปํ, อุปาทารูปํ. (1)}

\prose{อุปาทินฺนํ ธมฺมํ ปฏิจฺจ อนุปาทินฺโน ธมฺโม อุปฺปชฺชติ น-อารมฺมณปจฺจยา.\csromanspacer{} อุปาทินฺเน ขนฺเธ ปฏิจฺจ จิตฺตสมุฏฺฐานํ รูปํ.}

\prose{อนุปาทินฺนํ ธมฺมํ ปฏิจฺจ อนุปาทินฺโน ธมฺโม อุปฺปชฺชติ น-อารมฺมณปจฺจยา.\csromanspacer{} อนุปาทินฺเน ขนฺเธ ปฏิจฺจ จิตฺตสมุฏฺฐานํ รูปํ.\csromanspacer{} เอกํ มหาภูตํ ฯเปฯ.\csromanspacer{} พาหิรํ.\csromanspacer{} อาหารสมุฏฺฐานํ.\csromanspacer{} อุตุสมุฏฺฐานํ เอกํ มหาภูตํ ฯเปฯ มหาภูเต ปฏิจฺจ อุปาทารูปํ. (1)}

\prose{อุปาทินฺนญฺจ อนุปาทินฺนญฺจ ธมฺมํ ปฏิจฺจ อนุปาทินฺโน ธมฺโม อุปฺปชฺชติ น-อารมฺมณปจฺจยา.\csromanspacer{} อุปาทินฺเน ขนฺเธ จ มหาภูเต จ ปฏิจฺจ จิตฺตสมุฏฺฐานํ รูปํ. (1)}

\csromanheader{2}{\textbf{น-อธิปตฺยาทิ}}
\proseitem{423}{อุปาทินฺนํ ธมฺมํ ปฏิจฺจ อุปาทินฺโน ธมฺโม อุปฺปชฺชติ น-อธิปติปจฺจยา.\csromanspacer{} น-อนนฺตรปจฺจยา.\csromanspacer{} นสมนนฺตรปจฺจยา.\csromanspacer{} น-อญฺญมญฺญปจฺจยา.\csromanspacer{} น-อุปนิสฺสยปจฺจยา.}

\csromanheader{2}{\textbf{นปุเรชาตาทิ}}
\proseitem{424}{อุปาทินฺนํ ธมฺมํ ปฏิจฺจ อุปาทินฺโน ธมฺโม อุปฺปชฺชติ นปุเรชาตปจฺจยา.\csromanspacer{} อรูเป อุปาทินฺนํ เอกํ ขนฺธํ ปฏิจฺจ ตโย ขนฺธา ฯเปฯ เทฺว ขนฺเธ ฯเปฯ. \csromanfolio{540}ปฏิสนฺธิกฺขเณ อุปาทินฺนํ เอกํ ขนฺธํ ปฏิจฺจ ตโย ขนฺธา กฏตฺตา จ รูปํ ฯเปฯ เทฺว ขนฺเธ ฯเปฯ. (ยาว อสญฺญสตฺตา.) (1)}

\prose{อุปาทินฺนํ ธมฺมํ ปฏิจฺจ อนุปาทินฺโน ธมฺโม อุปฺปชฺชติ นปุเรชาตปจฺจยา.\csromanspacer{} อุปาทินฺเน ขนฺเธ ปฏิจฺจ จิตฺตสมุฏฺฐานํ รูปํ. (2)}

\prose{อนุปาทินฺนํ ธมฺมํ ปฏิจฺจ อนุปาทินฺโน ธมฺโม อุปฺปชฺชติ นปุเรชาตปจฺจยา.\csromanspacer{} อรูเป อนุปาทินฺนํ เอกํ ขนฺธํ ปฏิจฺจ ตโย ขนฺธา ฯเปฯ เทฺว ขนฺเธ ฯเปฯ.\csromanspacer{} อนุปาทินฺเน ขนฺเธ ปฏิจฺจ จิตฺตสมุฏฺฐานํ รูปํ.\csromanspacer{} เอกํ มหาภูตํ ฯเปฯ. (ยาว อสญฺญสตฺตา.) (1)}

\prose{อุปาทินฺนญฺจ อนุปาทินฺนญฺจ ธมฺมํ ปฏิจฺจ อนุปาทินฺโน ธมฺโม อุปฺปชฺชติ นปุเรชาตปจฺจยา.\csromanspacer{} อุปาทินฺเน ขนฺเธ จ มหาภูเต จ ปฏิจฺจ จิตฺตสมุฏฺฐานํ รูปํ.\csromanspacer{} นปจฺฉาชาตปจฺจยา.\csromanspacer{} น-อาเสวนปจฺจยา. (1)}

\csromanheader{2}{\textbf{นกมฺม}}
\proseitem{425}{อนุปาทินฺนํ ธมฺมํ ปฏิจฺจ อนุปาทินฺโน ธมฺโม อุปฺปชฺชติ นกมฺมปจฺจยา.\csromanspacer{} อนุปาทินฺเน ขนฺเธ ปฏิจฺจ อนุปาทินฺนา เจตนา.\csromanspacer{} พาหิรํ.\csromanspacer{} อาหารสมุฏฺฐานํ อุตุสมุฏฺฐานํ ฯเปฯ.\csromanspacer{} มหาภูเต ปฏิจฺจ อุปาทารูปํ. (1)}

\csromanheader{2}{\textbf{นวิปาก}}
\proseitem{426}{อุปาทินฺนํ ธมฺมํ ปฏิจฺจ อุปาทินฺโน ธมฺโม อุปฺปชฺชติ นวิปากปจฺจยา.\csromanspacer{} อสญฺญสตฺตานํ เอกํ มหาภูตํ ฯเปฯ มหาภูเต ปฏิจฺจ กฏตฺตารูปํ, อุปาทารูปํ. (1)}

\prose{อนุปาทินฺนํ ธมฺมํ ปฏิจฺจ อนุปาทินฺโน ธมฺโม อุปฺปชฺชติ นวิปากปจฺจยา.\csromanspacer{} อนุปาทินฺนํ เอกํ ขนฺธํ ปฏิจฺจ ตโย ขนฺธา จิตฺตสมุฏฺฐานญฺจ รูปํ ฯเปฯ เทฺว ขนฺเธ ฯเปฯ.\csromanspacer{} เอกํ มหาภูตํ ฯเปฯ. (ยาว อุตุสมุฏฺฐานํ.)}

\csromanheader{2}{\textbf{น-อาหาร}}
\proseitem{427}{อุปาทินฺนํ ธมฺมํ ปฏิจฺจ อุปาทินฺโน ธมฺโม อุปฺปชฺชติ น-อาหารปจฺจยา.\csromanspacer{} อสญฺญสตฺตานํ เอกํ มหาภูตํ. (1)}

\prose{อนุปาทินฺนํ ธมฺมํ ปฏิจฺจ อนุปาทินฺโน ธมฺโม อุปฺปชฺชติ น-อาหารปจฺจยา.\csromanspacer{} พาหิรํ.\csromanspacer{} อุตุสมุฏฺฐานํ ฯเปฯ. (1)}

\csromanheader{2}{68. อุปาทินฺนทุก \textbf{น-อินฺทฺริย}}
\proseitem{428}{\csromanfolio{541}อุปาทินฺนํ ธมฺมํ ปฏิจฺจ อุปาทินฺโน ธมฺโม อุปฺปชฺชติ น-อินฺทฺริยปจฺจยา.\csromanspacer{} อสญฺญสตฺตานํ มหาภูเต ปฏิจฺจ รูปชีวิตินฺทฺริยํ. (1)}

\prose{อนุปาทินฺนํ ธมฺมํ ปฏิจฺจ อนุปาทินฺโน ธมฺโม อุปฺปชฺชติ น-อินฺทฺริยปจฺจยา.\csromanspacer{} พาหิรํ.\csromanspacer{} อาหารสมุฏฺฐานํ, อุตุสมุฏฺฐานํ ฯเปฯ. (1)}

\csromanheader{2}{\textbf{นฌานาทิ}}
\proseitem{429}{อุปาทินฺนํ ธมฺมํ ปฏิจฺจ อุปาทินฺโน ธมฺโม อุปฺปชฺชติ นฌานปจฺจยา.\csromanspacer{} ปญฺจวิญฺญาณสหคตํ เอกํ ขนฺธํ ปฏิจฺจ ตโย ขนฺธา ฯเปฯ เทฺว ขนฺเธ ฯเปฯ.\csromanspacer{} อสญฺญสตฺตานํ ฯเปฯ.}

\prose{อนุปาทินฺนํ ธมฺมํ ปฏิจฺจ อนุปาทินฺโน ธมฺโม อุปฺปชฺชติ นฌานปจฺจยา.\csromanspacer{} พาหิรํ.\csromanspacer{} อาหารสมุฏฺฐานํ.\csromanspacer{} อุตุสมุฏฺฐานํ ฯเปฯ.}

\prose{อุปาทินฺนํ ธมฺมํ ปฏิจฺจ อุปาทินฺโน ธมฺโม อุปฺปชฺชติ นมคฺคปจฺจยา. (นเหตุสทิสํ.\csromanspacer{} โมโห นตฺถิ.) นสมฺปยุตฺตปจฺจยา.}

\csromanheader{2}{\textbf{นวิปฺปยุตฺตาทิ}}
\proseitem{430}{อุปาทินฺนํ ธมฺมํ ปฏิจฺจ อุปาทินฺโน ธมฺโม อุปฺปชฺชติ นวิปฺปยุตฺตปจฺจยา.\csromanspacer{} อรูเป อุปาทินฺนํ เอกํ ขนฺธํ ปฏิจฺจ ตโย ขนฺธา ฯเปฯ เทฺว ขนฺเธ ฯเปฯ.\csromanspacer{} อสญฺญสตฺตานํ เอกํ มหาภูตํ ปฏิจฺจ ฯเปฯ.}

\prose{อนุปาทินฺนํ ธมฺมํ ปฏิจฺจ อนุปาทินฺโน ธมฺโม อุปฺปชฺชติ นวิปฺปยุตฺตปจฺจยา.\csromanspacer{} อรูเป อนุปาทินฺนํ เอกํ ขนฺธํ ปฏิจฺจ ตโย ขนฺธา ฯเปฯ เทฺว ขนฺเธ ฯเปฯ.\csromanspacer{} พาหิรํ.\csromanspacer{} อาหารสมุฏฺฐานํ.\csromanspacer{} อุตุสมุฏฺฐานํ ฯเปฯ.\csromanspacer{} โนนตฺถิปจฺจยา.\csromanspacer{} โนวิคตปจฺจยา.}

\csromanheader{2}{2. ปจฺจยปจฺจนีย 2. สงฺขฺยาวาร}
\csromanheader{2}{\textbf{สุทฺธ}}
\proseitem{431}{นเหตุยา ปญฺจ, น-อารมฺมเณ จตฺตาริ, น-อธิปติยา ปญฺจ, น-อนนฺตเร จตฺตาริ, นสมนนฺตเร จตฺตาริ, น-อญฺญมญฺเญ จตฺตาริ, น-อุปนิสฺสเย จตฺตาริ, นปุเรชาเต จตฺตาริ, นปจฺฉาชาเต ปญฺจ, น-อาเสวเน ปญฺจ, \csromanfolio{542}นกมฺเม เอกํ, นวิปาเก เทฺว, น-อาหาเร เทฺว, น-อินฺทฺริเย เทฺว, นฌาเน เทฺว, นมคฺเค ปญฺจ, นสมฺปยุตฺเต จตฺตาริ, นวิปฺปยุตฺเต เทฺว, โนนตฺถิยา จตฺตาริ, โนวิคเต จตฺตาริ.}

\csromanheader{2}{3. ปจฺจยานุโลมปจฺจนีย}
\proseitem{432}{\textbf{เหตุ}ปจฺจยา น-อารมฺมเณ จตฺตาริ, น-อธิปติยา ปญฺจ ฯเปฯ นปุเรชาเต จตฺตาริ, นปจฺฉาชาเต ปญฺจ, น-อาเสวเน ปญฺจ, นกมฺเม เอกํ, นวิปาเก เอกํ ฯเปฯ นสมฺปยุตฺเต จตฺตาริ, นวิปฺปยุตฺเต เทฺว, โนนตฺถิยา จตฺตาริ, โนวิคเต จตฺตาริ.}

\csromanheader{2}{4. ปจฺจยปจฺจนียานุโลม}
\proseitem{433}{นเหตุปจฺจยา อารมฺมเณ เทฺว, อนนฺตเร เทฺว, สมนนฺตเร เทฺว, สหชาเต ปญฺจ, อญฺญมญฺเญ เทฺว, นิสฺสเย ปญฺจ, อุปนิสฺสเย เทฺว, ปุเรชาเต เทฺว, อาเสวเน เอกํ, กมฺเม ปญฺจ, วิปาเก ปญฺจ ฯเปฯ มคฺเค เอกํ, สมฺปยุตฺเต เทฺว ฯเปฯ อวิคเต ปญฺจ.}

\vspace{\tipitakaparskip}
\csromancenter{(สหชาตวาโร ปฏิจฺจวารสทิโส.)}
\csromanheader{2}{68. อุปาทินฺนทุก 3. ปจฺจยวาร}
\csromanheader{2}{1. ปจฺจยานุโลม 1. วิภงฺควาร}
\csromanheader{2}{\textbf{เหตุ}}
\proseitem{434}{อุปาทินฺนํ ธมฺมํ ปจฺจยา อุปาทินฺโน ธมฺโม อุปฺปชฺชติ เหตุปจฺจยา.\csromanspacer{} อุปาทินฺนํ เอขํ ขนฺธํ ปจฺจยา ตโย ขนฺธา ฯเปฯ เทฺว ขนฺเธ ฯเปฯ.\csromanspacer{} ปฏิสนฺธิกฺขเณ อุปาทินฺนํ เอกํ ขนฺธํ ปจฺจยา ตโย ขนฺธา กฏตฺตา จ รูปํ ฯเปฯ เทฺว ขนฺเธ ฯเปฯ.\csromanspacer{} ขนฺเธ ปจฺจยา วตฺถุ, วตฺถุํ ปจฺจยา ขนฺธา.\csromanspacer{} เอกํ มหาภูตํ ฯเปฯ มหาภูเต ปจฺจยา กฏตฺตารูปํ, อุปาทารูปํ.\csromanspacer{} วตฺถุํ ปจฺจยา อุปาทินฺนา ขนฺธา. (1)}

\prose{อุปาทินฺนํ ธมฺมํ ปจฺจยา อนุปาทินฺโน ธมฺโม อุปฺปชฺชติ เหตุปจฺจยา.\csromanspacer{} อุปาทินฺเน ขนฺเธ ปจฺจยา จิตฺตสมุฏฺฐานํ รูปํ.\csromanspacer{} วตฺถุํ ปจฺจยา อนุปาทินฺนา ขนฺธา. (2)}

\csromanheader{2}{68. อุปาทินฺนทุก}
\prose{\csromanfolio{543}อุปาทินฺนํ ธมฺมํ ปจฺจยา อุปาทินฺโน จ อนุปาทินฺโน จ ธมฺมา อุปฺปชฺชนฺติ เหตุปจฺจยา.\csromanspacer{} อุปาทินฺนํ เอกํ ขนฺธํ ปจฺจยา ตโย ขนฺธา จิตฺตสมุฏฺฐานญฺจ รูปํ ฯเปฯ เทฺว ขนฺเธ ฯเปฯ. (3)}

\prose{อนุปาทินฺนํ ธมฺมํ ปจฺจยา อนุปาทินฺโน ธมฺโม อุปฺปชฺชติ เหตุปจฺจยา.\csromanspacer{} อนุปาทินฺนํ เอกํ ขนฺธํ ปจฺจยา ตโย ขนฺธา จิตฺตสมุฏฺฐานญฺจ รูปํ ฯเปฯ เทฺว ขนฺเธ ฯเปฯ.\csromanspacer{} เอกํ มหาภูตํ ฯเปฯ มหาภูเต ปจฺจยา จิตฺตสมุฏฺฐานํ รูปํ, อุปาทารูปํ. (1)}

\prose{อุปาทินฺนญฺจ อนุปาทินฺนญฺจ ธมฺมํ ปจฺจยา อนุปาทินฺโน ธมฺโม อุปฺปชฺชติ เหตุปจฺจยา.\csromanspacer{} อุปาทินฺเน ขนฺเธ จ มหาภูเต จ ปจฺจยา จิตฺตสมุฏฺฐานํ รูปํ.\csromanspacer{} อนุปาทินฺนํ เอกํ ขนฺธญฺจ วตฺถุญฺจ ปจฺจยา ตโย ขนฺธา ฯเปฯ เทฺว ขนฺเธ จ ฯเปฯ. (1)}

\csromanheader{2}{\textbf{อารมฺมณ}}
\proseitem{435}{อุปาทินฺนํ ธมฺมํ ปจฺจยา อุปาทินฺโน ธมฺโม อุปฺปชฺชติ อารมฺมณปจฺจยา.\csromanspacer{} อุปาทินฺนํ เอกํ ขนฺธํ ปจฺจยา ตโย ขนฺธา ฯเปฯ เทฺว ขนฺเธ ฯเปฯ.\csromanspacer{} ปฏิสนฺธิกฺขเณ ฯเปฯ.\csromanspacer{} วตฺถุํ ปจฺจยา ขนฺธา.\csromanspacer{} จกฺขายตนํ ปจฺจยา จกฺขุวิญฺญาณํ ฯเปฯ.\csromanspacer{} กายายตนํ ฯเปฯ.\csromanspacer{} วตฺถุํ ปจฺจยา อุปาทินฺนา ขนฺธา. (1)}

\prose{อุปาทินฺนํ ธมฺมํ ปจฺจยา อนุปาทินฺโน ธมฺโม อุปฺปชฺชติ อารมฺมณปจฺจยา.\csromanspacer{} วตฺถุํ ปจฺจยา อนุปาทินฺนา ขนฺธา. (2)}

\prose{อนุปาทินฺนํ ธมฺมํ ปจฺจยา อนุปาทินฺโน ธมฺโม อุปฺปชฺชติ อารมฺมณปจฺจยา.\csromanspacer{} อนุปาทินฺนํ เอกํ ขนฺธํ ปจฺจยา ตโย ขนฺธา ฯเปฯ เทฺว ขนฺเธ ฯเปฯ. (1)}

\prose{อุปาทินฺนญฺจ อนุปาทินฺนญฺจ ธมฺมํ ปจฺจยา อนุปาทินฺโน ธมฺโม อุปฺปชฺชติ อารมฺมณปจฺจยา.\csromanspacer{} อนุปาทินฺนํ เอกํ ขนฺธญฺจ วตฺถุญฺจ ปจฺจยา ตโย ขนฺธา ฯเปฯ เทฺว ขนฺเธ จ ฯเปฯ. (1)}

\csromanheader{2}{\textbf{อธิปติ}}
\proseitem{436}{อุปาทินฺนํ ธมฺมํ ปจฺจยา อนุปาทินฺโน ธมฺโม อุปฺปชฺชติ อธิปติปจฺจยา.\csromanspacer{} วตฺถุํ ปจฺจยา อนุปาทินฺนา ขนฺธา. (1)}

\prose{อนุปาทินฺนํ ธมฺมํ ปจฺจยา อนุปาทินฺโน ธมฺโม อุปฺปชฺชติ อธิปติปจฺจยา.\csromanspacer{} อนุปาทินฺนํ เอกํ ขนฺธํ ปจฺจยา ตโย ขนฺธา จิตฺตสมุฏฺฐานญฺจ รูปํ ฯเปฯ เทฺว ขนฺเธ ฯเปฯ.\csromanspacer{} เอกํ มหาภูตํ ฯเปฯ มหาภูเต ปจฺจยา จิตฺตสมุฏฺฐานํ รูปํ, อุปาทารูปํ. (1)}

\prose{\csromanfolio{544}อุปาทินฺนญฺจ อนุปาทินฺนญฺจ ธมฺมํ ปจฺจยา อนุปาทินฺโน ธมฺโม อุปฺปชฺชติ อธิปติปจฺจยา.\csromanspacer{} อนุปาทินฺนํ เอกํ ขนฺธญฺจ วตฺถุญฺจ ปจฺจยา ตโย ขนฺธา ฯเปฯ เทฺว ขนฺเธ จ ฯเปฯ. (1) (สํขิตฺตํ.)}

\csromanheader{2}{1. ปจฺจยานุโลม 2. สงฺขฺยาวาร}
\csromanheader{2}{\textbf{สุทฺธ}}
\proseitem{437}{เหตุยา ปญฺจ, อารมฺมเณ จตฺตาริ, อธิปติยา ตีณิ, อนนฺตเร จตฺตาริ, สมนนฺตเร จตฺตาริ, สหชาเต ปญฺจ, อญฺญมญฺเญ จตฺตาริ, นิสฺสเย ปญฺจ, อุปนิสฺสเย จตฺตาริ, ปุเรชาเต จตฺตาริ, อาเสวเน ตีณิ, กมฺเม ปญฺจ, วิปาเก ปญฺจ ฯเปฯ อวิคเต ปญฺจ.}

\csromanheader{2}{2. ปจฺจยปจฺจนีย 1. วิภงฺควาร}
\csromanheader{2}{\textbf{นเหตุ}}
\proseitem{438}{อุปาทินฺนํ ธมฺมํ ปจฺจยา อุปาทินฺโน ธมฺโม อุปฺปชฺชติ นเหตุปจฺจยา.\csromanspacer{} อเหตุกํ อุปาทินฺนํ เอกํ ขนฺธํ ปจฺจยา ตโย ขนฺธา ฯเปฯ เทฺว ขนฺเธ ฯเปฯ.\csromanspacer{} อเหตุกปฏิสนฺธิกฺขเณ ฯเปฯ.\csromanspacer{} ขนฺเธ ปจฺจยา วตฺถุ, วตฺถุํ ปจฺจยา ขนฺธา.\csromanspacer{} เอกํ มหาภูตํ ฯเปฯ มหาภูเต ปจฺจยา กฏตฺตารูปํ, อุปาทารูปํ.\csromanspacer{} อสญฺญสตฺตานํ เอกํ มหาภูตํ ฯเปฯ.\csromanspacer{} จกฺขายตนํ ปจฺจยา จกฺขุวิญฺญาณํ ฯเปฯ.\csromanspacer{} กายายตนํ ฯเปฯ.\csromanspacer{} วตฺถุํ ปจฺจยา อเหตุกา อุปาทินฺนา ขนฺธา. (1)}

\prose{อุปาทินฺนํ ธมฺมํ ปจฺจยา อนุปาทินฺโน ธมฺโม อุปฺปชฺชติ นเหตุปจฺจยา.\csromanspacer{} อเหตุเก อุปาทินฺเน ขนฺเธ ปจฺจยา จิตฺตสมุฏฺฐานํ รูปํ.\csromanspacer{} วตฺถุํ ปจฺจยา อเหตุกา อนุปาทินฺนา ขนฺธา.\csromanspacer{} วตฺถุํ ปจฺจยา วิจิกิจฺฉาสหคโต อุทฺธจฺจสหคโต โมโห. (2)}

\prose{อุปาทินฺนํ ธมฺมํ ปจฺจยา อุปาทินฺโน จ อนุปาทินฺโน จ ธมฺมา อุปฺปชฺชนฺติ นเหตุปจฺจยา.\csromanspacer{} อเหตุกํ อุปาทินฺนํ เอกํ ขนฺธํ ปจฺจยา ตโย ขนฺธา จิตฺตสมุฏฺฐานญฺจ รูปํ ฯเปฯ เทฺว ขนฺเธ ฯเปฯ. (3)}

\proseitem{439}{อนุปาทินฺนํ ธมฺมํ ปจฺจยา อนุปาทินฺโน ธมฺโม อุปฺปชฺชติ นเหตุปจฺจยา.\csromanspacer{} อเหตุกํ อนุปาทินฺนํ เอกํ ขนฺธํ ปจฺจยา ตโย ขนฺธา จิตฺตสมุฏฺฐานญฺจ รูปํ ฯเปฯ}

\vspace{\tipitakaparskip}
\csromancenter{อุปาทินฺนทุก เทฺว ขนฺเธ ฯเปฯ.\csromanspacer{} เอกํ มหาภูตํ ฯเปฯ.\csromanspacer{} พาหิรํ.\csromanspacer{} อาหารสมุฏฺฐานํ.\csromanspacer{} อุตุสมุฏฺฐานํ ฯเปฯ.\csromanspacer{} วิจิกิจฺฉาสหคเต อุทฺธจฺจสหคเต ขนฺเธ ปจฺจยา วิจิกิจฺฉาสหคโต อุทฺธจฺจสหคโต โมโห. (1)}
\prose{\csromanfolio{545}อุปาทินฺนญฺจ อนุปาทินญฺจ ธมฺมํ ปจฺจยา อนุปาทินฺโน ธมฺโม อุปฺปชฺชติ นเหตุปจฺจยา.\csromanspacer{} อเหตุเก อุปาทินฺเน ขนฺเธ จ มหาภูเต จ ปจฺจยา จิตฺตสมุฏฺฐานํ รูปํ.\csromanspacer{} อเหตุกํ อนุปาทินฺนํ เอกํ ขนฺธญฺจ วตฺถุญฺจ ปจฺจยา ตโย ขนฺธา ฯเปฯ เทฺว ขนฺเธ จ ฯเปฯ.\csromanspacer{} วิจิกิจฺฉาสหคเต อุทฺธจฺจสหคเต ขนฺเธ จ วตฺถุญฺจ ปจฺจยา วิจิกิจฺฉาสหคโต อุทฺธจฺจสหคโต โมโห. (1) (สํขิตฺตํ.)}

\csromanheader{2}{2. ปจฺจยปจฺจนีย 2. สงฺขฺยาวาร}
\csromanheader{2}{\textbf{สุทฺธ}}
\proseitem{440}{นเหตุยา ปญฺจ, น-อารมฺมเณ จตฺตาริ, น-อธิปติยา ปญฺจ, น-อนนฺตเร จตฺตาริ, นสมนนฺตเร จตฺตาริ, น-อญฺญมญฺเญ จตฺตาริ, น-อุปนิสฺสเย จตฺตาริ, นปุเรชาเต จตฺตาริ, นปจฺฉาชาเต ปญฺจ, น-อาเสวเน ปญฺจ, นกมฺเม ตีณิ, นวิปาเก จตฺตาริ, น-อาหาเร เทฺว, น-อินฺทฺริเย เทฺว, นฌาเน เทฺว, นมคฺเค ปญฺจ, นสมฺปยุตฺเต จตฺตาริ, นวิปฺปยุตฺเต เทฺว, โนนตฺถิยา จตฺตาริ, โนวิคเต จตฺตาริ.}

\csromanheader{2}{3. ปจฺจยานุโลมปจฺจนีย}
\proseitem{441}{เหตุปจฺจยา น-อารมฺมเณ จตฺตาริ, น-อธิปติยา ปญฺจ ฯเปฯ นปุเรชาเต จตฺตาริ, นปจฺฉาชาเต ปญฺจ, น-อาเสวเน ปญฺจ, นกมฺเม ตีณิ, นวิปาเก ตีณิ ฯเปฯ นสมฺปยุตฺเต จตฺตาริ, นวิปฺปยุตฺเต เทฺว, โนนตฺถิยา จตฺตาริ, โนวิคเต จตฺตาริ.}

\csromanheader{2}{4. ปจฺจยาปจฺจนียานุโลม}
\proseitem{442}{นเหตุปจฺจยา อารมฺมเณ จตฺตาริ, อนนฺตเร จตฺตาริ ฯเปฯ มคฺเค ตีณิ ฯเปฯ อวิคเต ปญฺจ.}

\vspace{\tipitakaparskip}
\csromancenter{(นิสฺสยวาโร ปจฺจยวารสทิโส.)}
\csromanheader{2}{68. อุปาทินฺนทุก 5. สํสฏฺฐวาร}
\csromanheader{2}{1. ปจฺจยานุโลมาทิ}
\proseitem{443}{\csromanfolio{546}อุปาทินฺนํ ธมฺมํ สํสฏฺโฐ อุปาทินฺโน ธมฺโม อุปฺปชฺชติ เหตุปจฺจยา.\csromanspacer{} อุปาทินฺนํ เอกํ ขนฺธํ สํสฏฺฐา ตโย ขนฺธา ฯเปฯ เทฺว ขนฺเธ ฯเปฯ.\csromanspacer{} ปฏิสนฺธิกฺขเณ ฯเปฯ.}

\prose{อนุปาทินฺนํ ธมฺมํ สํสฏฺโฐ อนุปาทินฺโน ธมฺโม อุปฺปชฺชติ เหตุปจฺจยา.\csromanspacer{} อนุปาทินฺนํ เอกํ ขนฺธํ สํสฏฺฐา ตโย ขนฺธา ฯเปฯ เทฺว ขนฺเธ ฯเปฯ. (1)}

\prose{เหตุยา เทฺว, อารมฺมเณ เทฺว, อธิปติยา เอกํ, อนนฺตเร เทฺว, สมนนฺตเร เทฺว, สหชาเต เทฺว, อญฺญมญฺเญ เทฺว, นิสฺสเย เทฺว, อุปนิสฺสเย เทฺว, ปุเรชาเต เทฺว, อาเสวเน เอกํ, กมฺเม เทฺว, วิปาเก เทฺว ฯเปฯ อวิคเต เทฺว.\csromanspacer{} อนุโลมํ.}

\proseitem{444}{อุปาทินฺนํ ธมฺมํ สํสฏฺโฐ อุปาทินฺโน ธมฺโม อุปฺปชฺชติ นเหตุปจฺจยา.\csromanspacer{} อเหตุกํ อุปาทินฺนํ เอกํ ขนฺธํ สํสฏฺฐา ตโย ขนฺธา ฯเปฯ เทฺว ขนฺเธ ฯเปฯ.\csromanspacer{} อเหตุกปฏิสนฺธิกฺขเณ ฯเปฯ. (1)}

\prose{อนุปาทินฺนํ ธมฺมํ สํสฏฺโฐ อนุปาทินฺโน ธมฺโม อุปฺปชฺชติ นเหตุปจฺจยา.\csromanspacer{} อเหตุกํ อนุปาทินฺนํ เอกํ ขนฺธํ สํสฏฺฐา ตโย ขนฺธา ฯเปฯ เทฺว ขนฺเธ ฯเปฯ.\csromanspacer{} วิจิกิจฺฉาสหคเต อุทฺธจฺจสหคเต ขนฺเธ สํสฏฺโฐ วิจิกิจฺฉาสหคโต อุทฺธจฺจสหคโต โมโห. (1)}

\prose{นเหตุยา เทฺว, น-อธิปติยา เทฺว, นปุเรชาเต เทฺว, นปจฺฉาชาเต เทฺว, น-อาเสวเน เทฺว, นกมฺเม เอกํ, นวิปาเก เอกํ, นฌาเน เอกํ, นมคฺเค เทฺว, นวิปฺปยุตฺเต เทฺว.}

\vspace{\tipitakaparskip}
\csromancenter{ปจฺจนียํ.}
\csromancenter{(เอวํ อิตเร เทฺว คณนาปิ สมฺปยุตฺตวาโรปิ กาตพฺโพ.)}
\csromanheader{2}{68. อุปาทินฺนทุก \textbf{68. อุปาทินฺนทุก 7. ปญฺหาวาร}}
\csromanheader{2}{1. ปจฺจยานุโลม 1. วิภงฺควาร}
\csromanheader{2}{\textbf{เหตุ}}
\proseitem{445}{\csromanfolio{547}อุปาทินฺโน ธมฺโม อุปาทินฺนสฺส ธมฺมสฺส เหตุปจฺจเยน ปจฺจโย.\csromanspacer{} อุปาทินฺนา เหตู สมฺปยุตฺตกานํ ขนฺธานํ เหตุปจฺจเยน ปจฺจโย.\csromanspacer{} ปฏิสนฺธิกฺขเณ อุปาทินฺนา เหตู สมฺปยุตฺตกานํ ขนฺธานํ กฏตฺตา จ รูปานํ เหตุปจฺจเยน ปจฺจโย. (1)}

\prose{อุปาทินฺโน ธมฺโม อนุปาทินฺนสฺส ธมฺมสฺส เหตุปจฺจเยน ปจฺจโย.\csromanspacer{} อุปาทินฺนา เหตู จิตฺตสมุฏฺฐานานํ รูปานํ เหตุปจฺจเยน ปจฺจโย. (2)}

\prose{อุปาทินฺโน ธมฺโม อุปาทินฺนสฺส จ อนุปาทินฺนสฺส จ ธมฺมสฺส เหตุปจฺจเยน ปจฺจโย.\csromanspacer{} อุปาทินฺนา เหตู สมฺปยุตฺตกานํ ขนฺธานํ จิตฺตสมุฏฺฐานานญฺจ รูปานํ เหตุปจฺจเยน ปจฺจโย. (3)}

\prose{อนุปาทินฺโน ธมฺโม อนุปาทินฺนสฺส ธมฺมสฺส เหตุปจฺจเยน ปจฺจโย.\csromanspacer{} อนุปาทินฺนา เหตู สมฺปยุตฺตกานํ ขนฺธานํ จิตฺตสมุฏฺฐานานญฺจ รูปานํ เหตุปจฺจเยน ปจฺจโย. (1)}

\csromanheader{2}{\textbf{อารมฺมณ}}
\proseitem{446}{อุปาทินฺโน ธมฺโม อุปาทินฺนสฺส ธมฺมสฺส อารมฺมณปจฺจเยน ปจฺจโย.\csromanspacer{} จกฺขุํ ฯเปฯ.\csromanspacer{} กายํ อุปาทินฺเน รูเป.\csromanspacer{} คนฺเธ.\csromanspacer{} รเส.\csromanspacer{} โผฏฺฐพฺเพ.\csromanspacer{} วตฺถุํ อุปาทินฺเน ขนฺเธ อนิจฺจโต ฯเปฯ โทมนสฺสํ อุปฺปชฺชติ.\csromanspacer{} กุสลากุสเล นิรุทฺเธ วิปาโก ตทารมฺมณตา อุปฺปชฺชติ.\csromanspacer{} อุปาทินฺนํ รูปายตนํ จกฺขุวิญฺญาณสฺส ฯเปฯ.\csromanspacer{} อุปาทินฺนํ คนฺธายตนํ ฆานวิญฺญาณสฺส ฯเปฯ.\csromanspacer{} ผฏฺฐพฺพายตนํ กายวิญฺญาณสฺส อารมฺมณปจฺจเยน ปจฺจโย. (1)}

\prose{อุปาทินฺโน ธมฺโม อนุปาทินฺนสฺส ธมฺมสฺส อารมฺมณปจฺจเยน ปจฺจโย.\csromanspacer{} จกฺขุํ ฯเปฯ.\csromanspacer{} กายํ อุปาทินฺเน รูเป.\csromanspacer{} คนฺเธ.\csromanspacer{} รเส.\csromanspacer{} โผฏฺฐพฺเพ.\csromanspacer{} วตฺถุํ อุปาทินฺเน ขนฺเธ อนิจฺจโต ฯเปฯ โทมนสฺสํ อุปฺปชฺชติ.\csromanspacer{} ทิพฺเพน จกฺขุนา อุปาทินฺนํ รูปํ ปสฺสติ.\csromanspacer{} เจโตปริยญาเณน อุปาทินฺนจิตฺตสมงฺคิสฺส จิตฺตํ ชานาติ. \csromanfolio{548}อุปาทินฺนา ขนฺธา อิทฺธิวิธญาณสฺส, เจโตปริยญาณสฺส, ปุพฺเพนิวาสานุสฺสติญาณสฺส, อนาคตํสญาณสฺส, อาวชฺชนาย อารมฺมณปจฺจเยน ปจฺจโย. (2)}

\proseitem{447}{อนุปาทินฺโน ธมฺโม อนุปาทินฺนสฺส ธมฺมสฺส อารมฺมณปจฺจเยน ปจฺจโย.\csromanspacer{} ทานํ ฯเปฯ สีลํ ฯเปฯ อุโปสถกมฺมํ กตฺวา ตํ ปจฺจเวกฺขติ อสฺสาเทติ อภินนฺทติ, ตํ อารพฺภ ราโค ฯเปฯ โทมนสฺสํ อุปฺปชฺชติ.\csromanspacer{} ปุพฺเพ สุจิณฺณานิ ฯเปฯ.\csromanspacer{} ฌานา ฯเปฯ.\csromanspacer{} อริยา มคฺคา วุฏฺฐหิตฺวา มคฺคํ ปจฺจเวกฺขนฺติ.\csromanspacer{} ผลํ ฯเปฯ.\csromanspacer{} นิพฺพานํ ฯเปฯ.\csromanspacer{} นิพฺพานํ โคตฺรภุสฺส, โวทานสฺส, มคฺคสฺส, ผลสฺส, อาวชฺชนาย อารมฺมณปจฺจเยน ปจฺจโย.\csromanspacer{} อริยา ปหีเน กิเลเส ฯเปฯ.\csromanspacer{} วิกฺขมฺภิเต กิเลเส ฯเปฯ.\csromanspacer{} ปุพฺเพ ฯเปฯ.\csromanspacer{} อนุปาทินฺเน รูเป.\csromanspacer{} สทฺเท.\csromanspacer{} คนฺเธ.\csromanspacer{} รเส.\csromanspacer{} โผฏฺฐพฺเพ อนุปาทินฺเน ขนฺเธ อนิจฺจโต ฯเปฯ โทมนสฺสํ อุปฺปชฺชติ.\csromanspacer{} ทิพฺเพน จกฺขุนา อนุปาทินฺนํ รูปํ ปสฺสติ.\csromanspacer{} ทิพฺพาย โสตธาตุยา สทฺทํ สุณาติ.\csromanspacer{} เจโตปริยญาเณน อนุปาทินฺนจิตฺตสมงฺคิสฺส จิตฺตํ ชานาติ.\csromanspacer{} อากาสานญฺจายตนํ ฯเปฯ.\csromanspacer{} อากิญฺจญฺญายตนํ ฯเปฯ.\csromanspacer{} อนุปาทินฺนา ขนฺธา อิทฺธิวิธญาณสฺส, เจโตปริยญาณสฺส, ปุพฺเพนิวาสานุสฺสติญาณสฺส, ยถากมฺมูปคญาณสฺส, อนาคตํสญาณสฺส, อาวชฺชนาย อารมฺมณปจฺจเยน ปจฺจโย. (1)}

\prose{อนุปาทินฺโน ธมฺโม อุปาทินฺนสฺส ธมฺมสฺส อารมฺมณปจฺจเยน ปจฺจโย.\csromanspacer{} อนุปาทินฺเน รูเป.\csromanspacer{} สทฺเท.\csromanspacer{} คนฺเธ.\csromanspacer{} รเส.\csromanspacer{} โผฏฺฐพฺเพ อนุปาทินฺเน ขนฺเธ อนิจฺจโต ฯเปฯ โทมนสฺสํ อุปฺปชฺชติ.\csromanspacer{} กุสลากุสเล นิรุทฺเธ วิปาโก ตทารมฺมณตา อุปฺปชฺชติ.\csromanspacer{} อากาสานญฺจายตนกุสลํ วิญฺญาณญฺจายตนวิปากสฺส ฯเปฯ.\csromanspacer{} อากิญฺจญฺญายตนกุสลํ ฯเปฯ.\csromanspacer{} อนุปาทินฺนํ รูปายตนํ จกฺขุวิญฺญาณสฺส ฯเปฯ.\csromanspacer{} สทฺทายตนํ ฯเปฯ.\csromanspacer{} โผฏฺฐพฺพายตนํ กายวิญฺญาณสฺส อารมฺมณปจฺจเยน ปจฺจโย. (2)}

\csromanheader{2}{\textbf{อธิปติ}}
\proseitem{448}{อุปาทินฺโน ธมฺโม อนุปาทินฺนสฺส ธมฺมสฺส อธิปติปจฺจเยน ปจฺจโย.\csromanspacer{} \textbf{อารมฺมณาธิปติ}\csromandash{}จกฺขุํ ฯเปฯ กายํ.\csromanspacer{} อุปาทินฺเน รูเป.\csromanspacer{} คนฺเธ.\csromanspacer{} รเส.\csromanspacer{} โผฏฺฐพฺเพ.\csromanspacer{} วตฺถุํ อุปาทินฺเน ขนฺเธ ครุํ กตฺวา อสฺสาเทติ อภินนฺทติ, ตํ ครุํ กตฺวา ราโค อุปฺปชฺชติ.\csromanspacer{} ทิฏฺฐิ อุปฺปชฺชติ. (1)}

\csromanheader{2}{68. อุปาทินฺนทุก}
\prose{\csromanfolio{549}อนุปาทินฺโน ธมฺโม อนุปาทินฺนสฺส ธมฺมสฺส อธิปติปจฺจเยน ปจฺจโย.\csromanspacer{} \textbf{อารมฺมณาธิปติ}, สหชาตาธิปติ.\csromanspacer{}\csromanspacer{}\textbf{อารมฺมณาธิปติ}\csromandash{}ทานํ ฯเปฯ สีลํ ฯเปฯ อุโปสถกมฺมํ กตฺวา ตํ ครุํ กตฺวา ปจฺจเวกฺขติ, อสฺสาเทติ อภินนฺทติ, ตํ ครุํ กตฺวา ราโค อุปฺปชฺชติ.\csromanspacer{} ทิฏฺฐิ อุปฺปชฺชติ.\csromanspacer{} ปุพฺเพ ฯเปฯ.\csromanspacer{} ฌานา ฯเปฯ.\csromanspacer{} อริยา มคฺคา วุฏฺฐหิตฺวา มคฺคํ ครุํ กตฺวา ปจฺจเวกฺขนฺติ, ผลํ ฯเปฯ.\csromanspacer{} นิพฺพานํ ครุํ กตฺวา ปจฺจเวกฺขนฺติ.\csromanspacer{} นิพฺพานํ โคตฺรภุสฺส, โวทานสฺส, มคฺคสฺส, ผลสฺส อธิปติปจฺจเยน ปจฺจโย.\csromanspacer{} อนุปาทินฺเน รูเป.\csromanspacer{} สทฺเท.\csromanspacer{} คนฺเธ.\csromanspacer{} รเส.\csromanspacer{} โผฏฺฐพฺเพ อนุปาทินฺเน ขนฺเธ ครุํ กตฺวา อสฺสาเทติ อภินนฺทติ, ตํ ครุํ กตฺวา ราโค อุปฺปชฺชติ.\csromanspacer{} ทิฏฺฐิ อุปฺปชฺชติ.\csromanspacer{} \textbf{สหชาตาธิปติ}\csromandash{}อนุปาทินฺนาธิปติ สมฺปยุตฺตกานํ ขนฺธานํ จิตฺตสมุฏฺฐานานญฺจ รูปานํ อธิปติปจฺจเยน ปจฺจโย. (1)}

\csromanheader{2}{\textbf{อนนฺตร}}
\proseitem{449}{อุปาทินฺโน ธมฺโม อุปาทินฺนสฺส ธมฺมสฺส อนนฺตรปจฺจเยน ปจฺจโย.\csromanspacer{} ปุริมา ปุริมา อุปาทินฺนา ขนฺธา ปจฺฉิมานํ ปจฺฉิมานํ อุปาทินฺนานํ ขนฺธานํ อนนฺตรปจฺจเยน ปจฺจโย.\csromanspacer{} ปญฺจวิญฺญาณํ วิปากมโนธาตุยา อนนฺตรปจฺจเยน ปจฺจโย.\csromanspacer{} วิปากมโนธาตุ วิปากมโนวิญฺญาณธาตุยา อนนฺตรปจฺจเยน ปจฺจโย. (1)}

\prose{อุปาทินฺโน ธมฺโม อนุปาทินฺนสฺส ธมฺมสฺส อนนฺตรปจฺจเยน ปจฺจโย.\csromanspacer{} ภวงฺคํ อาวชฺชนาย.\csromanspacer{} วิปากมโนวิญฺญาณธาตุ กิริยมโนวิญฺญาณธาตุยา อนนฺตรปจฺจเยน ปจฺจโย. (2)}

\proseitem{450}{อนุปาทินฺโน ธมฺโม อนุปาทินฺนสฺส ธมฺมสฺส อนนฺตรปจฺจเยน ปจฺจโย.\csromanspacer{} ปุริมา ปุริมา อนุปาทินฺนา ขนฺธา ปจฺฉิมานํ ปจฺฉิมานํ อนุปาทินฺนานํ ขนฺธานํ อนนฺตรปจฺจเยน ปจฺจโย.\csromanspacer{} อนุโลมํ โคตฺรภุสฺส ฯเปฯ.\csromanspacer{} ผลสมาปตฺติยา อนนฺตรปจฺจเยน ปจฺจโย. (1)}

\prose{อนุปาทินฺโน ธมฺโม อุปาทินฺนสฺส ธมฺมสฺส อนนฺตรปจฺจเยน ปจฺจโย.\csromanspacer{} อาวชฺชนา ปญฺจนฺนํ วิญฺญาณาณํ ฯเปฯ.\csromanspacer{} อนุปาทินฺนา ขนฺธา วุฏฺฐานสฺส อนนฺตรปจฺจเยน ปจฺจโย.\csromanspacer{} สมนนฺตรปจฺจเยน ปจฺจโย.\csromanspacer{} สหชาตปจฺจเยน ปจฺจโย.\csromanspacer{} ปญฺจ. (ปฏิจฺจสทิสา.) อญฺญมญฺญปจฺจเยน ปจฺจโย.\csromanspacer{} เทฺว. (ปฏิจฺจสทิสา.) นิสฺสยปจฺจเยน ปจฺจโย. (ปจฺจยวาเร นิสฺสยสทิสา.) ปญฺจ.}

\csromanheader{2}{\textbf{อุปนิสฺสย}}
\proseitem{451}{\csromanfolio{550}อุปาทินฺโน ธมฺโม อุปาทินฺนสฺส ธมฺมสฺส อุปนิสฺสยปจฺจเยน ปจฺจโย.\csromanspacer{} \textbf{อนนฺตรูปนิสฺสโย, ปกตูปนิสฺสโย ฯเปฯ.}\csromanspacer{} ปกตูปนิสฺสโย\csromandash{} กายิกํ สุขํ กายิกสฺส สุขสฺส.\csromanspacer{} กายิกสฺส ทุกฺขสฺส อุปนิสฺสยปจฺจเยน ปจฺจโย.\csromanspacer{} กายิกํ ทุกฺขํ.\csromanspacer{} อุปาทินฺนํ อุตุ.\csromanspacer{} โภชนํ กายิกสฺส สุขสฺส.\csromanspacer{} กายิกสฺส ทุกฺขสฺส อุปนิสฺสยปจฺจเยน ปจฺจโย.\csromanspacer{} กายิกํ สุขํ.\csromanspacer{} กายิกํ ทุกฺขํ.\csromanspacer{} อุตุ.\csromanspacer{} โภชนํ กายิกสฺส สุขสฺส.\csromanspacer{} กายิกสฺส ทุกฺขสฺส อุปนิสฺสยปจฺจเยน ปจฺจโย. (1)}

\prose{อุปาทินฺโน ธมฺโม อนุปาทินฺนสฺส ธมฺมสฺส อุปนิสฺสยปจฺจเยน ปจฺจโย.\csromanspacer{} \textbf{อารมฺมณูปนิสฺสโย, อนนฺตรูปนิสฺสโย, ปกตูปนิสฺสโย ฯเปฯ.}\csromanspacer{} ปกตูปนิสฺสโย\csromandash{}กายิกํ สุขํ อุปนิสฺสาย ทานํ เทติ ฯเปฯ สํฆํ ภินฺทติ.\csromanspacer{} กายิกํ ทุกฺขํ.\csromanspacer{} อุปาทินฺนํ อุตุํ.\csromanspacer{} โภชนํ อุปนิสฺสาย ทานํ เทติ ฯเปฯ สมาปตฺติํ อุปฺปาเทติ.\csromanspacer{} ปาณํ หนติ ฯเปฯ สํฆํ ภินฺทติ.\csromanspacer{} กายิกํ สุขํ.\csromanspacer{} กายิกํ ทุกฺขํ.\csromanspacer{} อุตุ.\csromanspacer{} โภชนํ สทฺธาย ฯเปฯ ปตฺถนาย มคฺคสฺส, ผลสมาปตฺติยา อุปนิสฺสยปจฺจเยน ปจฺจโย. (2)}

\proseitem{452}{อนุปาทินฺโน ธมฺโม อนุปาทินฺนสฺส ธมฺมสฺส อุปนิสฺสยปจฺจเยน ปจฺจโย.\csromanspacer{} \textbf{อารมฺมณูปนิสฺสโย, อนนฺตรูปนิสฺสโย, ปกตูปนิสฺสโย ฯเปฯ.}\csromanspacer{} ปกตูปนิสฺสโย\csromandash{}สทฺธํ อุปนิสฺสาย ทานํ เทติ.\csromanspacer{} มานํ ชปฺเปติ, ทิฏฺฐิํ คณฺหาติ.\csromanspacer{} สีลํ ฯเปฯ.\csromanspacer{} ปตฺถนํ.\csromanspacer{} อุตุํ.\csromanspacer{} โภชนํ.\csromanspacer{} เสนาสนํ อุปนิสฺสาย ทานํ เทติ ฯเปฯ สํฆํ ภินฺทติ.\csromanspacer{} สทฺธา ฯเปฯ.\csromanspacer{} ปตฺถนา.\csromanspacer{} อุตุ.\csromanspacer{} โภชนํ.\csromanspacer{} เสนาสนํ สทฺธาย ฯเปฯ ปตฺถนาย มคฺคสฺส, ผลสมาปตฺติยา อุปนิสฺสยปจฺจเยน ปจฺจโย. (1)}

\prose{อนุปาทินฺโน ธมฺโม อุปาทินฺนสฺส ธมฺมสฺส อุปนิสฺสยปจฺจเยน ปจฺจโย.\csromanspacer{} \textbf{อนนฺตรูปนิสฺสโย, ปกตูปนิสฺสโย ฯเปฯ.}\csromanspacer{} ปกตูปนิสฺสโย\csromandash{} สทฺธํ อุปนิสฺสาย อตฺตานํ อาตาเปติ ปริตาเปติ, ปริยิฏฺฐิมูลกํ ทุกฺขํ ปจฺจนุโภติ.\csromanspacer{} สีลํ ฯเปฯ.\csromanspacer{} ปตฺถนํ.\csromanspacer{} อุตุํ.\csromanspacer{} โภชนํ.\csromanspacer{} เสนาสนํ อุปนิสฺสาย อตฺตานํ อาตาเปติ ปริตาเปติ, ปริยิฏฺฐิมูลกํ ทุกฺขํ ปจฺจนุโภติ.\csromanspacer{} สทฺธา ฯเปฯ.\csromanspacer{} เสนาสนํ กายิกสฺส สุขสฺส.\csromanspacer{} กายิกสฺส ทุกฺขสฺส อุปนิสฺสยปจฺจเยน ปจฺจโย.\csromanspacer{} กุสลากุสลํ ธมฺมํ วิปากสฺส อุปนิสฺสยปจฺจเยน ปจฺจโย. (2)}

\csromanheader{2}{68. อุปาทินฺนทุก \textbf{ปุเรชาต}}
\proseitem{453}{\csromanfolio{551}อุปาทินฺโน ธมฺโม อุปาทินฺนสฺส ธมฺมสฺส ปุเรชาตปจฺจเยน ปจฺจโย.\csromanspacer{} \textbf{อารมฺมณปุเรชาตํ}, วตฺถุปุเรชาตํ.\csromanspacer{}\csromanspacer{}\textbf{อารมฺมณปุเรชาตํ}\csromandash{} จกฺขุํ ฯเปฯ กายํ.\csromanspacer{} อุปาทินฺเน รูเป.\csromanspacer{} คนฺเธ.\csromanspacer{} รเส.\csromanspacer{} โผฏฺฐพฺเพ.\csromanspacer{} วตฺถุํ อนิจฺจโต ฯเปฯ โทมนสฺสํ อุปฺปชฺชติ.\csromanspacer{} กุสลากุสเล นิรุทฺเธ วิปาโก ตทารมฺมณตา อุปฺปชฺชติ.\csromanspacer{} อุปาทินฺนํ รูปายตนํ จกฺขุวิญฺญาณสฺส ฯเปฯ คนฺธายตนํ ฯเปฯ โผฏฺฐพฺพายตนํ กายวิญฺญาณสฺส ฯเปฯ.\csromanspacer{} \textbf{วตฺถุปุเรชาตํ}\csromandash{}จกฺขายตนํ จกฺขุวิญฺญาณสฺส ฯเปฯ กายายตนํ กายวิญฺญาณสฺส ฯเปฯ.\csromanspacer{} วตฺถุ อุปาทินฺนานํ ขนฺธานํ ปุเรชาตปจฺจเยน ปจฺจโย. (1)}

\prose{อุปาทินฺโน ธมฺโม อนุปาทินฺนสฺส ธมฺมสฺส ปุเรชาตปจฺจเยน ปจฺจโย.\csromanspacer{} \textbf{อารมฺมณปุเรชาตํ}, วตฺถุปุเรชาตํ.\csromanspacer{}\csromanspacer{}\textbf{อารมฺมณปุเรชาตํ}\csromandash{} จกฺขุํ ฯเปฯ กายํ.\csromanspacer{} อุปาทินฺเน รูเป.\csromanspacer{} คนฺเธ.\csromanspacer{} รเส.\csromanspacer{} โผฏฺฐพฺเพ วตฺถุํ อนิจฺจโต ฯเปฯ โทมนสฺสํ อุปฺปชฺชติ.\csromanspacer{} ทิพฺเพน จกฺขุนา อุปาทินฺนํ รูปํ ปสฺสติ.\csromanspacer{} \textbf{วตฺถุปุเรชาตํ}\csromandash{}วตฺถุ อนุปาทินฺนานํ ขนฺธานํ ปุเรชาตปจฺจเยน ปจฺจโย. (2)}

\proseitem{454}{อนุปาทินฺโน ธมฺโม อนุปาทินฺนสฺส ธมฺมสฺส ปุเรชาตปจฺจเยน ปจฺจโย.\csromanspacer{} \textbf{อารมฺมณปุเรชาตํ}\csromandash{}อนุปาทินฺเน รูเป.\csromanspacer{} สทฺเท.\csromanspacer{} คนฺเธ.\csromanspacer{} รเส.\csromanspacer{} โผฏฺฐพฺเพ อนิจฺจโต ฯเปฯ โทมนสฺสํ อุปฺปชฺชติ.\csromanspacer{} ทิพฺเพน จกฺขุนา อนุปาทินฺนํ รูปํ ปสฺสติ.\csromanspacer{} ทิพฺพาย โสตธาตุยา สทฺทํ สุณาติ. (1)}

\prose{อนุปาทินฺโน ธมฺโม อุปาทินฺนสฺส ธมฺมสฺส ปุเรชาตปจฺจเยน ปจฺจโย.\csromanspacer{} \textbf{อารมฺมณปุเรชาตํ}\csromandash{}อนุปาทินฺเน รูเป.\csromanspacer{} สทฺเท.\csromanspacer{} คนฺเธ.\csromanspacer{} รเส.\csromanspacer{} โผฏฺฐพฺเพ อนิจฺจโต ฯเปฯ โทมนสฺสํ อุปฺปชฺชติ.\csromanspacer{} กุสลากุสเล นิรุทฺเธ วิปาโก ตทารมฺมณตา อุปฺปชฺชติ.\csromanspacer{} อนุปาทินฺนํ รูปายตนํ จกฺขุวิญฺญาณสฺส ฯเปฯ สทฺทายตนํ ฯเปฯ โผฏฺฐพฺพายตนํ กายวิญฺญาณสฺส ปุเรชาตปจฺจเยน ปจฺจโย. (2)}

\prose{อุปาทินฺโน จ อนุปาทินฺโน จ ธมฺมา อุปาทินฺนสฺส ธมฺมสฺส ปุเรชาตปจฺจเยน ปจฺจโย.\csromanspacer{} \textbf{อารมฺมณปุเรชาตํ}, วตฺถุปุเรชาตํ.\csromanspacer{}\csromanspacer{}อนุปาทินฺนํ รูปายตนญฺจ จกฺขายตนญฺจ จกฺขุวิญฺญาณสฺส ฯเปฯ สทฺทายตนญฺจ โสตายตนญฺจ ฯเปฯ โผฏฺฐพฺพายตนญฺจ กายายตนญฺจ กายวิญฺญาณสฺส ปุเรชาตปจฺจเยน ปจฺจโย.\csromanspacer{} อนุปาทินฺนํ รูปายตนญฺจ วตฺถุ จ ฯเปฯ โผฏฺฐพฺพายตนญฺจ วตฺถุ จ อุปาทินฺนานํ ขนฺธานํ ปุเรชาตปจฺจเยน ปจฺจโย. (1)}

\prose{\csromanfolio{552}อุปาทินฺโน จ อนุปาทินฺโน จ ธมฺมา อนุปาทินฺนสฺส ธมฺมสฺส ปุเรชาตปจฺจเยน ปจฺจโย.\csromanspacer{} \textbf{อารมฺมณปุเรชาตํ, วตฺถุปุเรชาตํ}.\csromanspacer{} อนุปาทินฺนํ รูปายตนญฺจ วตฺถุ จ ฯเปฯ โผฏฺฐพฺพายตนญฺจ วตฺถุ จ อนุปาทินฺนานํ ขนฺธานํ ปุเรชาตปจฺจเยน ปจฺจโย. (2)}

\csromanheader{2}{\textbf{ปจฺฉาชาตาเสวน}}
\proseitem{455}{อุปาทินฺโน ธมฺโม อุปาทินฺนสฺส ธมฺมสฺส ปจฺฉาชาตปจฺจเยน ปจฺจโย.\csromanspacer{} ปจฺฉาชาตา อุปาทินฺนา ขนฺธา ปุเรชาตสฺส อิมสฺส อุปาทินฺนสฺส กายสฺส ปจฺฉาชาตปจฺจเยน ปจฺจโย. (1)}

\prose{อุปาทินฺโน ธมฺโม อนุปาทินฺนสฺส ธมฺมสฺส ปจฺฉาชาตปจฺจเยน ปจฺจโย.\csromanspacer{} ปจฺฉาชาตา อุปาทินฺนา ขนฺธา ปุเรชาตสฺส อิมสฺส อนุปาทินฺนสฺส กายสฺส ปจฺฉาชาตปจฺจเยน ปจฺจโย. (2)}

\prose{อุปาทินฺโน ธมฺโม อุปาทินฺนสฺส จ อนุปาทินฺนสฺส จ ธมฺมสฺส ปจฺฉาชาตปจฺจเยน ปจฺจโย. (สํขิตฺตํ.) (3)}

\prose{อนุปาทินฺโน ธมฺโม อนุปาทินฺนสฺส ธมฺมสฺส ปจฺฉาชาตปจฺจเยน ปจฺจโย. (สํขิตฺตํ.) (1)}

\prose{อนุปาทินฺโน ธมฺโม อุปาทินฺนสฺส ธมฺมสฺส ปจฺฉาชาตปจฺจเยน ปจฺจโย. (สํขิตฺตํ.) (2)}

\prose{อนุปาทินฺโน ธมฺโม อุปาทินฺนสฺส จ อนุปาทินฺนสฺส จ ธมฺมสฺส ปจฺฉาชาตปจฺจเยน ปจฺจโย. (สํขิตฺตํ.) (3)}

\proseitem{456}{อนุปาทินฺโน ธมฺโม อนุปาทินฺนสฺส ธมฺมสฺส อาเสวนปจฺจเยน ปจฺจโย.\csromanspacer{} เอกํ. (สํขิตฺตํ.)}

\csromanheader{2}{\textbf{กมฺม}}
\proseitem{457}{อุปาทินฺโน ธมฺโม อุปาทินฺนสฺส ธมฺมสฺส กมฺมปจฺจเยน ปจฺจโย.\csromanspacer{} อุปาทินฺนา เจตนา สมฺปยุตฺตกานํ ขนฺธานํ กมฺมปจฺจเยน ปจฺจโย.\csromanspacer{} ปฏิสนฺธิกฺขเณ อุปาทินฺนา เจตนา สมฺปยุตฺตกานํ ขนฺธานํ กฏตฺตา จ รูปานํ กมฺมปจฺจเยน ปจฺจโย. (1)}

\csromanheader{2}{68. อุปาทินฺนทุก}
\prose{\csromanfolio{553}อุปาทินฺโน ธมฺโม อนุปาทินฺนสฺส ธมฺมสฺส กมฺมปจฺจเยน ปจฺจโย.\csromanspacer{} อุปาทินฺนา เจตนา จิตฺตสมุฏฺฐานานํ รูปานํ กมฺมปจฺจเยน ปจฺจโย. (2)}

\prose{อุปาทินฺโน ธมฺโม อุปาทินฺนสฺส จ อนุปาทินฺนสฺส จ ธมฺมสฺส กมฺมปจฺจเยน ปจฺจโย.\csromanspacer{} อุปาทินฺนา เจตนา สมฺปยุตฺตกานํ ขนฺธานํ จิตฺตสมุฏฺฐานานญฺจ รูปานํ กมฺมปจฺจเยน ปจฺจโย. (3)}

\prose{อนุปาทินฺโน ธมฺโม อนุปาทินฺนสฺส ธมฺมสฺส กมฺมปจฺจเยน ปจฺจโย.\csromanspacer{} อนุปาทินฺนา เจตนา สมฺปยุตฺตกานํ ขนฺธานํ จิตฺตสมุฏฺฐานานญฺจ รูปานํ กมฺมปจฺจเยน ปจฺจโย. (1)}

\prose{อนุปาทินฺโน ธมฺโม อุปาทินฺนสฺส ธมฺมสฺส กมฺมปจฺจเยน ปจฺจโย.\csromanspacer{} \textbf{นานากฺขณิกา} อนุปาทินฺนา เจตนา วิปากานํ ขนฺธานํ กฏตฺตา จ รูปานํ กมฺมปจฺจเยน ปจฺจโย. (2)}

\csromanheader{2}{\textbf{วิปาก}}
\proseitem{458}{อุปาทินฺโน ธมฺโม อุปาทินฺนสฺส ธมฺมสฺส วิปากปจฺจเยน ปจฺจโย.\csromanspacer{} อุปาทินฺโน เอโก ขนฺโธ. (สํขิตฺตํ.\csromanspacer{} อุปาทินฺนมูลเก ตีณิ.)}

\prose{อนุปาทินฺโน ธมฺโม อนุปาทินฺนสฺส ธมฺมสฺส วิปากปจฺจเยน ปจฺจโย.\csromanspacer{} วิปาโก อนุปาทินฺโน เอโก ขนฺโธ ติณฺณนฺนํ ขนฺธานํ จิตฺตสมุฏฺฐานาญฺจ รูปานํ วิปากปจฺจเยน ปจฺจโย ฯเปฯ เทฺว ขนฺธา ฯเปฯ. (1)}

\csromanheader{2}{\textbf{อาหาร}}
\proseitem{459}{อุปาทินฺโน ธมฺโม อุปาทินฺนสฺส ธมฺมสฺส อาหารปจฺจเยน ปจฺจโย.\csromanspacer{} อุปาทินฺนา อาหารา สมฺปยุตฺตกานํ ขนฺธานํ อาหารปจฺจเยน ปจฺจโย.\csromanspacer{} ปฏิสนฺธิกฺขเณ ฯเปฯ.\csromanspacer{} อุปาทินฺโน กพฬีกาโร อาหาโร อิมสฺส อุปาทินฺนสฺส กายสฺส อาหารปจฺจเยน ปจฺจโย. (1)}

\prose{อุปาทินฺโน ธมฺโม อนุปาทินฺนสฺส ธมฺมสฺส อาหารปจฺจเยน ปจฺจโย.\csromanspacer{} อุปาทินฺนา อาหารา จิตฺตสมุฏฺฐานานํ รูปานํ อาหารปจฺจเยน ปจฺจโย.\csromanspacer{} อุปาทินฺโน กพฬีกาโร อาหาโร อิมสฺส อนุปาทินฺนสฺส กายสฺส อาหารปจฺจเยน ปจฺจโย. (2)}

\prose{\csromanfolio{554}อุปาทินฺโน ธมฺโม อุปาทินฺนสฺส จ อนุปาทินฺนสฺส จ ธมฺมสฺส อาหารปจฺจเยน ปจฺจโย.\csromanspacer{} อุปาทินฺนา อาหารา สมฺปยุตฺตกานํ ขนฺธานํ จิตฺตสมุฏฺฐานานญฺจ รูปานํ อาหารปจฺจเยน ปจฺจโย.\csromanspacer{} อุปาทินฺโน กพฬีกาโร อาหาโร อิมสฺส อุปาทินฺนสฺส จ อนุปาทินฺนสฺส จ กายสฺส อาหารปจฺจเยน ปจฺจโย. (3)}

\proseitem{460}{อนุปาทินฺโน ธมฺโม อนุปาทินฺนสฺส ธมฺมสฺส อาหารปจฺจเยน ปจฺจโย.\csromanspacer{} อนุปาทินฺนา อาหารา สมฺปยุตฺตกานํ ขนฺธานํ จิตฺตสมุฏฺฐานานญฺจ รูปานํ อาหารปจฺจเยน ปจฺจโย.\csromanspacer{} อนุปาทินฺโน กพฬีกาโร อาหาโร อิมสฺส อนุปาทินฺนสฺส กายสฺส อาหารปจฺจเยน ปจฺจโย. (1)}

\prose{อนุปาทินฺโน ธมฺโม อุปาทินฺนสฺส ธมฺมสฺส อาหารปจฺจเยน ปจฺจโย.\csromanspacer{} อนุปาทินฺโน กพฬีกาโร อาหาโร อิมสฺส อุปาทินฺนสฺส กายสฺส อาหารปจฺจเยน ปจฺจโย. (2)}

\prose{อนุปาทินฺโน ธมฺโม อุปาทินฺนสฺส จ อนุปาทินฺนสฺส จ ธมฺมสฺส อาหารปจฺจเยน ปจฺจโย.\csromanspacer{} อนุปาทินฺโน กพฬีกาโร อาหาโร อิมสฺส อุปาทินฺนสฺส จ อนุปาทินฺนสฺส จ กายสฺส อาหารปจฺจเยน ปจฺจโย. (3)}

\prose{อุปาทินฺโน จ อนุปาทินฺโน จ ธมฺมา อุปาทินฺนสฺส ธมฺมสฺส อาหารปจฺจเยน ปจฺจโย.\csromanspacer{} อุปาทินฺโน จ อนุปาทินฺโน จ กพฬีกาโร อาหาโร อิมสฺส อุปาทินฺนสฺส กายสฺส อาหารปจฺจเยน ปจฺจโย. (1)}

\prose{อุปาทินฺโน จ อนุปาทินฺโน จ ธมฺมา อนุปาทินฺนสฺส ธมฺมสฺส อาหารปจฺจเยน ปจฺจโย.\csromanspacer{} อุปาทินฺโน จ อนุปาทินฺโน จ กพฬีกาโร อาหาโร อิมสฺส อนุปาทินฺนสฺส กายสฺส อาหารปจฺจเยน ปจฺจโย. (2)}

\prose{อุปาทินฺโน จ อนุปาทินฺโน จ ธมฺมา อุปาทินฺนสฺส จ อนุปาทินฺนสฺส จ ธมฺมสฺส อาหารปจฺจเยน ปจฺจโย.\csromanspacer{} อุปาทินฺโน จ อนุปาทินฺโน จ กพฬีกาโร อาหาโร อิมสฺส อุปาทินฺนสฺส จ อนุปาทินฺนสฺส จ กายสฺส อาหารปจฺจเยน ปจฺจโย. (3)}

\csromanheader{2}{\textbf{อินฺทฺริยาทิ}}
\proseitem{461}{อุปาทินฺโน ธมฺโม อุปาทินฺนสฺส ธมฺมสฺส อินฺทฺริยปจฺจเยน ปจฺจโย.\csromanspacer{} อุปาทินฺนา อินฺทฺริยา สมฺปยุตฺตกานํ ขนฺธานํ อินฺทฺริยปจฺจเยน ปจฺจโย. (1)}

\vspace{\tipitakaparskip}
\csromancenter{อุปาทินฺนทุก ปฏิสนฺธิกฺขเณ อุปาทินฺนา อินฺทฺริยา สมฺปยุตฺตกานํ ขนฺธานํ กฏตฺตา จ รูปานํ อินฺทฺริยปจฺจเยน ปจฺจโย.\csromanspacer{} จกฺขุนฺทฺริยํ จกฺขุวิญฺญาณสฺส ฯเปฯ.\csromanspacer{} กายินฺทฺริยํ กายวิญฺญาณสฺส ฯเปฯ.\csromanspacer{} รูปชีวิตินฺทฺริยํ กฏตฺตารูปานํ อินฺทฺริยปจฺจเยน ปจฺจโย. (2)}
\prose{\csromanfolio{555}อุปาทินฺโน ธมฺโม อนุปาทินฺนสฺส ธมฺมสฺส ฯเปฯ. (อุปาทินฺนมูลเก ตีณิ, ปฐมสฺเสว รูปชีวิตินฺทฺริยํ, อิตเรสุ นตฺถิ.) (3)}

\prose{อนุปาทินฺโน ธมฺโม อนุปาทินฺนสฺส ธมฺมสฺส อินฺทฺริยปจฺจเยน ปจฺจโย.\csromanspacer{} อนุปาทินฺนา อินฺทฺริยา สมฺปยุตฺตกานํ ขนฺธานํ จิตฺตสมุฏฺฐานานญฺจ รูปานํ อินฺทฺริยปจฺจเยน ปจฺจโย. (1)}

\prose{อุปาทินฺโน ธมฺโม อุปาทินฺนสฺส ธมฺมสฺส ฌานปจฺจเยน ปจฺจโย.\csromanspacer{} จตฺตาริ.\csromanspacer{} มคฺคปจฺจเยน ปจฺจโย.\csromanspacer{} จตฺตาริ.\csromanspacer{} สมฺปยุตฺตปจฺจเยน ปจฺจโย.\csromanspacer{} เทฺว.}

\csromanheader{2}{\textbf{วิปฺปยุตฺต}}
\proseitem{462}{อุปาทินฺโน ธมฺโม อุปาทินฺนสฺส ธมฺมสฺส วิปฺปยุตฺตปจฺจเยน ปจฺจโย.\csromanspacer{} สหชาตํ, ปุเรชาตํ, ปจฺฉาชาตํ.\csromanspacer{}\csromanspacer{}\textbf{สหชาตา} ปฏิสนฺธิกฺขเณ อุปาทินฺนา ขนฺธา กฏตฺตารูปานํ วิปฺปยุตฺตปจฺจเยน ปจฺจโย.\csromanspacer{} ขนฺธา วตฺถุสฺส วิปฺปยุตฺตปจฺจเยน ปจฺจโย.\csromanspacer{} วตฺถุ ขนฺธานํ วิปฺปยุตฺตปจฺจเยน ปจฺจโย.\csromanspacer{} \textbf{ปุเรชาตํ} จกฺขายตนํ จกฺขุวิญฺญาณสฺส ฯเปฯ กายายตนํ กายวิญฺญาณสฺส ฯเปฯ.\csromanspacer{} วตฺถุ อุปาทินฺนานํ ขนฺธานํ วิปฺปยุตฺตปจฺจเยน ปจฺจโย.\csromanspacer{} \textbf{ปจฺฉาชาตา} อุปาทินฺนา ขนฺธา ปุเรชาตสฺส อิมสฺส อุปาทินฺนสฺส กายสฺส วิปฺปยุตฺตปจฺจเยน ปจฺจโย. (1)}

\prose{อุปาทินฺโน ธมฺโม อนุปาทินฺนสฺส ธมฺมสฺส วิปฺปยุตฺตปจฺจเยน ปจฺจโย.\csromanspacer{} สหชาตํ, ปุเรชาตํ, ปจฺฉาชาตํ.\csromanspacer{}\csromanspacer{}\textbf{สหชาตา} อุปาทินฺนา ขนฺธา จิตฺตสมุฏฺฐานานํ รูปานํ วิปฺปยุตฺตปจฺจเยน ปจฺจโย.\csromanspacer{} \textbf{ปุเรชาตํ} วตฺถุ อนุปาทินฺนานํ ขนฺธานํ วิปฺปยุตฺตปจฺจเยน ปจฺจโย.\csromanspacer{} \textbf{ปจฺฉาชาตา} อุปาทินฺนา ขนฺธา ปุเรชาตสฺส อิมสฺส อนุปาทินฺนสฺส กายสฺส วิปฺปยุตฺตปจฺจเยน ปจฺจโย. (2)}

\prose{อุปาทินฺโน ธมฺโม อุปาทินฺนสฺส จ อนุปาทินฺนสฺส จ ธมฺมสฺส วิปฺปยุตฺตปจฺจเยน ปจฺจโย.\csromanspacer{} \textbf{ปจฺฉาชาตา} อุปาทินฺนา ขนฺธา ปุเรชาตสฺส อิมสฺส อุปาทินฺนสฺส จ อนุปาทินฺนสฺส จ กายสฺส วิปฺปยุตฺตปจฺจเยน ปจฺจโย. (3)}

\proseitem{463}{\csromanfolio{556}อนุปาทินฺโน ธมฺโม อนุปาทินฺนสฺส ธมฺมสฺส วิปฺปยุตฺตปจฺจเยน ปจฺจโย.\csromanspacer{} สหชาตํ, ปจฺฉาชาตํ.\csromanspacer{}\csromanspacer{}\textbf{สหชาตา} อนุปาทินฺนา ขนฺธา จิตฺตสมุฏฺฐานานํ รูปานํ วิปฺปยุตฺตปจฺจเยน ปจฺจโย.\csromanspacer{} \textbf{ปจฺฉาชาตา} อนุปาทินฺนา ขนฺธา ปุเรชาตสฺส อิมสฺส อนุปาทินฺนสฺส กายสฺส วิปฺปยุตฺตปจฺจเยน ปจฺจโย. (1)}

\prose{อนุปาทินฺโน ธมฺโม อุปาทินฺนสฺส ธมฺมสฺส วิปฺปยุตฺตปจฺจเยน ปจฺจโย.\csromanspacer{} \textbf{ปจฺฉาชาตา} อนุปาทินฺนา ขนฺธา ปุเรชาตสฺส อิมสฺส อุปาทินฺนสฺส กายสฺส วิปฺปยุตฺตปจฺจเยน ปจฺจโย. (2)}

\prose{อนุปาทินฺโน ธมฺโม อุปาทินฺนสฺส จ อนุปาทินฺนสฺส จ ธมฺมสฺส วิปฺปยุตฺตปจฺจเยน ปจฺจโย.\csromanspacer{} \textbf{ปจฺฉาชาตา} อนุปาทินฺนา ขนฺธา ปุเรชาตสฺส อิมสฺส อุปาทินฺนสฺส จ อนุปาทินฺนสฺส จ กายสฺส วิปฺปยุตฺตปจฺจเยน ปจฺจโย. (3)}

\csromanheader{2}{\textbf{อตฺถิ-อาทิ}}
\proseitem{464}{อุปาทินฺโน ธมฺโม อุปาทินฺนสฺส ธมฺมสฺส อตฺถิปจฺจเยน ปจฺจโย.\csromanspacer{} สหชาตํ, ปุเรชาตํ, ปจฺฉาชาตํ, อาหารํ, อินฺทฺริยํ. (สํขิตฺตํ.\csromanspacer{} ยถา นิกฺขิตฺตปทานิ วิภชิตพฺพานิ ปริปุณฺณานิ.) (1)}

\prose{อุปาทินฺโน ธมฺโม อนุปาทินฺนสฺส ธมฺมสฺส อตฺถิปจฺจเยน ปจฺจโย.\csromanspacer{} สหชาตํ, ปุเรชาตํ, ปจฺฉาชาตํ, อาหารํ. (สํขิตฺตํ.\csromanspacer{} ยถา นิกฺขิตฺตปทานิ วิตฺถาเรตพฺพานิ.) (2)}

\prose{อุปาทินฺโน ธมฺโม อุปาทินฺนสฺส จ อนุปาทินฺนสฺส จ ธมฺมสฺส อตฺถิปจฺจเยน ปจฺจโย.\csromanspacer{} สหชาตํ, ปจฺฉาชาตํ, อาหารํ. (สํขิตฺตํ.\csromanspacer{} ยถา นิกฺขิตฺตปทานิ วิตฺถาเรตพฺพานิ.) (3)}

\prose{อนุปาทินฺโน ธมฺโม อนุปาทินฺนสฺส ธมฺมสฺส อตฺถิปจฺจเยน ปจฺจโย.\csromanspacer{} สหชาตํ, ปุเรชาตํ, ปจฺฉาชาตํ, อาหารํ. (สํขิตฺตํ.\csromanspacer{} ยถา นิกฺขิตฺตปทานิ วิภชิตพฺพานิ.) (1)}

\prose{อนุปาทินฺโน ธมฺโม อุปาทินฺนสฺส ธมฺมสฺส อตฺถิปจฺจเยน ปจฺจโย.\csromanspacer{} \textbf{ปุเรชาตํ}, ปจฺฉาชาตํ, อาหารํ.\csromanspacer{}\csromanspacer{}\textbf{ปุเรชาตํ}\csromandash{}อนุปาทินฺเน รูเป.\csromanspacer{} สทฺเท.\csromanspacer{} คนฺเธ.\csromanspacer{} รเส.\csromanspacer{} โผฏฺฐพฺเพ อนิจฺจโต ฯเปฯ โทมนสฺสํ อุปฺปชฺชติ.\csromanspacer{} กุสลากุสเล นิรุทฺเธ วิปาโก ตทารมฺมณตา อุปฺปชฺชติ.\csromanspacer{} อนุปาทินฺนํ}

\vspace{\tipitakaparskip}
\csromancenter{อุปาทินฺนทุก รูปายตนํ จกฺขุวิญฺญาณสฺส ฯเปฯ โผฏฺฐพฺพายตนํ กายวิญฺญาณสฺส ฯเปฯ.\csromanspacer{} \textbf{ปจฺฉาชาตา} อนุปาทินฺนา ขนฺธา ปุเรชาตสฺส อิมสฺส อุปาทินฺนสฺส กายสฺส อตฺถิปจฺจเยน ปจฺจโย.\csromanspacer{} อนุปาทินฺโน \textbf{กพฬีกาโร อาหาโร} อิมสฺส อุปาทินฺนสฺส กายสฺส อตฺถิปจฺจเยน ปจฺจโย. (2)}
\prose{\csromanfolio{557}อนุปาทินฺโน ธมฺโม อุปาทินฺนสฺส จ อนุปาทินฺนสฺส จ ธมฺมสฺส อตฺถิปจฺจเยน ปจฺจโย.\csromanspacer{} ปจฺฉาชาตํ, อาหารํ.\csromanspacer{}\csromanspacer{}\textbf{ปจฺฉาชาตา} อนุปาทินฺนา ขนฺธา ปุเรชาตสฺส อิมสฺส อุปาทินฺนสฺส จ อนุปาทินฺนสฺส จ กายสฺส อตฺถิปจฺจเยน ปจฺจโย.\csromanspacer{} อนุปาทินฺโน \textbf{กพฬีกาโร อาหาโร} อิมสฺส อุปาทินฺนสฺส จ อนุปาทินฺนสฺส จ กายสฺส อตฺถิปจฺจเยน ปจฺจโย. (3)}

\proseitem{465}{อุปาทินฺโน จ อนุปาทินฺโน จ ธมฺมา อุปาทินฺนสฺส ธมฺมสฺส อตฺถิปจฺจเยน ปจฺจโย.\csromanspacer{} \textbf{ปุเรชาตํ}, ปจฺฉาชาตํ, อาหารํ, อินฺทฺริยํ.\csromanspacer{}\csromanspacer{}\textbf{ปุเรชาตํ} อนุปาทินฺนํ รูปายตนญฺจ จกฺขายตนญฺจ จกฺขุวิญฺญาณสฺส ฯเปฯ.\csromanspacer{} อนุปาทินฺนํ โผฏฺฐพฺพายตนญฺจ กายายตนญฺจ กายวิญฺญาณสฺส อตฺถิปจฺจเยน ปจฺจโย.\csromanspacer{} อนุปาทินฺนํ รูปายตนญฺจ วตฺถุ จ ฯเปฯ โผฏฺฐพฺพายตนญฺจ วตฺถุ จ อุปาทินฺนานํ ขนฺธานํ อตฺถิปจฺจเยน ปจฺจโย.\csromanspacer{} \textbf{ปจฺฉาชาตา} อุปาทินฺนา ขนฺธา จ อนุปาทินฺโน \textbf{กพฬีกาโร อาหาโร} จ อิมสฺส อุปาทินฺนสฺส กายสฺส อตฺถิปจฺจเยน ปจฺจโย.\csromanspacer{} \textbf{ปจฺฉาชาตา} อนุปาทินฺนา ขนฺธา จ รูปชีวิตินฺทฺริยญฺจ กฏตฺตารูปานํ อตฺถิปจฺจเยน ปจฺจโย. (1)}

\prose{อุปาทินฺโน จ อนุปาทินฺโน จ ธมฺมา อนุปาทินฺนสฺส ธมฺมสฺส อตฺถิปจฺจเยน ปจฺจโย.\csromanspacer{} สหชาตํ, ปุเรชาตํ, ปจฺฉาชาตํ, อาหารํ.\csromanspacer{}\csromanspacer{}\textbf{สหชาตา} อุปาทินฺนา ขนฺธา จ มหาภูตา จ จิตฺตสมุฏฺฐานานํ รูปานํ อตฺถิปจฺจเยน ปจฺจโย.\csromanspacer{} \textbf{สหชาโต} อนุปาทินฺโน เอโก ขนฺโธ จ วตฺถุ จ ติณฺณนฺนํ ขนฺธานํ อตฺถิปจฺจเยน ปจฺจโย ฯเปฯ เทฺว ขนฺธา จ ฯเปฯ.\csromanspacer{} \textbf{ปุเรชาตํ} อนุปาทินฺนํ รูปายตนญฺจ วตฺถุ จ อนุปาทินฺนานํ ขนฺธานํ อตฺถิปจฺจเยน ปจฺจโย ฯเปฯ โผฏฺฐพฺพายตนญฺจ วตฺถุ จ อนุปาทินฺนานํ ขนฺธานํ อตฺถิปจฺจเยน ปจฺจโย.\csromanspacer{} \textbf{ปจฺฉาชาตา} อุปาทินฺนา ขนฺธา จ อนุปาทินฺโน \textbf{กพฬีกาโร อาหาโร} จ อิมสฺส อนุปาทินฺนสฺส กายสฺส อตฺถิปจฺจเยน ปจฺจโย. (2)}

\prose{อุปาทินฺโน จ อนุปาทินฺโน จ ธมฺมา อุปาทินฺนสฺส จ อนุปาทินฺนสฺส จ ธมฺมสฺส อตฺถิปจฺจเยน ปจฺจโย.\csromanspacer{} \textbf{อาหารํ}\csromandash{}อุปาทินฺโน จ อนุปาทินฺโน จ \csromanfolio{558}กพฬีกาโร อาหาโร อิมสฺส อุปาทินฺนสฺส จ อนุปาทินฺนสฺส จ กายสฺส อตฺถิปจฺจเยน ปจฺจโย. (3)}

\prose{นตฺถิปจฺจเยน ปจฺจโย.\csromanspacer{} วิคตปจฺจเยน ปจฺจโย.\csromanspacer{} อวิคตปจฺจเยน ปจฺจโย.}

\csromanheader{2}{1. ปจฺจยานุโลม 2. สงฺขฺยาวาร}
\csromanheader{2}{\textbf{สุทฺธ}}
\proseitem{466}{เหตุยา จตฺตาริ, อารมฺมเณ จตฺตาริ, อธิปติยา เทฺว, อนนฺตเร จตฺตาริ, สมนนฺตเร จตฺตาริ, สหชาเต ปญฺจ, อญฺญมญฺเญ เทฺว, นิสฺสเย ปญฺจ, อุปนิสฺสเย จตฺตาริ, ปุเรชาเต ฉ, ปจฺฉาชาเต ฉ, อาเสวเน เอกํ, กมฺเม ปญฺจ, วิปาเก จตฺตาริ, อาหาเร นว, อินฺทฺริเย จตฺตาริ, ฌาเน จตฺตาริ, มคฺเค จตฺตาริ, สมฺปยุตฺเต เทฺว, วิปฺปยุตฺเต ฉ, อตฺถิยา นว, นตฺถิยา จตฺตาริ, วิคเต จตฺตาริ, อวิคเต นว.}

\csromanheader{2}{2. ปจฺจนียุทฺธาร}
\proseitem{467}{อุปาทินฺโน ธมฺโม อุปาทินฺนสฺส ธมฺมสฺส อารมฺมณปจฺจเยน ปจฺจโย.\csromanspacer{} สหชาตปจฺจเยน ปจฺจโย.\csromanspacer{} อุปนิสฺสยปจฺจเยน ปจฺจโย, ปุเรชาตปจฺจเยน ปจฺจโย.\csromanspacer{} ปจฺฉาชาตปจฺจเยน ปจฺจโย.\csromanspacer{} อาหารปจฺจเยน ปจฺจโย.\csromanspacer{} อินฺทฺริยปจฺจเยน ปจฺจโย. (1)}

\prose{อุปาทินฺโน ธมฺโม อนุปาทินฺนสฺส ธมฺมสฺส อารมฺมณปจฺจเยน ปจฺจโย.\csromanspacer{} สหชาตปจฺจเยน ปจฺจโย.\csromanspacer{} อุปนิสฺสยปจฺจเยน ปจฺจโย.\csromanspacer{} ปุเรชาตปจฺจเยน ปจฺจโย.\csromanspacer{} ปจฺฉาชาตปจฺจเยน ปจฺจโย.\csromanspacer{} อาหารปจฺจเยน ปจฺจโย. (2)}

\prose{อุปาทินฺโน ธมฺโม อุปาทินฺนสฺส จ อนุปาทินฺนสฺส จ ธมฺมสฺส สหชาตปจฺจเยน ปจฺจโย.\csromanspacer{} ปจฺฉาชาตปจฺจเยน ปจฺจโย.\csromanspacer{} อาหารปจฺจเยน ปจฺจโย. (3)}

\proseitem{468}{อนุปาทินฺโน ธมฺโม อนุปาทินฺนสฺส ธมฺมสฺส อารมฺมณปจฺจเยน ปจฺจโย.\csromanspacer{} สหชาตปจฺจเยน ปจฺจโย, อุปนิสฺสยปจฺจเยน ปจฺจโย.\csromanspacer{} ปจฺฉาชาตปจฺจเยน ปจฺจโย.\csromanspacer{} อาหารปจฺจเยน ปจฺจโย. (1)}

\csromanheader{2}{68. อุปาทินฺนทุก}
\prose{\csromanfolio{559}อนุปาทินฺโน ธมฺโม อุปาทินฺนสฺส ธมฺมสฺส อารมฺมณปจฺจเยน ปจฺจโย.\csromanspacer{} อุปนิสฺสยปจฺจเยน ปจฺจโย.\csromanspacer{} ปจฺฉาชาตปจฺจเยน ปจฺจโย.\csromanspacer{} กมฺมปจฺจเยน ปจฺจโย.\csromanspacer{} อาหารปจฺจเยน ปจฺจโย. (2)}

\prose{อนุปาทินฺโน ธมฺโม อุปาทินฺนสฺส จ อนุปาทินฺนสฺส จ ธมฺมสฺส ปจฺฉาชาตปจฺจเยน ปจฺจโย.\csromanspacer{} อาหารปจฺจเยน ปจฺจโย. (3)}

\prose{อุปาทินฺโน จ อนุปาทินฺโน จ ธมฺมา อุปาทินฺนสฺส ธมฺมสฺส ปุเรชาตํ, ปจฺฉาชาตํ, อาหารํ, อินฺทฺริยํ. (1)}

\prose{อุปาทินฺโน จ อนุปาทินฺโน จ ธมฺมา อนุปาทินฺนสฺส ธมฺมสฺส สหชาตํ, ปุเรชาตํ, ปจฺฉาชาตํ, อาหารํ. (2)}

\prose{อุปาทินฺโน จ อนุปาทินฺโน จ ธมฺมา อุปาทินฺนสฺส จ อนุปาทินฺนสฺส จ ธมฺมสฺส อาหารปจฺจเยน ปจฺจโย. (3)}

\csromanheader{2}{2. ปจฺจยปจฺจนีย 2. สงฺขฺยาวาร}
\csromanheader{2}{\textbf{สุทฺธ}}
\proseitem{469}{นเหตุยา นว, น-อารมฺมเณ นว, (สพฺพตฺถ นว,) น-อาหาเร อฏฺฐ ฯเปฯ นสมฺปยุตฺเต นว, นวิปฺปยุตฺเต นว, โน-อตฺถิยา จตฺตาริ, โนนตฺถิยา นว, โนวิคเต นว, โน-อวิคเต จตฺตาริ.}

\csromanheader{2}{3. ปจฺจยานุโลมปจฺจนีย}
\proseitem{470}{เหตุปจฺจยา น-อารมฺมเณ จตฺตาริ ฯเปฯ น-อญฺญมญฺเญ ตีณิ, น-อุปนิสฺสเย จตฺตาริ, นสมฺปยุตฺเต ตีณิ, นวิปฺปยุตฺเต เทฺว, โนนตฺถิยา จตฺตาริ, โนวิคเต จตฺตาริ.}

\csromanheader{2}{4. ปจฺจยปจฺจนียานุโลม}
\proseitem{471}{นเหตุปจฺจยา อารมฺมเณ จตฺตาริ, อธิปติยา เทฺว, (อนุโลมมาติกา กาตพฺพา.) ฯเปฯ อวิคเต นว.}

\nitthitam{อุปาทินฺนทุกํ นิฏฺฐิตํ.}
\nitthitam{มหนฺตรทุกํ นิฏฺฐิตํ.}
\csromanmatikaanchor{mtk.68}
\csromantocmarknopage{section}{11. อุปาทานโคจฺฉก}{mtk.68}
\bookmark[dest={mtk.68},level=1]{11. อุปาทานโคจฺฉก}
\csromanheader{1}{11. อุปาทานโคจฺฉก}
\csromanmatikaanchor{mtk.69}
\csromantocmarkref{subsection}{69. อุปาทานทุก}{mtk.69}
\bookmark[dest={mtk.69},level=2]{69. อุปาทานทุก}
\csromanheader{2}{69. อุปาทานทุก 1. ปฏิจฺจวาร}
\csromanheader{2}{1. ปจฺจยานุโลม 1. วิภงฺควาร}
\csromanheader{2}{\textbf{เหตุ}}
\proseitem{1}{\csromanfolio{560}อุปาทานํ ธมฺมํ ปฏิจฺจ อุปาทาโน ธมฺโม อุปฺปชฺชติ เหตุปจฺจยา.\csromanspacer{} ทิฏฺฐุปาทานํ ปฏิจฺจ กามุปาทานํ.\csromanspacer{} กามุปาทานํ ปฏิจฺจ ทิฏฺฐุปาทานํ.\csromanspacer{} สีลพฺพตุปาทานํ ปฏิจฺจ กามุปาทานํ.\csromanspacer{} กามุปาทานํ ปฏิจฺจ สีลพฺพตุปาทานํ.\csromanspacer{} อตฺตวาทุปาทานํ ปฏิจฺจ กามุปาทานํ.\csromanspacer{} กามุปาทานํ ปฏิจฺจ อตฺตวาทุปาทานํ. (1)}

\prose{อุปาทานํ ธมฺมํ ปฏิจฺจ โน-อุปาทาโน ธมฺโม อุปฺปชฺชติ เหตุปจฺจยา.\csromanspacer{} อุปาทาเน ปฏิจฺจ สมฺปยุตฺตกา ขนฺธา จิตฺตสมุฏฺฐานญฺจ รูปํ. (2)}

\prose{อุปาทานํ ธมฺมํ ปฏิจฺจ อุปาทาโน จ โน-อุปาทาโน จ ธมฺมา อุปฺปชฺชนฺติ เหตุปจฺจยา.\csromanspacer{} ทิฏฺฐุปาทานํ ปฏิจฺจ กามุปาทานํ สมฺปยุตฺตกา จ ขนฺธา จิตฺตสมุฏฺฐานญฺจ รูปํ.\csromanspacer{} กามุปาทานํ ฯเปฯ. (สพฺพํ จกฺกํ กาตพฺพํ.) (3)}

\proseitem{2}{โน-อุปาทานํ ธมฺมํ ปฏิจฺจ โน-อุปาทาโน ธมฺโม อุปฺปชฺชติ เหตุปจฺจยา.\csromanspacer{} โน-อุปาทานํ เอกํ ขนฺธํ ปฏิจฺจ ตโย ขนฺธา จิตฺตสมุฏฺฐานญฺจ รูปํ ฯเปฯ เทฺว ขนฺเธ ฯเปฯ.\csromanspacer{} ปฏิสนฺธิกฺขเณ โน-อุปาทานํ เอกํ ขนฺธํ ปฏิจฺจ ตโย ขนฺธา กฏตฺตา จ รูปํ ฯเปฯ เทฺว ขนฺเธ ฯเปฯ.\csromanspacer{} ขนฺเธ ปฏิจฺจ วตฺถุ, วตฺถุํ ปฏิจฺจ ขนฺธา.\csromanspacer{} เอกํ มหาภูตํ ฯเปฯ มหาภูเต ปฏิจฺจ จิตฺตสมุฏฺฐานํ รูปํ, กฏตฺตารูปํ, อุปาทารูปํ. (1)}

\prose{โน-อุปาทานํ ธมฺมํ ปฏิจฺจ อุปาทาโน ธมฺโม อุปฺปชฺชติ เหตุปจฺจยา.\csromanspacer{} โน-อุปาทาเน ขนฺเธ ปฏิจฺจ อุปาทานา. (2)}

\prose{โน-อุปาทานํ ธมฺมํ ปฏิจฺจ อุปาทาโน จ โน-อุปาทาโน จ ธมฺมา อุปฺปชฺชนฺติ เหตุปจฺจยา.\csromanspacer{} โน-อุปาทานํ เอกํ ขนฺธํ ปฏิจฺจ ตโย ขนฺธา อุปาทานา จ จิตฺตสมุฏฺฐานญฺจ รูปํ ฯเปฯ เทฺว ขนฺเธ ปฏิจฺจ เทฺว ขนฺธา อุปาทานา จ จิตฺตสมุฏฺฐานญฺจ รูปํ. (3)}

\csromanheader{2}{69. อุปาทานทุก}
\proseitem{3}{\csromanfolio{561}อุปาทานญฺจ โน-อุปาทานญฺจ ธมฺมํ ปฏิจฺจ อุปาทาโน ธมฺโม อุปฺปชฺชติ เหตุปจฺจยา.\csromanspacer{} ทิฏฺฐุปาทานญฺจ สมฺปยุตฺตเก จ ขนฺเธ ปฏิจฺจ กามุปาทานํ. (สพฺเพ จกฺกา กาตพฺพา.) (1)}

\prose{อุปาทานญฺจ โน-อุปาทานญฺจ ธมฺมํ ปฏิจฺจ โน-อุปาทาโน ธมฺโม อุปฺปชฺชติ เหตุปจฺจยา.\csromanspacer{} โน-อุปาทานํ เอกํ ขนฺธญฺจ อุปาทานญฺจ ปฏิจฺจ ตโย ขนฺธา จิตฺตสมุฏฺฐานญฺจ รูปํ ฯเปฯ เทฺว ขนฺเธ จ ฯเปฯ.\csromanspacer{} อุปาทานญฺจ มหาภูเต จ ปฏิจฺจ จิตฺตสมุฏฺฐานํ รูปํ. (2)}

\prose{อุปาทานญฺจ โน-อุปาทานญฺจ ธมฺมํ ปฏิจฺจ อุปาทาโน จ โน-อุปาทาโน จ ธมฺมา อุปฺปชฺชนฺติ เหตุปจฺจยา.\csromanspacer{} โน-อุปาทานํ เอกํ ขนฺธญฺจ ทิฏฺฐุปาทานญฺจ ปฏิจฺจ ตโย ขนฺธา กามุปาทานญฺจ จิตฺตสมุฏฺฐานํ รูปํ.\csromanspacer{} เทฺว ขนฺเธ จ ฯเปฯ. (จกฺกํ กาตพฺพํ.) (3)}

\csromanheader{2}{\textbf{อารมฺมณ}}
\proseitem{4}{อุปาทานํ ธมฺมํ ปฏิจฺจ อุปาทาโน ธมฺโม อุปฺปชฺชติ อารมฺมณปจฺจยา. (นวปิ ปญฺหา กาตพฺพา.\csromanspacer{} รูปํ ฉฑฺเฑตพฺพํ.)}

\csromanheader{2}{1. ปจฺจยานุโลม 2. สงฺขฺยาวาร}
\proseitem{5}{เหตุยา นว, อารมฺมเณ นว, อธิปติยา นว, (สพฺพตฺถ นว.) วิปาเก เอกํ ฯเปฯ อวิคเต นว.}

\csromanheader{2}{2. ปจฺจยปจฺจนีย 1. วิภงฺควาร}
\csromanheader{2}{\textbf{นเหตุ}}
\proseitem{6}{โน-อุปาทานํ ธมฺมํ ปฏิจฺจ โน-อุปาทาโน ธมฺโม อุปฺปชฺชติ นเหตุปจฺจยา.\csromanspacer{} อเหตุกํ โน-อุปาทานํ เอกํ ขนฺธํ ปฏิจฺจ ตโย ขนฺธา จิตฺตสมุฏฺฐานญฺจ รูปํ ฯเปฯ เทฺว ขนฺเธ ฯเปฯ.\csromanspacer{} อเหตุกปฏิสนฺธิกฺขเณ ฯเปฯ.\csromanspacer{} ขนฺเธ ปฏิจฺจ วตฺถุ, วตฺถุํ ปฏิจฺจ ขนฺธา.\csromanspacer{} เอกํ มหาภูตํ ฯเปฯ มหาภูเต ปฏิจฺจ จิตฺตสมุฏฺฐานํ รูปํ, กฏตฺตารูปํ, อุปาทารูปํ.\csromanspacer{} พาหิรํ.\csromanspacer{} อาหารสมุฏฺฐานํ.\csromanspacer{} อุตุสมุฏฺฐานํ.\csromanspacer{} อสญฺญสตฺตานํ ฯเปฯ.\csromanspacer{} วิจิกิจฺฉาสหคเต อุทฺธจฺจสหคเต ขนฺเธ ปฏิจฺจ วิจิกิจฺฉาสหคโต อุทฺธจฺจสหคโต โมโห. (1)}

\csromanheader{2}{\textbf{น-อารมฺมณาทิ}}
\proseitem{7}{\csromanfolio{562}อุปาทานํ ธมฺมํ ปฏิจฺจ โน-อุปาทาโน ธมฺโม อุปฺปชฺชติ น-อารมฺมณปจฺจยา.\csromanspacer{} อุปาทาเน ปฏิจฺจ จิตฺตสมุฏฺฐานํ รูปํ. (1)}

\prose{โน-อุปาทานํ ธมฺมํ ปฏิจฺจ โน-อุปาทาโน ธมฺโม อุปฺปชฺชติ น-อารมฺมณปจฺจยา.\csromanspacer{} โน-อุปาทาเน ขนฺเธ ปฏิจฺจ จิตฺตสมุฏฺฐานํ รูปํ.\csromanspacer{} ปฏิสนฺธิกฺขเณ โน-อุปาทาเน ขนฺเธ ปฏิจฺจ กฏตฺตารูปํ.\csromanspacer{} ขนฺเธ ปฏิจฺจ วตฺถุ.\csromanspacer{} เอกํ มหาภูตํ ฯเปฯ. (ยาว อสญฺญสตฺตา.) (1)}

\prose{อุปาทานญฺจ โน-อุปาทานญฺจ ธมฺมํ ปฏิจฺจ โน-อุปาทาโน ธมฺโม อุปฺปชฺชติ น-อารมฺมณปจฺจยา.\csromanspacer{} อุปาทาเน จ สมฺปยุตฺตเก จ ขนฺเธ ปฏิจฺจ จิตฺตสมุฏฺฐานํ รูปํ.\csromanspacer{} อุปาทาเน จ มหาภูเต จ ปฏิจฺจ จิตฺตสมุฏฺฐานํ รูปํ. (1)}

\prose{น-อธิปติปจฺจยา.\csromanspacer{} น-อนนฺตรปจฺจยา ฯเปฯ.\csromanspacer{} น-อุปนิสฺสยปจฺจยา.}

\csromanheader{2}{\textbf{นปุเรชาต}}
\proseitem{8}{อุปาทานํ ธมฺมํ ปฏิจฺจ อุปาทาโน ธมฺโม อุปฺปชฺชติ นปุเรชาตปจฺจยา.\csromanspacer{} อรูเป อตฺตวาทุปาทานํ ปฏิจฺจ กามุปาทานํ.\csromanspacer{} กามุปาทานํ ปฏิจฺจ อตฺตวาทุปาทานํ. (1)}

\prose{อุปาทานํ ธมฺมํ ปฏิจฺจ โน-อุปาทาโน ธมฺโม อุปฺปชฺชติ นปุเรชาตปจฺจยา.\csromanspacer{} อรูเป อุปาทาเน ปฏิจฺจ สมฺปยุตฺตกา ขนฺธา.\csromanspacer{} อุปาทาเน ปฏิจฺจ จิตฺตสมุฏฺฐานํ รูปํ. (สํขิตฺตํ.\csromanspacer{} นวปิ ปญฺหา อรูเป เทฺว อุปาทานา.)}

\csromanheader{2}{2. ปจฺจยปจฺจนีย 2. สงฺขฺยาวาร}
\csromanheader{2}{\textbf{สุทฺธ}}
\proseitem{9}{นเหตุยา เอกํ, น-อารมฺมเณ ตีณิ, น-อธิปติยา นว, น-อนนฺตเร ตีณิ ฯเปฯ น-อุปนิสฺสเย ตีณิ, นปุเรชาเต นว, นปจฺฉาชาเต นว, น-อาเสวเน นว, นกมฺเม ตีณิ, นวิปาเก นว, น-อาหาเร เอกํ, น-อินฺทฺริเย เอกํ, นฌาเน เอกํ, นมคฺเค เอกํ, นสมฺปยุตฺเต ตีณิ, นวิปฺปยุตฺเต นว, โนนตฺถิยา ตีณิ, โนวิคเต ตีณิ.}

\csromanheader{2}{69. อุปาทานทุก \textbf{3. ปจฺจยานุโลมปจฺจนีย}}
\proseitem{10}{\csromanfolio{563}\textbf{เหตุ}ปจฺจยา น-อารมฺมเณ ตีณิ, น-อธิปติยา นว ฯเปฯ น-อุปนิสฺสเย ตีณิ, นปุเรชาเต นว, นปจฺฉาชาเต นว, น-อาเสวเน นว, นกมฺเม ตีณิ, นวิปาเก นว, นสมฺปยุตฺเต ตีณิ, นวิปฺปยุตฺเต นว, โนนตฺถิยา ตีณิ, โนวิคเต ตีณิ.}

\csromanheader{2}{4. ปจฺจยปจฺจนียานุโลม}
\proseitem{11}{นเหตุปจฺจยา อารมฺมเณ เอกํ, (สพฺพตฺถ เอกํ.) มคฺเค เอกํ ฯเปฯ อวิคเต เอกํ.}

\prose{(สหชาตวาโร ปฏิจฺจวารสทิโส.\csromanspacer{} วิภชนฺเตน ทิฏฺฐุปาทานํ “สหชาตํ กามุปาทานนฺ”ติ กาตพฺพํ.)}

\csromanheader{2}{69. อุปาทานทุก 3. ปจฺจยวาร}
\csromanheader{2}{1. ปจฺจยานุโลมาทิ}
\csromanheader{2}{\textbf{เหตุ}}
\proseitem{12}{อุปาทานํ ธมฺมํ ปจฺจยา อุปาทาโน ธมฺโม อุปฺปชฺชติ เหตุปจฺจยา.\csromanspacer{} ทิฏฺฐุปาทานํ ปจฺจยา กามุปาทานํ.\csromanspacer{} ตีณิ. (ปฏิจฺจสทิสา.)}

\prose{โน-อุปาทานํ ธมฺมํ ปจฺจยา โน-อุปาทาโน ธมฺโม อุปฺปชฺชติ เหตุปจฺจยา.\csromanspacer{} โน-อุปาทานํ เอกํ ขนฺธํ ปจฺจยา ตโย ขนฺธา จิตฺตสมุฏฺฐานญฺจ รูปํ ฯเปฯ.\csromanspacer{} เทฺว ขนฺเธ ฯเปฯ.\csromanspacer{} ปฏิสนฺธิกฺขเณ ฯเปฯ.\csromanspacer{} ขนฺเธ ปจฺจยา วตฺถุ ฯเปฯ. (ยาว อชฺฌตฺติกา มหาภูตา.) วตฺถุํ ปจฺจยา โน-อุปาทานา ขนฺธา. (1)}

\prose{โน-อุปาทานํ ธมฺมํ ปจฺจยา อุปาทาโน ธมฺโม อุปฺปชฺชติ เหตุปจฺจยา.\csromanspacer{} โน-อุปาทาเน ขนฺเธ ปจฺจยา อุปาทานา.\csromanspacer{} วตฺถุํ ปจฺจยา อุปาทานา. (2)}

\prose{โน-อุปาทานํ ธมฺมํ ปจฺจยา อุปาทาโน จ โน-อุปาทาโน จ ธมฺมา อุปฺปชฺชนฺติ เหตุปจฺจยา.\csromanspacer{} โน-อุปาทานํ เอกํ ขนฺธํ ปจฺจยา ตโย ขนฺธา อุปาทานา จ จิตฺตสมุฏฺฐานญฺจ รูปํ ฯเปฯ เทฺว ขนฺเธ ฯเปฯ.\csromanspacer{} วตฺถุํ ปจฺจยา อุปาทานา.\csromanspacer{} มหาภูเต ปจฺจยา จิตฺตสมุฏฺฐานํ รูปํ.\csromanspacer{} วตฺถุํ ปจฺจยา อุปาทานา สมฺปยุตฺตกา จ ขนฺธา. (3)}

\proseitem{13}{\csromanfolio{564}อุปาทานญฺจ โน-อุปาทานญฺจ ธมฺมํ ปจฺจยา อุปาทาโน ธมฺโม อุปฺปชฺชติ เหตุปจฺจยา.\csromanspacer{} ทิฏฺฐุปาทานญฺจ สมฺปยุตฺตเก จ ขนฺเธ ปจฺจยา กามุปาทานํ.\csromanspacer{} กามุปาทานญฺจ สมฺปยุตฺตเก จ ขนฺเธ ปจฺจยา ทิฏฺฐุปาทานํ. (จกฺกํ กาตพฺพํ.) ทิฏฺฐุปาทานญฺจ วตฺถุญฺจ ปจฺจยา กามุปาทานํ. (จกฺกํ กาตพฺพํ.) (1)}

\prose{อุปาทานญฺจ โน-อุปาทานญฺจ ธมฺมํ ปจฺจยา โน-อุปาทาโน ธมฺโม อุปฺปชฺชติ เหตุปจฺจยา.\csromanspacer{} โน-อุปาทานํ เอกํ ขนฺธญฺจ อุปาทานญฺจ ปจฺจยา ตโย ขนฺธา จิตฺตสมุฏฺฐานญฺจ รูปํ ฯเปฯ เทฺว ขนฺเธ จ ฯเปฯ. (จกฺกํ กาตพฺพํ.) อุปาทาเน จ มหาภูเต จ ปจฺจยา จิตฺตสมุฏฺฐานํ รูปํ.\csromanspacer{} อุปาทานญฺจ วตฺถุญฺจ ปจฺจยา โน-อุปาทาโน ขนฺธา. (2)}

\prose{อุปาทานญฺจ โน-อุปาทานญฺจ ธมฺมํ ปจฺจยา อุปาทาโน จ โน-อุปาทาโน จ ธมฺมา อุปฺปชฺชนฺติ เหตุปจฺจยา.\csromanspacer{} โน-อุปาทานํ เอกํ ขนฺธญฺจ ทิฏฺฐุปาทานญฺจ ปจฺจยา ตโย ขนฺธา กามุปาทานํ จิตฺตสมุฏฺฐานญฺจ รูปํ ฯเปฯ. (จกฺกํ.) ทิฏฺฐุปาทานญฺจ วตฺถุญฺจ ปจฺจยา กามุปาทานํ สมฺปยุตฺตกา จ ขนฺธา ฯเปฯ. (จกฺกํ.) (3)}

\prose{อารมฺมณปจฺจยา. (อารมฺมเณ โน-อุปาทานมูลเก ปญฺจายตนญฺจ วตฺถุญฺจ กาตพฺพา.)}

\prose{เหตุยา นว, อารมฺมเณ นว, อธิปติยา นว, (สพฺพตฺถ นว,) วิปาเก เอกํ ฯเปฯ อวิคเต นว.\csromanspacer{} อนุโลมํ.}

\proseitem{14}{โน-อุปาทานํ ธมฺมํ ปจฺจยา โน-อุปาทาโน ธมฺโม อุปฺปชฺชติ นเหตุปจฺจยา.\csromanspacer{} อเหตุกํ โน-อุปาทานํ เอกํ ขนฺธํ ปจฺจยา ฯเปฯ.\csromanspacer{} อเหตุกปฏิสนฺธิกฺขเณ ฯเปฯ. (ยาว อสญฺญสตฺตา.) จกฺขายตนํ ปจฺจยา จกฺขุวิญฺญาณํ ฯเปฯ.\csromanspacer{} กายายตนํ ปจฺจยา กายวิญฺญาณํ.\csromanspacer{} วตฺถุํ ปจฺจยา อเหตุกา โน-อุปาทานา ขนฺธา.\csromanspacer{} วิจิกิจฺฉาสหคเต อุทฺธจฺจสหคเต ขนฺเธ จ วตฺถุญฺจ ปจฺจยา วิจิกิจฺฉาสหคโต อุทฺธจฺจสหคโต โมโห.}

\prose{นเหตุยา เอกํ, น-อารมฺมเณ ตีณิ, น-อธิปติยา นว, น-อนนฺตเร ตีณิ, นสมนนฺตเร ตีณิ ฯเปฯ นปุเรชาเต นว ฯเปฯ นกมฺเม ตีณิ, นวิปาเก นว, (ปฏิจฺจสทิสํ.) ฯเปฯ โนวิคเต ตีณิ.}

\vspace{\tipitakaparskip}
\csromancenter{ปจฺจนียํ.}
\csromancenter{(เอวํ อิตเร เทฺว คณนาปิ นิสฺสยวาโรปิ กาตพฺโพ.)}
\csromanheader{2}{69. อุปาทานทุก \textbf{69. อุปาทานทุก 5. สํสฏฺฐวาร}}
\csromanheader{2}{1. ปจฺจยานุโลมาทิ}
\csromanheader{2}{\textbf{เหตุ}}
\proseitem{15}{\csromanfolio{565}อุปาทานํ ธมฺมํ สํสฏฺโฐ อุปาทาโน ธมฺโม อุปฺปชฺชติ เหตุปจฺจยา.\csromanspacer{} ทิฏฺฐุปาทานํ สํสฏฺฐํ กามุปาทานํ.\csromanspacer{} กามุปาทานํ สํสฏฺฐํ ทิฏฺฐุปาทานํ. (จกฺกํ.\csromanspacer{} เอวํ นวปิ ปญฺหา กาตพฺพา.)}

\vspace{\tipitakaparskip}
\csromancenter{\textbf{เหตุ}ยา นว, อารมฺมเณ นว, อธิปติยา นว, (สพฺพตฺถ นว.) วิปาเก เอกํ ฯเปฯ วิคเต นว, อวิคเต นว.\csromanspacer{} อนุโลมํ.}
\proseitem{16}{โน-อุปาทานํ ธมฺมํ สํสฏฺโฐ โน-อุปาทาโน ธมฺโม อุปฺปชฺชติ นเหตุปจฺจยา.\csromanspacer{} อเหตุกํ โน-อุปาทานํ เอกํ ขนฺธํ สํสฏฺฐา ตโย ขนฺธา ฯเปฯ.\csromanspacer{} เทฺว ขนฺเธ ฯเปฯ.\csromanspacer{} อเหตุกปฏิสนฺธิกฺขเณ ฯเปฯ.\csromanspacer{} วิจิกิจฺฉาสหคเต อุทฺธจฺจสหคเต ขนฺเธ สํสฏฺโฐ วิจิกิจฺฉาสหคโต อุทฺธจฺจสหคโต โมโห.}

\prose{นเหตุยา เอกํ, น-อธิปติยา นว, นปุเรชาเต นว, นปจฺฉาชาเต นว, น-อาเสวเน นว, นกมฺเม ตีณิ, นวิปาเก นว, นฌาเน เอกํ, นมคฺเค เอกํ, นวิปฺปยุตฺเต นว.}

\vspace{\tipitakaparskip}
\csromancenter{ปจฺจนียํ.}
\csromancenter{(เอวํ อิตเร เทฺว คณนาปิ สมฺปยุตฺตวาโรปิ กาตพฺโพ.)}
\csromanheader{2}{69. อุปาทานทุก 7. ปญฺหาวาร}
\csromanheader{2}{1. ปจฺจยานุโลม 1. วิภงฺควาร}
\csromanheader{2}{\textbf{เหตุ}}
\proseitem{17}{อุปาทาโน ธมฺโม อุปาทานสฺส ธมฺมสฺส เหตุปจฺจเยน ปจฺจโย.\csromanspacer{} อุปาทานา เหตู สมฺปยุตฺตกานํ อุปาทานานํ เหตุปจฺจเยน ปจฺจโย. (มูลํ กาตพฺพํ.) อุปาทานา เหตู สมฺปยุตฺตกานํ ขนฺธานํ จิตฺตสมุฏฺฐานานญฺจ \csromanfolio{566}รูปานํ เหตุปจฺจเยน ปจฺจโย. (มูลํ กาตพฺพํ.) อุปาทานา เหตู สมฺปยุตฺตกานํ ขนฺธานํ อุปาทานานญฺจ จิตฺตสมุฏฺฐานานํ รูปานํ เหตุปจฺจเยน ปจฺจโย. (3)}

\proseitem{18}{โน-อุปาทาโน ธมฺโม โน-อุปาทานสฺส ธมฺมสฺส เหตุปจฺจเยน ปจฺจโย.\csromanspacer{} โน-อุปาทานา เหตู สมฺปยุตฺตกานํ ขนฺธานํ จิตฺตสมุฏฺฐานานญฺจ รูปานํ เหตุปจฺจเยน ปจฺจโย.\csromanspacer{} ปฏิสนฺธิกฺขเณ ฯเปฯ. (1)}

\prose{โน-อุปาทาโน ธมฺโม อุปาทานสฺส ธมฺมสฺส เหตุปจฺจเยน ปจฺจโย.\csromanspacer{} โน-อุปาทานา เหตู สมฺปยุตฺตกานํ อุปาทานานํ เหตุปจฺจเยน ปจฺจโย. (มูลํ กาตพฺพํ.) โน-อุปาทานา เหตู สมฺปยุตฺตกานํ ขนฺธานํ อุปาทานานญฺจ จิตฺตสมุฏฺฐานานํ รูปานํ เหตุปจฺจเยน ปจฺจโย. (3)}

\proseitem{19}{อุปาทาโน จ โน-อุปาทาโน จ ธมฺมา อุปาทานสฺส ธมฺมสฺส เหตุปจฺจเยน ปจฺจโย.\csromanspacer{} อุปาทานา จ โน-อุปาทานา จ เหตู สมฺปยุตฺตกานํ อุปาทานานํ เหตุปจฺจเยน ปจฺจโย. (มูลํ กาตพฺพํ.) อุปาทานา จ โน-อุปาทานา จ เหตู สมฺปยุตฺตกานํ ขนฺธานํ จิตฺตสมุฏฺฐานญฺจ รูปานํ เหตุปจฺจเยน ปจฺจโย. (มูลํ กาตพฺพํ.) อุปาทานา จ โน-อุปาทานา จ เหตู สมฺปยุตฺตกานํ ขนฺธานํ อุปาทานานญฺจ จิตฺตสมุฏฺฐานานํ รูปานํ เหตุปจฺจเยน ปจฺจโย. (3)}

\csromanheader{2}{\textbf{อารมฺมณ}}
\proseitem{20}{อุปาทาโน ธมฺโม อุปาทานสฺส ธมฺมสฺส อารมฺมณปจฺจเยน ปจฺจโย.\csromanspacer{} อุปาทาเน อารพฺภ อุปาทานา อุปฺปชฺชนฺติ.\csromanspacer{} ตีณิ. (อารพฺภ กาตพฺพา.)}

\prose{โน-อุปาทาโน ธมฺโม โน-อุปาทานสฺส ธมฺมสฺส อารมฺมณปจฺจเยน ปจฺจโย.\csromanspacer{} ทานํ ทตฺวา, สีลํ ฯเปฯ อุโปสถกมฺมํ กตฺวา ตํ ปจฺจเวกฺขติ, อสฺสาเทติ อภินนฺทติ, ตํ อารพฺภ ราโค อุปฺปชฺชติ.\csromanspacer{} ทิฏฺฐิ อุปฺปชฺชติ.\csromanspacer{} วิจิกิจฺฉา ฯเปฯ.\csromanspacer{} อุทฺธจฺจํ ฯเปฯ.\csromanspacer{} โทมนสฺสํ อุปฺปชฺชติ.\csromanspacer{} ปุพฺเพ สุจิณฺณานิ ฯเปฯ.\csromanspacer{} ฌานา วุฏฺฐหิตฺวา ฌานํ ปจฺจเวกฺขติ.\csromanspacer{} อริยา มคฺคา วุฏฺฐหิตฺวา มคฺคํ ปจฺจเวกฺขนฺติ.\csromanspacer{} ผลํ ฯเปฯ.\csromanspacer{} นิพฺพานํ ปจฺจเวกฺขนฺติ.\csromanspacer{} นิพฺพานํ โคตฺรภุสฺส, โวทานสฺส, มคฺคสฺส, ผลสฺส, อาวชฺชนาย อารมฺมณปจฺจเยน ปจฺจโย.\csromanspacer{} อริยา โน-อุปาทาเน ปหีเน}

\vspace{\tipitakaparskip}
\csromancenter{อุปาทานทุก กิเลเส ฯเปฯ.\csromanspacer{} วิกฺขมฺภิเต กิเลเส ปจฺจเวกฺขนฺติ, ปุพฺเพ สมุทาจิณฺเณ ฯเปฯ.\csromanspacer{} จกฺขุํ ฯเปฯ วตฺถุํ โน-อุปาทาเน ขนฺเธ อนิจฺจโต ฯเปฯ โทมนสฺสํ อุปฺปชฺชติ.\csromanspacer{} ทิพฺเพน จกฺขุนา รูปํ ปสฺสติ.\csromanspacer{} ทิพฺพาย โสตธาตุยา สทฺทํ สุณาติ.\csromanspacer{} เจโตปริยญาเณน โน-อุปาทานจิตฺตสมงฺคิสฺส จิตฺตํ ชานาติ.\csromanspacer{} อากาสานญฺจายตนํ วิญฺญาณญฺจายตนสฺส ฯเปฯ.\csromanspacer{} อากิญฺจญฺญายตนํ เนวสญฺญานาสญฺญายตนสฺส ฯเปฯ.\csromanspacer{} รูปายตนํ จกฺขุวิญฺญาณสฺส ฯเปฯ.\csromanspacer{} โผฏฺฐพฺพายตนํ กายวิญฺญาณสฺส ฯเปฯ.\csromanspacer{} โน-อุปาทานา ขนฺธา อิทฺธิวิธญาณสฺส, เจโตปริยญาณสฺส, ปุพฺเพนิวาสานุสฺสติญาณสฺส, ยถากมฺมูปคญาณสฺส, อนาคตํสญาณสฺส, อาวชฺชนาย อารมฺมณปจฺจเยน ปจฺจโย. (1)}
\prose{\csromanfolio{567}โน-อุปาทาโน ธมฺโม อุปาทานสฺส ธมฺมสฺส อารมฺมณปจฺจเยน ปจฺจโย.\csromanspacer{} ทานํ ฯเปฯ สีลํ ฯเปฯ อุโปสถกมฺมํ กตฺวา ตํ อสฺสาเทติ อภินนฺทติ, ตํ อารพฺภ ราโค อุปฺปชฺชติ.\csromanspacer{} ทิฏฺฐิ อุปฺปชฺชติ.\csromanspacer{} ปุพฺเพ สุจิณฺณานิ ฯเปฯ.\csromanspacer{} ฌานา ฯเปฯ.\csromanspacer{} จกฺขุํ ฯเปฯ วตฺถุํ โน-อุปาทาเน ขนฺเธ อสฺสาเทติ อภินนฺทติ, ตํ อารพฺภ ราโค อุปฺปชฺชติ.\csromanspacer{} ทิฏฺฐิ อุปฺปชฺชติ. (2)}

\prose{โน-อุปาทาโน ธมฺโม อุปาทานสฺส จ โน-อุปาทานสฺส จ ธมฺมสฺส อารมฺมณปจฺจเยน ปจฺจโย.\csromanspacer{} ทานํ ฯเปฯ สีลํ ฯเปฯ อุโปสถกมฺมํ กตฺวา ฯเปฯ.\csromanspacer{} ปุพฺเพ สุจิณฺณานิ ฯเปฯ.\csromanspacer{} ฌานา วุฏฺฐหิตฺวา ฌานํ ฯเปฯ.\csromanspacer{} จกฺขุํ ฯเปฯ วตฺถุํ โน-อุปาทาเน ขนฺเธ อสฺสาเทติ อภินนฺทติ, ตํ อารพฺภ อุปาทานา จ สมฺปยุตฺตกา จ ขนฺธา อุปฺปชฺชนฺติ. (3)}

\prose{อุปาทาโน จ โน-อุปาทาโน จ ธมฺมา อุปาทานสฺส ธมฺมสฺส อารมฺมณปจฺจเยน ปจฺจโย.\csromanspacer{} ตีณิ. (อารพฺภ กาตพฺพา.)}

\csromanheader{2}{\textbf{อธิปติ}}
\proseitem{21}{อุปาทาโน ธมฺโม อุปาทานสฺส ธมฺมสฺส อธิปติปจฺจเยน ปจฺจโย.\csromanspacer{} \textbf{อารมฺมณาธิปติ}\csromandash{}อุปาทาเน ครุํ กตฺวา อุปาทานา อุปฺปชฺชนฺติ.\csromanspacer{} ตีณิ. (\textbf{อารมฺมณาธิปติ}เยว.)}

\proseitem{22}{โน-อุปาทาโน ธมฺโม โน-อุปาทานสฺส ธมฺมสฺส อธิปติปจฺจเยน ปจฺจโย.\csromanspacer{} \textbf{อารมฺมณาธิปติ}, สหชาตาธิปติ.\csromanspacer{}\csromanspacer{}\textbf{อารมฺมณาธิปติ}\csromandash{}ทานํ ทตฺวา, สีลํ ฯเปฯ อุโปสถกมฺมํ ฯเปฯ.\csromanspacer{} ปุพฺเพ ฯเปฯ.\csromanspacer{} ฌานา วุฏฺฐหิตฺวา ฌานํ ครุํ กตฺวา ปจฺจเวกฺขติ, อสฺสาเทติ อภินนฺทติ, อริยา มคฺคา \csromanfolio{568}วุฏฺฐหิตฺวา มคฺคํ ครุํ กตฺวา ปจฺจเวกฺขนฺติ ฯเปฯ ผลสฺส อธิปติปจฺจเยน ปจฺจโย.\csromanspacer{} จกฺขุํ ฯเปฯ วตฺถุํ โน-อุปาทาเน ขนฺเธ ครุํ กตฺวา อสฺสาเทติ อภินนฺทติ, ตํ ครุํ กตฺวา ราโค อุปฺปชฺชติ.\csromanspacer{} ทิฏฺฐิ อุปฺปชฺชติ.\csromanspacer{} \textbf{สหชาตาธิปติ}\csromandash{}โน-อุปาทานาธิปติ สมฺปยุตฺตกานํ ขนฺธานํ จิตฺตสมุฏฺฐานานญฺจ รูปานํ อธิปติปจฺจเยน ปจฺจโย. (1)}

\prose{โน-อุปาทาโน ธมฺโม อุปาทานสฺส ธมฺมสฺส อธิปติปจฺจเยน ปจฺจโย.\csromanspacer{} \textbf{อารมฺมณาธิปติ}, สหชาตาธิปติ.\csromanspacer{}\csromanspacer{}\textbf{อารมฺมณาธิปติ\csromandash{}}ทานํ ฯเปฯ สีลํ ฯเปฯ อุโปสถกมฺมํ กตฺวา ตํ ครุํ กตฺวา อสฺสาเทติ อภินนฺทติ, ตํ ครุํ กตฺวา ราโค อุปฺปชฺชติ.\csromanspacer{} ทิฏฺฐิ อุปฺปชฺชติ.\csromanspacer{} ปุพฺเพ ฯเปฯ.\csromanspacer{} ฌานา ฯเปฯ.\csromanspacer{} จกฺขุํ ฯเปฯ วตฺถุํ โน-อุปาทาเน ขนฺเธ ครุํ กตฺวา อสฺสาเทติ อภินนฺทติ, ตํ ครุํ กตฺวา ราโค อุปฺปชฺชติ.\csromanspacer{} ทิฏฺฐิ อุปฺปชฺชติ.\csromanspacer{} \textbf{สหชาตาธิปติ}\csromandash{}โน-อุปาทานาธิปติ สมฺปยุตฺตกานํ อุปาทานานํ อธิปติปจฺจเยน ปจฺจโย. (2)}

\prose{โน-อุปาทาโน ธมฺโม อุปาทานสฺส จ โน-อุปาทานสฺส จ ธมฺมสฺส อธิปติปจฺจเยน ปจฺจโย.\csromanspacer{} \textbf{อารมฺมณาธิปติ}, สหชาตาธิปติ.\csromanspacer{}\csromanspacer{}\textbf{อารมฺมณาธิปติ\csromandash{}}ทานํ ฯเปฯ สีลํ ฯเปฯ อุโปสถกมฺมํ ฯเปฯ.\csromanspacer{} ปุพฺเพ ฯเปฯ.\csromanspacer{} ฌานํ ฯเปฯ.\csromanspacer{} จกฺขุํ ฯเปฯ วตฺถุํ โน-อุปาทาเน ขนฺเธ ครุํ กตฺวา อสฺสาเทติ อภินนฺทติ, ตํ ครุํ กตฺวา อุปาทานา จ สมฺปยุตฺตกา จ ขนฺธา อุปฺปชฺชนฺติ.\csromanspacer{} \textbf{สหชาตาธิปติ}\csromandash{}โน-อุปาทานาธิปติ สมฺปยุตฺตกานํ ขนฺธานํ อุปาทานานญฺจ จิตฺตสมุฏฺฐานานํ รูปานํ อธิปติปจฺจเยน ปจฺจโย. (3)}

\prose{อุปาทาโน จ โน-อุปาทาโน จ ธมฺมา อุปาทานสฺส ธมฺมสฺส อธิปติปจฺจเยน ปจฺจโย.\csromanspacer{} ตีณิ.\csromanspacer{} \textbf{อารมฺมณาธิปติ}.\csromanspacer{} ตีณิ. (อารพฺภ กาตพฺพา.\csromanspacer{} \textbf{อารมฺมณาธิปติ}เยว.)}

\csromanheader{2}{\textbf{อนนฺตราทิ}}
\proseitem{23}{อุปาทาโน ธมฺโม อุปาทานสฺส ธมฺมสฺส อนนฺตรปจฺจเยน ปจฺจโย.\csromanspacer{} ปุริมา ปุริมา อุปาทานา ปจฺฉิมานํ ปจฺฉิมานํ อุปาทานานํ อนนฺตรปจฺจเยน ปจฺจโย. (มูลํ กาตพฺพํ.) ปุริมา ปุริมา อุปาทานา ปจฺฉิมานํ ปจฺฉิมานํ โน-อุปาทานานํ ขนฺธานํ อนนฺตรปจฺจเยน ปจฺจโย.\csromanspacer{} อุปาทานํ วุฏฺฐานสฺส อนนฺตรปจฺจเยน ปจฺจโย. (มูลํ กาตพฺพํ.) ปุริมา ปุริมา อุปาทานา ปจฺฉิมานํ ปจฺฉิมานํ อุปาทานานํ สมฺปยุตฺตกานญฺจ ขนฺธานํ อนนฺตรปจฺจเยน ปจฺจโย. (3)}

\csromanheader{2}{69. อุปาทานทุก}
\proseitem{24}{\csromanfolio{569}โน-อุปาทาโน ธมฺโม โน-อุปาทานสฺส ธมฺมสฺส อนนฺตรปจฺจเยน ปจฺจโย.\csromanspacer{} ปุริมา ปุริมา โน-อุปาทานา ขนฺธา ปจฺฉิมานํ ปจฺฉิมานํ โน-อุปาทานานํ ขนฺธานํ อนนฺตรปจฺจเยน ปจฺจโย.\csromanspacer{} อนุโลมํ โคตฺรภุสฺส ฯเปฯ.\csromanspacer{} ผลสมาปตฺติยา อนนฺตรปจฺจเยน ปจฺจโย. (1)}

\prose{โน-อุปาทาโน ธมฺโม อุปาทานสฺส ธมฺมสฺส อนนฺตรปจฺจเยน ปจฺจโย.\csromanspacer{} ปุริมา ปุริมา โน-อุปาทานา ขนฺธา ปจฺฉิมานํ ปจฺฉิมานํ อุปาทานานํ ขนฺธานํ อนนฺตรปจฺจเยน ปจฺจโย.\csromanspacer{} อาวชฺชนา อุปาทานานํ อนนฺตรปจฺจเยน ปจฺจโย. (2)}

\prose{โน-อุปาทาโน ธมฺโม อุปาทานสฺส จ โน-อุปาทานสฺส จ ธมฺมสฺส อนนฺตรปจฺจเยน ปจฺจโย.\csromanspacer{} ปุริมา ปุริมา โน-อุปาทานา ขนฺธา ปจฺฉิมานํ ปจฺฉิมานํ อุปาทานานํ สมฺปยุตฺตกานญฺจ ขนฺธานํ อนนฺตรปจฺจเยน ปจฺจโย.\csromanspacer{} อาวชฺชนา อุปาทานานํ สมฺปยุตฺตกานญฺจ ขนฺธานํ อนนฺตรปจฺจเยน ปจฺจโย. (3)}

\proseitem{25}{อุปาทาโน จ โน-อุปาทาโน จ ธมฺมา อุปาทานสฺส ธมฺมสฺส อนนฺตรปจฺจเยน ปจฺจโย.\csromanspacer{} ปุริมา ปุริมา อุปาทานา จ สมฺปยุตฺตกา จ ขนฺธา ปจฺฉิมานํ ปจฺฉิมานํ อุปาทานานํ อนนฺตรปจฺจเยน ปจฺจโย. (มูลํ กาตพฺพํ.) ปุริมา ปุริมา อุปาทานา จ สมฺปยุตฺตกา จ ขนฺธา ปจฺฉิมานํ ปจฺฉิมานํ โน-อุปาทานานํ ขนฺธานํ อนนฺตรปจฺจเยน ปจฺจโย.\csromanspacer{} อุปาทานา จ สมฺปยุตฺตกา จ ขนฺธา วุฏฺฐานสฺส อนนฺตรปจฺจเยน ปจฺจโย. (มูลํ กาตพฺพํ.) ปุริมา ปุริมา อุปาทานา จ สมฺปยุตฺตกา จ ขนฺธา ปจฺฉิมานํ ปจฺฉิมานํ อุปาทานานํ สมฺปยุตฺตกานญฺจ ขนฺธานํ อนนฺตรปจฺจเยน ปจฺจโย. (3)}

\prose{สมนนฺตรปจฺจเยน ปจฺจโย.\csromanspacer{} สหชาตปจฺจเยน ปจฺจโย. (ปฏิจฺจสทิสํ.) อญฺญมญฺญปจฺจเยน ปจฺจโย. (ปฏิจฺจสทิสํ.) นิสฺสยปจฺจเยน ปจฺจโย. (ปจฺจยสทิสํ.)}

\csromanheader{2}{\textbf{อุปนิสฺสย}}
\proseitem{26}{อุปาทาโน ธมฺโม อุปาทานสฺส ธมฺมสฺส อุปนิสฺสยปจฺจเยน ปจฺจโย.\csromanspacer{} อารมฺมณูปนิสฺสโย, อนนฺตรูปนิสฺสโย, ปกตูปนิสฺสโย ฯเปฯ.\csromanspacer{} ปกตูปนิสฺสโย\csromandash{}อุปาทานา อุปาทานานํ อุปนิสฺสยปจฺจเยน ปจฺจโย.\csromanspacer{} ตีณิ.}

\proseitem{27}{\csromanfolio{570}โน-อุปาทาโน ธมฺโม โน-อุปาทานสฺส ธมฺมสฺส อุปนิสฺสยปจฺจเยน ปจฺจโย.\csromanspacer{} อารมฺมณูปนิสฺสโย, อนนฺตรูปนิสฺสโย, ปกตูปนิสฺสโย ฯเปฯ.\csromanspacer{} \textbf{ปกตูปนิสฺสโย}\csromandash{}สทฺธํ อุปนิสฺสาย ทานํ เทติ ฯเปฯ สมาปตฺติํ อุปฺปาเทติ.\csromanspacer{} มานํ ชปฺเปติ, ทิฏฺฐิํ คณฺหาติ.\csromanspacer{} สีลํ ฯเปฯ.\csromanspacer{} ปญฺญํ.\csromanspacer{} ราคํ ฯเปฯ.\csromanspacer{} ปตฺถนํ.\csromanspacer{} กายิกํ สุขํ ฯเปฯ.\csromanspacer{} เสนาสนํ อุปนิสฺสาย ทานํ เทติ ฯเปฯ สมาปตฺติํ อุปฺปาเทติ.\csromanspacer{} มานํ ชปฺเปติ.\csromanspacer{} ทิฏฺฐิํ คณฺหาติ.\csromanspacer{} ปาณํ หนติ ฯเปฯ สํฆํ ภินฺทติ.\csromanspacer{} สทฺธา ฯเปฯ.\csromanspacer{} เสนาสนํ สทฺธาย ฯเปฯ ผลสมาปตฺติยา อุปนิสฺสยปจฺจเยน ปจฺจโย. (1)}

\prose{โน-อุปาทาโน ธมฺโม อุปาทานสฺส ธมฺมสฺส อุปนิสฺสยปจฺจเยน ปจฺจโย.\csromanspacer{} \textbf{ปกตูปนิสฺสโย}\csromandash{}สทฺธํ อุปนิสฺสาย มานํ ชปฺเปติ, ทิฏฺฐิํ คณฺหาติ.\csromanspacer{} สีลํ ฯเปฯ.\csromanspacer{} เสนาสนํ อุปนิสฺสาย อทินฺนํ ฯเปฯ มุสา ฯเปฯ ปิสุณํ ฯเปฯ สมฺผํ ฯเปฯ สนฺธิํ ฯเปฯ นิลฺโลปํ ฯเปฯ เอกาคาริกํ ฯเปฯ ปริปนฺเถ ฯเปฯ ปรทารํ ฯเปฯ คามฆาตํ ฯเปฯ นิคมฆาตํ กโรติ.\csromanspacer{} สทฺธา ฯเปฯ.\csromanspacer{} เสนาสนํ อุปาทานานํ อุปนิสฺสยปจฺจเยน ปจฺจโย. (มูลํ กาตพฺพํ.) สทฺธํ อุปนิสฺสาย มานํ ชปฺเปติ, ทิฏฺฐิํ คณฺหาติ.\csromanspacer{} สีลํ ฯเปฯ.\csromanspacer{} เสนาสนํ อุปนิสฺสาย อทินฺนํ อาทิยติ, มุสา ฯเปฯ ปิสุณํ ฯเปฯ สมฺผํ ฯเปฯ สนฺธิํ ฯเปฯ นิลฺโลปํ ฯเปฯ เอกาคาริกํ ฯเปฯ ปริปนฺเถ ฯเปฯ ปรทารํ ฯเปฯ คามฆาตํ ฯเปฯ นิคมฆาตํ กโรติ.\csromanspacer{} สทฺธา ฯเปฯ.\csromanspacer{} เสนาสนํ อุปาทานานํ สมฺปยุตฺตกานญฺจ ขนฺธานํ อุปนิสฺสยปจฺจเยน ปจฺจโย. (3)}

\prose{อุปาทาโน จ โน-อุปาทาโน จ ธมฺมา อุปาทานสฺส ธมฺมสฺส อุปนิสฺสยปจฺจเยน ปจฺจโย.\csromanspacer{} อารมฺมณูปนิสฺสโย, อนนฺตรูปนิสฺสโย, ปกตูปนิสฺสโย ฯเปฯ.\csromanspacer{} \textbf{ปกตูปนิสฺสโย}\csromandash{}ตีณิ.}

\csromanheader{2}{\textbf{ปุเรชาต}}
\proseitem{28}{โน-อุปาทาโน ธมฺโม โน-อุปาทานสฺส ธมฺมสฺส ปุเรชาตปจฺจเยน ปจฺจโย.\csromanspacer{} \textbf{อารมฺมณปุเรชาตํ}, วตฺถุปุเรชาตํ.\csromanspacer{}\csromanspacer{}\textbf{อารมฺมณปุเรชาตํ}\csromandash{} จกฺขุํ ฯเปฯ วตฺถุํ อนิจฺจโต ฯเปฯ โทมนสฺสํ อุปฺปชฺชติ.\csromanspacer{} ทิพฺเพน จกฺขุนา รูปํ ปสฺสติ.\csromanspacer{} ทิพฺพาย โสตธาตุยา สทฺทํ สุณาติ.\csromanspacer{} รูปายตนํ จกฺขุวิญฺญาณสฺส ฯเปฯ.\csromanspacer{} โผฏฺฐพฺพายตนํ กายวิญฺญาณสฺส ฯเปฯ.\csromanspacer{} \textbf{วตฺถุปุเรชาตํ}\csromandash{}จกฺขายตนํ จกฺขุวิญฺญาณสฺส ฯเปฯ กายายตนํ กายวิญฺญาณสฺส ฯเปฯ.\csromanspacer{} วตฺถุ โน-อุปาทานานํ ขนฺธานํ ปุเรชาตปจฺจเยน ปจฺจโย. (1)}

\csromanheader{2}{69. อุปาทานทุก}
\prose{\csromanfolio{571}โน-อุปาทาโน ธมฺโม อุปาทานสฺส ธมฺมสฺส ปุเรชาตปจฺจเยน ปจฺจโย.\csromanspacer{} \textbf{อารมฺมณปุเรชาตํ}, วตฺถุปุเรชาตํ.\csromanspacer{}\csromanspacer{}\textbf{อารมฺมณปุเรชาตํ}\csromandash{} จกฺขุํ ฯเปฯ วตฺถุํ อสฺสาเทติ อภินนฺทติ, ตํ อารพฺภ ราโค อุปฺปชฺชติ.\csromanspacer{} ทิฏฺฐิ อุปฺปชฺชติ.\csromanspacer{} \textbf{วตฺถุปุเรชาตํ}\csromandash{}วตฺถุ อุปาทานานํ ปุเรชาตปจฺจเยน ปจฺจโย. (2)}

\prose{โน-อุปาทาโน ธมฺโม อุปาทานสฺส จ โน-อุปาทานสฺส จ ธมฺมสฺส ปุเรชาตปจฺจเยน ปจฺจโย.\csromanspacer{} \textbf{อารมฺมณปุเรชาตํ}, วตฺถุปุเรชาตํ.\csromanspacer{}\csromanspacer{}\textbf{อารมฺมณปุเรชาตํ}\csromandash{}จกฺขุํ ฯเปฯ วตฺถุํ อสฺสาเทติ อภินนฺทติ, ตํ อารพฺภ อุปาทานา จ สมฺปยุตฺตกา จ ขนฺธา อุปฺปชฺชนฺติ.\csromanspacer{} \textbf{วตฺถุปุเรชาตํ}\csromandash{}วตฺถุ อุปาทานานํ สมฺปยุตฺตกานญฺจ ขนฺธานํ ปุเรชาตปจฺจเยน ปจฺจโย. (3)}

\csromanheader{2}{\textbf{ปจฺฉาชาตาเสวน}}
\proseitem{29}{อุปาทาโน ธมฺโม โน-อุปาทานสฺส ธมฺมสฺส ปจฺฉาชาตปจฺจเยน ปจฺจโย. (สํขิตฺตํ.) (1)}

\prose{โน-อุปาทาโน ธมฺโม โน-อุปาทานสฺส ธมฺมสฺส ปจฺฉาชาตปจฺจเยน ปจฺจโย. (สํขิตฺตํ.) (1)}

\prose{อุปาทาโน จ โน-อุปาทาโน จ ธมฺมา โน-อุปาทานสฺส ธมฺมสฺส ปจฺฉาชาตปจฺจเยน ปจฺจโย. (สํขิตฺตํ.) (1) อาเสวนปจฺจเยน ปจฺจโย.}

\csromanheader{2}{\textbf{กมฺม}}
\proseitem{30}{โน-อุปาทาโน ธมฺโม โน-อุปาทานสฺส ธมฺมสฺส กมฺมปจฺจเยน ปจฺจโย.\csromanspacer{} \textbf{สหชาตา}, นานากฺขณิกา.\csromanspacer{}\csromanspacer{}\textbf{สหชาตา} โน-อุปาทานา เจตนา สมฺปยุตฺตกานํ ขนฺธานํ จิตฺตสมุฏฺฐานานญฺจ รูปานํ กมฺมปจฺจเยน ปจฺจโย.\csromanspacer{} ปฏิสนฺธิกฺขเณ ฯเปฯ.\csromanspacer{} \textbf{นานากฺขณิกา} โน-อุปาทานา เจตนา วิปากานํ ขนฺธานํ กฏตฺตา จ รูปานํ กมฺมปจฺจเยน ปจฺจโย. (1)}

\prose{โน-อุปาทาโน ธมฺโม อุปาทานสฺส ธมฺมสฺส กมฺมปจฺจเยน ปจฺจโย.\csromanspacer{} โน-อุปาทานา เจตนา สมฺปยุตฺตกานํ อุปาทานานํ กมฺมปจฺจเยน ปจฺจโย. (2)}

\prose{โน-อุปาทาโน ธมฺโม อุปาทานสฺส จ โน-อุปาทานสฺส จ ธมฺมสฺส กมฺมปจฺจเยน ปจฺจโย.\csromanspacer{} โน-อุปาทานา เจตนา สมฺปยุตฺตกานํ ขนฺธานํ อุปาทานานญฺจ จิตฺตสมุฏฺฐานานํ รูปานํ กมฺมปจฺจเยน ปจฺจโย. (3)}

\csromanheader{2}{\textbf{วิปากาทิ}}
\proseitem{31}{\csromanfolio{572}โน-อุปาทาโน ธมฺโม โน-อุปาทานสฺส ธมฺมสฺส วิปากปจฺจเยน ปจฺจโย.\csromanspacer{} วิปาโก โน-อุปาทาโน เอโก ขนฺโธ ฯเปฯ.\csromanspacer{} เอกํ.}

\prose{โน-อุปาทาโน ธมฺโม โน-อุปาทานสฺส ธมฺมสฺส อาหารปจฺจเยน ปจฺจโย.\csromanspacer{} ตีณิ. (เอโกเยว กพฬีกาโร อาหาโร.) อินฺทฺริยปจฺจเยน ปจฺจโย.\csromanspacer{} ตีณิ. (รูปชีวิตินฺทฺริย เอกํเยว.) ฌานปจฺจเยน ปจฺจโย.\csromanspacer{} ตีณิ.}

\prose{อุปาทาโน ธมฺโม อุปาทานสฺส ธมฺมสฺส มคฺคปจฺจเยน ปจฺจโย.\csromanspacer{} อุปาทานานิ มคฺคงฺคานิ สมฺปยุตฺตกานํ อุปาทานานํ มคฺคปจฺจเยน ปจฺจโย. (อิมินา การเณน นว ปญฺหา กาตพฺพา.) สมฺปยุตฺตปจฺจเยน ปจฺจโย.\csromanspacer{} นว.}

\csromanheader{2}{\textbf{วิปฺปยุตฺต}}
\proseitem{32}{อุปาทาโน ธมฺโม โน-อุปาทานสฺส ธมฺมสฺส วิปฺปยุตฺตปจฺจเยน ปจฺจโย.\csromanspacer{} สหชาตํ, ปจฺฉาชาตํ.\csromanspacer{}\csromanspacer{}\textbf{สหชาตา} อุปาทานา จิตฺตสมุฏฺฐานานํ รูปานํ วิปฺปยุตฺตปจฺจเยน ปจฺจโย.\csromanspacer{} \textbf{ปจฺฉาชาตํ}\csromandash{}ปจฺฉาชาตา อุปาทานา ปุเรชาตสฺส อิมสฺส กายสฺส วิปฺปยุตฺตปจฺจเยน ปจฺจโย. (1)}

\prose{โน-อุปาทาโน ธมฺโม โน-อุปาทานสฺส ธมฺมสฺส วิปฺปยุตฺตปจฺจเยน ปจฺจโย.\csromanspacer{} สหชาตํ, ปุเรชาตํ, ปจฺฉาชาตํ. (สํขิตฺตํ.) (1)}

\prose{โน-อุปาทาโน ธมฺโม อุปาทานสฺส ธมฺมสฺส วิปฺปยุตฺตปจฺจเยน ปจฺจโย.\csromanspacer{} \textbf{ปุเรชาตํ} วตฺถุ อุปาทานานํ วิปฺปยุตฺตปจฺจเยน ปจฺจโย. (2)}

\prose{โน-อุปาทาโน ธมฺโม อุปาทานสฺส จ โน-อุปาทานสฺส จ ธมฺมสฺส วิปฺปยุตฺตปจฺจเยน ปจฺจโย.\csromanspacer{} \textbf{ปุเรชาตํ} วตฺถุ อุปาทานานํ สมฺปยุตฺตกานญฺจ ขนฺธานํ วิปฺปยุตฺตปจฺจเยน ปจฺจโย. (3)}

\proseitem{33}{อุปาทาโน จ โน-อุปาทาโน จ ธมฺมา โน-อุปาทานสฺส ธมฺมสฺส วิปฺปยุตฺตปจฺจเยน ปจฺจโย.\csromanspacer{} สหชาตํ, ปจฺฉาชาตํ.\csromanspacer{}\csromanspacer{}\textbf{สหชาตา} อุปาทานา จ สมฺปยุตฺตกา จ ขนฺธา จิตฺตสมุฏฺฐานานํ รูปานํ วิปฺปยุตฺตปจฺจเยน ปจฺจโย.\csromanspacer{} \textbf{ปจฺฉาชาตา} อุปาทานา จ สมฺปยุตฺตกา จ ขนฺธา ปุเรชาตสฺส อิมสฺส กายสฺส วิปฺปยุตฺตปจฺจเยน ปจฺจโย. (1)}

\csromanheader{2}{69. อุปาทานทุก \textbf{อตฺถิ}}
\proseitem{34}{\csromanfolio{573}อุปาทาโน ธมฺโม อุปาทานสฺส ธมฺมสฺส อตฺถิปจฺจเยน ปจฺจโย.\csromanspacer{} ทิฏฺฐุปาทานํ กามุปาทานสฺส อตฺถิปจฺจเยน ปจฺจโย. (จกฺกํ.) (1)}

\prose{อุปาทาโน ธมฺโม โน-อุปาทานสฺส ธมฺมสฺส อตฺถิปจฺจเยน ปจฺจโย.\csromanspacer{} \textbf{สหชาตํ}, ปจฺฉาชาตํ.\csromanspacer{}\csromanspacer{}\textbf{สหชาตา} อุปาทานา สมฺปยุตฺตกานํ ขนฺธานํ จิตฺตสมุฏฺฐานานญฺจ รูปานํ อตฺถิปจฺจเยน ปจฺจโย.\csromanspacer{} \textbf{ปจฺฉาชาตา} อุปาทานา ปุเรชาตสฺส อิมสฺส กายสฺส อตฺถิปจฺจเยน ปจฺจโย. (2)}

\prose{อุปาทาโน ธมฺโม อุปาทานสฺส จ โน-อุปาทานสฺส จ ธมฺมสฺส อตฺถิปจฺจเยน ปจฺจโย. (สํขิตฺตํ.\csromanspacer{} ปฏิจฺจสทิสํ.) (3)}

\proseitem{35}{โน-อุปาทาโน ธมฺโม โน-อุปาทานสฺส ธมฺมสฺส อตฺถิปจฺจเยน ปจฺจโย.\csromanspacer{} \textbf{สหชาตํ}, ปุเรชาตํ, ปจฺฉาชาตํ, อาหารํ, อินฺทฺริยํ. (สํขิตฺตํ.\csromanspacer{} วิตฺถาเรตพฺพํ.) (1)}

\prose{โน-อุปาทาโน ธมฺโม อุปาทานสฺส ธมฺมสฺส อตฺถิปจฺจเยน ปจฺจโย.\csromanspacer{} \textbf{สหชาตํ}, ปุเรชาตํ.\csromanspacer{}\csromanspacer{}\textbf{สหชาตํ}\csromandash{}(สหชาตสทิสํ.) ปุเรชาตํ\csromandash{} (ปุเรชาตสทิสํ.) (2)}

\prose{โน-อุปาทาโน ธมฺโม อุปาทานสฺส จ โน-อุปาทานสฺส จ ธมฺมสฺส อตฺถิปจฺจเยน ปจฺจโย.\csromanspacer{} \textbf{สหชาตํ}, ปุเรชาตํ. (สํขิตฺตํ.\csromanspacer{} สหชาตสทิสํ.\csromanspacer{} \textbf{สหชาตํ} วิภชิตพฺพํ.\csromanspacer{} ปุเรชาตสทิสํ ปุเรชาตํ.) (3)}

\proseitem{36}{อุปาทาโน จ โน-อุปาทาโน จ ธมฺมา อุปาทานสฺส ธมฺมสฺส อตฺถิปจฺจเยน ปจฺจโย.\csromanspacer{} \textbf{สหชาตํ}, ปุเรชาตํ.\csromanspacer{}\csromanspacer{}\textbf{สหชาตํ} ทิฏฺฐุปาทานญฺจ สมฺปยุตฺตกา จ ขนฺธา กามุปาทานสฺส อตฺถิปจฺจเยน ปจฺจโย. (จกฺกํ.) \textbf{สหชาตํ} ทิฏฺฐุปาทานญฺจ วตฺถุ จ กามุปาทานสฺส อตฺถิปจฺจเยน ปจฺจโย. (จกฺกํ.) (1)}

\prose{อุปาทาโน จ โน-อุปาทาโน จ ธมฺมา โน-อุปาทานสฺส ธมฺมสฺส อตฺถิปจฺจเยน ปจฺจโย.\csromanspacer{} \textbf{สหชาตํ}, ปุเรชาตํ, ปจฺฉาชาตํ, อาหารํ, อินฺทฺริยํ.\csromanspacer{}\csromanspacer{}\textbf{สหชาโต} โน-อุปาทาโน เอโก ขนฺโธ จ อุปาทาโน จ ติณฺณนฺนํ ขนฺธานํ จิตฺตสมุฏฺฐานานญฺจ รูปานํ อตฺถิปจฺจเยน ปจฺจโย ฯเปฯ เทฺว ขนฺธา จ ฯเปฯ.\csromanspacer{} \textbf{สหชาตา} อุปาทานา จ มหาภูตา จ จิตฺตสมุฏฺฐานานํ \csromanfolio{574}รูปานํ อตฺถิปจฺจเยน ปจฺจโย.\csromanspacer{} \textbf{สหชาตา} อุปาทานา จ วตฺถุ จ โน-อุปาทานานํ ขนฺธานํ อตฺถิปจฺจเยน ปจฺจโย.\csromanspacer{} \textbf{ปจฺฉาชาตา} อุปาทานา จ สมฺปยุตฺตกา จ ขนฺธา ปุเรชาตสฺส อิมสฺส กายสฺส อตฺถิปจฺจเยน ปจฺจโย.\csromanspacer{} \textbf{ปจฺฉาชาตา} อุปาทานา จ สมฺปยุตฺตกา จ ขนฺธา กพฬีกาโร อาหาโร อิมสฺส กายสฺส อตฺถิปจฺจเยน ปจฺจโย.\csromanspacer{} \textbf{ปจฺฉาชาตา} อุปาทานา จ สมฺปยุตฺตกา จ ขนฺธา รูปชีวิตินฺทฺริยญฺจ กฏตฺตารูปานํ อตฺถิปจฺจเยน ปจฺจโย. (2)}

\prose{อุปาทาโน จ โน-อุปาทาโน จ ธมฺมา อุปาทานสฺส จ โน-อุปาทานสฺส จ ธมฺมสฺส อตฺถิปจฺจเยน ปจฺจโย.\csromanspacer{} \textbf{สหชาตํ}, ปุเรชาตํ.\csromanspacer{}\csromanspacer{}\textbf{สหชาโต} โน-อุปาทาโน เอโก ขนฺโธ จ ทิฏฺฐุปาทานญฺจ ติณฺณนฺนํ ขนฺธานํ กามุปาทานสฺส จ จิตฺตสมุฏฺฐานานํ รูปานํ อตฺถิปจฺจเยน ปจฺจโย ฯเปฯ เทฺว ขนฺธา ฯเปฯ. (จกฺกํ.) \textbf{สหชาตํ} ทิฏฺฐุปาทานญฺจ วตฺถุ จ กามุปาทานสฺส สมฺปยุตฺตกานญฺจ ขนฺธานํ อตฺถิปจฺจเยน ปจฺจโย. (จกฺกํ.) (3)}

\csromanheader{2}{1. ปจฺจยานุโลม 2. สงฺขฺยาวาร}
\csromanheader{2}{\textbf{สุทฺธ}}
\proseitem{37}{เหตุยา นว, อารมฺมเณ นว, อธิปติยา นว, อนนฺตเร นว, สมนนฺตเร นว, สหชาเต นว, อญฺญมญฺเญ นว, นิสฺสเย นว, อุปนิสฺสเย นว, ปุเรชาเต ตีณิ, ปจฺฉาชาเต ตีณิ, อาเสวเน นว, กมฺเม ตีณิ, วิปาเก เอกํ, อาหาเร ตีณิ, อินฺทฺริเย ตีณิ, ฌาเน ตีณิ, มคฺเค นว, สมฺปยุตฺเต นว, วิปฺปยุตฺเต ปญฺจ, อตฺถิยา นว, นตฺถิยา นว, วิคเต นว, อวิคเต นว.}

\vspace{\tipitakaparskip}
\csromancenter{อนุโลมํ.}
\csromanheader{2}{2. ปจฺจนียุทฺธาร}
\proseitem{38}{อุปาทาโน ธมฺโม อุปาทานสฺส ธมฺมสฺส อารมฺมณปจฺจเยน ปจฺจโย.\csromanspacer{} สหชาตปจฺจเยน ปจฺจโย.\csromanspacer{} อุปนิสฺสยปจฺจเยน ปจฺจโย. (1)}

\csromanheader{2}{69. อุปาทานทุก}
\prose{\csromanfolio{575}อุปาทาโน ธมฺโม โน-อุปาทานสฺส ธมฺมสฺส อารมฺมณปจฺจเยน ปจฺจโย.\csromanspacer{} สหชาตปจฺจเยน ปจฺจโย.\csromanspacer{} อุปนิสฺสยปจฺจเยน ปจฺจโย.\csromanspacer{} ปจฺฉาชาตปจฺจเยน ปจฺจโย. (2)}

\prose{อุปาทาโน ธมฺโม อุปาทานสฺส จ โน-อุปาทานสฺส จ ธมฺมสฺส อารมฺมณปจฺจเยน ปจฺจโย.\csromanspacer{} สหชาตปจฺจเยน ปจฺจโย.\csromanspacer{} อุปนิสฺสยปจฺจเยน ปจฺจโย. (3)}

\proseitem{39}{โน-อุปาทาโน ธมฺโม โน-อุปาทานสฺส ธมฺมสฺส อารมฺมณปจฺจเยน ปจฺจโย.\csromanspacer{} สหชาตปจฺจเยน ปจฺจโย.\csromanspacer{} อุปนิสฺสยปจฺจเยน ปจฺจโย.\csromanspacer{} ปุเรชาตปจฺจเยน ปจฺจโย.\csromanspacer{} ปจฺฉาชาตปจฺจเยน ปจฺจโย.\csromanspacer{} กมฺมปจฺจเยน ปจฺจโย.\csromanspacer{} อาหารปจฺจเยน ปจฺจโย.\csromanspacer{} อินฺทฺริยปจฺจเยน ปจฺจโย. (1)}

\prose{โน-อุปาทาโน ธมฺโม อุปาทานสฺส ธมฺมสฺส อารมฺมณปจฺจเยน ปจฺจโย.\csromanspacer{} สหชาตปจฺจเยน ปจฺจโย.\csromanspacer{} อุปนิสฺสยปจฺจเยน ปจฺจโย.\csromanspacer{} ปุเรชาตปจฺจเยน ปจฺจโย. (2)}

\prose{โน-อุปาทาโน ธมฺโม อุปาทานสฺส จ โน-อุปาทานสฺส จ ธมฺมสฺส อารมฺมณปจฺจเยน ปจฺจโย.\csromanspacer{} สหชาตปจฺจเยน ปจฺจโย.\csromanspacer{} อุปนิสฺสยปจฺจเยน ปจฺจโย.\csromanspacer{} ปุเรชาตปจฺจเยน ปจฺจโย. (3)}

\proseitem{40}{อุปาทาโน จ โน-อุปาทาโน จ ธมฺมา อุปาทานสฺส ธมฺมสฺส อารมฺมณปจฺจเยน ปจฺจโย.\csromanspacer{} สหชาตปจฺจเยน ปจฺจโย.\csromanspacer{} อุปนิสฺสยปจฺจเยน ปจฺจโย. (1)}

\prose{อุปาทาโน จ โน-อุปาทาโน จ ธมฺมา โน-อุปาทานสฺส ธมฺมสฺส อารมฺมณปจฺจเยน ปจฺจโย.\csromanspacer{} สหชาตปจฺจเยน ปจฺจโย.\csromanspacer{} อุปนิสฺสยปจฺจเยน ปจฺจโย.\csromanspacer{} ปจฺฉาชาตปจฺจเยน ปจฺจโย. (2)}

\prose{อุปาทาโน จ โน-อุปาทาโน จ ธมฺมา อุปาทานสฺส จ โน-อุปาทานสฺส จ ธมฺมสฺส อารมฺมณปจฺจเยน ปจฺจโย.\csromanspacer{} สหชาตปจฺจเยน ปจฺจโย.\csromanspacer{} อุปนิสฺสยปจฺจเยน ปจฺจโย. (3)}

\csromanheader{2}{2. ปจฺจยปจฺจนีย 2. สงฺขฺยาวาร}
\vspace{\tipitakaparskip}
\csromancenter{นเหตุยา นว, น-อารมฺมเณ นว, (สพฺพตฺถ นว.) โน-อวิคเต นว.}
\csromanheader{2}{3. ปจฺจยานุโลมปจฺจนีย}
\proseitem{42}{\csromanfolio{576}\textbf{เหตุ}ปจฺจยา น-อารมฺมเณ นว, น-อธิปติยา นว ฯเปฯ น-อญฺญมญฺเญ ตีณิ, น-อุปนิสฺสเย นว, (สพฺพตฺถ นว.) นสมฺปยุตฺเต ตีณิ, นวิปฺปยุตฺเต นว, โนนตฺถิยา นว, โนวิคเต นว.}

\csromanheader{2}{4. ปจฺจยปจฺจนียานุโลม}
\proseitem{43}{นเหตุปจฺจยา อารมฺมเณ นว, อธิปติยา นว, (อนุโลมมาติกา กาตพฺพา.) ฯเปฯ อวิคเต นว.}

\nitthitam{อุปาทานทุกํ นิฏฺฐิตํ.}
\csromanmatikaanchor{mtk.70}
\csromantocmarkref{subsection}{70. อุปาทานิยทุก}{mtk.70}
\bookmark[dest={mtk.70},level=2]{70. อุปาทานิยทุก}
\csromanheader{2}{70. อุปาทานิยทุก 1. ปฏิจฺจวาร}
\proseitem{44}{อุปาทานิยํ ธมฺมํ ปฏิจฺจ อุปาทานิโย ธมฺโม อุปฺปชฺชติ เหตุปจฺจยา.\csromanspacer{} อุปาทานิยํ เอกํ ขนฺธํ ปฏิจฺจ ตโย ขนฺธา จิตฺตสมุฏฺฐานญฺจ รูปํ ฯเปฯ เทฺว ขนฺเธ ฯเปฯ.\csromanspacer{} ปฏิสนฺธิกฺขเณ ฯเปฯ.\csromanspacer{} ขนฺเธ ปฏิจฺจ วตฺถุ, วตฺถุํ ปฏิจฺจ ขนฺธา.\csromanspacer{} เอกํ มหาภูตํ ฯเปฯ. (ยถา โลกิยทุกํ, เอวํ กาตพฺพํ.\csromanspacer{} นินฺนานากรณํ.)}

\nitthitam{อุปาทานิยทุกํ นิฏฺฐิตํ.}
\csromanmatikaanchor{mtk.71}
\csromantocmarkref{subsection}{71. อุปาทานสมฺปยุตฺตทุก}{mtk.71}
\bookmark[dest={mtk.71},level=2]{71. อุปาทานสมฺปยุตฺตทุก}
\csromanheader{2}{71. อุปาทานสมฺปยุตฺตทุก 1. ปฏิจฺจวาร}
\csromanheader{2}{1. ปจฺจยานุโลม 1. วิภงฺควาร}
\csromanheader{2}{\textbf{เหตุ}}
\proseitem{45}{อุปาทานสมฺปยุตฺตํ ธมฺมํ ปฏิจฺจ อุปาทานสมฺปยุตฺโต ธมฺโม อุปฺปชฺชติ เหตุปจฺจยา.\csromanspacer{} อุปาทานสมฺปยุตฺตํ เอกํ ขนฺธํ ปฏิจฺจ ตโย ขนฺธา ฯเปฯ เทฺว ขนฺเธ ฯเปฯ. (1)}

\prose{อุปาทานสมฺปยุตฺตํ ธมฺมํ ปฏิจฺจ อุปาทานวิปฺปยุตฺโต ธมฺโม อุปฺปชฺชติ เหตุปจฺจยา.\csromanspacer{} อุปาทานสมฺปยุตฺเต ขนฺเธ ปฏิจฺจ จิตฺตสมุฏฺฐานํ รูปํ.\csromanspacer{} ทิฏฺฐิคตวิปฺปยุตฺตโลภสหคเต ขนฺเธ ปฏิจฺจ โลโภ จิตฺตสมุฏฺฐานญฺจ รูปํ. (2)}

\csromanheader{2}{71. อุปาทานสมฺปยุตฺตทุก}
\prose{\csromanfolio{577}อุปาทานสมฺปยุตฺตํ ธมฺมํ ปฏิจฺจ อุปาทานสมฺปยุตฺโต จ อุปาทานวิปฺปยุตฺโต จ ธมฺมา อุปฺปชฺชนฺติ เหตุปจฺจยา.\csromanspacer{} อุปาทานสมฺปยุตฺตํ เอกํ ขนฺธํ ปฏิจฺจ ตโย ขนฺธา จิตฺตสมุฏฺฐานญฺจ รูปํ ฯเปฯ เทฺว ขนฺเธ ฯเปฯ.\csromanspacer{} ทิฏฺฐิคตวิปฺปยุตฺตโลภสหคตํ เอกํ ขนฺธํ ปฏิจฺจ ตโย ขนฺธา โลโภ จ จิตฺตสมุฏฺฐานญฺจ รูปํ ฯเปฯ เทฺว ขนฺเธ ฯเปฯ. (3)}

\proseitem{46}{อุปาทานวิปฺปยุตฺตํ ธมฺมํ ปฏิจฺจ อุปาทานวิปฺปยุตฺโต ธมฺโม อุปฺปชฺชติ เหตุปจฺจยา.\csromanspacer{} อุปาทานวิปฺปยุตฺตํ เอกํ ขนฺธํ ปฏิจฺจ ตโย ขนฺธา จิตฺตสมุฏฺฐานญฺจ รูปํ ฯเปฯ เทฺว ขนฺเธ ฯเปฯ.\csromanspacer{} ทิฏฺฐิคตวิปฺปยุตฺตํ โลภํ ปฏิจฺจ จิตฺตสมุฏฺฐานํ รูปํ.\csromanspacer{} ปฏิสนฺธิกฺขเณ อุปาทานวิปฺปยุตฺตํ เอกํ ขนฺธํ ปฏิจฺจ ตโย ขนฺธา กฏตฺตา จ รูปํ ฯเปฯ เทฺว ขนฺเธ ฯเปฯ.\csromanspacer{} ขนฺเธ ปฏิจฺจ วตฺถุ, วตฺถุํ ปฏิจฺจ ขนฺธา.\csromanspacer{} เอกํ มหาภูตํ ฯเปฯ. (1)}

\prose{อุปาทานวิปฺปยุตฺตํ ธมฺมํ ปฏิจฺจ อุปาทานสมฺปยุตฺโต ธมฺโม อุปฺปชฺชติ เหตุปจฺจยา.\csromanspacer{} ทิฏฺฐิคตวิปฺปยุตฺตํ โลภํ ปฏิจฺจ สมฺปยุตฺตกา ขนฺธา. (2)}

\prose{อุปาทานวิปฺปยุตฺตํ ธมฺมํ ปฏิจฺจ อุปาทานสมฺปยุตฺโต จ อุปาทานวิปฺปยุตฺโต จ ธมฺมา อุปฺปชฺชนฺติ เหตุปจฺจยา.\csromanspacer{} ทิฏฺฐิคตวิปฺปยุตฺตํ โลภํ ปฏิจฺจ สมฺปยุตฺตกา ขนฺธา จิตฺตสมุฏฺฐานญฺจ รูปํ. (3)}

\proseitem{47}{อุปาทานสมฺปยุตฺตญฺจ อุปาทานวิปฺปยุตฺตญฺจ ธมฺมํ ปฏิจฺจ อุปาทานสมฺปยุตฺโต ธมฺโม อุปฺปชฺชติ เหตุปจฺจยา.\csromanspacer{} ทิฏฺฐิคตวิปฺปยุตฺตโลภสหคตํ เอกํ ขนฺธญฺจ โลภญฺจ ปฏิจฺจ ตโย ขนฺธา ฯเปฯ เทฺว ขนฺเธ จ ฯเปฯ. (1)}

\prose{อุปาทานสมฺปยุตฺตญฺจ อุปาทานวิปฺปยุตฺตญฺจ ธมฺมํ ปฏิจฺจ อุปาทานวิปฺปยุตฺโต ธมฺโม อุปฺปชฺชติ เหตุปจฺจยา.\csromanspacer{} อุปาทานสมฺปยุตฺเต ขนฺเธ จ มหาภูเต จ ปฏิจฺจ จิตฺตสมุฏฺฐานํ รูปํ.\csromanspacer{} ทิฏฺฐิคตวิปฺปยุตฺตโลภสหคเต ขนฺเธ จ โลภญฺจ ปฏิจฺจ จิตฺตสมุฏฺฐานํ รูปํ. (2)}

\prose{อุปาทานสมฺปยุตฺตญฺจ อุปาทานวิปฺปยุตฺตญฺจ ธมฺมํ ปฏิจฺจ อุปาทานสมฺปยุตฺโต จ อุปาทานวิปฺปยุตฺโต จ ธมฺมา อุปฺปชฺชนฺติ เหตุปจฺจยา.\csromanspacer{} ทิฏฺฐิคตวิปฺปยุตฺตโลภสหคตํ เอกํ ขนฺธญฺจ โลภญฺจ ปฏิจฺจ ตโย ขนฺธา จิตฺตสมุฏฺฐานญฺจ รูปํ ฯเปฯ เทฺว ขนฺเธ จ ฯเปฯ. (3)}

\csromanheader{2}{\textbf{อารมฺมณ}}
\proseitem{48}{\csromanfolio{578}อุปาทานสมฺปยุตฺตํ ธมฺมํ ปฏิจฺจ อุปาทานสมฺปยุตฺโต ธมฺโม อุปฺปชฺชติ อารมฺมณปจฺจยา.\csromanspacer{} อุปาทานสมฺปยุตฺตํ เอกํ ขนฺธํ ปฏิจฺจ ตโย ขนฺธา ฯเปฯ เทฺว ขนฺเธ ฯเปฯ. (1)}

\prose{อุปาทานสมฺปยุตฺตํ ธมฺมํ ปฏิจฺจ อุปาทานวิปฺปยุตฺโต ธมฺโม อุปฺปชฺชติ อารมฺมณปจฺจยา.\csromanspacer{} ทิฏฺฐิคตวิปฺปยุตฺตโลภสหคเต ขนฺเธ ปฏิจฺจ โลโภ. (2)}

\prose{อุปาทานสมฺปยุตฺตํ ธมฺมํ ปฏิจฺจ อุปาทานสมฺปยุตฺโต จ อุปาทานวิปฺปยุตฺโต จ ธมฺมา อุปฺปชฺชนฺติ อารมฺมณปจฺจยา.\csromanspacer{} ทิฏฺฐิคตวิปฺปยุตฺตโลภสหคตํ เอกํ ขนฺธํ ปฏิจฺจ ตโย ขนฺธา โลโภ จ ฯเปฯ เทฺว ขนฺเธ ฯเปฯ. (3)}

\proseitem{49}{อุปาทานวิปฺปยุตฺตํ ธมฺมํ ปฏิจฺจ อุปาทานวิปฺปยุตฺโต ธมฺโม อุปฺปชฺชติ อารมฺมณปจฺจยา.\csromanspacer{} อุปาทานวิปฺปยุตฺตํ เอกํ ขนฺธํ ฯเปฯ เทฺว ขนฺเธ ฯเปฯ.\csromanspacer{} ปฏิสนฺธิกฺขเณ ฯเปฯ.\csromanspacer{} วตฺถุํ ปฏิจฺจ ขนฺธา. (1)}

\prose{อุปาทานวิปฺปยุตฺตํ ธมฺมํ ปฏิจฺจ อุปาทานสมฺปยุตฺโต ธมฺโม อุปฺปชฺชติ อารมฺมณปจฺจยา.\csromanspacer{} ทิฏฺฐิคตวิปฺปยุตฺตํ โลภํ ปฏิจฺจ สมฺปยุตฺตกา ขนฺธา. (2)}

\prose{อุปาทานสมฺปยุตฺตญฺจ อุปาทานวิปฺปยุตฺตญฺจ ธมฺมํ ปฏิจฺจ อุปาทานสมฺปยุตฺโต ธมฺโม อุปฺปชฺชติ อารมฺมณปจฺจยา.\csromanspacer{} ทิฏฺฐิคตวิปฺปยุตฺตโลภสหคตํ เอกํ ขนฺธญฺจ โลภญฺจ ปฏิจฺจ ตโย ขนฺธา ฯเปฯ เทฺว ขนฺเธ จ ฯเปฯ. (3) (สํขิตฺตํ.)}

\csromanheader{2}{1. ปจฺจยานุโลม 2. สงฺขฺยาวาร}
\csromanheader{2}{\textbf{สุทฺธ}}
\proseitem{50}{เหตุยา นว, อารมฺมเณ ฉ, อธิปติยา นว, อนนฺตเร ฉ, สมนนฺตเร ฉ, สหชาเต นว, อญฺญมญฺเญ ฉ, นิสฺสเย นว, อุปนิสฺสเย ฉ, ปุเรชาเต ฉ, อาเสวเน ฉ, กมฺเม นว, วิปาเก เอกํ, อาหาเร นว, (สพฺพตฺถ นว,) มคฺเค นว, สมฺปยุตฺเต ฉ, วิปฺปยุตฺเต นว, อตฺถิยา นว, นตฺถิยา ฉ, วิคเต ฉ, อวิคเต นว.}

\csromanheader{2}{71. อุปาทานสมฺปยุตฺตทุก \textbf{2. ปจฺจยปจฺจนีย 1. วิภงฺควาร}}
\csromanheader{2}{\textbf{นเหตุ}}
\proseitem{51}{\csromanfolio{579}อุปาทานวิปฺปยุตฺตํ ธมฺมํ ปฏิจฺจ อุปาทานวิปฺปยุตฺโต ธมฺโม อุปฺปชฺชติ นเหตุปจฺจยา.\csromanspacer{} อเหตุกํ อุปาทานวิปฺปยุตฺตํ เอกํ ขนฺธํ ปฏิจฺจ ตโย ขนฺธา จิตฺตสมุฏฺฐานญฺจ รูปํ ฯเปฯ เทฺว ขนฺเธ ฯเปฯ.\csromanspacer{} อเหตุกปฏิสนฺธิกฺขเณ ฯเปฯ. (ยาว อสญฺญสตฺตา.) วิจิกิจฺฉาสหคเต อุทฺธจฺจสหคเต ขนฺเธ ปฏิจฺจ วิจิกิจฺฉาสหคโต อุทฺธจฺจสหคโต โมโห. (1)}

\csromanheader{2}{\textbf{น-อารมฺมณาทิ}}
\proseitem{52}{อุปาทานสมฺปยุตฺตํ ธมฺมํ ปฏิจฺจ อุปาทานวิปฺปยุตฺโต ธมฺโม อุปฺปชฺชติ น-อารมฺมณปจฺจยา.\csromanspacer{} อุปาทานสมฺปยุตฺเต ขนฺเธ ปฏิจฺจ จิตฺตสมุฏฺฐานํ รูปํ. (1)}

\prose{อุปาทานวิปฺปยุตฺตํ ธมฺมํ ปฏิจฺจ อุปาทานวิปฺปยุตฺโต ธมฺโม อุปฺปชฺชติ น-อารมฺมณปจฺจยา.\csromanspacer{} อุปาทานวิปฺปยุตฺเต ขนฺเธ ปฏิจฺจ จิตฺตสมุฏฺฐานํ รูปํ.\csromanspacer{} ทิฏฺฐิคตวิปฺปยุตฺตํ โลภํ ปฏิจฺจ จิตฺตสมุฏฺฐานํ รูปํ.\csromanspacer{} ปฏิสนฺธิกฺขเณ ฯเปฯ. (ยาว อสญฺญสตฺตา.) (1)}

\prose{อุปาทานสมฺปยุตฺตญฺจ อุปาทานวิปฺปยุตฺตญฺจ ธมฺมํ ปฏิจฺจ อุปาทานวิปฺปยุตฺโต ธมฺโม อุปฺปชฺชติ น-อารมฺมณปจฺจยา.\csromanspacer{} อุปาทานสมฺปยุตฺเต ขนฺเธ จ มหาภูเต จ ปฏิจฺจ จิตฺตสมุฏฺฐานํ รูปํ.\csromanspacer{} ทิฏฺฐิคตวิปฺปยุตฺตโลภสหคเต ขนฺเธ จ โลภญฺจ ปฏิจฺจ จิตฺตสมุฏฺฐานํ รูปํ. (1)}

\prose{น-อธิปติปจฺจยา.\csromanspacer{} น-อนนฺตรปจฺจยา.\csromanspacer{} นสมนนฺตรปจฺจยา.\csromanspacer{} น-อุปนิสฺสยปจฺจยา.}

\csromanheader{2}{\textbf{นปุเรชาตาทิ}}
\proseitem{53}{อุปาทานสมฺปยุตฺตํ ธมฺมํ ปฏิจฺจ อุปาทานสมฺปยุตฺโต ธมฺโม อุปฺปชฺชติ นปุเรชาตปจฺจยา.\csromanspacer{} อรูเป อุปาทานสมฺปยุตฺตํ เอกํ ขนฺธํ ปฏิจฺจ ตโย ขนฺธา ฯเปฯ เทฺว ขนฺเธ ฯเปฯ. (1)}

\prose{อุปาทานสมฺปยุตฺตํ ธมฺมํ ปฏิจฺจ อุปาทานวิปฺปยุตฺโต ธมฺโม อุปฺปชฺชติ นปุเรชาตปจฺจยา.\csromanspacer{} อรูเป ทิฏฺฐิคตวิปฺปยุตฺตโลภสหคเต ขนฺเธ ปฏิจฺจ โลโภ.\csromanspacer{} อุปาทานสมฺปยุตฺเต ขนฺเธ ปฏิจฺจ จิตฺตสมุฏฺฐานํ รูปํ. (2)}

\prose{\csromanfolio{580}อุปาทานสมฺปยุตฺตํ ธมฺมํ ปฏิจฺจ อุปาทานสมฺปยุตฺโต จ อุปาทานวิปฺปยุตฺโต จ ธมฺมา อุปฺปชฺชนฺติ นปุเรชาตปจฺจยา.\csromanspacer{} อรูเป ทิฏฺฐิคตวิปฺปยุตฺตโลภสหคตํ เอกํ ขนฺทํ ปฏิจฺจ ตโย ขนฺธา โลโภ จ ฯเปฯ เทฺว ขนฺเธ ฯเปฯ. (3)}

\proseitem{54}{อุปาทานวิปฺปยุตฺตํ ธมฺมํ ปฏิจฺจ อุปาทานวิปฺปยุตฺโต ธมฺโม อุปฺปชฺชติ นปุเรชาตปจฺจยา.\csromanspacer{} อรูเป อุปาทานวิปฺปยุตฺตํ เอกํ ขนฺธํ ปฏิจฺจ ตโย ขนฺธา ฯเปฯ เทฺว ขนฺเธ ฯเปฯ.\csromanspacer{} อุปาทานวิปฺปยุตฺเต ขนฺเธ ปฏิจฺจ จิตฺตสมุฏฺฐานํ รูปํ.\csromanspacer{} ทิฏฺฐิคตวิปฺปยุตฺตํ โลภํ ปฏิจฺจ จิตฺตสมุฏฺฐานํ รูปํ.\csromanspacer{} ปฏิสนฺธิกฺขเณ ฯเปฯ. (ยาว อสญฺญสตฺตา.) (1)}

\prose{อุปาทานวิปฺปยุตฺตํ ธมฺมํ ปฏิจฺจ อุปาทานสมฺปยุตฺโต ธมฺโม อุปฺปชฺชติ นปุเรชาตปจฺจยา.\csromanspacer{} อรูเป ทิฏฺฐิคตวิปฺปยุตฺตํ โลภํ ปฏิจฺจ สมฺปยุตฺตกา ขนฺธา. (2)}

\prose{อุปาทานสมฺปยุตฺตญฺจ อุปาทานวิปฺปยุตฺตญฺจ ธมฺมํ ปฏิจฺจ อุปาทานสมฺปยุตฺโต ธมฺโม อุปฺปชฺชติ นปุเรชาตปจฺจยา.\csromanspacer{} อรูเป ทิฏฺฐิคตวิปฺปยุตฺตโลภสหคตํ เอกํ ขนฺธญฺจ โลภญฺจ ปฏิจฺจ ตโย ขนฺธา ฯเปฯ เทฺว ขนฺเธ จ ฯเปฯ. (1)}

\prose{อุปาทานสมฺปยุตฺตญฺจ อุปาทานวิปฺปยุตฺตญฺจ ธมฺมํ ปฏิจฺจ อุปาทานวิปฺปยุตฺโต ธมฺโม อุปฺปชฺชติ นปุเรชาตปจฺจยา.\csromanspacer{} อุปาทานสมฺปยุตฺเต ขนฺเธ จ มหาภูเต จ ปฏิจฺจ จิตฺตสมุฏฺฐานํ รูปํ.\csromanspacer{} ทิฏฺฐิคตวิปฺปยุตฺตโลภสหคเต ขนฺเธ จ โลภญฺจ ปฏิจฺจ จิตฺตสมุฏฺฐานํ รูปํ. (2)}

\prose{นปจฺฉาชาตปจฺจยา.\csromanspacer{} น-อาเสวนปจฺจยา.}

\csromanheader{2}{\textbf{นกมฺม}}
\proseitem{55}{อุปาทานสมฺปยุตฺตํ ธมฺมํ ปฏิจฺจ อุปาทานสมฺปยุตฺโต ธมฺโม อุปฺปชฺชติ นกมฺมปจฺจยา.\csromanspacer{} อุปาทานสมฺปยุตฺเต ขนฺเธ ปฏิจฺจ สมฺปยุตฺตกา เจตนา. (1)}

\prose{อุปาทานวิปฺปยุตฺตํ ธมฺมํ ปฏิจฺจ อุปาทานวิปฺปยุตฺโต ธมฺโม อุปฺปชฺชติ นกมฺมปจฺจยา.\csromanspacer{} อุปาทานวิปฺปยุตฺเต ขนฺเธ ปฏิจฺจ สมฺปยุตฺตกา เจตนา.\csromanspacer{} พาหิรํ.\csromanspacer{} อาหารสมุฏฺฐานํ.\csromanspacer{} อุตุสมุฏฺฐานํ ฯเปฯ. (1)}

\prose{อุปาทานวิปฺปยุตฺตํ ธมฺมํ ปฏิจฺจ อุปาทานสมฺปยุตฺโต ธมฺโม อุปฺปชฺชติ นกมฺมปจฺจยา.\csromanspacer{} ทิฏฺฐิคตวิปฺปยุตฺตํ โลภํ ปฏิจฺจ สมฺปยุตฺตกา เจตนา. (2)}

\csromanheader{2}{71. อุปาทานสมฺปยุตฺตทุก}
\prose{\csromanfolio{581}อุปาทานสมฺปยุตฺตญฺจ อุปาทานวิปฺปยุตฺตญฺจ ธมฺมํ ปฏิจฺจ อุปาทานสมฺปยุตฺโต ธมฺโม อุปฺปชฺชติ นกมฺมปจฺจยา.\csromanspacer{} ทิฏฺฐิคตวิปฺปยุตฺตโลภสหคเต ขนฺเธ จ โลภญฺจ ปฏิจฺจ สมฺปยุตฺตกา เจตนา. (1) (สํขิตฺตํ.)}

\csromanheader{2}{2. ปจฺจยปจฺจนีย 2. สงฺขฺยาวาร}
\csromanheader{2}{\textbf{สุทฺธ}}
\proseitem{56}{นเหตุยา เอกํ, น-อารมฺมเณ ตีณิ, น-อธิปติยา นว, น-อนนฺตเร ตีณิ, นสมนนฺตเร ตีณิ, น-อญฺญมญฺเญ ตีณิ, น-อุปนิสฺสเย ตีณิ, นปุเรชาเต สตฺต, นปจฺฉาชาเต นว, น-อาเสวเน นว, นกมฺเม จตฺตาริ, นวิปาเก นว, น-อาหาเร เอกํ, น-อินฺทฺริเย เอกํ, นฌาเน เอกํ, นมคฺเค เอกํ, นสมฺปยุตฺเต ตีณิ, นวิปฺปยุตฺเต ฉ, โนนตฺถิยา ตีณิ.\csromanspacer{} โนวิคเต ตีณิ.}

\prose{(อิตเร เทฺว คณนาปิ กาตพฺพา.\csromanspacer{} สหชาตวาโร ปฏิจฺจวารสทิโส.)}

\csromanheader{2}{71. อุปาทานสมฺปยุตฺตทุก 3. ปจฺจยวาร}
\csromanheader{2}{1. ปจฺจยานุโลม 1. วิภงฺควาร}
\csromanheader{2}{\textbf{เหตุ}}
\proseitem{57}{อุปาทานสมฺปยุตฺตํ ธมฺมํ ปจฺจยา อุปาทานสมฺปยุตฺโต ธมฺโม อุปฺปชฺชติ เหตุปจฺจยา.\csromanspacer{} ตีณิ. (ปฏิจฺจสทิสํ.)}

\prose{อุปาทานวิปฺปยุตฺตํ ธมฺมํ ปจฺจยา อุปาทานวิปฺปยุตฺโต ธมฺโม อุปฺปชฺชติ เหตุปจฺจยา.\csromanspacer{} อุปาทานวิปฺปยุตฺตํ เอกํ ขนฺธํ ปจฺจยา ฯเปฯ. (ยาว อชฺฌตฺติกา มหาภูตา.) วตฺถุํ ปจฺจยา อุปาทานวิปฺปยุตฺตา ขนฺธา.\csromanspacer{} วตฺถุํ ปจฺจยา ทิฏฺฐิคตวิปฺปยุตฺโต โลโภ. (1)}

\prose{อุปาทานวิปฺปยุตฺตํ ธมฺมํ ปจฺจยา อุปาทานสมฺปยุตฺโต ธมฺโม อุปฺปชฺชติ เหตุปจฺจยา.\csromanspacer{} วตฺถุํ ปจฺจยา อุปาทานสมฺปยุตฺตา ขนฺธา.\csromanspacer{} ทิฏฺฐิคตวิปฺปยุตฺตํ โลภํ ปจฺจยา สมฺปยุตฺตกา ขนฺธา. (2)}

\prose{\csromanfolio{582}อุปาทานวิปฺปยุตฺตํ ธมฺมํ ปจฺจยา อุปาทานสมฺปยุตฺโต จ อุปาทานวิปฺปยุตฺโต จ ธมฺมา อุปฺปชฺชนฺติ เหตุปจฺจยา.\csromanspacer{} วตฺถุํ ปจฺจยา อุปาทานสมฺปยุตฺตกา ขนฺธา.\csromanspacer{} มหาภูเต ปจฺจยา จิตฺตสมุฏฺฐานํ รูปํ.\csromanspacer{} ทิฏฺฐิคตวิปฺปยุตฺตํ โลภํ ปจฺจยา สมฺปยุตฺตกา ขนฺธา จิตฺตสมุฏฺฐานญฺจ รูปํ.\csromanspacer{} วตฺถุํ ปจฺจยา ทิฏฺฐิคตวิปฺปยุตฺตโลภสหคตา ขนฺธา จ โลโภ จ. (3)}

\proseitem{58}{อุปาทานสมฺปยุตฺตญฺจ อุปาทานวิปฺปยุตฺตญฺจ ธมฺมํ ปจฺจยา อุปาทานสมฺปยุตฺโต ธมฺโม อุปฺปชฺชติ เหตุปจฺจยา.\csromanspacer{} อุปาทานสมฺปยุตฺตํ เอกํ ขนฺธญฺจ วตฺถุญฺจ ปจฺจยา ตโย ขนฺธา ฯเปฯ เทฺว ขนฺเธ จ ฯเปฯ.\csromanspacer{} ทิฏฺฐิคตวิปฺปยุตฺตโลภสหคตํ เอกํ ขนฺธญฺจ โลภญฺจ ปจฺจยา ตโย ขนฺธา ฯเปฯ เทฺว ขนฺเธ จ ฯเปฯ. (1)}

\prose{อุปาทานสมฺปยุตฺตญฺจ อุปาทานวิปฺปยุตฺตญฺจ ธมฺมํ ปจฺจยา อุปาทานวิปฺปยุตฺโต ธมฺโม อุปฺปชฺชติ เหตุปจฺจยา.\csromanspacer{} อุปาทานสมฺปยุตฺเต ขนฺเธ จ มหาภูเต จ ปจฺจยา จิตฺตสมุฏฺฐานํ รูปํ.\csromanspacer{} ทิฏฺฐิคตวิปฺปยุตฺตโลภสหคเต ขนฺเธ จ โลภญฺจ ปจฺจยา จิตฺตสมุฏฺฐานํ รูปํ.\csromanspacer{} ทิฏฺฐิคตวิปฺปยุตฺตโลภสหคเต ขนฺเธ จ วตฺถุญฺจ ปจฺจยา โลโภ. (2)}

\prose{อุปาทานสมฺปยุตฺตญฺจ อุปาทานวิปฺปยุตฺตญฺจ ธมฺมํ ปจฺจยา อุปาทานสมฺปยุตฺโต จ อุปาทานวิปฺปยุตฺโต จ ธมฺมา อุปฺปชฺชนฺติ เหตุปจฺจยา.\csromanspacer{} อุปาทานสมฺปยุตฺตํ เอกํ ขนฺธญฺจ วตฺถุญฺจ ปจฺจยา ตโย ขนฺธา ฯเปฯ เทฺว ขนฺเธ จ ฯเปฯ.\csromanspacer{} อุปาทานสมฺปยุตฺเต ขนฺเธ จ มหาภูเต จ ปจฺจยา จิตฺตสมุฏฺฐานํ รูปํ.\csromanspacer{} ทิฏฺฐิคตวิปฺปยุตฺตโลภสหคตํ เอกํ ขนฺธญฺจ วตฺถุญฺจ ปจฺจยา ตโย ขนฺธา โลโภ จ ฯเปฯ เทฺว ขนฺเธ จ ฯเปฯ. (3) (สํขิตฺตํ.\csromanspacer{} อารมฺมณปจฺจยมฺหิ ปญฺจ วิญฺญาณา กาตพฺพา.)}

\csromanheader{2}{1. ปจฺจยานุโลม 2. สงฺขฺยาวาร}
\proseitem{59}{เหตุยา นว, อารมฺมเณ นว, อธิปติยา นว, อนนฺตเร นว, (สพฺพตฺถ นว,) วิปาเก เอกํ ฯเปฯ อวิคเต นว.}

\csromanheader{2}{2. ปจฺจยปจฺจนีย 1. วิภงฺควาร}
\csromanheader{2}{\textbf{นเหตุ}}
\proseitem{60}{อุปาทานวิปฺปยุตฺตํ ธมฺมํ ปจฺจยา อุปาทานวิปฺปยุตฺโต ธมฺโม อุปฺปชฺชติ นเหตุปจฺจยา.\csromanspacer{} อเหตุกํ อุปาทานวิปฺปยุตฺตํ เอกํ ขนฺธํ ปจฺจยา ฯเปฯ. (ยาว}

\vspace{\tipitakaparskip}
\csromancenter{อุปาทานสมฺปยุตฺตทุก อสญฺญสตฺตา.) จกฺขายตนํ ปจฺจยา จกฺขุวิญฺญาณํ ฯเปฯ.\csromanspacer{} กายายตนํ ปจฺจยา กายวิญฺญาณํ ฯเปฯ.\csromanspacer{} วตฺถุํ ปจฺจยา อเหตุกา อุปาทานวิปฺปยุตฺตา ขนฺธา.\csromanspacer{} วิจิกิจฺฉาสหคเต อุทฺธจฺจสหคเต ขนฺเธ จ วตฺถุญฺจ ปจฺจยา วิจิกิจฺฉาสหคโต อุทฺธจฺจสหคโต โมโห. (สํขิตฺตํ.)}
\csromanheader{2}{2. ปจฺจยปจฺจนีย 2. สงฺขฺยาวาร}
\proseitem{61}{\csromanfolio{583}นเหตุยา เอกํ, น-อารมฺมเณ ตีณิ, น-อธิปติยา นว, น-อนนฺตเร ตีณิ, นสมนนฺตเร ตีณิ, น-อญฺญมญฺเญ ตีณิ, น-อุปนิสฺสเย ตีณิ, นปุเรชาเต สตฺต, นปจฺฉาชาเต นว, น-อาเสวเน นว, นกมฺเม จตฺตาริ, นวิปาเก นว, น-อาหาเร เอกํ, น-อินฺทฺริเย เอกํ, นฌาเน เอกํ, นมคฺเค เอกํ, นสมฺปยุตฺเต ตีณิ, นวิปฺปยุตฺเต ฉ, โนนตฺถิยา ตีณิ, โนวิคเต ตีณิ.}

\vspace{\tipitakaparskip}
\csromancenter{(เอวํ อิตเร เทฺว คณนาปิ นิสฺสยวาโรปิ กาตพฺโพ.)}
\csromanheader{2}{71. อุปาทานสมฺปยุตฺตทุก 5. สํสฏฺฐวาร}
\csromanheader{2}{1. ปจฺจยานุโลมาทิ}
\proseitem{62}{อุปาทานสมฺปยุตฺตํ ธมฺมํ สํสฏฺโฐ อุปาทานสมฺปยุตฺโต ธมฺโม อุปฺปชฺชติ เหตุปจฺจยา.\csromanspacer{} อุปาทานสมฺปยุตฺตํ เอกํ ขนฺธํ สํสฏฺฐา.\csromanspacer{} ตีณิ.}

\prose{อุปาทานวิปฺปยุตฺตํ ธมฺมํ สํสฏฺโฐ อุปาทานวิปฺปยุตฺโต ธมฺโม อุปฺปชฺชติ เหตุปจฺจยา. (ปฏิจฺจสทิสํ.\csromanspacer{} อรูปํเยว กาตพฺพํ.)}

\prose{อุปาทานวิปฺปยุตฺตํ ธมฺมํ สํสฏฺโฐ อุปาทานสมฺปยุตฺโต ธมฺโม อุปฺปชฺชติ เหตุปจฺจยา. (ปฏิจฺจสทิสํ.\csromanspacer{} อรูปํเยว กาตพฺพํ.)}

\prose{อุปาทานสมฺปยุตฺตญฺจ อุปาทานวิปฺปยุตฺตญฺจ ธมฺมํ สํสฏฺโฐ อุปาทานสมฺปยุตฺโต ธมฺโม อุปฺปชฺชติ เหตุปจฺจยา. (ปฏิจฺจสทิสํ.\csromanspacer{} อรูปํเยว กาตพฺพํ.)}

\prose{เหตุยา ฉ, อารมฺมเณ ฉ, อธิปติยา ฉ, (สพฺพตฺถ ฉ,) วิปาเก เอกํ ฯเปฯ อวิคเต ฉ.\csromanspacer{} อนุโลมํ.}

\prose{อุปาทานวิปฺปยุตฺตํ ธมฺมํ สํสฏฺโฐ อุปาทานวิปฺปยุตฺโต ธมฺโม อุปฺปชฺชติ นเหตุปจฺจยา. (สํขิตฺตํ.)}

\prose{\csromanfolio{584}นเหตุยา เอกํ, น-อธิปติยา ฉ, นปุเรชาเต ฉ, นปจฺฉาชาเต ฉ, น-อาเสวเน ฉ, นกมฺเม จตฺตาริ, นวิปาเก ฉ, นฌาเน เอกํ, นมคฺเค เอกํ, นวิปฺปยุตฺเต ฉ.}

\vspace{\tipitakaparskip}
\csromancenter{ปจฺจนียํ.}
\csromancenter{(เอวํ อิตเร เทฺว คณนาปิ สมฺปยุตฺตวาโรปิ กาตพฺโพ.)}
\csromanheader{2}{71. อุปาทานสมฺปยุตฺตทุก 7. ปญฺหาวาร}
\csromanheader{2}{1. ปจฺจยานุโลม 1. วิภงฺควาร}
\csromanheader{2}{\textbf{เหตุ}}
\proseitem{63}{อุปาทานสมฺปยุตฺโต ธมฺโม อุปาทานสมฺปยุตฺตสฺส ธมฺมสฺส เหตุปจฺจเยน ปจฺจโย.\csromanspacer{} อุปาทานสมฺปยุตฺตา เหตู สมฺปยุตฺตกานํ ขนฺธานํ เหตุปจฺจเยน ปจฺจโย. (1)}

\prose{อุปาทานสมฺปยุตฺโต ธมฺโม อุปาทานวิปฺปยุตฺตสฺส ธมฺมสฺส เหตุปจฺจเยน ปจฺจโย.\csromanspacer{} อุปาทานสมฺปยุตฺตา เหตู จิตฺตสมุฏฺฐานํ รูปานํ เหตุปจฺจเยน ปจฺจโย.\csromanspacer{} ทิฏฺฐิคตวิปฺปยุตฺตโลภสหคตา เหตู ทิฏฺฐิคตวิปฺปยุตฺตโลภสฺส จิตฺตสมุฏฺฐานานญฺจ รูปานํ เหตุปจฺจเยน ปจฺจโย. (มูลํ กาตพฺพํ.) อุปาทานสมฺปยุตฺตา เหตู สมฺปยุตฺตกานํ ขนฺธานํ จิตฺตสมุฏฺฐานานญฺจ รูปานํ เหตุปจฺจเยน ปจฺจโย.\csromanspacer{} ทิฏฺฐิคตวิปฺปยุตฺตโลภสหคตา เหตู สมฺปยุตฺตกานํ ขนฺธานํ โลภสฺส จ จิตฺตสมุฏฺฐานานญฺจ รูปานํ เหตุปจฺจเยน ปจฺจโย. (3)}

\proseitem{64}{อุปาทานวิปฺปยุตฺโต ธมฺโม อุปาทานวิปฺปยุตฺตสฺส ธมฺมสฺส เหตุปจฺจเยน ปจฺจโย.\csromanspacer{} อุปาทานวิปฺปยุตฺตา เหตู สมฺปยุตฺตกานํ ขนฺธานํ จิตฺตสมุฏฺฐานานญฺจ รูปานํ เหตุปจฺจเยน ปจฺจโย.\csromanspacer{} ทิฏฺฐิคตวิปฺปยุตฺโต โลโภ จิตฺตสมุฏฺฐานานํ รูปานํ เหตุปจฺจเยน ปจฺจโย.\csromanspacer{} ปฏิสนฺธิกฺขเณ ฯเปฯ. (มูลํ กาตพฺพํ.) ทิฏฺฐิคตวิปฺปยุตฺโต โลโภ สมฺปยุตฺตกานํ ขนฺธานํ เหตุปจฺจเยน ปจฺจโย. (มูลํ กาตพฺพํ.) ทิฏฺฐิคตวิปฺปยุตฺโต}

\vspace{\tipitakaparskip}
\csromancenter{อุปาทานสมฺปยุตฺตทุก โลโภ สมฺปยุตฺตกานํ ขนฺธานํ จิตฺตสมุฏฺฐานานญฺจ รูปานํ เหตุปจฺจเยน ปจฺจโย. (3)}
\proseitem{65}{\csromanfolio{585}อุปาทานสมฺปยุตฺโต จ อุปาทานวิปฺปยุตฺโตจ ธมฺมา อุปาทานสมฺปยุตฺตสฺส ธมฺมสฺส เหตุปจฺจเยน ปจฺจโย.\csromanspacer{} ทิฏฺฐิคตวิปฺปยุตฺตโลภสหคโต โมโห จ โลโภ จ สมฺปยุตฺตกานํ ขนฺธานํ เหตุปจฺจเยน ปจฺจโย. (มูลํ กาตพฺพํ.) ทิฏฺฐิคตวิปฺปยุตฺตโลภสหคโต โมโห จ โลโภ จ จิตฺตสมุฏฺฐานานํ รูปานํ เหตุปจฺจเยน ปจฺจโย. (2)}

\prose{อุปาทานสมฺปยุตฺโต จ อุปาทานวิปฺปยุตฺโต จ ธมฺมา อุปาทานสมฺปยุตฺตสฺส จ อุปาทานวิปฺปยุตฺตสฺส จ ธมฺมสฺส เหตุปจฺจเยน ปจฺจโย.\csromanspacer{} ทิฏฺฐิคตวิปฺปยุตฺตโลภสหคโต โมโห จ โลโภ จ สมฺปยุตฺตกานํ ขนฺธานํ จิตฺตสมุฏฺฐานานญฺจ รูปานํ เหตุปจฺจเยน ปจฺจโย. (3)}

\csromanheader{2}{\textbf{อารมฺมณ}}
\proseitem{66}{อุปาทานสมฺปยุตฺโต ธมฺโม อุปาทานสมฺปยุตฺตสฺส ธมฺมสฺส อารมฺมณปจฺจเยน ปจฺจโย.\csromanspacer{} อุปาทานสมฺปยุตฺเต ขนฺเธ อารพฺภ อุปาทานสมฺปยุตฺตา ขนฺธา อุปฺปชฺชนฺติ. (มูลํ กาตพฺพํ.) อุปาทานสมฺปยุตฺเต ขนฺเธ อารพฺภ อุปาทานวิปฺปยุตฺตา ขนฺธา จ โลโภ จ อุปฺปชฺชนฺติ. (มูลํ กาตพฺพํ.) อุปาทานสมฺปยุตฺเต ขนฺเธ อารพฺภ ทิฏฺฐิคตวิปฺปยุตฺตโลภสหคตา ขนฺธา จ โลโภ จ อุปฺปชฺชนฺติ. (3)}

\proseitem{67}{อุปาทานวิปฺปยุตฺโต ธมฺโม อุปาทานวิปฺปยุตฺตสฺส ธมฺมสฺส อารมฺมณปจฺจเยน ปจฺจโย.\csromanspacer{} ทานํ ฯเปฯ สีลํ ฯเปฯ อุโปสถกมฺมํ ฯเปฯ.\csromanspacer{} ปุพฺเพ สุจิณฺณานิ ฯเปฯ.\csromanspacer{} ฌานา วุฏฺฐหิตฺวา ฌานํ ปจฺจเวกฺขติ, อสฺสาเทติ อภินนฺทติ, ตํ อารพฺภ ทิฏฺฐิคตวิปฺปยุตฺโต ราโค ฯเปฯ.\csromanspacer{} วิจิกิจฺฉา ฯเปฯ อุทฺธจฺจํ อุปฺปชฺชติ.\csromanspacer{} ฌาเน ปริหีเน วิปฺปฏิสาริสฺส โทมนสฺสํ อุปฺปชฺชติ.\csromanspacer{} อริยา มคฺคา วุฏฺฐหิตฺวา มคฺคํ ปจฺจเวกฺขนฺติ.\csromanspacer{} ผลํ ฯเปฯ.\csromanspacer{} นิพฺพานํ ปจฺจเวกฺขนฺติ, นิพฺพานํ โคตฺรภุสฺส, โวทานสฺส, มคฺคสฺส, ผลสฺส, อาวชฺชนาย อารมฺมณปจฺจเยน ปจฺจโย.\csromanspacer{} อริยา อุปาทานวิปฺปยุตฺเต ปหีเน กิเลเส ฯเปฯ.\csromanspacer{} วิกฺขมฺภิเต กิเลเส ฯเปฯ.\csromanspacer{} ปุพฺเพ สมุทาจิณฺเณ กิเลเส ฯเปฯ.\csromanspacer{} จกฺขุํ ฯเปฯ วตฺถุํ อุปาทานวิปฺปยุตฺเต \csromanfolio{586}ขนฺเธ จ โลภญฺจ อนิจฺจโต ฯเปฯ วิปสฺสติ, อสฺสาเทติ อภินนฺทติ, ตํ อารพฺภ ทิฏฺฐิคตวิปฺปยุตฺโต ราโค ฯเปฯ วิจิกิจฺฉา ฯเปฯ อุทฺธจฺจํ ฯเปฯ โทมนสฺสํ อุปฺปชฺชติ ฯเปฯ. (สพฺพํ ปริปุณฺณํ.) ทิพฺเพน จกฺขุนา. (ยาว กายวิญฺญาณํ.) อุปาทานวิปฺปยุตฺตา ขนฺธา อิทฺธิวิธญาณสฺส, เจโตปริยญาณสฺส, ปุพฺเพนิวาสานุสฺสติญาณสฺส, ยถากมฺมูปคญาณสฺส, อนาคตํสญาณสฺส, อาวชฺชนาย อารมฺมณปจฺจเยน ปจฺจโย. (1)}

\prose{อุปาทานวิปฺปยุตฺโต ธมฺโม อุปาทานสมฺปยุตฺตสฺส ธมฺมสฺส อารมฺมณปจฺจเยน ปจฺจโย.\csromanspacer{} ทานํ ฯเปฯ.\csromanspacer{} ฌานํ ฯเปฯ.\csromanspacer{} จกฺขุํ ฯเปฯ.\csromanspacer{} วตฺถุํ อุปาทานวิปฺปยุตฺเต ขนฺเธ จ โลภญฺจ อสฺสาเทติ อภินนฺทติ, ตํ อารพฺภ ราโค อุปฺปชฺชติ.\csromanspacer{} ทิฏฺฐิ อุปฺปชฺชติ. (2)}

\prose{อุปาทานวิปฺปยุตฺโต ธมฺโม อุปาทานสมฺปยุตฺตสฺส จ อุปาทานวิปฺปยุตฺตสฺส จ ธมฺมสฺส อารมฺมณปจฺจเยน ปจฺจโย.\csromanspacer{} จกฺขุํ ฯเปฯ.\csromanspacer{} วตฺถุํ อุปาทานวิปฺปยุตฺเต ขนฺเธ จ โลภญฺจ อารพฺภ ทิฏฺฐิคตวิปฺปยุตฺตโลภสหคตา ขนฺธา จ โลโภ จ อุปฺปชฺชนฺติ. (3)}

\proseitem{68}{อุปาทานสมฺปยุตฺโต จ อุปาทานวิปฺปยุตฺโต จ ธมฺมา อุปาทานสมฺปยุตฺตสฺส ธมฺมสฺส อารมฺมณปจฺจเยน ปจฺจโย.\csromanspacer{} ทิฏฺฐิคตวิปฺปยุตฺตโลภ สหคเต ขนฺเธ จ โลภญฺจ อารพฺภ อุปาทานสมฺปยุตฺตา ขนฺธา อุปฺปชฺชนฺติ. (มูลํ กาตพฺพํ.) ทิฏฺฐิคตวิปฺปยุตฺตโลภสหคเต ขนฺเธ จ โลภญฺจ อารพฺภ อุปาทานวิปฺปยุตฺโต ขนฺธา จ โลโภ จ อุปฺปชฺชนฺติ. (มูลํ กาตพฺพํ.) ทิฏฺฐิคตวิปฺปยุตฺตโลภสหคเต ขนฺเธ จ โลภญฺจ อารพฺภ ทิฏฺฐิคตวิปฺปยุตฺตโลภสหคเต ขนฺธา จ โลโภ จ อุปฺปชฺชนฺติ. (3)}

\csromanheader{2}{\textbf{อธิปติ}}
\proseitem{69}{อุปาทานสมฺปยุตฺโต ธมฺโม อุปาทานสมฺปยุตฺตสฺส ธมฺมสฺส อธิปติปจฺจเยน ปจฺจโย.\csromanspacer{} อารมฺมณาธิปติ, สหชาตาธิปติ.\csromanspacer{}\csromanspacer{}\textbf{อารมฺมณาธิปติ\csromandash{}}อุปาทานสมฺปยุตฺเต ขนฺเธ ครุํ กตฺวา อุปาทานสมฺปยุตฺตา ขนฺธา อุปฺปชฺชนฺติ.\csromanspacer{} \textbf{สหชาตาธิปติ\csromandash{}} อุปาทานสมฺปยุตฺตาธิปติ สมฺปยุตฺตกานํ ขนฺธานํ}

\vspace{\tipitakaparskip}
\csromancenter{อุปาทานสมฺปยุตฺตทุก อธิปติปจฺจเยน ปจฺจโย. (มูลํ กาตพฺพํ.) \textbf{อารมฺมณาธิปติ}, สหชาตาธิปติ.\csromanspacer{}\csromanspacer{}\textbf{อารมฺมณาธิปติ\csromandash{}}อุปาทานสมฺปยุตฺเต ขนฺเธ ครุํ กตฺวา ทิฏฺฐิคตวิปฺปยุตฺโต โลโภ อุปฺปชฺชติ.\csromanspacer{} \textbf{สหชาตาธิปติ}\csromandash{} อุปาทานสมฺปยุตฺตาธิปติ จิตฺตสมุฏฺฐานานํ รูปานํ อธิปติปจฺจเยน ปจฺจโย.\csromanspacer{} ทิฏฺฐิคตวิปฺปยุตฺตโลภสหคตาธิปติ โลภสฺส จ จิตฺตสมุฏฺฐานานญฺจ รูปานํ อธิปติปจฺจเยน ปจฺจโย. (มูลํ กาตพฺพํ.) \textbf{อารมฺมณาธิปติ}, สหชาตาธิปติ.\csromanspacer{}\csromanspacer{}\textbf{อารมฺมณาธิปติ\csromandash{}} อุปาทานสมฺปยุตฺเต ขนฺเธ ครุํ กตฺวา ทิฏฺฐิคตวิปฺปยุตฺตโลภสหคตา ขนฺธา จ โลโภ จ อุปฺปชฺชนฺติ.\csromanspacer{} \textbf{สหชาตาธิปติ}\csromandash{}อุปาทานสมฺปยุตฺตาธิปติ สมฺปยุตฺตกานํ ขนฺธานํ จิตฺตสมุฏฺฐานานญฺจ รูปานํ อธิปติปจฺจเยน ปจฺจโย.\csromanspacer{} ทิฏฺฐิคตวิปฺปยุตฺตโลภสหคตาธิปติ สมฺปยุตฺตกานํ ขนฺธานํ โลภสฺส จ จิตฺตสมุฏฺฐานานญฺจ รูปานํ อธิปติปจฺจเยน ปจฺจโย. (3)}
\proseitem{70}{\csromanfolio{587}อุปาทานวิปฺปยุตฺโต ธมฺโม อุปาทานวิปฺปยุตฺตสฺส ธมฺมสฺส อธิปติปจฺจเยน ปจฺจโย.\csromanspacer{} \textbf{อารมฺมณาธิปติ}, สหชาตาธิปติ.\csromanspacer{}\csromanspacer{}\textbf{อารมฺมณาธิปติ\csromandash{}}ทานํ ฯเปฯ.\csromanspacer{} สีลํ ฯเปฯ.\csromanspacer{} ฌานา วุฏฺฐหิตฺวา ฌานํ ครุํ กตฺวา ปจฺจเวกฺขติ, อสฺสาเทติ อภินนฺทติ, ตํ ครุํ กตฺวา ทิฏฺฐิคตวิปฺปยุตฺโต ราโค อุปฺปชฺชติ ฯเปฯ.\csromanspacer{} อริยา มคฺคา วุฏฺฐหิตฺวา ฯเปฯ ผลสฺส อธิปติปจฺจเยน ปจฺจโย.\csromanspacer{} จกฺขุํ ฯเปฯ.\csromanspacer{} วตฺถุํ อุปาทานวิปฺปยุตฺเต ขนฺเธ จ โลภญฺจ ครุํ กตฺวา อสฺสาเทติ อภินนฺทติ, ตํ ครุํ กตฺวา ทิฏฺฐิคตวิปฺปยุตฺโต ราโค อุปฺปชฺชติ.\csromanspacer{} \textbf{สหชาตาธิปติ}\csromandash{} อุปาทานวิปฺปยุตฺตาธิปติ สมฺปยุตฺตกานํ ขนฺธานํ จิตฺตสมุฏฺฐานานญฺจ รูปานํ อธิปติปจฺจเยน ปจฺจโย. (1)}

\prose{อุปาทานวิปฺปยุตฺโต ธมฺโม อุปาทานสมฺปยุตฺตสฺส ธมฺมสฺส อธิปติปจฺจเยน ปจฺจโย.\csromanspacer{}\csromanspacer{}\textbf{อารมฺมณาธิปติ\csromandash{}}ทานํ ฯเปฯ.\csromanspacer{} ฌานํ ฯเปฯ.\csromanspacer{} จกฺขุํ ฯเปฯ.\csromanspacer{} วตฺถุํ อุปาทานวิปฺปยุตฺเต ขนฺเธ จ โลภญฺจ ครุํ กตฺวา อสฺสาเทติ อภินนฺทติ, ตํ ครุํ กตฺวา ราโค อุปฺปชฺชติ.\csromanspacer{} ทิฏฺฐิ อุปฺปชฺชติ.}

\prose{อุปาทานวิปฺปยุตฺโต ธมฺโม อุปาทานสมฺปยุตฺตสฺส จ อุปาทานวิปฺปยุตฺตสฺส จ ธมฺมสฺส อธิปติปจฺจเยน ปจฺจโย.\csromanspacer{}\csromanspacer{}\textbf{อารมฺมณาธิปติ\csromandash{}}จกฺขุํ ฯเปฯ วตฺถุํ อุปาทานวิปฺปยุตฺเต ขนฺเธ จ โลภญฺจ ครุํ กตฺวา ทิฏฺฐิคตวิปฺปยุตฺตโลภสหคตา ขนฺธา จ โลโภ จ อุปฺปชฺชนฺติ. (3)}

\proseitem{71}{\csromanfolio{588}อุปาทานสมฺปยุตฺโต จ อุปาทานวิปฺปยุตฺโต จ ธมฺมา อุปาทานสมฺปยุตฺตสฺส ธมฺมสฺส อธิปติปจฺจเยน ปจฺจโย.\csromanspacer{}\csromanspacer{}\textbf{อารมฺมณาธิปติ\csromandash{}}ทิฏฺฐิคตวิปฺปยุตฺตโลภสหคเต ขนฺเธ จ โลภญฺจ ครุํ กตฺวา อุปาทานสมฺปยุตฺตา ขนฺธา อุปฺปชฺชนฺติ. (มูลํ กาตพฺพํ.) ทิฏฺฐิคตวิปฺปยุตฺตโลภสหคเต ขนฺเธ จ โลภญฺจ ครุํ กตฺวา ทิฏฺฐิคตวิปฺปยุตฺโต โลโภ อุปฺปชฺชติ. (มูลํ กาตพฺพํ.) ทิฏฺฐิคตวิปฺปยุตฺตโลภสหคเต ขนฺเธ จ โลภญฺจ ครุํ กตฺวา ทิฏฺฐิคตวิปฺปยุตฺตโลภสหคตา ขนฺธา จ โลโภ จ อุปฺปชฺชนฺติ. (3)}

\csromanheader{2}{\textbf{อนนฺตร}}
\proseitem{72}{อุปาทานสมฺปยุตฺโต ธมฺโม อุปาทานสมฺปยุตฺตสฺส ธมฺมสฺส อนนฺตรปจฺจเยน ปจฺจโย.\csromanspacer{} ปุริมา ปุริมา อุปาทานสมฺปยุตฺตา ขนฺธา ปจฺฉิมานํ ปจฺฉิมานํ อุปาทานสมฺปยุตฺตกานํ ขนฺธานํ อนนฺตรปจฺจเยน ปจฺจโย. (มูลํ กาตพฺพํ.) ปุริมา ปุริมา ทิฏฺฐิคตวิปฺปยุตฺตโลภสหคตา ขนฺธา ปจฺฉิมสฺส ปจฺฉิมสฺส ทิฏฺฐิคตวิปฺปยุตฺตสฺส โลภสฺส อนนฺตรปจฺจเยน ปจฺจโย.\csromanspacer{} อุปาทานสมฺปยุตฺตา ขนฺธา วุฏฺฐานสฺส อนนฺตรปจฺจเยน ปจฺจโย. (มูลํ กาตพฺพํ.) ปุริมา ปุริมา ทิฏฺฐิคตวิปฺปยุตฺตโลภสหคตา ขนฺธา ปจฺฉิมานํ ปจฺฉิมานํ ทิฏฺฐิคตวิปฺปยุตฺตโลภสหคตานํ ขนฺธานํ โลภสฺส จ อนนฺตรปจฺจเยน ปจฺจโย. (3)}

\proseitem{73}{อุปาทานวิปฺปยุตฺโต ธมฺโม อุปาทานวิปฺปยุตฺตสฺส ธมฺมสฺส อนนฺตรปจฺจเยน ปจฺจโย.\csromanspacer{} ปุริโม ปุริโม ทิฏฺฐิคตวิปฺปยุตฺโต โลโภ ปจฺฉิมสฺส ปจฺฉิมสฺส ทิฏฺฐิคตวิปฺปยุตฺตสฺส โลภสฺส อนนฺตรปจฺจเยน ปจฺจโย.\csromanspacer{} ปุริมา ปุริมา อุปาทานวิปฺปยุตฺตา ขนฺธา ปจฺฉิมานํ ปจฺฉิมานํ อุปาทานวิปฺปยุตฺตานํ ขนฺธานํ อนนฺตรปจฺจเยน ปจฺจโย.\csromanspacer{} ทิฏฺฐิคตวิปฺปยุตฺโต โลโภ วุฏฺฐานสฺส อนนฺตรปจฺจเยน ปจฺจโย.\csromanspacer{} อนุโลมํ โคตฺรภุสฺส ฯเปฯ ผลสมาปตฺติยา อนนฺตรปจฺจเยน ปจฺจโย. (มูลํ กาตพฺพํ.) ปุริโม ปุริโม ทิฏฺฐิคตวิปฺปยุตฺโต โลโภ ปจฺฉิมานํ ปจฺฉิมานํ ทิฏฺฐิคตวิปฺปยุตฺตโลภสหคตานํ ขนฺธานํ อนนฺตรปจฺจเยน ปจฺจโย.\csromanspacer{} อาวชฺชนา อุปาทานสมฺปยุตฺตกานํ ขนฺธานํ อนนฺตรปจฺจเยน ปจฺจโย. (มูลํ กาตพฺพํ.) ปุริโม ปุริโม ทิฏฺฐิคตวิปฺปยุตฺโต โลโภ ปจฺฉิมานํ ปจฺฉิมานํ ทิฏฺฐิคตวิปฺปยุตฺตโลภสหคตานํ ขนฺธานํ โลภสฺส จ อนนฺตรปจฺจเยน ปจฺจโย.}

\vspace{\tipitakaparskip}
\csromancenter{อุปาทานสมฺปยุตฺตทุก อาวชฺชนา ทิฏฺฐิคตวิปฺปยุตฺตโลภสหคตานํ ขนฺธานํ โลภสฺส จ อนนฺตรปจฺจเยน ปจฺจโย. (3)}
\proseitem{74}{\csromanfolio{589}อุปาทานสมฺปยุตฺโต จ อุปาทานวิปฺปยุตฺโต จ ธมฺมา อุปาทานสมฺปยุตฺตสฺส ธมฺมสฺส อนนฺตรปจฺจเยน ปจฺจโย.\csromanspacer{} ปุริมา ปุริมา ทิฏฺฐิคตวิปฺปยุตฺตโลภสหคตา ขนฺธา จ โลโภ จ ปจฺฉิมานํ ปจฺฉิมานํ ทิฏฺฐิคตวิปฺปยุตฺตโลภสหคตานํ ขนฺธานํ อนนฺตรปจฺจเยน ปจฺจโย. (มูลํ กาตพฺพํ.) ปุริมา ปุริมา ทิฏฺฐิคตวิปฺปยุตฺตโลภสหคตา ขนฺธา จ โลโภ จ ปจฺฉิมสฺส ปจฺฉิมสฺส ทิฏฺฐิคตวิปฺปยุตฺตสฺส โลภสฺส อนนฺตรปจฺจเยน ปจฺจโย.\csromanspacer{} ทิฏฺฐิคตวิปฺปยุตฺตโลภสหคตา ขนฺธา จ โลโภ จ วุฏฺฐานสฺส อนนฺตรปจฺจเยน ปจฺจโย. (มูลํ กาตพฺพํ.) ปุริมา ปุริมา ทิฏฺฐิคตวิปฺปยุตฺตโลภสหคตา ขนฺธา จ โลโภ จ ปจฺฉิมานํ ปจฺฉิมานํ ทิฏฺฐิคตวิปฺปยุตฺตโลภสหคตานํ ขนฺธานํ โลภสฺส จ อนนฺตรปจฺจเยน ปจฺจโย. (3)}

\csromanheader{2}{\textbf{สมนนฺตราทิ}}
\proseitem{75}{อุปาทานสมฺปยุตฺโต ธมฺโม อุปาทานสมฺปยุตฺตสฺส ธมฺมสฺส สมนนฺตรปจฺจเยน ปจฺจโย.\csromanspacer{} สหชาตปจฺจเยน ปจฺจโย. (ปฏิจฺจสทิสํ.) นว.\csromanspacer{} อญฺญมญฺญปจฺจเยน ปจฺจโย. (ปฏิจฺจสทิสํ.) ฉ.\csromanspacer{} นิสฺสยปจฺจเยน ปจฺจโย. (ปจฺจยวารสทิสํ.) นว.}

\csromanheader{2}{\textbf{อุปนิสฺสย}}
\proseitem{76}{อุปาทานสมฺปยุตฺโต ธมฺโม อุปาทานสมฺปยุตฺตสฺส ธมฺมสฺส อุปนิสฺสยปจฺจเยน ปจฺจโย.\csromanspacer{} อารมฺมณูปนิสฺสโย, อนนฺตรูปนิสฺสโย, ปกตูปนิสฺสโย ฯเปฯ.\csromanspacer{} ปกตูปนิสฺสโย\csromandash{}อุปทานสมฺปยุตฺตา ขนฺธา อุปาทานสมฺปยุตฺตกานํ ขนฺธานํ อุปนิสฺสยปจฺจเยน ปจฺจโย. (มูลํ กาตพฺพํ.) อุปาทานสมฺปยุตฺตา ขนฺธา อุปาทานวิปฺปยุตฺตกานํ ขนฺธานํ โลภสฺส จ อุปนิสฺสยปจฺจเยน ปจฺจโย. (มูลํ กาตพฺพํ.) อุปาทานสมฺปยุตฺตา ขนฺธา ทิฏฺฐิคตวิปฺปยุตฺตโลภสหคตานํ ขนฺธานํ โลภสฺส จ อุปนิสฺสยปจฺจเยน ปจฺจโย. (3)}

\proseitem{77}{\csromanfolio{590}อุปาทานวิปฺปยุตฺโต ธมฺโม อุปาทานวิปฺปยุตฺตสฺส ธมฺมสฺส อุปนิสฺสยปจฺจเยน ปจฺจโย.\csromanspacer{} อารมฺมณูปนิสฺสโย, อนนฺตรูปนิสฺสโย, ปกตูปนิสฺสโย ฯเปฯ.\csromanspacer{} ปกตูปนิสฺสโย\csromandash{}สทฺธํ อุปนิสฺสาย ทานํ เทติ ฯเปฯ สมาปตฺติํ อุปฺปาเทติ, มานํ ชปฺเปติ.\csromanspacer{} สีลํ ฯเปฯ.\csromanspacer{} ปญฺญํ.\csromanspacer{} อุปาทานวิปฺปยุตฺตํ ราคํ.\csromanspacer{} มานํ.\csromanspacer{} ปตฺถนํ ฯเปฯ.\csromanspacer{} เสนาสนํ อุปนิสฺสาย ทานํ เทติ ฯเปฯ สํฆํ ภินฺทติ.\csromanspacer{} สทฺธา ฯเปฯ.\csromanspacer{} เสนาสนํ สทฺธาย ฯเปฯ ปญฺญาย อุปาทานวิปฺปยุตฺตสฺส ราคสฺส.\csromanspacer{} มานสฺส.\csromanspacer{} ปตฺถนาย.\csromanspacer{} กายิกสฺส สุขสฺส ฯเปฯ ผลสมาปตฺติยา อุปนิสฺสยปจฺจเยน ปจฺจโย. (1)}

\prose{อุปาทานวิปฺปยุตฺโต ธมฺโม อุปาทานสมฺปยุตฺตสฺส ธมฺมสฺส อุปนิสฺสยปจฺจเยน ปจฺจโย. (ตีณิ อุปนิสฺสยา.) สทฺธํ อุปนิสฺสาย มานํ ชปฺเปติ.\csromanspacer{} ทิฏฺฐิํ คณฺหาติ.\csromanspacer{} สีลํ ฯเปฯ.\csromanspacer{} ปญฺญํ.\csromanspacer{} อุปาทานวิปฺปยุตฺตํ ราคํ.\csromanspacer{} มานํ.\csromanspacer{} ปตฺถนํ.\csromanspacer{} เสนาสนํ อุปนิสฺสาย อทินฺนํ ฯเปฯ มุสา ฯเปฯ ปิสุณํ ฯเปฯ สมฺผํ ฯเปฯ สนฺธิํ ฯเปฯ นิลฺโลปํ ฯเปฯ เอกาคาริกํ ฯเปฯ ปริปนฺเถ ฯเปฯ ปรทารํ ฯเปฯ คามฆาตํ ฯเปฯ นิคฆาตํ กโรติ.\csromanspacer{} สทฺธา ฯเปฯ.\csromanspacer{} เสนาสนํ อุปาทานสมฺปยุตฺตสฺส ราคสฺส.\csromanspacer{} โมหสฺส.\csromanspacer{} มานสฺส.\csromanspacer{} ทิฏฺฐิยา.\csromanspacer{} ปตฺถนาย อุปนิสฺสยปจฺจเยน ปจฺจโย. (2)}

\prose{อุปาทานวิปฺปยุตฺโต ธมฺโม อุปาทานสมฺปยุตฺตสฺส จ อุปาทานวิปฺปยุตฺตสฺส จ ธมฺมสฺส อุปนิสฺสยปจฺจเยน ปจฺจโย.\csromanspacer{} อารมฺมณูปนิสฺสโย, อนนฺตรูปนิสฺสโย, ปกตูปนิสฺสโย ฯเปฯ.\csromanspacer{} ปกตูปนิสฺสโย\csromandash{}สทฺธํ อุปนิสฺสาย มานํ ชปฺเปติ.\csromanspacer{} สีลํ ฯเปฯ.\csromanspacer{} ปญฺญํ.\csromanspacer{} อุปาทานวิปฺปยุตฺตํ ราคํ ฯเปฯ.\csromanspacer{} เสนาสนํ อุปนิสฺสาย อทินฺนํ อาทิยติ ฯเปฯ.\csromanspacer{} คามฆาตํ กโรติ.\csromanspacer{} นิคมฆาตํ กโรติ.\csromanspacer{} สทฺธา ฯเปฯ.\csromanspacer{} เสนาสนํ ทิฏฺฐิคตวิปฺปยุตฺตโลภสหคตานํ ขนฺธานํ โลภสฺส จ อุปนิสฺสยปจฺจเยน ปจฺจโย. (3)}

\proseitem{78}{อุปาทานสมฺปยุตฺโต จ อุปาทานวิปฺปยุตฺโต จ ธมฺมา อุปาทานสมฺปยุตฺตสฺส ธมฺมสฺส อุปนิสฺสยปจฺจเยน ปจฺจโย.\csromanspacer{} ทิฏฺฐิคตวิปฺปยุตฺตโลภสหคตา ขนฺธา จ โลโภ จ อุปาทานสมฺปยุตฺตกานํ ขนฺธานํ อุปนิสฺสยปจฺจเยน ปจฺจโย. (มูลํ กาตพฺพํ.) ทิฏฺฐิคตวิปฺปยุตฺตโลภสหคตา ขนฺธา จ โลโภ จ อุปาทานวิปฺปยุตฺตานํ ขนฺธานํ โลภสฺส จ อุปนิสฺสยปจฺจเยน ปจฺจโย. (มูลํ กาตพฺพํ.) ทิฏฺฐิคตวิปฺปยุตฺตโลภสหคตา ขนฺธา จ โลโภ จ ทิฏฺฐิคตวิปฺปยุตฺตโลภสหคตานํ ขนฺธานํ โลภสฺส จ อุปนิสฺสยปจฺจเยน ปจฺจโย. (3)}

\csromanheader{2}{71. อุปาทานสมฺปยุตฺตทุก \textbf{ปุเรชาต}}
\proseitem{79}{\csromanfolio{591}อุปาทานวิปฺปยุตฺโต ธมฺโม อุปาทานวิปฺปยุตฺตสฺส ธมฺมสฺส ปุเรชาตปจฺจเยน ปจฺจโย.\csromanspacer{} \textbf{อารมฺมณปุเรชาตํ}, วตฺถุปุเรชาตํ.\csromanspacer{}\csromanspacer{}\textbf{อารมฺมณปุเรชาตํ}\csromandash{}จกฺขุํ ฯเปฯ วตฺถุํ อนิจฺจโต ฯเปฯ โทมนสฺสํ อุปฺปชฺชติ.\csromanspacer{} ทิพฺเพน จกฺขุนา รูปํ ปสฺสติ.\csromanspacer{} ทิพฺพาย โสตธาตุยา สทฺทํ สุณาติ.\csromanspacer{} รูปายตนํ จกฺขุวิญฺญาณสฺส ฯเปฯ.\csromanspacer{} โผฏฺฐพฺพายตนํ ฯเปฯ.\csromanspacer{} \textbf{วตฺถุปุเรชาตํ\csromandash{}}จกฺขายตนํ จกฺขุวิญฺญาณสฺส ฯเปฯ กายายตนํ ฯเปฯ.\csromanspacer{} วตฺถุ อุปาทานวิปฺปยุตฺตานํ ขนฺธานํ โลภสฺส จ ปุเรชาตปจฺจเยน ปจฺจโย. (1)}

\prose{อุปาทานวิปฺปยุตฺโต ธมฺโม อุปาทานสมฺปยุตฺตสฺส ธมฺมสฺส ปุเรชาตปจฺจเยน ปจฺจโย.\csromanspacer{} \textbf{อารมฺมณปุเรชาตํ}, วตฺถุปุเรชาตํ.\csromanspacer{}\csromanspacer{}\textbf{อารมฺมณปุเรชาตํ}\csromandash{}จกฺขุํ ฯเปฯ วตฺถุํ อสฺสาเทติ อภินนฺทติ, ตํ อารพฺภ ราโค อุปฺปชฺชติ.\csromanspacer{} ทิฏฺฐิ อุปฺปชฺชติ.\csromanspacer{} \textbf{วตฺถุปุเรชาตํ\csromandash{}}วตฺถุ อุปาทานสมฺปยุตฺตกานํ ขนฺธานํ ปุเรชาตปจฺจเยน ปจฺจโย. (2)}

\prose{อุปาทานวิปฺปยุตฺโต ธมฺโม อุปาทานสมฺปยุตฺตสฺส จ อุปาทานวิปฺปยุตฺตสฺส จ ธมฺมสฺส ปุเรชาตปจฺจเยน ปจฺจโย.\csromanspacer{} \textbf{อารมฺมณปุเรชาตํ}, วตฺถุปุเรชาตํ.\csromanspacer{}\csromanspacer{}\textbf{อารมฺมณปุเรชาตํ}\csromandash{}จกฺขุํ ฯเปฯ วตฺถุํ อารพฺภ ทิฏฺฐิคตวิปฺปยุตฺตโลภสหคตา ขนฺธา จ โลโภ จ อุปฺปชฺชนฺติ.\csromanspacer{} \textbf{วตฺถุปุเรชาตํ\csromandash{}}วตฺถุ ทิฏฺฐิคตวิปฺปยุตฺตโลภสหคตานํ ขนฺธานํ โลภสฺส จ ปุเรชาตปจฺจเยน ปจฺจโย. (3)}

\csromanheader{2}{\textbf{ปจฺฉาชาตาเสวน}}
\proseitem{80}{อุปาทานสมฺปยุตฺโต ธมฺโม อุปาทานวิปฺปยุตฺตสฺส ธมฺมสฺส ปจฺฉาชาตปจฺจเยน ปจฺจโย. (สํขิตฺตํ.) (1)}

\prose{อุปาทานวิปฺปยุตฺโต ธมฺโม อุปาทานวิปฺปยุตฺตสฺส ธมฺมสฺส ปจฺฉาชาตปจฺจเยน ปจฺจโย. (สํขิตฺตํ.) (1)}

\prose{อุปาทานสมฺปยุตฺโต จ อุปาทานวิปฺปยุตฺโต จ ธมฺมา อุปาทานวิปฺปยุตฺตสฺส ธมฺมสฺส ปจฺฉาชาตปจฺจเยน ปจฺจโย. (สํขิตฺตํ.) (1) อาเสวนปจฺจเยน ปจฺจโย.}

\csromanheader{2}{\textbf{กมฺมวิปาก}}
\proseitem{81}{\csromanfolio{592}อุปาทานสมฺปยุตฺโต ธมฺโม อุปาทานสมฺปยุตฺตสฺส ธมฺมสฺส กมฺมปจฺจเยน ปจฺจโย.\csromanspacer{} อุปาทานสมฺปยุตฺตา เจตนา สมฺปยุตฺตกานํ ขนฺธานํ กมฺมปจฺจเยน ปจฺจโย. (1)}

\prose{อุปาทานสมฺปยุตฺโต ธมฺโม อุปาทานวิปฺปยุตฺตสฺส ธมฺมสฺส กมฺมปจฺจเยน ปจฺจโย.\csromanspacer{} \textbf{สหชาตา}, นานากฺขณิกา.\csromanspacer{}\csromanspacer{}\textbf{สหชาตา} อุปาทานสมฺปยุตฺตา เจตนา จิตฺตสมุฏฺฐานานํ รูปานํ กมฺมปจฺจเยน ปจฺจโย.\csromanspacer{} ทิฏฺฐิคตวิปฺปยุตฺตโลภสหคตา เจตนา โลภสฺส จิตฺตสมุฏฺฐานานญฺจ รูปานํ กมฺมปจฺจเยน ปจฺจโย.\csromanspacer{} \textbf{นานากฺขณิกา} อุปาทานสมฺปยุตฺตา เจตนา วิปากานํ ขนฺธานํ กฏตฺตา จ รูปานํ กมฺมปจฺจเยน ปจฺจโย. (2)}

\prose{อุปาทานสมฺปยุตฺโต ธมฺโม อุปาทานสมฺปยุตฺตสฺส จ อุปาทานวิปฺปยุตฺตสฺส จ ธมฺมสฺส กมฺมปจฺจเยน ปจฺจโย.\csromanspacer{} อุปาทานสมฺปยุตฺตา เจตนา สมฺปยุตฺตกานํ ขนฺธานํ จิตฺตสมุฏฺฐานานญฺจ รูปานํ กมฺมปจฺจเยน ปจฺจโย.\csromanspacer{} ทิฏฺฐิคตวิปฺปยุตฺตโลภสหคตา เจตนา สมฺปยุตฺตกานํ ขนฺธานํ โลภสฺส จ จิตฺตสมุฏฺฐานานญฺจ รูปานํ กมฺมปจฺจเยน ปจฺจโย. (3)}

\proseitem{82}{อุปาทานวิปฺปยุตฺโต ธมฺโม อุปาทานวิปฺปยุตฺตสฺส ธมฺมสฺส กมฺมปจฺจเยน ปจฺจโย.\csromanspacer{} \textbf{สหชาตา}, นานากฺขณิกา.\csromanspacer{}\csromanspacer{}\textbf{สหชาตา} อุปาทานวิปฺปยุตฺตา เจตนา สมฺปยุตฺตกานํ ขนฺธานํ จิตฺตสมุฏฺฐานานญฺจ รูปานํ กมฺมปจฺจเยน ปจฺจโย.\csromanspacer{} ปฏิสนฺธิกฺขเณ ฯเปฯ.\csromanspacer{} \textbf{นานากฺขณิกา} อุปาทานวิปฺปยุตฺตา เจตนา วิปากานํ ขนฺธานํ กฏตฺตา จ รูปานํ กมฺมปจฺจเยน ปจฺจโย. (1)}

\prose{อุปาทานวิปฺปยุตฺโต ธมฺโม อุปาทานวิปฺปยุตฺตสฺส ธมฺมสฺส วิปากปจฺจเยน ปจฺจโย.\csromanspacer{} วิปาโก อุปาทานวิปฺปยุตฺโต เอโก ฯเปฯ.\csromanspacer{} เอกํ.}

\csromanheader{2}{\textbf{อาหาราทิ}}
\proseitem{83}{อุปาทานสมฺปยุตฺโต ธมฺโม อุปาทานสมฺปยุตฺตสฺส ธมฺมสฺส อาหารปจฺจเยน ปจฺจโย.\csromanspacer{} อินฺทฺริยปจฺจเยน ปจฺจโย.\csromanspacer{} ฌานปจฺจเยน ปจฺจโย.\csromanspacer{} มคฺคปจฺจเยน ปจฺจโย. (อิเมสุ จตูสุปิ ยถา กมฺมปจฺจเย ทิฏฺฐิคตวิปฺปยุตฺโต โลโภ ทสฺสิโต, เอวํ ทสฺเสตพฺโพ.\csromanspacer{} จตฺตาริ จตฺตาริ ปญฺหา.) สมฺปยุตฺตปจฺจเยน ปจฺจโย.\csromanspacer{} ฉ.}

\csromanheader{2}{71. อุปาทานสมฺปยุตฺตทุก \textbf{วิปฺปยุตฺต}}
\proseitem{84}{\csromanfolio{593}อุปาทานสมฺปยุตฺโต ธมฺโม อุปาทานวิปฺปยุตฺตสฺส ธมฺมสฺส วิปฺปยุตฺตปจฺจเยน ปจฺจโย.\csromanspacer{} สหชาตํ, ปจฺฉาชาตํ. (สํขิตฺตํ.) (1)}

\prose{อุปาทานวิปฺปยุตฺโต ธมฺโม อุปาทานวิปฺปยุตฺตสฺส ธมฺมสฺส วิปฺปยุตฺตปจฺจเยน ปจฺจโย.\csromanspacer{} สหชาตํ, ปุเรชาตํ, ปจฺฉาชาตํ. (สํขิตฺตํ.) (1)}

\prose{อุปาทานวิปฺปยุตฺโต ธมฺโม อุปาทานสมฺปยุตฺตสฺส ธมฺมสฺส วิปฺปยุตฺตปจฺจเยนปจฺจโย.\csromanspacer{} \textbf{ปุเรชาตํ} วตฺถุ อุปาทานสมฺปยุตฺตกานํ ขนฺธานํ วิปฺปยุตฺตปจฺจเยน ปจฺจโย. (2)}

\prose{อุปาทานวิปฺปยุตฺโต ธมฺโม อุปาทานสมฺปยุตฺตสฺส จ อุปาทานวิปฺปยุตฺตสฺส จ ธมฺมสฺส วิปฺปยุตฺตปจฺจเยน ปจฺจโย.\csromanspacer{} \textbf{ปุเรชาตํ} วตฺถุ ทิฏฺฐิคตวิปฺปยุตฺตโลภสหคตานํ ขนฺธานํ โลภสฺส จ วิปฺปยุตฺตปจฺจเยน ปจฺจโย. (3)}

\prose{อุปาทานสมฺปยุตฺโต จ อุปาทานวิปฺปยุตฺโต จ ธมฺมา อุปาทานวิปฺปยุตฺตสฺส ธมฺมสฺส วิปฺปยุตฺตปจฺจเยน ปจฺจโย.\csromanspacer{} สหชาตํ, ปจฺฉาชาตํ.\csromanspacer{}\csromanspacer{}\textbf{สหชาตา} ทิฏฺฐิคตวิปฺปยุตฺตโลภสหคตา ขนฺธา จ โลโภ จ จิตฺตสมุฏฺฐานานํ รูปานํ วิปฺปยุตฺตปจฺจเยน ปจฺจโย.\csromanspacer{} ปจฺฉาชาตา\csromandash{} ฯเปฯ. (1)}

\csromanheader{2}{\textbf{อตฺถิ-อาทิ}}
\proseitem{85}{อุปาทานสมฺปยุตฺโต ธมฺโม อุปาทานสมฺปยุตฺตสฺส ธมฺมสฺส อตฺถิปจฺจเยน ปจฺจโย.\csromanspacer{} เอกํ. (ปฏิจฺจสทิสํ.) (1)}

\prose{อุปาทานสมฺปยุตฺโต ธมฺโม อุปาทานวิปฺปยุตฺตสฺส ธมฺมสฺส อตฺถิปจฺจเยน ปจฺจโย.\csromanspacer{} สหชาตํ, ปจฺฉาชาตํ. (สํขิตฺตํ.) (2)}

\prose{อุปาทานสมฺปยุตฺโต ธมฺโม อุปาทานสมฺปยุตฺตสฺส จ อุปาทานวิปฺปยุตฺตสฺส จ ธมฺมสฺส อตฺถิปจฺจเยน ปจฺจโย. (ปฏิจฺจสทิสํ.) (3)}

\proseitem{86}{อุปาทานวิปฺปยุตฺโต ธมฺโม อุปาทานวิปฺปยุตฺตสฺส ธมฺมสฺส อตฺถิปจฺจเยน ปจฺจโย.\csromanspacer{} สหชาตํ, ปุเรชาตํ, ปจฺฉาชาตํ, อาหารํ.\csromanspacer{} อินฺทฺริยํ. (สํขิตฺตํ.) (1)}

\prose{อุปาทานวิปฺปยุตฺโต ธมฺโม อุปาทานสมฺปยุตฺตสฺส ธมฺมสฺส อตฺถิปจฺจเยน ปจฺจโย, สหชาตํ, ปุเรชาตํ. (สํขิตฺตํ.) (2)}

\prose{\csromanfolio{594}อุปาทานวิปฺปยุตฺโต ธมฺโม อุปาทานสมฺปยุตฺตสฺส จ อุปาทานวิปฺปยุตฺตสฺส จ ธมฺมสฺส อตฺถิปจฺจเยน ปจฺจโย.\csromanspacer{} \textbf{สหชาตํ, ปุเรชาตํ}. (อิเมสุ สหชาตํ สหชาตสทิสํ, ปุเรชาตํ ปุเรชาตสทิสํ.) (3)}

\proseitem{87}{อุปาทานสมฺปยุตฺโต จ อุปาทานวิปฺปยุตฺโต จ ธมฺมา อุปาทานสมฺปยุตฺตสฺส ธมฺมสฺส อตฺถิปจฺจเยน ปจฺจโย.\csromanspacer{} \textbf{สหชาตํ, ปุเรชาตํ}.\csromanspacer{}\csromanspacer{}\textbf{สหชาโต} อุปาทานสมฺปยุตฺโต เอโก ขนฺโธ จ วตฺถุ จ ติณฺณนฺนํ ขนฺธานํ อตฺถิปจฺจเยน ปจฺจโย ฯเปฯ.\csromanspacer{} เทฺว ขนฺธา จ ฯเปฯ.\csromanspacer{} \textbf{สหชาโต} ทิฏฺฐิคตวิปฺปยุตฺตโลภสหคโต เอโก ขนฺโธ จ โลโภ จ ติณฺณนฺนํ ขนฺธานํ อตฺถิปจฺจเยน ปจฺจโย ฯเปฯ.\csromanspacer{} เทฺว ขนฺธา จ ฯเปฯ. (1)}

\prose{อุปาทานสมฺปยุตฺโต จ อุปาทานวิปฺปยุตฺโต จ ธมฺมา อุปาทานวิปฺปยุตฺตสฺส ธมฺมสฺส อตฺถิปจฺจเยน ปจฺจโย.\csromanspacer{} \textbf{สหชาตํ, ปุเรชาตํ}, ปจฺฉาชาตํ, อาหารํ, อินฺทฺริยํ.\csromanspacer{}\csromanspacer{}\textbf{สหชาตา} อุปาทานสมฺปยุตฺตา ขนฺธา จ มหาภูตา จ จิตฺตสมุฏฺฐานานํ รูปานํ อตฺถิปจฺจเยน ปจฺจโย.\csromanspacer{} \textbf{สหชาตา} ทิฏฺฐิคตวิปฺปยุตฺตโลภสหคตา ขนฺธา จ โลโภ จ จิตฺตสมุฏฺฐานานํ รูปานํ อตฺถิปจฺจเยน ปจฺจโย.\csromanspacer{} \textbf{สหชาตา} ทิฏฺฐิคตวิปฺปยุตฺตโลภสหคตา ขนฺธา จ วตฺถุ จ โลภสฺส อตฺถิปจฺจเยน ปจฺจโย.\csromanspacer{} \textbf{ปจฺฉาชาตา} ทิฏฺฐิคตวิปฺปยุตฺตโลภสหคตา ขนฺธา จ โลโภ จ ปุเรชาตสฺส อิมสฺส กายสฺส อตฺถิปจฺจเยน ปจฺจโย.\csromanspacer{} \textbf{ปจฺฉาชาตา} ทิฏฺฐิคตวิปฺปยุตฺตโลภสหคตา ขนฺธา จ โลโภ จ กพฬีกาโร อาหาโร จ อิมสฺส กายสฺส อตฺถิปจฺจเยน ปจฺจโย.\csromanspacer{} \textbf{ปจฺฉาชาตา} ทิฏฺฐิคตวิปฺปยุตฺตโลภสหคตา ขนฺธา จ โลโภ จ รูปชีวิตินฺทฺริยญฺจ กฏตฺตารูปานํ อตฺถิปจฺจเยน ปจฺจโย. (2)}

\prose{อุปาทานสมฺปยุตฺโต จ อุปาทานวิปฺปยุตฺโต จ ธมฺมา อุปาทานสมฺปยุตฺตสฺส จ อุปาทานวิปฺปยุตฺตสฺส จ ธมฺมสฺส อตฺถิปจฺจเยน ปจฺจโย.\csromanspacer{} \textbf{สหชาตํ, ปุเรชาตํ}.\csromanspacer{}\csromanspacer{}\textbf{สหชาโต} ทิฏฺฐิคตวิปฺปยุตฺตโลภสหคโต เอโก ขนฺโธ จ โลโภ จ ติณฺณนฺนํ ขนฺธานํ จิตฺตสมุฏฺฐานานญฺจ รูปานํ อตฺถิปจฺจเยน ปจฺจโย ฯเปฯ.\csromanspacer{} เทฺว ขนฺธา จ ฯเปฯ.\csromanspacer{} \textbf{สหชาโต} ทิฏฺฐิคตวิปฺปยุตฺตโลภสหคโต เอโก ขนฺโธ จ วตฺถุ จ ติณฺณนฺนํ ขนฺธานํ โลภสฺส จ อตฺถิปจฺจเยน ปจฺจโย ฯเปฯ.\csromanspacer{} เทฺว ขนฺธา จ ฯเปฯ. (3)}

\csromanheader{2}{71. อุปาทานสมฺปยุตฺตทุก}
\prose{\csromanfolio{595}นตฺถิปจฺจเยน ปจฺจโย.\csromanspacer{} วิคตปจฺจเยน ปจฺจโย.\csromanspacer{} อวิคตปจฺจเยน ปจฺจโย.}

\csromanheader{2}{1. ปจฺจยานุโลม 2. สงฺขฺยาวาร}
\csromanheader{2}{\textbf{สุทฺธ}}
\proseitem{88}{เหตุยา นว, อารมฺมเณ นว, อธิปติยา นว, อนนฺตเร นว, สมนนฺตเร นว, สหชาเต นว, อญฺญมญฺเญ ฉ, นิสฺสเย นว, อุปนิสฺสเย นว, ปุเรชาเต ตีณิ, ปจฺฉาชาเต ตีณิ, อาเสวเน นว, กมฺเม จตฺตาริ, วิปาเก เอกํ, อาหาเร จตฺตาริ, อินฺทฺริเย จตฺตาริ, ฌาเน จตฺตาริ, มคฺเค จตฺตาริ, สมฺปยุตฺเต ฉ, วิปฺปยุตฺเต ปญฺจ, อตฺถิยา นว, นตฺถิยา นว, วิคเต นว, อวิคเต นว.}

\csromanheader{2}{2. ปจฺจนียุทฺธาร}
\proseitem{89}{อุปาทานสมฺปยุตฺโต ธมฺโม อุปาทานสมฺปยุตฺตสฺส ธมฺมสฺส อารมฺมณปจฺจเยน ปจฺจโย.\csromanspacer{} สหชาตปจฺจเยน ปจฺจโย.\csromanspacer{} อุปนิสฺสยปจฺจเยน ปจฺจโย. (1)}

\prose{อุปาทานสมฺปยุตฺโต ธมฺโม อุปาทานวิปฺปยุตฺตสฺส ธมฺมสฺส อารมฺมณปจฺจเยน ปจฺจโย.\csromanspacer{} สหชาตปจฺจเยน ปจฺจโย.\csromanspacer{} อุปนิสฺสยปจฺจเยน ปจฺจโย.\csromanspacer{} ปจฺฉาชาตปจฺจเยน ปจฺจโย.\csromanspacer{} กมฺมปจฺจเยน ปจฺจโย. (2)}

\prose{อุปาทานสมฺปยุตฺโต ธมฺโม อุปาทานสมฺปยุตฺตสฺส จ อุปาทานวิปฺปยุตฺตสฺส จ ธมฺมสฺส อารมฺมณปจฺจเยน ปจฺจโย.\csromanspacer{} สหชาตปจฺจเยน ปจฺจโย.\csromanspacer{} อุปนิสฺสยปจฺจเยน ปจฺจโย. (3)}

\proseitem{90}{อุปาทานวิปฺปยุตฺโต ธมฺโม อุปาทานวิปฺปยุตฺตสฺส ธมฺมสฺส อารมฺมณปจฺจเยน ปจฺจโย.\csromanspacer{} สหชาตปจฺจเยน ปจฺจโย.\csromanspacer{} อุปนิสฺสยปจฺจเยน ปจฺจโย.\csromanspacer{} ปุเรชาตปจฺจเยน ปจฺจโย.\csromanspacer{} ปจฺฉาชาตปจฺจเยน ปจฺจโย.\csromanspacer{} กมฺมปจฺจเยน ปจฺจโย.\csromanspacer{} อาหารปจฺจเยน ปจฺจโย.\csromanspacer{} อินฺทฺริยปจฺจเยน ปจฺจโย. (1)}

\prose{\csromanfolio{596}อุปาทานวิปฺปยุตฺโต ธมฺโม อุปาทานสมฺปยุตฺตสฺส ธมฺมสฺส อารมฺมณปจฺจเยน ปจฺจโย.\csromanspacer{} สหชาตปจฺจเยน ปจฺจโย.\csromanspacer{} อุปนิสฺสยปจฺจเยน ปจฺจโย.\csromanspacer{} ปุเรชาตปจฺจเยน ปจฺจโย. (2)}

\prose{อุปาทานวิปฺปยุตฺโต ธมฺโม อุปาทานสมฺปยุตฺตสฺส จ อุปาทานวิปฺปยุตฺตสฺส จ ธมฺมสฺส อารมฺมณปจฺจเยน ปจฺจโย.\csromanspacer{} สหชาตปจฺจเยน ปจฺจโย.\csromanspacer{} อุปนิสฺสยปจฺจเยน ปจฺจโย.\csromanspacer{} ปุเรชาตปจฺจเยน ปจฺจโย. (3)}

\proseitem{91}{อุปาทานสมฺปยุตฺโต จ อุปาทานวิปฺปยุตฺโต จ ธมฺมา อุปาทานสมฺปยุตฺตสฺส ธมฺมสฺส อารมฺมณปจฺจเยน ปจฺจโย.\csromanspacer{} สหชาตปจฺจเยน ปจฺจโย.\csromanspacer{} อุปนิสฺสยปจฺจเยน ปจฺจโย. (1)}

\prose{อุปาทานสมฺปยุตฺโต จ อุปาทานวิปฺปยุตฺโต จ ธมฺมา อุปาทานวิปฺปยุตฺตสฺส ธมฺมสฺส อารมฺมณปจฺจเยน ปจฺจโย.\csromanspacer{} สหชาตปจฺจเยน ปจฺจโย.\csromanspacer{} อุปนิสฺสยปจฺจเยน ปจฺจโย.\csromanspacer{} ปจฺฉาชาตปจฺจเยน ปจฺจโย. (2)}

\prose{อุปาทานสมฺปยุตฺโต จ อุปาทานวิปฺปยุตฺโต จ ธมฺมา อุปาทานสมฺปยุตฺตสฺส จ อุปาทานวิปฺปยุตฺตสฺส จ ธมฺมสฺส อารมฺมณปจฺจเยน ปจฺจโย.\csromanspacer{} สหชาตปจฺจเยน ปจฺจโย.\csromanspacer{} อุปนิสฺสยปจฺจเยน ปจฺจโย. (3)}

\csromanheader{2}{2. ปจฺจยปจฺจนีย 2. สงฺขฺยาวาร}
\proseitem{92}{นเหตุยา นว, น-อารมฺมเณ นว, น-อธิปติยา นว, (สพฺพตฺถ นว.) โน-อวิคเต นว.}

\csromanheader{2}{3. ปจฺจยานุโลมปจฺจนีย}
\proseitem{93}{เหตุปจฺจยา น-อารมฺมเณ นว, น-อธิปติยา นว, น-อนนฺตเร นว, นสมนนฺตเร นว, น-อญฺญมญฺเญ ตีณิ, น-อุปนิสฺสเย นว, (สพฺพตฺถ นว.) นสมฺปยุตฺเต ตีณิ, นวิปฺปยุตฺเต ฉ, โนนตฺถิยา นว, โนวิคเต นว.}

\csromanheader{2}{4. ปจฺจยปจฺจนียานุโลม}
\proseitem{94}{นเหตุปจฺจยา อารมฺมเณ นว, อธิปติยา นว, (อนุโลมมาติกา.) ฯเปฯ อวิคเต นว.}

\nitthitam{อุปาทานสมฺปยุตฺตทุกํ นิฏฺฐิตํ.}
\vspace{\tipitakaparskip}
\csromancenter{อุปาทาน-อุปาทานิยทุก \textbf{72.}\csromanspacer{}\textbf{ อุปาทาน-อุปาทานิยทุก 1.}\csromanspacer{}\textbf{ ปฏิจฺจวาร}}
\csromanheader{2}{1. ปจฺจยานุโลม 1. วิภงฺควาร}
\csromanheader{2}{\textbf{เหตุ}}
\csromanmatikaanchor{mtk.72}
\csromantocmarkref{subsection}{72. อุปาทาน-อุปาทานิยทุก}{mtk.72}
\bookmark[dest={mtk.72},level=2]{72. อุปาทาน-อุปาทานิยทุก}
\proseitem{95}{\csromanfolio{597}อุปาทานญฺเจว อุปาทานิยญฺจ ธมฺมํ ปฏิจฺจ อุปาทาโน เจว อุปาทานิโย จ ธมฺโม อุปฺปชฺชติ เหตุปจฺจยา.\csromanspacer{} ทิฏฺฐุปาทานํ ปฏิจฺจ กามุปาทานํ.\csromanspacer{} กามุปาทานํ ปฏิจฺจ ทิฏฺฐุปาทานํ. (จกฺกํ.) (1)}

\prose{อุปาทานญฺเจว อุปาทานิยญฺจ ธมฺมํ ปฏิจฺจ อุปาทานิโย เจว โน จ อุปาทาโน ธมฺโม อุปฺปชฺชติ เหตุปจฺจยา.\csromanspacer{} อุปาทาเน ปฏิจฺจ สมฺปยุตฺตกา ขนฺธา จิตฺตสมุฏฺฐานญฺจ รูปํ. (2)}

\prose{อุปาทานญฺเจว อุปาทานิยญฺจ ธมฺมํ ปฏิจฺจ อุปาทาโน เจว อุปาทานิโย จ อุปาทานิโย เจว โน จ อุปาทาโน จ ธมฺมา อุปฺปชฺชนฺติ เหตุปจฺจยา.\csromanspacer{} ทิฏฺฐุปาทานํ ปฏิจฺจ กามุปาทานํ สมฺปยุตฺตกา จ ขนฺธา จิตฺตสมุฏฺฐานญฺจ รูปํ. (จกฺกํ.) (3)}

\proseitem{96}{อุปาทานิยญฺเจว โน จ อุปาทานํ ธมฺมํ ปฏิจฺจ อุปาทานิโย เจว โน จ อุปาทาโน ธมฺโม อุปฺปชฺชติ เหตุปจฺจยา.\csromanspacer{} อุปาทานิยญฺเจว โน จ อุปาทานํ เอกํ ขนฺธํ ปฏิจฺจ ตโย ขนฺธา จิตฺตสมุฏฺฐานญฺจ รูปํ ฯเปฯ.\csromanspacer{} เทฺว ขนฺเธ ฯเปฯ.\csromanspacer{} ปฏิสนฺธิกฺขเณ ฯเปฯ. (ยาว อชฺฌตฺติกา มหาภูตา.) (1)}

\prose{อุปาทานิยญฺเจว โน จ อุปาทานํ ธมฺมํ ปฏิจฺจ อุปาทาโน เจว อุปาทานิโย จ ธมฺโม อุปฺปชฺชติ เหตุปจฺจยา.\csromanspacer{} อุปาทานิเย เจว โน จ อุปาทาเน ขนฺเธ ปฏิจฺจ อุปาทานา. (2)}

\prose{อุปาทานิยญฺเจว โน จ อุปาทานํ ธมฺมํ ปฏิจฺจ อุปาทาโน เจว อุปาทานิโย จ อุปาทานิโย เจว โน จ อุปาทาโน จ ธมฺมา อุปฺปชฺชนฺติ เหตุปจฺจยา.\csromanspacer{} อุปาทานิยญฺเจว โน จ อุปาทานํ เอกํ ขนฺธํ ปฏิจฺจ ตโย ขนฺธา อุปาทานา จ จิตฺตสมุฏฺฐานญฺจ รูปํ ฯเปฯ.\csromanspacer{} เทฺว ขนฺเธ ฯเปฯ. (3)}

\proseitem{97}{อุปาทานญฺเจว อุปาทานิยญฺจ อุปาทานิยญฺเจว โน จ อุปาทานญฺจ ธมฺมํ ปฏิจฺจ อุปาทาโน เจว อุปาทานิโย จ ธมฺโม อุปฺปชฺชติ เหตุปจฺจยา.\csromanspacer{} ทิฏฺฐุปาทานญฺจ สมฺปยุตฺตเก จ ขนฺเธ ปฏิจฺจ กามุปาทานํ. (จกฺกํ.) (1)}

\prose{\csromanfolio{598}อุปาทานญฺเจว อุปาทานิยญฺจ อุปาทานิยญฺเจว โน จ อุปาทานญฺจ ธมฺมํ ปฏิจฺจ อุปาทานิโย เจว โน จ อุปาทาโน ธมฺโม อุปฺปชฺชติ เหตุปจฺจยา.\csromanspacer{} อุปาทานิยญฺเจว โน จ อุปาทานํ เอกํ ขนฺธญฺจ อุปาทาเน จ ปฏิจฺจ ตโย ขนฺธา จิตฺตสมุฏฺฐานญฺจ รูปํ ฯเปฯ เทฺว ขนฺเธ จ ฯเปฯ.\csromanspacer{} อุปาทาเน จ มหาภูเต จ ปฏิจฺจ จิตฺตสมุฏฺฐานํ รูปํ. (2)}

\prose{อุปาทานญฺเจว อุปาทานิยญฺจ อุปาทานิยญฺเจว โน จ อุปาทานญฺจ ธมฺมํ ปฏิจฺจ อุปาทาโน เจว อุปาทานิโย จ อุปาทานิโย เจว โน จ อุปาทาโน จ ธมฺมา อุปฺปชฺชนฺติ เหตุปจฺจยา.\csromanspacer{} อุปาทานิยญฺเจว โน จ อุปาทานํ เอกํ ขนฺธญฺจ ทิฏฺฐุปาทานญฺจ ปฏิจฺจ ตโย ขนฺธา กามุปาทานญฺจ จิตฺตสมุฏฺฐานญฺจ รูปํ ฯเปฯ.\csromanspacer{} เทฺว ขนฺเธ จ ฯเปฯ. (จกฺกํ.) (3)}

\prose{(ยถา อุปาทานทุกํ.\csromanspacer{} เอวํ ปฏิจฺจวาโรปิ สหชาตวาโรปิ ปจฺจยวาโรปิ นิสฺสยวาโรปิ สํสฏฺฐวาโรปิ สมฺปยุตฺตวาโรปิ กาตพฺพา.\csromanspacer{} นินฺนานากรณา, อามสนํ นานากรณํ.)}

\csromanheader{2}{72. อุปาทาน-อุปาทานิยทุก 7. ปญฺหาวาร}
\csromanheader{2}{1. ปจฺจยานุโลม 1. วิภงฺควาร}
\csromanheader{2}{\textbf{เหตุ}}
\proseitem{98}{อุปาทาโน เจว อุปาทานิโย จ ธมฺโม อุปาทานสฺส เจว อุปาทานิยสฺส จ ธมฺมสฺส เหตุปจฺจเยน ปจฺจโย.\csromanspacer{} อุปาทานา เจว อุปาทานิยา จ เหตู สมฺปยุตฺตกานํ อุปาทานานํ เหตุปจฺจเยน ปจฺจโย. (1)}

\prose{อุปาทาโน เจว อุปาทานิโย จ ธมฺโม อุปาทานิยสฺส เจว โน จ อุปาทานสฺส ธมฺมสฺส เหตุปจฺจเยน ปจฺจโย.\csromanspacer{} อุปาทานา เจว อุปาทานิยา จ เหตู สมฺปยุตฺตกานํ ขนฺธานํ จิตฺตสมุฏฺฐานานญฺจ รูปานํ เหตุปจฺจเยน ปจฺจโย. (อุปาทานทุกสทิสา, นินฺนานา.\csromanspacer{} นว ปญฺหา.)}

\csromanheader{2}{72. อุปาทาน-อุปาทานิยทุก \textbf{อารมฺมณ}}
\proseitem{99}{\csromanfolio{599}อุปาทาโน เจว อุปาทานิโย จ ธมฺโม อุปาทานสฺส เจว อุปาทานิยสฺส จ ธมฺมสฺส อารมฺมณปจฺจเยน ปจฺจโย.\csromanspacer{} อุปาทาเน อารพฺภ อุปาทานา อุปฺปชฺชนฺติ.\csromanspacer{} ตีณิ.}

\prose{อุปาทานิโย เจว โน จ อุปาทาโน ธมฺโม อุปาทานิยสฺส เจว โน จ อุปาทานสฺส ธมฺมสฺส อารมฺมณปจฺจเยน ปจฺจโย.\csromanspacer{} ทานํ ฯเปฯ.\csromanspacer{} ฌานา วุฏฺฐหิตฺวา ฌานํ ปจฺจเวกฺขติ, อสฺสาเทติ อภินนฺทติ, ตํ อารพฺภ ราโค อุปฺปชฺชติ.\csromanspacer{} ทิฏฺฐิ อุปฺปชฺชติ.\csromanspacer{} วิจิกิจฺฉา.\csromanspacer{} อุทฺธจฺจํ.\csromanspacer{} โทมนสฺสํ อุปฺปชฺชติ.\csromanspacer{} อริยา โคตฺรภุํ ปจฺจเวกฺขนฺติ.\csromanspacer{} โวทานํ ปจฺจเวกฺขนฺติ.\csromanspacer{} ปหีเน กิเลเส ฯเปฯ.\csromanspacer{} วิกฺขมฺภิเต กิเลเส ฯเปฯ.\csromanspacer{} ปุพฺเพ สมุทาจิณฺเณ ฯเปฯ.\csromanspacer{} จกฺขุํ ฯเปฯ.\csromanspacer{} วตฺถุํ ฯเปฯ. (สํขิตฺตํ.) อนาคตํสญาณสฺส, อาวชฺชนาย อารมฺมณปจฺจเยน ปจฺจโย. (1)}

\prose{อุปาทานิโย เจว โน จ อุปาทาโน ธมฺโม อุปาทานสฺส เจว อุปาทานิยสฺส จ ธมฺมสฺส อารมฺมณปจฺจเยน ปจฺจโย. (สํขิตฺตํ.\csromanspacer{} อิตเร เทฺว อุปาทานทุกสทิสา.) (3)}

\prose{อุปาทาโน เจว อุปาทานิโย จ อุปาทานิโย เจว โน จ อุปาทาโน จ ธมฺมา อุปาทานสฺส เจว อุปาทานิยสฺส จ ธมฺมสฺส อารมฺมณปจฺจเยน ปจฺจโย.\csromanspacer{} ตีณิ. (เหฏฺฐา อธิปติ.\csromanspacer{} ตีณิ.\csromanspacer{} อุปาทานทุกสทิสา.)}

\csromanheader{2}{\textbf{อธิปติ}}
\proseitem{100}{อุปาทานิโย เจว โน จ อุปาทาโน ธมฺโม อุปาทานิยสฺส เจว โน จ อุปาทานสฺส จ ธมฺมสฺส อธิปติปจฺจเยน ปจฺจโย.\csromanspacer{} \textbf{อารมฺมณาธิปติ}, สหชาตาธิปติ.\csromanspacer{}\csromanspacer{}\textbf{อารมฺมณาธิปติ}\csromandash{}ทานํ.\csromanspacer{} ฌานา วุฏฺฐหิตฺวา ฌานํ ครุํ กตฺวา ปจฺจเวกฺขติ, อสฺสาเทติ อภินนฺทติ, ตํ ครุํ กตฺวา ราโค อุปฺปชฺชติ.\csromanspacer{} ทิฏฺฐิ อุปฺปชฺชติ.\csromanspacer{} เสกฺขา โคตฺรภุํ ครุํ กตฺวา ฯเปฯ.\csromanspacer{} โวทานํ ฯเปฯ.\csromanspacer{} จกฺขุํ ฯเปฯ วตฺถุํ อุปาทานิเย เจว โน จ อุปาทาเน ขนฺเธ ครุํ กตฺวา อุปาทานิยา เจว โน จ อุปาทานา ขนฺธา อุปฺปชฺชนฺติ.\csromanspacer{} \textbf{สหชาตาธิปติ}\csromandash{} อุปาทานิยา เจว โน จ อุปาทานาธิปติ สมฺปยุตฺตกานํ ขนฺธานํ จิตฺตสมุฏฺฐานานญฺจ รูปานํ อธิปติปจฺจเยน ปจฺจโย. (อวเสสา เทฺวปิ อารมฺมณาธิปติปิ สหชาตาธิปติปิ อุปาทานทุกสทิสา.) (3)}

\prose{\csromanfolio{600}(ฆฏนา อธิปติปิ ตีณิ.\csromanspacer{} อุปาทานทุกสทิสา.\csromanspacer{} สพฺเพ ปจฺจยา อุปาทานทุกสทิสา.\csromanspacer{} อุปาทานิเย โลกุตฺตรํ นตฺถิ, ปจฺจนียมฺปิ อิตเร เทฺว คณนาปิ อุปาทานทุกสทิสํ.)}

\nitthitam{อุปาทาน-อุปาทานิยทุกํ นิฏฺฐิตํ.}
\csromanmatikaanchor{mtk.73}
\csromantocmarkref{subsection}{73. อุปาทาน-อุปาทานสมฺปยุตฺตทุก}{mtk.73}
\bookmark[dest={mtk.73},level=2]{73. อุปาทาน-อุปาทานสมฺปยุตฺตทุก}
\csromanheader{2}{73. อุปาทาน-อุปาทานสมฺปยุตฺตทุก 1. ปฏิจฺจวาร}
\csromanheader{2}{1. ปจฺจยานุโลม 1. วิภงฺควาร}
\csromanheader{2}{\textbf{เหตุ}}
\proseitem{101}{อุปาทานญฺเจว อุปาทานสมฺปยุตฺตญฺจ ธมฺมํ ปฏิจฺจ อุปาทาโน เจว อุปาทานสมฺปยุตฺโต จ ธมฺโม อุปฺปชฺชติ เหตุปจฺจยา.\csromanspacer{} ทิฏฺฐุปาทานํ ปฏิจฺจ กามุปาทานํ. (จกฺกํ กาตพฺพํ.) (1)}

\prose{อุปาทานญฺเจว อุปาทานสมฺปยุตฺตญฺจ ธมฺมํ ปฏิจฺจ อุปาทานสมฺปยุตฺโต เจว โน จ อุปาทาโน ธมฺโม อุปฺปชฺชติ เหตุปจฺจยา.\csromanspacer{} อุปาทาเน ปฏิจฺจ สมฺปยุตฺตกา ขนฺธา. (2)}

\prose{อุปาทานญฺเจว อุปาทานสมฺปยุตฺตญฺจ ธมฺมํ ปฏิจฺจ อุปาทาโน เจว อุปาทานสมฺปยุตฺโต จ อุปาทานสมฺปยุตฺโต เจว โน จ อุปาทาโน จ ธมฺมา อุปฺปชฺชนฺติ เหตุปจฺจยา.\csromanspacer{} ทิฏฺฐุปาทานํ ปฏิจฺจ กามุปาทานํ สมฺปยุตฺตกา จ ขนฺธา. (จกฺกํ.) (3)}

\prose{อุปาทานสมฺปยุตฺตญฺเจว โน จ อุปาทานํ ธมฺมํ ปฏิจฺจ ฯเปฯ.}

\prose{(สํขิตฺตํ.\csromanspacer{} อามสนํ นานากรณํ.\csromanspacer{} อุปาทานทุกสทิสํ.\csromanspacer{} นว ปญฺหา.\csromanspacer{} รูปํ นตฺถิ.\csromanspacer{} เอวํ สพฺเพปิ วารา วิตฺถาเรตพฺพา.\csromanspacer{} อรูปํเยว.)}

\csromanheader{2}{73. อุปาทาน-อุปาทานสมฺปยุตฺตทุก 7. ปญฺหาวาร}
\csromanheader{2}{1. ปจฺจยานุโลม 1. วิภงฺควาร}
\csromanheader{2}{\textbf{เหตุ}}
\proseitem{102}{อุปาทาโน เจว อุปาทานสมฺปยุตฺโต จ ธมฺโม อุปาทานสฺส เจว อุปาทานสมฺปยุตฺตสฺส จ ธมฺมสฺส เหตุปจฺจเยน ปจฺจโย.\csromanspacer{} อุปาทานา เจว}

\vspace{\tipitakaparskip}
\csromancenter{อุปาทาน-อุปาทานสมฺปยุตฺตทุก อุปาทานสมฺปยุตฺตา จ เหตู สมฺปยุตฺตกานํ อุปาทานานํ เหตุปจฺจเยน ปจฺจโย. (มูลํ ปุจฺฉิตพฺพํ.) อุปาทานา เจว อุปาทานสมฺปยุตฺตา จ เหตู สมฺปยุตฺตกานํ ขนฺธานํ เหตุปจฺจเยน ปจฺจโย. (มูลํ ปุจฺฉิตพฺพํ.) อุปาทานา เจว อุปาทานสมฺปยุตฺตา จ เหตู สมฺปยุตฺตกานํ ขนฺธานํ อุปาทานานญฺจ เหตุปจฺจเยน ปจฺจโย. (3)}
\proseitem{103}{\csromanfolio{601}อุปาทานสมฺปยุตฺโต เจว โน จ อุปาทาโน ธมฺโม อุปาทานสมฺปยุตฺตสฺส เจว โน จ อุปาทานสฺส ธมฺมสฺส เหตุปจฺจเยน ปจฺจโย.\csromanspacer{} อุปาทานสมฺปยุตฺตา เจว โน จ อุปาทานา เหตู สมฺปยุตฺตกานํ ขนฺธานํ เหตุปจฺจเยน ปจฺจโย. (มูลํ ปุจฺฉิตพฺพํ.) อุปาทานสมฺปยุตฺตา เจว โน จ อุปาทานา เหตู สมฺปยุตฺตกานํ อุปาทานานํ เหตุปจฺจเยน ปจฺจโย. (มูลํ ปุจฺฉิตพฺพํ.) อุปาทานสมฺปยุตฺตา เจว โน จ อุปาทานา เหตู สมฺปยุตฺตกานํ ขนฺธานํ อุปาทานานญฺจ เหตุปจฺจเยน ปจฺจโย. (3)}

\proseitem{104}{อุปาทาโน เจว อุปาทานสมฺปยุตฺโต จ อุปาทานสมฺปยุตฺโต เจว โน จ อุปาทาโน จ ธมฺมา อุปาทานสฺส เจว อุปาทานสมฺปยุตฺตสฺส จ ธมฺมสฺส เหตุปจฺจเยน ปจฺจโย.\csromanspacer{} อุปาทานา เจว อุปาทานสมฺปยุตฺตา จ อุปาทานสมฺปยุตฺตา เจว โน จ อุปาทานา จ เหตู สมฺปยุตฺตกานํ อุปาทานานํ เหตุปจฺจเยน ปจฺจโย. (มูลํ ปุจฺฉิตพฺพํ.) อุปาทานา เจว อุปาทานสมฺปยุตฺตา จ อุปาทานสมฺปยุตฺตา เจว โน จ อุปาทานา จ เหตู สมฺปยุตฺตกานํ ขนฺธานํ เหตุปจฺจเยน ปจฺจโย. (มูลํ ปุจฺฉิตพฺพํ.) อุปาทานา เจว อุปาทานสมฺปยุตฺตา จ อุปาทานสมฺปยุตฺตา เจว โน จ อุปาทานา จ เหตู สมฺปยุตฺตกานํ ขนฺธานํ อุปาทานานญฺจ เหตุปจฺจเยน ปจฺจโย. (3)}

\csromanheader{2}{\textbf{อารมฺมณ}}
\proseitem{105}{อุปาทาโน เจว อุปาทานสมฺปยุตฺโต จ ธมฺโม อุปาทานสฺส เจว อุปาทานสมฺปยุตฺตสฺส จ ธมฺมสฺส อารมฺมณปจฺจเยน ปจฺจโย.\csromanspacer{} อุปาทาเน อารพฺภ อุปาทานา อุปฺปชฺชนฺติ. (มูลํ ปุจฺฉิตพฺพํ.) อุปาทาเน อารพฺภ อุปาทานสมฺปยุตฺตา เจว โน จ อุปาทานา ขนฺธา อุปฺปชฺชนฺติ. (มูลํ ปุจฺฉิตพฺพํ.) อุปาทาเน อารพฺภ อุปาทานา จ สมฺปยุตฺตกา จ ขนฺธา อุปฺปชฺชนฺติ. (3)}

\prose{\csromanfolio{602}อุปาทานสมฺปยุตฺโต เจว โน จ อุปาทาโน ธมฺโม อุปาทานสมฺปยุตฺตสฺส เจว โน จ อุปาทานสฺส ธมฺมสฺส อารมฺมณปจฺจเยน ปจฺจโย.\csromanspacer{} อุปาทานสมฺปยุตฺเต เจว โน จ อุปาทาเน ขนฺเธ อารพฺภ อุปาทานสมฺปยุตฺตา เจว โน จ อุปาทานา ขนฺธา อุปฺปชฺชนฺติ. (ตีณิปิ กาตพฺพา, ฆฏเน ตีณิปิ กาตพฺพา.)}

\csromanheader{2}{\textbf{อธิปติ}}
\proseitem{106}{อุปาทาโน เจว อุปาทานสมฺปยุตฺโต จ ธมฺโม อุปาทานสฺส เจว อุปาทานสมฺปยุตฺตสฺส จ ธมฺมสฺส อธิปติปจฺจเยน ปจฺจโย.\csromanspacer{} ตีณิ. (3)}

\prose{อุปาทานสมฺปยุตฺโต เจว โน จ อุปาทาโน ธมฺโม อุปาทานสมฺปยุตฺตสฺส เจว โน จ อุปาทานสฺส ธมฺมสฺส อธิปติปจฺจเยน ปจฺจโย.\csromanspacer{} อารมฺมณาธิปติ, สหชาตาธิปติ. (ตีณิปิ, ตีสุปิ เทฺวปิ อธิปตี กาตพฺพา, ฆฏนาธิปติปิ ตีณิ.)}

\csromanheader{2}{\textbf{อนนฺตร}}
\proseitem{107}{อุปาทาโน เจว อุปาทานสมฺปยุตฺโต จ ธมฺโม อุปาทานสฺส เจว อุปาทานสมฺปยุตฺตสฺส จ ธมฺมสฺส อนนฺตรปจฺจเยน ปจฺจโย.\csromanspacer{} ปุริมา ปุริมา อุปาทานา เจว อุปาทานสมฺปยุตฺตา จ ขนฺธา ปจฺฉิมานํ ปจฺฉิมานํ อุปาทานานํ อนนฺตรปจฺจเยน ปจฺจโย. (เอวํ นวปิ ปญฺหา กาตพฺพา, อาวชฺชนาปิ วุฏฺฐานมฺปิ นตฺถิ.}

\csromanheader{2}{\textbf{สมนนฺตราทิ}}
\proseitem{108}{อุปาทาโน เจว อุปาทานสมฺปยุตฺโต จ ธมฺโม อุปาทานสฺส เจว อุปาทานสมฺปยุตฺตสฺส จ ธมฺมสฺส สมนนฺตรปจฺจเยน ปจฺจโย.\csromanspacer{} นว.\csromanspacer{} สหชาตปจฺจเยน ปจฺจโย.\csromanspacer{} นว.\csromanspacer{} อญฺญมญฺญปจฺจเยน ปจฺจโย.\csromanspacer{} นว.\csromanspacer{} นิสฺสยปจฺจเยน ปจฺจโย.\csromanspacer{} นว.}

\csromanheader{2}{\textbf{อุปนิสฺสย}}
\proseitem{109}{อุปาทาโน เจว อุปาทานสมฺปยุตฺโต จ ธมฺโม อุปาทานสฺส เจว อุปาทานสมฺปยุตฺตสฺส จ ธมฺมสฺส อุปนิสฺสยปจฺจเยน ปจฺจโย ฯเปฯ.\csromanspacer{} ตีณิ.}

\csromanheader{2}{73. อุปาทาน-อุปาทานสมฺปยุตฺตทุก}
\prose{\csromanfolio{603}อุปาทานสมฺปยุตฺโต เจว โน จ อุปาทาโน ธมฺโม อุปาทานสมฺปยุตฺตสฺส เจว โน จ อุปาทานสฺส ธมฺมสฺส อุปนิสฺสยปจฺจเยน ปจฺจโย.\csromanspacer{} อารมฺมณูปนิสฺสโย, อนนฺตรูปนิสฺสโย, ปกตูปนิสฺสโย ฯเปฯ.\csromanspacer{} ปกตูปนิสฺสโย\csromandash{}อุปาทานสมฺปยุตฺตา เจว โน จ อุปาทานา ขนฺธา อุปาทานสมฺปยุตฺตานญฺเจว โน จ อุปาทานานํ ขนฺธานํ อุปนิสฺสยปจฺจเยน ปจฺจโย. (ตีณิ ฆฏนุปนิสฺสเยปิ ตีณิ.) อาเสวนปจฺจเยนปจฺจโย.\csromanspacer{} นว.}

\csromanheader{2}{\textbf{กมฺมาทิ}}
\proseitem{110}{อุปาทานสมฺปยุตฺโต เจว โน จ อุปาทาโน ธมฺโม อุปาทานสมฺปยุตฺตสฺส เจว โน จ อุปาทานสฺส ธมฺมสฺส กมฺมปจฺจเยน ปจฺจโย.\csromanspacer{} ตีณิ.\csromanspacer{} อาหารปจฺจเยน ปจฺจโย.\csromanspacer{} ตีณิ.\csromanspacer{} อินฺทฺริยปจฺจเยน ปจฺจโย.\csromanspacer{} ตีณิ.\csromanspacer{} ฌานปจฺจเยน ปจฺจโย.\csromanspacer{} ตีณิ.\csromanspacer{} มคฺคปจฺจเยน ปจฺจโย.\csromanspacer{} นว.\csromanspacer{} สมฺปยุตฺตปจฺจเยน ปจฺจโย.\csromanspacer{} นว.\csromanspacer{} อตฺถิปจฺจเยน ปจฺจโย.\csromanspacer{} นว.\csromanspacer{} นตฺถิปจฺจเยน ปจฺจโย.\csromanspacer{} นว.\csromanspacer{} วิคตปจฺจเยน ปจฺจโย.\csromanspacer{} นว.\csromanspacer{} อวิคตปจฺจเยน ปจฺจโย.\csromanspacer{} นว.}

\csromanheader{2}{1. ปจฺจยานุโลม 2. สงฺขฺยาวาร}
\csromanheader{2}{\textbf{สุทฺธ}}
\proseitem{111}{เหตุยา นว, อารมฺมเณ นว, อธิปติยา นว, อนนฺตเร นว, สมนนฺตเร นว, สหชาเต นว, อญฺญมญฺเญ นว, นิสฺสเย นว, อุปนิสฺสเย นว, อาเสวเน นว, กมฺเม ตีณิ, อาหาเร ตีณิ, อินฺทฺริเย ตีณิ, ฌาเน ตีณิ, มคฺเค นว, สมฺปยุตฺเต นว, อตฺถิยา นว, นตฺถิยา นว, วิคเต นว, อวิคเต นว.}

\csromanheader{2}{2. ปจฺจนียุทฺธาร}
\proseitem{112}{อุปาทาโน เจว อุปาทานสมฺปยุตฺโต จ ธมฺโม อุปาทานสฺส เจว อุปาทานสมฺปยุตฺตสฺส จ ธมฺมสฺส อารมฺมณปจฺจเยน ปจฺจโย.\csromanspacer{} สหชาตปจฺจเยน ปจฺจโย.\csromanspacer{} อุปนิสฺสยปจฺจเยน ปจฺจโย. (1)}

\prose{\csromanfolio{604}อุปาทาโน เจว อุปาทานสมฺปยุตฺโต จ ธมฺโม อุปาทานสมฺปยุตฺตสฺส เจว โน จ อุปาทานสฺส ธมฺมสฺส อารมฺมณปจฺจเยน ปจฺจโย.\csromanspacer{} สหชาตปจฺจเยน ปจฺจโย.\csromanspacer{} อุปนิสฺสยปจฺจเยน ปจฺจโย. (2)}

\prose{อุปาทาโน เจว อุปาทานสมฺปยุตฺโต จ ธมฺโม อุปาทานสฺส เจว อุปาทานสมฺปยุตฺตสฺส จ อุปาทานสมฺปยุตฺตสฺส เจว โน จ อุปาทานสฺส จ ธมฺมสฺส อารมฺมณปจฺจเยน ปจฺจโย.\csromanspacer{} สหชาตปจฺจเยน ปจฺจโย.\csromanspacer{} อุปนิสฺสยปจฺจเยน ปจฺจโย.}

\prose{(เอวํ นวปิ กาตพฺพา, เอเกกสฺส มูเล ตีณิ ตีณิ ปญฺหา.)}

\csromanheader{2}{2. ปจฺจยปจฺจนีย 2. สงฺขฺยาวาร}
\proseitem{113}{นเหตุยา นว, น-อารมฺมเณ นว, น-อธิปติยา นว, (สพฺพตฺถ นว.) โน-อวิคเต นว.}

\csromanheader{2}{3. ปจฺจยานุโลมปจฺจนีย}
\proseitem{114}{\textbf{เหตุ}ปจฺจยา น-อารมฺมเณ นว, น-อธิปติยา นว, น-อนนฺตเร นว, นสมนนฺตเร นว, น-อุปนิสฺสเย นว, (สพฺพตฺถ นว.) นมคฺเค นว, นสมฺปยุตฺเต นว, โนนตฺถิยา นว, โนวิคเต นว.}

\csromanheader{2}{4. ปจฺจยปจฺจนียานุโลม}
\csromanheader{2}{115. นเหตุปจฺจยา อารมฺมเณ นว, อธิปติยา นว ฯเปฯ. (อนุโลมมาติกา กาตพฺพา.)}
\nitthitam{อุปาทาน-อุปาทานสมฺปยุตฺตทุกํ นิฏฺฐิตํ.}
\csromanmatikaanchor{mtk.74}
\csromantocmarkref{subsection}{74. อุปาทานวิปฺปยุตฺต-อุปาทานิยทุก}{mtk.74}
\bookmark[dest={mtk.74},level=2]{74. อุปาทานวิปฺปยุตฺต-อุปาทานิยทุก}
\csromanheader{2}{74. อุปาทานวิปฺปยุตฺต-อุปาทานิยทุก 1. ปฏิจฺจวาร}
\csromanheader{2}{1. ปจฺจยานุโลมาทิ}
\csromanheader{2}{\textbf{เหตุ}}
\proseitem{116}{อุปาทานวิปฺปยุตฺตํ อุปาทานิยํ ธมฺมํ ปฏิจฺจ อุปาทานวิปฺปยุตฺโต อุปาทานิโย ธมฺโม อุปฺปชฺชติ เหตุปจฺจยา.\csromanspacer{} อุปาทานวิปฺปยุตฺตํ อุปาทานิยํ}

\vspace{\tipitakaparskip}
\csromancenter{อุปาทานวิปฺปยุตฺต-อุปาทานิยทุก เอกํ ขนฺธํ ปฏิจฺจ ตโย ขนฺธา จิตฺตสมุฏฺฐานญฺจ รูปํ ฯเปฯ เทฺว ขนฺเธ ฯเปฯ.\csromanspacer{} ปฏิสนฺธิกฺขเณ ฯเปฯ.\csromanspacer{} เอกํ มหาภูตํ ฯเปฯ มหาภูเต ปฏิจฺจ จิตฺตสมุฏฺฐานํ รูปํ, กฏตฺตารูปํ, อุปาทารูปํ. (1)}
\prose{\csromanfolio{605}อุปาทานวิปฺปยุตฺตํ อนุปาทานิยํ ธมฺมํ ปฏิจฺจ อุปาทานวิปฺปยุตฺโต อนุปาทานิโย ธมฺโม อุปฺปชฺชติ เหตุปจฺจยา.\csromanspacer{} อุปาทานวิปฺปยุตฺตํ อนุปาทานิยํ เอกํ ขนฺธํ ปฏิจฺจ ตโย ขนฺธา ฯเปฯ เทฺว ขนฺเธ ฯเปฯ. (1)}

\prose{อุปาทานวิปฺปยุตฺตํ อนุปาทานิยํ ธมฺมํ ปฏิจฺจ อุปาทานวิปฺปยุตฺโต อุปาทานิโย ธมฺโม อุปฺปชฺชติ เหตุปจฺจยา.\csromanspacer{} อุปาทานวิปฺปยุตฺเต อนุปาทานิเย ขนฺเธ ปฏิจฺจ จิตฺตสมุฏฺฐานํ รูปํ.}

\prose{อุปาทานวิปฺปยุตฺตํ อนุปาทานิยํ ธมฺมํ ปฏิจฺจ อุปาทานวิปฺปยุตฺโต อุปาทานิโย จ อุปาทานวิปฺปยุตฺโต อนุปาทานิโย จ ธมฺมา อุปฺปชฺชนฺติ เหตุปจฺจยา.\csromanspacer{} อุปาทานวิปฺปยุตฺตํ อนุปาทานิยํ เอกํ ขนฺธํ ปฏิจฺจ ตโย ขนฺธา จิตฺตสมุฏฺฐานญฺจ รูปํ ฯเปฯ เทฺว ขนฺเธ ฯเปฯ. (3)}

\prose{อุปาทานวิปฺปยุตฺตํ อุปาทานิยญฺจ อุปาทานวิปฺปยุตฺตํ อนุปาทานิยญฺจ ธมฺมํ ปฏิจฺจ อุปาทานวิปฺปยุตฺโต อุปาทานิโย ธมฺโม อุปฺปชฺชติ เหตุปจฺจยา.\csromanspacer{} อุปาทานวิปฺปยุตฺเต อนุปาทานิเย ขนฺเธ จ มหาภูเต จ ปฏิจฺจ จิตฺตสมุฏฺฐานํ รูปํ. (1)}

\prose{เหตุยา ปญฺจ, อารมฺมเณ เทฺว ฯเปฯ อวิคเต ปญฺจ.}

\prose{(อิมํ ทุกํ จูฬนฺตรทุเก โลกิยทุกสทิสํ.\csromanspacer{} นินฺนานากรณํ.)}

\nitthitam{อุปาทานวิปฺปยุตฺต-อุปาทานิยทุกํ นิฏฺฐิตํ.}
\nitthitam{อุปาทานโคจฺฉกํ นิฏฺฐิตํ.}
